% 本文件由 LaTeX 排版 Agent 根据有效 translation.json 自动生成。
% author=John Chrysostom; work_id=2001; title=玛窦福音讲道集

\chapter{\WorkTitle{玛窦福音讲道集}}

% structure: p0001 source=h1 target=omitted reason=page-title-already-rendered-by-parent

\Emph{讲道一}\newline{}\Emph{讲道二}\newline{}\Emph{讲道三}\newline{}\Emph{讲道四}\newline{}\Emph{讲道五}\newline{}\Emph{讲道六}\newline{}\Emph{讲道七}\newline{}\Emph{讲道八}\newline{}\Emph{讲道九}\newline{}\Emph{讲道十}\newline{}\Emph{讲道十一}\newline{}\Emph{讲道十二}\newline{}\Emph{讲道十三}\newline{}\Emph{讲道十四}\newline{}\Emph{讲道十五}\newline{}\Emph{讲道十六}\newline{}\Emph{讲道十七}\newline{}\Emph{讲道十八}\newline{}\Emph{讲道十九}\newline{}\Emph{讲道二十}\newline{}\Emph{讲道二十一}\newline{}\Emph{讲道二十二}\newline{}\Emph{讲道二十三}\newline{}\Emph{讲道二十四}\newline{}\Emph{讲道二十五}\newline{}\Emph{讲道二十六}\newline{}\Emph{讲道二十七}\newline{}\Emph{讲道二十八}\newline{}\Emph{讲道二十九}\newline{}\Emph{讲道三十}\newline{}\Emph{讲道三十一}\newline{}\Emph{讲道三十二}\newline{}\Emph{讲道三十三}\newline{}\Emph{讲道三十四}\newline{}\Emph{讲道三十五}\newline{}\Emph{讲道三十六}\newline{}\Emph{讲道三十七}\newline{}\Emph{讲道三十八}\newline{}\Emph{讲道三十九}\newline{}\Emph{讲道四十}\newline{}\Emph{讲道四十一}\newline{}\Emph{讲道四十二}\newline{}\Emph{讲道四十三}\newline{}\Emph{讲道四十四}\newline{}\Emph{讲道四十五}\newline{}\Emph{讲道四十六}\newline{}\Emph{讲道四十七}\newline{}\Emph{讲道四十八}\newline{}\Emph{讲道四十九}\newline{}\Emph{讲道五十}\newline{}\Emph{讲道五十一}\newline{}\Emph{讲道五十二}\newline{}\Emph{讲道五十三}\newline{}\Emph{讲道五十四}\newline{}\Emph{讲道五十五}\newline{}\Emph{讲道五十六}\newline{}\Emph{讲道五十七}\newline{}\Emph{讲道五十八}\newline{}\Emph{讲道五十九}\newline{}\Emph{讲道六十}\newline{}\Emph{讲道六十一}\newline{}\Emph{讲道六十二}\newline{}\Emph{讲道六十三}\newline{}\Emph{讲道六十四}\newline{}\Emph{讲道六十五}\newline{}\Emph{讲道六十六}\newline{}\Emph{讲道六十七}\newline{}\Emph{讲道六十八}\newline{}\Emph{讲道六十九}\newline{}\Emph{讲道七十}\newline{}\Emph{讲道七十一}\newline{}\Emph{讲道七十二}\newline{}\Emph{讲道七十三}\newline{}\Emph{讲道七十四}\newline{}\Emph{讲道七十五}\newline{}\Emph{讲道七十六}\newline{}\Emph{讲道七十七}\newline{}\Emph{讲道七十八}\newline{}\Emph{讲道七十九}\newline{}\Emph{讲道八十}\newline{}\Emph{讲道八十一}\newline{}\Emph{讲道八十二}\newline{}\Emph{讲道八十三}\newline{}\Emph{讲道八十四}\newline{}\Emph{讲道八十五}\newline{}\Emph{讲道八十六}\newline{}\Emph{讲道八十七}\newline{}\Emph{讲道八十八}\newline{}\Emph{讲道八十九}\newline{}\Emph{讲道九十}\newline{}\par

\SourceRecord{Translated by George Prevost and revised by M.B. Riddle.；Nicene and Post-Nicene Fathers, First Series；Vol. 10.；Edited by Philip Schaff.；Buffalo, NY: Christian Literature Publishing Co.,；1888.；<http://www.newadvent.org/fathers/2001.htm>.}

% structure: p0001 source=h1 target=section reason=page-title

\section{\WorkTitle{玛窦福音讲道一}}

实在，我们本不该需要成文的圣言，而应展现出如此纯洁的生活，以致圣神的恩宠为我们灵魂代替了书籍，并且正如这些书籍用墨水书写，同样我们的心也该用圣神书写。但是，既然我们已完全抛弃了这恩宠，来吧，至少让我们接受这次好的途径。\par

因为前者更好，天主已藉祂的言语和作为显明了。因为祂向诺厄、亚巴郎、他的后裔、约伯，以及梅瑟说话，不是藉着书写，而是亲自对亲自，因为他们心智纯洁。但全体希伯来人陷于邪恶的深渊之后，那时及此后才有了成文的圣言、法版，以及藉由这些而来的训诫。\par

我们也可以看出，不仅旧约的圣人们如此，新约的也是如此。因为天主并没有给宗徒们任何写下的东西，反而代替成文的圣言，祂应许赐给他们圣神的恩宠：因为我们的主说：\BibleQuote{祂要使你们想起一切}（\BibleRef{若 14:26}）。为使你们知道这要好得多，请听祂藉先知所说的话：\BibleQuote{我要与你们订立新盟约，将我的法律放在他们心中，写在他们的心版上}，并且\BibleQuote{他们都要蒙天主教导}。保禄也指出同样的优越性，说他们领受了法律，\BibleQuote{不是写在石版上，而是写在肉心版上}。\par

但随着时间的推移，有些人关于教义，另一些人关于生活与行为，都遭遇了沉船，于是他们再次需要藉着成文的圣言受到提醒。\par

2. 请反思，对我们而言这是何等重大的恶事：我们本应生活得如此纯洁，以致不需要成文的圣言，而将我们的心如同书籍献给圣神；如今我们失去了那荣耀，并且需要这些，却再次未能适当地运用这第二剂良药。因为如果需要成文的圣言，且未能将圣神的恩宠引到自己身上，已是责备；试想，在这帮助之后仍不肯受益，反倒轻忽所写的，仿佛它是漫无目的、随意抛出的，从而加重对我们的惩罚，这又是多么沉重的罪责。\par

但为避免产生这种后果，让我们严格留意所写的内容；让我们学习旧法律是如何赐下的，新盟约又是如何赐下的。\par

3. 那么，那法律在过去是如何赐下的？何时？何地？是在消灭埃及人之后，在旷野中，在西乃山上，当时烟和火从山中升起，号角吹响，雷霆闪电，梅瑟进入云彩深处。但在新盟约中并非如此——既不在旷野，也不在山上，没有烟、黑暗、云彩和风暴；而是在一天的开端，在一所房屋内，当众人一同坐着时，在极大的宁静中，一切发生了。因为对那些较不合理性、难以引导的人，需要外在的排场，如旷野、山、烟、号角声和其他类似事物；但那些品性更高、顺服、超越了纯肉体想象的人，则无此需要。事实上，新盟约所宣告的一切，是刑罚的除去、罪的赦免，以及\Quote{正义、圣化、救赎}，还有义子名分、天堂的产业、与天主圣子的关系——这都是祂来向众人宣告的：向仇敌、向悖逆者、向坐在黑暗中的人。那么，有什么能与此佳音相比呢？天主在地上，人在天上；一切都混合在一起，天使加入了人类的歌咏团，人类与天使及天上的其他权能共融；人们可以看到长期的战争结束了，天主与我们的本性之间达成了和好，魔鬼蒙羞，恶神逃窜，死亡被消灭，乐园被开启，诅咒被涂抹，罪恶被除去，错误被驱散，真理回归，虔诚的话语到处播种并生长茂盛，天上的政体栽植在地上，那些权能与我们有安全的交往，天使不断在地上显现，关于未来之事有丰厚的希望。\par

因此他称这历史为福音，因为其他一切事物确实只是空言，毫无实质；例如，巨额财富、强大权势、王国、荣耀、尊荣，以及人间被视为美好的一切事物；但那些由渔夫们所传扬的，完全合理且适当地被称为福音：不仅因为它是确定不移的祝福，超越我们的功德，而且因为它极其容易地赐给了我们。\par

因为我们不是藉着劳苦和流汗，不是藉着疲劳和痛苦，而单单因为蒙受天主的爱，我们领受了我们所领受的。\par

5. 那么，为什么当有如此多门徒时，只有两位宗徒写作，另有两位是他们的追随者？（因为一位是保禄的门徒，另一位是伯多禄的门徒，与玛窦和若望一同写了福音。）这是因为他们行事不为虚荣，而全为实用。\par

\Quote{那么怎样？难道一位福音作者不够讲述一切吗？} 一位的确足够；但如果四位写作，不在同一时间，不在同一地方，也未曾聚在一起彼此交谈，却仿佛出于同一口说出一切，这便成为真理的一个非常伟大的证明。\par

6. 或许有人会说：\Quote{然而相反的情形却发生了，因为许多地方他们都被查出有不一致之处。}其实，这一点正是他们真实性的极大证据。因为假如他们在所有事项上，甚至包括时间、地点乃至用词上都完全一致，那我们的仇人都会相信他们是聚在一起，按照某种人的约定写下了所写的内容；因为如此完全的吻合并不出于纯朴。但现在，即使在小事上看来存在的这种不一致，也使他们免除了所有怀疑，并清楚地为作者的性格作了辩护。\par

但若在时间或地点方面，他们有不同的记载，这丝毫不会损害他们所说之事的真理。对于这些事，我们也将在主允许的范围内，随着论述的进展，努力予以指明；同时，我们要求你们除了我们已经提到的事项之外，还要注意：在主要点——那些构成我们生命并为我们提供教义的事上——没有发现任何一处有丝毫分歧。\par

这些点是什么呢？即如下：天主降生成人，祂显了奇迹，祂被钉十字架，祂被埋葬，祂复活了，祂升天，祂将要审判，祂颁布了导向救恩的诫命，祂带来了一条与旧约不相违背的法律，祂是子，祂是独生子，祂是真子，祂与圣父同体，以及所有与此类似的事；因为关于这些，我们将发现它们之间完全一致。\par

若在奇迹中，他们并非全都提到了所有奇迹，而是一个提到这些，另一个提到那些，你们不要为此困扰。因为一方面，若有人讲述了所有奇迹，其余的人就会显得多余；另一方面，若所有人都写了新事且彼此不同，他们一致的证据就不会显明。为此，他们既共同处理了许多事，又各自接受并宣告了一些自己的东西：这样，一方面他不至显得多余，无目的地堆砌；另一方面，他能使我们对其肯定的真理的检验更加完善。\par

7. 现在，路加也告诉我们他着手写作的原因：\Quote{为叫你认清你所学到的道理是确实的}（\BibleRef{路 1:4}），即：通过不断被提醒，你能把握住确实性，并安于确实性。\par

至于若望，他本人对原因保持沉默；然而，按照从最初、从教父们那里传下来的传统，他也不是无缘无故地写作；而是因为那三位（对观福音作者）已致力于详述降生奥迹，而关于天主性的道理几乎被搁置，他受基督感动，于是才着手撰写自己的福音。这一点从历史本身以及他的福音开头便可显明。因为他不像其余福音作者那样从低处开始，而是从高处开始，从他所指向的那个点开始，并且他正是为此目的而撰写了整本书。不仅开头如此，整部福音他都比其余作者更为高超。\par

关于玛窦，据说那些从犹太人中来相信的人曾来到他面前，恳求他将他用口宣讲给他们的同样内容写下来留给他们，于是他也用希伯来语撰写了福音。关于马尔谷，据说在埃及，他应门徒们的请求，也做了同样的事。\par

因此，玛窦既然是为犹太人写作，便只力求显明祂是出于亚巴郎和达味；而路加是对所有人讲话，便将谱系追溯到更高，一直上溯到亚当。前者从祂的诞生开始，因为对犹太人来说，没有什么比被告知基督是亚巴郎和达味的后裔更令人安慰的了；后者并不如此，而是先提及许多其他事，然后再进行谱系。\par

8. 但我们将以接纳了他们陈述的全世界以及真理的仇敌自身来确立他们之间的和谐。因为自他们以来，产生了许多坚持反对他们言论的异端；其中一些异端接受了他们所说的一切，另一些则从其余部分中割下某些说法保留给自己。但若他们的陈述中存在敌意，那么持相反意见的异端就不会接受全部，而只会接受那些看似与自身一致的部分；那些割下一部分的异端也不会被那一部分彻底驳倒；因此，这些碎片也无法隐藏，而是大声宣告它们与整个身体的联系。正如你若从动物身上取下一部分，即使在那部分中，你也会发现构成整体的所有东西——神经、血管、骨骼、动脉、血液，以及可以说是整个团块的一个样本——同样，对于圣经，在其中的每一段陈述里，人们都可以看到与整体的联系清晰显现。然而，若它们不一致，这一点就无法指明，教义本身也早已归于无有，因为祂说：\Quote{凡一国自相纷争，必不能存立。}但如今，即使在这一点上也彰显了圣神的大能，即：祂使这些从事于更为必要和紧急之事的人，在这些小事上丝毫未受损害。\par

现在，关于他们各自写作时居住何处，我们不宜非常肯定地断言。\par

但他们彼此并不对立，这一点我们将在整部著作中努力证明。而你指控他们不一致，无异于坚持要他们使用同样的词句和表达方式。\par

9. 我尚未提及那些以修辞和哲学为荣的人，他们中许多人就同一主题写了许多著作，不仅表达方式不同，甚至彼此矛盾（因为表达不同是一回事，说法相左是另一回事）；这些我都不说。我绝不从那些人的疯狂中为我们的辩护找依据，也不愿借虚假的事来为真理作推荐。\par

但我很乐意查问：这些不同的记述是如何被相信的？如何得胜的？为何在说相反的话时，他们竟被钦佩、被相信、被普世传扬？\par

况且，他们所说之事的见证人众多，对手和敌人也众多。因为他们并非在一隅写下这些东西并埋藏起来，而是在海陆各处，在众人耳中展开；这些事在敌人面前被诵读，正如现今一样，他们所说的一切没有一件触犯任何人。这很自然，因为有一种神圣的力量贯穿一切，使它在众人中兴旺。\par

10. 若非如此，税吏、渔夫和未受教育的人如何能达到这样的哲学高度？那些外人甚至在梦中都未曾想象过的事物，被这些人极其确凿地宣扬并令人信服，不仅在生前，也在死后；不仅对两个人、二十个人、一百人、一千人、一万人，而是对城市、民族、人民，对海陆两地，对希腊人和野蛮人、有人居住和无人居住之地；而且都是关于远超我们本性的事。因为他们抛开大地，所谈论的全是关于天上的事，同时为我们带来另一种生命原则、另一种生活方式：财富与贫穷、自由与奴役、生命与死亡、我们的世界与我们的政体，全都改变了。\par

不像柏拉图，他写了那荒唐的\WorkTitle{共和国}，也不像芝诺，或任何其他写过政体或制定法律的人。因为关于这些人，他们自己已表明，有一个邪恶的灵、一个与我们种族为敌的残酷恶魔、一个与谦逊为敌、与良好秩序为敌、颠覆一切者，在他们灵魂中发声。例如，当他们使妇女成为共有、将裸体少女暴露在体育场中、引入男人目光时，当他们建立秘密婚姻、混淆一切、颠倒自然界限时，还有什么可说？因为这些言论全是魔鬼的发明，违背自然，甚至自然本身也会作证，不容忍我们所提及的事；尽管他们写作时没有迫害、危险或斗争，而是在完全安全和自由中，从多种来源用许多装饰来粉饰。但渔夫们的这些教义，尽管他们被追逐、鞭打、身处险境，却有学问和无学问者、仆人和自由人、国王和普通士兵、野蛮人和希腊人，都满怀善意地接受了。\par

11. 你不能说，因为这些事微不足道、卑下，所以容易被众人接受：不，因为这些教义远比那些更高。因为关于童贞，他们连其名称都从未曾在梦中想象过，更不用说自愿贫穷、禁食或任何其他崇高的事了。\par

但我们的同道不仅根除情欲，不仅惩罚行为，甚至惩罚不洁的目光、侮辱性的言语、放荡的欢笑、衣着、步伐、喧闹，并将严谨推行至最微小的事，使大地充满童贞的植物。关于天主和天上的事，他们劝人用那样的智慧去理解，那是那些学者中任何人在任何时候都未能心中领悟的。因为他们怎能，那些以野兽、地上爬行的怪物、甚至更卑贱之物为神的形象的人？\par

然而，这些高深的教义被人接受并相信，且每日兴旺增长；而其他那些却已消逝、灭亡，比蜘蛛网更容易消失。\par

这很自然，因为那是魔鬼所宣扬的；因此，除了它们的不洁之外，其晦涩也很大，所需的劳苦也更大。因为还有什么比那个\WorkTitle{共和国}更荒唐？在其中，除了我所提到的，哲学家花费无数篇幅才能说明什么是正义，除了这冗长之外，他的言谈还充满许多含糊不清。即使其中包含有益之物，也必然对人生毫无用处。因为如果农夫、铁匠、建筑者、舵手，以及每个靠双手劳动为生的人，都要放下自己的手艺和诚实的劳作，花费这样那样数年时间来学习什么是正义；那么在他学会之前，往往已被饥饿彻底摧毁，因这正义而灭亡，没有学到其他有用之物，以残酷之死结束生命。\par

12. 然而，我们的教训并非如此；基督教导了我们什么是正义、什么是合宜、什么是益处，以及一切德行，概括在简短明了的言语中：有时祂说，\BibleQuote{这两条诫命是法律和先知的基础}（\BibleRef{玛 22:40}），即爱天主和爱近人；有时祂说，\BibleQuote{凡你们愿意人给你们做的，你们也要照样给人做：这就是法律和先知。}（\BibleRef{玛 7:12}）\par

这些事，就连劳工、仆人、寡妇、幼童，以及好似极迟钝的人，都能明了，易于学习。因为真理的教训正是如此，实际的结果也为此作证。至少所有人都学会了当做之事，不仅学会，而且见贤思齐；不仅在城市里、市场中，也在高山之巅。\par

是的，因为在那里你要看见真正的智慧充盈，天使的歌咏团在人的身体里闪耀，\OriginalTerm{天上邦国}{commonwealth of Heaven}在此尘世彰显。因为这些渔夫也为我们书写了一个邦国，没有像那些其他的那样命令人从孩童时期就接受它，也没有立法说德行之人必须年长多少岁，而是将他们的言论针对每一个年龄。因为那些教训是孩童的玩具，这些却是事物的真理。\par

作为这个邦国的居所，他们指定了天堂，并引入了天主作为其缔造者，以及其中设立的法令的立法者；这实在是他们的本分。这个邦国的奖赏不是月桂叶或橄榄枝，不是公宴上的肉食，也不是铜像，这些冰冷平凡的东西，而是无尽的生命，成为天主的子女，加入天使的歌咏团，站在王座旁，永远与基督同在。这个邦国的民魁是税吏、渔夫、制帐幕者，不是那些短暂活着的人，而是现在永远活着的人。因此，即使他们死后，仍然可能对被治理者施以最大的善。\par

这个邦国不是与人作战，而是与魔鬼及那些无形体的权势作战。因此，他们的统帅不是人，也不是天使，而是天主自己。这些战士的铠甲也适合战争的性质，因为它不是由皮革和钢铁制成，而是由真理、正义、信德和一切对智慧的真爱组成。\par

13. 既然上述的邦国是本书所论述的主题，现在我们也要对此发言，让我们专心聆听玛窦，他清楚地谈论此事：因为他的话不是他自己的，而全是基督的，基督制定了这城的法律。我说，让我们专心，好使我们能够被登记于其中，并在那些已成为其公民、正等候不朽冠冕的人中发光。然而，对许多人来说，这番话似乎容易，而先知著作却难。但这又是那些不知道其中所藏思想深度的看法。因此，我恳求你们极其勤勉地跟随我们，以便进入所写之事的汪洋深海，以基督为我们在这次进入时的向导。\par

但是为使话语更容易学习，我们祈祷并恳求你们，如同我们对其他圣经所行的，请事先拿起我们即将解释的那部分圣经，以使你们的阅读为理解铺路（如同太监的情况，\BibleRef{宗 8:28}），从而大大便利我们的工作。\par

14. 这是因为问题又多又频繁。请看，就在他的福音书的开头，有多少困难接踵而至。首先，为什么追溯若瑟的家谱，他并不是基督的父亲？其次，怎样显明祂出自达味，而生祂的玛利亚的祖先却不知道，因为童贞女的谱系没有被追溯？第三，为何追溯若瑟的家谱，他与诞生毫无关系；而对于童贞女，她正是母亲，却没有表明她出自哪些父亲、祖父或祖先？\par

此外，这一点也值得探究：既然他在追溯男人的家谱时也提到了女人，为何他决定这样做后，却没有提及所有女人，而是越过更杰出的，如撒辣、黎贝加以及类似者，只引出了那些以恶事闻名的女人？例如，若有人是娼妓、奸妇、因非法婚姻成为母亲，或是外邦人或蛮夷。因为他提到了乌黎雅的妻子、塔玛尔、辣哈布和卢德，其中一人是异族人，一人是娼妓，一人被近亲玷污，且与他并非婚姻形式，而是通过偷情，她假装成娼妓；至于乌黎雅的妻子，无人不知，因为那罪行的恶名昭彰。然而，这位福音作者越过了所有其余，只在族谱中插入了这些。既然要提女人，就该提所有；若不是全部，而是某些，就该提那些以美德著称的，而非恶行。\par

你看，在第一个开端我们就需要多少留意？然而开端似乎比其余部分更浅显；对许多人来说甚至多余，不过是列举名字而已。\par

此后，另一点又值得探究：为何他省略了三位君王。因为如果他们因极其不虔敬而被隐没其名，那么他也不应提到其他像他们的人。\par

这又是另一个疑问：为何在谈及十四代之后，他在第三组中没有保持这个数目。\par

为何路加提到了其他名字，不仅不全相同，而且更多，而玛窦既少又不同，尽管他也以若瑟结尾，路加也同样以若瑟结束。\par

你看我们多么需要警醒的注意，不仅是为了解释，甚至为了学习我们必须解释的事。因为能够找出困难也不是小事；此外还有另一个难题，即依撒伯尔属肋未支派，如何是玛利亚的亲戚。\par

15. 但为免因堆砌过多而加重你们的记忆负担，我们暂时在此停住我们的讲论。因为要使你们彻底醒悟，只需知道这些疑问就够了。但若你们也渴望其解答，这又取决于你们自己，在我们开口之前。因为若我见你们彻底醒悟，且渴望学习，我将努力添加解答；但若你们打哈欠而不留心，我将隐藏这些难题及其解答，以顺从一项神圣的律法。因为主说：\BibleQuote{不要把圣物给狗，也不要把你们的珍珠投在猪前，免得它们用脚践踏了它们。}\par

但谁是用脚践踏它们的人呢？就是那些不把这些事视为宝贵和可敬的人。也许有人会问，谁如此可怜，竟不把这些事视为可敬，比一切更宝贵？就是那不肯在这事上花时间的，就像在撒殚的剧场中对待娼妓一样。因为在那里，群众整日消磨时光，为这不合时宜的消遣放弃不少家务，并且精确地记住所听到的一切，虽然这对他们的灵魂有害，他们仍然保持。但在这里，当天主发言时，他们却不肯稍作停留。\par

因此，我警告你们，我们与天堂毫无共同之处，我们的公民身份只限于言语。然而正因如此，天主甚至以地狱来威胁，并非为将我们投入其中，而是为说服我们逃避这痛苦的暴政。但我们却反其道而行之，每天奔向通往那里的道路，而且当天主命令我们不仅要听，还要实行祂所说的话时，我们甚至连听都不肯听从。\par

那么，我请问你们，我们何时才会遵行所命令的事，并着手实行，若我们甚至连与它们相关的话语都忍受不了听，却对我们在此逗留的时间（尽管极其短暂）感到不耐烦和不安呢？\par

16. 此外，当我们谈论无关紧要的事时，若见同座的人不注意，我们便称他们的行为为侮辱；但我们是否想到，当天主讲论如此重要的事时，我们若轻视所言，并看向别处，我们是在激怒天主呢？\par

何以一个年迈的人，游历过许多地方，能精确地告诉我们斯达狄亚的数目、城市的位置、规划、港口和市场；但我们自己却连离天上的城有多远都不知道。因为若我们知道距离，我们一定会努力缩短空间。那城不仅像天离地那样远，若我们疏忽，甚至更远得多；相反，若我们尽力，即使在一瞬间，我们也必能到达其城门。因为这些距离不是由地理空间，而是由道德心性来决定的。\par

但你们对世事，无论是新的、旧的，甚至很古老的，都了如指掌；你们能数算出过去所服侍过的王侯、竞赛的裁判、得胜者、军队的领袖——这些与你们毫不相干的事；但在这城中谁成了统治者，第一任、第二任或第三任，各任职多久，每人完成了什么，成就了什么，你们连做梦都未曾想到。而且你们不肯听这城中的律法，甚至当别人告诉你们时，你们也不留心。那么，我请问你们，当你们甚至不听所言时，怎能期望得到所应许的福分呢？\par

17. 但即使以前从未如此，现在至少让我们这样做吧。是的，因为我们正进入一座城（若天主许可），一座由黄金构成，比任何黄金更珍贵的城。\par

让我们留意她的基础，她的城门由蓝宝石和珍珠构成；因为我们确有\Person{玛窦}作为卓越的向导。因为通过他的门我们即将进入，因此需要我们十分勤勉。因为若祂见有人不注意，便将他逐出城外。\par

是的，因为这城极其王权而荣耀；不像我们这里的城，分有市场和宫廷；因为那里一切都是君王的宫廷。因此，让我们打开我们心灵的门，打开我们的耳朵，并怀着极大的战栗，在即将踏上门槛时，朝拜那在内的君王。因为确实，初次的接近立刻有能力使观看者震惊。\par

目前我们见城门紧闭；但当我们见它们敞开时（这便是难题的解答），那时我们将领悟内在光辉的伟大。因为在那里，有一位领你们以神目观看的人，他愿意向你们展示一切，甚至这位税吏；君王坐在何处，祂的扈从谁侍立于旁；天使在哪里，总领天使在哪里；这城中为新公民预备了什么位置，通往那里的路是怎样的，首先成为公民的人获得了什么样的份，以及他们之后的人和随之而来的人。这些部族的秩序有多少，元老院的有多少，尊严的等级有多少。\par

因此，让我们不要带着喧哗或骚动进入，而是带着一种神秘的静默。\par

因为在剧场中，当一片肃静时，君王的诏书被宣读，那么在这城中，众人更必须安静，以心灵和耳朵直立站立。因为即将被宣读的，不是任何地上主人的书信，而是天使之主的书信。\par

若我们如此安排自己，圣神本身的恩宠将带领我们达于圆满，我们将到达君王御座前，并获得一切美善，因我们的主耶稣基督的恩宠和爱人，愿光荣与权能归于祂，偕同圣父及圣神，自今至永远，永世无穷。阿们。\par

\SourceRecord{Translated by George Prevost and revised by M.B. Riddle.；Nicene and Post-Nicene Fathers, First Series；Vol. 10.；Edited by Philip Schaff.；Buffalo, NY: Christian Literature Publishing Co.,；1888.；<http://www.newadvent.org/fathers/200101.htm>.}

% structure: p0001 source=h1 target=section reason=page-title

\section{玛窦福音讲道集 第二篇}

\BibleRef{玛 1:1}\par

\Quote{亚巴郎之子，达味之子耶稣基督的族谱。}\par

你们果真记得我们最近向你们提出的要求，恳请你们怀着完全的静默和神秘的宁静，聆听一切所说的话吗？因为今天我们就要踏入圣洁的门廊，为此我也提醒你们这个要求。\par

如果犹太人，当他们要接近\Quote{一座燃烧的山，有火、黑暗、幽暗和风暴}——或者更确切地说，他们甚至不能接近，只能远远看见和听见这些——就奉命在三天前禁绝与妻子同房，洗自己的衣服，并且他们自己和梅瑟一同战栗恐惧；那么，我们这些将要聆听如此话语的人，不是站在冒烟的山远处，而是要进入天堂本身，岂不是应当表现出更大的克己吗？不是洗我们的衣服，而是擦净我们灵魂的外袍，摆脱所有世俗之物的混杂。因为你们将要看到的不是黑暗，不是烟，也不是风暴，而是那君王自己坐在那不可言喻的荣耀的宝座上，以及天使和总领天使侍立在他身旁，还有圣人一族，伴随着无数不可胜数的天军。\par

因为天主的城正是如此，拥有\Quote{首生者的教会，义人的灵魂，天使的集会，以及所洒的血}，藉此一切合而为一，天上接纳了地上的事物，地上接纳了天上的事物，那由天使和圣人们长久以来所渴望的和平已经来临。\par

在这里矗立着十字架的凯旋，荣耀而显赫，那是基督赢得的战利品，我们本性的初果，我们君王所获的掠物；这一切，我说，我们将从福音中完全知晓。如果你们安静地跟随，我们就能领你们到处周游，向你们展示死亡被钉在何处，罪恶被悬在何处，以及这一场战争、这一场战斗中许多奇妙的贡品又在何处。\par

你们同样会看见那暴君被捆绑在此，一大群俘虏跟在后面，以及那城堡，从那里那邪恶的恶魔在过去曾肆虐一切。你们将看到那强盗的藏身之处和洞穴，如今已被攻破并敞开，因为连那里也有我们的君王亲临。\par

但是，亲爱的，不要厌倦。如果有人描述一场可见的战争、战利品和胜利，你们绝不会感到满足；相反，你们不会为了这叙述而放弃饮食。如果那种叙述受欢迎，这叙述更是如此。因为请想想，这是何等之事：一边是天主从天上的\BibleQuote{王座中跃下}\BibleRef{智 18:15}到地上，甚至到阴间，置身于战阵之中；另一边是魔鬼列阵与他相对——或者更确切地说，不是与揭去面纱的天主，而是与隐藏在人性内的天主相对。\par

奇妙的是，你们将看见死亡被死亡摧毁，诅咒被诅咒消除，魔鬼的统治被他自己得胜的那些事物所推翻。因此，让我们彻底振作起来，不要沉睡，因为看哪，我见门已向我们敞开；但让我们以一切合宜的秩序和战栗进入，立刻踏上这门廊本身。\par

2. 这门廊是什么？\Quote{亚巴郎之子，达味之子耶稣基督的族谱。}\par

\Quote{你说什么？你不是答应要讲论天主的独生子吗？怎么却提及一个出生在一千代之后的达味，并说他既是父又是祖先？}且慢，不要想一下子学会一切，要温和地、一点一点地来。怎么，你正站在门廊里，就在门口；为什么要急着往内殿去？你还没有好好看清外面的一切。因为我暂时并不向你宣布那另一种生；或者更确切地说，连这后来的生也不宣布，因为它是无法言说、不可表述的。在我之前，先知依撒意亚已经告诉了你这事；当他宣告他的受难和他对世界的极大关怀，惊叹他是谁，他成了什么，他降到了何处时，他大声清楚地喊说：\Quote{谁能述说他的世代？}\par

我们现在不是要讲那一种生，而是讲这一种，即在地上发生的、有成千上万见证的。关于这一种，我们也要在蒙受圣神恩宠的情况下，尽可能地讲述。因为连这一种，也没有人能完全清楚地阐述，因为它也是最令人敬畏的。所以，当你们听到这诞生时，不要以为你们听到的是小事，而要振奋心灵，立刻战栗，因为你们被告知天主降临到了地上。这事如此奇妙、如此超乎意料，以致天使们为此歌唱，为世界献上他们的赞美，先知们从一开始就对此感到惊异，因为\BibleQuote{他出现在地上，与人同在} \BibleRef{巴 3:37}。是的，因为听到那位不可言说、不可名状、不可理解、与父平等的天主，竟通过了童贞女的子宫，屈尊由女人而生，以亚巴郎和达味为祖先，这远超一切思想。但我为什么说亚巴郎和达味呢？因为更令人惊叹的是，还有那些我们最近提到的妇女。\par

3. 听到这些，站起来，不要想任何卑微的事；反而正因为这一点，你们更应当惊叹：他是无始天主的圣子，他的真子，却允许自己也被称为达味之子，为的是使你成为天主的儿子。他让一个奴隶作他的父亲，为的是使主成为你这个奴隶的父亲。\par

你岂非从开头就看见福音属于何种性质？倘若你对自己有关的事心存疑虑，那么从他所属的事上也当相信这些。因为按人的理性判断，天主降生成人，远比人被宣告为天主子更难。因此，当你听说天主子是达味之子和亚巴郎之子时，不要再怀疑你——亚当之子——也将成为天主子。他并非随意或徒然地如此自谦，只愿高举我们。他按肉身诞生，为叫你按圣神而生；他由女人而生，为叫你不再作女人的儿子。\par

因此，这诞生具有双重性：既肖似我们，又超越我们。由女人而生，诚然是我们的命运；但\BibleQuote{不被血统、也不被肉身的意愿、也不被人所生，而是被圣神所生}\BibleRef{若 1:13}，则是预告那超越我们的诞生，即他将要藉圣神白白赐予我们的将来诞生。其余一切也是如此。他的圣洗也是同一方式：它分有旧，也分有新。由先知受洗，标记旧；圣神降临，预兆新。如同有人站在两个相距的人中间，伸出双手抓住两边，将他们联结起来；他也如此行，将旧盟约与新盟约、天主的本性与人的本性、他的事物与我们的事物结合起来。\par

你看见这城的闪烁光华了吗？它以何等大的辉煌从起初就眩惑了你？它如何立即以你的形象显明君王？如同在军营中：君王并不常以本有威仪出现，而是脱下紫袍与冠冕，常扮作普通士兵。但在那里，为免被认出而招致敌人；在此却相反，为免他若被认出，就使敌人避战，并使所有他自己的子民惊惶：因为他的目的是拯救，而非恐吓。\par

4. 因此，他也立即用这名号称呼他，命名为耶稣。因为耶稣这名并非希腊语，而在希伯来语中如此称为耶稣；译成希腊语，即\Quote{救主}。他之所以被称为救主，是由于他拯救他的子民。\par

你看见他如何给了听众翅膀，同时说着熟悉的事，又借此指示我们超越一切希望的事吗？我的意思是，这两个名号犹太人都很熟悉。因为所发生的事超出期待，名号的预象先行出现，为使从起初就除去一切新奇的不安。那在梅瑟之后带领百姓进入预许之地的，被称为耶稣。你看见预象吗？请看真理。那一位领进预许之地，这一位领进天堂及天上的美善；那一位在梅瑟死后，这一位在法律终止之后；那一位作领袖，这一位作君王。\par

然而，免得你因听见耶稣这个字，因同名而困惑，他补充说：\Quote{耶稣基督，达味之子}。但另一位并非出于达味，而是出于另一支派。\par

5. 但他为何称此书为\Quote{耶稣基督的世代之书}，而此书不仅包含诞生，还包含整个救恩计划？因为这是整个救恩计划的总纲，并成为我们一切福分的根源。正如梅瑟称书为天地之书\BibleRef{创 2:4}，虽他不仅论及天地，也论及其中万物；同样，此人以其所成就一切伟大之事的总和命名其书。因为令人惊叹、超乎希望与期待的事，就是天主成为人。此事既成，一切后来的事便顺理成章。\par

6. 但他为何不说\Quote{亚巴郎之子}，然后\Quote{达味之子}？这不是像有些人以为的，他要从较低点向上追溯（因为那样他会与路加一样，但现在他相反）。那么为何提及达味？此人因他的卓越和时代而在众人口中：他去世没有很久，不像亚巴郎。虽然天主对两者都作了许诺，但那一位因古老而默默无闻，达味却因新鲜和近时而被众人传诵。例如，他们说：\BibleQuote{基督岂非出自达味的后裔，出自白冷，就是达味所住的村镇吗？}\BibleRef{若 7:42}。无人称他为亚巴郎之子，众人皆称达味之子；因为达味更常被众人记起，既因时代（如我已说），也因他的王权。同理，此后所有他们尊重的君王，都从达味得名，无论百姓还是天主都如此称呼。因为厄则克耳及其他先知论到达味时，说他要来临、要复活；不是指已死的达味，而是指那些效法他德行的人。至于希则克雅，他说：\BibleQuote{我必要护卫这城，为我自己和我仆人达味的缘故。}\BibleRef{列下 19:34}。至于撒罗满，他也说，因达味的缘故，他在撒罗满有生之日不将王国夺回。因为此人在天主和人面前都大有荣耀。\par

为此，他立即从更知名者开始，然后上溯到他的父亲；认为对犹太人而言，将族谱推得更远是多余的。因为这些人尤其受钦佩：一位是先知和君王，另一位是族长和先知。\par

7. \Quote{然而，从哪里可以证明祂是达味的后裔？}有人可能会说。因为如果祂不是从男人所生，而是仅从女人所生，并且童贞女的族谱没有被追溯，我们如何知道祂是达味的后裔？因此，有两个问题被追问：一是为什么祂母亲的族谱没有被记述，二是为什么他们提到了若瑟，而他在生育上并无分涉？因为后者似乎是多余的，而前者则是一种缺陷。\par

那么，必须先谈哪一个？童贞女如何是达味的后裔。我们如何知道她是达味的后裔？请听天主吩咐加俾额尔去\Quote{一位童贞女，已许配给一个名叫若瑟的男子，是达味家室的人。}当你听到童贞女是达味家室的人时，还有什么比这更清楚的呢？\par

由此可知，若瑟也属于同一家族。是的，因为有一条法律，规定不可从其他支派娶妻，而只能从同一支派。先祖雅各伯也曾预言祂将出自犹大支派，如此说：\Quote{权杖不离犹大，柄杖不离他腿间，直到那应得的人来到，祂是外邦人所期待的。}\BibleRef{创 49:10}\par

\Quote{好了，这个预言确实表明祂出自犹大支派，但并没有表明祂出自达味的家族。那么，在犹大支派中难道只有达味这一个家族吗？不也有很多其他的吗？一个人可能属于犹大支派，却不属于达味的家族，难道不是这样吗？}\par

不，免得你这样说，福音作者已经消除了你的疑虑，他说祂是\Quote{达味家室的人。}\par

如果你还希望从另一个理由了解这一点，我们也不缺乏另一个证据。因为不仅在法律上不允许从其他支派娶妻，甚至也不允许从其他家族，即从其他亲属关系娶妻。所以，如果我们将\Quote{达味家室的人}这句话与童贞女联系起来，那么所说的话就成立；或者如果与若瑟联系起来，那么由此也证明了这一点。因为如果若瑟是达味家室的人，他就不会从不同于自己出身的家族娶妻。\par

\Quote{那么，}有人可能会说，\Quote{如果他违反了法律呢？}为此，作者预先作证若瑟是义人，正是为了让你不这样说；相反，得知他的德行后，你也会确信他不会违反法律。因为他是如此仁慈、无私，甚至不愿在猜疑之下对童贞女施加惩罚，他又怎会因私欲而违反法律呢？那位表现出超越法律的智慧和自制的（因为暗中休退她，是超越法律的自制行为），他又怎会做出任何违反法律的事呢？更何况没有任何原因驱使他这样做。\par

8. 童贞女出自达味家族，从这些事已清楚可见；但为什么他没有给出她的族谱，而是给出了若瑟的，这需要解释。那么，原因是什么呢？在犹太人当中，并没有追溯女性族谱的惯例。因此，为了遵守习俗，并避免一开始就显得革新，同时又让我们认识童贞女，他保持了沉默而不提她的祖先，却追溯了若瑟的族谱。因为如果他关于童贞女这样做了，就会显得引入新奇；而如果他对若瑟保持沉默，我们就不会知道童贞女的祖先。因此，为了让我们学习关于玛利亚的事，她是谁，出于什么家族，同时法律不受干扰，他追溯了她未婚夫的族谱，并表明他是达味家室的人。因为当这一点被清楚证明时，另一事实也随之显明，即童贞女也由此而出，因为这位义人，正如我已经说过的，不会容忍娶另一个家族的女子为妻。\par

还有另一个原因，一个更具奥义的原因，可以提及，因为童贞女的祖先被保持了沉默；但此刻不宜宣布，因为已经说了这么多。\par

9. 因此，让我们就此打住对问题的讨论，同时准确地记住已经启示给我们的事：例如，为什么他首先提到达味；为什么他称这本书为\Quote{族谱}；他说\Quote{耶稣基督}的原因是什么；生育如何是共同的，又不是共同的；玛利亚如何被证明是达味的后裔；以及为什么追溯若瑟的族谱，而她的祖先却被保持沉默。\par

因为如果你们记住这些，就会更加鼓励我们继续未来的事；但如果你们拒绝并将它们从心中丢弃，我们就会对余下的事更加迟疑。就如同没有农夫会愿意关心一块毁掉了先前种子的土地。\par

因此，我恳求你们反复思考这些事。因为思考这些事，灵魂中会生发出某种巨大的善，指向救恩。因为通过这些默想，我们将能够取悦天主自己；我们的口将免于侮辱、污言秽语和辱骂，而当它们操练在属灵的言语中时；我们将使魔鬼感到可怕，因为用这样的话武装我们的舌头；我们将更加吸引天主的恩宠到自己身上，它将使我们的眼睛更加锐利。因为确实，眼睛、口和耳朵都是为此目的而置于我们身上的，使我们的所有肢体都能事奉祂，让我们说祂的话，行祂的事，向祂不断歌唱赞歌，献上感恩的祭献，并借此彻底净化我们的良心。\par

因为正如身体在有洁净空气时会更加健康，同样，当灵魂在这样的操练中滋养时，也会更富有实践智慧。你们没有看到身体的眼目吗？当它们停留在烟雾中时，总是流泪；但当它们处于清澈的空气中、在草地上、在泉源和花园中时，它们变得目光敏锐且更加健康？灵魂的眼目也如此：如果它在属灵神谕的草地上喂养，就会清澈、敏锐、目光明锐；但如果它进入今生之事的烟雾中，就会无止境地哭泣，现今和将来都哀号。今生之事确实如同烟雾。因此有人曾说：\Quote{我的岁月如烟消逝。}他固然是指它们的短暂和虚幻，但我要说，我们不仅应从这意义上理解所说的话，也当从其浑浊的性质来理解。\par

因为没有什么比世俗的忧虑和欲望的丛集更能伤害和蒙蔽灵魂的眼目了。这烟雾正是由这些柴薪所喂养。正如火，当它沾上潮湿而浸透的燃料时，会生出许多烟雾；同样，这强烈燃烧的情欲，当它沾上一个（可以这样说）潮湿而放荡的灵魂时，也会在其方式上产生大量的烟雾。为此，我们需要圣神的甘露和那空气，以熄灭火焰、驱散烟雾，并给我们的思想以翅膀。因为一个被如此重大邪恶所压重的人，不可能，绝不可能翱翔到天上；即使毫无阻碍，我们能开出一条路来，已是万幸；或者更确切说，即使如此也不可能，除非我们获得圣神的翅膀。\par

既然要达到那样高度，既需要不受阻碍的头脑，也需要属灵的恩宠；那么，如果这些都没有，而我们反而吸引与它们相反的一切，即一种\OriginalTerm{撒殚般的重负}{satanical weight}，那又怎样呢？我们如何能被如此重负拖住还能向上翱翔呢？的确，如果有人试图用精准的天平来称量我们的话语；在一万塔兰特世俗的言谈中，他几乎找不到一百个德纳的属灵言语，或者更确切说，我该说，连十个法寻都没有。那么，这难道不是一种耻辱和极端的嘲弄吗？如果我们有一个仆人，我们大多在必要的事上使用他；但拥有舌头，我们却没有像对待奴隶那样善用我们的器官，反而用它来作不益的事和仅仅无足轻重的事？但愿只是无足轻重的事！但现在却是用于相反、有害且对我们毫无益处的事。因为我们所说的话如果对我们有益，肯定也会中悦天主。但事实上，无论魔鬼暗示什么，我们都说了出来：时而大笑，时而说俏皮话；时而诅咒辱骂，时而发誓、说谎、发假誓；时而抱怨，时而说虚妄的闲话，谈论比老妇更无聊的琐事；吐出一切与我们无关的话。\par

因为，请告诉我，你们中间站着的人，如果被要求，谁能背诵一篇圣咏或圣经的任何其他部分？没有一个人。\par

这不仅令人悲痛，而且你们对属灵的事如此退步，但在属于撒殚的事上，却比火更猛烈。例如，如果有人想要问你们魔鬼的歌曲和不洁的柔弱旋律，他会发现许多人完全知晓，且以极大的乐趣重复它们。\par

10. 但对于这些指控，有什么回答呢？你们会说：\Quote{我不是隐修士，我有妻子儿女和家庭的挂虑。}这恰恰是败坏一切的原因：你们以为阅读圣经只属于那些人，而你们比他们更需要它。因为生活在世俗中、每日受创伤的人，最需要药物。所以，不读经固然很糟，但认为这件事\Quote{多余}则更糟：因为这是魔鬼的诡计。你们没有听见保禄说：\Quote{这一切都写下来是为劝诫我们}吗？\BibleRef{格前 10:11}\par

而且，如果你们要拿起一本福音，不会不洗手就拿；但在其中所存的内容，你们难道不认为极其必要吗？正因为此，一切都颠倒了。\par

因为如果你们愿意学习圣经有多大益处，请审视自己：你们听圣咏时变成了什么，听撒殚的歌曲时又变成了什么；你们留在圣堂时的心态，和坐在剧场时的心态如何；你们会发现这灵魂与那灵魂之间有着巨大差别，尽管两者是同一个。因此保禄说：\Quote{坏同伴会败坏好习惯。}\BibleRef{格前 15:33}为此我们需要那些\OriginalTerm{来自圣神的护符}{charms from the Spirit}不断吟唱。是的，这正是我们优于无灵之物的原因，因为在一切其他方面，我们甚至远不及它们。\par

这是灵魂的食物，这是它的装饰，这是它的保障；而不听则是饥荒和消耗，因为他说：\Quote{我要给他们，不是食物的饥荒，也不是口渴的干渴，而是聆听上主言语的饥荒。}\BibleRef{亚 8:11}\par

还有什么比这更可悲的呢？天主以惩罚方式威胁的灾祸，你们却主动引到自己头上，给灵魂带来一种可悲的饥荒，使它成为世上最软弱的东西？因为它的本性既因言语而消耗，也因言语而得救。是的，言语使它走向愤怒；同样的话又使它温良：肮脏的言辞常会燃起它的情欲，而庄重的言语则使它接受节制。\par

但如果单单言语就有如此大的力量，告诉我，你们为何藐视圣经？如果劝诫能成就如此大事，何况带着圣神的劝诫呢？是的，因为从圣经而来的言语，在耳中回响，比火更能软化刚硬的心灵，使之适合一切美善之事。\par

11. 同样，\Person{保禄}见\Place{格林多人}骄傲自大、意气用事，便使他们平静下来，变得更加体贴。因为他们正以那些本应感到羞耻、掩面不及的事而自夸。但他们收到书信之后，请听他们有何改变——师傅亲自为他们作证说：\Quote{正是这事，你们依着天主的意思忧愁，在你们身上产生了何等殷勤，甚至何等辩白，何等愤慨，何等热忱，何等惩罚。}同样，我们也如此管教仆人和儿女、妻子和朋友，并使我们的仇敌成为朋友。\par

同样，那些伟大人物，即那些天主所爱的人，也变得更好。例如\Person{达味}，在他犯罪之后，因着某些话语的裨益，便达成了那最卓越的悔改；宗徒们也是如此，藉此方法成了他们所成的，并吸引了整个世界跟随他们。\par

\Quote{那有什么益处呢？}有人可能会说，\Quote{当一个人听了却不照所说的去做时？}即使仅是聆听，也将获益匪浅。因为他会继续自责，内心叹息，并终将去做那些所说的事。但那连自己犯了罪也不知道的人，何时才会停止怠惰？何时才会自责？\par

因此，我们不可轻视聆听圣经。因为这是\Person{撒殚}的诡计；他不让我们看见宝藏，免得我们获得财富。因此，他说聆听神圣的法律算不了什么，免得他看见我们从聆听中也能付诸实践。\par

我们既知他的这种邪恶伎俩，就当四面加强防御，好让我们以这铠甲护身，自己既能保持不被征服，又能击打他的头；如此，我们以光荣的胜利冠冕装饰自己，就能获得未来的福乐，藉着我们的主耶稣基督的恩宠和慈爱，愿光荣与权能归于他，直到永远。阿们。\par

\SourceRecord{Translated by George Prevost and revised by M.B. Riddle.；Nicene and Post-Nicene Fathers, First Series；Vol. 10.；Edited by Philip Schaff.；Buffalo, NY: Christian Literature Publishing Co.,；1888.；<http://www.newadvent.org/fathers/200102.htm>.}

% structure: p0001 source=h1 target=section reason=page-title

\section{玛窦福音第三篇讲道}

\Emph{\BibleRef{玛 1:1}}\par

\BibleQuote{耶稣基督，达味之子，亚巴郎之子的族谱。}\par

看，第三篇讲道，我们还没有结束前言。所以我说这些思想具有很大的深度，并非徒然。\par

来，让我们今天讲剩下的。那么现在需要什么？为何要追溯若瑟的族谱，而他在诞生中并无分？我们已经提到一个原因，但还需要提到另一个，比第一个更奥秘、更隐秘。那么这是什么？他不愿在诞生时让犹太人知道基督是由童贞女所生。\par

不，不要为这话的奇异而困扰。因为这不是我的主张，而是我们那些杰出而显赫的教父们的主张。因为如果他从起初隐藏了许多事，称自己为人子，甚至没有处处清楚地显示他与圣父的平等；那么，他暂时也隐藏了这件事，安排了一个伟大而奇妙的目的，你又为何惊奇呢？而且，他们会以奸淫罪定她的罪。因为，对于其他事——他们在旧约时代也经常有这样的先例——他们厚颜无耻地固执己见（例如，因为他驱魔，他们说他附魔；因为他在安息日治病，他们认为他是天主的敌人；然而，在以前安息日也多次被打破），如果这件事告诉他们，他们还会说什么呢？尤其是，在这之前所有时间都站在他们一边，因为从未发生过这样的事。因为，如果经过了这么多奇迹，他们仍称他为若瑟的儿子，那么在奇迹之前，他们怎么会相信他是由童贞女所生呢？\par

正是为此，若瑟的族谱被追溯，童贞女许配给他。因为，即使他——一个正义而奇妙的人——也需要许多事（天使、梦中的异象、先知的见证）才能接受所发生的事；那么，犹太人既愚钝又堕落，对他怀有如此敌意，又怎能将这种想法纳入心中呢？因为这件事的奇异和新鲜必定会极大地困扰他们，而且在他们祖先的时代，从未听说过这样的事。因为，一旦有人相信他是天主子，之后他就不会再怀疑这一点；但那些认为他是骗子、是天主的敌人的人，又怎能不被这事更冒犯，而被引向相反的想法呢？因此，宗徒们最初并没有直接说出这一点，而是经常谈论他的复活（因为在前时代有这方面的先例，尽管不似如此）；他们并不总是说他是童贞女所生；甚至他的母亲本人也不敢说出这话。请看，童贞女本人也对他说了什么： \BibleQuote{看，你的父亲和我一直在找你。} \BibleRef{路 2:48} 因为，如果这种怀疑被接受，他就不会再被视为达味之子，而这种意见不被坚持，还会产生许多其他祸患。因此，天使也没有向所有人说这些事，而只告诉了玛利亚和若瑟；但当向牧羊人宣告所发生之事的喜讯时，他们不再添加这一点。\par

2. 但为什么，在提到亚巴郎并说\Quote{他生了以撒，以撒生了雅各伯}之后，却没有提到他的兄弟；当来到雅各伯时，他却同时提到了\Quote{犹大和他的兄弟们}？现在有些人说，这是因为厄撒乌和之前其他兄弟的悖逆。但我不这么说；因为如果是这样，他稍后为何提到这样的妇女？恰恰相反，在此处，他的荣耀正是通过对比显明；不是因为有伟大的祖先，而是卑微和微不足道的。因为对崇高者而言，能够极度贬抑自己是极大的荣耀。那么，他为何没有提到他们？因为撒拉逊人、依市玛耳人、阿拉伯人以及所有从那些祖先衍生的人，与以色列民族毫无共同之处。因此，他跳过那些，急忙转向他的祖先，即犹太民族的祖先。所以他说：\BibleQuote{雅各伯生了犹大和他的兄弟们。} 因为从此时起，犹太民族开始有了其独特标记。\par

3. \BibleQuote{犹大由塔玛尔生了培勒兹和则辣黑。} \BibleRef{玛 1:3} \Quote{你做什么，人啊，让我们想起一段包含非法结合的历史？} 但为什么这样说？因为，如果我们是在讲述一个普通人的族谱，自然会对这些事保持沉默；但如果是在讲述降生成人的天主，非但不能沉默，反而应当以此夸耀，显示他的温柔眷顾和他的大能。是的，他正是为此而来，不是要逃避我们的耻辱，而是要将它们带走。因此，正如他更受赞叹，不仅因为他死了，而且甚至被钉十字架（尽管这是可耻的，但越可耻，越显示他对人的爱），同样地，我们也可以谈论他的诞生；我们惊异于他，不仅因为他取了肉体成为人，而且因为他屈尊有了这样的亲属，丝毫不以我们的恶为耻。他从诞生的最初就宣布，他不以任何属于我们的事物为耻；同时也教导我们，永远不要因祖先的邪恶而掩面，而要只追求一件事，即德行。因为这样的人，即使他的祖先是一个外邦人，即使他的母亲是个娼妓或任何什么样的人，也不会因此受到伤害。因为，如果连那淫乱者本人，一旦改变，也不因过去的生活而蒙羞，那么，从娼妓或奸妇所生的人，只要他有德行，祖先的邪恶更不能使他蒙羞。\par

然而祂这样做，不仅是为了教导我们，也是为了挫败犹太人的傲慢。因为他们忽视自己灵魂中的德性，却炫耀\Person{亚巴郎}的名号，以为先祖的功德可以为他们辩护；祂从一开始就表明，人不应在这些事上夸耀，而应夸耀自己的善行。\par

除此之外，祂也确立另一点：表明众人都处于罪中，连他们的祖先也不例外。至少他们的族长和同名人被证明犯了大罪，因为\Person{塔玛尔}站在他对面，控告他的淫乱。\Person{达味}也因他所玷污的妻子生下了\Person{撒罗满}。如果连伟人都没有遵行法律，那么更渺小的人更无法遵行。既然法律未被遵行，众人都有罪，基督的来临便成为必要。\par

为此，祂也提及了十二位族长，以此再次贬抑他们对先祖高贵出身的骄傲。因为他们中许多人也是由女奴所生；然而，父母的不同并未造成子女的差异。因为他们全都同样地是族长和支派的首领。因为这就是教会的优越性，这就是我们中间那从起初就预表的高贵特权。因此，无论你是为奴的还是自由的，你从那里得到的不多也不少；唯一的衡量标准就是心灵和灵魂的倾向。\par

4. 然而，除了我们所说的，还有另一个原因，祂为何也提及了这个历史；因为，当然，\Person{则辣黑}的名字并非偶然出现在\Person{培勒兹}的记载中。（因为当祂提到\Person{培勒兹}——基督的族谱由此追溯——之后，再提及\Person{则辣黑}确实是多余且无关的。）那么为何祂要提及他呢？当\Person{塔玛尔}即将生产时，分娩的阵痛临到她，\Person{则辣黑}首先伸出了手。\BibleRef{创 38:27} 接生婆看见后，为了认出谁是长子，用朱红线绑了他的手；但婴儿被绑住后，将手缩了回去；收回去后，\Person{培勒兹}先出来，然后才是\Person{则辣黑}。接生婆看见后说：\BibleQuote{你为什么为你破开了篱笆？}\par

你看到奥秘隐晦的表达了吗？因为这些事并非无意义地记录给我们；因为我们不值得费心去学习接生婆说了什么，也不值得叙述得知：那个后出来的孩子先伸出了手。那么，这奥秘的教训是什么？首先，从孩子的名字我们得知所探寻的事，因为\Person{培勒兹}的意思是\OriginalTerm{分裂}{division}和\OriginalTerm{破口}{breach}。而且，从所发生的事本身来看；因为伸出了手又把它缩回去，并非出于自然规律；这些事既不属于理性的运动，也不是按自然顺序发生的。因为手伸出后，另一孩子先出来或许并非不自然；但收回手让另一个先通过，这就不符合出生婴儿的方式了，而是天主的恩宠与这两个孩子同在，安排这些事，并藉此为我们描绘出未来之事的影像。\par

那么这样？有些仔细研究这些事的人说，这两个孩子是两族人民的预像。所以为了让你知道，后一族人的政体在先一族出现之前就已照亮，那伸手的孩子并未完全显现，而是将手缩回；在他的兄弟完全滑出后，他才完全显现。这发生在两族人身上也是如此。我的意思是，在\Person{亚巴郎}时代，教会的政体已经显现，然后在其进程中隐退，随后犹太民族和法律的政体出现，然后新的人民带着自己的法律完全显现。因此接生婆说：\BibleQuote{你为什么为你破开了篱笆？} 因为法律的进入，打破了政体的自由。的确，圣经常将法律称为篱笆；正如先知所说：\BibleQuote{你拆毁了它的篱笆，使所有过路的人都摘取它的葡萄。} 以及：\BibleQuote{我给它筑了篱笆。} 还有\Person{保禄}：\BibleQuote{拆毁了中间隔断的篱笆。} 但也有人说，\BibleQuote{你为什么为你破开了篱笆？} 这句话是针对新的人民说的：因为它的到来废除了法律。\par

5. 你看到吗？祂并非为了轻微或微小的原因而让我们记起关于\Person{犹大}的全部历史。为此，祂也提及了\Person{卢德}和\Person{辣哈布}，一个是外邦人，一个是妓女，为让你知道祂来是要消除我们的一切恶疾。因为祂来是作医生，而非审判官。因此，正如古时的人娶妓女为妻，同样天主也娶了那行淫的本性为妻；从起初先知们也宣布这在会堂身上发生了。但那配偶对那位作她丈夫的忘恩负义，而教会，一旦从祖先那里继承的恶中获得解放，便继续拥抱新郎。\par

试看\Person{卢德}的事，它与我们所经历的事何其相似。她出身异族，又极其贫穷，但当\Person{波阿次}看见她时，既不因她的贫穷而轻视她，也不因她的卑微出身而厌恶她；同样，\Person{基督}接纳了教会——教会本是外邦人，又极其贫穷——使她共享伟大的恩赐。然而正如\Person{卢德}，若她先前没有离开她的父亲，没有弃绝家庭、种族、国家和亲属，就不能获得这种联姻；同样，教会也舍弃了从祖先领受的习俗，然后（且仅在此之后）才成为新郎所爱慕的。为此先知对她说：\Quote{忘记你的民族，和你的父家，君王就会喜爱你的美貌。}\Person{卢德}也这样做了，因此她成为君王的母亲，教会也是如此。因为\Person{达味}本人就是出自她。所以，为了借此一切使他们感到羞愧，并说服他们不可骄傲，他编排了族谱，并引出了这些妇女。是的，因为后者，藉着中间的人，成了伟大君王的祖先，而\Person{达味}并不以此为耻。因为一个人不可能因其祖先的德行或恶行而成为好或坏、卑微或光荣；即使说些看似矛盾的话，那并非出自高贵祖先却成为卓越的人，反而更加闪耀。\par

6. 因此，谁也不可因这些事而骄傲，而应思想主的祖先，摒弃一切傲慢；要以善行为荣，或者更确切地说，连善行也不可自夸。因为正是这样，\Person{法利塞人}成了低于税吏的人。因此，你若想显示善行的伟大，就不要存高傲之心，这样你便证明善行更加伟大。你当自认一无所行，这样你就行了一切。因为如果我们这些罪人，当我们自认自己所是时，便成为义人，正如税吏那样；那么，当我们是义人时，若自认是罪人，岂不更是如此？既然罪人因谦卑之心（虽然这不算是谦卑，而是正直之心）而成义；那么，如果正直之心在罪人身上有如此大的效力，试想谦卑之心对义人又当如何呢？\par

不要糟蹋你的劳苦，也不要丢弃你辛劳的果实，更不要徒然奔跑，使你在许多赛程之后的全部努力化为乌有。因为你的主比你更了解你的善行。即使你只给一杯凉水，他也不会忽视；哪怕你只献出一文钱，或只是叹息一声，他都欣然接纳，记念在心，并为此赐予丰厚的赏报。\par

但你为何要查究自己的行为，并不断在我们面前炫耀呢？难道你不知道，你若赞扬自己，天主就不再赞扬你了吗？正如你若哀叹\Emph{自己}，他必会不停地在众人面前称赞你。因为他决不愿你的辛劳受到贬低。何止是不愿贬低？他反倒行事安排一切，为要因微小的事而给你加冕；他四处寻找借口，好使你免于地狱。为此，即使你在第十一时辰工作，他也给你全天的工资；即使你毫无得救的理由，他也说：\Quote{我为了我自己的缘故而行这事，免得我的圣名被亵渎：}\BibleRef{则 36:22}即使你只叹息一声，只哭泣一次，这一切他都迅速抓住，作为拯救你的契机。\par

因此，我们不可自高自大，而要自认无用，这样我们才能成为有用的。因为如果你自称可嘉，你就成了无用的，即使你确实可嘉；但如果你自称无用，你就成了有用的，即使你原是弃绝的。\par

7. 所以必须忘记我们的善行。\Quote{然而怎么可能呢？}有人会说，\Quote{不知道自己熟悉的事？}你怎么说？你不断地得罪你的主，却生活安逸，嬉笑度日，甚至不知道自己犯了罪，把一切遗忘；而对于你的善行，难道就不能忘却吗？然而敬畏是一种更强有力的事。可是我们恰恰相反：一方面，我们每天都在犯罪，却从不放在心上；另一方面，如果我们给了穷人一点钱，就时时惦记。这种行为完全是疯狂，对这样算计的人来说是极大的损失。因为善行的安全仓库就是忘记善行。如同衣物和金子，若把它们放在市场上展示，就会招来许多心怀恶意的人；但若把它们收藏在家里隐藏起来，便能安全存放。同样，对于我们的善行，若我们不断记念，就会激怒主，武装仇敌，让他来偷窃；但如果除了那位唯一应当知道的主之外，无人知晓，它们就会安全存留。\par

因此，不要永远炫耀它们，免得有人偷走它们。正如\Person{法利塞人}所遭遇的，因为他将善行挂在唇边，魔鬼便夺去了它们。然而他是以感恩之心提及它们，并将一切归于天主。但即便如此也不够。因为感恩不是侮辱别人，不是在许多面前炫耀，不是自高自大、蔑视有过失的人。你若是感谢天主，就当只满足于他，不要向人宣扬，也不要谴责你的邻人；因为这不是感恩。你愿学习感恩的话语吗？请听\Person{三青年}的话：\Quote{我们犯了罪，我们违了法。上主，你对我们所做的一切都是公义的，因为你用真实的判断使这一切临于我们。}因为承认自己的罪，就是向天主献上感恩的忏悔；这种事意味着人虽犯下无数过犯，却没有受到应有的惩罚。这样的人才是最感恩的。\par

8. 让我们谨慎的是，不要谈论自己，因为这会使我们在人面前可憎，在天主面前可恶。因此，我们行的善事越大，就越少谈论自己；这是在天主和人面前获得最大荣耀的方式。或者更确切地说，不仅是来自天主的荣耀，还有赏报，是的，巨大的酬报。所以不要要求赏报，以便你得到赏报。要承认自己是因恩宠得救，好使他承认自己是你的债务人；不仅因你的善行，也因你这种正直的心态。因为当我们行善时，我们只凭善行使他成为我们的债务人；但当我们甚至不认为自己行了任何善事时，那么也因这种心态本身而使他成为我们的债务人；而且因这种心态比因其他事更多：所以这等同于我们的善行。因为如果缺少这种心态，善行也不会显得伟大。同样，我们也是如此，当我们有仆人时（\BibleRef{路 17:10}），当他们善意地完成了所有服务，却不认为自己做成了什么大事时，我们最认可他们。因此，你若要使你的善行伟大，就不要认为它们伟大，这样它们才会伟大。\par

百夫长也如此说：\Quote{我不配你进到我的屋顶下}；因此他成为配得的，并且被\Quote{感到惊奇}\BibleRef{玛 8:8}超过所有犹太人。同样，保罗说：\Quote{我不配称为宗徒}\BibleRef{格前 15:9}；因此他竟成了第一。同样，若翰：\Quote{我不配解开祂的鞋带}；因此他成了\Quote{新郎的朋友}，并且他承认不配触摸祂鞋子的那只手，基督却把它按到自己头上。同样，伯多禄说：\Quote{离开我，因为我是个罪人}\BibleRef{路 5:8}；因此他成了教会的基石。\par

因为没有什么比自居末位更蒙天主悦纳。这是一切实用智慧的首要原则。因为谦卑且心灵破碎的人，不会贪图虚荣，不会发怒，不会嫉妒邻人，不会怀有任何其他激情。因为即使一只手受了伤，纵然我们千万次努力，也无法将它高高举起。同样，如果我们这样破碎我们的心，即使它被千万种膨胀的激情搅动，它也无法被举起，一点也不。因为如果有人为今生之事哀伤，驱除灵魂的一切疾病，那么为罪哀伤的人，将更享受节制的福分。\par

9. \Quote{但是谁能，}有人会说，\Quote{如此破碎自己的心呢？}请听达味，他因这一点而变得卓越，看他心灵的痛悔。在行了成千上万的善功之后，当他即将失去国家、家园和生命本身，在他灾难的时刻，看见一个卑贱的、被遗弃的普通士兵践踏他命运的转变并辱骂他；他非但没有回骂，反而坚决制止一位想要杀他的将领，说：\Quote{由他吧，因为主命他这样。}\BibleRef{撒下 16:10} 另外，当司祭们想抬着\Emph{天主}的约柜与他同行时，他不允许；但他说什么？\Quote{让我把它放在圣殿里，如果天主救我脱离眼前的危险，我必看到它的华美；但如果他对我说：我不喜悦你，看，我在这里，让他照他看为好的待我。} 关于撒乌耳的事，一次又一次，多次，表现出何等卓越的节制！是的，他甚至超越了旧法律，接近宗徒的训诫。为此，他默默地承受了从主手而来的一切；不反抗所遭遇的事，只专注于一个目标，即事事服从并遵行祂所定的法律。当他在做了这么多高尚的事之后，看见那个暴君、弑父者、杀害自己兄弟的人，那个有害的、疯狂的人，占据了他的王国，他也没有生气。但他却说：\Quote{如果天主喜欢这样，}他说，\Quote{使我被追赶、漂泊、逃跑，而使他受尊荣，我默然接受，并感谢天主诸多的磨难。} 不像许多无耻厚颜的人，他们连达味善行的最小部分也没有行过，但如果他们看见某人兴旺，而自己遭受一点挫折，就用千万次亵渎毁灭自己的灵魂。但达味不是这样的人；相反，他显示了一切谦和。因此天主也说：\Quote{我找到了叶瑟的儿子达味，一个合我心意的男子。}\par

我们也应获得这种精神，无论遭受什么，我们都能轻松承受，在进入天国之前，我们将在今世收获谦卑带来的益处。因此，\Quote{你们跟我学习，}他说，\Quote{因为我是良善心谦的，你们心里就必得享安息。}\BibleRef{玛 11:29} 因此，为了在今生和来世都享受安息，让我们勤勉地在灵魂中植入万善之母，我指的是谦卑。因为这样我们就能平静地渡过这生命的海洋，并在那宁静的港湾结束航行；赖我们的主耶稣基督的恩宠和慈爱，愿光荣和权能归于他，直到永远。阿们。\par

\SourceRecord{Translated by George Prevost and revised by M.B. Riddle.；Nicene and Post-Nicene Fathers, First Series；Vol. 10.；Edited by Philip Schaff.；Buffalo, NY: Christian Literature Publishing Co.,；1888.；<http://www.newadvent.org/fathers/200103.htm>.}

% structure: p0001 source=h1 target=section reason=page-title

\section{玛窦福音第四篇讲道}

\BibleRef{玛 1:17}\par

\BibleQuote{所以，从亚巴郎到达味共十四代；从达味到被掳往巴比伦共十四代；从被掳往巴比伦直到基督共十四代。}\par

他将所有的世代分成三部分，表明即使在他们的政体改变之后，他们也没有变得更好，而是在贵族政治、君主政治和寡头政治下都同样走在邪恶的道路上；无论是民众领袖、司祭还是君王管辖他们，对他们的德行都毫无益处。\par

但是为什么他在中间部分略过了三位君王，而在最后部分列出了十二代却声称是十四代呢？前一个问题我留给你们去查考；因为我不需要将所有的事都向你们解释，免得你们变得懈怠。但我们来解释第二个问题。在我看来，他在这里似乎是将被掳的时期和基督本身当作一代，千方百计地将祂与我们联系起来。他特意让我们想起那次被掳，表明即使他们下到那里，也没有变得更加清醒；这样，从一切事上就可以看出祂的来临是必要的。\par

\Quote{那么，}有人会说，\Quote{为什么玛尔谷没有这样做，也没有追溯基督的族谱，而是一切都说得简短呢？}在我看来，玛窦是第一个进入主题的（因此他精确地列出了族谱，并在需要之处停笔）；而玛尔谷在他之后，这就是为什么他采取简短的方式，因为他着手处理已经说过和显明的事。\par

那么路加不但追溯族谱，而且用了更多的名字，这是怎么回事？正如所料，玛窦带领之后，路加试图在先前陈述的基础上教导我们一些东西。而且每个人也都仿效了自己的老师：一位是\Person{保禄}，涌流超过任何河流；另一位是\Person{伯多禄}，力求简洁。\par

2. 那么，玛窦为什么不像先知那样在开头说\Quote{我所见的异象}和\Quote{那临到我的话}呢？因为他写信给善意的人，且非常留意他。因为所行的奇迹大声宣告，领受这话的人极其忠信。但在先知的情况下，没有那么多奇迹来宣告他们；此外，假先知的一伙不在少数，恣意闯入他们中间；犹太人民甚至更听从他们。因此在这种情况下，这种开头是必要的。\par

如果曾有奇迹发生，那是为了外邦人的缘故，以增加皈依者的人数；并为了彰显天主的权能，唯恐他们的敌人俘掳了他们之后，以为自己得胜了，因为他们的神很强大。如同在\Place{埃及}，有许多\Quote{闲杂人}出来；之后在\Place{巴比伦}，关于火窑和梦境的事。当他们独自在旷野时，奇迹也行过。在我们这里也是如此：因为在我们当中，当我们刚脱离谬误时，许多奇事显了出来；但后来当真宗教在各国扎根时，它们就停止了。\par

后期所发生的事稀少且有间隔；例如，当太阳在运行中停住，并倒转方向。这件事也可以看到发生在我们身上。因为在我们这一代，在一位超越一切不虔诚的人身上，我指的是\Person{犹利安}，发生了许多奇怪的事。当犹太人试图重建\Place{耶路撒冷}圣殿时，火从地基喷出，完全阻止了他们；当他的司库和他的叔叔兼同名者公开侮辱圣器时，一人\Quote{被虫咬而死}，另一人\Quote{肚腹崩裂}。此外，当在那里献祭时，泉水干涸，饥荒连同皇帝本人进入城市，这是一个非常大的征兆。因为天主通常做这样的事：当罪恶增多，祂看到自己的子民受折磨，而他们的仇敌因统治他们而极度陶醉时，就彰显祂自己的权能。祂在\Place{波斯}对犹太人也这样做了。\par

3. 因此，他将基督的祖先分成三部分并非无目的或偶然，从以上所说就很明显了。且注意他从哪里开始，在哪里结束。从\Person{亚巴郎}到\Person{达味}；从\Person{达味}到\Place{巴比伦}被掳；从那里直到基督自己。因为在开头他就把\Person{达味}和\Person{亚巴郎}紧密地放在一起，在总结时他也同样提到两者。这是因为，正如我已经说过的，应许是赐给他们的。\par

但是为什么当他提到\Place{巴比伦}被掳时，却没有提到下\Place{埃及}呢？因为他们不再害怕\Place{埃及}人，但仍然畏惧\Place{巴比伦}人。一件事是古老的，另一件事是新鲜的，是最近发生的。他们被带下去到\Place{埃及}不是因为罪恶，但被掳到\Place{巴比伦}则是因过犯而被迁移。\par

此外，就这些名字本身而言，如果有人试图翻译其词源，甚至可以从那里得出神圣思辨的重大材料，并且对于新约非常重要；例如，从\Person{亚巴郎}的名字，从\Person{雅各伯}的名字，从\Person{撒罗满}的名字，从\Person{则鲁巴贝耳}的名字。因为这些名字不是无缘无故给他们的。但恐怕我们因篇幅过长而显得令人厌倦，让我们略过这些，继续处理紧急的事。\par

4. 既然列举了祂所有的先祖，并止于若瑟，他没有就此停止，反而加上说：\Quote{若瑟是玛利亚的丈夫}，暗示他是为了她而追溯自己的族谱。然后，免得你听到\Quote{玛利亚的丈夫}时，以为基督是按普通自然律诞生的，请注意他是如何以随后的话来纠正的。\Quote{你听到了，}他说，\Quote{一个丈夫，你听到了一个母亲，你听到了一个被指定给孩子的名字，因此也要听听诞生方式。}\Quote{耶稣基督的诞生是这样的。}\BibleRef{玛 1:18}\Quote{请问，你既然已经提到了祂的祖先，你又告诉我什么诞生呢？}\Quote{我还是愿意告诉你祂诞生的方式。}你看到了吗？他是怎样唤醒听者？因为好像他就要讲述一件异乎寻常的事，他许诺也要告诉其方式。\par

请注意他所提及之事中最令人赞叹的次序。因为他没有直接进入诞生，而是首先提醒我们：从亚巴郎算起有多少代，从达味算起有多少代，从巴比伦流徙算起有多少代；这样，他使细心的听者思考这些时代，以显示这就是先知们所预言的基督。因为当你数算了这些世代，并从时间得知这就是祂，你也会乐意接受祂诞生时所发生的奇迹。因此，当他要讲述一件伟大的事——祂由童贞女诞生时，他首先将这陈述遮盖起来，直到他数完了世代，仅仅说\Quote{玛利亚的丈夫}；或者更准确地说，他甚至在短句中包含了诞生本身的叙述，然后继续数算年代，提醒听者：这就是先祖雅各伯曾说过的，要在犹太统治者终结时来临的那一位；这就是达尼尔先知预先宣告的，要在那些星期之后来临的那一位。如果有人数算天使对达尼尔所说的星期数中的年数，从建城之时直到祂的诞生，通过计算他就会发现两者相符。\par

5. 我请问你，祂是怎样诞生的？\Quote{他的母亲玛利亚已许配}\BibleRef{玛 1:18}他没有说\Quote{童贞女}，而只说\Quote{母亲}；这样他的叙述容易被人接受。因此，他事先预备听者期待一些普通的信息，借此抓住了他，然后通过加上那奇妙的事实使他惊讶，说：\Quote{他们在成婚以前，她就因圣神怀了孕。}他没有说\Quote{在她被带到新郎家以前}；因为她确实已经在他家中。古人大都习惯将许配的妻子安置在自己家中\BibleRef{创 19:8}，至少在那些地方，如今仍可见同样做法。因此，罗特的女婿也和他同在家中。玛利亚自己当时也同样与若瑟同住。\par

那么她为什么没有在许配之前怀孕呢？正如我起初所说的，是为了使所发生的事暂时隐藏起来，使童贞女免遭一切恶意的猜疑。因为那个最有理由嫉妒的人，不但没有将她示众或羞辱她，反而在她怀孕后仍接待并照顾她；很明显，除非他完全确信所发生的事是圣神的运作，否则他不会将她留在身边，并在一切事上服侍她。而且他非常恰当地说\Quote{她\InnerQuote{被发现}怀了孕}，这种表达常用于那些奇怪、出乎一切预期、始料未及的事。\par

因此，不要再进一步，也不要要求比已经说过更多的东西；不要说：\Quote{但圣神是如何使童贞女怀孕的呢？}因为，当本性在运作时，都无法解释形成的样式；那么，当圣神施行奇迹时，我们又如何能够表达呢？而且，为避免你不断询问这些事而使福音作者厌倦或打扰他，他已经说了施行奇迹的那一位，并以此隐退。\Quote{因为我知道，}他说，\Quote{没有更多，只有所发生的是圣神的工作。}\par

6. 那些忙于探究天上的受生的人真该羞惭。因为如果这个诞生——有无数见证者，并且早在许久以前就已宣告，且已显现并由手摸过——尚且没有人能解释；那么，那些对那不可言明的受生忙碌好奇的人，岂不几乎达到疯狂的程度？因为加俾额尔和玛窦都未能多说任何话，只说它是出于圣神；但如何出于圣神，或以何种方式，他们两人都没有解释；因为那也是不可能的。\par

不要以为你们因听到 \Quote{论及圣神} 就学会了一切；不，即使知道了这一点，我们还有许多无知的事；例如，无限者如何居于子宫之中，包含万有的那位如何像未出生者一样被女人怀抱，童贞玛利亚如何生育并仍然是童贞玛利亚。请告诉我，圣神如何建造了那圣殿？祂如何不是从子宫中取出全部血肉，而只取一部分，使之增长并成形？因为祂确实出自童贞玛利亚的肉身，祂通过说 \BibleQuote{在她内受孕的} \BibleRef{迦 4:4}，而保禄通过说 \BibleQuote{由女人所生}，来阻止那些说基督如同通过某条管道来到我们中间的人的口。因为如果真是这样，子宫有何必要？如果真是这样，祂与我们就没有丝毫共同之处，那肉躯是别样的，不属于我们所属的质料。那么，祂如何是叶瑟的根？如何是枝条？如何是人之子？如何玛利亚是祂的母亲？如何祂是达味的后裔？祂如何 \BibleQuote{取了奴仆的形体} \BibleRef{斐 2:7}？\BibleQuote{圣言如何成了血肉} \BibleRef{若 1:14}？而保禄在致罗马人书中说：\BibleQuote{按肉身说，基督也是从他们来的，祂是在万有之上，永远受赞美的天主。} \BibleRef{罗 9:5} 因此，祂来自我们、属于我们的实体、出自童贞玛利亚的子宫，从这些事和别的事上是明显的；但如何发生却不明显。所以你也不要追问；要接受所启示的，不要好奇于所隐藏的。\par

\BibleQuote{她的丈夫若瑟，} 他说，\BibleQuote{因是义人，不愿公开羞辱她，有意暗暗地休退她。} \BibleRef{玛 1:19}\par

说了那是因圣神，没有同居之后，他以另一种方式再次确立自己的陈述。免得有人会说：\Quote{这从何可见？谁曾听说、谁曾见过这样的事发生？}——或者免得你怀疑门徒为偏袒他的师傅而编造这些事；——他介绍了若瑟，以他亲身经历的事，来证实所述之事；他的叙事几乎是这样说的：\Quote{如果你怀疑我，如果你怀疑我的作证，就相信她的丈夫吧。} 因为他说：\Quote{若瑟，她的丈夫，因是义人。} 此处\Quote{义人}意指在一切事上有德的人。因为既不贪财是义，普遍的德行也是义；而圣经大多以后一种意思使用\Quote{义}这一名称；如经上说：\BibleQuote{那人原是纯朴正直的} \BibleRef{约 1:1}；又说：\BibleQuote{二人在天主前都是义人} \BibleRef{路 1:6}。他既是\Quote{义人}，即是良善而体贴的，\Quote{就愿意暗暗地休退她。} 他以此道出若瑟完全知情以前发生的事，使你不致怀疑他知情以后所做的事。然而，这样的人不仅不会被公开羞辱，而且法律也命令她应受惩罚。但若瑟不仅免去了那更大的惩罚，也免去了较小的，即羞辱。因为他非但未惩罚，甚至不愿公开羞辱她。你看一个自制的人，从最暴虐的情欲中解脱出来。因为你知道嫉妒是多么强烈的事：因此，那清楚知道这些事的主说：\BibleQuote{因为丈夫的嫉妒充满忿怒，在报复之日必不留情} \BibleRef{箴 6:34}；\BibleQuote{嫉妒如坟墓之残忍} \BibleRef{歌 8:6}。我们也都知道许多人宁愿舍弃性命，也不愿堕入受嫉妒的嫌疑。但在此事中，情况不止是嫌疑，子宮的重负完全证实了她。然而他如此没有激情，以至不愿在小事上使童贞玛利亚悲伤。因此，留住她在家里似乎违反法律，而暴露她、把她带上法庭又必会逼他交出她来处死；他两样都不做，现在遵循一种高于法律的准则。因为恩宠既已来临，从此必须有那崇高公民身份的许多标记。正如太阳，虽然尚未展现光线，却已从远处照亮了大半个世界；同样，基督正要从那子宫升起，甚至在祂出来之前，就已经照耀了全世界。因此，甚至在祂分娩之前，先知们就已欢跃，妇人们预言了将要发生的事，若翰尚未从母腹中出来，就从子宫中跳跃。因此，这人也表现出极大的自制，既不控告也不责骂，只打算休退她。\par

事情既然如此，所有人都不知所措，天使前来解决所有困难。但值得追问，为何天使不更早说话，在丈夫有这样的念头之前？而是直到 \Quote{他思念这事的时候}，他才来；因为经上说：\BibleQuote{他思念这事的时候，天使} 就来了。然而，天使在她怀孕之前就向她宣告了喜讯。这又包含另一个困难；因为即使天使没有说话，那被天使告知的童贞玛利亚为何沉默？当她看见她的未婚夫困扰时，为何不消除他的困惑？\par

那么，为什么天使不在若瑟烦乱以前说话？因为我们必须先解释前面的疑难。那么，天使为何不早说？免得若瑟不信，遭遇与匝加利亚相同的事。因为当事情显而易见时，相信就容易了；但事情尚未开始，接受他的话就不那么容易了。为此，天使起初没有发言，同样，童贞玛利亚也保持沉默。因为她不认为能取得未婚夫的信任，向他宣告闻所未闻的事，反而会更激怒他，仿佛她是在掩饰已犯的罪。既然她自己——这位领受了如此大恩宠的人——也多少带有人性的反应，说道：\Quote{这事如何成就？因为我不认识男人。}\BibleRef{路 1:34}那么他更会怀疑了，尤其是从被怀疑的女人口中听到这事。因此，童贞玛利亚什么也不对他说，而天使在时机需要时，向他显现。\par

9. 那么，或许有人问：为何天使不在童贞女身上也同样行事，而在她怀孕以后才向她宣告喜讯？免得她焦虑不安。因为她若不知实情，很可能对自己做出错误的判断，甚至可能上吊或自刺，无法忍受羞辱。因为那位童贞女实在奇妙，路加指出了她的卓越之处，说当她听到问候时，并未立刻倾倒自己，也没有接受那话，反而\Quote{惊慌不安}，思索\Quote{这问候是什么意思。}\BibleRef{路 1:29}这样一位如此完美细腻的女子，在想到自己的羞辱时，定会惊慌失措，因为她不指望凭自己所说的话能说服任何听到这事的人，反而会以为所发生的是奸淫。因此，为防止这些事，天使在怀孕之前就来了。此外，万物的创造者要进入的胎宫，理应不受烦扰；那被认为配得为如此奥迹服务的灵魂，也应当摆脱一切扰动。因此，出于这些原因，祂在怀孕前对童贞女说话，而在分娩的时刻对若瑟说话。\par

许多较为单纯的人因不理解而说这里有矛盾；因为路加说是向玛利亚传报喜讯，玛窦却说是向若瑟；他们不知道两者都发生了。在整个历史中，必须牢记这一点；因为这样，我们就能解决许多看似矛盾的地方。\par

10. 于是，当若瑟烦乱时，天使来了。因为除了上述原因外，也是为了显示他的自制，天使延迟了来临。但当事情即将发生时，他终于显现了。\BibleQuote{当他在思虑这些事时，一位天使在梦中向若瑟显现。}\BibleRef{玛 1:20}\par

你看见这丈夫的温和吗？他不但没有惩罚，甚至没有向任何人宣告，连他所怀疑的人也没有，独自在心里思忖，似乎连童贞女本人也想隐瞒原因。因为经上没有说他打算\Quote{休了她}，而是\Quote{暗暗地休退她}，这人真是极其温和良善。\BibleQuote{当他正在思虑这些事时，天使在梦中显现。}\par

为什么不像对牧人们、对匝加利亚和对童贞女那样公开显现呢？因为这人非常充满信德，不需要这样的异象。而童贞女，由于被宣告了极伟大的喜讯（比给匝加利亚的更大），且是在事情发生之前，所以也需要一个奇妙的异象；牧人们则因天性较为迟钝粗鲁。但此人是在怀孕之后，\newline{}两次显现之间的间隔很大；因此也不需要责备。\par

但通过说\BibleQuote{不要害怕}，他表明若瑟曾害怕，恐怕因收留一个淫妇而得罪天主；因为若非如此，他甚至不会想到要休退她。于是，天使从各方面指明自己来自天主，将他所想和心中所感受的一切都摆在他面前。\par

天使提到她的名字后，并未停留于此，还补充说\BibleQuote{你的妻子}；如果她已被玷污，他就不会这样称呼了。在这里，他称那已订婚的为\Quote{妻子}；正如圣经惯常称未结婚前的订婚女婿为女婿。\par

但\BibleQuote{娶过来}是什么意思？就是把她留在自己家中，因为她在意念上已被休弃了。天使说：\Quote{她已经休了，你就把她娶过来，当作是天主托付给你的，而不是她父母给托付的。天主托付她，不是为了婚姻，而是为了与你同住；借着我的声音，天主托付她。}正如基督后来将祂的母亲托付给祂的门徒，如今也同样托付给若瑟。\par

12. 然后，天使隐约提及了此事，没有提到那邪恶的猜疑；反而以更恭敬合宜的方式，通过告知分娩的原因消除了这猜疑；暗示那使他害怕并想要休弃她的原因——正是这个原因成了他娶她并留她在家的正当理由。这样，天使完全消除了他的痛苦。他说：\Quote{因为她不仅没有非法结合，而且她的怀孕甚至超越一切自然。因此，你不仅要消除恐惧，更要大大喜乐，\InnerQuote{因为那在她内受孕的是出于圣神。}}\par

天使所说的事是奇异的，超越人的理性，高于一切自然律。那么，这样全新的喜讯，他怎能相信呢？天使说：\Quote{借着过去的事，}天使说，\Quote{借着启示。}为此，天使将他心中一切——他所感觉的、所害怕的、所决定要做的——都显明出来，使他借此得以确信这一点。\par

更确切地说，他不仅借过去之事，也同样借未来之事赢得他的心。他说：\Quote{她将要生一个儿子，你要给祂起名叫耶稣。}（\BibleRef{玛 1:21}）因为你不要因祂是由圣神所生，就自以为与这一救恩计划的职分无关。因为虽然在生产上你无分，童贞女仍然没有接触，然而，凡属于父亲的事，不损伤童贞的尊荣，我就赐给你，就是为那所生的起一个名字：因为\Quote{你要称祂}。因为那后代虽不是你的，你仍要对他显出父亲的关怀。因此，我甚至从命名开始，就直接将你与那所生的联系起来。\par

然后，免得另一方面有人因此怀疑他是父亲，请听他如何精确谨慎地陈述此事。他说：\Quote{她将要生一个儿子，}他不是说\Quote{为你生一个儿子，}而只是说\Quote{她将要生，}说得不明指人：因为她不是为他而生，而是为整个世界而生的。\par

13. 为此缘故，天使也从天上带来祂的名字，再次表明这是一次奇妙的诞生：是天主自己从天上藉天使将名字送给若瑟。因为这不是没有目的的，而是万福的宝藏。因此天使也解释它，并暗示美好的希望，这样再次引导他相信。因为这些事我们通常更倾向，所以也更乐意相信它们。\par

因此，在借着一切——过去之事、未来之事、现在之事、以及给予他自己的尊荣——确立了信心之后，他适时地请先知也来，为这一切作证。但在引出先知之前，他预先宣告了那将要借着祂临于世界的福乐。这些是什么？就是罪被移除和消除。\Quote{因为祂要把自己的子民从他们的罪中拯救出来。}\par

这里再次表明这事是超乎一切期望的。因为他所宣告的救恩喜讯，不是来自可见的战争，也不是来自野蛮人，而是来自比这些远为重大的罪恶；这是以前从未有人能成就的工作。\par

但有人也许会问，他为什么说\Quote{祂的子民}，而没有加上外邦人呢？那是为了暂时不惊动听者。因为对于有理解的听者，他也隐约地暗示了外邦人。因为\Quote{祂的子民}不仅是犹太人，也是所有亲近并接受从祂而来的知识的人。\par

并且请注意，他如何顺便向我们揭示了祂的尊威，称犹太民族为\Quote{祂的子民}。因为这话无他，只是暗示那所生的是天主的儿子，他所谈论的是至高者的君王。同样，赦免罪过不属于任何别的权能，而只属于那唯一的神体。\par

14. 既然我们分享了如此伟大的恩赐，让我们尽力不做任何有辱这一恩惠的事。因为如果在获得这尊荣之前，所做的事已应受惩罚，那么在这无可言喻的恩惠之后，更应如此。我现在这样说不是无缘无故，而是因为我看到许多人在领受圣洗之后，生活得比未领洗的人更加漫不经心，在生活方式上没有任何特别之处来区别于他们。你看，正因为此，无论在市场还是在教堂，都无法很快认出谁是信徒谁是不信者；除非在举行奥秘礼仪时在场，看见一类人被赶出，另一类人留在了里面。然而，他们不应该由所处的位置来区分，而应由生活方式来区分。因为正如人的外在尊荣自然由他们披戴的外在标志所发现，同样，我们的尊荣应由灵魂辨别出来。也就是说，信徒不仅应借恩赐显明，也应由新生活显明。信徒应成为世界的光和盐。但当你连自己也不照亮，也不包扎自己的毒疮，还有什么留下让我们认识你呢？因为你进入了圣水？不，这对你成了惩罚的仓库。因为尊荣的伟大，对于那些不选择活出与尊荣相称生活的人，是惩罚的增加。是的，信徒不仅应由他从天主所领受的发光，也应由他自己所贡献的发光；并应在一切事上可辨别——由他的步伐，他的面容，他的衣着，他的声音。我说这些话，不是为了炫耀，而是为了观者的益处，以此作为我们塑造自己的准则。\par

15. 但现在，不论我试图借什么来认识你，我发现你在所有点上都被相反的事所区分。因为若我想按你的处所来辨别你，我看见你在赛马、剧场和不法场所中消磨时光，在市场中的邪恶集会里，在堕落的人群中；或者按你面容的样式，我看见你不断地过度大笑，像嬉笑和放荡的妓女那样不检点；或者按你的衣服，我看见你并不比舞台上的人更好；或者按你的随从，你带着谄媚者和奉承者；或者按你的言语，我听你说不出什么有益的、必要的、对我们生命重要的话；或者按你的餐桌，从那里更重的指控将临到你身上。\par

那么，告诉我，让我借什么来认出你身上的信徒？而我所提到的这一切都给你下了相反的判决。我为什么说\Quote{信徒}？因为我甚至不能清楚认出你是否是一个人。因为当你像驴一样踢人，像公牛一样放荡，像马一样向女人嘶鸣；当你像熊一样贪吃，像骡一样娇惯肉体，像骆驼一样怀恨；\newline{}还有吗？\par

再者，如果我要求你使另一个人变得温和，即使这样我也不应显得是在命令不可能的事；尽管如此，你可能那时会反驳说，你无法控制别人的性情，且这并不完全取决于你。但现在，这野兽是你自己的，并且完全依赖于你。那么你有什么借口呢？或者你能提出什么合理的托词，当你将一头狮子变成一个人，而你却不在乎自己身为一个人却变成了一头狮子？你将超乎本性之物施予野兽，却连本性的东西也不为自己保留？是的，当野兽因你的努力而提升到我们高贵的地位时，你自己却从天国的宝座上跌落，被驱逐到它们的疯狂中。因此，如果你愿意，想象你的愤怒是一种野兽，并且如同别人在狮子身上所显示的那么大热情，你也为了你自己而显出来，使那种态度变得温和与驯良。因为愤怒也有可怕的牙齿和利爪，如果你不驯服它，它就会毁灭一切。因为连狮子或毒蛇也没有如此力量撕裂内脏，像愤怒用它的铁爪不断这样做那样。我们看到，它不仅毁坏身体，连灵魂的健康也被它败坏，它吞噬、撕裂、扯碎一切力量，使之对一切都无用处。因为如果一个人在脏腑里养虫，就不能好好呼吸，他的内脏全都腐烂；那么我们心中有这样一条大蛇（我指的是愤怒）在里面吞噬一切，我们怎能产生任何高贵的事物呢？\par

那么，我们如何摆脱这祸害呢？如果我们能喝一种药水，能够杀死我们里面的虫子和蛇。\Quote{这是什么性质的？}有人会问，\Quote{这药水能有如此功效吗？}如果能以充分的信心领受\OriginalTerm{基督的宝血}{precious Blood of Christ}（因为这将有能力消灭一切疾病）；与此并行的是仔细聆听\WorkTitle{圣经}，并在聆听之外加上施舍；因为藉这一切，我们必能克制那玷污我们灵魂的情欲。然后我们才能生活；因为眼下我们确实不比死者强；因为当那些情欲活着时，我们不可能活着，而我们必须必然灭亡。除非我们先在此地杀死它们，它们在来世必会杀死我们；或者不如说，在那死亡之前，它们甚至在此地就要向我们索取极重的惩罚。是的，因为每一种这样的情欲既残忍又专横且不知足，永不停止地每天吞噬我们。因为\Quote{他们的牙齿是狮子的牙齿，}\BibleRef{岳 1:6}或者更甚，远更凶猛。因为狮子一旦满足，通常会丢下倒在路上的尸体；但这些情欲既不知足，也不离开它们所抓住的人，直到将他置于魔鬼近旁。因为它们的权势如此之大，以至于保罗向基督所显示的侍奉——\BibleRef{罗 8:38}——为祂的缘故既轻视地狱也轻视天国，这同样的事，它们也要求被它们所抓住的人做到。因为无论一个人是被女人的爱，还是财富，还是荣耀所缠住，他从此就嘲笑地狱，轻视天国，以便能成就这些情欲的意愿。那么，当保罗说他如此爱基督时，我们不要怀疑。因为当发现有些人如此服侍自己的情欲时，那另一件事以后怎么显得不可信呢？是的，这也是我们渴望基督更软弱的原因，因为我们所有的力量都消耗在这种爱上，我们抢夺、欺诈，并成为虚荣的奴隶；比虚荣更无价值的是什么？\par

因为即使你变得极其显赫，你也将比卑贱者好不了多少；反而正因为这个原因，你甚至会更卑贱。因为当那些愿意给你荣耀、使你显赫的人，正因为这个原因而嘲笑你——即你渴望来自他们的荣耀——这样的例子怎能不反过来对你不利呢？因为这件事确实属于招致谴责之事。因此，正如在想要犯奸淫或行淫的人身上，如果有人赞美或奉承他，这行为本身便使他成了指控者，而非称赞那纵容这种欲望的人；同样，对于渴望荣耀的人，当我们都赞美时，我们给予那些希望变得荣耀的人的是控诉而非赞美。\par

那么为何要把那通常会使你遭遇相反结果的事加在自己身上呢？是的，如果你愿意被荣耀，就轻视荣耀；这样你必比任何人更显赫。为何要像\Person{拿步高}那样感觉？因为他也立了一座像，以为凭木头和没有感觉的形象能为自己获取更大名声，而活人竟愿藉无生命之物显得更荣耀。你看到他疯狂的极端吗？他以为在表达尊敬，反而侮辱了自己？因为当这显明他更依靠无生命之物，而非依靠他自己和活在他内的灵魂，并且因为这个原因，他将那木块提升到如此高的地位，他岂能不显得可笑呢？他努力装饰自己，不是靠他的生活方式，而是靠木版。正如一个人因他房屋的铺地和美丽的楼梯而自命不凡，而非因他是一个人。如今我们中许多人也在模仿他。因为正如他为他的像，所以有些人要求因他们的衣服而受钦佩，另一些因他们的房屋，或他们的骡子和战车，以及他们房屋的柱子。因为既然他们失落了作为人的本质，他们便四处从别处收集那充满极端可笑的荣耀。\par

至于天主高贵而伟大的仆人们，并非藉这些方式，而是藉最合他们身份的方式，甚至藉此方式，他们闪耀发光。因为虽身为俘虏、奴隶、少年、异乡人，并被剥夺了自己的一切资源，他们当时却比那位拥有一切的人更令人敬畏。当\Person{拿步高}发现既无如此巨大的像，也无总督、将领、无数军团、大量黄金或其他足以满足其欲望、显示其伟大的排场时；对这些一无所有的人而言，他们崇高的自制已足够，并显示那位戴冠冕、穿紫袍的人，在荣耀上远不及那些没有这些的人，正如太阳比珍珠更光辉。因为他们被带到全世界面前，同时是少年、俘虏、奴隶；他们一出现，国王眼中就喷出怒火，总督、副官、省长以及魔鬼的整个圆形剧场都围在四周；四面八方的笛声、号角声和一切乐声直冲云霄，在他们耳边回响；火炉烧到无边的高度，火焰直达云端，到处充满恐怖和惊慌。但这一切都没有使他们惊慌，他们反而嘲笑这一切，如同嘲笑那些戏弄他们的孩子，并展现出他们的勇气和温和，发出比那些号角更清晰的声音，说：\Quote{大王！请您知道。}\BibleRef{达 3:18}因为他们不想冒犯国王，连一句话也不愿，只想表明自己的信仰。为此，他们也没有长篇大论，而是简明扼要地说：\Quote{因为有一位天主在天上，他能拯救我们，}\BibleRef{达 3:17}\Quote{你为何向我显示群众？为何显示火炉？为何显示磨利的剑？为何显示可怕的卫兵？我们的主比这一切更高、更强大。}\par

然后，当他们考虑到天主可能愿意甚至允许他们被烧死时，唯恐若这事发生，他们会被视为说谎，便又加上一句说：\Quote{若不发生这事，大王！请您知道，我们不侍奉你的神。}\BibleRef{达 3:18}因为若他们说：\Quote{他不拯救我们，是因为我们的罪过，}万一他不拯救，他们就不会被相信。为此，他们在此处对这点保持沉默，尽管在火炉中他们一再提到自己的罪过。但在国王面前，他们不说这样的话；只说即使被烧死，他们也不会放弃信仰。\par

因为他们行这事并非为了奖赏和回报，而是出于爱；然而他们当时也处于被掳、为奴的境况，未曾享受任何好处。是的，他们失去了自己的国家、自由和一切财产。请别向我提及他们在王宫里的尊荣，因为身为圣洁和正义的人，他们宁愿千次在家乡作乞丐，并分享圣殿中的福乐。经上说：\Quote{我宁愿在我天主殿中做看门人，也不愿在恶人的帐幕中居住。}又说：\Quote{在你的庭院中一日，胜过千日。}因此，他们宁愿千次在家乡作看门人，也不愿在巴比伦作君王。这从他们在火炉中所说的话可以明显看出，他们因滞留异乡而悲伤。因为虽然他们自己享有极大尊荣，但看到其他人的灾难，他们非常烦恼；尤其是圣人们的特征，即没有荣耀、尊荣或任何东西比邻人的福祉更珍贵。请看，例如，即使他们在火炉中，他们也为全体人民祈祷。但我们即使在宽裕时，也不想念我们的弟兄。再者，当他们询问异梦时，\BibleRef{达 2:17}他们关注的是\Quote{不是自己的利益，而是公共福祉}；因为他们用许多事证明了他们轻视死亡。但处处他们挺身而出，仿佛想以恳切祈求感动天主。接着，他们不认为自己足够，便投靠祖先；但对自己，他们只说他们所奉献的不过是\Quote{一颗痛悔的心}。\par

19. 让我们也效法这些人。因为如今也有一座金像被竖立起来，就是玛门的暴政。但让我们不要留意手鼓、笛子、竖琴和其余财富的排场；即使我们必须落入贫穷的火炉中，我们也要选择它，而不去朝拜那偶像，那么\Quote{在火中必有湿润的呼啸风}。所以，我们不要一听\Quote{贫穷的火炉}就战栗。因为那时掉进火炉的人显得更荣耀，而朝拜的人却被毁灭。只是那时一切都同时发生，但在现今，一部分将在此世成就，一部分在来世，一部分既在此世也在将来的日子。因为那些为了不朝拜玛门而选择贫穷的人，将在此世和那时都更荣耀；而那些在此世不义地富有的人，将在那时受到最严厉的惩罚。\par

拉匝禄也从这火窑中走了出來，并不比那些孩童逊色；但那拜了像的富人却受了地狱的罚。因为我们刚才所提的，正是这件事的预像。所以，正如在此，那些掉进火窑的毫发无损，而那些坐在外面的却被猛烈地抓住；同样，那时也是如此。圣人们走过火河时不会受苦，反而会显得喜乐；但那拜了像的，将看见火像凶猛的野兽一样扑向他们，把他们拖进去。因此，如果有人不信地狱，当他看见这火窑时，就让他从眼前的事相信将来的事，不要害怕贫穷的火窑，而应害怕罪的火窑。因为这火是火焰与折磨，而那火却是甘露与凉快；站在火旁的是魔鬼，站在火中的却是天使，他们扇开火焰。\par

20. 这些话让那些富有的人听吧，他们正在点燃贫穷的火窑。因为即使他们伤害不了别人——\Quote{甘露}会前来相助——但他们自己却会成为火焰的易猎之物，这火焰是他们亲手点燃的。\par

那时，一位天使与那些孩童一同下去；现在，让我们与那些在贫穷火窑中的人一同下去，借着施舍，制造一片\Quote{甘露的空气}，把火焰完全扇开，好使我们也能分享他们的冠冕；这样，地狱的火焰也必因基督的话而消散，祂说：\BibleQuote{你们看见我饿了，便给我吃的。}\BibleRef{玛 25:35}因为那时，那句话将代替\Quote{一阵湿润的风在火焰中呼啸}陪伴我们。让我们带着施舍下到贫穷的火窑里去吧；让我们看见那些节制的人在其中行走，脚踩火炭；让我们看见这奇事，不可思议：一个人在火窑中歌唱赞美，一个人在火焰中感恩，被极度的贫穷束缚，却对基督献上许多赞美。因为那些以感恩之心忍受贫穷的人，实在与那些孩童平等了。没有一种火焰像贫穷那样可怕，那样容易点燃我们。但那些孩童并没有着火；反而，当他们感谢主时，他们的绳索也立刻解开了。同样，现在，如果你陷入贫穷而感恩，那么绳索就解开了，火焰也熄灭了；即使它没有熄灭（这更奇妙），它也会变成泉源而不是火焰：那时同样的事发生了，在火窑中间他们享受了纯净的甘露。那是因为火确实没有熄灭，但它完全阻止了烈火灼烧那些被扔进去的人。这在那些按智慧准则生活的人身上也能看到，因为他们即使在贫穷中也比富人更安全。\par

因此，我们不要坐在火窑外面，对穷人无动于衷，免得我们遭遇那时那些刽子手的下场。因为如果你下到他们那里，与孩童们站在一起，火就不再伤害你；但如果你坐在上面，忽视他们在贫穷火焰中的处境，火焰就会烧死你。所以，下到火里去，免得被火烧死；不要坐在火外面，免得火焰抓住你。因为如果它发现你在穷人中间，它就会离开你；但如果你与他们疏远，它就会迅速扑向你，抓住你。因此，不要站在那些被扔进去的人之外，而是当魔鬼命令将那些没有朝拜金子的人扔进贫穷的火窑时，你不要成为扔别人的人，而应成为被扔进去的人，这样你就能成为得救者之一，而不是被烧者之一。因为最有效的\Emph{甘露}就是不被财富的欲望所束缚，与穷人为伍。这些人比所有人都富有，因为他们践踏了财富的欲望。正如那些孩童当时因藐视国王而变得比国王更荣耀。所以，如果你藐视世俗之物，你就会变得比全世界更可敬，就像那些圣人们，\BibleQuote{世界配不上他们}\BibleRef{希 11:38}。\par

为了能配得上天上的事物，我劝你嘲笑眼前的一切。因为这样你在此世会更荣耀，并享受未来的福乐，借着我们的主耶稣基督的恩宠和慈爱；愿光荣与权能归于祂，直到永远。阿们。\par

\SourceRecord{Translated by George Prevost and revised by M.B. Riddle.；Nicene and Post-Nicene Fathers, First Series；Vol. 10.；Edited by Philip Schaff.；Buffalo, NY: Christian Literature Publishing Co.,；1888.；<http://www.newadvent.org/fathers/200104.htm>.}

% structure: p0001 source=h1 target=section reason=page-title

\section{\WorkTitle{玛窦福音讲道集 第五篇}}

\BibleRef{玛 1:22-23}\par

\Quote{这一切事的发生，是为应验主藉先知所说的话：\InnerQuote{看，一位\OriginalTerm{童贞女}{Virgin}将怀孕生子，人要称祂的名字为\OriginalTerm{厄玛奴耳}{Emmanuel}。}}\par

我听见许多人说：\Quote{我们在这里聆听时，感到敬畏；但一出去，又变成另一个人，热心的火焰熄灭了。}那么，怎样才能不发生这事呢？让我们观察其原因。这变化究竟从何而来？是由于我们不善运用时间，以及和恶人交往。因为我们从\OriginalTerm{圣体}{Communion}退下后，不应立即从事与圣体不相称的事务，而应在回家后，立即拿起圣经，召集妻子儿女，共同整理所听见的，然后再去处理生活中的事务。\par

因为，如果你在沐浴后不愿匆忙进入市场，以免市场事务破坏沐浴带来的舒畅，那么领受圣体后更应如此。但事实上，我们却反其道而行，就这样浪费了一切。因为我们所听的道理尚未稳固，外界事务的强大冲击便将其一扫而空。\par

为免如此，你们从圣体退下后，必须认为没有比整理所听见的道理更重要的事了。是的，我们竟将五六天的时间用于今世的事务，却不肯在属灵的事上花费一天，甚至一小部分的时间，这实在是极大的愚昧。你们没看见自己的孩子们吗？无论给他们什么课业，他们整天都在学习。我们也当如此行，否则我们来这里将一无所得，就像每天把水倒进漏水的容器，而不像对待金银那样认真保存所听的。谁若得到几个钱，便会将其放进钱袋并封上印章；但我们得到比金银宝石更珍贵的圣言，领受圣神的宝藏，却不将其存入灵魂的仓库，反而漫不经心地让它们从心中流失。这样，还有谁会怜悯我们呢？我们竟如此自害，使自己陷入极深的贫穷！因此，为使这事不致发生，让我们为自己、为妻子、为儿女立下一条不可更改的律法：把每周的这一天完全用于聆听和回忆所听的道。这样，我们学习后面的内容时会更善于理解；我们的辛劳会更少，而你们的收获会更大，因为你们记着先前所听的，再听后面的就会更受益。这对理解所讲的道也有很大帮助，因为你们能准确知道我们为你们编织的思想脉络。既然不可能一天之内讲完所有内容，你们必须借着不断的记忆，将多日所听的内容连成一条链子，缠绕在灵魂上，使圣经的整体得以呈现。\par

因此，今天我们也不应在未先回忆近日所讲的之前，就继续下面的主题。\par

2. 但今天摆在我们面前的是什么？\Quote{这一切事的发生，是为应验主藉先知所说的话。}他以一种与奇迹相称的语气，尽力高声说道：\Quote{这一切事的发生。}因为他看见天主对人的爱的海洋和深渊，看见那从未被期望的事成为现实，自然律被打破，和好得以实现，那高于万有的降为低于万有的，\Quote{中间的隔墙被拆毁}\BibleRef{弗 2:14}，障碍被除去，此外还有许多更大的事；他用一句话概括了这奇迹，说：\Quote{这一切事的发生，是为应验主所说的话。}因为他说：\Quote{不要以为这些事是现在才决定的；它们早已被预表了。}保禄也处处努力证明这一点。\par

天使又把若瑟引向依撒意亚，为的是即使他醒来后忘记天使刚说的话，也能因着先知的话——他素常所滋养的——而记住天使所说内容的实质。天使没有向玛利亚提这些事，因为她是个年轻女子，不熟悉这些；但他向作为义人且研究先知书的丈夫，用先知的话来推理。在此之前，天使说：\Quote{玛利亚，你的妻子；}但现在，当他将先知引到若瑟面前时，他才把\OriginalTerm{童贞}{virginity}的名字托付给他；因为若瑟若没有先从依撒意亚那里听到关于童贞女的事，就不会在听天使提到童贞女时依然坚定。因为他从先知那里听到的并非新奇之事，而是他熟悉并长久默想的。为此，天使为使所说的话易于接受，引用了依撒意亚。但他并没有停在这里，而是将话语与天主联系起来。因为他没有称那话是依撒意亚的，而是说那是万有之主的。因此他没有说：\Quote{为应验依撒意亚所说的话，}而是\Quote{为应验主所说的话。}因为嘴是依撒意亚的，但神谕是从上而来的。\par

3. 那么，这神谕说了什么？\Quote{看，一位\OriginalTerm{童贞女}{Virgin}将怀孕生子，人要称祂的名字为\OriginalTerm{厄玛奴耳}{Emmanuel}。}\par

有人会说：那么，祂的名字怎么不叫\OriginalTerm{厄玛奴耳}{Emmanuel}，却叫耶稣基督呢？这是因为祂并没有说\Quote{你要称}，而是说\Quote{人要称}，即民众和事件的结果。因为在此祂将事件作为名字；这在圣经中是惯常的，即以发生的事件代替名字。\par

因此，说\Quote{人要称}祂为\OriginalTerm{厄玛奴耳}{Emmanuel}，无非就是说他们看见天主与人同在。因为祂确实一直与人同在，但从未如此明显地显露。\par

但如果犹太人固执，我们就要问他们：那孩子什么时候被称为\Quote{急速掠夺，赶快抢夺}？他们根本无法回答。那么，先知为什么说\Quote{给他起名叫\OriginalTerm{玛黑尔沙拉哈施巴兹}{Maher-shalal-hash-baz}}？因为当他出生时，发生了掠夺和分赃，所以在他时代发生的事件就成了他的名字。同样，经上还说，这城将被称作\Quote{正义之城，忠信的城熙雍}\BibleRef{依 1:26}。然而，我们从未发现这城被称为\Quote{正义}，它仍被称为耶路撒冷。但是，因为当这城变得更好时，这事确实发生了，所以经上说它被称为这样。因为当某一事件比名字更清楚地表明那成就它的人或受益者时，圣经就把事件的事实当作他的名字。\par

4. 但即使他们在这点上哑口无言，他们还会寻求另一个借口，即关于玛利亚童贞的经文，并向我们提出其他译者，说他们用的不是\Quote{童贞女}，而是\Quote{年轻女子}；首先我们要说，\OriginalTerm{七十贤士}{the Seventy}比所有其他译者更值得信赖。因为\OriginalTerm{他们}{these}是在基督来临之后才翻译的，仍然是犹太人，很可能出于敌意故意歪曲预言；而\OriginalTerm{七十贤士}{the Seventy}在基督来临前一百多年就开始了这项工作，完全消除了这类嫌疑，而且因为其年代、人数和一致性，更应值得信赖。\par

但即使他们引用那些人的见证，胜利的记号仍然在我们这边。因为圣经惯常以\OriginalTerm{青年}{youth}一词指\OriginalTerm{童贞}{virginity}；这不仅指女子，也指男子。正如经上说\Quote{少年和童贞女，老年与幼童}。又说到被侵犯的少女时，经上说\Quote{若那年轻女子呼救}，意思就是童贞女。\par

而且前面的经文也确立了这种解释。因为祂不仅说\Quote{看，童贞女要怀孕}，而且首先说\Quote{看，上主自己要给你们一个征兆}，然后接着说\Quote{看，童贞女要怀孕}\BibleRef{依 7:14}。然而，如果那要生产的不是童贞女，而是按婚姻方式发生，那么这事件算什么征兆？因为征兆当然必须超出寻常事件，必须是奇异非凡的；否则怎能是征兆？\par

5. \Quote{若瑟从睡梦中醒来，就照主的天使所吩咐的做了。}你看到他的顺从和谦卑的心了吗？你看到真正的醒悟和不玷污的灵魂了吗？因为当他怀疑某些痛苦或不当的事时，他不能容忍留住童贞女；但在消除这怀疑之后，他也不能把她赶走，反而留住她，并为整个救恩计划服务。\par

\Quote{遂娶了他的妻子玛利亚。}你看到福音作者多么持续地使用这个词吗？他不愿过早揭露那奥秘，并消除那邪恶的猜疑。\par

\Quote{若瑟娶了玛利亚，但与她同房以前，她生了头胎儿子。}他在这里用了\OriginalTerm{以前}{till}这个词，不是让你怀疑以后他认识了她，而是告诉你，在出生前，童贞女完全未被男人接触。但也许有人会说，为什么用\OriginalTerm{以前}{till}这个词？因为圣经中经常这样做，用这个表达而不指有限的时间。例如，关于方舟，经上说\Quote{乌鸦没有回来，直到地上的水干了}\BibleRef{创 8:7}，然而在那之后它也没有回来。又论及天主时，圣经说\Quote{从永远直到永远，你是天主}，这并不是限定时间。又当预言福音时说\Quote{在他的日子里，正义要兴盛，和平富足，直到月亮被除去}，并没有给这美丽的受造物设定界限。同样，这里也用了\OriginalTerm{以前}{till}一词，以确定出生前的情况，至于之后的事，则留给你推断。因此，祂自己说了你必须学习的：童贞女在出生前未被男人接触；但随后的事，既是从前一句推断出的，也是公认的，就留给你去理解：即在此之后，她既然成为了母亲，并获得了如此奇异的产痛和如此奇特的生育，那义人怎会再认识她呢？因为如果他认识了她，并将她作为妻子保留，主耶稣又怎会把她作为无助、无亲之人交给祂的门徒，并命令他接她到自己家里呢？\BibleRef{若 19:27}\par

那么，有人会说，雅各伯和其他人怎么被称为祂的兄弟呢？正如若瑟本人被认为是玛利亚的丈夫一样。因为有许多遮蔽，为的是使这如此出生的奥秘暂时隐藏。因此，连若望也这样称呼他们，说\Quote{连祂的兄弟也不信祂}。\par

6. 然而，那些起初不信的人，后来却成了可敬而显赫的。至少当保禄和他同行的人为了谕令上耶路撒冷时，他们径直去见雅各伯。因为他如此受人敬仰，以至于首先被托付了主教的职务。据说他过着极刻苦的生活，以致他的肢体都如同死了一般；由于不断的祈祷和与地面不断的接触，他的额头变得如此坚硬，以至于不比骆驼的膝盖好多少，只因他如此频繁地叩头于地。这个人后来当保禄再次上耶路撒冷时，亲自指示保禄说：\Quote{兄弟，你看，有多少万信众聚集了。}他的智慧和热忱如此之大，或者更确切地说，基督的德能如此之大。因为那些在他活着时嘲笑他的人，在他死后却充满敬畏，以至于极其乐意为他而死。这些事最有力地显示了他复活的德能。你看，这就是更光荣的事留待后来的原因，即为了使这个证据成为无可争议的。因为即使那些在我们中间活着时受人敬佩的人，当他们去世后，也容易被我们遗忘；那么，如果这个人只是众人中的一个，那些在他活着时轻视他的人，怎么后来竟认为他是天主呢？他们怎么会同意为他被杀，除非他们清楚无误地领受了他的复活？\par

7. 我们告诉你们这些事，是叫你们不仅听，还要效法他那种刚毅的严厉、坦率的言辞和凡事上的正义；好使无人对自己绝望，即使他至今仍是疏忽的，使他除了天主的仁慈之外，不把希望寄托在任何事上，而只寄托于自己的德行。因为如果这些人虽与基督同家同族，但直到他们证明了自己的德行之前，这种亲属关系并未使他们受益；那么当我们以正义的亲属和兄弟为理由恳求时，除非我们极其顺服并活于德行，否则我们能得到什么恩惠呢？正如先知也说了同样的事：\Quote{兄弟不能赎买，人岂能赎买？}不，即使他是梅瑟、撒慕尔、耶肋米亚。例如，请听天主对这位先知说的话：\Quote{你不要为这百姓祈祷，因为我不听你。}\BibleRef{耶 11:14} 你为什么惊奇，如果我不听你呢？\Quote{即使梅瑟和撒慕尔站在我面前，}\BibleRef{耶 15:1} 我也不会为他们接纳这些人的恳求。是的，如果是厄则克耳恳求，他也会被告知：\Quote{即使诺厄、约伯和达尼尔站出来，他们也不能救自己的儿女。}即使先祖亚巴郎为那些病入膏肓、不肯悔改的人求情，天主也会离开他继续前行，\BibleRef{创 18:33} 以致不接受他为他们的呼求。同样，如果是撒慕尔这样做，天主对他说：\Quote{你不要为撒乌耳悲伤。}\BibleRef{撒上 16:1} 即使有人为自己的姊妹求情，但时机不合，他也会得到与梅瑟相同的回答：\Quote{假如她父亲在她脸上吐唾沫。}\BibleRef{户 12:14}\par

因此，我们不要张口仰慕他人。的确，圣人们的祈祷具有极大的力量，但条件是我们悔改并改善自己。因为梅瑟曾解救了自己的兄弟和六十万人脱离那时从天主而来的愤怒，却无力解救自己的姊妹；况且罪过并不相等：她只是冒犯了梅瑟，而在另一件事中，他们所犯的是明显的亵渎。但我把这个难题留给你们；而更难的，我将试着解释。\par

因为何必提及他的姊妹呢？既然那为如此大的民族作辩护的人，却不能为自己获得恩惠；相反，在他无数的劳苦、磨难和四十年的勤劳之后，他被禁止踏上那片曾有如此多宣告和应许的土地。那么原因是什么呢？因为赐予这个恩惠并无益处，反而会带来许多害处，并肯定成为许多犹太人的绊脚石。因为如果他们只是从埃及被解救出来就离弃了天主，去追随梅瑟，并把一切都归功于他；那么如果他们看见他也带领他们进入应许之地，他们会陷于何等的不敬虔呢？为此，我也补充一点，他的坟墓也没有被显明。\par

同样，撒慕尔也不能拯救撒乌耳脱离上天的愤怒，尽管他时常保护以色列人。耶肋米亚不能拯救犹太人，但他或许曾借预言使某人免于灾祸。达尼尔拯救了巴比伦人免受屠杀，\BibleRef{达 2:24} 但他并没有解救犹太人脱离充军。\par

在福音中，我们也看到这两件事发生，不是在不同的人身上，而是在同一个人身上；同一个人有时为自己求得恩惠，有时却被舍弃。因为那欠一万塔冷通的仆人，虽然曾借恳求使自己脱离危险，但后来却又未能求得恩惠；\BibleRef{玛 18:26} 而另一个人，先前曾自暴自弃，后来却能最大程度地帮助自己。\BibleRef{路 15:13} 这是谁呢？就是那挥霍父亲家产的人。\par

因此，一方面，如果我们疏忽大意，即使有他人的帮助，我们也不能获得救恩；另一方面，如果我们警醒，我们就能靠自己做到，而且靠自己比靠他人更容易。是的，因为天主更愿意将祂的恩宠赐给我们自己，而不是借他人赐给我们；这样，我们亲自努力平息祂的愤怒，就能对祂充满信心，并成为更好的人。例如，祂怜悯了客纳罕妇人，就这样拯救了妓女，就这样拯救了强盗，当时没有人为他们作中介或辩护。\par

8. 我说这话，并非要我们停止恳求圣人，而是要防止我们疏忽，并将我们的挂虑仅托付给别人，而自己却后退沉睡。因为当祂说：\Quote{要为自己结交朋友}，并非仅止于此，反而加上：\Quote{用不义的钱财}，好使善工又成为你们自己的；因为祂在这里所指的无非是施舍。更奇妙的是，若我们远离不义，祂也并不与我们严格算账。因为祂所说的是这样的：\Quote{你得不义之财吗？好好花费。你以不义聚敛吗？以正义分散。}然而，用这样的所得来施舍，算是什么德行呢？但天主满有对人的爱，竟屈尊俯就至此，而如果我们这样行，祂应许我们许多美物。但我们如此麻木不仁，连从我们的不义之财中施舍也不肯，反而不停地掠夺，若我们贡献最微小的一部分，就以为已尽了所有。你们没有听见保禄说：\Quote{那少量播种的，也要少量收获}吗？\BibleRef{格后 19:6}那么你为何吝啬呢？什么，这行为是支出吗？是花费吗？不，这是收益和好买卖。哪里有买卖，哪里就有增长；哪里有播种，哪里就有收割。但你们，若你们必须耕种一片肥沃深厚、能接受许多种子的田地，你们会既花费已有的，还会向别人借贷，认为在这种情况下吝啬是损失；但是，当你们要耕种的是天堂，它不受任何气候影响，并必以丰厚的增长回报你们的支出时，你们却迟延退后，且不认为可能因吝啬而损失，因不吝啬而得益。\par

9. 因此，要分散，免得丧失；不要收藏，为要保存；要支出，为要积蓄；要花费，为要获得。如果你的财宝需要囤积，你不要囤积它们，因为你必定会丢弃它们；但要把它们托付给天主，因为在那里无人能夺走。不要做生意，因为你根本不知道如何获利；但要借给那一位，祂所给的利息大于本金。要借给那没有嫉妒、没有控告、没有恶谋、没有恐惧的地方。要借给那一无所缺、却为你的缘故而有所需要的那一位；祂养活众人，却忍饥挨饿，使你不致遭受饥荒；祂贫穷，为要你富足。在那里借出，你的回报不会是死亡，而是生命取代死亡。因为这种高利贷是王国的前兆，那种高利贷却是地狱的前兆；一种出于贪得无厌，另一种出于克己；一种出于残忍，另一种出于仁爱。那么，当我们有能力获得更多、而且有保障、在适当的时候、在极大的自由中、没有指责、没有恐惧、没有危险，却放弃这些收益，去追求那另一种卑劣下贱、不安全、易朽坏、并大大为我们增加火炉之刑罚的收益时，我们还有什么借口呢？因为世上没有什么，没有什么比这世界的高利贷更卑贱，更残忍的了。为什么？别人的灾难正是这种人的买卖；他从别人的痛苦中获利，为仁慈索取报酬，好像他怕显得慈悲似的，在仁慈的掩盖下挖掘更深的陷阱，借着帮助的行为刺痛人的贫穷，在伸出手时把他推倒，在接纳他如入港口时，使他遭遇海难，如同撞在岩石、浅滩或暗礁上。\par

有人会说：\Quote{那你要求什么呢？}\Quote{难道我应该把我积蓄的、对我有用的金钱给别人使用，而不要求回报吗？}绝非如此：我并非说这话；是的，我诚恳地盼望你们得到回报，但不是卑微或微小的，而是远为更大的；因为以金子为代价，我愿你们为高利贷而接受天堂。那么，为什么把自己封锁在贫穷中，在地上爬行，以小的换取大的呢？不，这是不懂如何致富的人的行为。因为当天主以少量金钱为代价，应许你们天堂的美好事物时，你们却说：\Quote{不要给我天堂，而要以天堂换取会朽坏的金子}，这是希望继续贫穷的人。因为那渴望财富和丰裕的人，必会选择持久而非短暂的事物；取之不尽而非消耗殆尽；以多换少，以不朽换可朽。因为这样，那另一类也会随之而来。因为那先寻求地上而非天堂的人，也必丧失地上；而那以天堂优先于地上的人，必能极好地同时享受两者。为使这成为我们的现实，让我们鄙视这里的一切，选择将来的美物。因为这样，我们必能借着我们的主耶稣基督的恩宠和慈爱，同时获得两者；愿光荣和权能归于祂，直到永远。阿们。\par

\SourceRecord{Translated by George Prevost and revised by M.B. Riddle.；Nicene and Post-Nicene Fathers, First Series；Vol. 10.；Edited by Philip Schaff.；Buffalo, NY: Christian Literature Publishing Co.,；1888.；<http://www.newadvent.org/fathers/200105.htm>.}

% structure: p0001 source=h1 target=section reason=page-title

\section{玛窦福音讲道集 第六篇}

\BibleRef{玛 2:1-2}.\par

\BibleQuote{耶稣在犹太的伯利恒诞生时，黑落德王在位，忽然有贤士从东方来到耶路撒冷，说：才诞生的犹太人的王在哪里？我们在东方看见了他的星，特来朝拜他。}\par

我们需要大量的警醒和许多祈祷，才能解释眼前的这段经文，并且认识这些贤士是谁，他们从哪里来，如何来，受谁的劝说，以及那颗星是什么。或者，如果你们愿意，让我们先提出真理的仇敌所说的话。因为魔鬼猛烈地吹向他们，甚至试图从这段经文武装他们来反对真理的话语。\par

那么，他们说了什么？他们说：\Quote{请看，}他们说，\Quote{就连基督诞生时也出现了一颗星；这是星象学可以依赖的记号。}那么，如果他的诞生是按照那定律，他怎能废除星象学、除去命运、堵住魔鬼的口、驱逐谬误、推翻一切这样的巫术呢？\par

再者，贤士们从星本身学到了什么？他是犹太人的王？然而他并不是这国的王；正如他也对比拉多说：\BibleQuote{我的国不属于这世界。}无论如何，他没有显示这种王权，因为他身边既没有持矛或盾的卫兵，也没有马匹、骡车，或任何其他这类东西；他反而过着卑微贫穷的生活，身边带着十二个出身微贱的人。\par

即使他们知道他是王，他们来的目的是什么？因为星象学的工作，据说，不是从星知道谁诞生，而是从人诞生的时辰预测他们将遭遇什么。然而这些人既没有在产妇的痛苦中陪伴，也不知道他诞生的时辰，也没有从那一刻开始，根据星辰的运行推算出将要发生的事；反而相反，他们在自己的国家早就看见一颗星出现，就来看那诞生的一位。\par

这种情况本身比前者更难解释。因为什么理由促使他们，或期望什么好处，去朝拜一个如此遥远的王？如果他要统治他们自己，即使如此，这件事也无法给出合理的解释。当然，如果他诞生在王宫里，他的父亲又是国王，在他身边，任何人都会自然地说，他们出于讨好父亲的愿望，朝拜了诞生的孩子，从而为自己预先积攒了大量的庇护关系。但现在他们甚至没有期望他是自己的王，而是一个陌生民族的王，远离他们的国家，又没有看到他长大成人；为什么他们要走上如此漫长的旅程，献上礼物，而且整个行动充满危险？因为黑落德一听到这事，就极其不安；众人听到他们讲述这些事，也惊慌失措。\par

\Quote{但是这些人没有预见到这个。}不，这不合理。因为他们即使极其愚笨，也不可能不知道，当他们来到一个有王的城市，宣告这样的事，并树立另一个王与现任国王对立时，他们必然会招致千百次死亡。\par

2. 他们为何去朝拜一个躺在襁褓中的婴儿？如果他是一个成人，也许可以说，他们预料到将来会得到他的援助，因而投入他们预见的危险；但波斯人、野蛮人、与犹太民族毫无关系的人，愿意离开家乡、放弃祖国、亲戚和朋友，并且服从另一个王国，这实在是极其不合理的。\par

如果这是愚笨的，那么下文就更愚笨了。是什么呢？他们踏上如此漫长的旅程，朝拜了，搅动了所有人，然后立刻离开。他们看到棚屋、马槽、襁褓中的婴儿和贫穷的母亲，到底看到了什么王权的迹象？此外，他们向谁献上礼物，为了什么目的？难道到处都有这样的惯例，向每一位诞生的国王献殷勤？他们是不是总是走遍天下，朝拜那些他们知道将从卑微低贱的地位成为国王的人，在他们登上王位之前？没有人能这样说。\par

他们到底为何朝拜他？如果是为了眼前的事物，他们期望从一个婴儿和一个卑微的母亲那里得到什么？如果是为了将来的事物，他们如何知道这个他们在襁褓中朝拜的孩子会记得当时所做的事？但如果他的母亲要提醒他，他们不但不应受尊敬，反而应受惩罚，因为他们把他带入他们必然预见的危险中。因为正是由此，黑落德才不安、搜索、查问，并着手杀害他。事实上，任何宣告未来的国王的人，在他幼年微贱时，无非是把他出卖给杀戮，并给他挑起无尽的战争。\par

你们看，如果我们按照人间事物的常规来考察这些事，会出现多少荒谬之处？因为不仅这些主题，还可以提到更多，比我们所说的包含着更多的问题。但为了避免将问题堆叠起来使你们困惑，现在让我们进入所询问题的解答，以星本身作为解答的开始。\par

3. 因為你若能知道那星是什麼、屬何種類、是否為普通星辰之一、或為新星而與眾不同、以及它是由本性為星辰、還是只是外觀上像星辰，我們便容易知道其他事情。那麼，這些要點從何而顯明呢？從所記載的本身。因此，這星不是普通之類，或者說根本不是星，至少在我看來是這樣，而是一種無形的能力轉變為此種外觀，首先從它的運行路線便清楚可見。因為沒有任何星辰是按這種方式運行的；不論你提到太陽、月亮，還是其他所有星辰，我們都看見它們從東向西運行；但這星卻是從北向南飄移；因為巴勒斯坦相對於波斯正是如此位置。\par

其次，從時間上也可見此。因為它不在夜間出現，而是在正午，太陽照耀之時；這不是星辰的能力所能及，甚至月亮也不能；因為月亮雖然遠超眾星，但當太陽的光線出現時，它立刻隱藏並消失。然而這星卻以其自身光輝的強烈勝過了太陽的光線，比它們更明亮，在如此強的光中更加輝煌地照耀。\par

第三，從它的出現與再次隱藏可見。因為當他們一路前往巴勒斯坦時，它在前引領他們；但當他們進入耶路撒冷後，它卻隱藏了；然後，當他們離開黑落德，告訴他前來的原因，正要出發時，它又顯現；這一切並不像是星辰的運動，而像是一種具有高度理性的能力。因為它本身甚至沒有任何固定的運行路線；當他們要前行時，它就行動；當他們要停留時，它就停留，按需要安排一切；正如雲柱那樣，有時停留，有時喚醒猶太人的營地，在需要的時候。\par

第四，從它指明祂的方式可以清楚察覺。因為它並沒有在高處指明那地方；因為他們不可能這樣確定；而是降下來履行這項職責。因為你們知道，一個如此狹小的地方，只有一個棚屋的大小，或者說只有一個嬰孩身體所佔據的空間，不可能由一顆星來標明。由於星的高度極大，它無法充分區分如此狹小的地方，並向那些渴望看見的人顯露。這一點任何人都可以從月亮看出，月亮雖然遠超眾星，但對所有居住在世界、分散在如此廣闊土地上的人來說——我說是——似乎每個人都覺得它很近。那麼，請告訴我，星如何能指出如此狹小的地方，僅僅是馬槽和棚屋的空間，除非它離開那高度，降下來，站在幼童的頭頂之上？福音作者暗示了這一點，當他說：\BibleQuote{看，那星在他們前頭走，直到來到小孩子所在的地方，就停在上頭。}\par

4. 你們看見了嗎？藉著多少證據表明這星不是眾星之一，也不是按照外在受造界的秩序顯現的？它出現是為了什麼目的？為了責備猶太人的麻木，並斷絕他們為故意無知所找的一切藉口。因為，既然那來臨者要終結古時的政體，並召喚世界來朝拜祂自己，且在一切陸地和海洋受朝拜，那麼祂一開始就為外邦人打開了大門，願意藉著陌生人來告誡自己的百姓。因此，因為先知們不斷被聽見談論祂的來臨，而他們卻不怎麼留意，祂竟使野蠻人從遠方前來，尋找在他們中間的君王。他們從波斯人的口中首先學到了他們不願從先知那裡學習的；這樣，如果他們秉性正直，他們便有最強烈的服從動機；如果他們好爭辯，他們從此便沒有任何藉口。因為他們還有什麼可說的呢？他們在這麼多先知之後不接受基督，而看見賢士因見到一顆星就接受了祂，並朝拜了顯現的祂？祂的行動與對待尼尼微人的方式很相似，那時祂派遣了約納；也與對待撒瑪黎雅婦人和客納罕婦人的方式相似。為此，祂也說：\BibleQuote{尼尼微人要起來，且要定這一代的罪；} \BibleQuote{南方的女王要起來，且要定這一代的罪。}\BibleRef{瑪 12:41} 因為這些人相信了較小的事，而猶太人卻連更大的事也不相信。\par

\Quote{有人會說：祂為什麼用這樣的異象吸引他們呢？} 祂該怎樣做呢？派遣先知？但賢士不會順服先知。從天上發出聲音？不，他們不會留心。派遣天使？連天使他們也會忽略。因此，為了這一切，天主放下所有這些方式，以極大的俯就，用他們熟悉的事物呼喚他們；並顯示一顆巨大非凡的星，使他們因它外觀的宏偉與美麗，以及它運行的方式而驚奇。\par

在效法此事上，保禄也曾从一个外邦人的祭台与希腊人辩论，并引用诗人的见证。他向犹太人宣讲时，并非没有割礼。他以祭献作为向生活在法律之下者讲论的开端。因为，既然对于每一个人，熟悉的事物是可爱的，那么天主和祂所差遣的人，都本着这个原则行事，以谋求世界的得救。因此，不要以为祂用星来召唤他们是有失身份；因为按照同样的理由，你也会指责所有犹太人的礼仪：祭献、取洁、新月、约柜，甚至连圣殿本身。因为连这些也源自外邦人的粗鄙。然而，为了拯救那些迷途的人，天主竟容忍自己藉着这些事物受人侍奉，而外邦人原本是用这些来侍奉魔鬼的；只是祂稍加改变，好逐步将他们从习俗中引开，导向最高的智慧。祂对待贤士们也是如此，并不鄙视用星象来召唤他们，以便日后将他们提升得更高。因此，在祂引导他们，牵着他们的手，将他们带到马槽旁之后，祂不再用星，而是用天使向他们说话。这样，他们就渐渐成了更好的人。\par

祂对待阿市刻隆和迦萨的人也是如此。因为那五座城（当约柜来到时，他们遭到致命的瘟疫打击，却无法摆脱所患的灾祸）——城中的人召叫了他们的先知，聚集开会，试图找到摆脱这神圣灾祸的方法。那时，他们的先知说，他们应将从未负轭、而且头胎的母牛犊套在约柜上，让它们自己走，无人引导；因为这样就会清楚看出瘟疫是来自天主，还是某种偶然带来了疾病（有人说：\Quote{因为如果，}他们说，\Quote{它们因不习惯而挣断轭，或是转向牛犊鸣叫的地方，\InnerQuote{那是偶然发生的事；}\BibleRef{撒上 6:9}但如果它们直行，不偏离道路，无论是牛犊的鸣叫还是对道路的陌生都不影响它们，那就很清楚是天主的手降临了那些城市：}）。我说，当那些城市的居民听从他们先知的话，按吩咐做了以后，天主也顺应了先知的建议，在那次也显示了迁就，并不认为实现先知的预言、使他们当时所说的话显得可信，是有失身份的。因为这样成就的善事更大，连祂的敌人自己也见证了天主的能力；是的，他们自己的教师也替祂说了话。人们可以看到天主成就了许多其他类似的事。关于那女巫所发生的事，\EditorialNote{撒慕尔纪上第二十八章}又像是这种安排；现在你也可以从以上所说的话来解释这件事了。\par

关于星，我们已经说了这些，也许你们还能说得更多；因为经上说：\Quote{劝导智慧人，他就更加智慧；}\BibleRef{箴 9:9}但我们现在必须回到所读经文的开头。\par

5. 开头是什么呢？\Quote{当耶稣在犹太的伯利恒诞生时，在黑落德王的的日子里，看，有贤士从东方来到耶路撒冷。}贤士在星的引导下跟随，这些人却不信，尽管先知的声音在他们耳中回响。但为什么他要向我们提到时间和地点，说\Quote{在伯利恒}和\Quote{在黑落德王的的日子里}？他又为什么加上他的身份？加上身份，是因为还有另一个黑落德，就是杀害若翰的那一个；但那一个是分封王，这一个却是王。他记下地点和时间，是为了让我们想起古代的预言；其中一个是米该亚所说的：\BibleQuote{而你，犹大地的伯利恒，在犹大的首领中决不是最小的；}\BibleRef{米 5:2}另一个是圣祖雅各伯所说的，明确向我们指出时间，并展示祂来临的重大标记。因为他说：\BibleQuote{权杖决不会离开犹大，领袖也不会离开他的腰间，直到那应得的人来到，祂是万民的期待。}\BibleRef{创 49:10}\par

这又值得追问：他们从何产生这样的想法？是谁激动他们做这事？因为在我看来，这不只是星的工作，也是天主的工作，祂感动了他们的灵魂；祂在居鲁士身上也做了同样的事，促使他释放犹太人。但祂这样做并没有摧毁他们的自由意志，因为即使祂从天上用声音召唤保禄时，祂也既显示了自己的恩宠，又显示了保禄的服从。\par

有人也许会问：为什么祂不向东方所有的贤士启示这事？因为不是所有人都会相信，只有这些人比其他人准备得更好；就像有无数民族灭亡，但只有尼尼微人被派去先知；有两个强盗在十字架上，但只有一个得救。你至少可以看到这些人的美德，不仅在于他们的到来，也在于他们说话的大胆。为了不让他们看起来像是某种骗子，他们告诉他们是谁指引了道路和旅途的长度；他们来了以后，说话很大胆：因为\Quote{我们来了，}他们说，\Quote{来朝拜祂；}他们既不惧怕百姓的愤怒，也不惧怕国王的暴政。因此，在我看来，他们在本国也应该是同乡的教师。因为那些在这里毫不退缩地说出这话的人，在自己的国家更会大胆发言，因为他们既接受了天使的神谕，又得到了先知的见证。\par

6. 但圣经记载：\BibleQuote{黑落德听了，就惊慌起来，全耶路撒冷也同他一起惊慌。}\Person{黑落德}自然惊慌，因为他是国王，为自己也为自己的子女害怕；但\Place{耶路撒冷}为什么惊慌呢？的确，先知们早已预言祂是救主、恩主和从上而来的解救者。那么，\Place{耶路撒冷}为什么惊慌呢？是由于同样一种情绪，这情绪从前也使他们，当天主向他们倾注恩惠时，背弃天主，并在享受极大自由之际，仍念念不忘\Place{埃及}的肉锅。\par

但请你留意先知们的准确性。因为恰恰是这事，先知从一开始就预言了，他说：\Quote{他们若被火烧，倒会欢喜；因为有一个婴儿为我们诞生了，有一个儿子赐给了我们。}\par

然而，尽管惊慌，他们并不想去查看发生了什么事，也不跟随贤士，也不作任何特别的询问；他们真是既好争吵又漫不经心，超过了所有人。因为当理应自豪——这王诞生在他们中间，并吸引了\Place{波斯}之地归向祂，他们几乎要让所有人都臣服于他们，好像他们的境况已迈向改善，且从起初祂的帝国就已如此光荣——的时候，他们竟不因此而变得更好。而且他们刚从那里的囚禁中获得释放；他们自然可以这样想（即使他们不知道那些高超和奥秘的事，只从眼前事物判断）：\Quote{如果他们在我们君王诞生时就这样战栗，那么祂长大后岂不更令他们畏惧和服从，我们的地位岂不比异邦人更光荣？}\par

7. 但这一切丝毫未能唤醒他们，他们的迟钝如此严重，还有他们的嫉妒：这两样我们都必须仔细从心中根除；而要在这样的阵列中站立得住，必须比火更炽热。因此基督也说：\BibleQuote{我来是把火投在地上，我多么希望它已经燃烧起来！}与此相同，虔诚的泪水是永久不衰的喜乐的种子。就这样，那位妓女被这火抓住时，变得比贞女更尊贵。也就是说，她被悔改彻底温暖后，遂被对基督的渴望带出自我；她解开头发，用眼泪沾湿祂的圣足，用自己发辫擦拭，并倾倒香液。这些都是外在的结果，但她内心所做的远比这些更炽热；这些唯有天主自己看见。因此，每个人听见时，都和她一同欢喜，赞赏她的善工，免去她的一切责备。但我们这些邪恶的人若作此判断，请想想她从那爱人的天主获得了怎样的判决；并且，甚至在领受天主的恩赐之前，她的悔改为她收获了怎样的祝福。\par

因为正如大雨过后，天空晴朗明净；同样，当泪水倾泻时，平静和安宁升起，我们罪孽后的黑暗全然消失。又如藉着水和圣神，我们也藉着泪水和告解得第二次洁净；除非我们这样做是为了炫耀和虚荣：因为对于一个眼泪属于那类别的女人，我该公正地称她为可责难的，比她用画线和颜色装饰自己更甚。因为我寻找的眼泪不是为炫耀而流，而是出于痛悔；那些秘密地、在斗室中、无人看见时、轻轻地、无声地滴下的眼泪；那些从灵魂深处涌出、在痛苦和忧伤中流下的、只为天主而流之泪；正如\Person{亚纳}的眼泪，经上记载：\BibleQuote{她的嘴唇在动，却听不到她的声音；}然而，她的眼泪独自发出比任何号角更清晰的呼喊。正因为此，天主也开了她的子宫，使那硬石变为沃土。\par

如果你也这样哭泣，你便成了你主的门徒。是的，因为祂也曾哭泣，为\Person{拉匝禄}，也为那座城；关于\Person{犹达斯}，祂极为忧伤。这一点我们常常看见祂做，但从未见祂大笑，甚至从未微笑；至少没有一位福音作者提过这事。因此关于\Person{保禄}，他哭泣，他三年来日夜流泪，他自己说过，别人也这样论他；但他大笑，他自己从未提过，也没有任何其他圣人提过他或任何像他这样的事；但这事只提到\Person{撒辣}，\BibleRef{创 18:12}当她被责备时，以及\Person{诺厄}的儿子，当他从自由人变作奴隶时。\BibleRef{创 9:25}\par

9. 我说这些话，并不是要抑制一切欢笑，而是要除去心灵的涣散。因为请问，你为什么还奢侈放荡？你仍背负如此重的指控，将站在可畏的审判台前，为在这里所做的一切作严格的交账。是的：我们要为我们故意犯的罪和我们不情愿犯的罪交账——因为祂说：\BibleQuote{凡在人前否认我的，我在我父前也必否认他；}\BibleRef{玛 10:33}——这样的否认确实不情愿，但它不能免于惩罚，我们也要为此交账：既为我们所知的，也为我们所不知的；有人说：\BibleQuote{我虽然不觉得自己有错，却也不能因此成义；}\BibleRef{格前 4:4}——既为我们无知所作的，也为我们明知而作的；经上说：\BibleQuote{我为他们作证，他们对天主有热心，但不按着真知识；}\BibleRef{罗 10:2}但这并不足以成为他们的借口。当他写信给格林多人时也说：\BibleQuote{我怕你们的心意，或许被蛇的诡诈诱惑，像\Person{厄娃}被诱惑一样，以致偏离了在基督内的纯朴。}\par

既然如此重大的事，你都要为此交账，你还能坐着谈笑风生、言辞俏皮，沉溺于享乐吗？\Quote{为什么？}有人会说，\Quote{如果我不这样做，而是悲伤，又有什么益处呢？}益处确实极大，甚至大到无法用言语表达。因为在世俗的法庭前，无论你如何痛哭流涕，判决之后也无法逃脱惩罚；但在这里，只要你叹息一声，就能撤销判决，获得宽恕。因此，基督向我们大谈哀恸，称哀恸的人为有福，而宣称欢笑的人为可怜。因为这里不是欢笑的剧场，我们聚集也并非为此目的，好让我们放纵无度的欢乐，而是为了叹息，并借着这叹息继承天国。然而，当你站在君王面前时，你连微微一笑都难以容忍；那么，当众天使之主居住在你内时，你难道不该战战兢兢、满怀应有的节制吗？你却反而发笑，甚至常常在祂不悦时也是如此？难道你不明白，你这样做比犯罪更惹怒祂吗？因为天主通常不会因人们犯罪而远离他们，却会因他们犯罪后不感到敬畏而远离他们。\newline{}\par

尽管如此，有些人仍愚钝到如此地步，竟在听完这些话后说：\Quote{不，愿我永远不会哭泣，而是愿天主赐我终生欢笑和嬉戏。}还有什么比这种心态更幼稚呢？因为赐予嬉戏的并非天主，而是魔鬼。至少听听那些嬉戏之人的结局是什么。经上记载：\BibleQuote{百姓坐下吃喝，起来玩耍。}索多玛的人也是如此，洪水时代的人也是如此。论到索多玛，经上同样说：\BibleQuote{她们在骄傲、丰足和安逸中放纵情欲。} \BibleRef{则 16:49} 诺厄时代的人也是如此，他们看见方舟建造了许多年，却仍活在愚昧的欢乐中，对即将来临的事毫无预见。因此，洪水袭来，将他们全部冲走，一瞬间造成了普世的沉船之灾。\newline{}\par

因此，不要向天主祈求这些你从魔鬼那里得到的事。因为天主所赐的，是痛悔谦卑的心，是清醒、自持、敬畏、充满悔悟和痛悔的心。这些才是祂的恩赐，因为我们最需要的就是这些。是的，因为一场严峻的战斗就在眼前，我们的搏斗是对抗看不见的权势，对抗\BibleQuote{邪恶的鬼神} \BibleRef{弗 6:12}，我们的战争是对抗\BibleQuote{率领者}、对抗\BibleQuote{掌权者}。只有当我们认真、清醒、彻底警醒时，我们才能抵挡住那野蛮的方阵。但如果我们欢笑嬉戏，总是轻松度日，那么甚至在战斗开始之前，我们就会被自己的松懈所打倒。\newline{}\par

10. 因此，我们不应总是欢笑、放荡、奢华，这属于那些舞台上的人，属于娼妓、为此目的打扮的男人、寄生虫和谄媚者；不属于那些蒙召进入天国的人，不属于那些登记在上城的人，不属于那些佩戴属神武器的人，而属于那些被征募在魔鬼一方的人。因为正是他，正是他，甚至把这件事变成了一门技艺，好削弱基督的士兵，软化他们热忱的锐气。为此，他在城市中建造剧场，训练那些小丑，借着他们有害的影响，使那种瘟疫降临到整个城市，劝人追随保禄命令我们逃避的那些事：\BibleQuote{愚妄的谈论和戏谑。} \BibleRef{弗 5:4} 而比这些更严重的是笑声的对象。因为当那些表演荒唐事的人说出任何亵渎或污秽的话时，许多较为轻率的人就发笑并欢喜，为他们本应投石击打的事鼓掌，并借着这娱乐为自己招致火狱。因为赞美说这种话的人，正是他们最促使人们这样说话；因此，他们更应为此类事所应得的惩罚负责。因为如果没有人观看这类事，也就没有人表演；但当他们看见你为了继续留在那里而放弃你的作坊、你的手艺、你的收入，总之放弃一切时，他们就因此更加大胆，更加努力地从事这些事。\newline{}\par

我说这话，并非为他们开脱责备，而是让你明白，主要是你为这种不法行为提供了原则和根源；你整日消耗在这些事上，亵渎地展示婚姻的神圣事物，公然嘲弄这伟大的奥秘。因为表演这些事的人固然有罪，但你比他更有罪；你命令他拿这些事开玩笑，或者更确切地说，你不仅命令他，甚至热心于此，取乐、欢笑、称赞所做的事，并以各种方式为这些魔鬼的工场增添力量。\newline{}\par

那么请问，在此之后，你用什么眼睛看家里的妻子，既然你亲眼看见她在那里受侮辱？或者，当你看到自然本身公然受辱时，你怎能不脸红，想起家里的伴侣？不，请不要对我说所做的事是演戏；因为这种演戏已经使许多人成了奸夫，颠覆了许多家庭。而最令我悲伤的是，所做的事甚至不显得邪恶，反而有掌声、喊叫和许多欢笑，当如此可耻的奸淫被犯下时。你说什么？所做的事是演戏？正是为此，他们应受一万次死亡，因为所有法律命令人逃避的事，他们却费力模仿。因为如果事情本身是恶的，那么对它的模仿也是恶的。我还没有说，那些表演奸淫场景的人制造了多少奸夫，如何使观看此类事的观众变得大胆无耻；因为没有比那忍受观看这类事的眼睛更充满淫乱和胆怯的了。\newline{}\par

你在市場上不願見到一個女子赤身露體，或者說即使在屋內也不願，還稱這樣的事為暴行。而你卻上劇院去，侮辱那男女共有的人性，玷污你自己的眼睛？不要說那被剝去衣服的是個娼妓；但本性相同，身體一樣，無論是娼妓的身體還是自由婦女的身體。如果這不算什麼錯，那麼為什麼你若在市場上見到這事，你自己會趕快走開，還會把那個行為不端的女子趕走？還是說，當我們各自獨處時，這樣的事是暴行，但當我們聚集一堂、所有人坐在一起時，就不再那麼可恥了？不，這是荒謬和羞辱，是極度瘋狂的話語；與其觀看這樣的過犯，倒不如用泥巴和污物塗滿眼睛。因為污物對眼睛的傷害，遠不及不潔的景象和觀看一個赤身露體的女子。例如，請聽起初是什麼導致了赤身露體，並閱讀這種恥辱的起因。那麼是什麼導致了赤身露體呢？我們的悖逆和魔鬼的計謀。如此，從一開始，甚至從起初，這就是他的詭計。然而他們赤身時至少感到羞恥，但你卻以此為榮，\Quote{正如宗徒所說的：\InnerQuote{以你的羞辱為光榮。}} \BibleRef{斐 3:19}\par

那麼，你的妻子從此以後將如何看待你？當你從如此邪惡的行徑歸來時，她將如何接待你？如何與你說話？在你如此公開羞辱了那共同的女性本性，並因這種景象成了娼妓的俘虜和奴隸之後？\par

如果你因聽見這些事而憂傷，我就大大感謝你，因為\BibleQuote{誰能使我快樂呢？不是那因我而憂傷的人嗎？} \BibleRef{格后 2:2} 因此，不要停止為這些事憂傷和煩惱，因為這種憂傷將成為你向善轉變的開端。為此，我也使我的言語更加有力，好藉著更深的切入，使你們脫離那陶醉你們之人的毒液；使你們恢復靈魂純潔的健康。願天主恩賜我們眾人藉著一切方法享受這健康，並達到為這些善行所預備的賞報；賴我們的主耶穌基督的恩寵和慈愛，願光榮與權能歸於祂，直到永遠。阿們。\par

\SourceRecord{Translated by George Prevost and revised by M.B. Riddle.；Nicene and Post-Nicene Fathers, First Series；Vol. 10.；Edited by Philip Schaff.；Buffalo, NY: Christian Literature Publishing Co.,；1888.；<http://www.newadvent.org/fathers/200106.htm>.}

% structure: p0001 source=h1 target=section reason=page-title

\section{Homily 7 on Matthew}

\BibleRef{玛 2:4-5}.\par

\BibleQuote{当他把所有司祭长和民间的经师都聚集起来时，便问他们基督应在何处诞生。他们回答他说：在犹太的白冷。}\par

你看，这一切如何是为了使犹太人知罪？只要他不在他们的视野中，嫉妒尚未抓住他们，他们便真实地复述有关\Emph{他}的证言；但一旦他们看见从奇迹中升起的荣耀，一种吝啬的精神就占据了他们，从此他们便出卖了真理。\par

然而，真理藉着一切被高举，甚至通过它的仇敌也越发得力。请看在这件事上，暗中预备的结果是多么奇妙、超乎期望。因为蛮族人和犹太人同时彼此更多地了解，也彼此教导。犹太人从贤士那里听说，有一颗星也在波斯地宣告了他；而贤士则从犹太人那里得知，这颗星所宣告的这个人，先知们从远古以来就已经预知了。他们询问的缘由，为双方提供了一个阐明更清楚、更完美教导的机会；真理的仇敌被迫违背自己的意愿阅读有利于真理的著作，并解释预言；尽管不是全部。他们说了白冷，以及他将出来统治以色列之后，为了谄媚国王，就没有继续添加下文。下文是什么？那就是\BibleQuote{他的根源从亘古、从太初就有。}\par

2. \Quote{但为什么，}有人会说，\Quote{如果他要从那里来，为什么他出生后住在纳匝肋，使预言模糊不清呢？} 不，他没有使它模糊，反而更加显明了。因为他的母亲常住在一个地方，而他却在另一个地方出生，这件事表明这是出于天主的安排。\par

我还要补充说，为此原因，他出生后并没有立即从那里离开，而是停留了四十天，给那些愿意探究的人提供机会，让他们仔细考察一切。因为事实上，有许多事情足以促使他们进行这样的探究，至少如果他们愿意留心的话。例如，当贤士来到时，全城轰动，连同王也一起，先知被提出来，并召集了最高权威的法庭；还有许多其他事情也在那里发生，路加都详细叙述了。比如有关亚纳、西默盎、匝加利亚、天使和牧童的事；所有这些，对于留心的人，都足以提供线索，以查明所发生的事。因为那些从波斯来的贤士都不至于不知道这地方，何况那些住在那里的人，他们更应该知道这些事。\par

他起初藉着许多奇迹显明自己，但当他们不愿看见时，他便隐藏了一段时间，之后再次从另一个更荣耀的开端被启示。因为不再是贤士，也不是星，而是天上的父在约旦河畔宣告了他；圣神也降临在他身上，将那个声音引向刚受洗者的头上；若翰用极其明畅的话语，在犹太各地呼喊，直到有人居住的和荒芜之地都充满了这种教训；还有奇迹、地、海和整个受造界的见证，都为他发出清晰的声音。但在诞生之时，只发生了恰好足以安静地指明他已来临的那些事。因此，为了避免犹太人会说：\Quote{我们不知道他何时出生，也不知道他在哪里}，所有这些涉及贤士的事件以及其他我们提到的事情，都藉着天主的眷顾发生了；这样，他们就无可推诿，没有为自己不去查究所发生的事而辩解。\par

但是也要注意预言的精确性。它没有说\BibleQuote{他将\Emph{住}在白冷}，而是说\BibleQuote{他将从那里\Emph{出来}}。所以，他的单纯诞生在那里，也是一项预言的对象。\par

然而，有些人厚颜无耻地说，这些话是指则鲁巴贝耳说的。但他们怎能正确呢？因为他的\BibleQuote{根源}当然不是\BibleQuote{从亘古、从太初就有}。\BibleRef{米 5:2} 又怎能使开头所说的\BibleQuote{你将出来}应验在他身上呢？则鲁巴贝耳不是在犹太出生，而是在巴比伦出生，因此他被称为则鲁巴贝耳，因为他的根源在那里？凡懂叙利亚话的人都知道我在说什么。\par

与所说的话一起，从这些事以来的一切时间也足以确立这见证。因为他说了什么？\BibleQuote{你在犹大的首领中决不是最小的，} 他又加上这卓越的原因说：\BibleQuote{将有一位从你那里出来。} 但除了祂以外，没有别人使那地方显赫或著名。例如：自从那诞生以来，人们从地极前来观看马槽和畜棚的所在地。先知从起初就高声预言了这一点，说：\BibleQuote{你在犹大的首领中决不是最小的；} 意思是在各支派的首领中。通过这句话，他甚至包括了耶路撒冷。但即便如此，他们也没有留心，尽管益处归于他们自己。是的，正因为这个原因，先知们在开始时很少谈论祂的尊威，而是谈论祂给他们带来的好处。因为如此，当童贞玛利亚怀着孩子时，他说：\BibleQuote{你要给祂起名叫耶稣；} \BibleRef{玛 1:21} 并给出理由说：\BibleQuote{因为祂要把自己的百姓从他们的罪恶中拯救出来。} 贤士们也没有说：\BibleQuote{天主子在哪里？} 而是说：\BibleQuote{那生下来作犹太人的王的在哪里？} 这里也没有断言：\BibleQuote{将有一位从你那里出来} 天主子，而是\BibleQuote{一位领袖，将牧养我的百姓以色列。} 因为起初必须以极其谦卑的语气与他们交谈，以免他们跌倒；并特别宣讲与他们的救恩有关的事，这样他们反而更容易被争取。无论如何，所有首先引用的、恰好在诞生时适时的见证，都没有说什么关于祂的伟大或崇高的事，也没有像神迹显现之后那些后来的见证那样；因为这些见证更清晰地谈论祂的尊威。例如，在许多神迹之后，孩子们向祂歌唱时，请听先知所说的：\BibleQuote{从婴孩和吃奶的口中，你完善了赞美。} 又说：\BibleQuote{我要观看你手指所造的诸天；} 这表示祂是宇宙的创造者。而升天之后提出的见证也显明了祂与父的平等，如此说：\BibleQuote{主对我主说：你坐在我右边。} 依撒意亚也说：\BibleQuote{那兴起统治外邦人的，外邦人要仰望祂。}\par

但他说白冷\BibleQuote{在犹大的首领中决不是最小的}是什么意思？因为不仅是在巴勒斯坦，而是在全世界，这村庄变得显赫了。为什么，到目前为止，他是在对犹太人说话；因此他也加上了：\BibleQuote{祂将牧养我的百姓以色列。} 然而祂牧养了全世界；但正如我所说，祂不愿意现在因揭示关于外邦人的话而冒犯他们。\par

但有人会说：祂怎么没有牧养犹太百姓呢？我回答：首先，这也成就了；因为在这里，他用\Quote{以色列}一词比喻性地指那些从犹太人中相信祂的人。保禄解释这一点说：\BibleQuote{因为从以色列生的不都是以色列人，} \BibleRef{罗 9:6} 而是所有因信德和恩许而生的人。如果祂没有牧养他们所有人，这是他们自己的错误和过失。因为当他们应该与贤士们一同朝拜，并荣耀天主，因为这样的时刻已经来到，除去了他们所有的罪恶（因为对他们说的不是审判或交账的话，而是一位温和良善的牧人）；他们却恰恰相反，烦恼不安，不断地制造阴谋。\par

3. \BibleQuote{那时，黑落德暗暗地叫了贤士来，仔细查问那星出现的时间。} \BibleRef{玛 2:7}\par

企图杀害那出生的婴孩——这是极端的愚蠢，不仅是疯狂；因为所说和所发生的事足以阻止他做出任何这样的企图。我的意思是，那些事件不像人所为。一颗星从天召唤贤士；贤士们长途跋涉来朝拜那躺在襁褓和马槽里的婴孩；先知们自古就预告了这一切——这些以及其他一切都不是人力所能及的；然而，这一切都没有阻止他。因为邪恶就是如此。它自相矛盾，总是企图不可能的事。请注意他完全的愚昧。如果一方面他相信预言，并认为它不可改变，那么他显然是在企图不可能的事；如果另一方面他不相信，也不指望那些话会实现，他就没有必要恐惧惊慌，也不必为此设谋。所以无论哪种方式，他的狡诈都是多余的。\par

这也出于极端的愚昧，认为贤士们会更重视他，而不是那出生的婴孩，为了这婴孩他们跋涉了如此遥远的路程。因为如果在看见祂之前，他们就如此渴望祂；在亲眼看见之后，又被预言所确认，他还指望说服他们把那幼童出卖给他吗？\par

然而，尽管有许多理由阻止他，他仍尝试了；他\Quote{私下召了贤士来，向他们询问}。因为他认为犹太人会眷顾这婴孩，而他万万想不到他们竟会疯狂到如此地步，竟愿意将他们的保护者和救主，就是那来拯救他们民族的那一位，交给祂的仇敌。为此他不仅私下召见他们，而且询问的不是婴孩的时辰，而是星辰的时辰：这样他就把追踪的目标定得远远超过实际所需。因为我想，那星必然早已出现。贤士们旅途所需时间很长。因此，为使他们在祂出生后立刻出现（祂理应在襁褓中就被朝拜，以显此事的神奇与非凡），那星早早地就显现了。倘若不是提前，而是在祂于巴勒斯坦出生那一刻，那星才在东方被他们看见，那么他们旅途耗时已久，到达时便见不到襁褓中的祂了。至于他杀害\Quote{两岁及以下的}婴孩，我们不必惊奇；因为他的愤怒与恐惧，为确保万无一失，将时间大大延长，以致无人能逃脱。\par

于是他召了他们来，说：\Quote{你们去仔细寻访那婴孩，找到了，就给我报信，好让我也去朝拜祂。}\BibleRef{玛 2:8}\par

你看见他极端愚蠢吗？你若是真心说这话，为何私下询问？但你若想谋害祂，难道你看不出，因秘密盘问，贤士们会察觉你的诡计吗？但如我所说，灵魂一旦被某种恶行俘获，就变得比任何东西更加迟钝。\par

他没有说\Quote{你们去了解关于这君王的事}，而是\Quote{关于这婴孩的事}；因为他甚至不愿用祂统治的名号来称呼祂。\par

4. 但贤士们因极大的敬畏，对此毫无察觉（因为他们万万想不到他会堕落到如此邪恶，竟试图图谋反对如此奇妙的安排）；他们离去时，对这些事毫无怀疑，而是凭自己的心意推测全人类都会如此。\par

\Quote{看，他们在东方所见的那星，在他们前头走。}\BibleRef{玛 2:9}\par

因为那星隐藏起来，只是为了让他们失去引导，从而不得不向犹太人询问，这样事情就向众人显明了。因为他们询问之后，有祂的仇敌作为报信者，那星又向他们显现。请注意这顺序何等美妙：先是那星，然后犹太民族接待他们，还有君王，他们引用预言来解释所显现的事；然后，在先知之后，又有一位天使接引他们，教导他们一切；但那时他们从耶路撒冷到白冷，是靠那星的引导，那星也从那里与他们同行；好使你由此知道，这不是普通星辰之一，因为没有一颗星辰有这种本性。它不只是移动，而是\Quote{在他们前头走}，引领他们在正午前行。\par

\Quote{但既然地点已查明，何必再需要这颗星？}有人会问：\Quote{为了也能看见那婴孩。}因为没有什么能显明祂：房屋并不显眼，祂的母亲也不荣耀或出众。因此需要那颗星，指示他们地方。所以当他们出了耶路撒冷，它又重新出现，一直走到马槽才停住。\par

奇事接连发生；因为两样都是奇事：贤士的朝拜和那星在前引领；足以吸引即使铁石心肠的人。因为如果贤士们说，他们听先知说过这些事，或有天使私下与他们交谈，他们可能不被相信；但现在，当那星在高处显现时，就连最无耻的人也哑口无言。\par

此外，那星停在婴孩上方时，再次停住了行程；这本身就比星辰的力量更大：时而隐藏，时而显现，显现后又停住。因此他们也增加了信德。为此他们也欢喜，因为他们找到了所寻求的，证明自己是真理的使者，没有空跑如此漫长的旅程；他们对基督有如此大的渴望（可以这样说）。因为那星先来停在祂的头顶，显示所生的是神；接着停在那里，引导他们朝拜祂；他们不单是蛮族人，而是他们中的智者。\par

你看，那星显现得多么合宜？因为即使在先知预言和司祭长与经师的解释之后，他们仍心系那颗星。\par

5. 愿马西昂感到羞耻，愿撒摩撒塔的保禄感到羞耻，因为他们拒绝看见那些贤士——教会的先祖——所看见的；因为我不耻这样称呼他们。愿马西昂羞惭，看见天主在肉身中被朝拜。愿保禄羞惭，看见祂被朝拜不是作为纯粹的人。至于祂在肉身中，首先由襁褓和马槽所表明；至于他们不是作为纯粹的人朝拜祂，他们在祂如此年幼时就献上适宜献给天主的礼物，以此宣告。愿犹太人也同他们一起羞惭，看见外邦人和贤士抢先，而他们甚至不愿跟随其后。因为那时所发生的事，是将来之事的预像，从起初就显示外邦人将先于他们的民族。\par

\Quote{但这是怎么回事？}有人可能会问，\Quote{不是在一开始，而是在后来，祂说了：\InnerQuote{你们去，使万民成为门徒}}？因为所发生的事是未来的预像，如我所说，是一种预先的宣告。因为自然秩序是犹太人应先到祂这里来；但因他们自愿放弃了自身的益处，秩序就颠倒了。因为即使在此事上，贤士们也不应比犹太人先来，来自如此遥远地方的人也不应比住在耶路撒冷附近的人先到，毫无听闻的人也不应阻挡在众多预言中受培育的人。但由于他们极其无知于自己的福祉，波斯人反而比耶路撒冷人先到。而这正是\Person{保禄}所说的：\Quote{天主的圣言应当先传给你们，但你们却自认为不配，看，我们转向外邦人。}\BibleRef{宗 13:46}因为即使他们先前没有听从，至少当从贤士们那里听闻时，他们本应赶紧前来；但他们不愿意。因此，当那些人沉睡时，这些人奔跑在前。\par

6. 让我们也跟随贤士们，让我们远离我们野蛮的习俗，并与之保持巨大距离，好使我们能看见基督，因为如果他们不是远离自己的国家，就会错过看见祂。让我们远离地上的事物。因为贤士们身在波斯时，只看见星；但他们离开波斯后，就看见了正义的太阳。或者不如说，若不是他们乐意从那里起身，就连星也看不见。让我们也起身吧；即令所有人都骚动，让我们奔向婴孩的房屋；即令君王、万民、暴君阻挡这条路，愿我们的渴望不消逝。因为这样我们才能彻底击退所有围攻我们的危险。因为这些人也是，除非他们看见了婴孩，就不能逃脱来自君王的危险。在看见婴孩之前，恐惧、危险和困扰从四面压迫他们；但朝拜之后，就是平静与安全；不再是星，而是天使接待他们，他们因朝拜而成了司铎；因为我们看见他们也献上了礼物。\par

因此，也要离开犹太人民、骚动的城市、嗜血的暴君、世界的浮华，奔向白冷，那里是\BibleRef{宗 13:46}属灵食粮的房屋。因为即使你是牧人，来到此地，你必看见马槽中的婴孩；即使你是君王，若不来到此地，你的紫袍也毫无益处；即使你是贤士之一，这也不妨碍你；只愿你的来临是为尊敬与朝拜，而非轻蔑天主子；只愿你怀着战栗与喜乐这样做：因为这两者可以并存于一身。\par

但要当心，不要像\Person{黑落德}那样，说：\Quote{我要来朝拜祂}，但来了之后却存心杀害祂。因为那些不相称地领受奥迹的人，正是仿效他；经上说：\BibleQuote{这样的人要得罪主的圣体圣血}\BibleRef{格前 11:27}。是的，因为他们自身怀着那因基督的王国而忧伤的暴君，就是比古时的\Person{黑落德}更邪恶的\Person{玛门}。因为他想要掌权，差遣自己的人表面上朝拜，实则朝拜时杀害。因此，让我们害怕，免得我们外表像是恳求者和朝拜者，实际上却显示出相反的事。\par

当我们朝拜时，让我们将一切从手中抛下；即使我们拥有金子，也当献给祂，而非埋藏。因为如果那些外邦人那时为尊敬而献上金子，你若不将祂所需之物献给祂，又将如何？如果那些人长途跋涉来朝拜刚诞生的祂，你却连一条小巷的长度都不愿多走，去探望病中或被束缚的祂，又有什么理由呢？然而，当他们生病或被束缚时，连我们的仇人都得我们怜悯；你对你的恩主和主的怜悯却被拒绝。他们献上金子，你几乎连面饼都不给。他们看见星就欢喜，你却看见基督亲自成了陌生人和赤身露体，却毫不动心。\par

你们中哪一个曾为基督的缘故作过如此长途的朝圣？你们这些领受了无数恩惠的人，还不如这些外邦人——或者不如说，这些比最明智的哲学家更明智的人。我为什么说长途旅程呢？不，我们的许多妇女如此娇弱，连一条街的交叉口都不愿多走，去属灵的马槽边朝拜祂，除非有骡子拉她们。而其他能行走的人，却宁愿将他们的参与托付给世俗事务的拥挤，或戏院，也不愿到这里来。然而，那些外邦人为祂的缘故完成了如此伟大的旅程，是在看见祂之前；你们在看见祂之后，却不效法他们，反而在看见祂之后离弃祂，跑去看戏子。（因为我再次触及同样的问题，正如我最近所做的。）你们看见基督躺在马槽中，却离开祂，去看舞台上的妇女。\par

7. 这些事难道不值得雷霆轰击吗？请告诉我，如果有人领你进入王宫，让你看见国王坐在宝座上，你是否真的会弃此而去看戏院？然而，在王宫里也毫无所得；但在此处，属灵的火泉从这祭台中涌出。你却离开它，跑到戏院里，去看妇女游泳，公开羞辱本性，而留下基督坐在水井旁？是的，因为现在，正如古时，祂坐在水井旁，不是与一个撒玛黎雅妇人谈论，而是与整个城。或者也许现在也只有一个撒玛黎雅妇人。因为现在也没有人与祂同在；有的只是用身体，有的连身体也没有。但祂仍不离开，仍在那里，向我们要求喝，不是水，而是圣善，因为\Quote{祂将祂的圣物赐给圣者}。因为祂从这泉源赐给我们的不是水，而是活的血；这确实是死亡的象征，却成了生命的原因。\par

但你们却舍弃血泉、那可畏之杯，径自走向魔鬼的泉源，去看妓女游泳，遭受灵魂的沉船。因为那水是淫欲的海洋，不淹没身体，却使灵魂沉船。她裸身游泳时，你注视之，便沉入淫欲的深渊。魔鬼的网罗正是如此：它使那些并非下水的人，而是坐在上面的人，比在其中翻滚的人更甚，如同昔日在海中与车马一同沉沦的\Person{法郎}，更痛苦地窒息他们。如果灵魂能被看见，我能指给你们看许多漂浮在这水上的人，如同当时埃及人的尸首。更可悲的是，他们甚至称这完全的毁灭为乐事，称这沉沦之海为愉快航行的水道。的确，人要安全渡过爱琴海或第勒尼安海，比观看这场景更容易。首先，整夜魔鬼以期盼预先占据他们的灵魂；然后，展示他们所期望的对象时，立刻束缚他们，使他们成为俘虏。不要以为因为你没有与那妓女结合，就免于罪；因为在你内心的意愿中，你已经完全行了。如果你被情欲所抓住，你便使火焰更高涨；如果你对所见的无动于衷，你便应受更重的控告，因为你成了他人的绊脚石，鼓励他们观看这些场景，又玷污了自己的眼目，连同眼目和灵魂。\par

然而，不只是一味指责，让我们也想出一种纠正的方法吧。那么这方法是什么呢？我要把你们托付给你们自己的妻子，使她们教导你们。诚然，按照 \Person{保禄} 的律法 \BibleRef{格前 14:34}，你们应当是教导者。但因那秩序被罪颠倒，身体居上，头部居下，就让我们采取这种方式吧。\par

但如果你们羞愧于让女人作你们的导师，那就逃避罪恶吧，你们很快就能登上天主赐给你们的宝座。因为你们犯罪时，圣经不仅把你们送到女人那里，甚至送到无理性的受造物那里，且是较低级的；是的，它毫不羞愧地把你们这些享有理性尊荣的人，当作门徒送到蚂蚁那里。\BibleRef{箴 6:6} 显然，这不是对圣经的控告，而是对那背叛自身高贵血统之人的控告。因此，我们也要这样做；眼下我们把你们托付给你们的妻子；但如果你们藐视她，我们将把你们送到野兽的学校，向你们指出多少飞鸟、鱼类、四足走兽和爬虫都比你们更尊贵、更贞洁。\par

如果现在你们羞愧并因这比较而脸红，那就登上你们的尊贵，逃离地狱之海和火焰之流，我是指剧场中的水池。因为这水池引向那海洋，点燃那火焰的深渊。既然\BibleQuote{凡注视妇女而对她起淫念的，已经在她心里犯了奸淫}\BibleRef{玛 5:28}，那被迫甚至见到她裸体的人，岂不成了万倍的俘虏？\Person{诺厄}时代的大洪水没有像这些游泳的女人那样，以极大的耻辱淹没所有在场的人。因为那场雨，虽然确实带来了身体的死亡，却压制了灵魂的邪恶；但这里却有相反的效果：身体存留，灵魂却被毁灭。而你们，当有关优先权的问题时，你们声称超过全世界，因为我们的城市首先以基督徒之名加冕；但在贞洁的竞争中，你们却毫不羞愧地落后于最粗野的城市。\par

8. \Quote{好吧}，有人说，\Quote{那么你要求我们做什么？去占据山林，成为隐修士吗？}这正是我叹息的原因，因为你们以为只有他们才适宜关心端庄和贞洁；然而基督的确使祂的法律为众人共通。因此，当祂说\BibleQuote{凡注视妇女而对她起淫念的}时，祂不是对独修者说的，也是对有妻子的人说的；因为事实上，那山上当时充满了各种这类人。那么，就在你心中形成那圆形剧场的情景，而憎恨这魔鬼的剧场吧。也不要谴责我话语的严厉。因为我既不\BibleQuote{禁止嫁娶}\BibleRef{弟前 4:2}，也不妨碍你们享受快乐；但我希望这事在贞洁中完成，而非羞耻、责备和无休止的指责。我没有立法要求你们去占据山林和旷野，而是要在城市中居住，做良善、体贴和贞洁的人。因为事实上，我们所有的法律对隐修士也是共通的，除了婚姻；甚至在这方面，\Person{保禄}命令我们完全与他们平等，说：\BibleQuote{因为这世界的容貌正在消逝：使有妻子的要像没有妻子一样。}\par

\Quote{所以}（他这样说）\Quote{我并没有命令你们占据山顶；我固然愿意这样，因为城市效法了索多玛所做的事；然而我并不强迫这样。你们可以安居，有房屋、有儿女、有妻子；但不要侮辱你们的妻子，不要使你们的儿女蒙羞，也不要把剧院的传染带进你们的家。}你们没有听见保禄说：\BibleQuote{丈夫对自己的身体没有主权，而是妻子，}\BibleRef{格前 7:4}，并且制定了双方共同的法律吗？但是你们，如果你们的妻子不断地往公共集会里钻，你们就严厉地责备她；而你们自己整天在公共表演上消磨时光，却不认为该受责备。是的，关于妻子的端庄，你们严格到甚至超过了必要或限度，连不可或缺的缺席都不允许她；但对自己却认为一切都可以。然而保禄不允许你们，他也给了妻子同样的权柄，因为他这样说：\Quote{丈夫当对妻子尽应尽的本分。}那么这是什么样的本分呢？当你在最重要的事上侮辱她，把她的身体交给娼妓（因为她的身体就是你的身体），当你把骚乱和战争带进家中，当你在市场上做出那些事，由你自己在家中向妻子讲说时，使她羞愧难堪，也使在场的女儿羞愧，更使你自己比她们更羞愧呢？因为你必然要么沉默，要么行为如此不雅，以至于连你的仆人也该为此受鞭打。请问您还有什么可说的呢？您这样热切地观看那些连提说都是可耻的事，宁愿选择这些连叙述都无法忍受的景象，超过其他一切景象？\par

为了不使你们负担过重，我现在暂且在这里结束我的讲道。但如果你们继续同样行径，我将把刀磨得更锋利，切口更深；我不停止，直到我驱散魔鬼的剧场，从而净化教会的集会。因为这样我们既会脱离当前的耻辱，也将收获来世生命的果实，是因我们的主耶稣基督的恩宠和慈爱，愿光荣和权能归于他，至于无穷之世。阿们。\par

\SourceRecord{Translated by George Prevost and revised by M.B. Riddle.；Nicene and Post-Nicene Fathers, First Series；Vol. 10.；Edited by Philip Schaff.；Buffalo, NY: Christian Literature Publishing Co.,；1888.；<http://www.newadvent.org/fathers/200107.htm>.}

% structure: p0001 source=h1 target=section reason=page-title

\section{《玛窦福音》第八篇讲道}

\BibleRef{玛 2:2}\par

\Quote{他们进了屋子，看见婴孩和祂的母亲玛利亚。}\par

那么，路加怎么说到祂躺在马槽里呢？这是因为，在出生时，她确实立即把祂放在那里（因为，在那征税的大集会中，不出所料，他们找不到房子；路加也以此表明，说\Quote{因为没有地方，她就把他放在那里}），但后来她抱起祂，抱在膝上。因为她一到白冷，就结束了产痛，好叫你从此也学习整个\OriginalTerm{救恩计划}{dispensation}，并知道这些事不是随意或偶然发生的，而是照某种神圣的预知和先知的次序，一切都在实现之中。\par

但什么促使他们朝拜呢？因为童贞女并不显眼，屋子也不出众，他们所见的其他任何事物都没有令人惊奇或吸引人的地方。然而他们不仅朝拜，还\Quote{打开他们的宝盒}，\Quote{献上礼物}；而且这些礼物不是献给一个人，而是献给天主。因为乳香和没药正是此意的象征。那么，是什么促使他们呢？那促使他们离家远行的事，就是那颗星，以及天主在他们心中所启迪的光照，逐步引导他们达到更完善的认识。因为，若不是这样，所有外在所见都很平凡，他们就不会显出如此大的敬意。因此，在那情景中，没有一样外在环境是伟大的，而是一个马槽、一个棚屋，和一位贫穷的母亲；这是要在你眼前，赤裸\Emph{且无掩饰}地摆出那些贤士的智慧之爱，并向你证明，他们不是作为普通人接近祂，而是作为天主和恩主。所以，他们没有被任何外在所见的事物所绊倒，反而朝拜并献上礼物；这些礼物不仅没有犹太教的粗俗（因为他们没有献上羊和牛犊），而且接近教会的自我奉献，因为他们向祂献上的是知识、服从和爱。\par

\Quote{他们在梦中得到指示，不要回到黑落德那里，就由另一条路返回自己的地方去了。}\BibleRef{玛 2:12}\par

由此也看到他们的信心，他们如何不被冒犯，而是顺服、明理；他们既不困扰，也不与自己理论说：\Quote{然而，如果这孩子是伟大的，并且有任何能力，何需逃跑和秘密撤退呢？为什么我们公开、勇敢地前来，对抗了那么多人和君王的愤怒，天使却像赶走逃亡者和逃犯一样把我们赶出城呢？}但他们既没有说，也没有想这些事。因为信心最特别的一点，就是不去寻求所命令之事的原由，而只是服从加在我们身上的诫命。\par

2. \Quote{他们离去后，看，上主的天使在梦中显现给若瑟说：\InnerQuote{起来，带着婴孩和祂的母亲，逃往埃及。}}\BibleRef{玛 2:13}\par

这里有些值得探究的问题，涉及贤士和婴孩两方面；因为如果他们甚至没有被困扰，而是以信心接受一切，那么我们就有理由考察，为什么他们和婴孩没有继续留在那里，而是贤士作为逃亡者去了波斯，祂和母亲去了埃及。但有什么办法呢？祂本应落入黑落德手中，落入后不被除掉吗？不，那样祂就不会被认为是取了肉身；\OriginalTerm{救恩计划}{Economy}的伟大就不会被相信。\par

因为，如果当这些事正在进行，许多情况在按人的方式神秘安排时，有些人竟敢说祂取得我们肉身是一个虚构；那么，如果祂以完全符合天主性和祂自己权能的方式行事，他们会陷于何等的不敬呢？\par

至于贤士，祂迅速打发他们离开，一方面派遣他们作老师前往波斯地，另一方面同时阻断君王的疯狂，好让他知道他在图谋不可能的事，并熄灭他的愤怒，放弃这徒劳的努力。因为不仅公开制伏敌人，而且轻易地欺骗他们，也是祂权能可配得的事。例如，祂在犹太人事件上也欺骗了埃及人，并且有能力公开将他们的财富转交希伯来人手中，却命令他们秘密、用计谋去做；这无疑不亚于其他奇迹，使祂在敌人眼中成为可畏之物。至少，阿市刻隆人和其余的人，当他们取了约柜，遭受打击后，就指示他们的同胞不要作战，不要反对祂，并连同其他奇迹也提出这一点，说：\Quote{你们为什么硬着心，像埃及和法郎硬着心一样？祂愚弄了他们之后，岂不随后就领出祂的子民，他们就走了吗？}他们这样说，是认为这件\Emph{新事}不亚于那些公开显出的其他神迹，足以显示祂的权能和伟大。同样的事在此次也发生了；这件事足以使暴君惊骇。因为黑落德被贤士欺骗，如此被嗤笑，他自然会有何等感受，他的呼吸几乎要停止呢？因为，即使他没有变得更好，那也不是奇妙安排这一切者的过错，而是黑落德疯狂的过度，甚至不屈服于那些足以说服他并阻止他行恶的事，反而继续前进，要因如此愚蠢而受到更严厉的惩罚。\par

3. 但或许有人会说：为何将婴孩送往埃及？首先，福音作者本人已提及原因，说：\BibleQuote{为应验经上所言：我从埃及召回了我的儿子。}同时，从那时起，世界就已预闻美好希望的肇始。也就是说，既然巴比伦和埃及，这两个世上最广大的地区，都曾为不虔敬的烈焰所焚燬，祂从一开始便示意要予以纠正和改善，并借此引导人们期盼祂同样惠及整个世界。为此，祂遣贤士往其一，而另一地则由祂亲自偕同母亲造访。\par

此外，除我所说之外，我们还可从中习得另一项道理，对培养我们的真正自制力不无裨益。是怎样的道理呢？即从一开始就应期待试探和阴谋。请看，从祂的襁褓之时起，事情便是如此。你看，祂一出生，先是暴君发怒，继而逃亡，远走他方；祂的母亲虽无罪，却被放逐至异邦之地。你听到这些，若自负有资格从事任何属灵事务，却见自己遭受不治之症、经历无数危险，切莫大为困扰，不要说：\Quote{这究竟是怎么回事？我本当受加冕、受称颂、因遵行主的诫命而荣耀显赫才对。}相反，你当以这个榜样为鉴，勇敢地承担一切，须知这正是所有属灵事物的秩序：试探总是与它们同命运。至少你看，这不仅发生在婴孩的母亲身上，也发生在那些异邦人身上：他们如同逃亡者一般秘密离去；而她，从未跨过自家门槛，却因这奇妙的诞生和她的属灵产痛，被迫承受如此漫长的痛苦旅程。\par

又见一奇事：巴勒斯坦图谋不轨，埃及却收留并保护那被谋害的对象。原来，不仅先祖们的事例中出现了\OriginalTerm{预像}{types}，在我主自己身上也是如此。无疑，在许多实例中，祂当时的作为都是对未来之事的预言性宣告，例如关于驴驹的实例就是如此。\par

4. 天使既已显现，并未与玛利亚交谈，而是与若瑟说话；他说了什么？\Quote{起来，带着婴孩和祂的母亲。}此处，他不再说\Quote{你的妻子}，而是说\Quote{祂的母亲}。因为生产既已完成，疑虑也已消除，丈夫得到了安抚，从此天使公开说话，不再称那孩子或妻子为若瑟的，而是说：\Quote{带着婴孩和祂的母亲，逃往埃及去。}他又提及逃亡的原因：\Quote{因为黑落德将要搜寻这婴孩，要杀害祂。}\par

若瑟听了这些话，并未动怒，也没有说：\Quote{这事难以理解：你方才不是说过，祂要拯救祂的子民吗？如今祂连自己都救不了，反倒必须逃跑，远离家乡，长久流离。事实与应许正好相反。}不，他并没有说任何这些话（因为那人是有信德的）；他也不追问何时可以返回，尽管天使这样不确定地说：\Quote{留在那里，直到我通知你。}然而即便如此，他也没有战栗，而是顺服听从，欣然经历了一切试炼。\par

这是因为满有慈爱的主以艰辛与甘甜相间，这正是祂对待一切圣者的方式：既不让危险连绵不断，也不让慰藉永无止境，而是将义人的人生交织成苦乐参半。此处祂也正是如此而行：请看，若瑟见童贞女怀孕，这使他不安至极，因为他怀疑这女子犯了奸淫。但天使立刻出现，消除了他的猜疑，驱散了他的恐惧；及至见婴孩降生，他获得了最大的喜乐。然而这喜乐过后，又发生了不小的危险：全城骚动，暴君疯狂搜寻那初生的婴孩。但这忧患之后，又有另一种喜乐接踵而至：异星和贤士的朝拜。在这喜乐之后，又有恐惧和危险：\Quote{因为黑落德将要搜寻这婴孩，要杀害祂。}祂必须像凡人一样逃跑、隐蔽；此时行奇事的时机尚不成熟。因为倘若祂自婴孩时期就显奇迹，就不会被当作人了。\par

为此，圣殿也并非立即造成；而是有正常的受孕、九个月的时间、产痛、分娩、哺乳，以及如此长时间的缄默，并等待成人的适当年龄；这样，祂\OriginalTerm{救恩计划}{Economy}的奥秘才能为众人所接纳。\par

但有人会说：\Quote{那么，为何一开始就行这些征兆呢？}为祂母亲的缘故；为若瑟和即将离世的西默盎的缘故；为牧羊人和贤士的缘故；也为犹太人的缘故。因为他们若愿意留心观察所发生的事，也能从这些事为将来的事情获得不小的助益。\par

若是先知们没有提及关于贤士的事，不要为此困扰；因为他们既没有预言所有的事，也没有对一切保持沉默。因为如果毫无预兆地看见这些事发生，自然会引发极大的惊奇和困扰；但如果一切都事先告知，又会使听者懈怠，并且让福音作者没有可补充的余地。\par

5. 倘若犹太人对预言提出疑问，说：\Quote{我从埃及召叫了我的儿子，}这句话是针对他们自己说的；我们要告诉他们，这是预言的规律：在许多情况下，为某一类人所说的话，却在另一类人身上应验；就如关于\Person{西默盎}和\Person{肋未}所说的话：\Quote{我要使他们分居，}祂说，\Quote{在\Place{雅各伯}，散居在\Place{以色列}}。\BibleRef{创 49:7}但这并没有在他们身上应验，而是在他们的后裔身上；\Person{诺厄}关于\Person{客纳罕}的说话，又在客纳罕的后裔基贝红人身上应验。关于\Person{雅各伯}，也可看到这样应验：因为那些祝福说：\Quote{你要做你兄弟的主人，你父亲的儿子要向你叩拜}，\BibleRef{创 27:19}在他身上并没有应验（如何能应验呢？因为他战战兢兢，屡次向兄弟下拜？\BibleRef{创 33:3}），却在他的后裔身上应验了。同样的话也可用在这件事上。因为谁可称为更真实的天主子？是那个拜金牛、依附巴耳培敖尔、将自己的儿子祭献给魔鬼的人？还是那位按本性为子、尊崇生祂的那一位的人？所以，若这个人没有来，那预言就不会得到应有的应验。值得留意的是，福音作者也用\Quote{好使预言应验。}这句话暗示同一件事，暗示若祂没有来，就不会应验。\par

这也使\Person{童贞玛利亚}获得非凡的光荣和尊贵：那曾是以色列全体人民在赞美方面所特有的恩赐，从此她也得以拥有。我的意思是，他们以离开\Place{埃及}为荣，并以此夸耀（这其实是先知暗示的，当他说：\Quote{我岂不是从\Place{卡帕多细亚}领来了外邦人，从深渊领来了亚述人吗？}），祂也使这一殊荣同样属于\Person{童贞玛利亚}。\par

然而，不如说，子民和族长下到那里，又从那地上来，共同完成了祂这次回来的预象，而这预象是\Emph{祂的}回来。因为正如他们下到\Place{埃及}是为躲避饥荒而死，祂也是为躲避阴谋而死。但他们在到达后暂时脱离了饥荒，而这个人，当祂下去时，却因脚踏其地而圣化了整个土地。\par

至少可以看出，在祂受辱的当中，如何显露出祂天主性的标记。首先，天使说：\Quote{逃往\Place{埃及}去。}却没有应许与他们同行，无论是下去还是回来；暗示他们有一位伟大的同行者，即那已诞生的婴孩；祂一出现就改变了一切，并使祂的仇敌在各方面为这救恩计划服务。于是，\Person{贤士}和外邦人离开祖传的迷信，前来朝拜；\Person{奥古斯都}通过人口普查的法令，为\Place{白冷}的诞生服务；\Place{埃及}接待并保护了被驱逐、遭谋害的祂，从而获得了与祂联合的初步推动；以致日后当她听到宗徒们传扬祂时，至少可以以此夸耀：她是最先接待祂的。然而，这特权本来只属于\Place{巴勒斯坦}；但后者（\Place{埃及}）比前者更为热切。\par

6. 如今，你若来到\Place{埃及}的旷野，你会看到这旷野变得比任何乐园更美，有成千上万的天使般的人群，无数殉道者团体，成群结队的贞女，而魔鬼的一切暴虐都被摧毁，基督的王国灿烂辉煌。诗人之母、贤士、术士都不过是愚蠢老妇的虚构，但真正的、堪配天国的哲学，却是渔夫们所传扬的。正因如此，连同他们在教义上的极其精确，他们也通过生活表现出极大的严谨。因为他们舍弃了自己的一切，对世界被钉在十字架上，又进一步向前奔跑，用自己身体的劳动来供养有需要的人。他们不因禁食和守夜而认为白天可以闲散；而是用夜晚诵念圣歌和守夜，白天祈祷，同时亲手劳作，效法宗徒的热忱。因为当整个世界都看着他，为了供养有需要的人时，他既占用工场，又操持手艺，而且如此忙碌以至于夜不安眠；那么，他们说，我们这些居住在旷野、与城市的纷扰无关的人，岂不更应当用这安闲的时光从事灵修工作吗？\par

那么，我们大家，无论是富足的还是贫穷的，都应当惭愧：那些人除了身体和双手之外一无所有，却奋力向前，热切地从中寻找供养穷人的方法；而我们拥有无尽的财富，却连自己的多余之物都不肯为这些目的动分毫。请问，那时我们有什么理由？又有什么借口？\par

再者，请思考，从前这些\Place{埃及}人如何既贪婪又贪食，还有其他恶习。因为以色列人还记得的肉锅（\BibleRef{出 16:3}），那里是肚腹的极大专制。然而，他们有了愿意的心，就改变了：他们从基督那里得了火，立刻启航驶向天堂；尽管他们比其余人类更热烈、更倔强，无论在愤怒还是肉欲方面，他们却在温良上效法无形体的天使，并在其余属于克己的离情去欲方面也如此。\par

7. 如今，若有人曾在乡间，他便知道我所言。但若他从未进入过那些帐幕，就让他想起那一位直到如今仍在众人口中的——那一位在宗徒之后，由\Place{埃及}所生的人——有福的伟大的圣\Person{安当}；并让他自我省思：\Quote{此人亦与法老同生于一国，然而他并未因此受损，反倒获得天主赐予的神视，并展现出基督的法律所要求的生活。}任何人若读了记载他生平的那本书，便会完全明了此事，其中也将看到许多预言。我指的是他关于受\Person{亚略}错误感染之人的预言，以及他陈述这些错误将造成的祸患；天主甚至在那时已将它们显示给他，并将一切将要发生的事在他眼前描绘出来。此事（在其余事中）尤其有助于证明真理，即外部的异端中无一人能提及如此人物。但为不依赖我们得知此事，请仔细研读那本书中所记载的，你将完全知晓一切，并从中学习到许多克己之道。\par

我给出此劝告：不仅要阅读书中所写，也要效法它，且莫以地域、教育或祖先的邪恶为借口。因为若我们留意自身，这一切都不会成为我们的阻碍，因为连\Person{亚巴郎}也有一个不虔敬的父亲\BibleRef{苏 24:2}，但他并未继承他的邪恶；\Person{希则克雅}的父亲是\Person{阿哈次}，然而他仍成为天主所喜爱的人。同样，\Person{若瑟}在\Place{埃及}中间，以节德的冠冕装饰自己；三位青年在\Place{巴比伦}和宫廷中，面对如\Place{西巴里斯}般的筵席，显示了最高的克己；\Person{梅瑟}在\Place{埃及}，\Person{保禄}在普世中，但这一切都没有成为任何人德行赛跑中的障碍。\par

让我们记住这一切，摒弃这些无用的借口和托词，专心致力于德行事业所需的劳苦。因为这样，我们既会从天主那里吸引更多的恩宠，又说服祂助我们奋战，并获享永恒的福祉；愿天主恩赐我们众人皆能达到，赖我们的主\Person{耶稣基督}的恩宠和爱人之心，愿光荣与胜利归于祂，世世无穷。亚孟。\par

\SourceRecord{Translated by George Prevost and revised by M.B. Riddle.；Nicene and Post-Nicene Fathers, First Series；Vol. 10.；Edited by Philip Schaff.；Buffalo, NY: Christian Literature Publishing Co.,；1888.；<http://www.newadvent.org/fathers/200108.htm>.}

% structure: p0001 source=h1 target=section reason=page-title

\section{\WorkTitle{玛窦福音讲道集 第九篇}}

\BibleRef{玛 2:16}\par

\BibleQuote{那时黑落德，见自己受了贤士的愚弄，就大发忿怒。}\par

然而，这实在不是动怒的时候，而是应当恐惧和敬畏；他本该明白自己是在企图做不可能的事。但他并未被阻止。因为当灵魂变得麻木不仁、无药可救时，它便不会接受天主所施的任何药剂。请看此人，他紧追之前的努力，在一件谋杀上又增添了多件，疯狂地冲向任何地方。因为被这愤怒和嫉妒如同被某种恶魔驱使，他什么都不顾及，甚至向本性本身发怒，并将对愚弄了他的贤士的怒气发泄在无辜的婴孩身上：当时在\Place{巴勒斯坦}做了一件与\Place{埃及}曾发生的事相似的事。因为经上说，他\BibleQuote{差人将\Place{白冷城}及其四境所有两岁及两岁以下的婴儿杀死，照着他由贤士们细心探询的时期。}\par

这里请仔细听我。因为许多人关于这些婴孩说了许多极其空洞的话，事件的进程被指责为不义，有些人以较温和的方式表达他们的困惑，另一些人则更加放肆和疯狂。为了使我们能解脱这些人的疯狂和那些人的困惑，请允许我略微讨论一下这个话题。显然，如果他们的控诉是这些婴孩被杀，那么他们也应当指责杀害看守\Person{伯多禄}的士兵一事。\BibleRef{宗 12:19}因为在这里，当幼童逃走时，其他婴孩代替被搜寻的祂被屠杀；同样在那时，\Person{伯多禄}被天使从监狱和锁链中救出后，那个与这位暴君同名且同样性情的人，当他寻找他而未找到时，就杀了看守他的士兵来替代他。\par

\Quote{这是什么？}可能有人说，\Quote{这不是解答，反而加深了我们的困惑。}我也知道这个，正是为此我提出所有这些情况，为的是对所有情况提出同一个解答。那么这些事的解答是什么？或者我们能给出什么合理的解释？那就是\Person{基督}并非他们被杀的原因，而是国王的残酷；同样，\Person{伯多禄}也不是那些人的原因，而是\Person{黑落德}的疯狂。因为如果他看到墙被打破，或门被推倒，他或许有理由指责看守宗徒的士兵疏忽；但现在，当一切如常，门大敞，锁链仍缚在看守他们的人手上（事实上他们是与他绑在一起的），他本可以从这些事推断（也就是说，如果他严格履行法官职责的话），这件事不是人的力量或技巧，而是某种神圣的、显奇迹的力量；他本应崇拜做这些事的那位，而不是与哨兵作战。因为天主所做的一切都是为了行善，使看守者不致暴露，反而藉着他们引导国王走向真理。但如果他证明自己麻木不仁，那么对于那善用一切事物来行善的灵魂的巧匠来说，病人的不服从又算什么呢？\par

2. \Quote{是的，}有人说，\Quote{你使\Person{黑落德}完全失去了借口，证明了他嗜血成性；但你还没有解决关于所发生之事的不义的问题。因为他若行事不义，天主为什么容许这事呢？}对此我们该说什么呢？那就是我在教堂里、在市场上、在每一处不断说的话，也是我希望你们牢记的，因为它对我们来说是一种规则，适用于所有这样的困惑。那么我们的规则是什么，我们的话是什么？就是：虽然有许多人作恶，却没有一个人是受害的。为了不让这谜语太困扰你们，我也赶紧加上解释。我的意思是，我们从任何人那里不义地遭受的，要么用于消除我们的罪过——天主将这错误归于我们，要么用于赏报。\par

2. \Quote{是的，}有人说，\Quote{你使\Person{黑落德}完全失去了借口，证明了他嗜血成性；但你还没有解决关于所发生之事的不义的问题。因为他若行事不义，天主为什么容许这事呢？}对此我们该说什么呢？那就是我在教堂里、在市场上、在每一处不断说的话，也是我希望你们牢记的，因为它对我们来说是一种规则，适用于所有这样的困惑。那么我们的规则是什么，我们的话是什么？就是：虽然有许多人作恶，却没有一个人是受害的。为了不让这谜语太困扰你们，我也赶紧加上解释。我的意思是，我们从任何人那里不义地遭受的，要么用于消除我们的罪过——天主将这错误归于我们，要么用于赏报。\newline{}ewline{} 即使你被贤士们欺骗，这与那些无辜的婴孩有什么关系呢？\par

为了使我的话更清楚，让我们用比喻来论证。例如：假设某个仆人欠他主人很多钱，然后这个仆人受到不义之人的虐待，被抢了一些财物。如果主人本有能力阻止掠夺者和作恶者，却没有归还那财产，而是将所夺去的算作仆人欠他的债，那么仆人受损了吗？决不会。但若主人甚至偿还他更多，难道他不是甚至比失去的获得更多吗？我想人人都明白这一点。\par

如今，我们对于自身的苦难也当如此推算。因为就事实而论，鉴于我们所可能无辜遭受的苦难，我们或得赦免罪过，或得领受更荣耀的冠冕（若我们的罪不那么深重的话）。请听\Person{保禄}如何论及那犯奸淫的人说：\BibleQuote{将这样的人交给撒殚，败坏他的肉体，为使他的灵魂得救。} \BibleRef{格前 5:5} \Quote{但这算什么？}你或许会说，\Quote{因为所论的乃是受他人伤害的人，而非受师长纠正的人。}我可以回答：并无区别；因为问题在于：受苦是否对受苦者构成侮辱？然而，为更切近所问之点，请忆起\Person{达味}：某次当他看见\Person{史米}攻击他、践踏他的苦难、无尽地辱骂他时，他的将领们想要杀死\Person{史米}，他却严加禁止，说：\Quote{由他咒骂吧，但愿上主垂看我的卑辱，并在今日因这咒骂而赐我好报。}在\BibleBook{}{圣咏}中他也歌咏说：\Quote{求你看顾我的仇敌，他们何其众多，且以无理的仇恨恨我，}并\Quote{赦免我的一切罪过。}此外，\Person{拉匝禄}也因同样的原由获得赦免，因他在此生已饱受无数苦难。因此，那些受冤屈的人并未受冤屈，若他们能高尚地忍受一切所受的痛苦；反而他们获益更丰，无论所受的是来自天主的打击，还是来自魔鬼的鞭笞。\newline{}\par

3. \Quote{但这些孩童有何种罪过，}或有人说，\Quote{以致能藉死亡消除罪过呢？因为论及成年人，他们犯了许多疏忽，如此说或许有理；但这些过早夭折的孩童，他们藉所受的苦难消除了何种罪过？}你们岂未听见我说，即使没有罪过，在此生受苦的人将来仍有赏报？那么，这些幼童因这般缘故被杀，被迅速带入那无浪的港口，他们又受了什么损害？\Quote{因为，}你说，\Quote{他们若活着，本可成就许多伟大的善行。}为此，天主预先为他们储积了不小的赏报，即为这般缘故而结束他们的生命。此外，若这些孩童将成为伟人，天主必不容许他们提前被夺走。因为，即使那些终将生活在持续邪恶中的人，天主尚且以如此宽容忍耐他们，那么，若祂预知这些孩童将成就任何大事，岂会容许他们如此被杀害？\newline{}\par

这些是我们所能给出的理由；然而这些并非全部，还有比这些更奥秘的理由，只有那亲自安排万事的主完全知晓。因此，让我们将对此事更圆满的理解交托给祂，并致力于下文所述，在别人的灾难中学习高尚地忍受一切。的确，当时\Place{白冷}所发生的悲剧绝非微小：孩童们从母亲的怀中夺走，被拉去遭受这不义的屠杀。\newline{}\par

若你仍心怯，不能在这些事上自持，请思量那胆敢行此一切者的结局，稍作振作。因他很快便为这些事受惩罚，为如此可憎的行为付出了应有的代价：以惨痛的死亡结束生命，比他当时所敢施行的更为可怜；此外还遭受了无数额外的灾祸，你们可从\Person{若瑟夫}的记载中知晓这些事。但为避免使我们的讲论冗长而打断连贯性，我们认为无需在此叙述那段记载。\newline{}\par

4. \BibleQuote{于是应验了\Person{耶肋米亚}先知所说的话：在\Place{辣玛}听到了声音，\Person{辣黑耳}痛哭她的子女，不愿受安慰，因为他们不在了。} \BibleRef{耶 31:15} \BibleRef{玛 2:17}\newline{}\par

如此，叙述了这些事——那暴烈而不义的杀戮，极其残酷和无法无天——使听者充满恐惧之后，他又安慰他们，说：这一切发生并非因天主无力阻止，也非祂不知情；乃是祂既知道，又藉先知大声预言了。因此，不要困扰，也不要沮丧，当仰望祂那不可言喻的眷顾；从祂所行和所容许的事上，人最清楚看见这一点。祂在另一处向门徒讲论时也暗示了这一点：即当祂预言审判台、处刑、世上的战争以及不休止的战斗时，为振奋他们的精神并安慰他们，祂说：\BibleQuote{两只麻雀不是卖一个铜钱吗？但你们在天之父若不许可，一只也不会掉在地上。} \BibleRef{玛 10:29} 祂说这些话，表明没有一件事是在祂不知情下发生的；但虽然祂知晓一切，却并不在一切事上都行动。\Quote{因此，不要困扰，}祂说，\Quote{也不要惊慌。}因为，若祂知道你们所受的苦，并有能力阻止，那么显然，祂正是出于对你们的眷顾和关怀才不阻止。在我们的考验中，我们也当铭记这一点，从中获得极大的安慰。\newline{}\par

但或许有人会说：\Person{辣黑耳}与\Place{白冷}有何关系？因为经文说：\BibleQuote{\Person{辣黑耳}痛哭她的子女。}\Place{辣玛}与\Person{辣黑耳}又有何关系？\Person{辣黑耳}是\Person{本雅明}的母亲，她死后被葬在此地附近的跑马场。因此，坟墓既近，且那地属于她的幼子\Person{本雅明}（因为\Place{辣玛}是\Person{本雅明}支派的），祂便先以支派之首，次以埋葬之地，自然地称呼那些被屠杀的幼童。为表明那创伤是无可医治且残酷的，祂说：\BibleQuote{她不愿受安慰，因为他们不在了。}\newline{}\par

这里再次教导我们先前提到的事：当天主所应许的与所发生的事相反时，切莫惊慌。请看，当祂来为百姓——不，是为世界——带来救恩时，祂的开端是怎样的呢？祂的母亲首先逃亡；祂的出生地陷入无法补救的灾祸中，发生了一场最惨烈的屠杀，到处是哀哭、大恸和悲号。但不要烦恼；因为祂惯于用相反的方式完成祂的救恩计划，由此给我们显明了祂的大能。\par

祂也这样引导自己的门徒，预备他们尽一切本分，以相反的方式成就万事，使奇迹更加伟大。他们尽管遭受鞭打、迫害和无穷的恐惧，却就这样胜过了那些鞭打和迫害他们的人。\par

5. \Quote{黑落德死后，有一位上主的天使在梦中显现给若瑟说：起来，带着婴孩和祂的母亲，到以色列地去。} \BibleRef{玛 2:19}\par

祂不再说\Quote{逃跑}，而是说\Quote{去吧}。你看见了吗？试炼之后，又是安息；安息之后，又是危险？他确实脱离了放逐，回到本国；看见那杀害婴儿的凶手被处死；但当他踏上本国土地时，又遇到先前危险残余——暴君的儿子还活着，并且作了国王。\par

但是，当般雀·比拉多任总督时，阿尔赫劳怎能统治犹太呢？黑落德刚死不久，王国尚未分成多份；他刚去世，儿子暂时继承王国\Quote{代替他的父亲黑落德}；他的兄弟也名叫黑落德，这就是为什么福音作者补充说：\Quote{代替\Emph{他的父亲}黑落德}。\par

但也许有人说：\Quote{若他因阿尔赫劳而惧怕定居犹太，他也有理由因黑落德而恐惧加里肋亚。}我回答：由于他更换了地方，整个事件从此被掩盖；因为全部攻击是针对\Quote{白冷及其四境}。因此，既然屠杀已经发生，青年阿尔赫劳没有别的想法，只以为一切已经结束，在所杀的人中，那所寻的已被除灭。再者，他亲眼目睹父亲如此结局，今后对进一步行动和继续那邪恶的道路更加谨慎。\par

因此若瑟来到纳匝肋，部分是为了躲避危险，部分也是由于喜爱住在故乡。为了给他更多的勇气，天使也赐给他关于此事的启示。然而，路加并没有说他因天主的警告而到那里，而是说他们满了取洁礼后，就回到了纳匝肋。\BibleRef{路 2:39} 那么我们可以说什么呢？路加讲这些话时，是在叙述下埃及之前的事。因为祂不会在取洁之前带他们下去，以免违反法律，而是等待她取洁完毕，再到纳匝肋，然后才下到埃及。之后，在他们回来时，祂吩咐他们到纳匝肋去。但在此之前，他们没有受到天主的警告要去那里，而是出于对故乡的思念，自行这样做了。因为他们上去（耶路撒冷）别无原因，只是为了登记，连住宿的地方都没有，所以当他们完成了此行的目的后，就下到了纳匝肋。\par

6. 这里我们看到天使为何让他们安心，并将他们送回故乡的原因。不仅如此，他还加上了一句预言，说：\Quote{这是为应验先知所说的话：祂将被称为纳匝肋人。} \BibleRef{玛 2:23}\par

那是什么样的先知说这话呢？不要好奇，也不要多事。因为许多先知著作已经失传；这一点可以从编年史（历代志）的历史看出。由于他们疏忽，不断陷入不敬虔，有些书卷任其消失，另一些被他们亲手焚烧或割碎。后一件事实耶肋米亚记述了；前一件事实，那编写列王纪第四卷的人说道，过了很久，申命纪才勉强被找到，埋在某处而丢失了。但如果当时没有外邦人在那里，他们尚且这样遗弃自己的书卷，何况外邦人入侵之后呢？至于先知确实预言了这件事，宗徒们自己在许多地方称祂为纳匝肋人。\BibleRef{耶 36:23}\par

也许有人说：\Quote{这岂不是给关于白冷的预言投下阴影吗？}绝不然；相反，这件事必能大大激动人，唤醒他们去查考关于祂所说的话。例如，纳塔乃耳也开始查问祂说：\Quote{纳匝肋能出什么好事吗？} \BibleRef{若 1:46} 因为那地方微不足道；不，不仅是那地方，就连整个加里肋亚地区也是如此。因此法利塞人说：\Quote{你们查考一下，就知道从加里肋亚不会出先知。}然而，祂不以从那里被命名而羞耻，表明祂不需要任何属人的事物；祂的门徒也是从加里肋亚拣选的；处处断绝那些懈怠之人的借口，并给我们凭据，表明如果我们修德行善，就不需要外在的事物。为此，祂甚至不为自己选择一所房屋，因为祂说：\Quote{人子没有枕头的地方。} \BibleRef{玛 8:20}\par

7. 為何你要以你的祖國自豪，當我命令你成為整個世界的陌生人？（祂這樣說）；當你被允許成為這樣的人，以至於整個宇宙都不配你？因為這些事情如此可鄙，甚至在希臘的哲學家們中間也不被認為值得任何考慮，而是被稱為\Emph{外在之物}，並佔據最低的位置。\par

\Quote{可是保祿}，有人說，\Quote{允許他們，如此說：\InnerQuote{就揀選而論，他們因祖宗的緣故是蒙愛的。}} 但請告訴我，他何時、論及何事、以及向誰說話？是對那些外邦人，他們因信德而自高自大，並高抬自己反對猶太人，從而更加使他們斷絕；為要壓制一方的驕傲，贏得另一方，並徹底激發他們相同的熱忱。因為當他論及那些高尚偉大的人物時，請聽他如何說：\BibleQuote{說這些話的人，明明表示他們在尋求一個祖國；他們若想念所離開的家鄉，還有可以回去的機會。但他們現在渴慕另一個、更美的家鄉。} \BibleRef{希 11:14} 又說：\BibleQuote{這些人都是存著信心死的，並沒有得著所應許的，卻從遠處看見，且歡喜迎接。}  若翰也對那些來到他那裡的人說：\BibleQuote{不要自己心裡說：我們有亞巴郎為父。} \BibleRef{瑪 3:9} 保祿又說：\BibleQuote{因為從以色列生的，不都是以色列人；也不因為是亞巴郎的後裔，就都作他的兒女。} \BibleRef{羅 9:6} 請告訴我，撒慕爾的兒子們因父親的高貴得了什麼益處？他們並沒有繼承父親的德行。梅瑟的兒子們，沒有效法他的成全，又有什麼益處？因此他們也沒有繼承統治權；他們雖認他為父，但百姓的治理卻轉給了別人，即藉德行成為他兒子的人。弟茂德是希臘父親所生，對他有什麼損害？反之，諾厄的兒子因父親的德行得了什麼益處？他竟成了奴隸而非自由人。你看，父親的高貴在代求方面對兒女有何等渺小的幫助？因為含的邪惡性情勝過了自然的律，不僅使他失去了因父親而有的高貴，也失去了自由的身份。厄撒烏呢？他不是依撒格的兒子嗎？他的父親不是為他作朋友嗎？是的，他的父親也努力並渴望他分享祝福，他自己也為此做了所吩咐的一切。然而，因為他乖謬，這些事沒有一樣使他獲益；雖然他按出生是長子，又有父親為他做一切，但因天主沒有與他同在，他失去了一切。\par

但我何必提人？猶太人曾是天主的兒子，卻因這高貴出身毫無所得。如今，若一個人成了天主的兒子，卻未能顯出與這高貴出身相稱的卓越，甚至會受到更重的懲罰；那你為何要向我提出遠近祖先的高貴？因為不只舊約之下，就是新約之下，也可見這法則成立。因為經上說：\BibleQuote{凡接受祂的，祂賜他們權能，成為天主的子女。} \BibleRef{若 1:12} 然而保祿斷言這些子女中有許多並未因他們的父親得益；他說：\BibleQuote{如果你們受割損，基督對你們就毫無益處。} \BibleRef{迦 5:2} 如果基督對那些不願謹慎自守的人毫無幫助，人又怎能為他們辯護？\par

8. 所以，我們不要因高貴出身或財富而自豪，反而要輕視那些這樣思想的人；也不要因貧窮而沮喪。但我們要追求那由善工構成的財富；讓我們逃避那引人行惡的貧窮，因那財主也正是因此成了貧窮的 \BibleRef{路 16:24}。所以他連一滴水也不得支配，儘管他多多懇求。然而，在我們中間，誰能貧窮到\Emph{連一點安慰的水也缺乏}呢？沒有人。因為即使是極度飢餓的人，也可能得到一滴水的安慰；不僅是一滴水，還有更豐富的舒暢。那財主卻不是這樣，他竟貧窮到如此地步：更可悲的是，他無法從任何來源緩解他的貧窮。那麼，我們為何還要貪戀財富，既然它們不能帶我們進入天堂？\par

请告诉我，如果世间某位君王说：富有的人不可能在朝中显赫或享受任何尊荣，难道你们不会轻蔑地抛弃自己的财富吗？既然如此，若有人将我们从地上王宫那样的尊荣中驱逐出去，他们理应受尽轻视；但当天国之王每日高声呼喊说：\Quote{他们难以踏足那神圣的门槛；}难道我们不该放弃一切，舍弃财产，以便能坦然进入天国吗？我们如此费心积攒那些阻碍我们进入天国的事物，不仅把它们藏在箱柜里，甚至埋在地下，而不愿托付给上天守护，这还配得上什么顾念呢？如今你们所做的，正像农夫得了种子要播种沃土，却弃田不顾，把所有种子埋进坑里，结果自己无法享用，种子也腐烂浪费，一无所成。当我们指责他们这些事时，他们常提出的借口是什么？他们说：知道一切安放在家中是极大的安慰。其实，不知其安放才是安慰。因为即使你不怕饥荒，但为了这储积，其他更可怕的事必然令你恐惧：死亡、战争、针对你的阴谋。若饥荒临到，百姓因饥饿所迫，就会执戈攻击你们的房屋。或者，你们这样做，首先是把饥荒引进我们的城市，其次是为自己的房屋挖出比饥荒更可怕的深渊。因为因饥荒而速死的人，我不知有多少；实际上，在多方有许多方法可以缓解那灾祸。但为了财产和财富及其相关的追求，我可以指出许多人走向了毁灭，有的暗中，有的公开。许多这样的实例充斥在大路、法庭和市场上。但我何必提大路、法庭和市场？看哪，连海洋都充满了他们的血。看来，这暴政不仅在地上得势，也在海洋上铺张扬厉。有人为黄金航海，另有人为同一目的被刺；同一暴虐的力量使一人成为商人，另一人成为杀人凶手。\par

还有什么比玛门更不可靠呢？因为为了它，人旅行、冒险、被杀。经上说：\Quote{谁怜悯被蛇咬的巫士？} \BibleRef{德 12:13} 我们既知道它残忍的暴政，就该逃离那奴役，消灭那痛苦的欲望。有人说：\Quote{但这怎么可能？} 只要引入另一种渴望——对天国的渴望。因为渴望天国的人必嘲笑贪财；成为基督奴仆的人不再是玛门的奴隶，而是它的主人；因为玛门惯于追随逃离它的人，而逃避追求它的人。它所尊重的，不是追求它的人，而是蔑视它的人；它最嘲笑的，正是那些渴望它的人；它不仅嘲笑他们，还用无数锁链捆绑他们。\par

因此，但愿我们，即使迟了，也要解开这些痛苦的锁链。为什么要把你理性的灵魂捆绑在无理性的物质上，那无数恶事的母亲呢？但，何其荒谬！我们口头上与它作战，它却以行动向我们开战，带领我们、驱策我们四处奔走，像对待花钱买来的、该受鞭挞的奴隶一样凌辱我们；还有什么比这更可耻、更不光彩的呢？\par

再者，如果我们不能胜过无感官的物质形态，又怎能胜过无形体的权势呢？如果我们不轻视卑贱的泥土和低微的石头，又怎能制伏那些率领者和掌权者呢？我们又如何修习节制？我的意思是，如果银子的光芒令我们目眩神迷，我们何时才能匆匆走过一张俊美的面孔呢？事实上，有些人如此被这暴政奴役，甚至见到金币的显现就被某种程度地感动，开玩笑说，眼睛看见金币反而更舒服。但不论你是谁，不可开这样的玩笑；因为没有什么比贪恋这些事物更伤害眼睛，无论是肉身的眼睛还是灵魂的眼睛。例如，正是这痛苦的欲望熄灭了那些童贞女的灯，把她们赶出了婚房。这（如你所说）\Quote{有益于眼睛}的视觉，竟不让可怜的犹大倾听主的声音，反而领他到绞架，使他肚腹崩裂；之后又领他下地狱。\par

还有什么比这更无法无天？更可怕？我不是指财富的实质，而是指对它不合时宜、疯狂的欲望。它甚至滴着人血，目光充满杀意，比任何野兽更凶猛，撕裂落在它面前的人，更糟的是，它甚至不让那些被撕裂者感到自己如此被撕裂。按理说，受此待遇的人应当向过路者伸手呼救，但他们竟为这种撕裂肉体而感恩，还有什么比这更可悲呢？\par

让我们铭记这一切，逃离这不治之症；医治它造成的伤口，远离这样的瘟疫；好使我们在此世能过安稳无忧的生活，并达到将来的宝藏；\Emph{愿天主恩赐我们众人达到那宝藏}，因我们的主耶稣基督的恩宠和慈爱，祂与圣父及圣神，一同享受荣耀、威能、尊崇，自今至永远，世世无尽。阿们。\par

\SourceRecord{Translated by George Prevost and revised by M.B. Riddle.；Nicene and Post-Nicene Fathers, First Series；Vol. 10.；Edited by Philip Schaff.；Buffalo, NY: Christian Literature Publishing Co.,；1888.；<http://www.newadvent.org/fathers/200109.htm>.}

% structure: p0001 source=h1 target=section reason=page-title

\section{玛窦福音讲道第十篇}

\BibleRef{玛 3:1-2}\par

\Quote{那时，洗者若翰出现在犹太旷野宣讲说：\InnerQuote{你们悔改罢！因为天国临近了。}}\par

何以说\Quote{那时}？并非指祂童年时，那时祂到了纳匝肋，而是三十年之后，若翰才出现，路加也为此作证。那么为何说\Quote{那时}？圣经常惯用这种说法，不仅指紧接着发生的事，也指多年后要发生的事。例如，当门徒们到橄榄山上坐在祂面前，问及祂的来临和耶路撒冷的陷落时，祂的回答涵盖了许多情事；\BibleRef{玛 24:3}而你们知道那几个时期之间相隔甚久。我的意思是，在祂预言了母城的颠覆、完成了那部分的讲论，并要过渡到末世论时，祂插入了\Quote{\Emph{那时}这些事也要发生；}\BibleRef{玛 24:23}并非用\Emph{那时}一词把时间拉近，而是指示这些事将要发生的那段时间。祂现在也做了类似的事，说\Quote{那时}。这不是指紧接着的日子，而是指祂将要叙述的那些事发生的时间。\par

\Quote{但为何过了三十年，耶稣才来领受圣洗？}有人或许会说。在此圣洗之后，祂将从此废除法律；因此直到这个能容纳所有罪过的年龄，祂继续遵守全部法律，免得有人以为祂自己不能遵守才废除。因为各种情欲并非在所有年龄都同样侵袭我们；在生命的初期，多有轻率和胆怯；随后到来时，享乐更为强烈；再之后，则是财富的欲望。为此，祂等待成年期圆满，并在整个时期遵守法律，然后才来领受圣洗，将此作为完全遵守所有其他诫命之后的追加行为。\par

为了证明这是祂在法律所规定的一切善行中的最后一项，请听祂自己的话：\Quote{因为这样才合乎我们全部的义务，尽诸般的义。}\BibleRef{玛 3:15}祂的意思是这样：我们已履行了法律的一切义务，没有违反一条诫命。既然只剩下这一项，也必须加上，这样我们才能\Quote{尽诸般的义}。祂在这里把\Quote{义}称为完全遵守所有诫命。\par

2. 由此显然可知，基督为此缘故来领受圣洗。但为何为祂预备了这个圣洗？因为匝加利亚的儿子并非出于自己，而是受天主推动——路加也如此声明，当他说：\Quote{天主的话传到了他，}\BibleRef{路 3:2}即祂的命令。若翰自己也说：\Quote{那派遣我来以水施圣洗的，对我说：你看见圣神像鸽子一样降下，停在谁身上，谁就是那要以圣神施圣洗的人。}\BibleRef{若 1:33}那么，他为何被派来施圣洗？洗者若翰又向我们讲明，说：\Quote{我不认识祂，但为使祂显示给以色列，因此我来以水施圣洗。}\BibleRef{若 1:31}\par

如果这是唯一的原因，那么路加如何说：\Quote{他来到约旦河一带地方，宣讲悔改的圣洗，为得罪之赦？}\BibleRef{路 3:3}但此圣洗并无赦罪之恩，因为这项恩赐属于后来所施的圣洗；在此圣洗中\Quote{我们与祂同葬}，我们的旧人也与祂同钉十字架，而在十字架之前，任何地方都找不到赦罪；因为赦罪处处归功于祂的血。保禄也说：\Quote{但你们已经洗净了，已经圣化了，}不是藉着若翰的圣洗，而是\Quote{因主耶稣基督之名，并因我们天主的圣神。}\BibleRef{格前 6:11}他在别处也说：\Quote{若翰固然宣讲了悔改的圣洗，}（他没有说\Quote{赦罪的}）\Quote{叫百姓信那在他以后要来的。}\BibleRef{宗 19:4}因为当祭献尚未举行，圣神尚未降临，罪过尚未除去，仇恨尚未消除，咒诅尚未摧毁时，赦罪如何能发生？\par

那么，\Quote{为得罪之赦}是什么意思？\par

犹太人是愚昧的，从未感受到自己的罪，尽管他们应当为最严重的恶行负责，却在各方面自义；这比任何事更导致了他们的毁灭，使他们远离信仰。例如，保禄本人曾指责他们说：\Quote{因为他们不认识天主的义，而企图建立自己的义，就不服从天主的义。}\BibleRef{罗 10:3}又说：\Quote{那么，我们说什么呢？那追求义的外邦人，反而得到了义，但追求法律之义的以色列，却没有达到法律之义。为什么？因为他们不凭信德，而凭行为去追求。}\par

由此可见，这正是他们祸患的根源，若翰来了，所做的无非是使他们认识到自己的罪。他的服装本身，即悔改和认罪的标记，也表明这点。他所宣讲的也指示这点，因为他所说的无非是\Quote{结出与悔改相称的果实。}\BibleRef{玛 3:8}既然他们不谴责自己的罪（如保禄也解释的），这使他们远离基督；而认识自己的罪会使他们渴望寻求救主，并希望得到赦罪；因此若翰来了，是要劝导他们悔改，不是为了让他们受罚，而是为了让他们因悔改变得谦卑，谴责自己，从而急忙去领受赦罪。\par

但是让我们看看他是如何精确地表达的；在他说了\BibleQuote{他来到犹太旷野宣讲悔改的洗礼}之后，他又加上\BibleQuote{为赦罪}，仿佛是说：为此他劝诫他们承认并悔改自己的罪，不是要他们受罚，而是使他们更容易领受随后的赦免。因为如果他们不先谴责自己，就不能寻求祂的恩宠；不寻求，就不能获得赦免。\par

这样，那洗礼就为这洗礼预备了道路；所以他（若翰）也说，他们\BibleQuote{当信那在他以后要来的}\BibleRef{宗 19:4}，连同前面所提到的，说明了他（耶稣）受洗的另一原因。因为，若翰去各家走动，拉着基督的手，说\Quote{信这个人}，比那有福的声音被发出，以及所有那些事在众人面前和眼前举行，都远不如。\par

因此，祂前来受洗。因为事实上，施洗者的声望和这事本身的意旨，吸引了全城，召唤他们到约旦河去，成了一场盛大的景象。\par

因此，当他们来到时，他（若翰）也使他们谦卑，劝他们不要自高自大；向他们表明，若不悔改，他们将面临最大的灾祸，并要离开他们的祖先和所有对他们的夸耀，迎接那将要来的祂。\par

因为事实上，关于基督的事直到那时还是隐藏的，许多人以为祂已经死了，由于发生在白冷的大屠杀。虽然祂在十二岁时显示了自己，但祂很快又再次隐藏了自己。为此，需要那辉煌的开场和更崇高的开始。所以，那时他（若翰）才第一次用清楚的声音宣告了犹太人从未听过的事，既不是从先知，也不是从任何其他人；他提到了天堂和那里的国度，不再说任何关于地上的事。\par

但在此处，他（若翰）所说的国度是指祂（基督）的第一次和末次来临。\par

3. \Quote{但这对犹太人有什么意义呢？}有人可能会说，\Quote{因为他们甚至不明白你说什么。}\Quote{正是为此，}他说，\Quote{我才这样说，好让他们因我的话隐晦而被激发，进而寻求我所宣讲的祂。}事实上，当人们靠近时，他这样以美好的希望激励他们，以至于连许多税吏和兵士都询问他们该做什么，该如何引导自己的生活；这标志他们从此摆脱一切世俗之事，并瞩目于其他更大的目标，预感到将要来的事。是的，因为当时所有的景象和言语，都引领他们进入崇高的思想。\par

想一想，例如，看到一个人在三十年后从旷野出来，他是大司祭的儿子，从未知道人类的一般需要，各方面都令人尊敬，并有依撒意亚与他同在，这是何等伟大的事。因为依撒意亚也显现出来，为他作证，说：\Quote{这就是我曾说要求的那位，他要在整个旷野中大声呼喊、宣讲。}因为先知们对这事的热忱是如此之大，以至于他们不仅早就宣告了他们的主，也宣告了那要侍奉祂的人；他们不仅提到他，还提到他居住的地点，他到来时要教授的教义的方式，以及他产生的良好效果。\par

至少，请看先知和洗者若翰是如何基于同样的观念，尽管用语不同。\par

因此，先知说他要来说：\BibleQuote{预备主的道，修直他的路。}\BibleRef{依 40:3}他自己来时则说：\BibleQuote{结出与悔改相称的果实。}\BibleRef{玛 3:8}这与\Quote{预备主的道}是相应的。你看到吗，无论是先知的言语，还是他自己的宣讲，都只显明了这一件事：他来了，来铺路和预备，并非赐予恩赐即赦免，而是及时调整那些将要接受万有之天主的灵魂。\par

但是路加表达得更为详尽：不仅没有重复开场白就继续下去，而且还记下了全部预言。他说：\BibleQuote{凡山谷都要填满，一切山岳丘陵都要铲平，弯曲的要修直，崎岖的要铺平；凡有血肉的，都要看见天主的救恩。} \BibleRef{路 3:5} 你看见先知如何以他的话预示了这一切吗？群众的聚集、事物的更好转变、所传之道的平易、以及所发生的一切的首要原因，尽管他以比喻的方式表达，因为他实际上是在说预言。因此，当他说：\BibleQuote{凡山谷都要填满，一切山岳丘陵都要铲平，崎岖的要修直；} 他是在指示卑微者的高升、倔强者的降卑，以及法律的严苛转变为信仰的平易。因为他说，不再有辛劳和苦工，而是恩宠和罪过的赦免，为救恩提供了极大的便利。接着，他指出了这些事的原因，说：\BibleQuote{凡有血肉的，都要看见天主的救恩；} 不再只是犹太人和皈依者，而是整个大地和海洋，以及全人类。因为藉着\Quote{弯曲的事}，他预指了我们整个腐朽的生活——税吏、娼妓、强盗、术士，以及所有原先悖逆、后来却行走在正道的人；正如主自己也曾说：\BibleQuote{税吏和娼妓要在你们之前进入天主的国，} \BibleRef{玛 21:31} 因为他们相信了。先知又以别的话宣告了同一件事，这样说道：\BibleQuote{那时豺狼和羔羊将一起牧放。} \BibleRef{依 11:6} 因为正如此处他以山岳和山谷来意指性格的不一致被融合为同一的节制之平坦，同样在那里，他也以野兽的性格来指明人们的不同性情，并说到它们被联结在同一的虔诚和谐之中。此处也如前所述，指出了原因。那原因是：\Quote{必有那一位兴起统治外邦人，外邦人将在祂内寄望；} 这与此处他所说的：\BibleQuote{凡有血肉的，都要看见天主的救恩，} 大致相同，处处表明这些福音的权能和知识将倾注到地极，使人类从野蛮的气质和凶暴的性情转变为极其温和与柔顺。\par

\BibleQuote{这若翰身穿骆驼毛的衣服，腰束皮带。} \BibleRef{玛 3:4}\par

请注意，先知们预言了一些事，另一些则留给福音作者。因此玛窦既记载了预言，又加上自己的部分，甚至认为谈论义人的服饰也不算多余。\par

因为确实，在人的形体上见到如此严峻的克苦，是一件奇妙而不可思议的事；这件事特别吸引了犹太人，因为他们在他身上看到了伟大的厄里亚，并因他们当时所见而想起那位有福之人的回忆；或者不如说，甚至更加惊奇。因为厄里亚曾在城市和房屋中被抚养长大，而若翰从襁褓起就完全居住在旷野中。因为那将要消除一切古代苦难（如劳苦、诅咒、忧愁、汗水）之主的先驱，他自己也应当具有这种恩赐的某些标志，并立即超越那判决。因此，他不耕地，不开犁沟，也不汗流满面才得糊口，而是他的餐桌迅速得到供应，他的衣服比餐桌更容易获得，他的住所比衣服更不麻烦。因为他既不需要屋顶，也不需要床、桌子或任何其他东西，而是在这血肉之躯中展示了一种天使般的生活。因此，他的衣服是用骆驼毛做的，好藉着他的服装教导人们脱离一切属人的事物，与大地毫无共同之处，并急速返回他们最初的尊贵，即亚当在需要衣服或长袍之前的状况。因为那装束所象征的，无非是王权与悔改。\par

且不要对我说：\Quote{他居住在旷野，从何处得到骆驼毛的衣服和皮带呢？} 因为如果你对此吹毛求疵，你还会追问更多的事：在冬季和夏季的酷热中，他如何继续留在旷野，而且身体娇嫩，年龄尚幼？他幼小的肉体如何承受如此剧变的天气、如此不寻常的饮食，以及旷野中所有其他的艰辛？\par

如今希腊的哲学家在哪里？他们随意且徒然地效仿犬儒学派的厚颜无耻（因为被关在木桶里，然后陷入如此放荡，有什么益处呢？）他们佩戴戒指、杯盏，有男仆女仆，还伴有大量排场，陷入两个极端。但若翰并非如此；他住在旷野如同住在天堂，显示出全备的节制之严格。从那里，他如同一位从天堂来的天使，下到城市，成为虔诚的斗士、战胜世界的冠冕得胜者，以及一位配得天国的哲学家的哲学。而这些都是发生在罪尚未除去、法律尚未废止、死亡尚未被束缚、铜门尚未被打破之时，而古老的政体仍然有效之时。\par

这就是高贵且完全警醒的灵魂的本性，因为它处处向前跃进，并超越为他所设定的界限；保禄对于新约政体也是如此。\par

但有人可能会问：为什么他的衣裳上束着腰带？这是古人的习俗，当时的人尚未进入这种柔软宽松的衣着。例如，\Person{伯多禄}在\BibleRef{若 21:7}中似乎也\Quote{束着腰}，\Person{保禄}也是如此；因为经上说：\Quote{那拥有这条腰带的人。}（\BibleRef{宗 21:11}）\Person{厄里亚}（\BibleRef{列下 1:8}）也是这样穿戴，所有圣人都如此，因为他们不断地劳作、奔波，或旅行，或处理其他必要事务；不仅为此，也为了践踏一切装饰，实践一切刻苦。因此，基督将这件事称为德行最大的赞美，说：\Quote{你们出去看什么？一个穿细软衣服的人吗？看！那穿细软衣服的人是在王宫里。}（\BibleRef{玛 11:8}）\par

但若连他——如此纯洁、比天更荣耀、超过众先知、凡妇人所生的没有比他更大的、且拥有如此大胆直言的人——都这样刻苦自己，如此蔑视一切放荡的安逸，训练自己过这种艰苦的生活；那么，我们这些领受了如此大恩、背负着无数罪债的人，却连他补赎的最小部分都没有表现出来，反而饮酒作乐、满身香水、打扮得如同舞台上的妓女，处处软化自己，使自己成为魔鬼易得的猎物，这还有什么借口可言呢？\par

5. \Quote{那时，全犹太和耶路撒冷，以及约但河一带的人，都出来到他那里，承认自己的罪，在他那里受洗。}（\BibleRef{玛 3:5}）\par

你看到先知来临的巨大力量了吗？他如何激动了全体人民？如何引导他们反省自己的罪？因为看到一个人形，却显露出如此的事，用如此大的自由言论，起来谴责所有的人如同孩童，而且他脸上闪耀着极大的恩宠，这确实是值得惊奇的。此外，在长时间没有先知之后，先知的显现也增加了他们的惊愕，因为这项恩赐已经中断了许久，如今又回到他们中间。他所宣讲的道理也是陌生而不寻常的。他们听到的不是那些习以为常的事，诸如战争、争战、地上的胜利、饥荒、瘟疫、巴比伦人、波斯人、城池的攻取，以及他们所熟悉的其他事情，而是关于天堂、天上的国度、地狱的惩罚。我还可以补充说，这正是为什么尽管那些在旷野中叛乱的人，就是跟随\Person{犹达斯}和\Person{特乌达}的人（\BibleRef{宗 5:36}），在不久之前全都被杀，他们却并没有因此退缩而不去那里。因为他召唤他们并非为了统治、叛乱或革命等目标，而是为了亲手引领他们到天上的国度。因此，他也没有把他们留在旷野，带着他们到处走，而是给他们施洗，教导他们克己的规则，然后打发他们回去；以此教导他们要轻视一切地上的事物，提升自己向往未来的事，并每日前进。\par

6. 因此，让我们也效法这个人，放弃奢华和醉酒，转向节制的生活。因为现在确实是告解的时刻，无论是对未领洗者还是已受洗者：对前者，使他们悔改后能领受神圣奥迹；对后者，使他们洗去领洗后的玷污，以清洁的良心走近祭台。因此，让我们放弃这种柔软娇纵的生活方式。因为不可能同时做补赎又过奢华的生活。让\Person{若翰}通过他的衣服、食物和居所来教导你这点。那么，你会说，你要求我们实践这样的节制吗？我不要求，但我劝告并推荐。但如果这对你不可能，至少让我们在城市中也显露出悔改，因为审判确已临近。即使它还很遥远，我们也不应该胆大妄为，因为对每一个被召唤的人来说，他生命的终结实质上就是世界的终结。但说它已经临近，请听\Person{保禄}说：\Quote{黑夜深了，白日将至。}（\BibleRef{罗 13:12}）又说：\Quote{那要来的就来，决不迟延。}（\BibleRef{希 10:37}）\par

因为宣告那日的征兆现在已经完成。祂说：\Quote{这天国的福音，必传遍天下，对万民作见证；然后终结才来到。}（\BibleRef{玛 24:14}）请仔细注意所说的话。祂没有说\Quote{当所有人都信了}，而是\Quote{当它传给了所有人}。为此，祂也说\Quote{对万民作见证}，以表明祂并不等所有人都相信了才来。因为\Quote{作见证}这句话的意思是\Quote{为控告}、\Quote{为责备}、\Quote{为定那不信之人的罪}。\par

但我们听见并看见这些事，却昏昏欲睡，如同做梦，沉沦在麻木中，仿佛处于最深的黑夜。因为现在的事，无论是顺境还是逆境，都不比梦更好。因此，我恳求你们现在终于要醒悟，转向另一边，朝向正义的太阳。因为没有人能在睡觉时看见太阳，也不能以它的光芒之美愉悦双眼；无论他看到什么，都如同在梦中看见。因此，我们需要大量的补赎和许多眼泪，因为我们迷失时处于麻木状态，而且我们的罪大且无可推诿。我并非说谎，听我讲道的人中大部分可以作证。然而，虽然它们无可推诿，让我们悔改，我们必将获得冠冕。\par

7. 然而我所说的悔改，不仅是要离弃从前的恶行，而且要显出比那些更大的善行。因为他说：\BibleQuote{结出与悔改相称的果实来}（\BibleRef{玛 3:8}）。但我们怎样结出呢？如果我们做相反的事：例如，你曾强夺他人的财物吗？从此以后，连你自己的也要施舍。你长期犯奸淫吗？在规定的日子里，连你的妻子也要禁欲，操练节制。你曾侮辱和殴打过往的人吗？从此以后，祝福那侮辱你的人，善待那打击你的人。因为拔出箭矢并不足以使我们痊愈，还必须把药敷在伤口上。你从前放纵无度、醉酒吗？要守斋，并注意饮水，以摧毁那在你内已长大的恶。你曾以不洁的眼目看别人的美貌吗？从此以后，连妇人也不要看，这样你才更安全。因为经上说：\BibleQuote{离恶行善}；又说：\BibleQuote{使你的舌头停止作恶，使你的嘴唇不说诡诈的话。}\Quote{但请告诉我那善事。}\BibleQuote{寻求和平，并追随它：}我所说的和平不仅是与人的和平，也是与天主的和平。他说得对：\Quote{追随}她：因为她被赶走、被驱逐；她离开了大地，到天上居住。但我们如果除去骄傲和自夸，以及一切阻碍她的事物，并追随这有节制和简朴的生活，就能把她带回来。因为没有什么比忿怒和暴怒更令人痛苦。这使人既自高又卑贱，前者令人可笑，后者令人可憎，同时带来相反的恶习——骄傲和谄媚。但我们若能斩断这情欲的贪求，就既能彻底谦卑，又能安全地高升。因为在我们身体中，一切疾病也来自过度；当身体的元素离开它们适当的界限，超过适度时，就会产生无数疾病和可怕的死亡。在灵魂方面，人们也可以看到类似的情形。\par

8. 因此，让我们割断过度，饮服节制的有益良药，安于适当的平衡，并留心我们的祈祷。即使我们未蒙垂允，也要坚持，好使我们能蒙垂允；若蒙垂允，则因已蒙垂允而坚持。因为祂绝不愿延迟施予，而是藉此延迟安排我们坚持。祂这样做，是比我们更知道什么对我们有益，祂爱我们比生育我们的父母更热烈。让这两个思虑成为我们的护身符，在每一次遇到的恐惧中，我们对自己吟诵它们，好能抑制沮丧，并在一切事上光荣为我们行事并安排一切的那一位——天主。\par

因为这样，我们便能轻易击退一切敌人的阴谋，并获得不朽的冠冕：藉着我们的主耶稣基督的恩宠与爱人之情，愿光荣、权能和尊荣，偕同圣神，归于圣父，从今时至永远，世世无尽。阿们。\par

\SourceRecord{Translated by George Prevost and revised by M.B. Riddle.；Nicene and Post-Nicene Fathers, First Series；Vol. 10.；Edited by Philip Schaff.；Buffalo, NY: Christian Literature Publishing Co.,；1888.；<http://www.newadvent.org/fathers/200110.htm>.}

% structure: p0001 source=h1 target=section reason=page-title

\section{\WorkTitle{玛窦福音讲道集 第十一讲}}

\Emph{\BibleRef{玛 3:7}}\par

\BibleQuote{但若翰看见许多法利塞人和撒杜塞人来受他的洗，就对他们说：\InnerQuote{毒蛇的种类！谁指教你们逃避那将要来到的忿怒呢？}}\par

那么，基督怎么说他们不信若翰呢？\BibleRef{路 20:5} 因为若他们不接受他所预许的那一位，这便是不信。他们自以为尊重先知和立法者，但主却说他们并不尊重他们，因为他们没有接受先知所预言的那一位。\BibleQuote{因为如果你们信了梅瑟，你们也会信我。}\BibleRef{若 5:46} 此后，基督又问他们：\BibleQuote{若翰的洗礼是从哪里来的？}\BibleRef{玛 21:25} 他们却说：\BibleQuote{如果我们说\InnerQuote{从人来的}，我们害怕群众；如果我们说\InnerQuote{从天上来}，他会对我们说：\InnerQuote{那么你们为什么没有相信他？}}\par

由此可见，他们确实来了并受了洗，却未坚持所宣讲的信仰。若翰也指出了他们的邪恶，因为他们派人到洗者若翰那里说：\BibleQuote{你是厄里亚吗？你是基督吗？}因此他又加了一句：\BibleQuote{那些被派来的是法利塞人。}\BibleRef{若 1:24}\par

\Quote{那么，群众不也是同样想法吗？}有人会说。不，群众是出于单纯的心而有此怀疑，但法利塞人却是想要抓住他。因为既然大家都承认基督是出自达味的村庄，而若翰却是肋未支派的人，他们就借这个问题设下陷阱，若他说出任何类似的话，他们便可立刻捉拿他。实际上，这由下文便可看出；因为当他不承认他们所期待的任何事情时，他们仍然抓住他，说：\BibleQuote{你若不是基督，为什么施洗呢？}\BibleRef{若 1:25}\par

为要让你确信法利塞人来的心意与群众不同，请听福音作者如何说明这一点：关于群众，他说\BibleQuote{他们来到他那里，承认自己的罪，在约旦河里受他的洗；}\BibleRef{玛 3:6} 但关于法利塞人，却不再如此，而是说\BibleQuote{他看见许多法利塞人和撒杜塞人来受洗，就对他们说：\InnerQuote{毒蛇的种类！谁指教你们逃避那将要来到的忿怒呢？}} 何等的伟大！他竟对那些渴慕先知之血、心性不比蛇好的人们说话！他以如此坦白的话贬低他们自己和他们的祖先！\par

2. \Quote{是的，}有人说；\Quote{他说话足够直白，但问题在于这种直白是否有道理。因为他没有看见他们犯罪，而是看见他们正在改变；因此他们不应受责备，反而应受赞扬和赞许，因为他们离开了城市和房屋，急忙来听他的讲道。}\par

那么我们该说什么？他并非着眼于眼前正在发生的事，而是由于天主的启示，他知道他们心中的秘密。由于他们以祖先为荣，这很可能成为他们灭亡的原因，并使他们陷入疏忽的状态，他便切断他们骄傲的根源。为此，依撒意亚也称他们为\BibleQuote{索多玛的官长}和\BibleQuote{哈摩辣的人民，}\BibleRef{依 1:10} 另一位先知说：\BibleQuote{你们对我岂不是如雇士人的儿子？}\BibleRef{亚 9:7} 所有先知都使他们脱离这种想法，制服他们的骄傲，这骄傲曾给他们带来无数灾祸。\par

\Quote{你将会说，先知们这样做是自然的，因为他们看见他们犯罪；但在此情况下，看见他们服从他，他为什么、出于什么原因这样做？}为使他们的心更加柔软。\par

但如果仔细分析他的话，就会发现他也用称赞来调和了责备。他说这些话，是因为他惊奇他们竟然能够（虽然为时已晚）做似乎是他们几乎不可能的事。你看，他的责备更像是要拉拢他们、促使他们觉醒。因为当他显得惊讶时，就暗示了先前的邪恶是巨大的，而他们的悔改是奇妙的、出乎意料的。因此他说：\Quote{发生了什么事，使得那些人的子孙、受如此恶劣教养的他们竟然悔改了？这巨大的变化从何而来？谁软化了他们心灵的刚硬？谁矫正了那不治之症？}\par

请注意，他一开始就惊醒了他们，以关于地狱的话作为基础。他没有谈论通常的话题：\Quote{谁指教你们逃避战争、蛮族入侵、被掳、饥荒、瘟疫？}而是关于另一种刑罚，从未向他们显明的刑罚，他敲响了序曲，这样说：\BibleQuote{谁指教你们逃避那将要来到的忿怒？}\par

他称他们为\BibleQuote{毒蛇的种类}也甚恰当。因为据说那种动物会吞噬正在分娩的母亲，咬穿她的肚子，从而来到光明中；这些人所做的也是如此，他们是\BibleQuote{杀父母的}\BibleRef{弟前 1:9}，亲手杀害了自己的导师。\par

3. 然而，他并不止于责备，也引入了劝告。他说：\BibleQuote{你们要结出与悔改相称的果实。}\par

因为仅仅逃避邪恶还不够，还必须显出伟大的德行。不要让我看到那种自相矛盾却又常见的情形：你们节制了一小会儿，又回到同样的邪恶中去。因为我们来到这里的对象，与先前的先知们不同。如今，一切都已改变，且更加崇高，因为审判者即将亲自来临，那正是君王本身、王国的主宰，祂引领我们走向更大的节制，召唤我们升天，将我们向上提至那些居所。为此，我也阐明了关于地狱的教义，因为美好的事物和痛苦都是永恒的。因此，不要停留在你们现在的状态，也不要提出那些习以为常的托辞：\Person{亚巴郎}、\Person{依撒格}、\Person{雅各伯}，你们祖先的高贵血统。\par

他说这些话，并非禁止他们说自己是那些圣人的后裔，而是禁止他们因忽视灵魂的德行而依恃这一点；他一面公开揭示他们心中的想法，一面预言未来的事。因为此后，他们果然会说：\BibleQuote{我们有亚巴郎为父，从未做过任何人的奴隶。}\BibleRef{若 8:33}既然这一点最使他们骄傲自大并毁灭他们，他首先将其压制下去。\par

请看，他在向族长表示敬意时，如何将此与对这些事的纠正相结合。即是说，他说了\BibleQuote{不要心想：我们有亚巴郎为父，}之后，并没有说\Quote{因为族长不能给你们带来任何益处，}而是以更温和、更可接纳的方式暗示了同样的事，说道：\par

\BibleQuote{天主能从这些石头给亚巴郎兴起子孙来。}\BibleRef{玛 3:9}\par

有些人说，他这些话是指着外邦人说的，以\Emph{石头}作比喻；但我说，这句话还有另一层含义。这是什么含义呢？他说：\Quote{不要以为，如果你们灭亡，就会使族长绝嗣。}不是这样，根本不是这样。因为在天主，既能从石头中造出人来，也能使他们与那关系相连；因为在起初也是如此做的。因为当孩童从那个硬化了的子宫中出来时，这就像是从石头中生人一样。\par

先知也正是暗示这一点，当他说：\BibleQuote{你们要望向那坚硬的磐石，你们是从那里被凿出来的；望向那坑穴的洞口，你们是从那里被挖出来的：你们要望向你们的父亲亚巴郎，望向那生你们的撒辣。}\BibleRef{依 51:1}请看，他提醒他们这段预言，表明如果起初祂使他成为父亲，其奇妙如同从石头中造出他来一样，那么现在这也可能发生。而且请看，他如何既恐吓他们，又切断他们的托词：他没有说\Quote{祂已经兴起了，}以免他们对自己绝望，而是说祂\Quote{能兴起；}他也没有说\Quote{祂能从石头中造出人来，}而是说了一件更伟大的事：\Quote{亚巴郎的亲属和子孙。}\par

你看见他如何适时地使他们脱离对肉身之事的虚妄想象，以及对他们祖先的依赖，从而使他们将自己救恩的希望寄托于自己的悔改和节制？你看见他如何通过抛弃血统关系，引进那凭信德而来的关系？\par

4. 那么请注意，他在随后的话中如何进一步增加他们的恐惧，加剧他们令人痛苦的畏惧。\par

在说了\BibleQuote{天主能从这些石头给亚巴郎兴起子孙来}之后，他又加上：\BibleQuote{现在斧子已放在树根上，}这无疑使他的话令人惊恐。因为正如他因自己的生活方式而拥有很大的直言不讳，他们也因长期荒芜而需要他严厉的斥责。因为\Quote{我为什么说}（他的话如此）你们将要失去与族长的关系，并要看见别人——甚至那些从石头中出来的——被引进你们的尊位？不，你们的刑罚不仅达到这一点，而且惩罚还要继续发展。\BibleQuote{因为现在，}他说，\BibleQuote{斧子已放在树根上。}没有比他话语的这番转折更可怕的事了。因为不再是\BibleQuote{飞镰}，也不是\BibleQuote{拆掉篱笆}，也不是\BibleQuote{践踏葡萄园}；\BibleRef{依 5:5}而是一把极其锋利的斧子，更糟的是，它已到了门口。因为他们不断不信先知，常说：\BibleQuote{上主的日子在哪里？}以及\BibleQuote{愿以色列圣者的计谋来临，好让我们知道，}\BibleRef{依 5:19}由于他们的话要过许多年才应验，为了消除他们从这一点获得的鼓励，他把恐怖放在他们附近。他通过说\Quote{现在}和放在\Quote{树根}上表明了这一点。\Quote{因为现在中间没有间隔，}他说，\Quote{它已放在树根上。}他没有说\Quote{树枝上}，也没有说\Quote{果实上}，而是\Quote{树根上}。这表示，如果他们疏忽大意，就要承受无法治愈的恐怖，连补救的希望也没有。因为现在来的不是仆人，像祂以前的那几位，而是万有的主，带着祂猛烈而极其有效的惩罚临到他们身上。\par

然而，儘管他再次使他們驚恐，他卻不讓他們陷入絕望；正如先前他沒有說\Quote{他已興起了}，而是說\Quote{天主能從這些石頭中給亞巴郎興起子孫來}（既警醒又安慰他們）；同樣，此處他也沒有說\Quote{它已觸及了根}，而是說\Quote{它已被放在根上，並且已臨近，顯示毫無延遲。}然而，即使他已將它帶得如此之近，他卻使砍伐取決於你們。因為若你們悔改並成為更好的人，這斧頭將離去而一無所成；但若你們繼續故我，他會將樹連根拔起。因此，請注意，它既沒有從根上移開，也沒有在近處砍伐：前者是為使你們不致懈怠，後者是為讓你們知道，即使在短時間內也能悔改而得救。為此，他也從各個方面增加他們的恐懼，徹底喚醒並催促他們悔改。如此，首先是他們脫離了祖先；其次，引進他人取而代之；最後，這些恐怖已臨門，必將遭受無可救藥的禍患（他藉著根和斧頭同時表明了這兩點），這足以徹底喚醒那些即使極其懈怠的人，使他們充滿憂慮。我可以補充說，保祿也提出了同樣的道理，當他說：\Quote{主將在普世完成一句話。}\BibleRef{羅 9:28}\par

但不要害怕；或者說，要害怕，卻不要絕望。因為你們仍有改變的希望；判決並非完全絕對，斧頭也沒有來砍伐（否則它既已臨近根，為何不砍？）；而是特意藉著這恐懼使你們變得更好，預備你們結出果實。為此，他補充說：\Quote{因此，凡不結好果子的樹，就砍下來，丟在火裡。}\BibleRef{瑪 3:10} 藉著\Quote{凡}這個詞，他再次剝奪了他們因高貴出身而享有的特權；他說：\Quote{是的，即使你是亞巴郎的後裔，即使你有成千上萬的族長可以列舉，你若結不出果實，只會遭受雙倍的懲罰。}\par

這些話使稅吏也感到驚恐，兵士的心靈被他震撼，既沒有使他們陷入絕望，卻擺脫了所有的安全感。因為在他所說的話中，除了恐懼，也包含著許多鼓勵；因為藉著\Quote{不結好果子的}這句話，他表明結出果實的樹將免於一切懲罰。\par

5. 有人說：\Quote{當斧刃已近，時間如此緊迫，期限縮短，我們怎能結出果實呢？}他說：\Quote{你們能夠，因為這種果實不同於普通樹木的果實，後者需要漫長的時間，受季節限制，還需許多其他管理；而你們只需願意，樹木便立即結出果實。因為不但根的本性，而且農人的技巧也大大有助於這種結果實。}\par

因為（讓我補充）為此——免得他們說：\Quote{你警醒我們、催促我們、強迫我們，舉著斧頭，威脅要砍我們，卻要求在受罰的時刻結出果實}——他為了表明結那果實的容易，補充說：\Quote{我固然以水洗你們，但在我以後要來的那位，比我更強，我連解祂的鞋帶也不配；祂要以聖神和火洗你們。}這意味著只需要決心和信德，而不是勞苦和辛勞；正如受洗容易，悔改和成為更好的人也容易。因此，他藉著對\Emph{天主}審判的恐懼、對\Emph{祂}懲罰的期望、以及斧頭的提及、祖先的失落、其他子孫的引進、砍伐和焚燒的雙重懲罰，以各種方式軟化了他們的剛硬，使他們渴望脫離如此巨大的邪惡；然後他才引出關於基督的話，並且不是簡單地，而是宣告祂的偉大優越。然後，在陳述他自己與基督之間的差異時，免得他似乎是在出於偏袒而說這話，他藉著比較二者所賜的恩賜來證實此事。因為他沒有立刻說\Quote{我不配解祂的鞋帶}，而是在首先表明自己洗禮的微小價值，並顯示它除了引人悔改之外別無他物（因為他沒有說以水洗禮的赦免，而是悔改的洗禮）之後，他才陳述基督的洗禮，其中充滿了不可言喻的恩賜。如此，他似乎是在說：免得當你們聽說祂在我以後來時，你們因祂來得較晚而輕視祂；要認識祂恩賜的德能，你們就會清楚知道，當我說\Quote{我不配解祂的鞋帶}時，我所說的並不相稱或偉大。同樣，當你們聽說\Quote{祂比我更強}時，不要以為我是以比較的方式說這話。因為我連被列入祂的僕人之中都不配，不，連祂最低微的僕人也不配，也不能接受祂事奉中最卑微的部分。因此，他不僅僅說\Quote{他的鞋子}，甚至連\Quote{鞋帶}也沒有說，而這種服務被認為是眾事中最末的。然後，為了防止你們將所說的話歸於謙卑，他還從事實中補充了證據，說：\Quote{因為祂要以聖神和火洗你們。}\par

6. 你看洗者若翰的智慧何其伟大？当主自己宣讲时，他说尽一切令人警醒、使他们焦虑的话；但当他差人往主那里去时，就只说那些柔和、能挽回他们的话：不提及斧头、被砍下焚烧的树、投入火中、将来的忿怒，而是罪过的赦免、刑罚的消除、正义、圣化、救赎、义子名分、兄弟情谊、承受产业的分沾、圣神的丰沛供应。这一切他都隐约预示，当他说：\BibleQuote{他要以圣神洗你们}，藉此比喻本身，即刻表明恩宠的丰盈（因为他没有说：\BibleQuote{他要把圣神赐给你们}，而是说\BibleQuote{他要以圣神洗你们}）；另一方面，藉着指明火，他指出了祂恩宠的强烈而不可抑制的性质。\par

试想，当听众想到自己将立刻像先知和那些伟人一样时，他们应当成为怎样的人。你看，正因如此，他才提及火，为引导他们思念那些人的往事。因为在向他们显现的一切异象中，我几乎可以说，大部分都出现在火中：天主在荆棘丛中向梅瑟说话，在西乃山上向全体百姓说话，在革鲁宾上向厄则克耳说话。\BibleRef{则 1:27}\par

再看，他如何激发听众，把最后才发生的事放在最前面。因为羔羊将被宰杀，罪过将被抹去，仇恨将被摧毁，埋葬将要发生，复活要发生，然后圣神才要来。但他目前一件都没有提及，而是把最后的事放在最前面，且为这一切前事所成就的，正是最适于彰显祂尊威的事；这样，当听众被告知他将领受如此伟大的圣神时，他会在心里思索：罪既如此盛行，这事如何、以何种方式发生？发现他心中充满思虑，准备好接受那教训，他便在那时引入他关于苦难要说的内容，这样在期待如此恩赐之下，再无人感到不快。\par

因此他又呼喊说：\BibleQuote{请看，天主的羔羊，背负世罪者！}他没有说\BibleQuote{除去}，而是说含有更细心照顾之意的\BibleQuote{背负}。因为单纯除去是一回事，而把它担负在自己身上是另一回事。前者可无危险而行，后者则需死亡。\par

他又说：\BibleQuote{祂是天主子。}\BibleRef{若 1:34}但这甚至没有向听众公开宣告祂的地位（因为他们还不懂得视祂为真子）：而是藉着如此伟大的圣神恩赐，那地位也得以确立。因此，圣父在派遣若翰时，\Emph{如你所知}，给了若翰这个作为临到者尊威的第一个标志，说：\BibleQuote{你看见圣神降下来，停在谁身上，谁就是那要以圣神洗人的那位。}因此他自己也说：\BibleQuote{我看见了，便作证：这就是天主子。}好像前者对后者永远是明确的证据。\par

7. 然后，既已说出任务中较温和的部分，安抚并缓和了听众，他又约束他们，免得他们变得松懈。因为犹太民族的特性就是这样：容易被一切鼓励的事所自满、败坏。因此他又提出他的恐吓，说：\BibleQuote{祂手里拿着簸箕。}\BibleRef{玛 3:12}\par

这样，正如他先前论及惩罚，此处他也指明判官，并引入永罚。因为他说：\BibleQuote{祂要把糠秕用不灭的火烧尽。}你看，祂是万有的主，祂自己就是农夫；虽然他在别处称圣父也是如此。因为祂说：\BibleQuote{我父是农夫。}\BibleRef{若 15:1}既然祂提到了斧头，免得你以为这事需要劳苦，分离难以做到，他便藉另一个比喻暗示其容易，意指全世界都属于祂；因为祂不能惩罚不属于祂的人。诚然，目前一切都混杂在一起（因为虽然麦子闪耀可见，却与糠秕同在，如在禾场上，而非在仓里），但到那时，分离将是巨大的。\par

现在，那些不信地狱之火的人在哪里？因为这里确实设定了两点：一是祂将用圣神施洗，二是祂将焚毁悖逆者。如果前者可信，那么后者也必然可信。是的，正因如此，他将这两个预言紧密连接，以便通过已经发生的事来证明那尚未成就的另一预言。因为基督自己也常在许多地方这样做，常对同一件事或相反的事，设立两个预言；祂在此世应验一个，在将来应许另一个；以便那些过于好辩的人，可以从已经成就的一个，相信那尚未成就的另一个。例如，祂应许那些为祂的缘故舍弃一切的人，在今世要得到百倍，并在来世得到永生；通过已经赐予的，使将来的也应可信。正如我们所见的，若翰同样在此处做了：设立了两件事，即祂将用圣神施洗，并用不灭的火焚烧。那么，如果祂没有用圣神为宗徒和每天所有愿意的人施洗，你或许会对那些其他事产生怀疑；但如果那看似更大更难、超越一切理性的事已经成就，并且每天都在成就，你怎能否认那容易的、合乎理性的事是真实的呢？因此，说了\Quote{祂要用圣神和火施洗}，并由此应许了巨大的祝福；免得你完全脱离了先前的事而变得懈怠，他加上了簸箕和由此宣告的审判。因此他说：\Quote{你们绝不要以为，如果以后成为凡俗之人，你们的圣洗就足够了}；因为我们既需要德行，也需要那众所周知的克己。因此，正如他用斧头催迫他们走向恩宠和圣洗池，同样，在恩宠之后，他用簸箕和不灭的火恐吓他们。对于那尚未受洗的一类，他没有区别，而是概括地说：\BibleQuote{凡不结好果子的树，就砍下来}，\BibleRef{玛 3:10}，惩罚所有不信者。而在圣洗之后，他进行了一种区分，因为许多相信的人会活出与信仰不相称的生活。\par

那么，谁也不要成为糠秕，谁也不要被抛来抛去，不要暴露于邪恶的欲望之下，被它们轻易地四处吹动。因为如果你保持为麦子，即使试探临到你，你也不会遭受任何可怕的事；不，因为在打谷场上，那车轮如锯子的齿并不会把麦子切碎；但如果你堕落成糠秕的软弱，你在此世就会遭受不治的恶事，被所有人击打，而在那里你将承受永罚。因为所有这样的人，在那火炉之前，在此世就已成为无理性情欲的食物，如同糠秕成为牲畜的食物；而在那里，他们又成为火焰的材料和食物。\par

如果直接说祂将审判人的行为，不会如此有效地使祂的教导被人接受；但若将比喻与之混合，从而完全确立它，则更易说服听者，并以更充分的鼓励吸引他。因此，基督自己也大多这样与他们谈论；打谷场、收割、葡萄园、榨酒池、田地、网、捕鱼，以及一切熟悉的事物，以及他们忙碌的事物，祂都用作祂讲论中的素材。那么，洗者若翰在此也做了同样的事，并提出了一个极其有力的证据来支持他的话，即赐予圣神。因为\Quote{祂既有如此大的权能，既能赦罪，又能赐予圣神，那么这些事也必定在祂的权能之内：}他这样说。\par

你看见现在奥秘如何按适当顺序被奠立根基，在复活和审判之前吗？\BibleRef{希 6:1}\par

\Quote{那么为什么，}也许有人说，\Quote{他没有提及祂立刻要行的征兆和奇迹？}因为这事比一切更大，并为了这事，所有那些都做了。因此，在他提到主要之事时，他包含了所有：死亡被消除，罪被废除，咒诅被抹去，那些长久的战争被终止；我们进入乐园，升天，与天使同为公民，分享未来的美善：因为事实上这就是这一切的凭据。所以，在提到这一点时，他也提到了我们身体的复活，祂奇迹在此的显示，我们分享祂的王国，以及那\BibleQuote{眼所未见，耳所未闻，人心所未想到的}美善。\BibleRef{格前 2:9}。因为所有这些美善，祂都藉着那恩赐赐予了我们。因此，谈论那立刻要发生、且视觉可以判断的征兆是多余的；但应当谈论那些他们所怀疑的事；例如，祂是天主子；祂远远超过若翰；祂\Quote{承担世人的罪}；祂将要求我们对所行的一切交账；我们的利益不限于现世，而是每个人在别处都要承受应得的惩罚。因为这些事还未能凭视觉证明。\par

8. 因此，知道这些事，让我们在打谷场上时，要极其勤勉；因为当我们在这里时，即有可能从糠秕变成麦子，同样，另一方面，许多人也从麦子变成了糠秕。那么，我们不要懈怠，也不要被各种风吹动；也不要与我们的弟兄分离，尽管他们似乎微小而卑下；因为麦子与糠秕相比，虽然量小，但本性更优。因此，不要看外在炫耀的形式，因为它们是为火预备的；而要注视这种属天的谦卑，如此坚固、不可分解，不能被砍断，也不能被火烧毁。正是为了他们的缘故，祂长久忍耐那糠秕本身，好使他们通过与他们的交往而变得更好。因此，审判尚未来到，好使我们能一同被加冕，好使许多人从邪恶归向德行。\par

那么，我们当因听到这个比喻而战栗。因为那火确实是不灭的。\Quote{怎么会不灭呢？}或许有人说。\Quote{它怎么会不灭？}难道你没有看见这太阳永远燃烧，却从未熄灭吗？难道你没有看见荆棘燃烧，却没有被烧毁吗？那么，如果你也渴望逃脱那火焰，就预先积攒施舍，这样你甚至不会尝到那火。因为，如果在这里你相信所告诉你的，那么在你离开到那地方之后，你将不会看见那火炉；但如果你现在不信，那么你将在那里通过经验充分认识它，那时任何逃脱都无可能。因为事实上，没有恳求能为那些没有活出正直生活的人免去惩罚。因为相信当然不够，就连魔鬼也向天主战栗，但他们仍将受罚。\par

9. 因此，我们对行为举止的关怀必须极大。这就是我们不断在此召集你们的原因；不是单要你们进来，而是也要从你们的停留中收获一些果实。但如果你确实常常来，却由此毫无收获地离去，那么你进入和出席这地方就没有任何益处。\par

因为我们送孩子去老师那里时，如果看到他们没有因此获益，就开始严厉责备老师，并常把他们转到别的老师那里；那么我们对自己的美德，难道有什么借口可以不付出甚至像对这些世俗事物那样多的勤勉，却总是把空白的写字板带回家吗？然而，我们这里的老师更多且更伟大。因为先知、宗徒、圣祖和所有义人，人数不少，都被我们安置在每间教会中作为你们的老师。但即便如此，也没有任何益处；如果你随意地、马马虎虎地唱了两三篇圣咏，做了惯常的祈祷，就这样被遣散，你认为这足以得救吗？难道你们没有听见先知（或者更确切地说，天主藉先知）说：\BibleQuote{这百姓用嘴唇尊敬我，他们的心却远离我？}\par

因此，为了避免我们也发生同样的情况，请擦去那文字——或者更确切地说，那印记——就是魔鬼刻在你灵魂上的；并给我一颗摆脱了世俗纷扰的心，好让我能毫无恐惧地写上我所愿意的。因为至少现在，除了他的文字——抢劫、贪婪、嫉妒、忌恨——之外，其他什么都分辨不出来。因此，当然，当我收到你们的写字板时，我甚至无法阅读它们。因为我找不到我们每个主日刻写在你们身上，然后让你们离开的文字；而是其他取而代之的、无法辨识且畸形的文字。那么，当我们擦去它们，并写上了属圣神的文字后，你们离去，把心交到魔鬼的作为中，又给他能力在你们身上替换他自己的文字。这一切的结局将会如何，即使不用我说，每个人的良心也知道。因为我确实不会停止尽我的本分，在你们身上写正确的文字。但如果你破坏了我们的勤勉，就我们而言，我们的赏报不变，而你们的危险却不小。\par

现在，虽然我极不愿说任何使你们不悦的话，但我再次恳求和规劝你们，至少要在这事上效法小孩子们的勤勉。因为他们先学习字母的形状，然后练习辨认扭曲的字母，最后才通过它们有序地阅读。让我们也这样做；让我们划分美德，先学习不发誓、不发假誓、不说坏话；然后进入另一行，不嫉妒、不放纵情欲、不贪食、不醉酒、不暴戾、不懒惰；这样我们可以从这些再进入属圣神的事，操练节制、克己、节欲、正义、超越荣耀、温柔和谦卑的心；让我们把这些彼此连接，写在我们的灵魂上。\par

10. 所有这些，让我们在家里，在自己朋友、妻子、儿女中间操练。目前，让我们从那些首要且更容易的开始；例如，从不发誓开始；让我们不断在家里操练这一字母。因为确实，家里有许多事妨碍我们操练：有时仆人激怒人，有时妻子烦扰并惹人发怒，有时不听话、任性的孩子催逼人威胁和发誓。如果现在在家里，当你不断被刺痛时，却能做到不被引诱而发誓，那么你在市场上也将容易保持不败。\par

同样地，你将通过不对妻子、仆人、或家中任何其他人辱骂，而做到保持自己不辱骂任何人。因为人的妻子也常常称赞这个或那个人，或自怨自艾，激起他说别人的坏话。但不要让自己被迫说那被称赞之人的坏话，而要高尚地忍受一切。如果你察觉到仆人在称赞其他主人，不要心烦意乱，而要高尚地站立。让你的家成为某种竞技场，一个操练美德的场所，好使你在那里好好训练后，能以充分的技巧面对外面的一切。\par

在虚荣方面也要这样做。因为如果你训练自己在妻子和仆人面前不虚荣，那么此后你在别人面前也很难轻易被这种情欲抓住。因为虽然这种病在任何情况下都是严重而专横的，但尤其当有女人在场时更是如此。因此，如果我们在这情形下压制了它的力量，我们便能轻易地在其他情形下也胜过它。\par

对于其他情欲，我们也当这样做，在家里操练自己对抗它们，并每天为自己傅油。\par

为使我们的操练更为容易，让我们进一步为自己制定一种处罚，用于我们违背任何决心之时。且让这处罚本身又带来不是损失，而是赏报——带来某些极大的收益。若我们判处自己更严格的禁食、常睡光地以及诸如此类的刻苦，便是如此。因为这样，从各方面都会有大量益处临于我们：我们将在此世度德行甘美之生活，并将获得未来的美善，且永远成为天主的朋友。\par

但为避免同样的事再次发生——免得你们在此赞赏所讲的话后离去，随意将你们心志的板片丢在任何地方，让魔鬼涂抹这些事——愿每个人回家后，叫自己的妻子，将这事告诉她，并让她协助自己；从这天起，让他进入那高贵的操练学校，以圣神的供应为膏油。即便你在训练中跌倒一次、两次、许多次，也不要失望，而要站起来再斗；不要放弃，直到你戴上了战胜魔鬼的光荣凯旋之冠，并为将来积累了德行财富于一个不可侵犯的宝库中。\par

因为若你使自身习惯于这高贵克己的习操，那么，即使懈怠时，你也不能违犯任何诫命，因为习惯模仿了本性的坚固。是的，如同睡眠、饮食、呼吸是容易的，德行的事工对我们也将是容易的，我们将为自己收获那纯净的快乐，安息于无浪的港湾，享受持续的宁静，带着巨大的负载将我们的船驶入港口，进入那城，在那日子；我们将获得不朽的荣冠，愿我们众人获得它，藉着我们的主耶稣基督的恩宠和爱人，愿光荣与权能归于祂，从现在直到永远，永无穷尽。阿们。\par

\SourceRecord{Translated by George Prevost and revised by M.B. Riddle.；Nicene and Post-Nicene Fathers, First Series；Vol. 10.；Edited by Philip Schaff.；Buffalo, NY: Christian Literature Publishing Co.,；1888.；<http://www.newadvent.org/fathers/200111.htm>.}

% structure: p0001 source=h1 target=section reason=page-title

\section{玛窦福音第十二篇讲道}

\BibleRef{玛 3:13}\par

\Quote{那时，耶稣由加里肋亚来到约但河，}等等。\par

主与仆役同列，审判者与罪犯同伍，前来受洗。但不要困扰；因为正是在这些卑微中，祂的崇高最显光辉。祂曾屈尊在童贞女的胎中怀胎如此之久，带着我们的本性从那里出来，又受鞭打，被钉十字架，并忍受了祂所忍受的一切——为何希奇祂也屈尊受洗，并与其余的人一起来到祂的仆人那里呢？因为惊异在于这一件事：祂是天主，却甘愿成为人；其余的一切都顺理成章。\par

为此，容我补充，若翰也预先说了他以前所说的全部话，即他\Quote{不配解祂的鞋带}，以及其余的话，例如祂是审判者，按各人的行为施行报应，并将丰沛地赐予众人祂的圣神；这样，当你看见祂来到洗礼时，就不会怀疑任何卑微的事。因此，祂来了，若翰也阻拦祂，说道：\par

\BibleQuote{我本来需要受祢的洗，而祢却来就我吗？} \BibleRef{玛 3:14}因为，既然这洗礼是\Quote{悔改的}，并引导人承认自己的过犯，免得有人以为祂也怀着这种心意\Quote{来到约但河}，若翰事先纠正了这一点，称祂为羔羊和救赎者，赎去世上的一切罪恶。既然祂能除去全人类的罪，祂自己更是无罪的。为此，他没有说：\Quote{看，那无罪的，}而是说了更重的：祂\Quote{除免世罪者}，好叫你连同这真理一同确信地领受那另一真理，并领受了便能觉察，祂是出于某种更深远的安排才来受洗的。因此，祂来的时候，若翰也对祂说：\Quote{我本来需要受祢的洗，而祢却来就我吗？}\par

他没有说：\Quote{而祢却由我受洗吗？}不，因为他不敢这样说；而是说什么？\Quote{而祢却来就我吗？}那么基督做了什么？祂后来对伯多禄所做的，当时也做了。因为伯多禄也拦阻祂洗脚，但听到\Quote{我所作的，你现在不明白，后来必要明白}，以及\BibleQuote{你与我无分} \BibleRef{若 13:7}时，便迅速放弃了自己的主张，转向了相反的一面。同样，若翰听到\Quote{你暂且容许吧，因为我们应当这样以履行全部义德}后，立即服从了。因为他们并非固执地争辩，而是显出了爱与服从，并致力于在一切事上受主的管束。\par

请留意祂如何恰恰以那最使他怀疑所发生之事的那一点来说服他；祂没有说：\Quote{这样是合理的，}而是说：\Quote{这样是相宜的。}因为，在他看来，最不相称的一点就是祂由仆人受洗，祂说出了这一点而非其他，恰恰与那印象相对：仿佛祂说：\Quote{你所避免和阻止的，难道不是不相宜吗？不，正是为了这个缘故，我命你容许，因为这是相宜的，而且是极相宜的。}\par

祂不仅说：\Quote{容许，}还加上了\Quote{现在。} \Quote{因为不会永远这样，}祂说，\Quote{但你将看见我如你所愿，只是目前，忍耐这个吧。}接着祂也表明这如何是\Quote{相宜的}。那么如何呢？\Quote{因为我们履行全部法律；}为此祂说：\Quote{全部义德。}因为义德就是遵守全部诫命。\Quote{既然我们履行了其余一切诫命，}祂说，\Quote{只剩下这一条了，也必须加上；因为我来了，是要解除因违犯法律而定的诅咒。所以我必须先全部履行，然后将你们从法律的谴责中释放出来，这样才结束它。因此，我履行全部法律是合宜的，正如同我解除那在律法中对你们所写的诅咒是合宜的一样；这正是我取得肉身、来到此间的目的。}\par

2. \Quote{于是若翰就容许了祂。耶稣受洗后，立时从水里上来，忽然天为祂开了，他看见天主的圣神有如鸽子降下，来到祂上面。}\par

因为许多人以为若翰比主更大，因为若翰一生在旷野长大，又是大司祭之子，穿着那样的衣服，召唤众人受他的圣洗，且由一位不育的母亲所生；而耶稣，首先，出自一位普通女子（因为童贞女生子尚未显明给众人）；此外，祂在家中长大，与众人交谈，穿着平常的衣服；他们便以为祂不如若翰，对那些隐秘之事尚一无所知——而且祂又受若翰的圣洗，这件事更助长了这种猜测，即使先前那些事都不存在；因为他们会想，此人不过是众人之一（因为祂若非众人之一，就不会与众人一同来受圣洗），而若翰比祂更大，更值得钦佩：——因此，为使这种看法不流行于群众之中，当祂受圣洗时，天就开了，圣神降下，并有声音与圣神同来，宣告独生子的尊威。因为那说\Quote{这是我所爱的儿子}的声音，对群众而言，似乎更属于若翰，因为并未加上\Quote{这受圣洗的}，而仅说\Emph{这}，每位听者都会以为是指施洗者，而非领受圣洗者，部分由于施洗者本身的尊威，部分由于上述的一切；圣神便以鸽子的形状降下，将声音引向耶稣，使众人明确知道，这\Emph{这}不是指施洗的若翰，而是指受圣洗的耶稣。\par

或许有人问：发生了这些事，他们为何仍不相信？因为在梅瑟时代，也曾有许多奇事发生，虽不如这些；而经过那些事、声音、号角和闪电之后，他们仍铸造了牛犊，并\Quote{与巴耳培敖尔连结}。而那些亲眼看见拉匝禄复活的人，不仅未相信行这事的主，甚至多次企图杀害祂。既然亲眼看见死人复活，他们尚且如此邪恶，又何必惊讶他们不接受从天上传来的声音呢？因为当灵魂不坦诚、乖僻，且被嫉妒之病缠扰时，它不会屈服于这些事；而灵魂坦诚时，则用信德接受一切，并不需要这些。\par

因此不要这样说：\Quote{他们不相信}，而要问：\Quote{难道所有应当使他们相信的事没有发生吗？}因为天主也藉先知为祂自己的道路普遍作出这种辩护。也就是说，犹太人濒临灭亡，并要被交付于极刑；免得有人因他们的邪恶而诬蔑天主眷顾，祂说：\Quote{对我的葡萄园，我应当做什么而没有做？}同样，在此也应考虑：\Quote{应当做什么却没有做？}的确，每当关于天主眷顾的争论出现时，就要对那些因许多人的邪恶而试图对之抱有偏见的人，使用这种辩护。请看，例如，发生了何等令人惊奇的事，是将来之事的预兆；因为此时不再是乐园，而是天开了。\par

但让我们将与犹太人的辩论留到另一时间；目前，在天主与我们一同工作时，我们要将我们的讲论转向眼前的事。\par

3. \Quote{耶稣受了洗，随即从水里上来；看！天为祂开了。}\BibleRef{玛 3:16}\par

为何天开了？为告诉你，在你受圣洗时，这事也照样发生，天主召叫你往高处之乡，并劝你不要与地有任何牵连。即使你看不见，也绝不要怀疑。因为在所有奇异和属神行事的开端，常有可感觉的异象和此类征兆出现，为那些性情有些迟钝、需要外在视觉、完全不能领会无形体本性、只被所见之物激动的人；这样，即使日后不再有此类事发生，那最初一次所宣告的事，可由你的信德领受。\par

因为在宗徒们身上，也有\Quote{大风的响声}\BibleRef{宗 2:2}和火舌的异象出现，但不是为宗徒们，而是因为当时在场的犹太人。尽管如此，即使没有可感觉的记号，我们也接受这些记号所曾启示的事。因为当时鸽子之所以显现，是为了像手指一样（可以这样说）向在场的人和若翰指出天主子。但不仅如此，也为教训你：在你受圣洗时，圣神也同样降临。然而，既然我们不再需要感觉的异象，信德就足以代替一切。因为\Quote{奇迹不是为信的人，而是为不信的人。}\BibleRef{格前 14:22}\par

但是爲甚麼以鴿子的形狀呢？這生物是溫良而純潔的。既然聖神也是\Quote{溫良之神}，所以祂以這種形狀出現。此外，祂令我們想起一段古老的歷史。從前，當一場普遍的洪水滅頂覆蓋了全世界，我們的種族瀕臨滅亡時，這生物出現，指明了從風暴中的救援，並銜著一根橄欖枝，\EditorialNote{創世紀第八章}傳報了全世界普遍平靜的佳音；這一切都是未來之事的預像。實在，那時人類的處境更爲惡劣，應受更嚴厲的懲罰。爲使你不致失望，祂便提醒你那段歷史。因爲那時，當事情絕望時，也有一種救援和革新；但那時是藉著懲罰，如今卻相反，是藉著恩寵和不可言喻的恩賜。\BibleRef{格後 9:15}因此鴿子也出現，不是銜著橄欖枝，而是爲我們指出從一切邪惡中救援我們的那一位，並提示恩寵的希望。因爲她不是帶領一個人從方舟出來，而是在她出現時帶領整個世界升天，並代替從橄欖枝而來的和平枝條，將義子之恩傳遞給全世界所有的後裔。\par

現在請思想這恩賜的偉大，不要因祂以這樣的形狀出現而輕看祂的尊嚴。因爲我確實聽見有些人說：\Quote{正如人與鴿子之間的差別，基督與聖神之間的差別也是這樣：因爲一位以我們的性體出現，另一位卻以鴿子的形狀出現。}對於這些話，我們該說甚麼呢？天主子確實取了人的性體，但聖神並沒有取鴿子的性體。因此，福音作者也沒有說\Quote{以鴿子的性體}，而是說\Quote{以鴿子的形狀}。所以，祂後來再也沒有以這種方式出現，僅在那一刻而已。如果你因此斷定祂的尊嚴較低，那麼革魯賓也將因這理由被認爲遠超過祂，正如鷹超過鴿子一樣：因爲他們也被塑造成那可見的形狀。天使們也同樣超過祂，因爲他們也多次以人的形狀出現。但事情並非如此，絕非如此。因爲救恩計劃的真理是一回事，臨時異象的俯就是另一回事。\par

現在，我懇求你，不要成爲忘恩負義的人，對待你的恩人，也不要以恰恰相反的事回報那將福樂之源賜給你的人。因爲在哪裏義子之恩被賜予，那裏就有邪惡的消除，並有一切美善的賜予。\par

4. 正因如此，猶太人的圣洗終止，而我們的圣洗開始。正如有關逾越節所發生的，同樣的事也發生在圣洗上。因爲在那件事上，祂同時顧及兩者，使前者結束，爲後者開端；同樣，在此處，祂既完成了猶太人的圣洗，同時也開了教會圣洗的門；如同從前在一張桌子上，如今在一條河流中，祂既勾勒了陰影，如今又加上真理。因爲唯獨這圣洗擁有聖神的恩寵，而若翰的圣洗卻缺乏這恩賜。正因如此，在那些受圣洗者身上並沒有發生這樣的事，唯獨在將要傳授這圣洗的那一位身上發生；這樣，除了我們所說的之外，你也可得知：產生這效果的，不是施圣洗者的純潔，而是受圣洗者的權能。在此之前，天門確實未曾敞開，聖神也未降臨。從此，祂帶領我們脫離舊的，進入新的政體，既爲我們敞開高天的門，又從那裏遣下聖神，召喚我們到天上的家鄉；不僅召喚，而且以最尊貴的標誌。因爲祂沒有使我們成爲天使或總領天使，而是使我們成爲\Quote{天主的子女}和\Quote{蒙愛者}，這樣帶領我們趨向我們的產業。\par

既然心中懷有這一切，就當彰顯與那召喚你的愛、與你在那世界的公民身份、以及與你所受的尊榮相稱的生活。你既已與世界同釘在十字架上，也把世界釘在自己身上，就當以嚴格的態度展示自己是天朝之公民。不要因你的身體沒有被提升到天上，就以爲你與大地還有任何關係；因爲你的頭已在天上。主正是以這個目的，先來到這裏，帶來了祂的天使，然後帶著你離開那裏；甚至在你上升到那地方之前，你就能明白，你有可能如同在天上一般居住在地上。\par

那麼，讓我們守護那從起初就領受的高貴出身；讓我們每天都更多地追求天上的宮殿，把這裏的一切視爲陰影和夢境。假設地上某個國王，見你貧窮而討飯，突然使你成爲他的兒子，你絕不會再記起你的小屋及其簡陋的陳設。然而在這種情況下，差別並不大。那麼，在這種情況下，也不要考慮任何從前的事，因爲你被召叫去承受更大的事。因爲召叫你的那一位是天使的君主，所賜予的美善超越一切言詞和思想。因爲祂不是像地上的國王那樣，將你從地上遷到地上，而是從地上遷到天上，從可朽的本性遷到不朽的，並遷到不可言喻的光榮，這光榮只有當我們實際享用它時才能恰當地彰顯出來。\par

如今，既然你将要分享如此大的恩赐，我岂能看到你挂念金钱，贪恋此世的浮华？你岂不认为一切所见之物比乞丐的破衣更为卑贱？你怎能显得配得上这尊荣？你又有何借口可陈？或者更确切地说，你受了如此大的恩赐却跑回你从前的呕吐物，你岂不要受惩罚？因为今后你受罚不再仅仅是人，而是犯了罪的天主儿子；你尊荣的伟大成了使你受更严厉惩罚的手段。因为我们惩罚做错事的奴隶和犯同样过错的儿子时也不一样；尤其当他们从我们这里领受了一些大恩惠时。\par

因为如果那以乐园为份的人，为了一次不顺服，在受尊荣后遭受了如此可怕的事；那么我们这些接受了天堂、并成为独生者共同继承人的，还有什么借口，在鸽子之后跑向蛇呢？因为不再是对我们说：\Quote{你是尘土，将归于尘土}\BibleRef{创 3:19}并且你\Quote{耕种田地}以及那些先前的话；而是比这些更严重的：\Quote{外黑暗}\BibleRef{玛 25:30}、不可挣开的锁链、毒蛇的虫子、\Quote{切齿}；这十分合理。因为连如此大的恩惠都不能改善的人，公正地会遭受最极端、更严厉的惩罚。厄里亚曾一度开阖天，那是为降雨和止雨；但你的天并非如此打开，而是为让你升到那里；不仅如此，不只是升到那里，还能引导别人也上去，只要你愿意；祂已在你一切事上赐予你如此大的信心和能力。\par

5. 既然我们的家在那里，我们就应将一切储存于那里，什么也不留在这里，以免失去。因为在这里，即使你装上门锁、房门、门闩，又派遣成千的仆人看守；即使你胜过了所有狡猾的人，逃过了嫉妒者的眼目、虫蚁、时间的侵蚀——这不可能——你无论如何也逃不过死亡，而会在顷刻之间被剥夺这一切；不仅被剥夺，还要常常将它们转移到你的仇敌手中。反之，你若能将它们转移到那家中，你将超乎万有之上。因为那里无需钥匙、门户或门闩；那座城有如此美德，这地方如此不可侵犯，本性上不受腐败和一切邪恶的侵害。\par

那么，在一切储藏之物注定毁坏和消耗的地方，我们将一切堆积于那里，而在东西恒存不坏、增长的地方，我们却连最小的一份也不储存，这岂非极端的愚妄？尤其当我们将要永远住在那里时。正因如此，外邦人都不信我们所说的话，因为他们愿意从我们接受的证明是我们的行为，而非言语；当他们看见我们为自己建造华屋、开辟花园、浴场、购置田地时，他们就不愿相信我们正在为另一处离开城市的居所做准备。\par

\Quote{如果是这样，}他们说，\Quote{他们就会把这里的一切换为金钱，预先储存到那里；}他们从今世的事推测出这一点。因为我们看到那些非常富有的人，主要在他们要居住的城市里为自己购置房屋、田地和其余一切。但我们做的却相反；我们以全部的热心占据这大地，却很快就要离开；我们不仅放弃金钱，甚至为几亩地和几座房屋付出我们的血；而为购买天堂，我们却连超出所需的部分也不肯付出，尽管我们本可廉价购得，并且一旦购得就永远拥有。\par

因此我说，我们将受极重的惩罚，赤身贫苦地离开那里；或者更准确地说，我们遭受这些无可补救的灾祸，不仅因我们自己的贫穷，也因我们使别人变得像我们自己一样。因为当外邦人看见那分享了如此大奥秘的人对这些事热衷时，他们自己会更加紧紧抓住这些东西，将多多的火堆在我们头上。因为我们本应教导他们轻视一切可见之物，自己却最有力地怂恿他们贪爱这些事物；那时我们怎能得救？我们还要为别人的丧亡负责。难道你们没有听见基督说，祂留我们在这世上作盐和光，为使我们既振作那些沉溺于奢华的人，又光照那些被财富挂虑所笼罩的人？因此，当我们甚至将他们推入更深的黑暗，并使他们更加放荡时，我们还有什么救恩的希望？绝无希望；我们将哀哭切齿，手脚被捆，进入地狱的火中，此前已被财富挂虑折磨得筋疲力尽。\par

既思虑这一切，让我们松开这般欺骗的束缚，使我们绝不陷入那些将我们交付于不灭之火的事物。因为那作金钱奴隶的人，此世和来世的锁链将使他一直受捆；但那摆脱这种欲望的人，将从两者获得自由。为使我们也能达到这自由，让我们打碎贪吝的重轭，为自己装上飞向天堂的翅膀；赖我们的主耶稣基督的恩宠和慈爱，愿光荣与权能归于祂，直到永远。阿们。\par

\SourceRecord{Translated by George Prevost and revised by M.B. Riddle.；Nicene and Post-Nicene Fathers, First Series；Vol. 10.；Edited by Philip Schaff.；Buffalo, NY: Christian Literature Publishing Co.,；1888.；<http://www.newadvent.org/fathers/200112.htm>.}

% structure: p0001 source=h1 target=section reason=page-title

\section{玛窦福音第十三篇讲道}

\BibleRef{玛 4:1}\par

\BibleQuote{那时，耶稣被圣神领到旷野，受魔鬼的试探。}\par

那时。是什么时候？在圣神降临之后，在从上面传来的声音之后，那声音说：\BibleQuote{这是我的爱子，我所喜悦的。} 更奇妙的是，这出于圣神；因为，他在这里说，是圣神领他上去。既然他为了我们的教导，既行了又受了所有的事，他也忍受被领到那里，与魔鬼争斗：为使每个领受圣洗的人，如果在他受洗后必须忍受更大的试探，不会因结果出乎意料而感到困扰，反而能继续高贵地忍受一切，仿佛事情是自然而然发生的。\par

是的，因此你拿起武器，不是为了闲散，而是为了战斗。为此，天主并不阻止试探的来临：首先，是为了教导你，你已变得更为强壮；其次，使你保持谦逊，不因恩赐的宏大而自高自大，因试探有能力抑制你；还有，为了使那邪恶的魔鬼，虽然暂时对你的离弃他感到怀疑，但通过试探的试金石，他能够确信你已经彻底弃绝并离开了他；第四，使你借此变得更坚强，比任何钢铁更精炼；第五，使你获得对所托付给你的宝藏的明确证明。\par

因为魔鬼不会攻击你，除非他看到你被提升到更大的尊荣。因此，例如，从一开始，他就攻击亚当，因为他看见亚当享有极大的尊位。为此，他又与约伯为敌，因为他看见约伯被万有的天主所加冕和宣告。\par

那么，祂如何说：\BibleQuote{你们要祈祷，免得陷于诱惑。}\BibleRef{玛 26:41} 为此，圣史并不简单地显示耶稣上去，却\Quote{被领上去}，这是按照救恩计划的原则；借此隐约地表示，我们不应自行冲入试探，而是被拖向它时，要英勇地站立。\par

请看：圣神领祂上到哪里去？祂取了祂，不是到城里和广场，而是到旷野。也就是说，祂有意吸引魔鬼，不仅藉着饥饿，也藉着这地方给魔鬼一个把柄。因为魔鬼最常攻击，是在他看见人独处、独自一人的时候。同样，在起初，他也这样攻击女人，趁她独自一人，且远离她的丈夫时抓住她。正如当他看见我们与别人在一起并团结在一起时，他就不那么自信，也不会发动攻击。因此，正因如此，我们极其需要不断聚集在一起，以免向魔鬼的攻击敞开。\par

2. 于是，他们在旷野找到了祂，在一条无路的旷野（因为旷野是那样的，马尔谷曾声明说，祂\BibleQuote{与野兽在一起}\BibleRef{谷 1:13}），请看，他是带着何等的狡猾和恶意靠近的，他是如何寻找机会的。因为不是在祂禁食时，而是在祂饥饿时接近祂；为教导你禁食是何等大的善，它是抵抗魔鬼的最强大的盾牌，并且，在（领洗的）水泉之后，人们应该委身于禁食，而非奢侈、醉酒和丰盛的筵席。因为，为此缘故，连祂也禁食，不是祂自己需要，而是为了教导我们。因此，既然我们在水泉之前的罪是由服侍肚腹带来的：就像一个治愈了病人的人，会禁止他做那些导致疾病的事一样；我们在这里也看到，祂自己在水泉之后引入了禁食。因为，的确，亚当因肚腹的不节制被逐出乐园；诺厄时代的洪水也是由此产生；这（不节制）也招来了对索多玛的雷霆。因为，虽然也有行淫的指控，但每一种惩罚的根源都是从这（不节制）成长起来的；厄则克耳也指明了这一点，他说：\BibleQuote{这就是索多玛的罪恶：她骄傲自大，粮食丰足，奢侈安逸。}\BibleRef{则 16:49} 同样，犹太人因醉酒和奢靡而被驱赶过犯，犯下了最大的邪恶。\BibleRef{依 5:11}\par

为此，连祂也禁食四十天，向我们指出我们救恩的良药；然而祂没有更进一步，免得另一方面，因这奇迹的极其伟大而使祂救恩计划的真理受到质疑。因为事实上，这并不是不可能的，既然梅瑟和厄里亚在祂以前，曾靠着天主的能力，得以持续到那么长的时间。而如果祂持续更久，那么从这一点及其他事情上，祂对我们血肉的取得，对许多人来说就会显得不可信了。\par

祂于是禁食了四十昼夜，\par

\BibleQuote{后来祂就饿了。}\BibleRef{玛 4:2} 给了祂一个可抓着来靠近的把柄，以便通过实际的交锋，祂能显示如何得胜和获胜。同样，摔跤手们也是这样：在教导门徒如何得胜和克服时，他们自愿在竞技场上与别人交手，使这些门徒在对手身上看到和学习征服的方式。同样的事情那时也发生了。因为祂有意把他引到如此地步，祂既让他知道自己的饥饿，也等待他的靠近，当祂等待他时，祂就将他摔倒在地，一次、两次、三次，以与祂相称的如此轻易。\par

3. 但为了不因我们匆忙掠过这些胜利而损害你们的益处，让我们从第一次进攻开始，仔细审查每一个。\par

因此，在祂饿了之后，经上说：\BibleQuote{试探者前来，对祂说：你若是天主子，就命这些石头变成饼。}\BibleRef{玛 4:3}\par

因为（魔鬼）曾听见从天上来的声音说：\Quote{这是我的爱子；}也听见若翰为他作了如此伟大的见证，之后又看见他饥饿，于是便陷入困惑，既不能相信他仅是人，因为他所听见论及他的话；另一方面，也不能接受他是天主子，因为他看见他饥饿。因此，在困惑中，他说出含糊的话。正如起初他来到亚当面前时，他捏造不实之事，以探知实情；同样，在此他因不清楚救恩计划那不可言喻的奥秘，以及来者是谁，便试图编织别的网罗，以为可以探知那隐藏而模糊的事。他说了什么？\Quote{你若是天主子，就命这些石头变成饼。}他并没有说\Quote{因为你饿了}，而是说\Quote{你若是天主子}，以为可以用谄媚来欺骗他。因此他也对饥饿闭口不提，以免显得是在提及并指责他。因为他不知道正进行的救恩计划之伟大，便以为这是对他的责难。故此他狡猾地奉承他，只提及他的尊位。\par

基督如何回答？为了挫败他的骄傲，并指明所发生的事并无任何可耻之处，也不有损他的智慧；那另一者（魔鬼）为奉承他而保持沉默的事，基督提出并阐明，说：\BibleQuote{人生活不只靠饼。}\BibleRef{玛 4:4}\par

因此，他从肚腹的必须开始。但请注意，我祈求你们，那恶者的狡猾，他从何处开始他的搏斗，以及他如何没有忘记自己的伎俩。因为他用何种手段驱出了第一个人，并用千万种其他罪恶包围了他，他在这里也用同样的手段编织他的欺骗；我指的是肚腹的不节制。同样，如今也可以听见许多愚昧人因为肚腹而说出千万句坏话。但基督，为了显示有德行的人甚至不被这种暴虐所迫使去做任何不当之事，首先饥饿，却不服从所命令的事；教导我们在任何事上都不听从魔鬼。如此，因为第一个人正是因此冒犯了天主并违反了法律，基督教导我们：即使魔鬼命令的并非违反法律的事，也不可听从。\par

为什么我说\Quote{违反法律}？\Quote{即使魔鬼提出什么有益的事，你也不要}，他说，\Quote{照样听从他们。}例如，他也堵住了那些魔鬼的口，他们宣称他是天主子。同样，保禄也曾斥责他们，喊着同样的话\BibleRef{宗 16:18}；尽管他们所说的有益，但保禄更加羞辱他们，阻止他们对我们的阴谋，即使当他们宣讲救恩的道理时也赶走他们，封住他们的口，命他们静默。\par

因此，在这件事上，基督也没有同意所说的话。但他说了什么？\BibleQuote{人生活不只靠饼。}他的意思是：\Quote{天主甚至能用一句话使饥饿的人得饱饫。}并引用古经的见证，教导我们，即使我们饥饿，甚至无论受什么苦，也永不离开我们的主。\par

但若有人说：\Quote{他仍应显示自己。}我会问他，有何意图和理由？因为那另一者如此说，并非要使他相信，而是按他所想，要辩倒他使他陷入不信。因为人类起初就是如此被他迷惑和辩倒，没有对天主怀有真诚的信德。因为他许诺了他们与天主所说相反的事，用虚妄的希望使他们自高自大，导致他们不信，从而将他们从实际拥有的福分中逐出。但基督表明自己并未同意，无论是那时还是后来对他的党羽犹太人要求征兆时；他始终教导我们，无论我们有何能力，都不要无谓地或随意行事；甚至当有所缺乏时，也不服从魔鬼。\par

4. 那么这个可咒骂者做了什么？他被战胜了，无法说服基督听从他的命令，尤其是在受到如此剧烈的饥饿压迫时，他转而进行另一件事，说：\par

\BibleQuote{你若是天主子，就跳下去，因为经上记载：祂会为你吩咐祂的天使，用手托着你。}\BibleRef{玛 4:6}\par

什么理由使得他在每次诱惑中都加上\Quote{你若是天主子？}这与他在之前的情况中所做的非常相似。也就是说，正如他那时诽谤天主，说：\BibleQuote{在你们吃的那天，你们的眼睛会开了}\BibleRef{创 3:5}，意在暗示他们被迷惑和欺骗，没有获得任何益处；同样，在此他也暗示同一件事，说：\Quote{天主白白称你为子，并用他的恩赐欺骗了你；因为，若非如此，就给我们一些明确的证据，证明你拥有那样的能力。}然后，因为基督曾用圣经与他辩论，他也引用了先知的一个见证。\par

基督如何回应？他并不愤怒，也不激动，而是以极大的温和再次用圣经与他辩论，说：\BibleQuote{你不可试探主你的天主。}\BibleRef{玛 4:7}教导我们，我们必须胜过魔鬼，不是靠奇迹，而是靠忍耐和宽容，并且我们绝不可为了炫耀和虚荣做任何事。\par

但请注意他的愚拙，甚至由他所提出的见证本身即可看出。因为主所引用的两处圣经都极其恰当；而他这边却是偶然、随意的言语，他也没有提出那适用于当前问题的经文。因为经上记载：\BibleQuote{祂必为你吩咐祂的天使}，这决不是劝人纵身跳下；而且，这句话甚至不是论及主的。但当时祂并未揭露这一点，尽管他的说话方式既带侮辱，又极不一致。因为对天主圣子，没有人要求这些事；但纵身跳下是魔鬼和恶魔的行为。天主的事，却是连那些跌倒的人也扶起来。如果祂本当显示自己的大能，那也不会是随意跳下和翻滚，而是拯救他人。但跳下悬崖、落入坑中，正是属于他的队伍。例如，他们中间的变戏法者到处如此。\par

但基督，即使这些话被说出，仍未显示自己，而是暂时像人一样与他交谈。因为所说的话：\BibleQuote{人生活不只靠饼}，以及：\BibleQuote{你不可试探主你的天主}，适合一位不十分显示自己，而是将自己表现为众人中一员的人。\par

但不要惊讶，如果他在与基督辩论时常常转身。因为正如拳击手受到致命打击后，满身是血、眼目昏花而摇晃；同样，他也因第一和第二次打击而昏暗，信口说出想到的话，并继续第三次攻击。\par

5.\BibleQuote{他带祂上了一座极高的山，把普世万国指给祂看，说：你若俯伏朝拜我，这一切我都给你。那时耶稣对他说：退到我后面去，撒殚！因为经上记载：你要朝拜主，你的天主，惟独事奉祂。}\BibleRef{玛 4:8}\par

因为他现在已到了冒犯圣父的地步，说凡属圣父的都是他的，并试图使自己被当作神，作为宇宙的创造者；那时祂斥责了他，但即使如此也不严厉，只是简单地说：\Quote{退去，撒殚！}这话本身带有命令甚于斥责的成分。因为祂刚对他说了\Quote{退去}，就使他逃跑了；因为他没有对祂提出其他诱惑。\par

路加怎么说，说\BibleQuote{他结束了所有试探}。在我看来，他提到主要的试探时，就已说了所有，好像其余都包含在这些里。因为构成无数邪恶实质的就是这些：作肚腹的奴隶、为虚荣行事、被财富的疯狂所辖制。那被咒骂者考虑到这些，便将最强大的一种，即贪求更多的欲望，放在最后；虽然起初他早已渴望达到这一点，但他仍保留它到最后，以为它比其余更有力。事实上，这就是他摔跤的方式：最后才用那些似乎更容易推翻人的东西。他对\Person{约伯}也是如此。因此，在这个例子中，他也从那些似乎更卑下、更弱的动机开始，然后转向更强大的。\par

那么我们如何胜过他？用基督所教导我们的方式，逃避到天主那里去避难；在饥荒中不沮丧，因为相信天主能用一句话喂养人；在领受任何恩惠时，不试探那赐予者，反而满足于从上来的荣耀，不看重人的荣耀，并时常轻看我们所需之外的事物。因为没有比贪求更多和爱慕贪婪更使我们落在魔鬼权势之下的了。这一点甚至可以从现在所做的事看出。因为现在也有那些人说：\Quote{这一切我们都要给你，只要你俯伏朝拜我们；}他们本性是人，却成了他的工具。因为那时他也曾接近祂，不单自己，也藉着他人。路加也说明这一点，当他说，\BibleQuote{他暂时离开了祂}，显示此后他藉着他合用的工具接近祂。\par

\BibleQuote{看，天使前来伺候祂。}\BibleRef{玛 4:11}因为当攻击进行时，祂不容许他们显现，以免因此驱逐猎物；但在祂在各方面谴责了他，并使他逃跑之后，他们才显现：好叫你也知道，在你们效法祂的胜利之后，天使也要迎接你们，为你们鼓掌，并在一切事上守护你们。例如，天使将\Person{拉匝禄}\BibleRef{路 16:22}带走，在经过贫困、饥饿和一切患难的熔炉之后。因为正如我已经说过的，基督在此场合展示了许多我们自己也当享用的事物。\par

6. 既然这一切都是为你做的，你当效法并模仿祂的胜利。如果那群恶神仆役中有人接近你，并品味属乎他的事，责备你说：\Quote{如果你奇妙而伟大，就挪开这座山；}不要烦恼，也不要惊慌，而要温柔地回答，说一些像你听你主所说的话：\BibleQuote{你不可试探主你的天主。}\par

或者，倘若他献上光荣与权柄、无尽的财富，吩咐你朝拜他，你要再次勇敢地站立。因为魔鬼不仅这样对待我们共同的主，而且每天也将这些诡计加诸于祂的每位仆人身上，不只在山上旷野，也不单凭他自己，同样在城市里、在市场里、在法庭里，以及通过我们的亲属，甚至人类。那么我们必须做什么？全然不信他，捂住耳朵不听他，当他谄媚时恨他，当他提供更多时，更要躲避他。因为\Person{厄娃}的情况也是如此，当他用希望高举她时，然后把她打倒，给她造成了最大的祸害。是的，因为他是一个不共戴天的仇敌，对我们发动了排除一切和约的战争。而我们对自己的救恩并不像他对我们的毁灭那样热切。那么让我们躲避他，不单用言语，也用行动；不单在心里，也在行为上；我们不要做任何他所赞同的事，这样我们就会做所有天主所赞同的事。是的，因为他也会作很多许诺，不是为了给予，而是为了夺取。他许诺通过掠夺，以剥夺我们的天国和义德；并把财宝放在地上，作为某种罗网或陷阱，以剥夺我们这两样以及天上的财宝；他要我们在地上富有，好叫我们不在那里富有。\par

如果他不能通过财富将我们从那里的份中赶出去，他就走另一条路——贫穷之路，如同他对\Person{约伯}所做的那样。也就是说，当他看到财富对他没有伤害时，他就用贫穷编织罗网，指望从那方面战胜他。但有什么比这更愚蠢呢？因为那能适度承受财富的人，更能勇敢地承受贫穷；那在财富存在时不贪恋的人，在财富不在时也不会寻求；正如那位有福的人没有那样做，反而因他的贫穷变得更加光荣。因为那邪恶的魔鬼固然有能力夺去他的财产，却不能夺去他对天主的爱，反而使之更加强烈；当他剥夺了他的一切时，却使他更丰盛地拥有祝福；因此魔鬼也陷入困惑。因为他越是给他带来灾祸，就越看见他变得强大。因此，如你们所知，当他试尽一切，彻底考验了他的品质，却毫无进展时，他就跑回他的老武器——女人，并装出一副关心的面具，以最可怜的音调描绘他的灾难，伪装出要除去他的祸害而引出那致命的劝告。但即使如此，他也没有得逞；因为那非凡的人识破了他的诱饵，并以极大的智慧堵住了那在魔鬼唆使下说话的女人的口。\par

同样，我们也必须这样做：即使是一个兄弟、一位经过考验的朋友、一个妻子，或任何我们最亲近的人，被魔鬼进入，说出不合适的话，我们也不该因说话的人而接受劝告，却要因为致命的劝告而远离说话者。因为事实上现在他也做许多这样的事，并在他面前摆出同情的面具，当他看似友善时，正在灌输他那些比毒药更致命的言语。因此，正如谄媚以行恶是他的本分，同样，责罚以使我们得益，是天主的本分。\par

7. 那么，我们不要被欺骗，也不要以任何手段追求安逸的生活。因为经上说：\BibleQuote{主所爱的，祂必管教}（\BibleRef{希 12:6}），因此，当我们生活在邪恶中却享受顺境时，我们更应悲伤。因为我们在犯罪时应该常常害怕，尤其是在没有遭受恶事的时候。因为当天主逐渐追讨我们的罚债时，他使我们的偿还变得容易；但当他对我们的每一项疏忽长期忍耐时，如果我们继续这样，他就是在为我们积蓄巨大的惩罚。既然对于行善者，磨难是必要的，对于犯罪者更是如此。\par

请看，例如\Person{法郎}遇到了多么大的容忍，后来却遭受了最极端的惩罚；\Person{拿步高}犯了多少罪，最终却为一切赎罪；而那个富翁，因为没有在这里遭受什么大苦，正因如此他主要成了可怜的，因为他在现世过奢侈生活后，去世到了那里为这一切受罚，在那里他完全不能得到任何东西来减轻他的灾祸。\par

尽管如此，有些人还是如此冷漠迟钝，总是只寻求眼前的事物，说出那些荒谬的话：\Quote{让我暂时享受眼前的一切，然后再考虑看不见的事：我要满足我的肚腹，我要成为欲望的奴隶，我要充分利用现世生活；给我今天，明天拿去。} 哦！愚昧至极！为什么？那些如此说话的人与山羊和猪有何区别？因为先知（\BibleRef{耶 5:8}）不允许他们被称为人，就是那些\BibleQuote{向邻舍的妻子嘶叫}的人，谁又能责备我们把这些人视为山羊和猪，比驴还要迟钝呢？因为他们把那些比我们所见更明显的事物视为不确定。为什么？如果你们什么也不信，就请看那些在受鞭打中的魔鬼，他们的一切行为，言语和行动，都是以伤害我们为目的。因为我相信你们不会否认这一点：他们做这一切是为了增加我们的安逸，消除对地狱的恐惧，并滋生对那世界审判的不信。然而，那些如此想法的人，常常通过哭号和哀叹来宣告那里的痛苦。那么，他们为什么这样说，说出与自己意愿相反的话呢？没有别的原因，只是因为他们在更强大的强制力下。因为他们不会自愿承认他们被死人折磨，或者他们遭受任何可怕的事。\par

如今我为何说了这话？因为那恶鬼承认地狱，它们倒愿地狱被人不信；而你，既享受如此大荣，又曾领受那不可言喻的奥迹，却连它们也不如，反比它们更为刚硬。\par

8. \Quote{但是，}有人会说，\Quote{谁曾从地狱中上来，把那些事告诉了我们？}其实，谁又从天堂来到这里，告诉我们有一位创造万有的天主？又怎能清楚我们有一个灵魂？因为显然，你若只信眼所能见的事，那么天主、天使、理智和灵魂对你都成了可疑，这样你便发现一切真理之道都已丧失。\par

然而，你若愿意相信那显而易见的事，那么不可见的事比那可见的事更应为你所信。虽我所说的是奇论，却仍然是真实的，并在明理之人中完全被承认。因为眼睛常被欺骗，不仅在不可见的事上（它们对这些事根本不觉察），甚至在人以为确实看见的事上，距离、空气、心不在焉、愤怒、挂虑以及无数其他事物都妨碍其准确性；而灵魂的推理能力，若接受圣经的光照，则将成为更准确、无误的现实标准。\par

我们不要徒然自欺，也不要在我们生活的懈怠之外（这懈怠正是此类学说的产物），为这些学说本身为自己堆积更炽烈的火。因为若没有审判，我们不必为行为交账，我们也就不为劳苦得酬报。请看，当你说天主是公义、仁慈、温和的，竟忽略如此大的劳苦和辛劳时，你的亵渎是何等趋向？这又怎能合理？因为，即使你不凭其他任何事，至少凭你自己家中的情况，我请你衡量这些事，然后你便会看见其荒谬。因为即使你自己比野兽更野蛮、更残忍、更狂野，你也不会在临死时选择不尊重那曾对你忠心的仆人，反而以自由和金钱酬报他；并且因为此后你已离去，不能再亲自为他行善，你便委托你财产的将来继承人，恳求、劝勉、竭尽全力，使他不致受不到酬报。\par

如此，你这邪恶之人，对你的仆人尚且如此仁慈和爱护；而那无限的美善，即天主，那不可言喻的对人类之爱，那如此浩瀚的仁慈，岂会忽略并让祂的仆人们，\Person{伯多禄}和\Person{保禄}、\Person{雅各伯}和\Person{若望}，那些为祂的缘故每日忍饥、被捆缚、被鞭打、被投入海中、被交与野兽、死亡、遭受我们无法数算的如此巨大苦难的人，得不到冠冕？奥林匹克裁判宣告并给胜者加冕，主人酬报仆人，君王酬报士兵，每人一般地以他所能给予的好事酬报那服事他的人；难道唯独天主，在那些如此巨大的辛劳和劳苦之后，不以任何大小好事酬报他们？那些正义而虔诚、行走在一切德行中的人，岂能与奸淫者、弑亲者、杀人者、掘墓者同列？这怎能合理？因为，若我们离世后一无所有，我们的利益不及于现世，那么这些人与那些人便处于相同情况，甚至不如说并非相同。因为，照你所说，将来他们固然命运相同，但在此世，整个时间中，恶人安逸，义人受罚。这岂是哪种暴君、哪种野蛮无情的人曾为他的仆人和臣民如此设想的？\par

你注意到这荒谬的巨大和这推论导致的结果了吗？因此，你若不由其他方式，至少由这些推理受教，以摆脱这邪恶思想，逃离恶习，并紧握那以德行终结的辛劳，那时你必确实知道，我们的事不限于现世生命。若有人问你：\Quote{谁曾从那里上来，把那里的事告诉了我们？}你便对他说：没有一个人；因为那人必常被不信，认为他自夸夸大其事；但天使之主却精确地告知了所有那些事。我们何需任何人？那将要向我们要求交账的祂，每日高呼，说祂已预备了地狱，也预备了天国，并给我们提供了这些事的清楚证明。因为祂若将来不审判，就不会在此世施罚。\par

9. \Quote{然而，关于这一点的合理性又如何？恶人中有些受罚，有些则不然？我是说，若天主不偏待人（祂确实不偏待），为何一个祂施罚，另一个祂却任其不受罚而去？这比前者更难解释。}\par

然而，若你们愿意坦诚听我们所说，我们也将解决这一难题。\par

那么，解决是什么？祂在此世并非对所有人施罚，免得你绝望于复活，失去对审判的期望，仿佛所有人都在此交账；祂也并非让所有人不受罚而去，免得你反以为一切都无天主眷顾；祂既施罚，也免于施罚：藉那些祂施罚的人，表明在来世祂也将向在此未受罚的人追讨；藉那些祂不罚的人，促使你相信我们离世后还有可怕的审判。\par

如果祂对我们的先前行为完全漠不关心，祂就不会在此惩罚任何人，也不会施予恩惠。但现在你看见祂为了你伸展天空，点燃太阳，奠定大地，倾泻海洋，扩展空气，为月亮设定轨道，为年岁季节制定不变的法则，以及所有其他事物都按照祂的号令准确运行。因为我们的本性，以及无理性受造物的本性——那些爬行、行走、飞翔、游泳的，在沼泽、泉源、河流、山脉、森林、房屋、空中、平原中的；还有植物、种子、树木，野生的和栽培的，结果子的和不结果子的；总之，万物都由那不知疲倦的手所推动，为我们的生命提供供应，从自身向我们献上服役，不仅为了我们的需要，也为了我们崇高地位的感受。\par

因此，看见如此伟大而美好的秩序（然而我们连其中最微小的部分都还没有提及），你竟敢说，那位为了你成就了如此多、如此伟大之事的祂，会在最关键的事情上忽略你，让你死后与驴和猪躺卧在一起？而且，祂以如此伟大的恩赐——即虔敬的恩赐——荣耀了你，借此甚至使你与天使同等，难道在你不计其数的劳苦和辛劳之后，祂会忽略你吗？\par

这怎能合理呢？看，这些事，如果我们缄默不言，\BibleQuote{石头就会立刻喊叫起来；}\BibleRef{路 19:40}，它们是这样清楚、显明，比阳光本身更明亮。\par

因此，在思考了这一切并说服我们自己的灵魂之后，知道在离开此世后，我们将站在那可畏的审判台前，为我们所做的一切交账，如果我们继续疏忽，我们将承受惩罚，接受判决；如果我们愿意稍加注意自己，我们将接受冠冕和无法言喻的福乐；让我们既堵住那些反对这些事之人的口，自己也选择美德之路；这样，以应有的信心前往那审判庭，我们就能获得为我们所预许的美物，藉着我们的主耶稣基督的恩宠和爱人之情，愿光荣与权能归于祂，自今而始，以至永远，世世无尽。阿们。\par

\SourceRecord{Translated by George Prevost and revised by M.B. Riddle.；Nicene and Post-Nicene Fathers, First Series；Vol. 10.；Edited by Philip Schaff.；Buffalo, NY: Christian Literature Publishing Co.,；1888.；<http://www.newadvent.org/fathers/200113.htm>.}

% structure: p0001 source=h1 target=section reason=page-title

\section{论玛窦福音第十四篇讲道}

\Emph{\BibleRef{玛 4:12}}\par

\Quote{耶稣听见若翰被交付了，就退到加里肋亚去。}\par

祂为何离去？再次教导我们不要主动迎向诱惑，而要退让和躲避。因为不将自己投入危险并不算耻辱，但在陷入危险后不能勇敢地站立才是耻辱。为了教导我们这些，并平息犹太人的嫉妒，祂退到葛法翁；既成全了预言，又急忙去捕捉世人的教师；因为正如你们所知，他们正留在那里，从事自己的手艺。\par

但请你留意，每当祂要前往外邦人那里时，总是犹太人给祂提供了机会。就像在此例，他们设计陷害祂的前驱，把他投入监狱，从而把基督赶到了外邦人的加里肋亚。因为为了说明祂既不是指一部分犹太民族，也不是含糊地指所有支派，留意先知如何区分那地方，说：\Quote{纳斐塔里地，沿海之路，约但河外，外邦人的加里肋亚，那坐在黑暗中的百姓看见了大光；}这里的黑暗不是指可感觉的黑暗，而是指人的错谬和不敬。因此他又加上：\Quote{坐在死影之地的人，有光照耀他们。}为了使你明白他所说的光与黑暗都不是可感觉的，在论及这光时，他称它不仅是一道光，而是一道\Quote{大光}，这光在别处他称为\Quote{真光}，见：\BibleRef{若 1:9}；在描述黑暗时，他称之为\Quote{死影}。\par

然后暗示他们不是自己寻求和找到，而是天主从上方向他们显示自己，他对他们说：\Quote{光照耀了；}也就是说，光自己照耀出来：并不是他们先奔向光。事实上，在基督来临之前，人的状况是最糟的。因为他们不仅是\Quote{行在黑暗中}，而且是\Quote{坐在黑暗中}；这象征着他们甚至不期望被拯救。因为像那些不知往哪里迈步的人一样，他们就这样坐着，被黑暗笼罩，甚至无法再站立。\par

\Quote{从那时起，耶稣开始宣讲说：你们悔改吧！因为天国临近了。}\par

\Quote{从那时起：}什么时间？若翰被下狱之后。为什么祂不从一开始就向他们宣讲？既然祂的事迹已经见证并宣告了祂，若翰还有什么必要？\par

这是为了让你也能认识祂的尊威；即，正如圣祖们一样，祂也有先知；为此匝加利亚也说：\Quote{至于你，孩子，你要称为至高者的先知。}\BibleRef{路 1:76}而且，为了不给无耻的犹太人留下任何借口，祂亲自提出的理由说：\Quote{若翰来了，也不吃，也不喝，他们就说：他有魔鬼。人子来了，也吃也喝，他们却说：看哪，一个贪食好酒的人，税吏和罪人的朋友。但智慧是由她的子女成义。}\par

此外，关于祂的事必须先由他人而非祂自己来说。因为即使在如此多、如此伟大的见证和显示之后，他们还说：\Quote{你为自己作证，你的证据不真。}\BibleRef{若 8:13}如果祂在若翰未说任何话之前就来到中间，首先为自己作证，他们还会说什么？因此，祂既不在若翰之前宣讲，也不在若翰被下狱之前行奇迹，以免群众因此分裂。因此若翰也没有行任何奇迹，这样做也是为了让群众归向耶稣，因为祂的奇迹吸引他们。\par

再者，即使在如此多的神圣预防措施之后，若翰的门徒，无论在他被囚之前还是之后，仍然嫉妒祂，人们也怀疑祂不是基督，而是以为若翰是基督；那么，如果这些都没有发生，结果会怎样？因此玛窦明确注明：\Quote{从那时起，祂开始宣讲；}当祂开始宣讲时，祂自己也教导了这同样的教义，即那人（若翰）所宣讲的；并且祂所宣讲的教义还没有关于祂自己的任何话。因为在当时，即使接受这一点也是一件大事，因为他们还没有对祂有正确的看法。因此，祂一开始并没有像那人那样提出严厉或苛刻的事，如斧头、砍下的树、簸箕、打谷场和不灭的火；但\Emph{祂}的前奏是仁慈的：天堂和天国是祂向听众宣告的喜讯。\par

\Quote{耶稣沿着加里肋亚海行走，看见两个兄弟：称为伯多禄的西满和他的兄弟安德肋，在海里撒网；因为他们是渔夫。祂对他们说：来跟随我，我要使你们成为渔人的渔夫。他们立刻舍了网，跟随了祂。}\BibleRef{玛 4:18}\par

然而，若翰却说他们是以另一种方式被召叫的。由此可知，这是第二次的召叫；并且从许多事上可以察觉这一点。因为那里说，他们来到祂跟前时\BibleQuote{若翰还没有被投进监里}；但这里，却是在他被监禁之后。并且在那里是安德肋召叫伯多禄，但在这里是耶稣召叫两者。若望说，耶稣看见西满前来，就说：\BibleQuote{你是约纳的儿子西满，你要称为\OriginalTerm{刻法}{Cephas}，翻译出来就是磐石。}\BibleRef{若 1:42}但玛窦说他已经被称为那个名字了；因为他写道：\BibleQuote{看见那称为伯多禄的西满}。并且从他们被召叫的地方，以及许多别的事上，可以察觉这一点；也可见于他们即时的服从与舍弃一切。因为如今他们已事先受教良多。因此，在另一种情况下，我们看到安德肋进了祂的房屋，听了许多话；但在这里，他们一听到一句赤裸裸的话，便立即跟随。因为他们起初跟随祂，然后看见若翰被投监、祂自己离开，又离开祂重返本业，并非不合理。因此你看见祂确实发现他们在打鱼。但祂既没有在他们起初想要退开时禁止他们，也没有在他们退出后完全放手；而是当他们离开祂时，祂让开，并再次来赢回他们；这正是捕鱼的要诀。\par

但请注意他们的信德和服从。因为他们虽在工作中（你知道捕鱼是多么贪婪的事），一听到祂的命令，他们不迟延，不拖延，没有说\Quote{让我们回家，与亲属交谈}，而是\BibleQuote{他们舍弃一切跟随了祂}，正如厄里叟跟随厄里亚一样。\BibleRef{列上 19:20}因为基督向我们要求的正是这样的服从：即使有绝对最必要的事紧迫催促我们，我们也不延迟片刻。因此，当另一个人来到祂跟前，请求许可埋葬自己的父亲时，\BibleRef{玛 8:21}连这事祂也不许他去做，以表明我们应将跟随祂置于一切之上。\par

但若你说道\Quote{这应许非常伟大}；正是为此我最钦佩他们，因为当他们尚未见到任何征兆时，就相信了如此伟大的应许，并将一切置于跟随祂之后。这是因为他们相信，借着什么话被捕获，也就能借同样的话捕获别人。\par

对这些人，祂的应许就是如此；但对雅各伯和若望，祂没说这样的话。因为先前之人的服从已经为这些人铺平了道路。此外，他们先前也听说过许多关于祂的事。\par

并且请看他是如何精确细致地向我们透露他们的贫困：祂发现他们在缝补渔网。他们的贫困如此巨大，以至于他们修补破旧之物，无力购买新的。而这一点在当时也是德行不小的证明：他们从容战胜贫困，以诚实的劳动维持生计，因爱的力量彼此相连，拥有他们的父亲与他们同在，并照顾他们。\par

4. 当祂捕获他们后，祂就在他们面前开始行奇迹，以祂的作为证实若翰关于祂的话。祂不断进他们的会堂，以此教导他们，祂并非某种天主的对手和骗子，而是按照圣父而来的。\par

并且，当祂进出会堂时，祂不仅宣讲，也显示奇迹。这是因为，每当有奇异和新奇的事发生，以及任何政制被引入时，天主惯于行奇迹作为祂权能的保证，并将之赐予要接受祂法律的人。例如，当祂要造人时，祂创造了整个世界，然后赐给他那在乐园中的法律。当祂要赐法律给诺厄时，祂又重新显了大奇迹，使整个受造界重归原质，使那可怕的海水泛滥了一年之久；并且，在那大风暴中，祂保全了那位义人。在亚巴郎的时代，祂也赐予了许多征兆，如他在战争中的胜利、法郎身上的瘟疫、他从危险中的得救。当祂要为犹太人立法时，祂显示了那些奇妙伟大的异事，然后赐下了法律。同样在此情形中，祂要引入某种崇高的政制，并告诉他们从未听过的事，便藉奇迹的展示来证实祂所说的话。\par

因此，因为祂所宣讲的国并未显现，祂便藉可见之事使那虽不可见的国彰显出来。\par

并且注意福音作者的谨慎，避免赘言；他并未告诉我们每一个被治愈的人，而是用寥寥数语掠过如雨的奇迹。\par

因为他说：\BibleQuote{他们带到他跟前的一切患各种疾病和痛苦的人、附魔的人、癫痫的人、瘫痪的人，祂都治好了他们。}\par

但我们的问题是：为什么祂没有要求他们中任何人的信德？因为祂没有说我们后来发现祂所说的\BibleQuote{你们信我能做这事吗？}\BibleRef{玛 9:28}，因为祂尚未给出祂权能的证明。此外，接近祂以及带别人到祂跟前的行为本身，就显示了非同寻常的信德。因为他们甚至从远处带来他们；若非他们已确信关于祂的伟事，他们绝不会带来。\par

那么，现在我们也当跟随祂；因为我们灵魂也有许多疾病，而祂尤其愿意治愈这些。因为祂纠正那种疾病，正是为了将我们灵魂中的这些疾病驱除。\par

5. 让我们因此到祂那里去，不要祈求任何属于此世的事物，而只求罪的赦免。因为祂确实现在也赐予这赦免，只要我们真心诚意。因为正如那时\BibleQuote{祂的名声传遍了叙利亚}，现在也传遍了全世界。他们一听说祂医治附魔的人，就都跑来聚集；而你，在经历了祂更多更大的权能之后，难道还不振奋自己，跑去吗？\par

而他们既离开了家乡、朋友和亲属；你难道不能为了亲近并获得远为更大的事物，而忍受哪怕离开你的屋子吗？或者我们并不要求你这么多，而只要求你离开你的恶习，那么你就能轻易得到痊愈，与朋友们一起留在家中。\par

但事实上，如果我们有任何身体的病痛，我们会想尽一切办法来摆脱使我们痛苦的事；但当我们的灵魂不适时，我们却拖延、退缩。正因如此，我们也没有从另一种病痛中得到解脱：因为我们把不可缺少的事变成了次要的，而把次要的事变成了不可缺少的；我们放开了我们疾病的源头，却想净化支流。\par

因为我们的身体疾病是由灵魂的邪恶引起的，这一点可以从瘫痪三十八年的人、从屋顶被缒下的人、以及之前的加音身上看出来；从许多其他事上也可以看出这一点。因此，让我们除掉我们邪恶的源头，那么我们一切疾病的渠道都将止息。因为疾病不仅是瘫痪，还有我们的罪；而这比那更严重，正如灵魂比身体更贵重一样。\par

因此，现在也让我们到祂那里去；让我们恳求祂坚固我们瘫痪的灵魂，并放下一切属于此世的事物，只专注于属灵的事。或者，如果你也依恋这些事物，也要在属灵之事之后才想到它们。\par

你也不可轻视这一点，因为你在犯罪时没有痛苦；反而正因如此你最应当哀悼，因为你感觉不到你过犯的痛苦。因为这不是因为罪不咬人，而是因为犯罪的灵魂麻木不仁。要这样看待那些对自己罪有感觉的人：他们如何比那些被切割或焚烧的人更痛苦地哀号；他们做了多少事，遭受了多少事，如何大大地悲伤和哀悼，以便从他们的邪恶良心中得到解脱。他们不会做任何这样的事，除非他们在灵魂中极其痛苦。\par

因此，最好的事是一开始就避免犯罪；其次，是感到我们犯了罪，并彻底改正自己。但如果我们没有这个，我们怎能向天主祈祷，祈求赦免我们的罪，我们这些不把这些事放在心上的人呢？因为当你自己——你冒犯了天主——甚至不愿意知道这个事实，即你犯了罪，那么你将为哪些过犯祈求天主宽恕呢？因为那些你不知道的过犯？你又怎能知道恩惠的伟大呢？因此，要具体说出你的过犯，好让你知道你得到了什么宽恕，从而能对你的恩主心怀感激。\par

但你，当你冒犯了一个人时，你会恳求朋友、邻居和看门人，花费金钱，花费许多天去拜访和恳求，尽管那被冒犯的人一次、两次、千万次地完全拒绝你，你也不沮丧，反而更加恳切地再三请求；但当万有的天主被冒犯时，我们却张着嘴，向后躺倒，过着奢侈和醉酒的生活，做一切平常的事。那么，我们何时才能平息祂的怒火呢？我们这样岂不是更加激怒祂吗？因为不仅犯罪，而且连犯罪也不痛苦，更激起祂的义怒和愤怒。因此，在这些事之后，我们应该沉入地底，甚至不再见这太阳，不再呼吸，因为我们有这样一位仁慈的主，我们首先激怒了祂，然后对激怒祂毫无悔意。然而，祂确实，即使发怒时，也不是出于仇恨和远离我们，而是为了借此方式赢得我们归向祂。因为如果祂在受辱后继续善待你，你会更加轻视祂。因此，为了不发生这种情况，祂暂时转面不顾，为要永远拥有你与祂同在。\par

6. 现在，我恳请你们，因祂对人的爱而鼓起勇气，并要表现出热切的悔改，趁着那使我们无法再从中获益的日子尚未来临。因为目前一切都取决于我们，但到那时，审判者独自控制着判决。\Quote{让我们因此带着忏悔来到祂面前；}让我们哀恸，让我们悲伤。因为如果我们能在指定的日子之前说服审判者赦免我们的罪，那么我们就不必进入法庭；反之，如果不这样做，祂将在世界面前公开听取我们的案件，我们将不再有任何获得赦免的希望。因为那些没有在此世消除自己罪过的人，当他们离开这里后，将无法逃避对这些罪过的交代；但正如那些从地上监狱中被带出来的人，带着锁链被带到审判台前，同样，所有灵魂，当他们带着自己罪过的多重锁链从这个世界离开时，被带到可怕的审判台前。因为事实上，我们现在的生命并不比监狱好多少。但是，就像我们进入那房间时，看到所有人都被锁链捆着一样，现在如果我们从外在的表象中抽身出来，进入每个人的生命、每个人的灵魂，我们会看到它被比铁更重的锁链捆绑着：尤其如果你进入那些富人的灵魂中。因为他们拥有的越多，就被捆绑得越紧。因此，关于囚犯，当你看到他背上、手上、常常还有脚上都戴着铁镣时，你最觉得他可怜；同样，关于富人，当你看到他周围环绕着无数的事务时，不要因此认为他富有，反而要因这些事情而认为他可怜。因为连同这些锁链，他还有一个残忍的狱卒——对财富的邪恶之爱；这狱卒不让他离开这个监狱，反而为他准备了成千上万的脚镣、守卫、门和门闩；并且当他把他投入内监时，还劝他甚至在这些锁链中感到快乐，这样他就无法找到任何从压在他身上的邪恶中得到释放的希望。\par

如果你在思想中剖开那人的灵魂，你不仅会看到它被捆绑，而且污秽、肮脏、充满蛆虫。因为奢侈的快乐并不比蛆虫好，甚至更可憎，并且更严重地摧毁身体，连同灵魂一起；它们给身体和灵魂带来成千上万的疾病鞭打。\par

因此，为这一切，让我们恳求我们灵魂的救赎者，愿祂断开我们的锁链，除去我们这残忍的狱卒，使我们脱离那些铁链的重担，使我们的精神比任何翅膀更轻快。当我们恳求祂时，让我们也贡献自己的一份：热忱、反省和卓越的热心。因为这样，我们就能在短时间内摆脱现在压迫我们的邪恶，并了解我们先前的状况，抓住属于我们的自由；愿天主使我们所有人都能达到这一点，赖我们的主耶稣基督的恩宠和祂对人的爱，愿光荣和权能归于祂，世世无穷。阿们。\par

\SourceRecord{Translated by George Prevost and revised by M.B. Riddle.；Nicene and Post-Nicene Fathers, First Series；Vol. 10.；Edited by Philip Schaff.；Buffalo, NY: Christian Literature Publishing Co.,；1888.；<http://www.newadvent.org/fathers/200114.htm>.}

% structure: p0001 source=h1 target=section reason=page-title

\section{\WorkTitle{玛窦福音讲道集 第十五讲}}

\Emph{玛 5:1-2}\par

\BibleQuote{耶稣看见群众，就上了山，祂坐了下来，门徒们到祂跟前。祂遂开口教训他们说：有福的，}等等。\par

请看祂是多么淡泊，毫无夸耀：祂并没有一路带领众人，而是当需要治病时，祂亲自四处奔走，遍访城镇和乡村；如今群众甚多，祂便坐于一处，且不在任何城市或广场当中，而是在山上和旷野里；教训我们做事不为显扬，当远离日常生活的喧哗，尤其在我们要研习智慧、谈论应行之事时，更当如此。\par

但当祂上了山，\BibleQuote{祂坐下后，门徒们来到祂跟前。}你看到他们德性上的长进了吗？又如何一时之间便成了更好的人？因为群众不过观看奇迹，但这些人从那时起也渴望听到一些伟大高深的事理。的确，正是这一点促使祂施教，并开始了这番讲论。\par

因为祂不仅医治人的身体，也在修正他们的灵魂；而再从对灵魂的照顾，转而关照身体。如此，祂同时变换所赐的救助，并将祂言语所予的教训与祂工作所显的荣耀交融；此外，祂也堵住了异端者的无耻之口，借着祂对我们存有两部分的照料，表明祂自己是整个受造界的造主。因此，祂对每一本性都赐予了丰厚的眷顾，时而修正其一，时而修正其二。\par

当时祂正是这样做的。因为经上记载，\BibleQuote{祂开口教训他们。}为何加上\BibleQuote{祂开口}这一句？是要告诉你，祂在沉默时也给予教训，并非只在发言时：有时借着\BibleQuote{开口}，有时借着所行的工作发声。\par

但当你听到祂教训他们时，不要以为祂只是对门徒讲话，而是通过他们向众人讲话。\par

由于群众如同群众常有的那样，且由匍匐于地的人组成，祂便将门徒的群体抽离，专对他们发言：在与他们谈话时，使其余尚远远不及祂言语水准的人，也觉得祂关于克己的教训不再令他们反感。这一点路加也暗示了，他说祂将话直接对他们说；玛窦也同样清楚表明，写道：\BibleQuote{门徒们来到祂跟前，祂就教训他们。}因为这样，其余之人也必定比祂对众人讲话时更加热切地留意祂。\par

那么祂从哪里开始？祂为我们新体制的根基立下何等样式？\par

让我们用心聆听所说的话。因为这话虽是对他们说的，却也是为了后来所有人的缘故而记载的。因此，尽管祂公开讲道时心中想到门徒，却不将祂的话局限于他们，而是毫无限制地将所有祝福的话应用于众人。祂并没有说：\BibleQuote{你们若成为贫穷的，就有福了}，而是说：\BibleQuote{贫穷的人是有福的。}我还可以补充说，即使祂是对他们说的，这劝言仍然对一切人普遍适用。正如祂说：\BibleQuote{看！我同你们天天在一起，直到今世的终结}\BibleRef{玛 28:20}，这不只是对他们讲话，也通过他们向全世界讲话。而祂称那些受迫害、被驱逐、忍受一切难堪的人为有福，这冠冕不仅为他们编织，也为一切达到同样卓越的人编织。\par

不过，为使这一点更加清晰，并告诉你们，你们和全人类一样，对祂的言语有极大的兴趣，如果有人愿意留心听的话，请听祂如何开始这些奇妙的话语。\par

\BibleQuote{神贫的人是有福的，因为天国是他们的。}\BibleRef{玛 5:3}\par

所谓\Quote{神贫的人}是什么意思？就是心灵谦卑和痛悔的人。因为这里祂用\Quote{精神}指灵魂和选择的官能。既然很多人并非自愿谦卑，而是被环境所迫；祂略过这些人（因为这并无可称赞之处），首先祝福那些自愿谦卑、自我收敛的人。\par

但祂为何不说\Quote{谦卑的人}，而说\Quote{贫穷的人}？因为这比那个更进一步。祂这里指的是那些敬畏并战栗于天主诫命的人。天主也曾借着先知依撒意亚热切悦纳这样的人，说：\BibleQuote{我要垂顾谁呢？除非是温良安静、对我的话战栗的人。}因为谦卑有多种：一种人是按照自己的尺度谦卑，另一种人是极度地谦卑。正是这最后一种谦卑，那位有福的先知赞许，向我们描述那不仅是受制、且是完全破碎的心境，他说：\BibleQuote{天主的祭献就是痛悔的精神，痛悔和谦卑的心，天主必不轻看。}那三位青年也将此作为大祭献给天主，说：\BibleQuote{但愿我们在痛悔的灵魂和谦卑的精神蒙受悦纳。}这也就是基督现在祝福的。\par

3. 因为最大的恶事，以及那些使全世界陷入混乱的恶事，都是从骄傲进来的。——魔鬼原本并非如此，却因此变成了魔鬼；正如保禄明确宣告的，说：\Quote{免得他被骄傲自大所蒙蔽，而陷于魔鬼所受的判决。}\BibleRef{弟前 3:6}——第一个人也受了魔鬼用这些希望所鼓动，成了鉴戒，变成了必死的（因为他期望成为神，竟连他所拥有的也失去了；天主也以此责备他，嘲笑他的愚昧，说：\Quote{看，亚当已成了我们中的一个}\BibleRef{创 3:22}）；以后每一个来到世上的人，都因妄想与天主平等，而在不敬中毁了自己。——既然，我说，这是我们邪恶的堡垒，一切恶事的根源和泉源，祂为这病症准备了一个合适的疗法，首先立下这条法律，作为坚固而安全的基础。因为若以此为根基，建造者便可安全地在其上建造一切。但如果这根基被除去，即使人在其人生旅途中达到天上，也容易全部崩毁，归于悲惨的结局。即使守斋、祈祷、施舍、节制，以及任何其他善事都聚集在你身上；若没有谦卑，一切都将消逝败坏。\newline{}\par

这正是在那个法利塞人的身上所发生的事。因为即使他到达了至高的顶峰，他却\Quote{下去}\BibleRef{路 18:14}，失去了一切，因为他没有那众德之母；因为正如骄傲是一切邪恶的泉源，谦卑则是所有自制之德的开端。因此祂也从这里开始，从祂听众的灵魂中连根拔起自夸的毒草。\newline{}\par

或许有人会问：\Quote{这与祂的门徒有什么关系？他们在各方面都是谦卑的。的确，他们并没有什么可骄傲的，因为他们是渔夫、贫穷、卑贱、没有学问。}\newline{}即使这些事与祂的门徒无关，也必然与当时在场的人有关，也与以后将要接待门徒的人有关，免得他们因此而轻视门徒。但更确切地说，这些事也与祂的门徒有关。因为即使当时不需要，后来当他们行了神迹奇事，得了世人的尊敬，有了对天主的信心之后，他们必定需要这帮助。因为财富、权势、甚至王权本身，都不及他们所充分拥有的事物更能使人骄傲。此外，甚至在行神迹之前，他们也可能在那个时候，当他们看到群众和所有围绕他们师傅的听众时，感到一些人性的软弱。因此祂立即抑制他们的骄傲。\newline{}\par

祂不是以劝告或命令的方式提出祂所说的，而是以祝福的方式，使祂的话不那么沉重，并向所有人敞开祂训诲的道路。因为祂不是说：\Quote{这个或那个人}，而是：\Quote{凡是这样做的人，都是\Emph{有福的}。}因此，即使你是奴隶、乞丐、贫穷、外方人、没有学问，只要你效法这美德，就没有什么能阻止你成为有福的。\newline{}\par

4. 现在，正如你所看到的，祂在最有必要的地方开始，然后进入另一条诫命，这条诫命似乎与全世界的判断相反。因为所有人都认为快乐的人是令人羡慕的，而沮丧、贫穷、哀恸的人是可怜的，但祂却称这些人为有福的，而不是那些人，祂这样说：\newline{}\par

\Quote{哀恸的人是有福的。}\BibleRef{玛 5:4}\newline{}\par

然而，所有人都称他们为可怜的人。因此祂事先行了神迹，以便在颁布这样的法令时，祂的言语得到信任。\newline{}\par

在这里，祂再次不仅指一切哀恸的人，而且指那些为罪哀恸的人；因为那另一种哀恸，即有关现世生活的哀恸，是被禁止的，而且是严厉禁止的。保禄也清楚表明了这一点，他说：\Quote{世俗的忧愁致死亡，但依照天主的忧愁，却能产生悔改，以致得救，不令人后悔。}\newline{}\par

祂自己也称这些人为有福的，他们的忧愁是那种忧愁；但祂不仅指忧愁的人，而是指深切忧愁的人。因此祂没有说\Quote{他们忧愁}，而是说\Quote{他们哀恸}。因为这条诫命再次适合教导我们完全的自制。因为那些为子女、妻子或任何其他离世的亲属而悲伤的人，在他们悲伤的那段时间里，不会贪恋利益或快乐；如果他们不追求荣耀，不受侮辱所激怒，不被嫉妒所俘虏，也不被任何其他情欲所困，他们的悲伤本身就完全占据了他们；那么，那些为自己的罪而哀恸，如同应当哀恸的人，将表现出比这更大的克己。\newline{}\par

接下来，这些人的赏报是什么？\Quote{因为他们要受安慰。}祂说。\newline{}\par

他们要在哪里受安慰？请告诉我。既在此世，也在来世。因为所吩咐的事是极其沉重和令人痛苦的，祂便应许赐予那最能使之轻省的事。所以，你若愿意受安慰，就要哀恸；不要以为这是难懂的话。因为当天主安慰时，即使忧伤像雪片一样成千上万地临到你，你也能超越它们。事实上，正如天主的报偿总是远远超过我们的劳苦，祂在这事上也是如此，称哀恸的人为有福，不是按他们所做的价值，而是按祂对人的爱。因为哀恸的人为过犯哀恸，这样的人得到宽恕并为自己辩护就足够了。但既然祂充满对人的爱，祂并没有将赏报限于免除我们的惩罚或赦免我们的罪，而是使他们甚至成为有福的，并赐予他们丰厚的安慰。\newline{}\par

但祂命令我们，不仅为自己的过犯哀恸，也要为别人的过犯哀恸。圣人们的灵魂就有这样的气质：梅瑟、保禄、达味都是如此；是的，所有这些人都曾多次为并非自己的过犯哀恸。\newline{}\par

5. \Quote{温良的人是有福的，因为他们要承受土地。}请告诉我，是什么样的土地？有人说是一种寓意的土地，但并非如此，因为我们在圣经中从未找到任何纯粹寓意的土地。那么这句话是什么意思呢？祂提出了一个可感的赏报；正如\Person{保禄}所做的，当他说：\Quote{应孝敬你的父亲和你的母亲，} \BibleRef{弗 6:2} 他加上：\Quote{如此你便在地上长寿。} 而主自己对右盗又说：\Quote{今天你就要与我同在乐园里。} \BibleRef{路 23:43}\par

因此，祂不仅以未来的恩赐激励我们，也以现世的恩赐，为了祂那些较为粗俗的听众，以及那些在追求未来之前先寻求现世恩赐的人。\par

例如，在后面祂又说：\Quote{你要与你的对头和息。} \BibleRef{玛 5:25} 然后祂为这种自制指定了赏报，说：\Quote{恐怕对头随时把你交给审判官，审判官交给差役。} \BibleRef{玛 5:25} 你看到祂是如何恐吓我们的吗？借着感官的事物，借着眼前发生的事。又说：\Quote{凡向弟兄说\OriginalTerm{拉加}{Raca}的，应受议会的裁判。} \BibleRef{玛 5:22}\par

\Person{保禄}也详细陈述了可感的赏报，并在劝勉中使用现世的事物；例如当他论及童贞时。他在那里没有提及天上之事，而是暂时借着现世之事来劝勉，说：\Quote{为了现今的艰难，} 和 \Quote{但我宽待你们，} 以及 \Quote{我愿你们无所挂虑。}\par

因此，基督也是这样将可感的事物与属灵的事物混合。因为温良的人被认为会失去自己的一切，祂却应许相反的事，说：\Quote{不，正是这个人安全地拥有他的财产，即那个不鲁莽、不夸耀的人；而那一种人常常会失去他的祖业，甚至他的性命。}\par

此外，由于在旧约中先知常这样说：\Quote{温良的人要继承土地；} 祂就将他们惯用的言词织入祂的讲论中，以便不至于到处说陌生的言语。\par

祂这样说，并非将赏报限于现世之物，而是将另一类恩赐也与之结合。因为祂在论及任何属灵之事时，并不排斥现世之物；同样，在应许现世之物时，也不将应许限于那一种。因为祂说：\Quote{寻求天主的国，这一切自会加给你们。} \BibleRef{玛 6:33} 又说：\Quote{凡舍弃了房屋或兄弟的，要在今世得到百倍，并在来世承受永生。}\par

6. \Quote{饥渴慕义的人是有福的。} \BibleRef{玛 5:6}\par

是什么样的义德呢？祂指的是全部的德行，或是与贪婪相对的那种特定德行。因为祂正要颁赐关于怜悯的诫命，为了表明我们必须如何显示怜悯，例如，不来自抢夺或贪婪，祂就祝福那些持守义德的人。\par

请看祂以何等强烈的语气说出这话。祂并没有说：\Quote{坚守义德的人是有福的，} 而是说：\Quote{饥渴慕义的人是有福的：} 为使我们不仅泛泛地，而是以全部的热忱去追求它。因为贪财最特有的属性是，我们对饮食的贪爱不及对获取和囤积越来越多的贪爱，祂就命令我们将这种欲望转向一个新的对象：脱离贪财。\par

然后祂指定了赏报，又是来自可感之物，说：\Quote{因为他们要得饱饫。} 因此，既然人们认为富人通常是因为贪婪而变得富裕，祂说：\Quote{不，正好相反：因为是义德成就了这事。所以，只要你行事正义，就不要害怕贫穷，也不要在饥饿面前战栗。因为勒索者正是那失去一切的人，正如那爱慕义德的人，反而安全地拥有众人的财物。}\par

但如果那些不贪图他人财物的人享有如此大的富足，那么那些放弃自己财物的人就更不用说了。\par

\Quote{怜悯人的人是有福的。} \BibleRef{玛 5:7}\par

在我看来，祂此处不仅指那些在施舍钱财上显示怜悯的人，也指那些在行为上怜悯的人。因为显示怜悯的方式多种多样，这条诫命是广泛的。那么其赏报是什么呢？\Quote{因为他们要受怜悯。}\par

这似乎确实是一种平等的回报，但这远超善行本身。因为他们作为人显示怜悯，却从万有的天主那里得到怜悯；人的怜悯与天主的怜悯并不相同；正如邪恶与良善之间的差距有多大，这两者之间的距离也有多大。\par

\Quote{心里洁净的人是有福的，因为他们要看见天主。} \BibleRef{玛 5:8}\par

看，赏报又是属灵的。祂在此称呼\Quote{洁净}，要么指那些达到一切德行、自己不觉任何邪恶的人；要么指那些生活在节制中的人。因为要看见天主，没有比这最后一种德行更需要的了。因此\Person{保禄}也说：\Quote{你们要追求与众人和睦，并要追求圣德，没有圣德，没有人能看见主。} \BibleRef{希 12:14} 祂在此说的是人可能拥有的那种看见。\par

因为有许多人显示怜悯，不做抢夺之事，也不贪婪，却犯奸淫和不洁；为表明单有前者并不足够，祂加上了这一点，与\Person{保禄}写信给格林多人时，为马其顿人作证的意思相同，即他们不仅在施舍上富足，也在一切其他德行上富足。因为论到他们在财物方面所表现的高尚精神，他说：\Quote{他们也把自己献给主，也献给了我们。} \BibleRef{格后 8:5}\par

7. \Quote{缔造和平的人是有福的。} \BibleRef{玛 5:9}\par

在此，祂不仅完全消除了我们之间的争斗与仇恨，更额外要求我们：还应使那些彼此争斗的人和好。\par

并且，祂所附加的赏报是属灵的。那么，它究竟如何呢？\par

\Quote{因为他们将被称为天主的子女。}\par

是的，因为这正是独生子的工作：使分裂者合一，使疏远者和好。\par

然后，免得你以为平安在任何情况下都是祝福，祂又补充道：\par

\Quote{为义德而受迫害的人是有福的。} \BibleRef{玛 5:10}\par

也就是说，为德行、为救助他人、为虔诚的缘故；祂常以\Quote{义德}之名来称呼灵魂的全部实践智慧。\par

\Quote{几时人为了我的缘故而辱骂你们、迫害你们，捏造一切坏话毁谤你们，你们是有福的。你们欢喜踊跃吧！} \BibleRef{玛 5:11}\par

祂仿佛说：\Quote{即使他们称你们为术士、欺骗者、瘟疫般的人，或任何其他，你们是有福的。}祂如此说。还有什么比这些训导更新奇的呢？其中，正是所有人都避之唯恐不及的事——贫穷、哀恸、迫害、恶名——祂却宣称是可取的。然而，祂不仅肯定了这一点，还说服了不是两、十、二十、一百或一千人，而是整个世界。听到这些如此沉重刺耳、与人类惯常方式如此相悖的话，群众\Quote{都惊奇不已。}说话者的能力何等伟大。\par

然而，免得你以为仅仅被人恶言相向就能使人有福，祂设定了两个条件：一是为祂的缘故，二是所说的都是假的；因为若无这些，那被恶评的人非但不是有福的，反而是不幸的。\par

接着，再看那赏报：\Quote{因为你们的赏报在天上是丰厚的。}但是，尽管你在每一端真福中没有听到天国的赐予，也不必沮丧。因为祂虽然给赏报冠以不同的名称，却将一切引向祂的天国。因此，当祂说：\Quote{哀恸的人要受安慰；} \Quote{怜悯人的人要受怜悯；} \Quote{心里洁净的人要看见天主；}缔造和平的人\Quote{要称为天主的子女；}祂通过这些话语所预示的，无非是天国。因为凡享有这些恩惠的人，必然达到那里。所以，不要以为这赏报只属于神贫的人，也属于饥渴慕义的人、温良的人，以及所有其他毫无例外的人。\par

正因如此，祂在一切真福上施予祝福，使你不再寻求任何感官的事物：因为那被与此生一同终结、比影子还快的事物所加冕的人，不能是有福的。\par

8. 但当祂说\Quote{你们的赏报是天上的丰厚的}之后，又附加了另一安慰：\Quote{因为他们也这样迫害了在你们以前的先知。}\par

如此，既然第一个应许——天国的应许——尚未来到，一切尚在期待中，祂便从今世给予他们安慰：从他们与那些先前受迫害之人的共融中得到安慰。\par

\Quote{不要以为，}祂说，\Quote{你们是因为言论和法令中的矛盾而受这些苦；或是因教导恶的教义而被他们迫害。那些阴谋和危险并非源于你们言论的邪恶，而是源于听众的恶意。因此，你们受冤屈并非你们的错，乃是行不义者的错。而这一切的真理，整个过去的时间都作证。因为反对先知，他们甚至没有提出任何违反法律和亵渎思想的指控，就掷石砸死一些人，赶走另一些人，用无数苦难围绕他们。因此，不要让这困扰你们，因为他们现在所做的一切，都是出于同样的心思。}你看见祂如何提升他们的精神，将他们置于梅瑟和厄里亚的行列之中吗？\par

同样，保禄在致得撒洛尼人书中说：\Quote{你们成了在犹太地为天主的各教会的效法者，因为你们也由自己的同胞遭受了同一痛苦，正如他们由犹太人遭受的一样；那些犹太人不但杀害了主耶稣和先知，也驱逐了我们；他们不使天主喜悦，且与众人为敌。} \BibleRef{得前 2:14}基督在此也确立了这一点。\par

而在其他真福中，祂说：\Quote{贫穷的人是有福的}、\Quote{怜悯的人是有福的}；但在此，祂没有一般地说，而是直接对他们说话：\Quote{几时人们辱骂你们、迫害你们，捏造一切坏话毁谤你们，你们是有福的。}这表明这是他们特别的恩惠，并且超越一切，教师将此据为己有。\par

与此同时，祂也在此隐晦地表明了自己的尊威，以及祂与生祂者同等的荣耀。因为\Quote{如同他们为了圣父的缘故，}祂说，\Quote{你们也要为了我而受这些苦。}但祂说\Quote{在你们以前的先知}时，暗示他们如今也已成了先知。\par

接着，为了表明这最有益于他们，并使他们荣耀，祂没有说：\Quote{他们会诽谤迫害你们，但我将阻止它。}因为祂要他们的安全不在于逃避恶名，而在于勇敢地忍受，并用行动驳斥那些诽谤者：这比前者伟大得多；正如被打而不受伤害，比躲避打击要伟大得多。\par

9. 现在在此处祂说：\Quote{你们的赏报在天堂是大的。}但路加记载祂说这话时，既恳切又带着更完整的安慰；因为祂不仅（如你所知）宣告那些为天主而受恶言的人是蒙福的，也同样宣告那些被众人称赞的人是有祸的。因为祂说：\Quote{祸哉，你们！当众人都称赞你们的时候。}然而宗徒们也曾被人称赞，但不是被所有人。因此祂没有说：\Quote{祸哉，你们！当人们称赞你们的时候}，而是\Quote{当众人}如此行的时候：因为那些在德行中生活的人甚至不可能被所有人称赞。\par

祂又说：\BibleQuote{几时人们因我的缘故弃绝你们的名，你们要欢喜踊跃。}\BibleRef{路 6:22}因为祂不仅为他们所经历的危险，也为毁谤，设定了莫大的赏报。因此祂没有说：\Quote{当人们迫害并杀害你们时}，而是\Quote{当人们辱骂你们，并说一切坏话时}。因为毫无疑问，人们的恶言比他们的行为更刺人。因为在我们遭遇危险时，有许多事能减轻辛劳，例如被众人鼓励，有许多人称赞、加冕、宣扬我们的赞美；而在受辱时，这种安慰也被夺走了。因为我们似乎没有成就任何大事；而这比一切危险更折磨战斗者：至少许多人甚至因忍受不了恶言而自缢。你为什么对别人感到惊奇呢？因为那个叛徒，那个无耻又可诅的，那个对任何事都不再羞耻的人，也正是因为这原因主要促使他匆忙去上吊。还有约伯，他全然坚如磐石，比岩石更牢固；当他被夺去一切所有，遭受那些不治之症，突然丧子，并看到自己的身体像泉水一样流出蛆虫，妻子攻击他时，他轻松地抵挡了一切；但当他看到朋友们责备、践踏他，对他怀有恶念，说他因某些罪而遭受这些，为邪恶付出代价时：那时这位伟大而高尚的人就有了困扰，有了动荡。\par

达味也忽略所受的一切，只向天主求那毁谤的报应。因为他说：\BibleQuote{由他咒骂吧，因为上主吩咐了他；愿上主看到我的卑微，就在今天偿还他的这咒骂。}\BibleRef{撒下 16:11}\par

保禄也宣告胜利，不仅属于那些遭受危险或丧失财物的人，也属于这些人，如此说：\Quote{你们要回忆过去的时日，那时你们才受光照，就忍受了艰苦的极大斗争：一面因辱骂与磨难而成了戏景。}因此，基督也为他们设立了莫大的赏报。\par

此后，免得有人说：\Quote{你在这里不给任何补偿，也不堵住人的口；只在那里给予赏报吗？}祂把先知放在我们面前，以表明就是在他们的时代，天主也没有给予矫正。如果在那时赏报临近时，祂用未来的事物安慰他们；那么如今，当这希望变得更加清晰，克己之心也增长了的时候，更是如此。\par

且留意，祂在多少条诫命之后才说了这话，因为祂一定不是无目的地这样做，而是要表明，一个没有准备好、没有用所有其他美德武装自己的人，不可能出去面对这些战斗。因此，你看，祂在每一实例中，以前一条诫命为后一条铺路，为我们编织了一条金链。这样，首先，那\Quote{谦卑}的，必也会为自己的罪\Quote{哀恸}；那这样\Quote{哀恸}的，必会是\Quote{温良}、\Quote{正义}和\Quote{怜悯}的；那\Quote{怜悯}和\Quote{正义}并\Quote{痛悔}的，当然也会是\Quote{心里洁净}的；这样的人也会是\Quote{缔造和平的人}；而那达到了这一切的人，更会武装起来面对危险，当人说他坏话、忍受无数严峻考验时，他不会烦乱。\par

10. 现在，在给予他们应有的劝勉后，祂再次用赞美来振奋他们。如下：诫命崇高，远超旧约中的；免得他们不安和困惑，说：\Quote{我们怎能做到这些？}请听祂说什么：\BibleQuote{你们是地上的盐。}\BibleRef{玛 5:13}这意味着祂极其必要地命令这一切。因为\Quote{不是为了你们自己的生命}，祂说，\Quote{而是为整个世界，你们要交账。因为我不是派你们到两座城、或十座、二十座、或一个民族那里，像派遣先知那样；而是到地、海和全世界，而且是在恶劣的情况下。}因为祂说\Quote{你们是地上的盐}，表明整个人性已\Quote{失了味}，被我们的罪腐蚀了。所以，你看，祂要求他们具备那些对管理大众最为必要和有益的美德。因为首先，那温良、谦让、怜悯、正义的人，不仅将善行封存在自己内，也提供这些善泉涌流出来，以益他人。而心里洁净、缔造和平并为真理受迫害的人，同样安排他的生活方式以利于公共福祉。\Quote{那么，不要以为}，祂说，\Quote{你们被引向普通的战斗，或要为些小事交账。}\Quote{你们是地上的盐。}\par

那么，他们恢复了败坏的吗？绝非如此；因为向已经败坏之物撒盐，也是不可能成任何善事的。因此他们并未这样做。反之，凡先前已得到恢复、交托给他们、并且摆脱了那恶臭的，他们就撒上盐，保存并保持它们在从主所领受的新鲜中。因为使人从罪的腐朽中获得释放，是基督的善工；而使他们不再返回罪中，则是这些人勤勉劳苦的目标。\par

你看见祂怎样逐渐指出他们比先知更优越吗？祂说他们不只是巴勒斯坦的教师，而是全世界的教师；并且不只是教师，而且是令人敬畏的教师。因为奇妙之处在于：他们不是靠谄媚或柔顺，而是靠像盐那样严厉地振奋人心，甚至这样成为众人所爱的。\par

\Quote{如今不要惊奇，}祂说，\Quote{如果我离开了所有其他人，向你们讲话，并吸引你们进入如此巨大的危险。因为想一想我将派你们去管辖多少城市、部落和民族。因此，我愿意你们不仅自己明智，而且也要使他人同样明智。这样的人极需明达，因为其余人的救恩掌握在他们手中：他们应当如此富有德行，甚至能使别人也受益。因为如果你们不成为这样的人，你们连自己都救不了。}\par

\Quote{那么，不要不耐烦，好像我的话太沉重。因为别的失去咸味的人还有可能通过你们得回，但如果你们自己落到这地步，你们就会连同自己一起毁灭别人。所以，你们手中掌管的事务越重大，你们就越需要更大的勤勉。}因此祂说：\par

\Quote{但如果盐失了味，用什么来使它再咸呢？它以后毫无用处，只好被丢在外面，被人践踏。}\BibleRef{玛 5:13}\par

因为别的人即使常常跌倒，或许还能得到宽恕；但教师若遭遇此事，便失去了一切借口，并要受最严厉的惩罚。因此，免得他们听到\Quote{当人辱骂你们，迫害你们，说一切坏话攻击你们时，}就胆怯不敢出去，祂告诉他们：\Quote{除非你们准备好与这一切搏斗，你们就是被白选了。}因为你们不应害怕恶名，而应害怕自己成为装假的同伙。因为那时\Quote{你们会失味，被人践踏；}但如果你们继续严厉地振奋他们，然后被人说坏话，就应当喜乐；因为盐的功用正是刺痛腐败者，使他们感到刺痛。因此他们的指责自会随之而来，丝毫不伤害你们，反而见证你们的坚定。但如果因害怕指责而放弃应有的热忱，你们将遭受更严重的痛苦，既被人说坏话，又被众人轻视。因为这就是\Quote{被人践踏}的意思。\par

11. 此后，祂又引向另一个更高的形象。\par

\Quote{你们是世界的光。}\BibleRef{玛 5:14}\par

\Quote{世界的}——再一次，不是指一个民族，也不是指二十个城邦，而是指整个可居住的大地。而且\Quote{光}是照亮心智的，远胜于这太阳光；他们先前是属灵的\Emph{盐}，如今是\Emph{光}；为叫你明白这些严格训诫的巨大收益和那严肃纪律的益处：它如何约束人，不容人放荡；并使人眼光清晰，引导人走向德行。\par

\Quote{建在山上的城是不能隐藏的；人点灯，也不放在斗底下。}\par

同样，这些话祂训练他们过严格的生活，教导他们要热切努力，如同在众人眼前，在世界这竞技场中争战。因为祂说：\Quote{不看我们现在坐在这里，在一个小角落里。因为你们将像建在山脊上的城、像家中灯台上的灯那样照亮众人，显明给所有的人。}\par

如今那些坚持不信基督德能的人在哪里呢？让他们听这些事，并朝拜祂的大能，惊叹预言的威力。因为想一想，祂向他们应许了多么伟大的事，他们连在自己的家乡都不为人知：地上海洋都要认识他们，他们的名声将传到可居住世界的边缘；或者，不只是名声，而是他们所行善事的实效。因为不是名声将他们带到各处，使他们显扬，而是他们工作的实际证明。因为他们仿佛长了翅膀，比阳光更猛烈地跑遍全地，播种虔诚之光。\par

但在这里，祂似乎也在训练他们放胆说话。因为说\Quote{建在山上的城是不能隐藏的}，是在宣告祂自己的能力。正如那座城无论如何都不能隐藏，同样，他们所宣讲的也绝不能沉入沉默和暗昧之中。因此，既然祂提到了迫害、毁谤、阴谋和战争，恐怕他们认为这些会堵住他们的口，祂就鼓励他们说：不仅不会被隐藏，反而要照亮全世界；并且正因如此，他们会变得显赫和著名。\par

这样祂就宣告了自己的能力。在接下来的话中，祂要求他们应有那份放胆说话的态度，就这样说：\par

\Quote{人点灯，也不放在斗底下，而是放在灯台上，它就照亮屋里所有的人。照样，你们的光也当在人前照耀，叫他们看见你们的善行，光荣你们在天之父。}\par

\Quote{因为我，}祂说，确实，我点燃了这光，但其持续燃烧，要靠你们的勤勉：不单为你们自己，也为那些将要因这些光芒而受益并被引向真理的人。因为毁谤断不能遮蔽你们的光辉，只要你们仍度着严谨的生活，如同那将要改变整个世界的人所当行的。因此，活出一个配得上祂恩宠的生活；好使这光，正如福音被到处传扬，也处处伴随福音。\par

接着，祂向他们提出另一种益处，除了人类的救恩，足以使他们热切奋斗，并引向一切勤勉。正如这样，祂说：\Quote{你们不仅，}祂说，\Quote{如果你们生活端正，将改善世界，而且还将使天主得光荣；反之，如果你们行事相反，就既毁灭人，又使天主的名字受亵渎。}\par

或许有人问：天主怎能因我们得光荣，若人们要毁谤我们？不，并非所有人；甚至那些因嫉妒而行此毁谤的人，良心上也会钦佩并赞同你们，正如那些对恶人表面奉承的人，内心里却谴责他们。\par

那么怎样？你们命令我们为炫耀和虚荣而活？绝不是！我并非说这话；因为我并没有说：\Quote{你们要努力彰显自己的善行}，也没有说：\Quote{显示它们}，而是说：\Quote{你们的光当照耀}。也就是说：\Quote{愿你们的德行伟大，火炽烈，光不可言喻}。因为当德行如此伟大，它不能隐藏，即使追求者将其遮蔽万倍。向人展示无可指摘的生活，使他们没有真实毁谤的理由；那么，即使有成千的毁谤者，没有人能向你投下阴影。祂说\Quote{你们的光}，说得很好，因为没有什么能使人如此卓越，无论他多么愿意隐藏自己，如德行的彰显。因为他仿佛披着阳光本身，他照耀，甚至比阳光更亮；他的光芒不限于地上，也超越天堂本身。\par

因此祂更丰厚地安慰他们。因为，\Quote{虽然毁谤使你痛苦，}祂说；你仍将有许多人因你而尊敬天主。在两方面你都积聚赏报，既因天主因你受光荣，又因你为天主的缘故受辱。这样，免得我们故意寻求被辱骂，当听到有赏报时：首先，祂并非简单地表达那情感，而是有两个限制，即当所说的是虚假的，且是为天主的缘故时；其次祂指明，不仅那样，而且好名声也有其大利，其光荣归给天主。并且祂向我们提出那些恩惠的希望。\Quote{因为，}祂说，恶人的毁谤不足以使所有其他人无法看见你们的光。因为只有当你们\Quote{失了味}时，他们才践踏你们；而不是当你们被妄证、行义时。是的，那时将有许多人钦佩，不仅是你们，而且为你们的缘故，你们的父也受钦佩。祂没有说\Quote{天主，}而是说\Quote{你们的父；}早已预先播下那即将赐予他们的高贵出身的种子。此外，表明祂在尊荣上的平等，正如祂上面所说：\Quote{你们被毁谤时不要忧愁，因为你们为我的缘故被这样毁谤，就够了；}所以在这里祂提到父：处处显明祂的平等。\par

12. 既然我们知道从这热忱而来的益处，以及懈怠的危险（因为若我们的主因我们受亵渎，那比我们的丧亡更糟），让我们\BibleQuote{不给人任何妨碍，无论是犹太人、外邦人，还是天主的教会。}\BibleRef{格前 10:32} 并且，虽然我们呈现在他们面前的生活比太阳还明亮，然而若有人毁谤我们，让我们不要因被辱骂而忧愁，只有当被公正地责骂时才忧愁。\par

因为，一方面，如果我们生活在邪恶中，即使无人说我们坏话，我们也是最可怜的人；另一方面，如果我们致力于德行，即使全世界都说我们坏话，那时我们也比任何人都更令人羡慕。并且我们将吸引所有选择得救的人跟随我们，因为不是恶人的毁谤，而是我们的善行，会吸引他们的注意。确实，号角也不如我们行为所提供的证明那样清晰；光明本身也不如纯洁的生活那样透明，即使毁谤者不计其数。\par

我说，如果上述所有品质属我们所有；如果我们温良、谦卑、慈悲；如果我们纯洁、缔造和平；如果听到辱骂不再还口，反而喜乐；那么我们将吸引所有观察我们的人，不少于奇迹所行的。并且所有人都将善待我们，即使他是野兽、魔鬼，或你所能想象的任何东西。\par

或者即使有人毁谤你们，也不要为此烦恼，也不要因他们在公开场合辱骂你们而介意；但要探查他们的良心，你们会看到他们在鼓掌、钦佩你们，并列举出千万的赞美。\par

请看，例如，\Person{拿步高}怎样赞美在火窑中的青年；诚然他是对手和敌人。但一看见他们高贵地站立，他就宣告他们的胜利，并给他们加冕；这并非因别的，而是因为他们违抗了他，听从了天主的法律。因为魔鬼，当它看见自己一无所成时，就从此离开，害怕它成为我们赢得更多冠冕的原因。当它离开后，即使一个可憎、堕落的恶人也会认明德行，那时迷雾已被除去。或者如果人们仍然固执地争辩，你们将从天主那里得到更大的赞美和钦佩。\par

现在请不要忧伤，我恳求你们，也不要沮丧；因为宗徒们对某些人是\BibleQuote{死亡的芬芳}，\BibleRef{格前 2:16} 对另一些人则是\BibleQuote{生命的芬芳}。若在你身上无可指摘，你便免于他们的一切指控；或者更确切地说，你变得更加有福。因此，要在你的生命中发光，不要理会那些毁谤你的人。因为一个关心德行的人不可能、绝不可能没有许多仇敌。然而，这对有德行的人算不得什么。因为藉此，他的光辉将越发丰盛地增加。\par

那么，让我们牢记这些事，只专注于一件事：如何严谨地管理我们自己的生活。因为这样，我们也将引导那些如今坐在黑暗中的人，走向那里的生命。因为那光明的德能，不仅在此照耀，也能引导其跟随者到达那里。当人们看见我们轻视一切现世之物，并为那将来之事预备自己时，我们的行为比任何言谈都更能说服他们。有谁如此愚昧，看见一个人，一两天前还生活在奢华与财富中，如今却舍弃一切，生出翅膀，整装以待饥荒、贫穷、一切艰辛、危险、流血、屠杀，以及一切被视为可怕的事物，岂不会从这景象中得出关于未来之事的明确证明吗？\par

但若我们纠缠于现世之事，并且越发沉溺其中，又怎能说服他们我们正奔向另一个旅程呢？\par

此后我们还有什么借口呢，如果对天主的敬畏在我们身上不如人的荣耀在希腊哲人身上那样有效？因为他们中有些人确实放弃了财富，并且轻视死亡，为要在人前炫耀；因此他们的希望也成了虚空。既然如此伟大的事物摆在我们面前，如此崇高的克己准则向我们敞开，我们却连他们所行的也不能做到，反倒毁灭了自己和他人，那还有什么托辞能解救我们呢？因为外邦人犯罪造成的损害，远不如基督徒犯同样的罪那样大。当然不是；因为他们的名誉已经丧失，而我们的名誉，因天主的恩宠，即使在恶人眼中也是可敬而荣耀的。因此，当他们最想辱骂我们、加重其恶言时，他们会加上这样的讥讽：\Quote{你这基督徒！}这个讥讽他们本不会说出，除非他们私下对我们教义有很高的评价。\par

难道你没有听见基督颁下了多少伟大的诫命吗？如今你何时才能遵守其中一条诫命，而你却抛弃一切，四处收敛利息，堆砌高利贷，着手经营买卖，购买成群的奴隶，置办银器，购置房屋、田地、无尽的财物？我希望仅此而已。但是当你在这些不合时宜的追求上再加上不义，挪移地界，强夺房屋，加剧贫穷，增加饥饿，你又何时才能踏上这些门槛呢？\par

13. 但有时你怜悯穷人。这我知道，和你知道的一样清楚。但即使在这件事上，祸害也很大。因为你这样做若不是出于骄傲，就是出于虚荣，以至于连你的善行也无法获益。还有什么比这更悲惨的，在港湾里就遭遇船难？为防此患，当你行了任何善事后，不要从我这里寻求感谢，好使天主成为你的债务人。因为，祂说：\Quote{借给那些你不指望偿还的人。}\par

你有你的债务人；为何离开祂，却向我这可怜的、有死的受造物讨债？怎么？那位债务人，当人向祂讨债时，祂会不悦吗？怎么？祂贫穷吗？祂不愿偿还吗？你没看见祂说不尽的宝藏吗？你没看见祂难以描述的慷慨吗？那么，抓住祂，向祂提出要求；因为祂喜欢人们这样向祂讨债。因为，若祂看见另一个人被要求偿还祂自己所欠的债，祂会觉得自己受了侮辱，不再偿还你；不仅如此，祂还会公允地责备说：\Quote{怎么，你定我什么忘恩负义之罪？你知道我有什么贫穷，以至于你从我身边跑开，去投靠别人？你借给了一位，却向另一个人讨债吗？}\par

因为虽然人接受了施舍，却是天主命令你施予；祂的旨意是让自己，按原初的意义，成为债务人和担保人，为你提供无数机会，从各方各面向祂讨债。因此，不要放弃如此大的便利和丰富，转而向我这什么都没有的人索取。怎么？你向我显示你对穷人的怜悯，究竟是何目的？怎么？难道是命令你施舍吗？难道你是从我这里听到这话，以致要向我讨回？祂亲自说过：\BibleQuote{怜悯贫穷的，就是借给上主。} \BibleRef{箴 19:17} 你借给了上主：记在祂的账上吧。\par

\Quote{但祂现在并不全部偿还。}没错，这也是祂为了你的益处而行的。因为祂是这样的一个债务人：不像许多人，只急于偿还所借的；祂却经营万事，为将那已给祂的也稳妥地投资。因此，你看，有些祂在此偿还，有些祂留待他处。\par

14. 因此，我们既然知道这些事，就让我们使自己的慈悲丰盛，以金钱的使用和我们的行为，证明我们极大的爱人之心。若我们看见任何人在市集上受虐待被打，我们若能付钱，就做吧；或者我们能用言语分开他们，就不要退缩。因为就连一句话也有它的赏报，叹息更是如此。有福的\Person{约伯}曾这样说：\BibleQuote{我为一切无助者哭泣，我见人困苦就叹息。} \BibleRef{约 30:25} 但如果眼泪和叹息有赏报，那么再加上言语、焦虑的努力及许多其他事物时，想想赏报会变得多么大。是的，因为我们也曾是天主的仇敌，而独生子和解了我们，祂置身其间，为我们受鞭打，为我们忍受死亡。\par

让我们也同样努力，将那些陷入无数邪恶的人拯救出来；而不是像我们现在这样，当我们看到有人互相殴打、撕裂时，我们却习惯站在一旁，以他人的耻辱为乐，围成一个魔鬼般的竞技场：还有什么比这更残忍的呢？你看到人们被辱骂，互相撕扯，撕裂衣服，打脸，你还能忍受静静地站在一旁吗？\par

什么！难道是熊在争斗？是野兽？是蛇？是人，是一个在各方面都与你有共融的人：一个兄弟，一个肢体。\BibleRef{弗 4:25} 不要旁观，要分开他们。不要以此为乐，而要改正那恶。不要煽动别人来看这可耻的景象，反而要驱散并分开那些聚集的人。这种对灾难的乐趣属于无耻的人和天生的奴隶，属于那些渣滓，无理性的驴。\par

你看到一个人举止不当，难道你不把这不当视为你自己的不当吗？你不介入，驱散魔鬼的队伍，终结人们的痛苦吗？\par

有人说：\Quote{我可能自己挨打}；\Quote{这也是你的命令吗？} 你连这一点也不会遭受；但如果你遭受了，这事对你来说就是一种殉道；因为你为天主而受苦。如果你对挨打反应迟钝，要想到你的主为你忍受十字架并不迟钝。\par

因为他们那边的人已在黑暗中沉醉；愤怒是他们的暴君和指挥官；他们需要某个健全的人来帮助他们，无论是作恶者还是受伤害者；一个是使他免于受苦，另一个是使他停止作恶。所以，你这位清醒的人，靠近并伸手给那个沉醉的人。因为也有一种愤怒的沉醉，比酒的沉醉更严重。\par

你没有看到海员们，当他们看到有人遭遇海难时，是如何张开船帆，全速出发，去把同行的船员从海浪中救出来吗？如果同行业的人如此彼此关心，那么同受同一本性的人更该做这一切！因为事实上，这里也有海难，比那更严重；因为人要么在激怒之下亵渎天主，从而抛弃一切；要么在愤怒的支配下发假誓，从而堕入地狱；要么打人杀人，从而再次遭受同样的海难。去吧，制止这恶；把那些溺水的人拉出来，即使你下到波浪的最深处；在拆散魔鬼的剧场之后，把每个人分开，劝告他平息火焰，平息波浪。\par

但如果燃烧的柴堆变大，火炉更猛烈，不要害怕；因为你有很多人帮助你，只要你提供一个开始，就伸出援手；最重要的是，你肯定有和平的天主与你同在。如果你首先扑灭了火焰，许多其他人也会跟随，对于他们所行的善事，你将亲自领受赏报。\par

请听基督给那些在地上爬行的犹太人颁布了什么诫命：祂说：\Quote{如果你看到你敌人的驮兽跌倒，不要匆忙走过，而要扶起它。} 你必须看到，分开并和好正在争斗的人比扶起跌倒的野兽要轻松得多。如果我们应当帮助扶起我们敌人的驴，那么更应帮助扶起我们朋友的灵魂；尤其当跌倒更严重时；因为这些跌倒不是陷入泥潭，而是陷入地狱之火，因为他们无法承受愤怒的重担。而你，当你看到你的兄弟躺在重担下，魔鬼站在一旁点燃火堆，你却残忍无情地跑过去；这种事即使对野兽也是不安全的。\par

而那撒玛黎雅人看见一个受伤的人，不认识，也与他毫无关系，却停下来，把他扶上牲口，带到客栈，请了医生，付了一些钱，还答应更多；而你，看见一个人跌倒，不是落在强盗手中，而是落在魔鬼的群中，被愤怒包围；这也不是在旷野，而是在广场中央；不需要花钱，不需要雇牲口，不需要长途运送，只需说几句话——你竟迟迟不做？你退缩，残忍无情地匆匆走过？你呼求天主，又怎能期望找到祂的仁慈呢？\par

15. 但让我也对你们说话，你们这些公开羞辱自己的人：对那恶意行凶、作恶的人。你是在施加打击吗？告诉我；还有踢、咬？你变成了野猪和野驴了吗？你不羞愧吗？你这样变成了野兽，背叛了自己的高贵，难道不脸红吗？因为即使你是穷人，你是自由的；即使你是工人，你是基督徒。\par

不，正因为你贫穷，你更应该安静。因为争斗属于富人，不属于穷人；属于那些有许多理由迫使他们征战的人。但你，没有财富的乐趣，却把财富的邪恶——仇恨、纷争、争斗——聚集到自己身上；你掐住你兄弟的喉咙，想勒死他，公开把他摔倒在众人面前：难道你不认为你自己更加丢脸吗？你模仿野兽的暴怒，甚至变得比它们更坏？因为它们一切都共有；它们互相成群，一起行动；而我们却毫无共有，一切都混乱：争斗、纷争、辱骂、仇恨、侮辱。我们既不敬畏我们共同蒙召的天堂，也不敬畏祂留给我们共享的大地，也不敬畏我们自己的本性；愤怒和贪爱钱财却席卷一切。\par

你们没有看见那欠一万塔冷通的，后来他的债被豁免了，却抓住同仆的喉咙，为了一百银钱，遭受了何等大的祸患，并被交付于无穷的惩罚吗？你们没有因这榜样而战栗吗？难道不害怕你们也遭受同样的事吗？因为我们也欠我们的主很多很大的债；然而，他宽容忍耐，既不催逼我们，像我们对待同仆那样，也不掐住我们的喉咙；然而，倘若他定意向我们追讨哪怕最小的一部分，我们早已灭亡了。\par

所以，诸位可爱的，我们要记住这些事，谦卑自下，并对那些欠我们债的人心怀感激。因为他们对我们，如果我们约束自己，就成了获得极其丰富宽恕的契机；我们付出少许，却将获得许多。那么，为什么还要强索呢？即使对方愿意偿还，你也本应主动豁免他，好从天主那里获得全部。但现在你反而做尽一切，强暴好争，好使你的债无一被豁免；当你想着凌辱邻人时，你正把刀刺向自己，从而加重你在地狱中的惩罚；然而，如果你在此稍微自制，你就能使你自己的账目变得容易。的确，天主之所以愿意我们率先行此慷慨，正是为了他有机会加倍地回报我们。\par

因此，凡欠你们债的，无论是金钱还是过犯，都让他们自由离去，并向天主索取你们如此宽宏大量的赏报。因为只要他们继续欠你们债，你们就不能使天主成为你们的债务人。但如果你们释放他们，你们就能留住你们的天主，并要求他丰厚地偿还如此伟大自制的赏报。设想有一个人走来，看见你逮捕你的债务人，便恳求你释放他，并将你对那人的债权转移给自己：在如此豁免之后，他怎能不公平呢？既然他已将全部要求转到自己身上：那么，当为了天主的诫命的缘故，若有人欠我们债，不论大小，我们都不对他们提出任何控告，反而使他们免于一切责任，天主岂能不加倍地、甚至一万倍地报答我们呢？因此，让我们不要想那因向债务人逼债而暂时升起的快乐，而要想那由此将来我们将承受的损失，是何等巨大！这使我们严重损害自己，在永恒的事上。因此，超越一切，让我们豁免那些必须向我们交账的人，无论债务还是冒犯，好使我们自己的账目得以宽容，并且我们藉所有其他德行所不能达到的，借着不对邻人怀怨而获得；如此，藉我们的主\Person{耶稣基督}的恩宠和爱人之情，享受永恒的福乐，愿光荣和权能归于他，从今世直到永远。阿们。\par

\SourceRecord{Translated by George Prevost and revised by M.B. Riddle.；Nicene and Post-Nicene Fathers, First Series；Vol. 10.；Edited by Philip Schaff.；Buffalo, NY: Christian Literature Publishing Co.,；1888.；<http://www.newadvent.org/fathers/200115.htm>.}

% structure: p0001 source=h1 target=section reason=page-title

\section{玛窦福音第十六篇讲道}

\BibleRef{玛 5:17}\par

\BibleQuote{不要以为我来是废除法律或先知。}\par

谁会怀疑这事呢？谁会控告祂，使祂必须为这一指控辩护呢？因为从前面所说的话，绝不会产生这样的怀疑。因为命令人温和、良善、慈悲、心地纯洁、追求正义，并不表示这样的意图，反而完全相反。\par

那么祂为什么要说这话呢？不是随意或空洞地说，而是因为祂即将颁布比旧约更伟大的诫命，说：\BibleQuote{你们一向听过对古人说：不可杀人；我却对你们说：不可发怒；}并指出一条神圣和天上生活的方式；为了使祂所说的话的陌生性不致扰乱听众的心灵，也不致使他们起来反对祂的话，祂就用这种预先纠正的方法。\par

因为他们虽然不遵守法律，但对法律仍怀有极大的尊敬；他们每天都用行为废除法律，却希望法律的字句保持不动，不准任何人再添加任何东西。或者更确切地说，他们容忍自己的首领添加东西，但这不是为了更好，而是为了更坏。因为这样他们就用自己添加的东西来废除对父母的尊敬，并且也以这些不合时宜的添加来摆脱其他许多诫命的约束。\par

因此，因为基督首先不属于司祭支派，其次祂要引入的是一种添加，但并不削弱德行，反而增强它；祂预先知道这两件事会困扰他们，所以在把这些奇妙的诫命写进他们心中之前，先消除了他们心中肯定存在的障碍。那么，他们心中存在的障碍是什么？\par

2. 他们认为祂这样说，是为了废除古代的体制。因此祂医治了这种怀疑；不仅在这里，而且在别处也这样做。例如，因为他们认为祂是天主的敌人，理由之一是祂不守安息日；为了医治他们的怀疑，祂在那里也提出了辩解，其中一些适合祂自己，如当祂说：\BibleQuote{我父工作至今，我也工作；}\BibleRef{若 5:17}有些则充满屈尊，如当祂提出在安息日丢失的羊，\BibleRef{玛 12:11}指出为了救羊而扰乱法律，并再次提到割礼也有同样的效果。\BibleRef{若 7:23}\par

因此，我们看到祂经常说一些比祂身份低微的话，以消除祂是神敌的外观。\par

为此，祂虽曾用一句话叫成千的死人复活，但在召唤\Person{拉匝禄}时，也加上了祈祷；然后，为了不使人们以为祂比生祂的那位更卑微，祂纠正了这种怀疑，说：\BibleQuote{我说这话，是为了站着的群众，使他们相信是你派遣了我。}祂并非一切都凭自己的权能行事，为彻底纠正他们的软弱；祂也并非一切都用祈祷，以免给后来者留下邪恶怀疑的把柄，以为祂无力无能：而是将后者与前者混合，又将前者与后者混合。祂这样做并非不加区别，而是凭自己特有的智慧。因为祂在行更大的事时用权威，在行较小的事时却仰望上天。例如，在赦罪、启示秘密、开启乐园、驱逐魔鬼、洁净癞病人、制服死亡、叫成千死人复活时，祂都是以命令行事；但在行比这些小得多的事，即用少数饼变出许多饼时，祂却仰望上天：表明这样做并非由于软弱。因为能用权威行更大的事，怎么会在小事上需要祈祷呢？但正如我所说，祂这样做是为了堵住他们的无耻。因此，我请你们在听到祂说卑微的话时，也要这样理解。因为那些言行有许多原因：例如，为了不显得与天主无关；为了教导和服侍众人；为了教导谦卑；为了祂被肉体包围；因为犹太人不能一下子完全领会；为了教导我们不要自高。因此，许多时候，祂亲自说了许多关于自己的卑微话，而大事却留给别人说。例如，祂亲自与\Person{犹太人}辩论时说：\BibleQuote{在\Person{亚巴郎}出现以前，我就有。}\BibleRef{若 8:58}但祂的门徒却不这样说，而是说：\BibleQuote{在起初已有圣言，圣言与天主同在，圣言就是天主。}\BibleRef{若 1:1}\par

再者，祂亲自创造了天、地、海和一切有形无形的万物，但祂本人没有在任何地方明确说过；而祂的门徒却直言不讳，毫不隐瞒地肯定这一点，一次、两次，甚至多次，写道：\BibleQuote{万物是藉着祂造成的；}和\BibleQuote{没有祂，什么也没有造成；}以及\BibleQuote{祂在世界内，世界是藉着祂造成的。}\par

如果别人说了比祂自己说更大的事，有什么可惊奇的？因为（更甚的是）在许多情况下，祂所行的祂没有公开说出来。例如，祂曾藉着那个瞎子清楚地表明祂创造了人类；但当祂在起初谈到造人时，祂没有说：\Quote{我造了，}而是\BibleQuote{那创造他们的，造了他们一男一女。}\BibleRef{玛 19:4}再者，祂藉着鱼、酒、饼、海上平静、十字架上遮暗的太阳光，以及许多其他事物，证明祂创造了世界和其中的万物；但在言语上，祂从来没有公开说过这一点，而祂的门徒——\Person{若望}、\Person{保禄}、\Person{伯多禄}——却不断地宣告。\par

因为那些日夜听祂讲论、见祂行奇迹的人，祂私下向他们解释了许多事，并赐给他们如此大的权能，甚至能复活死人；祂使他们如此成全，以致为祂舍弃一切：如果连这些人在领受圣神之恩以前，经过如此伟大的德性和克己，尚且不能完全承受这一切，那么犹太民众，既缺乏悟性，又远不及那样的卓越，只是偶然在祂行事或说话时在场，他们又怎能不被说服，认为祂与万有的天主无关——除非祂自始至终施行了如此伟大的屈尊俯就？\par

因此我们看到，即使当祂废除安息日时，祂也不是故意引入这样\Emph{祂的}律法，而是提出了许多不同的辩护理由。如果祂在使一条诫命失效时，尚且如此谨慎措辞，以免惊吓听众；那么，当祂要在整个律法之外再加上另一整套律法时，祂当然需要更多的处理与关注，以免惊动当时听祂的人。\par

同样，我们也没有发现祂处处清楚地教导关于祂自己的天主性。因为如果增添律法已经足以使他们困惑，那么宣称自己是天主就更会如此了。\par

3. 因此，祂所说的许多话远低于祂应有的尊严；而在这里，当祂要着手增添律法时，祂预先大量使用了矫正之辞。因为祂不仅一次说\Quote{我不废除律法}，而且重复了多次，并添上了另一件更大的事：在\Quote{你们不要以为我来是废除}之后，祂接着又说\Quote{我不是来废除，而是来成全。}\par

这不仅阻止了犹太人的固执，也堵住了那些说旧约盟约来自魔鬼的异端者的口。因为如果基督来是为了摧毁他的暴政，那么这盟约怎么不仅没有被摧毁，反而被祂成全了呢？因为祂不只说\Quote{我不废除它}——尽管这已经足够——而是说\Quote{我甚至要成全它}：这些话表明祂远非反对自己，反而是要建立它。\par

有人可能会问：祂如何没有废除它？祂在哪些方面反而成全了律法或先知？祂成全了先知，因为祂以自己的行动证实了所有关于祂的预言；因此，福音作者在每处都常说：\Quote{这就应验了先知所说的话。}当祂出生时（\BibleRef{玛 1:22}），当孩童向祂唱那奇妙的赞歌时，当祂骑在驴上时（\BibleRef{玛 21:5}），以及在许多其他场合，祂都成就了这同一应验；如果祂没有来，这一切都必不得应验。\par

但祂成全律法，不仅是一种方式，而是第二种和第三种。一种方式：祂没有违犯律法的任何一条。因为祂确实成全了整个律法，请听祂对若翰说的话：\Quote{因为这样，我们应当这样行，以成全各种义德。}（\BibleRef{玛 3:15}）祂也对犹太人说：\Quote{你们中谁能指证我有罪？}（\BibleRef{若 8:46}）又对祂的门徒说：\Quote{这世界的首领来了，他在我身上一无所获。}（\BibleRef{若 14:30}）先知从一开始就曾说祂\Quote{没有行过罪恶}（\BibleRef{依 53:9}）。\par

这是成全律法的第一种含义。另一种，祂也借着我们成全了律法；因为令人惊奇的是，祂不仅自己成全了律法，也将此赐给了我们。保禄在宣告这事时也说：\Quote{基督是律法的终向，使所有信者获得成义。}（\BibleRef{罗 10:4}）他还说：\Quote{祂定了罪恶在肉身中的罪，为使律法的成义，成就在我们这些不随从肉身的人身上。}（\BibleRef{罗 8:3}）又说：\Quote{这样，我们因信德废了律法吗？绝对不可！我们反而坚固了律法。}（\BibleRef{罗 3:31}）因为律法致力于使人成义，却没有能力，祂来了，带来了因信德而获得成义的道路，从而坚固了律法所渴望的；律法无法藉着文字成就的，祂藉着信德成就了。因此祂说：\Quote{我不是来废除律法。}\par

4. 但如果有人仔细探究，还会发现另一种、第三种意义，即这种事已经成就了。这又是怎样的呢？就是指祂将要交给他们的那套未来的律法。\par

因为祂的话并非废除前者，而是将其引申并充实。例如：\Quote{不可杀人}并未因\Quote{不可发怒}的话而被废除，反而被充实并更加稳固；其他一切也是如此。\par

因此，你看，正如祂先前已经不知不觉地播下了这一教导的种子；所以当祂通过比较新旧诫命，会更容易被怀疑为彼此对立时，祂预先用了矫正之辞。因为祂先前所说的话，实际上已经暗中播下了那些种子。例如：\Quote{神贫的人是有福的}，这与我们不可发怒相同；\Quote{心里洁净的人是有福的}，即不可\Quote{怀着贪恋去看女人}；\Quote{不要在地上积蓄财宝}，与\Quote{怜悯人的人是有福的}相和谐；\Quote{哀恸}、\Quote{受迫害}和\Quote{受辱骂}，与\Quote{从窄门进入}一致；\Quote{饥渴慕义}，正是祂后来所说的\Quote{凡你们愿意人给你们做的，你们也要照样给人做}。祂宣告\Quote{缔造和平的人是有福的}之后，又几乎说了同样的话，命令\Quote{留下礼物}，赶紧与那受委屈的人和好，以及关于\Quote{与仇人和好}的事。\par

但在那里祂列出了行善者的赏报，在这里则着重列出了忽视实践者的惩罚。因此，正如祂在那处说：\Quote{温良的人将继承土地；}同样，在这里：\Quote{称自己的兄弟为\InnerQuote{愚笨}的人，应受地狱的罚。}那里说：\Quote{心里洁净的人要看见天主；}这里说，凡以淫荡目光注视妇女的，就已经犯了奸淫。祂在那里称\Quote{缔造和平的人为天主的子女；}在这里祂又从另一方面警告我们，说：\Quote{恐怕仇人把你交给判官。}同样，前一部分祂祝福哀恸的人和受迫害的人；在后一部分，祂确立同一观点，威胁那些不走那条路的人要遭毁灭；因为祂说：\Quote{走\InnerQuote{宽路}的人，他们的结局就在那里。}而\Quote{你们不能事奉天主而又事奉钱财}在我看来与\Quote{怜悯人的人是有福的}以及\Quote{饥渴慕义的人}是一致的。\par

但正如我所说，由于祂将更清楚地说明这些事，不仅更清楚，而且还要比已经说的增加更多（因为祂不再只要求一个怜悯人的人，而是命令我们连外衣也给人；不仅仅是一个温良的人，而且还要转另一面颊给打我们的人）：因此祂首先消除了表面的矛盾。\par

因此，正如我已经说过的，祂不止一次这样说，而是一次又一次；在\Quote{不要以为我来是废除法律}之后，祂又加上\Quote{我不是来废除，而是来成全。}\par

\Quote{我实在告诉你们：即使天地过去了，一撇或一画也决不会从法律上过去，直到一切都实现。}\par

祂所说的大意是：不可能不实现，而是其中最小的事也必须应验。祂自己就做到了这一点，因为祂极其精确地成全了法律。\par

祂在此暗中向我们表明，整个世界的形象也正在改变。祂并非无缘无故地这样写，而是为了惊醒听者，并表明祂有充分的理由引入另一种规范，因为至少受造物的一切工作都将被转变，人类将被召唤到另一个国度，并以更高的方式实践如何生活。\par

5. \Quote{所以，谁若废除这些诫命中最小的一条，并这样教训人，他在天国里将称为最小的。}\BibleRef{玛 5:19}\par

这样，祂消除了邪恶的猜疑，堵住了那些想要反驳之人的口，然后祂终于开始发出警告，并为了支持祂即将颁布的法令而提出了严厉的宣告。\par

\Quote{我告诉你们：}祂说，\Quote{除非你们的义德超过经师和法利塞人的义德，你们决进不了天国。}\par

因为如果祂是针对古代法律发出威胁，祂怎能说\Quote{除非超过}？因为那些与古人同样行事的人，在义德上不可能超过他们。\par

但所需要的超过是什么呢？不轻易发怒，甚至不以淫荡的目光注视妇女。\par

那么，祂为什么称这些诫命为\Quote{最小的}，尽管它们如此伟大崇高？因为祂自己将要引入这些法令；正如祂谦卑自己，常常克制地谈论自己，同样也克制地谈论自己的法令，借此再次教导我们在一切事上要谦逊。此外，由于似乎存在着对新事物的猜疑，祂暂时有保留地安排祂的教导。\par

当你听到\Quote{在天国里称为最小的}时，不要猜想别的，只想到地狱和惩罚。因为祂惯常以\Quote{天国}所指的不仅是其享乐，也是复活的时候和那可怕的来临。如果称自己的兄弟为\Quote{愚笨}并违反一条诫命的人落入地狱，那么违反所有诫命并引诱他人同样行的人，怎能进入天国呢？所以祂的意思并非如此，而是说这样的人在那时将是最小的，即被抛弃的、最后的。而最后的人必定会落入地狱。因为祂是天主，祂预知许多人的松懈，预知有些人会认为这些话只是夸张，会辩论这些法律，说：如果一个人称另一个为\Quote{愚笨}，他会受惩罚吗？如果一个人只是注视妇女，他就算犯奸淫吗？正是为此，祂预先摧毁了这样的傲慢，对两种人——违犯者和引诱他人这样做的人——都设定了最严厉的警告。\par

既然我们知道了祂的威胁，就让我们既不自己违犯，也不劝阻那些倾向于遵守这些诫命的人。\par

\Quote{但谁若实行并教训，}祂说，\Quote{将称为大的。}\par

因为我们不仅应当有益于自己，也当有益于他人；因为引导自己正确的人的赏报，不如既引导自己又引导他人的人大。因为教导而不实行，定教师的罪（\Quote{你教导别人，}经上说，\Quote{却不教导自己吗？}\BibleRef{罗 2:21}）；同样，实行而不引导他人，也会减少我们的赏报。因此，人应当在两方面都卓越，先端正自己，然后才去关心其余的人。因为为此祂自己把实行放在教导之前，暗示这样最能够教导，否则不能。因为有人会对他说：\Quote{医生，医治你自己吧！}\BibleRef{路 4:23}因为那不能教导自己却企图纠正别人的人，会受到许多人的嘲笑。或者，这样的人根本就没有能力教导，因为他的行为在反对他。但如果他在两方面都完善，\Quote{他在天国里将称为大的。}\par

6. \BibleQuote{因为我告诉你们：除非你们的义德超过经师和法利塞人的义德，你们决不能进入天国。}\BibleRef{玛 5:20}\par

此处祂以义德指一切德行；正如论及\Person{约伯}时，祂说：\Quote{他是无瑕的人，是义人。} 根据同一词义，\Person{保禄}也称那人为\Quote{义人}，因为他说，没有法律为这样的人而设。\Quote{因为，他说：\InnerQuote{法律不是为义人设立的。}}\BibleRef{弟前 1:9} 在许多其他地方也可见到这名代表一般的德行。\par

然而，请注意恩宠的增加；祂要祂新来的门徒比旧约的师傅更好。因为此处祂所说的经师和法利塞人，不仅指不守法者，也指行善者。如果他们不行善，祂就不会说他们有义德；祂也不会将不真实者与真实者相比。\par

还要注意，祂如何称赞旧法律，通过将它与新法律比较；这暗示它们是同一种类族系。因为\Emph{更多}和\Emph{更少}属于同一种类。你看，祂并非指责旧法律，而是要它更严格。然而，如果它是邪恶的，祂就不会要求更多，也不会使它更完善，而是会弃绝它。\par

也许有人说：如果它如此，为何不带领我们进入天国？它现在不带领那些在基督来临后生活的人，他们蒙受更多力量，必须为更大的事奋斗。至于它自己的养子女，它却带他们全部进入。因为，祂说：\Quote{从东从西将有许多人来，在\Person{亚巴郎}、\Person{依撒格}和\Person{雅各伯}的怀抱中安卧。}\BibleRef{玛 8:11} \Person{拉匝禄}也获得大奖，被显示住在\Person{亚巴郎}的怀抱中。所有在旧约时代以卓越表现闪耀的人，都是藉着它而闪耀的。如果旧法律有任何邪恶或与祂相异之处，基督自己就不会在来临时完全履行它。因为如果祂这样做只是为了吸引犹太人，而不是为了证明它与新法律相似且一致，那么祂为何不也履行外邦人的法律和习俗，以便也吸引外邦人呢？\par

因此，从所有考虑来看，显然，它未能带领我们进入，不是因为它本身有任何不好，而是因为现在乃是更高规诫的时节。\par

如果它比新法律更不完善，这也并不暗示它是邪恶的；因为按照这个原则，新法律本身也会处于同样的情形。因为事实上，我们对此的知识，与将来那相比，是一种部分和不完全的东西，并在那另一个来临后被废除。\Quote{因为当那完全的来到时，那部分的就必被废除。}\BibleRef{格前 13:10} 正如旧法律通过新法律所遭遇的那样。但我们不能为此责备新法律，尽管当我们达到天国时，它也让位；因为祂说：\Quote{那时那部分的就必被废除。} 但尽管如此，我们仍称它是伟大的。\par

既然它的赏报更大，圣神赐予的能力也更丰富，按理它也要求我们的恩宠更大。因为它不再是流奶流蜜的地方，也不是舒适的晚年，也不是多子多孙，也不是五谷、美酒、羊群和牛群；而是天堂，和天堂中的美物，以及与独生子的义子身份和兄弟情谊，分享嗣业，受光荣，与祂一同为王，以及那些数不清的赏报。至于我们领受了更丰富的帮助，请听\Person{保禄}说：\Quote{如今那些在基督耶稣内的，不随从肉体，而随从圣神的，就不再定罪了；因为生命之圣神的法律，在基督耶稣内使我脱离了罪与死的法律。}\BibleRef{罗 8:1}\par

7. 如今在警告了违犯者，并为行善者设定了重大赏报，并表明祂合理地要求我们比先前的标准更多之后，祂从此开始立法，不是简单地，而是通过与古代法令的比较，想要暗示这两件事：第一，祂制定这些法令，不是要与先前的法令争竞，而是与之非常和谐；第二，祂为它们添加这些第二规诫是合宜且非常合时宜的。\par

为使这一点更加清楚，让我们听从立法者的话。那么祂自己说了什么呢？\par

\BibleQuote{你们听过对古人说：不可杀人。}\BibleRef{玛 5:21}\par

然而，正是祂自己也曾赐下那些法律，但至今祂不具名地陈述它们。因为一方面，如果祂说：\Quote{你们听过我对古人说，} 这句话就会难以接受，并会阻碍所有听众。另一方面，如果在说\Quote{你们听过我父对古人说}之后，祂又补充说\Quote{但我告诉你们}，祂就会显得更加自以为是。\par

因此祂只是简单地陈述了它，从而只得出一个论点：证明祂在合宜的时节来说这些事。因为藉着\Quote{对古人说}这句话，祂指出了时间的长度，自从他们接受这条诫命以来。祂这样做是为了使听众知耻，他们畏缩不前，不愿意进步到祂诫命的更高类别；就像一个老师对一个懒惰的孩子说：\Quote{你不知道你已经花了多少时间学音节吗？} 这样，祂也通过\Quote{古人}这个词含蓄地暗示了这一点，并由此在将来召唤他们到祂教训的更高秩序：好像祂说过：\Quote{你们学习这些功课已经够久了，从此你们必须努力追求比这些更高的。}\par

很好，祂没有打乱诫命的次序，而是从更早的诫命开始，法律也以这诫命开始。是的，这也适合一位在它们之间显示和谐的人。\par

\BibleQuote{但是我告诉你们：凡向自己弟兄发怒的，就要受裁判。}\par

你看见完全权威吗？你看见适合立法者的态度吗？为什么，哪一位先知曾这样说？哪一位义人？哪一位族长？没有一位，而是\Quote{主这样说。}但圣子却不这样。因为他们宣告他们主人的命令，祂宣告祂父的命令。当我说\Quote{祂父的}时，我是指祂自己的。\Quote{因为凡是我的，}祂说，\Quote{都是你的；你的，也是我的。}\BibleRef{若 17:10} 他们有同等身份的同伴来为之立法，而祂是为祂自己的仆人。\par

现在让我们问那些否认法律的人：\Quote{不可发怒}与\Quote{不可杀人}是否相反？或者一条诫命岂不是另一条的成全和发展？显然一条是另一条的成全，而且正因为如此它更伟大。因为那不被激动发怒的人，更会避免杀人；那控制怒气的人，更会约束自己的手。因为愤怒是杀人的根。你看见，祂砍断根，就更是除去枝子；或者更确切地说，根本不容许它们发芽。所以祂订立这些法令，不是为了废除法律，而是为了更完全地遵守它。因为法律规定这些事的目的是什么？岂不是为了不让任何人杀害他的邻人？由此可见，那反对法律的人就会命令杀人。因为杀人是与不杀人相反的。但如果祂连动怒都不容许，那么法律的精神就被祂更完全地确立了。因为那努力避免杀人的人，与那甚至抛弃愤怒的人相比，不会同样地远离杀人；后者离罪行更远。\par

8. 但为了也在另一方面定他们的罪，让我们提出他们所有的指控。那么他们断言什么？他们宣称那创造世界、\BibleQuote{使太阳上升，光照恶人，也光照善人；降雨给义人，也给不义的人}的天主，在某种意义上是一个恶者。但（毕竟）他们中间较温和的人虽然回避这一点，却断言祂是公义的，同时剥夺祂的良善。而另一个人，不是祂，也没有创造任何存在之物，他们将其立为基督的父。他们说那不善良的留在自己的范围内，保护自己的所有；但那善良的却寻求别人的，并突然想要成为那些祂不是其创造者的拯救者。你看见魔鬼的儿女吗？他们从他们父的泉源说话，将创造的工程从天主疏离，而若望呼喊说：\BibleQuote{祂来到自己的领域，}又说：\BibleQuote{世界是藉着祂造成的？}\par

其次，他们批评旧盟约中的法律，它命令\BibleQuote{以眼还眼，}和\BibleQuote{以牙还牙；}他们立刻侮辱并说：\Quote{怎么，祂这样说话怎能是善良的呢？}\par

那么我们对这个问题说什么？这是最高的仁爱。因为祂订立这法律，不是要我们彼此挖眼，而是恐怕因害怕被他人所害，我们就不再对他们做这样的事。就如祂恐吓尼尼微人要倾覆，不是为了毁灭他们（因为如果那是祂的旨意，祂应该沉默），而是为了借恐惧使他们改善，从而平息祂的愤怒；同样，祂也为那些肆意攻击他人眼睛的人规定了惩罚，以便如果善意不使他们避免这种残忍，恐惧可以抑制他们伤害邻人的视力。\par

如果这是残忍，那么制止杀人者和遏制奸淫者也是残忍。但这是无知之人的言语，是疯狂到极点的癫狂之人的言语。因为我非但不认为这是出于残忍，反而应该说，按照人的判断，与此相反才是非法的。而你说，因为祂命令\BibleQuote{以眼还眼，}所以祂是残忍；我说，如果祂没有给予这命令，那么在大众的判断中，祂就会显得是你所说的那样。\par

因为让我们假设这法律完全被废除，没有人害怕随之而来的惩罚，而所有恶人都被允许安然顺从自己的性情，奸淫者、杀人者、伪证者、弑亲者，难道一切不会颠倒吗？难道城市、市场、房屋、海洋、陆地以及整个世界不会充满无数的污秽和谋杀吗？每个人都看到了。因为如果当有法律、恐惧和威胁时，我们的恶劣习性尚且难以抑制；那么即使这安全感被夺走，有什么能阻止人选择邪恶？那时将会有多么大的祸患泛滥于整个人类生活呢？\par

更确切地说，因为残忍不仅在于允许恶人做他们想做的事，也在于同样地忽视、不关心那没有做错事，却无缘无故受苦的人。因为请告诉我：假如有人从各处聚集恶人，给他们武装刀剑，命令他们走遍全城，屠杀所有遇到的人，难道还有比这更野蛮的吗？而如果另一个人将那人所武装的人捆绑、严加拘禁，并从那些无法无天的爪下夺回那些即将被屠杀的人，难道还有比这更仁慈的吗？\par

我现在吩咐你们将这些比喻也应用在法律上；因为那命令拔除\Quote{以眼还眼}的，是将恐惧如同坚固的锁链置于恶人的灵魂上，仿佛那将凶手拘禁在监狱中的人；而那不为他们指定惩罚的，几乎是以这种安全武装他们，扮演着那将刀剑交在他们手中、放纵他们遍及全城之人的角色。\par

你们岂不见，这些诫命远非来自残忍，反而来自丰盛的慈悲？若因这些你们称立法者为严苛难忍，请告诉我，那一种命令更劳苦难当：\Quote{不可杀人}，还是\Quote{连生气也不可}？谁更极端，是为杀人索取刑罚者，还是仅为怒气？谁是在事后惩罚犯奸淫者，还是连淫念也处以刑罚，且是永罚？你们岂不见他们的推理完全颠倒？旧约的天主（他们称为残忍的）将被发现是温和良善的；而新约的那一位（他们承认为良善的）依他们的疯狂将成为严厉难堪的？但我们说，两个盟约只有同一位立法者，祂合宜地安排一切，并根据时代的差异调整两种法律体系之间的差异。因此，前约的诫命并不残忍，新约的也不严厉难堪，而是出于同一眷顾。\par

因为祂自己也赐下了旧约，请听先知的肯定——或者更确切地说（我们必须这样说），那同时是两者的一位所说的：\BibleQuote{我要与你们订立盟约，不像我与你们祖先所立的盟约。}（\BibleRef{耶 31:31}）\par

但如果那患摩尼教\OriginalTerm{摩尼教}{Manichæan}教义的人不接受这个，让他听保禄在另一处所说的同样的话：\BibleQuote{亚巴郎有两个儿子，一个由婢女所生，一个由自由妇人所生；这两个妇女代表两个盟约。}（\BibleRef{迦 4:22}）因此，正如在那例子里妻子不同而丈夫相同，同样这里也有两个盟约，但立法者是一位。\par

为向你们证明这出于同一温柔，在一处祂说：\Quote{以眼还眼}，但在另一处，\par

\BibleQuote{若有人打你的右颊，把另一面也转给他。}（\BibleRef{玛 5:39}）\par

因为正如在那一例子里，祂以受此苦的恐惧来制止那行恶者，在此也一样。\Quote{怎会如此？}有人会说，\Quote{祂竟命令把另一面颊也转给他？}不，这又何妨？祂这样命令并非要除去他的恐惧，而是命你让他完全满足。祂并未说那人不被惩罚，而是说：\Quote{你不要惩罚；}同时既增强了打人者的恐惧（若他继续的话），又安慰了被打者。\par

9. 但这些话我们像是顺便对所有诫命说的。现在我们必须回到我们面前的问题，并保持所肯定之事的线索。祂说：\Quote{向弟兄发怒而无故的，应受裁判的处罚。}因此祂并未完全除去发怒：首先，因为作为人，摆脱情欲是不可能的：我们固然可以驾驭它们，但完全无欲是不可能的。\par

其次，因为如果我们知道在合适的时间使用它，这种情欲甚至是有益的。例如，请看保禄对格林多人发怒的那次著名事件带来了多大的好处；以及这如何既使他们脱离了一场严重的灾祸，又同样借此使迦拉达人（他们也曾偏离）回头；还有其他许多人。那么发怒的适当时间是什么？当我们不是为自己报仇，而是制止他人无法无天的胡作非为，或强迫他们注意自己的疏忽时。\par

那么不合适的时间是什么？当我们是为自己报仇时：这连保禄也禁止，他说：\Quote{亲爱的诸位，你们不可为自己复仇，但应给天主的义怒留有余地。}（\BibleRef{罗 12:19}）当我们为财富而争时：是的，因为这一点他也除去，他说：\Quote{你们为什么不情愿受欺？为什么不情愿吃亏？}（\BibleRef{格前 6:7}）因为正如后一种是多余的，前一种则是必要而有益的。但大多数人却相反：自己受伤害时变得像野兽，但看见别人受欺辱时却懈怠怯懦：这两者都恰与福音的律法相反。\par

因此发怒并不是过犯，但不适时地发怒则是。为此先知也说：\Quote{你们发怒，但不可犯罪。}\par

10. \Quote{凡向弟兄说\InnerQuote{辣加}\OriginalTerm{辣加}{Raca}的，应受议会的裁判。}\par

在此处祂所说的议会是指犹太人的法庭；祂现在提及这一点，是为了使自己不至于处处显得像外来者和标新立异者。\par

但这个字\Quote{辣加}并非表示极大的无礼，而是说话者某种轻蔑和轻视。因为我们吩咐仆人或任何地位低下的人时，说：\Quote{走开；你，来告诉某人：}同样，使用叙利亚语的人说\Quote{辣加}，以此词代替\Quote{你}。但是爱人者天主连最小的过错也拔除，命令我们以体面且合宜的态度彼此相待，并由此也消灭更大的过错。\par

\Quote{但凡说\InnerQuote{愚妄人}的，应受地狱之火的处罚。}\par

对许多人来说，这条诫命似乎严厉难堪，如果仅仅为一个字就要受如此大的惩罚。有些人甚至说这是夸张的说法。但我害怕的是，如果我们在此用言语自欺，到那时在行为上或许会遭受那极重的惩罚。\par

因为，请你告诉我，这条诫命为何显得过于沉重？难道你不知道，大多数惩罚和大多数罪恶都是从言语开始的？是的，因为言语产生亵渎、否认、辱骂、斥责、伪证和作假见证。所以，不要以为它只是一句话，而要查问它是否有极大的危险。你不知道在敌意之时，当愤怒燃烧、灵魂被点燃时，最小的事也显得重大，并不算太羞辱的事也被视为不可忍受？常常这些小事甚至导致了谋杀，颠覆了整座城池。因为正如哪里有友谊，连严重的事也轻省，同样哪里潜伏着敌意，极琐碎的事也显得无法忍受。而且无论多么简单的一句话，都会被推测是怀着恶意说的。如同在火中：如果只有一点火星，即使旁边有成千的木板，也不容易点燃它们；但如果火焰变得强烈而高，它就不但轻易抓住木板，还抓住石头和所有挡路的东西；通常用来灭火的东西，反而使它烧得更旺（因为有人说，那时不仅是木头、麻絮和其他易燃物，就连泼上去的水也只能更加助长其力量）；愤怒也是如此；无论人说什么，顷刻间都成了这邪恶火焰的燃料。基督预先制止了这一切邪恶，首先将无故发怒的人定为受审判的危险（这正是祂说\Quote{发怒的人要受审判的危险}的原因）；然后，将说\Quote{\OriginalTerm{拉卡}{Raca}}的人定为受公议会。但至今这些还不是大事；因为惩罚是在此世。因此，对于称人\Quote{愚妄}的，祂加上了地狱的火，现在首次提到地狱之名。因为祂之前说了许多关于天国的事，直到那时才提到这地方；这表明前者是出于祂自己对人的爱与宽容，后者则是出于我们的疏忽。\par

11. 你看祂如何一步步地施行惩罚，几乎像是在向你们道歉，表明祂的本意并不想威胁这样的事，而是我们迫使祂发出这样的警告。请看：\Quote{我吩咐你们，}祂说，\Quote{不要无故发怒，因为你们要受审判的危险。你们轻视了前一条诫命：看愤怒产生了什么；它立刻引你们到侮辱，因为你们称自己的兄弟为\InnerQuote{\OriginalTerm{拉卡}{Raca}}。我又设了另一个惩罚——\InnerQuote{公议会}。如果你们连这个也忽视，进而做出更严重的事，我就不再用这些有限的惩罚处罚你们，而是用地狱永罚的刑罚，免得此后你们甚至犯下谋杀。}因为世界上没有什么比傲慢更令人难以忍受；它最能刺痛人的灵魂。但若所说的话本身比傲慢更具伤害性，火焰就加倍了。不要以为称别人为\Quote{愚妄}是小事。因为当您剥夺了兄弟那使我们与禽兽区别、尤其使我们成为人的东西——即心智与理智时，您就剥夺了他所有的高贵。\par

因此，我们不要只看言语，而要体认事情本身及其感受，思量这个词造成了多大的创伤，又引向多少邪恶。为此，保禄也把不仅\Quote{奸淫的}和\Quote{柔弱的}，而且连\Quote{辱骂的}都逐出天国\BibleRef{格前 6:9}。这很有道理：因为傲慢的人破坏了爱德的一切美善，给邻人带来无数祸患，制造持久的仇恨，撕裂基督的肢体，天天驱散天主所愿的和平；他借其伤害人的方式给魔鬼提供有利条件，增强它的力量。因此基督亲自砍断魔鬼力量的肌腱，制定了这条法律。\par

因为祂十分重视爱德：这爱德超越一切，是诸善之母，是祂门徒的标志，也是维系我们整个状况的纽带。所以祂极其认真地铲除那完全破坏爱德的仇恨的根源与源头。\par

因此不要以为这些话有任何夸张之处，而要思想它们所带来的益处，并赞叹这些法律的温和。因为天主如此费心，没有比这事更甚的：我们要彼此合一、紧密相连。因此，无论在旧约还是新约中，祂亲自并通过祂的门徒如此重视这条诫命，并严厉惩罚和报复那些忽视这一职责的人。因为事实上，没有比消除爱德更能给一切邪恶开辟入口并扎根的了。因此祂也说：\Quote{当罪恶增多，许多人的爱德必要冷淡。}\BibleRef{玛 24:12}这样，加音成了杀害自己兄弟的人；厄撒乌也是如此；若瑟的兄弟们也是如此；我们无数的罪恶也因这纽带断裂而蜂拥而来。你明白为什么祂自己也从各方面极其精确地根除一切伤害爱德的事物。\par

12. 祂不只停留在已提及的那些诫命上，还添加了更多：由此表明祂对此多么重视。也就是，在威胁了\Quote{公议会}、\Quote{审判}和\Quote{地狱}之后，祂又加上与先前一致的另外的话，这样说：\par

\Quote{若你到祭台前献礼物，在那里想起你的兄弟有什么对不住你的事，就当把你的礼物留在祭台前，先去与你的兄弟和好，然后再来献你的礼物。}\BibleRef{玛 5:23}\par

哦仁慈！哦卓越的愛人之情！祂不計較當受的尊榮，為著我們愛近人的緣故；表明祂先前發出那些威脅，並非出於任何敵意，也非出於懲罰的慾望，而是出於極深的慈愛。因為還有什麼比這些話更溫和呢？祂說：\BibleQuote{讓我的服務中斷，好使你的愛得以持續；因為這一點也是祭獻：你與你的弟兄和好。}是的，為此祂沒有說\BibleQuote{在奉獻之後}或\BibleQuote{在奉獻之前}；而是當禮物正放在那裡、祭獻已經開始時，祂差遣你去與你的弟兄和好；既不移開擺在我們面前的東西，也不在呈上禮物之前，而是當它放在中間時，祂命你快去那裡。\par

那麼祂為何如此命令，目的何在？依我看，祂藉此暗示並預先提供了兩個目的。第一，如我所說，祂願意指出祂高度重視愛德，視其為最大的祭獻；沒有它，祂甚至不接受另一種祭獻；其次，祂強加了無可推諉的和好義務。因為凡被命令在和好之前不得奉獻的人，若不為愛鄰人，也為使這禮物不致未祝聖，就會趕緊跑到那位受委屈的人那裡，消除敵意。為此，祂也以最顯著的方式表達了一切，以警醒並徹底喚醒他。因此，當祂說\BibleQuote{留下你的禮物}時，祂沒有停留在這裡，而是加上\BibleQuote{在祭台前}（再次以這地方使他戰慄）；\BibleQuote{且離去。}祂不僅說\BibleQuote{離去}，還加上\BibleQuote{先去，然後來奉獻你的禮物。}祂以這一切表明，這祭台不接受彼此為敵的人。\par

讓已受洗的人聽吧，凡是懷著敵意接近的；也讓未受洗的人聽吧；是的，因為這話也與他們有關。因為他們也奉獻禮物和祭獻：我指的是祈禱和施捨。因為這也是祭獻，聽先知怎麼說：\BibleQuote{讚頌的祭獻將光榮我}；又說：\BibleQuote{向天主奉獻讚頌的祭獻}；以及\BibleQuote{舉起我的手是晚上的祭獻。}因此，即使你以這樣的心境奉獻的只是一次祈禱，你最好留下你的祈禱，先去與你的弟兄和好，然後再奉獻你的祈禱。\par

因為這一切都是為了這個目的：為了這個目的，甚至天主也降生成人，並安排了所有那些工作，為了使我們和好。\par

在此處，祂將犯錯者打發到受委屈者那裡；在祂的祈禱中，祂引導受委屈者走向犯錯者，並使他們和好。因為正如那裡祂說\BibleQuote{寬免人們的債}，這裡祂說：\BibleQuote{如果他對你有什麼不滿，你就去他那裡。}\par

或者更確切地說，在我看來，這裡祂似乎也在打發受委屈的人：為此緣由，祂沒有說\BibleQuote{使你與你的弟兄和好}，而是說\BibleQuote{你要被和好。}雖然這話似乎指向侵略者，但實際上全部指向受委屈之人。如此，\Person{基督}說：\BibleQuote{如果你與他和好，通過你對他的愛，你將使我也可親近，並能極有信心地奉獻你的祭獻。但你若仍然生氣，要想到：我甚至樂於命令我的事物被輕視，好使你們成為朋友；願這些思想平息你的怒氣。}\par

祂沒有說：當你遭受了較大的不義時，才和好；而是說：\BibleQuote{即使他對你有微不足道的不滿。}祂也沒有加上\Quote{無論有理無理}，而只是說\BibleQuote{如果他對你有什麼不滿。}因為即使有理，你也不應延長敵意；因為\Person{基督}也公正地對我們發怒，然而祂仍將自己交出，為我們被殺，\BibleQuote{不算那些過犯。}見：\BibleRef{格後 5:19}\par

為此，\Person{保祿}在另一個地方勸導我們和好時，也說：\BibleQuote{不要讓太陽在你們的憤怒中落下。}見：\BibleRef{弗 4:26}。因為正如\Person{基督}用祭獻的論證，\Person{保祿}也用日子的論證，催促我們達到同一個點。因為實際上他害怕黑夜，怕它單獨襲擊受傷的人，使傷口更大。因為白天有許多人分散注意力、拉開他；夜裡，當他獨自一人、自己思考時，波濤洶湧，風暴變得更猛烈。因此，你看，\Person{保祿}為了防止這種情況，寧願把他已經和好地交給黑夜，免得魔鬼之後有機會利用他的孤獨重新點燃他憤怒的火爐，使其更猛烈。同樣，\Person{基督}也不允許任何延遲，即使時間很短，免得祭獻完成後，這樣的人變得更加鬆懈，日復一日地拖延；因為祂知道這種情況需要非常迅速的處理。正如一位高明的醫生，不僅展示我們疾病的預防措施，也展示其糾正措施，祂也是如此。同樣，禁止我們稱呼\BibleQuote{傻子}是敵意的預防措施；而命令和好則是消除敵意所產生病症的方法。\par

請注意兩條命令是如何帶著迫切提出的。因為正如在前者中祂威脅地獄，這裡在和好之前祂不接受禮物，表示極大的不悅，並以所有這些方法摧毀根和果實。\par

首先祂說：\BibleQuote{不要生氣}；之後：\BibleQuote{不要辱罵。}因為的確兩者互相增強：敵意產生辱罵，辱罵產生敵意。為此，祂現在醫治根，現在醫治果實；首先阻止邪惡最初萌芽；但如果它可能發芽並結出最邪惡的果實，那麼祂必以各種方式將其燒得更徹底。\par

13. 所以，你看，祂先提到了审判，然后公议会，然后地狱，并谈及祂自己的祭献，然后又添加了其他主题，这样说：\par

\Quote{赶快与你的对头和好，当你还在路上与他同行的时候。} \BibleRef{玛 5:25}\par

也就是说，免得你们说：\Quote{那么，如果我受了伤害呢？} \Quote{如果我被抢劫，还被拖到法庭前呢？} 连这个借口和托词祂也除去了，因为祂命令我们甚至在这种情况下也不可怀有敌意。然后，既然这个命令是重大的，祂从眼前的事物中引出劝告，这些事物比未来的事物更能约束较为粗鲁的人。\Quote{为什么，你说什么？}祂说。\Quote{你的对头比你强，并且冤枉了你？当然，如果你不和解，反而被迫上法庭，他会更加冤枉你。因为在第一种情况下，放弃一些钱财，你就能保住人身自由；但当你受到法官的判决时，你将被捆绑，并支付最重的罚金。但如果你避免那里的争斗，你将收获两个好处：首先，不必遭受任何痛苦；其次，此后所行的善事将是你自己的行为，不再是他强迫的结果。但如果你不听从这些话，你伤害的不是他，而是你自己。}\par

请看看祂如何催促他；因为说过\Quote{与你的对头和好}之后，祂又加上\Quote{快点}；祂还不满足于此，甚至在这迅速上祂还要求更进一步，说\Quote{当你还在路上与他同行的时候}；以此大力催迫和催促他。因为没有什么比迟延和拖延行善更使我们的生命颠倒的了。不，这常常使我们丧失一切。因此，正如保禄所说的：\Quote{在日落之前，消除敌意}；又如祂自己在上文所说的：\Quote{在献礼完成之前，先与他和好}；所以祂在此处也说：\Quote{快点，当你还在路上与他同行的时候}，在你到达法庭大门之前；在你站在被告席上，从此受审者管辖之前。因为在你进去之前，一切都在你自己的掌控之中；但如果你踏上那道门槛，即使你再怎么努力，也无法按自己的意愿安排事务，因为你已受制于他人。\par

但什么是\Quote{达成协议}？祂的意思要么是\Quote{宁可同意受冤枉}？要么是\Quote{如此辩护你的案件，就好像你处在对方的位置}；这样你就不会因自私而败坏正义，反而要如同考虑自己的事一样考虑别人的事，然后才对此事投出你的票。如果这是一件大事，不必惊奇；因为祂正是为了这个目的，才先展示了这一切恩赐，预先抚慰和预备听众的灵魂，使他更能接受祂的一切诫命。\par

现在有些人说，祂用对头的名号隐约指代魔鬼自己；并吩咐我们不可拥有他的任何东西（他们认为这就是与他\Quote{达成协议}）：因为离开此世之后，不可能有任何和解，也没有任何事物等待我们，只有那惩罚，任何祈祷都无法解救。但在我看来，祂似乎是在谈论这世上的判官，以及通往法庭的路，还有这监狱。\par

因为在以更高的事和未来的事令人羞愧之后，祂也以今生的事来警醒他们。保禄也同样行事，用将来和现在来影响他的听众：例如，当他劝人远离邪恶时，他给那倾向恶的人指出武装的掌权者，如此说：\Quote{如果你作恶，就应害怕；因为他不是徒然佩剑；他是天主的仆役。} \BibleRef{罗 13:4} 此外，他又吩咐我们应服从他，所提出的不仅是对天主的敬畏，也有对方的威胁和他的警戒。\Quote{因为你们必须服从，不单是为了愤怒，也是为了良心。} \BibleRef{罗 5:5} 因为正如我所说的，较愚钝的人更容易被这些显现于眼前的事物所纠正。因此基督也提到了，不仅是地狱，还有法庭，以及被拖到那里，还有监狱，和那里的一切痛苦；借此除去了杀戮的根源。因为那既不辱骂，也不提起诉讼，也不延长敌意的人，怎么会犯杀人罪呢？所以由此可见，我们邻人的利益正是我们自己的利益。因为那与对头和好的人，将更大地造福自己；他因自己的行动，从法庭、监狱和那里的悲惨中得获自由。\par

14. 那么，让我们听从祂的话；不要反对自己，也不要争论；因为首先，甚至在它们回报之前，这些诫命本身就含有快乐和益处。即使对大多数人来说，它们似乎是沉重的，所引起的麻烦也很大；但你要想着，你是在为基督而做，那么痛苦就会变得愉快。因为如果我们始终持定这种算计，我们将不会经历任何沉重，反而从各处收获巨大的快乐；因为我们的辛劳将不再显得是辛劳，反而越是增加，就越变得甘美和愉快。\par

因此，当恶习的惯性以及对财富的贪欲不断迷惑你时，你当以那告诉我们\Quote{我们将因轻视短暂之乐而获得巨大赏报}的思维与之作战；并对你的灵魂说：\Quote{你因我夺去你的快乐而沮丧吗？不，要喜乐，因为我正将你引入天堂。你这样做不是为人的缘故，而是为天主的缘故。所以，请忍耐片刻，你将看到获益何等巨大。为今生忍耐，你将获得无可言喻的勇气。}因为如果我们这样与自己的灵魂交谈，不仅考虑德性中艰难的一面，也思量随之而来的荣冠，我们就能迅速使灵魂远离一切邪恶。\par

因为如果魔鬼用暂时的快乐和永久的痛苦尚且如此强大并得胜，而我们的情况正好相反：劳苦是暂时的，快乐与利益是永生的，那么，在如此巨大的鼓励之后，我们若仍不追随德性，还有什么借口呢？因为我们劳苦的目标足以抵挡一切，而且我们清楚地知道，我们是为天主的缘故忍受这一切。因为一个人若使君王成为他的债务人，他便会认为这足以保障他的一生；更何况那使仁慈而永恒的天主成为自己的债务人、因大小善行而欠债的人，他的福气将是何等巨大呢？所以，不要向我提劳苦和汗水；因为天主不仅借着对未来的希望，也通过另一种方式使德性变得容易：他在各方面帮助我们，亲手扶助我们的工作。只要你贡献一点热心，其余一切都会随之而来。为此，他要你也稍作努力，甚至为了使胜利也属于你。正如一位君王要他的儿子亲自出现在阵列中，要他张弓射箭并展示自己，好让胜利的锦标归他所有，而实际上君王亲自完成了全部；同样，天主在与魔鬼的战争中，只要求你一件事：你要对敌人表现出真诚的憎恨。如果你将此献给他，他便亲自结束整场战争。尽管你被愤怒、贪财或任何暴虐的私欲所燃烧，只要他看到你只身脱去这些，并准备与之对抗，他就会迅速来到你身边，使一切变得容易，使你凌驾于火焰之上，如同古时在巴比伦火窑中的那些孩童；因为他们进入时除了善意之外一无所有。\par

为了使我们也能在此熄灭那混乱快感的一切火焰，从而逃脱那里的地狱，让我们每天以此为思考、挂虑和实践，借着我们对善工的完全决心和不断的祈祷，吸引天主的恩宠。因为这样，现在看似不堪忍受的事，将变得最容易、最轻松、最可爱。因为我们还在私欲中时，便认为德性崎岖、阴郁、艰巨，而恶习则令人渴望、最讨人喜欢；但如果稍微摆脱这些，那么恶习就会显得可憎、丑陋，德性则会变得容易、温柔、非常令人向往。从那些已经达致德性的人身上，你可以清楚地学到这一点。例如，听听保禄在脱离那些私欲之后如何以之为耻，他说：\BibleQuote{你们那时在你们现今所羞耻的事上有什么果实呢？}\BibleRef{罗 6:21}但保禄在劳苦之后仍宣称德性是轻松的，称我们患难中的劳苦是暂时的、\Quote{轻的}，并因自己的苦难而喜乐，因患难而夸耀，因自己为基督的缘故所受的印记而自豪。\par

为了使我们也能养成这种习惯，让我们每天按所说的话来规范自己，\BibleQuote{忘记背后，努力面前的，向着标竿直跑，要得天主在基督耶稣里从上面召我来得的奖赏。}\BibleRef{斐 3:13}愿天主恩赐我们众人达致这一目标，借着我们的主耶稣基督的恩宠与对人的慈爱，愿光荣与权能归于他，世世无穷。阿们。\par

\SourceRecord{Translated by George Prevost and revised by M.B. Riddle.；Nicene and Post-Nicene Fathers, First Series；Vol. 10.；Edited by Philip Schaff.；Buffalo, NY: Christian Literature Publishing Co.,；1888.；<http://www.newadvent.org/fathers/200116.htm>.}

% structure: p0001 source=h1 target=section reason=page-title

\section{玛窦福音第十七篇讲道}

\BibleRef{玛 5:27-28}.\par

\Quote{ \Emph{你们听过古人的话说：}不可奸淫；\Emph{但我对你们说：凡注视妇女，有意贪恋她的，他已在心里奸淫了她。} }\par

现在，祂完成了前一条诫命，并将其提升到克己的高度；祂继续前进，按顺序进入第二条，也在此遵守法律。\par

可是，也许有人说，这不是第二条，而是第三条；因为第一条不是\Quote{不可杀人}，而是\Quote{上主你的天主是唯一的主}。\par

因此，也值得探究，祂为何不从那一条开始。那是为什么呢？因为，如果祂从那里开始，就必须也扩充它，并把自己与圣父一同引入。但当时还不是教导关于祂自己的这种事情的时候。\par

此外，祂有一段时间只宣讲道德教义，旨在通过这最初的和祂的奇迹，使听众相信祂是天主子。现在，如果祂在说话或做任何事之前就立刻说：\Quote{你们听过古人的话说：\InnerQuote{我是上主你的天主，除我之外没有别的神}，但我对你们说：要敬拜我如同敬拜祂一样。} 这会使所有人都视祂为疯子。因为，即使在祂教导和如此大的奇迹之后，当祂还没有公开说这话时，他们尚且说祂附有魔鬼；\BibleRef{若 8:48} 如果祂在这一切之前就尝试说任何这样的话，他们还会说什么呢？他们还会想什么呢？\par

但是，在适当的时机保留关于这些主题的教导，祂使得这道为大众所接受。因此，现在祂快速略过它，但当祂通过奇迹和最卓越的教导到处确立了它之后，祂后来也在言语上揭示了它。\par

但目前，通过祂奇迹的显示和祂教导的方式，祂偶尔、逐渐而安静地展开它。因为祂如此制定法律，并如此权威地修正法律，会逐步引导专注而理解的听众走向祂教义之言。因为经上记载：\Quote{他们都惊奇祂的教训，因为祂教训他们，不像他们的经师。} \BibleRef{玛 7:28}\par

2. 因为从那些最属于我们全人类的激情开始，我指的是愤怒和欲望（因为这两者主要在我们内绝对支配，且比其他更自然）；祂以立法者应有的极大权威，严厉地纠正并规范了它们。因为祂没有说只有奸淫者受惩罚；而是祂对杀人者所做的，此处也做了，甚至惩罚不洁的注视：为了教导你，祂比经师多出了什么。因此，祂说：\Quote{凡注视妇女，有意贪恋她的，他已在心里奸淫了她。} 也就是说，那些以好奇明亮身形、搜寻优雅容貌、以目光宴饮灵魂、将双眼固定在美丽面容上的人。因为祂来不仅要拯救身体脱离一切恶行，也要拯救灵魂先于身体。因此，由于我们在心里接受圣神的恩宠，祂首先洁净它。\par

也许有人说：\Quote{如何可能摆脱欲望呢？} 我回答，首先，如果我们愿意，甚至这也可以被削弱并保持不活跃。\par

其次，祂在这里并没有绝对地除去欲望，而是除去因注视而产生的欲望。因为那好奇注视漂亮面容的人，他自己主要是那激情的火炉的点燃者，并使自己的灵魂成为俘虏，且很快也会走向行为。\par

因此我们看到为何祂没有说\Quote{凡有意奸淫的}，而是说\Quote{凡注视以致贪恋的}。在愤怒的情况下，祂作出了某种区分，说\Quote{无缘故的}和\Quote{无故的}；但此处不是这样；反而祂一次而永远地除去了欲望。然而，两者都是自然植入的，且两者都为了我们的益处安置在我们内：愤怒和欲望；愤怒是为了我们可以惩罚邪恶，纠正那些行为不端的人；欲望是为了我们可以有孩子，并且我们的种族可以通过这种延续得到补充。\par

那么，为何祂在此处也没有做出区分呢？不，如果你留意，此处也包括了一个很大的区分。因为祂没有简单地说\Quote{凡欲望的}，因为即使坐在山上也可能有欲望；而是说\Quote{凡注视以致贪恋的}；也就是说，那些将贪恋聚集到自己身上的人；那些在没有强迫时，将野兽带入平静的思绪中的人。因为这不再出于本性，而是出于放纵。甚至连古圣经从起初就纠正这一点，说：\Quote{不要注视别人的美貌。} \BibleRef{德 9:8} 然后，免得有人说：\Quote{那么，如果我注视而不被俘虏呢？} 祂惩罚注视，免得你信赖这种安全而某时陷入罪恶。\Quote{那么，}也许有人说：\Quote{如果我注视，并且确实有欲望，但没有行恶呢？} 即便如此，你也被列在奸淫者之中。因为立法者已如此宣告，你不能再问更多。因为如此注视一次、两次或三次，你或许能克制；但如果你持续这样做，并点燃火炉，你必然会被俘虏；因为你的地位并没有超越人类共有的本性。正如我们，如果看见一个孩子拿着刀，即使没有看到他受伤，我们也会打他，并禁止他再拿刀；同样，天主甚至在行为之前就除去不洁的注视，免得你何时也陷入行为。因为那一次点燃火焰的人，即使他注视的女子不在场，也会不断在脑海中形成羞耻的影像，并常常从这些影像走向行为。因此，基督甚至除去了仅存于心中的拥抱。\par

如今，那些拥有童贞女伴的人能说什么呢？照此法律，他们每天以欲望注视她们，便犯了千万次奸淫。为此，圣约伯 \BibleRef{约 31:1} 从一开始就订立了这一法律，从各方面阻挡自己这种注视。\par

因为实际上，注视渴望的对象而不拥有它，是一场更大的斗争；我们从注视中获得的快乐，并不如我们因加强欲望而遭受的伤害大；这样，我们使对手强壮，给魔鬼更多空间，再也无法击退他，既然我们已经将他带入内里，将我们的心灵向他敞开。因此祂说：\Quote{你的眼目不可犯奸淫，你的心思就不会犯奸淫。}\par

因为人也可以有另一种注视，如贞洁者的目光；所以祂并非完全禁止观看，而是禁止那伴随着欲望的观看。如果祂不是这个意思，祂就会简单地说：\Quote{凡注视妇女的。}但祂现在没有这样说，而是说：\Quote{凡注视以生贪念的}，\Quote{凡注视以满足眼目的。}\par

因为天主造你的眼睛，绝不是为了让你们由此引入奸淫，而是为了使你们在观看祂的受造物时，赞叹造物主。\par

正如一个人可以无故发怒，同样也可以无故投以目光；那就是当你出于情欲时。反之，如果你想要观看并享受乐趣，就看你自己的妻子，不断爱她；没有法律禁止。但如果你好奇于别人的美貌，你就既伤害了你的妻子（让你的眼睛游荡到别处），又伤害了你所注视的她（因为非法接触了她）。虽然你没有用手触摸她，却用眼睛爱抚了她；因此这也算作奸淫，并且在大罚之前，本身就有不小的惩罚。因为他内心充满了不安和动荡，风暴巨大，痛苦最甚，没有俘虏或锁链中的人比处于这种心态的人更糟糕。常常是，射箭者已经飞走了，伤口却仍然存在。或者不如说，不是她射了箭，而是你因不贞的目光给了自己致命的一击。我说这话是为了让端庄的妇女免于责任：因为实际上，如果一个人打扮自己，引诱遇到她的人的目光，即使她没有击伤遇见她的人，她也应受极重的惩罚：因为她调了毒药，准备了毒芹，即使她没有递上杯子。或者不如说，她也递上了杯子，尽管没有人喝。\par

3. \Quote{那么，祂为何不也与她们谈论呢？}有人会说。因为祂所规定的法律在每种情况下都是共同的，尽管祂似乎只对男人说话。因为在与头说话时，祂也使祂的劝诫对整个身体有效。因为祂视女人和男人为一个活物，从不在其中区分性别。\par

但如果你也想听祂对她们特别谴责的话，请听依撒意亚 \BibleRef{依 3:16} 用许多话斥责她们，嘲笑她们的服饰、容貌、步态、拖地的衣裳、急促的脚步、低垂的脖子。也请听他同圣保禄一起为她们制定许多法律；这些法律既\Emph{关于服装和金饰}，又关于编发、奢华生活以及所有类似之事，都严厉责备这一性别。基督在接下来的话中也隐约显示了这一点；因为祂说：\Quote{剜出并砍下那冒犯你的眼睛}，这表示祂对她们的愤怒。\par

3. 因此祂接着说：\par

\Quote{若你的右眼冒犯你，就剜出来，丢掉它。}\par

这样，免得你说：\Quote{但她若是我亲戚呢？若她因其他方式属于我呢？}因此祂给出了这些命令；不是谈论我们的肢体——远非如此——因为祂从未说过我们的肉身应受责备，而是处处指责邪恶的思想。因为看的不是眼睛，而是心思和意念。例如，我们常常完全转向别处，眼睛就看不见在场的人。所以这件事不完全取决于眼睛的工作。再者，如果祂是在谈论身体的肢体，祂就不会只说一只眼睛，而且只说右眼，而是两只眼睛。因为被右眼冒犯的人，显然也会因左眼遭受同样的邪恶。那么，祂为何提到右眼并加上手呢？为了向你表明，祂不是在谈论肢体，而是在谈论那些亲近我们的人。因此，祂说：\Quote{你若如此爱一个人，如同他是你的右眼；你若认为他对你如此有益，如同你的一只手，而他却伤害你的灵魂，就连这些也要砍掉。}并且请注意强调：祂没有说\Quote{离开他}，而是为了显示分离的完全，祂说：\Quote{剜出来，丢掉它。}\par

然后，因为祂的命令严厉，祂也藉着继续比喻，从益处和祸患两方面指出得失。\par

\Quote{因为这对你有利，}祂说：\Quote{你的一个肢体丧失，而不是你全身被投入地狱。}\par

因为如果他既不能救自己，也不能不毁灭你，那么两人都沉沦有什么好处呢？如果分开，至少一个可能得救。\par

但保禄那时为何选择成为诅咒 \BibleRef{罗 9:3}呢？不是因为自己一无所获，而是为了他人的得救。但在此案中，祸害涉及双方。因此祂不仅说\Quote{剜出来}，还说\Quote{丢掉}：不再接纳他，如果他仍保持原样。因为这样你既使他免于更重的责任，也使你自己免受毁灭。\par

但是，为使你们更清楚地看到这法律的益处，若你们同意，我们不妨就以身体本身为例来试一试。我说，如果你面临选择：要么保留你的眼睛，却掉进坑里丧命；要么剜出眼睛，保全身体其余部分。你难道不自然会选择后者吗？这是人人都明白的。因为这并非恨你的眼睛，而是爱你的身体其余部分。对他人，男人和女人，你们也应作同样的估计：如果那因友谊而伤害你的人持续不可救药，那么将他这样割除，既能使你摆脱一切祸害，他自己也能免于更重的指控，不再因你的毁灭连同他自己的恶行而负责。\par

你们看见这法律充满了温良和关怀，而一般人以为严酷之处，却显露出何等深切的仁爱吗？\par

让那些匆匆奔向剧场、每日自陷奸淫的人听听这些吧。如果法律命令割除那与我们交往而伤害我们的人，那么那些经常出入那些场所、每日吸引甚至素不相识的人、给自己招来无数毁灭机会的人，还能有什么借口呢？\par

因为从此以后，祂不仅禁止我们不洁地观看，而且在指明由此产生的祸害后，祂甚至进一步严格法律，命令割除、分离、远远丢弃。祂——曾说出无数关于爱的话语——这样制定一切，为要使你在两方面都学到祂的眷顾，以及祂如何从各种来源寻求你的益处。\par

4. \Quote{又有话说：\InnerQuote{凡休妻的，当给她休书。}我告诉你们，凡休妻的，除非为淫乱的缘故，就是使她犯奸淫；人若娶被休的妇人，也是犯奸淫。} \BibleRef{玛 6:31}\par

祂没有立即转向面前的话题，直到彻底阐明前面的论点。因为，看，祂又向我们显示了另一种奸淫。这是什么呢？有一条古代法律规定，\BibleRef{申 24:1}凡憎恶自己妻子的，无论出于何种原因，不应禁止他休妻，并另娶一妻。然而法律并非简单地命令他这样做，而是规定给妻子休书，使她不能再回到他那里；这样至少可以保存婚姻的形式。\par

因为如果祂没有命令这样做，而是允许先休妻再娶，然后再把前妻接回，混乱必然很大，人们会不断互相娶妻；此后事情就直接成了奸淫。为此，祂设计了休书作为不小的缓和措施。\par

但这些事是因另一件更严重的恶行而行的；我的意思是，如果祂命令必须将即使被憎恨的妻子留在家里，那丈夫由于憎恨会杀死她。因为犹太民族就是如此。他们连孩子都不饶恕，杀死先知，\Quote{倾流人血如水}，更不会对妇女手下留情。因此祂允许较小的恶，以消除更大的恶。因为这不是一条首要的法律，请听祂说：\Quote{梅瑟因你们心硬，才写了这条诫命}，免得你们在家里杀死她们，而宁愿休弃她们。然而，既然祂已经除去了所有愤怒，不仅禁止杀人，甚至禁止愤怒的情绪，祂便轻易地引入了这条法律。为此，祂不断提到先前的言语，以表明祂的话与它们并不矛盾，而是和谐一致：祂是在加强，而非推翻它们；成全，而非废除它们。\par

请注意祂到处向男人说话。例如，祂说：\Quote{凡休妻的，就是使她犯奸淫；凡娶被休的妇人的，就是犯奸淫。}这就是说，前者即使不另娶妻，仅此一举便已使自己有罪，因为他使前妻成了淫妇；后者又因娶了别人的妻子而成了奸夫。请不要对我说：\InnerQuote{别人已经休了她。}不，因为被休后，她仍然是休她者的妻子。然而，为了不让妻子因将一切归咎于休她者而更加任性，祂也关闭了后来接纳她者的门，因为祂说：\Quote{凡娶被休的妇人的，就是犯奸淫；}这样使妻子即使不情愿也保持贞洁，完全堵住了她接近一切人的路，不让她给嫉妒提供借口。因为那得知自己要么必须保留最初配给她的丈夫，要么被赶出家门后绝无其他避难所的女人——她即使不情愿，也被迫与配偶好好相处。\par

如果祂没有对女人谈论这些事，请不要惊讶；因为女人是软弱的受造物。为此，祂放过她，而在对男人的威胁中完全纠正了她的懈怠。正如一个人有一个放荡的孩子，他离开孩子，却斥责那些使孩子变成这样的人，禁止他们与孩子交往或接近。如果这令人痛苦，请你们回想祂先前说的话，祂是在什么条件下祝福听众的；你就会看到这是完全可能且容易的。因为那温柔、和平、贫穷、慈悲的人，怎么会休妻呢？那习惯于使别人和好的人，怎能与自己妻子不睦呢？\par

非但如此，祂还以另一种方式减轻了这项规定：因为祂甚至为男人留下一种休妻的例外，当祂说：\Quote{除非为淫乱的缘故}；否则事情又会回到同样的结局。因为如果祂命令将即使犯了许多淫乱的妻子留在家里，那就会再次导致奸淫。\par

你看见这些话如何与先前所说的相合吗？因为凡不以不洁的目光注视其他妇女的，就不会行淫；不行淫的，就不会给丈夫休妻的借口。\par

因此，你看，此后祂毫无保留地强调这一点，并以这种恐惧作为壁垒，提醒丈夫若休妻，他自身便要为她的奸淫负责，这是极大的危险。这样，免得你听到\Quote{剜出眼睛}时，以为这是针对妻子说的，祂及时添加了这项纠正，只允许休妻一种情况，绝无其他。\par

5. \Quote{你们又听见有对古人说：不可背誓，所起的誓，总要向主谨守。只是我告诉你们，什么誓都不可起。} \BibleRef{玛 5:33}\par

为什么祂没有立即论到偷窃，而是越过那条诫命，先论妄证呢？因为偷窃的人有时也会起誓；但一个既不起誓也不说谎言的人，更不会选择偷窃。所以借此祂也推翻了另一罪行：因为谎言源自偷窃。\par

但\Quote{要向主谨守你的誓言}是什么意思？就是\Quote{你起誓要真实。}\Quote{只是我告诉你们，什么誓都不可起。}\par

接着，为了引导他们远离指着天主起誓，祂说：\Quote{不可指着天起誓，因为天是天主的宝座；不可指着地起誓，因为地是祂的脚凳；也不可指着耶路撒冷起誓，因为它是大君王的京城。}——这些话仍是出自先知著作，表明祂自己并不与古人对立。这是因为他们有指着这些对象起誓的习俗，祂在福音的末尾暗示了这个习俗。\par

但请注意，我祈求你，祂是以什么理由来尊崇这些元素；不是从它们自身的本性，而是从天主与它们的关系，这关系是出于迁就而启示的。因为偶像崇拜的暴政甚大，免得元素因自身而被认为值得尊敬，祂便提出了我们提到的这个原因，这原因又归向天主的荣耀。因为祂既没有说\Quote{因为天是美丽而伟大的}，也没有说\Quote{因为地是有益的}；而是说\Quote{因为一是天主的宝座，另一是祂的脚凳}；从各方面催促他们归向他们的主。\par

\Quote{也不可指着你的头起誓，}祂说，\Quote{因为你不能使一根头发变白或变黑。} \BibleRef{玛 5:36}\par

在这里，祂不是因惊叹于人而阻止他指着自己的头起誓（因为那样人自己就会被崇拜），而是将荣耀归于天主，并表明你连自己也不是主人，因此当然也不是指着头所起的誓的主人。因为假如没有人会把自己的孩子交给别人，那么天主更不会把自己的作品交给你。因为虽然那是你的头，但它却属于另一位；你不但不能做它的主人，连最小的事也不能做。因为祂没有说\Quote{你不能使一根头发生长}，而是说\Quote{连改变它的颜色也不能}。\par

\Quote{但是，}也许有人说，\Quote{如果有人要求起誓并施加压力呢？}让对天主的敬畏胜过那压力吧：因为如果你搬出这样的借口，你将一事无成。\par

是的，首先关于你的妻子，你会说：\Quote{如果她好争吵又奢侈呢？}然后关于右眼：\Quote{如果我爱它，并且十分热心呢？}关于不洁的目光：\Quote{那么，如果我无法避免看见呢？}关于对弟兄发怒：\Quote{如果我性情急躁，不能控制舌头呢？}总而言之，你都可以这样践踏祂所有的教训。然而，关于人的法律，你绝不敢在任何情况下使用这种借口，也不敢说：\Quote{那么如果这或那情况发生呢？}而是，无论愿意与否，你都要接受所写的。\par

此外，你根本不会受到任何强迫。因为那先前听从了那些真福，并使自己成为基督所命令那样的人，不会受到任何人的这种强迫，反而被众人尊敬和崇敬。\par

\Quote{你们的话，是就说是，不是就说不是；那多余的，便是出于邪恶。}\par

那么什么是\Quote{超过是}和\Quote{不是}的呢？就是起誓，而不是背誓。因为后者是众所周知的，没有人需要学习它是出于邪恶；它并非多余，而是相反；而多余意味着更多的东西，是额外添加的；起誓就是这类东西。\par

\Quote{那么，}有人说，\Quote{它是出于邪恶的吗？如果它是出于邪恶的，它怎会成为法律呢？}同样，关于妻子你也会说：为何现在被视为奸淫的事，先前却是允许的？\par

对此可以回答什么呢？就是那时颁布的诫命是针对领受法律之人的软弱；因为用祭品的烟云来被崇拜非常不相称于天主，就像咿呀学语不相称于哲学家一样。因此，现在既然德行原理已经进步，那种事就被定为奸淫，起誓被定为出于邪恶。但如果这些事从一开始就是魔鬼的法律，它们就不会达到如此大的美善。是的，因为若不是那些先有先行者，我们现在所拥有的就不会如此容易被接受。所以，不要现在要求它们的卓越，既然它们的使用已经过去；而要那时，当时代召唤它们的时候。或者，如果你愿意，现在也可以：是的，因为现在它们的美德仍然显明；而且正是由于我们指责它们的原因，这一点最为突出。因为它们现在显得如此，正是对它们最大的称赞。因为如果它们没有好好抚养我们，使我们适合接受更大的诫命，它们就不会显得如此。\par

因此，正如乳房，当它完成了自己的一切职责，将婴儿打发到更为成人的饮食之后，便显得无用；而父母先前认为它对于婴儿是必需的，如今却用千百种嘲笑来糟蹋它（许多人甚至不满足于辱骂之词，还涂抹苦药在其上；好使当言语无力消除孩子不合时宜的倾向时，实际的事物能扑灭他们的渴望）：同样，基督也说，他们属于恶者，并不是要指明旧约法律属于魔鬼，而是要以最大的热忱将他们从昔日的贫乏中引领出来。祂对他们说这些话；至于那些麻木不仁、固守旧路的犹太人，祂以被掳的恐怖涂抹绕他们的城，如同某种苦药，使其无法接近。但由于连这也不能约束他们，他们仍渴望再见它，奔向它，犹如婴儿奔向乳房，祂便将它完全隐藏起来；既拆毁它，又将他们中的大部分从它那里远远带走：正如我们的牲口一样，许多人通过将牛犊隔开，久而久之诱导它们放弃对乳汁的旧习。\par

但是，如果旧约法律属于魔鬼，它就不会引领人们远离偶像崇拜，反而会牵引他们并投进其中；因为这是魔鬼所愿的。但现在我们看到旧约法律产生了相反的效果。事实上，起誓这件事本身，古时正是为此而设立的，好使他们不指着偶像起誓。因为祂说：\Quote{你要指着}，祂说，\Quote{真天主起誓}。那么，法律所产生的益处并非微小，而是极大。因为他们得以享用\Quote{固体食物}，乃是其关爱的成果。\par

\Quote{那么，}有人可能会说，\Quote{起誓不是出于恶者吗？}是的，它确实完全出于恶者；也就是说，如今，在如此高的克己准则之后；但那时并非如此。\par

\Quote{但是，}有人会说，\Quote{同一件事怎么会时而成为善的，时而成为不善呢？}不，我说恰恰相反：当万物都大声宣告它们如此时：技艺、地上的果实以及其余一切，它又怎能不成为善的和不善的呢？\par

请看，例如，这首先发生在我们同类身上。因此，在生命的最初阶段，被人抱着是好的，但后来却有害；在我们生命最先的时期，吃在嘴里软化过的食物是好的，但后来却令人厌恶；以乳汁喂养并扑向乳房，起初是有益且健康的，但后来却趋于败坏和损害。你看到同样的行为，因着时间，如何显得好，后来又显得不好吗？是的，当你还是孩童时，穿着童袍是好的；但相反，当你成人后，它就变成丢脸的。你也愿学习相反的情况吗？即对孩童来说，成年人的事物又是如何不适宜呢？给男孩一件成年人的袍子，将会引起大笑；而且危险更大，因为那样走路常常会跌倒。让他处理公务、经商、播种、收割，同样会引起大笑。\par

我为什么提到这些？当杀人，在所有恶行中被承认为恶者的发明，杀人，我说，找到了合适的时机，使行这事的丕乃哈斯获享司祭职的尊荣。\BibleRef{户 25:8} 因为那杀人是出自我刚才提到的那一位的工作，听基督说什么：\Quote{你们要行你们父亲的工作；他从起初就是杀人者。}\BibleRef{若 8:44} 但是丕乃哈斯成了杀人者，\Quote{这算为他}（祂如此说）\Quote{的正义}，而亚巴郎不仅成为杀人者，而且（更糟的是）成为杀子者，却赢得了更多的赞许。伯多禄也行了双重的杀戮，然而他所行的是出于圣神。\EditorialNote{宗徒大事录第五章}\par

因此，让我们不要只审视行为，也要审视时节、原因、心意、人的差异以及一切可能伴随它们的东西，让我们精确地探查这些；因为否则无法达致真理。\par

如果我们想要获得天国，就让我们殷勤地显出比旧诫命更多的东西；因为否则我们无法把握天上的事。因为如果我们仅仅达到古人的同等程度，我们就会站在门槛之外；因为\BibleQuote{除非你们的义德超过经师和法利塞人的义德，你们决不能进入天国。}\BibleRef{玛 5:20}\par

6. 然而，虽然如此重的警告已经记下，有些人却不仅没有超越这义德，反而连它也不及；不仅不避免起誓，甚至发假誓；不仅避免不贞的眼目，甚至堕入邪恶的作为本身。凡禁止的其他一切事，他们都敢于去做，如同失去感觉一般：只等待一件事，就是惩罚的日子，以及他们要为恶行付出最极端代价的时刻。而这只属于那些在邪恶中结束自己生命的人。因为这些人有理由绝望，从今以后除了惩罚别无期待；而那些仍然在此的人，却有力量重新战斗、得胜并轻易地接受冠冕。\par

因此，不要灰心，人啊，也不要放弃你高贵的热情；因为事实上所命令的事并不沉重。请问，避免起誓有什么困难呢？难道它需要花钱吗？需要流汗和辛苦吗？只要愿意，就大功告成了。\par

但如果你向我提出你的习惯；正是为此，我尤其说，你行善是容易的。因为如果你使自己养成另一种习惯，你就做成了一切。\par

请思考，例如，在希腊人中，有许多实例，口齿不清的人通过大量练习完全治好了他们结巴的舌头；而另一些人，习惯用不雅的方式耸肩并不断移动它们，通过在肩上放一把剑，使自己戒除了这个习惯。\par

因为你们不受圣经的劝导，我被迫用外邦人来羞辱你们。天主也曾对犹太人这样做，当祂说：\Quote{你们往\Place{基提海岛}去，到\Place{基达}那里，看有没有外邦人更换他们的神；其实那些并不是神。}祂也常打发我们到畜生那里去，这样说：\Quote{懒惰的人哪，你去察看蚂蚁，效法它的作为；}又\Quote{到蜜蜂那里去。}\par

因此，我现在也对你们说：请思考\Place{希腊}的哲学家们；那么你们就会知道，我们这些不服从天主法律的人，该受何等大的惩罚：因为他们为了在人前显得端庄，已费尽心力，而你们却不肯为天上的事付出同样的努力。\par

如果你们回答说：\Quote{习惯有奇妙的力量，甚至能迷惑那些非常认真的人；}我也承认这一点；然而，我还要加上另一件事：既然它能迷惑人，它也容易改正。因为如果你们在家里为自己设置许多监督者，例如你们的仆人、妻子、朋友，你们就会轻易打破坏习惯，因为受到众人的督促和严密约束。如果你们连续十天这样做，之后就不再需要更多时间，一切都会稳妥，在新的优秀习惯的坚固中扎根。\par

因此，当你们开始改正时，即使你们违反自己的准则一次、两次、三次、甚至二十次，也不要失望，而要重新站起来，恢复同样的努力，你们必定能够成功。\par

因为发假誓确实不是小事。如果发誓是出于恶者，那么假誓会带来多大的惩罚！你们对所讲的话表示赞赏了吗？不，我不要掌声，不要喧闹，不要噪音。我只希望你们安静而明智地聆听，并按所说的去做。这就是我的掌声，我的赞辞。但如果你们赞美我所说的，却不践行你们所称赞的，那么惩罚更大，控告更严厉；而对我则是耻辱和嘲笑。因为这里所呈现的不是戏剧表演；你们也不是坐在这里观看演员，仅仅为了鼓掌。这个地方是灵性的学校。因此，目标只有一个：适当地实行所讲的话，用行动表现我们的服从。因为只有这样，我们才算获得了一切。但实际上，说实话，我们已经相当绝望了。因为我不仅在私下遇见你们时，也在公开与你们全体交谈时，不断地给予这些劝诫。可是我没有看到任何成效，你们仍然停留在初级的粗陋阶段，这足以使老师感到厌倦。\par

例如，请看\Person{保禄}本人，他几乎无法忍受，因为他的弟子们长期停留在初级阶段：\Quote{因为按时间来说，你们本该做导师了，却还需要有人再把天主圣言的基础教导你们。}\par

因此，我们也在哀伤悲痛。如果我看到你们执迷不悟，我将禁止你们今后踏进这神圣的门槛，并分享不朽的奥迹，如同我们对待淫乱者、奸淫者和杀人犯一样。是的，因为与其聚集一大批违法者和败坏他人的人，不如与两三个遵守天主法律的人一起献上我们惯常的祈祷。\par

不要让我在这里看到任何富人、权贵对我吹胡子瞪眼；所有这一切对我来说都是神话、幻影、梦境。因为现在富裕的人，没有一个会在那里为我辩护，当我被传唤、被控告，说我未能用应有的热忱完全维护天主的法律时。因为正是这一点，毁了那位可敬的老人，\BibleRef{撒上 3:13}虽然他本人在自己的生活中无可指摘；然而，因为他容忍天主的法律被践踏，他和他的儿子们一起受罚，付出了惨重的代价。如果，出于自然如此强大的绝对权威，那个未能对自己的儿子们表现适当坚定的人竟受了如此重的惩罚；那么，我们摆脱了那种权威，却用谄媚毁掉一切，情何以堪？\par

因此，为使你们不毁掉我们和你们自己，我恳求你们，听我劝说：给自己设置许多监督者，他们会督促你们、要求你们交代，从而使自己摆脱发誓的习惯；这样，从那里开始循序渐进，你们就能轻而易举地达到一切美德，并享受将来的福乐；愿天主赐予我们众人获得这一切，因我们的主\Person{耶稣基督}的恩宠和慈爱，愿光荣和权能归于祂，从今世直到永远。阿们。\par

\SourceRecord{Translated by George Prevost and revised by M.B. Riddle.；Nicene and Post-Nicene Fathers, First Series；Vol. 10.；Edited by Philip Schaff.；Buffalo, NY: Christian Literature Publishing Co.,；1888.；<http://www.newadvent.org/fathers/200117.htm>.}

% structure: p0001 source=h1 target=section reason=page-title

\section{\BibleBook{玛窦}{福音}第十八篇讲道}

\BibleRef{玛 5:38-40}\par

\BibleQuote{\Emph{你们听见有话说：以眼还眼，以牙还牙。但我告诉你们：不要抵抗恶人：} \Emph{而且，凡打你右颊的，转给他另一面也；若有人愿在法庭上控告你，拿你的内衣，你连外衣也让他拿去。}}\par

你们看见了吗？祂从前制定律法，要剜出那得罪人的眼睛时，并非指眼睛本身，而是指那借着友谊伤害我们、将我们投进毁灭深渊的人。因为那在这处用了如此强烈言辞的祂，甚至当别人剜你的眼睛时，也不许你剜他的；祂又怎能制定律法要人剜自己的眼睛呢？\par

但若有人指责古代法律，因为它命令这样的报复，在我看来，那人极不擅长一位立法者应有的智慧，也不懂得时机的美德和屈尊俯就的益处。因为如果他考虑这些话语的听众是谁，他们性情如何，以及他们在何时接受了这部法典，他就会完全承认立法者的智慧，并看见是同一位制定了那些法律和这些法律，且按最有益的方式和适当的时机写下每一条。是的，因为若在开始时祂就引入这些崇高而最重要的诫命，人们既不会接受这些，也不会接受那些；但如今，祂在适当的时候分别颁布它们，借着这两者纠正了整个世界。\par

此外，祂命令这些，不是要我们互相剜出眼睛，而是要我们管住自己的手。因为遭受报复的威胁有效地约束了我们行恶的倾向。\par

事实上，祂正悄悄撒下大量自律的种子，至少在于祂命令以同样行为报复。然而，那开始这种侵犯的人确实应受更重的惩罚，这是抽象的正义所要求的。但因祂愿意将仁慈与正义调和，祂定那犯下极大过犯的人受比其应得更轻的惩罚：教导我们即使在受苦时也要显出极大的体谅。\par

因此，祂既提及了古代法律，并完全承认它，便再次表明：行这些事的不是我们的兄弟，而是恶者。为此，祂又附加上：\BibleQuote{但我告诉你们，不要抵抗恶者。} 祂没有说：\Quote{不要抵抗你的兄弟，} 而是说\Quote{恶者，} 指明是受他的推动，人才敢如此行；这样，祂将罪责转移到另一处，从而缓和并暗中消除了我们对侵犯者的大部分愤怒。\par

\Quote{那么怎样？} 有人问，\Quote{我们不该抵抗恶者吗？} 当然应该，但不是这样，而是如祂所命令的，甘愿忍受不义；因为这样你将胜过他。因为火不能被火熄灭，只能被水熄灭。为了向你指明，即使在旧法律下，受苦的人反而得胜，他才是赢得桂冠的那位；考察所发生的事，你就会看见他的优势很大。因为那开始行不义的人，会毁掉双方的眼睛，邻人的和自己的（因此他受到所有人的憎恨，千万控告对准他）；而受伤害的人，即使在同等报复之后，也没有做任何可怕的事。因此，许多人同情他，因为他在报复之后仍清白无辜。虽然灾祸对双方相同，但在天主和人面前的判决却不相同。如此看来，末了灾祸也不相等。\par

然而，在开始祂说：\BibleQuote{凡无缘无故向兄弟发怒的，} 和\BibleQuote{凡称兄弟为傻子的，将受地狱之火的危险，} 这里祂要求更完全的自制，命令那受害的人不仅要安静，还要更热切地转过另一边脸。\par

祂说这话，不只是为这样的耳光立法，更是教导我们在其他一切考验中应如何忍耐。正如祂说：\BibleQuote{凡无缘无故向兄弟发怒的，} 和\BibleQuote{凡称兄弟为傻子的，将受地狱之火的危险，} 这里祂要求更完全的自制，命令那受害的人不仅要安静，还要更热切地转过另一边脸。\par

2. \BibleQuote{若有人愿告你，拿你的外衣，连斗篷也让他拿去。} \BibleRef{玛 5:40}\par

因为祂不仅是在挨打的事上，就是在我们的财物上，也愿意我们表现出这样的容忍。因此祂又用了同样的强烈说法。也就是说，正如在前一事例中祂命令以受苦获胜，同样在此处，祂命令我们让自己被剥夺得比恶人所期望的更多。然而，祂并非仅仅这样说，还加上了增强语气的话：没有说\Quote{把斗篷给那求你的人}，而是\Quote{给那愿告你的人}，也就是说，\Quote{如果他把你拖进法庭，给你添麻烦。}\par

正如祂吩咐不可称弟兄为傻子，不可无故发怒之后，更进一步，命令连右脸也转过来；同样，在这里说了\Quote{与你的对头和息}之后，祂又扩大了诫命。因为现在祂命令我们不仅要给他人所想要的，甚至要表现出更大的慷慨。\par

\Quote{那么，}有人会说，\Quote{难道我要赤身露体行走吗？}如果我们准确遵守这些训言，我们不会赤身露体；相反，我们会比任何人都穿得更充足。因为首先，没有人会攻击具有这种心性的人；其次，即使万一有人如此野蛮无情，竟至于走到那一步，仍会有许多人找到机会给那个如此克己的人穿衣，不仅用衣裳，甚至如果可能，用自己的血肉。\par

再者：即使有人因这种克己而不得不赤身露体，那也不是什么耻辱。因为亚当在乐园里也是\BibleQuote{赤身露体} \BibleRef{创 2:25}，\BibleQuote{并不羞耻}；依撒意亚也曾\BibleQuote{赤身露体，赤脚而行}，比所有犹太人更荣耀；\BibleRef{依 20:2} 而且若瑟 \BibleRef{创 39:12} 在脱去衣服时，那时比任何时候都更显光辉。因为这样赤身露体并非坏事，反而像我们现在这样穿着昂贵的衣服，才是可耻和可笑的。因此，你看，那些人得了天主的称赞，而这些人却受祂借先知和宗徒的责备。\par

因此，我们不要认为祂的诫命是不可能的。不，因为它们不仅有益，而且如果我们清醒，它们是非常容易的；它们的益处如此之大，以致不仅对我们自己，甚至对那些虐待我们的人，都大有帮助。它们的主要优越之处就在于此：它们引导我们忍受不义，同时借此也教导那些行不义的人管束自己。因为恶人自以为夺取别人的东西是了不起的事，但你却向他表明，给你甚至他没有要求的东西也是容易的；当你用慷慨来平衡他的卑鄙，用明智的节制来平衡他的贪婪时，你想他会受到怎样的教训？他不仅靠言语，更靠实际行动，学会鄙视恶行，追求美德。\par

因为天主愿意我们不仅有益于自己，也有益于所有的近人。如果你只给予而不诉讼，你只寻求了自己的利益；但如果你还给他别的东西，你就也使他变得更好，然后打发他离开。盐就是如此，祂愿意他们成为盐；因为盐既保存自己，也保存所有与它接触的其他物体；光也是如此，因为它既为自己也向所有人显示物体。既然祂已把你放在这些东西的等级中，你就同样要帮助那坐在黑暗中的人，教导他以前也不是强取任何东西的；说服他，他没有行过任何强暴。是的，因为如果你表明你是自愿给予而非被抢夺，你自己也会更加受尊敬和敬畏。因此，通过你的节制，使他的罪成为你慷慨的实例。\par

3. 如果你认为这很了不起，那么请等一等，你会清楚地看到，你还没有达到成全。因为那位为忍耐制定律法者，并没有在这里止步，而是更进一步，这样说：\par

\BibleQuote{若有人强迫你走一里路，你就同他走两里。} \BibleRef{玛 5:41}\par

你看到克己的高度了吗？至少在这里，在给了外衣和斗篷之后，即使你的敌人想要利用你的赤身躯体从事艰苦和劳累的工作，甚至那样（祂说），你也不可禁止他。因为祂愿意我们与那些需要的人，也与那些侮辱我们的人，共有我们的一切，包括身体和财物：后者出于刚毅，前者出于怜悯。\par

因此，祂说：\BibleQuote{若有人强迫你走一里路，你就同他走两里：}再次引导你更高，命令你展现同样的好胜心。\par

因为如果祂开头所说的事，远小于这些，却已蒙受如此大的祝福；那么请想想，那些妥善履行这些事的人，等待着他们的是怎样的福分，甚至在得到赏报之前，他们在可朽坏的身体中赢得了完全的自由，不受情欲的困扰。因为当侮辱、打击、财物被抢都不能激怒他们时，当他们对这类事毫不让步，反而大大增加自己的忍耐时，请想想他们的灵魂在接受怎样的训练。\par

因此，正如在挨打的事上、在我们的财物上一样，在这个事例中祂也命令我们行事。\Quote{为什么？}祂说，\Quote{我提到侮辱和财物呢？即使他想要利用你的肢体从事劳苦和疲乏的工作，而且是不义地这样行，你也要再次胜过并超越他不义的愿望。}\par

因为\Quote{强迫}就是指不义地、毫无理由地、出于侮辱地拉人。但尽管如此，你也要预备好站定，以便承受比对方想对你做的更多。\par

\BibleQuote{求你的，就给他；向你借贷的，不可转开。} \BibleRef{玛 5:42}\par

末后者小于前者；但不要惊奇，因为主常这样做，将小的与大的相混。如果这些与那些相比是小的，那么那些夺取他人财物、将自己的家产分给娼妓、既因不义之财又因有害之支出为自己点燃双重火焰的人，应当倾听。\par

但这里主所说的\Quote{借贷}，不是指高利贷的契约，而是单纯的使用。在别处祂甚至加以扩展，说我们应当给那些我们不指望归还的人。\par

4.\BibleQuote{你们听过这话：\InnerQuote{要爱你的近人，恨你的仇人。}我却对你们说：爱你们的仇人，为那迫害你们的人祈祷；祝福那诅咒你们的人，对那恨你们的人行善。这样，你们便可以成为你们在天之父的儿女，因为祂使太阳上升，光照恶人也光照善人，降雨给义人也给不义的人。}\par

请看主如何将我们善行的顶峰设置得如此之高。正因如此，祂教导我们不仅忍受打击，还要转脸由他打；不仅加上外衣，还要陪那强迫你走一里路的人走两里；这样你们便可以轻易领受比这些更大的事。有人也许会说：\Quote{那么，比这些更大的是什么呢？}甚至不把做这些事的人当作仇人；或者更甚于此。因为祂没有说\Quote{不要恨}，而是说\Quote{要爱}；祂没有说\Quote{不要伤害}，而是说\Quote{要行善}。\par

若有人仔细考察，便会看见甚至在这一切之上还添加了更大的事。因为祂不仅命令爱，还命令祈祷。\par

你看见主攀登了多少阶梯，将我们置于德行的顶峰吗？不，请从开始数算。第一步：不先作不义；第二步：在别人开始不义后，以相等的报复自卫；第三步：不对那激怒我们的人作我们所受的，而是安静；第四步：甚至甘愿受冤枉；第五步：放弃比那行不义者所希望的更多；第六步：不恨那做这事的人；第七步：甚至爱他；第八步：也向他行善；第九步：亲自为他恳求天主。你看见何等高超的自制吗？因此，它的赏报也是荣耀的，如同我们所见的。这是因为所命令的事是伟大的，需要一个热切的灵魂和许多热诚，主为此所定的赏报也超过了先前的一切。因为祂在这里没有提及土地，如同对温良的人那样；也没有提及安慰和怜悯，如同对哀恸和怜悯的人那样；也没有提及天国；而是提及比一切都更令人振奋的：我们成为像天主一样，就人所能成为的程度。因为祂说：\Quote{这样，你们便可以成为像你们在天之父一样。}\par

并且请你留意，无论在这里还是在前面的部分，祂都没有称祂为自己的父，而是在那论及誓言的段落中称为\Quote{天主}和\Quote{大君王}，在这里则称为\Quote{他们的父}。祂这样做，是为将关于这些点要说的话保留到适当的时机。\par

5. 然后，为了使相似更加贴切，祂说：\par

\BibleQuote{因为祂使太阳上升，光照恶人也光照善人，降雨给义人也给不义的人。}\BibleRef{玛 5:45}\par

祂说：\Quote{因为祂不但不恨，反而厚待那些侮辱祂的人。}然而，这二者绝不可相提并论，不仅因为祂的恩惠无比崇高，也因祂的尊贵卓越。因为你确实被你的同仆轻视，而祂却被祂的奴隶轻视，这奴隶还从祂领受了数不尽的恩惠；你确实以言语祈祷祝福他，而祂则以行动——极大而奇妙的事——点燃太阳，赐下年雨。\Quote{然而，即便如此，我仍准许你与我同等，就人所能达到的程度。}\par

那么，不要恨那亏负你的人，他正为你赢得如此美善之事，并带你到如此伟大的尊荣。不要诅咒那凌辱你的人；因为你若这样做，你就承受了劳苦，却失去了果实；你承受了损失，却失去了赏报；这是最愚蠢的，既然承受了更重的，却不愿承受较轻的。有人问：\Quote{但这怎么可能发生呢？}你看见天主降生成人，如此降卑，为你受这么多的苦，你还在询问和怀疑，如何能宽恕同仆的伤害吗？难道你没有听见祂在十字架上说：\BibleQuote{父啊，宽恕他们吧，因为他们不知道他们做的是什么？}\BibleRef{路 23:34}？难道你没有听见\Person{保禄}说：\BibleQuote{那升到至高者之上，坐在右边的那位，为我们转求？}\BibleRef{罗 8:34}？你没有看见吗？即使在十字架上和升天以后，祂还派遣宗徒到杀害祂的犹太人那里，给他们带来祂无尽的祝福，尽管他们要在这些人手中受无尽的恐怖。\par

6. 但你受到了极大的冤屈吗？不，你所忍受的，哪能比得上你的主？祂被捆绑、被鞭打、被棍击、被仆人吐唾沫，忍受死亡，而且是万死之中最可耻的死亡，在施予了无数恩惠之后。即使你受了极大的冤屈，正因如此，你更应当向他行善，好使你自己的冠冕更加荣耀，也使你的兄弟从最严重的病弱中解脱。因为医生们也是这样：当他们被疯子踢打、凌辱时，他们反而最怜悯他们，并设法彻底医治他们，知道这侮辱来自他们疾病的极端。我也劝你以同样的心意对待那些图谋伤害你的人，并如此对待那些损害你的人。因为他们才是患病最重的人，他们才经受着一切的暴力。所以，你要将他从这可耻的侮辱中解救出来，允许他放下愤怒，并把他从那可怖的恶魔——忿怒——中释放。是的，因为如果我们看到被魔鬼附身的人，我们为他们哭泣；我们不会寻求自己也附身。\par

现在，让我们也同样对待那些发怒的人；因为的确，愤怒的人就像被附身的人；不，他们比被附身的人更可怜，因为他们清醒地发狂。因此，他们的狂乱更是不可原谅。所以，不要践踏跌倒的人，反而要怜悯他。例如，如果我们看到有人胆汁上涌，眼睛昏花，眩晕，挣扎着要吐出这有害的体液，我们就会伸出手，在挣扎中持续扶持他，即使弄脏了衣服也不在乎，只求一事：如何使他摆脱这痛苦的折磨。那么，让我们也这样对待愤怒的人：当他们呕吐挣扎时，持续扶持他们，不要放手，直到他吐出所有的苦毒。当他平静下来时，他就会对你怀着极大的感激；那时他就会清楚地知道你将他从多大的痛苦中释放了。\par

但我为何说起他的感谢呢？因为天主会立即给你加冕，并用万千荣耀赏报你，因为你使你的兄弟摆脱了严重的疾病；而且那位兄弟也会像尊敬主人一样尊敬你，永远敬畏你的忍耐。\par

你没有看到那些临产的妇人吗？她们如何咬旁边的人，而她们并不感到疼痛？或者更确切地说，她们感到疼痛，但勇敢地忍受，并同情那些因产痛而受苦的人。你也要效法她们，不要比女人更软弱。因为在这些妇人生育之后（因为这些男人比女人更软弱），那时他们才会知道你是一个真正的男子汉。\par

9. 那么，我们这些被命令效法天主，却或许连税吏都不如的人，有什么可配得的呢？因为如果\Quote{爱那些爱我们的人}是税吏、罪人和外邦人的行为，那么当我们连这都没有做到（而当我们嫉妒那受尊荣的弟兄时，我们就确实没有做到）时，我们这些被命令超越经师，却落在外邦人之后的人，将受什么惩罚呢？那么，请告诉我，我们怎能见到天国？我们怎能踏上那神圣的门槛，我们连税吏都未能超越？因为当祂说\Quote{税吏不也是这样做吗？}时，祂暗中暗示了这一点。\par

在祂的教导中，最值得我们钦佩的是这一点：虽然祂在每个例子中都极其详尽地列出了战斗的奖赏，例如\Quote{看见天主}、\Quote{继承天国}、\Quote{成为天主的子女}、\Quote{如同天主}、\Quote{获得怜悯}、\Quote{受安慰}和\Quote{大赏报}；但如果必须提及令人悲伤的事，祂则用压抑的语气说出。例如，首先，\Quote{地狱}这个词在如此多的句子中只出现了一次；在其他一些例子中，祂也克制地纠正听者，仿佛祂的讲论更倾向于羞辱而非威胁；祂说：\Quote{税吏不也是这样做吗？}、\Quote{盐若失了味}、\Quote{他在天国里要称为最小的}。\par

在某些地方，祂把罪本身当作惩罚提出来，让听者自己推断惩罚的严重性：例如祂说：\Quote{他在心里已经与她犯了奸淫}、\Quote{那休妻的，使她犯了奸淫}、\Quote{此外的，都是出于邪恶}。\par

因此，祂在这里也引述外邦人和税吏，以人物的身份来羞辱门徒。保禄也这样做过，他说：\BibleQuote{不要忧伤，像那些没有望德的其他人一样}\BibleRef{得前 4:13}；又说：\BibleQuote{像那些不认识天主的外邦人一样}\BibleRef{得前 4:5}。\par

为表明祂所要求的并不过分，只是比习惯的稍多一点，祂说：\par

\BibleQuote{外邦人不也是这样做吗？}\BibleRef{玛 5:47} 然而祂并未在此止住讲论，而是以奖赏和那些美好的希望作结，说：\par

\Quote{所以你们当是成全的，如同你们的天父。}\par

祂处处丰富地提到诸天的名号，以这地方本身彻底提升他们的心智。因为，我不知为何，那时他们还有些软弱迟钝。\par

10. 那么，我们既记住所说的一切，就要对我们的仇敌显出大爱；并且要抛弃那荒谬的习俗，许多轻率的人常随从它，即等待遇见他们的人先打招呼。他们对于那含有大祝福的事没有热心，却追求那荒谬的事。\par

那么，你现在为什么不肯先对他说话？答复是：\Quote{因为他在等候这个，}不，正因如此，你更应该迎向他，好赢得冠冕。\Quote{不，}他说，\Quote{因为这就是他的目的。}还有什么比这种愚蠢更糟糕呢？就是说，\Quote{因为这就是他的目的——成为我获得赏报的牵线人——我就不愿动手去做他所提示的事。}现在，如果他先对你说话，即使你招呼他，你也不会得到什么。但如果你先迎上前去对他说话，你就从他那骄傲中为自己谋得利益，并在某种程度上从他固执中收获了丰硕的果实。那么，当我们从空话中就能收获如此大的果实时，却放弃赢利，并谴责他，却在同一件事上跌倒，这岂非极端的愚昧？因为你若为此事责怪他，即他先等候别人先对他说话，那么你又为何效法你所指控的同一件事？你曾说那是邪恶的，为何你要模仿它为善？你看见没有，没有什么比一个与恶为伴的人更无知了？因此，我恳求，让我们逃避这种邪恶而可笑的作风。确实，这瘟疫已经摧毁了成千上万的友谊，制造了许多敌意。\par

因此，让我们抢先向对方致意。既然我们受命要忍受打击、被迫旅行、被仇敌剥夺、并忍受这一切，我们若在单纯的问候形式上表现出如此大的好争论，那该当得什么样的宽容呢？\par

11. \Quote{为什么，}有人说，\Quote{我们一放弃这个，就被人轻视和唾弃。}为了避免人轻视你，你就冒犯天主吗？为了避免你那狂躁的同仆轻视你，你就轻视施予你如此大恩惠的主吗？不，如果同你平等的人轻视你是错的，那么你轻视那创造你的天主，更当如何呢？\par

与此一起，也请思考另一点：当他轻视你时，他正在那一刻为你获得更大的赏报。因为你为了天主而忍受它，因为你听从了祂的律法。这等荣誉，有什么不能相比？相当于多少冠冕？但愿我的份份就是为天主而受辱被轻，胜于受所有君王的尊崇；因为没有什么，没有什么能与这光荣相等。\par

那么，我们就当这样追求，如同祂自己所命令的，不理会人的事物，而在一切事上表现出完美的自制，如此引导我们自己的生活。因为这样，就在此刻，从这时候起，我们就要享受天上的福乐和那里的冠冕，在人间像天使般行走，在世上像天使的能力一样活动，并远离一切情欲、一切纷扰。\par

再加上这一切，我们还要领受那不可言喻的福分：愿我们都能达到，借着我们的主耶稣基督的恩宠与对人的慈爱，愿光荣、权能、钦崇归于祂，与无始的圣父及圣而善的圣神，自今至永远，世世无尽。阿们。\par

\SourceRecord{Translated by George Prevost and revised by M.B. Riddle.；Nicene and Post-Nicene Fathers, First Series；Vol. 10.；Edited by Philip Schaff.；Buffalo, NY: Christian Literature Publishing Co.,；1888.；<http://www.newadvent.org/fathers/200118.htm>.}

% structure: p0001 source=h1 target=section reason=page-title

\section{玛窦福音第十九讲}

\Emph{\BibleRef{玛 6:1}}\par

\BibleQuote{\Emph{你们应当小心，不可在他人面前行你们的施舍}，\Emph{为得被人看见。}}\par

他在剩下部分根除了最专横的情欲，即关于虚荣的愤怒和疯狂，这在行义的人中滋生。因为起初他完全没有谈论它；在说服他们做应该做的事之前，教导他们应如何实践和追求这些东西，确实是多余的。\par

但在他引导他们达到自制之后，他进而清除其中潜藏的渣滓。因为这病绝不是偶然发生的；而是当我们恰当地履行了许多诫命之后。\par

所以应当先植入德行，然后除去败坏其果实的激情。\par

且看他以什么开始：以禁食、祈祷和施舍；因为这些善行最易成为虚荣的巢穴。例如，法利塞人因此而骄傲，他说：\BibleQuote{我每周禁食两次，我捐什一之税。}\BibleRef{路 18:12}他在祈祷中也是虚荣的，为了展示。因为当时没有别人在场，他指着自己向税吏说：\BibleQuote{我不像别人，也不像这个税吏。}\BibleRef{路 18:11}\par

且注意基督如何开始，好像谈论某种野兽，难以捕捉，狡猾地欺骗那些不很警醒的人。因此，他说：\Quote{你们小心，关于你们的施舍。}保禄也对斐理伯人说：\Quote{你们要提防狗。}这是有理由的，因为这恶兽秘密地进入我们里面，无声无息地把一切都吹走，悄悄地拿走里面的一切。\par

因为他已经对施舍说了很多，并且引出了天主，\BibleQuote{祂使太阳升起，光照恶人和善人，}\BibleRef{玛 5:45}并从各方面以\Emph{动机}敦促他们这样做，并说服他们因慷慨施舍而欢欣；最后他也清除了所有妨碍这美丽橄榄树的东西。为此他说：\Quote{你们应当小心，不可在他人面前行你们的施舍，}因为先前所说的，是\Quote{天主的}施舍。\par

2. 当他说了\Quote{不可在他人面前行}之后，他又补充了\Quote{为得被人看见}。虽然看起来像是重复，但如果仔细注意，它并非同一件事，而是一件不同于另一件；并且它有极大的保障，以及说不尽的关怀和温柔。因为可能有人在他人面前施舍，却不是为得被人看见；又可能有人不在他人面前施舍，却是为得被人看见。因此，他所惩罚和奖赏的不是单纯的行为，而是意图。若非如此严格，这会使许多人对施舍更加退缩，因为完全秘密地做并不总是可能的。为此，他解除了这种约束，将惩罚和奖赏都不以行为的结果，而以行者的意图来界定。\par

也就是说，免得你说：\Quote{什么？难道若他人看见，我就更糟了？}——他说：\Quote{我所寻求的不是这个，而是你里面的心思和你所作所为的基调。}因为他愿意完全调整我们的灵魂，并医治它的一切疾病。现在如你所见，他禁止人为展示而行，并教导他们由此而来的惩罚，即徒劳无益之后，他再次唤起他们的精神，提醒他们父和天，不仅以损失刺痛他们，而且以给予他们存在的那位的回忆来羞辱他们。\par

他说：\BibleQuote{因为你们在天上的父那里没有赏报了。}\BibleRef{玛 6:1}\par

即使在此他也没有停止，而是进一步用其他动机增加他们的厌恶。因为正如上面他以税吏和外邦人为例，以身份羞辱效仿者，此处他以假善人为例。\par

他说：\BibleQuote{因此，当你施舍时，不可在你前面吹号，像假善人那样。}\BibleRef{玛 6:2}\par

并不是他们真有号，而是他以这种修辞手法显示他们疯狂的程度，借此嘲弄并展示他们。\par

他很好地称他们为\Quote{假善人}，因为面具是怜悯，而精神却是残酷和不人道。因为他们这样做，不是出于怜悯邻人，而是为了自己获得声誉；这出于极端的残酷：当另一个人快要饿死时，却寻求虚荣，而不终止他的痛苦。\par

所以，所要求的不是施舍本身，而是按应该的方式施舍，为这样的目的施舍。\par

在充分嘲笑了那些人，并如此对付他们，使听者甚至应为他们感到羞耻之后，他再次彻底纠正了如此失调的意念；说了我们不应如何行动之后，他另一方面表明我们应如何行动。那么，我们应该如何施舍呢？\par

他说：\BibleQuote{不要叫左手知道右手所作的。}\BibleRef{玛 6:3}\par

这里他的隐晦意思不是关于手，而是他夸张地表达了此事。意思是：\Quote{如果可以，连你自己也不知道，这应该是你的目标；如果可能，甚至对行施舍的手也隐藏。}这不是像一些人说的，要躲避执拗的人，因为他在这里命令要躲避所有人。\par

然后也要思考赏报，它是何其大。因为在他谈论了来自某一方的惩罚之后，他也指出了来自另一方的荣耀；从两方面推动他们，引导他们走向更高的教训。是的，因为他正在说服他们知道：天主无所不在，我们的利益不限于今生，当我们离开这里时，还有一个更可畏的法庭将接我们，以及我们所有行为的账目、荣耀和惩罚；而且没有任何人能在做大事或小事时隐藏起来，尽管他似乎对人隐藏。因为这一切，当他说话时，他隐约地表明了：\par

\BibleQuote{你的圣父在暗中看见，必要公开报答你。}\par

为他设立了一个伟大而庄严的观众集合，并且他所渴望的，正是那件事本身，大大地赐予他。\Quote{因为，}他说，\Quote{你想要什么？岂不是要有一些观看所发生之事的人吗？看，那么你有一些：不是天使，不是总领天使，而是万有的天主。}如果你也希望有人作为观众，他也不会在适当的时刻剥夺你这个愿望，反而更丰富地赐予你。因为，如果你现在炫耀，只能向十人、二十人或（我们暂且说）一百人炫耀；但如果你现在努力隐藏，天主自己将在整个宇宙面前公开宣布你。因此，最重要的是，如果你愿意人看见你的善行，现在就把它们隐藏起来，以便那时所有人都能以更高的荣耀观看它们，天主使它们显明、赞扬它们，并在众人面前宣告它们。再者，现在观看的人反而会谴责你为虚荣；当他们看见你被加冕时，他们非但不谴责，反而都会钦佩你。因此，当你稍等片刻，既能获得赏报，又能收获更大的钦佩时，想想将自己逐出这两者是多么愚蠢；并且当你在寻求天主的赏报、当天主正在观看时，却召唤人来观看所发生之事。是的，如果必须展示我们的爱，我们首先应向我们的圣父展示；尤其当圣父有能力加冕和惩罚时。\par

让我补充，即使没有惩罚，那渴望荣耀的人也不应放弃我们这个剧场，而用人的剧场来交换。因为谁如此可怜，以至于当君王急忙前来观看他的成就时，他会让君王离开，而集合穷人乞丐作为观众？为此缘故，他不仅命令不要炫耀，甚至要努力隐藏：不追求公开与追求隐藏根本不是一回事。\par

3. \BibleQuote{当你们祈祷时，}他说，\BibleQuote{不要像假善人一样，因为他们喜爱站在会堂里和十字路口祈祷。我实在告诉你们，他们已获得了他们的赏报。}\par

\BibleQuote{至于你，当你祈祷时，要进入你的内室，关上门，向你在暗中的圣父祈祷。}\par

这些他又称为\BibleQuote{假善人}，非常恰当；因为他们一边假装向天主祈祷，一边环顾四周看人；穿着不是恳求者的外衣，而是可笑之人的外衣。因为那要行恳求之职的人，应放下一切其他，只仰望那能赐予他请求的一位。但如果你离开这一位，到处游荡、四处张望，你将空手而归。因为这是你自己的意愿。因此他没有说\Quote{这样的人不会得到赏报}，而是\BibleQuote{他们已得到了赏报}：也就是说，他们确实会得到，但来自他们自己渴望的那些人。因为天主不愿意这样；他反而愿意将出于他自己的赏报赐予人；但他们寻求来自人的赏报，便不能再正当地从他那里领受，因为他们没有为他做什么。\par

但请注意，我请求你，天主的慈爱，因为他应许甚至为我们向他祈求的善事赐予赏报。\par

在贬低了那些不按应当的方式安排这职责的人——既从地方也从他们的心态——并表明他们非常可笑之后，他介绍了祈祷的最佳方式，并再次赐予赏报，说：\BibleQuote{进入你的内室。}\par

\Quote{那么，}也许有人说，\Quote{我们不应该在圣堂里祈祷吗？}当然我们应该，但要有这样的精神。因为天主无处不在寻求一切行为的心意。即使你进入内室，关上门，却为炫耀而做，门对你毫无益处。\par

在此也值得注意，他的定义多么精确，当他说：\BibleQuote{为叫人看见。}所以即使你关上门，他更希望你好好履行这一点，而非仅仅关门，即关上心灵的门。因为在一切事上，摆脱虚荣都是好的，尤其在祈祷中。因为即使没有这虚荣，我们也会游荡、分心，如果我们又带着这种疾病进入，何时才能留意我们所说的话呢？如果我们这些祈祷恳求的人都不留意，我们怎能期望天主留意呢？\par

4. 然而，有些人，在如此认真严厉的告诫之后，在祈祷中表现得如此不当，以至于即使他们的身子被隐藏，他们也通过声音向所有人显明自己，杂乱地喊叫，既凭姿势又凭声音使自己成为笑柄。你难道没有看见，即使在市场上，如果有人像这样上前，喧闹地乞求，他会赶走他所求的人；但如果安静地、以适当的姿势，他反而能赢得那能施恩的人吗？\par

那么，让我们不要以身体的姿势、也不以声音的响亮来祈祷，而是以心灵的恳切；不要以喧闹、喊叫和炫耀，以致打扰我们周围的人，而是要以完全的谦逊、心灵的痛悔和内心的眼泪。\par

但你的心痛苦，无法不呼喊？然而，那极其痛苦的人，正应以我所说的方式祈祷和恳求。因为梅瑟也痛苦，并这样祈祷，且蒙垂听；为此天主也对他说：\BibleQuote{你为什么向我呼喊？}\BibleRef{出 14:15}。亚纳也是如此，她的声音虽未被听见，却成就了她所愿的一切，因为她的心在呼喊。\BibleRef{撒上 1:13}但亚伯尔不仅沉默时祈祷，甚至在临终时也祈祷，他的血发出比号角更清晰的呼声。\BibleRef{创 4:10}\par

因此，你也应当像那位圣者一样呻吟，我并不禁止。\BibleQuote{你们要撕裂你们的心，而不是你们的衣服，}如同先知所命，\BibleRef{岳 2:13}从深渊中呼求天主，因为经上说：\BibleQuote{上主，我从深渊向你呼求。}从深处，从心中，发出声音，使你的祈祷成为奥秘。你是否看到，即使在君王的宫殿里，一切喧哗都被除去，四周一片寂静？因此，你也当进入宫殿——不是地上的，而是远为可畏的，即天上的宫殿——显出极大的端庄。是的，因为你与天使的唱诗班联合，与总领天使共融，与色辣芬一同歌颂。所有这些群体都显示出极好的秩序，以极大的敬畏歌唱那奥妙的曲调，以及他们献给万王之王天主的圣歌。因此，当你祈祷时，要与他们联合，效法他们的奥秘秩序。\par

因为你不是向人祈祷，而是向无处不在的天主，祂在声音发出之前已垂听，祂知道内心的秘密。你若这样祈祷，你所获得的赏报是大的。\par

\BibleQuote{因为你们的父，}祂说，\BibleQuote{在暗中看见，必要公开报答你们。}\par

祂没有说\Quote{要白白赐给你们，}而是说\Quote{要报答你们；}是的，因为祂使自己成为了你们的债务人，并由此赐予了你们极大的尊荣。因为祂本身是不可见的，祂也愿你们的祈祷是如此。\par

5. 然后，祂说出了祈祷的言辞本身。\par

\Quote{你们祈祷时，}祂说，\BibleQuote{不要重复无意义的言语，如同外邦人那样。}\par

你们看到，当祂谈论施舍时，祂只消除了那来自虚荣的祸害，并未增添别的；祂既未说应从何处施舍——如由诚实的劳动，而非由抢夺或贪财——因为这在众人中已是充分承认的。而且在此以前，祂已彻底阐明了这一点，当祂祝福\BibleQuote{那些饥渴慕义的人}。\newline{}\par

但关于祈祷，祂还补充了一些：\Quote{不要重复无意义的言语。}如同在那里祂嘲笑伪君子，同样在这里祂嘲笑外邦人；祂处处以人物的卑贱来使听者羞愧。因为最令人刺痛和刺伤的，莫过于我们被比作卑贱的人；藉着这个主题，祂劝阻他们；祂在这里将轻浮称为\Quote{无意义的重复}：即当我们向天主祈求不相宜的事物，如王国、荣耀、战胜敌人、财富丰裕，以及一般与我们无关的事。\par

\BibleQuote{因为祂知道，}祂说，\BibleQuote{你们需要什么。}\BibleRef{玛 6:8}\par

在此，祂似乎要我命令，我们不应使祈祷过长；我说长，不是指时间，而是指所提及之事的数量和长度。因为坚持不懈地重复同样的请求是我们的责任：祂的话是\BibleQuote{恒心祈祷。}\BibleRef{罗 12:12}\par

祂自己，也以那寡妇为例，她因不断恳求，说服了那无情残忍的官长；\BibleRef{路 18:1}又以那朋友为例，他在深夜前来，从床上唤醒那睡觉的人，\BibleRef{路 11:5}并非因友谊，而是因他的固执；祂除了订立一条法律，要大家不断向祂恳求，还做了什么？祂并没有命我们编造一个有万条条款的祈祷，然后来到祂面前，仅仅重复它。祂隐约地指出了这一点，当祂说：\BibleQuote{他们以为因多言必蒙垂听。}\par

\BibleQuote{因为祂知道，}祂说，\BibleQuote{你们需要什么。}若祂知道，有人会说，我们还需要什么祈祷？不是为了教导祂，而是为了赢得祂；为了因不断恳求而变得与祂亲密；为了变得谦卑；为了被提醒你的罪过。\par

6. \BibleQuote{所以，你们要这样祈祷，}祂说：\BibleQuote{我们的天父，你在天上。}\par

请看祂如何立刻振奋听者，并提醒他天主在起初的一切恩惠。因为称天主为父的人，在这一称号中，就承认了罪过的赦免、刑罚的除去、正义、圣化、救赎、义子名分、继承、与独生子的兄弟情谊，以及圣神的赐予。因为一个人若不获得所有这些恩惠，就不能称天主为父。因此，祂双重地激发了他们的精神：既因所呼求者的尊贵，也因他们所享受的恩惠的伟大。但当祂说\Quote{在天上}时，祂并非说天主被局限在那里，而是要使祈祷者脱离尘世，定于高处和天上的居所。\par

祂还教导我们要为我们的弟兄共同祈祷，因为祂并没有说：\Quote{我的圣父，你在天上}，而是说：\Quote{我们的圣父}，为整个身体共同献上祈求，丝毫不顾自己的事，而处处顾及近人的益处。借此，祂一下子消除了仇恨，压服了骄傲，驱除了嫉妒，带来了诸善之母——爱德，消除了人间的不平等，并显示出国王与穷人之间的平等能延伸到多远，至少在那些最大、最不可或缺的事上，我们都是同袍。因为我们在下等亲属关系上有何损害呢？在上等关系中，我们都是结合在一起的，没有人比另一个人拥有更多：富人比穷人，主人比仆人，统治者比臣民，国王比普通士兵，哲学家比野蛮人，有学问者比无学识者，有何更多呢？因为祂赐予所有人同一份高贵，屈尊被称为众人相同的圣父。\par

7. 因此，当祂提醒了我们这份高贵、这来自高天的恩赐、我们与弟兄的平等以及爱德；并且当祂将我们移离了大地，安置在天上之后，让我们看看祂命令我们之后祈求什么。不仅如此，首先，仅仅那句话本身就足以在所有美德中植入教导。因为那称天主为圣父、且为共同的圣父的人，理当显露出相称的行为，以免显得不配这份高贵，并表现出与恩赐相称的勤奋。然而祂并不满足于此，还加上了另一句，这样说：\par

\Quote{愿你的名被尊为圣。}\par

那称天主为圣父的人，其祈祷理当不求任何先于祂的圣父的荣耀的事，而视一切事物次于赞美祂的工作。因为\Quote{被尊为圣}就是\Emph{受光荣}。因为祂自己的荣耀是圆满的，永远如一，但祂命令祈祷者寻求祂也藉我们的生命受光荣。这恰恰是祂之前也说过的话：\Quote{你们的光也当在人前照耀，好使他们看见你们的善行，光荣你们在天的圣父。}\BibleRef{玛 5:16}是的，连色辣芬也这样颂扬说：\Quote{圣，圣，圣。}所以\Quote{被尊为圣}的意思就是，即\Quote{受光荣}。也就是说，他说：\Quote{求你赐予我们如此纯洁地生活，使众人藉着我们都能光荣你。}这又属于完美的自制，即向所有人呈现无可指责的生命，使每一位观者都能为此向主献上应得的赞美。\par

\Quote{愿你的国来临。}\BibleRef{玛 6:10}\par

这又是一个心地正直的儿女的语言：不执着于看得见的事物，也不把眼前的事看为重大，而是奔向我们的圣父，渴慕将来的事。这出自良好的良心和摆脱尘世事物的灵魂。例如，保禄本人每天都渴慕这事：因此他也说，\Quote{连我们这些拥有圣神初果的，也在叹息，等待义子的身份，就是我们身体的救赎。}因为凡有这种爱慕的人，既不会被今生的好事得意忘形，也不会被它的悲伤所羞辱，而如同住在天上一样，摆脱了各种纷乱。\par

\Quote{愿你的旨意奉行在人间，如同在天上。}\par

请看何等卓越的思路！因为祂命令我们确实渴慕将来的事，奔向那居所；并且直到那时之前，即使我们尚在此地，也要努力表现出与天上者同样的行为。因为祂说：你们必须渴慕天堂和天上的事；然而，甚至在到达天堂之前，祂已命令我们使大地成为天堂，并且即使我们仍在其中，也要如同在那里有我们生活一般行事说话；因此这些也应成为我们向主的祈求对象。因为没有什么能阻止我们达到天上权势的完美，因为我们居住在地上；但即使在此地，也可能做到一切，就像已经身居高处一样。因此祂所说的是：\Quote{如同那里万事毫无阻碍地成就，天使不是部分服从部分不服从，而是事事顺从和听命（因为经上说：\InnerQuote{大能者，执行祂的话语}），同样求祢赐予我们人类不半心半意地行祢的旨意，而是按祢的意愿成就一切事。}\par

你看到祂如何教导我们谦虚，因为祂清楚说明德行不仅来自我们的努力，也来自上天的恩宠？并且，祂又命令我们每一个祈祷的人要承担起对全世界的关怀。因为祂根本没有说：\Quote{愿你的旨意奉行} \Emph{在我内}，或\Emph{在我们内}，而是说\Quote{在人间各处}；好使错误被消灭，真理被植入，一切邪恶被驱逐，德行回归，从此天上人间不再有区别。\Quote{因为如果这事成就，}祂说，\Quote{那么下界和上界之间，尽管本性分离，也不再有区别；大地将向我们展示另一批天使。}\par

8. \Quote{求你今天赐给我们日用的食粮。}\BibleRef{玛 6:11}\par

什么是\Quote{\Emph{日用的食粮}}？就是一天的食粮。\par

因为祂既然这样说了：\Quote{愿祢的旨意奉行在人间，如同在天上}，但对那些被肉体包围、受本性需要束缚、不能像天使那样不动情的人讲话——祂既命令我们也当实行这些诫命，如同他们实行一样；祂在接下来的话中也俯就我们本性的软弱。因此，祂说：\Quote{行为成全，我要求同样大，却不要你们毫无情感；不，因为本性的暴政不允许这一点：它需要必要的食物。}但请注意，我恳求你，即使在属身体的事上，属灵的事也何等丰富。因为祂命令我们祈求的，既不是为财富，也不是为美食，也不是为昂贵衣服，也不是为任何其他这类事物，而仅仅是为饼。并且是\Quote{每日的饼}，为了不为明天忧虑。因这一点，祂又加上：\Quote{\Emph{每日的}饼}，即一天的饼。\BibleRef{玛 6:34}\par

祂甚至不满意这说法，后来又加上另一句，说：\Quote{求祢\Emph{今天}赐给我们}；这样我们就不至于为此而耗尽心力为次日操心。因为那日子，你不知是否能见到它的到来，为何要为它担忧呢？\par

祂接下来更进一步吩咐说：\Quote{不要为明天忧虑。}祂愿意我们在各方面毫无牵挂，如鸟展翅飞翔，只给本性以必要所要求的那么多。\par

9. 既然我们在重生之洗后仍会犯罪，祂在此也显示祂对人的大爱，命令我们为得罪之赦而来到爱人者天主面前，并这样说：\par

\Quote{宽恕我们的罪债，如同我们也宽恕我们的债户。}\par

你见到超越的仁慈了吗？在消除了如此大的邪恶、并在祂恩赐的不可言喻的伟大之后，如果人再犯罪，祂仍视他们为可以宽恕的。因为这篇祈祷属于信徒，这由教会的法律和祈祷的开端教导我们。因为未入门者不能称天主为父。那么，如果祈祷属于信徒，而他们祈求罪过得赦免，那么显然即使在洗濯之后，悔改的益处也没有被除去。因为，如果祂不愿表明这一点，祂就不会立法要我们如此祈祷。如今，祂既使人记起罪恶，又命令我们求赦免，并教导我们如何获得宽恕，从而使道路容易；显然，祂引入这祈求的规则，是因为知道并表明，即使在洗礼之泉后，我们仍能从过犯中洗净自己；祂提醒我们的罪，劝我们谦卑；祂命令我们宽恕他人，使我们脱离一切报复之情；同时，祂应许以此也赦免我们，给予我们美好的希望，并教导我们对天主向人不可言喻的仁慈怀有崇高的见解。\par

但我们最应注意的是：尽管在每一条中，祂都提到了整个德行，并以此也包括了忘怀伤害（因为这样，\Quote{愿祢的圣名被尊为圣}是完美言行的精确要求；而\Quote{愿祢的旨意奉行}再次宣告同一件事；能称天主为\Quote{父}是无可指责生活的宣认；所有这些也包含了平息我们对冒犯者的愤怒的义务），但祂仍不满足于此，而为了表明祂在这件事上多么恳切，祂特别将它写下来，并在祈祷之后，不再提及别的诫命，只说了这话：\par

\Quote{因为如果你们宽恕别人的过犯，你们的天父也必宽恕你们。}\BibleRef{玛 6:14}\par

所以开端在我们，我们自己掌握着将要临到我们的审判。因为为使任何一个人，即使愚蠢的人，在被带到审判时，不能有任何或大或小的抱怨；祂使判决取决于你，你这要交账的人；祂说：\Quote{你怎样审判自己，我也怎样审判你。}如果你宽恕你的同仆，你也会从我获得同样的恩惠；虽然两者并不相等。因为你是在需要中宽恕，而天主一无所缺；你宽恕你的同奴，天主宽恕祂的奴仆；你负有无数罪责，天主却无罪。但即使这样，祂仍显示祂对人的慈爱。\par

因为即使没有这些，祂本也可以赦免你所有的过犯；但祂愿意你由此也获得益处；从各方面赐给你无数温和与爱人的机会，驱逐你内里的兽性，熄灭愤怒，用一切方式将你与你的肢体紧密连结。\par

你能说什么呢？说你无辜遭受了邻人的某种伤害？（因为只有这些是过犯，若行为是公义的，就不是过犯。）但你也正接近为这类事获得宽恕，甚至是为更大的事。甚至在宽恕之前，你已经接受了不小的恩赐，即被教导要有一颗人的心灵，并被训练成全然温和。同时，将来也有一个大赏赐为你存留，即不必为你的任何过犯交账。\par

那么，当我们领受了这特权却又背叛自己的救恩时，我们该受怎样的惩罚？如果我们连那些取决于我们的事上都不留情，又怎能要求在其他事上被垂听呢？\par

10. \Quote{不要让我们陷入诱惑，但救我们脱离那邪恶者：因为国度、权能和光荣，都是祢的，直到永远。阿们。}\par

这里祂清楚教导我们自己的卑贱，并抑制我们的骄傲，训诲我们祈求避免一切冲突，而不是贸然冲入。因为如此，我们的胜利将更光荣，魔鬼的败亡更可耻。我的意思是：当我们被拖出时，必须勇敢站立；但当我们未被召时，应当安静，等待冲突的时刻；使我们能显出既无虚荣，又有高尚精神。\par

祂在此称魔鬼为 \Quote{\OriginalTerm{恶者}{the wicked one}}，命令我们向他发动一场无休止的战争，并暗示他并非本性如此。因为邪恶并非出于本性，而是出于我们自己的选择。他之所以被特别如此称呼，是因为他邪恶至极，并且他在未受我们任何伤害的情况下向我们发动无可和解的战争。因此，祂并未说\Quote{救我们脱离恶者们}，而是说\Quote{救我们脱离恶者}；教导我们绝不要因邻人可能加给我们的任何伤害而对他们心怀怨怼，而是要将敌意从他们身上转移到他身上，因为他是一切伤害的根源。\par

祂既这样藉着提醒我们敌人而使我们如同在战斗前焦虑，又斩除了我们一切的懈怠之后，祂再次激励并振奋我们的精神，让我们记起我们在其麾下受训的君王，并表明祂比一切更有能力。祂说：\Quote{因为王国、权能和荣耀都是祢的。}\par

难道这不是说，如果王国属于祂，我们就无需惧怕任何人，因为没有人能与祂抗衡并分割帝国？因当祂说\Quote{王国是属于祢的}，祂甚至将向我们作战的那一位也置于我们面前，表明他已受制伏，尽管他看似反抗，天主暂时容许如此。事实上，他也在天主的仆人之列，虽然属于堕落、有罪的那一类；若未先从上天获得许可，他不敢攻击任何一位同仆。而且我说\Quote{他的同仆}何指？他甚至不敢向猪群施暴，直到祂亲自允许他（\BibleRef{路 8:32}）；也不敢攻击羊群或牛群，直到他从上天获得许可（\BibleRef{约 1:12}）。\par

祂说：\Quote{还有权能。}因此，无论你的软弱多么多样，你都有理由充满信心，因为有这样一位统治你，祂能完全成就一切，并且轻而易举，甚至藉着你。\par

\Quote{荣耀，直到永远。阿们。}因此，祂不仅救你脱离临近的危险，还能使你变得光荣显赫。因为祂的权能伟大，祂的荣耀也不可言喻，它们都是无限的，没有穷尽。你看见祂如何用各种方式膏立了祂的斗士，并使他充满信心吗？\par

11. 然后，正如我先前所说，祂要表明祂最厌恶、最憎恨的是记仇，而最悦纳的是与那恶习对立的德行；因此在祈祷之后，祂又再次提醒我们同一件好事；既以设定的惩罚，又以指定的奖赏，敦促听者遵守这道命令。\par

祂说：\Quote{因为如果你们宽恕别人，你们的天父也会宽恕你们。但如果你们不宽恕，祂也不会宽恕你们。}\par

为这缘故，祂又提及了天和他们的父；同样用这一话题来使听者羞愧：身为这样一位天父的儿子，竟会变成一头野兽；被召上天，却怀着一副属世的凡俗心思。因为你看，我们不仅应当因恩宠而成为祂的儿女，也当因我们的行为。没有什么比愿意宽恕恶人和作恶者更使我们像天主；正如祂先前教导的，当祂说到\Quote{祂使太阳照耀恶人和善人}（\BibleRef{玛 5:45}）。\par

为此原因，祂又在每一条句中命令我们作共同的祈祷，说：\Quote{我们的天父}，\Quote{愿祢的旨意行在地上如同行在天上}，\Quote{求祢赐给\Emph{我们}日用的食粮，并宽免\Emph{我们}的罪债}，\Quote{不要让\Emph{我们}陷入诱惑}，\Quote{救\Emph{我们}脱离凶恶}；处处命令我们用这个复数词，使我们连一丝一毫对邻人的愤怒都不保留。\par

那么，做了这一切之后，他们非但不自己宽恕，反而甚至求天主向他们的敌人复仇，简直公然违反这条律法，他们要受多大的惩罚呢？而天主正在做并安排一切来阻止我们彼此不和。因为爱是万善之根，祂从各处移除一切破坏爱的事物，将我们聚拢并相互结合。因为没有，绝对没有一个人，无论是父亲、母亲、朋友或任何其他人，像创造我们的天主那样爱我们。这一点，祂每日的恩赐和祂的诫命都最清楚地显明了。但如果你对我说到苦难、悲伤和生活中的恶事；想一想你每日在多少事上得罪祂，你就不会再惊讶，即使更多这样的恶事临到你；但如果你享有任何好事，那时你就会惊讶和诧异。然而事实上，我们关注临到我们的灾祸，却不想我们每日得罪祂的过犯；因此我们困惑。因为如果我们严格地计算仅一天的罪过，那么我们就该知道我们必须承受何等的恶。\par

至于我们每人所犯的其他恶行暂且不提，只谈今天所犯的；虽然我当然不知道我们每人可能犯了什么罪，但我们恶行如此之多，即使知道一切的人也无法仅从这些中挑选出来。例如，我们中有谁没有在祈祷中疏忽？谁没有傲慢或虚荣？谁没有说过弟兄的坏话，谁没有接纳过邪恶的欲望，谁没有用不贞洁的眼睛看过，谁没有怀着敌意回忆事情，甚至使自己的心膨胀？\par

如果我们身在教会，短短时间内就犯下如此严重的恶行，那么当我们离开这里时又会怎样呢？如果港湾中波浪已如此汹涌，当我们驶入罪恶的河道——我指的是广场、公共事务以及家中的挂虑——我们还能认得自己吗？\par

然而，从如此众多且重大的罪恶中，天主赐给我们一条短而容易的解脱之道，且毫无劳苦。因为原谅那得罪我们的人有何劳苦呢？不，若不原谅，保持敌意，才是劳苦。就如摆脱愤怒，既在我们心中产生极大的舒畅，对于愿意的人也十分容易。因为无需渡海，无需长途跋涉，无需翻山越岭，无需花费金钱，无需折磨自己的身体；只需愿意，一切罪恶便被清除。\par

但是，如果你非但不自己原谅他，反而向天主求情反对他，那么你还有什么得救的希望呢？正当你本该平息天主之怒时，你却惹怒祂；你披着哀求的外衣，却发出野兽的嚎叫，并朝自己射出恶者的毒箭吗？因此，保禄在提到祈祷时，最强调的莫过于遵守这条诫命；因为他说：\BibleQuote{举起圣洁的手，没有愤怒和疑虑。} 如果在你需要怜悯的时候，你仍不肯放下愤怒，反而极其记恨，尽管你知道这是在将剑刺向自己；那么，你何时才能变得仁慈，吐出这邪恶的毒液呢？\par

但是，如果你还没有充分看到这种蛮横，那么设想这发生在人中间，你就会看出这傲慢的过分。比如：有一个人来到你——一个人——面前寻求怜悯，但在他俯伏在地时，看见一个敌人，便停止哀求，开始殴打他；难道你不会更加生他的气吗？你要想，这同样发生在天主面前。因为同样地，当你向天主祈求时，你在祈求中停了下来，用言语打击你的敌人，并侮辱天主的法律。你召唤那位制定法律消除一切愤怒的天主，去反对那些得罪了你的人，并要求祂做违背自己诫命的事。你在报复方面，不仅自己违背了天主的法律，难道还要恳求祂也这样做吗？什么？祂忘记了自己所命令的吗？什么？说这些话的是人吗？祂是天主，知道一切，并愿意自己的法律被最严格地遵守，祂非但不会做你所要求的事，反而会因你说这些话而厌恶和憎恨你，并要你受最严厉的惩罚。那么，你怎能从祂那里求得祂严厉禁止你去做的事呢？\par

然而，有些人竟到了这样兽性的地步，不仅向天主求情反对敌人，甚至诅咒自己的孩子，并且如果可以的话，还要尝他们的肉；更确切地说，他们甚至正在尝。因为你不要对我说，你没有把牙齿咬进那得罪你的人的身体；因为就你而言，你做了更严重的事：你要求上天的愤怒降到他身上，使他被交付于永罚，并连同他的全家被推翻。\par

哎呀，什么样的咬啮比这更凶猛？什么样的武器比这更恶毒？基督并没有这样教导你；祂并没有命令你以血污口。不，因人的血肉而变得血腥的嘴，也不如这样的舌头令人震惊。\par

那么，你怎能问候你的弟兄呢？你怎能触摸祭献呢？你怎能尝主的血呢？当你心中充满这等毒害时？因为当你说：\Quote{把他撕碎，推翻他的房屋，毁灭一切}，当你对他咒骂万千种死亡时，你与凶手毫无分别，更恰当地说，与吃人的野兽无异。\par

那么，让我们停止这种疾病和疯狂，并要显示出祂所命令的仁慈，对待那些得罪了我们的人：使我们能成为像\BibleQuote{我们在天的父}一样。我们将停止这样做，如果我们回想起自己的罪过；如果我们严格地检查自己在家、在外、在市场和在教会中的所有恶行。\par

12. 因为即使不因为别的，单凭我们在此处的不敬，就值得我们受最严厉的惩罚。因为当先知们吟咏、宗徒们歌唱、天主发言时，我们却在外游荡，招来世俗事务的纷扰。我们甚至没有给天主的法律以像戏院观众给皇帝谕旨那样的静默——他们为谕旨保持安静。因为在那里，当这些谕旨被宣读时，议员、总督、元老院和人民，都直立静听。如果在那种深沉的静默中有人突然跳起来喊叫，他必受最严厉的惩罚，因为对皇帝不敬。但在这里，当来自天上的书信被宣读时，周围却一片混乱。然而，那发送书信的远比我们的君王更伟大，集会也更庄严：因为不仅有人，还有天使在其中；这些书信所报的凯旋，远比地上的更为可畏。因此，不仅人，而且天使和总领天使，以及天上的万军和我们地上所有的人，都被命令要赞美。因为经上说：\BibleQuote{请众赞美主，祂的一切工程。} 是的，因为祂的成就并非小事，而是超越一切言词、思想和人的理解。\par

先知们天天宣报这些事，各自以不同的方式宣扬这光荣的胜利。因为一位说：\BibleQuote{祢上升至高天，掳掠了俘虏，在人间接受了恩赐。}又说：\Quote{主，大能而英勇的。}另一位说：\Quote{祂要分取强者的战利品。}\BibleRef{依 53:12} 祂来正是为此，为\Quote{向俘虏宣告释放，向盲者宣告复明。}\par

祂高声向死亡发出胜利的呼喊，说：\Quote{死亡，你的胜利在哪里？坟墓，你的刺在哪里？}另一位又宣告最深远的和平喜讯，说：\Quote{他们要将刀剑打成犁头，将枪矛打成镰刀。}一位呼召耶路撒冷说：\BibleQuote{熙雍女子，你应尽情欢乐，因为你的君王到你这里来，祂是谦逊的，骑在驴驹上。}\BibleRef{匝 9:9} 另一位也预言祂的第二次来临，这样说：\Quote{你们所寻求的主，必要来到。\BibleRef{拉 3:1} 你们要像挣脱束缚的小牛跳跃。}另一位又对这些事感到惊奇，说：\BibleQuote{这是我们的天主，没有别的可比拟祂。}\BibleRef{巴 3:35}\par

然而，虽然这些和许多更甚的言论都被宣说出来，我们本该战栗，甚至不该自认为仍在地上；但我们却好像在广场之中，喧哗骚动，将整个庄严集会的时间耗费在谈论与我们无关的事上。\par

因此，无论在大小事上，无论在听与行上，无论在外与在家，在教会中，我们都如此疏忽；加上这一切，还祈祷反对我们的仇敌：在如此多的严重罪过上又加上另一个严重的罪，且与它们相等，即这不法的祈祷，我们还能有什么救恩的希望呢？\par

那么，若有什么出乎意料或痛苦的事发生在我们身上，我们后来还能惊奇吗？我们反倒该在无事发生时惊奇。因为前者是出于常理，而后者却超越一切理性和期望。因为那些成为天主的仇敌、激怒祂的人，竟能享受阳光、雨水以及一切，这实在不合情理；这些人比野兽更野蛮，彼此敌对，用咬邻人的血玷污自己的舌头，在参与神宴和领受如此大恩、承受无数诫命之后。\par

因此，思虑这些事，让我们吐出那毒汁，结束我们的仇视，并献上相称于我们身份的祈祷。让我们舍弃魔鬼般的残忍，取用天使般的温和；无论我们受了什么伤害，让我们顾念自己的处境，并反思为这条诫命所预备的赏报，从而平息愤怒，平息怒涛，以便既安然度过今生，又在我们离世后，发现我们的主对待我们，如同我们对待同仆一样。如果这是沉重可畏的事，让我们使它变得轻松可爱；让我们开启向祂坦然无惧的光荣之门；我们未能通过远离罪恶成就的事，就通过善待那些得罪我们的人来完成（这的确不是困难或沉重的）；让我们通过善待仇敌，为自己预先积蓄许多仁慈。\par

因为这样，今生所有人都会爱我们，而天主更超越一切地恩待我们、为我们加冕，并使我们堪得将来的一切美善；愿我们都能藉着我们的主耶稣基督的恩宠和慈爱，达到这一切，愿光荣和权能归于祂，直到永远。阿们。\par

\SourceRecord{Translated by George Prevost and revised by M.B. Riddle.；Nicene and Post-Nicene Fathers, First Series；Vol. 10.；Edited by Philip Schaff.；Buffalo, NY: Christian Literature Publishing Co.,；1888.；<http://www.newadvent.org/fathers/200119.htm>.}

% structure: p0001 source=h1 target=section reason=page-title

\section{《玛窦福音》第二十篇讲道}

\BibleRef{玛 6:16}\par

\BibleQuote{你们禁食的时候，不要像假善人那样面带愁容；因为他们把脸弄得难看，叫人看出他们禁食。}\par

在这里理应大声叹息，痛哭流涕：因为我们不仅效仿假善人，甚至已经超过了他们。我知道，是的，我知道有许多人，不仅禁食并炫耀，甚至忽略禁食，却戴着禁食者的面具，用比他们的罪更糟糕的借口来掩饰自己。\par

他们说：\Quote{我这样做，是为了不得罪许多人。} 你说什么？有天主的命令规定这些事，你还谈论得罪人？你以为遵守它就得罪人，违反它反而使众人免于得罪？还有什么比这更愚蠢的呢？\par

你岂不要停止变得比假善人更坏，使你的伪善加倍吗？当你思量这邪恶的极大泛滥时，岂不因眼前这句话的力量而羞愧？祂没有只说\Quote{他们演戏}，更愿意深深触动他们，便说：\BibleQuote{因为他们把脸弄得难看}，即他们毁损、糟蹋自己的面容。\par

如果因虚荣而使面色苍白是毁容，那么对于以涂脂抹粉毁坏面容、导致不贞节青年堕落的妇女，我们该说什么呢？因为那些男人只伤害自己，而这些妇女既伤害自己也伤害看她们的人。所以我们应该远远躲避这两种祸害，保持足够甚至更多的距离。因为祂不仅命令不要炫耀，甚至要寻求隐藏。这在前面祂已经做过类似的事。\par

关于施舍，祂没有简单地说，而是先说\BibleQuote{小心不要在人面前行施舍}，又加上\BibleQuote{为叫人看见}；然而关于禁食和祈祷，祂却没有这样限制。这是为什么呢？因为施舍完全隐藏是不可能的，而祈祷和禁食则可能。\par

因此，当祂说\BibleQuote{不要叫左手知道右手所作的}时，祂不是指手，而是指严格向所有人隐藏的责任；同样，当祂命令我们进入内室时，并非绝对只在那里祈祷，也不是首要指那里，而是暗示同样的事；同样在这里，祂命令我们\BibleQuote{要抹油}，并非规定我们一定要抹油；否则我们都将成为这法律的违犯者，尤其是那些最努力遵守它的人——住在山上的隐修士团体。祂所吩咐的不是这个，而是因为古人有习惯在寻欢作乐时经常抹油（这一点从\BibleRef{撒下 12:20}的达味和\BibleRef{达 10:3}可以清楚看出），祂说我们要抹油，并非一定这样做，而是尽一切努力，极其严格地隐藏这项善工。为使你确信如此，祂自己在行动中履行了祂用言语所吩咐的：祂禁食了四十天，且是秘密地禁食，既没有抹油，也没有沐浴。然而，尽管祂没有做这些事，祂却毫无疑问地完全做到了没有虚荣。所以，祂也同样吩咐我们，既把假善人摆在我们面前，又通过两次重复的告诫劝阻听众。\par

祂还以\Quote{假善人}这个名称暗示了别的意思。也就是说，祂不仅通过这件事的荒唐和它带来的极端惩罚，而且通过表明这种欺骗只是暂时的，来使我们脱离那邪恶的欲望。因为演员只在观众坐着的时候才显得荣耀；或者即使在观众面前也并非完全如此，因为大多数观众知道他是谁，在演什么角色。然而，一旦散场，他就更清楚地被所有人识破。你看，虚荣的人必然要经历这些。因为即使在这里，大多数人也看出他们并非自己所显示的那样，只是戴着面具；但在将来，当一切事物都\Quote{赤裸敞开}时，他们将更清楚地被揭露。\par

祂还通过另一个动机使听众远离假善人，即表明祂的诫命是轻省的。因为祂并没有使禁食更严格，也没有命令我们更多的禁食，而是不要失去其奖赏。所以，似乎难以承受的事，是我们和假善人共同的，因为他们也禁食；而最轻省的事，即劳苦之后不失去报酬，祂说：\Quote{我所吩咐的就是这个}，没有给我们的劳苦增加任何东西，而是稳妥地为我们积聚工资，不让我们像他们那样空手而归。他们甚至不愿意效仿在奥林匹克运动会上摔跤的人，尽管有如此多的观众和王子坐在那里，他们却只求取悦一个人，即在他们中间判定胜利的那个人；即使那人比他们低下得多。但是你呢，虽然你有双重的理由向祂展示胜利：第一，祂是判定者；第二，祂远超剧场中所有在座的人——你却向那些人展示，他们非但无益，反而暗中给你造成最大的损害。\par

然而，连这一点我也没有禁止，祂说。只是，如果你也渴望在人前炫耀，且等候，我也要更丰盛地、大大有益地赐予你。因为目前，这完全使你与那属于我的荣耀分离，就像轻视这些事使你紧密联合一样；但那时，你将完全安全地享受一切；甚至在最后的赏报之前，你就在这个世界里收获不少的果实，即：你已践踏了所有人间的荣耀，脱离了世人沉重的奴役，成为真正德行的工人。而现在，只要你这样存心，即使你处在旷野，你的所有德行也会被你抛弃，因为无人看见你。这就像是在侮辱德行本身，如果你追求它不是为它本身，而是为了制绳匠、铜匠和下等平民，使坏人和远离德行的人赞赏你。你竟召来德行的敌人来展示和观看它，如同有人选择自律，不是为自律的卓越，而是要在娼妓面前炫耀。你也一样，看来你不是选择德行，而是为了德行的敌人；而你本应正是为此而钦佩她，因为她甚至使敌人称赞她——但要恰当钦佩她，不是为别人，而是为她本身。因为我们自己，当我们被爱不是为我们自己，而是为别人时，就认为这是一种侮辱。同样，我吩咐你在德行上也这样算，不要为别人而追随她，也不要为人的缘故而服从天主，而要为人而服从天主。因为如果你反其道而行之，虽然你看起来追随德行，但你同样激怒了她，如同不追随她的人一样。正如他不服从是因不做，你则是因不合法地做。\newline{}\par

2. \Quote{不要在地上为自己积蓄财宝。}\par

这样，祂在驱除了虚荣心的病患之后（而不是之前），恰当地引入了关于自愿贫穷的论述。因为没有什么比贪求荣耀更能训练人贪爱财富。正是因此，人们想出那些成群的奴隶、那群太监、那些金饰的马匹、银桌以及所有其他更加可笑的东西；不是为了满足任何需要，也不是为了享受任何乐趣，而是为了在众人面前炫耀。\par

此前，祂只说了必须行仁慈；但这里祂也指出了必须行多大的仁慈，当祂说\Quote{不要积蓄财宝}。因为一开始不可能一下子全部引入轻视财富的论述，由于那情欲的暴虐，祂将其分成小块，并在释放听者的心意后，灌输进去，使之变得可接受。因此，你看，祂首先说\Quote{怜悯人的人是有福的}；此后\Quote{与你的对头和好}；然后又\Quote{若有人要控告你，拿你的内衣，连外衣也给他}；但这里，是比所有这些都大得多的事。因为那里祂的意思是说：\Quote{你若看见诉讼临头，就这样做；因为缺乏而脱离争斗，比拥有而争斗更好}；但这里，假设既无对头，也无人与你诉讼，且完全不提任何其他这样的人，祂独自教导轻视财富本身，暗示这法律不仅是为那些蒙怜悯的人，更是为施予者的缘故，所以即使没有人在伤害我们，或拉我们上法庭，我们也可以轻视我们的财产，把它们施舍给需要的人。\par

而且在这里，祂也没有把全部都说出来，即使在这处也是温和地说；尽管祂在旷野中已经极其充分地显示了祂在这方面的争战。\BibleRef{玛 4:9} 然而祂没有说出这一点，也没有提出来；因为揭示的时机尚未到来；但祂暂时寻找理由，在谈论此主题时，保持顾问而非立法者的地位。\par

因为祂说了\BibleQuote{不要在地上积蓄财宝}之后，又加上了\BibleQuote{那里有虫蛀、锈蚀，也有贼挖洞偷窃}。\par

目前，祂从地方和毁坏财宝之物两方面，指明这里的财宝之害和那里的财宝之益。并且祂没有在这里停止，还加上另一个论证。\par

首先，祂从他们最害怕的事来敦促他们。因为祂说：\Quote{你们怕什么？} \Quote{怕你们的东西若施舍就耗尽了吗？不，你们施舍吧，这样它们就不会耗尽；而且更重要的是，它们非但不会耗尽，反而会得到更大的增加；是的，因为天上的东西也要加给它们。}\par

然而，祂暂时没有说出这一点，而是放在后面。但就目前而言，祂提出了最能说服他们的一点，即财宝会为他们保留下来而不耗尽。\par

而且祂从两方面吸引他们。因为祂不仅说\Quote{你若施舍，它就得以保存}；还威胁了相反的事，即若不施舍，它就毁灭。\par

看看祂说不出的明智。因为祂也没有说\Quote{你只是把它们留给别人}；因为这对人也是愉快的；然而祂用一个新理由恐吓他们，指出连这一点他们也得不到：因为即使没有人的欺诈，那些必定会欺诈的东西，\BibleQuote{飞蛾}和\BibleQuote{锈蚀}也在那里。因为虽然这种损害似乎很容易防止，但它却是不可抵抗和无法控制的，无论你如何设计，你都无法制止这种损害。\par

\Quote{那么，飞蛾会毁掉金子吗？} 即使不是飞蛾，也有贼。\Quote{那么，所有人都被掠夺了吗？} 即使不是所有人，也是大部分人。\par

3. 因此，祂又加上另一个我已提过的论证，说：\par

\BibleQuote{因为你的财宝在哪里，你的心也必在那里。}\par

因为即使这些事一件也不发生，祂说，你们也必受不小的损害：被钉在下面的事物上，沦为奴隶而非自由人，将自己逐出天上的事物，无力思考任何崇高之事，而只想着金钱、利息、借贷、收益和卑贱的交易。还有什么比这更悲惨呢？因为事实上，这样的人比任何奴隶更糟，给自己带来最沉重的暴政，放弃了最首要的东西，即人的尊贵与自由。因为无论谁向你们讲论，你们都无法听见任何与你们有关的事，因为你们的心思被钉在金钱上；反而像一条拴在坟墓上的狗，被财富的暴政捆绑，比任何锁链更沉重，向所有靠近的人狂吠，你们只有这一项持续的工作：为他人看守你们所积蓄的。还有什么比这更悲惨呢？\par

然而，因这话对于听众的心智过于高深，且一般人既不易看清其害处，也不易看清其益处，反而需要一种更自制的精神才能领悟两者，所以祂先将此点放在那些显而易见的题目之后，说：\BibleQuote{因为你的财宝在哪里，你的心也必在哪里。}\BibleRef{玛 6:21} 接着，祂又清楚地阐释，将祂的言论从理智层面转移到可感觉的层面，说：\par

\BibleQuote{身体的眼睛就是灯。}\BibleRef{玛 6:22}\par

祂所说的是这样的意思：不要将金子埋在地下，也不要作任何这类事，因为你不过是给蛾子、锈蚀和盗贼积蓄。即使你完全逃脱了这些祸患，你仍无法逃脱心灵的奴役、以及将它钉在一切下面的事物上：\BibleQuote{因为你的财宝在哪里，你的心也必在哪里。}\BibleRef{玛 6:21} 因此，如果你在天上积蓄财宝，你不仅收获这一果实——获得这些事的赏报——而且在此世就已得到报偿：你在那里找到了港湾，将你的情感寄托在那里，并关心那里的事（因为你在哪里积蓄了财宝，显然你也会把你的心转移过去）；同样，如果你在地上这样做，你将会经历相反的事。\par

但如果这话对你来说晦涩难懂，请听紧接着的话。\BibleQuote{身体的眼睛就是灯；所以，如果你的眼睛单纯，你全身就充满光明。但如果你眼睛邪恶，你全身就充满黑暗。那么，如果你里面的光明成了黑暗，那黑暗该是多么大啊！}\BibleRef{玛 6:22-23}\par

祂将祂的言论引向更接近我们感官的事物。我是说，因为祂已谈到心灵被奴役、被俘虏，而许多人不易辨识这一点，所以祂将教训转至外在的事物——那些摆在人眼前的事物——好让那些外在的事也能使内在的事达于人的理解。因此祂说：\Quote{如果你们不知道，}祂说，\Quote{心思受到损害是怎样一回事，那就从身体的事上来学习吧；因为眼睛之于身体，正如心思之于灵魂。} 因此，正如你不愿穿戴黄金、身着丝绸，而眼睛却被弄瞎，反而认为眼睛的健康远比这一切多余之物更可贵（因为如果你丧失或败坏了这健康，你其余的生命便毫无益处）：正如当眼睛被弄瞎时，其他肢体的能力大多随之丧失，因为它们的光明被熄灭了；同样，当心思败坏时，你的生命将被无数邪恶所充满。因此，正如在身体上我们的目标是保持眼睛健康，同样在灵魂上也要保持心思健康。但如果我们使那本应光照其余部分的心思残废，那么我们又靠什么来清晰看见呢？正如毁掉源头的人也弄干了河流，同样，熄灭悟性的人也混乱了他此生的一切行为。因此祂说：\BibleQuote{如果你里面的光明成了黑暗，那黑暗该是多么大啊！}\BibleRef{玛 6:23}\par

因为当舵手淹死、蜡烛熄灭、将军被俘时，此后那些受指挥的人还有什么希望呢？\par

因此，祂在此不再谈论财富所引发的阴谋、争讼、诉讼（这些祂以上已暗示过，当祂说：\BibleQuote{对头把你交给审判官，审判官交给衙役。}\BibleRef{玛 5:25}），而是指出这一切中更严重的事——那必然发生的事——如此将我们从邪恶的欲望中引开。因为坐监牢的痛苦远不及心灵被这种病态所奴役；前者并不必然发生，而后者却是贪财的直接后果。这就是为什么祂将这一条放在第一条之后，作为更严重且必然发生的事。\par

因为天主，祂说，赐给我们理智，好使我们驱除一切无知，对事物有正确的判断，并使用这理智如同武器和光明，对抗一切令人痛苦或有害的事，从而保持安全。但我们却为了一些多余无用的事物而背叛了这份恩赐。\par

因为当将军被俘时，再多的金饰兵士有何用？当舵手沉入波涛之下时，一艘装备华丽的船有何益？当眼睛的视线被挖去时，一具比例匀称的身体有何利？因此，正如有人使那本应健康以消除我们疾病的医生染病，然后让他躺在银床上、金屋内，这对病人毫无益处；同样，如果你败坏了那能制伏我们情欲的心思，即使你将心思置于财宝旁，你非但未给它带来任何益处，反而造成了最大的损失，并损害了你整个灵魂。\par

4. 你岂不见祂藉着那些最常引人作恶的事，竟也藉此最有效地劝阻他们，并将他们引回德行吗？\Quote{你们贪求财富，是怀着什么意向呢？}祂说；\Quote{岂不是为了享乐和奢华吗？然而，藉此你们恰恰得不到这些，反会适得其反。}因为，倘若我们双目被剜，因这灾祸便不能感受任何愉快的事；那么，在心灵的邪僻和残废上，更是如此。\par

再者，你们把它埋在地里，是怀着什么意向呢？为了安全保存吗？但在这里，祂说，又是恰恰相反。\par

同样，对于那为了虚荣而禁食、施舍和祈祷的人，祂藉着那些他最渴望的事，劝诫他不要虚荣：——\Quote{你们怀着什么意向呢？}祂说，\Quote{你们如此祈祷和施舍，岂非为了从人而来的荣耀吗？那么，不要这样祈祷，}祂说，\Quote{这样你们将在来日获得它。}——同样，祂也藉着那贪财之人最热衷的事，俘虏了他。\Quote{你们要怎样呢？}祂说，\Quote{要使你们的财富得以保存，并享受快乐吗？这两样我都将丰盛地赐给你们，只要把你们的金子存放在我所吩咐的地方。}\par

的确，后来祂更清楚地揭示了此事对心灵的恶果，就是当祂提到荆棘时；\BibleRef{玛 13:22}，但此刻，即使在这里，祂也已显著地暗示了同样的事，将如此丧心病狂的人描述为昏暗无光。\par

如同在黑暗中的人看不见任何清晰之物，若看见一根绳子，便以为是蛇；若看见山岭和深谷，便吓得要死；同样，这些人也是如此：对于有视力之人毫不惊恐的事，他们却疑虑重重。因此，他们惧怕贫穷；不，不仅贫穷，甚至任何微小的损失。是的，若他们失去一点小东西，那些缺乏必需食物的人也不及他们那样悲伤哀痛。至少许多富人已至于自缢，不能忍受那样的厄运；而且，受人侮辱和恶待，他们觉得如此难堪，以致为此又有许多人竟自绝于今生。因为财富已使他们对一切变得软弱，唯有对服侍财富却不然。当财富命令他们服侍自己时，他们敢于谋杀、鞭打、辱骂和一切羞耻之事。这正出于极端的悲惨：在应当自制之处，却是众人中最柔弱的；而在更需要谨慎之处，反而变得更无耻、更固执。实在，他们所遭遇的，就如同一个将全部家产挥霍于不当之物的人所必须忍受的。因为这样的人，当必须开支之时到来，因一无所有而无法供给，便遭受不治之祸，因他所有的一切早已被滥用殆尽。\par

如同那些在舞台上、精于那些邪术的人，在其中经历许多奇怪而危险的事，但在其他必要而有用的事上，无人比他们更可笑；这些人的情形也是如此。因为那些走绳索的人，展现了如此大的勇气，但若遇某种紧急情况需要勇敢或冒险，他们却不能，甚至不能忍受思考这样的事。同样，那些富人也是如此：为钱财敢于一切，但为了节制之故，却不肯忍受任何事，无论大小。如同前者从事既危险而无益的事务；同样，这些人经历许多危险和跌倒，却达不到任何有益的目的。是的，他们经历双重的黑暗：既因心灵的邪僻而眼睛被盲，又因忧虑的欺骗而陷入浓雾。因此，他们甚至不能轻易看透。因为在黑暗中的，一旦太阳出现，便脱离黑暗；但眼睛被损坏的，即使太阳照耀也不看见；这正是这些人的情形：如今公义的太阳已经照耀，且正在劝诫，他们却不听，因为财富闭上了他们的眼睛。所以他们要遭受双重的黑暗：部分来自他们自己，部分来自对师傅的漠视。\par

5. 让我们于是留心听祂，好使我们虽迟却终能恢复视力。人如何能恢复视力呢？倘若你明白你是如何失明的。那么，你是如何失明的呢？因你的邪欲。因为贪爱钱财，如同一种邪恶的体液聚集在明亮的眼球上，使云翳变得浓厚。\par

但这云翳若我们愿意接受基督教训的光束，若我们愿意听祂劝诫并说：\Quote{不要在地上为自己积蓄财宝}，便易于驱散和打破。\par

\Quote{但是，}有人说，\Quote{聆听于我何益，只要我仍被欲望所困？}首先，持续聆听本身就有力量摧毁欲望。其次，即使欲望继续纠缠你，你要想，这件事其实根本算不上欲望。因为这是什么样的欲望呢？——陷入痛苦的束缚，屈服于暴政，四面被困，居于黑暗，充满混乱，忍受无益的劳苦，为别人保存财富，且常为仇敌保存？这些事与何种欲望相符？或者，它们岂不更应被逃避和厌恶？什么样的欲望会在盗贼中间储藏财宝？相反，如果你真渴望财富，就把它移到安全无扰的地方。因为你现在的所作所为，并非渴望财富，而是渴望束缚、侮辱、损失和不断的烦恼。然而，假如世上有人向你显示一个不受侵扰的地方，哪怕他带你进入旷野，保证你能安全保存财富——你绝不会迟疑或退缩；你信任他，把你的财物放在那里；但当应许你的是天主而非人时，当祂在你面前摆出的不是旷野，而是天堂时，你却接受相反的事。然而，无论他们在地上的保障有多少，你永远无法摆脱对它们的挂虑。我的意思是，即使你未失去它们，你也永远无法摆脱失去的焦虑。但在那里，你将不会经历这些事；还有，更值得注意的是，你不仅埋葬了金子，还将它播种。因为金子既是财宝，又是种子；甚至比两者更胜。因为种子不会永远存留，但金子却永远长存。再者，财宝不会发芽，但这金子却为你结出永不凋零的果实。\par

6. 但如果你向我提及时间和赏报的延迟，我也能指出并告诉你，即使在此世你已获得多少回报；此外，我还要从今世本身的事物中，证明你以此借口毫无道理。我的意思是，即使在今生，你也会预备许多你自己无法享用的事物；倘若有人责备你，你便以子女和孙辈为理由，以为这样便足以掩饰你多余的劳苦。当你在极度年老时建造华丽的房屋，这些房屋常常在你离世之前尚未完工；当你栽种树木，它们多年后方才结果；当你购置产业和遗产，所有权需经很长时间才归你所有，并且你热切忙于许多其他类似之事，却无法收获其果实——你这样做，究竟是为自己，还是为后代？那么，在此世，你丝毫不因时间延迟而犹豫，尽管这种延迟会使你失去一切劳苦的报酬；而在那里，却因这种等待而完全麻木，这岂不是极度的愚昧？更何况，那等待会带来更大的收益，且不会将你的好处转给别人，反而为你自己赢得恩赐。\par

但此外，这延迟本身并不长久；不，那些事物就在门口，我们不知道在我们这一代，一切与我们有关的事是否都会完成，那可怕的日子是否来临，将那威严而不朽的审判台摆在我们面前。是的，因为大部分预兆已经应验，福音也已在天下宣讲，关于战争、地震和饥荒的预言已经发生，间隔已不大了。\par

你难道没有看见任何征兆吗？为什么，这件事本身就是一个极大的征兆。因为在\Person{诺厄}的时代，他们也没有看见那场普世毁灭的预兆，却在玩耍、吃喝、婚嫁、做他们惯常的事时，被那可怕的审判所吞没。同样，在\Place{索多玛}，他们生活在享乐中，丝毫不怀疑将临的事，就被那时从天而降的闪电烧死了。\par

因此，想到这一切，让我们致力于准备离开此世。\par

因为即使普遍的终结之日不会临到我们，每个人的结局也已在门口，无论他是老年还是青年；人一旦离开此世，便不能再买油，也不能藉祈祷获得宽恕，即使求情的是\Person{亚巴郎}、\Person{诺厄}、\Person{约伯}或\Person{达尼尔}。\BibleRef{路 16:24}\BibleRef{则 14:14}\par

所以，趁还有机会，让我们预先为自己积蓄许多信心，大量收集油，将一切迁往天堂，以便在适当的时候，在我们最需要的时候，享受这一切：藉着我们的主耶稣基督的恩宠和对人的爱，愿光荣与权能归于祂，从现在直到永远。阿们。\par

\SourceRecord{Translated by George Prevost and revised by M.B. Riddle.；Nicene and Post-Nicene Fathers, First Series；Vol. 10.；Edited by Philip Schaff.；Buffalo, NY: Christian Literature Publishing Co.,；1888.；<http://www.newadvent.org/fathers/200120.htm>.}

% structure: p0001 source=h1 target=section reason=page-title

\section{玛窦福音释义 第二十一篇}

\BibleRef{玛 6:24}\par

\Quote{没有人能事奉两个主人：他或是要恨这一个，而爱那一个，或是要依附这一个，而轻忽那一个。}\par

难道你没有看见祂如何逐步地引导我们远离现世的事物，并更详尽地介绍祂关于自愿贫穷的教导，从而推翻贪财的统治吗？\par

因为祂并不满足于先前那些众多而伟大的言论，又增添了其他更令人警醒的话。\par

因为，如果我们要因财富而失去事奉基督的机会，还有什么比祂现在所说的更令人警醒呢？反之，如果我们因轻视财富，而对祂的爱戴和爱德得以成全，还有什么比这更令人渴望呢？因为我不断重复的，现在也同样要说：祂以双重方式敦促听众服从祂的言论——既通过有益的一面，也通过有害的一面；恰如一位高明的医生，既指出疏忽导致的疾病，也指出顺从带来的健康。\par

例如，请看祂指明这是何种裨益，又如何借着使人摆脱相反的事物来确立其益处。祂说：\Quote{财富不仅在这一方面伤害你——它使你招来强盗，又使你心志极度昏昧——而且也在使你被排除在天主的服务之外，成为无生命的财富的俘虏；它在两方面都加害于你：一方面使你成为本应管理的事物的奴隶，另一方面使你被排除在天主的服务之外，而事奉天主是你最不可少的。}因为正如在另一处，祂指出祸患是双重的：既在此处积蓄，\Quote{那里有虫蛀}，又不积蓄在那里，那里的守护是坚不可摧的；同样，在此处，祂也指出损失是双重的：既使人远离天主，又使人受制于玛门。\par

但祂并非直接陈述，而是先以一般原则来确立，如此说：\Quote{没有人能事奉两个主人}，这里指的是两个命令相反的主人，因为若非如此，他们甚至不是两个。因为\BibleQuote{众多信友都是一心一意}\BibleRef{宗 4:32}，但他们却分属许多身体；然而他们的一致使多人成为一体。\par

然后，为加强力量，祂说：\Quote{他不但不事奉，反而会恨恶和厌弃：}祂说：\Quote{因为不是恨这个，爱那个，就是依附这个，轻视那个。}这似乎是把同样的话说了两次；但祂并非无缘无故选择这种形式，而是为了表明向善的改变是容易的。我是说，免得你说：\Quote{我一旦成为奴隶；我被财富的暴政所辖制，}祂指出人有可能转移自己，正如从第一个转向第二个，也同样可以从第二个转向第一个。\par

你看，祂如此一般地说话，为要说服听众成为祂话语的公正审判官，并按事物本身的性质来判决；当祂确认了听众的同意后，才启示自己。因此祂立即加上：\Quote{你们不能事奉天主而又事奉玛门。}让我们战栗地思考，我们使基督说出了什么：在天主的名号旁，竟然提到了黄金。但如果这令人震惊，那么在实际行为中，我们宁愿服从黄金的暴政而非敬畏天主，就更令人震惊了。\par

\Quote{那么，难道这在古人中不可能吗？}绝不可能。\Quote{那么，}有人说，\Quote{亚巴郎和约伯怎么得到称赞呢？}请不要向我提那些富有的人，而是那些事奉财富的人。因为约伯也是富有的，但他并不事奉玛门，而是拥有它并统治它，他是主人，而不是奴隶。因此，他拥有这一切，就像管理别人的财物一样；不仅不剥夺他人，甚至把自己所有的也给予需要的人。更有甚者，当他拥有它们时，它们并不使他快乐：所以他也宣告说：\BibleQuote{我若因财物丰盛而欢喜，}\BibleRef{约 31:25}因此当它们失去时，他也不悲伤。但现在的富人不像他那样，反而比任何奴隶更糟糕，仿佛向某个暴君进贡一般。因为他们的心思好像被贪财所占据的堡垒，从那里每天向他们发出充满一切不义的命令，无人敢违抗。所以不要过于狡辩。因为天主已经明确宣布并断言，事奉这一位和另一位是不可能一致的。因此，你不要说\Quote{这是可能的。}当一位主人命令你强夺暴力，另一位却要你放弃财产；一位要你贞洁，另一位却要你犯奸淫；一位要你醉酒奢华，另一位却要你节制口腹；一位又要你轻视现世之物，另一位却要你紧紧抓住当下；一位要你欣赏大理石、墙壁和屋顶，另一位却要你蔑视这些，而尊重自制：这些怎么可能一致呢？\par

祂在这里称玛门为\Quote{主人}，不是因为它的本性，而是因为向它屈膝者的悲惨。同样，祂也称\Quote{肚腹为神}\BibleRef{斐 3:19}，不是因为这位女主人的尊贵，而是因为那些受奴役者的悲惨：这比任何刑罚更恶劣，在惩罚之前，对于陷入其中的人来说，就足以成为报应。因为那些被定罪的犯人，有谁像这些人这样悲惨呢？他们本有上主作为他们的主，却从那温和的统治逃离，投奔这严苛的暴政，而且他们的行为甚至在此生就带来如此大的损害？的确，他们因此而遭受的损失难以言表：诉讼、骚扰、纷争、劳苦，以及灵魂的盲目；而比这一切更可悲的是，他们从最高的福乐中坠落；因为成为天主的仆人就是这样一种福乐。\par

3. 如你们所见，祂现已从各方面教导了轻视财富的益处，既是为了保全财富本身，也是为了灵魂的喜乐、获得自制以及确保虔诚；接着祂便着手确立这条诫命的可行性。因为这对于最好的立法尤其重要：不仅要指示何为有益，还要使其成为可能。为此祂接着说：\par

\BibleQuote{不要为你们的生命忧虑吃什么。}\par

也就是说，免得他们说：\Quote{那么，如果我们抛弃一切，我们怎能活下去？}对此异议，祂随后非常适时地予以驳斥。因为正如一开始若祂就说\Quote{不要忧虑}，这话会显得沉重；同样，既然祂已显明贪婪所产生的恶果，这劝诫随后便易于接受了。因此祂也不是简单地说\Quote{不要忧虑}，而是增加了理由，然后这样吩咐。在说了\BibleQuote{你们不能事奉天主又事奉玛门}之后，祂又加上\BibleQuote{所以我告诉你们，不要忧虑。}\BibleQuote{所以}是为了什么？因为那无法言喻的损失。因为你们所受的伤害不仅在于财富，更在于最要害的部分，在于那倾覆你们救恩的事；因为它将你们从创造你们、眷顾你们、爱你们的天主那里抛离。\par

\BibleQuote{所以我告诉你们，不要忧虑。}这样，在祂表明伤害无法言喻之后，此时才使诫命更加严格；祂不仅命令我们抛弃所有，甚至禁止为必需的饮食忧虑，说：\BibleQuote{不要为你们的灵魂忧虑吃什么。}这并非因为灵魂需要食物，因为它是无形的；而是依照通常的说法。因为灵魂虽不需要食物，但除非身体得到喂养，它便不能留在身体中。祂这样说并非简单直接，而是又从我们已经拥有的事物以及其他例证中提出论据。\par

从我们已经拥有的事物出发，这样说：\par

\BibleQuote{难道灵魂不是贵于食物，身体不是贵于衣服吗？}\par

因此，那赐予更大者的，岂不也要赐予更小者？那创造了被喂养的肉身的，岂能不赐予食物？所以祂不是简单地说\BibleQuote{不要忧虑吃什么}或\BibleQuote{穿什么}，而是说\Quote{为了身体}和\Quote{为了灵魂}；因为祂要从这些出发加以证明，以比较的方式推进论述。灵魂祂已一次赐下，且保持不变；而身体却每日增长。因此，祂既指出前者的不朽，又指出后者的脆弱，接着说：\par

\BibleQuote{你们中谁能增加一肘身高？}\BibleRef{玛 6:27}\par

这样，祂对灵魂不再多说，因为它不增长，只论及身体；也由此显明：增长的并非食物，而是天主的眷顾。\Person{保禄}以别的方式表明这点时说：\BibleQuote{可见，栽种的不算什么，浇灌的也不算什么，只在那使之生长的天主。}\BibleRef{格前 3:7}\par

祂用这种方式从我们已经拥有的事物来敦促我们；又从其他事物的榜样，说：\BibleQuote{你们看那天上的飞鸟。}\BibleRef{玛 6:26}以免有人说：\Quote{我们通过忧虑来做好事。}祂既以更大的事（即灵魂和身体）来劝阻他们，又以更小的事（即飞鸟）来劝阻。因为如果祂对非常低级的事物如此眷顾，岂不更要赐予你们？祂对群众这样说，因为当时他们还是普通人；但对魔鬼却不是这样，而是如何？\BibleQuote{人生活不只靠饼，而也靠天主口中所发的一切言语。}\BibleRef{玛 4:4}但在这里祂提到飞鸟，而且以一种大大羞辱他们的方式；这对劝诫极有价值。\par

4. 然而，有些不敬神的人竟疯狂到攻击祂的比喻。他们说：因为那加强道德原则的事，不应当用天生的优势来激励。对于那些动物，他们补充说，这是出于本性。对此我们该说什么呢？即使于它们是天生的，但我们仍可能透过选择达到。因为祂并没有说\Quote{看飞鸟如何飞翔}，这对人是不可能的；而是说它们不忧虑而被喂养，这是如果我们愿意，我们也能做到的事。那些在行动上完成此事的人已经证明了这一点。\par

因此，我们极当赞叹我们立法者的深思远虑，因为祂本可以人间的例子来作比喻，也可以提及\Person{梅瑟}、\Person{厄里亚}、\Person{若翰}以及类似的、不为生计操心的人，但为了更深刻地触动他们，祂却提到了无理性的受造物。若祂提到那些义人，他们或许能说：\Quote{我们尚未达到他们那样的境界。}但如今，祂默而不提他们，反而举出天空的飞鸟，从而断绝了他们的一切借口。在此处，祂也效法了旧法律。诚然，旧盟约同样引导人去看蜜蜂、蚂蚁、斑鸠和燕子。\BibleRef{耶 8:7} 这并非不光彩的标记：因为那些动物凭本性所拥有的，我们也能藉自己的选择行为来完成。既然祂为我们而存在的受造物如此关怀，更何况对我们呢？既然祂眷顾仆人，更何为主人？因此祂说：\Quote{你们看那天上的飞鸟}，祂没有说：\Quote{因为它们不从事买卖、不作交易}——这些原是严加禁止的事。祂说的是什么呢？\Quote{它们不播种，也不收割。}\Quote{那么，}有人也许会说：\Quote{我们就不必播种了吗？}祂并没有说：\Quote{我们不必播种}，而是说：\Quote{我们不必忧虑}；祂既不是说人不该工作，而是说人不可自卑，不可用忧虑折磨自己。因为祂也命令我们得食物，但不是\Quote{通过忧虑}来获得。\par

关于这一教训，\Person{达味}从古时便奠定了基础，以隐晦的方式说：\Quote{祢张开祢的手，满足一切生物的需求}；又说：\Quote{祂赐食物给牲畜，向啼叫的小乌鸦施恩。}\par

\Quote{那么，谁没有忧虑呢？}有人会说。难道你没有听见我举出那么多义人吗？难道你没有看见\Person{雅各伯}，他离开父家时一无所有吗？难道你没有听见他祈祷说：\Quote{若主赐我食物吃，衣服穿}？\BibleRef{创 28:20} 这不是忧虑之人的行为，而是凡事求靠天主之人的行为。宗徒们也达到了这境界，他们舍弃一切，毫不忧虑；此外，还有\Quote{五千人}和\Quote{三千人}。\par

5. 但若你听了这些高超的言语，仍不能摆脱这些沉重的束缚，那么请思虑此事之无益，从而平息你的忧虑。因为\par

\Quote{你们中间谁能用忧虑使自己的身量增加一肘呢？}\BibleRef{玛 6:27}\par

你看见祂如何以显而易见的事，揭示了隐晦的事吗？祂说：\Quote{正如关乎你们的身体，你们无法藉忧虑稍稍增加它的尺寸，同样，你们也无法藉忧虑获得食物。你们无论如何思虑，也是如此。}由此可见，即使在我们看似行动之时，成就一切的并非我们的勤勉，而是天主的眷顾。因此，若祂舍弃我们，任何挂虑、焦虑、辛劳或其他类似之事，终将一无所成，一切都将完全消逝。\par

因此，我们不要以为祂的命令是不可能的：因为有许多人正在忠实地履行它们。若你不认识他们，那并不足为奇，因为连\Person{厄里亚}也曾认为自己是孤单的，但却被告知：\Quote{我为自己保留了七千人。}由此可见，现今也有许多人活出这种生活，就像当时的\Quote{三千人}和\Quote{五千人}一样。若我们不信，并非因为没有人行善，而是因为我们离行善甚远。正如嗜酒之徒不易相信有人甚至不沾水（但在我们时代，许多独修者确实做到了）；又如与无数女人纠缠者，不易相信守贞生活是易行的；又如掠夺他人财物者，不易相信有人甘愿舍弃自己的财物：同样，那些每日为无尽忧虑而消磨自己的人，也不易接受此事。\par

至于确实有许多人已达此境界的事实，我们甚至可以从我们这世代中实践这种克己的许多人身上证明。\par

但对你而言，眼下仅需学习不可贪恋，并知道施舍是一善事；且要知道你当分施你所有之物。亲爱的，若你忠实地履行这些事，你必会迅速进展到其余那些事。\par

6. 因此，目前让我们放下过分的奢侈，学习节制，并学习以诚实的劳动获取我们应有的一切。因为就连真福\Person{若翰}，当他与收税的人和士兵交谈时，也命令他们\Quote{当以薪饷为满足。}\BibleRef{路 3:14} 他虽渴望引导他们达到另一种更高的自制，但因他们尚不适合，便谈及了次要的事。因为若他提及比这些更高的事，他们将无法投身其间，也会从其他的事上坠落。\par

正因为此，我们也在这些次要的义务上训练你们。是的，因为我们知道，目前自愿贫穷的重担对你们来说过于沉重，而天堂与大地之间的距离，也不及这种自我舍弃离你们之远。那么，让我们至少抓住那些最低的诫命吧，因为即使这一点也不是小小的鼓励。然而，在外邦人中，有些人甚至已经做到了这一点，尽管精神并不正确，他们舍弃了自己所有的财产。不过，我们满足于你们的情况，只要你们慷慨地施舍；因为如果我们这样前进，很快也会达到那些别的义务。但如果我们连这一步都没有做到，我们这些奉命超越旧约下的人们，却显得不如外邦哲学家，那还配得什么恩宠呢？我们这些应当成为天使和天主子的人，连做人的本性都不能完全持守，还有何言可说？因为掠夺和贪图不属于人的温良，而属于野兽的凶残；是的，侵犯邻人财物的人比野兽更坏。因为对野兽来说，这是出于本性，但我们这些蒙受理性尊荣的人，却堕落至那种违反本性的卑劣，还能得到什么宽恕呢？\par

那么，让我们考虑摆在我们面前的这训诫的尺度，至少奋力达到中等的地位，这样我们既能免于将来的惩罚，又能按部就班地抵达一切美善的顶峰；愿我们众人，靠着我们的主\Person{耶稣基督}的恩宠和爱人，达到这顶峰，愿光荣和权能归于他，至于无穷之世。阿们。\par

\SourceRecord{Translated by George Prevost and revised by M.B. Riddle.；Nicene and Post-Nicene Fathers, First Series；Vol. 10.；Edited by Philip Schaff.；Buffalo, NY: Christian Literature Publishing Co.,；1888.；<http://www.newadvent.org/fathers/200121.htm>.}

% structure: p0001 source=h1 target=section reason=page-title

\section{《玛窦福音》第二十二篇讲道}

\Emph{玛 5:28-29。}\par

\BibleQuote{你们看看田间的百合花怎样生长：它们不劳作，也不纺织。可是我告诉你们：连撒罗满在他极盛的荣华时所披戴的，也不如这些花中的一朵。}\par

在谈到我们必需的饮食，并表明连这些也不应挂虑之后，祂接下来转向更容易的话题。因为衣服不如食物那样必需。\par

那么，为什么祂在此处没有使用同样的例子，即飞鸟，也没有向我们提及孔雀、天鹅和绵羊呢？因为那里有许多这样的例子可供借鉴。因为祂要指出这论证可以双管齐下：既从享有如此优美之物的卑贱者，也从赐予百合花装饰的慷慨。为此，当祂为它们装饰之后，便不再称它们为百合花，而是\BibleQuote{田野的草}（\BibleRef{玛 6:30}）。而且祂还不满足于这个名称，又加上了另一个卑贱的境况，说\BibleQuote{今天还在}。祂没有说\BibleQuote{明天就不在了}，而是说了更卑贱的：\BibleQuote{扔在炉里}。祂也没有说\BibleQuote{穿戴}，而是\BibleQuote{这样穿戴}。\par

你看见祂如何在各处大量运用强化和递进吗？祂这样做，是为了深深触动他们。所以祂还加上了\BibleQuote{何况你们呢，祂岂不更要给你们穿戴？}因为这也有很大的强调：\Quote{你们}一词的力量无非是暗示我们族类被赋予的巨大价值及关怀；好像祂在说：\Quote{你们，祂赐给了灵魂，为你们塑造了身体，为你们创造了所有可见之物，为你们派遣了先知，赐予了法律，行了那无数善工；为你们舍弃了祂的独生子。}\par

直到祂清楚证明之后，祂才进而责备他们，说：\BibleQuote{小信德的人哪！}因为这就是劝告者的特质：祂不仅劝诫，也责备，以更加唤醒人们听从祂话语的说服力。\par

祂由此教导我们不仅不要挂虑，也不要被人们衣着的昂贵所迷惑。因为这样的美丽是草的美丽，这样的荣华是青草的荣华；甚至草比这种衣饰更珍贵。那么，你为何要以那些连植物都占尽优势的事物而自傲呢？\par

请看祂如何从一开始就表明这条命令是容易的；通过对比以及他们所惧怕的事物，带领他们远离这些挂虑。当祂说\BibleQuote{你们看看田间的百合花怎样生长}时，祂加上了\BibleQuote{它们不劳作}：可见祂颁布这些命令，是渴望使我们摆脱劳苦。事实上，劳苦不在于不挂虑，而在于为这些事情挂虑。正如祂说\BibleQuote{它们不播种}时，所废除的不是播种，而是焦虑；同样，祂说\BibleQuote{它们不劳作也不纺织}时，所废除的不是工作，而是挂虑。\par

但是，如果撒罗满也被它们的美所超越，而且不是一次两次，而是整个统治期间——因为不能说他在某个时候穿着这样的衣服，之后就不再如此；其实他连一天也没有装饰得如此美丽：因为基督通过说\BibleQuote{在他极盛的荣华时}表明了这一点——并且如果撒罗满并未超越这朵花而嫉妒另一朵，而是向所有花朵都让步（所以祂也说\BibleQuote{像这些中的一朵}：因为如同真品与仿品之间，这些衣袍与那些花朵之间的差距也是如此之大）——那么，如果那位比所有古代君王都更荣耀的撒罗满承认自己不如花朵，你何时才能超越，或者即使接近那完美的形态呢？\par

此后，祂教导我们根本不要追求这样的装饰。请看其结局：在它的荣耀之后，\BibleQuote{被扔在炉里}。如果祂对卑微、无价值、没有大用处的事物都显示了如此大的关怀，祂又怎会放弃你——所有受造物中最重要的一位呢？\par

那么，祂为何使它们如此美丽呢？为的是显示祂的智慧和祂大能的卓越，使我们从万物中认识祂的荣耀。因为不仅\BibleQuote{诸天宣扬天主的荣耀}，大地也同样；达味在说\BibleQuote{结果实的树木和所有香柏木，请赞美主}时也宣告了这一点。因为有些借着果实，有些借着高大，有些借着美丽，向那创造者发出赞美：这也是智慧极大卓越的一个标志，即在甚至非常卑贱的事物上（有什么比今天还在、明天就不在的更卑贱呢？）祂也倾注了如此大的美丽。那么，如果祂赐给草它所不需要的东西（因为美丽对于烧火有何帮助？），祂岂不更要赐给你所需要的东西吗？如果那一切中最卑贱的事物，祂慷慨地装饰了，而且不是出于需要，而是出于慷慨，那么祂岂不更要尊荣你——一切中最尊贵的事物——在必需的事上吗？\par

2. 现在，如你所见，祂已经证明了天主眷顾的伟大，接下来本应责备他们，但即使在这一点上，祂也节省了言辞，没有责备他们缺乏信德，只说是信德薄弱。因为祂说：\BibleQuote{如果天主这样装饰田野的草，何况你们呢！小信德的人哪！}（\BibleRef{玛 6:30}）\par

然而，所有这些事确实是祂亲自做的。因为\BibleQuote{万物是借着祂而造成的；凡受造的，没有一样不是借着祂而造成的。}（\BibleRef{若 1:3}）但祂至今还没有任何一处提到自己；因为在那时，只要祂在每条诫命中说\Quote{你们一向听古人说，我却对你们说}，就足以表明祂的圆满权能了。\par

既然如此，当祂在随后的事例中也隐藏自己，或谈论一些关于祂自己的卑微之事时，不必惊奇。因为当时祂只有一个目标，就是使祂的话成为他们容易接受的，并在各方面证明祂并非天主的某种对手，而是与圣父同心合意。\par

因此，祂在这里也是如此。因为祂用这么多的话，不断将祂呈现在我们面前，赞叹祂的智慧、祂的眷顾、祂对一切大小事物的温柔关怀。因此，当祂谈到耶路撒冷时，祂称它为\BibleQuote{大王的城}，\BibleRef{玛 5:35}；当祂提到天时，祂又称它为\BibleQuote{天主的宝座}，\BibleRef{玛 5:34}；当祂论及祂在世界的安排时，祂又将这一切归于祂，说：\BibleQuote{祂使太阳升起，照耀恶人也照耀善人；降雨给义人也给不义的人}，\BibleRef{玛 5:45}。在祈祷中祂也教我们说，\BibleQuote{国度、权柄和荣耀都是祂的。}在这里，当祂论及祂的眷顾，并说明祂即使在微小事物中也是最杰出的艺术家时，祂说：\BibleQuote{祂给田野的青草穿上衣服}。祂没有称祂为自己的父，而是称为他们的父；为的是借这尊称责备他们，并使祂称祂为父时，他们不再不悦。\par

如今，若连基本必需品都不可挂虑，我们这些为奢侈之物挂虑的人，还能得到什么宽恕呢？或者说，那些甚至不睡觉，为要夺取他人之物的人，又能得到什么宽恕呢？\par

3. \BibleQuote{所以，不要忧虑说：我们吃什么？喝什么？穿什么？因为这一切都是外邦人所寻求的。}\par

你看见祂如何再次更加羞辱他们，并在途中表明祂所命令的并非严苛或沉重的事吗？因此，正如当祂说：\BibleQuote{如果你们爱那些爱你们的人}，你们所行的并非什么大事，因为外邦人也同样行；祂提到外邦人，是为了激励他们走向更伟大的事。同样，现在祂也提出他们来责备我们，并表明这是祂向我们索要的必要债务。因为如果我们必须显出比经师和法利塞人更多的事，我们这些远未超越他们，反而停留于外邦人的卑微境地，且效法他们心灵狭隘的人，又能配得什么呢？\par

然而祂并没有止于斥责，而是借此责备并唤醒了他们，用尽一切言辞羞辱他们之后，又用另一个论点安慰他们，说：\BibleQuote{因为你们的天父知道你们需要这一切。}祂没有说\Quote{天主知道}，而是说\Quote{你们的父知道}，为引导他们得到更大的希望。因为如果祂是父，且是这样一位父，祂一定不会在极端痛苦中忽略祂的儿女；因为即使是人，作为父亲，也忍受不了这样做。\par

祂对此还添加了另一个论点。是什么样的论点呢？那就是\Quote{你们需要}这些东西。祂的话是这样说的：什么！难道这些是多余的东西，以至于祂会忽略它们吗？然而，即使在多余的事上，祂也没有在野草的例子中表现出缺乏关怀；但现在这些甚至是必需之物。因此，你认为使你焦虑的原因，我说正是这原因足以使你脱离这种焦虑。我的意思是，如果你说：\Quote{因此我必须要挂虑，因为它们是必需的；}相反，我说：\Quote{不，正因如此，不要挂虑，因为它们是必需的。}因为如果它们是多余之物，我们也不应该绝望，而应确信它们会得到供应；但现在它们是必需的，我们更不应怀疑。因为哪个父亲能忍受连必需品也不供给他的儿女呢？因此，天主必定会最确切地赐予这些东西。\par

因为祂确实是我们本性的创造者，祂完全知道其需要。因此你也不能说：\Quote{祂确实是我们的父，我们所寻求的是必需的，但祂不知道我们需要它们。}因为那认识我们本性本身、创造它、并形成它如此的那位，显然比你这处于缺乏之中的人更知道它的需要。祂按照祂的旨意使我们的本性处于如此大的需要之中，因此祂不会自相矛盾：先使它必然处于如此大的缺乏，另一方面又剥夺它所需要的东西，甚至连绝对必需品也不给。\par

因此，我们不要焦虑，因为焦虑毫无益处，只会折磨自己。因为无论我们挂虑与否，祂都赐予；而且我们越不挂虑，祂赐予的反而越多。你的焦虑能得着什么，不过是自我施加多余的惩罚罢了？因为一个即将赴丰富筵席的人，肯定不会让自己为食物挂虑；一个走向泉水的人，也不为饮水焦虑。因此，既然我们有比任何泉水或无数预备好的宴席更丰富的供应，即天主的眷顾，就让我们不要作乞丐，也不要有卑微的心志。\par

4. 因为除了已经说过的，祂还提出另一个让我们对这些事充满信心的理由，说：\par

\BibleQuote{你们先寻求天国，这一切自会加给你们。}\par

这样，当祂使灵魂免于挂虑之后，祂也提及了天堂。因为祂确实来是为废除旧事，召叫我们迈向更伟大的家乡。因此祂所做的一切，为使我们脱离无益的事物和对尘世的依恋。为此祂也提到了外邦人，说\Quote{外邦人寻求这一切}——他们毕生劳碌只为今生，不顾将来之事，也无思想天堂。但你们，这些并非首要之事，而是别样。因为我们并非为此而生，即吃喝穿衣，而是为能中悦天主，获得将来的福乐。因此，正如这些事物在我们的劳作中是次要的，在我们的祈祷中也应是次要的。所以祂也说：\Quote{寻求天国，这一切自会加给你们。}\par

祂说的不是\Quote{必赐给你们}，而是\Quote{必加给你们}，为使你明白，与将来的伟大事物相比，现今的事物绝非祂恩赐的重大部分。因此，祂甚至没有吩咐我们为这些事祈求，而是当我们为其他事祈求时，要怀有信心，好像这些也附加在那事上。那么，寻求将来的事，你也必获得现今的事；不要寻求看得见的事，你必确实得到它们。是的，因为你为这些事接近你的主，是配不上你的。你本应将全部热忱和关怀用于那些不可言喻的福分，却将其消耗在对易逝之物的欲望上，大大羞辱了自己。\par

\Quote{那么怎样？}有人说，\InnerQuote{祂不是吩咐我们求食粮吗？}不，祂加上了\Quote{日用}，又加上了\Quote{今天}，事实上祂在这里也做了同样的事。因为祂没说\Quote{不要忧虑}，而是\Quote{不要为明天忧虑}，同时给我们自由，又使我们的灵魂专注于那些对我们更为必要的事。\par

为此缘故，祂甚至吩咐我们祈求这些，并非因为天主需要我们的提醒，而是为使我们知道，我们所成就的一切都是借助祂的帮助，并使我们因不断为这些事祈祷而更属于祂。\par

你们看见祂如何再次说服他们，他们必会得到现今的事物吗？那赐予更大的，岂不更将较小的赐予？祂说：\Quote{我告诉你们不要忧虑、不要祈求，并非为使你们受苦、赤身裸体，而是为使你们在这些事上也丰足有馀。}你们看，这比什么都能吸引他们归向祂。同样，在施舍的事上，当祂劝阻他们不要在人前炫耀时，祂主要是借应许更慷慨地供给他们而赢得他们——因为祂说：\BibleQuote{你们的圣父在暗中看见，必要在明处报答你们；}\BibleRef{玛 6:4}——同样在这里，在引导他们脱离对这些事的寻求时，祂说服的主题就是祂应许在他们不寻求的情况下，更丰盛地赐予他们。因此，祂说，我吩咐你们不要寻求，并非为使你们不得，而是为使你们丰盛地得；为使你们以适合你们的方式、带着你们应得的益处而得；为使你们不因忧虑和挂心这些事，而使自己不配得到这些和属灵的事；为使你们不经无谓的痛苦，又不从摆在你们面前的事上失落。\par

5. \BibleQuote{所以不要为明天忧虑：一天的恶足够一天的了。}\BibleRef{玛 6:34} 也就是说，苦难和它的折磨。\BibleRef{玛 5:34} 你们以脸上的汗水吃食粮，难道还不够吗？当你即将甚至从先前的劳苦中获得解脱时，为什么还要加上忧虑带来的额外痛苦？\par

祂在这里所说的\OriginalTerm{恶}{evil}，并非指邪恶，绝非如此，而是指苦难、困苦、灾祸；就像在另一处祂也说：\BibleQuote{城中若有祸患，岂不是上主所为？}也不是指这些，而是指从上而来的鞭打。又说：\BibleQuote{我造平安，也降灾祸。}\BibleRef{依 45:7} 因为在这处，祂也不是说邪恶，而是饥荒、瘟疫，这些被多数人视为灾祸的事；一般人惯于称这些事为灾祸。例如，那五个首长的司祭和先知，当把牛套在约柜上，让它们带着幼犊离去时，\BibleRef{撒上 6:9} 就把那些天降的瘟疫，以及由此在他们内心产生的惊慌和痛苦，称为\Quote{灾祸}。\par

这也是祂在这里的意思，当祂说\Quote{一天的恶足够一天的了。}因为没有什么比担忧和焦虑更使灵魂痛苦。同样，保禄在劝勉独身时也提出忠告说：\Quote{我愿意你们无忧无虑。}\par

但当祂说\Quote{明天自有明天的忧虑}时，祂这样说，并非以为日子会为这些事忧虑，而是因为祂必须对一群尚不完美的民众说话，为使祂的话更具表现力，祂将时间拟人化，按照一般人的习惯向他们讲话。\par

祂在此确实劝告，但随后更使之成为一条律法，说：\BibleQuote{你们不要带金、银，也不要在腰带里带铜钱；路上不要带口袋。}\BibleRef{玛 10:9} 这样，祂先在行动中充分显示了这些，然后更确定地引入关于这些事的言语规定，而诫命也因此变得更容易接受，因为先前已由祂自己的行动所证实。那么祂在哪里以行动证实了这一切呢？听祂说：\BibleQuote{人子没有枕头的地方。}\BibleRef{玛 8:20} 祂不仅满足于此，还在祂的门徒身上充分证明这些，同样塑造他们，却不使他们缺乏任何东西。\par

但也要留意祂的慈爱关怀，祂如何超越任何父亲的情感。如此，祂说：\Quote{我命令这事，别无其他，只是为了要救你们脱离多余的焦虑。因为即使你们今天为明天思虑，明天你们还得再思虑。那么，什么才是过分的？为何要强行让一天承受比它应得的忧愁更多的东西，并且连同它自己的烦恼，还加上次日的重担？而这样做，你根本不可能减轻另一天的负担，反而只是显明你在贪求多余的烦恼。}这样，为了更严厉地责备他们，祂几乎赋予时间以生命，让它作为受委屈者出现，并抗议他们无端的恶待。你们既已领受了这一天，就该为它的事操心。那么，为何还要把另一天的事加给它呢？难道它自己的忧虑还不够重吗？请问，你们为什么还要使它更沉重呢？当立法者，那要审判我们的那一位，说出这些话时，请想想祂所启示给我们的希望是何等美好：祂亲自作证，此生是悲惨而劳苦的，以至于仅仅一天的忧虑就足以伤害和折磨我们。\par

6. 尽管如此，在听到这么多如此严肃的话语之后，我们仍然为这些事思虑，却不再为天上的事思虑了：反而我们颠倒了祂的次序，从两方面对抗祂的话。请看：祂说\Quote{你们一点也不要寻求眼前的事物}；我们却一直寻求这些事物：祂说\Quote{你们要寻求天上的事}；我们却连短短的片刻也不寻求，反而根据我们为世俗之事所显示的忧虑程度，我们对属灵之事就更加漠不关心；甚至远胜于此。但这不会永远顺利，也不能持久。即便我们藐视十天？二十天？一百天？难道我们不是必定要离开，并落入审判者的手中吗？\par

\Quote{但迟延带来了安慰。}什么样的安慰呢？就是每天等待着惩罚和报应吗？不，如果你想要从这迟延中得到一些安慰，就要借着悔改为自己积蓄改过自新的果实。因为如果单单惩罚的迟延在你看来是一种缓解，那么不落入惩罚岂不更是得利？让我们善用这迟延，为的是能完全脱离逼在眉睫的危险。因为所吩咐的事没有一样是沉重或难担的，全都是那么轻松容易，只要我们怀有真诚的决心，即使我们犯下无数的罪过，也能完成一切。例如，\Person{默纳舍}曾经犯下无数不洁之事，他向圣人伸出毒手，在圣殿中引入可憎之物，使全城充满凶杀，并做了许多不可原谅的事；尽管如此，在经历了如此长久和巨大的邪恶之后，他洗净了自己这一切。如何洗净？以什么方式？借着悔改和反省。\par

因为没有，是的，没有任何罪恶不肯屈服于悔改的力量，或者更确切地说，不肯屈服于基督的恩宠。因为只要我们愿意改变，祂就会帮助我们。如果你渴望成为善人，没有人能阻止我们；或者说，魔鬼会阻止我们，但它没有力量，只要你选择最好的，并吸引天主来帮助你。但如果你自己不愿意，反而退缩，祂又如何保护你呢？因为祂不愿意你被迫或被强迫得救，而是自愿得救。因为如果你自己有一个充满仇恨和厌恶你的仆人，他不断地离开你、逃避你，你也不愿意留他，即使你需要他的服务；更何况天主，祂做一切事不是为了自己的利益，而是为了你的得救，祂岂会选择强迫你留下？另一方面，如果你只显示出正确的意向，祂绝不会放弃你，无论魔鬼做什么。因此，我们自己应为我们自己的灭亡负责。因为我们不亲近祂，不恳求祂，不按照应有的方式祈求祂；即使我们亲近，也不是作为需要领受的人，也没有应有的信心，也没有像提出要求那样，而是在一种懈怠和漫不经心的态度中做一切。\par

7. 然而，天主要求我们向祂求取事物，并因此将祂自己大大地看作亏欠于你。因为祂是所有债主中唯一一个，当索求发出时，反视为恩情，并给予我们未曾借给祂的。若祂见那索求者恳切地催迫，祂甚至付清我们未曾给祂的；但若懒散，祂也拖延；并非不愿施予，而是因祂乐于我们向祂提出请求。为此，祂也告诉你们那个夜间来借饼的朋友的比喻；\BibleRef{路 11:5}以及那不敬畏天主、也不尊重人的审判官的比喻。\BibleRef{路 18:1}祂并未停留在比喻上，更在亲自的行动中表明，当祂让那位腓尼基妇人离去时，以伟大恩赐充满她。因藉着她，祂表明：祂赐给恳切祈求者，甚至那本不属于他们的事物。\Quote{因为不合适}，祂说，\Quote{拿儿女的饼扔给狗。}然而祂仍给了，因她恳切地向他求。但藉着犹太人，祂显示：对漫不经心的人，祂连他们自己的也不给。他们因此一无所获，反倒失去了自己的。当这些人因不求而连自己的也未得着时，她却因以恳切进逼而获得了属于他人的事物；狗得了儿女的饼。\OriginalTerm{强求}{importunity}是如此大的善！纵然你是狗，若强求，你将胜过那疏忽的孩子：因为情感所不能成就的，一切皆由强求成就。因此不要说：\Quote{天主是我的仇敌，不会俯听我。}你若不断烦扰祂，祂立刻应允你——即使不是因为你是祂的朋友，也因你的强求。无论是敌意、不合时宜，或任何事都不能成为障碍。不要说：\Quote{我不配，所以不祈祷。}那位叙利亚-腓尼基妇人也是如此。不要说：\Quote{我犯了许多罪，不能向我所惹怒的那位恳求。}因为天主不看功绩，只看心态。若那不怕天主、不羞于人的长官被寡妇所胜，那良善的一位岂不更要被不断的恳求所征服？\par

因此，即使你不是朋友，即使你索求的并非你应得的，即使你挥霍了父亲的财产、长久不见踪影，即使你卑微、最末，即使你怒气冲冲地接近祂、极为不悦；只要愿意祈祷并回头，你必得着一切，且必迅速熄灭愤怒与定罪。\par

但是，\Quote{看，我祈祷}，有人说，\Quote{却没有结果。}为什么？你并不像那些人一样祈祷；我是指如叙利亚-腓尼基妇人、深夜来借饼的朋友、不断烦扰审判官的寡妇、以及耗尽父亲财物的儿子。你若如此祈祷，必速速得着。因为即使人对祂行了侮辱，祂仍是父亲；即使被激怒，祂仍疼爱儿女；祂只寻求一事：不报复我们的冒犯，而是见你悔改并恳求祂。但愿我们也能如此炽热，如同那心肠被激起对我们的爱一般。但这火只求一个开端；你若给它一点火星，就能点燃全然的恩惠火焰。因为祂并非因被侮辱而恼怒，而是因你——侮辱祂的人——如此疯狂。若我们这些恶人，当儿女令我们烦恼时，为他们忧伤；何况天主——祂根本不会受侮辱——岂不更因犯罪者你而大大发怒？若我们按本性爱，何况祂，那位超乎本性而仁慈的？\Quote{因为即便}，祂说，\Quote{妇人能忘记她孕育的果实，我也不会忘记你。}\BibleRef{依 49:15}\par

8. 因此，让我们亲近祂，并说：\Quote{主，是的；因为小狗也吃主人桌子上掉下来的碎屑。}\BibleRef{玛 15:27}让我们\Quote{无论得时不得时}靠近祂；更确切地说，靠近祂永不会不合时宜，因为不持续亲近才是失时。对那愿意给予的一位，祈求永远是合时的；是的，正如呼吸永不合时宜，祈祷也非失时，而不祈祷才是。因为我们需要这气息，同样也需要从祂而来的帮助；我们若愿意，必轻易吸引祂到我们这里来。先知为表明这点并指明祂恩惠的恒常准备，说：\Quote{我们必寻见祂预备好如同清晨。}因我们无论何时靠近，必见祂等着我们的行动。若我们未能从祂永不枯竭的良善中汲取，责任全在我们。例如，祂曾向某些犹太人抱怨说：\Quote{我的仁慈如早晨的云雾，又如早上的露水，转瞬即逝。}祂的意思是这样：\Quote{我固然尽了本分，但你们如同升起的烈日驱散云雾和露水，使之消失；同样，你们也因着极大的邪恶遏制了那不可言喻的恩惠。}\par

这本身又是天主眷顾的一个实例：即使祂见我们不配领受恩惠，祂仍暂缓施恩，免得我们变得漫不经心。但若我们稍有改变，哪怕仅是意识到自己犯了罪，祂便涌出超过泉源，倾注超过海洋；你领受越多，祂越是欢喜；并以此方式更被激发而赐予更多。因为祂视我们的得救和祂向求告者慷慨施予为自己的财富。这似乎是保禄所宣告的，当他说：祂\Quote{对于所有呼求祂的人，祂是丰富的，而且优于一切。}因为我们不祈祷时，祂就发怒；我们不祈祷时，祂就远离我们。为此，\Quote{祂成了贫穷，为使我们成为富足}；为此，祂承受了这一切苦难，为激励我们祈求。\par

因此，我们不要失望，反而有如此多的动机和美好的希望，即使我们每天都犯罪，让我们接近祂，恳求、祈求、寻求罪过的赦免。因为这样，我们将来会更不易犯罪；这样，我们将驱走魔鬼，并呼求天主的慈爱，获得将来的美好事物，藉着我们的主耶稣基督的恩宠和祂对人的爱，愿光荣和权能归于祂，直到永远。阿们。\par

\SourceRecord{Translated by George Prevost and revised by M.B. Riddle.；Nicene and Post-Nicene Fathers, First Series；Vol. 10.；Edited by Philip Schaff.；Buffalo, NY: Christian Literature Publishing Co.,；1888.；<http://www.newadvent.org/fathers/200122.htm>.}

% structure: p0001 source=h1 target=section reason=page-title

\section{《玛窦福音》讲道集第二十三篇}

\BibleRef{玛 7:1}\par

\BibleQuote{你们不要判断，免得你们受判断。}\par

那么，难道我们不该责备那些犯罪的人吗？因为保禄也说同样的话；或者更确切地说，在保禄那里也是基督藉着保禄说话，并且说：\BibleQuote{你为什么判断你的弟兄？……你，为什么轻视你的弟兄？}（\BibleRef{罗 14:10}）以及\BibleQuote{你是谁，竟判断别人的仆人？}（\BibleRef{罗 14:4}）又说：\BibleQuote{所以，时候未到，你们不要判断，只等主来。}（\BibleRef{格前 4:5}）\par

那么，祂为何在别处说：\BibleQuote{斥责、劝勉、鼓励}（\BibleRef{弟后 4:2}），以及\Quote{在众人前责备那些犯罪的人？}还有基督对伯多禄说：\Quote{你去，只在你和他之间指出他的过错；}如果他听不理，再带一两个人；如果仍不听从，便告诉教会。祂又怎样委派了这么多人来责备我们，不仅责备，还要惩罚？对不听任何这些劝告的人，祂命令\BibleQuote{看作外邦人或税吏。}（\BibleRef{玛 18:17}）祂又怎样赐给他们钥匙？如果他们不判断，就在任何事情上没有权柄，领受捆绑和释放的能力也是枉然。\par

此外，如果这样实行，一切都要败坏，无论是在教会、国家还是家庭中。因为如果主人不判断仆人，主母不判断使女，父亲不判断儿子，朋友不互相判断，各种邪恶就会增多。我何必说朋友？如果我们不判断敌人，就永远不能结束仇恨，一切都要颠倒。\par

那么这话是什么意思？让我们仔细留意，免得救恩的药物和平安的法律被人看作颠覆和混乱的法律。首先，祂在下面的话中已经向有理智的人指明了这条法律的卓越性，说：\BibleQuote{为什么你看见你兄弟眼中的木屑，而不理会自己眼中的大梁？}（\BibleRef{玛 7:3}）\par

但若对许多不大注意的人仍显得晦涩，我将设法从头解释。我以为，在这里，祂并非简单地命令我们不要判断任何人的罪，也并非绝对禁止这样做，而是针对那些充满无数罪恶，却因小事践踏他人的人。我想这里也暗示了某些犹太人，因为他们对邻人的小过、无关紧要的事严厉指控，而自己却不知不觉地犯下致命的罪。在结尾处祂也这样责备他们，当祂说：\BibleQuote{你们把沉重难担的担子捆起来，放在别人的肩上，自己却连一个指头也不肯动}（\BibleRef{玛 23:4}），以及\BibleQuote{你们捐献十分之一的薄荷、茴香，却忽略了法律上更重要的事，即公义、怜悯和信德。}（\BibleRef{玛 23:23}）\par

所以，我认为这些人在祂的控诉之内；祂预先阻止他们将来指控祂的门徒的那些事。虽然祂的门徒没有犯过这样的罪，但他们却被认为有过错；例如，不守安息日，不洗手吃饭，与税吏一同坐席。关于这些，祂在别处也说：\Quote{你们滤出蠓虫，却吞下骆驼。}但与此同时，祂也在这些事上制定了一条普遍的法律。\par

保禄对格林多人（\BibleRef{格前 4:5}）也没有绝对禁止判断，而是不要判断自己的上司，并且不要根据不认可的理由；不是绝对禁止纠正犯罪的人。同样，当时祂也不是一概责备所有人，而是责备那些如此行的门徒对老师，以及那些自己犯下无数罪却诬蔑无辜者的人。\par

这就是基督在这里所暗示的意思；不仅仅是暗示，而且还用极大的恐怖和无法以祈祷解脱的惩罚来加以保护。\par

2. \BibleQuote{因为你们用什么判断来判断，你们也要受什么判断。}（\BibleRef{玛 7:2}）\par

就是说，\Quote{你谴责的不是别人，}基督说，\Quote{而是你自己，你使审判席对自己变得可怕，使账目变得严格。}正如在赦免我们的罪时，起点在于我们，同样在判断中，是我们自己定下被定罪的标准。你看，我们不应该责备或践踏他们，而应该劝告；不应该辱骂，而应该建议；不应该傲慢攻击，而应该温柔纠正。因为你不是不宽恕他，而是把你自己交给极端的报复，当需要对他的过犯作出判决时，你却不宽恕他。\par

你看见吗？这两条诫命既容易遵守，又给服从的人带来极大的祝福，反之则给不关心的人带来灾难。因为那宽恕邻人的人，首先使自己从抱怨的根源中解脱，而且不费劳力；那以温柔和宽容调查别人过犯的人，藉着自己的判断为自己预先储存了宽恕的大量余地。\par

\Quote{那么！}你说：\Quote{若有人行淫，难道我不能说行淫是坏事，也不能纠正那行淫的人吗？}不，要纠正，但不要像敌人或对手那样索取惩罚，而要像医生提供药物。因为基督并没有说：\Quote{不要阻止犯罪的人，}而是说：\Quote{不要判断，}即不要刻薄地作出判决。\par

此外，正如我已经指出的，这话不是指重大的事或被禁止的事，而是指那些甚至不被视为过犯的事。所以祂也说：\par

\BibleQuote{为什么你看见你兄弟眼中的木屑？}（\BibleRef{玛 7:3}）\par

是的，如今许多人就这样做：他们若看见一个隐修士穿着不必要的衣服，就拿出我们主的法律来指责他，\newline{}\BibleRef{玛 10:10}，而他们自己却无休止地勒索，天天欺诈他人。若看见他饮食稍丰，便严厉控告，而他们自己却每日纵酒暴食，不知晓他们除了自己的罪过外，还因此为自己积聚更大的火焰，并剥夺了自己的一切辩护。因为关于这一点，即你的行为必须受到严格审查，你首先已制定了法律，通过如此审判邻人的行为。因此，不要觉得你也要经历同样的审判是件痛苦的事。\par

\BibleQuote{假善人！先从你眼中取出木梁吧。} \BibleRef{玛 7:5}\par

在此，祂愿意表明祂对行这样事的人所怀的极大愤怒。因为每当祂要指出罪恶重大，以及随之而来的惩罚和愤怒可怕时，祂总是以斥责开始。正如祂对那个逼债一百银币的人深感不悦地说：\BibleQuote{恶仆！我免了你一切债务，} \BibleRef{玛 18:32}；同样在这里，祂也说：\BibleQuote{假善人！} 因为这样的判断并非出于关怀，而是出于对人的恶意；人一面戴上仁慈的面具，一面却做出极恶之事，使无端的指责和控告附着于邻人，并且篡夺了教师的地位，而他自己甚至不配做个门徒。为此，祂称他为\BibleQuote{假善人}。因为你在别人的事上如此苛刻，连微小之事都看见；何以在自己事上如此松懈，连大事都被你忽略了呢？\par

\BibleQuote{先从你眼中取出木梁吧。}\par

你看见吗？祂并非禁止判断，而是命令先取出你眼中的木梁，然后再纠正其余世人的行为。因为每个人对自己的事比对他人的事更为了解；看见大的比小的更清楚；爱自己胜过爱邻人。因此，你若出于关怀之心而行，我劝你先照顾自己，因为你的罪既更确实也更大。但若你忽略自己，显然你判断弟兄也不是出于关心他，而是出于仇恨，并希望揭露他。即使他应当受判断，也应由一个不犯这样罪的人来判断，而不是由你。\par

这样，因为祂引入了关于克己的伟大崇高教义，以免有人说：如此在言语上实践是容易的；祂愿意表明祂的完全确信，以及祂对上述诸事都无可指摘，并且已完全履行了一切，于是说了这个比喻。这也是因为祂后来也要审判，说：\BibleQuote{祸哉，你们经师和法利塞假善人！} \BibleRef{玛 23:1} 但祂自己却并不受上述罪过的指摘；因为祂既没有取出木屑，眼中也没有木梁，而是清洁无瑕地纠正了所有人的过犯。\Quote{因为一个人若自己犯有同样的过失，就绝不适合审判他人。} 祂说。为何惊奇祂制定这条法律呢？就连那个强盗在十字架上也明白这一点，对另一个强盗说：\Quote{你连天主也不怕吗？我们受的是同样的刑罚，} 表达了与基督相同的情感。\par

但是你，非但未曾取出自己的木梁，甚至没有看见它；而别人的木屑你不仅看见，还加以审判，并试图取出它。如同一个人患了严重的水肿，或任何其他不治之症，却忽视此病，反而责备另一个忽略了轻微肿胀的人一样。若看不见自己的罪是恶，那么坐在审判他人的位置上，而自己却仿佛麻木不仁，眼中带着木梁，便是双倍或三倍的恶了。因为没有一根木梁比罪更沉重。\par

因此，祂这样命令的意思是：一个犯有无数恶行的人，不应成为他人过犯的严厉指责者，尤其当这些过犯微不足道时。祂并非推翻责备或纠正，而是禁止人忽略自己的过失，却幸灾乐祸于他人的过犯。\par

因为这正是人走向大恶的原因，带来了双重的邪恶。因为那惯于轻视自己的重大过失，却苛刻地追究他人轻微过错的人，会从两方面败坏自己：第一，因轻视自己的过失；第二，因与众人结下仇恨和纷争，并每天训练自己变得极其凶狠和缺乏同情心。\par

3. 祂以这卓越的立法排除了这一切之后，又加上另一条命令，说：\par

\BibleQuote{不要把圣物给狗，也不要把你们的珍珠丢在猪前。}\par

\Quote{然而往后，}有人会说，祂命令说：\BibleQuote{你们在暗中所听到的，要在屋顶上宣讲出来。}\BibleRef{玛 10:27}但这与先前的话绝不矛盾。因为祂在那里并非简单地命令向所有人传讲，而是告诉祂应当对谁说话，就叫他们自由地宣讲。祂用\Quote{狗}比喻那些生活在无可救药的不敬虔中、不提供任何改过自新希望的人；用\Quote{猪}比喻那些一直生活在不洁生活中的人，祂判定这一切人不配听这样的教训。\Person{保禄}也显明了这一点，当他说：\BibleQuote{但属血气的人不接受圣神的事，因为为他乃是愚妄。}在许多其他经文中，祂也说生活的败坏是人不能接受更完美道理的原因。因此祂命令不要向他们开门；因为他们学了后反而更加傲慢。因为对于善良和聪明的人，事情显明时是可敬的；对于麻木的人，事情未知时反而更可敬。既然从他们的本性不能学习，祂说：\Quote{让这事隐藏着，}祂说，\Quote{至少因无知而敬畏。因为猪根本不知道珍珠是什么。既然它不知道，也不要让它看见，免得它践踏它所不知道的。}\par

因为除了给那些如此倾向的人带来更大损害外，没有其他结果；因为圣物被他们亵渎，他们不知道是什么；并且他们更加自高自大，与我们作对。因为\BibleQuote{免得他们践踏在脚下，转过来撕裂你们}就是指这个。\par

\Quote{当然，}有人说：\Quote{他们应当坚强到在人们学习后依然同样坚不可摧，并且不因别人给我们提供机会。}但这不是事情本身造成的，而是因为这些人就是猪；就像珍珠被践踏，不是因为珍珠真的可鄙，而是因为落在猪中间。\par

祂说\Quote{转过来撕裂你们}很恰当：因为他们假装温顺以便受教；然后学了之后，完全变成另一种人，嘲笑、讥讽我们，视我们为受骗者。因此\Person{保禄}也对\Person{弟茂德}说（见：\BibleRef{弟后 4:15}）：\BibleQuote{你也要提防他，因为他极力反对我们的话；}又在另一处说：\BibleQuote{你要避开这样的人，}并且\BibleQuote{对异端者，一次两次警告后，就该弃绝。}\par

你们看，并不是那些真理给他们提供武器，而是他们自己如此成为愚人，充满更大的固执。因此，让他们停留在无知中并非小益，因为这样他们就不会成为完全的藐视者。但如果他们学了，祸害加倍。因为他们自己丝毫不会因此获益，反而更加受损，而你们将陷入无尽的麻烦。\par

让那些无耻地与所有人交往、并使可畏之事变得可鄙的人听吧。因为我们的奥秘也因此在闭门后举行，不让未入门者进入，不是因为我们认为我们的礼仪有任何软弱，而是因为许多人尚未准备好接受它们。正因为这个原因，祂自己也用比喻向犹太人讲了许多，\BibleQuote{因为他们看却看不见。}为此，\Person{保禄}同样命令\BibleQuote{要知道应当怎样回答每个人。}（见：\BibleRef{哥 4:6}）\par

4. \BibleQuote{你们求，就给你们；你们找，就要找到；你们敲，就给你们打开。}\BibleRef{玛 7:6}\par

因为祂既然吩咐了伟大和奇妙的事，命令人超越一切情感，引领他们直达天堂本身，命令他们追求相似的对象，不是天使和总领天使，而是（尽可能）万有的主；并吩咐祂的门徒不仅自己要妥善完成这一切，还要纠正他人，区别恶人与非恶人、狗与非狗（尽管人心中隐藏着许多）：——免得他们说：\Quote{这些事是沉重和难以忍受的}（果然，后来\Person{伯多禄}说了类似的话：\Quote{谁能得救？}又说：\Quote{如果那人的情形如此，倒不如不结婚}）：因此，为了使他们现在不这样说；正如先前祂已证明一切容易，列举了许多理由来说服人：所以在这一切之后，祂又加上了一切容易的顶峰，为我们劳苦提供了来自恒心祈祷的帮助，这并非小事。这样，祂说，我们不只是自己奋斗，也要呼求上面来的帮助；那帮助必定会来，与我们同在，帮助我们争斗，使一切变得容易。因此祂既命令我们祈求，又答应给予。\par

然而，祂命令我们不仅要简单地祈求，还要带着极大的勤勉和热忱。因为这就是\Quote{寻找}的意思。因为寻找的人，把一切其他事都放在一边，只专注于所寻找的，对在场的人毫无概念。我说的这件事，那些丢失了金子或仆人并努力寻找的人都知道。\par

通过\Quote{寻找}，祂表明了这一点；通过\Quote{敲}，表示我们以热忱和火热的心接近。\par

所以，人啊，不要失望，不要对德行表现出比他们对财富的欲望更少的热情。因为那种东西你常常寻找却找不到，但尽管你知道不一定能找到，你仍动用各种搜索方式；但在这里，既然有必定会得到的应许，你却没有表现出那种热忱的最小部分。即使你没有立刻得到，也不要绝望。因为为此祂说\Quote{敲}，表明即使祂不立刻开门，我们也应留在那里。\par

5. 你若怀疑我的肯定，至少当相信祂的榜样。\par

祂说：\BibleQuote{你们中间有哪个人，儿子向他求饼，反而给他石头呢？}\par

因为，在人中间，你若如此继续，甚至被视为烦扰和可厌；同样，在天主面前，你若不然，反而更激怒祂。你若继续祈求，即使没有立刻获得，也必定会获得。因为门关闭正是为了促使你敲；祂不立即应允，正是为了要你祈求。那么，继续做这些事，你必会获得。为使你不致说：\Quote{那么，我若祈求而不得，怎么办呢？}，祂用那比喻堵住了你的借口，再次设论，用那些人间之事激励我们对此满怀信心；暗示我们不仅应当祈求，而且应当祈求合宜之事。\par

\BibleQuote{你们中间有哪一位父亲，儿子向他求饼，反而给他石头呢？}所以，你若不得，是因为你求石头，这正是不得的原因。因为即使你是儿子，这也不足以使你获得；反而正是这本身阻碍了你获得，即身为儿子，却求无益之事。\par

因此，你也当不求任何世俗之物，只求属神的一切，就必获得。请看撒罗满，因他求了当求之事，何等地迅速获得！可见，祈祷者应有两点：恳切祈求，并求当求之事。祂说：\Quote{你们虽为父亲，也等待儿子们来求；若他们向你们求任何不当之事，你们拒绝给予；反之，若为有益，你们就应允并赐予。}你也要思考这些，未获之前不退却，未寻见之前不离去，未开门之前不松懈。因为，你若带着这样的心意说：\Quote{我若不获得，绝不离去，}你必会获得，只要你所求之事既适合施予者施予，又对祈求者有益。那么，这些事是什么？是寻求一切属神的事；宽恕得罪你的人，并如此近前求赦免；\Quote{举起圣洁的手，不发怒，不怀疑。}我们若如此祈求，必会获得。然而，实际上，我们的祈求简直是嘲弄，是醉汉而非清醒之人的行为。\par

有人说：\Quote{那么，我若求属神之事而不得呢？}你定然没有恳切地敲，或使自己不配获得，或很快就放弃了。\par

或许有人问：\Quote{为什么祂没有说，我们当求什么？}不，祂在之前已经全部提及，并指明了我们应当为何事而近前。那么，不要再说：\Quote{我近前了，却没有获得。}因为，我们不得，绝非由于天主——这位爱护我们甚至超过父亲，超过他们正如良善超越这邪恶的本性。\par

\BibleQuote{你们纵然不善，尚且知道把好东西给你们的儿女，何况你们的在天之父，岂不更要把好东西赐给求祂的人？}\par

祂这样说，并非要贬损人的本性，也不是要谴责我们的种族为恶；而是为了与祂自己的良善相对照，称父爱为不善，以显明祂对人的慈爱何等超越。\par

你看到这无可言喻的论证了吗？它足以唤醒甚至彻底绝望之人，令他怀有良好的希望。\par

在此，祂藉我们的父亲来表明祂的良善，但在之前，祂藉那首要的恩赐——\Quote{灵魂}和身体来表明。祂在任何地方都没有列出万善之最，也没有提出祂自己的降临。——因为那位如此急忙献出祂的圣子受死，\Quote{岂不也把万有与祂一同白白赐给我们？}——这尚未成就。但保禄确实如此宣告说：\BibleQuote{祂没有怜惜自己的儿子，岂不也把万有与祂一同白白赐给我们？}\BibleRef{罗 8:32}然而，祂与他们讲论，仍然基于人间之事。\par

6. 之后，为表明我们既不可在祈祷中自信而疏忽自己的行为，也不可只依靠自己的努力，而既要寻求上面的助佑，又要同时尽己本分，祂将两者相辅相成地陈明。因为，在长时间劝勉之后，祂又教导如何祈祷；教导完如何祈祷，祂又继续劝勉关于我们当行之事；然后，又从那里再次转到不断祈祷的必要性，说：\Quote{祈求}、\Quote{寻找}、\Quote{敲门}。接着，又转到我们自己也须勤勉的必要性。\par

祂说：\BibleQuote{所以，凡你们愿意人给你们做的，你们也要照样给人做。}\par

祂以此概括一切，表明德行是简明的、容易的，且为众人所熟知。\par

祂不仅说\Quote{凡你们愿意的}，而是说\Quote{所以，凡你们愿意的}。祂加上\Quote{所以}并非无意，而是含有深意：祂说：\Quote{你们若愿意蒙垂听，除了我所说的，还要做这些事。}那么，这些事是什么？\Quote{凡你们愿意人给你们做的。}你看见祂如何也借此表明，祈祷之外还需要端正的行为。祂没有说：\Quote{凡你们愿意天主为你们做的，你们也要对邻人做，}免得你说：\Quote{这怎么可能？祂是天主，我是人，}而是说：\Quote{凡你们愿意同伴为你们做的，你们自己也当向邻人显示。}还有什么比这更轻松？还有什么比这更公平？\par

随后，酬报之前的赞美也是极大的。\par

\BibleQuote{因为这就是法律和先知。}由此可见，德行是合乎我们本性的；我们众人自然知晓自己的本分；我们永远不可能以无知为借口。\par

7. \Quote{你们要从窄门进去，因为通向灭亡的门是宽的，路是广的，进去的人很多；但通向生命的门是窄的，路是狭的，找到它的人很少。}\par

然而在这之后祂说，\BibleQuote{我的轭是柔和的，我的担子是轻松的。}\BibleRef{玛 11:30} 而且在祂刚说的话中，祂也暗示了同样的事：那么祂在这里怎么说它是窄而受限的呢？首先，如果你留意，即使在这里祂也指出它是非常轻、容易且可接近的。\Quote{那么，}也许有人会说，\InnerQuote{这窄而受限的路怎么容易呢？}因为它是一条路和一扇门；正如另一条路，虽然宽、虽然宽阔，也同样是一条路和一扇门。这些都没有永恒性，一切都在消逝，无论是痛苦还是今世的美好事物。\par

不仅德行的这部分是容易的，而且它的终点也使它更容易。因为我们的辛劳与劳苦的消逝，以及它们指向一个美好结局（因为它们以生命为终结）的事实，足以安慰那些战斗的人。因此，我们劳苦的暂时性、我们冠冕的永恒性，以及劳苦在先、冠冕在后的次序，必定在辛劳中带来极大的安慰。所以保禄也称他们的苦难为\Quote{轻的}；不是出于事情的本质，而是出于战斗者的心态和对未来的希望。\BibleQuote{因为我们这现时轻微的苦难，}他说，\BibleQuote{正给我们造成一个永远的荣耀分量，我们不是注目那看得见的，而是注目那看不见的。}\BibleRef{格后 4:17}\par

无论如何，看祂在另一方面也使它容易，命令不要与狗交往，也不要将自己交给猪，并要\Quote{提防假先知；}如此从各方面使人们感觉像是在真正的冲突中。而且，称之为窄这个事实本身也大大有助于使它容易；因为它促使他们警醒。正如保禄那时说：\Quote{我们不是对抗血和肉，}他这样说不是为了沮丧，而是为了振奋士兵的精神：同样，祂也为了摇醒旅行者，称这条道路崎岖。祂不仅以此方式使人们警醒，而且还补充说，这条路上也有很多绊脚石；更严重的是，他们甚至不公开攻击，而是隐藏自己；因为假先知就是这类。\Quote{但是不要看这个，}祂说，\InnerQuote{就是它是崎岖而狭窄的，而要看它通向何处；也不要看相反的路是宽而广的，而要看它通向哪里。}\par

祂说这一切，都是为了彻底唤醒我们的热忱；正如祂在别处也说：\Quote{暴力的人以暴力夺取它。}因为凡是在争战中的，当他看到竞赛的裁判惊叹于他努力的艰苦时，就更加振奋。\par

那么，当许多令人烦恼的事由此产生时，不要让我们困惑。因为路是窄的，门是狭的，但城不是。因此，我们不应该在这里寻找安息，也不要在那里期待任何痛苦的事。\par

现在祂说：\Quote{找到它的人很少，}在此祂再次既宣告了众人的大意，又教训祂的听众不要看重多数人的幸福，而要看重少数人的劳苦。因为大多数人，祂说，不仅不走这条路，甚至不选择它：这是极其罪恶的事。但我们不应看重多数人，也不应为此烦恼，而应效法少数人；并且要千方百计装备自己，好行走在其上。\par

因为除了它狭窄之外，在通往那里的路上还有许多要绊倒我们的事。因此祂又加上，\par

8. \Quote{你们要提防假先知，他们来到你们这里，外面披着羊皮，里面却是贪婪的狼。}看哪，与狗和猪一起，还有另一种埋伏和阴谋，比那更严重。因为那些是公认的、公开的，但这些是隐蔽的。因此，对于那些人，祂命令避开；对于这些，祂吩咐人仔细警惕，好像不能一开始就看出它们似的。因此祂也说：\Quote{要提防}；使我们更精确地分辨它们。\par

然后，免得他们听到这是窄而狭的，他们必须走与众人相反的路，必须提防猪和狗，并同这些一起还有另一种更邪恶的种类，就是狼；免得，我说，他们因这许多烦恼而消沉，既要走与大多数人相反的路，又因此要如此焦虑这些事：祂用\Quote{假先知}这个词提醒他们发生在他们祖先的日子的事，因为那时也同样发生这样的事。你们现在不要烦恼（祂这样说），因为没有什么新的或奇怪的事要临到你们。因为对于一切真理，魔鬼总是暗中用相应的欺骗来替代。\par

而且，借着\Quote{假先知}这个形象，我认为祂在这里暗示的不是异端者，而是那些生活败坏却戴着德行面具的人；一般人通常称他们为骗子。因此祂又进一步说，\par

\BibleQuote{凭他们的果实，你们可认出他们来。}\BibleRef{玛 7:16}\par

因为在异端者中，人们常能找到实际的美善，但在那些我所提到的人中，却绝对找不到。\par

也许有人说：\Quote{那么，倘若他们也在这些事上伪装，又当如何？}\Quote{不然，他们很容易被识破；因为这道路的本性如此——我曾吩咐人走这条路，是痛苦而令人厌烦的；但伪善者不会选择受苦，只会做表面工夫；因此他也很容易被揭穿。}既然祂说过\Quote{找到它的人很少}，祂便将那些尚未找到却假装找到的人，从其中清除出去，吩咐我们不要只看那些戴着面具的人，而要看那些真正追求它的人。\par

也许有人说：\Quote{但祂为何不使他们显明，却要我们去寻找他们？}这是为了让我们警醒，时刻准备应战，既要防备伪装的敌人，也要防备公开的敌人；\Person{保禄}所指的正是这种情形，他说他们\Quote{以甜言蜜语欺骗了老实人的心}。因此，当我们如今看到许多这样的人时，不要感到困扰。因为基督从一开始就预言了这事。\par

请看祂的温和：祂没有说\Quote{惩罚他们}，而是说\Quote{不要被他们伤害}，\Quote{不要毫无防备地落在他们中间}。然后，免得你说\Quote{不可能分辨那类人}，祂又用一个人间的例子来论证，这样说：\par

\BibleQuote{人岂能从荆棘上摘葡萄，从蒺藜里摘无花果呢？照样，凡好树都结好果子，而坏树结坏果子。好树不能结坏果子，坏树也不能结好果子。}\BibleRef{玛 7:16}\par

祂所说的是这样：他们毫无温柔甜美之处；只是外表像羊；因此也容易辨认。为避免你丝毫怀疑，祂将此比作某些自然必然性，在这些事上只有一种结果。\Person{保禄}也以此意义说：\Quote{属血肉的思念是死亡；因为它不服从天主的法律，也的确不能服从。}\par

祂重复同一件事，并不是重复啰嗦。但若有人说：\Quote{坏树虽结坏果子，但也结好果子，使分辨变得困难，因为收成有两种。}祂说：\Quote{并非如此，它只结坏果子，绝不可能结好果子；正如相反的情况也是如此。}\par

\Quote{那么怎么说？难道没有好人变成坏人的事吗？相反的情况也发生，生活中这样的例子很多。}\par

但基督并非说，恶人无法改变，或者善人不会堕落，而是说，当他生活在邪恶中时，他不能结出好果子。因为恶人固然可以转变为善，但若继续留在邪恶中，他不会结出好果子。\par

那么，\Person{达味}原是善人，岂没结出坏果子吗？但他不是持续为善，而是改变了；因为毫无疑问，如果他始终如一，就不会犯下所犯的罪。\par

祂也借此堵住那些信口开河之人的口，给所有毁谤者的嘴唇加上缰绳。因为许多人因恶人而怀疑善人，但祂用这句话使他们失去了一切借口。\Quote{因为你不能说：\InnerQuote{我被欺骗、被蒙蔽了；}因为我已给了你这种以他们的果实分辨他们的方法，并吩咐你去看他们的行为，而非随意混淆一切。}\par

9. 既然祂没有吩咐惩罚他们，只是要提防他们，祂为了安慰那些被他们困扰的人，同时也为了警告并改变他们，便以祂手中将施与他们的惩罚作为屏障，说：\par

\BibleQuote{凡不结好果子的树，就砍下来，丢在火里。}\BibleRef{玛 7:19}\par

然后，为使这话不那么严厉，祂又加了：\par

\Quote{所以，你们凭他们的果实，可以认出他们来。}\par

这样，祂不是以恐吓为主导，而是以劝诫和忠告来激发他们的心。\par

这里在我看来，祂也暗示了犹太人，他们正结出这样的果实。因此，祂也以同样的言词提醒他们想起\Person{若翰}的话，描绘他们的惩罚。因为\Person{若翰}也说过同样的话，向他们提到\Quote{斧头}、\Quote{砍下的树}和\Quote{不灭的火}。\par

虽然被焚烧看似是单一的刑罚，但若仔细考察，这其实是两种惩罚。因为被焚烧的人当然也被逐出天主的国；而这后一种惩罚比前者更重。我知道许多人只对地狱感到战栗，但我敢说，失去那荣耀比地狱可怕得多。若无法用言语充分表达，不足为奇。因为我们也不知道那些好事的福乐，因此也不能清楚地领会失去它们的悲惨；\Person{保禄}既清楚知道这些，便明白从基督的荣耀中跌落比一切都更痛苦。那时我们将会知道，当我们实际经历它的时候。\par

但愿这永远不会临到我们，天主的独生子啊，但愿我们永远不会经历这无法补救的惩罚。因为从那些好事中跌落是多么大的灾祸，实在无法准确描述；但我将尽力，通过一个例子，哪怕在微小程度上，向你们阐明这一点。\par

让我们设想一个奇异的孩子，他除了自己的德行外，还拥有对全世界的统治权，并且在各方面都如此有德，以致能使所有人都产生父爱的渴望。你认为这个孩子的父亲会不愿意忍受任何偷窃，只为了不被逐出他的交往吗？他会不欢迎任何大小邪恶，只要能看到并享受他吗？现在，让我们同样就那荣耀来推理。因为没有一个孩子，无论他多么有德，对父亲而言，比我们拥有那些美善的一份，并\BibleQuote{离世与基督同在}更令人渴望和可爱了。\BibleRef{斐 1:23}\par

无疑地，地狱和那种惩罚是无法忍受的。然而，即使设想一万个地狱，也远远比不上失去那蒙福的荣耀、被基督憎恨、听到\BibleQuote{我不认识你们}\BibleRef{玛 25:12}、并因当我们看见他饥饿时没有喂养他而被控告\BibleRef{玛 25:42}。是的，宁愿忍受一千个闪电，也比看见那温和的面容转离我们，那和平之眼不愿注视我们更好。因为如果他在我仍是仇敌、憎恨他、并转离他时，竟如此追随我，甚至不惜牺牲自己，舍己至死：那么，当我这一切之后，竟不肯在他饥饿时施舍一块面包，我将以怎样的眼睛再次看见他？\par

但请注意他在这件事上的温和；他根本不提自己的恩惠，也不说：\Quote{你藐视了那对你行了如此多好事的主}；也不说：\Quote{我，那从无中把你带出，向你吹入灵魂，使你管理地上万物，为你创造了地、天、海、空气及一切存在，曾被藐视，甚至被看得比魔鬼还不值，却仍未撤退，在你之后更有无数思念；那选择了成为奴隶，被鞭打、被唾弃、被杀害，死于最耻辱的死亡，在天上仍为你代祷，无偿赐予你他的圣神，赐你王国，许你如此承诺，愿意成为你的头、新郎、衣裳、房屋、根、肉、饮、牧者、君王，并把你收为兄弟、嗣子、与他同嗣；那把你从黑暗带入光明的国度}——这些事，我说，以及比这些更多的话他本可说，但他一件也未提；而是什么？只提出了罪本身。\par

在此处他仍显出他的爱，并表明他对你的渴望：他不说\Quote{你们进入那为你们预备的火里去}，而是说\Quote{为魔鬼预备的}。而且在告诉他们做错事之前，他也没有耐心提及全部，只提少数。而且在此之前，他先召唤那些做了好事的人，这处也表明他责备他们是公正的。\par

那么，什么惩罚与这些话同样令人痛苦呢？因为如果有人看见自己的恩人饥饿，却没有照顾他，一旦被指责，他宁愿钻进地底，也不愿在两三个朋友面前听到这件事；那么，当我们当着全世界听到这些话时，我们会有何感受？即使到了那时，若非他认真为自己辩护，他也不会说出这些。因为他提出这些并非为了责备，而是为了辩护，并表明他说\Quote{离开我}不是无缘无故或随意说的；这从他不可言喻的恩惠中可以明显看出。因为如果他有意责备，他会提出所有这一切，但现在他只提到他所受到的待遇。\par

10. 因此，亲爱的，让我们害怕听到这些话。生命不是儿戏：或者说我们现在的生命是儿戏，但未来的事不是；或许我们的生命不仅是儿戏，甚至更糟。因为它不以欢笑结束，反而给那些不决心严格安排自己道路的人带来极大的损害。请问，盖房子玩的孩子和我们建造华丽房屋的人之间有什么区别？他们假装吃饭和我们享受美味佳肴之间又有什么区别？没有，只是我们做这些事时有受惩罚的危险。如果我们还没有完全意识到所发生事情的虚空，那并不奇怪，因为我们还没有成为大人；但当我们成为大人时，就会知道这一切都是幼稚的。\par

因为当我们长大成人时，那些其他事也被我们嘲笑；但当我们是孩子时，我们认为它们值得焦虑；当我们收集碎陶片和泥巴时，我们并不比那些建造高大城墙的人看低自己。然而它们很快就毁坏倒塌，即使站立时也对我们毫无用处，正如那些华丽房屋一样。因为天国的公民不能接受它们，他也不能忍受住在其中，那拥有天上家乡的人；但当我们用脚踢倒孩子们的东西时，他也用他的崇高精神踢倒那些东西。当我们为它们的倒塌哭泣时，他们不仅嘲笑，也哭泣：因为他们的心肠是怜悯的，由此产生的祸害巨大。\par

因此，让我们成为有智慧的人。我们要在地上爬行到何时，以石头和木块自夸？我们要嬉戏到何时？但愿我们只是嬉戏！但现在我们甚至连自己的救恩也出卖了；正如儿童在荒废学业时，在空闲时以这些事自娱，却遭受极其严厉的鞭打；同样，我们也将全部精力耗费在这些事上，然后当我们在行为上被要求交出属灵的功课，却无法交出时，就必须承受最重的惩罚。而且没有人能救我们；无论是父亲、兄弟，或任何你所想的。但当这一切都要过去时，随之而来的折磨却是永恒不息的；这正如儿童的事一样，父亲因他们的懒惰而完全毁掉他们幼稚的玩具，使他们不停地哭泣。\par

11. 为使你相信这些事确实如此，让我们把财富——那似乎比任何事都更值得我们劳苦的——摆在我们面前，并与灵魂的一项德行（随便你选哪一项）相对照，然后你就能最清楚地看到它的卑贱。让我们假设有两个人（我现在不说伤害他人的事，只说诚实的财富）；在这两人中，让一个人积聚钱财，航海经商，耕种田地，并寻找许多其他经商之道（虽然我不完全知道，他这样做是否能获得诚实的收益）；不过，就让我们这样假设，并允许他的收益是诚实的；他购买田地、奴隶以及所有这类东西，并假设其中没有不义。但让另一个人，拥有同样多的财产，却出售田地、房屋、金银器皿，并施舍给穷人；让他供给贫困者，医治病人，释放受困者，有些人他从捆绑中释放，另一些人他从矿坑中解救，这些人他从绞索中带回，那些被掳掠的人他解救他们脱离刑罚。那么，你愿意站在哪一边？我们还没有说到将来，只说到目前的事。那么，你们愿意站在哪一边？是站在积聚金子的人一边，还是站在消除灾祸的人一边？是站在购买田地的人一边，还是站在使自己成为人类避难港的人一边？是站在身穿许多金子的人一边，还是站在头戴无数祝福之冠的人一边？前者难道不像一位从天降下为改善其余人类的天使吗？而后者却不像一个人，倒像个胡乱积聚一切的小孩子？\par

但如果诚实地赚钱尚且如此荒谬，近乎极端疯狂；那么，当连诚实都不存在时，这样的人怎能不比其他任何人更悲惨呢？我说，如果荒谬如此之大；再加上地狱和天国的丧失，那么无论他活着还是死了，都该为他发出多么巨大的哀号呢？\par

12. 或者，你愿意我们着手探讨另一项德行？那么，让我们再引入另一个人，他握有大权，命令众人，身负崇高尊严，有华丽的前驱、腰带、侍从官和大批随从。这难道不显得伟大，并配称为有福吗？那么，让我们再拿另一个忍受冤屈、温良、谦卑、忍耐的人来对比；让这个人受凌辱、被殴打，而他安静地承受，并祝福做这些事的人。\par

请问，哪一个人更值得钦佩？是那个自高自大、怒气冲冲的人，还是那个自我克制的人？前者岂不又像那毫无情欲的天上者，后者却像个鼓胀的膀胱，或一个患水肿、严重发炎的人？前者像属灵的医生，后者像个可笑的孩子在鼓着腮帮子？\par

因为，人啊，你为什么自夸？因为你高高在上驾着车？因为一对骡子在拉你？这算什么呢？这连木料和石头都能发生。是因为你穿着美丽的衣服吗？但请看那位以德行为衣的人，你就会看到自己如同枯萎的干草，而他像一棵结出奇妙果实的树，给观看者带来许多喜悦。因为你身上带着喂养虫子和蛀虫的食物，它们若攻击你，便会迅速剥去你这装饰（因为衣服、金银，一是虫子的织品，一是泥土和灰尘，最终又成为泥土，一无所剩）；但那以德行为衣的人，他的衣服不仅虫子不能损伤，连死亡本身也不能。这很自然；因为灵魂的这些德行并非源于大地，而是圣神的果实；因此它们也不受虫子的侵蚀。这些衣服是在天堂织成的，那里没有蛀虫、虫子，也没有任何这类东西。\par

那么，告诉我，哪个更好？是富有还是贫穷？是掌权还是受辱？是奢侈还是饥饿？这很清楚：是尊荣、享受和财富。因此，如果你想要事物本身而不是名号，就离开大地和这里的事物，在天堂为自己找个锚地：因为这里的事是影子，而那里的一切都坚定不移，牢不可破，不受任何攻击。\par

因此，让我们以勤奋和谨慎选择它们，好使我们摆脱世间事物的纷扰，航入那平静的港湾，带着丰富的货物和那无可言喻的施舍财富；愿天主恩赐我们都能达到这一切，借着我们的主\Person{耶稣基督}的恩宠和祂对人的慈爱，愿光荣与权能归于祂，至于无穷之世。阿们。\par

\SourceRecord{Translated by George Prevost and revised by M.B. Riddle.；Nicene and Post-Nicene Fathers, First Series；Vol. 10.；Edited by Philip Schaff.；Buffalo, NY: Christian Literature Publishing Co.,；1888.；<http://www.newadvent.org/fathers/200123.htm>.}

% structure: p0001 source=h1 target=section reason=page-title

\section{玛窦福音讲道第二十四篇}

\BibleRef{玛 7:21}\par

\Quote{不是凡向我说\InnerQuote{主啊，主啊}的人，就能进入天国；惟独遵行我在天之父的旨意的人才能进入。}\par

为何祂不说：\Quote{惟独遵行我的旨意？}因为，为当时的人，能接受这第一点已是一大收获；是的，鉴于他们的软弱，这已是极大的收获。而且，祂也借此暗示了另一层意思。此外还可以说，事实上圣子的旨意与圣父的旨意并无不同。\par

在此，祂似乎主要是责备犹太人，因为他们把全部重心放在教义上，而不关心实践。为此，保禄也曾责备他们说：\BibleQuote{看，你称为犹太人，又依靠法律，且以天主自夸，并知道祂的旨意。}那么，\Quote{这些人是哪些？}你问。许多相信了的人接受了恩赐，例如那个赶鬼却不跟随祂的人；又如犹达斯，他虽邪恶，也得了恩赐。在旧约中也可以找到这样的例子：恩宠常常运行在不配的人身上，为的是有益于他人。也就是说，既然并非所有人都适宜做所有事，有些人生活纯洁却没有那么大的信德，另一些人则恰恰相反，那么，藉着这些话，祂一方面敦促前者彰显大信德，另一方面也以这不可言喻的恩赐召唤后者成为更好的人。因此，祂极其丰盛地赐予了那恩宠。因为经上说：\BibleQuote{我们行了许多异能。}但\Quote{那时我要向他们声明：我从不认识你们。}因为\Quote{如今他们自以为是我的朋友；但那时他们将会知道，我赐予他们恩赐并非因为他们是朋友。}\par

如果祂将恩赐赐给了那些相信祂却没有相称生活的人，这又何足为奇呢？因为即使对那些从信德和善行中都堕落的人，祂也显然照样运作。例如，巴郎既远离信德，也远离真正善良的生活，但恩宠仍为服侍他人而运行在他身上。法郎也属同类，但祂仍向他启示了未来的事。拿步高充满不义，但祂仍向他启示了多代以后的事。\EditorialNote{达尼尔书第三章}又对他儿子，虽比他更不义，祂也启示了未来的事，安排了一项奇妙而伟大的计划。\EditorialNote{达尼尔书第五章}因此，当时福音的初生时期，需要彰显其大能，所以许多不配的人也领受了恩赐。然而，那些奇迹对他们毫无益处，反倒使他们受更重的惩罚。所以祂对他们说出了那可怕的话：\BibleQuote{我从不认识你们。}有许多人，祂的忿怒甚至在此世就已开始；甚至在审判之前，祂就已弃绝他们。\par

为此，亲爱的诸位，我们要恐惧，要十分谨慎地对待我们的生活。不要因为我们现在不行奇迹，就以为自己更差。因为行奇迹对我们永远没有任何益处，同样，若我们注重一切德行，不行奇迹也毫无损害。因为行奇迹时，我们自己是欠债者；但关乎生活和行为，我们却有天主作我们的债务人。\par

看，现在祂已讲完了一切，精确地论述了所有德行，并指出了各种伪善者：那些为了炫耀而禁食、祈祷的人；那些披着羊皮来的；还有那些败坏德行的人，祂称他们为猪和狗。然后祂继续说明，德行在此世有多么大的益处，邪恶有多么大的祸害，祂说：\par

\Quote{所以，凡听了我这些话而实行的，就好比一个聪明人。}\par

意思是：那些不行这些话的人（即使行奇迹）将受什么，你们已经听见了；但你们也当知道，遵守这一切话的人将享受什么，不仅是在来世，就是在此世也是如此。因为祂说：\Quote{凡听了我这些话而实行的，就好比一个聪明人。}\par

你看见祂如何改变话题吗？祂时而说：\BibleQuote{不是凡向我说\InnerQuote{主啊，主啊}的，}并启示了自己；时而说：\BibleQuote{惟独遵行我父旨意的；}又将自己设为审判者：\BibleQuote{因为到那天，有许多人要向我说：主啊，主啊，我们不是因你的名说过预言吗？而我要说：我不认识你们。}在此，祂又表明自己对万有权柄，因此祂说：\Quote{凡听了我这些话。}\par

因此，虽然祂的一切论述都涉及将来——关于天国、不可言喻的赏报和安慰等等——但祂也愿意从现世的事物中给予他们果实，并表明德行的力量在今生何其强大。那么，这力量是什么呢？就是生活在平安中，不被任何恐怖轻易征服，超越一切虐待我们的人。有什么能与此相比呢？因为这一点，连戴王冠的人也做不到，只有追随德行的人才能做到。因为只有他完全拥有它：在现世事物的涨落中，他享有极大的平静。真正令人惊奇的是，不是在好天气中，而是在暴风雨猛烈、动荡巨大、诱惑不断之时，他丝毫不会动摇。\par

\BibleQuote{因为雨淋了，}祂说，\BibleQuote{水冲了，风吹了，撞着那房子，它却不倒塌，因为建基在磐石上。}\BibleRef{玛 7:25}\par

这里，祂用\Quote{雨}、\Quote{洪水}和\Quote{风}来比喻降临于人身上的灾难与痛苦，例如诬告、阴谋、丧亲、死亡、失去朋友、来自陌生人的烦恼——所有人在生活中能想到的不幸。\BibleQuote{但这样的人心，}祂说，\BibleQuote{绝不向其中任何一个屈服；原因在于它建立在磐石上。}祂将自己的教义的坚定性称为磐石，因为事实上，祂的诫命比任何磐石都坚固，使人超乎一切人间风浪之上。因为严格恪守这些诫命的人，不仅会在人们欺压他们时占上风，甚至在魔鬼算计他们时也能如此。这样说并非虚妄的自夸，\Person{约伯}就是我们的见证，他承受了\Person{魔鬼}所有的攻击，屹立不动；\Person{宗徒}也是我们的见证，因为当整个世界的波浪向他们袭来，当各国和君王、本国人和外邦人、恶灵和恶魔、以及一切器械都发动起来时，他们比磐石更坚定，将这一切都驱散了。\par

如今，还有什么比这种生活更幸福的呢？因为能保障这种生活的，不是财富，不是体力，不是荣耀，不是权势，也不是其他任何东西，而唯独是德行的持守。因为我们找不到、也找不到另一种远离一切邪恶的生活，唯独这一种。你们就是见证，你们知道王宫中的阴谋、富人家中的骚乱与烦恼。但在\Person{宗徒}们中间却没有这样的事。\par

那么怎样？难道这样的事没有发生在他们身上吗？难道他们没有从任何人手中受过苦吗？不，最奇妙的正在于此：他们确实成为许多阴谋的对象，许多风暴向他们袭来，但他们的灵魂并没有被这些所动摇，也没有陷入绝望，反而赤身裸体地搏斗、得胜、凯旋。\par

那么，同样地，如果你愿意严格执行这些事，你将会耻笑一切灾祸。是的，因为只要你以这些劝诫中所含的这样一种哲学得以坚固，就没有什么能伤害你。因为一个存心设谋伤害你的人，能在何事上害你呢？他会夺走你的钱财吗？然而，早在他们威胁之前，你就已被命令轻视钱财，并且如此彻底地戒绝它，以至于甚至不向你的主求告这样的事。他会把你投入监牢吗？为什么，在被监禁之前，你已被命令要如此生活，以至于甚至对世界被钉在十字架上。他出言毁谤吗？不，\Person{基督}甚至已将你从这种痛苦中解救出来，祂应许你在忍耐恶事之后不费力地得到大赏报，并使你如此远离由此产生的愤怒和烦恼，以至于甚至命令你为他们祈祷。他流放你，使你陷入无数灾祸吗？他是在为你使荣冠更辉煌。他毁灭你、杀害你吗？即使如此，他也大大地有益于你：为你获得殉道者的赏报，更快地引领你进入平静的港湾，为你提供更丰厚报偿的素材，并设法使你从普世的惩罚中获利。这一切中最奇妙的是，那些设谋者非但毫无伤害，反而使他们所蔑视的对象更加得到赞许。还有什么能与这样的生活方式相比呢？我是说，选择这样的生活方式，而非其他。\par

因此，既然祂曾称那条路是窄而狭的；为了在这一方面也安慰我们的劳苦，祂表明这条路的安全是多么大，快乐是多么大；正如相反的道路有多大的不稳与损害。因为，正如祂从今世之事已表明德行有其赏报，同样罪恶也有其惩罚。因为我常说的，现在也要说：祂在两方面处处促成听众的救恩：一方面借着对德行的热忱，另一方面借着对罪恶的憎恨。因此，因为会有人钦佩祂所说的话，却不用行为证明它，祂就预先唤醒他们的恐惧，说：虽然所说的是好的，但只听是不够得到安全的，还需要在行动上服从，并且全部的关键在于此。祂在这里结束祂的讲论，将恐惧留在他们心中高涨。\par

因为关于德行，祂不仅以将来的事敦促他们（谈论天国、天堂、说不尽的赏报、安慰和无数福乐），也以眼前的事，指出磐石的坚固与不动摇的性质；同样，关于邪恶，祂不仅以期待的事（如树被砍下、不灭的火、不能进入天国，以及祂所说的\BibleQuote{我不认识你们}）激起他们的恐惧，也以眼前的事——我指的是论到房屋时所說的倒塌。\par

4. 因此，祂也借助比喻来试验其力量，使祂的论证更具表现力；因为说\Quote{有德之人将坚不可摧，而恶人则易于制服}与设想磐石、房屋、河流、雨水、风等等，是不同的。\par

祂说：\BibleQuote{凡听了我这些话而不实行的，就好像一个愚昧人，把房屋盖在沙土上。}\par

祂称这人为\Quote{愚昧}，实在有理：因为还有什么比一个人把房屋盖在沙土上，既承担了劳苦，又剥夺了自己的果实和安息，反而受罚更愚蠢的呢？因为那些追随邪恶的人也劳苦，这是人人都明白的：勒索者、奸淫者、诬告者，都辛劳疲倦，以成就他们的恶事；但他们从这些劳苦中非但得不到任何益处，反而遭受大损失。因为\Person{保禄}也暗示了这点，他说：\BibleQuote{那撒种在自己肉身上的，必将由肉身收获败坏。}\BibleRef{迦 6:8}那些在沙土上建造的人，也与这人类似，如同那些沉溺于淫乱、放纵、醉酒、愤怒及一切其他事物的人。\par

阿哈布就是这样的人，但厄里亚却不是这样（因为我们把德行与恶习并列时，会更准确地知道区别）：因为一个人建在磐石上，另一个人建在沙土上；因此，尽管他是国王，却害怕并颤抖于先知，那个只有羊皮衣的人。犹太人就是这样，但宗徒们却不是；所以尽管他们人数少且被捆绑，却显出了磐石的坚定；而那些人，虽然人数众多且全副武装，却显出了沙土的软弱。因为他们说：\BibleQuote{我们怎样处置这些人？}\BibleRef{宗 4:16} 你看那些困惑的人，不是那些在别人手中被捆绑的人，而是那些积极压制和捆绑别人的人？还有什么比这更奇怪呢？你抓住了别人，却还完全困惑？是的，而且很自然。因为他们把一切都建在沙土上，所以也比所有人都软弱。为此他们又说：\BibleQuote{你们为什么想要把这人的血归到我们身上？}\BibleRef{宗 5:28} 他说什么？你鞭打别人，却害怕？你凌辱别人，却惊慌？你审判别人，却颤抖？邪恶如此软弱。\par

但宗徒们却不是这样，而是怎样？\BibleQuote{我们不能不说我们所见所闻的事。}\BibleRef{宗 4:20} 你看到高贵的灵性了吗？你看到磐石嘲笑波浪了吗？你看到不动摇的房子了吗？更奇妙的是，他们不仅没有因针对他们的阴谋而变得胆怯，反而更加勇敢，使别人陷入更大的焦虑。因为打击金刚石的人，自己反被击伤；踢刺棒的人，自己反被刺伤，受重伤的正是他自己；对付义人的人，自己反处于危险中。因为邪恶越是对抗德行，就越变得软弱。正如用衣服包火的人，不能熄灭火焰，反而烧毁了衣服；同样，凌辱义人、压迫他们、捆绑他们的人，使他们更加光荣，却毁灭了自己。因为你越是正义地生活而受苦，你就变得越强；因为我们越是尊重克己，就越不需要什么；越不需要什么，就越强壮，越超乎一切。若翰就是这样的人；因此没有人能使他痛苦，他却使黑落德痛苦；所以一无所有的人胜过掌权者；戴王冠、穿紫袍、拥有无尽排场的人，却在赤身露体的人面前颤抖害怕，甚至被斩首后也未能无惧地看到他的头。因为即使在死后，他仍对他充满恐惧，听他说：\Quote{这是若翰，我所斩首的。} 而\Quote{我所斩首的}这句话，并不是得意之言，而是平息自己的恐惧，说服自己不安的灵魂回想：是他自己杀了他。德行如此强大，甚至死后比活人更有力量。同样，当他活着时，许多富有的人到他那里说：\Quote{我们该做什么？} 你拥有那么多，却想从一无所有的人那里学习致富之道吗？富人向穷人学习？士兵向连房子都没有的人学习？\par

厄里亚也是如此：因此他也以同样的自由向民众讲话。因为前者说：\BibleQuote{毒蛇的种类！}\BibleRef{玛 3:7} 后者说：\Quote{你们两脚踌躇要到几时呢？} 前者说：\Quote{你杀了人，又得了产业吗？} 后者说：\BibleQuote{你不可拥有你兄弟斐理伯的妻子。}\BibleRef{谷 6:18}\par

你看到磐石了吗？你看到沙土了吗？它多么容易下沉，多么容易向灾祸屈服？尽管有王权、人数、贵族的支持，它却被推翻了？对于追求它的人，它使他们比所有人都更愚昧。\par

它不仅倒塌了，而且惨重地倒塌：因为祂说：\Quote{那倒塌是大的。} 危险不是小事，而是灵魂的、天堂的失落，以及那些不朽的祝福。或者说，甚至在那种失落之前，没有比追随邪恶更悲惨的生活了：他们生活在不断的沮丧、惊恐、忧虑、焦虑中；一位智者也曾暗示说：\BibleQuote{恶人无人追赶，却逃跑。}\BibleRef{箴 28:1} 这样的人害怕自己的影子，怀疑朋友、敌人、仆人、认识他们的人、不认识他们的人；在受罚之前，就在这里遭受极重的惩罚。基督说这一切，说：\Quote{那倒塌是大的。} 用这个合适的结尾来结束这些美好的诫命，甚至通过当前的事情说服最不信的人逃离邪恶。\par

因为虽然来自未来的论证更为宏大，但这对于约束粗俗之人、使他们远离邪恶更为有力。因此祂也以此结束，使它的益处能常驻在他们心中。\par

因此，意识到这一切，无论是现在还是将来，让我们逃离恶习，效法德行，使我们不致徒劳无益地劳作，而是既能享受这里的平安，又能分享那里的荣耀：愿天主恩赐我们众人获得这荣耀，因我们的主耶稣基督的恩宠和祂对人的爱，愿光荣和权能归于祂，直到永远。阿们。\par

\SourceRecord{Translated by George Prevost and revised by M.B. Riddle.；Nicene and Post-Nicene Fathers, First Series；Vol. 10.；Edited by Philip Schaff.；Buffalo, NY: Christian Literature Publishing Co.,；1888.；<http://www.newadvent.org/fathers/200124.htm>.}

% structure: p0001 source=h1 target=section reason=page-title

\section{\WorkTitle{玛窦福音释义第二十五讲}}

\BibleRef{玛 7:28}\par

\Quote{耶稣讲完了这些话，众人都惊奇祂的教训。}\par

然而，他们倒更应为祂话语的严苛而忧伤，为祂诫命的崇高而战栗；但如今，这位导师的能力如此伟大，以致许多人甚至被祂的话语所吸引，陷入极大的惊叹，并且因祂话语的甘甜而被说服，甚至在祂停止讲论之后，也再不愿离开祂。因为听众并没有在祂下山后离去，而是全体听众都跟随了祂；祂在他们心中激起了如此大的爱慕之情。\par

但他们最感惊奇的是祂的权威。因为祂说话不像先知和梅瑟那样引用他人，而是处处表明自己拥有定断之权。例如，在颁布法律时，祂不断加上：\Quote{我却对你们说。}在提醒他们那日子时，祂以惩罚和奖赏表明自己是审判者。\par

然而，这也可能使他们困扰。因为如果当他们看到祂以作为显示权威时，经师们就想要用石头砸祂并迫害祂；而现在只有言语证明这一点，他们怎会不冒犯呢？尤其是在开始时说这些话，在祂证明自己的能力之前。不过，他们对此毫无感觉；因为当心灵诚实坦率时，就容易被真理的话语说服。正因如此，有些人即使奇迹宣扬了祂的权能，仍然冒犯；而另一些人仅凭言语就被说服并跟随了祂。我要补充说，福音作者也在暗示这一点，当他说：\Quote{有许多群众跟随了祂，}\BibleRef{玛 8:1} 其中没有统治者或经师，而是那些没有邪恶、判断力未被败坏的人。你看到整部福音中，这样的人紧紧依附祂。因为祂讲话时，他们静静地听，不插话，不打断祂的话语，也不像法利塞人那样试探祂，想找把柄；在祂教训之后，他们又跟随祂，感到惊奇。\par

但是，请你留意主的体贴，祂如何变换使听众受益的方式：在奇迹之后开始讲论，又再次从言语训导转向奇迹。因为在上山之前，祂治愈了许多人，为祂的言论铺路；在向群众讲完那长篇道理之后，祂又回到奇迹，以行动证实所说的话。因此，既然祂教导如同\Quote{有权威的人}，以免祂这样教导显得自夸和傲慢，祂在行动中也同样以权威行事，好使他们在看到祂以同样方式行奇迹时，不再困惑。\par

2. \Quote{耶稣下山后，有一个癞病人前来，说：主，祢若愿意，就能使我洁净。}这个前来的人有极大的理解和信德。因为他没有打断教训，也没有冲进听众中，而是等待适当时机，在\Quote{祂下山后}接近祂。而且不是随意地，而是极其恳切地，在祂膝前恳求，正如另一位福音作者所说，带着真诚的信德和正确的见解。因为他没有说：\Quote{祢若祈求天主}，或\Quote{祢若祈祷}，而是说：\Quote{祢若愿意，就能使我洁净。}他没有说：\Quote{主，洁净我，}而是把一切交给祂，使自己的痊愈依赖于祂，并见证一切权柄属于祂。\par

\Quote{那么，}有人说，\Quote{如果癞病人的看法错了呢？}就应该消除它，责备并改正它。祂这样做了吗？完全没有；相反，祂确认并坚固了所说的话。因此，你看，祂没有说：\Quote{你洁净了，}而是说：\Quote{我愿意，你洁净吧；}好使这道理不再是他个人的猜测，而是祂的认可。\par

但宗徒们却不如此：他们怎样？全体百姓都很惊讶，他们说：\Quote{为什么注视我们，好像凭我们的能力和虔诚使这人行走呢？}然而主虽然常常谦卑地说话，低于祂的荣耀，但在这里祂说了什么来坚固那些因祂的权威而惊讶的人的道理呢？\Quote{我愿意，你洁净吧。}尽管在祂所行的许多伟大奇迹中，祂似乎从未说出这句话。但在这里，为了确认众百姓和癞病人对祂权威的猜测，祂特意加上了\Quote{我愿意。}\par

祂不仅说了这话，也没有不做；事功也立刻随之而来。如果他说得不对，这说法是亵渎，那么事功就应该中断。但现在，自然本身在祂的命令下屈服了，而且迅速，正如所应许的，甚至比福音作者所说的更快。因为\Quote{立刻}一词远不及那事工中实际的速度。\par

但祂不仅说了\Quote{我愿意，你洁净了吧！}而且还\Quote{伸手抚摸了他}；这是一件尤其值得探究的事。因为，当祂以旨意和言语洁净那人时，为何还要加上手的触摸呢？依我看来，别无他意，只是为了藉此表明：祂不受法律管辖，而是超越法律；而且，对洁净的人，从此以后，没有不洁之物。\BibleRef{铎 1:15} 正因如此，我们看到厄里叟连看都没有看纳阿曼，虽然他知道那人对祂不出来触摸自己感到不悦，但厄里叟仍遵守法律的严格规定，待在家中，只打发他去约但河洗浴。而主却触摸了那人，为要表明祂并非以仆人的身份，而是以绝对的主权施行医治。因为祂的手并未因癞病而不洁，反而是那癞病的身体被祂圣洁的手洁净了。\par

因为，正如我们所知，祂来不仅是为医治身体，更是要引导灵魂走向自我克制。因此，从那时起祂便不再禁止不洗手吃饭，并引入了那条高妙的关于食物无分别的诫命；同样，在此情形中，祂也教训我们今后应当注重灵魂——应当摒弃外在的洁净，而洁净灵魂，并唯独畏惧灵魂的癞病，即罪恶（因为成为癞病患者并不阻碍德行）——祂自己首先触摸了那癞病人，却无人指责。因为那法庭并非腐败，旁观者也未受嫉妒支配。因此，他们非但没有责难，反而对奇迹感到惊异，并顺服于此；他们因祂所说的和所做的，都崇拜祂那不可制伏的权能。\par

3. 祂既然医治了他的身体，就吩咐他说：\par

\Quote{不可告诉任何人，但要去让司铎察看，并献上梅瑟所命令的祭品，作为对他们的证据。}\BibleRef{玛 8:4}\par

现在有人认为，祂这样吩咐他不可告诉别人，是为了使他们对辨别祂的治愈毫无诡计可施；这真是他们愚蠢的猜疑。因为祂并非如此医治，以致治愈尚有疑问，而是祂吩咐他\Quote{不可告诉任何人}，教导我们避免夸耀和虚荣。然而祂深知那人不会听从，反而会宣扬他的恩主；但祂仍尽了自己的本分。\par

有人可能会问：\Quote{那么，为什么祂在别处却叫人宣扬这事呢？}这不是因为祂与自己冲突或矛盾，而是为了教导人要知恩。因为在那里，祂也没有命令人宣扬祂自己，而是要人\Quote{归光荣于天主}；藉着这癞病人，祂训练我们远离骄傲和虚荣；藉着另一件事，则训练我们感恩和知恩；并且事事教训我们，要将临到我们的一切赞美都归于主。也就是说，因为人们大多在患病时记念天主，痊愈后却变得松懈；所以祂吩咐他们，无论患病还是健康，都要常常注视主，祂说：\Quote{归光荣于天主。}\par

然而，祂为什么又吩咐他去让司铎察看，并献上祭品呢？这又是为了满全法律。\BibleRef{肋 14:1} 因为祂并非在每一件事上都废除法律，也并非在每一件事上都遵守法律，而是有时这样做，有时那样做；一方面为将要来临的高超生活准则预备道路，另一方面暂时堵住犹太人狂妄的言论，并迁就他们的软弱。若说在起初祂自己就这样做，又何足为奇呢？就连宗徒们，在他们受命往外邦人那里去、普世教导之门已敞开、法律被废止、诫命被更新、一切古旧之事都已止息之后，有时仍可见他们遵守法律，有时却忽略法律。\par

但是，也许有人会说：\Quote{你去让司铎察看}这话对遵守法律有什么贡献呢？这贡献不小。因为古时有一条法律：癞病人得洁净后，不可将洁净的判断交给自己，而要去让司铎察看，将证明呈现在他眼前，并藉着那判决被算作洁净之人。因为如果司铎没有说\Quote{这癞病人洁净了}，那人就仍留在营外与不洁的人在一起。所以祂说：\Quote{你去让司铎察看，并献上梅瑟所命令的祭品。}祂没有说\Quote{我所命令的}，而是暂时将他遣送到法律那里，用尽一切方法堵住他们的口。这样，免得他们说祂夺取了司铎的尊荣；虽然祂自己行了那工作，却将证明之事托付给他们，并让他们作为祂自己奇迹的审判者。祂说：\Quote{我远远不是要与梅瑟或司铎争斗，相反，我引导受我恩惠的人去顺服他们。}\par

但\Quote{作为对他们的证据}是什么意思？是为责备、为证明、为控告，如果他们不知恩的话。因为他们说，我们迫害祂，如同迫害骗子与冒充者，如同迫害天主的仇敌和法律的违犯者；祂说：\Quote{那时你们要为我作证，我并非法律的违犯者。因为，我医治了你之后，就让你去找法律和司铎的认可；}这是尊敬法律、钦佩梅瑟、并且不与古教义对立的行为。\par

即使他们实际上不会因此变得更好，但由此最可以看见祂对法律的尊重：尽管祂预知他们不会受益，祂仍完全履行了祂的本分。因为祂确实预知此事，并预言了它：祂没有说\Quote{为改正他们}，也没有说\Quote{为教导他们}，而是说\Quote{作为对他们的证据}，也就是说，作为控告、谴责和见证，证明我这一方已尽了一切；虽然我预知他们会继续执迷不悟，但我仍没有省略应做的事；只是他们自始至终坚持自己的邪恶。\par

我们可以注意到，祂在别处也这样说：\BibleQuote{这福音必传遍天下，对万民作见证，然后末期才来到。}\BibleRef{玛 24:14} 对万民，对那不服从的，对那不信仰的。如此，免得有人说：\Quote{既然不是所有的人都信，为什么还要向所有人宣讲呢？}——这是为了使我发现自己已尽了自己的一份，往后无人能指责说未曾听见。因为那宣讲本身要作证反对他们，他们往后不能再声称：\Quote{我们没有听见}；因为虔敬之言已\Quote{传到地极}。\par

4. 因此，我们当铭记这些事，也当向近人尽一切本分，并不断向天主感恩。因为我们每日实际上享受祂的恩惠，却连口头上都不承认这恩惠，这实在太荒谬了；何况这承认又能将一切益处归给我们。当然，祂并不需要我们的任何东西，但我们却需要从祂那里得到一切。如此，感恩本身并未为祂增添什么，却使我们更亲近祂。因为如果人的恩惠，当我们记起时，更能以其独特的爱品温暖我们；那么，当我们不断思念我们的主对我们的高贵作为时，岂不更应谨守祂的诫命吗？\par

为此，保禄也说过：\Quote{你们要心怀感恩。}\BibleRef{哥 3:15} 因为任何恩惠最好的保存，就是对此恩惠的纪念和不断的感恩。\par

为此，那在每个共融中举行的、充满伟大救恩的可敬畏的奥迹，被称为感恩的祭献，因为它们是对许多恩惠的纪念，并表明天主对我们眷顾的总和，且以各种方式促使我们感恩。因为若祂由童贞女诞生是一个伟大的奇迹，而福音作者惊异地写道：\Quote{这一切事就这样成就了}；那么，祂被杀害，我们又当如何称呼呢？请告诉我。我的意思是，如果诞生被称为\Quote{这一切事}，那么被钉十字架、倾流祂的血、并将祂自己赐予我们作为精神的盛宴和筵席——那又该称为什么呢？因此，让我们不断感谢祂，并让这感恩先于我们的言语和行动。\par

然而，我们感恩不仅当为自己的祝福，也当为别人的祝福；因为这样我们就能除去嫉妒，并坚固我们的爱德，使之更真诚。因为你若为某人感谢主，便不可能再嫉妒他了。\par

因此，正如你们所知，当那祭献被呈上时，司铎也命令我们为世界、为先前的事、为现在的事、为以前为我们所做的一切、为将来要发生在我们身上的事而感恩。\par

因为这是使我们脱离尘世、升入天堂、并使我们由人变为天使的事。因为天使也组成歌咏团，为天主赐予我们的恩惠而感谢天主，说：\Quote{光荣归于至高之天主，在地上平安，良善之人得蒙喜悦。} \Quote{这与我们有何关系呢？我们既不在世上，也不是人。} \Quote{不，这于我们关系重大，因为我们被教导要如此爱我们的同仆，甚至将他们的祝福视为我们的。}\par

因此，保禄在他的书信中也处处为天主对世界的恩惠行动而感恩。\par

因此，我们也当不断感恩，为自己的祝福，也为别人的祝福，无论大小。因为尽管礼物很小，但因是天主的恩赐而变得伟大；或者更确切地说，凡出自祂的，没有一样是小的，不仅因为是由祂所赐，而且按其本性也是如此。\par

至于其他一切，多如海沙，难以尽述；有什么能比为我们而施行的这救恩计画更伟大呢？祂将比一切更宝贵的，即祂的独生子，赐给了我们这些仇敌；不仅赐下，而且在赐下之后，还将祂置于我们面前作为食物；祂自己为了我们的益处做了这一切，既在赐予祂时，也在使我们为此感恩时。因为人大多忘恩，祂便亲自处处着手并促成我们得益。就如同祂对待犹太人，通过地方、时间和节期提醒他们祂的恩惠，同样，祂在此事上也是如此，通过祭献的方式使我们不断记念祂在这方面的恩赐。\par

没有人能像创造我们的天主那样，如此努力地使我们被悦纳、伟大、并在一切事上心存正直。因此，祂常常违背我们的意愿而善待我们，而且往往是在我们不知不觉中。若你对我所说的感到惊奇，我指出这并非发生在普通人身上，而是发生在真福保禄身上。因为就连那位真福之人，在极大的危险和痛苦中，也曾屡次祈求天主使试探离开他；然而，天主并未顾及他的请求，而是顾及他的益处，为表明这一点，祂说：\Quote{我的恩宠够你用，因为我的德能在软弱中才全显出来。}\BibleRef{格后 12:9} 因此，在祂向他说明原因之前，祂已违背他的意愿、在他不知不觉中施恩于他。\par

5. 如今，祂要求我们以感恩回报如此温柔的眷顾，这算是什么大事呢？那么，让我们服从，并处处坚持这样做。因为犹太人从未被任何事如此摧毁，如同被忘恩负义所摧毁；那接连而来的许多打击，正是由此而来；或者，甚至在这些打击之前，忘恩负义就已摧毁并败坏了他们的灵魂。因为经上说：\Quote{忘恩者的希望，像冬天的霜雪}；\BibleRef{智 16:29} 它使灵魂麻木和死寂，就像霜雪对我们的身体一样。\par

这源于骄傲，以及认为自己配得什么。但痛悔的人会因一切事向天主谢恩，不仅为好事，也为看似逆境的事；他无论遭受多少苦难，都不会认为自己的痛苦是白受的。因此，我们越在德行上进步，就越应使自己痛悔；因为实际上，这比任何事都更是德行。正如我们的视力越敏锐，就越清楚地知道自己离天空有多遥远；同样，我们在德行上越进步，就越明白自己与天主之间的差距。而能认清自己的本分，正是真正智慧中不小的部分。因为那认为自己算不了什么的人，才是最认识自己的人。我们看到，达味和亚巴郎在他们达到德行最高峰时，最完美地履行了这一点；他们一个称自己为\Quote{尘土和灰烬}，\BibleRef{创 18:27}，另一个称自己为\Quote{虫}；所有圣人也如同他们一样，承认自己的卑微。因此，那骄傲自大的人，恰恰是最不认识自己的人。所以在我们的日常用语中，我们也习惯对骄傲的人说：\Quote{他不认识自己}，\Quote{他对自己无知}。而不认识自己的人，又能认识谁呢？因为正如认识自己的人会认识一切，不认识自己的人也不会认识其余的事。\par

那人就是说出\Quote{我要将我的宝座高举于诸天之上}的人，而他竟不认为自己配得上宗徒的头衔，尽管行了那么多伟大的善事。\par

因此，让我们效法他，跟随他。如果我们摆脱尘世和地上的事物，我们就会跟随他。因为没有什么比迷恋世俗事务更使人不认识自己的了；同样，也没有什么比不认识自己更使人迷恋世俗事务的了：这两者是相互依存的。我的意思是，那喜爱外在荣耀、高度评价当下事物的人，如果永远奋斗，就不被允许认识自己；而那轻视这些事物的人，则会轻易认识自己；认识自己之后，他就会依次进入德行的所有其他部分。\par

因此，为了学习这美好的知识，让我们脱离所有在我们心中燃起如此大火焰的易朽事物，认清它们的卑劣，彰显出所有的谦卑和自制：这样我们才能获得现世和未来的福乐，借着我们的主耶稣基督的恩宠和对人的爱，愿光荣、权能和尊荣归于圣父，偕同圣神及善良之灵，从今时直到永远，世世无尽。阿们。\par

\SourceRecord{Translated by George Prevost and revised by M.B. Riddle.；Nicene and Post-Nicene Fathers, First Series；Vol. 10.；Edited by Philip Schaff.；Buffalo, NY: Christian Literature Publishing Co.,；1888.；<http://www.newadvent.org/fathers/200125.htm>.}

% structure: p0001 source=h1 target=section reason=page-title

\section{玛窦福音释义 第26篇}

\BibleRef{玛 8:5}\par

\Quote{耶稣进了葛法翁，有一个百夫长前来恳求祂说：主，我的仆人瘫痪在家里，痛苦难当。}\par

癞病人是在耶稣\Quote{下山后}来到祂面前的；但这位百夫长，是在\Quote{耶稣进了葛法翁}时来的。那么，为什么两人都没有上山呢？并非出于懈怠，因为两人的信德都十分热切，而是为了不打断祂的教训。\par

但他来到祂面前说：\Quote{我的仆人瘫痪在家里，痛苦难当。}有些人说，他提到原因是为自己辩护，为何没有把仆人带来。\Quote{因为那瘫痪者，}他说，\Quote{痛苦万分，奄奄一息，无法抬来。}至于他即将断气，路加说：\Quote{他快要死了。}但我说，这显示他有极大的信德，甚至比那些从屋顶缒下病人的人更大。\BibleRef{路 5:19}\par

那么耶稣做了什么？祂做了以前从未做的事。因为在所有情况下，祂总是顺从求告者的意愿，但在这里，祂反而主动趋向于这意愿，不仅赐予医治，还主动提出要到家里去。祂这样做，是为了让我们认识到百夫长的美德。如果祂没有提出这个邀请，而是说：\Quote{你回去吧，你的仆人必得痊愈；}我们便不会知道这些事了。\par

至少祂在相反的情况下，对待那位腓尼基妇人也是如此。在这里，当没有人邀请祂到家里时，祂主动说要去，好让你认识百夫长的信德和极大的谦卑；但在腓尼基妇人的情况下，祂不仅拒绝了请求，还让坚持祈求的她陷入极大的困惑。\par

因为祂是一位智慧且满有策略的医生，知道如何用相反的事物来达成相反的结果。正如这里通过祂主动提出前往，同样在那里通过祂断然拖延和拒绝，彰显了那位妇人的信德。祂在亚巴郎的事上也是如此，说：\Quote{我决不能向我的仆人亚巴郎隐藏；}好让你知道那人的慈爱和对索多玛的关心。至于罗特的例子，\BibleRef{创 19:2}那些被派遣的人拒绝进入他的房屋，好让你知道那义人慷慨好客的伟大。\par

那么百夫长怎么说？\Quote{我不敢当，祢到我舍下来。}\BibleRef{玛 8:8}让我们倾听，凡是即将接待基督的人：因为现在仍然可能接待祂。让我们倾听、效法，并以同样的热忱接待祂；因为事实上，当你接待一个饥饿赤贫的穷人时，你就接待并照顾了祂。\par

2. \Quote{只要祢说一句话，我的仆人就必痊愈。}\par

请看这人，也如那癞病人一样，对祂持有正确的看法。因为他既没有说\Quote{恳求，}也没有说\Quote{祈祷，恳求，}而是说\Quote{只要命令。}然后出于害怕祂因谦逊而拒绝，他说：\par

\Quote{因为我也在人的权下，也有兵在我以下；我对这个说：去，他就去；对那个说：来，他就来；对我的仆人说：你做这事，他就去做。}\BibleRef{玛 8:9}\par

\Quote{那又如何？}有人说，\Quote{若是百夫长这样揣测呢？问题是基督是否肯定并认可了这一点。}你说得很好，很有道理。让我们来看这件事本身；我们会发现，在癞病人身上发生的事，在这里也同样发生了。因为正如癞病人说\Quote{如果你愿意}（我们不仅从癞病人得知祂的权柄，也从基督的声音得知；因为祂非但没有消除这种揣测，反而以更多的话确认了它，祂说\Quote{我愿意，你洁净了吧，}本来不必多说，但为了确立那人的道理）：同样在这里，也应该看是否发生了这样的事。事实上，我们会发现同样的事再次发生。因为当百夫长说了这些话，并为祂的如此大权柄作了见证后，祂非但没有责备，反而赞同了它，并且比赞同进一步。因为福音作者并没有说祂只是称赞了那话，而是说祂的称赞带有某种热忱，祂甚至\Quote{惊讶；}而且祂不只是简单地惊讶，而是在全体民众面前，并把他作为榜样给其余的人，使他们效法他。\par

你看见吗？每一个为祂的权柄作证的人，都受到\Quote{惊讶？群众都惊奇祂的教训，因为祂教训他们，正像有权威的人；}\BibleRef{玛 7:29}祂非但没有责备他们，反而在下山时带着他们，并用祂对癞病人所说洁净的话，证实了他们的判断。再者，那个癞病人说：\Quote{如果你愿意，你就能洁净我；}\BibleRef{玛 8:2}祂非但没有斥责，反而照他所说的话洁净了他。同样，这位百夫长说：\Quote{只要祢说一句话，我的仆人就必痊愈；}\BibleRef{玛 8:8}祂惊讶地看着他，说：\Quote{在以色列中，我从未见过这样大的信德。}\par

现在，为从反面也说服你这一点：玛尔大没有说这一类话，反而说\Quote{祢无论向天主求什么，天主必赐给祢}\BibleRef{若 11:22}，她不仅没有被称赞，虽然她是祂的熟人，为祂所亲爱，且是曾向祂表现极大热忱的人之一，反而被祂责备和纠正，因为她没有说得好；当祂对她说：\Quote{我不是对你说过：你若信，就必看见天主的荣耀吗？}\BibleRef{若 11:40}责怪她好像她还不相信。再者，因为她曾说：\Quote{祢无论向天主求什么，天主必赐给祢}，为引导她离开这种猜测，并教导她不需要从别人领受，而祂自己就是一切美善的泉源，祂说：\Quote{我是复活，是生命}\BibleRef{若 11:25}，这就是说：\Quote{我不是等待接受能力，而是凭自己行一切事。}\par

因此，对于那位百夫长，祂既惊讶，又将他置于所有人之上，以天国的恩赐尊荣他，并激发其余人有同样的热忱。为了让你知道祂如此说正是为此目的，\Emph{viz.}，为了教导其余的人同样相信，请听福音作者的精确记载，看他如何暗示了这一点。因为，\par

\Quote{耶稣，}他说，\Quote{转过身来，对跟随祂的人说：\InnerQuote{在以色列中，我连这样大的信德也没有见过。}}\par

由此可见，对祂怀有崇高的想象，这尤其是信德的行为，并且有助于获得天国和祂的其他恩赐。因为祂的赞美不仅止于言语，而且祂使那病人痊愈，作为对他信德的酬报，并为他编织荣耀的冠冕，还应许了伟大的恩赐，这样说：\par

\Quote{将有从东从西来的人，与亚巴郎、依撒格、雅各伯一同坐席；但本国的子民，反被驱逐到外边。}\par

这样，既然祂显示了诸多奇迹，祂便更无保留地与他们交谈。\par

然后，为使无人以为祂的话是出于奉承，而使众人知道百夫长的心意正是如此，祂说：\par

\Quote{你去吧！照你所信的，给你成就。}\BibleRef{玛 8:13}\par

果不其然，事工立刻随之而来，证明了他的品格。毫无问题；因为问题是，他们中的每一位是否都将此人的热忱以及他对基督的正确看法摆在我们面前。但很可能，在派出他的朋友之后，他本人也来了并说了这些话。如果路加没有提到一件事，同样玛窦也没有提到另一件；这不是他们彼此矛盾，而是互相补充对方遗漏之处。但请从另一件事上看路加如何宣扬了他的信德，说他的仆人\Quote{快要死了}\BibleRef{路 7:2}。然而，这并没有使他沮丧，也没有使他放弃；他仍然相信他能成功。如果玛窦肯定基督说了\Quote{在以色列中，我连这样大的信德也没有见过}，从而清楚地表明他不是以色列人；而路加却说\Quote{他为我们建造了会堂}；这也不是矛盾。因为一个人，即使不是犹太人，也可能建造会堂并热爱这个民族。\par

4. 但愿你，我恳求你，不仅查问他所说的话，还要加上他的身份，然后你就会看到此人的卓越。因为的确，在权位之人的骄傲是大的，即使在患难中他们也不降低姿态。例如，若望福音中记载的那人，他想要把祂拉到自己的家里，并说：\Quote{下来吧，因为我的孩子快要死了}\BibleRef{若 4:49}。但此人不同；他远胜于那人，也胜于那些从房顶缒下床的人。因为他并不寻求祂的身体临在，也没有将病人带到医生跟前；这暗含了对祂不平凡的想象，而是对祂神性尊威的揣测。他说：\Quote{祢只要说一句话。}起初他甚至没有说\Quote{说一句话}，而只描述他的痛苦：因为他极其谦卑，并不期望基督立即同意并询问他的家。因此，当他听到祂说\Quote{我去治好他}时，然后，而不是在此之前，他说\Quote{说一句话}。患难也没有使他慌乱，反而在灾难中他冷静地推理，与其说关注仆人的健康，不如说避免任何不恭敬的行为。\par

然而，并不是他催促，而是基督主动提出：但即便如此，他仍害怕，恐怕被人认为超出自己的本分，招致力不能及的事。你看到他的智慧了吗？注意犹太人的愚昧，他们说：\Quote{这个人值得祢为他做这事}\BibleRef{路 7:4}。因为当他们本应投靠耶稣对人的爱时，反而提出此人的配得；他们甚至不知道以什么理由提出。但他不然，他宣称自己极其不配，不仅不配得到恩惠，甚至不配在家里接待主。因此，即使他说\Quote{我的仆人瘫痪在床}，他也没有加上\Quote{说一句话}，因为害怕自己不配获得恩赐；他只是陈明自己的痛苦。当他看到基督反过来热切时，他也没有冲上前，而是始终保持自己的适当分寸。\par

如果有人会说：\Quote{基督为何不反过来尊荣他呢？}我们要说，祂确实回报了他尊荣，而且是极大的：首先，通过显明他的心，这主要显现在祂没有进入他的家；其次，通过将他引入天国，并让他比整个犹太民族更受喜爱。因为，因为他自称不配甚至接待基督到自己家里，他反而配得王国，并获享亚巴郎所享有的那些福分。\par

\Quote{但有人会说}，一人可说，\Quote{为什么那癞病人没有受到称赞，他显出了比这些更大的事呢？}因为他并没有说\Quote{说一句话}，而说\Quote{只要你愿意}，这远远更大，这正是先知论及圣父所说的：\Quote{祂随意而行}。但他也受到了称赞。因为当祂说：\BibleQuote{献上梅瑟所命令的礼物，作为对他们的证据}（见：\BibleRef{玛 8:4}），祂的意思无非是：\Quote{你要成为他们的控告者，因为你相信了}。此外，一个犹太人相信与一个外邦人相信并不相同。因为百夫长不是犹太人，这从他是百夫长以及经上说\Quote{在以色列我从未见过这样大的信德}可以明显看出。对于一个不属于犹太民族的人能有如此伟大的思想，这是非常了不起的事。因为依我看来，他所想像的并不亚于天军；或者疾病、死亡以及其他一切都服从于祂，如同他的士兵服从于他一样。\par

因此他也说：\BibleQuote{因为我也是一个在权下的人}，这就是说，你是天主，我是人；我在权下，你不在权下。因此，如果我是人且在权下，尚且能行这样多的事；更何况祂，既是天主，又不在权下，岂不更能？这样，他以最强烈的表达想要说服祂，他说這話，不是给出一个相似例子，而是一个远远超越的例子。因为如果我（他说），与我命令的人在荣誉上平等，且在我权下，却因我职位的微小优越能行这样大的事，没有人反驳我，凡我所命令的，都成就了，尽管命令是各样的（\BibleQuote{我对这个说：去，他就去；对另一个说：来，他就来}），见：\BibleRef{玛 8:9}，那么你自身岂不更能？\par

有些人实际上这样读这段经文：\Quote{因为如果我是一个人}，然后插入一个停顿，他们加上\Quote{在我权下拥有士兵}。\par

但请你注意，他如何表明基督能够像战胜奴隶一样战胜死亡，像主人一样命令它。因为藉着说：\Quote{来，他就来}，和\Quote{去，他就去}，他表达了这个意思：\Quote{如果你命令他的结局不要临到他，它就不会来临}。\par

你看见他多么有信心吗？因为那以后要显明给众人的事，这里已经有一个人使之显明了：祂对死亡和生命都有权柄，并且\Quote{下到地狱的门，又领上来}。\EditorialNote{撒慕尔纪上 2:6}。他不仅说到士兵，也说到奴隶；这关涉到更完全的服从。\par

5. 然而，尽管他有如此大的信德，他仍自认不配。但基督表明他配得祂进入他的家，做了更大的事，惊叹他，宣扬他，并赐予他比所求更多的。因为他来确实是为他的仆人求身体的健康，但离开时却得到了天国。你看见这话已经应验了吗？\Quote{寻求天上的国，这一切都要加给你们}。因为由于他显示了大的信德和谦卑的心，祂既赐给他天国，也加给了他健康。\par

祂不只以此尊荣他，还藉此表明在谁被驱逐之后，他被引入。因为从此时起，祂开始向所有人宣示：救恩是因信德，而非因法律的行为。这就是为什么这恩赐不只赐给犹太人，也赐给外邦人，而且给后者胜于前者。因为祂说：\Quote{千万不要以为}这事只发生在这人身上；不，这将要发生在全世界。祂说这话，是在预言外邦人，并给予他们美好的希望。因为事实上，有些人从外邦人的加里肋亚跟随祂。祂说这话，一方面不让外邦人失望，另一方面压下犹太人的骄傲。\par

但为使祂的话不冒犯听众，也不给他们任何把柄；祂没有突出地提出祂关于外邦人的话，而是藉着百夫长的机会；也没有赤裸裸地使用\Quote{外邦人}一词：祂没有说\Quote{许多外邦人}，而是说\BibleQuote{从东从西有许多人来}（见：\BibleRef{玛 8:11}）。这是指出外邦人的说法，但没有太冒犯听众，因为祂的意思隐藏着。\par

祂不仅以这种方式缓和祂教义表面上的新颖，还通过提及\Quote{亚巴郎的怀中}代替\Quote{天国}。因为那个术语对他们并不熟悉；而且提到亚巴郎会给他们更尖锐的刺痛。因此若翰起初也没有提及地狱，而是说了最能使他们伤心的话：\BibleQuote{不要想自己说：我们有亚巴郎为父}（见：\BibleRef{玛 3:9}）。\par

祂还顾念另一点：不使自己在任何意义上显得与旧的政体对立。因为那尊敬圣祖们、以他们的怀抱作为福分产业的人，也远远足以消除这个疑虑。\par

所以，不要有人以为威胁只有一个，因为一方的惩罚和另一方的喜乐都是双重的：对一方，不仅他们失落了，而且是从他们自己的失落；对另一方，不仅他们得到了，而且得到了他们所没有期望的；还有第三点：一方得到了属于另一方的东西。祂称他们为\Quote{天国之子}，天国是为他们预备的：这最能使他们恼火；因为祂曾通过祂的提供和应许指出他们在祂怀中，最后却把他们赶出去。\par

6. 然后，因祂所说的只是单纯的肯定，祂就以奇迹加以证实；正如祂也反过来，通过后来预言的应验，来显明奇迹。因此，那不信当时仆人所获痊愈的人，不妨从今日已应验的预言，也相信那另一件事。因为这样，那预言甚至还未发生，就已借着当时发生的征兆向众人显明了。为此，你看，祂先说出那个预言，然后才在那时，而不是之前，使瘫子起来；祂要以眼前的事使将来可信，以较小的事使更大的事可信。因为善人享受祂的恩惠，恶人受祂的惩罚，并没什么不合理的，而是合理的事，符合法律的精神；但使软弱的人刚强，使死人复活，却是超乎本性的事。\par

但是，对于这伟大而奇妙的工作，百夫长也贡献不小；我们看到，基督也这样声明，说：\Quote{你回去吧，就照你所信的，给你成就了。}你看，仆人的痊愈如何同时宣扬了基督的大能和百夫长的信德，并且成为未来的保证？或者更确切地说，这全是基督大能的宣扬。因为祂不仅彻底治愈了仆人的身体，也亲自借着祂的奇迹，将百夫长的灵魂引领到信仰里。\par

并且，你不仅应当看这一点：一人相信，另一人痊愈，而且还要惊叹何等的迅速。因为福音者也这样声明说：\Quote{他的仆人就在那个时辰痊愈了}；就如关于癞病人，他也说：\Quote{他立刻洁净了。}因为祂不是单靠医治，而是以奇妙的方式和在转瞬之间行这事，来彰显祂的大能。祂不仅以这种方式使我们获益，而且在显奇迹时，祂也一贯趁机开讲关于祂的国度的道理，吸引众人归向它。因为，连那些祂威胁要赶出去的人，祂威胁他们也不是为了赶出他们，而是为了通过这样的恐惧，用祂的话将他们吸引进去。如果连这样他们也不得益，那完全是他们的过错了，正如所有患同样病症的人一样。\par

因为不仅是在犹太人中间可以看到这事发生，也在那些已经相信的人中间。因为犹达斯也曾是天国之子，并有人与门徒们一起对他说：\Quote{你们要坐在十二个宝座上}；\BibleRef{玛 19:28}但他却成了地狱之子；而厄提约丕雅人，虽是蛮族，并且是从东从西的人，却要同亚巴郎、依撒格和雅各伯享受荣冠。这事如今也发生在你我中间。\Quote{因为有许多}，祂说，\Quote{在先的将成为在后的，而在后的将成为在先的。}\BibleRef{玛 19:30}祂说这话，为叫前者不因认为无法回头而松懈，后者也不因稳固而自负。若翰也从起初就预先宣告了这事，他说：\Quote{天主能从这些石头给亚巴郎兴起子孙来。}\BibleRef{玛 3:9}所以，既然这事必然发生，就早已宣告了，好使无人因事件的奇异而惊慌。但若翰说的是可能之事（因为他是最先的），而基督却说是必将发生之事，并用自己的工作提供证明。\par

7. 那么，我们这些站立的人不要自信，而要对自已说：\Quote{凡自以为站得稳的，务要小心，免得跌倒}；\BibleRef{格前 10:12}我们这些跌倒的人也不要失望，而要对自已说：\Quote{跌倒的，岂能不起来吗？}\BibleRef{耶 8:4}因为许多人甚至已经登上了天顶，表现了极度的刻苦，在旷野中隐居，连梦中也没见过女人，只因稍微松懈，便失足跌入罪恶的深渊。而另一些人则从那里升上了天堂，从舞台和乐队转入了天使的纪律，显出了如此伟大的德行，以致驱赶邪魔，并行了其他许多这样的奇迹。这些例子，圣经里比比皆是，我们的生活中也充满了。连嫖客和柔弱之徒也封住了摩尼教徒的口，这些人说邪恶无法改变，把自己列在魔鬼一边，削弱那些愿意认真行事之人的手，颠覆了我们全部的生活。\par

因为那些教导这些事的人，不但伤害人有关未来的事，而且在此世也把一切颠倒了，至少在他们自己方面。因为当一个人认为回归和改过是不可能的，那在邪恶中生活的人，什么时候会尊重德行呢？因为如果现在，既有法律存在，又有刑罚恐吓，还有一般的舆论能挽回普通人，人们期望地狱，应许天国，邪恶受谴责，善行受称赞，却几乎没有人愿意为德行付出艰苦的努力；倘若你把这些都拿走，还有什么能阻止普遍的毁灭与腐败呢？\par

因此，知道魔鬼的诡计，并且知道外邦人的立法者、天主的圣言、自然的理性、众人的普遍意见，甚至野蛮人、斯基泰人、色雷斯人，以及一般所有人，都直接反对这些人和那些力图确立命运教义的人：让我们保持清醒，亲爱的，向这一切告别，沿着窄路行走，既充满信心又怀着畏惧：畏惧是因为两边都是悬崖，信心是因为我们的向导耶稣。让我们清醒警醒地前行。因为只要稍一打盹，他就迅速被冲走。\par

8. 因为我们并不比\Person{达味}更成全，他因一点疏忽便坠入了罪的深渊。但他又迅速站了起来。所以，不要只看他犯了罪，也要看到他洗去了自己的罪。因为天主记载那历史，正是为此：不是要你看见他跌倒，而是要你钦佩他站起来；为教导你，当你跌倒时，应当如何起来。就如医生选出最严重的疾病，写在他们的医书上，并教导在类似情况下的治疗方法；如果人们在较重的病症上操练，就能轻易掌握较轻的；同样，天主也提出了最大的罪，使得那些在小事上犯罪的人，也能借着这些大罪轻易找到治疗的方法：因为如果那些大罪可以得到医治，那么小罪更不待言。\par

让我们来看那有福之人患病的样式和他迅速康复的样式。那么他患病的方式是什么呢？他犯了奸淫和杀人。因为我不怕大声宣扬这些事。既然圣神不以为耻地记录了这全部历史，我们更不应该掩盖它。因此，我不仅宣扬它，还要添加另一个情况。事实上，凡隐藏这些事的人，正是最使他的德行黯然失色。正如那些对与\Person{哥肋雅}的战斗只字不提的人，剥夺了他不小的荣冠；同样，那些匆匆略过这段历史的人也是如此。我的话难道看似矛盾吗？不，稍等一下，然后你们就会知道，我们有理由这样说。因为我夸大罪过，使我的陈述更显奇异，正是为了更充分地提供药物。\par

那么我所添加的是什么呢？这人的德行；它也使罪过更大。因为并非所有人在所有事上都受到同样的审判。经上说：\Quote{有权势的人}（经上说）\Quote{必受严厉的拷问：}并且：\BibleQuote{那知道主的旨意而不实行的，必受重重鞭打}\BibleRef{路 12:47}所以更多的知识是更重惩罚的根据。同样，司铎如果犯了与平民相同的罪，他不会遭受相同的惩罚，而是要受更重得多的惩罚。\par

或许，看到对他的指控被放大，你们战栗害怕，对我感到惊奇，好像我正在走下悬崖。但我对那位义人如此自信，以至于我甚至要更进一步；因为我越加重指控，就越能彰显对\Person{达味}的赞颂。\par

\Quote{此外还有什么呢？}你们要说，\Quote{可以说出来吗？}多得很。因为就如\Person{加音}的事，他所做的不仅是杀人，甚至比许多杀人更坏；因为他杀的不是外人，而是兄弟；而且是受了冤枉、不曾行不义的兄弟；不是在许多杀人者之后，而是首先创始这可怕罪行的人：同样，在这里所犯的也不仅是杀人。因为做这事的人不是普通人，而是一位先知：他杀的不是行不义的人，而是受不义的人；因为那人确实受了致命的冤屈，由于他的妻子被抢走；然而，在此之后他还加上了这事。\par

9. 你们看，我是如何没有姑息那位义人？我是如何毫无保留地提到了他的过犯？然而，我对他辩护的信心如此之大，以至于在他罪孽如此沉重之后，我宁愿摩尼教徒——他们最嘲笑这一切——和在\Person{玛尔基翁}方式上有病的人都在场，好让我彻底堵住他们的嘴。因为他们确实说\Quote{他犯了杀人和奸淫}；但我不仅说这些，还证明了杀人是双重的：首先，由于那受冤屈的人；其次，由于犯罪者的身份。因为对于蒙受圣神眷顾、领受如此大恩、得以如此自由发言、又处在这样的年纪的人来说，去犯那种罪行，与没有这一切的人犯同样的罪行，并非一回事。然而，即使在这方面，这位杰出的人也最值得钦佩：当他落入邪恶的深坑时，他没有沉沦、没有绝望，也没有因松懈而自暴自弃，在受到魔鬼如此致命的一击之后；反而迅速，更确切地说，立刻并强有力地，给了比他所受的更致命的一击。\par

同样的事发生了，好似在战争和战斗中，某个野蛮人将他的长矛刺入首领的心脏，或射箭击中他的肝脏，并在第一次创伤之上又加上了更致命的一击；而遭受这些重创的人，倒在地上，周围血流遍地，却能迅速站起来，然后向伤他的人投掷长矛，并立即将他击倒在地。同样，在这个事例中，你宣称伤口越重，你就越能暗示受伤者的灵魂值得钦佩，因为他竟能在如此沉重的创伤之后，再次站起来，站在战线的最前沿，并击倒那伤他的人。\par

这是多么伟大的事，凡是陷入重罪的人最清楚。因为正直地行走、一路奔跑，固然是慷慨而活跃的灵魂的证明（因为这样的灵魂有美好的希望伴随，鼓励它、振奋它，使它坚强、更热切）；但更伟大的证明是，在获得无数荣冠、如此多的战利品和胜利之后，遭受了极大的损失，还能重新踏上同样的道路。为使我的话说得明白，我将努力向你们提出另一个例子，完全不亚于前一个。\par

请你们想一想，我求你们：假如有位舵手，航行了无数海域，驶过整个大洋，经历那许多风暴、礁石和波浪之后，却在港口的入口处连同满载的货物沉没了，几乎赤身逃出这场可怕的海难；他对于海洋、航行和这样的劳苦，会自然产生怎样的感受呢？这样的人，除非具有很高尚的灵魂，还会愿意再看到海滩、船只或港口吗？我想不会；他反倒会掩面躺卧，终日如同黑夜，躲避一切事物；他宁愿靠乞讨度日，也不愿再动手从事同样的劳作了。\par

然而这位有福之人并非如此；他虽然遭受了这样一场海难，经历了那无数患难与辛劳，却没有掩面停留，而是重新放下船只，张起帆，握起舵，投身于同样的劳苦，并使自己的财富再次更加丰盈。既然站着不动如此可赞，跌倒后不再永远躺卧；重新站起来并做出这样的事迹，这又配得上何等的冠冕呢？\par

诚然，有许多事足以使他绝望：第一，他罪过的深重；第二，这些事不是发生在人生之初——那时我们的希望也更丰沛——而是临近终了。因为一位刚出港就遇难的商人与一位经过无数次交易后触礁的商人，悲伤的程度并不相同。第三，他已经获得了巨大的财富，却遭遇此事。是的，因为到那时他已积攒了不小的货物：例如他青年早期牧羊时的事迹；关于\Person{哥肋雅}的事迹，他立下了光荣的战利品；以及有关他对\Person{撒乌耳}的自制。因为他甚至显出了福音般的忍耐，将他的敌人无数次交在他手中，却不断宽恕他；他宁愿成为本国、自由乃至生命的流放者，也不愿杀死那对他图谋不轨的人。同样，在他登基为王之后，他也不乏高尚的作为。\par

此外，除了我所说的，他在众人中的声望，以及从如此辉煌的荣耀中坠落，也会引起非同寻常的困惑。因为紫袍远未像他罪恶的污点那样玷污他。你们当然知道，恶行被揭露是何等大事，这样的人需要何等伟大的灵魂，才能在众议的谴责之下、在有许多自己过犯的见证人面前，不灰心丧志。\par

尽管如此，这位高贵的人却从自己灵魂中拔出了所有这些箭矢，此后如此光芒四射，如此抹去了污点，变得如此纯洁，以至于在他死后，他的后代还因他而罪过得减轻；我们发现天主也论及这人说了同样的话；或者不如说，对后者说的更多。因为关于圣祖，经上说：\Quote{我记起了我与亚巴郎所立的盟约；}\BibleRef{出 2:24}但这里天主不是说的\Quote{盟约}，而是怎样说的？\Quote{我要因我的仆人\Person{达味}的缘故保护这城。}\BibleRef{依 37:35}此外，因天主对他的宠爱，天主不容\Person{撒罗满}失去王国，尽管他犯的罪很大。这人的荣耀如此之大，以至于多年后，\Person{伯多禄}在劝勉犹太人时这样说：\Quote{诸位同人！容我坦白地对你们讲论圣祖\Person{达味}，他死了，也埋葬了。}\BibleRef{宗 2:29}基督也对犹太人说话，表明他犯罪之后仍蒙圣神赐予如此大的恩宠，以至于他配得再次预言，甚至关于天主性；藉此堵住他们的口，他说：\Quote{那么，达味怎么在圣神内称祂为主，说：\InnerQuote{主对吾主说：你坐在我右边？}？}而且，正如\Person{梅瑟}所遭遇的，\Person{达味}也遭遇了类似的事。因为\Person{米黎盎}，即使违背\Person{梅瑟}的意愿，也因对兄弟的傲慢而受到天主的惩罚，\BibleRef{户 12:13}因为天主极其爱这位圣人；同样，这人被他的儿子恶待，天主也迅速为他复仇，而且是违背他的意愿。\par

这些事便已足够，是的，比所有其他事都更能表明这人的卓越。因为当天主宣判时，我们不应再追问。但如果你特别想了解他的自制力，你可以通过细读他犯罪后的历史，看出他对天主的信赖、他的仁爱、他在德行上的成长、以及他至死不变的严谨。\par

因此，既然有了这些榜样，让我们保持谨慎，努力不灰心，如果我们有时跌倒，也不要卧倒不起。因为我提及\Person{达味}的罪过，不是要使你们陷入懈怠，而是要在你们心中激起更多的敬畏。因为如果那位义人因一点疏忽便受了这样的伤，我们这些每天都疏忽的人又将遭受什么呢？所以不要看他的跌倒而懈怠，而要想想他此后所做的伟大事迹，何等的哀恸，何等的悔改，他昼夜相加，泪流成泉，以眼泪洗床榻，且身穿苦衣。\par

如果他需要如此巨大的皈依，那么我们在犯下那么多罪后感觉麻木，何时才能得救呢？因为有许多善行的人，甚至轻易就能以此遮盖他的罪过；但未武装的人，无论何处中箭，都会受到致命伤。\par

因此，为使这事不发生，让我们用善行武装自己；如果我们有任何过犯，让我们洗去它：这样我们就能在今生为光荣天主而生活后，配得享受来世；愿我们都能藉着我们的主耶稣基督的恩宠和慈爱，到达那来世；愿光荣和权能归于祂，直到永远。阿们。\par

\SourceRecord{Translated by George Prevost and revised by M.B. Riddle.；Nicene and Post-Nicene Fathers, First Series；Vol. 10.；Edited by Philip Schaff.；Buffalo, NY: Christian Literature Publishing Co.,；1888.；<http://www.newadvent.org/fathers/200126.htm>.}

% structure: p0001 source=h1 target=section reason=page-title

\section{玛窦福音第廿七篇讲道}

\BibleRef{玛 8:14}\par

\BibleQuote{\Emph{耶稣进了伯多禄的家，看见他的岳母躺着发烧:}\Emph{祂摸了她的手，热就退了，她就起来伺候祂。}}\par

但是马尔谷还加上了\Quote{立刻}，意思是要同时宣布时间；而这位福音作者只记载了奇迹，没有说明时间。至于其他福音作者说，那生病的妇人也恳求了祂，这一点他也略过了。但这并非来自任何不一致，而是前者为了简洁，后者为了叙述准确。祂进入伯多禄的家是为了什么意图？依我看，是为了进食。至少当记载时，就表明了这一点。\par

\BibleQuote{她就起来伺候祂。}\BibleRef{玛 8:15}\par

因为祂时常探访祂的门徒（正如玛窦被召时一样），以此尊重他们并使他们更加热诚。\par

但请你注意，伯多禄对祂的尊敬。虽然他岳母在家卧病，且高烧很重，他却未将祂领进屋内，而是先等训导结束，再等所有其他人被治好；然后当祂进入时，才恳求祂。这样，从一开始，他就被教导要将众人的事置于自己的事之上。\par

因此，他自己并没有领祂进来，而是祂自动进入（在百夫长说过\BibleQuote{我不配祢进我的屋}\BibleRef{玛 8:8}之后）：为显示祂赐予门徒多大的恩宠。然而，请想想这些渔夫的房屋是怎样的；尽管如此，祂却不嫌进入他们简陋的小屋，以此教导你处处践踏人类傲慢。\par

有时祂只用言语治愈，有时祂甚至伸手，有时祂两样都做，以显明祂治愈的方式。因为祂不愿总是以更超越的方式行奇迹；祂需要被隐藏一段时间，尤其涉及到祂的门徒时；因为他们出于极大的喜悦会宣扬一切。这一点从以下事实可以明显看出，即使上了山，也需要嘱咐他们不可告诉任何人。\par

因此，祂触摸了她的身体，不仅退烧，还使她恢复完全的健康。如此，既然疾病是普通的，祂便以治愈的方式显示祂的能力；这是任何医术都无法做到的。因为你知道，即使在退烧之后，病人还需要很长时间才能恢复以前的健康。但那时，一切都同时发生了。\par

不仅在此处，在海上的情况也一样。因为祂在那里不仅平息了风和风暴，也立刻止住了浪涛的翻腾；这也是一件奇异的事。因为即使风暴停息，海浪仍会翻腾很长一段时间。\par

但在基督那里并非如此，而是一切同时结束；这妇人也经历了同样的事。因此，福音作者为表明这一点，说\BibleQuote{她就起来伺候祂；}\BibleRef{玛 8:15}这既是基督大能的标记，也是这妇人对基督所怀的态度的标记。\par

此外，我们从此还可观察到另一件事：基督将治疗某些人的恩宠赐予他人的信德。因为在此处，别人替她恳求了祂，正如在百夫长的仆人的例子中一样。祂赐予这恩宠，是在那将要被治愈的人没有不信，或由于疾病不能到祂面前，或由于无知对祂没有伟大概念，或由于年龄幼小的时候。\par

2.\BibleQuote{到了晚上，他们带了许多附魔的人到祂跟前，祂用一句话赶出恶灵，并治好了所有患病的人：这是为应验先知依撒意亚所说的话：祂承担了我们的软弱，背负了我们的疾病。}\par

你看见群众此时在信德上增长吗？因为即使时间紧迫，他们也不愿离开，也不认为在傍晚带病人到祂那里是不合时宜的。\par

但是，我求你注意，福音作者迅速略过了被治愈的庞大群众，没有逐一提到，也没有一一叙述，而是用一句话穿越了不可言喻的奇迹之海。然后，免得这奇事之伟大又使我们陷入不信，即如此多的人和各种各样的疾病竟在祂的一瞬间被解除和治愈，祂又引入了先知来为所发生的事作证：表明我们在每种情况下都有圣经的充分证据，其丰富程度甚至超过奇迹本身；祂说，\BibleQuote{祂承担了我们的软弱，背负了我们的疾病。}祂没有说\Quote{祂消除了}，而是\Quote{祂承担和背负}；在我看来，先知所说的话更是指罪，与若望的话一致，他说：\BibleQuote{看，天主的羔羊，除免世罪者！}\BibleRef{若 1:29}\par

那么，福音作者为何在此将经文应用于疾病呢？或者是以历史意义复述这段经文，或是为了表明我们的多数疾病源于灵魂的罪。因为正如万物的终结——死亡本身，其根源和基础来自罪，那么我们的多数疾病更是如此：因为我们自身受苦的能力本身也起源于那里。\par

3.\BibleQuote{耶稣看见许多群众围着祂，就吩咐渡到对岸去。}\BibleRef{玛 8:18}\par

你再一次看见祂如何不慕虚荣吗？正如其他福音作者说：\Quote{祂嘱咐魔鬼不要说是祂}，同样，这位作者说，祂遣散群众离开祂。祂这样做，一方面训练我们节制，同时消除犹太人的嫉妒，并教导我们做事不为炫耀。因为我们知道，祂不但是身体的医治者，也是灵魂的治疗者，和自我克制的导师；祂借着这两方面显示自己：既除去他们的疾病，又行事不为炫耀。因为他们确实依恋祂，爱慕祂，渴望瞻仰祂。因为谁会离开行这样奇迹的那一位呢？谁不渴望，哪怕只是看看祂的面容，和那说出如此话语的口呢？\par

因为祂不仅行奇迹时令人惊奇，即使只是显现自己时，也充满极大的恩宠；为此先知说：\Quote{你在世人中美丽超群}。如果依撒意亚说：\Quote{祂没有俊美，也没有华丽}，他或是相对于祂天主性的荣耀而言，那荣耀超乎一切言说和描述；或是说明祂受难时所发生的事，以及祂在十字架上所受的侮辱，和祂整个生活中处处所显示的卑微。\par

再者：祂并非先发命\Quote{渡到对岸}，直到医好他们之后。因为那时他们一定无法承受。因此，正如在山上，他们不仅当祂劝勉时与祂同在，而且在祂沉默时也跟随祂；同样在这里，他们不仅在祂行奇迹时服侍祂，而且在祂再次停止后，单单从祂的面容上就得到不小的益处。若梅瑟的脸发光，斯德望的脸像天使的脸，那么，你想想我们的共同主，在这样的时刻，祂的面容该是怎样的呢？\par

也许现在有许多人陷入热切渴望看到那面容的情感；但若我们愿意，我们将看到一个远胜于此的面容。因为我们若能以基督徒的勇气度过今生，我们必将在云中迎接祂，以不朽且不能腐朽的身体与祂相遇。\par

但请注意，祂并非简单地赶走他们，以免伤害他们。因为祂不是说\Quote{离开}，而是\Quote{发命到对岸去}，使他们期待祂一定会来那里。\par

4. 群众方面表现出这么大的爱情，怀着深厚的感情跟随祂；但有一个人，是财富的奴隶，充满极大的傲慢，走近祂说：\par

\Quote{师傅，祢无论往哪里去，我要跟随祢。} \BibleRef{玛 8:19}\par

你看见他的傲慢有多大吗？因为他不屑于与群众为伍，表示他在常人之上，所以他走近了。因为犹太人的性格就是这样，充满不合时宜的自信。同样，后来另一个人，在众人沉默时，自动起来说：\Quote{哪一条是第一条诫命？}\par

然而，主并没有责备他不合时宜的自信，教导我们要容忍这样的人。因此，祂并没有公开指责那些图谋不轨的人，而是回应他们内心的思想，只让他们自己知道被揭露了，从而双重地造福他们：首先，显示祂知道他们良心上的事；其次，在揭露之后，又允许他们隐藏，并让他们若愿意的话，可以重新振作。祂在这人身上也是如此做的。\par

因为他看到许多迹象，许多人被吸引跟随祂，就以为可以从这些奇迹中图利；所以他才急于跟随祂。这从何而知？从基督的回答，祂不是按字面回应问题，而是按其含义所表现的心态。因为祂说：\Quote{什么？你以为跟随我可以积聚财富吗？难道你没有看见我连住处都没有，甚至不如飞鸟吗？}\par

因为祂说：\Quote{狐狸有洞，天空的飞鸟有巢，但人子却没有枕头的地方。} \BibleRef{玛 8:20}\par

这些话不是拒绝的话，而是试探他邪恶的倾向，同时允许他（若愿意在这样的前景下）跟随祂。为证明他的邪恶，当他听了这些话并被考验后，他并没有说：\Quote{我准备好跟随祢。}\par

5. 在其它许多地方，基督也明显这样做；祂没有公开责备，而是借着回答显明前来之人的意图。就如对那说\Quote{善师}的人，他以为用这样的奉承可以赢得祂的欢心，祂就按他的意图回答，说：\Quote{你为什么称我为善？除了一个，就是天主，没有善的。}\par

当他们向祂说：\Quote{看，祢的母亲和祢的兄弟在外面找祢。} \BibleRef{玛 12:47} 因为他们受某种人性的软弱影响，不是想听有益的话，而是想炫耀与祂的亲属关系，并以此自夸；请听祂所说的话：\Quote{谁是我的母亲？谁是我的兄弟？}\par

又对祂的兄弟们，他们向祂说：\Quote{祢把自己显示给世人看罢，} \BibleRef{若 7:4} 并想借此满足他们的虚荣心，祂说：\Quote{你们的时机}（祂这样说）\Quote{常是现成的，但我的时机还没有来到。}\par

祂在相反的事上也如此行，例如对\Person{纳塔乃耳}说：\BibleQuote{看，这确是一个以色列人，心中毫无诡诈。}\BibleRef{若 1:47}又说：\BibleQuote{你们去，把你们所听见、所看见的告诉\Person{若翰}。}\BibleRef{玛 11:4}因为在此祂不是回答她们的话，而是回答那打发他们来的人的意向。祂又同样对民众说话，针对他们的良心说：\BibleQuote{你们出去到旷野里要看什么？}\BibleRef{玛 11:7}这是因为他们大概对\Person{若翰}有些想法，以为他是个轻浮易变的人；为纠正他们的猜疑，祂说：\BibleQuote{你们出去到旷野里要看什么？被风吹动的芦苇？}或者说：\BibleQuote{穿细软衣服的人？}藉这两种比喻说明，他既不是自己动摇的人，也不会被任何奢侈生活软化。因此，在目前的情况下，祂也针对他们的用意作答。\par

请看在此祂也显示出极大的节制：祂没有说：\Quote{我确实拥有，但轻视它}，而是说：\Quote{我没有。}你看到祂的屈尊降贵伴随着何等精确的关怀吗？正如祂在吃喝时，当祂似乎在与若翰相反的方式行事时，祂这样做也是为了犹太人的救恩，或者更确切地说，是为了整个世界的救恩，同时既堵住异端者的口，又渴望更丰盛地赢得当时的人归向祂自己。\par

6. 但是，我们读到，另一个人对祂说：\par

\BibleQuote{主，请许我先去埋葬我的父亲。}\BibleRef{玛 8:21}\par

你看出区别了吗？一个无礼地说：\Quote{无论祢往哪里去，我都要跟随祢；}但另一个，虽然请求一件神圣的义务，却说：\Quote{请允许我。}然而祂没有允许他，却说：\BibleQuote{让死人埋葬他们的死人，你只管跟随我。}因为在每一件事上，祂都关注人的意图。为什么祂没有允许他？有人会问。因为，一方面，有那些会履行那义务的人，死人不会不被埋葬；另一方面，这个人不适合被从更重要的事情中拉走。但祂说\Quote{他们自己的死人}，暗示这人不属于祂的死人。这是因为那死人，至少我认为，是不信者中的。\par

现在，如果你钦佩这个年轻人，为了一件如此必要的事，他恳求耶稣，并没有自行离去；那么你更应当钦佩他，因为他在被禁止时也留下了。\par

那么，有人会说，不参加父亲的葬礼，岂不是极端的忘恩负义吗？如果他这样做是由于疏忽，那确实是忘恩负义，但如果是为了不中断一件更必要的工作，那么他离开则肯定是极端的轻率。因为耶稣禁止他，并非命令我们轻视对父母应有的尊敬，而是表明，对我们而言，没有什么比天上的事更紧迫，我们应当全心全意地持守这些事，即使有极其必要和紧急的事务，也一刻不可拖延。因为有什么比埋葬父亲更必要呢？有什么更容易呢？因为它甚至不会耗费很长时间。\par

但如果连为埋葬父亲所需的时间都不应当花费，连与属神的事务分开这么长时间都不安全；那么想想我们这些终身远离与基督有关的事、把很平常的事置于必要的事之上、并在没有事催促我们时懈怠的人，该当受什么惩罚呢？\par

在此我们也应当钦佩祂教导的教益性：祂将他牢牢钉在祂的话上，并以此使他脱离那些无尽的祸患，如哀哭、悲伤以及随后的事。因为埋葬之后，他必须接着询问遗嘱，然后关于遗产的分割，以及随之而来的一切；如此，一波接一波的浪潮接连涌向他，将他远远冲离真理的港湾。因此，祂吸引他，并把他系在自己身上。\par

但如果你仍然感到惊奇和困惑，为什么他没有被允许参加他父亲的葬礼；那就想一想，许多人不会让病人知道，如果是父亲死，或母亲，或孩子，或任何其他亲属，也不会让他去送葬；我们并不因此指责他们残忍或无情：这是非常合理的。因为，相反地，将处于这种状态的人带出去参加葬礼，才是残忍。\par

但如果为亲属哀恸和心灵痛苦是邪恶的，那么我们从属神的讲论中被抽离更是如此。为此祂也在别处说：\Quote{手扶着犁向后看的人，不适合进天国。}当然，宣扬天国、将别人从死亡中拉回，远远好于埋葬那毫无益处的死尸；尤其是当有人能履行所有这些义务时。\par

7. 因此，我们学到的无他，就是我们决不可随意浪费任何时间，哪怕是最小的瞬间，即使有千万件事压迫我们；而是要把属神的事置于一切之上，甚至是最不可或缺的事，并要认识什么是生命，什么是死亡。因为许多看起来活着的人，其实并不比死人好；因为他们生活在邪恶中；或者不如说，他们比死人更糟；因为经上说：\BibleQuote{已死的人是脱离了罪}\BibleRef{罗 6:7}，但这个人却是罪的奴隶。因为你不要对我说，他没有被虫吃，也不躺在棺材里，也没有闭上眼睛，也没有被裹尸布缠住。不，因为他比死人更痛苦地承受这些事，虽然没有虫子咬他，但他灵魂的情欲比野兽更猛烈地撕裂他。\par

而如果他的眼睛是睁开的，这又比闭着眼睛更糟。因为死人的眼睛看不见任何坏事，但这个人在睁开眼睛时，却为自己聚集了无数的疾病。而另一个躺在棺材里，对任何事都无动于衷，这个人却被埋葬在他无数疾病的坟墓里。\par

但你所见的并非他的尸体已腐烂。那又怎样？因为在他的身体之前，他的灵魂已经败坏、毁灭，并遭受更大的腐朽。因为尸体只臭几天，但此人一生都散发出恶臭，他的口比阴沟更污秽。\par

因此，两者之间的差别恰恰在于：死者确实只经历来自自然的腐朽，但此人连同那自然腐朽，还带来了来自放纵的败坏，每日制造无数腐败的缘由。\par

但他骑马代步吗？那又怎样？死者不也躺在床榻上吗？更严重的是，死者腐朽败坏时无人看见，他有棺材作遮蔽，而此人却带着恶臭到处走动，将死去的灵魂装在身体内，如同在坟墓里。\par

若有人能看见那生活在奢侈和恶习中的人的灵魂，你就会发现，被捆绑在坟墓里远胜于被罪恶的锁链拴住；有石头压在身上，远胜于那厚重的麻木之盖。因此，最重要的是，这些死者的朋友们，既然他们已失去知觉，就当为他们到\Person{耶稣}跟前去，如同当时\Person{玛利亚}为\Person{拉匝禄}所做的那样。即使他\BibleQuote{臭了}，即使他\BibleQuote{死了四天}，也不要绝望，而要上前，先挪开石头。是的，因为你将看见他如同躺在坟墓里，被殓布缠裹着。\par

如果你愿意，就让我们将某个伟大显赫的人带到你面前。不必害怕，因为我会说出例子而不提名字；或者，即使我提了名字，也无需害怕：有谁会害怕一个死人呢？因为无论你做什么，他仍是死的，死人不能伤害活人分毫。\par

那么，让我们看看他们的头被缠住。的确，当他们永远醉酒时，如同死者被许多裹尸布和殓布包裹，他们所有感官都封闭捆绑。你若看他们的手，也会发现它们被绑在肚子上，如同死者的手，但捆绑他们的不是殓布，而是更可怕的贪婪之带：贪婪不允许他们伸手施舍或行其他善事；反而使他们的手比死者的手更无用。你还想看他们的脚被捆在一起吗？请看它们被忧虑缠住，因此永远不能跑到\Place{天主的殿宇}。\par

你见过死人吗？请看那涂香料的人。那么，这些人的涂香料者是谁？是\Person{魔鬼}，他仔细地将他们绑紧，不再让人显出人形，而是一根枯木。因为没有眼睛、没有手、没有脚，也没有其他器官，这样的人怎能显出人形？同样，我们也可以看见他们的灵魂被包裹着，与其说是灵魂，不如说是一个形象。\par

既然他们处于某种麻木状态，变成了死人，就让我们替他们到\Person{耶稣}跟前，恳求祂复活他们，让我们挪开石头，让我们解开殓布。因为如果你挪开石头，即他们对自身苦难的麻木，你就能迅速将他们带出坟墓；将他们带出后，你就能更容易地解下他们的束缚。那时，当你复活、被解开时，\Person{基督}将认识你；然后祂将召叫你赴祂的晚餐。因此，你们凡\Person{基督}的朋友、门徒、爱那死去之人的，都到\Person{耶稣}跟前祈祷吧。因为即使他的恶臭蔓延且极其强烈，我们这些他的朋友也不应抛弃他，反而更要靠近；正如当时\Person{拉匝禄}的姐妹们所做的那样；我们也不应停止代祷、恳求、祈请，直到我们看见他复活。\par

因为如果我们这样安排自己和邻人的事务，我们也将迅速获得永生；愿我们众人靠着我们的主\Person{耶稣基督}的恩宠和爱人，达到永生；愿光荣归于祂，直到永远。阿们。\par

\SourceRecord{Translated by George Prevost and revised by M.B. Riddle.；Nicene and Post-Nicene Fathers, First Series；Vol. 10.；Edited by Philip Schaff.；Buffalo, NY: Christian Literature Publishing Co.,；1888.；<http://www.newadvent.org/fathers/200127.htm>.}

% structure: p0001 source=h1 target=section reason=page-title

\section{论玛窦福音第二十八篇讲道}

\Emph{玛 8:23-24}\par

\Quote{耶稣上了船，祂的门徒跟随着祂。看，海上起了大震荡，以致那船为波浪所掩盖，祂却睡着了。}\par

至于路加（\BibleRef{路 8:22}），为摆脱时间顺序的约束，这样记载：\Quote{有一天，耶稣和门徒上了船；}马尔谷也是如此记载。但这位福音作者并非如此，他在此处也保持顺序。因为并非所有作者都以这种方式记载所有事。我先前已提及这些事，免得有人因省略而以为存在不一致。\par

于是祂遣散了群众，却把门徒留在身边：因为其他福音作者也提及此事。祂带领他们，并非无故，也非偶然，而是为使他们在即将发生的奇迹中成为见证。如同一位卓越的教练，祂为两个目的训练他们：既要在危险中无所畏惧，又要在尊荣中保持谦逊。为使他们在受到祂的保留——却遣散其余人——后不至于高傲，祂容许他们在风暴中颠簸；同时纠正他们，并训练他们勇敢地承受考验。\par

之前的奇迹固然伟大，但这次也包含了一种训练，而且是不小的训练，且与古时的征兆相似。为此祂只带领门徒同行。因为当显示奇迹时，祂也容许群众在场；但当考验和恐惧临到祂时，祂只带领那些将被训练的世界冠军们同行。\par

玛窦只提到祂\Quote{睡着了}，路加却说是在\Quote{枕头上}；这既表示祂的无傲，也教训我们一种高度的刻苦。\par

风暴极其猛烈，海涛汹涌，\Quote{他们唤醒祂说：主啊！救我们，我们丧亡了！}（\BibleRef{玛 8:25}）但祂先斥责他们，然后斥责海。因为如我所说，这些事是为训练而允许的，并且他们是将来要临到的考验的预像。因为此后祂又经常容许他们陷入更严重的命运风暴，并长久容忍他们。因此保禄也说：\Quote{弟兄们！我们不愿意你们不知道，我们在亚细亚所受的磨难：我们被压得超过力量，以致连活命的希望也没有了；}（\BibleRef{格后 1:8}）此后又说：\Quote{祂从这么大的死亡中救出了我们。}因此，为表明他们应当有信心，尽管浪涛汹涌，且祂安排一切为善，祂首先责备他们。因为他们的惊慌本身也是有益的事，使奇迹显得更伟大，并使他们对事件的记忆持久。\par

例如梅瑟首先惧怕蛇，且不仅是惧怕，甚至是极其痛苦；然后他看见那奇怪的事发生（\BibleRef{出 4:3}）。同样，这些人首先感到丧亡，然后得救，为使他们因承认危险，而认识奇迹的伟大。\par

因此祂也睡着了：因为倘若祂醒着发生此事，他们或许不会惧怕，或不会恳求祂，或甚至不会想到祂能行如此的事。因此祂睡着，为给他们的胆怯提供机会，并使他们更清楚地意识到所发生的事。因为人观看他人所经历的事，与观看自己所经历的事，眼光不同。因为他们曾看见别人获得益处，自己却未获得益处，且麻木（因为他们既不瘸腿，也无其他残障）；他们理应通过自己的亲身体验而享受祂的恩惠：于是祂容许风暴，使他们通过得救而更清楚地认识恩惠。\par

因此祂也不在群众面前行此事，免得他们因小信德而被定罪，而是将他们带到一边，纠正他们；并在水的风暴之前，先平息他们灵魂的风暴，斥责他们说：\par

\Quote{为什么你们害怕？你们这小信德的人啊！}如此教训他们：人的恐惧并非因考验的临近，而是因心灵的软弱。\par

若有人说，他们前来唤醒祂并非出于恐惧或小信德，我要说：这件事本身正是他们对祂缺乏正确认识的标志。即他们知道祂被唤醒后能斥责，但尚不知道祂在睡着时也能如此。\par

为何惊奇如今尚且如此，甚至在许多其他奇迹之后，他们的印象仍然不完全？因此他们常受责备，如当祂说：\Quote{你们还是不明白吗？}所以，当门徒处于如此不完全的状态时，群众对祂没有崇高想象，就更不足为奇了。因为\par

他们惊讶说：\Quote{这是怎样的一个人？竟连风和海也听从祂？} \BibleRef{玛 8:27}\par

但基督并未因他们称祂为人而责备他们，而是等待藉着征兆教训他们，使他们知道自己的猜想是错误的。他们为何以为祂是人？首先因祂的外貌，其次因祂睡觉，并使用了船。为此他们陷入困惑，说：\Quote{这是怎样的一个人？}因为睡觉和外貌显明祂是人，而海和风平浪静却宣告祂是天主。\par

因为梅瑟曾行过类似的事，主也在这事上显明自己的超越，并显明一个是作为奴仆行奇事，另一个是作为主行奇事。祂没有像梅瑟那样举起木杖，也没有向天伸手，更不需要任何祈祷，而是像主人命令婢女、或造物主命令受造物那样，仅凭言语和命令就平静并约束了大海；一切狂浪立即止息，没有留下一丝扰乱的痕迹。为此，福音作者记述说：\Quote{就大为平静了。} \BibleRef{玛 8:26} 那曾论及圣父作为大事所说的，祂也借作为显明出来。关于圣父曾说了什么？经上说：\Quote{祂一发言，暴风就止息。} 在此同样，祂一发言，\Quote{就大为平静了。} 为此，群众最惊奇祂；倘若祂像梅瑟那样行事，谁还会惊奇呢？\par

2. 当祂离开大海后，接着发生了另一件更可畏的奇迹。因为附魔的人如同邪恶的逃奴看见主人一样，说道：\par

\Quote{耶稣，天主子，我们与祢有什么相干？祢来这里是要在时期以前折磨我们吗？} \BibleRef{玛 8:29}\par

因为群众称祂为人，魔鬼却来宣告祂的天主性；那些听不见大海汹涌和平静的人，从魔鬼口中听到了同样的呼声，正如大海的平静大声宣告一样。\par

然后，恐怕这事看似出于谄媚，他们便按实际经验呼叫说：\Quote{祢来这里是要在时期以前折磨我们吗？} 为此，他们预先声明敌意，以免他们的恳求引起猜疑。因为他们正无形中受鞭打，大海不如他们那样狂暴；他们被祂的临在折磨、烧灼、忍受难堪的痛苦。因此，无人敢将他们带到祂面前，基督亲自走向他们。\par

玛窦固然记述他们说：\Quote{祢来这里是要在时期以前折磨我们吗？} 但其他福音作者补充说，他们也恳求和哀求主，不要把牠们赶入深渊。因为他们以为惩罚已近，害怕即将落入惩罚中。\par

虽然路加及其随从者说只有一个附魔的人，而这位福音作者说有两个，这并不显示任何不一致。假如他们说只有一个，再没有别人，那么他们似乎与玛窦不一致；但若一说一个，另一说两个，这并非出于不一致，而是叙述方式不同。也就是说，我个人认为，路加是为叙述而挑出其中最凶恶的一个，因此他也以更悲惨的方式报告他们的可怜情形，例如说他挣断锁链和绳索，在荒野中游荡。马尔谷也说，他常用石头砍自己。\par

他们的话语也充分暴露了他们固执和无耻的本性。因为他说：\Quote{祢来这里是要在时期以前折磨我们吗？} 你看，他们不能否认自己犯了罪，但要求不要在时期以前受罚。因为祂当场捉住他们正在实施那些无法治愈和无法无天的暴行，并且以各种方式败坏和惩罚祂的受造物；他们以为祂因他们罪大恶极，不会等到惩罚的时刻；因此他们恳求和哀求祂。那些连铁链也捆不住的，如今被捆绑；那些在山野间奔跑的，来到平原；那些阻挡别人经过的，一见祂挡住道路，便站住了。\par

3. 那么，他们为什么也爱住在坟墓中呢？他们想给群众灌输一种有害的见解，即认为死者的灵魂变成魔鬼，愿天主永不让我们接受这种观念。\Quote{但你要怎么说呢？} 有人可能会问，\Quote{既然许多巫师杀害儿童，为的是后来有灵魂帮助他们？} 这从何而知？因为杀害儿童的事，确有许多人告诉我们；但被杀者的灵魂是否与他们同在，请问你怎么知道？\Quote{附魔者自己，} 有人回答说，\Quote{呼叫说：我是某人的灵魂。} 但这也是一种戏剧表演和魔鬼的欺骗。因为那不是死者的灵魂在呼叫，而是邪灵假装这些事以欺骗听众。因为如果灵魂能进入邪灵的实质，那么更可能进入自己的躯体。\par

此外，受伤害的灵魂与作恶者合作，或人能改变无形体之物的本质，这不合情理。因为在身体上这是不可能的，人不能使人的身体变成驴的身体；在无形的灵魂上更不可能，也不能把它变成邪灵的实质。所以，这些都是昏聩老妇的胡言，是用来吓唬孩子的幻影。\par

确实，灵魂脱离身体后，不可能再在此处游荡。因为\BibleQuote{义人的灵魂在天主手里；} \BibleRef{智 3:1} 如果义人如此，那么这些儿童的灵魂也是如此，因为他们也不是恶人；罪人的灵魂也立即被领离此地。这从拉匝禄和富翁的事可以显明；此外，基督在别处也说：\BibleQuote{今夜就要索回你的灵魂。} \BibleRef{路 12:20} 灵魂离开身体后，不可能在此处游荡，理由也不难看出。因为如果我们行走在熟悉的地上，有身体包围，在陌生的路上行走时，若无人引导就不知该走哪条路；那么灵魂脱离身体后，离开所有习惯的区域，若无人指路，又如何知道往哪里走呢？\par

从许多其他事上也可以看出，一个脱离肉身的灵魂不可能留在此处。因为\Person{斯德望}说：\BibleQuote{主耶稣，接收我的灵魂！}\BibleRef{宗 7:59}；保禄也说：\BibleQuote{我渴望求解脱而与基督同在一起，这实在是再好不过了}\BibleRef{斐 1:23}。关于圣祖，圣经也说：他\BibleQuote{在安享高年后，归到他的祖先那里}。至于罪人的灵魂也不能继续留在此处的证明，请听那富人为此恳求却没有得到，因为如果可能的话，他早就来告诉那里所发生的事了\BibleRef{路 16:27}。由此可知，我们离开此世后，灵魂被领到某个地方，再没有能力自己回来，只等待那可怕的日子。\par

4. 如果有人问：\Quote{基督为什么成全了魔鬼的请求，允许它们进入猪群呢？}那么我们的回答是：祂这样做并非顺从它们，而是借此成就许多目的。第一，要教导那些从这些邪恶暴君手中被解救出来的人，认识他们阴险仇敌的恶意是何等巨大；第二，使众人得知，连猪它们都不敢侵犯，除非祂允许；第三，它们对这些人的迫害会比猪更厉害，除非在灾祸中他们也享受到天主眷顾的极大保护。因为魔鬼恨我们远胜过牲畜，这是人人显然可知的。所以，它们既然不放过猪，在一瞬间把猪全部推下悬崖，那么它们对附身的人更会如此，会把他们引向旷野，带走他们，除非在它们肆虐时，天主的守护之爱足够约束并遏制它们的暴力。由此可见，没有一个人不享受天主眷顾的恩惠。即使并非人人同等，也不是以同一方式，这本身就是眷顾的一个极大实例，因为按各人的益处，眷顾的工作也以不同方式显示出来。\par

此外，从这件事我们还学到另一件事：祂的眷顾不仅普遍地临于众人，也具体地临于每一个人；关于祂的门徒，祂也这样宣布说：\Quote{就是你们的头发也都一 一数过了}\BibleRef{玛 10:30}。从这些附魔者身上，我们也可以清楚看出这一点；他们若不是享有从上而来的许多温柔眷顾，早就被扼死了。\par

因此，祂容许它们进入猪群，也是为了使那些地方的居民认识祂的大能。因为在祂的名声广传的地方，祂没有特别彰显自己；但在没有人认识祂、他们仍处于麻木状态的地方，祂使奇迹发出光辉，以引导他们认识祂的天主性。因为从事件的结果明显可见，那城的居民是某种愚昧无知的人；因为他们本应朝拜祂、惊叹祂的大能，却把祂送走，并且\Quote{求祂离开他们的境界}。\par

但是魔鬼为什么毁灭猪群呢？它们处处致力把人驱入恐慌，处处以毁灭为乐。例如，魔鬼对约伯就是这样做的，虽然在那件事上，天主也容许了，但同样不是因为顺从魔鬼，而是愿意使祂的仆人更加荣耀，断绝那恶灵的无耻借口，并使对义人所做的恶事归到它自己头上。因为如今也发生了与它们愿望相反的事。基督的大能被光荣地宣扬，魔鬼的邪恶——祂从它们手中解救被附的人——得以更清楚地表明；以及它们是如何没有得到万有之天主的许可，连猪也不敢触碰。\par

如果有人愿意以隐微的意义理解这些事，也无不可。因为历史事实的确如此，但我们一定要确知，猪猡般的人尤其容易受魔鬼的运作。只要他们还是遭受这等事的人，就还能常常得胜；但若他们完全变成了猪，那么不但被附，还会被推下悬崖。此外，以免有人以为所发生的事只是演戏，而不明确相信魔鬼已经出来了，猪的死便使这一点显明出来。\par

也要留意祂的温良与祂的大能。当那地方的居民在接受如此恩惠后，却驱赶祂时，祂没有抗拒，而是退去，并离开那些表明自己不配受教导的人，却把那些从魔鬼手中被释放的人以及猪倌留给他们作教师，使他们可以从中得知所发生的一切；而祂自己退去，反使畏惧在他们心中保持活跃。因为那损失的巨大也传播了所行之事的声誉，那事件渗入他们的心。从四面八方传来宣告这奇迹奇异的声音：来自被治愈者、被淹死者、猪的主人、放猪的人。\par

5. 这些事如今也可以看到在发生，甚至有许多在坟墓里的人被恶灵附身，没有任何东西能阻止他们的疯狂：不是铁链，不是锁链，不是众多的人，不是劝告，不是告诫，不是恐怖，不是威胁，也不是任何其他这类事物。\par

因为当一个放荡的人贪求各种拥抱时，他与附魔者毫无区别，也像他一样赤身露体地行走，虽然穿着衣服，却被剥夺了真正的遮盖，失去了应有的荣耀；他不是用石头割伤自己，而是用比许多石头更伤人的罪。那么谁能捆绑这样的人呢？谁能阻止他的不雅和狂乱，阻止他永不醒悟、永远流连于坟墓之间呢？因为那些娼妓的场所正是如此，充满了许多恶臭、许多腐朽。\par

那贪财的人是怎样的？他岂不也像这样？因为谁能捆绑他呢？岂没有恐惧、每日的威胁、劝诫和劝告吗？然而，他挣断所有这些捆绑；若有人来释放他，他反而恳求不要被释放，认为不处于折磨之中才是最残酷的折磨——还有比这更可悲的吗？至于那恶鬼，虽然它藐视人，却顺从基督的命令，迅速从那人身上出来；但这个人连基督的命令也不顺从。请看，他每日听见祂说：\Quote{你不能同时事奉天主而又事奉玛门，}\BibleRef{玛 6:24}并听见祂警告地狱和无法治愈的刑罚，却不听从；这不是因为基督不够强大，而是因为基督不愿强迫我们。因此，这样的人虽住在城中，却如同生活在旷野。因为凡有理性的人，谁愿与这样的人为伍呢？我宁可和一万个附魔者同住，也不愿和一个患有这种病的人同住。\par

我说这话并非错误，从他们各自的情感即可看出。因为后者（贪财者）视那未伤害他们的人为仇敌，甚至想把自由人变成奴隶，并使他陷入无数灾祸；但附魔者并不如此，他们只是自己受疾病的折磨。而贪财者倾覆许多家庭，使天主的名受到亵渎，成为城市和全地的祸害；但被恶灵困扰的人，更值得我们怜悯和流泪。前者（贪财者）大多在麻木中行事，而后者（附魔者）虽在理智中却疯狂，在城中举行他们的狂欢，被某种新的癫狂所支配。因为所有附魔者所做的恶事，有哪一件比得上犹大所行的？他展示了极端的邪恶。所有效法他的人，如同从笼中逃出的猛兽，扰乱他们的城市，无人能制止。因为他们四周也有捆绑：法官的恐惧、法律的威胁、大众的谴责，以及其他种种；但他们甚至连这些也挣断，把一切颠倒混乱。若有人将所有这些捆绑从他们身上除去，他必会清楚地看出他们里面的魔鬼远比刚刚出去的那个更凶猛、更疯狂。\par

但由于这是不可能的，让我们暂且为论证而设想一下：让我们从他身上除去所有锁链，这样我们就能清楚认识他明显的疯狂。但当我们揭露这怪物时，不要害怕；因为这仅是言辞上的描绘，并非事实。那么，让我们设想一个人：眼中冒火，身体漆黑，双肩垂下的不是手，而是蛇；口中装有锋利的剑代替牙齿；舌头是涌出的毒泉和致命的毒药；肚子比任何熔炉更贪焚，吞噬一切投入之物；脚似飞翼，比火焰更猛烈；面容兼有狗和狼的形状；发出的声音不是人声，而是刺耳、不悦、可怕的声音；手中还拿着火把。也许我们所说的这些在你看来已经可怕，但我们还没有完全刻画出他的样子；因为除此之外，我们必须加上其他：他会杀死遇到的人，吞噬他们，咬住他们的肉。\par

然而，贪财者比这更凶猛，他如地狱般攻击所有人，吞没一切，成为人类共同的敌人。他巴不得没有人存在，好让他独揽一切。他甚至不停止于此，当他渴望毁灭所有人之后，他还想破坏大地的实体，使之全都变成金子；不单是大地，还有山丘、森林、泉源，总之，一切可见之物。\par

为使你确信我们尚未完全描绘出他的疯狂，让我们暂且假设没有一个人来控告或恐吓他，并暂时除去法律的恐惧，你就会看见他拿起剑，对所有人施暴，不饶恕任何人：既不饶恕朋友、亲属、兄弟，甚至也不饶恕自己的父母。其实，在这种情况下，甚至无需假设，我们只需问他：他是否不断在心中构想这样的图景，并意念中周旋于所有人之间，要毁灭他们——包括朋友、亲属、兄弟，甚至父母？其实无需询问，因为事实上所有人都知道，被这种疾病支配的人甚至对父亲的年老感到厌倦；而生育子女本是甜美且普遍渴望的，他们却视为痛苦和不受欢迎；许多人为此甚至花钱使自己不育，损伤自己的本性，不仅杀死已出生的子女，甚至不容他们出生。\par

6. 因此，不要惊奇我们如此描绘贪财者（事实上，他比我们所说的还要坏得多）；让我们思考如何将他从魔鬼手中解救出来。我们如何解救他呢？若他能清楚认识到，贪爱金钱正是他在获取金钱这件事上的巨大障碍，因为那些想在细小事上获利的人遭受重大损失；由此产生了一句相关谚语。例如，许多人屡次希望以高利贷放款，因贪图利益而不调查借款人的信誉，结果连本带利损失了全部本金。另一些人陷入危险，不愿放弃微小利益，结果连性命也失去了。\par

再者，当人们能以些许卑劣手段购得有利可图的职位或其他类似之物时，他们却已尽失一切。因为他们不知如何播种，却惯于收割，自然屡屡失收。无人能常收割，正如无人能常获利。既然他们不愿花费，也不知如何获利。倘若他们必须娶妻，同样的事又会临到他们；要么他们受骗，将贫女误作富女娶回；要么他们娶回富妇，却过失无数，此时损失已大于收益。因为使人致富的不是过剩，而是德性。若她挥霍无度、放荡不羁，比任何狂风更猛烈地将一切向外抛撒，她的财富又有何益？若她不贞，引来无数情人，又如何？若她酗酒，岂不很快使丈夫沦为最贫穷的人？但他们不仅以巨大风险娶妻，还因贪婪以同样风险购物，劳心寻找的不是好奴隶，而是廉价的奴隶。\par

请思量这一切（因为关于地狱和天国的话语，你们尚不能听受），并记念你们因贪爱钱财而在借贷、购物、婚姻、权位以及其他一切事上屡屡遭受的损失，从而戒除对钱财的迷恋。\par

这样，你们才能安全地度过今世，并在稍加进步后，得以聆听关于节制的言语，透视并仰望那义德的太阳，获得祂所应许的美善福乐；愿我们都能获此福乐，赖我们的主耶稣基督的恩宠和慈爱，愿光荣与权能归于祂，直到永远。阿们。\par

\SourceRecord{Translated by George Prevost and revised by M.B. Riddle.；Nicene and Post-Nicene Fathers, First Series；Vol. 10.；Edited by Philip Schaff.；Buffalo, NY: Christian Literature Publishing Co.,；1888.；<http://www.newadvent.org/fathers/200128.htm>.}

% structure: p0001 source=h1 target=section reason=page-title

\section{《玛窦福音》释义第二十九篇}

\BibleRef{玛 9:1-2}\par

\BibleQuote{耶稣上了船，渡过去，来到自己的城。看，有人给他送来一个躺在床上的瘫子；耶稣一见他们的信德，就对瘫子说：孩子，放心！你的罪赦了。}\par

祂所说的\Quote{自己的城}此处是指葛法翁。因为生祂的是白冷，养育祂的是纳匝肋，而祂常居住的是葛法翁。\par

然而这个瘫子与\BibleBook{若望}{福音}\BibleRef{若 5:1}所记述的那个不同。因为那个躺在池边，这个在葛法翁；那个患病三十八年，关于这个却没有提及此事；那个无人照顾，这个却有人看顾，他们将他抬起并带来。对这个祂说：\BibleQuote{孩子，你的罪赦了。}对那个祂说：\BibleQuote{你愿意痊愈吗？}\BibleRef{若 5:6}那个祂在安息日治愈，这个却不是安息日，否则犹太人也会以此责难祂；在这个人的事上他们沉默，在那个人的事上他们却不断迫害。\par

我这样说并非无的放矢，免得有人因为怀疑是同一个瘫子而认为有矛盾。\par

但请你注意我们主的谦卑和温良。祂此前也曾遣散群众，而且在加达拉被众人送走时，祂没有抗拒，反而退去，只是没有退到远处。\par

祂又上了船渡过去，其实祂本可以步行过去。因为祂不愿时常行奇迹，免得损害祂人性的道理。\par

\Person{玛窦}确实说\BibleQuote{他们带来他}，但其他福音作者说他们也拆开了屋顶，将他缒下去。他们将病人放在基督面前，一言不发，将一切交托给祂。因为虽然起初祂亲自四处走动，并不要求来者如此大的信德，但在此处，他们既来到祂面前，也需要他们的信德。因为经上说\BibleQuote{一见他们的信德}，即那些缒下病人之人的信德。因为祂并非在所有情况下都只要求病人的信德，例如当他们神志不清，或因疾病而无法自控时。或者，在此处病人也参与了信德，因为他若不信，就不会让人将自己缒下去。\par

既然他们显出了如此大的信德，祂也显出自己的权能，以完全权威赦免他的罪，从各方面表明祂与生祂者同等光荣。请看；祂从起初就藉着教导暗示了这一点，祂教导他们像有权威的人；藉着癞病人，祂说：\BibleQuote{我愿意，你洁净了吧}\BibleRef{玛 8:3}；藉着百夫长，当他说：\BibleQuote{只要你说一句话，我的仆人就会痊愈}，祂称赞他超过众人。\BibleRef{玛 8:8}；藉着海，祂仅用一句话就平息了它；藉着魔鬼，他们承认祂是审判者，祂以极大权威驱逐了他们。\par

在此处，祂又以另一种更大的方式迫使祂的仇敌承认祂的同等光荣，并藉他们的口使之显明。因为祂为表明祂不介意光荣（因有许多围观者堵住了入口，所以他们才从上面缒下病人），并没有立即急于治愈可见的身体，而是从他们身上找到机会；祂先治愈了那不可见的灵魂，赦免了他的罪；这固然救了那人，却未给祂自己带来多大光荣。他们自己反而因恶意困扰，想要攻击祂，结果甚至违背他们的意愿，使所行的事显明出来。事实上，祂以丰富的智慧，利用他们的嫉妒来彰显奇迹。\par

于是，当他们抱怨说：\BibleQuote{这人说了亵渎的话；除了天主一位，谁能赦罪呢？}我们来看看祂说什么。祂是否消除了这怀疑？然而，如果祂不是同等的，祂本应说：\Quote{为什么你们把不合宜的想法加在我身上？我远没有这种权能。}但如今祂一句这样的话也没有说，反而完全肯定并证实了，既藉自己的声音，也藉着所行的奇迹。这样，因为祂说自己的一些事会引起听众反感，祂就藉他人之口来肯定祂要说关于自己的事；而且真正奇妙的是，不单藉朋友，也藉敌人；这正是祂智慧的超卓。一方面藉朋友，当祂说：\BibleQuote{我愿意，你洁净了吧}\BibleRef{玛 8:3}，以及当祂说：\BibleQuote{在以色列我从未见过这么大的信心}\BibleRef{玛 8:10}；但如今藉敌人。因为他们曾说：\BibleQuote{除了天主一位，谁也不能赦罪，}祂接着说：\par

\BibleQuote{但为叫你们知道，人子在地上有权赦罪——于是对瘫子说——起来，拿起你的床，回家去吧。}\par

不仅在此处，而且在另一场合，当他们说：\BibleQuote{我们不是为了善事砸死你，而是为了亵渎的话，并且因为你是人，却把自己当作天主}\BibleRef{若 10:33}，那时祂也没有推翻这看法，反而再次肯定，说：\BibleQuote{如果我不做我父的工作，你们不要信我；但如果我做，纵然你们不信我，也要信这些工作。}\BibleRef{若 10:37}\par

2. 在此处，祂还显明了另一个记号，而且不小，就是祂的天主性和与父同等的光荣。因为他们说：\Quote{赦罪只属于天主，}祂不仅赦罪，而且在此之前还行了另一件只属于天主的事：揭露心中的隐秘。因为他们并未说出他们的想法。\par

因为经上说：\BibleQuote{看，有几个经师}心里说：\BibleQuote{这人说了亵渎的话。}耶稣知道他们的心意，就说：\BibleQuote{你们为什么心里思念恶事？}\BibleRef{玛 9:3}\par

但惟有天主才能知道人的秘密，请听先知的话：\BibleQuote{惟独你知道众人的心}；\BibleRef{编下 6:30}又说：\BibleQuote{天主考验人心和肺腑}；耶肋米亚也说：\BibleQuote{人心最狡猾，已极败坏，谁能识透？}以及\BibleQuote{人看外貌，上主却看人心。}\BibleRef{撒上 16:7}从许多事上可以看出，知道人心内的思念，惟独天主所有。\par

因此，祂暗示自己就是天主，与生祂的天主相等；他们自己在心里所推论的（因为害怕群众，他们不敢说出自己的心意），祂就揭示并显明了他们的这个意见，在此也显示祂极大的温良。\par

祂说：\BibleQuote{你们为什么心里思念恶事？}\BibleRef{玛 9:4}\par

其实，若有人该不悦，就应该是那病人，因为他完全受骗，应该说：\Quote{我来是为得医治，你却更改别的事？凭什么证明我的罪得赦免呢？}\par

但他此时却不说这样的话，而是将自己交托给医治者的权能；但这些人出于好奇和嫉妒，图谋破坏他人的善行。因此，祂责备他们，却十分温良。祂说：\Quote{如果你们不信先前的事，认为我的说话是夸口；看，我再加一件，就是揭露你们的隐秘；此后又加另一件。}这另一件是什么呢？就是使瘫子的身体复原。\par

当祂对瘫子说话时，祂没有清楚显明自己的权柄；因为祂不是说：\Quote{我赦免你的罪}，而是说：\BibleQuote{你的罪赦了}；因他们的逼迫，祂更清楚地显明自己的权柄，说：\BibleQuote{但为叫你们知道，人子在地上有赦罪的权柄——}\par

你看见了吗？祂多么不愿被认为与圣父不相等？因为祂绝没有说：\Quote{人子需要另一位}；或\Quote{祂给了祂权柄}，而是说：\Quote{祂有权柄。}祂说这话也不是为求荣耀，而是\Quote{为使你们相信}，祂说，\Quote{我没有因使自己与天主平等而亵渎。}\par

因此，祂处处愿意提供清楚且无可争议的证明；例如祂说：\BibleQuote{你去，把你自身给司祭看}；\BibleRef{玛 8:4}以及祂指出伯多禄的岳母服侍，并准许猪群冲下山崖。同样在此，首先，作为赦罪的凭证，祂使他的身体复原；随后又叫他拿起床，以防事实被当作幻象。祂在询问他们之前就这样做。因为祂说：\BibleQuote{说\InnerQuote{你的罪赦了}，或说\InnerQuote{拿起你的床，行走}哪一样更容易？}\BibleRef{玛 9:5}祂的话大意是这样：你们看，哪一样更容易：是治愈紊乱的身体，还是除去灵魂的罪？显然，治愈身体更容易。因为灵魂比身体更高贵，所以除罪是比这更大的事；但因为前者不可见，后者可见，我便把较次但更可感觉的事加上，好使那更大的、不可见的事藉此得到证明。这样，祂以工作预先印证了若翰所说的：\BibleQuote{祂除免世罪。}\par

于是，祂使他起来，送他回家；在此又显示祂的不自夸，并证明这事不是幻象；因为祂使同一批人成为他软弱和健康的见证人。祂说：\Quote{我实在愿意藉你的灾祸，也医治那些看似健康、却心灵患病的人；但他们不愿意，你就回家去吧，医治你家里的人。}\par

你看见祂如何显示自己是灵魂和身体的创造者吗？因此，祂治愈了两种实质上的瘫痪，以可见的事使不可见的事显明。但他们仍然匍匐在地。\par

\BibleQuote{群众看见，就惊讶光荣天主，}据说，\BibleQuote{因为赐给了人们这样的权柄}；因为肉体成了他们的绊脚石。但祂没有责备他们，而是以工作唤醒他们，提升他们的思想。因为当时，祂被认为超越众人、来自天主，已是一件大事。如果他们心中确立了这些事，循序渐进，他们就会知道祂甚至就是天主子。但他们没有清楚地持守这些，所以也不能接近祂。因为他们又说：\BibleQuote{这人不是从天主来的}；\BibleRef{若 9:16}\BibleQuote{这人怎会是来自天主的？}他们不断重复这些话，作为自己私欲的掩饰。\par

3. 如今许多人也是这样；他们认为自己在替天主报仇，却满足自己的私欲，而他们本该以节制处理一切。因为万有的天主，虽有力量向亵渎祂的人发出闪电，却仍使太阳升起，降下雨露，并丰厚地赐给他们一切；我们应当效法祂，以温柔相待、劝告、训诫，不发怒，不使自己成为野兽。\par

因为他们的亵渎对天主毫无损害，你何必发怒？那亵渎的人自己反而受了伤。为此你要叹息、哭泣，因为这灾祸确实该流泪。那受伤的人，再没有比温柔更能医治他的：我说温柔，它比任何力量都更强有力。\par

请看，例如，祂自己，那受侮辱的一位，如何与我们交谈，无论在旧约还是新约中；在旧约中说：\Quote{我的百姓啊，我向你们做了什么呢？}\BibleRef{米 6:3} 在新约中说：\Quote{扫禄，扫禄，你为什么迫害我？}\BibleRef{宗 9:4} 保禄也吩咐：\Quote{用温柔规劝那些反对的人。}\BibleRef{弟后 2:25} 基督再次，当祂的门徒来到祂面前，要求从天上降下火来时，严厉地斥责他们说：\Quote{你们不知道你们是出于什么精神。}\par

在这里祂又没有说：\Quote{你们这些可诅咒的、行邪术的；你们这些嫉妒的、敌视人类救恩的人；} 而是说：\Quote{你们为什么心里怀着恶念？}\par

我们必须，你们看，用温柔来根除疾病。因为那因畏惧人而变好的人，会很快再次回到邪恶。为此祂也命令留下莠子，给予指定的悔改日子。的确，他们中有许多人悔改了，从前是坏的，后来变好了；例如，保禄、税吏、盗贼；因为这些真正是莠子的，转变成了好麦子。因为虽然在种子中这不可能，但在人的意志中，这既是可掌控的，也是容易的；因为我们的意志不受任何本性的限制，而是有自由选择作为其特权。\par

因此，当你看到一个真理的敌人时，要等待他，照顾他，通过显露出色的生活，运用\Quote{无可指责的言语}\BibleRef{铎 2:8}，给予关注和温柔的关怀，尝试每一种改正的方法，来引领他回归德行，效法最好的医生。因为他们也不是只用一种方法医治，而是当他们看到伤口对第一种疗法没有反应时，就加上另一种，之后再加上另一种；有时用刀，有时包扎。你也要照样成为灵魂的医生，按照基督的法律运用每一种治疗方式，好使你获得拯救自己和益助他人的报酬，一切为了天主的荣耀，你自己也因此得荣耀。\Quote{因为荣耀我的，}祂说，\Quote{我必荣耀；藐视我的，必被轻视。}\par

我们要，我说，为祂的荣耀做一切事，好使我们能达到那蒙福的份，愿天主恩赐我们都能达到，藉着我们的主耶稣基督的恩宠和慈爱，愿荣耀和权能归于祂，直到永远。阿们。\par

\SourceRecord{Translated by George Prevost and revised by M.B. Riddle.；Nicene and Post-Nicene Fathers, First Series；Vol. 10.；Edited by Philip Schaff.；Buffalo, NY: Christian Literature Publishing Co.,；1888.；<http://www.newadvent.org/fathers/200129.htm>.}

% structure: p0001 source=h1 target=section reason=page-title

\section{玛窦福音讲道集第三十篇}

\Emph{\BibleRef{玛 9:9}}\par

\Quote{\Emph{祂从那里前行，看见一个人坐在税关上}，\Emph{名叫玛窦；祂就对他说：\InnerQuote{跟随我。}}}\par

因为祂行了奇迹以后，并没有停留，免得因在眼前而更激起他们的嫉妒；反而以退隐和缓和他们的情绪来迁就他们。所以我们也应这样做：不要对抗那些陷害我们的人，倒要缓和他们的创伤，让步并消减他们的激烈。\par

但祂为什么没有在召叫伯多禄、若望和其他人时一起召叫他呢？就如祂在召叫他们时，是在知道他们会服从的时刻；同样，祂召叫玛窦时，是在确信他会顺从自己的时候。同样，复活以后，祂抓住了保禄，如同渔夫捕鱼。因为那位洞悉人心、知道人内心隐秘的主，知道每人在何时会服从。因此，祂不是一开始就召叫他，那时他还相当刚硬，而是在祂行了无数奇迹、名声大噪之后，知道他已经实际预备好服从时，才召叫他。\par

我们也有理由钦佩圣史的无私，他不掩饰自己过去的生活，反而写上了自己的名字，而其他圣史是用另一个称呼将他隐藏起来的。\par

但他为什么说他\Quote{坐在税关上}呢？是为了说明召叫他的那一位的权能：他不是在放弃或离开这种邪恶的营生之后，而是在作恶中被祂提拔上来；正如祂也转化了蒙福的保禄，那时他正疯狂暴怒、口吐火焰；保禄自己以此作为祂权能的证明，对迦拉达人说：\Quote{你们听说过我从前在犹太教中的生活，我是怎样极度地迫害天主的教会。} \BibleRef{迦 1:13} 渔夫们也是在他们忙碌时蒙召的。但那是一种不算恶名的行业，只是粗鲁的人，不与别人交往，而且非常纯朴；而眼下这个行当却充满了一切傲慢和胆大妄为，是一种无法公正解释的获利方式，是一种无耻的买卖，是法律掩盖下的抢劫。然而，发出呼召的那一位并不因这些事而羞耻。\par

我说祂不因一个税吏而羞耻，这还算什么？因为即使对一个妓女，祂非但不羞于召叫她，反而允许她亲吻自己的脚，并用眼泪沾湿它们。\BibleRef{路 7:38} 是的，祂正是为此而来，不仅医治身体，也要治愈灵魂的邪恶。祂在瘫子身上也这样做了，清楚表明祂能赦免罪过，然后，不是之前，祂才来到我们此刻所谈论的那个人面前；为的是他们不再因看到一个税吏被选入门徒的行列而烦恼。因为那能消除我们一切过犯的，为什么惊讶祂甚至使这个人成为宗徒呢？\par

正如你看见了召叫者的权能，也看看被召者的服从：他既不抗拒，也不争论说：\Quote{这是什么？祂召叫我这样的人，难道不是一种欺骗？} 不；因为这种谦卑又是不合时宜的；他立刻服从，甚至没有要求回家和亲戚商量这事；渔夫们也没有；正如他们离开网、船和父亲，他也离开税关和收益，跟随了祂，展现出一个准备好一切的心志；并且立即与一切世俗的事物决裂，以完全的服从证实了召叫他的人选择了恰当的时机。\par

有人可能会说，为什么他没有告诉我们其他门徒是如何、以什么方式被召叫的，而只告诉了伯多禄、雅各伯、若望和斐理伯，却没有提到其他人？\par

因为这些人比其他人在生活上更为奇特和卑微。没有比税吏的行当更恶劣的，也没有比打鱼更普通的。而斐理伯也是卑贱中的一员，从他的家乡可以显明。因此，这些圣史特别向我们宣扬这些，连同他们的生活方式，以表明我们应当相信他们历史中光荣的部分。因为那些不放过任何被认为可耻的事情，反而比其余更精确地公布这些事情——无论是关于老师还是门徒——又怎么在他们理应被尊崇的部分被怀疑呢？尤其是许多征兆和奇迹被他们略过，而十字架的事件（被认为是一种耻辱）却仔细而响亮地表达出来；门徒们的追求和过失，以及他们老师祖先中因罪而臭名昭著的人，\BibleRef{玛 3:6} 他们用清晰的声音揭示出来。由此可见，他们极其重视真理，写作没有出于偏袒，也不是为了炫耀。\par

2. 因此，召叫他之后，祂以极大的尊荣加给他，立刻与他一同坐席；因为这样祂既给了他未来的美好希望，又引导他更加自信。因为不是长久的，而是立刻，祂治愈了他的恶习。祂不仅与他，还与许多其他人一同坐席；虽然这件事本身被当作对祂的指控，说祂没有赶走罪人。但他们也不隐瞒这一点，有人试图在祂的行动上定下什么样的罪过。\par

现在税吏们聚集起来，如同到一个同行的那里；因为他因\Person{基督}的到来而欢欣，便把他们全都召了来。事实上，\Person{基督}曾尝试各种方式；不但在讲论时，在医治时，在斥责祂的敌人时，甚至在早餐时，祂也常纠正那些行为不端的人；借此教导我们，每一个时节和每一件事都可能为我们带来益处。然而，当时摆在面前的确实来自不义和贪婪；但\Person{基督}并不拒绝分享，因为随后的收获是巨大的；祂甚至与犯了这些过犯的人同席共餐。因为医生的特质正是如此：他若不能忍受病人的腐坏，就无法使他们脱离病弱。\par

然而，祂无疑因此招致了恶评：首先是因为与那人一同进食，其次是在\Person{玛窦}的家里，第三是与许多税吏在一起。请看他们是怎样以此责备祂的。\BibleQuote{看，一个贪吃好酒的人，税吏和罪人的朋友。}\BibleRef{玛 11:19}\par

让那些力求以斋戒为自己博取巨大荣誉的人听吧，让他们想一想，我们的主曾被称作\Quote{一个贪吃好酒的人}，祂却并不羞愧，反而忽视了这一切，为要完成祂所定下的事；这的确随之成就了。因为那个税吏确实悔改了，并因此成了一个更好的人。\par

并且为教导你们，这件大事是通过与祂共餐而成就的，请听另一个税吏\Person{匝凯}所说的话。我是说，当他听到\Person{基督}说\Quote{今天我必须住在你家里}，喜乐给了他翅膀，他说\Quote{我把我一半的财物施舍给穷人，我若讹诈了谁，就偿还四倍。}\Person{耶稣}就对他说\Quote{今天救恩临到了这家。}所以，通过所有方式都有可能施教。\par

但是有人会说，\Person{保禄}怎么命令说\BibleQuote{若有人称为弟兄，是淫乱的或贪婪的，这样的人，连与他一同吃饭也不可？}\BibleRef{格前 5:11}首先，尚不明确，他是否也将这条命令给了教师们，而不仅仅是给弟兄们。其次，这些人还没有达到成全者的行列，也尚未成为弟兄。此外，\Person{保禄}甚至对那些已经成为弟兄的人，若是他们继续如故，才命令要远离他们；但这些人都已经停止了行恶，并且悔改了。\par

3. 但这些事没有一件使法利塞人羞愧，他们反而向祂的门徒控告祂，说：\par

\BibleQuote{为什么你们的师傅与税吏和罪人一同吃饭？}\BibleRef{玛 9:11}\par

当门徒似乎做错事时，他们就向祂求情，说\BibleQuote{看，你的门徒做了在安息日不合法的事。}\BibleRef{玛 12:2}；但在这里，他们却当着门徒的面毁谤祂。这一切都是出于诡诈之人所为，想要使门徒的队伍离开他们的师傅。那么，无限智慧说了什么？\par

\Quote{健康的人不需要医生，}祂说，\Quote{而是有病的人。}\par

请看祂如何将他们的推理转向相反的结论。也就是说，当他们以此指控祂与这些人在一起时，祂反而说，祂若不与他们同在，反而不配于祂自己，不配于祂对人类的爱；并且纠正这样的人不仅无可指责，而且是卓越的、必要的，值得各种赞美。\par

此后，为免祂所说\Quote{有病的人}这一说法使那些受邀者感到羞愧，请看祂又如何弥补，通过责备其他人并说：\par

\BibleQuote{你们去研究一下这是什么意思：我喜欢仁爱，而非祭献。}\BibleRef{玛 9:13}\par

祂说这话，是为了责备他们对圣经的无知。因此，祂也以更尖刻的方式安排祂的言论，并非出于愤怒，远非如此；而是为使那些税吏不至于全然困惑。\par

当然，祂本可以说，\Quote{你们没有注意到吗？我赦免了瘫痪病人的罪，我使他的身体恢复力气？}但祂并没说这样的话，而是先以人们普遍的推理与他们辩论，然后才引用圣经。因为祂说了\Quote{健康的人不需要医生，而是有病的人}，并隐性指出祂自己就是医生；之后祂说\Quote{你们去研究一下这是什么意思：我喜欢仁爱，而非祭献。}\Person{保禄}也是如此：当他先用普通事物的比喻确立了他的论证，并说了\BibleQuote{谁牧养羊群，而不吃羊奶呢？}\BibleRef{格前 9:7}然后他也引用了圣经，说\Quote{梅瑟法律上记载：牛在场上踹谷的时候，不可笼住它的嘴。}又说\Quote{主也这样规定：传福音的人应靠福音生活。}\par

但对祂的门徒，祂却不这样，而是以祂的奇事提醒他们，这样说：\BibleQuote{你们还不记得那五千人的五个饼，和你们收了多少篮子吗？}\BibleRef{玛 16:9}然而对这些人，祂却不这样，而是提醒他们我们共同的弱点，并指明他们无论如何都属于软弱之列；这些人甚至不认识圣经，却轻视其他德行，把一切重心都放在祭献上；当祂简要地写下所有先知所肯定的话时，祂也殷切地暗示他们这一点，说：\Quote{你们去学习这是什么意思：我喜欢仁爱，而非祭献。}\par

事实上，祂借此表示，不是祂违反法律，而是他们；好像祂说：\Quote{你们为什么指责我？因为我引领罪人悔改吗？那么你们也必须为此指责圣父。}正如祂在其他地方也为确立这一点而说：\BibleQuote{我父一直工作到现在，我也工作。}\BibleRef{若 5:17}所以在这里又说：\Quote{你们去研究一下\InnerQuote{我喜欢仁爱，而非祭献}是什么意思？}因为基督说：\Quote{既然是祂的旨意，也是如此。}你看见哪一个是多余的，哪一个是必要的吗？因为祂并没有说：\InnerQuote{我喜欢仁爱和祭献}，而是说：\Quote{我喜欢仁爱，而非祭献。}也就是说，祂允许了一件，而抛弃了另一件；并证明了他们所指责的，非但不是被禁止的，反而是法律所规定，甚至比祭献更重要；祂还引用了旧约，说出与自己和谐的话语，并制定法律。\par

既用一般的比喻和圣经责备了他们以后，祂又加上：\par

\Quote{我来不是为召义人，而是为召罪人悔改。}\par

祂对他们说这话，含讥讽的意味；正如祂说：\Quote{看，亚当成了我们中的一个；}以及\Quote{如果我饿了，我也不必告诉你。}因为保禄宣布，世上没有义人，说：\BibleQuote{因为众人都犯了罪，失掉了天主的光荣。}\BibleRef{罗 3:23}这也安慰了另一些人，即客人。祂说：\Quote{我远远没有厌恶罪人，甚至只为他们的缘故我才来了。}然后，免得使他们更加疏忽，祂没有停在\Quote{罪人}一词，而是加上了\Quote{悔改。}因为祂说：\Quote{我来不是要他们继续做罪人，而是要他们改变、改善。}\par

4. 祂既从圣经，又从事物的自然结果，多方堵住了他们的口；他们无话可说，被证明他们所加给祂的指控，反而落到了他们自己头上，而且他们成了法律和旧约的反对者；于是他们离开祂，又把控告转移到门徒身上。\par

路加确实肯定说法利塞人说了这话，但这位福音作者却说是若翰的门徒说的；但很可能两者都说了。也就是说，他们可以预料到，在完全困惑中，拉上另一派人；正如后来他们也与黑落德党人一起行动。事实上，若翰的门徒总是容易嫉妒祂，并反对祂；只有当若翰初被囚在狱中时，他们才谦卑下来。至少那时他们来了，\Quote{并告诉了耶稣；}但后来他们又回到先前的嫉妒。\par

他们说什么？\BibleQuote{为什么我们和法利塞人常常守斋，而你的门徒却不守斋？}\BibleRef{玛 9:14}\par

这是基督很久以前就根除的毛病，祂说：\BibleQuote{当你守斋时，要抹油梳头，洗脸；}\BibleRef{玛 6:17}预先知道由此产生的祸害。但即便如此，祂也没有责备他们，也没有说：\Quote{你们这些虚荣好管闲事的人；}而是极其温和地对他们说：\Quote{伴郎岂能当新郎与他们在一起的时候守斋？}如此，当别人（即税吏）被代祷，要抚慰他们受伤的心灵时，祂对责骂者更加严厉；但当他们嘲笑祂自己和祂的门徒时，祂极其温和地作答。\par

他们的意思是这样：他们说：\Quote{就算你这样做是作为医生，为什么你的门徒也放弃守斋，而依附这样的筵席？}然后，为使指控更重，他们先提自己，再提法利塞人；希望借比较加重罪责。因为经上说：\Quote{我们和法利塞人都常常守斋。}事实上他们确实守斋，一是从若翰学来的，一是从法律学来的；正如法利塞人也说：\BibleQuote{我每周守斋两次。}\BibleRef{路 18:12}\par

耶稣怎样回答？\Quote{伴郎岂能当新郎与他们在一起的时候守斋？}先前，祂自称医生，但这里自称新郎；用这些名号启示祂不可言喻的奥秘。当然，祂本可以更严厉地告诉他们：\Quote{这些事不取决于你们，以致你们制定这样的法律。因为当内心充满邪恶，当你们指责别人、定别人的罪，眼睛里有大梁，并且一切为了炫耀时，守斋有什么用？不仅如此，在这一切之先，你们应该抛弃虚荣，在其它一切义务上，即在爱德、温良、兄弟之爱上成为完善。}然而，祂没有说任何这类话，而是极其温和地说：\Quote{伴郎岂能当新郎与他们在一起的时候守斋？}使它们想起若翰的话，他说：\BibleQuote{有新娘的是新郎；新郎的朋友，站着听祂，因新郎的声音而喜乐。}\BibleRef{若 3:29}\par

祂的意思是这样：现在是喜乐欢庆的时期，所以不要带来忧郁的事情。因为守斋是忧郁的事，不是本质上，而是对尚在软弱状态的人而言；因为至少对那些愿意节制自己的人，守斋是十分愉快和可取的。正如身体健康时精神高昂，同样灵魂状态良好时喜乐更大。但祂是根据他们先前的印象说这话。同样依撒意亚在谈论它时也称它为\Quote{压制灵魂的事；}梅瑟也一样。\par

然而祂不仅以此堵住他们的口，还用另一个话题说：\par

\BibleQuote{日子将到，新郎要从他们中被带走，那时他们就要守斋。}\BibleRef{玛 9:15}\par

因为祂借此表明，他们所行的并非出于贪食，而是属于某种奇妙的安排。同时，祂在与他人争论时，预先为祂要说的关于祂苦难的话奠定基础，以此教导祂的门徒，并训练他们如今能通晓那些被视为令人忧伤的事。因为若直接对他们说这话，本会令他们痛苦和刺痛——我们知道后来说出时，的确困扰了他们——但向别人说出，对他们而言反而会变得不那么难以承受。\par

他们也自然以\Person{若翰}的灾难自夸，祂便借此话题抑制了他们这样的傲慢；然而祂尚未加上关于祂复活的教导，因为时机未到。因为一个人被认为是凡人而死亡，这固然是自然的，但另一件事却超乎自然。\par

5. 于是，祂先前所行的，在此再次行之。我的意思是，正如当那些人试图证明祂与罪人同食是可责备的时，祂反而向他们证明，祂的做法非但无可指摘，反而是祂完全的赞美；同样，如今当他们想显示祂不知如何管理门徒时，祂便表明，这样的话是那些不知如何推论、随意挑剔之人的言行。\par

祂说：\Quote{没有人把新布补在旧衣服上。}\par

祂再次通过日常生活的比喻来建立祂的论证。祂所说的大意是：\Quote{门徒们尚未变得坚强，仍需要极大的迁就。他们尚未被圣神更新，对处于这种状态的人，不应该施加任何诫命的负担。}\par

祂说这些话，是为祂自己的门徒制定法则和规则，以便当他们将来要从全世界接纳各式各样的人作为门徒时，他们能够极其温柔地对待他们。\par

\Quote{也没有人把新酒装在旧皮囊里。}\par

你看见祂的比喻如何与旧约相似吗？那衣服？那皮囊？因为\Person{耶肋米亚}也称百姓为\Quote{腰带}，又提到\Quote{皮囊}和\Quote{酒}。\BibleRef{耶 13:10} 这样，既然论题涉及贪食和筵席，祂便从同一事物中取比喻。\par

但\Person{路加}记载同样的话，两三次且常常重复；然而并非令人厌烦的方式，而是轻松愉快地。现在你转向她，时而奉承她，时而追求她。\par

难道你没有看见画师们，当他们绘制一幅美丽的肖像时，要涂抹多少、添加多少吗？那么，你不要比他们差。因为若他们在绘制身体的肖像时如此勤勉，我们在塑造灵魂时，更当用尽一切方法。因为若你妥善塑造这灵魂的形态，你将不会看到身体的容貌显得不雅，嘴唇玷污，口如熊嘴般染血，眉毛黑得像厨房器皿的煤灰，脸颊白得像坟茔墙壁的尘土。因为这一切都是煤烟、灰烬、尘土和极度丑陋的标志。\par

可是，且住！我不知怎地不知不觉地陷入这些话中；在劝诫他人要温柔教导时，我自己却冲动而发怒。因此，让我们再回到更温和的劝诫方式，容忍我们妻子的一切缺点，好使我们能达成所愿。难道你没有看见我们如何容忍孩子的哭闹，当我们想要给他们断奶时，我们如何忍受一切，只为了一个目的：说服他们鄙弃先前的食物？同样，在此事上我们也要如此，容忍其余一切，好能完成此事。因为当这一项得到改正后，你将看到其他一切也按部就班地进行，你又会重新谈论金饰，并以同样方式推理它们，这样一点一点地将你的妻子引向正确的规则，你将成为一个美丽的画师、一个忠信的仆人、一个杰出的农夫。\par

同时，也要提醒她古时的妇女：\Person{撒辣}、\Person{黎贝加}，无论是貌美的还是不那么貌美的，并指出她们如何同样实践了谦逊。因为连圣祖的妻子\Person{肋阿}虽不美貌，也并未被迫想出任何这类事，尽管她相貌平平，又不太为丈夫所爱，她既未想出任何这类事，也未毁损自己的面容，而是继续保存其不受歪曲的线条，虽则由外邦人抚养长大。\par

但你这有信德的妇人，你以基督为元首，难道要给我们引进一种撒殚的技艺吗？难道你不记得那曾泼在你面容上的水、那装饰你嘴唇的祭献、那染红你舌头的血吗？因为若你思考这一切，即使你爱打扮到一万倍，你也不敢也不愿将那尘土和灰烬放在自己身上。你要明白，你已经与基督结合，要戒除这种不雅。因为祂并不喜欢这些涂抹，而是寻求另一种美丽——祂极其喜爱的那一种——我是指灵魂中的美丽。先知也曾嘱咐你珍视这美丽，并说：\Quote{君王必爱慕你的美容。}\par

因此，我们不要好奇地使自己变得不雅。因为天主的任何作品都不是有瑕疵的，也不需要你来修正。因为即使有人对已树立的皇帝肖像试图添加自己的作品，这种尝试也不安全，反而会招致极大的危险。那么，人为的制品，你不添加；但当天主的作品，你却要修改它？难道你不考虑地狱之火吗？难道你不考虑你灵魂的贫乏吗？正因为你所有的关心都浪费在肉身上，灵魂才被忽略。\par

然而，我为何只谈论灵魂呢？因为就连肉身本身，一切结果也与你所求的相反。请思考一下：你希望显得美丽吗？这反倒使你显得不美。你希望取悦你的丈夫吗？这反而使他难堪；不仅使他，也使陌生人成为你的控告者。你希望显得年轻吗？这会迅速将你带入老年。你希望穿戴体面吗？这却使你蒙羞。因为这样的人不仅在同伴面前羞耻，就连知晓她秘密的婢女和仆人面前也羞耻；而首先是在她自己面前。\par

但我为何还要说这些呢？因为我现在还遗漏了比这一切更严重的事，那就是：你得罪了天主；你破坏了端庄，点燃了嫉妒的火焰，效仿了妓女在娼馆中的行径。\par

那么，你们这些妇女，请思考这一切，并嘲笑撒殚的炫耀和魔鬼的诡计；舍弃这种修饰——或者不如说是毁容——在你们自己的灵魂中培养那连天使也喜爱、天主所渴望、并令你们丈夫喜悦的美丽；这样，你们就能获得现世的荣耀，以及那将来的荣耀。愿天主因我们的主耶稣基督的恩宠和爱人之情，赐予我们众人得以达到这一目标。愿光荣与权能归于他，直到永远。阿们。\par

\SourceRecord{Translated by George Prevost and revised by M.B. Riddle.；Nicene and Post-Nicene Fathers, First Series；Vol. 10.；Edited by Philip Schaff.；Buffalo, NY: Christian Literature Publishing Co.,；1888.；<http://www.newadvent.org/fathers/200130.htm>.}

% structure: p0001 source=h1 target=section reason=page-title

\section{玛窦福音第三十一篇讲道}

\BibleRef{玛 9:18}\par

\Quote{\Emph{耶稣正对他们说这些话时，忽然来了一个} \Emph{会堂长，朝拜祂说：\InnerQuote{我的女儿刚才死了，但请祢来，把手按在她身上，她必会活。}}}\par

所行的事印证了所说的话，使法利塞人更加无话可说。因为来者是一个会堂长，他的痛苦是可怕的。那女孩是他独生的女儿，年方十二岁，正值青春年华；正因如此，主特别将她复活，并且是立即的。\par

如果路加说有人来告诉祂：\Quote{不要劳烦老师，因为她已经死了。}\BibleRef{路 8:49}，那么我们可以说，\newline{}\Quote{她刚才死了}这话是那人根据路程时间所作的推测，或是夸大其词。因为处于困境中的人通常会在报告时夸大自己的苦难，说一些超过实际情况的话，以更打动他们所恳求的人。\par

但请看他的迟钝：他向基督要求两件事，既需要祂亲临，又需要祂按手；这又表明他离开时女孩还在呼吸。那位叙利亚人纳阿曼也向先知如此要求。\Quote{因为我想，}他说，\Quote{他一定会出来，并用手抚摸患处。}的确，那些较为迟钝的人，往往需要看得见、感觉得到的东西。\par

至于马尔谷（\BibleRef{谷 5:37}）说祂带了三位门徒，路加也是如此（\BibleRef{路 8:51}），我们的福音作者只说\Quote{门徒}。那么祂为什么不带玛窦同去呢？尽管玛窦刚刚来到祂面前。这是为了激发他更强烈的渴望，也因为他尚处于不完美的状态。因为祂尊荣那些人，是为了让这些人也能成长为那样的人。至于玛窦，他当时足以看到患血漏妇人的遭遇，并因主的餐桌和祂分享的盐而得尊荣。\par

当祂起来时，许多人跟随祂，为要观看这伟大的奇迹，既因那来的人，也因为大多数人心思粗俗，寻求的不是灵魂的医治，而是身体的治愈；他们涌来，有的为自己的病痛所迫，有的急于观看他人如何得治愈。然而，习惯上主要为祂的话语和教导而来的人还很少。尽管如此，祂没有让他们进入屋里，只带了门徒，甚至连门徒也不是全部，处处教导我们要拒绝群众的称赞。\par

2. 经上说：\Quote{看，有一个患血漏十二年的妇人，从后面走来，摸了祂的衣边；因为她心里说：\InnerQuote{我只要摸到祂的衣裳，就会痊愈。}}\par

为什么她不勇敢地走上前去？她因自己的病而羞愧，自认是不洁的。如果行经的女人被认为不洁，那么患这种病的妇人更当自认为不洁；事实上，按法律，这病被视为大不洁（\BibleRef{肋 15:25}）。因此她隐藏自己，遮掩自己。因为她还没有对祂有正确和恰当的认识，否则不会以为能隐藏。这是第一个在公开场合来到祂面前的女人，她当然听说过祂也医治妇女，并且祂正前往那已死的小女儿那里。\par

她不敢邀请祂到自己家里，尽管她富有；她也没有公开上前，而是怀着信心悄悄摸了祂的衣裳。因为她没有怀疑，也没有心里说：\Quote{我真的会脱离疾病吗？我真的不会得痊愈吗？}但她确信自己会痊愈，就这样来到祂面前。经上记载：\Quote{她心里说：\InnerQuote{我只要摸到祂的衣裳，就会痊愈。}}因为她看到祂从什么样的房子里出来（税吏的家），以及谁跟随祂（罪人和税吏），这一切使她充满希望。\par

那么基督做了什么？祂不容她隐藏，反而把她带到中间，使她显明出来，为了多方面的目的。\par

确实，有些愚昧的人说：\Quote{祂这样做是出于虚荣。}他们说：\Quote{为什么祂不让她隐藏起来？}不洁的人，你说了什么？祂禁止张扬，祂略过无数的奇迹，难道祂是爱虚荣的吗？\par

那么祂为什么把她带出来？首先，祂消除了妇人的恐惧，免得她因良心自责（仿佛偷了恩赐）而痛苦。其次，她纠正了她自以为能隐藏的想法。第三，祂向众人显明她的信德，以激励其余的人也效仿；而祂止住她血漏的征兆，并不比祂显明祂的全知更为伟大。此外，会堂长正濒临彻底的不信，因而面临完全的毁灭，祂通过妇人纠正了他。因为那些来的人说：\Quote{不要劳烦老师，因为女孩死了；}家里的人当祂说\Quote{她睡着了}时都嘲笑祂；父亲也可能有类似的感受。因此，为了预先纠正这个弱点，祂把这平凡妇人带到台前。至于那会堂长是何等粗俗，请听祂对他说的话：\Quote{不要怕，只管信，她必会得救。}（\BibleRef{路 8:50}）\par

祂也特意等待死亡到来，然后祂才到达，以使复活的证据更加清晰。为此，祂走得较为缓慢，并与那妇人更多地交谈，好让女孩死去，并让那些来报告的人前来，他们说：\Quote{不要麻烦主。}这无疑又是福音作者隐晦地指出的，他说：\Quote{祂还在说话时，有人从屋里来，说：\InnerQuote{你的女儿死了，不要麻烦主。}}因为祂的旨意是要她的死被相信，这样她的复活就不会被怀疑。祂在每一件事上都如此。同样在拉匝禄的事上，祂等了第一、第二和第三天。\BibleRef{若 11:6}\par

因此，为了这一切，祂把她带上前来，说：\Quote{女儿，放心，}正如祂也对瘫痪者说过：\Quote{孩子，放心。}因为那妇人确实极其惊慌；因此祂说：\Quote{放心，}并称她为\Quote{女儿；}因为她的信德使她成为女儿。随后也有对她的称赞：\Quote{你的信德救了你。}\par

但路加还告诉我们关于这妇人的更多其他事。如此，他说，当她靠近祂，并得到了痊愈，基督并没有立刻叫她，而是先问：\Quote{谁摸了我？}然后，当伯多禄和同祂在一起的人说：主，群众拥挤着你，压着你，你说谁摸了我？\BibleRef{路 8:45}（这是一个非常可靠的迹象，表明祂身上有真实的身体，并且祂摒弃了一切虚荣，因为他们并非远远地跟随祂，而是从四面拥挤祂）；祂继续回答说：\Quote{有人摸了我，因为我感觉有能力从我身上出去了；}这是以较为粗俗的方式，顺应听者的感受。但祂说这些话，是为了促使那妇人自己承认。为此，祂并没有立即揭发她，而是为了在表明祂清楚知道一切之后，促使她自愿说出一切，并让她自己宣告所发生的事，免得祂因自己说出来而招致怀疑。\par

你看见那妇人胜过会堂长吗？她没有留住祂，没有抓住祂，只是用指尖触摸了祂，虽然她来得更晚，却先得痊愈而离去。那人确实要把医生请到他家里，但对这妇人，仅一触摸就够了。因为她虽被痛苦束缚，但她的信德给了她翅膀。且看祂如何安慰她，说：\Quote{你的信德救了你。}当然，若祂把她带出来是为了炫耀，祂就不会加上这句话；但祂说这话，部分是为了教导会堂长相信，部分是为了宣告对妇人的称赞，并用这些话赐予她与她身体健康相等的喜悦和益处。\par

因为祂这样做，是出于要光荣她，并改善他人，而不是要显扬自己，这从以下事实显而易见：即使没有这件事，祂本也同样是可钦羡的对象（因为奇迹如雪片般纷纷降临祂周围，祂已经做了并将要做比这些更大的事）；但这妇人，若没有发生这件事，就会隐藏而去，失去那些巨大的赞美。为此，祂把她带上前来，宣告对她的赞美，并驱除了她的恐惧（因为据说\Quote{她来了，}\Quote{战战兢兢}）；祂使她鼓起勇气，并在身体健康之外，还给她的旅程准备了其他所需，说：\Quote{平平安安地去吧。}\BibleRef{路 8:48}\par

3. \Quote{耶稣来到会堂长的家里，看见吹笛的和乱哄哄的群众，就对他们说：退出去，因为女孩不是死了，而是睡着了。他们就讥笑祂。}\par

这些无疑是会堂长的尊贵表现；在她死亡之际，笛子和铙钹奏起挽歌！基督做了什么？祂把所有其余的人都赶出去，却把父母领进来；以免留下别人说祂用别的方式治愈了她。而且在使她复活之前，祂用言语使她复活，说：\Quote{女孩不是死了，而是睡着了。}此外在许多事例中祂也这样做。正如在海上，祂从旁观者心中驱除骚动，同时表明祂复活死者是多么容易（祂对拉匝禄也这样做，说：\Quote{我们的朋友拉匝禄睡着了}；\BibleRef{若 11:11}同时也教导我们不要害怕死亡；因为死亡不再是死亡，而是从此成为睡眠）。因此，既然祂自己将要死亡，祂就在别人身上预先训练门徒，使他们鼓起勇气，并温和地承受结局。因为事实上，当祂来到时，死亡从此成为了睡眠。\par

但那些人还是讥笑祂；然而祂并不因那些祂稍后就要为他们行奇迹的人的不信而愤怒；祂也没有责备他们的讥笑，为的是让讥笑、笛子、铙钹以及所有其他事物，都能成为她死亡的确凿证据。因为通常在奇迹发生后，人们仍然不信，祂就预先用他们自己的话来弄住他们；这在拉匝禄和梅瑟的事上都如此。因为祂首先对梅瑟说：\Quote{你手里是什么？}\BibleRef{出 4:2}为的是当他看见它变成蛇时，他不会忘记它原本是棍杖，而是因自己的话语被提醒，对所发生的事感到惊奇。关于拉匝禄，祂说：\Quote{你们把他安放在哪里？}\BibleRef{若 11:34} 以便那些说过\Quote{你来看看，}和\Quote{他已经臭了，因为死了四天}的人，再也不能不信祂复活了死人。\par

于是，看见铙钹和群众，祂把他们全都赶出去，并在父母面前施行了奇迹；不是引入另一个灵魂，而是召回那已经离去的同一个灵魂，如同从睡梦中唤醒她。\par

祂握住她的手，以此向旁观者保证，好藉此景象为她的复活开辟信德之路。因为那父亲说：\Quote{请将手按在她身上。}\BibleRef{玛 9:18} 而祂却做得更多：祂没有按手在她身上，反倒抓住她、扶起她，表明对祂而言万事都已预备妥当。祂不仅使她复活，还吩咐给她食物，免得此事看似幻象。祂不是亲自给她，而是吩咐他们；对拉匝禄也是这样：祂说：\Quote{解开他，让他走。}\BibleRef{若 11:44} 后来还使他参与祂的宴席\BibleRef{若 12:2}。因为祂总是这样完全确立两点，以最圆满的方式同时证明死亡与复活。\par

但请你留意，我恳求你，不只是她的复活，还有祂命令\Quote{不要告诉任何人}；借此你尤其要学习祂的谦逊与不虚荣。此外还要学习这一点：祂把那些捶胸痛哭的人赶出屋子，宣布他们不配见到这样的景象；你不要跟着吹鼓手出去，而要留在伯多禄、若望和雅各伯身边。\par

因为若祂当时将他们赶出去，何况现在呢？因为那时死亡成为睡眠尚不明显，但现在这比太阳本身更清楚。也许祂现在没有使你的女儿复活？但祂必定会使她复活，而且带着更丰富的荣耀。因为那少女复活后还要再死；而你的孩子若复活，从此将永存不朽。\par

4. 所以，不要再有人捶胸痛哭，不要哀号，不要贬低基督的成就。因为祂确实战胜了死亡。你为什么徒然哀号？这事已成睡眠。你为什么悲叹哭泣？即使外教人这样做，也当被嘲笑；但当信徒在这些事上行为失当，他还有什么理由？犯了如此愚昧之罪的人，在这样长的时间、如此明确的复活证据之后，还有什么借口？\par

而你，仿佛在加重对你的控诉，竟还带进外教妇女唱挽歌，来点燃你的情感，彻底煽起炉火；你不听保禄的话说：\Quote{基督与贝里雅耳有什么和谐？信者与不信者有什么相干？}\par

当外教人的子女，他们对复活一无所知，却能找到安慰的话说：\Quote{要勇敢承受，因为既成事实无法挽回，悲叹也于事无补}；而你，听到比这更智慧、更好的话语，难道不羞愧于表现得更不像样吗？因为我们根本不说\Quote{要勇敢承受，因为既成事实无法挽回}，而是说\Quote{要勇敢承受，因为他必定会复活}；孩子睡着了，并没有死；他安息了，没有丧亡。因为复活将是他的终局，永生、不朽和天使的福分。你没有听圣咏说：\Quote{我的灵魂，回归你的安息吧，因为上主厚待了你}？天主称之为\Quote{厚待}，而你却哀号？\par

还有什么更多的，如果你是与死者为敌为仇？如果必须哀悼，那应当是魔鬼哀悼。他可以捶胸，可以痛哭，因为我们正走向更大的福乐。这哀哭属于他的邪恶，不属于你，你即将被加冕并安息。是的，因为死亡是美好的港湾。无论如何，请思量我们的今生充满了多少邪恶；请回想你多少次诅咒了今世。因为事情每况愈下，并且从一开始你就陷入了不小的定罪。因为祂说：\Quote{你将在痛苦中生子；}\Quote{你必汗流满面才得糊口；}\Quote{在世上你们将有苦难。}\BibleRef{若 16:33}\par

但在那里，我们的处境没有任何这样的话，反而全是相反的：即\Quote{忧伤、悲哀和叹息都已逃避。}\BibleRef{依 35:10} 以及\Quote{将有许多人从东方和西方来，并要在亚巴郎、依撒格和雅各伯的怀中安息。}\BibleRef{玛 8:11} 并且那里是属灵的洞房、明亮的灯，以及升到天上。\par

5. 那么为什么要辱没逝者？为什么要使其余的人惧怕死亡、战栗？为什么要使许多人指责天主，仿佛祂做了极可怕的事？或者，为什么此后要邀请穷人、恳求司铎祈祷？他说：\Quote{为使亡者进入安息；使他能寻得法官的慈悯。} 你既为此哀号痛哭，岂不是与自己交战，因他进入了港湾而自招风暴？\par

\Quote{但我能做什么呢？}他说：\Quote{这是本性。} 责任不在本性，也不在那事的必然结果；而是我们颠倒了一切，被软弱征服，放弃了应有的高贵，使不信者变得更糟。因为我们将如何与他人谈论不朽？我们将如何说服外教人，如果我们比他们更怕死，更因死而战栗？例如，许多希腊人虽然当然不知道不朽，却在子女夭折时给自己戴上花冠，穿着白衣，为得现世的荣耀；而你，即使为了将来的荣耀，也不停止你的妇人之态和哀号。\par

你没有继承人，没有人继承你的财产？你宁愿他成为你财产还是天上的继承人？你希望他继承消逝之物、不久后必放弃的，还是那存留不动之物？你未曾有他为继承人，而是天主替你得到了他；他没有成为与他兄弟的共同继承人，而是成了\Quote{与基督同为继承人}。\par

\Quote{但是要留给谁？}他说：\Quote{我们要把衣服留给谁？把房屋留给谁？把奴隶和土地留给谁？}要留给他，而且比他活着时更稳妥，因为没有任何阻碍。因为若外邦人将逝者的财物一起焚烧，那么你更应当将逝者所拥有的东西送与他同去：不是像那些人一样化为灰烬，而是使他获得更大的荣耀；并且若他去世时是罪人，这可以消除他的罪；若是义人，则能增加赏报和酬劳。\par

但你渴望见他吗？那么与他过同样的生活，你很快就会获得那神圣的目睹。\par

并且请考虑这一点：即使你不听我们的话，你也一定会屈服于时间。但那时对你没有赏报；因为安慰来自于日子的数量。然而，如果你现在愿意管束自己，你将获得两个极大的好处：首先，你将摆脱中间的痛苦；其次，你将从天主那里得到更明亮的荣冠。因为实际上，无论施舍还是其他任何事，都不如温顺地忍受苦难那样伟大。\par

请记住，连天主子也死了：他为你而死，而你为自己而死。当他说：\BibleQuote{如果可能，就让这杯从我面前过去吧，}\BibleRef{玛 26:39}并且忍受痛苦，处于痛苦之中，然而他并没有逃避结局，而是承受了它，并且伴随着极其悲惨的全过程。也就是说，他不仅仅简单地忍受了死亡，而是最耻辱的死亡；在死亡之前，有鞭打；在鞭打之前，有辱骂、嘲笑和毁谤；教导你要勇敢地承受一切。虽然他死了，脱下了身体，但又在更大的荣耀中重新取回，这同样给你带来了美好的希望。如果这些不是虚构的故事，就不要哀叹。如果你认为这些事是确实的，就不要哭泣；但如果你哭泣，你怎能说服外邦人相信呢？\par

6. 但即便如此，事情对你来说仍然不可忍受吗？那么，正因为如此，你不应该为他哀悼，因为他已脱离了许多这样的灾难。因此不要怨恨他，也不要嫉妒他：因为因他夭折而为自己求死，并为他未能活着忍受许多这样的事而哀悼，这更像是怨恨和嫉妒的表现。\par

并且不要想这个：他不会再回家来了；而是想：你自己过不多久也要到他那里去。不要注意这个：他不再回到这里，而是也要注意这些看得见的事物不会保持现状，它们也在被改变。是的，因为天、地、海和一切都在重新组合，那时你将更荣耀地找回你的孩子。\par

而且如果他去世时是罪人，他的邪恶就停止了；因为的确，如果天主知道他正在悔改，他必不会在他悔改之前就将他接走；但如果他结束生命时是义人，他现在就安全地拥有了一切美善。由此可见，你的眼泪不是出于亲情，而是出于无理性的激情。因为如果你爱那逝者，你应该为他已脱离现世的风浪而欢喜快乐。\par

因为还有更多，请问？有什么新的呢？我们不是每天都看到同样的事物在循环吗？白天和黑夜，黑夜和白天，冬天和夏天，夏天和冬天，没有更多了。这些事物始终如一；但我们的邪恶却是新的，更新奇。那么你愿意他每天都获取更多的这些，并留在这里，生病、哀哭、恐惧战栗，忍受现世的一些痛苦，又害怕将来会忍受另一些吗？因为肯定你不能说，航行在这大海之上可能没有沮丧和忧虑，以及所有其他类似的事。\par

此外，也要考虑这一点：你并没有生他不死；而且如果他现在不死，他不久也必须死。但是你还没有与他相处够吗？但你一定会在那里享受他。但你也想在这里见到他吗？有什么阻止你呢？因为即使在这里，如果你警醒，你也被允许；因为未来事物的希望比看见更清晰。\par

但如果你，若他在某个君王的宫廷里，你只要听到他的好消息，就不想见他；现在看到他去了更美好的地方，你却为一点点时间而沮丧？而且，当你有一个人代替他与你同住时？\par

但你没有丈夫吗？然而你有一个安慰，就是孤儿的父亲和寡妇的审判者。请听\Person{保禄}宣布这种寡妇是有福的，并说：\Quote{真正无依无靠的寡妇，孤苦伶仃，信赖主。}因为这样的人将显得更加蒙悦纳，显出了更大的忍耐。因此不要为你的荣冠而哀哭，那正是你要求赏报的缘由。\par

因为你已归还了他的寄存物，如果你展示了所托付给你的东西。不要再多挂虑，因为你已将财产存放在一个不能侵犯的宝库里。\par

但如果你真的想学习，我们现在的存在是什么，以及我们未来的生命是什么；并且知道一个是蜘蛛网和阴影，而那里的所有事物都是不动和不朽的；那么你之后就不需要其他论据了。因为现在你的孩子已脱离了一切变化；如果他在这里，也许他会继续好，也许不会。你看不到有多少人公开抛弃自己的子女吗？有多少人被迫将他们留在家里，尽管他们比公开被抛弃的更糟糕吗？\par

让我们考虑所有这一切并操练自制；因为这样我们既关心逝者，也从人那里获得许多赞美，并从天主那里接受忍耐的大赏报，达到永恒的美善；愿我们都能达到，藉着我们的主\Person{耶稣基督}的恩宠和爱人之心，愿光荣和权能归于他，直到永远。阿们。\par

\SourceRecord{Translated by George Prevost and revised by M.B. Riddle.；Nicene and Post-Nicene Fathers, First Series；Vol. 10.；Edited by Philip Schaff.；Buffalo, NY: Christian Literature Publishing Co.,；1888.；<http://www.newadvent.org/fathers/200131.htm>.}

% structure: p0001 source=h1 target=section reason=page-title

\section{\BibleBook{玛窦}{福音}第三十二篇讲道}

\BibleRef{玛 9:27-30}\par

\BibleQuote{\Emph{耶稣从那里前行，有两个瞎子跟着祂喊叫说：\InnerQuote{达味之子，可怜我们罢！}祂来到家里，瞎子们便走到祂跟前；耶稣对他们说：\InnerQuote{你们信我能作这事吗？}他们回答说：\InnerQuote{是，主。}于是祂摸了摸他们的眼说：\InnerQuote{照你们的信德，给你们成就罢！}他们的眼就开了。}}\par

为什么祂要拖延他们，任他们喊叫呢？这里再次教训我们完全抛弃来自群众的荣耀。因为房子就在附近，祂领他们到那里，私下医治他们。这事由祂随后还嘱咐他们不可告诉人而明显可见。\par

但这对犹太人不是轻微的指控：这些人虽然眼被蒙蔽，却单凭听觉接受了信仰，而他们目睹神迹，亲眼见证所发生的事，却完全相反。请看他们的热忱，既通过喊叫，也通过祈祷本身。因为他们不仅走近祂，而且大声喊叫，只提到\Quote{怜悯}。\par

他们称祂为\Quote{\InnerQuote{达味之子}}，因为这称号被认为是光荣的。例如在好多处，先知们也这样称呼他们想要尊崇、显其伟大的君王。\par

把他们领进屋里后，祂进一步询问他们。因为在许多情况下，祂坚持在祈求后施行医治，以免有人以为祂是出于虚荣而匆忙行神迹；不仅为此，也为了表明他们应得医治，并且没有人会说：\Quote{如果祂是出于纯粹的怜悯而拯救，所有人都应当得救。}因为祂的爱人慈心也有一种比例，取决于被治愈者的信仰。但祂要求他们信仰不仅为此，而且因为他们称祂为\Quote{\InnerQuote{达味之子}}，祂提升他们到更高处，教导他们应当对祂本人怀有怎样的想法，便说：\Quote{你们信我能作这事吗？}祂没有说：\Quote{你们信我能恳求我的父、我能祈祷吗？}而是\Quote{我能作这事吗？}\par

那么他们的回答是什么？\Quote{是，主。}他们不再称祂为达味之子，而是飞升更高，承认祂的主权。\par

于是最后祂按手在他们身上说：\Quote{照你们的信德，给你们成就罢。}祂这样做是为了坚定他们的信德，表明他们参与了这件善工，并见证他们的话不是谄媚之辞。因为祂没有说\Quote{让你们的眼睁开}，而是\Quote{照你们的信德，给你们成就罢}；祂对许多来到祂面前的人都这样说，在医治他们身体之前，就急于宣告他们灵魂中的信德，既使他们更被认可，也使其他人更加认真。\par

对于瘫子也是如此：因为在那里，在使身体有力之前，祂先扶起堕落的心灵，说：\Quote{孩子，放心！你的罪赦了。}还有那少女，当祂使她起来后，留住她，用食物使她认识她的恩主；百夫长的事上祂也照样行，将一切交托给他的信德；至于祂的门徒们，当祂在海上救他们脱离风暴时，先救他们脱于不信。同样在这里：祂在喊叫之前，就已知道他们内心的秘密；但为了也引导其他人到同样的热忱，祂将他们的信德向其余人公开，藉着治愈的结果宣告他们隐藏的信德。\par

在他们被治愈后，祂命令他们不可告诉人；不单是命令，而且非常严格。\par

\BibleQuote{因为耶稣，严厉地嘱咐他们说：\InnerQuote{小心，不可让人知道。}但他们出去后，却将祂的名声传遍了那整个地区。}\BibleRef{玛 9:30}\par

然而他们没有忍受这事，反而成了宣讲者和传福音者；当被吩咐隐藏所行的事时，他们不能忍受。\par

如果我们在另一处发现祂说：\Quote{你们去，宣扬天主的荣耀}，那与此并不矛盾，反而是非常一致的。因为祂教导我们，对于我们自己，不要说什么，甚至禁止那些赞扬我们的人；但如果荣耀归于天主，那么不仅不禁止，反而命令人这样做。\par

2. \Quote{他们出来的时候，看，有人带了一个附魔的哑巴来到祂跟前。}\par

因为这病症不是自然的，而是恶灵的伎俩；因此也需要别人带他来。因为他自己不能恳求（因为他哑了），也不能恳求别人，当恶灵捆绑了他的舌头，连同舌头捆绑了他的心灵。\par

为此原因，祂也不要求他的信德，而是直接治愈疾病。\par

\BibleQuote{魔鬼被赶出去后，哑巴就说出话来。群众惊奇说：\InnerQuote{在以色列从来没有见过这样的事。}}\BibleRef{玛 9:33}\par

这件事尤其激怒了法利塞人，因为众人把祂置于一切人之上，不仅当时的人，而且所有以往的人。他们推崇祂，不是因为祂的医治，而是因为祂轻而易举、迅速地治疗了无数不治之症。\par

群众如此；但法利塞人恰恰相反，不但诋毁这些事，而且自相矛盾，毫不羞愧。邪恶就是这样。因为他们说什么？\par

\Quote{祂是靠着魔王驱魔。}\par

有什么比这更愚蠢的呢？因为首先，正如祂随后所说的，魔鬼不可能驱逐魔鬼，因为那存在惯于修复属于自己的，而不是拆毁它。但祂不仅驱逐魔鬼，还洁净癞病人，复活死人，平静大海，赦免罪过，宣讲天国，并引领人归向圣父；这些事是魔鬼绝不会选择，也任何时候都不能做到的。因为魔鬼引人崇拜偶像，使人远离天主，并说服他们不信未来的生命。魔鬼在被侮辱时不会施恩；因为即使不被侮辱，它也伤害那些讨好和尊敬它的人。\par

但祂却做相反的事。因为在他们这些侮辱和谩骂之后，\par

\Quote{祂周游}，经上说，\Quote{各城各村，在他们的会堂里施教，宣讲天国的福音，并治好各种疾病和各种灾殃。}\par

而且祂并未因他们的麻木而惩罚他们，甚至没有简单地责备他们；既显示祂的温良，从而驳斥了诽谤；同时祂也意在通过随后发生的征兆更充分地展示祂的证明：然后还以言语提出反驳。因此祂周游各城各乡，并在他们的会堂里；教导我们，要以更大的善行回应诽谤我们的人，而不是以新的诽谤。既然，若不是为了人的缘故，而是为了天主的缘故，你善待同仆；无论他们做什么，你不可停止善待他们，好使你的赏报更大；因为那因他们的诽谤而停止行善的人，表明他追求这种美德是为了人的称赞，而不是为了天主。\par

为此原因，基督为了教导我们，祂纯粹出于良善而着手此事，祂非但没有等待病人到祂这里来，反而亲自赶快到他们那里去，带给他们两个最大的祝福：一是天国的福音，二是完全治愈他们的一切疾病。祂没有忽略一座城，没有匆忙经过一个村庄，而是访问了每一个地方。\par

4. 祂甚至不止于此，还展示了祂预见的另一个实例。即，\par

\Quote{祂看见}，经上说，\Quote{群众，就对他们动了怜悯的心，因为他们困苦流离，像没有牧人的羊。于是祂对门徒说：庄稼实在多，工人却少；所以你们应当求庄稼的主，派遣工人来收割祂的庄稼。}\par

再看祂的不求虚荣。为免吸引众人归向祂自己，祂派遣门徒出去。\par

而且不仅为此目的，也为了教导他们，在巴勒斯坦操练之后，如同在某种训练学校里，使他们卸下装备，以便与世俗搏斗。为此目的，祂使这些操练甚至比实际的战斗更严肃，只要关乎他们自己的德性；以便他们更容易应付即将来临的斗争；就像一群幼雏，祂最终要带领它们飞翔。目前，祂使他们成为身体的医生，以后赐予他们灵魂的医治，这是主要的事。\par

且看祂如何指出此事的容易与必要。因为祂说什么？\Quote{庄稼实在多，工人却少。} 即，祂说，\Quote{不是去播种，} \Quote{而是去收割，我派遣你们。} 这一点在若望中由祂表达为：\Quote{别人劳苦，你们却享受他们的劳苦。} \BibleRef{若 4:38}\par

祂说这些话，既抑制他们的骄傲，又鼓励他们鼓起勇气，并表明主要的部分工作已经预先完成了。\par

且看祂在此也是从爱世人开始，没有任何回报。因为祂怜悯他们，因为他们困苦流离，像没有牧人的羊。这是祂对犹太首领的控诉，他们身为牧人，却行狼的作为。因为他们非但没有改善群众，反而妨碍了他们的进步。例如，当他们惊讶地说：\Quote{在以色列从未见过这样的事}时，这些人却肯定相反的事：\Quote{祂是靠魔王驱魔。}\par

但祂在这里说的是什么工人？是指十二门徒。那么怎样？祂既然说：\Quote{工人却少，} 祂增加了他们的人数吗？绝不，但祂单独派遣他们出去。那么祂为什么说：\Quote{你们应当求庄稼的主，派遣工人来收割祂的庄稼，} 却没有增加他们的人数呢？因为虽然他们只是十二人，但从那时起祂使他们成为众多，不是增加人数，而是赐予他们权能。\par

然后为表示这恩赐有多大，祂说：\Quote{你们应当求庄稼的主；}并间接宣称这是祂的特权。因为在说了\Quote{你们应当求庄稼的主}之后，他们还没有任何恳求或祈祷，祂自己就立刻任命了他们，同时提醒他们若翰关于打谷场、簸扬的人、糠秕和麦子的话（\BibleRef{玛 3:12}）。由此显然可知，祂自己就是农夫，祂自己就是庄稼的主，祂自己就是先知的主人和拥有者。因为如果祂派遣他们去收割，显然不是收割别人的，而是收割祂自己藉先知所播种的。\par

但祂不仅用这种方式间接鼓励他们，称他们的职务为收割；而且还使他们有能力胜任这职务。\par

\Quote{祂叫了}，经上说，\Quote{祂的十二门徒来，赐给他们权柄制服邪魔，以驱逐它们，并治好各种疾病和各种灾殃。}\par

然而圣神尚未赐下。因为经上说：\Quote{还没有赐下圣神，因为耶稣尚未受到光荣。}那么他们如何驱逐魔鬼呢？藉着祂的命令，藉着祂的权柄。\par

也请留意，我恳求你，这派遣的时机何等恰当。因为祂并非一开始就派遣他们，而是当他们充分享受了跟随祂的益处，并看见死人复活，大海受斥责，魔鬼被驱逐，瘫痪者重新得力，罪过得赦免，癞病人被洁净，并且藉着言行对祂的能力有了足够证明之后，祂才派遣他们出去；而且不是去危险的事，因为那时巴勒斯坦尚无危险，他们只需抵挡恶言。然而，祂甚至预先警告他们这些事，我是指他们的危险；在时候之前就预备他们，并藉着不断作那类预言，使他们如同在冲突中一样。\par

5. 然后，既然祂曾向我们提及两对宗徒，即伯多禄的那对和若望的那对，并在那之后指出了玛窦的蒙召，但对其他宗徒的蒙召和名字却对我们只字未提；此时祂必然列出了他们的名单和数目，并使人知道他们的名字，这样说：\par

\Quote{十二宗徒的名字是这些：首先，\Person{西满}，又称\Person{伯多禄}。} \BibleRef{玛 10:2}\par

因为还有另一个\Person{西满}，加纳人；又有\Person{依斯加略的犹达斯}，以及\Person{雅各伯}的兄弟\Person{犹达斯}；又有\Person{阿耳斐}的儿子\Person{雅各伯}，和\Person{载伯德}的儿子\Person{雅各伯}。\par

现在，\Person{马尔谷}确实按他们的尊严排列；因为在两位首领之后，祂将\Person{安德肋}列在其次；但我们的福音作者却不是这样，而是不加区别；或者更好说，祂甚至将那远不如祂的\Person{多默}放在祂自己之前。\par

但让我们从开始看他们的名单。\par

\Quote{首先，\Person{西满}，又称\Person{伯多禄}，和他的兄弟\Person{安德肋}。}\par

连这也不是小的赞美。因为祂给前者命名是出于他的德行，给后者命名是出于他的高贵亲属，这符合祂的性情。\par

\Quote{\Person{载伯德}的儿子\Person{雅各伯}，和他的兄弟\Person{若望}。}\par

你们看见祂如何不按他们的尊严排列吗？因为在我看来，\Person{若望}似乎更大，不仅超过其他人，甚至超过他的兄弟。\par

此后，当祂说了\Quote{\Person{斐理伯}和\Person{巴尔多禄茂}}之后，祂又加上\BibleQuote{\Person{多默}和税吏\Person{玛窦}}。\BibleRef{玛 10:3}\par

但\Person{路加}却不是这样，而是以相反的顺序，并且把他放在\Person{多默}之前。\par

其次，\Quote{\Person{阿耳斐}的儿子\Person{雅各伯}}。因为还有\Person{载伯德}的儿子，正如我已经说过的。然后，在提及了\Quote{\Person{勒拜}，其别号是\Person{达陡}}和\Quote{\Person{西满}}热诚者（他也称之为\Quote{加纳人}）之后，祂才来到叛徒那里。祂并非像对待仇敌或敌人那样，而是如同写历史的人那样描述了他。祂没有说\Quote{不圣洁的、极不圣洁的}，而是以他的城市命名他，\Quote{\Person{依斯加略的犹达斯}}。因为还有另一个\Person{犹达斯}，\Quote{\Person{勒拜}，其别号是\Person{达陡}}，\Person{路加}说他是\Person{雅各伯}的兄弟，说\BibleQuote{\Person{雅各伯}的兄弟\Person{犹达斯}}。\BibleRef{路 6:16}因此，为了将他与这个人区分，经上说\BibleQuote{\Person{依斯加略的犹达斯}，祂的背叛者。}\BibleRef{玛 10:4}祂不以说\Quote{祂的背叛者}为耻。\par

在所有中为首的和歌咏团的领唱者，是那\BibleQuote{没有学问的、无知的人。}\BibleRef{宗 4:13}\par

但让我们看祂派遣他们往何处、到谁那里去。\par

\BibleQuote{这十二人，}据说，\BibleQuote{\Person{耶稣}派遣出去。}\BibleRef{玛 10:5}\par

这些是什么样的人？渔夫、税吏：因为确实有四个渔夫和两个税吏，\Person{玛窦}和\Person{雅各伯}，还有一个甚至是叛徒。祂对他们说什么？祂当下命令他们说：\par

\BibleQuote{外邦人的路，你们不要走；撒玛黎雅人的城，你们不要进；宁可往以色列家迷失的羊那里去。}\BibleRef{玛 10:5}\par

\Quote{你们千万不要以为，}祂说，\Quote{因为他们侮辱我，称我附魔者，我就恨他们并远离他们。不，正如我起初热切寻求改正他们一样，所以使你们远离其余的人，派遣你们到他们那里作教师和医生。我不仅禁止你们在这些事之前向别人宣讲，我甚至不允许你们踏上通往那里的路，也不许进入那样的城。}因为撒玛黎雅人也与犹太人处于敌对状态。然而，与他们打交道是一件更容易的事，因为他们对信仰更怀好意；但犹太人的情况更困难。尽管如此，祂仍派遣他们去完成更艰难的任务，表明祂对他们的照顾保护，堵住犹太人的口，并为宗徒的教导预备道路，以免日后人们因他们\BibleQuote{进入未受割礼的人那里}而责备他们，\BibleRef{宗 11:3}并认为他们有正当理由回避和厌恶他们。祂称他们为\Quote{迷失的}，不是\Quote{走散的}，\Quote{羊}，千方百计设法为他们开脱，并赢得他们的心归向自己。\par

6. \BibleQuote{你们随走随传，}祂说，\BibleQuote{说：天国近了。}\BibleRef{玛 10:7}\par

你们看见他们职务的伟大吗？你们看见宗徒的尊荣吗？他们受命不谈任何感官对象的事，也不谈\Person{梅瑟}和先于他们的先知所说的事，而是谈一些新奇的事。因为前者宣讲的不是这样的事，而是地上和地上的美好事物；这些宗徒却宣讲天国和其中所有的一切。\par

不仅因这情形他们更为伟大，也因他们的服从：他们不像古人那样退缩或落后；而是，尽管被警告有危险、战争和那些无法承受的邪恶，他们仍以极大的服从接受祂的命令，作为天国的先驱。\par

有人却说：\Quote{他们所传讲的，没有可悲或痛苦的事，所以他们欣然听从，这有什么稀奇呢？}你说什么？没有严命他们什么痛苦的事吗？难道你没有听见监狱、死刑、内战、众人的仇恨吗？这些事，主在不久后就说他们必须承受。的确，对于其他人，主派遣他们去成为无数祝福的\OriginalTerm{获取者与传报者}{procurers and heralds}，但为他们自己，主预先宣告他们要忍受可怕且难治的祸患。\par

此后，要使他们可信，主说：\par

\Quote{医治病人，洁净癞病人，驱逐魔鬼：你们白白得来的，也要白白施舍。}\par

请看主多么关注他们的品行，丝毫不亚于关注他们的奇迹，暗示没有品行，奇迹算不得什么。主既藉着说\Quote{你们白白得来的，也要白白施舍}来压伏他们的骄傲，又安排他们远离贪婪。此外，恐怕他们以为这是自己的功劳，因所行的奇迹而自高，主就说\Quote{白白得来的}。主说：\Quote{你们并未施恩于接待你们的人，因为你们没有为这些付出代价，也没有劳苦：恩宠属于我。同样，你们也要白白施舍给他人，因为找不到相称的代价。}\par

7. 此后，主立即拔除恶行的根 \BibleRef{弟前 6:10}，说：\par

\Quote{不要在你们的腰带里预备金、银、铜钱，也不要有行囊、两件内衣、鞋，也不要带棍杖。}\par

主没有说\Quote{不要随身携带}，而是说\Quote{即使能从别人那里得到，也要逃避这恶疾}。你看，主这样安排成全了许多好事：首先，让门徒不受猜疑；其次，解除他们的一切挂虑，使他们能全心投入圣言；第三，教导他们主自己的大能。后来主果然向他们明说：\Quote{我派遣你们时，赤身露体、无鞋无杖，你们缺少过什么吗？}\par

主没有立刻说\Quote{不要预备}，而是先说\Quote{洁净癞病人，驱逐魔鬼}，然后说\Quote{什么也不要预备；白白得来的，白白施舍}；主这样安排，既顾及了他们的利益、信誉，也顾及了他们的能力。\par

但或许有人会说：其他诫命可以理解，但\Quote{不要有行囊、两件内衣、棍杖、鞋}，为什么这样吩咐？因为主要训练他们趋于完全；而且更早的时候，主已经不准他们为第二天忧虑。因为主要派遣他们到全世界作导师。因此，主使他们（可以这样说）甚至成为天使，解除他们一切世俗挂虑，使他们只专注于教导一事；或者不如说，连这事主也释放了他们，说：\Quote{不要思虑怎样说话，或说什么。}\par

这样，那看似非常沉重伤人的事，主却显明为极其轻省容易。因为没有什么比摆脱焦虑和挂虑更使人喜乐；尤其当他们这样得释放，一无所缺，天主与他们同在，成为他们的一切。\par

接着，恐怕他们问：\Quote{那么我们从哪里获得日用的食物呢？}主没有对他们说：\Quote{你们听过我从前说的：\InnerQuote{你们看天空的飞鸟。}}\BibleRef{玛 6:26}（因为他们还不能在实际中实践这一诫命），而是加上远不及此的话说：\Quote{工人配得他的食物}，表明他们必须由门徒供养；这样，他们既不会向受教者高傲，仿佛全给予而一无所取；后者也不会因被导师轻看而背离。\par

此后，恐怕他们说：\Quote{那么你命令我们靠乞讨为生吗？}并以此为耻，主表明这理所当然：既称他们为\Quote{工人}，又称所得的为\Quote{报酬}。主说：\Quote{不要以为因为劳苦在于言语，你们施予的福分微小；不，这事有许多辛劳；无论受教者给予什么，都不是他们施舍的恩赐，而是偿还的酬报：\InnerQuote{工人配得他的食物。}}主这样说，并非宣告宗徒的劳苦值那么多，绝非如此；而是既以此立法，叫他们不再求更多，又使给予者相信，他们所行的并非慷慨之举，而是债务。\par

8. \Quote{你们无论进哪一城或哪一村，要查问谁是可堪当的，就住在那里，直到离去。}\par

即是说：\Quote{我虽说了\InnerQuote{工人配得他的食物}，但这并不意味着我向你们敞开所有人的门；在此事上，我也要求你们十分谨慎。因为这将有益于你们的信誉，也有益于你们的生计。若那人堪当，他必然给你们食物；尤其当你们只求必须之物，不求更多。}\par

主不仅要求他们寻找堪当的人，还不可换家居住；这样，既不使接待者烦恼，也不使自己落得贪食放纵的名声。主藉着说\Quote{住在那里，直到离去}表明此意。从其他福音记载也可看出这一点。\BibleRef{路 10:7}\par

你看，主又藉此使他们尊荣，也使接待他们的人谨慎，表明后者在尊荣和利益上更是受益者。\par

然后，主继续同一话题说：\par

\Quote{你们进了那家，要请安。如果那家堪当，你们的平安就临到它；如果不堪当，你们的平安仍归于你们。}\BibleRef{玛 10:12}\par

你可看见祂如何毫不退让地推行自己的训令？这是极为恰当的。因为祂正培育他们成为敬畏天主的斗士，并向普世宣讲者。为此，祂教导他们操练节制，使他们成为可爱的人，祂说：\par

\BibleQuote{凡不接待你们，不听你们话的，你们离开那家或那城的时候，就把脚上的尘土跺下去。我实在告诉你们：在审判之日，索多玛和蛾摩拉地所受的惩罚比那城还要容易忍受。} \BibleRef{玛 10:14}\par

也就是说，祂说：\Quote{不要因为你们是老师，就等着别人向你们致敬，而要先表示敬意。}接着，祂暗示这不仅是问候，更是一种祝福，说：\BibleQuote{如果这家配得，你们的平安就必临到这家；}但如果他们无礼对待，首先的惩罚将是得不到你们平安的益处；其次是他们必受索多玛的罚。有人会说：\Quote{那么，他们的惩罚与我们何干？}你们将拥有配得之人的家。\par

但\BibleQuote{跺去脚上的尘土}是什么意思？要么表示他们从那些人那里一无所获，要么是向他们见证为了他们的缘故所走的漫长路途。\par

但请看，我恳求你们，祂甚至还没有把全部交托给他们。因为祂尚未赐予他们预知的能力，以分辨谁配得、谁不配得；而是吩咐他们查问，等待考验。那么，祂自己如何与一个税吏同住呢？因为那人因悔改而成为配得的。\par

请看，我恳求你们，当祂剥夺了他们的一切后，又赐给他们一切，允许他们住在那些成为门徒者的家中，并且空手进入。因为这样既使他们自己免于焦虑，又会使其他人相信，他们来纯粹是为了他们的得救：首先，因为他们不携带任何东西进入；其次，除了必需品，他们不要求更多；最后，他们不无区别地进入所有人的家。\par

因为祂不仅希望他们因奇迹而出众，更希望在奇迹之前，因他们自身的德行而出众。因为没有什么比摆脱多余之物，并尽可能地摆脱需求更能体现生活的严谨。这一点连假宗徒也知道。因此保禄也说：\BibleQuote{好叫他们在所夸的事上，也能像我们一样。} \BibleRef{格后 11:12}\par

但若我们在异乡，前往素不相识的人那里，除了当天的食物外不可寻求更多，那么在家时更当如此。\par

9. 这些话我们不仅应聆听，还应效仿。因为这不单是对宗徒们说的，也是对后来的圣人们说的。因此，我们要配得款待他们。因为这份平安的来临和离去，取决于接待者的心意。这不仅取决于教导者的勇敢宣讲，也取决于接受者的配得。\par

我们切不可轻视未能享有这份平安的损失，因为这份平安正是先知自古以来所宣告的：\BibleQuote{那传报和平福音者的脚步何等美丽！}然后为解释其价值，他补充说：\BibleQuote{那传报美善喜讯者的脚步！}\par

这份平安，基督也宣告为伟大的，当祂说：\BibleQuote{我把平安留给你们，我将我的平安赐给你们。} \BibleRef{若 14:27}我们应当做一切事以享受这份平安，无论是在家还是在教会。因为在教会中，主持者也赐予平安。我们现在所说的平安是那平安的预像。你们应当以全部的热忱来接受它，在心中先于实际的共融。因为若在共融之后不将平安传递是可厌的，那么驱逐宣布平安的人是何等可厌！\par

为了你们，司铎坐着；为了你们，教师站着，劳苦工作。若你们不予以足够的欢迎来聆听他，你们将有什么借口呢？因为教会确实是众人共同的家，当你们先占据了它，我们进入时严格遵循他们所显示的预像。为此，我们一进入就向众人宣告\Quote{平安}，按照那条法则。\par

因此，没有人应当漫不经心，没有人应当不留意，当司铎们进入并教导时；因为为此确实设有不小的惩罚。是的，我个人宁愿进入你们的家中千万次而遭拒绝，也不愿在这里说话时不被聆听。后者对我而言比前者更难忍受，因为这地方的尊贵程度更高；我们伟大的财富确实存放在这里，我们所有的希望都在这里。因为这里的一切，哪样不是伟大而可敬畏的呢？因此，这张桌子比其他的桌子更珍贵、更令人喜悦；这灯台也比那里的灯台更美好。那些凭着信德并适时用油傅抹而除去疾病的人都知道这一点。这箱子也比其他的箱子更好、更不可或缺；因为它关着的不是衣服，而是施舍；即使拥有它们的人很少。这里还有比其他的床更好的床；因为神圣圣经的安息比任何床榻更令人愉悦。\par

我们若在和睦上已达到卓越，便别无他家。而此言无所过重，以\BibleQuote{三千人}\BibleRef{宗 2:41}为证，又以\BibleQuote{五千人}\BibleRef{宗 4:4}为证，他们共有一家、一桌、一心；因为\BibleQuote{众信者都是一心一意}\BibleRef{宗 4:32}。但既远逊于彼等之德，散居各宅，则至少在此集会时，当竭力如此。因虽于他事皆贫穷匮乏，于此则富足。故至少在此处，当以爱心接纳我们进来。当我言\Quote{愿你们平安}，你答\Quote{也与你的心灵同在}时，莫仅以口言，亦当以心言；非仅唇舌，亦须悟性。但若在此处你言\Quote{平安也与你的心灵}，出门外你却为仇敌，唾骂诽谤，暗地予以无数羞辱，此为何等平安？\par

因为我固然，即使你千万次毁谤我，也以清洁的心、真诚的意念给你那平安，且任何时候都不会说你一句坏话；因我有慈父的心肠。我若有时责备你，是出于关心。但至于你，暗中抨击我，不在这主的家里接纳我，我怕你反而会增加我的忧愁；并非因你侮辱我，或因你驱逐我，而是因你拒绝我们的平安，从而为自己招致那严厉的惩罚。\par

因为我虽不跺下脚上的尘土，虽不转身离开，所警告的仍不改变。因为我常频繁向你们宣讲平安，并不会停止；如果你们除了侮辱我之外，仍不接纳我，我甚至也不跺下尘土；并非我不听命于我们的主，而是我热切渴望你们。此外，我根本没有为你们受苦；我既没有远道而来，也没有穿那样的衣服和自愿的贫穷而来（因此我们先自责），也并非赤足或只穿一件外衣；也许这也是你们未能尽本分的原因。然而，这对你们并非充分的辩护；因为我们的谴责虽更重，对你们却无解释余地。\par

10. 那时候，房屋就是圣堂，但现在圣堂却成了房屋。那时在房屋里可说不出任何世俗的事，如今在圣堂里却不能说出属灵的事；甚至在这里，你们也带来市场上的生意；当天主在讲话时，你们不静心聆听他的话语，反而引入相反的事，制造不和谐。我但愿是你们自己的事，但如今却是与你们毫无关系的事，你们既说也听。\par

为此我哀哭，且不会停止哀哭。因为我无力离开这房屋，但在此处我们必须存留，直到离开今生。因此，就如\Person{保禄}所吩咐的，请\BibleQuote{接纳我们}\BibleRef{格后 7:2}。因为他在那处所说的话，并非关乎饮食，而是关乎心志与心灵。我们向你们所求的也是这个，就是爱，即热烈而真挚的情感。但若你们连这个也不能忍受，至少爱你们自己，且放下当前的懈怠。这足以为我们的安慰，如果我们看到你们自勉并成为更好的人。我也将同样显示更大的爱，即使\Quote{我越发爱你们，却越少被爱}。\par

因为确实有许多事将我们联结在一起。同一张桌子摆在众人面前，同一位圣父生了我们，我们都是同一阵痛所生，同一饮品已赐给众人；更不止是同一饮品，而且是饮自同一杯。因为我们的圣父渴望引导我们进入友善的情感，也设计了这一点，即我们应当饮自同一杯；这属于深切之爱的事。\par

但\Quote{宗徒与我们之间无可比拟}。我也承认，且永不否认。因为我不仅说与他们自己，就连他们的影子我们也无法相比。\par

然而，尽管你们应尽的部分仍需完成。这并不会使你们蒙羞，反而使你们更有益处。因为即使对不配之人，你们仍显出如此的爱与服从，那时你们将获得更大的赏报。\par

因为我们所说的并非我们自己的话，因为你们在地上本无师傅；而是我们领受了什么，也给予什么；在给予时，我们向你们别无所求，只求被爱。即使我们连这个也不配，但借着爱你们，我们很快便会成为可配的。虽然我们受命不仅爱那爱我们的人，也爱仇敌。那么，谁如此心硬，谁如此野蛮，以至于受了这样的律法之后，竟厌恶甚至憎恨那爱他的人，尽管他可能充满无数罪恶？\par

我们已分享了属灵的筵席，让我们也分享属灵的爱。因为，如果强盗在共食盐之后便忘记了他们的本性，那持续分享主的圣体却连他们的温和也不仿效的我们，还有什么借口呢？然而，对许多人而言，不只是一张桌子，甚至同属一城便足以成为友谊；但我们，拥有同一座城、同一所屋、同一张桌子、同一条路、同一道门、同一根源、同一生命、同一头、同一牧者、同一君王、同一师傅、同一审判者、同一创造者、同一父亲，且一切都为共有；如果我们彼此分裂，那该受何等宽恕？\par

11. 但是，你们所追寻的，或许是奇迹，诸如他们进来时所行的：癞病人被洁净，魔鬼被驱逐，死人复活？不，这恰恰是你们高贵出身和你们爱情的伟大标志，即你们应无条件地相信天主。事实上，正是这一点，再加上另一件事，乃是天主使奇迹止息的原因。我的意思是，如果当奇迹不再施行时，那些在其他优势上——例如，或因智慧之言，或因虔诚的外表——自夸的人变得虚荣、骄傲、彼此分离，那么，如果奇迹仍然发生，怎能不产生激烈的分裂呢？我所说的并非纯属猜测，格林多人可以作证，他们正是因此而分裂成许多派别。\par

所以，你不要寻求神迹，而要寻求灵魂的健康。不要寻求看见一个死人复活；不，因为你已经知道整个世界正在复活。不要寻求看见一个盲人痊愈，但要看见所有人现在都恢复了那更美好、更有益的视觉；并且你也要学习贞洁地观看，改正你的眼睛。\par

因为事实上，如果我们都按应当的方式生活，行奇迹的人不会比我们更受外邦人的钦佩。至于神迹，它们往往带有纯粹幻想或另一种恶意猜疑的概念，尽管我们的神迹并非如此。但纯洁的生活不能接受任何这样的指责；是的，所有人的口都被德行的获得堵住了。\par

那么，让德行成为我们的追求：因为她的财富丰沛，在她内所行的奇迹何等伟大。她赐予真正的自由，并使这自由甚至在奴隶制中也能被辨识；不是从奴隶制中释放，而是在人继续为奴时，显出他们比自由人更尊贵；这远胜于给他们自由：不是使穷人变富，而是在他仍然贫穷时，显出他比富人更富有。\par

但是，如果你也想行奇迹，摆脱过犯，你就完全成就了它。是的，因为罪是一个大魔鬼，亲爱的；如果你消灭了它，你就行了比那些驱逐一万个魔鬼的人更伟大的事。你要听保禄如何说话，他怎样更看重德行而非奇迹。他说：\Quote{你们要切慕最好的恩赐；但我还要显示一条更卓越的道路。} \BibleRef{格前 12:31} 当他宣布这条\Quote{道路}时，他没有谈论复活死人，没有谈论洁净癞病人，也没有谈论任何这类事；而是以爱德代替了这一切。也要听基督的话，祂说：\Quote{不要因为魔鬼服从你们而欢喜，但要因你们的名字已录在天上而欢喜。} \BibleRef{路 10:20} 在此之前又说：\Quote{在那一天，许多人会向我说：\InnerQuote{主啊，主啊，我们不是因祢的名字说过预言，因祢的名字驱过魔鬼，因祢的名字行过许多大能的事吗？} 那时我要向他们声明说：\InnerQuote{我从来不认识你们。}} \BibleRef{玛 7:22} 当祂即将被钉十字架时，祂召叫了门徒，对他们说：\Quote{众人因此就知道你们是我的门徒，} 不是 \Quote{如果你们驱魔，} 而是 \Quote{如果你们彼此相爱。} \BibleRef{若 13:35}\par

因为，至于奇迹，它们常常在有益于他人的同时，伤害了拥有这能力的人，使他陷入骄傲和虚荣，或可能以其他方式；但在我们的善工中，没有任何这类猜疑的余地，它们反而有益于遵循它们的人以及许多其他人。\par

那么，让我们以极大的勤勉实行这些事。因为如果你从不人道变为施舍，你就伸出了那枯干的手。如果你离开剧院前往教会，你就治愈了那跛脚的。如果你从娼妓和不是自己的美色前收回目光，你就睁开了瞎眼。如果你不再唱撒旦的歌曲，而学会了属灵的圣咏，你原来哑巴，现在却说话了。\par

这些是最伟大的奇迹，这些是奇妙的神迹。如果我们继续行这些神迹，我们既会通过这些使自己成为伟大而可钦佩的人，也会赢得所有邪恶的人归向德行，并享受未来的生命；愿我们众人都因我们的主耶稣基督对人的恩宠与爱情而获得这生命，愿光荣与权能归于祂，直到永远。阿们。\par

\SourceRecord{Translated by George Prevost and revised by M.B. Riddle.；Nicene and Post-Nicene Fathers, First Series；Vol. 10.；Edited by Philip Schaff.；Buffalo, NY: Christian Literature Publishing Co.,；1888.；<http://www.newadvent.org/fathers/200132.htm>.}

% structure: p0001 source=h1 target=section reason=page-title

\section{玛窦福音讲道集 第三十三篇}

\BibleRef{玛 10:16}\par

\BibleQuote{看，我派遣你们好像羊进入狼群中，所以你们要机警如同蛇，纯朴如同鸽子。}\par

主先使他们对于必需的食物感到安心，为他们敞开所有人的家门，并使他们的进入带有一种令人起敬的外表，吩咐他们不要像流浪者和乞丐那样进来，而要远比接待他们的人更可敬（因为祂藉着说\BibleQuote{工人自当得他的工资}来表明这一点；并藉着命令他们查问谁是配得的，就住在那家，并吩咐他们向接待他们的人请安；又藉着用那些不可救药的灾祸来恐吓不接待他们的人）：我说，就这样消除了他们的焦虑，并以奇迹的彰显武装了他们，使他们好像全成了铁和金刚石，藉着免除他们一切世物，释放他们脱离一切暂时的挂虑：接下来祂又讲到将要临到他们身上的恶事；不仅是那些不久就要发生的，还有那些在漫长岁月中将要发生的；从一开始，甚至很久以前，就为他们对魔鬼的战争作准备。是的，由此获得了许多好处：首先，他们学到了祂预知的能力；其次，没有人会怀疑这些恶事是因他们导师的软弱而临到他们身上；第三，遭受这些事的人不会因它们出乎意料、有违期望而惊慌；第四，他们不会在十字架的时刻因听到这些事而困扰。事实上，他们当时正是如此受影响；那时祂也责备他们说：\BibleQuote{只因我给你们说了这话，你们就满心忧愁，并且没有人问我说：\InnerQuote{你往那里去？}} \BibleRef{若 16:6} 然而祂还没有说到关于祂自己的事，比如祂要被捆绑、鞭打和处死，以免也扰乱他们的心；但当时祂只是预先宣布他们自己要遇到的事。\par

然后，为叫他们知道这场战争体系是崭新的，战斗阵式是奇特的；正如祂差遣他们赤着脚，只穿一件外衣，不穿鞋，不带棍杖，不束腰带，不带行囊，并吩咐他们由接待他们的人供养；同样，祂在这里也没有停止祂的讲话，而是为表明祂那不可言喻的大能，说：即使这样出发，也要显出\Quote{羊}的温良，而且你们是要往\Quote{狼}那里去；不只是往狼那里，而是\Quote{往狼群中}。\par

祂不仅要求他们拥有羊的温良，还要有鸽子的纯朴。\Quote{因为这样我就要最好地显示我的大能，当羊胜过狼时，当羊处在狼群中，忍受千般咬啮时，不仅不被消灭，反而使狼改变——一件比杀死他们更伟大、更奇妙的事——改变他们的精神，矫正他们的心灵；而完成这事的是仅仅十二个人，当时整个世界都充满了狼。}\par

那么，我们这些行事相反的人——像狼一样扑向我们的敌人——应该羞愧。因为只要我们作羊，我们就得胜：纵然万匹狼围绕四周，我们仍能克服并获胜。但若我们变成狼，我们就被击败，因为我们的牧者不再帮助我们：祂不牧养狼，只牧养羊；祂会离开你，退去，因为你不容许祂的大能显扬。因为，正如你若在受虐待时显示温良，祂便将整个胜利归于自己；同样，你若还手并还击，你就遮盖了祂的胜利。\par

2. 但是请你思量，我恳求你，听到这些如此艰难而劳苦的命令的是谁：是那些胆怯无知、目不识丁、未受教育的人；是那些在各样事上都默默无闻、从未受过外邦人律法训练、不轻易出现在公共场所的人；是渔夫、税吏，充满无数缺点的人。因为如果这些事足以使那些高尚伟大的人都感到惶恐，那么它们又何尝不足以使那些全未经历、从未有过任何高尚想象的人沮丧惊惶呢？然而它们并未使他们沮丧。\par

\Quote{这很自然，} 也许有人会说；\Quote{因为祂赐给了他们洁净癞病人、驱逐魔鬼的能力。} 我的回答如下：不，这件事本身尤其使他们困惑，那就是尽管他们能复活死人，却要忍受这些无法忍受的恶事——包括审判、处决、所有人向他们发动的战争以及世人的普遍仇恨——并且当他们行奇蹟时，竟有如此恐怖等待着他们。\par

3. 那么，对于这一切，他们的安慰是什么呢？是派遣他们的那一位的大能。因此祂也将这一点放在一切之前，说：\BibleQuote{看，我派遣你们。} 这对你们的鼓励已经足够，这足以建立你们的信心，并叫你们不怕任何攻击你们的人。\par

你看见权能吗？你看见特权吗？你看见无敌的大能吗？现在祂的意思是这样：\Quote{不要烦恼}（祂如此说），\Quote{因我派遣你们到狼群中，却吩咐你们如同羊和鸽子。因我本可以反其道而行，使你们不经历任何可怕的事，也不像羊一样暴露于狼；我可以使你们比狮子更可怕；但这样是适宜的。这使你们更加光荣；这也彰显我的大能。}\par

祂也对保禄说了这话：\BibleQuote{我的恩宠够你用，因为我的德能在软弱中才全显出来。} \BibleRef{格后 12:9} \Quote{是我，请注意，我使你这样。} 因为祂说\BibleQuote{我派遣你们好像羊}，祂暗示了这一点。\Quote{所以不要沮丧，因为我知道，我确实知道，如此你们将比任何其他方式都更能坚不可摧。}\par

此后，为使他们在自己一方也做出一些贡献，且使一切似乎不全出于祂的恩宠，也不致被认为他们是随意、枉然地得冠冕，祂说：\BibleQuote{因此，你们要像蛇那样机敏，像鸽子那样温和。}\Quote{可是，}有人可能会说，\Quote{我们的智慧在如此大的危险中能有什么用处？不仅如此，当那么多波浪将我们淹没时，我们怎能拥有智慧呢？因为一只羊即使再聪明，当它处在狼群中，且是那么多狼时，它又能做什么呢？鸽子即使再温和，当那么多鹰攻击它时，又有什么益处呢？}在禽兽身上，确实毫无益处；但在你们身上，却大有助益。\par

但让我们看祂在此要求怎样的智慧。祂说，就是蛇的智慧。因为正如那动物放弃一切，即使身体必须被砍断，也不极力保护它，为要保全头部；同样，祂说，你也当放弃一切，唯独信德不可放弃；即使财产、身体、生命本身，都应舍弃。因为那是头和根；若它得以保全，即使你失去一切，也将以更大的荣光重新获得一切。\par

因此，祂既没有命令单纯做一个天真纯朴的人，也没有命令仅仅做聪明人，而是将两者结合，使它们成为德行；采纳蛇的智慧，好使我们不在要害处受伤；以及鸽子的温和，好使我们不报复得罪我们的人，也不向设圈套的人复仇；因为若不加上这个，智慧也是无用的。我问，有什么比这些命令更严格呢？怎么，遭受冤枉还不够吗？祂说，不，我甚至不许你们心生愤慨。因为这正是\Quote{鸽子}。好比将一根芦苇投入火中，命令它不被火焚烧，反而要灭火。\par

然而，我们不要烦恼；因为这些事情已经发生，并且得以成就，在实质行为中显明了，人变得像蛇那样机敏，像鸽子那样温和；他们并非具有另一种本性，而是与我们相同。\par

那么，不要让任何人认为祂的命令不可实行。因为祂比任何人都更知道事物的本性；祂知道凶狠不能由凶狠消灭，而是由温和消灭。如果你也愿意在人的实际行为中看到这个结果，请阅读\WorkTitle{宗徒大事录}这本书，你将看到多少次，当犹太百姓起来反对他们，磨砺牙齿时，这些人效法鸽子，以合宜的温顺回答，消除了他们的怒气，平息了他们的疯狂，挫败了他们的冲动。例如当他们说：\BibleQuote{我们不是严格命令你们，不可再因这名发言吗？}\BibleRef{宗 5:28} 虽然能行无数奇迹，他们却没有说或做任何粗暴的事，而是以极温顺的态度为自己辩护，说：\BibleQuote{听从你们比听从天主在你们眼中是否正当，请你们判断。}\BibleRef{宗 4:19}\par

你看到了鸽子的温和吗？请看蛇的智慧。\BibleQuote{因为我们不得不说我们所看见所听见的事。}\BibleRef{宗 4:20} 你看到我们必须在各方面完美，既不被危险所贬抑，也不被愤怒所激怒吗？\par

4. 因此，祂又说了这话：\BibleRef{玛 10:17}\par

\BibleQuote{你们要提防人，因为他们要把你们交给公议会，也要在他们的会堂里鞭打你们；并且你们要为我的缘故被带到总督和君王面前，对他们和外邦人作证。}\par

祂又这样预备他们保持警醒，在每种情况中，都使他们承受冤枉，而允许他人施加冤枉；为教导你们：胜利在于忍受困厄，祂光荣的战利品因此设立。因为祂根本没有说：\Quote{你们也要争战，抵抗那些困扰你们的人，}而只说：\Quote{你们将遭受极端的祸患。}\par

啊，说话者的能力何其大！听者的自制力何其大！因为我们非常有理由惊奇：他们听到这些话后，竟然没有立刻从祂身边逃走，尽管他们容易因每个声音而受惊，而且从未去过那鱼塘以外的更远地方；他们也没有思考，对自己说：\Quote{这一切之后，我们往哪里逃呢？法庭反对我们，君王反对我们，总督、犹太人的会堂、外邦人的民族、统治者、被统治者。}（因为借此祂不仅预先警告他们有关巴勒斯坦及其中的祸患，还揭示了全世界的战争，说：\BibleQuote{你们将站在君王和总督面前；}表明之后祂也要派遣他们到外邦人那里作先驱。）\Quote{你使我们与世界为敌，你武装了地上所有的居民、民族、暴君、君王来反对我们。}\par

而接下来的事又更可怕，因为人们要因我们而成为杀害兄弟、杀害子女、杀害父亲的凶手。\par

\BibleQuote{因为兄弟，}祂说，\BibleQuote{要把兄弟、父亲要把儿女置于死地；儿女要起来反对父母，并把他们置于死地。}\BibleRef{玛 10:21}\par

\Quote{那么，}有人会说，\Quote{其余的人如何能相信呢？当他们因我们而看到父亲杀害子女、兄弟杀害兄弟，一切都充满可憎之事时？}什么？难道他们不会像对待毁灭性的魔鬼一样，像对待该诅咒者和世界之害一样，从各处驱逐我们吗？当他们看到大地充满亲人的血和如此多的凶杀时？当然，我们带给人们家庭并赐予他们的和平是美好的（不是吗？），同时我们却用如此多的杀戮充满那些家庭。为什么，如果我们人数众多，而不是十二个；如果我们不是\Quote{不学无术的}，而是有智慧、精通修辞、口才出众；不仅如此，如果我们甚至是君王，拥有军队和大量财富，我们怎么能说服任何人呢？当我们挑起内战，甚至比内战更糟的战争时？为什么，即使我们不顾自己的安全，其他所有人中谁会留意我们呢？\par

但他们对这些事既不想也不说，也没有要求对祂的吩咐作任何解释，只是简单地服从并遵从。这不仅是出于他们自身的德行，也出于他们导师的智慧。因为请看，祂如何为每一件可怕的事附上鼓励；正如对于那些不接受他们的人，祂说：\Quote{在审判之日，索多玛和哈摩辣地所受的惩罚比那城还要容易忍受。}同样在这里，当祂说：\Quote{你们将被带到总督和君王面前}时，祂又加上：\Quote{为了我的缘故，为向他们和外邦人作证。}这并非微不足道的安慰，因为他们为基督并为使外邦人信服而受苦。这样，纵然无人留意，天主却显明无处不在，行祂自己的工。如今这些事给他们安慰，并非因为他们渴望惩罚别人，而是他们可以有信心，确信那位预言并预知这些事的天主必常与他们同在；而且他们遭受这一切，不是因为他们是恶人和祸害。\par

除此之外，祂又加上另一个并非微不足道的安慰，说：\par

\BibleQuote{但当他们交付你们时，不要思虑如何或说什么，因为在那时刻，必会赐给你们当说的话。因为不是你们说话，而是你们圣父的圣神在你们内说话。}\BibleRef{玛 10:19}\par

恐怕他们会说：\Quote{当这些事发生时，我们怎能说服人呢？}祂也吩咐他们对辩护充满信心。确实，在别处祂说：\BibleQuote{我要赐给你们口才和智慧，}\BibleRef{路 21:15}但在这里，祂说：\BibleQuote{是你们圣父的圣神在你们内说话，}将他们提升到先知的尊严。因此，当祂提及所赐予的能力后，祂也加上了恐怖、凶杀和杀戮。\par

\BibleQuote{因为兄弟要把兄弟，}祂说，\BibleQuote{交付死亡，父亲交付孩子，孩子要起来反对父母，并害死他们。}\BibleRef{玛 10:21}\par

\BibleQuote{你们要为所有的人所恨。}而这里安慰又近在眼前，因为，\BibleQuote{为了我的名字的缘故，}祂说，\BibleQuote{你们要忍受这些事。}接着又加上另一个：\BibleQuote{但那坚持到底的，必然得救。}\BibleRef{玛 10:22}\par

这些事从另一个角度也同样足以振奋他们的精神；因为无论如何，他们的福音能力将如此高涨，以至于本性被轻视、亲属被弃绝、圣言被置于一切之上，有力地驱逐一切。因为如果没有任何本性的暴政足以抵挡你们的话，反而被瓦解和践踏，那么还有什么能胜过你们呢？然而，并非你们的生命将安全，因为这些事将要发生；相反，你们将拥有全世界的居民作为共同的敌人和仇敌。\par

5. 如今\Person{柏拉图}在哪里？\Person{毕达哥拉斯}在哪里？斯多葛派的漫长谱系呢？因为第一个，在享受了极大的荣誉之后，实际上被驳斥到甚至被卖出国境，没有达到任何目标，甚至连一个暴君也没有成功；是的，他背叛了门徒，悲惨地结束了生命。而犬儒学派，不过是污秽，都像梦和影子一样过去了。然而，确实没有这样的事发生在他们身上，反而他们因外邦哲学而被视为光荣，\Place{雅典}人公开纪念\Person{柏拉图}的书信，这些信是\Person{狄翁}寄给他的；他们安闲地度过所有时间，并拥有不少财富。例如，\Person{亚里斯提卜}常购买昂贵的娼妓；另一个立了遗嘱，留下非同寻常的遗产；另一个，当他的门徒像桥一样躺下时，他从他们身上走过；而\Place{西诺普}的那位，据说甚至在市场里做出不体面的事。\par

是的，这些是他们的光荣之事。但这里没有这样的事，而是严格的节制、完美的端庄、一场为真理和虔诚反对全世界的战争、每天被杀，直到此后他们获得荣耀的胜利。\par

但有人会说，他们中间也有一些善于作战的人，如\Person{地米斯托克利}、\Person{伯里克利}。但这些与渔夫们的事迹相比，也只是儿童玩具。因为你能说什么呢？他使\Place{雅典}人登上船只，当\Person{薛西斯}进军\Place{希腊}时？为什么在这个事例中，当不是\Person{薛西斯}进军，而是魔鬼与整个世界，以及无数恶灵攻击这十二个人时，不仅在一个危机中，而是在他们整个一生中，他们得胜并征服了；而且真正奇妙的是，不是通过杀死敌人，而是通过转化和改造他们。\par

为此，你尤其应当留意整个过程中，他们没有杀害或毁灭那些图谋反对他们的人，反而在发现他们如同魔鬼般邪恶时，使他们成为天使的竞争者，将人性从这邪恶暴政中解放出来。至于那些扰乱万物的可憎恶魔，他们将其驱逐出市场、房屋，甚至旷野之中。对此，隐修士的歌咏团可以作证，他们到处安置了隐修士，不仅洁净了可居住之地，连不可居住之地也清理干净。更奇妙的是，他们并非在公平战斗中完成这一切，而是在忍受苦难中成就了一切。因为人们实际上把他们——十二个无学问的人——围在中间，捆绑、鞭打、拖来拖去，却无法堵住他们的口；正如不能束缚阳光一样，他们的舌头也无法被束缚。原因是，\Quote{不是他们}自己\Quote{说话，}而是圣神的能力。例如，\Person{保禄}就这样战胜了\Person{阿格黎帕}和\Person{尼禄}，后者在邪恶上超越众人。\BibleQuote{主}他说，\BibleQuote{站在我身旁，坚固了我，并把我从狮子口中救了出来。}\BibleRef{弟后 4:17}\par

但是，你也要钦佩他们：当对他们说\Quote{不要忧虑}时，他们仍然相信并接受了，没有恐惧能惊吓他们。如果你说，祂通过说\Quote{将是你们父的圣神在说话}给了他们足够的鼓励；即便如此，我最钦佩的是，他们没有怀疑，也没有寻求从危险中解脱；而且，他们不是两三年要忍受这些事，而是终生如此。因为\Quote{坚持到底的，必然得救}这句话正暗示了这一点。\par

因为祂的旨意是，不仅祂的部分要贡献出来，而且善行也要由他们来完成。例如，注意，从一开始，一部分是祂的，一部分是祂门徒的。行奇迹是祂的，但什么都不准备是他们的。同样，打开所有人的家门，是来自上主的恩宠；但只求所需，不多求，是他们自己的克己。\Quote{工人配得他的工资。}赐予平安是天主的恩赐，寻找配得之人，不加区别地进入所有人，是他们自己的命令。同样，惩罚那些不接受他们的人，是祂的；但温和地离开他们，不辱骂或侮辱他们，是宗徒们的温良。赐予圣神并使他们不用忧虑，是那派遣他们者；但成为绵羊和鸽子，并高贵地承受一切，是他们的平静和明智。被憎恨而不沮丧，并忍受，是他们自己的；拯救那些忍受的人，是那派遣他们者。\par

因此祂也说：\Quote{坚持到底的，必然得救。}这是因为大多数人起初确实热情，但后来就衰颓，所以祂说：\Quote{我要求的是终点。}种子起初茂盛，但稍后枯萎，有什么用呢？因此，祂要求他们的是持续的忍耐。我的意思是，免得有人说，祂独自成就了一切，他们之所以如此，是因为没有忍受什么不可忍受的事，所以祂对他们说：\Quote{你们也需要忍耐。因为我虽会救你们脱离最初的危险，但我正为你们保留更严重的危险，此后还会有别的危险接踵而至；只要你们还有一口气，就不会停止为你们设下陷阱。}因为祂在说\Quote{但坚持到底的，必然得救}时暗示了这一点。\par

为此，尽管祂说：\Quote{不要忧虑你们要说什么；}但别处祂说：\BibleQuote{你们要时常准备好，对每一个询问你们心中所怀希望的理由的人做出答复。}\BibleRef{伯前 3:15}\par

因为这确实是一件非常伟大的事：一个忙于湖泊、皮革和收税的人，当暴君坐在宝座上，总督和卫兵站在他们旁边，刀剑出鞘，所有人站在他们一边时，独自进入，被捆绑，低着头，却能够开口。因为他们既不允许他们为自己的教义辩护或辩护，反而折磨他们至死，视他们为世界的公敌。因为\BibleQuote{他们}据说\BibleQuote{那些搅乱天下的人，也到这里来了。}和\BibleQuote{他们宣传与凯撒谕令相违背的事，说耶稣基督是王。}\BibleRef{宗 17:6}到处法庭都充斥着这样的猜疑，需要许多来自上天的帮助，以表明他们所传教义的真实性，以及他们没有违反普通法律；这样，他们既不会在热切谈论教义时被怀疑推翻法律，也不会在热切表明他们没有推翻普通政府时损害教义的完整性：所有这些你都会看到在\Person{伯多禄}、\Person{保禄}以及所有其他人身上都得到了周全成就。是的，他们被全世界指控为叛乱者、革新者和革命者；然而他们既抵挡了这种印象，又为自己披上了相反的形象，所有人都赞美他们为救主、守护者和恩人。而这一切都是通过他们巨大的耐心达成的。因此\Person{保禄}也说：\Quote{我天天死。}他继续\Quote{身处危险}直到终点。\par

6. 那么，我们拥有如此崇高的榜样，却在平安中沉溺于柔靡与松懈，我们该配得什么？没有敌人（这太明显了）我们就被击杀；无人追赶我们却发昏，在平安中我们被要求得救，连这一点我们也做不到。而他们，当世界着火，火堆在全地被点燃，进入并从火焰中、从里面抢出那些正在燃烧的人；但你甚至不能保全自己。\par

那么，我们还有什么信心？什么恩惠？没有鞭打，没有监狱，没有统治者，没有会堂，也没有任何这类东西加在我们身上；恰恰相反，我们统治并得胜。因为国王们都是虔诚的，基督徒有诸多尊荣、优先地位、尊贵和豁免，即使如此我们仍不能得胜。而他们，每日被带去处决，无论是教师还是门徒，忍受无数鞭打和不断烙印，比住在乐园的人更享安乐；我们这些从未受过丝毫苦楚、连梦里也没有的人，却比任何蜡还软。\Quote{但他们}，有人会说，\Quote{行了奇迹。}这难道使他们免于鞭打吗？使他们免受迫害吗？不，奇怪的是，他们常常甚至从他们帮助的人手中遭受这些，却毫不慌乱，以善报恶只换来恶。但你若对任何人施予一点小惠，然后被回报任何不愉快之事，就慌乱、烦恼，并后悔自己所行之事。\par

如果现在发生，正如我祈求不会发生、也永不发生，一场反对教会的战争和迫害，试想嘲笑何等巨大，责难何等严厉。这很自然；因为当没有人角力学校操练，如何在比赛中脱颖而出？哪个冠军，不习惯教练，当被奥林匹克竞赛召唤时，能够对抗对手展现出伟大崇高之事？我们岂不应当每天角力、搏斗和奔跑吗？你们岂不见那些被称为\OriginalTerm{五项全能运动员}{Pentathli}的人，在没有对手时，如何装满一袋沙，挂起来尽力击打？而那些更年轻的人，则在同伴身上练习对敌搏斗吗？\par

你也当效法这些，操练克己的角力。因为确实有许多人引你发怒，煽动情欲，燃起大火。因此，要抵挡你的情欲，高尚地忍受心灵中痛苦，使你也能够忍受身体的痛苦。\par

7. 因为蒙福的\Person{约伯}，若没有在争战之前好好操练自己，就不会在争战中如此闪耀。除非他操练了脱离一切沮丧，他会在儿女死时说些鲁莽的话。但事实上，他抵挡了所有攻击：家产破产，如此巨大财富的毁灭；儿女丧失，妻子的同情，身体的灾祸，朋友的责难，仆人的辱骂。\par

若你也要看他的操练方式，请听他怎样说，他怎样轻视财富：他说：\Quote{我若因财富丰盈而欢喜，若把黄金堆成堆，若依靠宝石。}因此，当财富被夺去时他也不慌乱，因为当财富在场时他并不贪恋。\par

请听他如何管理关乎儿女的事务，不像我们那样失之过软，而是要求他们凡事谨慎。因为他甚至为他们的隐罪献祭，试想他对那些明显的罪是何等严格的审判者。\BibleRef{约 1:5}\par

若你也愿听他追求贞洁的努力，请听他所说的：\Quote{我与眼立了约，决不顾恋处女。}\BibleRef{约 31:1}为此，他的妻子并未使他丧志，因为他在此以前就爱她，然而并非过分，而是如同对妻子所当爱的。\par

因此我甚至感到惊奇，魔鬼明明知道约伯先前的操练，怎么竟会想到挑起争战？那么它是怎么想到的呢？那\OriginalTerm{怪物}{monster}是邪恶的，从不绝望：这成了对我们极大的谴责，因为它从不放弃毁灭我们的希望，而我们却对自己的救恩绝望。\par

至于身体的残缺和凌辱，请注意他如何操练自己。何以如此？因为他自己从未经历过任何这类事，而是一直生活在财富、奢华和一切荣华中，他常常一件一件地揣摩别人的灾祸。当他这样说时，他宣告了这点：\Quote{我所恐惧的临到了我；我所惧怕的来到了我身上。}\BibleRef{约 3:25}又说：\Quote{我哭泣每一个无助的人，我见到困苦的人就叹息。}\BibleRef{约 30:25}\par

因此，所发生的一切没有一样使他慌乱，那些巨大且难以忍受的灾祸也没有一样。因为我叫你们不看他的家产破产，不看儿女丧失，不看那不治的灾病，也不看他妻子对他的计谋；而要看那些远比这些更痛苦的事。\par

\Quote{那么，}有人说，\Quote{约伯所受的比这些更苦吗？}因为从他的历史中，我们知道的不过这些。因为我们睡着了，我们没有学到，因为那急切并仔细寻找珍珠的人，必会知道比这些更多的细节。因为更痛苦、更易引起更大困惑的事是别的。\par

首先，他对天国和复活一无所知，这也正是他所哀叹的。他说：\Quote{因为我不会永远活着，好让我长久忍耐。}其次，他自知有许多善行。第三，他自知没有邪恶之事。第四，他认为这是从天主手中所受的；或者若来自魔鬼，这也足以令他跌倒。第五，他听到朋友们指责他邪恶，他们说：\Quote{你受的鞭打，还没有按你的罪孽该受的。}\BibleRef{约 11:6}第六，他看见作恶的人反倒亨通，并在上面欺压他。第七，他没有其他可以仰望的、曾遭受同样苦难的人。\par

8. 若你愿知道这些事何等重大，请想想我们现今的状况。因为如今，当我们期待天国、盼望复活和说不尽的福乐，自知无数恶行，又拥有许多榜样，并分享如此高超的哲思时；倘若有人丢失了一点金子，且往往是强取之后，便认为生无可恋——他们没有妻子加害，没有丧失儿女，没有受朋友责备，没有被仆人侮辱，反而有许多人用言语和行为安慰他们——那么，那位看见自己诚实劳动所得无故被夺，然后遭遇无数试炼如冰雹般降下，却始终不动摇，并为这一切向主献上应有感恩的人，岂不该得到何等光荣的冠冕？\par

即使没有人说出其他讥讽的话，单凭他妻子的话，就足以完全动摇磐石。请看她的狡猾。她闭口不谈金钱、骆驼、羊群和牛群（因为她知道丈夫在这些事上的自制），却谈论比这一切更难承受的事，即他们的孩子；而且她加深了悲剧，并加上自己的影响力。\par

倘若人在富贵无虞时，尚且在诸多事上多次听从妇人；那么，该是何等勇敢的灵魂，才能力拒她用如此强大的武器来攻击，并将两种最暴虐的激情——欲望与怜悯——践踏在脚下！然而，许多人战胜了欲望，却屈服于怜悯。例如，高贵的\Person{若瑟}制服了最暴虐的快乐，拒绝了那用无数诡计纠缠他的异族妇人；但他却未能忍住眼泪，当他看见曾亏待他的兄弟们时，他被那激情燃遍全身，迅速卸下面具，揭穿了他一直扮演的角色。但这里，她首先是他的妻子，她的话语凄切，时机对她有利，加上他的创伤和鞭伤，以及那无数灾祸的波浪；若非如此，又怎能不恰当地宣告那在如此风暴中无动于衷的灵魂比任何金刚石更坚固呢？\par

请容我直言：即使是宗徒们，即使不逊于这位有福之人，也至少不比他更伟大。因为他们因基督而受苦得到了安慰；这剂良药每日足以使他们康复，以致主处处说：\Quote{为我，因我的缘故，}以及\Quote{他们若称家主为贝耳则步。}\BibleRef{玛 10:25}但他却缺乏这样的鼓励，也缺乏来自奇迹和恩宠的鼓励；因为他没有如此大的圣神之能。\par

更伟大的是，他是在极度优渥中长大的，并非出自渔夫、税吏和俭朴度日之人，而是在享受如此荣耀之后，遭受了他所遭受的一切。在宗徒们身上最难忍受的，他也同样经历了：被朋友、仆人、敌人以及曾受他恩惠的人所憎恨；而那神圣的锚、无浪的港湾，即对宗徒们所说的\Quote{因我的缘故}，他却没有见到。\par

我再次钦佩那三青年，因为他们敢入火窑，敢于对抗暴君。但请听他们说：\Quote{我们不事奉你的神，也不敬拜你所立的金像。}\BibleRef{达 3:18}这对他们是最大的鼓励：确知他们为天主而受一切苦。但\Person{约伯}却不知这一切是较量与搏斗；因为若他知道，就不会对所发生的事感到痛苦。至少，当他听到\Quote{你以为我向你宣告我的神谕是徒然的吗？或者是为了使你被称为义人？}时，请看，他如何立即因一句话而重获呼吸，如何自谦，如何认为自己连所受的苦也当不得，如此说：\Quote{我何必再争辩，因受主的训诫和责备，听见这些事，我算什么？}又说：\Quote{我以前只风闻有你，现在亲眼看见了你；因此我厌恶自己，在尘土和炉灰中懊悔。}\par

这位在法律与恩宠之前之人的如此刚毅与节制，也让我们这些在法律与恩宠之后的人来效法，好使我们也能与他同享永恒的帐幕。愿我们众人赖我们的主耶稣基督的恩宠与慈爱，得以到达那里；愿荣耀与胜利归于他，至于无穷之世。阿们。\par

\SourceRecord{Translated by George Prevost and revised by M.B. Riddle.；Nicene and Post-Nicene Fathers, First Series；Vol. 10.；Edited by Philip Schaff.；Buffalo, NY: Christian Literature Publishing Co.,；1888.；<http://www.newadvent.org/fathers/200133.htm>.}

% structure: p0001 source=h1 target=section reason=page-title

\section{玛窦福音第三十四篇讲道}

\Emph{\BibleRef{玛 10:23}}\par

\Quote{但是，几时人们在这城迫害你们，你们就逃往另一城去；我实在告诉你们：你们还没有走完以色列的城邑，人子就来了。}\par

\Quote{但是，几时人们在这城迫害你们，你们就逃往另一城去；我实在告诉你们：你们还没有走完以色列的城邑，人子就来了。} 祂在讲了那些可怕又可怖的事，足以熔化最坚硬的心——这些事是在祂的十字架、复活和升天之后要临到他们身上的——之后，又将祂的言论转向更为平和的内容，让祂所训练的人得以喘息，并给予他们充分的安全。因为祂根本没有命令他们在受迫害时与敌人对抗，而是逃跑。也就是说，既然现在还只是开端和序幕，祂便将祂的言论转向非常迁就的态度。因为祂现在所说的并非后来的迫害，而是十字架和受难之前的迫害。祂用这话表明：\InnerQuote{你们还没有走完以色列的城邑，人子就来了。} 也就是说，免得他们说：\InnerQuote{那么，如果我们受迫害而逃跑，在那里他们又追上我们，赶走我们，该怎么办呢？} 为消除这种恐惧，祂说：\InnerQuote{你们还没有走遍巴勒斯坦，我就立刻来到你们跟前。}\par

请注意，祂在这里并非消除恐怖，而是在他们遇险时与他们同在。因为祂没有说：\Quote{我要把你们救出来，结束迫害；} 而是什么？\Quote{你们还没有走完以色列的城邑，人子就来了。} 是的，因为只要见到祂，就足以安慰他们了。\par

但请你观察，我恳求你，祂并非在每一个场合都把一切留给恩宠，而是也要求他们有所贡献。\Quote{因为如果你们害怕，} 祂说，\Quote{就逃跑，} 因为祂用这话表明：\Quote{你们就逃} 和 \Quote{不要怕。} \BibleRef{玛 10:26} 祂并没有命令他们一开始就逃跑，而是在受迫害时才撤离；祂允许他们走的距离也不远，只是让他们走遍以色列的城邑。\par

然后，祂又训练他们另一种自制：首先，消除他们对食物的所有挂虑；其次，消除对危险的所有恐惧；现在，则消除对毁谤的恐惧。因为祂藉着说：\Quote{工人配得他的工资} \BibleRef{玛 10:10} 并表明会有许多人接待他们，解除了他们第一种焦虑；又藉着说：\Quote{不要思虑怎样或说什么，} 以及 \Quote{凡坚持到底的，必要得救，} 解除了他们对危险的忧虑。\par

然而，此外，他们很可能也会招致恶名，这似乎比一切都更难忍受；请看，祂甚至在这种情况下也安慰他们，从祂自己和一切关于祂所说的话中汲取鼓励；没有别的东西能与之相比。因为正如祂在另一处所说：\Quote{你们要为所有的人所恨，} 并加上\Quote{因我的名字，} 同样地，在此处也是如此。\par

祂又用另一种方式减轻了这种痛苦，将一个新的主题与先前那个联系起来。那么，这是什么主题呢？\par

\Quote{门徒，} 祂说，\Quote{不能高于他的老师，仆人也不能高于他的主人。门徒能像他的老师一样，仆人能像他的主人一样，就够了。如果人家称家主为贝尔则步，何况他的家人呢？所以不要怕他们。} \BibleRef{玛 10:24}\par

请看，祂如何显明自己是万有的主、天主和创造者。那么怎样呢？难道门徒高于他的老师，或仆人高于他的主人吗？只要他是门徒或仆人，就那尊荣的本性而言，他并非如此。因为请不必在这里向我讲述罕见的事例，而要从事物的普遍规律出发。祂并没有说：\Quote{何况祂的仆人，} 而是说 \Quote{祂的家人，} 以表明祂感到他们与祂多么亲近。在另一处祂也说：\Quote{今后我不再称你们为仆人，你们是我的朋友。} 祂并没有只说：他们侮辱了家主并毁谤了祂，还陈述了侮辱的具体形式，即他们 \Quote{称祂为贝尔则步。}\par

然后，祂又给予了另一种安慰，并不亚于这个：因为这一种固然是最大的，但由于那些尚未严格度生的人还需要另一种能够特别有力地使他们振奋的安慰，祂也陈述了这种安慰。这段话在形式上似乎是一句普遍性的断言，但并非对所有事物，而只是针对眼前的事物。因为祂说的是什么？\par

\Quote{没有掩盖的事，将来不被揭露的；也没有隐藏的事，将来不被知道的。} \BibleRef{玛 10:16} 现在祂所说的好像是这样：你们受鼓励，因为我也与你们分享了同样的羞辱；我是你们的主人和老师。但如果你们听到这些话仍然忧伤，那么就再考虑另一件事：你们很快就要从这种怀疑中解脱出来。因为你们为什么忧伤呢？因为他们称你们为巫师和骗子？但稍等片刻，所有的人都会称你们为世界的救主和恩人。是的，因为时间会揭示一切隐藏的事；它既要驳斥他们的诬告，也要显明你们的德行。因为当事实显示出你们是救主、恩人和一切德行的榜样时，人们将不会在意他们的话，而会关注事情的真相；他们将显为诬告者、说谎者和诽谤者，而你们将比太阳更明亮，时间的长度揭示并宣告你们，发出比号角更清晰的声音，使所有的人都成为你们德行的见证人。因此，不要让现在所说的话使你们自卑，而要让未来之事的希望提升你们。因为与你们有关的事是不可能隐藏的。\par

2. 于是，在使他们摆脱了一切苦恼、恐惧、焦虑，并使他们超越人的责难之后，然后，也只有到那时，祂才适时地向他们谈论勇敢宣讲的事。\par

\Quote{我在暗中告诉你们的，你们要在光天化日之下说出来；你们在耳边听到的，要在屋顶上宣讲出来。} \BibleRef{玛 10:27}\par

然而，当祂说这些话时，并非全然在暗处，也不是在耳旁低语；祂用了强烈的修辞这样说。这是因为祂只与他们独自交谈，且在\Place{巴勒斯坦}的一个小角落，因此祂说\Quote{在暗处}和\Quote{在耳旁}，使那即将赐予他们的宣讲的胆量，与当时谈话的口吻形成对比。\Quote{因为你们不是向一个、两个或三个城市宣讲，而是向全世界，}祂说，\Quote{走遍陆地与海洋，有人居住之地和旷野；对君王与民族，对哲学家与演说家，直言无讳，完全自由地宣讲。}因此，祂说\Quote{在房顶上}和\Quote{在光明中}，毫不退缩，完全自由。\par

为什么祂不只说\Quote{在房顶上宣讲}和\Quote{在光明中说话}，还加上\Quote{我在暗中告诉你们的}和\Quote{你们在耳中所听的}？这是为了提升他们的精神。正如祂曾说：\Quote{信我的人，我所做的工程，他也要做，并且还要做比这更大的工程；}\BibleRef{若 14:12}同样在这里，为了表明祂要藉着他们做一切，并且超过祂自己所做的，祂加了这句话。\Quote{因为开端，}祂说，\Quote{我已赐下，并作了前奏；但更大部分，我愿藉着你们完成。}这不只是命令，也是预先宣告将要发生的事，并用祂的话鼓励他们，暗示他们将战胜一切，同时也悄悄消除他们对恶评的忧虑。因为正如这教义在隐藏一段时间后将遍及一切，同样，犹太人的恶意猜疑也将迅速消失。\par

然后，因为祂已将他们高举，祂又警告他们危险，给他们心智添上翅膀，将他们高高提升于万有之上。因为祂说什么？\Quote{你们不要害怕那些杀害身体，却不能杀害灵魂的。}\BibleRef{玛 10:28}你看见祂如何将他们置于万有之上，劝他们不仅蔑视焦虑、诽谤、危险和阴谋，甚至那被视为最可怕的死亡？不仅是死亡，而且是暴力死亡？祂没有说\Quote{你们将被杀}，而是以与祂相称的尊严，将这话摆在他们面前，说：\Quote{你们不要害怕那些杀害身体，却不能杀害灵魂的；但更要害怕那能将灵魂和身体都毁灭在地狱里的。}将论证转向其对立面，正如祂一贯所做的。为什么？祂说，你们害怕死亡吗？因此你们迟迟不宣讲吗？不，正因为如此，我命令你们宣讲，好让你们害怕死亡：因为这会使你们从真正的死亡中得救。即使他们杀了你们，他们也不能胜过那更好的部分，尽管他们千方百计。因此祂没有说\Quote{不杀灵魂的}，而是说\Quote{不能杀灵魂的}。因为他们即使想，也不能得逞。所以，如果你们害怕惩罚，就害怕那更严厉得多的惩罚吧。\par

你看见祂如何再次不承诺他们免于死亡，而是允许他们死亡，赐予他们比不准他们受苦更大的恩赐？因为免于死亡远不如劝人蔑视死亡伟大。你现在看到，祂不是将他们推入危险，而是将他们置于危险之上，并在短短一句话中，将他们灵魂不朽的教义铭刻在他们的心中；用两三个词植入了拯救的教义后，祂又用其他考虑来安慰他们。\par

因此，免得他们以为，当他们被杀被害时，是像被遗弃的人一样受苦，祂再次引入天主眷顾的论证，这样说：\Quote{两只麻雀不是卖一个铜钱吗？但若没有你们父的许可，一只也不会掉在网罗里。就是你们的头发也都数过了。}\BibleRef{玛 10:29}\Quote{还有什么比它们更微贱的呢？}祂说，\Quote{然而，连这些也没有一只会在天主不知情下被取去。}因为祂的意思不是\Quote{因祂的作为它们坠落}，因为这不符合天主；而是\Quote{没有一件事瞒得过祂。}那么，如果祂对我们的遭遇一无所不知，并且爱我们胜过父亲，如此爱我们，以至于数过我们的头发，我们就不该害怕。祂说这话，不是因为天主数我们的头发，而是为了表明祂的完全知识和对他们的伟大眷顾。因此，如果祂既知道所发生的一切，又能拯救你们，也愿意拯救，那么无论你们可能遭受什么，都不要以为你们是像被遗弃的人一样受苦。因为祂的旨意不是要救你们脱离恐惧，而是劝你们蔑视恐惧，因为这比任何事情都更能使人脱离恐惧。\par

3. \Quote{所以，你们不要害怕；你们比许多麻雀还要贵重。}\BibleRef{玛 10:31}你看见恐惧已经胜过他们了吗？是的，因为祂知道心里的秘密，所以祂又说：\Quote{因此，不要害怕他们。}因为即使他们得胜，也只是战胜那低劣的部分，我指的是身体；即使他们不杀它，本性也定会将它取去而逝去。所以连这件事也不取决于他们，而是人从本性得来的。如果你们害怕这个，就更应该害怕那更大的，畏惧\Quote{那能将灵魂和身体都毁灭在地狱里的。}祂现在没有明说那就是祂自己\Quote{能将灵魂和身体毁灭的}，但祂先前曾宣告自己是审判者，这已是显明的了。\par

但现在情况相反：我们不怕那能毁灭灵魂（即惩罚灵魂）的祂，却对那些杀害身体的人战栗。然而，既然祂连同灵魂也惩罚身体，他们甚至连身体也不能惩戒，更不用说灵魂了；即使他们严厉地惩戒它，反而使它更加光荣。\par

你看见祂如何表示这些争战是容易的吗？因为事实上，死亡曾极大地搅扰他们的灵魂，一时激起恐惧，因为那时它尚未被轻易征服，那些将要蔑视它的人也尚未领受圣神的恩宠。\par

你看，祂驱散了那搅扰他们灵魂的恐惧和忧苦；接着祂又鼓励他们，以恐惧驱除恐惧；不单以恐惧，还以巨大赏报的望德；祂满有权威地发出威胁，从两方面敦促他们为真理大胆宣讲；并接着说：\par

\Quote{凡在人面前宣认我的，我也必在我天上的圣父面前宣认他；凡在人面前否认我的，我也必在我天上的圣父面前否认他。} \BibleRef{玛 10:32}\par

这样，祂不仅从善的一面，也从相反的一面敦促他们；并以那可怕的部分作为结束。\par

要留意祂何等精细：祂说的不是\Quote{我}，而是\Quote{在我内}，暗示宣认者并非凭自己的能力，而是靠从上而来的恩宠的助佑，才作出自己的宣认。但对于否认的人，祂说的不是\Quote{在我内}，而是\Quote{我}；因为那人既失去了恩赐，他的否认便随之而来。\par

\Quote{那么，他为何受责难，}有人会说，\Quote{如果他被舍弃，他否认了呢？}\par

\Quote{你得着了好处？}祂说，\Quote{因你首先在这里宣认了我？我也要在你身上得着好处，将更大、且不可言喻更大之物赐给你；因为我要在那里宣认你。}难道你没看见，善与恶都是在那里施报的吗？那么，你为何急忙催促自己？为何在此世寻求赏报，你这\Quote{因望德而得救}的人？所以，无论你行了什么善事而未在此世得到报酬，不要烦恼（因为那赏报在将来会连同增加等待你）；或者你行了什么恶事而未受罚，不要心安；因为那里必有惩罚迎候你，除非你悔改归正。\par

但若你不信，可由此世之事推测将来之事。为何？如果在争战之时，那些宣认的人如此荣耀，试想他们在冠冕之时将如何。如果仇敌在此鼓掌，那最慈爱的父亲怎能不赞赏并宣告你？是的，那时我们必得着善的赏赐和恶的惩罚。因此，那些否认的人将在今世与来世都受损害：在此世，他们虽永不死，却怀着恶的良心而活，确实已死；在那里，则受最后的刑罚。而另一等人则在此世与来世都得利益：在此世，他们从死亡中获益，因而比活着的人更荣耀；在那里，则享受那些不可言喻的福乐。\par

天主绝不只急于惩罚，也急于施恩；而且后者甚于前者。但祂为何只一次设置赏报，却两次设置惩罚？祂知道这样更能纠正我们。为此，当祂说了\Quote{要畏惧那能把灵魂和身体都投入地狱毁灭的祂}之后，祂又说\Quote{我也要否认他。}保禄也是如此，不断提及地狱。\par

这样，我们看到祂以各种方式培育祂的门徒（既为他打开天堂，又将那可畏的审判台摆在他面前，并指向天使的圆形剧场，以及在它们中间如何宣告冠冕，这事从此便为虔诚的话语易于接受预备了道路）；在接下来的部分，为免他们胆怯而使话语受阻，祂吩咐他们甚至要为被杀本身做好准备；使他们知道，那些继续陷于错误的人，将因谋害他们而受苦（此外还有其他事）。\par

4. 因此，让我们蔑视死亡，尽管要求我们这样做的时刻尚未到来；因为它确实会将我们迁往更美好的生命。\Quote{但身体朽坏。}事实上，正因如此，我们最应当欢喜，因为死亡朽坏，必死性消灭，而非身体的本体。因为，倘若你看见一座雕像正在铸造，你不会称这过程为毁灭，而是改进的形成。对于身体，你也这样想，不要哀哭。那时，若它仍在惩罚中，才当哀哭。\par

\Quote{但是，}有人说，\Quote{这事应当不发生我们身体的朽坏；它们应当保持完整。}这对生者或逝者有何益处？你们这些爱惜身体的人要到几时呢？你们禁锢于大地、渴望影儿要到几时呢？这曾有什么好处？或者，岂非反而有什么害处？因为，倘若我们的身体不腐朽，首先，最大的恶——骄傲——就会在许多人中持续。因为即使在这事发生、蛆虫涌出之时，还有许多人对成为神热切追求；若身体持续，结果会怎样呢？\par

其次，人们就不会相信它是出于泥土；因为，如果它的结局见证了这一点，而有些人仍然怀疑，那么如果他们没看到这一点，又会怀疑什么呢？第三，身体会受到过度的喜爱；大多数人会变得更加肉性和粗鄙；如果现在还有人依附于人的坟墓和棺材，在它们本身已经朽坏之后，他们还会做什么呢，如果连他们的形象都保存下来？第四，他们不会热切渴望将来的事物。第五，那些说世界是永恒的人会更加确信，并否认天主是造物主。第六，他们不会知道灵魂的卓越，以及灵魂在身体内临在是何等伟大的事。第七，许多失去亲人的人会离开他们的城市，住在坟墓里，变得疯狂，不断与自己的死者交谈。因为如果现在人们还为自己塑造形象，既然他们不能保留身体（因为不可能，无论愿意与否，它都会从他们身上滑落和消逝），并且被钉在木板上，那么他们那时会想出什么怪诞的事呢？依我看来，一般人甚至会为这样的身体建造庙宇，那些精通此类巫术的人会说服邪魔通过它们说话；因为至少现在，那些冒险使用招魂术的人尝试了许多比这些更离奇的事。从这会产生多少偶像崇拜呢？当人们在尘土和灰烬之后，仍然热衷于那些做法。\par

因此，天主为了除去我们所有的过分行为，并教导我们远离一切尘世事物，就在我们眼前毁灭身体。因为即使是一个迷恋身体、看到美丽少女就深受影响的人，如果他不通过言谈学习那物质的丑陋，也会通过亲眼所见而知道。是的，许多与他所爱之人同龄、常常也更美丽的女子，死后第一天或第二天就散发出恶臭、腐烂物，并生虫。那么想想你所爱的是何种美丽，是何等优雅有力量如此扰乱你。但如果身体不腐烂，这一点就不会被清楚知道；就如邪魔跑到人的坟墓，我们许多情人也是如此，不断坐在坟墓旁，会在灵魂中接受邪魔，并很快在这可悲的疯狂中灭亡。\par

但事实上，与其他一切事物一起，这一点也安慰灵魂：看不见形体，使人忘记自己的痛苦。的确，如果不是这样，就根本不会有坟墓，但你会看到我们的城市有尸体代替雕像，每个人都渴望看到自己的死者。由此会产生许多混乱，普通人不会关心自己的灵魂，也不会给不朽的教义留出位置；还会有许多比这些更可怕的事，甚至说出来也不合适。因此，它立刻腐烂，好让你能看到灵魂之美的揭示。因为如果她是所有美和生命的促成者，那么她本身必然更加卓越。如果她保存了那如此畸形和难看的东西，那么她本身更甚。\par

5. 因为美不在于身体，而在于表情，以及灵魂倾注于其物质上的光彩。现在我吩咐你爱那使身体也成为其样子的东西。我为什么说到死亡？即使在生命本身中，我也要你注意，凡是美的都是属于她的。因为无论她是否喜悦，她都在脸颊上洒落玫瑰；无论她是否痛苦，她都会拿走那美，并用暗袍将其全部包裹。如果她持续欢乐，身体状况就会改善；如果悲伤，她会使它比蜘蛛网更薄更弱；如果发怒，她又使它变得可憎和污秽；如果她使眼睛平静，她就赋予极大的美；如果她表达嫉妒，她撒在我们身上的色调非常苍白和发青；如果爱，她立即赋予丰富的优雅。因此，事实上许多女人面貌不美，却从灵魂中获得了许多恩宠；另一些面色光鲜，却因有不优雅的灵魂而破坏了她们的美。想想一张苍白的脸如何变红，通过颜色的变化产生极大的愉悦，当需要羞耻和脸红时。另一方面，如果它无耻，它使面容比任何怪物更令人不快。\par

因为没有比美丽的灵魂更美丽、更甜蜜的了。因为对身体的渴望伴随着痛苦，而在灵魂的情况下，快乐是纯净和平静的。那么为什么要放弃国王，而为传令官疯狂呢？为什么要离开哲学家，而盯着他的解释者呢？你看到一只美丽的眼睛了吗？了解内在的东西；如果那并不美丽，那么也轻视这个。因为的确，如果你看到一个丑陋的女人戴着美丽的面具，她不会给你留下印象；正如另一方面，你也不会让一位美丽的人被面具所伪装，而是会拿掉它，因为选择看到她的美貌显露。\par

那么我吩咐你们也对灵魂这样做，并首先了解它；因为灵魂是用身体代替面具来穿戴的；因此，身体保持其原样；但另一个，尽管它可能是畸形的，可以很快变得美丽。即使它有一只难看、粗暴、凶猛的眼睛，它也可以变得美丽、温和、平静、性情温和、温柔。\par

因此，让我们寻求这美，装饰这面容；好使天主也\Quote{对我们的美感到喜悦}，并借着我们的主耶稣基督的恩宠和爱人之心，将祂永久的祝福分赐给我们，愿光荣和权能归于祂，直到永远。阿们。\par

\SourceRecord{Translated by George Prevost and revised by M.B. Riddle.；Nicene and Post-Nicene Fathers, First Series；Vol. 10.；Edited by Philip Schaff.；Buffalo, NY: Christian Literature Publishing Co.,；1888.；<http://www.newadvent.org/fathers/200134.htm>.}

% structure: p0001 source=h1 target=section reason=page-title

\section{\WorkTitle{玛窦福音第35篇讲道}}

\BibleRef{玛 10:34}\par

\BibleQuote{\Emph{不要以为我来是给大地送和平；我来不是送和平，而是送刀剑。}}\par

他再次提出更难忍受的事，而且加重语气。他预先说出他们所必然遇到的反对。我的意思是，免得他们听见这话后会说：\Quote{那么，你来的目的，就是要毁灭我们和服从我们的人，使大地充满战争吗？}他首先自己说：\BibleQuote{我来不是送和平到大地。}\par

那么，他如何嘱咐他们进人家时要先祝福和平？再者，天使又如何说：\BibleQuote{天主受享光荣于高天，主爱的人在世享平安}？\BibleRef{路 2:14} 所有先知又怎能传出这喜讯？因为这正是和平的最大表现，当患病的肢体被切除，当叛乱的分子被除去。这样，天上就能与大地结合。同样，医生切除不治的部分，以保全身体其他部分；将军使那些共谋作恶的人分离。在著名的巴贝耳塔事件中也是如此：他们恶意的和平\BibleRef{创 11:7}被善意的纷争所终止，从而缔造了和平。保禄也是这样分裂了那些合谋反对他的人。\BibleRef{宗 23:6} 在纳波特的事上，那种和睦比任何战争更令人痛心。\EditorialNote{列王纪上 21章} 因为一致并非总是好事，连强盗也彼此同心。\par

战争不是出于他的意愿，而是出于他们的性情。他的意愿本是所有人都同心合意于虔敬的话，但因他们陷入纷争，战争便发生了。然而他并未如此说，而是说：\BibleQuote{我来不是送和平}，以此安慰他们。好像是说：不要以为这些事要归咎于你们；是我如此安排，因为人倾向如此。因此，你们不要惊惶，好像事情出于意外。我为此而来，要在人间引发战争；这是我的意愿。所以，当大地处于战争时，你们不要烦恼，好像中了敌人的诡计。因为当较坏的部分被割裂后，天上才能与较好的部分结合。\par

他说这些话，是为坚固他们，以应对群众的恶意猜疑。\par

他说的不是\Quote{战争}，而是比战争更可怖的\Quote{刀剑}。若这些言辞带有痛苦和惊吓，不要奇怪。因为他要用严厉的话语训练他们的耳朵，使他们在艰难环境中不致动摇，他特意如此讲论，免得有人说他是用谄媚和隐瞒艰难来劝服他们。因此，他甚至将那些本该用温和语气表达的事，用更刺激、更痛苦的话语说出。因为看一个人的温柔体现在实际上，而非言语上。\par

2. 他对此尚不满足，还揭露了战争的本质，表明它比内战更可怕。他说：\BibleQuote{我来是为了叫人与父亲对立，女儿与母亲对立，媳妇与婆婆对立。}\par

不仅是朋友或同胞，连至亲也要彼此为敌，天性要分裂自身。他说：\BibleQuote{我来是为了叫人与父亲对立，女儿与母亲对立，媳妇与婆婆对立。}也就是说，战争不仅发生在同家之人中间，而且在最亲近、极亲密的人之间。这尤其显示了他的能力：他们听见这些事，不但接受了他，还去劝服众人。\par

这并非他造成的——当然不是——而是另一方的邪恶所为；但他却说是自己的作为。因为圣经有这样的惯用语。别处他也说：\BibleQuote{天主给了他们眼睛，叫他们看不见}\BibleRef{罗 11:8}；这里他这样说，是为了让门徒借着这些话事先得到锻炼，当遭受辱骂和侮辱时，不致慌乱。\par

若有人认为这些话难以忍受，让他们记起一段古代历史。因为在古时也发生过这样的事，这尤其显示旧约与新约同出一源，说这些话的与颁布那些诫命的是同一位。我是指犹太人：当各人杀了自己的邻人时，天主才止息了对他们的愤怒；既在他们造牛犊时，又当他们依附巴耳培敖尔时。那么，那些说\Quote{那个神是恶的，这个神是善的}人在哪里？因为看，他使世界充满了亲族相残的血。但我们仍认为这是极爱人作为。\par

因此，他暗示自己也赞同那些古时的事，又提及一段预言，虽非为此目的而说，却包含同一意义。这预言是什么？\par

\BibleQuote{人的仇敌就是自己家里的人。}\BibleRef{玛 10:36}\par

因为在犹太人中也发生类似的事。有先知和假先知，人民分裂，家庭纷争；有人信这个，有人信那个。故此先知告诫说：\BibleQuote{不要信赖朋友，不要依靠首领；甚至对怀中的妻子也要谨慎，不可向她吐露任何事}；并且说：\BibleQuote{人的仇敌就是自己家里的人。}\BibleRef{米 7:5}\par

祂这样说，是为预备那应接受这话的人超乎一切之上。因为死并非恶事，但恶死才是恶事。为此祂又说：\BibleQuote{我来是把火投在地上。}\BibleRef{路 12:49} 祂这样说，是为表明祂所要求的爱的热烈与温暖。因为祂既然非常爱我们，也愿意我们这样爱祂。这些话也坚固当时在场的人，并提升他们到更高境界。\Quote{因为如果那些人，}祂说，\Quote{要轻看亲属、儿女和父母，想想你们身为他们的老师，该当是怎样的人。既然苦难不止于你们，还将延续到其余的人。因为我既已带来极大的恩赐，也要求极大的顺服和内心的决意。}\par

3. \BibleQuote{爱父亲或母亲超过我，不配属于我；爱儿子或女儿超过我，不配属于我；不背起自己的十字架跟随我，也不配属于我。}\BibleRef{玛 10:37}\par

你看到老师的尊荣吗？你看到祂如何表明自己是生祂者的真儿子，命令我们放下一切下界的事物，而优先选择爱祂吗？\par

\Quote{何必说，}祂说，\Quote{朋友和亲属？即使你把自己的生命置于我的爱之上，你也离我的门徒很远。}那么，这些话是否与旧约相悖？绝不然，倒是非常和谐。因为那里也命令不仅要恨拜偶像的人，甚至要用石头砸死他们；在\WorkTitle{申命记}中，祂赞扬这些人说：\BibleQuote{他对父亲和母亲说：我没有见过你们；他不认他的弟兄，也不认自己的儿子；他遵守了祢的训言。}\BibleRef{申 33:9} 如果\Person{保禄}关于父母有许多指示，命令我们在凡事上顺从他们，请不要惊讶；祂的意思是只在那些不妨碍虔敬的事上顺从。因为孝敬他们，给予他们一切其他尊荣，本是神圣的义务；但当他们要求过分时，就不该顺从。为此\Person{路加}说：\BibleQuote{若有人到我这里来，不恨自己的父亲、母亲、妻子、儿女、兄弟、姐妹，甚至自己的性命，不能做我的门徒；}\BibleRef{路 14:26} 这不是简单地命令恨他们，因为这样甚至与法律完全相反；而是\Quote{当有人希望被爱超过我时，就在这一点上恨他。因为这样既毁了所爱者，也毁了爱者。}祂说这些话，既为使儿女们更坚定，也为使可能阻碍他们的父母更温和。因为当他们看到祂有如此大的力量和权能，能将他们的儿女从他们身边夺走时，他们就会放弃，因为尝试不可能的事。因此祂也离开父母，将祂的话转向儿女，教训前者不要作此尝试，因为这是不可能的事。\par

然后，为了避免他们愤慨或认为苛刻，请看祂如何使祂的论点转向：在说了\Quote{谁不恨父亲和母亲}之后，祂加上\Quote{和自己的性命}。因为，祂说，你为什么向我提父母、兄弟、姐妹、妻子呢？对任何人来说，没有比性命更亲近的了；然而，如果你连这个也不恨，你就必须在一切事上承受与爱我之人相反的命运。\par

祂的命令不单是简单地恨它，而是要把它暴露于战争、战斗、杀戮和流血之中。\Quote{因为那不背起自己的十字架跟随我的，不能做我的门徒。}祂这样说，不仅是我们必须对抗死亡，还要对抗惨死；不仅是惨死，还有耻辱的死。\par

祂尚未谈及自己的受难，为的是当他们经过一段时间的这些教导后，能更容易接受祂关于受难的话。那么，岂不令人惊讶吗？当他们听到这些话时，他们的灵魂没有从身体飞走，因为苦难处处当前，而美善之物尚在期望中？那么它为什么没有飞走呢？因为说话者的能力极大，听者的爱也极大。因此，虽然听到远比那些伟人\Person{梅瑟}和\Person{耶肋米亚}所忍受的更难以忍受、更刺耳的话，他们仍然顺从，没有反对。\par

\BibleQuote{找到性命的，}祂说，\BibleQuote{必要丧失性命；为我丧失性命的，必要找到性命。}\BibleRef{玛 10:39} 你看到那些过分爱惜性命的人所受的损害何其大吗？你看到那些恨恶性命的人所得的利益何其大吗？我是说，因为这些命令令人不快，当祂吩咐他们反对父母、儿女、本性、亲属、世界，甚至自己的灵魂时，祂也摆出了极大的利益。因此，\Quote{这些事，}祂说，\Quote{不但无害，反而大有裨益；其反面则有害；}祂常藉他们所渴望的事物来激励他们。因为你为什么愿意轻看你的性命呢？因为你爱它吗？那么正是为此轻看它，这样你就能最大限度地使它受益，并尽到爱它者的本分。\par

请看一种无以言表的体贴。因为祂不但就我们的父母实行这种推理，也涉及我们的儿女，更涉及我们的性命，这比一切更亲近；这样，那一点从此就无可置疑，他们也就学会，他们也将这样尽可能多地造福他们的亲属；因为对待我们的性命也是这样，它比一切对我们更为本质。\par

4. 这些话足以劝人接纳他们，即他们所指定的医治者。是的，谁不情愿地接纳那些如此高贵、如此真正的英雄，像狮子一样奔跑在地上，轻看自己的一切，以便别人得救呢？然而祂还提出另一种赏报，表明祂照顾款待者比照顾客人更多。\par

祂赐予的第一种尊荣是说：\par

\BibleQuote{接待你们的，就是接待我；接待我的，就是接待那派遣我来的。}\BibleRef{玛 10:40}\par

有甚么能与这相比呢？就是人接纳圣父和圣子！但祂于此还提出了另一个赏报。\par

祂说：\BibleQuote{谁因先知的名接纳先知，必得先知的赏报；谁因义人的名接纳义人，必得义人的赏报。}\BibleRef{玛 10:41}\par

祂先前如何恐吓那些不接纳他们的人要受惩罚，如今也为善人规定了某种慰藉。为教训你们祂对他们更大的眷顾，祂没有简单地说\Quote{谁接纳先知}或\Quote{谁接纳义人}，而是加上\Quote{因先知的名}和\Quote{因义人的名}；也就是说，如果他接纳他不是为了任何世俗的利禄，也不是为任何其他暂时的事物，而是因为他或是先知或是义人，他必得先知的赏报和义人的赏报；即适合那接纳先知或义人者所得的，或者如同那人本身所要领受的。保罗也曾说过类似的话：\BibleQuote{你们的富裕可补足他们的缺乏，好使他们的富裕也补足你们的缺乏。}\BibleRef{格后 8:14}\par

然后，免得有人推说贫穷，祂说：\par

\BibleQuote{无论谁因门徒的名，只把一杯凉水给这些小子中的一个喝，我实在告诉你们，他决不会失去他的赏报。}\BibleRef{玛 10:43}\par

\Quote{即使你的礼物只是一杯凉水，毫无花费，你们也将为此积存赏报。因为我做这一切都是为了你们这些接纳者。}\par

你们看见祂用了何等有力的说服，如何为他们在普天下打开了门户？是的，祂表明人们是他们的债务人：第一，祂说\Quote{工人当得起他的工资}；第二，祂差遣他们出去时一无所有；第三，祂让他们为那些接纳他们的人经历战争和战斗；第四，祂又把奇迹托付给他们；第五，祂藉他们的口将平安——一切福泽的缘由——引入接纳者的家中；第六，祂威胁那些不接纳他们的人，要遭受比索多玛更重的灾祸；第七，祂表明凡欢迎他们的人就是接纳祂自己和圣父；第八，祂应许先知的赏报和义人的赏报；第九，祂承诺即使是一杯凉水也会得到丰厚的回报。这其中每一件，即使单独而言，都足以吸引他们。因为，请告诉我，有哪位将军，在无数伤口染满鲜血，经过多次凯旋从战争和冲突中归来，出现在眼前，人们不会开门接纳他，打开家中每一扇门呢？\par

5. 但如今谁是这样的人？有人会说。因此祂加上\Quote{因门徒的名，因先知的名，因义人的名}，为教训你：赏报的指定，不是基于来访者的功绩，而是基于欢迎者的意向。因为虽然祂在此说的是先知、义人和门徒，但在别处祂却命令人接纳最被遗弃的人，并惩罚那些不这样做的人。因为\BibleQuote{你们没有给这些最小中的一个做，就是没有给我做}\BibleRef{玛 25:45}；祂又对同样的人肯定了相反的情形。\par

因为即使他不行什么大事，他也是人，与你们同住一个世界，看见同一个太阳，具有同一个灵魂，同一位主，与你们同受一样的奥迹，被召进同一个天国；他有一个强有力的理由，就是他的贫穷和缺乏必需的食物。但现在，那些在冬天用笛子和箫声吵醒你们，无目的无果实地打扰你们的人，从你们那里得到许多赠礼后便离去。而那些带着燕子、弄脏自己、辱骂任何人的人，也因这种咒术得到赏报。但如果有一个穷人来到你们面前需要面包，就难免有责骂、斥责、懒惰的指控、责备、侮辱、嘲讽；你们不思想自己也是懒惰的，但天主仍将祂的恩赐给予你们。因为不要对我说：你也在做什么事；倒要告诉我你所做的和忙碌的事是否有任何必要。但如果你说的是赚钱、买卖、照管和增加你的货物，我也要对你说：这些不是真正的工，而是施舍、祈祷、保护受屈者以及诸如此类的事——在这些事上我们完全懒散。然而天主从未对我们说：\Quote{因为你们懒惰，我不为你们点亮太阳；因为你们不做真正要紧的事，我熄灭月亮，使大地荒芜，制止湖泊、泉源、河流，涂抹大气，停止岁雨；}而是丰丰富富地赐给我们一切。祂甚至将这些东西白白赐给那些不仅懒惰，而且作恶的人。\par

因此，当你看见一个穷人，并说：\Quote{我气愤不过：这个年轻力壮的人，一无所靠，宁可懒惰度日；他一定是个奴隶或逃奴，抛弃了他的主人。}我吩咐你对你自己说这些话；或者更好，让他自由地对你说，他将更公义地说：\Quote{我气愤不过：你身强力壮，却懒惰，不实行天主吩咐的任何事，反倒逃离你主人的诫命，在邪恶中游荡，如同在异地，在醉酒、贪吃、偷窃、勒索、倾覆别人家中生活。}你确实归咎于懒惰，但我归咎于恶行：在你的阴谋、你的诅咒、你的谎言、你的掠夺、你做无数这类事中。\par

我说这话，并非制定法律赞成懒惰，绝不是；而是极其热切地希望人人都忙碌；因为懒惰是万恶的导师；但我恳求你们不要无情，不要残忍。因为保禄虽然无数次抱怨，且说：\BibleQuote{若有人不肯工作，就不应当吃饭}，但并未止于此，而是加上：\BibleQuote{至于你们，行善不可厌倦}。\Quote{不，这些话是矛盾的。因为你既然命令他们不可吃饭，怎么又劝我们给予呢？}他说：我这样做，因为我也命令要躲避他们，\BibleQuote{不要与他们来往}；我又说：\BibleQuote{不要以他们为仇敌，而要劝戒他们}；（\BibleRef{得后 3:14}）我并非制定矛盾的法律，而是完全和谐一致的。因为，如果你急于施恩，那么他，穷人，很快就会摆脱懒惰，而你也会摆脱残忍。\par

你回答说：\Quote{他有许多谎言和诡计。}那么，正因此他更可怜，因为他陷入如此困境，以致连这种行为都麻木了。但我们非但不怜悯，反而加上这些残忍的话：\Quote{你不是已经接受过一次又一次了吗？}我们就这样说话。那么怎样呢？因为他曾吃饱一次，他就无需再吃吗？你为什么不也为自己的肚腹制定这些法律，也向它说：\Quote{你昨天和前天都吃饱了，今天不要寻求了？}但你却无度地填满它，甚至到破裂，而对他，当他只求一点合理的东西时，你却转身离去；而你正因此应当怜悯他，因为他被迫每天来求你。是的，如果没有别的感动你，你也应当为此怜悯他；因为贫困的迫使他不得不这样做。你不怜悯他，因为他听到这样的话而不感到羞耻：原因是他的需要太过强大。\par

你非但不怜悯，反而羞辱他；天主命令要暗中施舍，你却公开暴露那向你乞讨的人，并谴责他本应使你怜悯的事。如果你不想施舍，为何还要加以斥责，伤害那疲惫可怜的灵魂？他来到你跟前，如同进入港口，寻求帮助；你为何要掀起波涛，使风暴更剧烈？你为何谴责他卑鄙？什么？如果他想到会受到这些话，他会来找你吗？或者，如果他实际来了并预见到了这些，你就更应该怜悯他，并战栗于你的残忍，因为即使如此，当你看到他迫于不可抗拒的需要时，你也不变得温和，也不认为他因饥饿的恐惧而有足够的理由来纠缠，反而指责他无耻；然而你自己常常犯下更大的无耻，甚至是严重的事情。因为在这里，无耻本身就带来了原谅的理由，而我们，常常犯下应受惩罚的事，却厚颜无耻；当我们应当记住这一切并谦卑时，我们甚至践踏那些可怜的人，当他们请求药膏时，我们却增加他们的伤口。我说，如果你不愿意给，为什么还要打击？如果你不愿意慷慨，为什么还要傲慢？\par

\Quote{但他不接受其他方式的拒绝。}那么，如那位智者所命令的，去做吧。\BibleQuote{你要以温和回答他平安的话。}（\BibleRef{德 4:8}）因为，他如此纠缠不休，当然不是出于自愿。没有，也不可能有人愿意为了羞耻本身而蒙羞。无论有多少人争辩，我永远无法被说服：一个生活富足的人会选择乞讨。\par

6. 那么，不要让人用论证欺骗我们。尽管保禄说：\BibleQuote{若有人不肯工作，就不应当吃饭}，（\BibleRef{得后 3:10}）他是对他们说的；但对我们，他说的不是这个，相反，他说：\BibleQuote{行善不可厌倦}。（\BibleRef{得后 3:13}）我们在家里也这样做；当两个人互相争吵时，我们把每个人拉到一边，给他们相反的建议。天主和梅瑟也这样做了。因为，他（梅瑟）对天主说：\Quote{如果你愿意赦免他们的罪，就赦免；否则，也请把我涂抹。}他反而命令他们互相残杀，并毁灭所有属于他们的人。然而这些事是相反的；但两者都指向同一个目的。\par

再者，天主在犹太人听得到的情况下对梅瑟说：\BibleQuote{你且由我，让我消灭这人民。}（\BibleRef{出 32:10}）（因为当他们听到天主这样说时，并不在场，但后来他们会听到的）：但祂私下给他相反指示。梅瑟后来被迫透露这事，这样说：\Quote{什么？难道我怀了他们，你竟对我说：你要把他们抱在怀里，像乳母怀抱乳儿？}\par

这些事在家里也发生；常常一个父亲私下责备教师，因为他用责骂对待孩子，说：\Quote{不要粗暴，也不要严厉，}但对年轻人却说相反的话：\Quote{即使你遭受不公正的斥责，也要忍受。}从这些相反的事中造就某种健全的结果。同样，保禄对健康而乞讨的人说：\BibleQuote{若有人不肯工作，就不应当吃饭，}（\BibleRef{得后 3:10}）以促使他们就业；但对能施恩的人说：\BibleQuote{你们，行善不可厌倦。}（\BibleRef{得后 3:13}）以引导他们施舍。\par

同样，当他告诫外邦人，在\WorkTitle{罗马书}里，不要对犹太人骄傲时，他也引用了野橄榄树，他似乎对这些人说一件事，对那些人说另一件事。（\BibleRef{罗 11:17}）\par

因此，让我们不要陷入残忍，而要听从保禄说：\BibleQuote{行善不可厌倦}；听从主说：\BibleQuote{凡求你的，就给他}，（\BibleRef{路 6:30}）和\BibleQuote{你们应当慈悲，像你们的父那样慈悲。}（\BibleRef{路 6:36}）虽然主谈了许多事，但从未使用过这个措辞，唯独关于我们慈悲的行为。因为没有什么比行善更使我们与天主相似。\par

\Quote{但没有什么比穷人更无耻的，}有一个人说，\Quote{比一个穷人。}请问，为什么呢？因为他跑上来，跟在您后面喊叫？那么让我指出，我们如何比他们更纠缠不休，且极其无耻。请回想，我说，在守斋的季节，傍晚时分当您摆好餐桌，叫了服侍的仆人；他稍一迟缓，您便把一切推翻，脚踢、辱骂、呵斥，仅仅因为一点耽搁；虽然您十分确信，若不是立刻，稍后片刻您也能享用食物。那时您并不称自己为无耻，竟为了一点小事变成了野兽；而穷人，为了更重大的事惊慌战栗（因为他的恐惧全不是关于耽搁，而是关于饥荒），您却称他鲁莽、无耻、厚颜，以及所有最辱骂的名称？这不是极端的厚颜无耻又是什么呢？\par

但我们不考虑这些事：因此我们觉得这样的人麻烦；如果我们稍微省察自己的行为，并与他们的相比较，我们便不会认为他们不可容忍了。\par

那么，不要作严厉的审判者。即使您没有一切罪过，天主的法律也不允许您苛刻地追究别人的罪过。如果法利塞人因此而丧亡，我们又能找到什么辩护呢？如果祂不允许那些行善的人苛求别人的行为，更何况那些犯罪的人呢？\par

7. 那么，让我们不要残忍、不要残酷、不要没有自然情感、不要不和解、不要比野兽更坏。因为我知道许多人甚至走向了那样野蛮的地步，为了一点麻烦，就轻视饥饿的人，并且说出这样的话：\Quote{我现在身边没有仆人；我们离家很远；我不认识任何兑换银钱的人。}哦，残酷！您应许了更大的事，却不履行那较小的吗？为了节省您走一小段路，他就该饿死吗？哦，傲慢！哦，骄傲！即使要走十斯塔狄亚，难道您就退缩吗？您难道没有想到，这样您的赏报会更伟大吗？因为，您施舍时，只因施舍而得到赏报；而当您自己也亲自前往时，为此又为您预备了赏报。\par

是的，我们钦佩族长本人，因他亲自跑到牛群中，抓起牛犊，\BibleRef{创 18:7}——那时他家里有三百一十八个生在家中的仆人。但现在有些人充满了如此多的骄傲，竟让仆人去做这些事，并不以为耻。有人说：\Quote{但您要我自己做这些事吗？}有人可以说，\Quote{那样我岂不显得虚荣？}不，事实上，您是被另一种虚荣心驱使去做这事，因羞于被人看见和穷人说话。\par

但我在这方面毫不严格；只管施舍，无论您愿意亲自做还是借他人之手；不要控告、不要打击、不要辱骂。因为到您这里来的人需要的是药物，而不是伤口；需要怜悯，而不是刀剑。请告诉我，如果有人被石头击中，头上受了伤，抛开所有的人，满身是血地跑到您的膝前，您难道要用另一块石头击打他，给他加上另一个伤口吗？我确实认为不会；相反，您会尽力医治。那么，您为什么对穷人做相反的事呢？您不知道言语有多大力量，既能扶起，也能推倒吗？\Quote{因为一句话，}经上说，\Quote{胜过馈赠。}\BibleRef{德 18:16}\par

您岂不考虑，您正在把剑刺向自己，并承受更重的伤，当穷人被辱骂后，默然退去，唉声叹气，泪流满面？因为他的确是天主派到您这里来的。那么，您要想到，侮辱他，您是将侮辱转移到谁身上？天主派他到您这里来，命令您施舍，但您不但不施舍，反而在他到来时侮辱他。\par

如果您不清楚这是多么严重的错误，请以人间的事来看，然后您就会完全知道这罪过的严重性。比方说：如果您的仆人受了您的命令，到另一个欠您钱的仆人那里去收钱，却不仅空手而回，还受到侮辱，您会怎样对待那施行侮辱的人？您会如何惩罚他？仿佛此后是您自己受到了虐待。\par

请您也这样对待天主；因为确实是祂把穷人送到我们这里，我们施舍的也是属于祂的，如果我们确实施舍的话。但如果不施舍，还让他们受侮辱离去，请想想，我们这样做该遭多少雷霆、多少闪电呢？\par

那么，我们充分思考这一切之后，既要勒住我们的舌头，又要除去无情，伸出施舍的手；不仅要借着金钱，也要借着言语，救济那些有需要的人；这样我们既能逃避辱骂的惩罚，又能因祝福和施舍而承继天国，赖我们的主耶稣基督的恩宠和慈爱，愿光荣与权能归于祂，至于无穷之世。阿们。\par

\SourceRecord{Translated by George Prevost and revised by M.B. Riddle.；Nicene and Post-Nicene Fathers, First Series；Vol. 10.；Edited by Philip Schaff.；Buffalo, NY: Christian Literature Publishing Co.,；1888.；<http://www.newadvent.org/fathers/200135.htm>.}

% structure: p0001 source=h1 target=section reason=page-title

\section{\WorkTitle{玛窦福音第三十六篇讲道}}

\BibleRef{玛 11:1}\par

\BibleQuote{耶稣嘱咐完了祂的十二门徒，就从那里走了，为在他们的城里施教宣讲。}\par

也就是说，在祂派遣他们之后，祂便退开，好给他们空间和机会去执行祂所吩咐的。因为当祂在场并施行医治时，没有人愿意接近他们。\par

\Quote{若翰在狱中听到基督所行的，就派自己的门徒去问祂说：\InnerQuote{你就是要来的那一位，或是我们还要等候另一位？}} \BibleRef{玛 11:2}\par

但路加说，他们也把奇迹告诉了若翰，然后若翰派了他们去。\BibleRef{路 7:18} 然而，这并不包含什么难题，只是值得思考；因为这件事，除了别的，表明他们对祂的嫉妒。\par

但下面的事完全属于有争议之点。这性质如何？他们说：\BibleQuote{你就是要来的那一位，或是我们还要等候另一位？} 也就是说，那位在祂行奇迹之前就认识祂的人，那位从圣神那里学到这事的人，那位从父那里听到的人，那位在众人面前宣告祂的人；他现在却派人去问祂，是否为祂自己，或不是？如果你还不知道那确实是祂，你怎么能认为自己可信，而断言你所不知道的事？因为要为别人作证的人，必须先自己可信。你不是说：\BibleQuote{我是不配解祂的鞋带的？}\BibleRef{若 1:27} 你不是说：\BibleQuote{我不认识祂，但那位派遣我来以水施洗的，对我说：你看见圣神降下并停留在谁身上，谁就是那以圣神施洗的？}\BibleRef{若 1:33} 你不是看见圣神以鸽子的形状？你不是听见声音？你不是完全阻止祂说：\BibleQuote{我需要受你的洗？}\BibleRef{玛 3:14} 你不是甚至对你的门徒说：\BibleQuote{祂必兴盛，我必衰微？}\BibleRef{若 3:30} 你不是教训众人说：\BibleQuote{祂要以圣神和火给你们施洗；}\BibleRef{玛 3:11} 以及\BibleQuote{祂是天主的羔羊，除免世罪者？} 在祂的征兆和奇迹之前，你不是已经宣告了这一切？那么现在，当祂已向众人显明，祂的名声传遍各处，死人复活，魔鬼被驱逐，如此大的奇迹已展示，你怎么之后还派人去问祂？\par

那么事实是什么？难道所有这些话都是某种欺诈、戏耍和寓言？不，任何一个有理智的人会这样说吗？我且不说若翰，他在母胎中跳跃，在他出生之前就宣告了祂，那旷野的居民，那展示天使般生活的人；即使他是一个常人，一个完全被遗弃的人，在如此多的证言之后——既有他自己的，也有别人的——他也不会犹豫。\par

因此，显然他派人去并非因为自己怀疑，也不是出于无知而问。因为没有人能说，虽然他完全知道，但由于牢狱之灾，他变得胆怯了；因为他并不指望从中获释，即使他指望，也不会背叛对天主的职责，他早已武装起来面对各种死亡。除非他没有准备，否则他不会对那习惯于流先知血的全体人民表现出如此大的勇气；也不会在城中、在广场上，当着所有人的面，以如此大的胆量斥责那野蛮的暴君，严厉地责备他，好像他是一个孩子。即使他变得更胆怯，他又怎能不在自己的门徒面前羞愧呢？他曾在他们面前多次为祂作证，却通过他们提问，这本该由别人去做。然而他清楚地知道，他们也嫉妒基督，想找把柄攻击祂。他又怎能不在犹太人面前羞愧呢？他曾在他们面前宣告如此崇高的事？这对他从捆绑中解脱有什么好处？因为他不是因基督的缘故被投入监狱，也不是因宣告祂的能力，而是因他自己关于那不合法的婚姻的责备。哪个如此愚笨的孩子，哪个如此疯狂的人，会这样装模作样呢？\par

2. 那么他到底在做什么？因为从我们所说的，显然若翰也不能对此怀疑，任何普通人，甚至极愚蠢和疯狂的人也不能。现在我们要加上解释。\par

他派去问的目的是什么？若翰的门徒正在远离耶稣，这一点任何人都能看到，他们一直对祂怀有嫉妒。这从他们对他们师傅的话可以清楚看出：\BibleQuote{曾同你在约但河对岸、你曾为祂作证的那位，看，祂也施洗，众人都到祂那里去了。}\BibleRef{若 3:26} 又说：\Quote{若翰的门徒和犹太人之间发生了一场关于洁净的争论。} 他们又来到祂跟前说：\BibleQuote{我们和法利塞人常常禁食，而你的门徒却不禁食？}\BibleRef{玛 9:14} 因为他们还不知道基督是谁，以为耶稣只是一个凡人，而若翰是超越凡人的伟大者，就因前者受重视而烦恼，而后者正如他所说的，现在已停止。这阻碍了他们到祂那里去，他们的嫉妒完全阻塞了通道。当若翰还和他们在一起时，他不断劝诫和教导他们，但即使这样也没有说服他们；当他即将去世时，他更加努力：因为他害怕留下不良教义的基础，他们继续与基督分离。正如他在开始时殷勤把一切属于他的人带到基督那里；既然他未能说服他们，现在临近终结，他更加热心。\par

假如他说：\Quote{你们到他那里去吧，他比我更好，}他不会说服他们，因为他们本来就不愿轻易离开他；他反而会被认为是出于谦逊而这样说，他们会更加依附于他。或者如果他保持沉默，那么同样无济于事。那么他怎么做呢？他等到他们听到基督行奇迹的消息，即使如此，他也没有劝诫他们，也没有差遣所有人，而是差遣了两个人（他或许知道这两人比其他门徒更易受教）；这样询问就可不引起怀疑，以便他们通过他的作为了解耶稣和他之间的区别。他说：你们去说：\BibleQuote{你就是那要来的那一位，还是我们应当期待另一位？}\BibleRef{玛 11:3}\par

但是基督知道若翰的意图，并没有说\Quote{我就是那一位}；因为这会再次触怒听者，尽管这按常理是他应当说的，但他让他们从他的作为中学习。因为经上说，当这些人来到他那里时，他就\BibleQuote{治愈了许多人。}\BibleRef{路 7:21}然而，当被问\BibleQuote{你就是那一位吗？}，他却对此不置一词，而是立即医治那些患病的人，这有何相称呢？除非他有意要确立我刚才提到的这一点？因为他们自然认为他行为的见证比他的言语更可靠、更无可怀疑。\par

因此，身为天主，他知道若翰差遣他们来的心思，便立即治愈了瞎眼的、瘸腿的以及许多其他人；不是为了教导他（因为怎能教导一个已确信的人呢？），而是为了这些心存疑惑的人。治愈他们之后，他说：\par

\BibleQuote{你们去，再告诉若翰你们所听见、所看见的事：瞎子看见，瘸子行走，癞病人洁净，聋子听见，死人复活，穷人听到福音。}\BibleRef{玛 4:5}他又加了一句：\BibleQuote{凡不因我而绊倒的，是有福的。}暗示他知道他们甚至未说出口的心思。因为假如他说了\Quote{我就是那一位}，这就会像我已说过的，使得他们绊倒；而且即使他们不说出来，也会像犹太人对他说的那样：\BibleQuote{你为自己作证。}因此他自己不说这话，而是让他们从奇迹中得知一切，使他所教导的免受怀疑，而且变得更加清楚。因此他也暗中加入了他们对他的责备。也就是说，因为他们\Quote{因他而绊倒}，他通过陈述他们的情形并将此事留给他们的良心，且不为此控告召叫任何证人，只有他们自己知道一切，这样也越加吸引他们归向他，说：\BibleQuote{凡不因我而绊倒的，是有福的。}因为他说这话时，暗指的就是他们。\par

3. 但是，为了使真理通过不同陈述的比较更清楚地向你们显明，我们不仅要提出我们自己的话，也要提出别人所说的，所以我们有必要对它们做一些说明。\par

那么有人主张什么呢？他们断言，上面所说的并不是原因，而是若翰处于无知之中，但并非完全无知；他知道他是基督，但不知道他是否也要为人类而死，因此他说：\BibleQuote{你就是那要来的那一位吗？}即要下降到地狱的那一位。但这并不成立，因为若翰对此也并非无知。至少在众人之前他就宣告了此事，并首先为此作证：\BibleQuote{看，天主的羔羊，除免世罪者。}\BibleRef{若 1:29}他称他为羔羊，以此宣布十字架；而他说：\BibleQuote{除免世罪者}时，也宣告了同一件事。因为除非藉着十字架，他不能成就此事；正如保禄同样说：\BibleQuote{那反对我们的罪状，他把它除去，钉在十字架上。}而他说：\BibleQuote{他要以圣神洗你们，}\BibleRef{玛 3:11}这是预言复活后事件的人所说的话。\par

对啊：他们说他（若翰）知道耶稣要复活，并且要赐予圣神；但他不知道他也要被钉十字架。那么，一个没有受难、没有被钉死的人如何能复活呢？而一个连先知所知道的事都不知道的人，如何比先知更大呢？因为基督亲自作证说他比先知更大，\BibleRef{玛 11:9}但先知们知道受难的事，这是每个人都清楚的。因为依撒意亚这样说：\BibleQuote{他如同一只被牵到宰杀之地的羔羊，又如同母羊在剪毛的人前默不作声。}\BibleRef{依 53:7}在这段见证之前，他也说：\BibleQuote{从叶瑟的根茎要生出一枝，那将要兴起统治外邦人的，外邦人都将寄望于他。}\BibleRef{依 11:10}然后讲到他的受难和随之而来的荣耀，他加上：\BibleQuote{他的安息乃是荣耀。}这位先知不仅预言了他要被钉十字架，还预言了与谁同钉。\BibleQuote{因为，}他说，\BibleQuote{他被列于罪犯之中。}\BibleRef{依 53:12}不仅如此，还预言他绝不为自己辩护；\BibleQuote{因为这人，}他说，\BibleQuote{不开口；}以及他将被不义地定罪；\BibleQuote{因为在他受屈辱时，}他说，\BibleQuote{他的审判被夺去。}在这之前，达味也说了这事，并描述了审判厅。\BibleQuote{为什么，}他说，\BibleQuote{外邦人争吵，万民谋划虚妄的事？地上的君王起来，首领聚集在一起，反对上主，反对他的受傅者。}在别处，他也提到了十字架的形象，这样说：\BibleQuote{他们刺穿了我的手和我的脚，}并以所有精确性补充了兵士们胆大妄为所做的事：\BibleQuote{他们分了我的衣服，}他说，\BibleQuote{为我的长衣拈阄。}在别处他又说，他们还给他酸醋喝；\BibleQuote{因为他们给了我，}他说，\BibleQuote{苦胆作食物，我渴时，他们给我酸醋喝。}\par

那么，先知们这么多年前就谈到审判厅、定罪、与祂同钉十字架的人、分衣服、拈阄，以及许多别的事（因为现在实在不必一一列举，免得拖长我们的讲论）：难道这个人，比他们所有人都伟大，竟对这些事一无所知吗？不，这怎么合理呢？\par

\Quote{为什么祂没有说：\InnerQuote{祢是要到地狱里来的那一位吗？}} 而只是说：\InnerQuote{那将要来的那一位？}虽然这比其他说法更荒谬，我是指他们说：\InnerQuote{祂说这些话，是为了在离去后也在那里宣讲。}对于这些人，适当地说：\Quote{弟兄们，在心意上不要作小孩子，但在恶事上却要作小孩子。}\BibleRef{格前 14:20} 因为现世生命确实是正当言行的时机，但死后是审判和惩罚。经上说：\InnerQuote{在地狱里，谁会承认祢？}\par

\Quote{铜门怎样被击破，铁闩怎样被折断？} 通过祂的身体；因为那时首次显示了一个不朽的身体，摧毁了死亡的暴政。此外，这表明死亡权势的毁灭，而不是解除那些在祂来临之前已死之人的罪。如果不是这样，而是祂将所有在祂之前的人从地狱里释放出来，那么祂怎么说：\InnerQuote{索多玛和哥摩辣地方所受的惩罚要比那城容易忍受呢？}\BibleRef{玛 10:15} 因为这句话假定那些人也要受惩罚；虽较轻，但仍要受惩罚。然而他们在此世也受了最严厉的惩罚，但即便如此也不能拯救他们。如果对他们如此，对那些未曾受苦的人更将如何呢？\par

\Quote{那么，} 有人会说，\InnerQuote{那些在祂来临之前生活的人受到了冤屈吗？}绝不，因为那时人们即使没有宣认基督也可以得救。因为这并不要求他们，而是要他们不拜偶像，并认识真天主。经上说：\InnerQuote{上主你的天主是唯一的主。}\BibleRef{申 6:4} 因此玛加伯们受到钦佩，因为他们为遵守法律而受苦；还有三个青年，以及犹太人中许多其他人，表现了极有德行的生活，并保持了这种知识的标准，对他们没有更多要求。因为那时，正如我已经说过的，只认识天主就足以得救；但现在不再是这样，还需要认识基督。因此祂说：\Quote{我若没有来向他们说话，他们就没有罪；但现在他们对自己的罪没有借口了。}\par

同样，关于实践准则。那时，杀人毁掉杀人者，但现在甚至连发怒也是。那时，犯奸淫和与别人妻子同房带来惩罚，但现在甚至连用不洁的眼睛看也是。因为正如知识一样，生活准则现在也变得更严格。所以那里不需要一个先驱。\par

此外，如果不信者在死后因相信而得救，那就没有人会灭亡了。因为那时所有人都会悔改和朝拜。作为这点的证明，请听保禄说：\InnerQuote{一切唇舌都要承认，一切膝盖都要跪拜，在天上的、在地上的、在地下的。}\BibleRef{斐 2:10} 又说：\InnerQuote{最后被消灭的仇敌是死亡。}\BibleRef{格前 15:26} 但那种顺服没有益处，因为它不是出于正确抉择的意向，而是因事物必然性，可以这么说，从那时起就发生了。\par

我们不要再引进这种老妇的谬论和犹太人的传说。至少听听保禄关于这些事所说的话。\InnerQuote{凡没有法律而犯了罪的，也必不按法律而丧亡；}\BibleRef{罗 2:12} 这里他指的是在法律时代以前生活的人；又说：\InnerQuote{凡在法律之下犯了罪的，必按法律受审判，}\BibleRef{罗 2:12} 指的是梅瑟之后的所有人。又说：\InnerQuote{天主的忿怒从天上显示出来，反对一切不虔敬和不义的人，}\BibleRef{罗 1:18} 又说：\InnerQuote{忿恨和愤怒，患难和困苦，加给一切作恶的人，先是犹太人，后是外邦人。}\BibleRef{罗 2:8} 然而外邦人在此世遭受了无数的灾祸，这既由异教徒的历史，也由我们手中的圣经所宣告。因为谁能数算巴比伦人或埃及人的悲剧灾难呢？但是为了证明那些在基督降生成人之前没有认识基督，却远离偶像崇拜，只敬拜天主，并表现出卓越生活的人，将享受一切福乐；请听所说的话：\InnerQuote{荣耀、尊贵、平安加给一切行善的人，先是犹太人，后是外邦人。}你看到吗，他们的善行有许多赏报，而作恶的人又有惩治和刑罚？\par

4. 那么，请问，那些完全不信地狱的人在哪里？如果那些在基督来临之前的人，他们甚至没有听过地狱的名字，也没有听过复活，并且在此世受了罚，将来也要在那里受罚；那么，我们这些在如此多严格德行的教训中培育的人，将如何呢？\par

有人问，那些从未听过地狱的人，怎么会落入地狱呢？因为他们会说：\Quote{如果你以地狱来威胁我们，我们就会更加畏惧，并清醒过来。}的确如此（难道不是吗？）看看我们现在的生活方式，我们每天都听到关于地狱的言论，却毫不在意。\par

此外，还有一点可说：那不被眼前审判所约束的人，更不会为那些未来的审判所约束。因为较不理性的人和性情粗野的人，往往更容易被眼前立即发生的事所警醒，而非那些许久以后才会发生的事。\Quote{但加在我们身上的恐惧更大，}有人会说，\Quote{他们在这方面受了亏待。}绝非如此。首先，给我们和他们的度量并不相同，而是给我们大得多。如今那些承担更重大劳苦的人，理应享受更大的帮助。我们的恐惧增加了，这并不是小小的帮助。如果我们比他们在预知未来上有优势，那么他们比我们在严厉惩罚即刻加诸于他们身上也有优势。\par

但与此相关还有一事，是众人所说的。因为他们说：\Quote{天主的公义在哪里，若有人在这里犯罪，却同时在这里和那里受罚？}那么，我是否要你们记起自己的话，好使你们不再麻烦我们，而从你们自己内部提供解答？我听过我们中的许多人，如果偶然有人告诉他们一个凶杀犯在法庭上被处决，他们如何愤慨，并这样谈论：\Quote{这邪恶可咒的恶徒，犯了三十起凶杀，甚至更多，自己只受一次死；公义在哪里呢？}所以你们自己承认，一次死亡不足以惩罚；那么现在你们为何作出相反的判决呢？因为你们审判的对象不是别人，而是你们自己：自爱对我们觉察什么是正义是这么大的障碍。正因为此，当我们审判别人时，我们严格地探究一切，但当我们在自己身上审判时，我们就被蒙蔽了。因为如果我们也在自己的事上像对待其他人那样探究这些事，我们就会作出公正的判断。因为我们也有罪，该当的不是两三个死，而是千万次死。且略过其余，让我们反省自己，我们中凡是不相称地分享奥秘的人，这样的人就是干犯基督的圣体圣血的人。因此，当你谈论那个凶杀犯时，也要考虑你自己。因为他确实杀了一个人，而你却有杀害主的罪债；他未分享奥秘，而我们却在享受圣筵的恩惠。\par

那些咬噬并吞食自己的弟兄、倾泻如此多毒液的，是些什么人？那抢夺穷人食物的，是些什么人？因为若那不分施自己财物的人，即如我所说的，\Emph{那夺取他人财物的人更是如此}。贪婪者在邪恶上超过多少盗贼！强暴者超过多少凶杀犯和盗墓者！又有多少人在掠夺他人之后甚至贪求他们的血！\par

\Quote{不，}他说，\Quote{天主不许。}现在你说，天主不许。当你有一个仇人时，你就说天主不许，并记起所说的话，活出一个充满极大严谨的生命；免得\Place{索多玛}的分也等待我们，免得我们遭受\Place{哈摩辣}的命运，免得我们经历\Place{提洛}人和\Place{漆冬}人的灾祸；或者更确切地说，免得我们得罪基督，那是比一切更痛苦\Emph{和更可畏惧}的事。\par

因为虽然许多人认为地狱是可怕的事，但我本人却不会停止不断地这样说：这比任何地狱更痛苦和更可畏惧；我也恳求你们拥有同样的心意。因为这样我们就能既脱离地狱，又享受基督所赐的光荣；愿我们众人能藉着我们的主耶稣基督的恩宠和慈爱达到这光荣，愿光荣和权能归于他，永远无穷世。阿们。\par

\SourceRecord{Translated by George Prevost and revised by M.B. Riddle.；Nicene and Post-Nicene Fathers, First Series；Vol. 10.；Edited by Philip Schaff.；Buffalo, NY: Christian Literature Publishing Co.,；1888.；<http://www.newadvent.org/fathers/200136.htm>.}

% structure: p0001 source=h1 target=section reason=page-title

\section{\WorkTitle{玛窦福音释义 第三十七篇}}

\BibleRef{玛 10:7-9}.\par

\Quote{\Emph{你们出去到旷野里要看什么呢？风摇动的芦苇吗？你们出去到底要看什么？穿细软衣服的人吗？看哪，穿细软衣服的人是在王宫里。但你们出去要看什么？一位先知吗？}\Emph{是的，我告诉你们，而且比先知还大。}}\par

因为若翰门徒的事确实已经料理得很好，他们因刚行的奇迹而得到确证；但之后还需要对民众进行补救。虽然他们对自己的师傅不会有任何猜疑，但普通人可能从若翰的询问中产生许多奇怪的怀疑，不知道他派门徒来的用意。他们自然会心里议论说：\Quote{那位曾如此作证的人，现在难道改变了看法，怀疑这位或另一位是不是要来者？难道他与耶稣有分歧才这样说？难道坐牢使他变得胆怯？难道他先前的话是虚妄、漫不经心说的？}既然他们很容易这样猜测，你看祂如何纠正他们的软弱，消除这些猜疑。因为\Quote{他们离开后，祂开始对群众讲话。}为什么\Quote{他们离开后？}为了不显得祂在奉承那人。\par

在纠正民众时，祂没有公开他们的猜疑，只补充了解决他们心中困扰思想的方法：表明祂知道所有人的秘密。因为祂不说，如同对犹太人那样，\Quote{你们为什么心里怀着恶念？}（\BibleRef{玛 9:4}）因为他们心中这样想，并非出于恶意，而是对所说之事无知。因此，祂也不以责备的方式对他们说话，只是纠正他们的理解，为若翰辩护，表示他并未偏离先前看法，也没有改变，他根本不是一个容易动摇反复的人，而是坚定可靠，绝不可能背叛所交付给他的事。\par

在确立这一点时，祂首先不凭自己的判断，而是用他们以前的见证，指出他们如何用言语和行为证明他的坚定。\par

因此祂说：\Quote{你们出去到旷野里要看什么呢？}如同祂说，你们为什么离开你们的城市和房屋，全都一起到旷野来？要看一个可怜而灵活的人吗？不，这完全不合理；这不是那种热忱和众人聚集到旷野所表明的。那么多百姓和城市不会在那个时代怀着如此大的热忱涌向旷野和约旦河，除非你们期望看到某个伟大而奇妙的人，一个比任何岩石都更坚定的人。是的，你们出去看的绝不是\InnerQuote{被风摇动的芦苇}，因为灵活而容易改变的人，今天说这个明天说那个，什么都不坚定，最像那样。\par

还有，祂省略了所有邪恶行为，只提及当时最困扰他们的事，消除了轻率的猜疑。\par

\Quote{但你们出去要看什么？一个穿细软衣服的人吗？看哪，穿细软衣服的人是在王宫里。} \BibleRef{玛 11:8}\par

祂的意思是这样：他本人并非动摇者；你们的热忱也证明了这一点。更不能说，他确实坚定，但后来因耽于奢侈而变得懒散。因为在人中，有些人是天生如此，另一些人则是后天变成；例如，一个人天生易怒，另一个人因长期患病而有了这个弱点。同样，有些人天生灵活善变，另一些人则因成为奢侈的奴隶、生活方式柔弱而变成这样。祂说：\Quote{但是若翰，}祂说，\Quote{既不是天生这种性格，因为你们出去看的不是芦苇；也不是因耽于奢侈而失去他原有的优点。}因为他的衣服、旷野和监狱都表明他没有成为奢侈的奴隶。要是他打算穿细软衣服，就不会住在旷野或监狱，而是在王宫里；他只需保持沉默，就能享受无限的尊荣。既然希律如此尊敬他，即使他责备希律并被囚禁，那么他若保持沉默，希律岂不更加讨好他？你看，他确实证明了自己的坚定和刚毅；他怎么可能理应受到那种猜疑呢？\par

2. 因此，当祂通过地点、衣服和众人的聚集描绘了若翰的品格后，祂接着引出了先知。因为说了\Quote{你们出去要看什么？一位先知吗？是的，我告诉你们，而且比先知还大；}祂接着说：\Quote{这人就是经上所记载的：看，我派遣我的使者在你面前，他要在你前面预备你的道路。} \BibleRef{玛 11:10} 祂先前已列出犹太人的见证，然后应用先知的见证；或者更确切地说，祂首先列出犹太人的论断（这必然是非常有力的证明，因为见证出自祂的仇敌），其次是若翰的生活，第三是祂自己的判断，第四是先知，用各种方法堵住他们的口。\par

然后，免得他们说：\Quote{但如果那时他确实是那样的人，现在却改变了呢？}祂还加上了以下的事：他的衣服、他的监狱，以及随之而来的预言。\par

然后，在说了他比先知更大之后，祂也指明他在哪方面更大。他在哪方面更大呢？在于接近那已经来临的那一位。因为祂说：\Quote{我派遣我的使者在你面前。}即临近你。因为如同在君王面前，那些靠近战车骑行的人比其余的人更显赫，照样，若翰也在他的行程中临近那降临本身。请看，祂也借此指明了若翰的卓越；即使如此祂也没有停止，随后又加上祂自己的见证，说：\Quote{我实在告诉你们：在女人所生者中，没有兴起比洗者若翰更大的。} \BibleRef{玛 11:11}\par

祂所说的就是这样：\Quote{女人没有生过比这人更大的。}祂这句话本身已经足够；但如果你也愿意从事实学习，就请想想他的餐桌、他的生活方式、他灵魂的高超。因为他生活得如同在天上；他超越了本性的需要，走了一条新路，将所有时间用在赞美诗和祈祷中，不与任何人交往，而只与天主不断来往。因为他甚至没有见过任何一个同伴，也没有被他中间任何人看见；他不靠奶食生活，不享受床铺、屋顶、市场或任何人世事物的安慰；然而他同时既温和又刚毅。例如，请听他如何审慎地与自己的门徒论理，如何勇敢地与犹太民众论理，如何公开地与君王论理。为此祂也说：\Quote{在女人所生者中，没有兴起比洗者若翰更大的。}\par

3. 但是，唯恐对若翰的过度赞扬会使犹太人产生一种过分的感情，尊敬若翰超过基督；请看祂又如何纠正这一点。正如那些建树祂门徒的话语对群众造成了伤害，因为他们以为祂是个软弱的人；同样，为众人所用的补救措施可能更有害，他们从基督的话语中对若翰产生了比祂更尊敬的看法。\par

因此，祂以不引人怀疑的方式纠正了这一点，说道：\Quote{在诸天之国里，较小的比他还大。}较小的是指年龄和按群众的看法，因为他们甚至称祂为\Quote{一个贪吃嗜酒的人}，\BibleRef{玛 11:19} 并且\Quote{这不是木匠的儿子吗？} \BibleRef{玛 13:55} 而且每次他们都轻视祂。\par

\Quote{那么，}也许有人说，\Quote{祂是通过比较才比若翰大吗？}远非如此。因为当若翰说\Quote{祂比我更有力量}时，\BibleRef{玛 3:11} 他并不是在比较他们；保禄在提及梅瑟时写道\Quote{因为这人被算为比梅瑟配得更大的荣耀}，\BibleRef{希 3:3} 他也不是用比较的方式写的；并且祂自己说\Quote{看，这里有一位比撒罗满更大的}时，\BibleRef{玛 12:42} 也不是在作比较。\par

或者即使我们承认这是祂用比较的方式说的，那也是出于迁就，因为听者的软弱。因为这些人确实非常注目于若翰；而且他的监禁和对君王的直言使他更显卓越；目前重要的是，即使这样也能被群众接受。同样，旧约也以同样的方式纠正犯错者的灵魂，用比较的方式将不可比较的事物放在一起；例如，当它说：\Quote{在诸神中，没有像你一样的，上主；}又\Quote{没有神像我们的天主。}\par

现在有些人断言基督是指宗徒说的，另有些人又是指天使说的。当人们偏离真理时，他们往往会迷失在许多道路上。因为说天使或宗徒有什么联系呢？此外，如果祂是论宗徒，什么妨碍祂指名道姓地提出他们呢？然而，当祂论及自己时，祂自然隐藏了自己的位格，因为当时仍然存在的猜疑，并且为了使自己似乎不说什么伟大的事；是的，我们常发现祂这样行。\par

但\Quote{在诸天之国里}是什么意思？在属灵的存在中，以及在诸天中的一切中。\par

而且祂说\Quote{在女人所生者中，没有兴起比若翰更大的}，这是适合将自己与若翰对比的人说的，并以此默默地将自己除外。因为虽然祂也是女人所生，却不像若翰那样；因为祂不是纯粹的人，也不是像人一样出生的，而是以一种奇异而奇妙的出生方式。\par

4. \Quote{从洗者若翰的日子直到现在，诸天之国是受强力的，强力的人以暴力夺取它。}\par

这与前面所说的有什么关系呢？确实有很大关系，并且完全一致。因为祂也借此话题敦促并催促他们信从祂自己；同时，祂所说的也与若翰先前所说的相符。因为如果直到若翰一切事都应验了，我就是\Quote{那一位要来者。}\par

\Quote{因为所有的先知和法律都预言直到若翰。}祂说，\BibleRef{玛 11:13}\par

因为先知们不会停止，除非我来了。因此不要再期待什么，也不要等候任何人。因为我是那一位，这从先知们的停止，以及每天\Quote{以暴力夺取}对我的信仰的人就可以显明。因为如此明显和确实，甚至有许多人以暴力夺取它。请问，谁这样做了？告诉我。所有用心认真接近它的人。\par

然后祂又提出了另一个无误的记号，说：\Quote{如果你们愿意接受，他就是那将要来的厄里亚。}因为经上说：\Quote{我要派遣提斯贝人厄里亚到你们那里，他必使父亲的心转向儿子。}所以这个人就是厄里亚，如果你们仔细听，祂说。因为\Quote{我要派遣我的使者在你面前。}祂说。\par

祂说得很好：\Quote{你们若愿意接受}，以表明没有强迫。因为祂说：我不勉强。祂这样说，是要求一颗坦诚的心，并指明若翰就是厄里亚，厄里亚就是若翰。因为二者接受了同一使命，二者都成了前驱。因此祂不是简单地说：\Quote{这人是厄里亚}，而是说：\Quote{你们若愿意接受，这人就是厄里亚}，也就是说，如果你们以坦诚的心留意所发生的事。祂甚至不止于此，在说了\Quote{这人是将要来的厄里亚}之后，祂又加上，为表明需要理解：\Quote{有耳听的，听罢！} \BibleRef{玛 11:15}\par

祂用了这么多隐晦的话，是为了激发他们去探究。如果连这样都不能唤醒他们，那么若一切都是明白清楚的，就更不必说了。因为肯定没人会说他们不敢问祂，或祂难以接近。那些问祂问题、在普通事上试探祂、无数次被堵住嘴的人，仍未离开祂；如果他们真渴望学习，怎能不去追问祂、恳求祂关于那些必不可少的事呢？因为关于法律的事，他们曾问\Quote{哪条是第一条诫命}以及诸如此类的问题，虽然当然不需要祂告诉他们；那么对于祂自己所说的话，他们岂不该问其含义，而祂也必须在回答中给予说明吗？尤其当祂自己正鼓励和吸引他们去这样做时。因为祂说\Quote{暴力的人夺取了它}，是激发他们认真；祂又说\Quote{有耳听的，听罢}，也是在同样做。\par

5. \Quote{我可把这一代比作什么呢？}祂说，\Quote{它像坐在街市上的孩童，向同伴呼喊说：我们给你们吹笛，你们却不跳舞；我们向你们举哀，你们却不捶胸。}这似乎与前面不连贯，但却是最自然的引申。祂仍坚持同一个要点：证明若翰的行动与祂自己一致，尽管结果相反；正如关于他的询问也是如此。祂暗示，凡是为他们的得救所应做的事，没有一件被忽略；这正是先知论葡萄园所说的：\Quote{我为我葡萄园所应做的，还有什么没有做到？}因为祂说：\Quote{我可把这一代比作什么呢？它像坐在街市上的孩童，向同伴呼喊说：我们给你们吹笛，你们却不跳舞；我们向你们举哀，你们却不捶胸。因为若翰来了，也不吃，也不喝，他们便说：他附了魔。人子来了，也吃也喝，他们便说：看哪，一个贪吃嗜酒的人，税吏和罪人的朋友。}\par

祂的意思是这样：我和若翰各走了一条相反的路；我们就像猎人们对付一头难以捕捉的野兽，可能通过两条路落入陷阱；好像两人各从一条路截断它并驱赶它，彼此相对而立；这样它必然落入其中一个陷阱。请看，整个人类如何对禁食的奇迹和这种克苦自制的生活感到惊奇。正因如此，若翰才被安排从幼年起就这样受培养，以便借此（以及其他原因）使他的言语获得信任。\par

但有人会问：为什么祂自己不走那条路？首先，祂自己也走过这条路，当祂禁食四十天，四处教导，没有枕头的地方时。尽管如此，祂也用另一种方式实现了同一目标，并获得了由此而来的益处。因为被走这条路的人（若翰）所见证，与亲自走这条路是同一回事，甚至更伟大。\par

此外，若翰确实只展现了生活与为人；因为经上说：\Quote{若翰没有行过神迹} \BibleRef{若 10:41}；但祂自己也有神迹和奇迹的见证。因此，祂让若翰以禁食著称，而祂自己则走相反的路，甚至赴税吏的宴席，又吃又喝。\par

让我们问犹太人：\Quote{禁食是好事，值得钦佩吗？那么你们应听从若翰，接纳他，相信他的话。因为那些话会引导你们走向耶稣。禁食是痛苦、沉重的事吗？那么你们应听从耶稣，相信那位走相反道路而来的。这样，无论如何，你们都会发现自己进入了天国。}但你们像一头倔强的野兽，诋毁两者。因此，过错不在于不被相信的人，而在于不信的人。因为没有人会同时诋毁相反的事，也不会同时称赞它们。我的意思是：一个喜欢开朗自由性格的人，不会喜欢忧郁严肃的人；一个称赞忧郁面容的人，不会称赞开朗的人。因为你不可能同时赞成两者。所以祂说：\Quote{我们给你们吹笛，你们却不跳舞；}意思是：\Quote{我展现了更自由的生活，你们却不听从；} \Quote{我们向你们举哀，你们却不捶胸；}意思是：\Quote{若翰走的是严峻克苦的生活，你们却不留意。}祂不是说：\Quote{他这样做，我那样做，}而是因为两者的目的虽生活方式相反却是一致，所以祂将他们的行为说成是共同的。是的，因为他们走相反的路正是出于最精确的一致，始终着眼于同一个目标。既然如此，你们还有什么可推诿的呢？\par

为此祂接着说：\Quote{智慧由其子女得称为义。}这就是说，即使你们不被说服，以后也不能再指责我。正如先知论及圣父所说：\Quote{祢在祢的言语上得称为义。}因为天主，即使祂对我们的眷顾不再产生任何效果，仍尽祂的一切本分，以致给那些寡廉鲜耻之人连一丝不诚实的怀疑的借口都不留下。\par

若这些比喻卑微、听来不佳，不必惊讶，因为祂是顾及听众的软弱而发言的。既然厄则克耳也提及许多类似的比喻，且不配天主的威严，但这尤其适合于祂的慈父关怀。\par

且注意他们，在另一方面又如何被牵着走向自相矛盾的意见。因为他们曾论若翰说：\Quote{他有魔鬼。}\BibleRef{玛 11:18}他们仍不满足，又对祂说了同样的话，尽管祂所行的完全相反；如此，他们永远被牵着走，陷入冲突的意见。\par

但路加在此还记下另一条更重的指控反对他们，说：\Quote{因为税吏已领受了若翰的洗礼，就称天主为义。}\par

6. 然后，在智慧已得称为义，祂已显示一切全然成就之后，祂就开始责骂那些城市。也就是说，祂既未能说服他们，如今只是为他们哀叹；这比恐吓更甚。因为祂已藉祂的言语展示了祂的教导，藉祂的奇迹展示了祂行奇事的能力。但因他们固守自己的不信，祂如今只是责骂。\par

因为经记载说：\BibleQuote{那时，耶稣就开始责骂那些祂行过大部分异能的城市，因为他们没有悔改，说：\InnerQuote{苛辣匝因，你有祸了！贝特赛达，你有祸了！}}\BibleRef{玛 11:20}\par

然后，为向你们表明他们并非生来如此，祂还提及那城的名号，那城曾出了五位宗徒。因为斐理伯，以及两对首席宗徒，都出自那里。\BibleRef{若 1:44}\par

\BibleQuote{因为，如果在你们中所行的异能，行在提洛和漆冬，他们早已披麻蒙灰悔改了。但我告诉你们：在审判之日，提洛和漆冬所受的，比你们还要容易忍受。至于你，葛法翁，你本被高举到天上，却要被打入地狱；因为如果在你们中所行的异能，行在索多玛，它必存留至今。但我告诉你们：在审判之日，索多玛地所受的，比你们还要容易忍受。}\par

祂并非无故将索多玛与别的城市并列，而是为了加重对他们的指控。是的，因为这是一个极大的邪恶证明：不仅当今之人，就连所有历史上曾行恶的人之中，也找不到像他们那样坏的。\par

同样，祂在其他地方也作比较，用尼尼微人和南方女王来谴责他们；但那里是藉行义的人，这里却是藉犯罪的人，这是一件更为严重的事。厄则克耳也知道这个谴责的法则，所以他对耶路撒冷说：\BibleQuote{你在你的一切罪恶中，使你的姊妹显得正义。}因此，祂处处喜欢停留在旧约中，如同在喜爱的地方。即使如此，祂也没有停止发言，反而使他们的恐惧更加剧烈，说他们所应受的苦比索多玛人和提洛人更重，好藉各种方式把他们聚集起来，既用哀叹，又用警吓。\par

7. 我们也要听同样的事：因为祂不仅为不信者，也为我们设立了比索多玛人更重的惩罚，若我们不接待来到我们这里的陌生人；我是指，当祂命令要跺掉脚上的尘土时，这是很合宜的。因为索多玛人虽然犯了重罪，却是在法律和恩宠之前；但我们，在领受了如此多的眷顾之后，若显出如此的不接待，将需要帮助的人关在门外，并在门前堵住耳朵，我们又配得什么宽恕呢？或者，不仅仅是对穷人，而是对宗徒们本身？因为我们这样对待穷人，正是因为这样对待宗徒们。因为保禄被诵读，你们却不留心；若望宣讲，你们却不听；那你们何时会接待一个穷人呢？你们连宗徒都不接待。\par

为使我们的房屋能持续向穷人敞开，我们的耳朵向宗徒敞开，让我们从灵魂的耳朵中清除污秽。因为正如污秽和泥土堵塞我们肉身的耳朵，同样，娼妓的歌谣、世俗的新闻、债务、高利贷和借贷的生意，更堵塞心灵的耳朵，且比任何污秽更甚；不，它们不仅堵塞，还使之不洁。那些向你们讲述这些事的人，正是在把粪便放入你们的耳朵。而外邦人所恐吓的，说：\BibleQuote{你要吃自己的粪便，}以及后续的话；\BibleRef{依 36:12}这些人也叫你们在行为上，而不是言语上，经历同样的事；或者，甚至更糟。因为那些歌谣确实比这些更可憎；更糟的是，你们听到它们时，非但不感到烦恼，反而大笑，而你们本该厌恶并逃避它们。\par

但若它们并非可憎，你就下到戏台，模仿你所称赞的，或者，你只需和那引人发笑的人同行。不，你承受不了。那为何给他如此大的荣誉？是的，外邦人所制定的法律要羞辱他们，你却全城出动，像接待大使和将军一样接待他们，并召集众人，在耳朵里接收粪便。你的仆人，若在你面前说任何污秽的话，必受许多鞭打；无论是儿子、妻子，或任何人，若做了我所说的事，你会称之为冒犯；但若那些该受鞭打的卑劣之人，邀请你去听那些污言秽语，你非但不愤怒，反而喜悦喝彩。还有什么比这更愚昧的呢？\par

但是你自己难道从来没有说过这些卑贱的话吗？这有什么益处呢？或者更确切地说，这件事本身从何显明呢？因为如果你不说这些事，你听到它们时就根本不会笑，也不会这样热切地奔向那使你羞愧的声音。\par

请告诉我，你听到别人亵渎时，会感到高兴吗？你难道不战栗，并捂住耳朵吗？我想你一定会。为什么呢？因为你自己不亵渎。同样，对于污言秽语，你也应当如此；如果你想清楚表明你并不喜欢污言秽语，那就连听也不要听。因为当你耳中充满这些声音而长大时，你何时才能变得良善？你何时才敢承担贞洁所要求的劳苦？既然你现在正通过这笑声、这些歌曲和污词秽语渐渐滑落。是的，一个灵魂若能保持自身远离这一切，便能成为庄重贞洁的人，这是了不起的事；更何况一个在这种听闻中滋养长大的灵魂呢？你难道不知道，我们天性更倾向于邪恶吗？那么，当我们甚至把这当成一种技艺和事业时，我们何时才能逃出那火炉呢？\par

8. 难道你没有听到\Person{保禄}说，\BibleQuote{你们要喜乐于主？}\BibleRef{斐 4:4} 他没有说，\Quote{在魔鬼内。} 那么你何时才能听\Person{保禄}的话呢？何时才能意识到你的错误行为呢？你因我所谈论的场面而酩酊大醉，永无休止。因为你来这里并无什么奇妙或伟大之处，或者更确切地说，是奇妙的。因为你来这里只是随便应付一下，以满足一点顾虑，但到那里却是勤勉、迅速且十分乐意。从你回家时所带回的东西就明显可见。\par

因为那里所倾倒给你的一切污泥——通过言语、歌曲、笑声——你都收集起来，每个人都带回家，或者更确切地说，不只是带回家，更是带进自己的心灵。\par

对于不应憎恶的事，你反而回避；而对于应憎恶的事，你非但不恨，反而追求。例如，许多人从坟墓回来，习惯洗身；但从剧场回来，他们从未叹息，也未曾流出泪泉。然而，死人并不是不洁之物，而罪却造成如此污点，即使用千万道泉水也不能洗净，只能用眼泪和告解。但没有人意识到这个污点。因此，因为我们不惧怕应当惧怕的，所以我们躲避不应当躲避的。\par

再者，那鼓掌是什么？那喧哗、那属撒殚的呼喊、那魔鬼般的姿态是什么？首先，一个青年男子，蓄着长发，改变本性成女子，在面容、姿态、衣着以及一切方面，竭力模仿柔弱的少女。另一个年长者，则相反，剃发，束腰，在头发之前先割去他的羞耻，准备好被杖打，准备说任何话、做任何事。女人们又是不戴头巾，站立而毫不脸红，与众人交谈，如此习于无耻；于是将一切厚颜和不洁倾注到听众的灵魂中。她们唯一的目的是从根本上铲除一切贞洁，羞辱我们的本性，满足邪恶魔鬼的欲望。是的，还有污秽的言辞，更污秽的姿态；发型、步态、服装、声音、肢体的弯曲，都朝向那个方向；还有眼波的转动、笛子、风笛、戏剧、阴谋；总之，一切都充满最极端的不洁。那么，我请问，你何时才会清醒呢？现在魔鬼正为你大量倒出淫乱的烈酒，调制那么多不洁的杯？因为那里有奸淫和私奔，有妓女卖淫，男子卖身，少年自甘堕落：那里充满了一切不义、一切巫术、一切羞耻。因此，旁观者不应嘲笑这些事，而应痛哭和悲叹。\par

\Quote{那么，我们要关闭舞台吗？} 有人会说，\Quote{难道因你一句话，一切都要天翻地覆吗？} 不，但现在一切都已天翻地覆。因为请告诉我，那些谋害我们婚姻的人从何而来？难道不是来自这个剧场吗？那些挖墙进入内室的人从何而来？难道不是来自那个舞台吗？难道不是因此，丈夫变得使妻子无法忍受？因此，妻子受到丈夫的藐视？因此，大多数人成了奸夫？所以，颠覆一切的人就是去剧场的人；正是他带来了严酷的暴政。\Quote{不然，} 你将会说，\Quote{这是由法律的良好秩序所规定的。} 怎么，抢夺别人的妻子、侮辱少年、推翻家庭，那是占据城堡的人干的事。\Quote{哪一个奸夫，} 你会说，\Quote{是由这些表演造成的呢？} 不，谁没有成为奸夫呢？如果现在可以指名道姓，我可以指出那些娼妓使多少丈夫与妻子分离，掳掠了多少人，甚至将一些人从婚床上拉走，不让另一些人活在婚姻之中。\par

\Quote{那么，请问，我们要推翻所有法律吗？}不，而是推翻无法无天，如果我们废除这些表演。因为由此产生那些破坏我们城市的人；由此产生，例如，叛乱和骚乱。因为那些靠舞者供养、将自己的嗓子卖给肚腹、以喊叫和做尽一切怪诞之事为业的人，这些人正是煽动民众、在我们的城市制造骚乱的人。因为青年若与闲懒携手，在如此重大的恶中成长，就变得比任何野兽更凶猛。还有那些召唤亡灵的人，请问，他们从何而来？岂不正是因此，为了刺激那些无所事事闲荡的人，让舞者获得大量响亮的掌声，并加固妓女以对抗贞洁者，他们在邪术中走得如此之远，甚至不惜惊扰死者的骸骨？岂不正是因此，当人们被迫无限制地花费在那邪恶的魔鬼合唱团上？还有淫荡，以及它无数的祸害，从何而来？你看，是你通过吸引人们做这些事在颠覆我们的生活，而我通过制止它们来修复生活。\par

\Quote{那么，让我们拆毁舞台，}他们说。但愿能拆毁它；或者，如果你愿意，就我们而言，它已被拆毁并挖起。然而，我并不断言这样的事。这些场所既然存在，我要求你使它们失效；这比拆毁它们更值得称赞。\par

至少请效仿外邦人，如果没有别人；因为他们确实完全远离寻求这样的景象。那么，我们这些天国的公民、与革鲁宾的合唱团共舞、与天使共融的人，若在这方面甚至比外邦人更差，而当无数其他更优的乐趣在我们伸手可及时，我们还有什么借口？\par

因为，如果你希望你的灵魂得到快乐，就去游乐场，到流过的河边，到湖边，观赏花园，聆听蟋蟀歌唱，常去殉道者的墓旁，那里有身体的健康和灵魂的益处，没有伤害，没有像这里一样在快乐后的懊悔。\par

你有妻子，你有儿女；什么能与这快乐相比？你有房屋，你有朋友，这些才是真正的乐趣：除了它们的纯洁，它们还带来极大的益处。因为，请问，有什么比儿女更甜蜜？有什么比妻子更甜蜜，对于内心贞洁的人？\par

为此，我们听说，外邦人在某个场合说出了一句充满智慧严厉的话。我的意思是，他们听说了这些邪恶的表演和它们不合时宜的快乐后，\Quote{为何罗马人，}他们说，\Quote{设计了这些娱乐，好像他们没有妻子和儿女似的；}暗示没有什么比儿女和妻子更甜蜜，如果你愿意正直地生活。\par

\Quote{那么，}有人会说，\Quote{如果我指出有些人，他们一点也没有被那里的消遣所伤害呢？}首先，即使这也是一种伤害，即无目的地浪费时间并成为他人的绊脚石。因为即使你没有受伤，你也使别人对此更加热衷。而你自己怎么可能不受伤呢？你给正在发生的事提供了借口？是的，占卜者、男妓、妓女，以及所有那些魔鬼的合唱团，都把他们所作所为的罪责归到你头上。因为如果没有观众，就不会有人从事这些行当；同样，既然有观众，他们也要分担这些行为所应得的火。所以，即使在贞洁上你完全没有受伤（这是不可能的），但你仍要因他人的毁灭而作出严重的交代：包括观众，以及那些召集他们的人。\par

在贞洁上你也会获得更多益处，如果你不去那里。因为即使现在你是贞洁的，避开这些景象会使你更贞洁。那么，我们不要沉溺于无用的争论，也不要设计无益的辩解：只有一个辩解，就是逃离巴比伦的火窑，远离埃及的妓女，即使必须赤身逃脱她的手。\BibleRef{创 39:12}\par

这样，我们就必享受许多快乐，良心不责备我们，并且我们要以贞洁度此现世生活，达到将来的善，藉着我们的主耶稣基督的恩宠和对人的爱；愿光荣和权能归于他，从现在直到永远，永无穷尽。阿们。\par

\SourceRecord{Translated by George Prevost and revised by M.B. Riddle.；Nicene and Post-Nicene Fathers, First Series；Vol. 10.；Edited by Philip Schaff.；Buffalo, NY: Christian Literature Publishing Co.,；1888.；<http://www.newadvent.org/fathers/200137.htm>.}

% structure: p0001 source=h1 target=section reason=page-title

\section{玛窦福音讲道集第三十八篇}

\BibleRef{玛 11:25-26}\par

\Quote{\Emph{在那时，耶稣回答说：我承认祢，} \Emph{圣父，天地的主宰；因为祢将这些事瞒住了智慧明达的人，而启示给婴孩。是的，圣父，因为祢喜欢这样。}}\par

你看见祂用多少方式引导他们进入信德吗？第一，\BibleRef{玛 11:7}藉着称赞若翰。因为指出他是伟大而奇妙的人，祂也证明了他的一切话都是可信的，而他常用这些话引导他们认识祂。第二，\BibleRef{玛 11:12}藉着说：\Quote{天国遭受暴力，暴力的人以武力夺取它。}因为这是督促和催促他们的话。第三，\BibleRef{玛 6:13}藉着指明先知的时代已经终结；因为这也显示祂就是那预传的人。第四，\BibleRef{玛 6:14}藉着指出凡由祂所行的事都已成就；那时祂也提及了孩童的比喻。第五，藉着责备那些不信的人，并用极大的警戒和威胁恐吓他们。\BibleRef{玛 11:20}第六，藉着为那些相信的人感谢。因为\Quote{我承认祢}这句话在此就是\Quote{我感谢祢。}\Quote{我感谢祢，}祂说，\Quote{因为祢将这些事瞒住了智慧明达的人。}\par

那么，祂喜悦人的灭亡吗？喜悦别人没有获得这认识吗？断然不是；但这是祂拯救人的最卓越的方式：不强迫那些完全拒绝、不愿接受祂话的人；因为他们没有因祂的召唤而变得更好，反倒退后并轻蔑它，那么祂的抛弃他们能使他们渴求这些事。同样，留心的人也会变得更热切。\par

而祂向这些人启示是喜乐的事，向那些人隐瞒则不再是喜乐，而是眼泪。至少祂为城哀哭时就是如此行。因此祂并非为此喜乐，而是因为智慧人所不知道的，这些人却知道了。正如保禄所说：\Quote{我感谢天主，你们曾是罪的奴仆，却从心里服从了那传给你们的教义规范。}你看，保禄也不是因他们是\Quote{罪的奴仆}而喜乐，而是因为既是这样的人，却蒙受了如此大恩。\par

此处祂所称的\Quote{智慧人}是指经师和法利塞人。祂说这些话，为使门徒更热切，并显明当其他人都错过了时，渔夫们获得了何等恩典。称他们为\Quote{智慧人}，祂并非指真实可赞的智慧，而是他们因天生的机敏而看似拥有的。因此祂也没有说\Quote{祢启示给愚者}，而是\Quote{启示给婴孩}；即天真单纯的人；祂暗示他们并未不配地错过这些特权，反倒是理所当然的。祂也教导我们始终远离骄傲，追求纯朴。为此，保禄也更恳切地写道：\Quote{若有人在这世上自以为有智慧，就让他成为愚者，好成为有智慧。}\BibleRef{格前 3:18}因为天主的恩宠就这样彰显出来。\par

但祂为何感谢圣父，尽管这本来是祂自己成就的？正如祂祈祷并恳求天主，显示祂对我们的极大爱德，同样祂也这样做：因为这也是出于大爱。祂表明，他们不仅离弃了祂，也离弃了圣父。因此，祂对门徒所说的\Quote{不要把圣物给狗}\BibleRef{玛 7:6}，祂自己就先行了。\par

此外，祂在此表明祂自己的主要旨意和圣父的旨意：祂的旨意，是因祂对所发生的事感谢和喜乐；圣父的旨意，是暗示圣父也并非受恳求而做，而是出于自己的旨意自己做的；\Quote{因为祢喜欢这样}，祂说，即\Quote{祢乐意这样。}\par

为什么这事向他们隐瞒？听保禄说：\Quote{因为他们企图建立自己的义，就不服从天主的义。}\BibleRef{罗 10:3}\par

现在想想，门徒听到这些话会有什么感觉：智慧人所不知道的，这些人知道了，而且作为婴孩继续知道，并藉着天主的启示而知道。但路加说，\Quote{就在那时}，当七十位回来向祂报告关于魔鬼的事时，祂\Quote{欢喜}并说了这些话，\BibleRef{路 10:21}这不但增加他们的勤勉，也使他们谦逊。意思是，因为他们自然以驱魔为傲，祂就借此抑制他们：无论做了什么，都是启示，并非他们的努力。因此，经师和智慧人自以为聪明，却因自己的虚荣而失落。那么，如果因此事向他们隐藏，\Quote{你们也当恐惧，并继续作婴孩。}因为这让你们获得了启示的益处，正如相反的事使他们失去了它。因为当祂说\Quote{祢隐藏了}时，绝非意指全是天主的作为；而是如同保禄说\Quote{任他们陷于堕落的心思}\BibleRef{罗 1:28}和\Quote{祂蒙蔽了他们的心思}时，并非指祂是行这事者，而是指那些给与机会的人；同样，这里祂也用了\Quote{祢隐藏了}这句话。\par

因为祂曾说：\Quote{我感谢祢，因为祢隐藏了这些事，却启示给婴孩。}为使你不以为祂自己缺乏这种能力，无法成就这事，才如此感谢，祂说：\par

\Quote{我父将一切交给了我。} \BibleRef{玛 11:27} 又对那些因魔鬼服从他们而欢喜的人，祂说：\Quote{不，你们为什么惊奇？} \BibleRef{路 10:22} 魔鬼屈服于你们？一切都是我的；\Quote{一切都被交给了我。}\par

但当你听到\Quote{它们被交付了}时，不要猜测任何属人的事。因为祂使用这个表达，是为了防止你们想象两个\OriginalTerm{非受生的神}{unoriginate Gods}。因为祂同时是受生的和万有的主，祂在许多方面和其他地方也宣告了这一点。\par

2. 然后祂说了比这更伟大的事，提升你们的思念；\Quote{除了父外，没有人认识子；除了子外，也没有人认识父。} 这在无知者看来似乎与前面无关，但实际上完全一致。就是说，祂说\Quote{我父将一切交给了我}之后，又加上\Quote{这有什么可惊奇的？} 祂如此说，\Quote{如果我是万有的主？我还有一个更大的特权，就是认识父并与祂\OriginalTerm{同体}{same substance}。} 是的，祂也通过自己是唯一这样认识祂的人来隐秘地表示这一点。因为当祂说\Quote{除了子，没有人认识父}时，就是这个意思。\par

且看祂在什么时候说这话。当他们因祂的作为获得了祂能力的确实证明时，不仅看见祂行奇迹，而且因祂的名被赋予如此大的权能。然后，因为祂曾说过\Quote{你将这些事启示给了小孩子}，祂表明这同样与祂自己有关；因为\Quote{也没有人认识父}，祂说，除了子，以及子愿意启示给谁的人；\BibleRef{玛 11:27} 不是\Quote{他可能被命令指示给谁}，\Quote{他可能被命令给谁}。但如果祂启示了祂，那么祂自己也是如此。但祂把这当作公认的而略过，却记下了另一项。祂在每一处都证实了这一点；正如祂所说：\Quote{除非藉着我，没有人能到父那里去。} \BibleRef{若 14:6}\par

并且由此祂确立了另一点，即祂与父和谐一致，同心同德。\Quote{为什么？}祂说，\Quote{我如此远离与祂争战冲突，以至除非藉着我，没有人能到祂那里去。} 因为最令他们冒犯的是祂似乎是一位敌对的神，所以祂用尽一切方法消除这一点；祂对此的热心丝毫不亚于甚至超过对祂奇迹的热心。\par

但当祂说\Quote{除了子，没有人认识父}时，祂的意思不是众人都对祂无知，而是说没有人以祂认识父的那种认识来认识父；这也可以对子这么说。因为祂所说的并不是像\Person{马西昂}所主张的那样，是一位无人认识、从未启示给人的神；而是祂在这里暗示知识的圆满，因为我们也不认识子如同祂所应被认识的那样；而且这一点，不用再多说，\Person{保禄}也宣告过，当他说\Quote{我们现在所知道的，只是局部的；我们作先知所讲的，也只是局部的。} \BibleRef{格前 13:9}\par

3. 接着，借着祂的言语使他们生出热切的渴望，并表明了祂不可言喻的权能之后，祂就邀请他们说：\Quote{凡劳苦和负重担的，你们都到我这里来，我要使你们安息。} \BibleRef{玛 11:28} 不是这个或那个人，而是所有在焦虑中、在悲伤中、在罪中的人。来吧，不是我要追究你们，而是我要消除你们的罪；来吧，不是我要你们的尊荣，而是我要你们的救恩。\Quote{因为我，}祂说，\Quote{要使你们安息。} 祂没有只说\Quote{我要拯救你们}，而是说了更重要的：\Quote{我要使你们完全安稳。}\par

\Quote{你们背起我的轭，跟我学吧！因为我是良善心谦的：这样你们必要找得你们灵魂的安息。因为我的轭是柔和的，我的担子是轻松的。} \BibleRef{玛 11:29} 所以，\Quote{不要害怕，}祂说，一听到说轭，因为它是柔和的：不要怕，因为我说了\Quote{担子}，因为它是轻松的。\par

而祂先前怎么说？\Quote{门是窄的，路是狭窄的？} \BibleRef{玛 7:13} 当你们疏忽时，当你们懈怠时；反之，如果你们妥善实行祂的话，担子就会轻松；因此祂现在也这样称呼它。\par

但怎样妥善实行呢？如果你们成为谦卑、温良、和善的人。因为这种美德是生活一切严谨之母。所以，在开始那些神圣诫命时，祂也是从这一点开始的。\BibleRef{玛 5:3} 而这里祂又做了同样的事，祂指定的赏报是极其伟大的。\Quote{因为你不只是有用于别人；也首先使你自己得到安息，}祂说。\Quote{因为你们必要找到灵魂的安息。}\par

甚至在将来的事物之前，祂就在这里赐予你们赏报，并预先给了你们报酬，使这话可接受，既藉此，也藉着把祂自己作为榜样。因为，\Quote{你们害怕什么？}祂说，\Quote{怕因你们的卑微而成为失败者吗？注视我，注视我所有的一切；向我学习，然后你们就会清楚地知道你们的祝福何其大。} 你们看祂是如何在各方面带领他们走向谦卑？藉着祂自己的所作所为：\Quote{向我学习，因为我是良善心谦的。} 藉着他们自己将获得的：因为祂说\Quote{你们必要找到}灵魂的安息。藉着祂赐予他们的：因为祂说\Quote{我也要使你们安息。} 藉着使它变得轻松：\Quote{因为我的轭是柔和的，我的担子是轻松的。} 同样地，\Person{保禄}也说：\Quote{因为我们这现时轻微的苦难，正分外无比地给我们造就永远的光荣厚报。} \BibleRef{格后 4:17}\par

有人或许会说：这重担怎能是轻的，当祂说：\Quote{除非人恼恨父亲和母亲；}又说：\Quote{不论谁，若不背起自己的十字架跟随我，就不配是我的；}还说：\Quote{不论谁，若不舍弃他所有的一切，不能做我的门徒。}祂甚至命令我们舍弃自己的生命？\BibleRef{玛 16:25} 让保禄教导你，他说：\Quote{谁能使我们与基督的爱隔绝？是困苦吗？是窘迫吗？是迫害吗？是饥饿吗？是赤身露体吗？是危险吗？是刀剑吗？\BibleRef{罗 8:35} } 又说：\Quote{现今的苦楚，与将来在我们身上要显示的光荣，是不能比拟的。\BibleRef{罗 8:18} } 让那些从犹太公会回来、满身鞭痕的人教导你，他们\Quote{因被算是配为这名受辱而欢喜。\BibleRef{宗 5:41} } 如果你仍然害怕，一听到轭和重担就发抖，那恐惧并非来自事情的本性，而是来自你的懈怠；因为只要你准备好并认真，一切对你都会容易而轻省。为此，基督也要我们亲自劳苦，祂不仅提到恩惠的事就闭口不言，也不仅提到痛苦的事，而是把两者都写下来。因此，祂既提到了\Quote{轭}，又称它为\Quote{轻省的}；既命名了重担，又加上它是\Quote{轻的}；好叫你们既不因困难而逃避，也不因过于容易而轻视。\par

但即使如此，若德行对你仍显得烦难，就当想到恶习更为烦难。这一点祂也暗示了，因为祂先不说\Quote{背起我的轭在你们身上}，而是先说\Quote{凡是劳苦和负重担的，你们来吧}，从而表明罪也有劳苦，有沉重而难担的重担。因为祂不仅说\Quote{劳苦}，还说\Quote{负重担}。先知论及罪的这种性质时也说：\Quote{如同沉重的担子，重重地压在我身上。}匝加利亚描述她时，称她是\Quote{一块铅}。\BibleRef{匝 5:7}\par

经验本身也证实了这一点。因为没有什么像罪恶感那样沉重地压迫和压伏灵魂，没有什么像获得义德和德行那样扶助灵魂、提升它向上。\par

请看吧：请问，有什么比没有财产更令人痛苦？转过另一边面颊，被人击打而不还手？死于横死？然而，如果我们克己，这一切都是轻省、容易并且令人愉悦的。\par

但不要烦乱；让我们逐一考察，认真查问；如果你愿意，先看看那件被许多人视为最痛苦的事。告诉我，哪一件更痛苦、更沉重？是只操劳一个肚腹，还是为成千上万的东西操心？是只穿一件外衣，不再寻求别的，还是家里有许多衣服，却日夜在恐惧、战栗中为保存它们而哀叹，为失去它们而忧虑、窒息？怕一件被虫蛀，怕仆人偷窃卷走？\par

4. 但无论我说什么，我的言辞都不如实际经验那样能提供证明。因此，我愿在此有那已达到克己顶峰的人与我们同在，这样你就能确知其中的喜乐；并且，那些喜爱自愿贫穷的人，即使有人拿出万贯财富，也不会接受。\par

你或许说：这些人会愿意变得贫穷，抛弃他们拥有的焦虑吗？那又怎样？这只能证明他们的疯狂和严重疾病，而不是证明这件事有什么令人愉快之处。这一点连他们自己也会向我们作证，他们每日为这些焦虑而哀叹，认为自己的生命不值得活。但那些克己的人并非如此；相反，他们欢笑、雀跃，戴王冠的人也不如他们在贫穷中那样荣耀。\par

同样，转过面颊，对于留心的人来说，比击打别人更不痛苦；因为争斗由此开始，由此终结；前者使你点燃了对方的火焰，后者却甚至扑灭了你自己的火焰。不被烧比被烧更愉快，这对每个人都是显然的。如果这在身体上成立，那么在灵魂上更是如此。\par

哪一样更轻省：是争斗还是加冠？是战斗还是得奖赏？是忍受波浪还是进入港湾？因此，死也比生更好。因为死使我们脱离波浪和危险，而生却增添它们，使人置身于无数的阴谋和痛苦之中，这些已使生命在你们眼中不值得活。\par

如果你们不信我们的话，请听那些曾目睹殉道者在争斗时面容的人：他们如何被鞭打、剥皮，却极其喜乐欢欣，被放在烙铁上时，仍满心欢喜，胜过躺在玫瑰花床上的人。因此，保禄在即将离世、以横死结束生命时，也说：\Quote{我喜乐，也与你们一同喜乐；同样，你们也要喜乐，与我一同喜乐。}你们看，他以何等强有力的言辞邀请普世分享他的喜乐！他深知自己的离世是何等美善、可爱、值得祈求，那可怕的事——死亡。\par

5. 但此德行的轭是甘饴且轻省的，这从许多其他方面也显然可见；但为了总结，如果你愿意，也让我们看看罪的负担。那么就让我们举出贪婪的人、无耻交易的零售商和二手商贩。还有什么比这些交易更沉重的负担呢？从这些收益中，每天滋生出多少忧愁、多少焦虑、多少失望、多少危险、多少阴谋和战争？多少烦恼和扰动？因为正如人永远不会看到无浪的海洋，同样也永远不会看到没有焦虑、沮丧、恐惧和扰乱的灵魂；是的，第二个接踵而至，又有其他涌来，当这些尚未止息时，又有其他的兴起。\par

或者你愿看见辱骂者和暴躁者的灵魂吗？什么比这种折磨更痛苦呢？什么比他们内心的创伤更甚呢？什么比那不断燃烧的熔炉和永不熄灭的火焰更甚呢？\par

或者你愿看见好色者和留恋此生者吗？什么比这种束缚更痛苦呢？他们过着加音的生活，在每一次死亡来临时都生活在不断的颤抖和恐惧中；死者的亲属不像这些人为自己的结局那样哀伤。\par

再者，什么比骄傲自满者更加充满动荡和疯狂呢？\BibleQuote{因为你们该向我学习，}祂说，\BibleQuote{我是良善心谦的，这样你们必要找得灵魂的安息。}\BibleRef{玛 11:29}因为坚忍是所有美善之事的母亲。\par

因此，你不要害怕，也不要躲避那使你从这一切中得轻省的轭，而要尽心竭力地服在它之下，这样你就会切实尝到其中的喜乐。因为它绝不压伤你的颈项，只是为秩序的缘故而加在你身上，劝你循规蹈矩地行走，引领你进入王道，从两边的悬崖中拯救你，使你在窄路上轻松行走。\par

既然它的益处如此之大，它的安全如此之大，它的喜乐如此之大，就让我们倾尽全灵、竭尽所能地负起这轭；好使我们在此世\BibleQuote{找得灵魂的安息}\BibleRef{玛 11:29}，并将来世的美善，赖我们的主耶稣基督的恩宠和慈爱，愿光荣和权能归于祂，自今而往，及于万世。阿们。\par

\SourceRecord{Translated by George Prevost and revised by M.B. Riddle.；Nicene and Post-Nicene Fathers, First Series；Vol. 10.；Edited by Philip Schaff.；Buffalo, NY: Christian Literature Publishing Co.,；1888.；<http://www.newadvent.org/fathers/200138.htm>.}

% structure: p0001 source=h1 target=section reason=page-title

\section{玛窦福音释义 第三十九篇}

\BibleRef{玛 12:1}\par

\BibleQuote{\Emph{那时，耶稣在安息日由麦田中经过；祂的门徒饿了，开始掐麦穗吃。}}\par

但路加却说：\Quote{在\OriginalTerm{双重安息日}{double Sabbath}。}什么是双重安息日？即停止劳苦是双重的：既因常规的安息日，又因适逢另一个庆日。因为他们称一切停止劳苦为安息日。\par

然而，那预知万事的主，为何不领他们避开呢？除非祂愿意安息日被违犯。祂的确愿意，但并非无理由；因此祂从无故违犯安息日，而是给出合理的托辞：以便既结束法律，又不惊动他们。但有些场合祂甚至直接废除安息日，而不加条件：例如当祂用泥土涂抹瞎子的眼睛时（\BibleRef{若 9:6}）；又如当祂说：\BibleQuote{我父作工直到现在，我也作工。}祂这样做，一方面为光荣自己的父，另一方面为安抚犹太人的软弱。在此祂正致力于后者，以自然的迫切需要作为借口；虽然在公认的罪行中，这当然绝不能成为托词。因为杀人者不能以他的愤怒为托词，通奸者不能以他的情欲为托词，不，任何其他理由也不行；但在这里，祂提到他们的饥饿，就使他们完全免责。\par

但我愿你钦佩门徒们，他们如何完全自制，不将身体之事放在心上，将肉体的餐桌视为次要之事，虽然不得不与持续不断的饥饿作斗争，却甚至不因此而退缩。因为若非饥饿严厉地逼迫他们，他们连这一点也不会做。\par

法利塞人于是怎样呢？经上记着：\BibleQuote{他们看见了，就对祂说：\InnerQuote{看，祢的门徒做了安息日不许做的事。}}（\BibleRef{玛 12:2}）\par

在此他们确无大的激动（然而这在他们本是合理的），但他们并未被激怒，而只是责备。但当祂伸出枯干的手并治好了它时，他们便如此狂怒，竟至商议要杀害并消灭祂。因为当没有伟大高贵的事发生时，他们平静；但当他们看到有人痊愈，便野蛮、烦恼，没有人比他们更难容忍：他们乃是人类得救如此之敌。\par

耶稣如何为门徒辩护呢？祂说：\BibleQuote{你们没有念过达味在圣殿里，当他和他的人饥饿时做了什么吗？他怎样进了天主的殿，吃了陈设饼，这饼不是他和他的人可以吃的，惟独司祭才可以吃。}（\BibleRef{玛 12:3}）\par

如此，为门徒辩护时，祂援引达味；为祂自己，则援引圣父。（\BibleRef{若 5:17}）\par

注意祂责备的口吻：\BibleQuote{你们没有念过达味做了什么吗？}因为那位先知实在伟大，以致伯多禄后来向犹太人辩护时也这样说：\BibleQuote{诸位仁人弟兄，容我放胆对你们论及\OriginalTerm{圣祖达味}{the patriarch David}，他死了，也埋葬了。}（\BibleRef{宗 2:29}）\par

但祂为何无论在此或以后都不以他的爵位称呼他呢？也许因为祂出自他的后裔。\par

倘若他们是坦诚之人，祂便会转向门徒饥饿受苦的话题；但他们既可恶又无人性，祂便向他们重述一段历史。\par

但马尔谷说：\BibleQuote{在\OriginalTerm{大司祭}{High Priest}阿彼雅塔尔的日子里。}他并非说与历史矛盾的话，而是暗示他有双名；并加上\BibleQuote{他给了他，}表明在这事上达味更有话说，因为连司祭本人也容许他；非但容许，甚至服侍了他。因为你们不要说达味是先知，即便如此也不合法，这特权是属于司祭的；因此祂又加上\BibleQuote{惟独司祭才可以。}因为即使他是万倍先知，也不是司祭；即使他本人是先知，与他同来的人却不然；我们知道他也给了他们。\par

\Quote{那么，}或许有人说，\Quote{他们与达味都是一样吗？}何必对我谈尊严，那里似乎有法律上的违犯，即使是自然的需要？是的，如此祂更彻底地解除了对他们的指控，因为更大的人被发现做了同样的事。\par

\Quote{这与问题何干？}有人可说，\Quote{难道达味违犯的是安息日吗？}你告诉我更大的事，这尤其显出基督的智慧：祂放下安息日，提出了比安息日更大的例子。因为违犯一日，与触摸那任何人不得触摸的圣桌，绝不相同。因为安息日确实曾被违犯，而且屡次；不，它不断地被违犯，无论是因割礼还是许多其他工作；在耶里哥（\BibleRef{苏 6:15}）也能看到同样的事发生；但这事却只发生在那时。所以祂更完全获得了胜利。那么，为何无人责备达味呢？尽管还有比这更重的指控，即司祭的谋杀，其根源就在此事？但祂不提它，只针对当前问题。\par

2. 随后祂又以另一种方式驳斥他们。因为正如祂起初引述达味，以人物的尊贵压下他们的骄傲；同样，当祂堵住他们的口、压下他们的夸耀之后，便进而加上更为贴切的驳斥。这驳斥是什么呢？\BibleQuote{你们不知道，在圣殿里司铎们亵渎了安息日，而无罪吗？}\BibleRef{玛 12:5} 因为前一个例子，祂说，是紧急情况造成了放宽，但这里却是毫无紧急情况下的放宽。然而祂并没有立刻就这样驳斥他们，而是先以让步的方式，然后才坚持自己的论点。因为理当最后才引出更强的推论，尽管先前的论证当然也有其应有的分量。\par

因为你们不要对我说，提出另一个犯了同样罪的人并不能使自己免责。因为当行此事者不受指责时，他所冒险的行为便成为他人辩解的惯例。\par

尽管如此，祂并不满足于此，还加上了更有决定性的论述，说这行为根本不是罪；而这比任何事都更标志着光荣的胜利，即指出法律自我废止，并且以两种方式如此行：首先通过地点，其次通过安息日；或者甚至可以说三种方式：因为所做的工作是双重的，而且还有另一件事，即这事是由司铎们做的；更有甚者，它甚至不被当作指控。\Quote{因为，}\ 祂说，\Quote{他们是无罪的。}\par

你们看祂陈述了多少要点？地点，因为祂说：\Quote{在圣殿里；}\ 人物，因为他们是\Quote{司铎们；}\ 时间，因为祂说：\Quote{安息日；}\ 行为本身，因为\Quote{他们亵渎；}（祂没有说\Quote{他们违犯，}\ 而是更严重的\Quote{他们亵渎；}）他们不仅逃脱惩罚，而且甚至免于责备，\Quote{因为，}\ 祂说，\Quote{他们是无罪的。}\par

所以，祂说，你们不要把这当作前一个例子。因为前一个例子既只发生一次，又非由司铎所行，而且是出于必要；因此他们理应得到原谅；但后一个例子则是每个安息日都做，由司铎在圣殿里，并且合乎法律。因此，他们再次不是出于恩惠，而是以合法的方式被免于指控。因为我这样说，绝非指责他们，祂说，也不是以宽恕的方式免其罪责，而是按照正义的原则。\par

祂似乎确实在为他们辩护，但实际上是在为祂的门徒洗脱所指控的过错。因为当祂说，\Quote{那些是无罪的}\ 时，祂的意思是，\Quote{这些人更无罪。}\par

\Quote{但他们不是司铎。}\ 不，他们比司铎更大。因为圣殿的主亲自在此：这是真理，而非预象。因此祂还说：\par

\BibleQuote{但我告诉你们，在这里有一位比圣殿更大的。}\par

然而，尽管他们听到了如此伟大的言论，他们却没有回答，因为救赎人类并不是他们的目的。\par

然后，因为这对听众来说似乎过于严厉，祂迅速掩盖了这一点，像先前一样使祂的讲论转向温和，然而即便如此，祂仍以责备的方式表达自己。\BibleQuote{但你们若明白这话的意思：\InnerQuote{我喜欢仁爱胜过祭献}，你们就不会定罪无罪的人了。}\BibleRef{玛 12:7}\par

你们看到祂如何再次将祂的话引向温和，却又表明他们已超出温和的范围？\Quote{因为你们不会定罪，}\ 祂说，\Quote{无罪的人。}\ 先前祂从关于司铎的话中推断出同样的事，说\Quote{他们是无罪的；}\ 但在这里祂以自己的权威陈述；或许这也是出自法律，因为祂引用了一句先知的话。\BibleRef{欧 6:6}\par

3. 此后祂又提及另一个理由；\BibleQuote{因为人子，}\ 祂说，\BibleQuote{是安息日的主；}\BibleRef{玛 12:8} 这是指着祂自己说的。但马尔谷记述祂也说过这涉及我们共同的本性；因为祂说：\BibleQuote{安息日是为了人而设立的，并不是人为了安息日。}\BibleRef{谷 2:27}\par

那么，为何那个拾柴的人受了惩罚？\BibleRef{户 15:32} 因为如果法律从一开始就被藐视，那么以后当然几乎不会被遵守。\par

的确，起初安息日带来了许多莫大的益处；例如，它使他们对自己家中的人温和仁慈；它教导他们天主的眷顾和创造，正如厄则克耳所说；\BibleRef{则 20:12} 它逐步训练他们远离邪恶，并使他们倾向于关注属神的事。\par

因为当他们不能承受时，如果祂在制定安息日法律时说：\Quote{你们在安息日做善事，但不要做恶事，}\ 因此祂禁止他们一切工作，说：\Quote{你们不可做任何工作；}\ 祂说：即使如此，他们也没有被约束住。但祂自己在制定安息日法律的同时，也藉着这句话隐约表明祂只希望他们远离恶事：\Quote{你们不可做工，除非是为了你们的性命。}\ 而且在圣殿里一切照常进行，并且更加殷勤、加倍辛劳。\BibleRef{户 28:9} 就这样，祂甚至藉着影子悄悄地向他们开启了真理。\par

或许有人会说：基督难道废除了一项如此有益的事吗？远非如此；祂反而大大提升了它。因为那时他们应当藉更高的规则在一切事上受到训练，而且那从邪恶中得释放、为一切善工展翅高飞的人，他的手无需被捆绑；也无需由此让人知道天主创造了万物；也无需让人如此变得温良，因为他们是蒙召效法天主对人类自身的慈爱（因为祂说：\Quote{慈悲，如同你们的天父一样}）；\BibleRef{路 6:36} 也无需让人把一天当作庆节，因为他们奉命终身守庆节；（因为\Quote{你们应当守庆节}，经上说，\Quote{不可用旧酵，也不可用奸恶和邪毒的酵母，但要用真诚和真理的无酵饼}）；\BibleRef{格前 5:8} 同样，那些拥有万有的主住在他们里面、并在一切事上与祂相通（藉祈祷、奉献、圣经、施舍，并让祂在他们内）的人，也无需站在约柜和金祭台旁。看，如今，对于那常常守庆节、他的思念在天上的人，为什么还需要安息日呢？\par

4. 那么，让我们不断守庆节，不做任何恶事；因为这就是庆节。让我们的属灵之事变得强烈，而属世之事退避。让我们安息于属灵的安息，节制双手不贪财，使我们的身体远离多余和无益的劳苦，远离古时希伯来人在埃及所忍受的那种劳苦。因为聚集金子的人，与那些被捆绑在泥浆中、制作砖块、收集碎秸、挨打的人，并无区别。是的，因为如今魔鬼也像当时的法郎一样，命令我们制作砖块。金不是泥浆吗？银不是碎秸吗？至少，碎秸燃起欲望的火焰；金也像泥浆一样玷污拥有它的人。\par

因此，祂没有从旷野派遣梅瑟，而是从天上派遣了祂的圣子。那么，如果祂来了之后，你仍留在埃及，你将与埃及人一同受苦；但如果你离开那地，与属灵的以色列人一同上去，你将看到一切奇迹。\par

然而，连这也不足以获得救恩。因为我们不仅必须被救出埃及，还必须进入应许之地。因为正如保禄所说，犹太人也曾穿越红海，吃玛纳，喝属灵的饮料，但他们全都灭亡了。\par

所以，免得我们同样遭遇，让我们不要迟延，也不要退缩。但当你甚至现在听到邪恶的探子对狭窄和艰难的道路提出恶意的报告，说出与古时那些探子同类的话时，不要让群众、而要若苏厄和耶孚乃的儿子加肋布成为我们的榜样。你不要放弃，直到你获得了应许，进入了天上。\par

也不要认为旅途艰难。\Quote{因为我们作仇敌的时候，尚且藉着圣子的死得与天主和好；既已和好，就更要因祂的生命得救了。} \BibleRef{罗 5:10} 或许有人说：\Quote{这条路是狭窄而艰难的。} 然而，你来的那条路不仅狭窄艰难，而且不可通行，充满了凶猛的野兽。正如若没有那个奇迹，就无法穿越红海；同样，我们若停留在过去的生活中，也无法升到天上，除非藉着圣洗圣事介入。如果不可能的事变为可能，那么艰难的事必定更容易。\par

或许有人说：\Quote{但那只是恩宠。} 正因为如此，你更有理由鼓起勇气。因为如果在那只是恩宠的地方，祂与你一同工作，那么当你们也显出劳苦的功绩时，祂岂不更加帮助你们？如果祂在你一无所作时拯救了你，那么当你工作时，祂岂不更加帮助你？\par

上文我确实说过，你应当从不可能的事中对困难的事也鼓起勇气；但现在我补充说，如果我们警觉，这些事甚至不会如此困难。请看：死亡被践踏，魔鬼堕落了，罪恶的法律被熄灭，圣神的恩宠被赐予，生命被浓缩为一个狭小的空间，重担被减轻了。\par

为了藉实际的结果使你们相信这一点，请看有多少人已超越了基督的诫命，而你却害怕那仅仅是他们分寸的事？那么，当别人跃过界限，而你自己对所规定的事却过于懒惰时，你还有什么借口？\par

例如，我们劝你们从你们所有的中施舍，但另一个人却已舍弃了全部财产。我们要求你们与妻子贞洁地生活，但另一个人却根本没有进入婚姻。我们恳求你们不要嫉妒，但另一个人甚至为了爱德舍弃了自己的生命。我们再次恳求你们在判断时要宽厚，不要苛刻对待犯罪的人，但另一个人甚至挨打时也转过另一边脸颊。\par

那么，请问我们该说什么？我们有什么借口，连这些事都不做，而别人却远超过我们？他们不可能超过我们，除非这事非常容易。因为嫉妒他人福乐的人与与他们同乐同喜的人，谁更憔悴？怀着猜疑和不断战栗看待一切的人，是贞洁的人还是奸夫？被善的希望所鼓舞的人，是强暴抢夺的人，还是显示仁慈、将自己的所有分施给穷人的？\par

因此，让我们牢记这些事，不要在德行的赛程中麻木迟钝；而要带着充分的准备投身于这些光荣的角力，稍事劳苦，以便赢得永久的、不朽的冠冕。愿我们所有人都能藉着我们的主耶稣基督的恩宠和慈爱达到这一点，愿光荣和权能归于祂，直到永远。阿们。\par

\SourceRecord{Translated by George Prevost and revised by M.B. Riddle.；Nicene and Post-Nicene Fathers, First Series；Vol. 10.；Edited by Philip Schaff.；Buffalo, NY: Christian Literature Publishing Co.,；1888.；<http://www.newadvent.org/fathers/200139.htm>.}

% structure: p0001 source=h1 target=section reason=page-title

\section{玛窦福音讲道集 40}

\BibleRef{玛 12:9-10}。\par

\BibleQuote{耶稣离开那地方，进了他们的会堂；看，那里有一个人，他的一只手枯萎了。}\par

\newline{}祂又在安息日治病，为祂门徒所行的辩护。其他福音记载，祂将那人\Quote{立在中间}，问他们：\Quote{在安息日行善是否合法。}\par

请看主的柔情。\Quote{祂将他立在中间}，为藉着这一景象制服他们；使他们因目睹而抛弃恶念，并由于对那人的某种羞愧而停止残忍的行径。但他们既不温柔也不人道，宁愿损害基督的名声，也不愿看到这人痊愈；在两方面都暴露了他们的邪恶：既与基督为敌，又如此好斗，甚至藐视祂对他人的仁慈。\par

其他福音记载是祂提问，而这位则说问题是向祂提出的。\BibleQuote{他们就问祂说：安息日治病可以不可以？为的是要控告祂。}\BibleRef{玛 12:10}很可能两者都发生了。因为他们是不圣洁的恶徒，而且确信祂必定会继续治病，就急于用问题来抢先阻止祂。所以他们问：\Quote{安息日治病可以不可以？}不是为了求知，而是为了\Quote{要控告祂}。然而，如果他们真要控告祂，这工作本身就足够了；但他们还想在祂的话中也找到把柄，为自己预先准备大量的论据。\par

但祂由于对人的爱，也做了这事：祂回答他们，教导自己的温良，并将一切反击到他们身上，指出他们的不人道。祂将那人\Quote{立在中间}，并非害怕他们，而是力图使他们得益，并激发怜悯。\par

但即使如此，祂仍未能说服他们，于是祂悲愤，因他们心硬而发怒，就说：\par

\BibleQuote{你们中间谁有一只羊，如果在安息日掉进坑里，而不把它抓住拉上来呢？人比羊贵重得多了！所以，在安息日行善是可以的。}\BibleRef{玛 12:3}\par

这样，为避免他们执拗，并再次控告祂违犯律法，祂就用这个例子来驳斥他们。请你注意，祂如何针对每种情况，以多样而恰当的方式提出祂为守安息日辩护的理由。首先，在瞎子的事上（\BibleRef{若 9:6}），当祂和泥时，祂甚至没有为自己辩护；尽管那时他们也责难祂，但创造的方式足以表明律法的主宰和所有者。其次，在瘫痪者的事上，当他拿着床走而他们挑剔时（\BibleRef{若 5:9}），祂为自己辩护，时而如天主，时而如人；如人时祂说：\BibleQuote{若人在安息日受割礼，免得律法被违背}（祂没有说\Quote{使人得益}）；\BibleQuote{你们就对我发怒，因为我在安息日使一个人完全康复吗？}\BibleRef{若 7:23}又如天主时祂说：\BibleQuote{我父到现在一直工作，我也工作。}\BibleRef{若 5:17}\par

但为门徒受责难时，祂说：\BibleQuote{你们没有读到达味做了什么，当他饥饿时，他和随从的人进了天主的殿，吃了供饼？}\BibleRef{玛 12:3}祂也提到了司祭。\par

而在此处：\BibleQuote{在安息日行善或作恶，哪样可以？你们谁有一只羊？}因为祂知道他们爱财，他们全被财富占据，而非热爱人类。实际上，另一位福音记载说，祂在问这些问题时还环视他们，想用目光赢得他们；但即使如此，他们也没有变得更好。\par

然而此处祂只说话；而在别处许多情况下，祂也藉按手来医治。但这一切都不能使他们温良；反而，那人被治愈了，而他们因他的健康变得更坏。\par

祂的愿望确实是在医治那人之前先医治他们，祂尝试了无数种医治方法，无论是藉着在他们面前做的事，还是藉着说的话；但既然他们的痼疾终究无法治愈，祂就着手工作。\BibleQuote{于是祂对那人说：伸出手来。他一伸，手就复原了，像另一只手一样。}\BibleRef{玛 12:13}\par

2. 他们做了什么？他们出去，商议要杀死祂。因为法利塞人，经上说：\BibleQuote{出去商议怎样除掉祂。}他们没有受到任何伤害，却图谋杀害祂。嫉妒是多么大的恶行！因为它不仅与外人争战，也与自己人争战。马尔谷说，他们与黑落德党人一同商议。\BibleRef{谷 3:6}\par

那么，温柔和良善的那一位做了什么？祂知道了这事就退避。\BibleQuote{耶稣知道了他们的计谋，就离开那里。}\BibleRef{玛 12:15}如今那些说\Quote{应该行奇迹}的人在哪里？看，这些事表明，不诚实的心灵甚至不被此说服；祂也清楚表明，祂的门徒也被他们无故责难。但我们应该注意，他们特别在看到邻人得益时变得凶猛；当他们看见有人从疾病或邪恶中被解救时，正是他们挑剔、变成野兽的时候。他们曾如此毁谤祂，当祂要拯救娼妓时，当祂与税吏一起吃饭时，现在又当他们看见那只手复原时。\par

但请你观察，我恳求你，祂如何既不停止对软弱者的温柔关怀，又平息了他们的嫉妒。\Quote{有许多群众跟着祂，祂治好了他们所有的人，并且吩咐那些被治好的人，不可把祂显扬出去。}因为虽然群众到处钦佩并跟随祂，但他们并没有停止他们的邪恶。\par

然后，免得你对所发生的事和他们的奇怪狂乱感到困惑，祂也引用了先知，预言这一切。因为先知的精确度如此之大，他们甚至没有忽略这些事，而是预言了祂的行程、地点的变动以及祂在其中行动的目的；这样你就可以学习，他们是如何完全藉着圣神说话的。因为人的秘密无法通过任何技巧被知晓，那么要了解基督的目的就更加不可能，除非圣神启示了它。\BibleRef{格前 2:11}\par

那么先知说了什么呢？接着就记录下来：\Quote{这是为应验那藉依撒意亚先知所说的话：\InnerQuote{看，我的仆人，我所拣选的，我所爱的，我心所喜悦的。我要将我的圣神赐给祂，祂要向外邦人宣讲正义。祂不争辩，也不呼喊，也没有人在街市上听见祂的声音。压伤的芦苇，祂不折断；将熄的灯心，祂不吹灭，直到祂使正义获胜。外邦人将因祂的名而信赖。}}\par

先知赞美祂的温和与不可言喻的能力，并为外邦人打开了\Quote{又大又有效的门}；他也预言了将要临到犹太人身上的灾难，并表明祂与圣父的同心合意。因为祂说：\Quote{看，我的仆人，我所拣选的，我所爱的，我心所喜悦的。}如果祂拣选了祂，那么基督并非作为敌对者而废除法律，也不是作为立法者的敌人，而是作为与祂有相同心意和相同目标者。\par

然后，宣告祂的温和，他说：\Quote{祂不争辩，也不呼喊。}因为祂的愿望确实是在他们面前治好他们；但由于他们推开了祂，祂甚至没有为此而争辩。\par

并暗示祂的力量和他们的软弱，他说：\Quote{压伤的芦苇，祂不折断。}因为确实很容易像芦苇一样把他们全都粉碎，而且不仅仅是芦苇，而是已经压伤的。\par

\Quote{将熄的灯心，祂不吹灭。}在这里他既指出了他们被点燃的愤怒，又指出了祂能够压制他们的愤怒、并轻易将其扑灭的力量；由此表明祂极大的温和。\par

那么怎样呢？这些事会一直持续吗？祂会永远容忍他们，让他们如此疯狂地密谋反对祂吗？远非如此；但当祂完成了祂的部分时，那么祂也将执行其他目的。因为祂通过说\Quote{直到祂使正义获胜：外邦人将因祂的名而信赖。}正如保禄同样所说：\Quote{我已经预备好了，等你们完全服从的时候，才责罚一切的不服从。}\par

但什么是\Quote{当祂使正义获胜时}？当祂完成了祂自己的一切部分时，我们就被告知，祂也会将祂的惩罚降在他们身上，而且是一个完全的惩罚。那时他们将遭受祂的恐怖，当祂的奖杯荣耀地竖立起来，从祂所发出的命令已经得胜，并且祂没有给他们留下任何无耻的辩驳借口时。因为祂习惯于称正义为\Quote{审判}。\par

但祂的救恩计划不仅仅限于惩罚不信者，而且祂也将赢得整个世界归向祂自己。因此祂补充说：\Quote{外邦人将因祂的名而信赖。}\par

然后，为了让你知道这也是按照圣父的旨意，先知从一开始就向我们保证了这一点，连同之前所说的；他说：\Quote{我的爱子，我心所喜悦的。}因为对于这爱子，很明显祂也按照那被爱者的心意做了这些事。\par

3.\Quote{那时，有人给祂带来一个附魔的、又瞎又哑的人，祂治好了他，以致那哑巴既会说话，瞎子也会看见。}\par

恶灵何等邪恶！它堵住了那人的两个通道，即视觉和听觉，藉此他本应相信；然而基督却把两者都打开了。\par

\Quote{所有的人都惊奇，说：\InnerQuote{这难道就是达味之子吗？}但法利塞人却说：\InnerQuote{这人驱魔，无非是靠着魔王贝耳则步。}}\par

然而，那到底说了什么大事呢？尽管如此，他们还是不能忍受；正如我已经指出的那样，他们总是被对邻人所做的好事刺痛，没有什么比人的得救更使他们痛苦的了。而实际上祂已经退避，给了时间让他们的激情平息；但因为又施行了一次恩惠，邪恶再次被点燃；恶灵不像他们那样愤怒。因为恶灵确实离开了身体，让位并逃跑了，没有发出任何声音；但这些人在试图先杀害祂，后毁谤祂。也就是说，他们的第一个目标没有成功，他们就想伤害祂的名誉。\par

嫉妒就是这样一种东西，没有比它更坏的邪恶了。因为奸淫者确实享受某种（尽管是那样的）快乐，并在短时间内完成他的特定罪；但嫉妒者惩罚自己，比对他所嫉妒的人更报复自己，并且从不停止他的罪，而是持续地犯这罪。因为正如猪喜欢在污泥中，恶灵喜欢伤害我们，同样他也以邻人的灾难为乐；如果有什么痛苦的事发生，他就得到安慰，松一口气；把别人的不幸当作自己的喜乐，把别人的祝福当作自己的灾祸；他不考虑自己可能得到什么快乐，而考虑邻居会有什么痛苦。那么，这些人难道不应该像疯狗、像毁灭性的魔鬼、像复仇女神一样被用石头打死、用棍棒打死吗？\par

因为正如甲虫以粪为食，这些人也以他人的灾祸为食，他们简直是我们本性的公敌和仇人。其余的人甚至怜悯被杀死的牲畜，而你看见一个人受益，却变得像野兽一样，战栗、脸色苍白？还有什么比这疯狂更糟糕的呢？所以，你看，淫乱者和税吏能进入天国，而嫉妒的人虽在里面却出去了：因为经上说：\Quote{\BibleQuote{天国的子民要被扔在外边。}}\BibleRef{玛 8:12}。前者一旦摆脱了当前的邪恶，就得到了他们从未指望的东西；而后者却连他们已经拥有的好处也失去了——这是非常合理的。因为这把人变成魔鬼，使其成为凶恶的邪魔。第一桩谋杀就是由此而来；本性就是这样被遗忘的；大地就是这样被玷污的；之后大地张开口，活生生地吞没了达堂、科辣黑、阿彼兰和那全体民众，彻底毁灭了他们。\par

4. 但有人会说，斥责嫉妒是容易的；我们还应当考虑如何使人摆脱这种病症。那么，我们怎样才能摆脱这种邪恶呢？如果我们谨记：正如犯奸淫的人不能合法地进入教会，嫉妒的人也不能——而且后者比前者更不能。因为就现状而言，嫉妒甚至被视为无所谓的事；因此也很少被重视；但若它的真正恶劣之处显明出来，我们就会轻易地避免它。\par

那么，哭泣、叹息吧！哀恸、恳求天主吧！要学会为此感到痛悔，如同犯了大罪一样。你若存此心，就会很快摆脱这病症。\par

有人会说：谁不知道嫉妒是坏事呢？确实没有人不知道；但人对这种情欲的看法与对奸淫和淫乱的看法并不相同。至少，有谁曾因嫉妒而严厉地谴责自己？有谁曾为这祸患恳求天主，求主怜悯他？从来没有人。但如果他禁食，并给穷人一点钱，即使他嫉妒到极点，他也认为自己没有做什么可怕的事，因为他被最可咒诅的情欲所束缚。例如，加音为何变成那样？厄撒乌为何那样？拉班的儿女呢？雅各伯的儿子们呢？科辣黑、达堂、阿彼兰及其同伙呢？米黎盎呢？亚郎呢？魔鬼自己呢？\par

与此同时，也要考虑这一点：你伤害的不是你嫉妒的人，而是把剑刺向你自己。因为加音伤害了亚伯尔什么呢？他不是反而违背自己的意愿，使亚伯尔更快地进入天国了吗？但加音自己却被无数灾祸刺透。厄撒乌怎样伤害了雅各伯？雅各伯不是发了财，享受了无数的祝福吗？而厄撒乌自己却成了父亲家里的流亡者，在诡计之后漂泊异乡。雅各伯的儿子们又怎样损害了若瑟呢？即使他们甚至动用了血。他们不是遭受了饥荒，面临极端的危险吗？而若瑟却成了全埃及的王。因为你越嫉妒，就越为你所嫉妒的对象谋取更大的祝福。因为有一位天主鉴察这些事；当他看见那受伤害的人并未伤害别人时，他就更加高举他，使他荣耀，却惩罚你。\par

因为如果他们因敌人跌倒而幸灾乐祸，天主不容他们不受惩罚（经上说：\Quote{\BibleQuote{你的仇敌跌倒，你不要高兴，免得主看见而不喜悦}}\BibleRef{箴 24:17-18}）；更何况那些嫉妒没有做错事的人呢。\par

那么，让我们铲除这多头野兽吧。因为嫉妒的种类确实很多。如果爱那爱自己的人只不过与税吏相等，那么恨那没有伤害自己的人又该站在哪里呢？他变得比外教人更坏，又如何逃脱地狱呢？因此，我极其忧伤，因为我们奉命效法天使，甚至效法天使的主，却效法了魔鬼。因为在教会里，甚至在领导者中间，嫉妒也确实很多。因此，我们甚至必须对自己讲道。\par

5. 那么请告诉我，你为什么嫉妒你的邻人？因为你看见他获得尊荣和好评？难道你不记得尊荣给不谨慎的人带来多少恶事吗？使他们陷入骄傲、虚荣、傲慢、轻蔑，并使他们更加粗心大意？除这些恶事之外，它们还轻易地使人凋零。最严重的是，由此产生的恶事永存不朽，而当时的快乐却一闪即逝。就为了这些事，你嫉妒吗？请告诉我。\par

\Quote{但他对统治者有很大影响力，随意引导和驱使一切，能伤害得罪他的人，善待奉承他的人，且有很大权力。}这些是世俗之人的话，是那些被钉在地上的人的话。因为属神的人，没有什么能伤害他。\par

因为别人能对他造成什么严重的伤害呢？罢免他的职务？那又怎样？如果做得公正，他甚至受益；因为没有什么比不配地担任司铎职务更惹天主发怒。但如果不公正，责任又落在别人身上，而非他身上；因为那无端受苦并高尚忍受的人，反而因此获得向着天主的更大信心。\par

我們不要追求如何處於有權勢、有榮譽、有權威的地位，而要追求如何在德行和克己中度生。因為權威的地位往往說服人行許多天主不喜歡的事，而需要極大的靈魂力量才能善用權威。正如被剝奪權威的人，無論情願與否，都得操練自制，同樣，享有權威的人就像與一位優雅美麗的少女同住，卻被規定不得以不潔的眼光看她。權威就是這樣的事。因此，它甚至使許多人不情願地變得傲慢；它喚醒憤怒，解除舌頭的韁繩，撕裂嘴唇的門戶，像風一樣煽動靈魂，使小船沉入罪惡的最深淵。那麼，處在如此大危險中的人，你竟羨慕他，說他值得嫉妒嗎？這豈不是極大的瘋狂！至少，除了我們所提及的，請想想這個人有多少敵人、控告者、諂媚者圍攻他。請問，這些是稱人為有福的理由嗎？誰能這樣說呢？\par

\Quote{但是人民，}你说，\Quote{很看重他。}这算什么呢？人民毕竟不是天主，他要向天主交账；所以你说人民，无非是在说其他破坏者、岩石、浅滩和暗礁。因为获得人民的喜爱，越是使人显赫，它所带来的危险、忧虑和沮丧就越大。因为这样的人民根本没有任何喘息的机会，有一个如此严厉的主人。我何必说\Quote{喘息和止步}呢？这样的人即使有再多的善工，也几乎不能进入天国。因为没有什么比来自群众的荣誉更容易使人败坏，使人胆怯、卑贱、谄媚、虚伪。\par

例如，法利塞人為什麼說基督被附魔？難道不是因為他們貪圖群眾的榮譽嗎？\par

群眾為什麼能對祂作出正確的判斷？難道不是因為他們沒有被這種病症所困擾嗎？因為沒有什麼比貪圖群眾的榮譽更使人無法無天、愚昧無知；也沒有什麼比輕視它更使人光榮和堅定不移。\par

因此，一個人若要抵抗如此強烈的衝擊和如此猛烈的風暴，就需要極大的靈魂力量。因為當事情順利時，他把自己置於一切之上；當遭遇相反時，他又恨不得活埋自己。這對他來說既是地獄，也是天國，當他被這種激情壓倒時。\par

請問，這一切是值得嫉妒的事，而不是哀傷和眼淚的事嗎？每個人都能看到。但是你嫉妒那種名望中的人，就像一個人看見另一個人被無數野獸捆綁、鞭打、撕裂，卻嫉妒他的傷口和鞭痕。事實上，群眾有多少人，他就有多少鎖鏈、多少暴君；更糟的是，每個人都有不同的想法；他們對奴役他們的人隨意下判斷，不經審查；而是根據這個人或那個人的決定來確認。\par

有什麼樣的波濤、什麼樣的暴風雨比這更痛苦呢？是的，這樣的人一會兒被快樂吹脹，一會兒又輕易沉沒，永遠在波動中，從不安寧。因此，在集會和演講競賽之前，他充滿焦慮和恐懼；集會之後，他要麼因沮喪而死，要麼相反地毫無節制地歡喜，這比憂傷更糟。因為那種快樂不比憂傷更小，從其對靈魂的影響就可以看出：它使靈魂輕浮、不穩定、飄忽不定。\par

這一點甚至可以從古代的例子看出。例如，達味何時值得欽佩：當他歡喜時，還是當他痛苦時？猶太人民呢：是當他們呻吟呼求天主時，還是在曠野中歡躍並敬拜牛犢時？因此，撒羅滿比任何人都更了解什麼是快樂，他說：\BibleQuote{往哀悼之家去，比往欢笑笑之家更好。}\BibleRef{训 7:2}因此，基督也祝福前者，說：\BibleQuote{哀慟的人是有福的，}\BibleRef{玛 5:4}卻為後者哀哭，說：\BibleQuote{你們現在歡笑的人是有禍的，因為你們將要哀哭。}\BibleRef{路 6:25}這很恰當。因為在快樂中，靈魂更加鬆弛和軟弱；而在哀傷中，靈魂變得堅強、清醒、脫離一切情慾，變得更高尚、更堅強。\par

既然知道這一切，讓我們逃避來自群眾的榮耀和由此而來的快樂，好使我們贏得真實而永恆的榮耀；願我們眾人賴我們的主耶穌基督的恩寵和慈愛，達到這一榮耀；願光榮和權能歸於祂，直到永遠。阿們。\par

\SourceRecord{Translated by George Prevost and revised by M.B. Riddle.；Nicene and Post-Nicene Fathers, First Series；Vol. 10.；Edited by Philip Schaff.；Buffalo, NY: Christian Literature Publishing Co.,；1888.；<http://www.newadvent.org/fathers/200140.htm>.}

% structure: p0001 source=h1 target=section reason=page-title

\section{玛窦福音讲道 第四十一篇}

\Emph{\BibleRef{玛 12:25-26}}.\par

\BibleQuote{凡一国自相纷争，必成废墟；凡一城或一家自相纷争，必不得存立。如果撒殚驱逐撒殚，是自相纷争，那么他的国如何能存立呢？}\par

甚至在此以前，他们已经用这事控告过祂，说：\Quote{祂是仗赖贝耳则步驱魔。}但是，那时祂并没有责备他们，而是任由他们藉着更多的奇迹认识祂的能力，并通过教导学习祂的威严；如今，既然他们继续说着同样的话，祂也开始责备他们，首先以此显示祂的天主性：祂公开了他们的秘密；其次，藉着轻易驱魔的行为本身。\par

事实上，这控告也是极其无耻的。因为，如同我所说的，嫉妒并不寻求该说什么，而只是寻求能说些什么。尽管如此，基督并没有因此轻视他们，而是以祂特有的忍耐为自己辩护，教导我们对仇敌要温良；即使他们说这样的话，而我们既不自知，也毫无可能性，也不应惊慌或困扰，而是要以极大的宽容向他们解释。这正是祂在那场合最特别做的，提供了最有力的证据，证明他们所说的都是假的。因为一个附魔者不会表现出如此的温良；一个附魔者不会知道人的秘密。\par

事实上，由于这种猜疑过于厚颜无耻，又因为害怕群众，他们不敢公开提出这些指控，而是在心中盘算。但是祂，为了表明祂也知道这一切，并没有陈述指控，也没有揭露他们的恶意；而是加上反驳，让那些说了这话的良心去定他们的罪。因为祂只专注于一件事：对犯罪者行善，而非揭露他们。\par

然而，如果祂有意延长祂的言论，使他们显得可笑，并同时向他们索取最严厉的惩罚，没有任何东西能阻止祂。然而祂将这些都放在一边，只着眼于一个目标：不使他们更加争辩，而是更坦诚，从而更好地使他们趋向改正。\par

那么祂如何与他们辩论呢？不是藉着引用圣经（因为他们不会听从，反而肯定会扭曲其意义），而是藉着日常生活中的事件。因为祂说：\BibleQuote{凡一国自相纷争，必不得存立；一城一家，若自相纷争，也必迅速败落。}\par

因为外来的战争并不像内战那样具有破坏性。是的，身体也是如此；万事万物也是如此；但这一次祂从那些更为众所周知的事物中取譬。\par

然而，世上有什么比一个王国更有力量呢？没有什么，但若内部分裂，它也会灭亡。如果在这情况下，有人归咎于事务的沉重负担，认为它因自身重量而垮塌；那么对一座城市，你能说什么呢？对一家呢？因此，无论是一件小事，还是一件大事，若与自己分裂，它就会灭亡。那么，如果我有一个魔鬼，并藉着他驱魔，魔鬼之间就有分裂和争斗，他们彼此敌对。但如果他们彼此敌对，他们的力量就会削弱和消灭。\Quote{因为如果撒殚驱逐撒殚}（祂没有说\Quote{魔鬼}，暗示魔鬼彼此之间极大的合一），\Quote{他便是自相纷争；}祂这样说。但若他分裂了，他就变弱了，并且败亡；若他败亡了，他又如何能驱逐另一个呢？\par

你看到这控告是多么荒谬，多么愚蠢，多么矛盾吗？因为不可能同样的人先说祂站立着并驱魔，然后又说祂藉着那很可能导致祂覆灭的东西站立。\par

2. 这是第一个反驳，接下来的第二个是关于门徒的。因为祂不仅用一种方式，而且用第二、第三种方式来解决他们的异议，旨在最充分地堵住他们的无耻。祂对安息日也做了类似的事，提出了达味、司祭以及引用了经书的话：\Quote{我要的是怜悯，而不是祭献，}以及安息日设立的原因——\Quote{安息日是为人设立的。}所以祂在目前的情况下也这样做：在第一个反驳之后，祂继续进行第二个更清楚的反驳。\par

\Quote{因为如果我，}祂说，\Quote{是仗赖贝耳则步驱魔，那么你们的子侄是靠谁驱魔呢？}\BibleRef{玛 12:27}\par

请看这里祂的温和。因为祂没有说\Quote{我的门徒}，也没有说\Quote{宗徒}，而是说\Quote{你们的子侄}；这样，如果他们愿意回归与他们相似的尊贵，便可由此获得强有力的动力；但如果他们不坦率，继续走同样的路，那么他们此后就无法提出任何辩解，即使再无耻也不能。\par

祂所说的话是这样的：\Quote{宗徒们是靠谁驱魔的呢？}事实上，他们已经在这样做，因为他们从祂那里得到了权柄，而这些人并没有控告他们；他们的争论不是针对行为，而只是针对人。既然祂想表明他们的言论只是出于对祂的嫉妒，祂便提出了宗徒们：如果我是这样驱魔，那么那些从我这里得到权柄的人更是如此。然而，你们并没有对他们说过这样的话。那么，你们为何控告我做这些事的创始者，却对那些人免于指控呢？但这并不能免去你们的惩罚，反而会更定你们的罪。因此祂又加上：\Quote{他们将要作你们的审判官。}因为当你们中间的人，又受过这些事训练的人，既相信我，又服从我，那么很明显，他们也将定那些在言行上反对我的人的罪。\par

\BibleQuote{但若我赖圣神驱逐魔鬼，那么天主的国已临到你们中间了。}\par

所谓\Quote{天主的国}是什么意思？是\Quote{我的来临}。请看祂又如何安抚他们，平息他们，引导他们认识祂自己，并表明他们是在攻击自身的利益，反对自己的救恩。祂说：\Quote{你们本该欢喜踊跃，因为那一位已来临，赐予那些伟大而不可言喻的祝福，这些祝福是古时先知们所歌颂的，并且你们福乐的时刻已经临近；你们却反其道而行之，不但不接受祝福，反而恶言毁谤，捏造不存在的指控。}\par

玛窦固然说：\Quote{若我赖天主的神驱魔}，但路加却说：\Quote{若我赖天主的手指驱魔}\BibleRef{路 11:20}，意思是驱魔是最伟大的能力之工，而非寻常的恩宠。祂的用意是，他们应该由这些事推断并说：若是如此，那么天主子已经来临。然而祂没有这样说，而是以含蓄的方式，以免刺激他们，祂隐晦地暗示说：\Quote{那么天主的国已临到你们中间了。}\par

你看到何等的智慧？借着他们所指责的同一件事，祂显示了自己的临在照耀着。\par

然后，为安抚他们，祂没有简单地说\Quote{天主的国已来临}，而是说\Quote{临到你们}\newline{}好像祂在说：福份已临到你们身上，那么你们为何对自己的祝福感到不悦？为何要反对自己的救恩？这就是先知们许久以前所预言的时期；这就是他们歌颂的那个来临的记号，即这些事借着天主的能力运行。因为这些事确实运作，你们知道；但这些事是借着天主的能力运作，是事实本身所呼号的。是的，撒殚现在不可能更强；相反，他必定是软弱的。但软弱者不可能像强者那样赶出强壮的魔鬼。\par

祂这样说，表明了爱德的力量，以及分裂与争执的软弱。因此，祂自己也时常在所有场合嘱咐门徒们关于爱德，并教导他们，魔鬼为了颠覆爱德，不遗余力。\par

3. 在说完第二次反驳之后，祂又加了第三次，这样说：\par

\BibleQuote{怎能进入壮士的家，夺取他的家产，除非先捆绑那壮士，然后才夺取他的家产？}\par

因为根据以上所述，撒殚不可能驱逐撒殚，这是显而易见的；但除非先胜过他，否则也不可能以任何其他方式驱逐他，这一点也是人人承认的。\par

那么由此确立了什么呢？先前的论点，有了更充分的证据。祂说：\Quote{我远远不像你们所说的那样利用魔鬼为盟友，我反而与他作战，捆绑他；而掠夺他的家产就是这事的明证。}请看，他们试图确立的论点，却被证明了相反。因为他们想表明，祂不是凭自己的能力赶鬼，祂却表明，不但是魔鬼，甚至他们的魁首也被祂以全部权威捆绑；并且在驱逐魔鬼之前，祂已凭自己的能力胜过了他。这由所成就的事显然可见。因为如果他是首领，他们是属下，那么除非他被击败并屈服，他们怎能被掠夺呢？\par

这里祂的话在我看来似乎也是一个预言。因为我不认为只有邪灵是魔鬼的家产，那些行他恶事的人也是。因此，为宣告祂不仅驱逐魔鬼，还要将世间的一切谬误驱除，压下他的巫术，使他的所有伎俩化为无用，祂说了这些话。\par

祂不是说\Quote{将取走}，而是说\Quote{将掠夺}，以表达借着权柄所为。但祂称他为\Quote{壮士}，不是因为他是天性如此（绝对不可这样说），而是表明他因我们的懈怠而造成的昔日暴政。\par

\BibleQuote{不与我相合的，就是反对我的；不与我一齐收集的，就是分散的。} \BibleRef{玛 12:30}\par

请看第四个反驳。祂说：我的愿望是什么？是引人归向天主，教导德行，宣讲天主的国。魔鬼与邪灵的愿望又是什么？与此相反。那么，那个不与我一齐收集、甚至根本不与我一齐的，怎能与我合作呢？我说合作做什么？不，恰恰相反，他的愿望是分散我的成果。那么，这个不仅不合作反而分散的人，怎能表现出与我如此一致，以至于与我一同驱逐魔鬼呢？\par

我们可以自然而然地推测，祂说这话不仅指魔鬼，也指祂自己，因为祂自己一方面反对魔鬼，分散他的家产。那么，也许有人会说：不与我相合的，怎么就是反对我呢？正因为他没有收集。但如果这是真的，那么反对他的就更加是了。因为不合作的人尚且是敌人，那么作战的人就更不用说了。\par

但祂说这一切，是为了指出祂对魔鬼的敌意是多么伟大而难以言喻。因为请告诉我，如果你必须与某人交战，不愿站在你那边作战的人，岂不是因此反对你吗？若在别处祂说：\Quote{不反对你们的，就是帮助你们的}，这与这里并不矛盾。因为这里祂指的就是那些实际反对他们的人，但那里祂指的是那些部分站在他们一边的人：因为据说\Quote{他们因祢的名字驱逐魔鬼}。\par

但我以为祂在这里还暗示了犹太人，将他们置于魔鬼一边。因为他们也反对祂，分散祂收集的。至于祂暗示了他们，祂这样说便表明了：\par

\Quote{因此，我告诉你们：一切罪过和亵渎，都可得赦免。}\par

5. 主为自己辩护之后，驳斥了他们的异议，证明了他们无耻行为之虚妄，遂转而恐吓他们。因为劝诫与纠正的一个重要部分，不仅是恳求和说服，也包括威胁；祂在立法和施教的多处经文中，正是这样做的。\par

虽然这话似乎相当隐晦，但若我们留心，其解答却很简易。\par

首先，最好来听这些话本身：\BibleQuote{一切罪过和亵渎，人都可得赦免；但是亵渎圣神的罪，必不得赦免。凡说话干犯人子的，可得赦免；但谁若说话干犯圣神，必不得赦免，无论是在今世，或是在来世。}\par

祂现在所肯定的是什么呢？你们说了许多反对我的话，说我是骗子，是天主的仇敌。这些事，只要你们悔改，我就宽恕你们，不再追究处罚。但亵渎圣神的罪，即令悔改，也不得赦免。这怎么合理呢？因为连这罪，也曾因悔改而得赦免。至少那些说过这话的人，后来许多都信了，一切都被赦免了。那么祂是什么意思呢？是说这罪是万不可赦的。为什么？因为虽然他们不知道我是谁，但对圣神却已有充分的认识。因为先知们所说的一切，都是借着圣神说的；事实上，旧约全体都对圣神有极高的认识。\par

祂说的就是这样：就算你们因我所披戴的肉体而对我反感；难道你们对圣神也能说\Quote{我们不认识祂}吗？因此你们的亵渎不可赦免，你们要在此世和来世都受惩罚。因为许多人只在此世受罚（例如犯了奸淫的人，以及在格林多人中不配领受奥秘的人）；但你们却是今世和来世都受罚。\par

现在，论到你们在十字架前对我的亵渎，我宽恕它们；连十字架本身的胆大恶行也一样；你们不会仅因不信而被定罪。（因为即使那些在十字架前相信的人，信心也不完全。而且祂多次嘱咐他们在受难之前，不要将祂显明给人；在十字架上祂也说这罪已被赦免。）但你们论及圣神的话，却毫无推诿的理由。因为为了证明祂说的是受难前有关祂的话，祂又加上：\BibleQuote{凡说话干犯人子的，可得赦免；但谁若说话干犯圣神，}就不再赦免了。为什么？因为你们对此是知道的，并且你们所硬心反抗的真理是昭彰的。\par

因为论到人，有些人在此世和来世都受罚，有些人只在此世受罚，有些人只在来世受罚，另有些人则此世和来世都不受罚。在此世和来世受罚的，就如这些人本身（因为他们在自己城邑被攻陷时，在此世已遭受了那些不治的灾祸，在来世还将受极重的惩罚），又如同索多玛的居民，以及许多其他人。只在来世受罚的，如同那在火焰中受苦、连一滴水也得不到的富翁。只在此世受罚的，如同格林多人中犯奸淫的那人。此世和来世都不受罚的，如同宗徒们、先知们、有福的约伯；因为他们的苦难决非惩罚，而是竞赛和搏斗。\par

因此，让我们努力成为与这些人同列的人；若不能，至少成为在此世洗清自己罪过的人。因为那另一个审判实在可怕，那报应毫不留情，那惩罚无药可救。\par

6. 但如果你们连在此世也不愿受罚，就审判自己，施行自己的惩罚。听保禄所说的：\BibleQuote{我们若自己省察，就不致受审判。}如果你们这样做，按部就班，你们甚至会得到冠冕。\par

但有人问：我们如何施行自己的惩罚呢？哀哭、痛切呻吟、自卑、克苦，仔细回想你的各项罪过。这对人的灵魂来说不是什么小的折磨。凡曾处于痛悔状态的人都知道，灵魂所受的折磨莫过于此。凡曾生活在记忆罪恶中的人，都知道由此产生的痛苦。因此，天主将成义作为这种悔改的赏报，说：\BibleQuote{你们先要陈述你们的罪过，为能成义。}因为把我们的所有罪过汇集起来，并不断反复细数它们的细节，这确实是走向改正的重要一步。因为这样做的人会心碎，以至于认为自己不配活着；而这样想的人会比蜡还柔软。因为你们不要只对我讲奸淫、通奸，以及那些显现于外、众人皆知的事；也要汇集你们暗中的计谋、诬告、毁谤、虚荣、嫉妒以及诸如此类的事。因为连这些也不会带来微小的惩罚。因为毁谤者也要堕入地狱；醉汉在天国无份；不爱近人者如此得罪天主，连自己的殉道也救不了他；忽略自己家人的人已背弃信仰；忽视穷人者被投入火中。\par

所以，不要把这些事看为细小，而要全都汇集起来，如同写在册上。因为你们若写下来，天主就涂抹它们；反之，你们若疏于书写，天主不但记录它们，还要追讨刑罚。那么，由我们写下而由上天涂抹，远胜于相反地，当我们忘记时，天主在那日将之陈列在我们眼前。\par

因此，为免如此，让我们严格计算一切，我们会发现自己对许多事负有责任。因为谁能远离贪恋？不，不要告诉我数量，因为即使在少量上，我们也要受同样的惩罚，要思考这点并悔改。谁能摆脱一切傲慢？然而这会将人投入地狱。谁没有暗中说过邻人的坏话？然而这使人丧失天国。谁没有固执己见？然而这样的人比所有人更不洁。谁没有用不贞的眼观看？然而这是完全的奸夫。谁没有\Quote{无缘无故向弟兄发怒}？这样的人是\Quote{公议会定罪的}。谁没有发过誓？然而这是出于恶者。谁没有发过假誓？然而这样的人比出于恶者更甚。谁没有事奉过\OriginalTerm{玛门}{mammon}？然而这样的人已经背离了对基督的真切事奉。\par

我还有其他比这些更大的事要提：但即使这些也足够，且能使人——若他并非石心，也非全然麻木——产生痛悔之心。因为如果每一条都足以投入地狱，当它们全部集合在一起时，还有什么不能成就呢？\par

那么人如何能得救？可能会有人问。通过运用相反的补救方法：施舍、祈祷、痛悔、悔改、谦卑、痛悔的心、轻视财物。因为天主为我们指出了无数得救的途径，只要我们愿意留心。让我们留心吧，让我们从各方面洁净我们的伤口，显示慈悲，饶恕那些得罪我们的人的愤怒，为一切事感谢天主，按我们的能力守斋，真诚祈祷，\Quote{用不义的钱财结交朋友}。如此我们就能获得罪过的赦免，赢得所应许的美善；愿我们众人因我们的主\Person{耶稣基督}的恩宠与对人的爱，而被认为配得这一切，愿荣耀与德能归于他，直到永远。阿们。\par

\SourceRecord{Translated by George Prevost and revised by M.B. Riddle.；Nicene and Post-Nicene Fathers, First Series；Vol. 10.；Edited by Philip Schaff.；Buffalo, NY: Christian Literature Publishing Co.,；1888.；<http://www.newadvent.org/fathers/200141.htm>.}

% structure: p0001 source=h1 target=section reason=page-title

\section{\WorkTitle{《玛窦福音》讲道集 42}}

\BibleRef{玛 12:33}\par

\BibleQuote{你们或认为树好，它的果子也好；或认为树坏，它的果子也坏；因为树是由果子认出来的。}\par

祂又以另一种方式羞辱他们，并不满足于先前的驳斥。祂这样做，并非为自己辩护（因为前面已经足够），而是希望改正他们。\par

祂的意思是这样的：你们中没有人指责所治愈的人并未得治愈，也没有人说驱魔是件坏事。即使他们再无耻，也不可能说出这话。\par

既然他们对祂的作为没有提出控告，却毁谤行这事的主，祂便指明，这控告既违背常理，也不合时宜。这是一种加重的无耻，不仅恶意曲解，而且编造出违背常人观念的攻击。\par

请看祂何等不争辩。祂并没有说：\Quote{要使树好，因为果子也是好的；}而是彻底堵住他们的口，显示自己的仁厚和他们的傲慢，祂说：即使你们想要指责我的作为，我毫不禁止，只要你们不提出自相矛盾的控诉。因为这样，他们必会最清楚地被揭露，与明显的事实对抗。因此，你们的恶意和自相矛盾的言论是徒劳的，因为实际上，树是由果子显出的，而不是果子由树显出的，但你们却反其道而行。即使树是果子的根源，但果子使树被人认识。合乎情理的做法是，要么指责我们的作为，要么赞美这些作为，从而免除我们这些行者的罪责。但现在你们却相反：对作为（即果子）无可指责，却对树作了相反的判断，称我为附魔者，这是极度疯狂。\par

此外，祂先前所说的（\BibleRef{玛 7:16}）如今也再次确立：好树不能结坏果子，反之亦然。因此，他们的指控全然违背逻辑和本性。\par

既然祂并非为自己辩护，而是为圣神辩护，祂便如洪流般倾泻斥责，说道：\BibleQuote{毒蛇的种类，你们既是恶的，怎能说出善来？} \BibleRef{玛 12:34}\par

这既是谴责，也是从他们自身的情况证明祂的话。看，祂说，你们既是坏树，就不能结出好果子。那么，你们这样说话就不足为奇了：因为你们既受恶祖宗的恶劣教养，又养成了邪恶的心。\par

请看祂如何小心谨慎、毫无例外地表达谴责：祂没有说\Quote{你们既是毒蛇的种类，怎能说出善来？}（因为这后者与前无关），而是说：\Quote{你们既是恶的，怎能说出善来？}\par

祂称他们为\Quote{毒蛇的种类}，因为他们以祖先为荣。为表明他们并未因此得益，祂将他们从与亚巴郎的关系中排除，并为他们指派了性情相似的祖先，剥夺了他们引以为傲的根基。\par

\BibleQuote{因为心里充满什么，口里就说什么。}这里祂再次显示自己的天主性，知道他们的隐秘；并且不仅为言语，也为了恶念，他们将受惩罚；祂作为天主知道一切。祂说，连人也能知道这些；因为这是自然的结果：当邪恶在内满溢时，话语便从嘴唇溢出。因此，当你听见人说出恶言时，不要以为他只有话语所显示的那些邪恶，而要猜测那源头更为丰富；因为外在所说的，不过是内在满溢的一部分。\par

请看祂何等严厉地斥责他们。如果他们所说的如此邪恶，且出于魔鬼的心意，那么他们话语的根源和源泉该达到何等地步。这自然如此：因为舌头常因羞耻而不立即倾尽所有邪恶，但心里没有人的见证，便无所畏惧地生出任何邪恶；对天主却不太在意。既然人的言语受到审查并摆在众人面前，而心却被隐藏，因此前者的邪恶减少，后者的却增多。但当内心邪恶倍增时，隐藏的一切便猛烈涌出。正如呕吐之人起初努力抑制涌出的体液，但若无法克制，便吐出大量污秽；同样，图谋恶事并诽谤邻人的人也是如此。\par

祂说：\BibleQuote{善人从善库中发出善来，恶人从恶库中发出恶来。}\par

绝不要以为，祂说，只在恶事上如此，在善事上也是这样：因为那里，内在的德行也多于外在的言语。借此祂表明，他们比言语所显示的更为邪恶，而祂自己则比言语所宣告的更为完美良善。祂称之为\Quote{库}，表示其丰富。\par

然后祂又以极大的恐惧震慑他们。绝不要以为，祂说，事情到此为止，即只受到群众的谴责；不，凡如此行恶的人，都将受最严厉的惩罚。祂没有说\Quote{你们}，部分是为了教训全人类，部分是为了使祂的话不那么沉重。\par

\BibleQuote{但我告诉你们：}祂说，\BibleQuote{凡人们所说的闲话，在审判之日，都要交账。} \BibleRef{玛 12:36}\par

所谓闲话，即是不符合事实、虚假、包含不公正指责的话；也有人认为，无聊的话也包括在内，例如引发不当大笑的话，或污秽、无礼、粗鄙的话。\par

\Quote{因为凭你的话，你要成义，也要凭你的话，你要被定罪。} \BibleRef{玛 12:37}\par

你看见这审判是多么不偏不倚吗？所要交的账又是多么合理？因为法官的判决，并非根据别人所说关于你的话，而是根据你自己所说的话；这乃是极其公平的事：因为说话与否，确实在于你自己。\par

因此，需要焦虑战栗的，不是那些被毁谤的人，而是毁谤者。因为前者不必为自己被别人所说的恶事答辩，而后者却要为自己所说的恶事答辩；全部的危机都悬于后者头上。所以，被指责的人应当无忧无虑，因为他们不必为别人所说的恶事交账；但指责别人的人却应当焦虑战栗，因为他们自己要为此被拖到审判台前。这实在是魔鬼的罗网，是一种毫无乐趣、只有伤害的罪。而且，这样的人是在自己的灵魂中积累邪恶的宝藏。如果一个人体内有恶毒的体液，他首先会尝到疾病的苦果，那么，一个在心中积蓄比胆汁更苦的东西（即邪恶）的人，岂不更要遭受极大的祸患，给自己招来沉重的疾病？从他口中所吐出的东西就足以证明。因为如果这些东西使别人如此痛苦，那么生出它们的灵魂所受的痛苦就更大了。\par

这样，设阴谋的人首先毁灭了自己；正如踩火的人烧伤自己，撞金刚石的人自讨苦吃，踢刺棒的人自刺流血。因为那懂得忍受冤枉、勇敢承担的人，就如同金刚石、刺棒和火；而惯于行恶的人，却比任何泥土都软弱。\par

所以，受冤枉不是恶，行冤枉以及不懂得忍受冤枉才是恶。例如，\Person{达味}受了多大的冤枉！\Person{撒乌耳}行了多大的冤枉！哪一个更强、更有福？哪一个更可怜、更悲惨？难道是行恶的那一个吗？请细想：撒乌耳曾许诺，如果达味杀死培肋舍特人，就招他为婿，并慷慨地将女儿赐给他。达味杀了培肋舍特人，撒乌耳却背约，不但不把女儿给他，反而图谋杀害他。那么，谁更光荣？一个因绝望和恶魔而窒息，另一个因胜利和忠于天主而比太阳更闪耀？再者，在妇女的歌咏队前，一个因嫉妒而窒息，另一个却默默忍受一切，赢得了众人，使他们都归向他。甚至当达味将撒乌耳交在手里却放了他时，谁又是幸福的？谁又是悲惨的？谁软弱？谁强大？难道是那个甚至没有为自己合法复仇的人吗？这很自然，因为撒乌耳有武装的士兵，而达味有正义——比万军更强大的盟友和助手。因此，虽然被人不义地谋害，他却连合法地杀死对方也不忍心。因为他从先前的事知道：不行恶而忍受恶，才能使人更为强大。身体如此，树木也如此。\par

\Person{雅各伯}怎样？他岂非被\Person{拉班}所亏待、受了他的害？那么，谁更强？是那个将对方控制在自己手中却不敢触动他，反而恐惧战栗的人（\BibleRef{创 31:29}），还是那个我们看见手无寸铁、无兵无卒，却比无数国王更令对方畏惧的人？\par

但为了让我从相反的一面再给你们一个更明显的证明，让我们再次以\Person{达味}本人为例，用相反的推理来验证。因为这位受冤枉而如此强大的人，后来因行了冤枉，反而变得软弱。当达味亏待了\Person{乌黎雅}后，他的地位改变了：软弱转到了行恶者身上，力量却归给了受害者；乌黎雅虽然死了，却使达味的家室荒废。一个是活着的君王，却无能为力；另一个不过是阵亡的士兵，却把一切有关他对手的事都搞得天翻地覆。\par

你们愿意我用另一种方式把我说的话讲得更明白吗？让我们观察那些为自己合法报复的人。至于行恶的人，他们是最卑劣的人，与自己灵魂作战，这自然是众所周知的。\par

但是，那一位合法报复自己的人，却点燃了无尽的祸害，用许多灾难和忧伤刺透了自己？这就是达味军队的队长。因为他挑起了惨烈的战争，遭受了数不清的祸患；如果他懂得自制，这些祸患一件也不会发生。\par

因此，让我们远离这罪，无论在言语上还是行为上，都不要亏负我们的邻人。因为祂没有说：如果你毁谤、召来法庭，而只是说：如果你说坏话，即使只是在心里，你也将受极大的惩罚。即使你所说的是真的，即使你是出于确凿的证据，仍将有惩罚临到你。因为天主并不按照别人所作所为，而是按照你所说的话来判决；祂说：\Quote{因为凭你的话，你要被定罪。}法利塞人岂不也说了真话，肯定了人所共知的事，并未揭露隐藏的事么？然而他受了极重的惩罚。\par

如果我们连对人所共知的事也不该指责，那么对尚有争论的事更不该指责；因为犯罪者自有一位审判者。我劝你们，不要僭夺独生子的特权，因为审判的宝座是为祂保留的。\par

3. 然而，你愿作审判官吗？你有一个大有裨益、毫无责难的审判法庭。请作为审判官，在你的良心上坐下来，把我所犯的一切过犯带到它面前，搜查你灵魂的罪愆，严格追究其账目，并说：\Quote{你为何胆敢做这事或那事？}如果它回避这些事，反而查问别人的事，你就对它说：\Quote{我不是就这些事审判你，你也不是为这些事来这里申辩的。就算某人是恶人，与你何干？你为何犯了这件或那件罪？为你自己答辩，而非控告；管你自己的事，不要管别人的事。}你要不断敦促它进行这种焦虑的审判。然后，如果它无话可说，反而退缩，你就用鞭子折磨它，如同对待一个不安分、不贞的使女。要使这个法庭每天开庭，并描绘火河、毒虫和其余刑罚的景象。\par

不容它再与魔鬼为伍，也不容忍它的无耻之言：\Quote{他来到我这里，他陷害我，他诱惑我。}反倒要告诉它：\Quote{如果你不愿意，这一切都毫无用处。}如果它又说：\Quote{我被肉体纠缠，我披戴着肉身，我住在世上，我居于地上。}你就告诉它：\Quote{这些都是借口和托词。因为某人也是披戴着肉身，另一个住在世上、居于地上的人却蒙嘉许；连你自己，在做得好时，也是披戴着肉身做的。}如果它因听见这话而伤痛，不要松开你的手；因为你打击它，它不会死，你反而要救它脱离死亡。如果它又说：\Quote{某人激怒了我。}你就告诉它：\Quote{但你可以选择不被激怒；你至少常常克制了怒气。}如果它说：\Quote{那女人的美貌打动了我。}你就告诉它：\Quote{你本可以克制自己。}提起那些得胜的人，提起第一个女人，她说：\BibleQuote{蛇欺骗了我。}\BibleRef{创 3:13}，却仍未逃脱罪责。\par

当你查究这些事时，不要让任何人介入，不要有人打扰你；正如审判官坐在帷幕下审判，你也当同样寻求安静的时间和地点，代替帷幕。当你晚餐后起身，准备躺下时，就举行这个审判；这是适合你的时间，地点是你的床和你的内室。先知也同样命令说：\BibleQuote{你们在床上要心中思想，并要安静。}对于轻微的过犯，要求重大的补偿，这样你或许永远不敢接近大罪。如果你每天这样做，你必能坦然无惧地站在那可怕的审判台前。\par

保禄就是这样成为洁净的；因此他也说：\Quote{因为我们若先审断自己，就不致受审判。}\BibleRef{格前 11:31} 约伯也是如此洁净他的儿子们。\BibleRef{约 1:5} 因为那为隐密的罪献祭的人，岂不更要追究明显的过犯吗？\par

4. 但我们不这样做，反而全然相反。因为一躺下休息，我们反而思想一切世俗的事；有人引入不洁的思念，有人引入高利贷、契约和暂时的挂虑。\par

如果我们有一个女儿，是童贞女，我们严格看守她；但那比女儿更宝贵的我们的灵魂，我们却任凭她行淫、玷污自己，将无数的邪念引到她内。无论是贪财之爱、奢侈之爱、美色之爱、愤怒之爱，还是别的什么想要进来的，我们都敞开大门，吸引并邀请它，帮助它从容地玷污我们的灵魂。还有什么比这更野蛮呢？忽视比一切都宝贵的灵魂，被如此多的奸夫凌辱，长期与他们为伴，直到他们满足？这绝不会发生。因此，当睡意袭来时，它们才离开她；甚至那时也不离开，因为我们的梦和幻想又给她呈现同样的形象。因此，当白天到来时，盛满这些形象的灵魂往往堕落去实际行出那些幻想。\par

而你，不容一粒灰尘进入你的眼珠，却对你那积聚了一堆如此多罪恶的灵魂视而不见吗？我们何时才能清除这每日积聚在我们内的污秽？何时才能砍掉荆棘？何时才能撒种？你岂不知道收割的时期就在眼前吗？但我们的田地甚至还没有耕过。如果农夫来责问，我们有什么话说？我们作何回答？说没有人给我们种子？可是种子天天都在撒播。那么没有人砍掉荆棘吗？我们每天都在磨快镰刀。但生活必要的牵挂使你分心吗？你为什么没有把自己钉在十字架上死于世界？因为那人仅仅偿还了所领受的，尚且被定为恶，因为他没有使银子加倍；那连自己的银子都浪费的人，将有什么话对他说呢？如果那人被捆绑，被丢在外面的黑暗里，那里有哀号和切齿，我们这些有无数理由被引向德行却退缩不愿的人，将受什么苦呢？\par

因为有什么事物缺少说服你的力量呢？你岂不见世界的卑贱、生命的不定、为现世事物所付出的劳苦和汗水吗？怎么？难道德行必须劳苦追求，而恶习却可不劳而获吗？如果两者都需要劳苦，为什么你不选择这个有如此大益处的呢？\par

或者，德行的某些部分甚至无需劳苦。因为不毁谤、不说谎、不发誓、对邻人息怒，这算是什么劳苦呢？相反，做这些事才是劳苦的，带来许多焦虑。\par

那么，我们有什么可说的，有什么借口，连这些事也不做正呢？由此明显可见，由于懈怠和懒惰，那些更劳苦的义务也完全离我们而去了。\par

我们当思虑这一切；逃避恶行，选择美德，以便藉着我们的主耶稣基督的恩宠与对人的慈爱，既能获得现世的美善，也能获得将来的美善，愿光荣与威能归于祂，至于永世无穷。阿们。\par

\SourceRecord{Translated by George Prevost and revised by M.B. Riddle.；Nicene and Post-Nicene Fathers, First Series；Vol. 10.；Edited by Philip Schaff.；Buffalo, NY: Christian Literature Publishing Co.,；1888.；<http://www.newadvent.org/fathers/200142.htm>.}

% structure: p0001 source=h1 target=section reason=page-title

\section{玛窦福音讲道 第四十三篇}

\BibleRef{玛 12:38-39}\par

\BibleQuote{\Emph{那时，有几个经师和法利塞人回答祂说：\InnerQuote{老师，我们愿意祢显一个征兆给我们看。}但祂回答说：}\Emph{\InnerQuote{邪恶淫乱的世代寻求征兆，除了约纳先知的征兆外，再没有征兆给它。}}}\par

还有什么比这些人更愚蠢的呢（不仅是更不虔敬）？在那么多奇迹之后，好像一个都没行过一样，他们说：\Quote{我们愿意祢显一个征兆给我们看？} 他们这样说话是什么意图呢？是为了再次抓住祂。因为祂的话语已经一次、两次、多次堵住他们的嘴，遏制了他们无耻的舌头，他们就又来攻击祂的作为。对于这一点，福音作者再次惊奇，说：\par

\Quote{那时，有几个经师回答祂，要求一个征兆。}\par

\Quote{那时}，是什么时候？当他们本该俯伏在祂面前，钦佩、惊奇、退让的时候，\Quote{那时}他们却没有停止他们的邪恶。\par

看他们的话也充满了谄媚和虚伪。因为他们认为这样能把祂吸引到他们那边去。现在他们侮辱祂，现在又谄媚祂；时而称祂为附魔的，时而又称祂为\Quote{老师}，两者皆出于邪恶的意念，尽管他们说的话多么相反。\par

因此祂严厉地责备他们。当他们粗暴地质问祂、侮辱祂时，祂温和地与他们辩论；当他们谄媚时，祂则责备地、极其严厉地对待他们，表明祂超脱了两种情感，既不会在一种情况下被激怒，也不会在另一种情况下被谄媚软化。看祂的责备，不仅是严厉的话语，并且包含了他们邪恶的证明。因为祂说什么？\Quote{邪恶淫乱的世代寻求征兆。} 现在祂所说的是这个意思：如果你们对我——至今对你们还是陌生的——这样行为，那有什么奇怪呢？既然你们对父也是如此，你们对祂有如此多的经验。你们离弃祂，投奔魔鬼，吸引邪恶的情人到自己身边。厄则克耳也常常用这些责备他们。\par

\Quote{邪恶淫乱的世代寻求征兆。} 现在祂所说的是这个意思：如果你们对我——至今对你们还是陌生的——这样行为，那有什么奇怪呢？既然你们对父也是如此，你们对祂有如此多的经验。你们离弃祂，投奔魔鬼，吸引邪恶的情人到自己身边。厄则克耳也常常用这些责备他们。\par

因此祂称他们为\Quote{邪恶的世代}，因为他们总是对他们的恩人忘恩负义；因为他们在恩惠之后变得更坏，这属于极端的邪恶。\par

祂称它为\Quote{淫乱的}，表明他们从前和现在的不信；由此祂再次暗示自己与父平等，如果至少不信祂会使它成为\Quote{淫乱的}。\par

\Quote{再没有征兆给它，除了约纳先知的征兆。}\par

那么，有人可能会说，难道没有给它征兆吗？没有给它求的征兆。因为祂行奇迹不是为了吸引他们（祂知道他们心硬），而是为了纠正别人。要么可以这样说，要么就是他们不会得到像那样的征兆。因为一个征兆确实降临到他们身上，当他们因自己的惩罚而认识祂的能力时。这里祂是带着威胁说话，并暗含着这个意思：好像祂说，我显示了无数的恩惠，没有一件能使你们归向我，你们也不愿意敬拜我的大能。因此你们要通过相反的记号认识我的能力，当你们看到你们的城被夷为平地，城墙也被拆毁，圣殿成为废墟；当你们被剥夺了从前的公民身份和自由，到处流离失所，无家可归，被流放时。（因为这一切都在十字架之后发生了。）这些事对你们将是极大的征兆。的确，这是一个极大的征兆，那就是他们的灾祸没有改变；虽然成千上万的人试图改变，却没有人能够推翻已经对他们发出的判决。\par

然而祂没有说这一切，而是留待以后的时间向他们阐明，但目前祂正在试探祂复活的道理，他们将来要通过他们将要遭受的事情来认识它。\par

\Quote{因为正如约纳}，祂说，\Quote{在大鱼腹中三天三夜，人子也要在地心里三天三夜。} \BibleRef{玛 22:40} 这样，祂并没有公开说祂将要复活，因为他们甚至会嘲笑祂，但祂以这样的方式暗示，使他们会相信祂已经预知了这事。至于他们是否意识到，他们对比拉多说：\Quote{那个骗子说过，}这是他们的话，\Quote{祂还活着的时候，\InnerQuote{三天后我要复活}。}然而我们知道祂的门徒们对此一无所知；他们之前比这些人更愚钝，因此这些人成了自我定罪的。\par

但请看他如何精确地表达，即使是在隐晦的言语中。因为他没有说\Quote{在地里}，而是说\Quote{在地心里}；这是为了指明他的坟墓本身，并且使人不会怀疑只是一个表象。为此他也允许了三天的时间，使他的死亡事实得以被相信。因为他不只通过十字架和众人的目睹使之确实，也通过那些日子的时候。因为对于复活，确实后来的整个时期都将为之作证；但十字架，除非当时有许多征兆为之作证，否则就会被怀疑；而这种怀疑也会导致对复活的完全怀疑。因此他也称它为\OriginalTerm{征兆}{sign}。但若他没有被钉十字架，这征兆就不会给出。为此他也提出了预象，使真理得以被相信。因为请告诉我：\Person{约纳}在鲸鱼腹中只是一个表象吗？不，你不能这样说。因此，基督在地心里也不是这样的。因为预象当然不是在真理中，而真理却是在表象中。为此，我们到处彰显他的死亡，无论是在奥迹中，还是在圣洗中，以及其余一切中。因此\Person{保禄}也清楚地呼喊说：\Quote{我绝不夸耀，除非夸耀我们主耶稣基督的十字架。}\par

由此清楚可见，那些受\Person{马尔西翁}式疾病之苦的人，是魔鬼的儿女，他们抹杀这些真理——而基督行了如此多的事正是为了避免它们被废弃，魔鬼却如此费尽心机要使它们被废弃：我指的是他的十字架和他的受难。\par

3. 因此他在别处也说：\Quote{拆毁这座圣殿，三天内我要把它重建起来}；\BibleRef{若 2:19} 又说：\Quote{日子将到，新郎要从他们中间被带走}；\BibleRef{玛 9:15} 而在这里说：\Quote{必不给他征兆，除非\Person{约纳}先知的征兆}；既宣告了他要死，又宣告了\Quote{牛认识主人，驴认识主人的槽}。\BibleRef{依 1:3} 同样在这里，当他通过一个比喻指出他们的顽固后，随后也谈论了他们的惩罚。\par

那么这话是什么意思？他说：正如附魔者，若从那种病中得到释放后稍有疏忽，就会给自己招来比以前更严重的迷惑；你们也是如此。因为从前你们也被魔鬼附身，当你们敬拜偶像，将儿子献给魔鬼时，显出了极大的疯狂；然而我没有舍弃你们，而是通过先知赶出了那个魔鬼；并且我自己又来，愿意更彻底地洁净你们。既然你们不愿听从，反而陷入了更大的邪恶（因为杀害先知远不如杀害他那样严重和可悲）；因此你们的苦难将比以前——我指的是在\Place{巴比伦}、在\Place{埃及}和最初\Person{安提约古}下的那些苦难——更甚。因为在\Person{维斯帕先}和\Person{提托}时期所遭遇的事，远比以前更严重。因此他也说：\Quote{那时要有大灾难，是从来没有过的，也不会再有。}\BibleRef{玛 24:21} 但比喻不仅宣告了这一点，也宣告了他们将完全丧失一切德行，比那时更容易被魔鬼的力量攻击。因为那时，即使他们犯了罪，他们中间也有行为正直的人，天主的眷顾与他们同在，圣神的恩宠照顾、纠正、尽一切本分；但现在，这种守护也将被完全剥夺；他这样告诉他们；因此，现在既有德行的更大匮乏，也有更严重的苦难，以及魔鬼更暴虐的运作。\par

因此，你们知道，甚至在我们这一代，当那个在邪恶上无人能及的人——我指\Person{朱利安}——被他的狂热冲昏头脑时，他们是如何与外邦人站在一起，如何讨好他们的党派。所以，即使现在他们似乎在小程度上受到了一些惩罚，那也是皇帝们的恐惧使他们安静；因为若不是为此，他们所胆敢做的会比从前更糟。因为在他们所有其他的恶行中，他们超过了前人；他们大肆表现出巫术、魔法、不洁。此外，尽管有强烈的约束抑制着他们，他们仍然常常制造叛乱，起来反抗君王，结果导致他们受到最坏灾难的刺穿。\par

现在，那些寻求征兆的人在哪里？让他们明白，需要的是明智的心，如果缺少了这一点，征兆就毫无用处。请看，例如，\Place{尼尼微}人没有征兆就相信了，而这些人在那么多奇迹之后，却变得更坏，使自己成为无数魔鬼的住所，给自己带来万千灾难；这是极自然的。因为当一个人一旦从自己的恶中获释后仍不改正，他将会比从前遭受更严重的苦难。是的，因此他说\Quote{他找不到安息}，以表明那人必然且毫无例外地会被魔鬼的埋伏所追上。因为确实，他本应通过这两件事——他先前的苦难和他的得救——而清醒；或者，还有第三件事也加上了，即还要遭受更糟的威胁。但这一切却没有使他们变好。\par

5. 这一切也适宜地不仅对他们，也对我们说，当我们已经蒙光照，\BibleRef{希 6:4} 并从我们先前的邪恶中得救后，又重蹈覆辙时，我们后来的罪的惩罚也将更严重。因此，基督也对那瘫痪病人说：\Quote{看，你已痊愈；不要再犯罪，以免有更糟的事临到你身上}；\BibleRef{若 5:14} 这是对一位患病三十八年的人说的。有人会问：他能遭受比这更糟的事吗？远更糟、更不可忍受的事。因为千万不要让我们遭受我们所能承受的一切。因为天主并不缺乏惩罚。因为正如他的仁慈浩大，他的义怒也同样浩大。\par

藉此，祂也藉\Person{厄则克耳}责备\Place{耶路撒冷}。祂说：\Quote{我看见你，}祂说，\Quote{你在血中染污；我洗净了你，给你傅了油；你因美丽而获名声；你向邻近的人，}祂说，\Quote{倾泻你的淫行。}因此，当你犯罪时，祂的警告也更加严厉。\par

但由此你不仅推断出你的惩罚，也要推断出天主的无限忍耐。我们多少次至少伸手行同样的恶事，而祂仍长久忍耐！但我们不要自恃，而要恐惧；因为连法郎，若他因第一灾受到教训，就不会经历后来的灾祸；他后来就不会和他的全军一同被淹死。\par

我说这话，是因为我知道有许多人，像法郎一样至今仍在说：\Quote{我不认识天主，}\BibleRef{出 5:2} 并使他们辖下的人紧黏在泥土和砖块上。有多少人，尽管天主命令他们平息他们的\Quote{恐吓，}\BibleRef{弗 6:9} 却连减轻劳苦也不肯！\par

\Quote{但我们如今没有红海可以随后渡过。}但我们有火海，一片无论在种类或大小上都不像那海的火海，而是更大更猛，它的波浪是火焰，是一种奇异可怖的火焰。那里有一个极大的深渊，充满最无法忍受的烈焰。因为到处可见火焰像某种凶猛的野兽飞快地游走。如果在这里这感觉到的、物质的火焰像野兽从火窑中跳出，扑向坐在外面的人，\BibleRef{达 3:22} 那么那另一团火对落入其中的人又会怎样呢？\par

论到那日子，听先知说：\Quote{主的日子是无可救药的，充满愤怒和怒火。}\BibleRef{依 13:9} 因为那时将无人站在一旁，无人救援，也看不见基督那如此温和宁静的面容。但正如那些在矿场工作的人被交给某些残忍的人，看不到任何朋友，只有那些辖管他们的人；那时也是如此：或者不如说并非如此，而是甚至更痛苦。因为在这里，有可能到国王那里去求情并释放被判刑的人；但那里却不再可能；因为祂不允许，他们要继续在灼热的折磨中，在极大的痛苦中，言语无法形容。因为若在这里有人被火烧，没有任何言语能描述他们剧烈的痛苦，更何况那些在那边受苦的人呢？因为在这里一切都很快结束，但在那里却确实在燃烧，但被烧的却不被耗尽。\par

那么我们在那里将做什么？因为我也对自己说这些话。\par

6. 有人说：\Quote{但如果你，}有人说，\Quote{你是我们的老师，对自己这样说，我就不再关心了；若我受罚，又有什么稀奇呢？}不，我恳求，不要让人寻求这种安慰；因为这根本不是慰藉。请告诉我：魔鬼不是无形体的能力吗？他不是比人优越吗？然而他堕落了。有谁会因与他一同受罚而得安慰呢？绝不可能。所有在埃及的人呢？他们不也看见那些居高位的人受罚，每个家庭都哀悼吗？他们当时因此而得释放和安慰吗？当然没有；从他们随后所做的可以明显看出，好像被某种火折磨的人，一同起来反抗国王，迫使他赶出希伯来人。\par

是的，这种说法非常无意义，以为与所有人一同受罚就带来安慰，说：\Quote{大家都这样，我也一样。} 为什么我要谈地狱呢？请想想那些患痛风的人，当他们被剧痛折磨时，即使你让他们看到上万个更痛苦的人，他们也根本不会在意。因为痛苦的剧烈不允许他们的理智有闲暇去想别人，从而找到安慰。所以我们不要用这些冰冷希望喂养自己。因为从邻人的不幸中得安慰，只发生在普通的痛苦中；但当折磨过大，我们内在完全混乱，灵魂甚至不能认识自己时，它从哪里得安慰呢？所以所有这些说法都是荒谬的，是愚昧儿童的幻想。因为你所谈论的这种情况，发生在沮丧中，且是适度的沮丧，当我们被告知\Quote{某人遭遇了同样的事；}但有时甚至在沮丧中也不发生：若在那种情况下它都没有力量，何况在那无法形容的痛苦和重负中呢，就是\Quote{咬牙切齿}所指的。\par

我知道我这些话刺痛你们，使你们痛苦；但我能做什么呢？因为我本不愿这样说，而是希望在自身和你们众人中都意识到德行；但由于我们多数人处在罪中，谁又能赐我能力真正刺痛你们，并深入听众的理解呢？那样我就可以安心了。但现在我怕有人轻视我的话，因为他们漫不经心地听，惩罚就会更大。因为当主人发出威胁时，若一个同仆听见并轻视他的警告，他定会受惩罚，被激怒的主人会匆匆经过他，甚至加重责罚。因此我恳求你们，当我们听到祂关于地狱的言论时，让我们刺透自己的心。因为没有什么比这言论更令人愉悦，正如没有什么比现实更痛苦。但有人会问：谈论地狱怎么会令人愉悦呢？因为落入地狱是远非愉悦的，而我们这些看似刺耳的话恰恰阻止了这种结果。在此之前它们还提供另一种快乐：它们振奋我们的灵魂，使我们更加虔诚，提升心智，给思想插上翅膀，并驱逐那恶毒地困扰我们的欲望；这事就成了医治。\par

7. 因此，接着让我也谈一下惩罚中的羞耻。因为正如犹太人那时将受尼尼微人的审判，同样我们也要受许多现在看来不如我们的人的审判。\par

那么让我们想象那是多么大的嘲笑，多么大的定罪；让我们想象，并最终打下一些基础，打开一扇悔改之门。\par

這些話我是對自己說的，我先給自己這勸告，但願沒有人因覺得受責備而惱怒。讓我們踏上窄路。奢侈要到幾時？怠惰要到幾時？懶散、嬉笑、拖延，我們還不夠嗎？豈不又是同樣的事：宴飲、飽食、花費、財富、積蓄、建築？結局是什麼？死亡。結局是什麼？灰燼、塵土、棺材、蟲子。\par

那麼，讓我們展現一種新的生活。讓我們使大地成為天堂；讓我們藉此向希臘人表明，他們被剝奪了何等大的福分。因為當他們在我們身上看到良好的品行時，他們就會看到天國的面貌。是的，當他們看到我們溫良、純潔，沒有憤怒、邪惡慾望、嫉妒、貪婪，正確履行所有其他職責時，他們會說：\Quote{如果基督徒在這裡成了天使，那麼他們離開這裡後會是怎樣的呢？如果他們在作為客旅的地方如此發光，那麼當他們贏得故鄉時，他們會變得何等偉大！}這樣，他們也會改正，而虔敬的話語\Quote{會順利奔馳}，\BibleRef{得後 3:1}，不亞於宗徒時代。因為如果他們十二個人就使整個城市和地區歸化；如果我們所有人都因謹慎的品行而成為導師，請想像我們的聖業將被高舉到何等地步。因為甚至一個死人復活，也不如一個克己的人那樣有力地吸引希臘人。對前者他會驚奇，但對後者他會受益。前者發生並過去了；但後者常存，並且是對他靈魂的持續培育。\par

因此，讓我們注意自己，好使我們也能贏得他們。我不說任何沉重的事。我不說：不要結婚。我不說：拋棄城市，退出公共事務；而是要投身其中，展現德行。是的，我寧願在城中忙碌的人比佔據山頭的人更受讚許。為什麼？因為由此產生的益處很大。\Quote{因為沒有人點燈，放在斗底下。}因此，我願意所有的燈都放在燈台上，好使光變大。\par

那麼，讓我們點燃祂的火；讓我們使那些坐在黑暗中的人脫離他們的錯誤。不要對我說：\Quote{我有妻子和兒女，是一家之主，不能妥善實行這一切。}因為即使你沒有這一切，如果你粗心大意，一切都會失去；即使你被這一切包圍，如果你熱心，你必能達到德行。因為所需要的只有一件事：一顆慷慨之心的準備；無論年齡、貧窮、財富、命運的逆轉，還是任何其他事物，都不能阻礙你。因為事實上，無論老少、有妻子並養育兒女的人、以手工為業的人、當兵的人，都妥善履行了所吩咐的一切。因為達尼爾年輕，若瑟是奴隸，阿桂拉從事手工，賣紫布的女人經營作坊，另一個是獄卒，另一個是百夫長，如科爾乃略；另一個身體虛弱，如弟茂德；另一個是逃奴，如敖尼息摩；但這些對他們任何人都不是障礙，而是所有人，無論男女、老少、奴隸、自由人、士兵、平民，都獲得了讚許。\par

因此，我們不要製造虛假的藉口，而是要預備一顆徹底良善的心，無論我們是誰，我們必能達到德行，並獲得將來的福分；藉著我們的主耶穌基督的恩寵和對人的愛，願光榮、權能、尊崇歸於祂，偕同聖父及聖神，現今及永遠，世世無盡。亞孟。\par

\SourceRecord{Translated by George Prevost and revised by M.B. Riddle.；Nicene and Post-Nicene Fathers, First Series；Vol. 10.；Edited by Philip Schaff.；Buffalo, NY: Christian Literature Publishing Co.,；1888.；<http://www.newadvent.org/fathers/200143.htm>.}

% structure: p0001 source=h1 target=section reason=page-title

\section{玛窦福音讲道 44}

\Emph{\BibleRef{玛 12:46-49}}\par

\BibleQuote{ \Emph{祂正同群众讲话时，看，祂的母亲和祂的兄弟站在外面，想要同祂说话。有人告诉祂说：看，祢的母亲和祢的兄弟站在外面，想要同祢说话。祂却回答那告诉祂的人说：谁是我的母亲？谁是我的兄弟？祂遂伸手指着门徒说：看，我的母亲和我的兄弟！} }\par

我最近所说的，即当缺乏德行时，一切都是虚妄的，这一点如今也极为充分地显明了。因为我确实说过，年龄、本性以及独居旷野等一切，若无良好的心智，均无益处；但今天我们还学到另一件事：即使在母腹中怀有基督并生出那奇妙的诞生，若没有德行，也毫无益处。\par

这一点由此尤其明显。\Quote{因为祂正同群众讲话时}，经上记载：\Quote{有人告诉祂说：祢的母亲和祢的兄弟在外面找祢。祂却说：谁是我的母亲？谁是我的兄弟？}\par

祂这样说，并非因祂的母亲而羞愧，也不是否认那生育祂的；因为若祂以她为耻，就不会经过那母胎了；而是表明她若不行一切所当行的事，便毫无益处。事实上，她所试图做的，是出于多余的虚荣；因为她想向群众显示她对祂儿子有权柄和权威，当时对祂尚未有伟大思想；因此她的行事不合时宜。无论如何，看她和她亲属的自信。因为他们本应进去与群众一同聆听，或者若不如此，也应等待祂讲完话后再近前；他们却叫祂出来，并在众人面前这样做，显出多余的虚荣，并想让人以为他们能以很大权柄命令祂。福音作者也表明他在责备这一点，因为他正是带着这个暗示这样表述：\Quote{祂正同群众讲话时}，好像说：什么？难道没有别的机会吗？难道不能私下与祂交谈吗？\par

他们究竟想说什么？若涉及真理的道理，他们理当公开提出，在众人面前陈述，使其余人也得到益处；但若是关于他们自己的其他事情，就不应如此急切。因为祂尚且不允许埋葬父亲，以免中断祂的服务，那么他们更不应为微不足道的事而打断祂向群众的讲道。由此可见，引领他们做此事的无非是虚荣；若望也如此声明，说：\BibleQuote{连祂的兄弟也不信从祂}\BibleRef{若 7:5}，他还记载了他们的一些充满极大愚昧的话语；告诉我们，他们想拉祂去耶路撒冷，别无目的，只为他们自己能因祂的奇迹而得荣耀。\BibleQuote{因为祢若行这些事}，据说，\BibleQuote{就当将自己显明给世界。因为没有人行事隐秘，却还想要显明自己；}当时祂也责备他们，将此归因于他们的血肉思想。就是说，因为犹太人指责祂说：\BibleQuote{这不是木匠的儿子吗？他的父亲和母亲我们不是知道吗？他的兄弟不是同我们在一起吗？}他们想要摆脱因祂出身而带来的轻视，便叫祂去显奇迹。\par

为此，祂毅然拒绝他们，有意治愈他们的软弱；因为倘若祂有意否认祂的母亲，祂会在犹太人指责她时否认她。但事实上，我们看见祂如此照顾她，甚至在十字架上，将祂托付给祂最爱的门徒，并交给他一项关于她的重大托付。\par

但如今祂不这样做，乃是出于对她和祂兄弟的关心。我的意思是，因为他们看待祂如同一个平常人，并且充满虚荣，祂驱除这疾病，并非侮辱，而是纠正他们。\par

但我恳求你，不仅要考察那含有温和责备的话语，也要考察祂兄弟的不当行为，以及他们所表现的狂妄，并那责备者是谁——不是凡人，而是天主的独生子；以及祂责备的目的——并非有意令他们困惑，而是为将他们从那最专制的激情中解救出来，逐步引导他们正确认识自己，并使她明白，祂不仅是她的儿子，也是她的主：这样你就会看出，这责备是最高度适合祂且有益于她的，同时也含有极大的温柔。因为祂并没有说：\Quote{你走吧，告诉我的母亲，你不是我的母亲，}而是对那告诉祂的人说话，说：\Quote{谁是我的母亲？}同时，如已提及的，也为另一个目的。那是什么呢？就是叫他们和任何信赖亲属关系的人，不要忽略德行。因为若她因是祂的母亲而无益处，除非她具有那品质，那么几乎没有别人能靠亲属关系得救。因为唯一的高贵，就是承行天主的旨意。这种高贵出身比另一种更高贵，也更真实。\par

2. 因此，既然知道这些事，我们不可因儿女有美名而自夸，除非我们拥有他们的德行；也不可因出身高贵而自夸，除非我们在气质上相似他们。因为那生养人的人，可能不是他的父亲；而那并未生养他的人，却可能是他的父亲。因此，在另一处，当一个妇人说：\Quote{怀过祢的胎和祢所吸过的乳房是有福的。}祂并没有说：\Quote{那胎没有怀过我，我也没有吸过那乳房，}而是说：\Quote{是的，但更有福的是那些遵行我父旨意的人。}你看见了吗？祂在每一次机会中都不否认本性上的亲属关系，而是添加德行上的亲属关系。祂的前驱也说：\Quote{毒蛇的种类！不要想这样说：\InnerQuote{我们有亚巴郎作我们的父亲。}}\BibleRef{玛 3:7}这并不是说他们不是出于亚巴郎的，而是说他们出于亚巴郎毫无益处，除非他们有品格上的亲属关系；基督也宣告了这一点，当祂说：\Quote{如果你们是亚巴郎的子女，你们就该作亚巴郎的作为；}\BibleRef{若 8:39}这并非剥夺他们按肉体的亲属关系，而是教导他们追求比那更大、更真实的亲属关系。\par

祂在这里也确立了这一点，但方式更温和、更有分寸，正如祂对母亲说话时应有的那样。祂根本没有说：\Quote{她不是我的母亲，那些人也不是我的兄弟，因为他们不遵行我的旨意；}祂也没有宣布和判定他们的罪，而是把选择权留在他们自己手中，以祂应有的温和说话。\par

\Quote{因为凡遵行，}祂说，\Quote{我父旨意的，这人就是我的兄弟、姊妹和母亲。}\par

因此，如果他们愿意如此，就让他们走这条路吧。当那妇人再次喊道：\Quote{怀过祢的胎是有福的，}祂并没有说：\Quote{她不是我的母亲，}而是说：\Quote{如果她愿意有福，就让她遵行我父的旨意吧。因为这样的人就是兄弟，也是姊妹，也是母亲。}\par

何等尊荣！何等德行！她把追随她的人领到何等的高度！有多少妇人祝福过那位圣童贞女和她的胎，并祈求自己也能成为那样的母亲，并放弃一切！那么，有什么阻碍呢？看，祂为我们开辟了一条宽阔的道路；不仅妇女，而且男人也被准许获得这一等次，或者说是更高一等的。因为这使她成为祂的母亲，远比那些产痛更甚。所以，如果那是一个蒙福的理由，那么这更是，因为这也更真实。因此，不要仅仅渴望，而且要在那引你达到渴望的道路上，以极大的勤奋前行。\par

3. 说了这些话之后，\Quote{祂从屋里出来。}你看见了吗？祂既责备了他们，又做了他们所愿的事。祂在加纳婚宴上也这样做了。\BibleRef{若 2:1}因为在那里，祂也同时责备她不合时宜的请求，却又没有拒绝她；前者纠正她的软弱，后者显示祂对母亲的慈爱。同样，在此场合，祂也医治了虚荣心，并向祂的母亲给予了应有的尊荣，尽管她的请求不合时宜。因为据说：\Quote{就在那一天，耶稣从屋里出来，坐在海边。}\par

为什么？如果你们愿意看见和听见，祂说，看，我出来讲道。这样，行了许多奇迹之后，祂再次提供祂教训的益处。祂\Quote{坐在海边，}垂钓，把岸上的人收进祂的网里。\par

但祂\Quote{坐在海边}并非没有目的；这一点，福音作者也隐含地表达了。因为为表示祂这样做的原因是希望精确地安置祂的听众，不把任何人留在祂背后，而是让所有人都面对祂，\par

\Quote{有许多群众，}祂说，\Quote{聚集到祂跟前，以致祂上了船坐下，群众都站在岸上。}\par

祂在那里坐下，用比喻讲话。\par

\Quote{祂就用比喻，}经上说，\Quote{对他们讲了许多道理。}\BibleRef{玛 13:3}\par

然而我们知道，在山上祂并没有这样做，没有用这么多比喻编织祂的讲论，因为当时只有群众和单纯的百姓；但在这里还有经师和法利塞人。\par

但请你注意，我祈求你，祂先讲的是哪一种比喻，以及玛窦是如何按顺序排列它们的。那么，祂先讲的是哪一个？那是最需要先讲的，那使听众更加留心的。因为祂要用隐晦的话向他们讲论，祂首先用比喻彻底唤起听众的心。因此，另一位福音作者说，祂责备他们，因为他们不明白；说：\Quote{你们怎么不明白这比喻呢？}但祂用比喻说话不仅为此，还要使祂的讲论更加生动，更完美地铭记在听众心中，并将事物呈现在他们眼前。先知们也是这样做的。\par

4. 那么，这比喻是什么？\Quote{看，}祂说，\Quote{有一个撒种的出去撒种。}那充满万有、无所不在的祂，从何处出去呢？或者说，祂怎样出去呢？不是地方上的，而是条件上和对我们的安排上的，通过披上肉身，更接近我们。因为我们不能进入，我们的罪把我们挡在入口之外，祂就出来到我们这里。祂出来做什么？要毁灭那长满荆棘的土地？要报复那农夫？绝不是；而是要耕种、照料它，并撒下虔敬的道。因为这里的种子祂指的是祂的教训，土地是人的灵魂，撒种者是祂自己。\par

那么，这种子结果如何？三份失丧，一份得救。\par

\Quote{当祂撒种的时候，有些种子落在，}祂说，\Quote{路旁，飞鸟来把它们吃尽了。}\par

祂没有说祂扔下它们，而是说它们\Quote{落下了。}\par

\Quote{有些种子落在石地上，那里没有多少土壤；因土壤不深，立刻发了芽；太阳升起后，就被晒焦了；又因为没有根，就干枯了。有些落在荆棘中，荆棘长起来，便把它窒息了。可是有的落在好地里，就结了果实：有一百倍的，有六十倍的，有三十倍的。有耳的，听罢！}\par

四分之一得救；而且这四分之一也不是完全一样，而是其中也有很大的差别。\par

祂说这些话，表明祂毫无保留地向众人讲论。因为正如播种者对所耕种的土地不加区别，只是单纯而平等地撒种；同样，祂也不分富人和穷人、智者和愚者、懒惰和勤勉、勇敢和懦弱，而是向众人讲论，尽自己的本分，尽管祂预知结果；这样祂可以说：\Quote{我当做什么，却未曾做呢？}\BibleRef{依 5:4} 先知们也论及人民如同葡萄树：经上说：\Quote{我所爱的，有一个葡萄园；}\BibleRef{依 5:1} 以及\Quote{从埃及移来一棵葡萄树；}但祂却像撒种一样。这要表明什么呢？就是现在的服从将是迅速而比较容易的，并且会立即结出果实。\par

当你听到\Quote{播种者出去撒种}时，不要以为是多余的重复。因为播种者常常也出去做别的事情，比如耕地、拔除杂草、拔除荆棘，或处理类似的事；但祂是出去撒种。\par

那么请问，为什么大部分种子丢失了呢？不是因为播种者，而是因为接受种子的土地，即那不听从的灵魂。\par

祂为什么不说：有些种子因漫不经心的人接收而失去；有些被富人窒息；有些被肤浅的人出卖？祂不愿严厉责备他们，以免使他们陷入绝望，而是将责备留给听众的良心。\par

不仅种子如此，网也是：因为网也打出了许多无用的东西。\par

5. 祂讲这个比喻，是为了鼓励祂的门徒，教导他们：即使迷失的人比接受圣言的人更多，他们也不应沮丧。因为他们的主尚且如此，祂完全预知这些事会发生，仍不停止撒种。\par

有人说：把种子撒在荆棘中、石头上、路旁，这怎么合理呢？对于种子和土地来说，确实不合理；但对于人的灵魂和教导来说，这却是值得赞扬的。因为农夫若这样做，确实应受责备：石头不可能变成土地，路旁也不可能不再是路旁，荆棘仍是荆棘。但在有理智的事物中则不然：石头可以改变，成为沃土；路旁不再被践踏，不再向所有过路人敞开，而能成为肥沃的田地；荆棘可以除掉，种子得到充分安全。如果这是不可能的，这位播种者就不会撒种。而如果并非所有人都改变，这不是播种者的过错，而是那些不愿意改变的人之过：祂已尽了自己的本分；如果他们辜负了从祂那里领受的，祂是无辜的，是展示了如此大爱的人。\par

但请你注意：灭亡之路不止一条，而是有不同且相差甚远的道路。因为像路旁的，是粗俗、冷漠、漫不经心的人；而像石头地的，是仅仅因软弱而跌倒的人。\par

因为祂说：\Quote{那撒在石地上的，是指人听了圣言，立刻欣然接受；但他心里没有根，只是暂时的；一旦为圣言遭遇患难或迫害，立刻就跌倒了。}祂又说：\Quote{凡听了圣言却不明白的，那恶者就来，把撒在他心里的夺去。这是撒在路旁的。}\par

但圣言枯萎，是在无人迫害或动摇其基础的情况下，与在试探紧逼时不同。那些比作荆棘的人，比这些更不可原谅。\par

6. 为此，为了不使我们遭遇这些事，让我们以热忱和不断的记忆来保守所听的道。因为即使魔鬼夺去它们，但被夺去与否取决于我们；即使植物枯萎，这不是因为炎热（祂没说因为炎热而枯萎，而是说\Quote{因为没有根}）；即使祂的话被窒息，也不是因为荆棘，而是因为那些让荆棘生长的人。因为如果你愿意，有一种方法可以抑制这种邪恶的生长，并善用我们的财富。所以祂没有说\Quote{世界}，而是说\Quote{世俗的焦虑}；没有说\Quote{财富}，而是说\Quote{财富的迷惑}。\par

因此，我们不要责怪事物，而要责怪败坏的思想。因为人可以富有而不被迷惑，可以在世界之中而不被其忧虑窒息。财富确实有两种相反的不利之处：一是忧虑，耗尽我们，给我们带来黑暗；二是奢侈，使我们软弱。\par

祂说得很好：\Quote{财富的迷惑。}因为与财富相关的一切都是欺骗；它们只是名称，并不附着于事物。因为快乐、荣耀、华丽的服饰以及所有这些事物，都是一种虚浮的炫耀，而非现实。\par

在讲完毁灭的道路之后，祂接着提到好地，不让他们绝望，而是给予悔改的希望，并指出从前面所说的情况转变到好地是可能的。\par

然而，即使土地是好的，播种者是一位，种子也是同样的，为何有的结一百倍，有的六十倍，有的三十倍呢？这里，差别再次源于土地的性质，因为即使土地是好的，其内也有很大差别。你看见了吗，受责备的不是农夫，也不是种子，而是接受种子的土地？不是就其本性，而是就其心态。在此，祂对人的仁慈也显为大，因祂不要求单一的德行尺度，而是接纳第一等，不抛弃第二等，也给第三等留一席之地。\par

祂说这些话，免得跟随祂的人以为单凭听讲就足以获得救恩。有人可能会说：祂为什么没有提到其他恶习，如情欲、虚荣？祂在谈到\Quote{今世的挂虑和财富的迷惑}时，已把一切都包括了。是的，虚荣及其他一切都属于这世界和财富的迷惑；诸如逸乐、贪食、嫉妒、虚荣等类。\par

祂还加上了\Quote{路旁}和\Quote{石地}，表明仅仅摆脱财富是不够的，还必须培养德行的其他部分。因为即使你确实摆脱了财富，却仍软弱怯懦，那又怎样呢？或者你并不怯懦，却对听道懈怠疏忽，那又如何？可见，单单一方面并不足以使我们得救，而是需要先有认真的听道和不断的忆念，然后有刚毅，最后是轻视财富并从一切世俗事物中解脱出来。\par

事实上，祂把这一点放在其他之前，是因为这点是首先需要的（因为\BibleQuote{若未曾听见，怎能信祂呢？}\BibleRef{罗 10:14}；同样，我们若不留意所说的话，就甚至无法学习自己该做什么）；之后便是刚毅和对现实的轻视。\par

7. 因此，听了这些话，让我们从各方面坚固自己，牢记祂的教训，扎下深根，从一切世俗事物中洁净自己。但若我们做了一方面，忽视了另一方面，就不会有任何益处；因为即使我们不因这种方式灭亡，也必因另一种方式灭亡。因为如果我们不因财富而毁掉，却因懈怠而毁掉呢？或不因懈怠，却因软弱而毁掉呢？正如农夫，无论以这种或那种方式损失收成，都同样哀叹。因此，让我们不要因自己没有以所有这些方式灭亡而自慰，而应因我们是以何种方式灭亡而悲伤。\par

让我们烧尽荆棘，因为它们窒息了圣言。这一点是那些富人所知道的，他们不仅在这些事上，而且在其他事上也证明是无用的。因为他们成了自己逸乐的奴隶和囚徒，不仅在世俗事务上无用，若在世俗上无用，在天国的事上就更无用了。是的，我们的思想由此在两方面被败坏：一方面因奢侈，另一方面也因焦虑。这两者中的任何一个都足以使小船倾覆；而当两者同时发生时，想象那波涛有多高。\par

不要奇怪祂称我们的奢侈为\Quote{荆棘}。因为，你虽因沉溺于情欲而不自知，但那些健康的人知道，它比任何荆棘刺得更痛，奢侈比焦虑更能耗费灵魂，给灵魂和肉体带来更痛苦的伤害。因为焦虑的打击不像暴饮暴食那样严重。当失眠、太阳穴跳动、头重、腹绞痛抓住这样的人时，你可以想象这些荆棘在痛苦上超过多少。正如荆棘无论从哪边被抓住，都会使抓它的手流血，同样，奢侈折磨脚、手、头、眼，以及我们所有的肢体；它也枯干、不结果，如同荆棘，且比它伤害更大，且在要害部位。是的，它带来早衰，迟钝感官，蒙蔽理性，使敏锐的头脑变瞎，使身体肿胀，导致废物过度沉积，积聚大量邪恶；它使负担过大，重担压垮；因此我们跌倒频繁，船难不断。\par

请告诉我，为什么要娇养你的身体？怎么？是要把你当作祭品宰杀，摆在桌子上吗？娇养鸟雀倒是可以，或者不如说，即使对它们也不好；因为当它们被养肥时，对健康饮食也无益。奢侈之害竟如此之大，甚至在不理性的动物身上也显明了。因为即使是它们，我们也通过奢侈使它们变得对自己和我们都无用。因为它们多余的肉难以消化，那种肥腻会引起一种较湿润的腐败。而那些不这样喂养，而是可以说生活在节制、适度饮食、劳作和艰辛中的动物，对自己和他人最为有益，无论作为食物还是其他方面。至少，吃它们的人身体更健康；但那些以另一种为食的人则像它们一样，变得迟钝多病，使自己的枷锁更沉重。因为没有比奢侈更敌视、更有害于身体的了；没有比放纵更能撕裂、压垮、败坏身体的了。\par

因此，最令人惊讶的是，这些人竟如此愚昧，连别人对待酒囊的那点小心，他们也不愿用在自己身上。因为酒商不让酒囊装得过多，以免破裂；但这些人对于自己可怜的肚腹，却连这点先见也不肯有，反而将它填满、撑胀，装到耳朵、鼻孔、直到咽喉，从而将精神和指导活物的力量挤压到只剩一半的空间。怎么？你的喉咙是为了这个目的而被赐予的吗？是为了让你把它灌满酸败的酒和各样腐败之物，直到嘴边吗？不是为这个，人啊，而是为了使你首先向天主歌唱，献上神圣的祈祷，宣读神圣的法律，并给邻人有益的劝告。但你却好像它被赐予你就是为此目的，连片刻也不让它有空从事那职务，反而一生都使它屈从这个邪恶的奴役。就好像有人得到一把里拉琴，琴弦是金的，结构精美，却不去用它奏出最和谐的乐曲，反而用大量粪土和泥巴把它盖满；这些人正是如此。我说\Quote{粪土}，不是指生活必需品，而是指奢侈的生活和极大的放纵。因为超出必需的东西不是滋养，而只是有害。事实上，肚腹被造只是为了接收食物；但口、喉咙和舌头，则是为了其他远更必要的事；或者不如说，连肚腹也不是单纯为了接收食物，而是为了接收适度的食物。这由它在我们因贪婪而刺激它时大声呼喊抗议便可证明；它不仅向我们呼喊，而且为那伤害报复我们，施加最严厉的惩罚。首先它惩罚那带领我们前往邪恶宴饮的脚，然后惩罚那为它服务的手，将它们捆绑起来，因为它们带给它如此多量和种类的食物；许多人的嘴、眼睛和头都因此扭曲。如同一个仆人接到超出自己能力的工作，常常因绝望而对抗主人；同样，肚腹也常连同这些肢体，因过度紧张而毁坏甚至破坏大脑本身。天主安排得极好：从过度中产生如此多的祸害，为的是即使你自愿不肯节制，至少因惧怕如此巨大的毁灭而不得不学会节制。\par

既然知道这些，让我们逃避奢侈，学习节制，好使我们既享有身体的健康，又使灵魂脱离一切软弱，得以获得将来的美善，借着我们的主耶稣基督的恩宠和对人的爱，愿光荣和权能归于他，直到永远。阿们。\par

\SourceRecord{Translated by George Prevost and revised by M.B. Riddle.；Nicene and Post-Nicene Fathers, First Series；Vol. 10.；Edited by Philip Schaff.；Buffalo, NY: Christian Literature Publishing Co.,；1888.；<http://www.newadvent.org/fathers/200144.htm>.}

% structure: p0001 source=h1 target=section reason=page-title

\section{玛窦福音讲道四十五}

\BibleRef{玛 13:10-11}\par

\Quote{\Emph{门徒前来对祂说：\InnerQuote{祢为什么用比喻对他们讲话？}祂回答他们说：\InnerQuote{因为天国的奥秘是给你们知道的}}\Emph{知道天国的奥秘，但不是给他们知道的。}}\par

我们有充分的理由钦佩门徒们，他们虽渴望学习，却知道何时该问。因为他们不在众人面前这样做；玛窦藉着说\Quote{他们前来}表明了这一点。而且，这说法并非推测，马尔谷更清楚地表达了这一点，说他们\Quote{私下到祂跟前}。（见：\BibleRef{谷 4:10}）那么，祂的兄弟和母亲也该这样做，而不该把祂叫出来，引人注目。\par

但也要留意他们的亲切关爱，他们对他人多么关心，首先寻求他们的利益，然后才求自己的。经上记载：\Quote{为什么祢用比喻对他们讲话？}他们并没有说：\Quote{为什么祢用比喻对我们讲话？}是的，在其他场合，他们对他人的慈爱也在多方面显明出来；例如他们说：\Quote{请遣散群众}（见：\BibleRef{路 9:12}），以及\Quote{祢知道他们被冒犯了吗？}（见：\BibleRef{玛 15:12}）。\par

那么基督说了什么？祂说：\Quote{因为天国的奥秘是给你们知道的，但不是给他们知道的。}但祂这样说，并非引入必然性，或任意武断的分配，而是暗示他们是自己一切罪恶的制造者，并愿意表明这是一件恩赐，是从上而来的\OriginalTerm{恩宠}{grace}。\par

然而，绝不能说因为这是一件恩赐，就因此取消了\OriginalTerm{自由意志}{free will}；这从后面的话中可以清楚看出。为此，为使前者不绝望，后者不粗心，在被告知\Quote{是给予的}之后，祂表明起点在于我们自己。\par

\BibleQuote{因为凡有的，还要给他，叫他有余；凡没有的，连他所有的也要夺去。}\par

虽然这话充满许多晦涩，却指明了说不尽的正义。因为祂的话大意如下：当一个人有热忱和热诚时，从天主方面也将赐给他一切；但如果他缺乏这些，也不付出自己的一份，那么天主的恩赐也不会赐予。因为甚至\Quote{他所有的}，祂说，\Quote{也要被夺去}；天主并非如此夺去，而是视他为不配领受祂的恩赐。我们也是如此做的：当我们看到有人漫不经心地听，尽管我们再三恳求也无法说服他留心时，我们只得沉默。因为如果我们继续下去，他的漫不经心只会加重。但对于那努力求学的人，我们引导他，并倾注许多。\par

祂说得很好：\Quote{连他所有的。}因为他实际上连这个也没有。\par

接着祂又把所说的话说得更清楚，指明了\Quote{凡有的，还要给他；凡没有的，连他所有的也要夺去}这句话的含义。\par

祂说：\Quote{因此，我用比喻对他们讲话，因为他们看却看不见。}（见：\BibleRef{玛 13:13}）\par

有人可以说：\Quote{那么，若他们看不见，就该开他们的眼睛。}不，如果瞎眼是自然的，就该开；但因为这是一种自愿、自选的瞎眼，所以祂没有简单地说\Quote{他们看不见}，而是说\Quote{看却看不见}；这样瞎眼乃是出于他们自己的邪恶。因为他们看见魔鬼被驱逐，就说：\Quote{祂是靠魔王贝耳则步驱魔。}（见：\BibleRef{玛 12:14}）他们听见祂引导他们归向天主，显示祂与天父极大的合一，就说：\Quote{这个人不是从天主来的。}（见：\BibleRef{若 9:16}）既然他们所作的判断既违背所见又违背所闻，因此祂说，连听闻本身我也从他们夺去。因为他们从中得不到任何益处，反而招致更大的定罪。因为他们不仅不相信，而且还挑剔、控告、设陷阱。然而，祂没有说这些，因为祂不愿以指控使他们反感。因此，祂起初也没有这样对他们讲论，而是非常直白地讲；但因他们自甘堕落，从此祂就用比喻讲话。\par

2. 在此之后，为避免有人认为祂的话纯属指控，也避免人们说：\Quote{祂是我们的仇敌，在控告、诽谤我们；}祂也引入了先知，宣告与祂自己相同的判决。\par

祂说：\Quote{因为在他们身上应验了依撒意亚的预言，说：\InnerQuote{你们听是听见，却不明白；看是看见，却不认辨。}}\par

你看见先知也同样精确地指控他们吗？因为祂没有说\Quote{你们看不见}，而是\Quote{你们看却认不清}；也没有说\Quote{你们听不见}，而是\Quote{你们听却不明白}。所以他们首先自己加害了自己，塞住耳朵，闭上眼睛，使心迟钝。因为他们不仅听不见，而且\Quote{听得发沉}，祂说，他们这样做，\par

\Quote{免得他们转回，我就医好他们；}描述了他们严重的邪恶和决意离弃祂的心。祂这样说，是为了吸引他们归向祂，刺激他们，并表明如果他们回转，祂会医好他们：好比有人说：\Quote{他不肯看我，我感谢他；因为他若垂顾我，我立刻就会屈服。}他说这话，表明他如何愿意和好。同样在此处也说：\Quote{免得他们转回，我就医好他们；}暗示他们的回转是可能的，并且他们若悔改就能得救，而且祂做这一切并非为自己的荣耀，而是为他们的得救。\par

因为如果那不是祂的意愿——即他们应该听见而得救——祂本应保持沉默，不藉比喻说话；但如今，正是藉着这件事，甚至藉着在面纱下说话，祂激励了他们。\Quote{因为天主不愿罪人丧亡，却愿他归依于祂而生活。}\par

为了证明我们的罪不属于本性，也不属于必然和强迫，请听祂对宗徒们说的话：\Quote{但你们的眼睛有福，因为看得见；你们的耳朵有福，因为听得见；}\BibleRef{玛 13:16}这里所说的看见和听见，不是指肉体的，而是指心灵的。事实上，这些人也是犹太人，并在相同的环境中长大；但他们从未因这预言而受伤害，因为他们在自己内牢固地确立了祂祝福的根，即\OriginalTerm{选择原则}{principle of choice}和他们的判断。\par

你看见了吗？\Quote{给你们知道了}并非出于必然。因为如果善行不是他们自己的，他们也不会被称作有福。不要对我说这话讲得隐晦；因为他们本来可以来问祂，如同门徒们那样；但他们不愿，因他们懈怠和疏忽。为什么我说他们不愿？不仅如此，他们还做了相反的事：不仅不信、不听，甚至还敌对，并准备对祂所说的一切怀有恶意——这正是先知归咎于他们的，说：\Quote{他们听得很沉重。}\par

但这些人却非如此；因此祂也祝福了他们。接着，祂以另一种方式再次向他们保证说：\par

\Quote{我实在告诉你们：许多先知和义人，想看你们所看见的，却没有看见；想听你们所听见的，却没有听见。}\BibleRef{玛 13:17}祂指的是我的来临，我的奇迹，我的声音，我的教导。因为在这里，祂不仅将他们置于这些堕落者之上，甚至置于那些行善的人之上；而且祂宣称他们比后者更有福。这是为什么呢？因为他们不仅看见了犹太人没有看见的，甚至看见了古人所渴望看见的。古人只藉信德看见，而这些门徒却用肉眼也看见了，而且更为清晰。\par

你看见了吗？祂再次将\OriginalTerm{旧的安排}{old dispensation}与新的联系起来，表明古人不仅知道未来的事，而且非常渴望它们。但如果这些事属于某个陌生敌对的神，他们绝不会渴望它们。\par

祂说：\Quote{所以你们要听这撒种的比喻。}\BibleRef{玛 13:18}接着祂提到了我们先前所讲的，关于疏忽和留意，关于怯懦和\OriginalTerm{刚毅}{fortitude}，关于财富和\OriginalTerm{自愿贫穷}{voluntary poverty}，指出前者带来的损害和后者带来的益处。\par

然后祂也提出了美德的不同形式。因为祂充满对人的爱，不仅指出一条路，也没有说：\Quote{除非有人结出一百倍，他就是被弃绝的。}而是说结出六十倍的也得救，并且不仅是他，就连结出三十倍的也得救。祂这样说，表明了救恩是容易的。\par

3. 那么，你无法修持\OriginalTerm{贞洁}{virginity}吗？就在婚姻中守洁。你无法舍弃你的财产吗？就捐出你的财物。你承担不了那个重担吗？就与基督分享你的财富。你不愿把一切献给祂吗？就把一半、三分之一给祂。祂是你的兄弟、\OriginalTerm{共同继承人}{joint-heir}，也要在这里使祂成为你的共同继承人。无论你给祂什么，都是给了你自己。你难道没有听见先知说的话吗？\Quote{对你后裔中的人，你不可忽视。}如果我们不能让我们的亲属被遗忘，更何况我们的主呢？祂不仅以主的权威对待你，还有亲属的名分，以及其他更多。因为祂也使你分享了祂的财富，没有从你那里收取任何东西，反而以这不可思议的恩赐开始。那么，若非极度的麻木不仁，怎能不因这恩赐而变得仁慈？怎能不因白白的恩赐而有所回报？怎能不以较小的回报更大的恩典？祂使你成为天堂的继承人，你难道连地上的事物也不分给祂吗？当你还没有行善，甚至是敌人时，祂就与你和好了；作为朋友和恩人，你难道不报答祂吗？\par

然而，诚然，即使在王国和其余一切之前，单凭祂赐予这一事实，我们就应当感到亏欠于祂。因为仆人们也是如此，当他们邀请主人赴宴时，并不认为自己是在给予，而是在接受；但这里的情况却相反：不是仆人邀请主人，而是主人首先邀请仆人赴祂的筵席；而你，难道连在这之后也不邀请祂吗？祂首先领你进入祂自己的屋檐下；你难道连其次也不接待祂吗？祂赤身裸体时给你穿上衣服；你难道连在这之后也不接待这位陌生人吗？祂首先用自己的杯给你喝；而你难道连冷水也不分给祂吗？祂使你饮于圣神，而你难道连祂肉身的干渴也不缓解吗？你本应受罚，祂却使你饮于圣神；而祂口渴时，你竟忽略祂，况且做这一切用的是祂自己的东西？难道你不认为，能拿着基督要喝的杯，并递到祂唇边，是一件大事吗？难道你没有看见，只有司铎才被允许递送祂血的杯？但祂说，我在这方面绝不严格；即使你自己递送，我也接受；即使你是平信徒，我也不拒绝。我不要求像我所赐予的那样；因为我寻求的不是血，而是冷水。想想你给谁喝，战栗吧。想想，你已成为基督的司铎，用自己的手给予，不是肉，而是饼；不是血，而是一杯冷水。祂给你穿上了救恩的衣服，并且亲自给你穿上；你至少要通过你的仆人来给祂穿上。祂使你在天上荣耀，你当使祂免于战栗、赤身和羞耻。祂使你成为天使的同乡，你至少要把你屋顶的遮盖分给祂，给祂如同给你自己的仆人一样的住处。\Quote{I refuse not this lodging and that, having opened to you the whole Heaven. I have delivered you from a most grievous prison; this I do not require again, nor do I say, deliver me; but if you would look upon me only, when I am bound, this suffices me for refreshment. When thou were dead, I raised you; I require not this again of you, but I say, visit me only when sick.}\par

既然祂的恩赐如此之大，祂的要求又极其容易，而我们连这些也不供应，那么我们该当受什么样的地狱深渊呢？我们理当进入那为魔鬼及其使者所预备的火里，因为我们比任何磐石还要麻木。请告诉我，我们这些领受了并将领受如此多恩惠的人，却做金钱的奴隶（我们不久将被迫与金钱分离），这是何等的麻木不仁？而别人甚至舍弃了自己的生命，流了自己的血；难道你连自己的多余之物也不肯为天国、为如此巨大的冠冕而舍弃吗？\par

那么你配得什么恩宠？什么成义？你在耕种土地时，慷慨地播撒一切；在向人放债取利时，毫不吝惜；但在通过穷人喂养你的主时，你却残忍、不近人情？\par

既然我们思考了这一切，计算了我们所领受的、将要领受的、以及对我们所要求的，就让我们在属灵的事上表现出我们全部的勤勉。让我们终于变得温和与仁慈，以免给自己招来无法忍受的惩罚。因为有什么不能定我们的罪呢？我们享受了如此众多而巨大的恩惠；我们对大事没有什么要求；我们所要求的，是我们将在此世不情愿地留下的；我们在世俗事务上表现出如此慷慨。这些事中的每一件，即使单独一件，也足以定我们的罪；但当它们聚集在一起时，还有什么救恩的希望呢？\par

因此，为了使我们能逃脱这一切定罪，让我们向那些需要的人表现一些慷慨。因为这样，我们将在此世和来世享受一切美善；愿我们都能达到这些，借着我们的主耶稣基督的恩宠和慈爱，愿光荣与权能归于祂，直到永远。阿们。\par

\SourceRecord{Translated by George Prevost and revised by M.B. Riddle.；Nicene and Post-Nicene Fathers, First Series；Vol. 10.；Edited by Philip Schaff.；Buffalo, NY: Christian Literature Publishing Co.,；1888.；<http://www.newadvent.org/fathers/200145.htm>.}

% structure: p0001 source=h1 target=section reason=page-title

\section{玛窦福音释义第四十六篇}

\Emph{\BibleRef{玛 13:24-30}}\par

\Emph{\Quote{祂又给他们设了一个比喻说：\InnerQuote{天国好比一个人在自己田里撒了好种子。\newline{}但在人睡觉的时候，他的仇人来，在麦子中间撒上莠子，就走了。\newline{}当苗长起来，结出果实的时候，莠子也显出来了。\newline{}家主的仆人就来对他说：主人！你不是在你田里撒了好种子吗？那么从哪里来的莠子呢？\newline{}他回答说：这是仇人做的。仆人对他说：那么你愿意我们去把莠子收集起来吗？\newline{}他却说：不，免得你们收集莠子时，连麦子也一起拔出来。\newline{}让它们一齐长到收割的时候吧。}}}\par

这个比喻和前面的比喻有什么区别呢？前面祂讲的是那些根本不持守祂，而偏离了、丢失了种子的人；但这里祂指的是异端者的团体。为了使祂的门徒连这事也不致不安，祂在教导他们为何用比喻说话之后，也预言了这事。前面的比喻意味着他们不接受祂；这个比喻则意味着他们接受了败坏者。因为这也是魔鬼诡计的一部分：总是在真理旁边引入错误，在其上涂上许多相似之处，以便容易欺骗易受骗的人。因此祂不称之为别的种子，而称之为\OriginalTerm{莠子}{tares}，外表上有点像麦子。\par

然后祂也提到了他诡计的方式。因为祂说：\Quote{当人睡觉时}。祂借此将不小的危险悬于我们的统治者头上，他们尤其是受托看守田地的人；不仅是统治者，也是臣民。\par

祂也表明错误是在真理之后出现的，实际的事件也证明了这一点。因为在先知之后，有假先知；在宗徒之后，有假宗徒；在基督之后，有假基督。因为魔鬼若看不见要模仿什么，或要谋害谁，他既不会尝试，也不知道如何做。现在他也看到\Quote{有一百倍的，有六十倍的，有三十倍的}之后，就另走别路。也就是说，他不能拔除已扎根的，也不能窒息或晒焦它，就藉另一种诡计阴谋对付它，偷偷放进自己的发明。\par

有人会说，睡觉的人与像路边的人之间有什么区别呢？在后一种情形中，他立刻夺去了种子，甚至不让它扎根；但在这里需要更多的诡计。\par

基督说这些话，教导我们要常常警醒。因为祂说，即使你完全免除了那些伤害，还有另一种伤害。因为就像那些例子中的\Quote{路边}、\Quote{石头地}和\Quote{荆棘}一样，在这里睡觉又造成了我们的覆灭；所以需要不断的警醒。因此祂也说：\BibleQuote{唯独坚持到底的，才可得救。}\BibleRef{玛 10:22}\par

类似的事情甚至在起初就发生了。我说的是，许多\OriginalTerm{主教}{prelates}把恶人，即乔装的\OriginalTerm{异端首领}{heresiarchs}带进教会，为安放那种网罗提供了极大的便利。因为魔鬼一旦把他们安插在我们中间，甚至不必费任何力气。\par

有人会说：\Quote{人怎能不睡觉呢？}的确，就自然的睡眠而言，是不可能的；但就道德官能的睡眠而言，却是可能的。因此保禄也说：\BibleQuote{你们要警醒，在信德上站稳。}\BibleRef{格前 16:13}\par

此后，祂指出这件事不仅是伤害性的，而且是多余的；因为，在土地已经耕种好、什么都不需要之后，这个仇敌才再撒种；正如异端者所做的，他们除了为了\OriginalTerm{虚荣}{vainglory}之外，并无其他原因而注入他们自己的毒液。\par

不仅凭这一点，而且凭随后的话，祂也准确地描绘了他们的全部行为。因为祂说：\Quote{当苗长起来，结出果实的时候，莠子也显出来了}，这些人所做的也正是这一类的事。起初他们伪装自己；但当他们获得了很大的信任，有人向他们传授圣言的教导时，他们就倒出他们的毒液。\par

那么为什么祂要让仆人来说明了所发生的事呢？是为了表明杀害他们是不对的。\par

祂称他为\Quote{仇人}，因为他加害于人。虽然这冒犯是针对我们，但其根源来自他对天主的敌意，而不是对我们的敌意。由此明显可见，天主爱我们胜过我们爱自己。\par

再从另一件事上看魔鬼的恶毒诡计。因为他之前没有撒种，因为他没有可毁坏的东西；但当一切都完成后，他才撒种，以挫败农夫的勤奋；他常常以这样的敌意对待祂。\par

也要注意仆人的情感。我是说，他们多么急于立刻拔除莠子，即使是轻率地去做；这表明他们对庄稼的关心，所关注的只有一件事：不是惩罚那个仇敌，而是保全所撒的种子。当然后者并不是紧迫的考虑。\par

因此，他们如何能除去当前的祸害，这是他们的目标。而且他们甚至不绝对追求这一点，因为他们不信任自己，而是等待主人的决定，说：\Quote{你愿意吗？}\par

那么主人如何做呢？祂禁止他们，说：\Quote{免得你们收集莠子时，连麦子也一起拔出来。}祂这样说，是为了阻止战争、流血和杀戮。因为处死一个异端者是不对的，否则会在世界上引起无情的战争。因此祂用这两个理由来约束他们：一是免得麦子受到伤害；二是如果他们患有不治之症，惩罚必定会临到他们。所以，如果你们要他们受惩罚，但又无害于麦子，我吩咐你们等待适当的季节。\par

但是\Quote{免得你们将麦子也一同拔出来}是什么意思？\newline{}或者祂的意思是说：如果你们拿起武器去杀害异端者，许多圣人也不得不跟他们一起被推翻；或者是指那些莠子中，很可能有许多会改变，变成麦子。因此，如果你们预先拔掉它们，就伤害了将要变成麦子的部分，杀死了那些还有改变和改善余地的人。所以祂并不禁止我们制止异端者、封住他们的口、剥夺他们的言论自由、解散他们的集会和同盟，但禁止我们杀害和屠戮他们。\par

但你当注意祂的温和：祂不仅判决并禁止，还列出了理由。\par

那么，如果莠子一直存留到末了呢？\newline{}\Quote{那时我要对收割的人说：你们先将莠子收集起来，捆成捆，好焚烧他们。}\BibleRef{玛 13:30}祂再次提醒他们若翰的话，\BibleRef{玛 3:12}介绍祂为审判者；并且祂说：只要他们站在麦子旁边，我们就该宽容他们，因为他们甚至有可能变成麦子；但当他们离开，毫无益处时，那无法逃避的严厉惩罚就必然临到他们。因为祂说：\Quote{我要对收割的人说：你们先将莠子收集起来。}为什么说\Quote{先}呢？\newline{}以免这些人惊慌，以为麦子也跟他们一起被带走。\Quote{捆成捆，好焚烧他们，但把麦子收入我的仓里。}\par

2. 祂又给他们设了一个比喻说：\Quote{天国好像一粒芥菜子。}\par

就是说，既然祂说过收成中三份损失，只有一份得救，而且在这得救的一份中又发生了这么大的损害；免得他们说：\Quote{那谁，又有多少将是信的人？}祂便又以芥菜子的比喻引导他们去相信，并指出无论如何福音要传遍四方。\par

因此祂摆出了这棵菜的比喻，它与当前的主题十分相似。祂说：\Quote{它原是百种中最小的，但长大后，却比一切菜蔬都大，竟成了一棵树，以致天上的飞鸟都来栖息在它的枝上。}\par

祂这样是要陈明它伟大最决定性的标志。祂说：\Quote{对于福音也是这样。}是的，因为祂的门徒曾是最软弱、最微小的，但由于在他们内的大能，它已在世界各地传扬开了。\par

此后，祂又在这个比喻上加上了酵，说：\par

\Quote{天国好像酵母，一个妇人取了，藏在三斗面里，直到全部发了酵。}\par

因为如同这酵把大量的面变成自己的性质，你们也必将转变整个世界。\par

请看祂的智慧：祂引用自然界的事物，表明正如前者不能不发生，后者也不能不实现。你们不要对我说：\Quote{我们十二个人，投身于这么多的人群中，能做什么呢？}不，正是这一点最彰显你们的大能：你们与群众混合而不被击退。因此，如同酵只有靠近面，并且不只是靠近，而是实际混合，才能发面——因为祂不是说\Quote{放}，而是说\Quote{藏}——同样，当你们亲近你们的敌人，并与他们成为一体时，你们就必将胜过他们。如同酵虽被掩埋，却不被消灭，而是逐渐将一切转变成本身的状况，这里关于福音的结果也将类似。因此，你们不要因为我说会有许多损害而害怕：因为即使如此，你们也必将闪耀，并胜过一切。\par

这里祂以\Quote{三斗}指许多，因为祂惯常用这个数字表示众多。\par

你们不要奇怪，祂在谈论天国时提到了小种子和酵；因为祂是在对没有经验、无知、需要藉这些方式引导的人讲论。因为他们是如此单纯，即使在这一切之后，他们仍需要不少解释。\par

外邦人的子孙如今在哪里？\newline{}让他们学习基督的大能，看见祂事迹的真实性，并且从两方面都来朝拜祂：因为祂预言了如此伟大的事，并且成就了它。是的，因为是祂把能力放进酵里。祂存此心意，将信祂的人与群众混合，好让我们把我们的智慧分施给其余的人。因此，不要有人因我们人数少而责备我们。因为福音的大能是伟大的，那已经发酵了的，又成为剩余部分的酵。如同一星之火落在木材上，使已燃的部分助长火焰，并如此继续烧向其余部分，福音也是这样。但祂没有说\Quote{火}，却说\Quote{酵}。这是为什么？\newline{}因为在火的情形中，整个效果不单是火产生的，部分也来自被点燃的木材；但这里，酵独自完成全部工作。\par

3. 如果十二个人使整个世界发酵，试想我们是如何地卑劣：我们这么多人，却不能改正那些余下的人；我们本应足以应付一万个世界，并成为它们的酵。但也许有人说：\Quote{但他们曾是宗徒。}那又怎样？\newline{}他们不是与你们同有份吗？他们不是在城里长大的吗？他们没有享受同样的恩惠吗？他们没有从事手艺吗？什么，他们是天使吗？什么，他们是从天上降下来的吗？\par

\Quote{但他们有标记}，人会这样说。\newline{}并不是标记使他们可钦佩。我们还要多久用那些奇迹作为我们自己懈怠的遮羞布？\Emph{请看圣人行列：他们并非因那许多奇迹而闪耀}。为什么？\newline{}许多实际赶走过魔鬼的人，因为他们行不义，不但没有受到敬佩，反而还受到惩罚。\par

他要说，那么是什么使他们显得伟大呢？就是轻视财富、藐视荣耀、摆脱世俗事物。因为如果他们缺乏这些品质，反而受自己情欲的奴役，就算他们使一万个死人复活，也非但没有任何益处，反而会被视为骗子。正是他们那处处闪耀光辉的生活，也吸引着圣神的恩宠降临。\par

若翰行了什么样的奇迹，竟使如此众多的城市关注他？至于他并未行过奇事，请听福音作者说：\Quote{若翰没有行过奇迹。}而厄里亚何以令人钦佩？难道不正是由于他在君王面前的勇敢？由于他对天主的热情？由于他甘愿的贫穷？由于他那件羊皮衣、他的山洞和他的山岭？他的奇迹是在这一切之后才行的。至于约伯，他在魔鬼眼前行了什么奇迹，使魔鬼对他感到惊异？其实没有奇迹，而是一种闪耀发光的生活，展现出比任何金刚石更坚定的忍耐。达味年轻时行了什么奇迹，竟使天主说：\BibleQuote{我找到了叶瑟的儿子达味，他是一个合我心意的人？}\BibleRef{宗 13:22}亚巴郎、依撒格和雅各伯，他们使谁复活？洁净了哪个癞病人？你们不知道吗？如果我们不保持清醒，奇迹在许多情况下甚至会造成伤害。比如，许多格林多人彼此分裂；许多罗马人变得骄傲；西满就被抛弃了。还有那个曾经想要跟随基督的人，当他对他说：\BibleQuote{狐狸有洞，天空的飞鸟有巢，}\BibleRef{玛 8:20}时，就被拒绝了。因为这些人中，一个贪图财富，另一个追求奇迹带来的荣耀，都堕落而灭亡了。然而，修行与对美德的热爱，不但不会产生这样的欲望，反而会消除它（如果存在的话）。\par

同样，主自己为门徒制定律法时，他说了什么？\Quote{行奇迹，好让人们看见你们？}绝对不是。而是说什么？\BibleQuote{你们的光当在人前照耀，使他们看见你们的善行，光荣你们在天之父。}\BibleRef{玛 5:16}他又对伯多禄说：\Quote{如果你爱我，}\Quote{行奇迹，}而是\BibleQuote{喂养我的羊群。}\BibleRef{若 21:16}而他在处处将伯多禄与雅各伯和若望一起高举于众人之上，请问，他为何高举他们？是因为他们的奇迹吗？绝非如此，因为大家都同样洁净癞病人、复活死人；他把同样的权柄赐给了所有人。\par

那么他们从何获得优势？来自他们灵魂中的美德。你们看到吗？处处都要求实践和行为的证明？\BibleQuote{凭他们的果实，你们可认出他们来，}\BibleRef{玛 7:16}他说。那么，什么使我们的生活值得称赞？是奇迹的展示，还是卓越言谈的圆满？显然，后者才是；而奇迹既源于此，也归结于此。因为一个展现卓越生活的人，会吸引这恩赐；而那接受恩赐的人，也是为了这个目的——修正他人的生活。因为基督自己也为此行了那些奇迹：藉此使自己成为可信的，吸引人们归向他，从而把美德带入我们的生活。因此，他更加重视后者。因为他绝不满足于仅仅行征兆，还以地狱威胁，应许天国，并颁布那些令人震惊的律法；他安排这一切，都是为了使我们能与天使同等。\par

为什么我说基督做这一切都是为了这个目标？因为，如果有人给你选择，是奉他的名复活死人，还是为他的名而死——请问，你更愿意接受哪一个？难道不是后者吗？然而前者是奇迹，后者只是行为。如果又有人让你选择，是把青草变成黄金，还是能够像对待青草一样轻视一切财富——你岂不更愿意接受后者？这是非常合理的。因为这会以更强大的方式吸引人类。如果他们看见青草变成了黄金，他们自己也会渴望获得那种能力（如同西满那样），从而使贪爱金钱在他们中间增长；但如果他们看见我们所有人都轻视、忽略黄金，如同对待青草一样，他们早就从这疾病中解脱了。\par

4. 你们看到吗？我们的行为更有效地行善。我所说的行为，不是指你们的禁食，也不是指你们铺上麻布和灰，而是指：你们轻视财富如同应当轻视的那样；你们心存仁爱；你们把面包给饥饿的人；你们克制愤怒；你们抛弃虚荣；你们除去嫉妒。主自己就是这样教导的：因为他说：\BibleQuote{你们跟我学习，因为我是良善心谦的。}\BibleRef{玛 11:29}他没有说：\Quote{因为我禁食了，}虽然他能提到那四十天，但他没有说这个，而是说：\Quote{我是良善心谦的。}还有，当他派遣门徒时，他没有说：\Quote{禁食，}而是说：\Quote{吃摆在你们面前的一切。}然而，关于财富，他要求他们非常严格，说：\BibleQuote{不要在你们的钱囊里预备金、银、铜钱。}\BibleRef{玛 10:9}\par

我讲这一切，并非要轻视禁食——绝无此意，反而要高度赞扬它。但我感到忧伤的是，当其他责任被忽略时，你们却以为只需禁食就够了，认为它在美德的歌咏团中只占最末的位置。因为最重要的乃是爱德、节制和施舍；这些甚至比童贞更崇高。\par

因此，如果你渴望与宗徒们平等，没有什么能阻止你。因为达到这种德行就足以使你丝毫不逊于他们。所以不要让任何人等待奇迹。因为恶神被赶出身体时会悲伤，但当它看到一个灵魂从罪恶中得救时，它更加悲伤。因为这真是它的大能。\BibleRef{宗 8:10} 这大能导致基督死亡，为要终结它。是的，因为这带来了死亡；由于这个，一切都被颠倒了。如果你除去这个，你就斩断了魔鬼的神经，你\Quote{打伤了它的头}，你终结了它所有的力量，你驱散了它的军队，你展示了一个比所有标记更大的标记。\par

这话不是我的，而是蒙福的保禄的。因为当他说：\Quote{你们应切求那最好的恩赐，但我还要指给你们一条更卓越的道路；} \BibleRef{格前 12:31} 他接下来没有说到标记，而是说爱德，我们一切好事的根源。如果我们实践这个，以及由此而来的所有克己，我们就不需要标记；反之，如果我们不实践它，我们通过标记也一无所获。\par

那么，记住这一切，让我们效法那些使宗徒们变得伟大的事。他们是如何变得伟大的？听伯多禄说：\Quote{看，我们舍弃了一切，跟随了你；那么，我们将得到什么呢？} \BibleRef{玛 19:27} 也听基督对他们说：\Quote{你们要坐在十二个宝座上，} 以及\Quote{凡舍弃了房屋、或兄弟、或父亲、或母亲的，在今世必要得到百倍，并承受永生。} 因此，让我们从一切世俗的事物中抽身，将自己奉献给基督，这样我们就能照祂的宣告与宗徒们平等，并享受永生；愿我们所有人都能达到，藉着我们的主耶稣基督的恩宠和对人的爱，愿光荣和权能归于祂，直到永远。阿们。\par

\SourceRecord{Translated by George Prevost and revised by M.B. Riddle.；Nicene and Post-Nicene Fathers, First Series；Vol. 10.；Edited by Philip Schaff.；Buffalo, NY: Christian Literature Publishing Co.,；1888.；<http://www.newadvent.org/fathers/200146.htm>.}

% structure: p0001 source=h1 target=section reason=page-title

\section{玛窦福音讲道 第四十七篇}

\BibleRef{玛 13:34-35}\par

\Quote{\Emph{耶稣用比喻对群众讲了这一切，不用比喻就不对他们讲什么}；\Emph{这是为应验那借先知所说的话：\InnerQuote{我要开口说比喻，要说出由创世以来隐藏的事。}}}\par

但是马尔谷说：\Quote{祂按照他们所能听懂的，用比喻给他们讲道。}\BibleRef{谷 4:33}\par

然后，指出祂不是在做什么新事，祂也引用了先知，预先宣告祂这种教导方式。为教导我们基督的目的，祂如何以这种方式讲话，不是要他们无知，而是要引导他们追问，祂接着说：\Quote{祂没有用比喻就不对他们讲什么。}然而，祂确实讲过许多没有比喻的话；但在此处却什么都没有。尽管如此，没有人问祂问题，而我们知道，先知们是常常追问的：例如\Person{厄则克耳}，还有许多其他人：但这些人却没有这样做。当然，祂的话足以使他们困惑，并激发他们追问；因为那些比喻确实威胁着非常严厉的惩罚；然而，即便如此，他们也不为所动。\par

因此，祂也离开他们走了。因为，\par

\Quote{然后，}祂说，\Quote{耶稣遣散了群众，回到自己的家里。}\par

没有一个经师跟随祂；由此可知他们跟随祂没有别的目的，只是为了抓住祂。但当他们不注意祂的话时，祂就任由他们去了。\par

门徒们来到祂跟前，问祂关于莠子的比喻；虽然有时他们想学，却害怕问（\BibleRef{谷 9:32}）。那么，这次他们的信心从何而来？他们曾被告知：\Quote{天国的奥秘是给你们知道的}；于是他们壮了胆。因此，他们私下问；不是出于嫉妒群众，而是遵守他们师傅的规范。因为，祂说：\Quote{这些人，是不给的。}\par

为什么他们放过了关于酵母和芥菜籽的比喻，而问这个呢？他们放过那些，因为更清楚；但对于这个，因为它与前面说的有关，而且展示了更多的东西，他们渴望学习（因为祂不会第二次对他们说同样的话）；因为他们确实看到其中发出的威胁是多么严厉。因此，祂也不责备他们，反而完善了祂之前的陈述。\par

正如我常说的，比喻不能逐一按字面解释，否则会引出许多荒谬；这一点，连祂自己也在解释这个比喻时教导了我们。因此，祂根本没有说那些来到祂面前的仆人是谁，但暗示祂把他们引入，是为了某种秩序和构成画面，祂省略了那一部分，只解释那些最迫切、最本质、也是比喻之所以说的部分；表明祂自己是所有人的审判者和主。\par

\Quote{祂回答说}，据说，\Quote{对他们说：撒好种子的就是人子；田地是世界；好种子是天国之子；莠子则是恶者之子；那撒莠子的仇敌是魔鬼；收获时期是世界的终结，收割者是天使。因此，如将莠子收集起来用火焚烧，在世界终结时也将如此。人子要派遣祂的天使，把一切使人跌倒的事和作恶的人收集起来，扔进火炉里，那里将有哀号和切齿。那时，义人在他们父的国度里，要像太阳一样发光。}\par

因为祂自己是撒种者，且是在自己的田里，并从自己的国度里收集，显然现今世界也是祂的。\par

但请注意祂对人不可言喻的爱，祂倾向于宽仁，不倾向惩罚；因为祂播种时亲自播种，但惩罚时却借他人之手，即天使。\par

\Quote{那时，义人在他们父的国度里，要像太阳一样发光。}不是因为只有那样，而是因为在我们所知中，没有星比这更明亮。祂使用我们熟悉的比较。\par

然而，祂在别处确实说，收获已经来到；正如祂论到撒玛黎雅人时说：\Quote{你们举目观看田地，庄稼已经发白，可以收割了。}\BibleRef{若 4:35}又说：\Quote{庄稼固多，工人却少。}那么，祂在那处说收获已经来到，而这里却说将来才有，是什么意思呢？依照另一种意义。\par

祂在别处说：\Quote{一人撒种，另一人收割}\BibleRef{若 4:37}，而在这里却说祂自己撒种，这是怎么回事呢？因为在那处，祂说话是为了区分宗徒们，不是区分祂自己，而是区分先知们，并且是在犹太人和撒玛黎雅人的情况下。因为当然也是祂通过先知们撒的种。\par

有时祂把同一件事既称为收获也称为撒种，根据不同的关联来称呼。因此，当祂谈到归化者的信服和服从时，祂称那事为\Quote{收获}，好像已完成了一切；但当祂寻求他们听道的果实时，祂称它为种子，而结束为收获。\par

祂在别处如何说：\Quote{义人先被提到云中？}\BibleRef{得前 4:17}因为他们确实先被提，但基督来后，其他人被交于惩罚，然后义人离开进入天国。因为他们必须在天上，但祂自己要来审判这里所有的人；对这些判了刑之后，祂如同一位君王与祂的朋友一同起来，带领他们进入那有福的境地。你看到惩罚是双重的：首先是烧毁，然后失去那荣耀吗？\par

2. 但为何在其他人退去后，祂仍继续向这些人讲比喻？他们因祂的话而变得更有智慧，以致能明白。无论如何，祂后来对他们说：\par

祂说：\BibleQuote{这一切你们都明白了吗？他们回答说：是的，主。}比喻连同其他目的，也完全实现了这一点，就是使他们更加明彻。那么，祂又说什么呢？\par

\BibleQuote{天国好像藏在田里的宝藏；人找到了，就把它藏起来，喜出望外，去变卖一切所有的，买了那块田。天国又好像一个商人寻找精美的珍珠；当他发现了一颗贵重的珍珠，就去变卖一切所有的，买了它。}\par

正如在别处，芥子和酵母彼此略有差异，同样，这里的两个比喻——宝藏和珍珠——也是如此。当然，两者都表明：我们应将福音置于万有之上。以前的比喻（关于酵母和芥子）是针对福音的能力以及它必定胜过世界而说的；但这些比喻则宣告它的价值和贵重。因为正如福音像芥子一样扩展，像酵母一样渗透，它也像珍珠一样宝贵，像宝藏一样带来丰盛。那么，我们不仅应当学习这一点：我们应当舍弃一切，紧紧抓住福音；而且也要明白：我们应当以喜乐这样做；当人放弃自己的财物时，他要知道这交易是收获，而非损失。\par

你看，福音如何藏在世界中，而福音中的美善也是如此？\par

除非你变卖一切，否则不能买；除非你拥有一个焦虑而求索的灵魂，否则不能找到。因此，两件事是必需的：远离世俗之事和警醒。因为祂说：\Quote{一个寻找精美珍珠的商人，当他发现了一颗贵重的珍珠，就变卖了一切，买了它。}因为真理是唯一的，而不是分散在许多碎片中。\par

正如拥有珍珠的人知道自己富有，但别人常常不知道他手中握着它（因为珍珠没有形体）；同样，对于福音，拥有它的人知道自己富有，而不信的人因不认识这宝藏，对我们的财富也一无所知。\par

3. 此后，为了叫我们不只是因听道而对福音有信心，也不以为单凭信德就足以得救，祂又讲了另一个令人敬畏的比喻。这是哪一个呢？就是网的比喻。\par

\BibleQuote{天国好像一个撒在海里的网，捕了各种鱼类。网满了，他们就拉上岸，坐下，拣好的收在器皿里，把坏的扔在外面。}\par

这比喻与莠子的比喻有何不同？因为那里也是一样的得救，一样的灭亡；但那里是因为选择了邪恶的道理；而此前的那些是因为不听祂的话，但这些是因为行为的邪恶；这些人是最悲惨的，他们得到了祂的知识，并被网住，却仍然不能得救。\par

然而，祂在别处说，牧者亲自将他们分开，但这里祂说天使这样做（\BibleRef{玛 25:32}）；对莠子的比喻也是如此。那么，这是怎么回事呢？祂有时用更适合他们迟钝的方式讲话，有时用更高深的语调。\par

这个比喻祂没有等人来问就解释了，而是主动解释了一部分，增加了他们的敬畏。因为怕你们听到\Quote{他们把坏的扔在外面}就以为那毁灭没有危险；祂藉着解释指明了惩罚，说：\BibleQuote{他们将把他们投进火炉。}（\BibleRef{玛 13:50}）祂又宣告了哀号和切齿，那痛苦是无法言喻的。\par

你看，灭亡的道路有多少？藉着石头，藉着荆棘，藉着路旁，藉着莠子，藉着网。因此，祂说：\BibleQuote{通往毁灭的门是宽敞的，走那条路的人很多。}（\BibleRef{玛 7:13}）并非没有理由。\par

4. 说完这一切，并以令人恐惧的语气结束祂的讲论，祂表明这些是大多数情况（因为祂更多地强调了它们）。祂说：\par

祂说：\BibleQuote{这一切你们都明白了吗？他们回答说：是的，主。}\par

于是，因为他们明白了，祂又称赞他们说：\par

\BibleQuote{因此，凡是受教于天国的经师，就好像一个家主，从他的宝库里拿出新旧的东西来。}（\BibleRef{玛 13:52}）\par

因此，祂在别处也说：\BibleQuote{我将派遣智者与经师到你们这里来。}（\BibleRef{玛 23:34}）你看，祂非但不排斥旧约，反而称赞它，公开为它说话，称它为\Quote{宝藏}？\par

因此，凡是不认识神圣圣经的人，就不能成为\Quote{家主}；这些人既不自己拥有，也不从别人接受，却忽略自己的处境，在饥荒中灭亡。不仅如此，异端者也与此福分无分。因为他们不能拿出新旧的东西。因为他们没有旧的，所以也没有新的；同样，没有新的，也没有旧的，两者都被剥夺了。因为这两者是相互联系、交织在一起的。\par

那么，我们这些忽略阅读圣经的人，让我们听听我们使自己遭受了怎样的损害，怎样的贫穷。因为当我们甚至不知道指导我们实践的律法时，我们何时才能致力于真正的德性实践呢？但是，那些富有、迷恋财富的人，不断地抖落衣服，以防虫蛀；而你们，看到遗忘比任何蛀虫更坏地侵蚀你们的灵魂，却忽略了与书籍交谈？你们难道不驱除这瘟疫，装饰你们的灵魂，持续注视德性的形象，并熟悉她的肢体和头部吗？因为她也拥有一个头部和肢体，比任何优雅美丽的身体更加得体。\par

\Quote{那么，有人会说，德行的首要是什麼？谦逊。因此基督也以此开始，说：\InnerQuote{神贫的人是有福的。}这首要并不卷发或垂鬈，却有得蒙天主喜悦的美丽。因为祂说：\InnerQuote{但我要垂顾谁呢？岂非顾念那谦卑、心神痛悔、对我言语战栗的人吗？}\BibleRef{依 66:2}又说：\InnerQuote{我的眼目眷顾地上的谦卑人。}以及：\InnerQuote{上主亲近心灵破碎的人。}这首要，代替鬈发与飘垂的头发，献上蒙天主悦纳的祭献。它是金祭台，是属神的祭坛；\InnerQuote{因为痛悔的精神是献给天主的祭献。}这是智慧之母。人若有此，也将拥有其余的一切。}\par

\Quote{你见过从未见过的头吗？你也要看见脸，或者更确切地说，注意它？那么请留意它目前的色泽：多么红润、鲜艳且迷人；并观察它的成分。\InnerQuote{那么，是哪些成分呢？}羞怯与脸红。因此也有人说过：\InnerQuote{羞怯的人面前，恩宠先行。}\BibleRef{德 32:10}这羞怯也为其他肢体增添许多美丽。纵然你调和万种颜色，也无法产生这样的红润。}\par

\Quote{你若也愿意看眼睛，请看它们如何以端庄和节制精确地勾勒。因此它们也变得如此美丽和锐利，以至能看见主自己。因为祂说：\InnerQuote{心里洁净的人是有福的，因为他们要看见天主。}\BibleRef{玛 5:8}}\par

\Quote{她的口是智慧与明智，以及对属神诗歌的认识。她的心是圣经的熟识、正统道理的持守、仁爱与良善。因为没有后者便无法生存，同样，没有前者也绝无救恩。是的，因为她的一切卓越皆由此而生。她也有足与手，即其善工的彰显。她还有一个灵魂，即虔诚。她还有一个金子的胸怀，比金刚石更坚固，即刚毅；万物皆可被俘虏，却比撕裂那胸怀更容易。而在头脑和心中的精神，乃是爱德。}\par

\Quote{5. 你愿意我以她实际的作为，也向你展示她的形象吗？请思考，我求你，这位圣史本人：虽然我们没有将他的一生全写下来，但即使从几件事实，也能看见他的形象闪耀。}\par

\Quote{首先，关于他的谦逊与痛悔，请听他在其福音后，自称为税吏；关于他的仁慈，看他舍弃一切跟随耶稣；关于他的虔诚，从其教义显而易见。他的智慧也容易从他撰写的福音中看出，以及他的爱德（因为他关心整个世界）；他的善工的彰显，从他将要坐的宝座；\BibleRef{路 22:30}还有他的勇气，\InnerQuote{由于他欢喜地离开公议会。}\BibleRef{宗 5:41}}\par

\Quote{让我们效法这德行，尤其是他的谦逊与施舍，没有这些无人能得救。这由五个贞女以及法利塞人一同表明。因为没有贞洁，固然可能见到天国，但没有施舍却不可能。因为这属于本质之事，维系一切。所以我们并非不恰当地称它为德行之核心。但这核心，若不向一切输送气息，便很快熄灭。同样，泉源若将其水流局限在自己内，就会腐败；富人也如此，当他们把财产据为己有时。因此在我们日常谈话中，我们也说：\InnerQuote{此人财富消耗巨大}；而不说：\InnerQuote{财富丰盛，宝藏巨大。}因为确实有一种消耗，不仅是拥有者的消耗，也是财富本身的消耗。因为存放的衣服会被虫蛀，金子会生锈，谷物会被虫吃，而主人的灵魂比这一切更加被忧虑锈蚀和败坏。}\par

\Quote{若你愿意，把一个吝啬者的灵魂放在面前；它就像一件被无数虫蛀的衣服，没有任何完好的部分，你会发现它如此——被忧虑从四面八方穿透；被罪恶腐蚀、锈坏。}\par

\Quote{但穷人的灵魂并非如此——我指的是自愿贫穷的人的灵魂——它如金子般灿烂，像珍珠一样闪耀，如玫瑰般绽放。因为没有蛾子，没有贼，没有世俗的忧虑，而是像天使一样交往。}\par

\Quote{你愿看见这灵魂的美吗？你愿了解贫穷的富足吗？他不命令人，却命令邪魔。他不站在君王身旁，却站在天主近旁。他不是人的同伴，而是天使的同伴。他没有两三个或二十个箱子，却有如此丰盛，以至视整个世界为无物。他没有宝藏，却有天堂。他不需要奴隶，更确切地说，他以情欲为奴隶，以统治君王的动机为奴隶。因为那命令穿紫袍之人的动机，在他面前退缩。王权、黄金以及一切此类事物，他看作孩童的玩具；像铁环、骰子、陀螺和球一样，他视这一切为可鄙。因为他有一种装饰，是玩这些玩具的人所不能见的。}\par

\Quote{那么，有什么能超过这穷人呢？他至少有天空为地板；若地板如此，请想像屋顶！他没有马匹和战车吗？他将被云彩承载，与基督同在，还需要这些做什么？}\par

\Quote{将这些印在我们心中，让我们——无论男女——追求那财富和不能掠取的丰盈，好使我们也能到达天国，藉我们的主耶稣基督的恩宠与对人的爱，愿光荣与权能归于祂，至于无穷之世。阿们。}\par

\SourceRecord{Translated by George Prevost and revised by M.B. Riddle.；Nicene and Post-Nicene Fathers, First Series；Vol. 10.；Edited by Philip Schaff.；Buffalo, NY: Christian Literature Publishing Co.,；1888.；<http://www.newadvent.org/fathers/200147.htm>.}

% structure: p0001 source=h1 target=section reason=page-title

\section{\WorkTitle{玛窦福音讲道集 48}}

\BibleRef{玛 13:53}\par

\Quote{耶稣讲完了这些比喻，就从那里走了。}\par

祂为什么说\Quote{这些}呢？因为祂还要讲别的比喻。再者，祂为什么离开？因为祂愿意将圣言播撒各处。\par

\Quote{祂来到自己的家乡，在他们的会堂里教训他们。} \BibleRef{玛 13:54}\par

祂如今称什么地方为家乡？在我看来，是\Place{纳匝肋}。\Quote{祂在那里没有行许多奇事，} \BibleRef{玛 13:58} 经上说，但在\Place{葛法翁}祂行了奇迹；因此祂也说：\Quote{至于你，\Place{葛法翁}，你被高举到天上，必被投到地狱里；因为若是在\Place{索多玛}所行的异能，行在你们那里，它必存留到今天。} \BibleRef{玛 11:23}\par

祂到了那里，在奇迹方面略有收敛，以免激起他们更大的嫉妒，也免得因他们不信的加重而更严厉地定他们的罪；然而祂仍提出一种教训，其奇妙不亚于奇迹。因为那些完全麻木的人，本应惊叹并惊讶于祂言语的权能，反而因那似乎是祂父亲的人而轻视祂；但我们知道，他们在古时有许多这样的例子，从卑微的父亲中见到杰出的儿子。例如，\Person{达味}是卑微农夫\Person{耶瑟}的儿子；\Person{亚毛斯}是牧羊人的孩子，且自己也是牧羊人；\BibleRef{亚 7:14} 而立法者\Person{梅瑟}也有一个远不如他的父亲。因此，他们本应为此而崇拜和惊讶，即祂出于这样的父母却说出如此的话，这清楚表明并非出于人的关怀，而是出于天主的恩宠；然而他们却因那些本应钦佩祂的事而轻视祂。\par

再者，祂不断出入会堂，免得若总是留在荒野，他们会更指责祂制造分裂，对抗他们的政体。因此他们惊讶而困惑地说：\Quote{这人从哪里得了这样的智慧和异能？} 或许称奇迹为异能，甚或智慧本身。\Quote{这不是木匠的儿子吗？} \BibleRef{玛 13:55} \Quote{祂的母亲不是叫\Person{玛利亚}吗？祂的弟兄不是\Person{雅各伯}、\Person{若瑟}、\Person{西满}和\Person{犹达}吗？祂的姊妹不是都在我们这里吗？那么这人从哪里得了这一切？} 他们就因祂而跌倒。\BibleRef{玛 13:56-57}\par

你看见祂讲道的地方是\Place{纳匝肋}吗？经上说：\Quote{祂的弟兄不是某某某吗？} 这又如何？正因如此，你们更应被引向信德。但嫉妒，你看到，是一种卑贱可鄙的情感，常常自相矛盾。因为那些令他们感到陌生、奇妙、足以使他们归向的事，反倒使他们反感。\par

那么\Person{基督}对他们说什么？祂说：\Quote{先知除了在自己的家乡和本家之外，没有不受尊敬的。} 经上又说：\Quote{祂没有多行异能，因为他们的不信。} 但\Person{路加}说：\Quote{祂在那里没有行许多异能。} 按理祂应该行异能。因为若对祂的惊奇在增长（其实连那里的人也惊讶祂），祂为何不行异能？因为祂不看自己的显扬，而看他们的益处。因此当这事没有成功时，祂就忽略了自己的事，以免加重他们的惩罚。\par

且看过了多久祂才来到他们那里，又行了多少奇迹；但即使这样他们仍不能忍受，反而再次被嫉妒点燃。\par

那么祂为何仍行几个奇迹？免得他们说：\Quote{医生，医治你自己吧！} \BibleRef{路 4:23} 免得他们说：\Quote{祂是我们的仇敌，忽视自己的家乡；} 免得他们说：\Quote{若行了奇迹，我们也会相信。} 因此祂既行了奇迹，又停了下来：行奇迹是为完成祂的部分，停下来是为不再定他们的罪。\par

且看你当思想祂话语的能力，至少在此：他们虽被嫉妒控制，仍不能不钦佩。正如对于祂的作为，他们不指责所行的，却捏造不存在的理由，说：\Quote{祂是靠着\OriginalTerm{贝耳则步}{Beelzebub}驱魔；} 同样，在这里他们也不指责教训，却躲进祂出身的低微中。\par

但请你留意师傅的温和，祂如何不责骂他们，而是以极大的温良说：\Quote{先知除了在自己的家乡，没有不受尊敬的。} 祂并未止于此，还加上：\Quote{并在自己的本家。} 在我看来，祂暗中暗指祂自己的弟兄而加了这句话。\par

但在\Person{路加}那里，祂还举了例子说：\Person{厄里亚}并没有到自己的家乡，而是到了外邦寡妇那里；\Person{厄里叟}也没有医治其他癞病人，而是外邦人\Person{纳阿曼}。\BibleRef{路 4:25-27} 以色列人既未受益，也未施益，而是外邦人得利。祂说这些话，处处表明他们心术不正，并且在祂身上并没有发生什么新事。\par

2. \Quote{那时，分封侯\Person{黑落德}听见\Person{耶稣}的名声，} 因为杀婴孩的\Person{黑落德}王，这人的父亲，已经死了。\par

但福音作者并非无目的地注明时间，而是要你注意暴君的傲慢和他的轻率，因为他不是在起初就打听\Person{基督}，而是在很长一段时间之后。因为掌权者、被众多排场包围的人，往往很晚才得知这些事，因为他们不看重这些。\par

但请你留意德行是多么伟大，即使\Person{黑落德}在\Person{基督}死后仍害怕祂，并且因恐惧而对复活之事说出明智的话。\par

\Quote{他说：}经上记着，\BibleQuote{他对自己的仆人说：这是我斩首的\Person{若翰}，他从死者中复活了，故此这些异能在他身上运行。} \BibleRef{玛 13:2} 你看见他恐惧的程度吗？因为那时他也不敢公开宣扬，仍只是对仆人们说。\par

然而，这个说法仍带有军人的粗鲁，且是荒谬的。因为许多别的人也曾从死者中复活，却没有行过这类奇事。他的话在我看来，既是虚荣，又是恐惧之言。因为无理性的灵魂往往同时容纳对立的激情。\par

但\Person{路加}肯定群众说：\Quote{这是\Person{厄里亚}，或是\Person{耶肋米亚}，或是古先知中的一位。}而他却自以为比别人高明，作出了这样的断言。\par

但可能在此之前，当有人对他说他是\Person{若翰}（因为许多人曾这样说），他否认了，并说：\Quote{我杀了他，}以此自夸和骄傲。因为\Person{玛尔谷}和\Person{路加}都记载他说：\Quote{我斩了\Person{若翰}的首级。}但当传闻占了上风时，他也随从了民众的说法。\par

然后，福音作者给我们讲述了这段历史。他为什么不把它单独作为主题叙述呢？因为他们的全部努力都只是讲述关于基督的事，除了有助于同一目的外，他们没有为自己安排次要的工作。因此，若非为了基督的缘故，并因\Person{黑落德}说了\Quote{\Person{若翰}复活了}，他们现在也不会提及这段历史。\par

但\Person{玛尔谷}说，\Person{黑落德}极其尊重那人，即使在受责斥时也是如此。\BibleRef{谷 6:20} 德行是一件多么伟大的事啊！\par

然后他的叙述这样继续：\BibleQuote{原来\Person{黑落德}曾派人拿住\Person{若翰}，把他捆绑在监狱里，是为了他的兄弟\Person{斐理伯}的妻子\Person{黑落狄雅}的原故。因为\Person{若翰}对他说：你不可占有她。\Person{黑落德}本想杀他，但害怕群众，因为他们都认为他是一个先知。} \BibleRef{玛 13:3}\par

为什么他完全不向她发言，而只向那男人说话呢？因为责任更在于他。\par

但请看他是如何无犯意地提出控告，像是叙述历史而非提出指控。\par

4. \Quote{但当\Person{黑落德}的生日到了，}他说，\BibleQuote{\Person{黑落狄雅}的女儿在众人面前跳舞，使\Person{黑落德}欢喜。} \BibleRef{玛 13:6} 哦，魔鬼的狂欢！哦，撒殚的表演！哦，无法无天的舞蹈！更无法无天的是对舞蹈的奖赏。因为一桩比所有谋杀更不虔敬的谋杀被执行了；那本应受冠冕和公众尊荣的人，在筵席中被杀害，魔鬼的战利品被摆上了餐桌。\par

而胜利的方式也与所行的事相称。因为，\par

\Quote{\Person{黑落狄雅}的女儿，}经上说，\BibleQuote{在中间跳舞，使\Person{黑落德}欢喜。于是他以誓言许诺，无论她求什么，都给她。她受了母亲的指使，就说：请把\Person{洗者若翰}的头放在盘子里给我。} \BibleRef{玛 13:6}\par

她的耻辱是双重的：第一，她跳舞了；第二，她取悦了他，并且是如此取悦，以至于甚至以谋杀作为她的酬报。\par

你看见他是多么野蛮、无感觉和愚蠢吗？他让自己受誓言的约束，却给了她提出要求的全权。但当邪恶实际发生时，经上说\Quote{他忧愁了}；然而起初他竟把他关押起来。那么他为何忧愁呢？这就是德行的本性：即使在恶人之中，它也应得钦佩和赞美。但唉，她的疯狂！她本应钦佩，是的，甚至向他屈膝，因为他试图纠正她的不义，她却反而协助策划阴谋，设下陷阱，并要求一个恶魔般的恩惠。\par

他因\Quote{誓言的缘故}，经上说，\Quote{并因同席的人}而害怕。你怎么不害怕那更严重的事呢？如果你害怕有见证人看到你背誓，那么你更应该害怕有这么多见证人看到如此无法无天的谋杀。\par

然而，我想许多人并不了解这不平之事，即谋杀源于何处，我必须也说明这一点，以便你们了解立法者的智慧。那么，那古老的法律是什么，\Person{黑落德}践踏它，而\Person{若翰}维护它？那无子女而死者之妻应归于其兄弟。\BibleRef{申 25:5} 因为死亡是无法医治的恶，而一切都为生命起见；他制定法律，让活着的兄弟娶她，并将所生的孩子以死者的名字命名，以免他的家完全灭绝。因为如果死者连留下孩子——这最大的死亡缓解——都不能，悲伤将无法弥补。因此，你们看到，立法者为那些天生无子女的人设计了这种安慰，并命令所生的子女归属于死者。\par

但当有子女时，这种婚姻就不再被允许。\Quote{为什么？}有人会说，\Quote{如果对另一个人是合法的，对兄弟就更应合法。}绝不是。因为他愿意扩大人们的亲属关系，并增加我们彼此关注的源泉。\par

那么，在没有子嗣的死亡情况下，为什么另一个男人不娶她呢？因为孩子就不会被认为是死者的；但现在由兄弟生育，这种假定就变得合理了。此外，任何其他男人都没有义务去建造死者的家，但兄弟由于亲属关系而负有这一责任。\par

因此，\Person{黑落德}娶了他兄弟的妻子，而她又已有了孩子，所以\Person{若翰}责备他，并且责备得很有节制，在显示他的勇敢的同时，也显示了他的谨慎。\par

但请你注意，我恳求你，整个剧场是多么魔鬼般的。第一，它由醉酒和奢华组成，从中不能产生任何健康的东西。第二，其中的观众是堕落的，而设宴者是所有罪人中最坏的。第三，有非理性的快感。第四，那少女——因她而婚姻是非法的，她本应甚至隐藏自己，因为她的母亲被她羞辱了——却出来表演，使所有娼妓都黯然失色，尽管她是个处女。\par

时间因素也对此滔天罪行的不无谴责。因为正当他应当感谢天主在那一天使他来到世间时，他却胆敢做出那些违法之举。当本该由他释放被囚之人时，他却在束缚之上又加了杀戮。\par

请听，你们这些贞女，或者不如说你们这些已婚妇女，凡是在他人的婚礼上同意这种不体面之事——跳跃、蹦跳、羞辱我们共同本性的人，都请听。请听，你们这些男人，凡追随那些充满耗费与酗酒的宴席的人，都当畏惧那恶者的深渊。因当时那恶人如此牢牢地抓住了那个可怜人，以至于他起誓甚至要给她半个王国：这是马尔谷的记述：\BibleQuote{无论你向我求什么，我必给你，甚至我王国的一半。}\BibleRef{谷 6:23}\par

他对自己王权的估价竟是如此；他既已为情欲所俘，以至于愿为一个舞蹈而舍弃\Emph{他的王国}。他贬抑、辱骂、侮辱。但圣人们并非如此；他们反倒为犯罪者哀哭，而非咒骂他们。\par

8. 因此，让我们也这样行，并为黑落狄雅和那些效法她的人哀哭。因为如今也常有这样的狂欢，虽然若翰未被杀害，但基督的肢体却被杀害，且以更痛苦的方式。因当代的舞者所求的，不是盘子里的头，而是赴宴者的灵魂。因为他们使这些人成为奴隶，引他们走向非法的爱情，用妓女包围他们，他们并非割下头颅，而是杀害灵魂，使他们成为奸夫、柔弱者与淫乱之徒。\par

你们总不能告诉我，当你们酒足饭饱、酩酊大醉、看着一个跳舞并口出秽言的女人时，你们对她毫无感觉，也没有被情欲所胜而奔向放荡？不，那种可怕的事发生在你们身上：你们 \BibleQuote{将基督的肢体作为娼妓的肢体}\BibleRef{格前 6:15}\par

因为即使黑落狄雅的女儿不在场，那当初在她身上跳舞的魔鬼，如今也在那些人身上举行它的歌咏队，并带着被掳的赴宴者的灵魂离去。\par

但即使你们能够避免酗酒，你们却参与了另一种最严重的罪恶；因为这样的狂欢也充满了许多掠夺。我请你们不要看摆在他们面前的肉食和糕点；而要思考这些是从哪里来的，你们就会看到它们来自烦恼、贪婪、暴力和掠夺。\par

\Quote{不，我们的宴席并非来自这样的来源，}有人会说。但愿不如此：因为我也不愿这样。然而，即使它们没有这些，我们奢华的宴席也并非无可指责。无论如何，请听，即使抛开这些，先知如何指责他们，这样说道：\Quote{祸哉，那些饮过滤之酒、用上等油膏抹自己的人。}你们看见他如何谴责奢侈吗？因为他在这里所责备的不是贪婪，而仅仅是挥霍。\par

而你饮食过度，基督甚至连必需之物也没有；你吃各种糕点，他连干面包也没有；你喝塔索斯葡萄酒，而他在口渴时，你连一杯冷水也没有给他。你躺在柔软绣花的床上，而他却在寒冷中濒死。\par

因此，即使宴席不涉及贪婪，它们仍是被诅咒的，因为当你自己事事过度时，你连他的必需也没有给他；而你却是靠他的东西过着奢侈的生活。假如你是一个孩子的监护人，占有他的财物，却置他于绝境而不顾，你会遭到千万人的控告，并受到法律规定的惩罚；而现在你占有基督的财物，这样无谓地消耗它们，你以为自己将不必交账吗？\par

9. 我说这些话，并不是针对那些把妓女带到宴席上的人（因为我对他们无话可说，如同我对狗也无话可说），也不是针对那些欺骗一些人、纵容另一些人的人（因为我也与他们无关，如同我与猪和狼无关）；而是针对那些确实享用自己财产、却不肯分施给他人的人；那些随意挥霍祖产的人。因为这些人也并非无可指摘。因为，请告诉我，你若不受责备和指责，怎能逃脱呢？当你的食客被养肥、你身旁的狗被喂养，而基督的价值在你眼中甚至不如它们时？当那食客因说了一句俏皮话而满载而归，而这一位教给我们——倘若我们没有学到，我们便与狗无异——却连那食客所得的最小一部分都不配得时？\par

听到这样说，你不战栗吗？那么，就在现实中战栗吧。赶走那些食客，让基督与你同席。如果他尝了你的盐，分享你的餐桌，他在审判你时就会温和；他知道如何尊重人的餐桌。是的，如果强盗知道这一点，主岂不更知道？例如，想想那个妓女，他在餐桌上如何称她为义，并指责西满说：\BibleQuote{你没有给我亲吻。}\BibleRef{路 7:54} 我说，如果他喂养你，而你不做这些事，那么你若做这些事，他会更赏报你。不要看那穷人来到你面前时污秽不堪，而要想到基督借着他正进入你的家，并停止你的粗暴和无情的言语——你甚至用这些言语辱骂到你面前的人，称他们为骗子、懒汉和其他更恶劣的名字。\par

并且，当你这样说话时，想想那些食客：他们成就了什么工作？他们在哪方面有益于你的家？他们真的使你的晚餐愉快吗？愉快——因他们挨打和说污言秽语？不，有什么比这更不愉快呢？当你击打那按天主的肖像所造的人，并从对他的傲慢中为自己寻找乐趣，把你的家变成剧场，用剃了头的戏子充满你的宴席——你这个出身高贵、自由的人，却模仿那些剃了头的戏子？因为在他们中间，也有欢笑和粗鲁的殴打。\par

这些事，我请问你，你竟称之为快乐吗？它们实在该当许多眼泪、许多哀恸和悲叹！当适于劝勉他们度善生、及时给予忠告时，你竟引导他们去发假誓、说无秩序的话，并称那事为享乐？那导致地狱的事，你竟视为快乐的题目？而且，当他们想不出机智的话时，便以发誓和假誓来付清全部代价。这些事难道该当欢笑，而不该当哀恸和流泪吗？谁会有理智这样说呢？\par

我说这话，并非禁止供养他们，而是不要为了这样的目的。相反，让他们的生计出于仁慈，而非残忍；出于同情，而非傲慢。因为他是穷人，就供养他；因为基督受到供养，就供养他；而不是为了引人说出撒殚的话，耻辱自己的生活。不要看他外表欢笑，而要审视他的良心，那时你便会看见他对自己发出千万咒骂、呻吟和哀号。即使他没有表露出来，这也是你的过错。\par

10. 那么，让你餐桌的同伴是贫穷而自由的人，不是发假誓者，也不是戏子。如果你必须向他们索取饮食的回报，就命令他们，若见到任何不当之事，要斥责、劝诫，并辅助你照顾家庭、管理仆人。你有孩子吗？让这些人成为他们共同的父亲，让他们与你分担责任，让他们为你产生天主所爱的利益。使他们从事属神的交易。若你看到有人需要保护，就命令他们援助、服务。藉着这些人寻找陌生人，藉着这些人给赤身者穿衣，藉着这些人送食物到监狱，终结他人的痛苦。\par

让他们为他们的食物给你这样的回报，这既有益于你也有益于他们，且不带有任何定罪。\par

藉此，友谊也更为牢固。因为现在，虽然他们似乎被人爱，但毕竟他们感到羞愧，如同在你的家中毫无目标地生活；但如果他们完成这些目的，他们自己就更加愉快，你在供养他们时也更加满足，因为没有白费钱财；而他们也将以勇敢和应有的自由与你同住，你的家将不再是戏院，而成为你的教会，魔鬼将逃跑，基督将进入，天使的合唱也将来到。因为哪里有基督，哪里也有天使；哪里有基督和天使，那里就是天堂，那里有比这太阳更明亮的光。\par

如果你还愿藉着他们得到另一种安慰，就命令他们，在你闲暇时，拿起他们的书阅读天主的法律。他们这样服侍你，会比另一种方式更令他们快乐。因为这些事为你和他们都增添尊重，但那些事却为所有人带来耻辱；为你，如同一个傲慢的人和一个酒徒；为他们，如同可怜和贪食的人。因为如果你供养他们是为了侮辱他们，那就比把他们处死更糟；但如果是为他们的好处和利益，那就比把他们从行刑路上带回来更有用。现在你确实辱没了他们，甚至超过你的仆人；而你的仆人所享有的言论自由和良心自由，比他们更多；但那时你将使他们如同天使一样。\par

因此，释放他们和你自己，除去\Quote{寄生虫}的称号，称他们为你的餐桌同伴；抛弃\Quote{谄媚者}的名称，赐予他们\Quote{朋友}的称呼。天主确实正是为此目的创造了我们的友谊，不是为了对爱者和被爱者造成伤害，而是为了他们的好处和利益。\par

但这些友谊比任何敌意更令人痛苦。因为藉着我们的敌人，如果我们愿意，我们甚至能获益；但藉着这些朋友，我们必然受害，这是毫无疑问的。所以不要结交教导你害处的朋友；不要结交那些迷恋你的餐桌胜过你的友谊的朋友。因为所有这样的人，如果你削减你的好生活，他们也会削减他们的友谊；但那为了美德而与你交往的人，却时常存留，忍受一切变化。\par

此外，那一班\Quote{寄生虫}常常会报复你，使你声名狼藉。因此，我认识许多可敬的人获得了坏名声，有些人因行邪术而被恶名传播，有些人因奸淫和败坏少年而受责难。因为他们无事可做，白白浪费自己的生命；他们的服务被众人怀疑为与败坏少年的行为相同。\par

因此，我们既要从恶名，更要从将来的地狱中拯救自己，并做那中悦天主的事，让我们终止这魔鬼的习俗，好使\BibleQuote{我们或吃或喝，都要为光荣天主而做一切事}\BibleRef{格前 10:31}，并享受那来自祂的光荣；愿我们众人藉着我们的主耶稣基督对人的恩宠和慈爱，都能达到这光荣，愿光荣和权能归于祂，从今世直到永远。阿们。\par

\SourceRecord{Translated by George Prevost and revised by M.B. Riddle.；Nicene and Post-Nicene Fathers, First Series；Vol. 10.；Edited by Philip Schaff.；Buffalo, NY: Christian Literature Publishing Co.,；1888.；<http://www.newadvent.org/fathers/200148.htm>.}

% structure: p0001 source=h1 target=section reason=page-title

\section{\WorkTitle{玛窦福音讲道集 第49讲}}

\Emph{\BibleRef{玛 14:13}}\par

\Emph{\Quote{但耶稣一听说，便从那里乘船独自退到荒野地方；群众听说后，就从各城步行跟随祂。}}\par

请看见祂在一切场合\Quote{退去}。因为不仅藉祂的外表，更藉祂的行动，祂要确认这一点，因祂知道魔鬼的诡计，知道牠会不遗余力地摧毁这个道理。\par

祂为此而退去，但群众即使如此也不离开祂，反而紧紧跟随祂，连若翰的悲惨结局也未能使他们惊惧。热忱是何等伟大，爱是何等伟大；它如此克服并驱散一切危险。\par

因此他们也立刻获得了赏报。因为经上说，耶稣\Quote{出来，看见一大群人，就对他们动了怜悯的心，治好了他们的病人。} \BibleRef{玛 14:14}\par

因为无论他们的勤勉有多大，然而祂的作为却超过任何勤勉所能赚取的。因此祂也表明祂如此医治他们的动机——祂的怜悯，深切的怜悯；并且祂治愈了所有人。\par

而且祂在此并不要求信德。因为他们既来到祂跟前，又离开自己的城市，且殷勤寻求祂，并在饥饿紧迫时仍与祂同在，这些都显明了他们的信德。\par

但祂也要喂养他们。祂并非凭自己这样做，而是等待被恳求；在每一种场合，正如我所说，祂保持这个原则：不抢先施行奇迹，不预先阻止他们，而是出于某种请求。\par

为什么群众中没有人上前为他们说话？他们极其尊敬祂，甚至因渴望与祂同在而感觉不到饥饿。连祂的门徒来到祂跟前时，也没有说\Quote{喂养他们}，因为那时他们还处于不完善的状态；但他们说了什么呢？\par

\Quote{到了傍晚，}经上说，\Quote{祂的门徒来到祂跟前说：\InnerQuote{这地方是荒野，时候已经过了；请遣散群众，叫他们自己去买食物。}} \BibleRef{玛 14:15}\par

因为即使在奇迹之后，他们忘记了所发生的事，在篮子之后，还以为祂在谈论饼，当祂称法利塞人的教训为\Quote{酵母}时；\BibleRef{玛 16:6} 更何况当他们还从未经历过这样的奇迹时，他们更不会期待这样的事。然而祂已经开始实际治愈了许多病人；但即便如此，他们也未从这件事预料到饼的奇迹；那时他们是多么不完善。\par

但是，我请求你注意，这位导师的技巧，祂如何清晰地将他们召唤到信仰中。因为祂没有立刻说\Quote{我喂养他们}——这确实不会轻易被接受——但祂说了什么呢？\par

\Quote{但是耶稣}，如经上所载，对他们说：\Quote{他们不必离开，你们给他们吃的吧。} \BibleRef{玛 14:16}\par

祂没有说\Quote{我给他们}，而是说\Quote{你们给他们}；因为那时他们仍以看待一个人的眼光看待祂。但他们即使这样也未被唤醒，仍然像与一个人推理一样，说：\par

\Quote{我们这里只有五个饼和两条鱼。} \BibleRef{玛 14:17}\par

因此马尔谷也说：\Quote{他们不明白这话，因为他们的心是迟钝的。}\par

因此他们仍然在地上爬行，然后祂终于介入自己的角色，说：\Quote{把饼拿到我这里来。} 因为虽然地方是荒野，但喂养世界的那位在这里；虽然时间已经过了，但不属于时间的那位正在与你们说话。\par

但是若望也说，那是\Quote{大麦饼}，\BibleRef{若 6:9} 祂并非无目的地提及此事，而是教导我们践踏奢侈生活的骄傲。先知们的饮食也是如此。\par

2. \Quote{于是耶稣拿起那五个饼和两条鱼，吩咐群众，}经上说，\Quote{坐在草地上，举目望天，祝福了，掰开，递给门徒，门徒再递给群众。众人都吃了，吃饱了，把剩下的碎块收集起来，装满了十二篮子。吃的人除妇女孩子外，约有五千。}\par

为什么祂举目望天并祝福呢？这是为了让人相信祂既出于圣父，又与圣父同等。但这些事的证明似乎彼此矛盾。因为祂的同等由祂凭权能行一切事表明，而祂出自圣父的起源，除非祂以极大的谦卑行事，并指向圣父，在祂的工作上呼求圣父，否则他们无法被说服。因此我们看到祂并非只做这些事或那些事，而是两者都做，使两者都得着确认；祂有时凭权能行奇迹，有时藉祈祷行奇迹。\par

此外，为了使祂的作为不显得矛盾，在较小的事上祂举目望天，但在更大的事上祂凭权能行事；为要教导你在小事上也是如此，祂这样做不是从别处接受能力，而是为了尊敬那生祂的。例如：当祂赦免罪过、开启乐园、引入盗贼、彻底废掉旧律、复活无数死人、平息大海、斥责人未说出的思想、创造了一只眼睛——这些唯独天主才能成就的事——我们在任何事例中都未见到祂祈祷；但当祂使饼自行增多时，这是远小于这一切的事，祂却举目望天；既确立了我所说的这些真理，又教导我们在未感谢那赐予我们这食物者之前，不可碰触饭食。\par

为什么祂不从无中造出饼呢？为要堵住马尔西翁和摩尼教徒的口，他们使祂的受造物与祂疏离，并藉祂的作为教导我们，连所有可见之物都是祂的作为和受造物，并表明正是祂自己赐予果实，祂在起初曾说：\Quote{让地生出青草}，以及\Quote{让水生出有生命的蠕动之物。}\par

因为这完全不比其他工作逊色。因为虽然那些是从无中造成的，但毕竟是由水而成；从地上生出果实，从水中生出有生命的活物，并不比从五个饼变出那么多，又从鱼变出更多（这是证明祂是天地的主宰）更伟大。\par

因此，既然病人常常是祂奇迹的对象，祂也施行一项普遍的恩惠，使众人不仅成为他人遭遇的旁观者，自己也分享这恩赐。\par

在旷野中令犹太人感到奇妙的事（他们至少说：\Quote{祂也能赐食吗？或在旷野摆设筵席吗？}）祂以行动显明了。为此，祂也领他们进入旷野，使奇迹远避嫌疑，无人以为附近某村庄为这餐饭贡献了什么。为此，祂不仅提到地点，也提到了时辰。\par

我们还学到另一件事：门徒在必要事物上的节制，以及他们如何轻视饮食。因为十二个人只有五个饼和两条鱼；肉体之事对他们如此次要，他们只依恋属灵之事。\par

就连那一点点，他们也没有紧紧抓住，而是被要求时就交出来了。由此我们应受教导：即使我们只有一点点，也应将它施予有需要的人。因此，当受命拿出五个饼时，他们没有说：\Quote{我们从哪里得食物呢？从哪里满足我们自己的饥饿呢？}而是立刻服从。\par

除了我所提到的，我至少认为，祂用他们现有的材料来行这事，是为引导他们到信德；因为那时他们仍处于软弱状态。\par

所以祂也\Quote{举目望天。}因为其他奇迹他们有许多例子，但这个却没有。\par

3. 因此，\Quote{祂拿起饼来，} 然后 \Quote{并且擘开，递给门徒，} 以此尊崇他们；不仅为尊崇他们，也为使奇迹发生后，他们不致怀疑或忘记，因自己的手为他们作证。\par

因此，祂也让群众先感到饥饿，等着他们先来求祂，并由他们使众人坐下，由他们分施；为要借他们自己的宣认和行动，预先感动每一个人。\par

因此，祂也从他们手中接过饼，使所行之事的见证增多，并使他们对奇迹有纪念。因为若连这些事之后他们都忘记了（\BibleRef{玛 16:9}），若祂省略那些预备，他们又会如何呢？\par

祂命令他们坐在被踏过的草地上，以此教导群众克己。因为祂的旨意不仅是喂养他们的身体，也是教导他们的灵魂。因此，借着地点，以及只给他们饼和鱼，且将同样的食物摆在众人面前，使之共享，不使人比他人多得，祂教导他们谦逊、节制、爱德，并彼此同心，视万物为公有。\par

\Quote{祂擘开递给门徒，门徒又递给群众。}五个饼被擘开分施，五个饼在门徒手中增多了。祂甚至未让奇迹止于此，还使它们有余，不是作为整饼，而是作为碎块，以表明这些饼有剩余，并使未在场的人知道所发生的事。\par

为此，祂也使群众饥饿，以免有人以为所发生的是幻觉。\par

为此，祂又使剩下的装满十二篮，使犹达斯也拿一篮。因为祂本可满足他们的饥饿，但门徒就不会知道祂的能力，正如厄里亚的情形所发生的（\BibleRef{列上 17:16}）。\par

无论如何，犹太人为此大大惊奇，甚至想立祂为王（\BibleRef{若 6:15}），虽在其他奇迹上他们从未如此。\par

什么推理能说明饼如何增多，如何在旷野中涌流，如何足够这么多人（因为\Quote{除了妇女孩子，约有五千男子；}这是对百姓极大的称赞，即妇女和男人都跟随祂）；剩余如何存在（这也不比前者小），且如此丰富，使篮子数目与门徒相等，不多不少？\par

拿起碎块后，祂没有给群众，而是给门徒，因为群众比门徒更不完美。\par

行了奇迹后，\Quote{随即祂催门徒上船，先渡到对岸，等祂叫众人散开。}\par

因为即使祂在场时看似制造幻觉，而非行真理；但不在场时当然不会。为此，祂将自己的作为置于精确的验证之下，命令那些持有纪念品和奇迹证明的人离开祂。\par

此外，当祂行大事时，祂将群众和门徒安置别处，教导我们在任何事上都不追求来自民众的荣耀，也不聚集人群环绕我们。\par

说\Quote{祂催他们，}表明门徒的紧紧跟随。\par

祂遣散他们的借口是群众，但祂自己却想上山；祂这样做，教导我们既不要总是与群众交往，也不要总是逃避人群，而要按需要各得其时，适当交换。\par

4. 因此，让我们也学习等候耶稣；但不要为感官的恩惠而等候，免得像犹太人一样受责备。因为祂说：\BibleQuote{你们找我，并不是因为见了神迹，而是因为吃饼吃饱了。}\BibleRef{若 6:26} 因此祂也没有一直行这个奇迹，只行了第二次；好叫他们学习不成为肚腹的奴隶，而是恒常依附于属神的事。\par

让我们也依附于这些事，寻求天上的食粮；领受之后，抛开一切世俗的挂虑。因为如果那些人撇下房屋、城市、亲属和一切，留在旷野，即使饥饿逼人也不退去；那么我们在接近这样的筵席时，更当显示更大的节制，并将爱情寄托于属神的事，而将感官的事视为次要。\par

因为犹太人受责备，并非因为他们为饼寻求祂，而是因为他们只为这个、并首先为这个寻求祂。因为如果有人轻视大恩赐，却紧紧抓住小的、施予者要他们轻视的东西，那么他也会失去这些后者；反之，如果我们爱那些大的，祂也会加上这些小恩赐。因为这些是那些恩赐的附属品；与那些相比，它们是如此卑贱和微不足道，即使它们是大的。因此，我们不要将操劳花费在它们上面，而将它们得失都同样视为无关紧要，正如约伯一样：他不紧抓它们当它们存在，也不寻求它们当它们不存在。因此，它们被称为\OriginalTerm{财物}{\GreekText{χρματα}}，不是要我们将它们埋在地里，而是为要善用它们。\par

正如工匠各有人其特有的技能，富人也当如此：他既不会治铜，也不会造船、纺织、建造房屋或任何此类事——那么让他学习善用他的财富，怜悯穷人；这样，他就学会一种比所有这些更好的技艺。\par

因为这技艺确实超越所有那些技艺。它的工坊建在天上。它的工具不是铁或铜，而是良善和正确的意志。基督和祂的父亲是这技艺的老师。因为祂说：\BibleQuote{你们应当慈悲，如同你们在天上的父那样慈悲。}\BibleRef{路 6:36}\par

而且令人惊奇的是，它既如此超越其他技艺，却不需要劳苦和时间来成全；只要愿意，一切都完成了。\par

但让我们也看它的结局是什么。那么它的结局是什么呢？天堂，天上的美物，那说不出的荣耀，属神的洞房，明亮的灯，与新郎同住；其他一切，言语甚至理智都无法表达的事。\par

因此在这方面，它与所有其他技艺也有很大不同。因为大多数技艺在今世使我们受益，但这技艺也为来世带来益处。\par

5. 但如果它如此超越那些为我们现世所必需的技艺，例如医学、建筑和所有其他类似的技艺：那么其他的就更不用说了，如果有人仔细审查，他甚至不会承认它们是技艺。因此，我至少不会称那些不必要的技艺为技艺。因为精致的烹调与制作酱料对我们有何益处？毫无益处；反而大大无益且有害，对灵魂和身体都造成伤害，因为它给我们带来所有疾病和痛苦之母——奢侈，连同巨大的挥霍。\par

不仅这些，连绘画或刺绣我也不会承认是技艺，因为它们只是把人投入无用的花费中。技艺应该关注我们生命中必要和重要的事，去提供并制作它们。因为天主赐给我们技能，正是为了让我们发明方法，用以装备我们的生活。但是墙壁上或衣服上的图像，请问有何用处？同样，制鞋匠和织布匠的技艺也应在许多方面大大削减。因为他们在其中大部分已流于庸俗的炫耀，败坏了其必要的用途，将诚实的技艺与邪恶的手艺混合；这在建筑技艺中也是如此。但只要建筑技艺建造房屋而不是戏院，从事必要之事而非多余之事，我就称之为技艺；同样，纺织之事，只要它制作衣服和被子，而不模仿蜘蛛，以许多荒谬和难言的娇柔压倒人，如此我就称之为技艺。\par

至于制鞋匠的行业，只要它制作鞋子，我不会剥夺其技艺的名称；但当它使人转向女性的姿态，并借由他们的鞋子变得放荡和娇柔时，我们将它归在有害和多余的事物中，甚至不称它为技艺。\par

我很清楚，在许多人看来，我过于琐细，忙于谈论这些；但我不会因此停止。因为我们一切邪恶的根源就在于这些过错常常被视为微不足道，因而被忽视。\par

但你们说，比这更小的罪是什么？就是拥有一双装饰华丽、闪闪发光、合脚的鞋；如果似乎还应该称它为罪的话。\par

那么，你们愿意我放开舌头谈论它，显示它的不雅有多少吗？你们不会生气吗？或者，即使你们生气，我也不太在乎。因为你们自己应对这愚妄负责，你们甚至不认为这是罪，从而迫使我们着手责备这种奢侈。来，让我们审视它，看看它是怎样的恶。因为当那些甚至不应编织在你们衣服中的丝线，却被你们缝在鞋子里时，这些事该受何等谴责，何等嘲笑？\par

若你们藐视我们的判断，请听保禄如何以极大的热忱禁止这些事，然后你们就会明白其中的荒谬。他说了什么？\Quote{不以编发、黄金、珍珠或华贵衣饰为妆饰。}\BibleRef{弟前 2:9} 那么，当保禄禁止已婚妇人穿华贵衣服时，你们竟将这柔靡之风延伸到鞋子上，并为了这种讥讽和羞辱而无休止地设计，你们还能配得什么恩宠？是的：先建造一艘船，然后召集桨手，配备舵手和掌舵者，扬起帆，横渡大洋，抛下妻子、儿女和祖国，商人将生命交给波涛，来到蛮夷之地，为了这些线，经历无数危险，好让你们最终取来，缝在鞋上，装饰皮革。还有什么比这愚妄更糟的呢？\par

但古时的风气并非如此，而是合于大丈夫的。因此，我预料随着时间推移，我们中间的年轻人甚至会穿女人的鞋子，并且不以为耻。更令人痛心的是，做父亲的看见这些事并不十分不悦，反倒视之为无所谓。\par

你们是否愿意我再加上更令人痛心的事：当有许多穷人时，这些事仍在发生？你们是否愿意我将基督带到你们面前：祂饥饿、赤身露体、到处漂泊、被锁链束缚？你们怎能不配遭受千雷击顶：忽视祂缺乏必需的饮食，却如此精心装饰这些皮革？当祂给门徒立法时，甚至不许他们穿鞋，但我们连赤脚行走都忍受不了，更不用说穿着合宜的鞋子了。\par

7. 那么，还有什么比这种不雅和荒谬更糟的呢？因为此事首先标志着一个柔弱的灵魂，其次是无情和残忍，然后是好奇和无聊。因为一个沉迷于这些多余之事的人，何时能专注任何必要的事？这样一个年轻人何时能忍受关注自己的灵魂，甚至思考自己有灵魂？是的，他必定是个轻浮的人，禁不住欣赏此类事物；他必定残忍，因这些事而忽视穷人；他必定没有德行，将全部精力花费在这些事上。\par

因为一个对线的美丽、颜色的鲜艳和编织的卷须好奇的人，何时能举目望天？何时能欣赏天上的美？他因皮革之美而激动，弯腰俯向地上，怎能欣赏那里的美？天主铺张穹苍，点亮太阳，引你向上看，你却强迫自己向下看，向着大地，如同猪一般，服从了魔鬼。因为恶毒魔鬼设计了这种不雅，要将你从那美中引开。为此他引你走这条路；而天主展示穹苍，却被一个魔鬼用某些皮革胜过——不，甚至不是皮革（因为连这些也是天主的化工），而是柔靡和恶俗的技艺。\par

这年轻人弯腰向地四处行走，他本应寻求天上的智慧，却为这些琐事沾沾自喜，仿佛完成了一件伟大的善工；他在广场上踮脚而行，为此给自己带来不必要的忧愁和烦恼：冬天怕弄脏鞋子，夏天怕蒙上灰尘。\par

人哪，你说什么？你因奢侈将整个灵魂抛入泥淖，任由它在地上拖行，却如此关心一双鞋？记住它们的用途，并尊重你加给它们的判断。鞋子是为践踏泥淖、污秽和路面所有斑点而造的。你若不能忍受，就把它们挂在脖子上，或顶在头上。\par

你们听到这些笑，但我却为这些人的疯狂和对此事的执着而流泪。因为实际上，他们宁愿用泥玷污身体，也不愿玷污那些皮革。\par

这样，他们成了轻浮之人，又在另一方面成了贪财者。因为习惯于为这类事狂热心急的人，对衣服和一切也需要巨大花费和丰厚收入。\par

若父亲慷慨，他的束缚更甚，荒谬嗜好更烈；若父亲节俭，他便被迫做其他不雅的事，以凑钱应付这类开销。\par

因此，许多年轻人甚至出卖了男子气概，成为富人的食客，从事其他奴役性工作，以此购买这些欲望的满足。\par

所以，此人必定既贪财又轻浮，且在重要事上最懒散，并被迫犯许多罪，这是显而易见的。说他残忍和虚荣，无人能反驳：残忍，因为他看见穷人时，因爱装饰而视若无睹，却用黄金装饰这些东西，任穷人饿死；虚荣，因为即使在琐事上，他也训练自己追逐旁观者的赞美。因为我认为，没有一位将军如此为自己的军团和战利品自豪，如这些放荡青年为装饰鞋子、曳地长袍和发型而自豪；然而，这一切都是他人手艺的作品。如果人们连在他人作品中都止不住虚荣，又何时能在自己的事上止住呢？\par

8. 我是否还应提其他更严重的事？或者这些已足够？那么，我必须在此结束讲论；因为我说这些也是因为那些好争辩的人，他们坚持认为此事并非如此恶劣。\par

虽然我知道许多年轻人不会留意我所说的，因为他们已完全沉醉于这种嗜好，但我仍不应因此缄默。因为那些有见识、尚未败坏的父亲，将能强迫他们，甚至违背其意愿，走向合宜的端庄。\par

那么，不要说：\Quote{这件事无关紧要，那件事无关紧要}；因为这件事、这件事已毁了一切。因为正应借此训练他们，并借着看似琐碎的事情，使他们成为庄重、心胸宽广、超越外在服饰的人；这样我们也会发现他们在大事上得到认可。因为还有什么比学习字母更平常呢？然而，人正是借此成为修辞学家、诡辩家和哲学家；如果他们不识字，也就永远不会拥有那种知识。\par

我们这样说，不仅是对青年男子，也是对妇女和年轻女子。因为她们也容易受到同样的指责，而且更加如此，因为端庄是适合贞女的美德。\par

因此，对其他人所说的，你们也要当作是对你们说的，这样我们就不必再重复同样的话了。\par

因为现在是用祈祷来结束我们讲论的时候了。你们众人都要与我们一同祈祷，使教会中的青年男子首先能够过有秩序的生活，并达到合宜的老年。对于那些不如此生活的人，确实最好不要活到老年。但对于那些在青年时期就已成熟的人，我祈祷他们也能达到白发的深处，成为合格子女的父亲，并给生养他们的人带来喜乐，而最重要的是给创造他们的天主带来喜乐；并且摧毁每一种不健康的想法，不仅关于他们的鞋子、衣服，也包括其他各种。\par

因为正如未开垦的土地，被忽视的青年也是如此，从四面八方生出许多荆棘。让我们向他们发出圣神之火，烧尽这些邪恶的欲望；让我们开垦我们的田地，为接受种子做好准备；让我们中间的年轻人比别处的老人更加清醒明智。因为这实在是奇妙的事：当节制在青年中闪耀时；因为老年时的节制确实没有大的赏报，因为年龄本身已提供了保障。而奇妙的是，在波涛中享受平静，在火窑中不被焚烧，在青年时期不放纵妄为。\par

因此，让我们将这些事放在心上，效法那有福的\Person{若瑟}，他在所有这些试炼中都闪耀着光芒，以便我们与他一起获得同样的荣冠；愿我们都能获得，藉着我们的主\Person{耶稣基督}的恩宠和对人的爱，愿光荣归于圣父，同圣神，从现在直到永远，永世无尽。阿们。\par

\SourceRecord{Translated by George Prevost and revised by M.B. Riddle.；Nicene and Post-Nicene Fathers, First Series；Vol. 10.；Edited by Philip Schaff.；Buffalo, NY: Christian Literature Publishing Co.,；1888.；<http://www.newadvent.org/fathers/200149.htm>.}

% structure: p0001 source=h1 target=section reason=page-title

\section{玛窦福音第50篇讲道}

\BibleRef{玛 14:23-24}.\par

\Quote{\Emph{祂遣散群众后，独自上山祈祷：到了晚上，祂独自一人在那里。但船已行至海中，被波浪摇撼：} \Emph{因为风不顺。}}\par

祂为何上山？为教导我们，当我们向天主祈祷时，孤独和退隐是有益的。你看，祂正是为此不断退入旷野，常常整夜祈祷，教导我们恳切寻求祈祷中的宁静，这是时间和地点所能赋予的。因为旷野是宁静之母，是风平浪静的港湾，将我们从一切纷扰中救出。\par

祂自己随即上山，门徒却又被波浪摇撼，遭遇风暴，与先前一般无二。但先前遇险时，祂在船上；如今他们却独自一人在海上。祂就这样温和地逐步激励他们，促使他们变得更好，甚至能高尚地承受一切。因此我们看见，当他们初次临近那危险时，祂虽在睡觉，却也在场，以便随时救助；但现在，祂带领他们达到更大的\OriginalTerm{坚忍}{endurance}，连这点也不做了，而是离开，在海上任凭风暴兴起，使他们甚至不从任何方面指望得救的希望；并且让他们整夜颠簸，我想，这是要彻底唤醒他们刚硬的心。\par

因为恐惧的本性如此，时间与恶劣天气相合而产生的恐惧。同时，祂也使他们\OriginalTerm{痛悔}{compunction}，并激发他们对祂更强烈的渴慕和不断的怀念。\par

因此，祂也没有立刻显现给他们。因为，如经上所载，\Quote{在第四更，}如此记载，\Quote{祂行走在海面上，到了他们那里}；\BibleRef{玛 14:25}教训他们不要急于从迫在眉睫的危险中寻求解救，而要勇敢地承担一切遭遇。\par

经上说：\Quote{门徒看见祂在海面上行走，就惊慌说：\InnerQuote{是个妖怪。}并且吓得喊叫起来。}\par

是的，祂常常这样做；当祂要消除我们的恐惧时，却给我们带来其他更坏、更可怕的事：我们看见当时也正是如此。因为除了风暴，那景象也困扰他们，不亚于风暴。因此，祂既没有消除黑暗，也没有立刻显示自己，正如我所说，祂借着这些恐惧的持续来训练他们，教导他们准备好忍受。在\Person{约伯}身上，祂也这样做了；当祂要消除约伯的恐怖和诱惑时，却让结局变得更加痛苦：我指的不是他子女的死亡或他妻子的话，而是他的仆人和朋友的责备。当祂要从异乡的苦难中解救\Person{雅各伯}时，祂让他的烦恼被唤醒和加剧：先是他的岳父追上他并威胁要杀他，接着他的兄长随即来到，将极端的危险悬在他头上。\par

因为人既不能长久受试探，又不能受严重的试探；当义人的争战即将结束时，祂愿意他们获得更多，便加强他们的斗争。祂在\Person{亚巴郎}身上也这样做了，指定以他的儿子作为他最后的争战。因为当这样的事临到我们，而使解救近在咫尺、就在门口时，甚至难以忍受的事也会变得可以忍受。\par

基督当时也是如此，直到他们喊叫之前，祂没有显露自己。因为他们越惊慌，就越欢迎祂的到来。后来当他们喊叫时，记载说，\par

\Quote{耶稣立刻对他们说：放心！是我，不要怕！}\BibleRef{玛 14:27}\par

这句话消除了他们的恐惧，使他们有了信心。因为他们因祂奇妙的行走方式和时间而未能凭视觉认出祂，祂便用声音显明自己。\par

2. 那么，处处热忱、总领先于他人的\Person{伯多禄}说什么呢？\par

\Quote{主，如果是祢，}他说，\Quote{请吩咐我走到祢那里去，在水面上。}\BibleRef{玛 14:28}\par

他没有说\Quote{祈求恳求}，而是说\Quote{吩咐。}你看见他的热忱和\OriginalTerm{信德}{faith}有多么大吗？然而，他确实常因寻求超出自己能力的事而陷入危险。因为在这里，他也要求了极其伟大的事，纯粹出于爱，而非炫耀。因为他没有说\Quote{吩咐我在水面上行走}，而是说什么？\Quote{吩咐我走到祢那里去。}因为没有人像他那样爱耶稣。\par

复活后，他也这样做了；他等不及与别人同来，而是纵身跳入水中。\BibleRef{若 21:7}他不仅显出爱，也显出\OriginalTerm{信德}{faith}。因为他不仅相信祂自己能行走在海面上，也相信祂能带领别人行在其上；他渴望尽快接近祂。\par

\Quote{耶稣说：\InnerQuote{来吧。}伯多禄就从船上下来，在水面上行走，向耶稣走去。但他一见风势猛烈，就害怕起来，开始下沉，便喊道：\InnerQuote{主，救我！}耶稣立刻伸手拉住他，对他说：\InnerQuote{小信德的人哪，你为什么怀疑？}}\BibleRef{玛 14:29}\par

这比先前更奇妙。因此，这事发生在之后。因为当祂显明祂掌管大海时，祂便将这记号推向更奇妙的事。那时祂只斥责风；但现在祂自己行走，也允许别人行走。如果祂一开始就要求这样做，伯多禄不会如此好地接受，因为他还没有获得如此大的\OriginalTerm{信德}{faith}。\par

那么，\Person{基督}为什么允许他呢？因为如果祂说\Quote{你不能}，热忱的伯多禄会再次反驳祂。因此，祂用事实说服他，使他今后保持清醒。\par

但即使如此，他也没有坚持住。因此，他一下船就头晕；因为他害怕了。这是海浪造成的，但他的恐惧是由风引起的。\par

但\Person{若望}说，\Quote{他们就欣然接祂上船；船立刻到了他们所要去的地方。}记述同一件事。如此，当他们即将抵达岸边时，祂上了船。\par

\Person{伯多禄}于是下船，到祂那里去，并不因在水上行走而多么喜乐，而是因来到祂跟前。当他得胜了那更大的事，却险些被那较小的事所害，即那风的猛烈，我是说，不是海的。因为人性如此；并非不常成就大事，却在小事上暴露自己；正如\Person{厄里亚}在\Person{依则贝耳}面前，\Person{梅瑟}在埃及人面前，\Person{达味}在\Person{巴特舍巴}面前。同样，这人也是如此；当他们的恐惧仍高涨时，他鼓起勇气在水上行走，却再也无法抵挡风的袭击，即使他离基督如此之近。所以，仅仅靠近基督，若不以信德靠近祂，毫无益处。\par

这也显示了师傅与门徒之间的区别，并安抚了其他人的情绪。因为若论那两位弟兄，他们曾心怀愤慨，何况在此呢？因为他们尚未蒙赐圣神。\par

但后来他们就不再如此。例如，每逢场合，他们都将首位荣誉让给\Person{伯多禄}，在向民众讲话时推举他，尽管他比他们任何人都更粗鲁。\par

为何祂没有命令风止息，而是亲自伸手抓住他？因为在他身上需要信德。当我们这方面有所欠缺时，天主那方面也止步不前。\par

因此，祂表明并不是风的袭击，而是他缺乏信德导致他的倾覆，祂说：\Quote{你为什么疑惑？小信德的人哪！}所以，若他的信德不是软弱的，他也能轻易抵御那风。正因如此，你看，即使祂已抓住他，仍让风吹着，表明当信德坚定时，风不会带来伤害。\par

正如雏鸟过早地离巢，行将坠落时，母鸟用翅膀托起它，带它回巢；基督也是如此行的。\par

\BibleQuote{他们一上了船，风就停了。}\BibleRef{玛 14:32}\par

然而在此之前他们曾说：\BibleQuote{这是什么人，连风和海也服从他！}\BibleRef{玛 8:27} 但现在却不是这样。因为经上说：\BibleQuote{\Quote{在船里的人}前来朝拜祂说：\InnerQuote{你真是天主子。}}\BibleRef{玛 14:33} 你看见祂如何逐渐引导他们越来越高升吗？因为既因祂在水上行走，又因祂命令别人也如此行，并在危难中保全他，他们的信德从此变得强大。因为那时祂确实斥责了海，但现在祂却没有斥责，以另一种方式更充分地显示祂的能力。因此他们也说：\BibleQuote{你真是天主子。}\par

那么怎样呢？祂责备他们这样说吗？不，恰恰相反，祂反而证实了他们所说的，以更大的权威治愈那些前来接近祂的人，且不像从前那样。\par

\Quote{他们渡到对岸，}经上说，\Quote{来到\Place{革乃撒勒}地方。那地方的人一认出是祂，就派人到周围整个地区，把一切有病的都带到祂跟前，恳求祂让他们只摸一摸祂外衣的繸子；凡摸着的，都完全痊愈了。}\par

因为他们不再像从前那样，拉祂进自己屋里，寻求摸祂的手，以及从祂口中得指示；而是以更高昂的热忱，更多的克己，更丰盛的信德，努力为自己赢得治愈；因为患血漏的女人教导他们所有人要热切寻求智慧。\par

圣史也暗示，祂间隔很久才访问各邻乡，说：\Quote{那地方的人认出祂，就派人到周围地区，把有病的人都带到祂跟前。}但间隔非但没有消除他们的信德，反而使之更大，并保持活力。\par

3. 那么我们也当触摸祂外衣的繸子，或者更确切地说，若我们愿意，我们拥有祂全人。因为如今祂的身体已摆在我们面前，不只是祂的外衣，而是祂的身体；不仅供我们触摸，也供我们吃且饱足。那么，现在让我们每一位有疾病的人都以信德前来亲近。因为若那些触摸祂外衣繸子的人从祂那里获得了如此大的能力，何况那些拥有祂全人呢？以信德亲近，不仅是指领受祭品，也是以纯洁的心触摸它；要存这样的心，如同亲近基督自己。因为，纵使你听不到声音，你仍看见祂被陈设；或者你更听见祂的声音，因祂正藉着圣史们说话。\par

所以，要相信，即使是现在，那也正是祂亲自坐席的晚餐。因为这与那一次毫无分别。因为人既不制作这个，祂自己也不制作那个；但两者都是祂的工程。因此，当你看见司铎递给你时，不要以为是司铎如此行，而是基督的手伸出。\par

正如当某人施洗时，不是他给你施洗，而是天主以无形的能力掌握你的头，没有任何天使或总领天使或任何其他受造物敢接近并触摸你；现在也是如此。因为当天主生育时，恩赐唯独属于祂。你们难道没有看见，那在此世收养儿子的人，如何不将此事委托给奴隶，而是亲自出席审判台吗？同样，天主也没有将祂的恩赐委托给天使，而是亲自临在，命令并说：\BibleQuote{不要称地上的人为父；}\BibleRef{玛 23:9}这不是要你们侮辱生育你们的人，而是要使你们在众人之上，偏爱那造了你们、并将你们登记在自己子女当中者。因为那赐予更大恩赐的，即将自己置于你们面前，祂岂会不屑于将祂的身体分施给你们呢？因此，让我们——无论是司铎还是信友——听一听我们被恩赐了什么；让我们听而战栗。祂已让我们饱饫了祂自己的圣体；祂已将自己祭献放在我们面前。\par

那么，当我们以这样的食物为食，却犯下这样的罪时，我们还有什么借口呢？当我们吃羔羊时，我们却变成了狼？当我们以羊为食时，我们却像狮子一样以暴力掠夺？\par

因为祂命令这奥迹始终保持纯洁，不仅免于暴力，甚至免于单纯的敌意。是的，因为这奥迹是和平的奥迹；它不允许我们执着于财富。因为如果祂为我们没有顾惜自己，那么，我们顾惜自己的财富，却浪费灵魂——祂为这灵魂没有顾惜自己——我们该当如何呢？\par

天主每年在犹太人的节期上，将他们特殊恩惠的纪念系于他们身上；但对你来说，我可以说，每天通过这奥迹，也是如此。\par

因此，不要以十字架为耻：因为这是我们的可敬之物，这是我们的奥迹；我们用这恩赐来装饰自己，我们因它而美丽。\par

如果我说，祂铺张了穹苍，祂展开了大地和海洋，祂派遣了先知和天使，我所说的都不足为道。因为祂恩惠的总和是：\BibleQuote{祂没有顾惜自己的儿子，}\BibleRef{罗 8:32}为要拯救祂那疏远的仆人。\par

4. 所以，不要让任何犹大接近这祭台，也不要让任何西满到来；因为此二人皆因贪财而灭亡。因此，让我们逃离这深渊；不要以为，如果我们掠夺了寡妇和孤儿，然后为这祭台献上金杯和宝石杯，就足以得救。不，如果你渴望尊崇这祭献，就献上你的灵魂——那祭献正是为它而宰杀的；使它成为金的；但如果它仍然比铅或陶土更坏，而器皿却是金的，那有什么益处呢？\par

因此，不要让这成为我们的目标：只献上金器皿，而是也要从诚实的收入中这样做。因为这些比金子更珍贵，这些是无害的。因为教会不是金库，也不是银作坊，而是天使的集会。因此，我们需要的是灵魂，因为事实上天主为了灵魂的缘故而接受这些。\par

那时，那桌子不是银的，那杯也不是金的，基督用它们将祂自己的血赐给祂的门徒；但那里的一切都是宝贵的，且令人敬畏，因为他们充满了圣神。\BibleRef{弗 5:18}\par

你愿意尊崇基督的身体吗？不要在他赤身露体时忽视他；不要在这里用丝织品尊崇他，却忽视他在外面因寒冷和赤身而濒死。因为那说\BibleQuote{这是我的身体}，并用祂的话证实了此事的那一位，也同样说：\BibleQuote{你们看见了我饿了，没有给我吃；}\BibleQuote{凡你们没有给这些最小中的一个做的，就是没有给我做。}\BibleRef{玛 25:42}因为祂确实不需要遮盖物，而需要纯洁的灵魂；但那一端却需要很多关注。\par

因此，让我们学习在生活上严谨，并照基督自己所愿的尊崇祂。因为对于被尊崇者，那最令祂喜悦的尊崇，是祂自己愿意得到的，而不是我们认为最好的。因为伯多禄也曾以为，通过阻止耶稣洗他的脚，是在尊崇祂，但这样做不是尊崇，而是相反。\par

同样，你也应用祂所命定的这尊崇来尊崇祂，将你的财富花费在穷人身上。因为天主完全不需要金器皿，而是需要金灵魂。\par

我说这些话，并不是禁止提供这样的奉献，而是要求你们连同它们、并在它们之前，施舍救济。因为祂确实接受前者，但更接受后者。因为在其中，只有奉献者受益；而在另一中，受益者也是接受者。在这里，行为似乎甚至是炫耀的缘由；但在那里，一切都是仁慈和爱人的。\par

因为有什么益处呢？当祂的桌子满是金杯，而祂却因饥饿而消亡？先使饥饿的祂得饱足，然后也丰盛地装饰祂的桌子。你给祂做一个金杯，却不给祂一杯凉水？有什么益处呢？你为祂的桌子铺上镶金布，却连必要的遮盖也不给祂自己？那又有什么好处呢？告诉我，如果你看见一个人缺乏必要的食物，却不先消除他的饥饿，却先用银子覆盖他的桌子，他真会感谢你吗？难道不会反而愤怒吗？又如，如果你看见一个人衣衫褴褛、冻得僵硬，却忽视给他衣服，反而建造金柱，说\Quote{你在尊崇他}，难道他不会说你是在嘲笑，并视为侮辱，而且是最严重的侮辱吗？\par

因此，你也当对基督这样想：当祂四处漂泊，无家可归，需要一个屋顶来遮盖祂时，你却忽略接待祂，反而装饰地板、墙壁、柱头和灯台，用银链悬挂灯具，但祂被捆锁在监狱里，你甚至不愿去看祂一眼。\par

5. 我说这些话，并不是禁止在这些事上慷慨大度，而是劝告你们要同时做那些事，甚至优先于这些。因为不履行这些事，从未有人被责备；但不履行那些事，则有地狱的威胁、不灭的火、以及和恶鬼一同受的刑罚。所以，不要一边装饰祂的房屋，一边忽视在困苦中的弟兄，因为他比房屋更是圣殿。\par

这些你的积蓄将会被不信的君王、暴君和盗贼所夺去；但无论你为饥饿、作客、赤身露体的弟兄所做的一切，连魔鬼也不能夺走，反而会积存于一个不可侵犯的宝库中。\par

那么为什么祂自己却说：\Quote{穷人你们常与你们同在，但你们不常与我同在？}这恰恰是我们最应该施舍的原因，因为我们现在生活中不常有饥饿的祂。但若你渴望也学习这话的完整意义，要明白这话不是针对祂的门徒说的，虽然看似如此，而是针对那妇人的软弱。也就是说，她的心性仍相当不完美，而他们对她怀疑；为了鼓励她，祂说了这些话。为了证明祂是为了安慰她而说这话，祂还添了一句：\Quote{你们为什么为难这妇人？}\BibleRef{玛 26:10}至于我们真正常与祂同在这一点，祂说：\Quote{看，我天天与你们同在，直到今世的终结。}\BibleRef{玛 28:20}由此显然，说这话的唯一目的，是不让门徒的责备挫伤那妇人刚刚萌芽的信德。\par

因此我们现在不要提出这些因某种计划而说出的话，而应阅读旧约和新约中关于施舍的所有法律，并认真对待此事。因为这能洁净罪恶。因为\BibleQuote{施舍，一切对你们就都洁净了。}\BibleRef{路 11:41}这比祭献更大。\BibleQuote{因为我喜欢仁慈，不喜欢祭献。}这能开启天门。因为\BibleQuote{你的祈祷和你的施舍已上达天主面前，获得纪念。}\BibleRef{宗 10:4}这比童贞更不可少：因为那些童女因此被赶出婚筵；而另一些因此被带入。\par

让我们思考这一切，并慷慨地播种，以便我们能更丰盛地收获，并赖我们主耶稣基督的恩宠和爱人之心，获得将来的美善，愿光荣归于祂，直到永远。阿们。\par

\SourceRecord{Translated by George Prevost and revised by M.B. Riddle.；Nicene and Post-Nicene Fathers, First Series；Vol. 10.；Edited by Philip Schaff.；Buffalo, NY: Christian Literature Publishing Co.,；1888.；<http://www.newadvent.org/fathers/200150.htm>.}

% structure: p0001 source=h1 target=section reason=page-title

\section{\WorkTitle{《玛窦福音》讲道集第五十一篇}}

\BibleRef{玛 15:1}\par

\Emph{\BibleQuote{那时，有来自耶路撒冷的经师和法利塞人来到耶稣跟前，说：\InnerQuote{祢的门徒为什么……}等等。}}\par

那时；何时？当祂行了无数的奇迹之后；当祂用衣繸的触碰治愈了病人之后。因为福音作者正是以此标记时间，为要表明他们那无法言说的邪恶，竟不被任何事物抑制。\par

但所谓\Quote{来自耶路撒冷的经师和法利塞人}是什么意思？他们在各支派中分散，分作十二份；但占据首府的那些人比别人更坏，因为他们既享有更多的荣誉，也积累了相当的傲慢。\par

但我请你们注意，他们如何被这问题本身揭露出来；他们不说：\Quote{为什么他们违犯梅瑟的法律？}而说：\Quote{为什么他们违犯先人的传统？}由此可见，司祭们正在发明许多新规，尽管梅瑟曾以极大的恐惧和威胁命令不可增添也不可删减。他说：\Quote{\BibleQuote{你不可增添你今日吩咐你的话，也不可从中删减。}}\BibleRef{申 4:2}\par

但他们仍然照样创新；例如：不可用未洗的手吃饭，必须清洗杯子和铜器，也必须清洗自身。这样，当人们本应从这些规条中逐渐获得解放的时候，他们却反而用更多的规条束缚他们，生怕有人夺去他们的权力，并希望引起更大的敬畏，好像他们自己也是立法者一样。这事已经发展到如此严重的地步，以至于他们自己的诫命被遵守，而天主的诫命却被违反；并且他们如此得势，以致这竟成了控告的理由。这双重指控：他们既发明新规，又对自己所定的严加执行，而对天主的诫命却毫不顾及。\par

他们省略了其他东西（如锅和铜器，因为太可笑），只提出那些看起来比其余更合理的事，我想，他们试图以此激怒祂。为此他们也提及先人，以便祂因轻视它们而给自己留下把柄。\par

但首先应当询问，门徒为何用未洗的手吃饭。那么他们为何如此吃呢？不是因为他们坚持这样做，而是因为他们从此忽略那些多余的事，只关注必要的事；他们并没有洗或不洗的规定，而是随意做。因为那些连自己必要的食物都轻视的人，又怎会把这些事看得很重要呢？这样的事往往不经心地发生——例如，当他们在旷野吃饭时，当他们掐麦穗时——现在却被那些人当作罪名提出来，而他们自己总是在大事上犯过，却对细节斤斤计较。\par

2. 那么基督说什么？祂没有反对，也没有作任何辩护，而是立即反过来责备他们，挫败他们的自信，并表明那犯大罪的人不应在细事上苛责他人。祂说：\Quote{\InnerQuote{你们本应受责备，反而责备别人吗？}}\par

但你当留意，每当祂决意废除律法中的某些规定时，祂总是以辩护的方式行事；就像在这件事上一样。祂决不立刻去违犯它，也不说：\Quote{这不算什么}；因为那样肯定会让他们更加放肆；而是首先彻底削去他们的胆量，提出更重的指控并把它们引到他们自己头上。祂既不说：\Quote{他们违犯它是好的}，免得他们抓住祂的把柄；也不谴责他们的做法，免得确认那条法律；另一方面，祂也不责备先人，视他们为无法无天的邪恶之人；因为无疑他们会避开祂，视之为侮辱者和伤害者；但祂放弃所有这些方式，走了另一条路。祂似乎是在责备那些来到祂跟前的人，但实际上是在斥责那些制定这些法律的人；祂没有一处提到先人，但藉着对经师的指控也推倒了他们，并表明他们的罪是双重的：第一，不服从天主；第二，为了人的缘故这样做；好像祂在说：\Quote{\InnerQuote{为什么？正是这事毁了你们：你们凡事顺从先人。}}\par

然而祂并不如此说，但这正是祂所暗示的，因为祂如此回答他们：\par

\Quote{你们为什么因你们的传统而违犯天主的诫命？因为天主命令说：\BibleQuote{应孝敬你的父亲和你的母亲；}又说：\BibleQuote{咒骂父亲或母亲的，应处死。}你们却说：\InnerQuote{谁若对父亲或母亲说：我所能供养你的，已成了献仪，}且不孝敬他的父亲或他的母亲——你们就因你们的传统而废除了天主的诫命。}\par

祂没有说\Quote{先人的传统}，而说\Quote{你们的}。又说：\Quote{你们说}；再次没有说\Quote{先人说}：为使祂的话不那么刺痛人。也就是说，因为他们想证明门徒违犯律法，祂指出他们自己正是如此，而门徒却是无可指责的。因为当然，那不是一条法律，而是人所制定的（所以祂也称它为\Quote{传统}），尤其是那些违犯法律的人所制定的。\par

既然这并不与法律有任何矛盾之处，即命令人洗手，祂又提出另一条与法律相对立的传统。祂所说大致如下：\Quote{他们以虔诚为外衣，教导年轻人轻视他们的父亲。}如何以及什么方式？\Quote{如果父母中一位对儿子说：\InnerQuote{把你拥有的这只羊，或这头牛，或任何这类东西给我}，他们就会说：\InnerQuote{这是献给天主的礼物，你可能会从我这里得到益处，你不能拥有它。}由此产生了两个弊端：一方面他们并没有将礼物献给天主，另一方面他们以献礼为名欺骗父母，既为了天主的缘故侮辱父母，又为了父母的缘故侮辱天主。}但祂没有立即说出这一点，而是先复述法律，以此表明祂热切希望父母受到尊敬。因为祂说：\Quote{尊敬你的父亲和你的母亲，好使你在世上长寿。}又说：\Quote{咒骂父亲或母亲的，应处死刑。}\BibleRef{出 21:17}\par

但祂省略了第一项，即为孝敬父母者所定的赏报，而提出更严厉的，即对不孝者所威胁的惩罚；意欲既令他们惊惧，又感化那些有悟性的人；并暗示他们因此当死。因为若有人在言语上侮辱父母尚且受罚，何况你们，不仅在实际行为上如此行，且教导别人这样做呢？\Quote{你们这些本不应当存活的人，又如何指责门徒呢？}\par

\Quote{又有什么可惊奇的呢？如果你们向我这个尚且不为人所知的人如此侮辱，况且对圣父你们也行了同样的事？}因为祂处处申明并暗示，他们的这种傲慢是从祂开始。\par

但有些人也作另一种解释：\Quote{这是礼物，凡你可能会从我这里得到益处的；}也就是说，我不欠你任何尊敬，但如果我确实尊敬你，那是我给你的无偿礼物。但基督不会提到那种侮辱。\par

而马尔谷又更清楚地说明了这一点，他说：\Quote{这是\OriginalTerm{科尔班}{Corban}\BibleRef{谷 7:11}，即凡你可能会从我这里得到益处的；}意思是，这不是礼物和馈赠，而是正当地奉献。\par

然后，既然表明了那些践踏法律的人无权正当指责别人违反某些长者的命令，祂又同样从先知那里指出同一件事。这样，祂一旦严厉地抓住他们，便进一步进行：正如祂每次所做的那样，引出圣经，从而证明自己与天主一致。\par

先知说了什么？\BibleQuote{这百姓用嘴唇尊敬我，心却远离我。他们恭敬我也是枉然，将人的规条当作道理教导人。}\par

你看见一个与祂的话完全一致的预言，并且从一开始就预先宣告了他们的邪恶吗？因为基督现在责斥他们的，依撒意亚也早已说过：他们藐视天主的话，\Quote{他们恭敬我也是枉然，}祂说；但对他们自己的话却很重视，\Quote{将人的规条，}祂说，\Quote{当作道理教导人。}因此，门徒不遵守它们是有道理的。\par

3. 你看，祂给了他们致命一击；首先从事实出发，然后从他们自己的认同出发，再通过先知加重了指控，祂的确不再与他们说话，因为他们现已无可救药，而是将话转向群众，以便引入祂那伟大、崇高且充满严厉的教诲；并利用前一个话题的机会，进一步插入更重要的内容，即废止对食物的遵守。\par

但要留意是什么时候。当祂洁净了癞病人，废除了安息日，显示自己是天地的主宰，制定了法律，赦免了罪过，复活了死人，给了他们许多天主性的证明之后，祂才谈论食物之事。\par

事实上，犹太人的全部宗教都包含于此；如果你废除了这一点，你就废除了全部。因为祂借此表明，割损也必须废止。但祂自己并没有突出地提出这一点（因为割损比其它诫命更古老，且更受尊重），而是通过祂的门徒来制定。因为此事如此重大，以至于门徒们在很长时间以后，想要废除它时，也先实行它，然后才废除了它。\BibleRef{宗 16:3}\par

但要留意祂如何引入祂的法律：\Quote{祂叫了群众来，对他们说：你们听，且要明白。}\BibleRef{玛 15:11}\par

因此，祂绝非简单地将其启示给他们，而是通过尊重和礼貌，首先使祂的话被人们接受（因为圣史说：\Quote{祂叫他们到祂跟前来}）；其次，也通过时机；在他们被驳倒、祂战胜他们、以及先知指控之后，祂才开始立法，那时他们也更容易接受祂的话。\par

祂不仅叫他们到祂跟前来，还使他们更加留心。因为祂说：\Quote{明白}，即\Quote{思考，振作起来；因为即将颁布的法律正是如此。如果他们甚至不合时宜地为自己的传统而搁置法律，而你们听从了他们；那么你们更应当听从祂，我在适当的时机正带领你们迈向更高的自律准则。}\par

祂并没有说\Quote{遵守食物的规定是无关紧要的，也不是梅瑟给出了错误的命令，也不是祂出于迁就而这样做}，而是以劝诫和规劝的方式，并从事物的本性中取出祂的见证，祂说：\BibleQuote{不是入口的污秽人，而是出口的污秽人} \BibleRef{玛 15:11}。祂在命令和证明中都诉诸自然本身。然而他们听了这一切，没有回答，也没有说\Quote{祢说什么？当天主关于食物的遵守给了无数命令时，祢却制定这样的法律？} 但祂不仅通过驳斥他们，而且通过揭露他们的诡计，暴露他们暗中做的事，揭示他们内心的秘密，彻底堵住了他们的嘴；他们的嘴被堵住，就离开了。\par

但是请注意，我恳求你们，祂还没有明确地大胆反对食物。因此祂既没有说\Quote{食物}，而是说\Quote{入口的不污秽人}；这也是他们自然会怀疑到未洗手的事。因为祂确实在谈论食物，但也会被理解为这些事。\par

这固然是出于对食物的顾虑之感如此强烈，以至于复活后伯多禄还说：\BibleQuote{主啊，万万不可，因为我从来没有吃过俗物或不洁净的东西。} \BibleRef{宗 10:14} 虽然他说这话是为了别人，并为自己留下辩护的理由以应对指责者，指出他实际上曾经抗议，即使这样也未被宽免，然而这也表明了他们对那点的印象之深。\par

因此你们看见，祂自己在开始时也没有公开谈论食物，而是说\Quote{入口的东西}；后来当祂似乎说得更明白时，又用结论将它掩盖，说：\BibleQuote{但用未洗手吃饭不污秽人。} \BibleRef{玛 15:20} 这样祂似乎是从那件事引出话题，并且仍是在谈论同一件事。因此祂并没有说\Quote{吃食物不污秽人}，而是仿佛在谈论另一话题，使他们无可辩驳。\par

4. 因此当他们听了这些事，经上记载说\Quote{法利塞人}就\Quote{起了反感}\BibleRef{玛 15:12}，不是群众。因为\Quote{祂的门徒}，经上记载说，\Quote{前来对祂说：祢知道法利塞人听见这话就起了反感吗？} 然而肯定没有话是对他们说的。\par

那么基督说什么？祂没有因他们而除去那绊脚石，反而责备他们，说：\BibleQuote{凡我圣父没有栽种的植物，必要拔出来。} \BibleRef{玛 15:13} 因为祂惯于轻视反感，又不轻视反感。例如，在别处祂说：\BibleQuote{但为了避免使他们反感，往海里去投钩。} \BibleRef{玛 17:27} 但在这里祂说：\Quote{任凭他们吧，他们是瞎子的领路人；若是瞎子领瞎子，两人都要掉进坑里。}\par

但祂的门徒说这些话，不仅是为那些人忧伤，也是因为他们自己也感到有些困惑。但因为他们不敢以自身名义说，就愿意通过告诉祂别人的事来了解。至于是否如此，请听此后热切且勇于直前的伯多禄来到祂面前说：\BibleQuote{请为我们解释这个比喻} \BibleRef{玛 15:15}，显露出他灵魂中的困扰，确实不敢公开说\Quote{我起了反感}，而是要求通过祂的解释来解除困惑；因此他也受了责备。\par

那么基督说什么？\BibleQuote{凡我圣父没有栽种的植物，必要拔出来。}\par

这话被那些染上摩尼教祸害的人断言是对法律说的；但他们的嘴被前面的话堵住了。因为如果祂是在谈论法律，那么祂又如何在后来为法律辩护，并为之争辩，说：\BibleQuote{你们为什么因你们的传统而违反天主的诫命？} 祂又怎样引出先知？但祂是对他们自己和他们的传统说这话。因为如果天主说：\BibleQuote{孝敬你的父亲和你的母亲}，那么由天主所说的话，怎么不是天主栽种的呢？\par

下文也表明，这话是对他们自己和他们传统的说的。因此祂补充说：\BibleQuote{他们是瞎子的领路人。} 然而，如果祂是对法律说的，祂会说：\Quote{它是瞎子的领路人。} 但祂没有这样说，而是说：\BibleQuote{他们是瞎子的领路人。} 从而免除法律的责任，将责难完全归到他们身上。\par

然后，为了将百姓也同他们分开，因为他们正因他们而快要掉进坑里，祂说：\Quote{若是瞎子领瞎子，两人都要掉进坑里。}\par

仅仅是瞎眼就是大恶，但处于这种情况且无人引领，并且还占据向导之位，便是双重和三重的责难理由。因为如果瞎子没有向导是危险的，那么他竟还想做别人的向导，就更危险得多。\par

那么伯多禄说什么？他没有说\Quote{祢说的这话是什么意思？} 而是装作充满晦涩似的提出问题。他也没有说\Quote{祢为什么说了违反法律的话？} 因为他害怕，怕被认为起了反感，而是断言它晦涩。然而，它并不晦涩，而是他起了反感，这是明显的，因为它没有任何晦涩之处。\par

因此祂也责斥他们说：\Quote{连你们也不明白吗？}\BibleRef{玛 15:16}因为对群众而言，也许他们并不太明白这话；但你们却是那被绊倒的人。因此，起初他们好像是为法利塞人求问，想要被告知；但当他们听见祂宣布一个大威胁，说：\Quote{凡我天父没有栽种的植物，必要被拔除}，以及\Quote{他们是瞎眼的领路者}时，他们就沉默了。但伯多禄总是热忱，即使如此也不能忍耐保持沉默，却说：\Quote{请给我们讲解这个比喻。}\BibleRef{玛 15:15}\par

那么基督说什么？祂以严厉的斥责回答说：\Quote{连你们也不明白吗？你们还不了解吗？}\par

祂说这些话并责备他们，是为了消除他们的偏见；但祂并未停留于此，还加上其他话说：\Quote{凡入口的，进入肚腹，然后排泄到茅厕里；但那从口里出来的，是从心里发出来的，才污秽人。因为从心里发出来的，有恶念、凶杀、奸淫、邪淫、盗窃、亵渎、妄证、毁谤：这些都是污秽人的；至于不洗手吃饭，并不污秽人。}\BibleRef{玛 15:17}\par

你们看见祂如何严厉地对待他们，并且是以斥责的方式吗？\par

然后祂藉我们共同的本性来确认祂的话，并着眼于他们的治疗。因为当祂说：\Quote{它进入肚腹，然后排泄到茅厕里}时，祂还是按照犹太人的浅薄见解回答。因为祂说：\Quote{它不存留，而是出去：}即使它存留，也不会使人不洁。但那时他们还不能听从这话。\par

可以注意到，正因为此，立法者只允许一段足够的时间，即它留在人体内的时间，一旦出去，就不再算为不洁。例如，祂命你在晚上洗身，然后成为洁净；这是衡量消化和排泄的时间。\BibleRef{肋 11:24}但心里的事，祂说，存留在内，当它们出去时就污秽人，而不仅仅是存留的时候。祂首先列出恶念，这是属于犹太人的一种东西；祂尚未从事物的本性来反驳，而是从肚腹和心各自的产生方式，以及从一种存留、另一种不存留的事实；一种从外面进入，又向外出去，而另一种在里面滋生，出去时就污秽人，并且出去时更是如此。因为正如我所说，他们还不能以应有的严格来学习这些事。\par

但玛尔谷说，祂说这话是\Quote{洁净食物}。然而祂并未明说，也完全没有说：\Quote{但吃这样那样的食物并不污秽人}，因为他们也不能忍受祂这样明确地告诉他们。因此祂的结论是：\Quote{但不洗手吃饭并不污秽人。}\BibleRef{玛 15:20}\par

5. 那么让我们学习什么污秽人；让我们学习，并逃避它们。因为在教会中，我们看见这样的习俗在众人中盛行：人们致力于穿着干净的衣服前来，并洗净双手；但如何向天主呈献一颗洁净的灵魂，他们却不加理会。\par

我说这话，并不是禁止他们洗手或漱口；而是愿意人如此洗涤，这洗涤是相宜的，不仅用水，而且以一切德行代替水。因为口的污秽是恶言、亵渎、辱骂、愤怒的言语、污秽的谈话、嬉笑、戏谑：如果你自觉没有说出其中任何一样，也没有被这种污秽玷污，就大胆地前来；但如果你无数次沾染了这些污渍，为何徒劳地洗你的舌头——用水，却带着如此致命有害的污秽呢？告诉我，如果你的手上沾了粪和污泥，你还会大胆祈祷吗？绝不会。然而这并无伤害；而那种却是毁灭。那么，你在不同的事上如此恭敬，在禁止的事上却掉以轻心吗？\par

那么，难道我们不应该祈祷吗？有人问。实在应该祈祷，但不可在污秽中，并带着那种污泥。\par

\Quote{那么，如果我一时不慎跌倒了？}有人说。洁净你自己。\Quote{如何，以什么方式？}哭泣、呻吟、施舍、向受冒犯的人道歉、藉此与他修和、擦净你的舌头，免得你更严重地激怒天主。因为如果有人两手沾满了粪，然后抓住你的脚恳求你，你非但不会听他，反而会一脚踢开他；那么，你怎么敢这样来到天主面前呢？事实上，舌头是祈祷者的手，我们藉它抓住天主的膝盖。所以不要玷污它，免得祂也对你说：\Quote{你们尽管多祈祷，我仍不俯听。}\BibleRef{依 1:15}是的，\Quote{生死在舌头的权下；}\BibleRef{箴 18:21}并且\Quote{凭你的话，你将被定无罪；凭你的话，你将被定罪。}\BibleRef{玛 12:37}\par

我命令你们要守护自己的舌头，胜过守护眼中的瞳仁。舌头是君王的坐骑。你若给它套上缰绳，教它步履有序，君王便会安坐其上；但你若任它无缰奔突、肆意跳跃，它就成了魔鬼和恶灵所骑的畜生。你刚从自己妻子的身旁过来，尚且不敢祈祷——尽管这并无丝毫可责之处；而你在辱骂和侮辱之后（这带来的不亚于地狱），未经彻底洗净，就举起双手？你怎么不战栗？告诉我。你没有听见保禄说：\Quote{婚姻当受尊重，床笫是无污的吗？}\BibleRef{希 13:4}但若从无污的床笫起来，你尚且不敢近前祈祷，那么你从魔鬼的床笫过来，又怎能呼求那可怕而可畏的名字？因为真正的魔鬼之床，就是沉溺于侮辱和辱骂。愤怒如同一个邪恶的奸夫，以极大的快乐与我们调情，将致命的种子播入我们内，使我们生出魔鬼般的仇恨，所做的一切都与婚姻相反。婚姻使二人成为一体，而愤怒却将合一的人分裂成许多部分，切割并粉碎灵魂。\par

为使你们能满怀信心地亲近天主，当愤怒临到时，不要接纳它，不要容它与你们同在，而要像赶走疯狗一样驱赶它。\par

因为保禄也如此命令；他说：\Quote{举起圣洁的手，没有愤怒和争议。}\BibleRef{弟前 2:8}所以，不要侮辱你的舌头，它若失去了应有的信心，又如何为你祈求呢？要以温柔和谦卑装饰它，使它配得上那被祈求的天主，以祝福和大量的施舍充满它。因为甚至可以用言语施行施舍。\Quote{因为言语胜过礼物，}\BibleRef{德 18:16}并\Quote{要以温和答复穷人。}\BibleRef{德 4:8}在其余的时间里，也要以讲述天主的法律来装饰它：\Quote{是的，愿你们的一切交谈都在至高者的法律中。}\BibleRef{德 9:15}\par

我们这样装饰自己之后，让我们来到我们的君王面前，俯伏在他膝前，不仅用身体，也用心灵。让我们思想我们亲近的是谁，是为了谁，以及我们要成就什么。我们正亲近天主，色辣芬看见他便转过脸去，无法承受他的光辉；大地一见他便战栗。我们亲近天主，\Quote{他住在不可接近的光中。}\BibleRef{弟前 6:16}我们亲近他，是为了从地狱中得救，为了罪的赦免，为了逃脱那些难以忍受的惩罚，为了达到天堂及那里的美善。让我们，我说，在他面前俯伏，身体和心灵一起，好叫他扶起我们这些倒地的人；让我们以全然的温柔和良善交谈。\par

有人会说：谁如此可怜和悲惨，竟会在祈祷中不变得温柔呢？那以诅咒祈祷、充满愤怒、并呼喊反对仇敌的人。\par

6. 不，你若想控告，就控告你自己吧。你若想磨砺并磨快你的舌头，就让它攻击你自己的罪过。不要述说别人对你行了什么恶，而要述说你自己对自己行了什么恶；因为这才是真正的恶；因为除非你伤害自己，否则没有人真正能伤害你。因此，你若想反对那些错待你的人，就先反对你自己吧；没有人会阻拦；因为如果你去法庭控告别人，你只会带着更大的伤害离去。\par

你究竟能提到什么真正的伤害呢？某人侮辱了你，用暴力抢夺了你，使你陷入危险？不，这不是受伤害，如果我们清醒，这反而是最大的益处；受伤的是行这事的人，而不是承受的人。这是我们一切邪恶的首要原因：我们根本不知道谁是受伤者，谁是害人者。因为如果我们明白这一点，我们就永远不会伤害自己，就不会祈祷反对别人，因为我们知道不可能从别人那里受苦。因为不是被抢夺，而是抢夺，才是恶。因此，如果你抢夺了，就控告你自己；但如果你被抢夺了，反而要为抢夺你的人祈祷，因为他给了你最大的好处。虽然行者的意图并非如此，但你若勇敢承受，便得到了最大的益处。因为对于他，人和天主的法律都宣告他是有祸的；但你，受损害的一方，他们为你加冕，并宣扬你的赞美。\par

就如一个患热病的人，若从别人那里强行夺走一个盛水的器皿，满足了他有害的欲望，我们不会说被夺者受了伤害，而是说抢夺者；因为他加剧了自己的热病，使疾病更重。同样，我命令你们也要这样思考那爱财富和金钱的人。因为他也患有比前者更严重的热病，通过这种掠夺，他煽起了自己内心的火焰。\par

又如一个疯子从某人那里夺走一把剑，并自杀，这又是谁受伤害呢？是被夺者，还是抢夺者？显然，是抢夺者。\par

那么，在强占财产的事上，让我们也做出同样的判断。因为财富之于贪财者，正如剑之于疯子；甚至更糟。因为疯子拿剑刺穿自己后，既脱离了疯狂，也不会再受第二刀；但爱财者每天受比这更重的万般创伤，却不脱离自己的疯狂，反而使其更加剧；他受的伤越多，就越引发其他更严重的打击。\par

因此，默想这些事，让我们逃避这把剑；逃避这疯狂；虽然晚了，让我们变得节制。因为这一德性也应被称为节制，不亚于所有人通常所称的节制。因为后者只需对抗一种情欲的统治，而在这里，我们必须驾驭许多种情欲，而且是各种各样的。\par

诚然，没有比财富的奴隶更愚蠢的了。他以为自己在征服，实则被征服；他以为自己是主人，实则是奴隶，给自己套上枷锁，还沾沾自喜；让野兽更加凶猛，他却高兴；成为俘虏，却自傲、雀跃；看见一条狂犬扑向他的灵魂，本该捆绑它、以饥饿削弱它，他却反而供应它丰盛的食物，让它更猛烈地扑向他，更加可怖。\par

既然思考了这一切，让我们解开束缚，杀死这怪兽，驱除这疾病，赶走这疯狂，好能享受宁静而纯净的健康，带着极大的喜悦驶入平静的港湾，获得永恒的福乐；愿我们都能因我们的主耶稣基督的恩宠与对人的慈爱而获得，愿光荣与权能归于祂，从现在直到永远，世世无穷。阿们。\par

\SourceRecord{Translated by George Prevost and revised by M.B. Riddle.；Nicene and Post-Nicene Fathers, First Series；Vol. 10.；Edited by Philip Schaff.；Buffalo, NY: Christian Literature Publishing Co.,；1888.；<http://www.newadvent.org/fathers/200151.htm>.}

% structure: p0001 source=h1 target=section reason=page-title

\section{论\WorkTitle{玛窦福音}第五十二篇讲道}

\BibleRef{玛 15:21-22}\par

\BibleQuote{耶稣从那里出去，退到提洛和漆冬一带。看，有一个客纳罕妇人从那一带出来，喊着说：\InnerQuote{主，达味之子，可怜我吧！我的女儿被魔鬼纠缠得好苦！}}\par

但马尔谷说，祂进了屋子，却\BibleQuote{不能隐藏}\BibleRef{谷 7:24}。祂为什么来到这些地方呢？当祂使门徒从食物的规条中得自由以后，便循序渐进，也为外邦人敞开大门；正如伯多禄先被指示废除这法律，然后被派往科尔乃略那里。\par

但若有人问：祂既对门徒说\BibleQuote{不要走外邦人的路}，自己却接待她？首先，我们可以这样回答：祂命令门徒的，并不等于祂自己也受约束；其次，祂离开并非为了宣讲；马尔谷也暗示这点，说祂甚至隐藏自己，却未被隐瞒。\par

因为祂不先赶到他们那里，是祂行事常理的步骤；但若赶走前来投靠的人，则有损祂对人的慈爱。因为若逃跑的应当追赶，那么追赶的就更不应躲避。\par

无论如何，请看这妇人是多么配得一切恩惠。她甚至不敢上耶路撒冷去，因为她畏惧并自觉不配。若非如此，她必会前去，这从她当前的恳切和她从本乡出来便可得知。\par

也有人以寓意解释：当基督从犹太出来时，教会才敢接近祂，她也从自己的边界出来。经上说：\BibleQuote{忘记你的民族和你的父家。}因为基督出了自己的边界，妇人也出了自己的边界，这样他们才得以相遇。因此祂说：\Quote{看，有一个客纳罕妇人从那一带出来。}\par

圣史提到这妇人，是为了彰显她的奇妙事迹，更加赞美她。因为你听到客纳罕妇人时，应记起那些邪恶的民族，他们甚至颠覆了自然律的基本法则。想到这些，还要思考基督来临的能力。因为那些被驱逐以免败坏犹太人的外邦人，竟表现得比犹太人更善良，甚至从自己的边界出来接近基督；而犹太人却在基督来到他们那里时驱赶祂。\par

2. 她来到祂跟前，没有说别的，只说\BibleQuote{可怜我吧}，她的喊声吸引了众多围观者。这确实是一幅令人怜悯的景象：一个妇女在极大的痛苦中高声哭喊，而且是一位母亲，为女儿恳求，女儿又处于如此悲惨的境地；她甚至不敢将附魔的女儿带到主面前，而是让她躺在家中，自己来恳求。\par

她只述说自己的痛苦，并不多说；她没有像那位大臣那样把医生拉到家里，说：\BibleQuote{求祢来按手在她身上}，\BibleQuote{趁我的孩子还没有死，求祢下去}。\par

她描述了自己的灾难和疾病的严重性，恳求主的仁慈，大声呼喊；她不说\BibleQuote{可怜我的女儿}，而是说\BibleQuote{可怜我吧}。因为女儿对自己的疾病无知觉，而我却承受着她无数的痛苦；我的疾病有意识，我的疯狂能自知。\par

2. \BibleQuote{耶稣却一言不答。}\BibleRef{玛 15:23}\par

这是何等新奇古怪的事？对顽固的犹太人，祂引导他们，辱骂时祂仍劝告他们，试探时祂不遣散他们；但对这跑来恳求祂、在法律和先知中未受教育、却表现出如此大虔敬的妇人，祂竟不屑回答。\par

这事与传闻如此相反，谁不会因此跌倒呢？因为他们曾听说祂走遍乡村治病，而这妇人来到祂跟前，祂却完全拒绝她。她的痛苦和她为女儿的恳求，谁不被感动呢？她前来不是自认为配得，也不是要求应得的，而是恳求怜悯，只是陈述自己的苦难；然而她连回答都不配得到。\par

也许许多听众会跌倒，但她却没有。我说听众，何必呢？我想连门徒们也会为这妇人的痛苦多少感到触动，大大困扰，失去信心。\par

尽管如此，他们在这困扰中也不敢说\Quote{请赐予她这个恩典}，而是\Quote{门徒前来恳求祂说：\InnerQuote{请打发她走吧，因为她在我们后面喊叫。}}\par

但基督说：\BibleQuote{我被派遣，只是为了以色列家迷失的羊。}\BibleRef{玛 15:24}\par

那么，妇人听了这话以后怎样呢？她沉默了吗？放弃了吗？减弱了恳切吗？绝不，她反而更加坚持。但我们不然，一旦得不到便放弃；其实这反而应当促使我们更加迫切。\par

然而，谁能不为所说的话感到困惑呢？祂的沉默已足以使她绝望，而祂的回答更是如此。因为不仅她自己，连为她求情的人也完全困惑，并听见这件事甚至不可能办到，这足以使她陷入难以言表的困惑。\par

然而，妇人并未困惑，反而见她的辩护人毫无成效，便以美好的厚颜让自己变得不知羞耻。\par

因为在这之前，她连露面都不敢（因为经上记着说：她\Quote{在我们后面喊叫}），当人可以预料她应在绝望中走得更远时，她却在那时候走上前来，朝拜说：\Quote{主，帮助我。}\BibleRef{玛 15:25}\par

这是怎么回事，妇人啊？难道你比宗徒们更有勇气？更有力量？她说：\Quote{勇气和力量，决不是；相反，我甚至充满羞愧。然而我还是把我的无耻当作恳求；祂会尊重我的勇气。}这是什么呢？你没有听过祂说：\Quote{我除了被派往以色列家迷失的羊那里，别处不去？}她说：\Quote{我听过，但祂本身就是主。}因此她也没有说\Quote{恳求并哀求}，而是说\Quote{帮助我}。\par

3. 那么基督说什么呢？即便如此祂仍不满足，反而再次加深她的困惑，说：\par

\Quote{拿儿女的饼扔给小狗，是不对的。}\BibleRef{玛 15:26}\par

祂一开口对她说话，就比沉默时更严厉地打击了她。祂不再将原因推给别人，也不说\Quote{我没有被派遣}，而是她越迫切恳求，祂也就越催促祂的拒绝。祂不再称他们为\Quote{羊}，而称\Quote{儿女}，称她为\Quote{狗}。\par

那么妇人说什么呢？她从祂自己的话中构造了她的请求。她说：\Quote{纵使我是一条狗，我也不是外人。}\par

基督公平地说：\Quote{我为审判来到这世界。}\BibleRef{若 9:32} 这妇人操练了高度的自制，并显示了全部的忍耐和信德，而且是在受到侮辱时；而他们，虽受奉承和尊荣，却报以相反的态度。\par

她说：\Quote{那食物对儿女是必需的，我也知道；然而我作为狗，也未被禁止。因为若接受是不合法的，那么即使碎屑也不可分享；但如果即使量少，他们也应分享，那么我作为狗，也未被禁止；甚至正因我是狗，我肯定有份。}\par

基督故意拖延她，因为祂知道她会这样说；祂拒绝恩赐，是为了彰显她高度的自制。\par

因为如果祂无意赐予，后来祂也不会赐予，也不会再堵住她的口。但正如祂在百夫长的事上所说：\Quote{我去治好他}（\BibleRef{玛 8:7}），让我们学习那人的敬畏，并听到他说：\Quote{主，我不堪当祢到我舍下来}（\BibleRef{玛 8:8}）；又如祂在患血漏的妇人身上所说：\Quote{我觉得有能力从我身上出去了}（\BibleRef{路 8:46}），以彰显她的信德；又如撒玛黎雅妇人的事，以表明即使受责备她也不停止（\BibleRef{若 4:18}）；同样在这里，祂不愿妇人如此伟大的德行被隐藏。所以祂的话不是出于侮辱，而是为了召唤她，并显明她内在的宝藏。\par

但我请你，连同她的信德也看她的谦卑。因为祂称犹太人为\Quote{儿女}，但她不满足于此，甚至称他们为\Quote{主人}；她毫不因别人的赞美而悲伤。\par

她说：\Quote{可是小狗也吃主人桌子上掉下来的碎屑。}\BibleRef{玛 15:27}\par

你看见这妇人的智慧吗？她甚至不敢反驳一句，也不因别人的赞美而刺痛，也不因责备而恼怒。你看见她的坚毅吗？祂说\Quote{不合适}，她说\Quote{主，是的}；祂称他们为\Quote{儿女}，她称\Quote{主人}；祂用了狗的称呼，她则加上了狗的动作。你看见这妇人的谦卑吗？\par

听听犹太人的骄傲言语：\Quote{我们是亚巴郎的后裔，从来没有给任何人做过奴仆}（\BibleRef{若 8:33}）；以及\Quote{我们是生自天主的}（\BibleRef{若 8:41}）。但这妇人却不这样，她称自己为狗，称他们为主人；因此她成为了儿女。那么基督说什么呢？\Quote{妇人，你的信德真大！}（\BibleRef{玛 15:28}）\par

是的，祂曾拖延她，正是为了大声宣告这句话，为了给妇人加冕。\par

\Quote{照你的愿望成就吧。}祂所说的话意思是这样：\Quote{你的信德确实能成就比这些更大的事；然而，照你的愿望成就吧。}\par

这类似于那声音说：\Quote{让天存在，就有了天。}\BibleRef{创 1:3}\par

\Quote{她的女儿从那时候就痊愈了。}\par

你看见这妇人也为她女儿的治疗贡献了不少吗？为此基督没有说\Quote{让你的小女儿痊愈}，而是说\Quote{你的信德真大，照你的愿望成就吧}；为教导你，这些话不是随便说的，也不是谄媚之言，而是她的信德力量伟大。\par

然而，事件的结果提供了确切的考验和证明。于是她的女儿立刻痊愈了。\par

但我请你注意，当宗徒们失败、没有成功时，这妇人却成功了。恒心祈祷是多么伟大的事啊！是的，祂甚至更愿意由我们这些有罪的人为属于我们的人向祂恳求，而不愿由别人代求。然而他们拥有更多的说话自由，但她表现了极大的忍耐。\par

通过结果，祂也向门徒们为迟延辩护，表明祂没有同意他们的请求是有理由的。\par

\Quote{\Person{耶稣}从那里出来，来到\Place{加里肋亚海}边，上了山，坐在那里。有许多群众来到祂跟前，带着瘸子、瞎子、残废的、哑巴，都放在祂脚前；祂治好了他们。以致群众看见哑巴说话，残废的痊愈，瘸子行走，瞎子看见，都很惊奇，并光荣了\Place{以色列}的天主。}\par

现在祂亲自四处行走，又坐下来等候病人，并让瘸子被带到山上。他们不再只是触摸祂的衣服，而是更上一层，被放在祂脚前：他们双倍地表现了信德，第一，虽然瘸腿却上了山；第二，不想要别的，只求被放在祂脚前。\par

神奇且非凡的是，看见那些被抬来的行走，瞎子不需要任何人引领。的确，被治愈的人群之多以及治疗的便捷都使众人惊讶。\par

你看见了吗？祂医治那妇人时如此拖延，但这些却立刻治愈？并非因为这些比她更好，而是因为她比他们更有信德。因此，在她身上，祂推迟和延迟，以显明她的恒心；对这些人，祂立即赐予恩赐，堵住不信的犹太人的口，并切断他们的一切借口。因为一个人领受的恩惠越大，若他麻木不仁，且这尊荣不能使他变得更好，他就越应受惩罚。因此，你看到富人也因邪恶而比穷人受更重的惩罚，因为他们甚至在顺境中也未软化。不要对我说他们施舍了。因为如果他们不按自己的财产比例施舍，他们也不能逃脱；我们的施舍不是按礼物的分量，而是按我们心地的宽厚来评判。但如果这些人受惩罚，那么那些贪图无用之事的人更甚，他们建造两三层楼房却忽视饥饿的人；他们贪图贪婪却忽视施舍。\par

5. 既然论述已涉及施舍，那么，让我们今天再次继续这个三天前关于仁爱的论证，当时我尚未完成。你们还记得，当我最近谈到你们对鞋子的虚荣、那些无谓的烦恼以及年轻人的奢靡时，正是从施舍的论述转到了对你们的责备。那时提出的问题是什么？即施舍是一种艺术，其工场在天堂，其教师不是人，而是天主。然后我们探讨什么是艺术，什么不是艺术，谈到了徒劳的劳动和邪恶的计谋，其中也提到了关于人鞋子的艺术。\par

你们还记得吗？那么，让我们今天也继续我们当时的论述，并表明施舍是一种艺术，且优于所有艺术。因为如果艺术的特性是产生有用的结果，并且没有什么比施舍更有用，那么显然这既是一种艺术，又优于所有艺术。因为它为我们制造的并非鞋子，也不编织衣服，也不建造泥土的房子；而是获得永生，从死亡的手中拯救我们，在今生和来世都使我们荣耀，并建造天堂的住所和永恒的帐幕。\par

这（施舍）不让我们的灯熄灭，也不让我们在婚礼上穿着污秽的衣服出现，而是洗净它们，使它们比雪更洁白。\BibleQuote{因为你们的罪虽似朱红，我将使它们洁白如雪。}\Emph{不要让他们空着肚子回去，免得他们在路上昏倒。}\par

\SourceRecord{Translated by George Prevost and revised by M.B. Riddle.；Nicene and Post-Nicene Fathers, First Series；Vol. 10.；Edited by Philip Schaff.；Buffalo, NY: Christian Literature Publishing Co.,；1888.；<http://www.newadvent.org/fathers/200152.htm>.}

% structure: p0001 source=h1 target=section reason=page-title

\section{\BibleBook{玛窦}{福音}讲道辞第五十三篇}

\BibleRef{玛 15:32} \BibleQuote{耶稣叫了祂的门徒来，说：\InnerQuote{我很怜悯这群众，因为他们同我一起已经三天，也没有吃的；我不愿遣散他们空腹回去，怕他们在路上晕倒。}}\par

在上面，当祂要行这奇迹时，祂先治愈了身体残缺的人，而在这里祂也做了同样的事；从治愈瞎子和瘸子，祂又进行了这个奇迹。\par

但是为什么呢，那时祂的门徒说\Quote{遣散群众}，而现在他们却不这样说；尽管已经过了三天？要么是他们自己此时已有所长进，要么是看到群众没有强烈的饥饿感；因为他们为所发生的事而光荣天主。\par

但请看，在此例中祂也并未立即施行奇迹，而是先召唤他们来。因为那些为治病而来的群众确实不敢求饼；但祂，这位仁慈而眷顾的主，甚至赐给那些没有求的人，并对祂的门徒说：\Quote{我很怜悯，不愿遣散他们空腹回去。}\par

因为唯恐他们会说他们来时带了路上的食物，祂就说：\Quote{他们同我一起已经三天了；}这样即使他们来时带了什么，也都用尽了。正因如此，祂自己未在第一、二天行这事，而是等他们所有的都消耗完了，为了让他们先感到匮乏，从而更热切地接受祂的作为。\par

因此祂说：\Quote{免得他们在路上晕倒；}暗示他们路途遥远，并且已一无所有。\par

\Quote{那么，若你不愿遣散他们空腹回去，为何不行奇迹呢？}藉这个问题和他们的回答，祂要使门徒更加留意，并使他们前来对祂说\Quote{做饼吧}以显示他们的信德。\par

但即便如此，他们也不明白祂提问的动机；因此后来祂如马尔谷所记，对他们说：\BibleQuote{你们的心仍然执拗吗？你们有眼看不见吗？有耳听不见吗？}\BibleRef{谷 8:17}\par

既然不是这样，祂为何要对门徒说话，表示群众配得恩惠，并加上祂自己感到的怜悯呢？\par

但玛窦说，此后祂也斥责他们，说：\BibleQuote{小信德的人哪，你们还不明白，也不记得那五千人的五个饼，和你们收了多少篮子吗？也不记得那四千人的七个饼，和你们收了多少篮子吗？}\BibleRef{玛 16:8}福音作者们如此完全地彼此和谐。\par

那么门徒说什么呢？他们仍在地上爬行，尽管祂做了许多事，为使那奇迹被记在心中；如通过祂的提问、回答、让他们在此服务、并分配篮子；但他们的心境仍是相当不完全的。\par

因此他们也问祂：\Quote{在荒野中我们从哪里得这么多饼呢？}\par

在此前和现在，他们都提到荒野；他们用一种软弱的论证方式如此说话，但即使如此，也藉此使奇迹无可置疑。也就是说，唯恐有人断言（正如我已说过的），他们是从某个邻近村庄获得食物，地点被确认，好使奇迹被相信。为这个目的，祂在荒野中、远离村庄的地方，行了从前和现在的奇迹。\par

门徒们毫不顾及这一切，说：\Quote{在荒野中我们从哪里得这么多饼呢？}因为他们确实以为祂这样说，是打算随后命令他们喂养群众；这是极其愚昧的，因为祂曾带着这个意图说过，而且就在不久前，\BibleQuote{你们给他们吃吧。}\BibleRef{玛 14:16}祂要使他们迫切需要恳求祂。\par

但现在祂没有说\Quote{你们给他们吃吧}，而是说什么呢？\Quote{我很怜悯他们，不愿遣散他们空腹回去；}使门徒更接近，并更激发他们，赐给他们更清楚的眼光，好向祂求这些事。因为这些实在是那表明祂有能力不遣散他们空腹回去者的言语；是彰显祂权柄者的言语。因为\Quote{我不愿}这个说法，表示祂有此意向。\par

2. 但由于他们仍只谈论群众、地方和荒野（因为据说：\Quote{在荒野中我们哪里能有这么多饼，来填饱这么一大群人？}），甚至仍未理解祂所说的，祂便继续尽祂自己的部分，对他们说：\par

\BibleQuote{你们有几块饼？}他们回答说：\BibleQuote{七个，还有几条小鱼。}\BibleRef{玛 15:34}\par

他们不再像先前那样说：\BibleQuote{但这点东西在这么多人中间算什么呢？}\BibleRef{若 6:9}所以，虽然他们仍未达至祂的全部意义，但终究逐渐提升了。因为祂也这样，藉此唤醒他们的心思，几乎如先前一样发问，为藉这询问的形式提醒他们已行的事。\par

但正如你由此看见了他们的不完全，也请观察他们精神的严格，并钦佩他们对真理的热爱，他们如何自书时，不隐瞒自己如此大的缺点。因为立刻忘记这刚发生的奇迹，并非小错；因此他们也被斥责。\par

并由此也请考虑他们在另一事上的严格，他们如何克服了食欲；如何训练有素，不重视饮食。因为在荒野中逗留三天，他们只有七个饼。\par

现在所有其余的事，祂都如先前一样做；祂既使他们坐在地上，也使得饼在门徒手中增多。\par

因为据说，\BibleQuote{祂命令群众坐在地上。祂拿起七个饼和鱼，祝谢了，擘开，递给门徒，门徒再递给群众。}\par

但当我们来到结尾，有一个不同之处。\par

因为，经上说：\BibleQuote{他们都吃了}，\BibleQuote{并且都饱了；他们把剩下的碎块收集起来，装满了七篮子。吃的人，除了妇女和孩子，约有四千人。} \BibleRef{玛 15:37}\par

但是为什么先前那次，有五千人时，剩下的碎块装满了十二篮子；而这里只有四千人时，却只有七篮子呢？请问，这是什么缘故，什么原因使剩下的较少，而客人并非那么多呢？\par

那么，或许可以这样说：这一次的篮子比先前所用的更大；或者，即使不是这样，为了避免奇迹的相等性再次使他们陷入遗忘，祂通过差异唤醒他们的记忆，好使他们因变化而同时想起两者。因此，在那次，祂使装满碎块的篮子数目与门徒的人数相等；这一次，则使篮子数目与饼的数目相等。这甚至借此指示了祂那不可言喻的能力，以及祂行使权柄时的轻易，因为祂既能以这种方式，也能以另一种方式行这样的奇迹。因为，无论从前还是现在，能够维持精确的数目——从前是五千人，现在是四千人——并且不使剩下的碎块超过一个场合或另一个场合所用的篮子，尽管客人的数目不同，这绝非小事。\par

结局又像从前一样。因为正如那时祂离开群众，乘船退去，现在也是如此；若望也记载了这事。\BibleRef{若 6:17} 因为没有什么神迹像增饼的奇迹那样促使他们跟随祂；他们不仅愿意跟随祂，还愿意立祂为王。\BibleRef{若 6:15} 为了避免任何篡夺王权的嫌疑，祂在行了这奇事后急忙离开；祂甚至不步行离去，免得他们跟随祂，而是上了船。\par

\BibleQuote{于是祂遣散了群众，} 经上这样说，\BibleQuote{就上船，来到玛加丹地区。}\par

3. \BibleQuote{法利塞人和撒杜塞人前来，要祂显给他们一个从天上来的征兆。祂却对他们说：到了晚上，你们说：天色发红，天要晴；早晨则说：今天天色发红而阴沉，要下雨。你们能分辨天色，却不能分辨时代的征兆吗？邪恶淫乱的世代寻求征兆，但除了约纳先知的征兆外，再没有征兆给它。祂就离开他们走了。}\par

但是马尔谷说，当他们来到祂跟前，与祂辩论时，\BibleQuote{祂从心里深深地叹息，说：为什么这世代寻求征兆？}\BibleRef{谷 8:12}\par

应该说，他们的询问是应该发怒和甚为不悦的；然而，仁慈而眷顾的那位并不发怒，反而怜悯他们，哀悼他们如同患了不治之症，在如此充分地显示了祂的能力之后，仍然试探祂。\par

因为他们寻求征兆不是为了相信，而是要抓住祂的把柄。倘若他们怀着相信的心来到祂跟前，祂必会赐给他们。因为祂对那妇人说：\BibleQuote{拿儿女的饼扔给小狗吃，是不对的} \BibleRef{玛 15:26} 但后来还是给了，更何况对这些人的慷慨呢？\par

可是，因为他们不是为相信而寻求，所以祂也称他们为伪善者，因为在别处他们口是心非。诚然，倘若他们相信，他们甚至不会问。从另一件事也可看出他们不信：当受到指责和揭露时，他们不留在祂身边，也不说：\Quote{我们无知，愿意学习。}\par

但他们所求的是什么样的天上征兆呢？是要祂命令太阳，或抑制月亮，或降下雷霆，或改变空气，或类似的事。\par

对此祂说了什么呢？\Quote{你们能分辨天色，却不能分辨时代的征兆吗？}请看祂的温柔和节制。因为祂不仅像先前那样只是拒绝，说：\Quote{不会再给他们征兆，}而且说明了不给的原因，尽管他们并非为求教而问。\par

那么原因是什么呢？祂说：\Quote{正如在天上，一件事是风暴的征兆，另一件是晴天的征兆，没有人见了恶劣天气的征兆还会寻求平静，同样在平静和晴天也不会寻求风暴；所以你们对我也应这样想。因为我来的这个现今时期，与将来的时期不同。现在需要的是这些在地上的征兆，那些天上的征兆则保留到那时。现在我来像一位医生，那时我将在那里如同审判官；现在我来寻找迷失的，那时我要查账。因此，我隐藏地到来，但那时将非常公开，卷起天穹，遮掩太阳，不让月亮发光。那时\BibleQuote{诸天的权势将要动摇，}\BibleRef{玛 24:29} \BibleQuote{而我显现的来临将如同闪电，同时显给众人。}\BibleRef{玛 24:27} 但现在不是这些征兆的时候；因为我来是要受死，并承受一切极端痛苦。}\par

你们没有听先知说：\BibleQuote{祂不争辩，也不呼喊，街上也听不到祂的声音吗？}\BibleRef{依 42:2} 另一位又说：\BibleQuote{祂像雨落在羊毛上？}\par

如果有人提到法郎时期的征兆，那时需要从敌人中解救出来，一切都按部就班发生。但来到朋友中间的那位，不需要那些征兆。\par

\Quote{况且，当小的不被人相信时，我怎能显大的征兆呢？}我说小，是指在显扬方面，因为后者在能力上比前者大得多。因为有什么能与赦罪、复活死人、驱逐魔鬼、创造身体、以及妥善安排一切其他事相比呢？\par

但请看他们硬化的心：被告知\Quote{不给他们任何征兆，除了\Person{约纳}先知的征兆}之后，他们并不询问。然而，既知道这位先知和他所遭遇的一切，又再次听到这样的话，他们原当探询并明了这话的意思；但，如我所说，他们这些行为中毫无求知的渴望。为此，\Quote{祂也离开了他们，便走了。}\par

4. \Quote{当祂的门徒，如经上记载，来到对岸，他们忘了带饼。耶稣便对他们说：你们应当谨慎，提防法利塞人和撒杜塞人的酵母。} \BibleRef{玛 16:5}\par

为什么祂不说得明明白白：\Quote{当心他们的教导}？祂的意愿是要提醒他们所经历过的事，因为祂知道他们已经忘记。但若立即指责他们，似乎没有合理依据；而借着他们自己的事来引出话题，再斥责他们，会使责备更可接受。\Quote{为何祂当时不责备他们，当他们说\InnerQuote{我们从哪里在荒野得到这么多饼呢？}因为那时似乎正是适宜说出这里的话的时机。} 这样祂就不会显得仓促地行奇迹。此外，祂不愿在群众面前责备他们，也不愿在他们面前寻求荣耀。而现在，这指责也更有理由，因为重复行奇迹之后他们仍然如此固守己见。\par

因此祂也行另一个奇迹，并且那时才责备他们；我是说，祂提出了他们心中所思虑的。但他们所思虑的是什么呢？\Quote{因为，如经上记载，我们没有带饼。}因为他们那时仍对犹太人的洁净礼和食物的规条充满恐惧。\par

因此，祂在各方面甚至严厉地责备他们，说：\Quote{小信德的人哪，你们为什么在内心思虑没有带饼呢？你们还不明白，也不理解吗？你们的心还是硬了吗？你们有眼看不见吗？有耳听不见吗？\BibleRef{谷 8:17}你们不记得那五千人的五个饼和你们取了多少篮子吗？也不记得那四千人的七个饼和你们取了多少篮子吗？\BibleRef{玛 16:9}}\par

你看到极度的不悦吗？因为祂在其他地方似乎从未如此责备过他们。那么祂为何这样做呢？为的是再次消除他们对食物的偏见。我是说，基于这个意图，当时祂只是说\Quote{你们还不明白，也不理解吗？}而在这里，带着强烈的斥责，祂说\Quote{小信德的人哪。}\par

因为并非处处宽厚都是好事。正如祂曾允许他们自由发言，祂也责备他们，通过这种多样方式为他们的救恩作准备。同时请注意祂的责备是多么强烈，而祂的温和也是如此。因为祂几乎是为自己严厉的责备向他们辩解，祂说：\Quote{你们还不思想那五个饼，和你们取了多少篮子；以及那七个饼，和你们取了多少篮子吗？}为此祂还列出了数目，包括用餐的人数以及碎块的数目，同时使他们想起过去，并使他们更加专注于未来。\par

为了教导你祂的责备有多大能力，以及它如何唤醒了他们昏沉的心，请听福音作者所说的话。耶稣没有再多说，只是责备了他们，并且只加了这句话：\Quote{你们怎么不明白，我不是对你们说饼的事要你们防备，而是要防备法利塞人和撒杜塞人的酵母；}接着祂继续说道：\Quote{那时他们才明白，祂不是叫他们防备饼的酵母，而是防备法利塞人和撒杜塞人的教导。} \BibleRef{玛 14:12}尽管祂并未说出这解释。\par

请看祂的责备成就了多少好处。因为它既使他们脱离犹太人的规条，又当他们懈怠时，使他们更加谨慎，并使他们脱离无信德的状态；\par

\SourceRecord{Translated by George Prevost and revised by M.B. Riddle.；Nicene and Post-Nicene Fathers, First Series；Vol. 10.；Edited by Philip Schaff.；Buffalo, NY: Christian Literature Publishing Co.,；1888.；<http://www.newadvent.org/fathers/200153.htm>.}

% structure: p0001 source=h1 target=section reason=page-title

\section{《玛窦福音》讲道集 54}

\BibleQuote{耶稣来到了斐理伯的凯撒勒雅境内，就问门徒说：\InnerQuote{人们说人子是谁？}}\BibleRef{玛 14:13}\par

为什么他提到这座城的创建者？因为还有另一座凯撒勒雅，即斯特拉托的凯撒勒雅。但他不是在那一座，而是在这一座向他们发问，带领他们远离犹太人，好使他们摆脱一切恐惧，能坦然说出心中所想的一切。\par

为什么他不立刻问他们自己的意见，而是先问众人的意见？为的是当他们说出众人的意见后，再被问到\InnerQuote{但你们说我是谁？}，借着提问的方式，他们能被引导到更高超的观念，而不会陷于与群众同样的低级看法。因此他并非在传道之初就问他们，而是在行了许多奇迹，并与他们谈论了许多高深的道理，提供了如此多关于他天主性和与圣父合一的清晰证据之后，才提出这个问题。\par

他没有说：\Quote{经师和法利塞人说我是谁？}尽管这些人经常来到他面前与他交谈；而是说：\Quote{人们说人子是谁？}探询人民无偏见的看法。因为虽然这种看法远低于应有的水平，却没有恶意，而另一种则充满了许多邪恶。\par

为了表明他多么渴望他的\OriginalTerm{救恩计划}{Economy}被承认，他说：\InnerQuote{人子}，以此表明他的天主性，他在许多其他地方也这样做。因为他说：\BibleQuote{没有人上过天，除了那自天降下而仍在天上的人子。}\BibleRef{若 3:13}又说：\BibleQuote{但当你们看见人子升到祂以前所在之处。}\BibleRef{若 6:62}\par

然后，当他们说：\BibleQuote{有人说是洗者若翰，有人说是厄里亚，有人说是耶肋米亚，或先知中的一位}\BibleRef{玛 16:14}，说出了他们错误的看法，他接着又说：\BibleQuote{但你们说我是谁？}\BibleRef{玛 16:15}借着第二次提问，他召唤他们对他持有更高的想象，并表明他们先前的判断与他的尊荣相去甚远。因此他向他们寻求另一个判断，并再次发问，以免他们附和群众。群众因看到他的奇迹超越人性，以为他是一个人，不过是一个复活后显现的人，正如黑落德也认为的那样。\BibleRef{玛 14:2}但他为了将他们从这种观念中带出来，说：\InnerQuote{但你们说我是谁？}即\InnerQuote{你们这些常与我同在，看见我行奇迹，并且自己也借着我行了许多大能之事的人。}\par

2. 那么，宗徒的口、永远热忱的、宗徒歌咏团的领袖伯多禄说了什么？当众人都被问到时，他回答。而当祂问众人的意见时，所有人都回答了问题；但当祂问他们自己的意见时，伯多禄抢先一步，说：\BibleQuote{你是基督，永生天主之子。}\BibleRef{玛 16:16}\par

基督怎么说？\BibleQuote{约纳的儿子西满，你是有福的，因为不是血和肉启示了你。}\BibleRef{玛 16:17}\par

然而，除非他正确宣认祂为圣父所生，否则这并非启示的作为；若他认为我们的主是众人之一，他的话就不配得到祝福。因为此前他们也曾说：\BibleQuote{你真是天主子}\BibleRef{玛 14:33}，我指的是风暴后船上那些人，他们看见了并说了这话，却没有被祝福，当然他们说的是真话。因为他们所宣认的并非伯多禄所宣认的那种父子关系，他们以为祂是真天主子如同众人中的一位，虽然特别超越众人，却不是同体的。\par

纳塔乃耳也说：\BibleQuote{辣彼，你是天主子，你是以色列的君王}\BibleRef{若 1:49}，他非但没有被祝福，反而受到祂的责备，因为他所说的远不及真理。祂至少回答说：\BibleQuote{因为我告诉你说，我看见你在无花果树下，你就信吗？你要看见比这更大的事。}\BibleRef{若 1:50}\par

那么为什么这个人被祝福？因为他承认祂是真正的儿子。因此你看，在先前那些例子中祂没有说这样的话，而在此处祂也指明了是谁启示的。也就是说，免得他的话在众人看来（因为他是个热爱基督的人）是出于友情和奉承，以及一种偏爱祂的态度，祂就将使之在灵魂中回响的那位摆出来，让你知道伯多禄虽然说了，却是圣父提示的，并使你相信这话不再是人的意见，而是神圣的教训。\par

为什么祂自己不宣布，不说\InnerQuote{我是基督}，而是通过提问来确定这一点，引导他们去承认？因为这样做既更适合祂，在那时也是必要的，并且更能吸引他们相信所说的话。\par

你看见圣父如何启示圣子，圣子如何启示圣父吗？因为祂说：\BibleQuote{除了圣子，也没有人认识父，只有圣子愿意启示给谁，谁才能认识。}因此，不可能从别人那里认识圣子，只能从圣父；也不可能从别人那里认识圣父，只能从圣子。所以，由此也显明了他们的尊荣和本体是同一的。\par

3. 基督于是说什么？\Quote{你是若纳的儿子西满，你要叫刻法。}\Quote{既然你宣认了我的父，我也要指出生养你的那位；}几乎是说：\Quote{如同你是若纳的儿子，同样我也是我父的儿子。}\newline{}否则说\Quote{你是若纳的儿子}便是多余；但因他曾说\Quote{天主子}，为指明祂是那一位天主子，如同那一位是若纳的儿子，与生养祂者同体，故此祂加了这句话：\Quote{我对你说：你是伯多禄，在这磐石上，我要建立我的教会；}\BibleRef{玛 16:18}即建在他宣认的信德上。祂以此表明许多人在此即将相信，并振作他的精神，立他为牧者。\Quote{阴间的门不能胜过她。}\Quote{既然不能胜过她，更不能胜过我了。所以你们即将听到我将被出卖和钉死时，不要烦恼。}\par

然后祂又提到另一项尊荣：\Quote{我还要给你天国的钥匙。}\newline{}但这是什么意思？\Quote{我还要给你？}\Quote{如同父给了你认识我，我也要给你。}\par

祂没有说\Quote{我要恳求父}（虽然祂权柄的彰显是伟大的，恩赐的浩瀚是难以言喻的），而是说\Quote{我要给你。}你要给什么？请告诉我。\Quote{天国的钥匙：凡你在地上捆绑的，在天上也要被捆绑；凡你在地上释放的，在天上也要被释放。}\newline{}那么，当祂说\Quote{我要给你}时，\BibleRef{玛 20:23}哪里说\Quote{坐在祂右边和左边不是祂赐予的}呢？\par

你看见祂自己如何引导伯多禄对祂有高超的思想，并启示自己，藉这两项应许暗示祂是天主子吗？因为那些唯独属于天主的事（赦免罪恶、使教会在这般汹涌的浪涛中不可颠覆、使一个渔夫比任何磐石更坚固，尽管全世界与他为敌），祂都应许亲自赐予；正如父对耶肋米亚所说的，要使他成为\Quote{铜墙铁壁}\BibleRef{耶 1:18}；但只对一国，而这人却对全世界。\par

因此，我愿问问那些想要贬低圣子尊威的人：哪样恩赐更大？是父赐给伯多禄的，还是子赐给他的？因为父赐给伯多禄对子的启示；而子赐给他将父和自己播种在全世界，并将天上一切权柄托付给一个必死的人，赐给他钥匙；他使教会扩展到全世界，并宣告她比天更强固。\Quote{天地要过去，但我的话绝不会过去。}\BibleRef{玛 24:35}那么，赐予如此恩赐、成就如此大事的那位，如何更小呢？\par

我说这些，并非分隔父与子的工作（\Quote{万物是藉着祂造的；凡受造的，没有一样不是藉着祂造的}），而是为勒住那些胆敢如此说话的无耻舌头。\par

但请看，祂的权威贯穿一切：\Quote{我对你说：你是伯多禄；我要建立我的教会；我要给你天国的钥匙。}\par

4. 然后，祂说了这话后，\Quote{就嘱咐门徒，不要对任何人说祂是基督。}\BibleRef{玛 16:20}\par

祂为什么嘱咐他们？为的是当那些绊脚石被除去，十字架完成，其余的苦难成就，再没有任何事会妨碍和扰乱百姓对祂的信德时，对祂的正确意见才能纯洁而稳固地铭刻在听众心中。因为事实上，祂的权能尚未清楚地显示出来。所以祂愿意当他们宣讲时，事实的明显真理和祂德能的效力，都支持宗徒们的宣称。因为在巴勒斯坦看见祂，时而行奇迹，时而受侮辱和迫害（特别当十字架本身即将接踵而至），与后来在全世界看见祂受朝拜和相信，不再受任何祂曾受的苦，那是截然不同的。\par

因此祂吩咐他们\Quote{不要告诉任何人。}因为那曾生根而后被拔除的事，即使再栽种，也难以为大众所保持；但那些一经固定就不可动摇、未受任何损害的事，容易生长，并长进到更大的成长。\par

如果那些享有多项奇迹、参与如此多不可言喻奥迹的人，仅仅听到这事就跌倒；更准确地说，不仅是他们，连他们的领袖伯多禄也是如此；那么设想一般民众会有什么感觉？他们先被告知祂是天主子，然后看见祂甚至被钉死、被唾弃，并且对这些奥迹的秘密一无所知，也未参与圣神的恩赐。因为如果祂对门徒们说：\Quote{我还有许多事要告诉你们，但你们现在不能担当；}\BibleRef{若 16:12}那么，如果在适当时间之前将这些奥迹中最大的启示给他们，其余百姓更会完全失败。因此，祂禁止他们告诉人。\par

为了教导你，当他们后来在绊脚石过去后学习祂完整的道理，这是何等的大事，就从这个领袖身上学习吧。因为这个就是伯多禄，他在如此多的奇迹之后，竟软弱到否认祂，害怕一个卑微的使女；但在十字架之后，他接受了复活的确实证明，再也没有什么绊脚石和扰乱他的事，他如此坚定不移地持守圣神的教导，以至于比狮子更猛烈地扑向犹太百姓，尽管面临各种危险和无数的死亡威胁。\par

因此，祂在受难之前，有理由嘱咐他们不要告诉众人，因为就连那些将要教导的人，祂也未敢在受难之前将一切托付给他们。\Quote{我还有许多事要告诉你们，}祂说，\Quote{但你们现在不能承担。}\par

而且，对于祂所说的事，他们有许多也不明白，祂在受难之前并没有说清楚。至少当祂从死者中复活后，那时而非之前，他们才明白祂的一些话。\par

5. \Quote{从那时起，祂开始向他们表明祂必须受难。\BibleRef{玛 16:21} 从那时起。} 那时是什么时候？就是当祂把道理稳固在他们心中，并引进了外邦人的开端之时。\par

但即便如此，他们也不明白祂所说的话。\Quote{因为这话，}据说，\Quote{对他们隐藏了；}\BibleRef{路 18:34} 他们仿佛处于某种困惑之中，不知道祂必须复活。因此，祂更着重谈论艰难之处，并展开祂的论述，为要开启他们的心智，使他们能明白祂所讲的是什么。\par

\Quote{但他们不明白，这话对他们是隐藏的，他们也害怕问这事；}\BibleRef{路 9:45} 不是问祂是否要死，而是如何死，以何种方式，以及这奥秘究竟是什么。因为他们甚至不知道这复活是什么，并且认为最好不要死。因此，当其他人感到不安和困惑时，伯多禄又出于热忱，独自大胆谈论这些事；但他也不是公开地，而是把祂拉到一边，即与其余门徒分开，然后说：\Quote{主，千万不可！这事绝不会临到祢身上。} 这是怎么回事？那个获得启示、蒙受祝福的人，竟这么快就跌倒、被征服，以至于害怕祂的受难？何况一个人若在这些事上未得到任何启示，有这种感受又有什么可奇怪的呢？是的，为了告诉你，他说的那些别的话也不是出于自己，请看在这些没有启示给他的事上，他是如何被扰乱和推翻，即便被告知千万次，也不知道那话是什么意思。\par

因为他已经知道祂是天主子，但十字架和复活的奥秘尚未向他显明：因为\Quote{这话，}据说，\Quote{对他们隐藏了。}\par

你看见吗？祂有理由吩咐他们不要向其他人宣告。因为如果这事使那些必须知道的人也如此困惑，那么其他一切人会怎样感觉呢？\par

6. 然而，为了表明祂远非违背自己的意愿而走向受难，祂既责备了伯多禄，又称他为撒殚。\par

让那些以基督十字架的苦难为耻的人听吧！因为如果首席宗徒，甚至在完全明白之前，就因这种感觉而被称作撒殚，那么那些在如此充分的证明之后仍否认祂的\OriginalTerm{经世}{economy}的人，还有什么借口呢？我说，那个如此蒙福、作了如此宣认的人，竟受到这样的话；试想，那些在这之后仍否认十字架奥秘的人将要受什么苦呢？\par

祂没有说：\Quote{撒殚藉着你说话，} 而是说：\Quote{退到我后面去，撒殚。}\BibleRef{玛 16:23} 这确实是仇敌的愿望，愿基督不要受苦。因此，祂以如此严厉的责备斥责他，因为祂知道他和其余人尤其害怕这件事，不会轻易接受。\par

因此，祂也揭示了他心中的想法，说：\Quote{你所体会的，不是天主的事，而是人的事。}\par

\Quote{你所体会的，不是天主的事，而是人的事} 是什么意思？伯多禄以人性和世俗的理性来审视这件事，认为这对祂来说是可耻的、不合适的。因此，祂严厉地触动他说：\Quote{我的受难并非不合适的事，但你以肉性的心思下此判断；你若以属灵的方式听从我的话，脱离你肉性的理解，你就会知道这事最合于我。因为你确实以为受苦对我不配；但我告诉你，对我不受苦才是魔鬼的心意；} 藉着相反的表述来抑制他的惊惶。\par

正如若翰认为由他给基督受洗是基督不配的，却因基督说\Quote{你暂且容许吧，因为我们理当这样尽全义}\BibleRef{玛 3:15} 而被说服；同样，这位伯多禄也曾阻止祂洗脚，主却说\Quote{我若不洗你，你就与我无分}\BibleRef{若 13:8}；同样，在这里祂也藉着提到相反的事和严厉的斥责，抑制了他对受苦的恐惧。\par

7. 因此，没有人应以我们救恩的尊贵象征和一切美善的元首为耻。我们藉着它生活，藉着它存在；要像戴冠冕一样，将基督的十字架佩戴在身上。是的，因为藉着它，我们在中间所行的一切事都得以成就。无论一个人是要重生，十字架在那里；或要以那奥秘的食物得滋养，或要受圣职，或要做其他任何事，处处都有我们胜利的象征。因此，我们在房屋、墙壁、窗户上，并在我们的额上和心上，都精心刻划它。\par

因为这是为我们所成就的救恩、我们共同的自由和主的良善的标记。\Quote{祂像羔羊被牵去宰杀。}\BibleRef{依 53:7} 因此，当你划十字圣号时，要思想十字架的目的，熄灭愤怒和一切其他情欲。当你划十字圣号时，要使你的额充满一切勇气，使你的灵魂自由。你们确实知道什么事物给予自由。因此，保禄也带领我们，我意思是到那合宜的自由那里，他藉此方式带领我们，提醒我们主的十字架和血。他说：\Quote{你们是重价买来的，不要作人的奴隶。} 他说，请想一想为你们所付的代价，你们就不会作任何人的奴隶；这里代价指的是十字架。\par

既然不仅应当用手指镌刻它，更在此之前以心中的信德来镌刻它。如果你以这种方式将它标记在脸上，就没有一个不洁的灵能靠近你，因为它们看见了那使其受到创伤的利刃，看见了那给予它致命一击的剑。因为如果我们看见处决罪犯的场所就战栗，那么试想魔鬼看见那武器——基督藉以终结它一切权势、斩断那龙的头颅的武器——时，它要承受何等痛苦。\par

不要因如此大的恩宠而感到羞耻，免得基督在祂带着荣耀来临、那记号在祂面前出现、比太阳光芒更加耀眼时，为你感到羞耻。因为十字架那时确实来临，以它的出现发出声音，为我们的主向全世界辩护，表明属于祂的一切丝毫没有失落。\par

这记号，在我们祖先的时代和如今，打开了关闭的门；这记号熄灭了毒药；这记号解除了毒芹的效力；这记号医治了毒兽的咬伤。因为如果它打开了地狱之门，敞开了天堂的拱门，为乐园开辟了一条新的入口，并斩断了魔鬼的神经，那么它战胜毒药、毒兽和所有这一类事物，又有什么可惊奇的呢？\par

因此，你要将这记号刻在你的心中，拥抱我们灵魂的救恩。因为这十字架拯救并改变了世界，驱散了错谬，带回了真理，使大地成为天堂，将人塑造成天使。因着它，魔鬼不再可怕，反而可鄙；死亡也不再是死亡，而是睡眠；因着它，一切与我们为敌的都被击倒在地，被践踏在脚下。\par

如果有人对你说：\Quote{你敬拜那被钉十字架的吗？}你要以喜悦的声音和欣慰的面容说：\Quote{我确实敬拜祂，并且永不止息地敬拜祂。}如果他嘲笑，你要为他哭泣，因为他是癫狂的。感谢主，祂赐予了我们如此大的恩惠，甚至没有祂从上而来的启示，人便无法得知。这正是他嘲笑的原因，因为\Quote{属血气的人不接受天主圣神的事}。\BibleRef{格前 2:14} 因为我们的孩子看见任何伟大奇妙的事物时也会有这样的反应；如果你带一个孩子进入奥迹，他会发笑。如今外邦人就像这些孩子；甚至比他们更不成熟；因此他们更可怜，因为不是在未成熟的年龄，而是在成年时，却拥有婴儿的情感；因此他们也不配得到宽容。\par

但让我们以清晰的声音，高声呼喊说（即使所有外邦人都在场，也要更加自信）：十字架是我们的光荣，是一切恩惠的总和，是我们的依靠，是我们全部的花冠。我愿能像保禄那样说：\Quote{借着它，世界于我已被钉在十字架上，我于世界也被钉在十字架上。} \BibleRef{迦 6:14} 但我不能，因为被各种情欲所束缚。\par

8. 因此，我劝告你们，也首先劝告我自己，要向着世界被钉在十字架上，不与大地有任何共同之处，而要将你们的爱慕放在天上的家乡，以及从那里而来的荣耀和美好事物。因为我们确是天国君王的士兵，穿着属灵的武装。那么，我们为何要染上商人、骗子，甚至虫子的生活呢？因为君王在哪里，士兵也当在哪里。是的，我们成为了士兵，不是属于远方之王的，而是属于近处之王的。因为地上的君王尚且不允许所有人都待在王宫里、在他身边，但天上的君王却愿意所有人都靠近祂的御座。\par

也许有人说：我们在此世，怎么可能站在那御座旁呢？因为保禄虽然在地上，却与革鲁宾和色辣芬同在，而且比这些君王的近卫更接近基督。因为这些近卫频繁地转动脸庞，但没有什么能迷惑或分心保禄，他始终将整个心神专注于君王。所以，如果我们愿意，我们也能做到。\par

因为如果祂在空间上远离我们，你或许会怀疑；但祂无所不在，对于努力奋斗和真诚的人，祂是临近的。因此先知也说：\Quote{我不惧怕灾祸，因为你与我同在；}而天主自己又说：\Quote{我是近处的天主，不是远处的天主。} \BibleRef{耶 23:23} 正如我们的罪使祂远离我们，同样我们的义行使我们亲近祂。因为经上说：\Quote{你还在说话时，我就要说：我在这里。} 有哪个父亲会这样顺从自己的子女？哪个母亲如此随时准备，不断站立，等着孩子呼唤她？没有，没有父亲，也没有母亲；但天主一直站着，等待祂的仆人或许呼唤祂；当我们按应有的方式呼唤时，祂从不拒绝垂听。因此祂说：\Quote{你还在说话时，}我不等你说完，就立刻垂听。\par

9. 因此，让我们按祂愿意被呼唤的方式呼唤祂。但祂的旨意是什么呢？祂说：\Quote{解开一切不义的束缚，松开压迫性盟约的扭结，撕毁每一份不公正的契约。把你的饼分给饥饿的人，将无地容身的穷人领到你的家中。若见赤身露体者，给他遮盖，不要忽视你的骨肉。那时，你的光要在清晨升起，你的医治要迅速发生，你的公义要在你前面行，主的荣耀要作你的后盾。那时你呼求我，我必应允你；你还在说话时，我就要说：看哪，我在这里。}\par

也许有人问：谁能做到这一切呢？我倒要问：有谁做不到呢？因为在我所提及的事中，哪一件是困难的？哪一件是劳苦的？哪一件是不容易的？\par

何以如此？这些训言不仅是可行的，甚至容易至极，致使许多人实际上已经超越了这些话语的尺度，他们不仅撕毁了不义的契约，甚至舍弃了自己所有的财产；他们不仅欢迎穷人来到屋檐下和餐桌前，甚至迎接他们到自己身体汗水的劳作中，辛苦劳动以供养他们；他们不仅善待亲属，甚至还善待仇敌。\par

然而，这些训言中究竟有什么艰难之处呢？因为祂并没有说：\Quote{越过山岭，渡过大海，挖掘如此多亩的土地，不进食，披上麻衣；}而是说：\Quote{分施给穷人，分施你的面包，取消不义订立的契约}。\par

还有什么比这更容易的呢？请告诉我。但即使你认为它困难，我恳求你，也看看赏报，那么它对你就会容易了。\par

因为正如我们的皇帝在赛马场上，在竞赛者面前堆积冠冕、奖品和衣物，同样，基督也将祂的赏报放置在比赛途中，借着先知的话语，仿佛通过许多手递出来。而皇帝们，即使他们是千万次的皇帝，但由于是人，他们所拥有的财富处于消耗过程中，他们的慷慨也会枯竭，他们总想把小的显得大；因此他们分别将每件东西交到不同侍从的手中，再这样呈上前来。但我们的君王却相反，祂将一切都堆积在一起（因为祂极其富有，且不为炫耀而行任何事），就这样呈上前来，祂所递出的东西无限巨大，需要许多手才能捧住。为使你意识到这一点，请仔细探究其中的每一个细节。\par

祂说：\Quote{那时，} \Quote{你的光明要在清晨迸发。}这恩赐在你看来不是某一样东西吗？但它并非一样；相反，它包含了许多东西，既有奖品，又有冠冕，还有其他赏报。如果你愿意，让我们将它拆开，并尽可能地展示其全部财富；只是你们不要厌倦。\par

首先，让我们明白\Quote{它要迸发}的含义。因为祂根本没有说\Quote{要显现}，而是说\Quote{要迸发}；以此向我们表明其迅速与丰富，以及祂多么切望我们的救恩，又表明美好事物本身如何急切地要出来，并催促着；并且没有任何事物能遏制那难以言说的力量；祂用这一切来表明其丰富和祂无穷财富的无限。但\Quote{清晨}是什么意思？意思是\Quote{不是在经历人生的诱惑之后，也不是在我们的祸患降临之后}；不，它远比这些来得更早。正如在果实中，我们把在季节之前就显现的称为早熟；同样，在这里，为了表明其迅速，祂这样说话，正如上面祂所说的：\Quote{你还在说话时，我就要说：看，我在这里。}\par

但祂说的是什么样的光呢？这光可能是什么呢？不是这感官所及的光，而是另一种远胜于此的光，它向我们显示天堂、天使、总领天使、革鲁宾、色辣芬、宝座、宰制、率领、掌权、万军、王宫、帐幕。因为如果你堪当得到这光，你将会看见这些，并从地狱、毒虫、咬牙切齿、无法挣开的锁链、痛苦与磨难、无光之黑暗、被斩断、火河、诅咒和哀伤之所中被解救出来；你将前往\BibleQuote{那里没有忧愁和哀叹}（\BibleRef{依 35:10}），那里有大喜乐、平安、爱、愉悦和欢乐；那里有永生、不可言喻的光荣和无以言表的美；那里有永恒的帐幕、君王不可言喻的光荣，以及\BibleQuote{眼所未见，耳所未闻，人心所未想到的}美好事物（\BibleRef{格前 2:9}）；那里有属灵的洞房、天上的殿宇、持明亮灯火的童女，以及穿着婚宴礼服的人；那里有许多我主所拥有的产业和君王的库房。\par

你看见了吗？祂用一个表述就提出了多么伟大、多么众多的赏报，并且将一切汇集在一起？\par

同样，当我们逐一解开后面每个表述时，我们将会发现我们的恩惠浩大，海洋无垠。那么，我们还迟疑什么呢？我恳求你们，还迟迟不肯怜悯那些需要帮助的人吗？不，我恳求你们，即使我们必须舍弃一切，被投入火中，冒犯刀剑，跃向匕首，承受你所能想象的一切苦难；让我们轻松承受这一切，以便获得天国婚宴的礼服和那无与伦比的光荣；愿我们众人，因我主耶稣基督的恩宠和祂对人的爱，都能获得这光荣，愿光荣与权能归于祂，至于无穷之世。阿们。\par

\SourceRecord{Translated by George Prevost and revised by M.B. Riddle.；Nicene and Post-Nicene Fathers, First Series；Vol. 10.；Edited by Philip Schaff.；Buffalo, NY: Christian Literature Publishing Co.,；1888.；<http://www.newadvent.org/fathers/200154.htm>.}

% structure: p0001 source=h1 target=section reason=page-title

\section{\WorkTitle{讲道词第五十五篇 论玛窦福音}}

\Emph{\BibleRef{玛 16:24}}\par

\Quote{ \Emph{那时，耶稣对他的门徒说：谁若愿意跟随我，该弃绝自己}，\Emph{背起自己的十字架来跟随我。} }\par

然后呢？什么时候？就是当伯多禄说：\Quote{主，千万不可！这事绝不会临到你身上}，并被回答说：\Quote{撒殚，退到我后面去}（\BibleRef{玛 16:22}）的时候。因为祂绝不满足于仅仅斥责，反而更愿意充分显示伯多禄所说的话的过分之处以及祂苦难的益处，于是祂说：\Quote{你对我所说的话是：\InnerQuote{主，千万不可！这事绝不会临到你身上}，但我对你说：\InnerQuote{不但阻止我、对我的苦难不悦是有害的、毁灭性的，而且你也不可能得救，除非你自己也时常准备好赴死。}}\par

因此，为免他们以为祂的苦难与祂不相称，祂不仅借着先前的事，也借着将到来的事件，教导他们其中的益处。于是祂先引用了\BibleBook{若望}{福音}的话说：\BibleQuote{一粒麦子如果不落在地里死了，仍只是一粒；如果死了，才结出许多子粒来}（\BibleRef{若 12:24}）。但在这里更充分地发挥，不仅关于祂自己，也关于他们，祂提出了必须死去的声明。\Quote{因为这益处如此之大，以至于在你们身上，不愿死去是痛苦的，但做好准备却是好的。}\par

然而，祂借着随后的话使这事清楚明白，但暂时祂只在一方面加以发挥。请看祂如何使自己的言论无懈可击：祂绝不说：\Quote{不管你愿不愿意，你都必须受这些苦}，而是怎么说？\Quote{谁若愿意跟随我。} \Quote{我不强迫，我不勉强，而是使每人成为自己选择的主人；因此我也说：\InnerQuote{若有人愿意}。因为我召叫你们是为了美好的事物，而不是邪恶或沉重的事物；不是为了惩罚和报复，以致我必须强迫。不，事情的本质本身足以吸引你们。}\par

祂这样说时，更加吸引他们归向祂。因为使用强迫的人往往使人远离，但让听者选择的人更吸引人。因为安抚比强迫更有力。因此，连祂自己也说：\Quote{若有人愿意。}祂说：\Quote{因为我赐给你们的好处是伟大的，甚至使人主动跑向它们。因为若有人给金子、献上宝藏，他不会用强迫来邀请。如果那邀请是不强迫的，更何况这邀请——指向天上的美好事物呢？因为如果事物的本性不能说服你跑去，你根本不配得到它，即使你得到了，你也不会真正认识你所得到的。}\par

因此基督并不强迫，而是劝勉，体恤我们。因为他们似乎私下因这话而深感不安，祂说：\Quote{不需要惊慌或烦恼。如果你们不认为我所提及的——即使是落在你们自己身上——是无数祝福的原因，我不使用武力，也不强迫，但若有人愿意跟随，我就召叫那人。}\par

\Quote{千万不要以为这就是跟随我，我指的是你们现在所做的——陪伴我。你们需要许多辛劳，许多危险，如果你们要跟随我。因为，伯多禄，你绝不可因为我承认你是天主子就只期待冠冕，并认为这足够你的得救，而将来可以像完成了所有事一样享受安全。因为虽然我作为天主子，有能力阻止你遭受任何那些艰难；但为了你的缘故，我并不愿意这样，为的是你也能自己贡献些什么，并更得蒙悦纳。}\par

因为假若有人在竞技场上担任裁判，有一位朋友在参加竞赛，他不会只凭恩惠就给他戴上冠冕，而是也希望因他自己的劳苦而给他加冕；尤其因为这正是他爱他的缘故。同样，基督也是这样：祂最爱的人，祂最愿意让他们也凭自己的努力显明自己，而不仅靠祂的帮助。\par

但请看祂如何同时使祂的话不显得沉重。因为祂绝不只以恐怖包围他们，而是将道理普遍地推向世界，说：\Quote{若有人愿意，}无论是女是男，是统治者还是臣民，让他走这条路。\par

虽然祂似乎只说了一件事，但祂的话有三句：\Quote{弃绝自己}，\Quote{背起自己的十字架}，\Quote{跟随我}；其中两句是连在一起的，但一句是分开的。\par

但让我们先看看舍弃自己是什么意思。我们先学习舍弃别人是什么意思，然后就知道舍弃自己可能是什么。那么，舍弃别人是什么意思呢？那舍弃别人的人——例如，无论是弟兄、仆人或任何你愿意的人——若看见他被打、被绑、被处决，或受任何苦，他都不站在他旁边，不帮助他，不为所动，对他毫无感觉，因为已经与他完全隔绝了。因此，祂要我们这样不顾自己的身体，无论人鞭打、流放、焚烧，或做任何事，我们都不要顾惜它。因为这才是顾惜它。正如父亲们把孩子托付给老师时，命令不必顾惜他们，这恰恰是顾惜他们。\par

同样，基督也没有说：\Quote{不要顾惜自己}，而是非常严格地说：\Quote{弃绝自己}；那就是，与他自身无关，而是把自己交付给一切危险和战斗；并且要感觉好像这一切都发生在另一个人身上。\par

祂没有说\Quote{否认}，而是说\Quote{弃绝}；甚至通过这个小小的添加再次暗示，这有多么深远。因为后者比前者更进一步。\par

\Quote{也要背起他的十字架。}这是由前一句引申出来的。为了防止你以为言语、侮辱和谴责是我们克己的极限，祂也说人应当克己到什么地步，那就是至死，且是一种耻辱的死。因此祂没有说\Quote{让他克己至死}，而是说\Quote{让他背起他的十字架}，指明了这种耻辱的死；并且不是一次、两次，而是整个一生都该如此行。\Quote{是的，}祂说，\Quote{要不断带着这死，日日准备被宰杀。因为许多人确实轻视了财富、享乐和荣耀，但他们并不藐视死亡，反而惧怕危险；我，}祂说，\Quote{要我的战士甚至奋战至流血，他的赛程的终点应该是被杀；因此，即使必须经历死亡，带着耻辱的死亡，被咒诅的死亡，并且是在恶名之下，我们也要勇敢地忍受一切，反而因被怀疑而喜乐。}\par

\Quote{并且跟随我。}也就是说，一个人可能受苦却没有跟随祂，当他不是为祂受苦的时候（因为强盗、盗墓者和行邪术的人也常常受重苦）；为了防止你以为仅仅是苦难的性质就足够了，祂加上了这些苦难的缘由。\par

这是什么缘由呢？那就是：当你这样行、这样受苦时，你要跟随祂；你要为祂承受一切；你也要拥有其他的德行。因为这也由\Quote{让他跟随我}表达出来，以便不仅显明在苦难中所锻炼的刚毅，也显明节制、中和及一切克己。这正确地就是\Quote{跟随}，即留意其他的德行，并为祂的缘故承受一切。\par

因为有些人跟随魔鬼甚至到了忍受这一切的地步，并且为它舍弃自己的生命；但我们是为基督，或者更确切地说，是为我们自己：他们确实为了在此世和来世都伤害自己；但我们是为了赢得两方面的生命。\par

那么，我们连那些丧亡之人的同样刚毅都表现不出来，岂非极其迟钝？何况我们还要从中收获如此多的冠冕？然而，基督自己当然与我们同在，作我们的帮助，而他们却无人相助。\par

其实，祂在派遣他们时已经说过这个命令，说：\BibleQuote{不要往外邦人的路上去}（因为祂说：\BibleQuote{我派遣你们好像羊进入狼群}，以及\BibleQuote{你们要被带到君王和总督面前}），但现在更加激烈和严厉了。因为那时祂只说到死，但这里还提到了十字架，并且是持续的十字架。因为\Quote{让他背起，}祂说，\Quote{他的十字架}，也就是说，\Quote{让他不断地背负它、忍受它。}这是祂在一切事上惯常的做法：不是一开始，不是起初，而是安静地、逐步地引入更大的诫命，以免听众觉得奇怪。\par

3. 然后，因为这话似乎严厉，请看祂如何用接下来的话缓和它，并设立了超越我们劳苦的赏报；不仅是赏报，还有罪恶的惩罚：不，祂更着重于后者，因为赐予恩惠不如威吓严厉更能使平凡人醒悟。请看祂如何从这里开始，也以此结束。\par

\BibleQuote{因为谁若愿意救自己的性命，必要丧失性命；但谁若为我的缘故丧失自己的性命，必要找到性命。人纵然赚得了全世界，却赔上了自己的灵魂，为他有什么益处？或者，人能拿什么作为自己灵魂的代价？}\par

祂所说的是这个意思：\Quote{我命令这些，不是对你们无情，反而是极其顾惜你们。因为那怜惜自己儿子的，反而毁了他；那不怜惜的，却保全了他。}一位智者也曾说：\BibleQuote{你若用棍杖打你的儿子，他必不至死，你却要救他的灵魂脱离死亡。}\BibleRef{箴 23:13}又说：\BibleQuote{娇惯自己儿子的，必包扎他的创伤。}\BibleRef{德 30:7}\par

在军营中也是如此。如果将军顾惜士兵，命令他们总是留在营内，他连同城中的居民一起毁灭。\par

\Quote{因此，为了不让这事也发生在你们身上，}祂说，\Quote{你们必须面对持续不断的死亡。因为如今一场惨烈的战争将要燃起。所以不要坐在营内，而要出去作战；即使你们倒在岗位上，那时你们便获得了生命。}因为在可见的战争中，那些在岗位上被杀的，比其他人更杰出、更不可战胜、更令敌人畏惧；虽然我们知道，死后他所效忠的君王不能使他复活。何况在这些战争中，除了有复活的希望，那将自己性命置于死地的人，必要找到它：一方面，因为他不会很快被捉住；另一方面，即使他倒下，天主也会将他的生命引领到更高的生命。\par

4. 然后，因为祂曾说：\Quote{那愿意救的，必要丧失；但那丧失的，必要救得，}并在一边设立了得救和丧亡，在另一边设立了丧亡和得救；为了防止有人以为这个丧亡和得救与另一个相同，并为了清楚地教导你，这个得救与那个得救之间的差别如同丧亡与得救之间的差别一样大，祂也从反面作了一次推论，一劳永逸地确立这些论点。\BibleQuote{人纵然赚得了全世界，却赔上了自己的灵魂，为他有什么益处？}祂说。\par

你看见了吗？错误的保存乃是丧亡，比一切丧亡更糟，因为无可补救，缺乏任何能赎回它的东西。因为\Quote{不要这样对我说，}祂说，\Quote{那逃脱了这些危险的人保存了他的性命；但即使把全世界连他的性命一起加给他，若灵魂丧亡，他有什么益处呢？}\par

请你告诉我，你若见你的仆人在享乐，而你自己却处于极大的灾难中，身为家主，你究竟能得什么益处？断然不能。对于你的灵魂，也当这样去衡量：当肉身处于安逸与富足之中，而灵魂却等候着将来的毁灭。\par

\BibleQuote{人还能拿什么作为自己灵魂的代价？}\par

主又重复这一点。他说：什么？你还有另一个灵魂可以代替这个灵魂吗？你失去金钱，尚能用金钱赎回；失去房屋、奴隶或任何其他财产也是如此。但若失去灵魂，你便没有另一个灵魂可以付出。纵然你拥有全世界，纵然你作全地的君王，你也不能用所有地上的财物并大地本身，赎回一个灵魂。\par

这有什么可惊奇的呢？即使是身体，我们也能看到同样的情况。你虽佩戴千万冠冕，但若身体天生孱弱且无药可医，你即使用整个王国，加上无数的人、城邑和财物，也不能恢复这身体的健康。\par

我也请你这样对待你的灵魂，甚至更应如此。你当放下一切，将全部精力倾注于它。不要只顾别人的事而忽略你自己和属于你的事物——这是现今众人常做的事，就像在矿坑里干活的人一样。他们既不能从劳动中得到任何益处，也不能从财富中获利，反而遭受大害：因为他们冒了无益的风险，且是为他人冒风险，从这样的劳苦和死亡中得不到任何回报。如今许多人仿效他们，为他人挖掘财富；或者不如说，我们比他们更悲惨，因为在我们这样的劳苦之后，地狱在等待着我们。他们因死亡而停止了那些劳苦，但死亡对我们却是无数祸患的开端。\par

但若你说，你在财富中获得了劳苦的果实，请向我显示你的灵魂快乐，这样我才相信。因为在我们内，灵魂是首要的。即使身体肥壮，而灵魂消瘦，这繁荣对你毫无益处（就像婢女欢乐，而主母却憔悴，婢女的幸福对主母无益；又像华服与孱弱的肉体相比，毫无价值）；但是基督要对你说：\BibleQuote{人还能拿什么作为自己灵魂的代价？} 处处命令你为此费心，并只以此为准。\par

5. 主既以此震慑他们，又用自己的美好事物安慰他们。\par

\Quote{因为人子将要在他父的光荣中，同他的圣天使一起降来，} 他说，\BibleQuote{那时，他要按照每人的行为予以报应。}\par

你看见父与子的光荣是何等一致吗？既然光荣是一，那么本体显然也是一。因为若在同一本体中有光荣的差别（\BibleQuote{因为太阳有一种光荣，月亮有另一种光荣，星辰有另一种光荣；而且星辰与星辰的光荣也有分别；}见：\BibleRef{格前 15:41}，虽然本体为一），那么光荣为一者，本体怎能不相同呢？因为他根本不是说\Quote{像父那样的光荣}，使你或许因此设想某种变化，而是说\Quote{在同样光荣中}，他说，\Quote{他将降临，} 表示视为同一。\par

\Quote{现在，你为什么害怕，哦，\Person{伯多禄}？}（他这样说）\Quote{那么，你将在父的光荣中看见我。我若在光荣中，你也在光荣中；你的利益绝不局限于现世，另一种更好的命运将接纳你。} 然而，在谈论美好事物之后，他并未停留于此，而是将可畏的事物也掺入其中，提出审判宝座、无情的清算、不可改变的判决和无法欺瞒的审判。\par

但他没有让演讲显得只有阴郁，而是用美好的希望使之缓和。因为他并没有说\Quote{那时他要惩罚犯罪的}，而是说\Quote{他要按照每人的行为予以报应}。这样说不仅提醒罪人惩罚，也提醒善人奖赏和荣冠。\par

6. 他说这话，部分是为了安慰善人，但我每次听到都战栗不已，因为我不属于那些受冠冕的人；我想其他人也与我一同恐惧和焦虑。有谁听到这句话而不深自反省良心，不感到惊骇，不相信自己更需要穿苦衣、长期禁食，甚至超过尼尼微人呢？因为我们的担忧不是关于一座城的倾覆和普遍的结局，而是关于永恒的惩罚和不灭的火。\par

因此，我也赞颂和钦佩那些居于旷野的隐修士，既为其余的一切，也为这句话。他们在用过晚餐（他们从不吃午餐，因为他们知道现今是哀恸和禁食的时候）之后，以某些感恩的圣歌赞美天主时，也提及了这句话。若你愿意听这些圣歌，以便你也能不断诵念，我将为你复述那整首神圣的歌曲。其文如下：\Quote{永受赞颂的天主，你从我的幼年养育了我，你赐食物给一切血肉；求你以喜乐和欢愉充满我们的心，使我们常常知足，在基督耶稣我们的主内，充足地行一切善事。愿光荣、尊威和权能，借着我们的主耶稣基督，偕同圣神，归于你，直到永远。阿们。光荣归于你，上主，光荣归于你，圣者，光荣归于你，君王，因为你赐给我们食物，使我们快乐。求你以圣神充满我们，使我们得以在你面前蒙悦纳，在你按每人的行为施报时，不致蒙羞。}\par

如今这赞美诗处处堪受赞叹，尤其是上述的结尾。这是因为，饮食往往使人逸散和沉重，他们便将这话语如同缰绳置于灵魂上，在纵情之时提醒它审判的时刻。因为他们已得知以色列人因奢侈的筵席所遭遇的事。\Quote{我的爱人，}祂说，\Quote{吃了，肥壮了，就踢跳。}因此\Person{梅瑟}也说：\BibleQuote{当你吃了、喝了、饱足了，要记念上主你的天主。}\BibleRef{申 6:11}\par

因为在这次宴会之后，他们才胆敢做出那些无法无天的行为。\par

因此你也要注意，免得同样的事临到你。因为你虽不向石头或金子祭献牛羊，但要注意，不要向忿怒祭献你自己的灵魂，不要向淫乱或其他类似的情欲祭献你自己的救恩。是的——正因如此，他们害怕这些堕落，在享受了他们的餐食之后，或者更确切地说，在禁食之后（因为他们的餐食实际上是禁食），便提醒自己那可怕的审判宝座和那一日。他们用禁食、席地而卧、守夜、苦衣以及千万种方法纠正自己，尚且需要这样的提醒，那么，对我们这些摆满无数残骸的餐桌、且根本不做祈祷（无论是餐前还是餐后）的人来说，何时才能活出美德呢？\par

7. 因此，为终结这些沉船之祸，让我们把那赞美诗摆在我们面前，并完全展开它，好使我们看见其益处，也常常在我们的餐桌上一遍遍吟唱它，平息肚腹的粗暴冲动，将那些天使的品行和规矩引入我们的家中。因为你本应去那里收获这些果实；但既然你不愿意，至少通过我们的话语，聆听这属灵的旋律，让每个人饭后说出这些话语，并这样开始。\par

\Quote{应受赞美的天主。}因为他们立刻履行了宗徒的训令，命令说：\BibleQuote{无论我们做什么，或在言语上，或在行为上，都要因主\Person{耶稣基督}的名而行，借着祂感谢天主父。}\BibleRef{哥 3:17}\par

其次，感恩不只是为了那一天，而是为了他们的一生。因为经上说：\Quote{从我的幼年至今，你喂养了我。}由此引出节制的教训：既然天主喂养我们，我们就不必忧虑。因为若有一位君王应许从自己的库中供应你每日的粮食，你就能对未来充满希望；何况当天主赐予，万物如泉源般涌向你时，你更应从一切挂虑中得自由。是的，正为此他们这样说，为要说服他们自己和他们的门徒，卸下一切世俗的忧虑。\par

然后，为使你不认为他们只是为自己献上这感恩，他们进一步说：\Quote{你赐食物给一切血肉之人，}代表全世界感恩；如同全地之父，他们为众人献上赞美，并训练自己真诚的兄弟之爱。因为他们不可能恨那些他们为其感谢天主得食物的人。\par

你看见了吗？感恩既引入了爱德，又藉着前面和后面的话语驱逐了世俗的忧虑。因为若祂喂养一切血肉之人，那么更喂养那些献身于祂的人；若祂喂养那些陷入世俗忧虑的人，那么更喂养那些摆脱了这些忧虑的人。\par

为确立这一点，\Person{基督}亲自说：\Quote{你们比许多麻雀贵重得多？}祂这样说，教导他们不要信赖财富、土地和种子；因为喂养我们的不是这些，而是天主的话。\par

藉此他们堵住了摩尼教徒和\Person{瓦伦提努}的门徒以及所有类似思想病态之人的口。因为这位存有者绝不是邪恶的，祂将自己的库藏敞开给众人，甚至给那些亵渎祂的人。\par

然后便是祈求：\Quote{请以喜乐和欢欣充满我们的心。}那么，这是怎样的喜乐呢？此世的喜乐吗？绝对不可！若他们意指此世的喜乐，就不会占据山顶和旷野，也不会披上苦衣；但他们所指的喜乐，与此生毫无共同之处，是天使的喜乐，是上天的喜乐。\par

他们不只是简单地求它，而且极其丰盛地求；因为他们不说\Quote{赐予，}而说\Quote{充满，}且不说\Quote{我们，}而说\Quote{我们的心。}因为这尤其是心灵的喜乐；\BibleQuote{圣神的果子就是爱、喜乐、平安。}\BibleRef{迦 5:22}\par

这样，因为罪恶带来了忧愁，他们祈求藉着喜乐，正义得以栽种在他们内，因为喜乐无法以其他方式产生。\par

\BibleQuote{好使我们时常有充足的一切，可以多多行善。}\BibleRef{格后 9:8} 看他们如何实践福音的那句话：\Quote{求你赐给我们今日的日用粮，}并且他们甚至为此寻求属灵的目的。因为他们的措辞是：\Quote{我们可以多多行善。}他们不是说\Quote{我们只尽本分，}而是\Quote{甚至超过所吩咐的}，因为\Quote{我们可以多多}正是指此。他们向天主祈求所需之物的充足，而他们自己不仅乐意在充足上服从，更乐意在大量和一切事上服从。这是好仆人的作为，这是严格行善之人的作为，常常在一切事上多多善行。\par

然后，他们再次提醒自己的软弱，以及若无上天的助佑，便不能做任何高尚的事；在说了\Quote{我们可以多多行善}之后，他们又加上\Quote{在我们的主\Person{耶稣基督}内，愿荣耀、尊威和权能归于你，直到永远。阿们。}他们如同以感恩之线编织这结尾，就像它的开始一样。\par

8. 此后，他们似乎又重新开始，但仍在坚持同一论点。正如\Person{保禄}在书信的开头，以光荣颂结束时所写的：\Quote{按照我们天主和圣父的旨意，愿光荣归于祂，直到永远。阿们。}\BibleRef{迦 1:4}，然后又重新开始他所论述的主题。同样，在另一处，当他说了\Quote{他们崇拜并事奉受造物超过造物主，祂是永远可赞美的，阿们}\BibleRef{罗 1:25}之后，并没有结束他的讲论，而是重新开始。\par

因此，我们也不要责怪这些天使们，认为他们行为无序，因为他们在以光荣颂结束之后，又重新开始神圣的赞美诗。因为他们遵循法则，从光荣颂开始，也在其中结束，而在结束之后又重新开始。\par

因此他们说：\Quote{光荣归于祢，上主；光荣归于祢，圣者；光荣归于祢，君王；因为祢赐给了我们食物，使我们欢乐。}\par

因为我们不仅应为大事，也应为小事感恩。他们也为这些小事感恩，从而羞惭了摩尼派的异端，以及所有主张现世生命是邪恶的人。因为为了不让你们因他们高度的自持和对肚腹的轻视，而怀疑他们厌恶食物（像前面提到的那些异端者，他们以窒息自尽的方式），他们通过祈祷教导你们：他们之所以戒绝大部分食物，并非出于对天主造物的厌恶，而是为了操练自制的德行。\par

请看，在感谢祂过去的恩赐之后，他们又为更大的事恳求，而不停留于今生的事物，而是升到诸天之上，说：\Quote{求祢以圣神充满我们。}因为若不被那恩宠充满，人甚至不可能适当地证明自己；因为若没有基督恩宠的影响，便无法成就任何高贵或伟大的事。\par

因此，正如当他们说\Quote{使我们多行各种善工}时，又加上了\Quote{在基督耶稣内}；同样，在这里他们也说：\Quote{求祢以圣神充满我们，好使我们能成为在祢面前蒙悦纳的人。}\par

你们看见了吗？他们不为今世的事物祈求，而只感恩；但对于属神的事物，他们既感恩又祈求。因为祂说：\Quote{你们先寻求天国，这一切自会加给你们。}\par

还要注意他们身上另一种严苛的良善：即他们说：\Quote{好使我们在祢面前蒙悦纳，不蒙羞愧。}因为他们说：\Quote{我们不介意来自多数人的羞辱，无论人们怎样嘲笑、责备我们，我们都不放在心上；我们的一切努力在于不在那时蒙羞。}在这些话中，他们也引入了火河、奖赏和报酬。\par

他们没有说\Quote{我们不受惩罚}，而是说\Quote{我们不蒙羞愧}。因为对我们而言，这比地狱更令人畏惧，即似乎得罪了我们的主。\par

但由于大多数人和较为粗俗的人并不畏惧这一点，他们又加上：\Quote{当祢按各人的行为施行报应的时候。}你们看见这些陌生人和旅客——这些旷野的公民，或者更确切地说是天上的公民——对我们有多大的益处吗？因为我们对天国是陌生的，却是地上的公民，而他们恰恰相反。\par

在这首赞美诗之后，他们怀着极大的痛悔和许多热泪，如此进入睡眠，只取用足以稍事休息的睡眠。而夜晚，他们则使之变为白天，用于感恩和唱诵圣咏。\par

不单是男子，连妇女也实践这种克己，以热忱之丰沛克服了本性的软弱。\par

那么，让我们因他们的热忱而感到羞愧吧，我们是男子，却仍被捆绑于现世的事物、影子、梦境和烟雾之中。因为我们生命的大部分都在麻木中度过。\par

因为人生的第一个阶段充满许多愚昧，而走向老年的阶段又使我们的所有感觉枯萎，中间的短暂时间则能够有感觉地享受快乐；或者说，即使在那段时间里，由于无数困扰我们的忧虑和劳苦，也无法纯粹地享受快乐。\par

因此，我恳求，让我们寻求那不动摇的永恒财富，以及永不衰老的生命。\par

因为即使是住在城里的人，也可以效法隐修士们的克己；是的，有妻子且忙于家务的人，也可以祈祷、斋戒并学习痛悔。因为那些最初由宗徒教导的人，虽然住在城里，却显出了旷野居住者的虔诚；另一些管理作坊的人，如\Person{普黎斯加}和\Person{阿桂拉}也是如此。\par

而且，所有先知也都有妻子和家庭，如\Person{依撒意亚}、\Person{厄则克耳}、伟大的\Person{梅瑟}，他们并未因此受到伤害而妨碍德行。\par

那么，让我们也效法他们，不断向天主献上感恩，不断向祂咏唱赞美诗；让我们专注于节制和所有其他德行，并把旷野中所实践的克己带入我们的城市；这样，我们便能在天主面前蒙悦纳，并在人前蒙赞许，借我们的主耶稣基督的恩宠和慈爱，获享未来的美善。愿光荣、尊威和大能，藉着祂，偕同祂，并联合圣神及生命之赋予者，归于圣父，自今而始，以至永远，永世无疆。阿们。\par

\SourceRecord{Translated by George Prevost and revised by M.B. Riddle.；Nicene and Post-Nicene Fathers, First Series；Vol. 10.；Edited by Philip Schaff.；Buffalo, NY: Christian Literature Publishing Co.,；1888.；<http://www.newadvent.org/fathers/200155.htm>.}

% structure: p0001 source=h1 target=section reason=page-title

\section{玛窦福音讲道 56}

\BibleRef{玛 16:28}\par

\BibleQuote{我实在告诉你们：站在这里的人中，就有些人在未尝到死味以前，必要看见人子来到自己的国内。}\par

因此，既然祂已谈论了许多关于危险和死亡，关于祂自己的受难，以及关于门徒的杀戮，并给他们加上了那些严厉的诫命；而这些都是现世生活且近在眼前的，但美好的事物却在希望和期待中：——例如，\Quote{凡丧失自己生命的，必将得着它} \Quote{祂将在他父的光荣中来临} \Quote{祂将给予报偿}——祂愿意使他们亲眼确信，并显示祂将要来临时的光荣是怎样的，只要他们能理解；祂在现世生活中就显示并启示了这一点，好使他们不再悲伤，无论是对自己的死亡，还是对主的死亡，尤其是伯多禄因祂的悲伤而特别难过。\par

并且看看祂做了什么。谈论了地狱和天国之后（因为祂既说了\BibleQuote{那找着自己生命的，必要丧失它；那为了我而丧失自己生命的，必要找到它；} \BibleRef{玛 16:25}又说了\BibleQuote{祂将按照每人的行为予以赏报}，祂已经将这两者都显明了）：我已经说过，论及两者之后，祂在异象中所显示的确实是天国，但地狱却还没有。\par

为何如此？因为如果他们属于另一种人，较为粗俗，那么这也将是必要的；但既然他们是经过考验而富有思想的，祂便以更温和的方式引导他们。但祂作出这个启示，不仅仅因此，也因为这对祂自己更为合适。\par

然而祂并不是忽略了这个主题，而是在有些地方，祂几乎将地狱的实情带到我们眼前；例如祂引入拉匝禄的比喻，提到那追讨一百银钱的仆人，以及那穿着污秽衣服的人，以及其他不少人。\par

2. \BibleQuote{六天后，耶稣带着伯多禄、雅各伯和若望} \BibleRef{玛 17:1}\par

另一位却说是\BibleQuote{八天后} \BibleRef{路 9:28}，这并不是与这位作者矛盾，而是完全一致。因为一位说出了祂说话的那一天以及祂带领他们上山的那一天；而另一位只说了之间的天数。\par

可是请你注意玛窦的严厉的善意，他没有隐瞒那些被偏爱的人。若望也常常这样做，极其真诚地记录对伯多禄的特别称赞。因为这群圣人的歌队处处都纯洁，没有嫉妒和虚荣。\par

因此，祂带着那几位首领，\BibleQuote{独自带他们上了一座高山，在他们面前变了容貌：祂的面貌发光如太阳，祂的衣服洁白如光。忽然，梅瑟和厄里亚显现给他们，正在同耶稣谈论。} \BibleRef{玛 17:2}\par

为何祂只带了他们？因为他们比其他门徒更优秀。伯多禄确实以其对祂的极大爱情显示了优越性；若望则因祂对他极大的爱；雅各伯又因他与他兄弟的回答，他们说\BibleQuote{我们能饮这杯}；并不仅因他的回答，也因他的行为；既因其余的事，也因他实践了自己所说的话。因为他如此热忱，使犹太人感到不快，以致黑落德自己认为他此举是对犹太人的极大恩惠，我是指杀害雅各伯。\par

但为何祂不立刻带他们上山？为了保全其他门徒，避免他们有任何人性的软弱感：因此祂也省略了将要上山之人的名字。这样做，是因为其余门徒会极其渴望跟随，为的是看见那光荣的典范；他们会因被忽视而痛苦。因为虽然祂是以某种有形的方式作出了启示，但这启示本身仍是一件非常值得渴慕的事。\par

那么，祂为何要预先告诉他们呢？为了使他们通过祂的预言更愿意领悟那高深的含义；并在这几日内被更强烈的渴望所充满，从而以完全清醒和充满专注的心智出现在那里。\par

3. 但为何祂也将梅瑟和厄里亚带来？可以提到许多理由。首先是这样：因为群众中有人说祂是厄里亚，有人说是耶肋米亚，有人说是古代先知中的一位，祂便带来歌队的领袖，使他们即使在此也能看出仆人与主之间的区别；并且知道伯多禄因为承认祂是天主子而受到称赞是正当的。\par

除此之外，还可以提到另一个理由：因为人们不断指控祂违犯法律，并认为祂是亵渎者，因祂将不属于自己的光荣——甚至父的光荣——归于自己，他们说\BibleQuote{这人不是从天主来的，因为他并不守安息日} \BibleRef{若 9:16}；又说\BibleQuote{为了善工，我们并不用石头砸死你，而是为了亵渎的话，并且因为你是人，却把自己当作天主} \BibleRef{若 10:33}。为了使这两种指控都显明源于嫉妒，并证明祂对两者都无可指摘；也为了证明祂的行为并非违犯法律，祂自称与父平等也不是将不属于自己的光荣归于自己，祂便带来那些在各方面都显赫的人：梅瑟，因为他颁布了法律，犹太人可能会推断他不会像他们所认为的那样，容忍法律被践踏，也不会尊重违犯法律者和立法者的仇敌；厄里亚，他则为天主的荣光而热心，若是有人与天主为敌，自称天主，使自己与父平等，而实际上他并非如他所说，也无权这样做；厄里亚不会站在一旁听他说话。\par

除了已经说过的理由，还可以提到另一个。那么是什么呢？就是告诉他们：祂对死亡和生命都有权能，是上下两界的主宰。为此祂带来了那已经死去的和那从未经历过死亡的人。\par

然而，第五个动机（因为除了已经提到的那些，还有第五个），连福音作者自己也都揭示了出来。那么，这动机是什么呢？就是显示十字架的光荣，并安慰伯多禄和其余的人，使他们不再害怕苦难，同时提升他们的心志。因为他们来到后，绝没有保持沉默，而是如经上所说，\Quote{谈论}祂在耶路撒冷将要完成的光荣，\BibleRef{路 9:31}也就是说，谈论苦难和十字架；因为他们总是这样称呼它。\par

祂不仅以此安慰他们，还借着这些人的卓越本身，就是祂特别要求他们具有的那种卓越。我的意思是，祂曾说过：\Quote{谁若愿意跟随我，该背起自己的十字架来跟随我。}现在，祂将那些曾为天主的诫命和托付给他们的百姓死过一万次的人，摆在他们的面前。因为这些人中的每一位，都失去了自己的生命，又找回了它。因为其中一位曾勇敢地向埃及法老说话，另一位向阿哈布说话；为了那些心硬和悖逆的人；并且正是被他们拯救的那些人，使他们陷入极端的危险；两人都希望引导人们离开偶像崇拜；两人都是没有学问的：因为前者是\Quote{拙口笨舌}的（\BibleRef{出 4:10}），言语迟钝；而后者在举止上也有几分粗鲁；两人都是自愿贫穷的严格奉行者：梅瑟没有谋取任何利益，厄里亚除了一张羊皮外别无所有；而且这是旧约时代，他们还没有领受如此大的神恩。因为梅瑟分开红海又如何？伯多禄却在水面上行走，能移山，能医治各种身体的疾病，驱赶凶猛的魔鬼，甚至用他的影子行出那些奇妙而伟大的奇迹，改变了整个世界。如果厄里亚也曾复活死人，这些人却复活了数以万计；而且是在圣神尚未赐给他们的时候。因此，祂也为此缘故把他们带出来。因为祂希望他们效法这些人对待百姓的亲和态度、镇定和坚定；希望他们像梅瑟那样温良，像厄里亚那样为天主热诚，并如他们一样充满慈爱。因为厄里亚为了犹太人民忍受了三年饥荒；而梅瑟说：\Quote{如果祢宽恕他们的罪，就宽恕；否则，也请把我从祢所写的书上抹去。}\BibleRef{出 32:32}这一切，祂都借着异象提醒他们。\par

因为祂也将那些人带来并显示在光荣中，并不是要这些门徒停留在原来的境界，而是使他们甚至超越那些界线。例如，当他们说：\Quote{我们要叫火从天降下来吗？}并提及厄里亚曾如此行时，祂说：\Quote{你们不知道你们有的是什么精神}；借着他们恩赐的优越来训练他们学会容忍。\par

谁也不要以为我们在谴责厄里亚为不完全；我们并非这样说；因为他确实是非常完全的，但在他自己的时代，那时人的心智在某种程度上尚属幼稚，需要这种管教。同样，梅瑟在这方面也是完全的；然而，这些人所被要求的，比他所被要求的更多。因为\Quote{除非你们的义德超过经师和法利塞人的义德，你们决不能进入天国。}\BibleRef{玛 5:20}因为他们不是进入埃及，而是进入整个世界，而这个世界比埃及更败坏；他们也不是与法老说话，而是与魔鬼本身，即邪恶的首领，进行肉搏。是的，他们被指定的战斗是：要捆绑魔鬼，并夺取他的一切财物；他们这样做不是分开红海，而是借着叶瑟的棍杖分开不敬神明的深渊——一个波涛更为凶险的深渊。试看有多少事情使这些人恐惧：死亡、贫穷、耻辱、数不尽的苦难；他们对这些的恐惧，远超过古时犹太人对那红海的恐惧。但尽管如此，祂说服他们勇敢地冒险去面对这一切，并像在干地上一样安全地渡过。\par

因此，为了训练他们应对这一切，祂将那些在旧约下光照人的圣者带了出来。\par

4. 那么，热情奔放的伯多禄说了什么呢？\Quote{我们在这里真好。}\BibleRef{玛 16:4}因为他曾听说基督要去耶路撒冷受难，即使在祂责备之后，他仍然害怕并为祂战栗，他不敢再上前说同样的话：\Quote{主，千万不可！}\BibleRef{玛 16:22}但他从那种恐惧中，用其他的话语隐约表达了同样的意思。也就是说，当他看见一座山，如此幽静和僻静时，他想：\Quote{祂在这里非常安全，即使从地点来看也是如此；而且不仅从地点来看，还因为祂不再去耶路撒冷。}因为他希望祂一直留在那里；所以他也提到了\Quote{帐幕}。因为他说：\Quote{如果可以这样，我们就不必上耶路撒冷；如果我们不去，祂就不会死，因为祂说经师们会在那里攻击祂。}\par

但他确实不敢这样说；然而，他渴望事情如此安排，便毫不怀疑地说：\Quote{我们在这里真好，}那里也有梅瑟和厄里亚；厄里亚曾在这座山上使火降下，梅瑟曾进入浓云中与天主交谈；没有人会知道我们在哪里。\par

你看到这位热爱基督的人有多么热切吗？现在不要看他劝说的方式是否经过深思熟虑，而要看他多么热切，他对基督的爱情有多么炽烈。因为为了证明他说这些话不是出于对自己的恐惧，请听当基督预言自己未来的死亡和将要受到的攻击时，他说了什么：\Quote{我要为祢舍掉我的性命。}\BibleRef{若 13:37}\Quote{即便我必须与祢同死，我也决不会否认祢。}\BibleRef{玛 26:35}\par

且看，即便在真正的危险之中，他如何为自己出了错误的劝告。我们知道，当如此众多的人群包围他们时，他非但没有逃跑，反而拔剑砍下了大司祭仆人的耳朵。他如此不顾自身利益，只为他的师傅担忧。然后，因为他说话如同陈述事实，他便制止自己，心想：若他再受责备又当如何？他说：\Quote{你若愿意，我们就在这里搭三座帐棚，一座为你，一座为梅瑟，一座为厄里亚。}\par

伯多禄，你说什么？你刚才不是把他与仆人区别开了吗？你又把他与仆人列在一起了吗？你看见他们在受难之前是多么不完备吗？因为虽然圣父已向他启示了这事，但他并没有始终持守这启示，反而被他的恐惧所困扰；不仅是我提到的这个恐惧，还有由那景象产生的另一个恐惧。事实上，其他福音作者为了说明这一点，并指出他说这些话时心神的混乱乃是由那恐惧所致，这样写道：注意，\BibleQuote{他不知道该说什么，因为他们都吓坏了；}\BibleRef{谷 9:6}但路加在他说了\BibleQuote{我们搭三座帐棚}之后，又补充了\BibleQuote{他不知道自己在说什么。}\BibleRef{路 9:33}接着，为表明他和其余的人都被极大的恐惧所抓住，他说：\Quote{他们被睡意沉重压着，醒来时就看见了祂的光荣；}此处所谓沉睡，是指那异象在他们心中产生的深深麻木。因为正如眼睛因过度明亮而昏暗，那时他们也有同样感受。我想，那不是黑夜，而是白天；那光的极度强烈压低了他们眼睛的软弱。\par

5. 那么怎样呢？祂自己什么也不说，梅瑟和厄里亚也不说，但那比一切都大、更值得相信的圣父，从云中发出了声音。\par

为何从云中呢？因为天主常常如此显现。\Quote{云和黑暗在祂四围；}并且\BibleQuote{祂乘驾轻云；}\BibleRef{依 19:1}又\Quote{祂以云彩为车辇；}以及\BibleQuote{有一朵云彩接祂去了，离开他们的视线；}\BibleRef{宗 1:9}还有\BibleQuote{人子乘着天上的云彩而来。}\BibleRef{达 7:13}\par

为了使门徒相信那声音来自天主，它便从那里发出。\par

而且那云彩是光亮的。因为\BibleQuote{他还在说话的时候，看，有一片光耀的云彩遮蔽了他们；并且，看，从云中有声音说：这是我的爱子，我所喜悦的，你们要听从他。}\BibleRef{玛 17:5}\par

因为，当祂发出警告时，祂显示一朵暗云——如同在西乃山；因为据说\Quote{梅瑟} \Quote{进入了云中和浓密的黑暗中；烟像蒸气一样上升；}先知在谈及祂的警告时也说：\Quote{天上的云中藏着黑暗的水；}——但在此处，因为祂的意愿不是要恐吓，而是要教导，所以云是光亮的。\par

伯多禄曾说过\Quote{让我们搭三座帐棚}，但主却显示了一座非人手所造的帐棚。因此，在西乃山那里是烟和炉子的蒸气，而在这里却是说不出的光和声音。\par

然后，为表明这声音并非指那三人中的某一位，而唯独是指基督，当声音发出时，他们就被移开了。因为若这声音是针对他们中的任何一位，那人就不会单独留下，而另外两人必会与他分离。\par

那么，为什么云彩没有单独接走基督，而是把他们都一起接去了呢？如果云彩只接走了基督，人们会以为那声音是祂自己发出的。故此，福音作者为确认这一点，指出那声音是从云中发出的，也就是说，来自天主。\par

那声音说了什么呢？\Quote{这是我的爱子。}既然祂是爱子，伯多禄，你不要害怕。因为你确实应当已经知道祂的能力，并对祂的复活深信不疑；但既然你不知道，至少从圣父的声音中获取勇气吧。因为如果天主是强大的——祂无疑强大——那么圣子显然也是如此。因此，不要害怕那些可怕的事。\par

如果你至今仍不能领受这一点，至少想一想另一事实：祂是圣子，且是爱子。因为经上说：\Quote{这是我的爱子。}既然祂是爱子，就不要害怕。因为没有人会放弃他所爱的人。因此，你不要慌乱；即使你爱祂超过一切，你爱祂也不如那生了祂的爱祂之深。\par

\Quote{我所喜悦的。}因为祂爱祂，不仅因为生了祂，也因为祂在一切方面与祂同等，并与祂同心。因此，爱的魅力是双重的，甚至可以说是三重的：因为祂是圣子，因为祂是爱子，因为在祂身上祂得了喜悦。\par

但\Quote{我所喜悦的}是什么意思？仿佛祂说：\Quote{在我内我得安息，在我内我得喜悦；}因为祂在一切方面与祂自己完全平等，在祂与圣父之间只有一个旨意，并且祂虽是圣子，却在一切方面与圣父为一。\par

\Quote{你们要听从他。}所以，即使祂选择被钉十字架，你也不可反对祂。\par

6. \BibleQuote{门徒听见了，就俯伏在地，非常害怕。耶稣前来，抚摸他们说：起来，不要害怕。他们举目一看，不见一人，只有耶稣独自一人。}\BibleRef{玛 17:6}\par

怎么他们听了这些话就惊惶了呢？然而先前在约旦河也曾有类似的声音发出，当时有许多人在场，却没有人感到那样；后来当人们说\BibleQuote{打雷了}时，\BibleRef{若 12:28}那时他们也没有经历这样的事。那么，为什么他们在山上俯伏在地呢？因为那里僻静、高耸、极其宁静，并有令人敬畏的显圣容、纯净的光芒和铺开的云彩——所有这些都使他们极度恐慌。惊骇从四面八方涌来，他们既因恐惧又因朝拜而俯伏在地。\par

但恐怕长久存留的恐惧会驱散他们的记忆，当下祂便止住了他们的惊慌，只让他们看见祂自己，并命令他们在这事上不可告诉任何人，直到祂从死者中复活。\par

因为\BibleQuote{当他们从山上下来时，祂嘱咐他们不要把所见的事告诉任何人，直到祂从死者中复活。}祂将要做什么。\par

7. 再没有比宗徒们更有福的了，尤其是那三位，他们甚至在云中获准与主同在一屋之下。\par

但我们若愿意，也将看见基督，不是像他们那时在山上，而是以更加明亮的光辉。因为将来祂来临的方式并非如此。当时为了怜恤祂的门徒，祂只显露出他们所能承受的光辉；将来祂要在圣父的光荣中来临，不仅与梅瑟和厄里亚同来，还要与无数的天使、总领天使、革鲁宾以及那无尽的队列同来，祂头上不再有云彩，连天本身也要卷起。\par

正如审判官的情况：当他们公开审判时，侍从拉开帷幕，将他们显给众人；同样，那时众人要看见祂坐着，全人类都要站在一旁，祂要亲自回答他们；对一些人祂要说：\BibleQuote{你们蒙我父降福的，来吧！因为我饿了，你们给了我吃的；}\BibleRef{玛 25:34}对另一些人祂要说：\BibleQuote{好！善良忠信的仆人，你既在少许事上忠信，我必委派你管理许多大事。}\BibleRef{玛 25:23}\par

然后又宣判相反的判决，对一些人祂要回答：\BibleQuote{离开我，到那给魔鬼和他的使者预备的永火里去吧！}\BibleRef{玛 25:41}对另一些人祂要说：\BibleQuote{可恶懒惰的仆人。}\BibleRef{玛 25:26}有些人祂要\Quote{割断}，\Quote{交给刑役；}但另一些人祂要命令\Quote{捆起他的手和脚，丢在外面的黑暗中。}\BibleRef{玛 22:13}斧头之后将有火炉；凡是从网里丢弃的，都要落在其中。\par

\BibleQuote{那时，义人要发光如同太阳。}\BibleRef{玛 13:43}或者更胜过太阳。但这样说，并非因为他们的光只有那么多，而是因为我们不知道比这更亮的星，祂便用已知的例子来说明圣人未来的光辉。\par

因为在山上，当祂说\BibleQuote{祂发光如同太阳}时，也是因同一缘故而如此说。因为宗徒们仆倒在地，表明那比较并未达到祂的光辉。倘若那光辉不是纯然的，且与太阳相比，他们就不会仆倒，而能轻易承受。\par

因此，义人在那时要像太阳一样发光，甚至比太阳更明亮；而罪人则要受尽极苦。那时将不需要记录、证据、证人。因为审判者本身就是一切，祂既是证人，又是证据，也是审判者。因为祂知道一切详细；\BibleQuote{万有在祂眼前都是袒露敞开的。}\BibleRef{希 4:13}\par

那里不再有富人或穷人，有权势的或软弱的，有智慧的或愚拙的，为奴的或自主的；这些面具都要被击碎，只审查他们的行为。因为在我们法庭上，若有人因篡权或杀人受审，无论他是总督、执政官还是什么人，所有这些尊荣都消失，被定罪者受极刑；在那里更是如此。\par

8. 因此，为使这事不发生，让我们脱去污秽的衣服，穿上光明的铠甲，天主的荣耀必将环绕我们。因为诫命中有什么是沉重的呢？有什么不是容易的呢？请听先知的话，你就晓得其容易。\BibleQuote{即使你低头如同芦苇，以苦衣和灰烬铺在下面，你也不能称此为悦纳的斋戒；但要解除不义的锁链，解开轭上的绳索。}\BibleRef{依 58:6}\par

请看先知的智慧，他首先指出令人厌烦的事，并将其除去，然后劝他们藉着容易的本分获得救恩；这表明天主不需要劳苦，而是顺从。\par

然后他暗示德行容易，而恶习沉重并令人难堪，仅凭名字就说明了：\Quote{因为，}他说，\Quote{恶习是锁链，}和\Quote{轭上的绳索，}而德行则是摆脱这一切。\par

\Quote{撕毁一切不义的契约；}这是指人们关于应付利息和所借款项的票据。\par

\Quote{释放受压迫的人；}即受苦难的人。因为债务人就是这样：当他看见债权人时，心灵破碎，比畏惧野兽更甚。\par

\BibleQuote{将流浪的穷人领到你家中；你若看见赤身露体的，给他衣服穿，不要忽略你的骨肉。}\BibleRef{依 58:7}\par

在我们上次对你们宣讲赏报的讲道中，我们指出了这些行为所生的财富；但现在让我们看看是否有任何诫命是沉重的，超越我们本性的。不，我们不会发现任何这类事，恰恰相反：这些道路非常容易，而恶习的道路则充满劳苦。因为有什么比放债、思虑利息和契约、要求担保、为证券、本金、文书、利息、保证人而恐惧战栗更令人烦恼呢？\par

因为世俗之事就是如此；确实，没有什么比那些被认为有保障并为这目的而设计的东西更不牢靠、更可疑的了；但施舍怜悯是容易的，且使人摆脱一切焦虑。\par

我们不可贩卖他人的灾祸，也不可将我们的慈惠当作买卖。我知道许多人听了这话会不悦；但沉默有何益处？因为我若缄默不言，不以言语烦扰你们，我也不能借着这沉默使你们免受惩罚；反倒适得其反，刑罚加重，不仅对你们，对我也是如此沉默招致惩罚。那么，当我们言语虽好，行为却不助益我们，反而有害时，这言语有何意义？言语讨人喜欢，行为却惹人恼怒，给耳朵带来快乐，给灵魂带来惩罚，这有什么好处？因此，我必须在此使你们忧伤，好使我们不在那里受罚。\par

9. 的确，亲爱的，一种可怕的疾病，可怕且需要精心照料，已降临在教会身上。那些被命令不可——甚至通过正当的劳力——积蓄财宝，而应开放家宅给有需要的人的人，却利用别人的贫穷牟利，制造出一种似是而非的抢劫，一种貌似合理的贪婪。\par

你们不要向我提及外邦人的法律；因为连税吏也遵守那外邦人的法律，但依然要受惩罚：我们也是如此，除非我们停止压迫穷人，停止利用他们的需要和贫乏作为无耻交易的借口。\par

因为你们拥有财富的目的，是为了纾解贫困，而不是为了利用贫困图利；但是你以救济为表，却使灾难更加严重，并出售慈惠以换取金钱。你尽管出售，我不禁止你，但要为了天国而出售。不要为如此善举收取低廉的报酬，即你每月一分的利息，而要收取那不朽的生命。你为何如此卑贱、贫穷、低下，用你伟大的事物换来微小的事物，甚至换来腐朽的财物，本应为永恒的国度而交换？你为何离弃天主，去获取人的收益？你为何越过那富有者，去打扰那贫穷者？离弃那可靠的付款人，与那忘恩负义者交易？那一位渴望偿还，而这一位甚至在偿还时也吝啬。这一位难以偿还百分之一，而另一位\BibleQuote{百倍和永生}。这一位带着侮辱和辱骂，而另一位带着赞美和吉言。这一位激起人对你的嫉妒，而另一位甚至为你编织冠冕。这一位难在此世，而另一位既在此世，也在来世。\par

那么，这难道不是极端的愚昧，甚至连如何获利都不知道吗？有多少人为了利息而失去了本金？有多少人因高利贷的收益而陷入危险？有多少人因他们那难以言表的贪婪，而使自己和他人陷入极度的贫困？\par

你们不要对我说，他乐于接受，并感谢贷款。这其实是你们残忍的结果。因为亚巴郎也曾策划如何使他的计划在野蛮人中成功，而将自己的妻子交给他们；然而并非出于自愿，而是出于对法老的恐惧。同样，穷人因为你认为他甚至不值那些钱，实际上被迫感激这种残忍。\par

在我看来，似乎你若将他从危险中解救出来，你也要他为这解救付出代价。他说：\Quote{走开，}\Quote{不要这样。}你说什么？你将他从更大的恶中解救出来时，不愿索取金钱，而为了较小的恶，你却表现出如此不人道吗？\par

你们难道看不见这行为将招致多大的惩罚吗？你们难道没有听见，即使在旧法律中这也是被禁止的吗？但许多人的托词是什么？有人告诉我：\Quote{我收到利息后，就施舍给穷人。}请说敬畏的话，人啊！天主不喜爱这样的祭献。不要巧妙地规避法律。与其那样施舍，不如不施舍给穷人；因为那通过诚实劳动积累的金钱，你常常因那邪恶的增加而使之成为不法；就好像强迫一个正常的子宫生出蝎子一样。\par

我为何要提天主的法律呢？连你们自己不也称之为\Quote{污秽}吗？但如果你们这些获利者如此表态，那么试想天主将对你们作出何等裁判。\par

如果你也询问外邦人的立法者，他们也会告诉你，连他们也认为这是最无耻的证明。例如，那些身居荣誉职位、属于他们称为元老院的伟大议会的人，依法不得以这类收益玷污自己；他们中间有一条法律禁止这样做。\par

那么，如果你们不肯把立法者给予罗马议会的同等尊荣给予天上的政体，反而使天上的比地上的得着更少，你们又不对这事的愚蠢感到羞愧，这难道不可怕吗？还有什么比这更愚蠢的呢？除非一个人没有土地、雨水或犁，却坚持要播种。因此，那些设计出这邪恶耕作的，所收割的是要投入火中的莠子。\par

难道没有许多正当的行业吗？在田间、羊群、牛群、畜牧、手工艺、财产管理方面？为何要疯狂、狂热地耕作无益的荆棘？即使地上的出产会遭遇不幸：冰雹、枯萎和过多的雨水，那又怎样？但不如金钱交易所遭受的。因为在那些情况下，损失当然关乎出产，但本金，即土地，仍存留。但在此处，许多人常常连本金也遭遇海难；而且在损失之前，他们也一直处于沮丧之中。放债人从不享受他的财富，也不在其中找到乐趣；当利息呈上时，他不因获得收益而欢喜，反而因利息尚未还清本金而忧伤。在这邪恶的产物完全诞生之前，他又迫使它生产，使利息成为本金，强迫它生出其不合时宜的、蛇蝎般早产的幼崽。高利贷的收益就是这种性质；它们比那些野兽更吞噬和撕裂不幸者的灵魂。这就是\Quote{不义的契约}；这就是\Quote{压迫性交易的扭曲的结}。\par

是的，\Quote{我给予，}祂似乎说，\Quote{不是为叫你们收受，而是使你们能偿还更多。}而天主甚至命令不收受所给予的（因为祂说：\BibleQuote{给予那些你们不指望从他们那里收受的人}），你们却要求多于所给予的，并且将你们未曾给予的，当作债务，强迫收受人偿付。\par

你确实以为你的资财由此增加，但代替资财，你是在点燃那永不熄灭的火。\par

为使这事不发生，让我们割除高利贷收益的邪恶母腹，让我们扼杀这些无法无天的分娩，让我们吸干这有害多产的场所，只追求真实而伟大的收益。\Quote{这些是什么？}请听保禄说：\BibleQuote{虔敬而知足乃是大利。}\BibleRef{弟前 6:6}\par

因此，让我们只在这财富上富足，好使我们既在此世享受安稳，又能达到来世的美好事物，藉着我们的主耶稣基督对人类的恩宠与爱，愿光荣与德能归于祂，连同圣父及圣神，自今至永远，直到无穷世。阿们。\par

\SourceRecord{Translated by George Prevost and revised by M.B. Riddle.；Nicene and Post-Nicene Fathers, First Series；Vol. 10.；Edited by Philip Schaff.；Buffalo, NY: Christian Literature Publishing Co.,；1888.；<http://www.newadvent.org/fathers/200156.htm>.}

% structure: p0001 source=h1 target=section reason=page-title

\section{玛窦福音讲道57}

\Emph{\BibleRef{玛 17:10}}\par

\BibleQuote{并且祂的门徒问祂说：那么，为什么经师说厄里亚必须先来？}\par

他们这并不是从圣经知道的，而是经师们自己解释，这说法在无知百姓中流传，如同关于基督的也一样。\par

为此，撒玛黎雅妇人也说：\BibleQuote{默西亚来了；当祂来了，祂会告诉我们一切} \BibleRef{若 4:25}，而他们自己问若翰说：\BibleQuote{你是厄里亚，还是那位先知？} \BibleRef{若 1:21}。因为，如我说，那说法，既关于基督也关于厄里亚，很盛行，不过没有被他们正确解释。\par

因为圣经讲述基督两次来临，一次是过去的，一次是未来的；保禄宣告这些说：\BibleQuote{天主的救恩恩宠已经显现，教训我们，弃绝不敬神和世俗的情欲，应当节节自制地、正义地、虔诚地生活}。看前一次，听他如何也宣告另一次；因为说了这些后，他加上：\BibleQuote{期待那有福的盼望和我们伟大的天主和救主耶稣基督的显现} \BibleRef{铎 2:13}。而且先知们也提到两者；不过其中之一，即第二次，他们说厄里亚将作前驱。因为第一次，若翰是前驱；基督也称他为厄里亚，不是因为他真是厄里亚，而是因为他完成了那位先知的职分。因为正如这位将是第二次来临的前驱，另一位也是第一次的前驱。但经师们混淆这些，误导百姓，只向百姓提到那第二次来临，并说：\Quote{如果这人是基督，厄里亚就该先来}。所以门徒们也这样说：\Quote{那么，为什么经师说厄里亚必须先来？}\par

因此法利塞人也派人到若翰那里，问他：\BibleQuote{你是厄里亚吗？} \BibleRef{若 1:21} 他们丝毫未提前一次来临。\par

基督所提出的解答是什么？\Quote{厄里亚确实在我第二次来临之前要来；而且如今厄里亚已经来了}；如此称呼若翰。\par

在这个意义上，厄里亚已经来了；但如果你寻找提斯贝人，他要求。所以祂也说：\BibleQuote{厄里亚确实要来，且将恢复一切}。一切什么？如玛拉基亚先知所说的；因为经上说：\BibleQuote{我要派遣你们厄里亚，提斯贝人，他将使父亲的心转向儿子，免得我来彻底打击大地}。\par

你看到预言用词的精确吗？因为基督称若翰为厄里亚，基于他们职务的共同点，免得你以为先知在此处也指这个意思，祂还加上他的故乡，说：\Quote{提斯贝人}；然而若翰不是提斯贝人。此外祂还提出另一个标记，说：\BibleQuote{免得我来彻底打击大地}，指向祂第二次可怕的来临。因为第一次祂来不是为了打击大地。因为祂说：\BibleQuote{我来不是为审判世界，而是为拯救世界} \BibleRef{若 12:47}。\par

为显示提斯贝人是在那另一个带有审判的来临之前来，祂说了这话。并且祂也教导祂来的原因。这原因是什么？就是当祂来了，可以说服犹太人相信基督，使他们不致在祂来临时全部灭亡。所以祂引导他们回想那话，说：\BibleQuote{他将恢复一切}；即纠正那时代犹太人的不信。\par

所以祂的说法极其精确；因为他没有说\BibleQuote{他将使儿子的心转向父亲}，而是\BibleQuote{使父亲的心转向儿子}。因为犹太人曾是宗徒们的父亲，他的意思是：他将使父亲们的心，即犹太民族的心，转向他们儿子的道理，即宗徒们的道理。\par

\BibleQuote{但我告诉你们：厄里亚已经来了，他们却不认识他，反而任意对待了他。同样，人子也要受他们的苦。} \BibleRef{玛 17:12}\par

然而，无论是经师们没有这样说，圣经也没有这样说；只是因为他们现在对祂的话更敏锐、更专注，就很快领会了祂的意思。\par

门徒们从哪里知道这事？祂已经告诉过他们：\BibleQuote{他就是那将要来的厄里亚} \BibleRef{玛 11:14}；但这里说，他已经来了；又说：\BibleQuote{厄里亚要来并恢复一切}。但你不要困扰，不要以为祂的话摇摆不定，虽然祂有时说\Quote{他要来}，有时说\Quote{他已经来了}。因为这些都是真的。因为当祂说\Quote{厄里亚确实要来并恢复一切}时，祂指的是厄里亚本人，以及那时即将发生的犹太人的归化；但当祂说\Quote{那将要来的}时，祂称若翰为厄里亚，按他管理的方式。是的，先知们也常称他们每个被认可的君王为达味，称犹太人为\BibleQuote{索多玛的官长} \BibleRef{依 1:10} 和\BibleQuote{雇士人的子孙} \BibleRef{亚 9:7}，因为他们的行为。因为正如那一位将是第二次来临的前驱，这位也是第一次的前驱。\par

2. 不仅如此，祂到处称他为厄里亚，还为了表示祂与旧约完全一致，以及这次来临也是符合预言的。\par

因此祂再次加上：\BibleQuote{他来了，他们却不认识他，反而任意对待了他。} \BibleRef{玛 17:12} \Quote{任意对待了他}是什么意思？他们把他投进监狱，凌辱他，杀了他，把他的头放在盘子里端来。\par

\Quote{同样人子也将被他们苦待。}你们看他又在适当时刻提醒他们关于自己的受难，从若翰的受难中为他们积蓄大量安慰。而且不止如此，他还通过当下行了大奇迹来加以安慰。确实，每当他谈及自己的受难时，他就当下行奇迹，既在这些话之后，也在这些话之前；在许多地方，可以看见他保持这规律。\par

\Quote{那时，}经上说，\Quote{他开始示意他必须往耶路撒冷去，被杀害，并受许多苦。}\BibleRef{玛 16:21}\Quote{那时：}何时？当他被承认为基督、天主的儿子时。\par

再者，在山上，当他向他们显示了明显的异现，先知们谈论了他的荣耀时，他提醒他们他的受难。因为谈论了关于若翰的事迹后，他又补充道：\Quote{同样人子也将被他们苦待。}\par

过了一小会儿，当他赶出了那个门徒们无法赶出的鬼时，因为在那时，\Quote{他们住在加利勒那时，}经上说，\Quote{耶稣对他们说，人子将被交在罪人手中，他们要杀害他，第三天他要复活。}\BibleRef{玛 17:23}\par

他这样做，是用巨大的奇迹减弱他们过度的悲伤，并在各方面安慰他们；就像这里，通过提及若翰的死，为他们带来了许多安慰。\par

但若有人说，\Quote{为何他连现在也不掀起厄里亚派遣他来，如同他如此见证他的到来是如此巨大的善事呢？}我们应回答，就是当时，当人们以为他是厄里亚时，也没有相信他。因为\Quote{有人说，}这样的话，\Quote{你是厄里亚，另有人说，你是耶肋米亚。}\BibleRef{玛 16:14}而且在若翰和厄里亚之间，除了时间之外没有区别。\Quote{那时他们怎么会相信呢？}可能有人说。为何？\Quote{他将恢复一切，}不是简单地因为被认出，而且也因为基督的荣耀将直到那日不断增大，在所有人中将比阳光更明亮。因此，当他在这样的观念和期望之后而来，发出与若翰同样的宣告，并且亲自也宣布耶稣时，他们将更容易接受他的话。但当他说，\Quote{他们不认识他，}他也在为自己身上所做的事开脱。\BibleRef{路 23:24}\par

不止如此，他还通过指出若翰在他们手中的苦难无论如何是不应得的，并且通过两个征兆——一个在山上，另一个即将发生——来掩盖他们可能担心的事，从而安慰他们。\par

但当他们听到这些话时，他们不问他厄里亚何时来；或被他的受难的悲伤所困，或是因恐惧。因为在许多场合，见他不愿明说某事，他们就保持沉默，事情就此了结。例如，当他们住在加利勒时，他说：\Quote{人子将被交出，他们要杀害他，}\BibleRef{玛 17:22}马尔谷接着说：\Quote{他们不明白这话，又不敢问他。}\BibleRef{谷 9:32}路加说：\Quote{这话向他们隐藏了，使他们无法明白，他们也不敢问他这话。}\BibleRef{路 9:45}\par

3. \Quote{当他们来到众人那里时，有一个人来到他身边，跪下来对他说：主啊，可怜我的儿子，因为他是月症病，很受折磨，因为他经常掉在火里，经常掉在水里。我带他到你的门徒那里，他们却不能治好他。}\BibleRef{玛 17:14}\par

经上表示，这个人的信德极其弱小；这点可以从许多方面看得出；从基督的话说：\Quote{在有信德的人，一切都是可能的，}\BibleRef{谷 9:23}从那个走过来的人自己的话说：\Quote{帮助我的不信，}\BibleRef{谷 9:24}从基督命令那鬼\Quote{再不要进入他}\BibleRef{谷 9:25}，以及那个人又对基督说：\Quote{你若能，}\BibleRef{谷 9:22}\Quote{但如果他的不信是原因，}可能有人说，\Quote{为什么鬼没有出去，他却责备门徒们？}这表示，即使没有人带病人带着信德前来，他们在许多情况下也可以治好人。因为，就像带人前来的人的信德经常可以使人从次级的侍奉者那里获得治疗；同样，行事者的能力经常也可以够用，即使前来的人没有信德，也能行奇迹。\par

在经上，这两种情况都有记载。因为，科尔乃略的人凭着他们的信德获得了圣神的恩宠；另外，在厄里叟的情况下，\BibleRef{列下 13:21}当没有人相信时，一个死人又复活了。因为，扔人的人不是因信德，而是因胆怯扔人，无意中和偶然，因为怕一伙强盗，他们就逃跑了；而被扔下的人本身已经死了，但仅凭圣者的遗体的德能，那死人就复活了。\par

由此可以看出，在这个案件中，门徒们也是弱小的；但不是所有的门徒，因为那些支柱\BibleRef{加 2:9}当时并不在场。另外，从另一个方面看这个人缺乏考虑，他在众人面前向耶稣控告他的门徒们说：\Quote{我带他到你的门徒那里，他们却不能治好他。}\par

但他在众人面前开脱了门徒们的责任，将大部分抵当归于他。因为，\Quote{喔，不信与悖逆的世代啊，}这是他的话，\Quote{我将与你们同在多久呢？}\BibleRef{玛 17:17}从中不仅针对他个人，以免他感到困惑，而且也是对所有的犹太人。因为当时在场的许多人可能会被得罪，对他们产生不当的看法。\par

但当他说：\Quote{我将与你们同在多久呢？}他再度表示死亡对他而言是可欢迎的，是所渴望的事，他切望离开，并且，不是十字架，而是与他们同在，才是苦难。\par

然而祂并未止于责难；祂说什么？\BibleQuote{把他带到我这里来。}并且祂亲自问他：\BibleQuote{这事临到他有多久了？}\BibleRef{谷 9:21}这样，既为祂的门徒辩护，又将那人引向美好的希望，好使他相信能从恶神手中获得解救。\par

祂任凭他被折磨，并非为了炫耀（因此，当人群开始聚集时，祂便前去斥责它），而是为了那父亲本人的益处，好叫他看到因基督的一声呼唤，恶神就受惊扰，这样，即使没有别的方式，至少也能使他相信将要发生的奇迹。\par

因他曾说：\Quote{自幼}，并说：\Quote{你若能帮助我}，基督便说：\BibleQuote{为相信的人，一切都可能。}\BibleRef{谷 9:23} 再次将抱怨转回他身上。至于那癞病人说：\BibleQuote{你若愿意，就能使我洁净，}\BibleRef{玛 8:2} 为祂的权柄作证，基督称赞他，并用祂的话证实说：\Quote{我愿意，你洁净了吧；}但在此人的事例中，当他说出与祂大能不相称的话——\Quote{你若能，帮助我}——请看祂如何纠正它，视之为说得不当。因为祂说什么？\Quote{你若能信，为相信的人，一切都可能。}祂的意思是这样：\Quote{我有如此充裕的大能，甚至能使别人也成就这些奇迹。所以，你若以应有的方式相信，连你自己也能，}祂说，\Quote{治好这个人和许多其他人。}说了这话，祂便释放了那附魔的人。\par

但你不仅由此要观察祂的眷顾及仁慈，也要从那另一时刻观察——那时祂允许魔鬼留在他身上。因为确然，即使在那时，若那人未蒙受诸多眷顾的恩宠，他早已丧命；因为据说：\Quote{它把他投入火中}，\Quote{和水中。}那胆敢如此行此事的，定会毁灭那人，除非天主在这极大的疯狂中，以强有力的缰绳约束它；正如那些在旷野中奔跑、用石头割伤自己的裸体之人的情况。\par

倘若他称他为\Quote{癫癇的}，你完全不必困扰，因为这是附魔之人的父亲所说的话。那么，为何福音作者也说：\Quote{祂治好了许多癫癇的人？}他依据众人的印象这样称呼他们。因为恶神为了污损受造的天性，便趁此机会，以病症的发作来遮掩；因为他所见到的这些患儿，并无缘由；但基督明确分辨了它们，使这骗术不能长久，因为祂立即斥责它，它便离开了那人。但在这里，我要向众位说一些关于醉酒的话。因为这病也使人癫癇，使人疯癫，使人像附魔一般。而那比任何魔鬼更凶残的，是它使人成为魔鬼的玩物。你所见的癫癇者，因受疾病之害，尚可同情；但醉酒者，却因自身意志而疯狂，理受憎恶，不应得任何宽恕。因为你患病，并非由于身体的本性，而是由于你自己的懈怠。况且那里有需要修补的船？因为那灵魂之船，即使尚未沉没，也已遭受巨大的船难。又是什么原因？是由于酒吗？因为船体越脆弱，船难就越彻底，无论她是自由的还是为奴的。因为自由妇女在奴隶作为观众面前失态，奴隶也同样在奴隶面前失态，她们使天主的恩赐被愚人亵渎地议论。\par

例如，我听见许多人当这些过犯发生时说：\Quote{但愿没有酒。} 哦，愚蠢！哦，疯狂！当他人犯罪时，你竟责难天主的恩赐？这是何等疯狂？什么？酒，人啊，产生了这邪恶？不是酒，而是那些在其中行邪恶之放纵者的无节制。那么你说：\Quote{但愿没有醉酒，没有奢侈；}但若你说：\Quote{但愿没有酒，}你便会逐步说下去：\Quote{但愿没有刀剑，因为凶杀者；没有黑夜，因为盗贼；没有光明，因为告密者；没有女人，因为奸淫；}总之，你会毁灭一切。\par

但你不要这样；因为这属于撒殚的心思；不要责难酒，而要责难醉酒；当你发现这同一个人清醒时，描绘出他的一切丑态，并对他说：酒被赐予，是为了使我们喜乐，而不是使我们行为不端；是为了使我们欢笑，而不是使我们成为笑柄；是为了使我们健康，而不是使我们生病；是为了矫正我们身体的软弱，而不是摧毁我们灵魂的力量。\par

天主以这恩赐尊荣了你，你为何以过分之事羞辱自己？请听保禄所说：\BibleQuote{为你的胃和你的屡次软弱，用一点酒吧。}\BibleRef{弟前 5:23} 但那圣人，即使被疾病折磨，忍受接连的疾病，直到他的导师允许他之前，都不曾饮酒；我们这些在健康中醉酒的人，将有何借口？因为祂对他说：\Quote{为你的胃，用一点酒；}但对你们每一位醉酒的人，祂会说：\Quote{用一点酒，为你的邪淫、你屡次的污言秽语，以及醉酒常产生的其他邪欲。}但即使你不愿为这些原因而戒酒，至少因随之而来的沮丧和烦恼，你们要戒酒。因为酒是为喜乐而赐予的，正如经上所说：\BibleQuote{是的，酒，使人的心喜乐。}但你甚至破坏了它的这一优点。因为什么样的喜乐是忘乎所以、有无穷的烦恼、看万物都在旋转、被头晕所压迫、如同发热病人需要有人用油浇头呢？\par

6. 这些话并非我向所有人说的；或者更确切地说，它们是对所有人说的，并非因为所有人都醉酒（愿天主阻止），而是因为那些不醉酒的人不关心醉酒的人。因此，我更是针对你们这些健康的人发言；因为医生也离开病人，转而向坐在他们旁边的人说话。因此我向你们发言，恳求你们永不被这情欲所胜，并且如同用绳索牵引那些已沉沦的人，使他们不致比禽兽更糟。因为他们实在只寻求所需之物，但这些甚至变得比它们更禽兽，超过了节制的界限。因为驴子比这些人好多少？狗好多少！因为这些动物以及所有其他动物，无论需要吃喝，都知足为限，不会超过所需；虽有无数的人强迫，它们也不愿过度。\par

在此方面，我们甚至比牲畜更不如，这不只是健康之人的判断，连我们自己也是如此。因为你们自认比狗和驴更低贱，这已向伯多禄启示，祂又借此再次确认。而且祂并未止步于此，更以祂的俯就本身宣告了这同一真理；这是极其智慧的一个实例。\par

\SourceRecord{Translated by George Prevost and revised by M.B. Riddle.；Nicene and Post-Nicene Fathers, First Series；Vol. 10.；Edited by Philip Schaff.；Buffalo, NY: Christian Literature Publishing Co.,；1888.；<http://www.newadvent.org/fathers/200157.htm>.}

% structure: p0001 source=h1 target=section reason=page-title

\section{\WorkTitle{玛窦福音讲道集 第五十八篇}}

\BibleQuote{当耶稣同门徒在加里肋亚时，耶稣对他们说：人子将被交于人之手，他们要杀害他，但第三天他要复活。他们就非常忧愁。}\par

这就是说，为防止他们说\Quote{我们为何一直待在这里}，祂再次对他们讲论受难的事；他们听到这话，就不愿再见到耶路撒冷。值得注意的是，当伯多禄被责斥后，梅瑟和厄里亚谈论了此事，并称之为光荣，圣父从天上发出声音，又行了如此多的奇迹，复活就在眼前（因为祂说祂绝不在死亡中久留，而是第三天复活），即使如此，他们仍不能忍受，反而悲伤；不仅是悲伤，而且是极其悲伤。\par

这种悲伤源于他们尚未明白祂话语的含义。马尔谷和路加间接表达此事，一个说：\BibleQuote{他们不明白这话，却害怕追问祂。}另一个说：\BibleQuote{这话为他们隐藏着，以致他们不能领悟，他们也怕问祂这话。}\par

然而，如果他们无知，又怎能悲伤？因为他们并非完全无知；他们知道祂要死，因为常听到，但这是什么死，以及会迅速解脱，并带来无数的祝福，他们尚未清楚明白；也不知道这复活是什么；但他们不明白，因此悲伤；因为他们确实非常热切地依恋他们的师傅。\par

\BibleQuote{当他们来到葛法翁，收\OriginalTerm{双德拉克马}{didrachma}的人来到伯多禄跟前说：\Quote{你们的师傅不纳双德拉克马吗？}}\par

这个\OriginalTerm{双德拉克马}{didrachma}是什么？当天主击杀了埃及人的长子后，祂就以肋未支派代替他们。后来，因为肋未支派的人数少于犹太人中长子的数目，为了使数目齐全，祂命令献纳一个\OriginalTerm{舍客勒}{shekel}；由此形成了一种习俗，即长子应缴纳这个税赋。\par

因为基督是长子，而伯多禄似乎是宗徒之首，所以他们来找他；我想，他们的习惯是在每个城市征收此税；因此在祂的家乡，他们也来找祂；因为葛法翁被认为是祂的家乡。\par

他们不敢直接找祂，而是找伯多禄；对伯多禄也不是强逼，而是温和地。因为他们并非责备，而是询问说：\Quote{你们的师傅不纳\OriginalTerm{双德拉克马}{didrachma}吗？}因为对祂的正确看法尚未形成，只如同看待一个人那样；但因着先前的神迹，他们还是向祂表示了一些尊敬和敬意。\par

那么伯多禄说什么？他说：\Quote{是的。}他对那些人说祂纳，但对耶稣却没有说，也许是因为羞于向祂提及这些事。因此那位温和者，祂明知一切，抢先说话：\Quote{\Person{西满}，你以为如何？地上的君王向谁征收关税或贡银？向自己的儿子，还是向外人？}当他说\InnerQuote{向外人}时，祂回答说：\Quote{那么，儿子就是自由的。}\par

因为免得伯多禄以为祂是从别人那里听说的，祂抢先说话，一方面表明所说的话，另一方面允许他自由发言，因他本不情愿先开口说这些事。\par

祂所说的是这样的：\Quote{我确实无需纳税。因为如果地上的君王不向自己的儿子收税，而向臣民收税，那么我更应该免于此要求，因为我不是地上君王的儿子，而是天上君王的儿子，并且我自己也是君王。}你看见祂如何区分儿子与非儿子吗？如果祂不是儿子，那么引用君王的例子也徒劳无益。有人会说：\Quote{是的，祂是儿子，但不是真\OriginalTerm{受生}{begotten}。}那么祂就不是儿子；如果不是儿子，也不是真受生的，那么祂就不属于天主，而属于其他。但如果祂属于其它，那么比较也就失去了应有的力量。因为祂所谈论的不是一般意义上的儿子，而是真正的儿子，是人们自己所生的；是与父母共享王国的儿子。\par

因此，作为对比，祂提到\Quote{外人}；所谓\Quote{外人}，是指不是他们生的，而\Quote{自己的儿子}，是指他们自己所生的。\par

我也请你们注意这一点：祂如何借此再次证实那启示给伯多禄的高超教义。祂并不止于此，反而借着祂的卑就本身，宣告了同一真理；这是极其智慧的实例。\par

\BibleQuote{但为了不冒犯他们，你去海边垂钓，拿起第一条上钩的鱼，打开它的口，就会找到一块钱；拿着去交给他们，当做我和你的税银。}\par

请看祂如何既不拒绝纳税，也不简单地命令缴纳，而是先证明自己不必纳税，然后才缴纳：一是为了避免百姓跌倒，二是为了避免祂周围的人跌倒。因为祂缴纳此税完全不是作为债务，而是为他们的软弱做最好的事。然而，在其他场合，当祂论及食物时，祂却不顾忌冒犯，\BibleRef{玛 15:11}教导我们知道何时应当顾及那些跌倒的人，何时可以不顾他们。\par

而且，藉着给予的方式本身，祂再次启示了自己。因为祂为何不命令他从他们所积存的东西中拿取呢？这是为了，如我所说，在此事上也表明祂是万有的天主，并且海洋也处于祂的统治之下。因为祂确实早已藉着斥责风浪，并藉着命令这同一位伯多禄在水面上行走，表明了这一点；\BibleRef{玛 8:26} \BibleRef{玛 14:29}但现在祂再次表明同一件事，尽管以另一种方式，却令人大为惊叹。因为预见到那条首先从深海中出来、会出现在他面前的鱼将是那纳贡的鱼，这并非小事；并将祂的命令像网一样抛入深渊，把那带钱币的鱼捉了上来；这实在是神圣而不可言喻的权能，竟能使海洋也献上礼物，并且使海洋在各方面都显示出对祂的服从——无论是它在狂怒中寂静时，还是它虽凶猛却接纳了它的同僚时，还是如今它为祂向那些索要的人纳贡时。\par

\Quote{拿给他们，}祂说，\Quote{为我和您。}你看见这极大尊荣了吗？也要看见伯多禄心灵的自主。因为这一点，这位宗徒的追随者马尔谷似乎没有记载，因为它表明了加给伯多禄的巨大荣誉；但对于否认之事，他却与其他作者一样记载了下来，而使他显赫的事他却略过不提，也许他的老师请求他不要提及关于自己的大事。祂用了\Quote{为我和您}这句话，因为伯多禄也是一个头胎的孩子。\par

正如你对基督的权能感到惊奇，我也劝你钦佩门徒的信德，因为他如此听从一件在自然上不可能的事。因为这事按本性来说确实远远超出可能。因此，为了回报他的信德，祂使他在纳贡之事上与祂自己联合。\par

3. \Quote{那时，门徒来到耶稣跟前，说：在天国里谁最大？}\par

门徒经历了一些人性软弱的感受；因此福音作者也附加了这句注释，说：\Quote{那时；}即在祂将他置于众人之上之后。因为雅各伯和若望中也有一个头胎的儿子，但祂并没有为他们行这样的事。\par

然后，他们羞于承认自己的感受，没有公开地说：\Quote{为什么祢将伯多禄置于我们之上？}或\Quote{他比我们大吗？}因为他们羞愧；而是含糊地问：\Quote{那么谁最大？}因为当他们看见那三位被拣选时，他们没有这样的感受；但如今荣誉归于一人，他们便烦恼了。不仅如此，还有许多其他的事汇集起来点燃了这种感受。因为祂曾对他说：\Quote{我要把钥匙给你；}\BibleRef{玛 16:19}对他说：\Quote{你有福了，西满巴约纳；}在此处对他说：\Quote{拿给他们，为我和您；}而且他们普遍看见他讲话多么自由，这使他们有些懊恼。\par

如果马尔谷说，\BibleRef{谷 9:34}他们不是问，而是在心中思量，这与此并不矛盾。因为很可能他们两样都做了，既问过，也彼此商议过；而在以前，在别的场合，他们曾有过这种感受，一次两次，如今他们既表明了感受，也彼此思量了。\par

但我对你们说：\Quote{不要只看他们的过犯，也要考虑这一点：首先，他们寻求的不是现世的事物；其次，这种情欲后来也被他们抛开了，并且彼此让出了首位。}但我们现在连他们的过失也达不到；我们也不寻求\Quote{在天国里谁最大}，而是寻求谁在世俗国度里最大，谁最富有，谁最有权势。\par

那么基督说什么？祂揭露了他们的良心，并回应了他们的感受，而不仅仅是他们的话。\Quote{因为祂叫了一个小孩子来，}经上说，\Quote{并且说：你们若不悔改，变成如同小孩子一样，决不能进入天国。}\BibleRef{玛 18:2}\Quote{你们，}祂说，\Quote{询问谁最大，并且为头等的荣誉争论不休；但我说，那没有变得最卑微的人，配不上进入那里。}\par

祂不仅引用了这个榜样，而且还把小孩子放在他们中间，以眼见的情景使他们羞愧，并说服他们也要同样谦卑和纯朴。因为小孩子既无嫉妒，也无虚荣，又无对首位的渴望；他拥有最大的德行——纯朴，以及一切天真和谦卑的品性。\par

因此，不仅需要勇气和智慧，还需要这种德行——我指的是谦卑和纯朴。是的，如果我们没有这些德行，那么属于我们救恩的事，甚至在最关键之处也会跛足。\par

小孩子，无论被侮辱挨打，还是被尊崇受赞美，都不会因前者而心生急躁或嫉妒，也不会因后者而骄傲自大。\par

你看见祂又如何召叫我们走向一切自然的优越，表明凭自由意志可以达致这些，从而平息了摩尼派\OriginalTerm{摩尼派}{Manichæans}邪恶的狂热吗？因为如果本性是邪恶的，为什么祂要从这里引出严峻美善的榜样呢？\par

我以为祂放在中间的那个小孩子是一个非常年幼的孩子，完全摆脱了所有这些情欲。因为这样的小孩子没有骄傲、虚荣的疯狂、嫉妒、好争以及所有这类情欲；而拥有许多美德——纯朴、谦卑、不恋世俗——却不为任何美德自夸；这是一种双重的严峻美善：拥有这些，却不因此而自大。\par

因此祂将小孩子带来，放在他们中间；并且祂不只以这一点结束祂的讲论，而是进一步发挥这个训诫，说：\Quote{无论谁因我的名而接待这样一个孩子，就是接待我。}\par

\Quote{因为知道吧，}祂说，\Quote{你们不但自己变成这样，会得到大赏报；而且若因我的缘故尊敬这样的人，我也因你们尊敬他们而给你们天国作为赏报。} 或者更确切地说，祂所设定的远更伟大，说：\Quote{他接纳我。凡谦卑而纯朴的一切，在我眼中如此珍贵。} 因为此处祂用\Quote{一个小孩子}，意指那些如此纯朴、谦卑、且在常人眼中卑微可鄙的人。\par

4. 此后，为使祂的话更被接纳，祂不仅藉尊荣，也藉惩罚来确立，继续说：\Quote{无论谁使这些小子中的一个跌倒，倒不如把一块磨石拴在他的颈上，沉在海的深处。}\par

\Quote{因为正如，}祂说，\Quote{凡为我而尊敬这些人的人，拥有天国，甚或拥有比天国更大的尊荣；同样，凡羞辱他们的人（因为这即是使他们跌倒），将受极重的惩罚。你们不必惊奇祂称侮辱为\InnerQuote{得罪}；因为许多软弱的人因遭受轻慢和侮辱而受到不小的伤害。因此，为加重和加剧责备，祂陈述了由此产生的祸害。}\par

祂并没有以同样的方式继续表达惩罚，而是从我们熟悉的事物中表明它是如何不可忍受。因为当祂要最尖锐地触动粗俗之人时，便引入可感知的形象。因此，在此处为指明他们将要受的惩罚之重大，并挫败轻蔑者的傲慢，祂提出了一种可感知的惩罚，即磨石和淹死。然而，按照前文，祂本应说：\Quote{凡不接纳这些小子中的一个的，就是不接纳我；}这比任何惩罚都更痛苦；但由于那麻木不仁、极其粗俗之人并不被此（尽管可怕）所触动，祂便设了\Quote{一块磨石}和\Quote{淹死}。祂并没有说：\Quote{一块磨石要拴在他的颈上，}而是\Quote{他倒不如}受此罚；暗示更严重的另一种惩罚等候着他；如果这已是不可忍受，那么那更甚者更是如此。\par

你看到祂如何从两方面使威胁可怕：首先通过与已知形象的比较使之更鲜明，然后藉其更甚之处使之在想象中远大于那可见的惩罚。你看到祂如何根除傲慢的精神；如何医治虚荣的溃疡；如何教导我们在任何事上不追求尊荣；如何劝勉那些贪求一切尊荣的人去寻求最低的位置？\par

5. 因为再没有比傲慢更坏的了。这甚至使人失去理智，使他们显出愚昧之相；或者不如说，它确实使他们完全像傻子一样。\par

因为正如有人身高三尺，却竭力想比山还高，或甚至以为如此，并挺起身来，以为超过了山顶，我们无需再寻其他证据证明他失去了理智；同样，当你看见一个傲慢的人，自以为超过众人，并以与常人共处为耻，你无需再寻求其他证据证明此人的疯狂。他比任何天生的傻子更可笑，因为他完全刻意制造这种疾病。不仅在此他是悲惨的，而且因为他不知不觉堕入了邪恶的深渊。\par

因为这样的人何时会正确认识任何罪？何时会察觉自己在得罪人？不，他更如一个卑贱被掳的奴隶，被魔鬼抓住带走，完全成为猎物，四面受到攻击，遭受无数侮辱。\par

因为魔鬼最终引他们进入如此大的愚昧，以至他们竟对自己的孩子、妻子，甚至自己的祖先傲慢。反之，他也使另一些人因祖先的显赫而骄傲。还有比这更愚蠢的吗？当人们因相反的原因同样自高自大——一些人因他们的父亲、祖父和祖先卑贱，另一些人因他们的祖先光荣而杰出——怎能分别减轻这种肿胀的毒疮呢？可以对后者说：\Quote{追溯到你的祖父和直系祖先之外，你会找到许多厨子、赶驴的和小贩；}而对那些因祖先卑贱而自傲的人，则相反：\Quote{而你，若在你祖先中继续上溯，你会发现许多比你更显赫的人。}\par

因為這本性有此歷程，且讓我甚至從聖經向你證明這一點。撒羅滿是君王之子，且是顯赫君王的兒子，但那位君王的父親卻出身卑賤且不名譽。而他的外祖父也是如此；否則他不會將女兒許配給一個普通的士兵。如果你再從這些卑微的人往更高處追溯，你會看到更顯赫、更王室的世系。在撒烏耳的例子也是如此，在許多其他人身上也是如此，人必會得出同樣的結論。因此，我們不可在此自誇。因為出身是什麼？告訴我。不過是一個名稱，沒有實質；你將在那一日知道這事。但因為那一日尚未來到，現在就讓我們從眼前的事物說服你，從中不會產生任何優越感。因為若戰爭臨到我們，若饑荒，或其他任何事，所有這些出身高貴的膨脹虛榮都將受到考驗：若疾病、瘟疫臨到我們，它無法區分富人和窮人、榮耀的和無榮耀的、出身高貴和並非如此的人；死亡也是如此，命運的其他逆轉也是如此，但它們都同樣攻擊所有人；若我可以說出一些奇妙的，更是針對富人。因為他們越是少受這些事的鍛鍊，當被這些事追上時就越發滅亡。而且恐懼在富人中也更大。因為沒有人像他們那樣在君王面前顫抖；在群眾面前，也不亞於在君王面前，甚至更多；事實上，許多這樣的家室曾因群眾的憤怒和君王的威脅而傾覆。但窮人卻免受這兩種動盪之水。\par

6. 因此，放下這貴族身份吧，你若願向我顯示你是高貴的，就顯示你靈魂的自由，正如那位有福之人（他是一個窮人）所擁有的，他對黑落德說：\Quote{你不可佔有你兄弟斐理伯的妻子；}\BibleRef{谷 6:18}也正如那在他之前、與他相似、此後也將如此的人所擁有的，他對阿哈布說：\Quote{我並非擾亂以色列，而是你和你父家；}\BibleRef{列上 18:18}正如先知們所擁有的，正如所有宗徒所擁有的。\par

但被財富奴役之人的靈魂並非如此，反而如同在千萬導師和工頭之下的人，他們甚至不敢抬起眼睛，為義理而勇敢發言。因為對財富、對榮耀、以及其他事物的貪愛，可怕地注視著他們，使他們成為奴顏婢膝的諂媚者；沒有什麼比糾纏於世俗事務、佩戴那些被視為尊榮標誌的東西更能剝奪自由了。因為這樣的人不是有一位主人，也不是兩位、三位，而是千萬位。\par

若你樂意甚至數算他們，讓我們引進一位在君王宮廷中受尊榮的人，讓他既擁有極多的財富，又擁有極大的權力，出生地超越他人，祖先顯赫，且被眾人仰望。現在讓我們看看，這人是否正是比所有人更為奴役的那位；讓我們將他與一個不僅是奴隸，而且是奴隸的奴隸相比，因為許多僕人也有奴隸。那麼這個奴隸的奴隸只有一位主人。雖然那一位不是自由人，卻只有一位，而這只討他的歡心。因為即使他主人的主人看似對他擁有權力，但目前他只服從一人；若他們兩人之間關係良好，他一生都將安居無憂。但我們這位先生不僅有一位或兩位主人，而是有許多位，且是更苛刻的主人。首先他關心君主本人。擁有一個卑微之人為主人，與擁有一個國王為主人是不同的，國王的耳朵被許多人嗡嗡作響，他時而成為這群人時而成為那群人的私有財產。\par

我們這位先生，雖然心中無愧，卻懷疑一切；既懷疑他的同僚和部下，也懷疑他的朋友和敵人。\par

但你可以說，另一個人也害怕他的主人。然而，擁有一位或許多位主人，使人膽怯，這豈是同一回事？或者，若人仔細探究，他將發現連一位也沒有。如何？在什麼意義上？因為那個奴隸沒有人想要把他從他的服務中趕走並取代自己（因此也沒有人陰謀反對他）；而這些人不僅沒有其他追求，反而要動搖那位更受統治者喜愛和愛戴的人。因此他也必須諂媚所有的人：他的上級、同僚、朋友。因為哪裡有嫉妒和愛榮耀，那裡甚至真誠的友誼也無法堅固。就像同行業的人不能以完全真摯的愛彼此相愛，那些在榮譽上競爭、以及追求世俗事物的人也如此。因此內戰也很大。\par

你看見一群主人，而且是苛刻的主人嗎？你願我向你顯示另一個比這更嚴重的嗎？那些在他後面的人，全都努力趕在他前面；那些在他前面的人，都阻止他靠近他們並超越他們。\par

7. 啊，奇蹟！我原本承諾向你們展示主人，但我們的言論，我們發現，在繼續並變得熱切時，所做的比我承諾的更多：指出了敵人而非主人；或者更確切地說，同一批人既是敵人又是主人。因為當他們像主人一樣被奉承時，他們像敵人一樣可怕，而且他們像仇敵一樣陰謀反對我們。那麼，當一個人擁有同一批人既是主人又是敵人時，還有什麼比這災難更糟糕的呢？奴隸雖然服從命令，但無論如何仍從發號施令者那裡得到關懷和善意；而這些人，雖然接受命令，卻成了敵人，彼此對立；而且比戰場上的人更嚴重，因為他們暗中傷害，在朋友的面具下對待他人如同敵人，並且常常利用他人的災難為自己贏得聲譽。\par

但我们的情况并非如此；相反，若有人遭遇不幸，便有许多人与他一同悲伤；若他获得尊荣，也有许多人与他一同喜乐。宗徒却不是这样：他说：\Quote{因为若是一个肢体受苦，所有的肢体都一同受苦；若是一个肢体得荣耀，所有的肢体都一同欢乐。}\BibleRef{格前 12:26} 而且，给予这些劝诫者的话，时而说：\Quote{我的希望或喜乐是什么呢？岂不是你们吗？}\BibleRef{得前 2:19}，时而说：\Quote{现在若你们在主内站稳，我们就活了。}\BibleRef{得前 3:8}，时而说：\Quote{我怀着极大的忧苦和痛苦的心给你们写了信。}\BibleRef{格后 2:4}，以及\Quote{谁软弱，我不软弱？谁跌倒，我不焦急呢？}\BibleRef{格后 11:29}。\par

那么，我们为什么还要忍受外界的风浪波涛，而不奔向这个平静的港湾，并离开那些美名的影子，去追求事物本身呢？因为对于他们而言，光荣、尊严、财富、声望以及诸如此类，都是空名；对于我们而言，却是现实；正如那些痛苦的事——死亡、羞辱、贫穷以及任何类似之物——在我们看来是空名，在他们看来却是现实。\par

如果你愿意，让我们先提出光荣，这是他们所有人都喜爱和渴望的。我不是说它短暂易逝，而是说当它在鼎盛之时，请你把它指给我看。不要除去那个娼妓的脂粉和彩线，而是把她打扮好带出来，展示给我们，好让我随后揭露她的丑陋。那么，你当然会说到她的排场、她众多的扈从、传令官的声音、各阶层的倾听、民众的肃静、对挡路者的鞭打以及普遍的注视。这些难道不是她的辉煌吗？那么，让我们来考察这些事物是否虚妄，是否只是无益的幻想。因为我们所谈论的那个人，凭借这些事物，在身体或灵魂上得到了什么好处？因为这才是构成人的要素。他会因此更高大、更强壮、更健康、更敏捷吗？或者他的感官会更敏锐、更透彻吗？不，没有人能这样说。那么，让我们转向灵魂，看看是否在那里可以发现任何由此而来的益处。那么，通过那种侍从，这样的人会更节制、更温和、更明智吗？绝不，反而完全相反。因为在这里，结果并不像在身体上那样。因为在身体方面，身体确实没有获得任何属于其自身卓越的东西；但在这里，祸害不仅在于灵魂没有结出任何善果，而且实际上还因此接受了诸多邪恶：因为这等事物使它陷入傲慢、虚荣、愚蠢、愤怒以及成千上万类似的过失。\par

你会说：\Quote{但他因这些事而欢乐、欣喜，这些事使他容光焕发。}你的这句话正是他邪恶的极点，也是疾病不可救药的部分。因为因这些事而欢乐的人，虽然很容易摆脱罪恶的根源，却不愿意；相反，他因这快乐而堵住了自己治愈的道路。所以，最糟糕的就在这里：当疾病在他身上增长时，他不但不痛苦，反而欢乐。\par

因为欢乐并不总是好事；就连盗贼也因偷窃而欢乐，奸夫因玷污邻人的婚床而欢乐，贪婪者因强夺而欢乐，杀人者因谋杀而欢乐。那么，我们不要只看他是否欢乐，而要看是否为了有益的事，以免我们发现他的欢乐如同奸夫和盗贼的欢乐一样。\par

因为，请告诉我，他为什么欢乐？为了在群众中的声望，因为他能自高自大，被人注视？不，有什么比这种欲望和这种不当的偏爱更糟糕呢？或者，如果这不是坏事，你们就必须停止嘲笑虚荣的人，并停止不断地以讥讽辱骂他们：你们也必须停止诅咒傲慢和轻蔑的人。但你们不会忍受这样的事。那么，他们也应受许多责备，尽管他们有许多扈从。我所说的这一切是关于那些较为可容忍的统治者；因为我们会发现大部分统治者比强盗、杀人犯、奸夫、盗墓者更严重地犯罪，因为他们没有善用权力。的确，他们的盗窃更无耻，他们的屠杀更冷酷，他们的不洁远比其他人更严重；他们挖掘的不是一道墙，而是无尽的庄园和房屋，他们的特权使他们非常容易做到这些。\par

他们服侍着最痛苦的奴役，既卑屈地屈服于自己的情欲，又惧怕所有的同谋者。因为只有摆脱了情欲的人，才是自由的，才是统治者，比所有君王更威严。\par

因此，知道了这些事，让我们追求真正的自由，从邪恶的奴役中解救自己，不要认为权力的显赫、财富的统治或任何其他类似的事物是幸福的；而只认为德行是幸福的。这样，我们在此世就能享有安全，并因我们的主耶稣基督的恩宠和慈爱而获得未来的美好事物。愿光荣和权能归于祂，与圣父及圣神，世世无尽。阿们。\par

\SourceRecord{Translated by George Prevost and revised by M.B. Riddle.；Nicene and Post-Nicene Fathers, First Series；Vol. 10.；Edited by Philip Schaff.；Buffalo, NY: Christian Literature Publishing Co.,；1888.；<http://www.newadvent.org/fathers/200158.htm>.}

% structure: p0001 source=h1 target=section reason=page-title

\section{关于玛窦福音的讲道集 59}

\Emph{\BibleRef{玛 18:7}}\par

\BibleQuote{祸哉，世界因恶表而受祸！因为恶表是免不了的，但使恶表临于那人的人有祸了。}\par

\BibleQuote{如果恶表是免不了的}（我们的对手或许会说）\Quote{那么祂为什么哀叹世界，而不更应伸出援手相助？因为这才是一位医生和保护者的本分，而那后者就连任何普通人都能做到。}\par

那么，对于如此无耻的舌头，我们能说些什么呢？不，你还要寻找什么与祂这种医治关怀相配的事呢？因为天主为你成为人，取了奴仆的形体，经历了一切极端，没有留下任何祂所应做的事不做。但既然忘恩的人对此毫无改善，祂就为他们哀叹，因为经过如此多的培育关爱，他们仍处于病态之中。\par

这就像有人为那个得到许多照顾却不愿听从医生嘱咐的病人哀叹说：\Quote{这人有祸了，因他的病弱因自己的疏忽而加重。}但在那种情况下，哀叹毫无益处；而在这里，预告将要发生的事并为之哀叹，本身也是一种治疗。因为许多人虽然被劝诫时毫无效果，但被哀悼时却会改正。\par

正因如此，祂最常用\Quote{祸哉}这个词，彻底唤醒他们，使他们认真起来，并促使他们保持警醒。同时，祂也显示了对那些人的善意和祂自己的温和，即使他们反对，祂仍为他们哀叹，不只是厌恶，而是通过哀叹和预言来纠正他们，以便赢得他们。\par

\Quote{但这怎么可能？}祂或许会说。\Quote{因为如果\BibleQuote{恶表是免不了的}，又怎能逃脱这些呢？}因为恶表固然必须来，但人灭亡并非必然。就像一位医生（我们不妨再用这个比喻）说：这病一定会来，但注意的人不一定就被病摧毁。祂这样说，如我所说，是为了唤醒众人，尤其是祂的门徒。因为他们被派往和平安宁的生活，不可懈怠，祂就预示他们内外都有许多争战。保禄宣告这事说：\BibleQuote{外有争战，内有恐惧。}\BibleRef{格后 7:5}又：\BibleQuote{在假弟兄中的危险。}\BibleRef{格后 11:26}祂在向米肋托人的讲话中也说：\BibleQuote{你们中也要有人起来讲说谬论。}\BibleRef{宗 20:30}祂自己也说：\BibleQuote{人的仇敌就是自己家里的人。}\BibleRef{玛 10:36}但当祂说\BibleQuote{恶表是免不了的}时，并非剥夺他们选择的能力，也不是否定道德原则的自由，也不是将人的生活置于任何绝对必然的约束之下，而是预告必将发生的事。路加用另一种方式表述：\BibleQuote{恶表是不可能不来的。}\BibleRef{路 17:1}\par

什么是恶表？就是正道上的障碍。舞台上的演员也这样称呼那些精通此道、扭曲身体的人。\par

所以，不是祂的预言带来了恶表；绝非如此。也不是因为祂预告了，因此就发生；而是因为事情必然发生，所以祂预告了。因为如果那些引进恶表的人不愿意作恶，恶表也不会来；如果恶表不会来，也就不会被预告。但因为那些人行恶，且病入膏肓，恶表就来了，祂预告了将要发生的事。\par

有人说，如果这些人持守正道，没有人引进恶表，这话岂不显为虚假？决不会，因为这也不会说出来。因为如果所有人都持守正道，祂就不会说\BibleQuote{恶表是免不了的}，而是因为祂预知他们自己不会改过，所以说恶表必定会来。\par

又有人说，祂为什么不除掉恶表呢？为什么要除掉呢？为了那些受害的人吗？但受害者的毁灭不是来自恶表，而是来自他们自己的懈怠。正直的人证明了这一点，他们不仅不受害，反而大大受益，例如约伯、若瑟、所有的义人和宗徒。如果许多人灭亡，那是在于他们自己的沉睡。如果不是这样，而是毁灭出于恶表，那么所有人都必灭亡。但既然有人逃脱，让那没有逃脱的人归咎自己吧。因为恶表，如我所说，不仅唤醒和使他更加警醒、敏锐，也唤醒那跌倒的人，如果他迅速站起来，恶表会使他更加安全，更难被打败。因此，如果我们警醒，就能从中获得巨大益处，即保持警醒。因为当我们有仇敌，当如此多的危险压迫我们时，我们还在睡觉，如果生活在安逸中，我们会怎样呢？如果你愿意，请看第一个人。他在乐园里住了很短时间，也许不到一整天，享受了快乐，却堕入如此邪恶的地步，甚至想与天主平等，把欺骗者当作恩人，连一条诫命也不遵守；如果他在余生中没有苦难，他会做出什么事呢？\par

2. 但我们说这些话时，他们又提出其他反对意见，问道：\Quote{那么天主为何把他造成这样？}天主并没有把他造成这样，绝非如此，否则祂也不会惩罚他。因为在那些我们自身是原因的事情上，我们尚且不责备自己的仆人，何况万有的天主呢？\Quote{但这又是如何发生的？}有人会说。\Quote{出于他自己和他的懈怠。}\Quote{所谓\Quote{出于他自己}是什么意思？}请问你自己。因为如果恶人不是出于自身而为恶，那么就不要惩罚你的仆人，也不要责备你的妻子所犯的过失，也不要打你的儿子，也不要责怪你的朋友，也不要恨那欺侮你的仇敌：因为所有这些人都值得怜悯，而非惩罚，除非他们乃是出于自身而犯罪。\Quote{但我无法做到自制。}有人会说。然而，当你认识到原因不在于他们，而在于另一种必然性时，你就能做到自制。至少当仆人患病，未能完成所吩咐的事时，你非但不责备，反而原谅他。可见你自己作证：一件事是出于自身，另一件事则不是出于自身。因此，在此处也是如此，如果你知道他生来就是邪恶的，你非但不会责备，反而会宽容他。因为既然你因他的疾病而原谅他，那么如果他从起初就被造成这样，你岂会不原谅天主的造化呢？\par

从另一个角度也很容易堵住这些人的口，因为真理的力量极其丰盛。为什么你从不责备你的仆人，因为他容貌不美，身材不高，不能飞翔？因为这些是自然的。所以，对于违反本性的缺点，他是无罪的，无人能反驳。因此，当你责备时，你表明过错不在于本性，而在于他的选择。因为对于那些我们不责备的事，我们承认完全是出于本性，那么显然，当我们责备时，我们宣称冒犯是出于选择。\par

因此，我恳求你，不要提出悖逆的推理，也不要提出比蜘蛛网更纤细的诡辩和网罗，而是再回答我这个问题：天主造了所有的人吗？这对每个人来说都很清楚。那么，为什么在德行和恶习方面，并非人人平等？善良、温良、谦和的人从何而来？无价值和邪恶的人从何而来？如果这些事情不需要任何意向，而是出于本性，那么为何一些人如此，另一些人如彼？因为如果所有人本性都是恶的，那么就不可能有任何人成为善的；但如果本性是善的，那么就没有人是恶的。因为如果所有人的本性相同，那么在这方面他们都必须是一样的，无论他们是善还是恶。\par

但如果我们说，一些人天生是善的，另一些天生是恶的，那将是不合理的（正如我们已表明的），因为这些东西必须是不可改变的，因为本性的事物是不可改变的。注意：所有会死的人都容易受苦；没有人能免除痛苦，尽管他无止境地努力。但现在我们看到许多好人变成坏人，许多坏人变成好人，前者因懈怠，后者因勤勉；这最有力地表明这些事情并非出于本性。\par

因为本性的事物既不会改变，也不需要勤奋去获得。就像对于看和听，我们不需要辛劳，同样，如果德行是自然分配的，我们也就不需要在德行上劳苦了。\par

\Quote{但祂为何要造出无价值的人呢？祂本可以造所有的人都是善的。那么恶事从何而来？}他说。请问你自己；因为我的责任是表明它们不出于本性，也不来自天主。\par

\Quote{那么它们是出于自身吗？}他说。绝不是。\Quote{那么它们是无始的吗？}说话要恭敬，人啊，远离这种疯狂，将同样的尊荣——而且是至高的尊荣——归于天主和恶事。因为它们若是无始的，便有力量，甚至不能被拔除，也不能归于虚无。因为无始之物是不朽坏的，这对所有人都是显然的。\par

3. 既然恶有如此大的力量，为什么又有如此多的善人？那有起源的，如何比那无始的更强呢？\par

\Quote{但天主会毁灭这些东西。}他说。何时？祂将如何毁灭那些与自己同等尊荣、同等力量、可说同样永恒的事物呢？\par

哦！魔鬼的恶毒！它发明了多大的恶！它说服人用何等的亵渎包围天主！它披着敬虔的外衣，又杜撰出另一种渎神的说法！因为要表明恶不是出于祂，他们引入了另一种邪恶的教义，说这些东西是无始的。\par

\Quote{那么恶从何而来？}有人会说。\Quote{来自愿意和不愿意。}\Quote{但我们愿意和不愿意这件事本身，又从何而来？}出于我们自己。你这样问，就好像当你问\Quote{看和看不见从何而来？}然后我说\Quote{来自闭眼或不闭眼}，你又问\Quote{闭眼或不闭眼本身从何而来？}然后你听说那是出于我们自己和我们的意志，你又去寻找另一个原因。\par

因为恶无非就是违背天主。\Quote{那么人如何发现这个？}有人会说。\Quote{怎么，发现这个很难吗？我请问。}\Quote{不，我并非说这件事困难；而是他为何变得想要违背？}\Quote{出于懈怠。因为他有能力作两者，却倾向了这边。}\par

但你若听了这话仍感困惑和头晕，我且问你一个不难也不复杂，而是简单明了的问题。你是否有时变坏？有时又变好？我的意思是这样的：你曾否克服了情欲，却又被情欲所制？你曾否被醉酒胜过，又胜过了醉酒？你曾一度动怒，又一度不动怒？你曾忽视穷人，又曾不忽视他？你曾行淫，又变得贞洁？那么，这一切是从何而来？告诉我，从何而来？你若不告诉我，我便说：这是因为你一度克制自己并努力，后来却变得懈怠和疏忽。因为对那些绝望的、常行恶的、麻木的、疯狂的、甚至不愿听那能改正他们的话的人，我连谈论克制也不愿；但对那些时而处于一种状态、时而处于另一种状态的人，我乐意说话。你曾一度强取不属于你的东西，之后又因怜悯，甚至将自己的东西分给有需要的人？那么这改变是从何而来？难道不是很明显来自理智和意志的选择吗？\par

这很明显，没有人不这么说。因此我恳求你认真起来，紧守美德，你就不需要这些问题了。因为如果我们愿意，我们的恶事仅仅是名称而已。所以不要再追问恶从何来，也不要困扰自己；既然已经发现它们仅来自懈怠，就要逃避恶行。\par

若有人说这些事不是出于我们，你看见他对仆人发怒，对妻子动气，责备孩子，谴责那些伤害他的人时，就对他说：那么你怎能说恶事不是出于我们呢？因为若不是出于我们，你为何指责呢？再说：你辱骂和侮辱是出于你自己吗？若不是出于你自己，就无人该对你发怒；但若是出于你自己，你的恶行便是出于你自己和你的懈怠。\par

然而怎样？你以为有一些好人吗？因为若真没有好人，你从何处得到这个词？赞美从何而来？但若有好人，显然他们也会责备恶人。然而，若没有人是自愿作恶，也不是出于自己，那么好人就会被发现是在不公正地责备恶人，而他们自己也因此变成恶人。因为还有什么比使无辜者遭受指责更坏的呢？但如果他们在我们看来仍是好人，虽然责备，而这恰恰是他们善的证明，那么即使是愚昧人也能由此明白，从来没有人是出于必然为恶的。\par

但若在这之后你仍要追问恶从何来？我要说，来自懈怠，来自懒惰，来自与恶人交往，来自轻视美德；因此既有恶事本身，也有人追问恶从何来的事实。因为行正直之事、立志公平节制生活的人，确实无人追问这些事；但那些敢于行恶、并想通过这类讨论为自己编造愚蠢安慰的人，是在织蜘蛛网。\par

但让我们不仅以言语，也以行动将这些撕碎。因为这些事情也不是必然的。因为若是必然的，他就不会说：\BibleQuote{那引入绊脚石的人有祸了。}（见：\BibleRef{玛 18:7}）因为他只哀叹那些因自己选择而作恶的人。\par

若他说\BibleQuote{凭着他}，不要惊奇。因为他这样说，并不是说别人藉着他带来绊脚石，而是视他本人为整个事件的起因。因为圣经常以\Quote{凭着他}代替\Quote{出于他}；例如经上说：\BibleQuote{我靠天主得了这个人}，这不是指第二因，而是指第一因；又说：\Quote{解梦不是出于天主吗？}（见：\BibleRef{创 40:8}）以及\BibleQuote{天主是信实的，你们是由他所召、与他的圣子耶稣基督我们的主契合的。}（见：\BibleRef{格前 1:9}）\par

4. 为使你们知道这不是必然的，也请听接下来的话。因为他哀叹了他们之后，说：\BibleQuote{倘若你的手或你的脚使你跌倒，砍下它来，从你身上扔掉；你残废或瘸腿进入生命，比有双手双脚被投入火里更好。倘若你的右眼使你跌倒，剜出它来；你独眼进入生命，比有双眼被投入地狱的火里更好。}他不是说肢体，绝不是；而是说朋友、亲戚，那些我们视同必要肢体的人。这话他以前说过，现在又说。因为没有什么比坏同伴更有害。强制不能做到的事，友谊常能做到，无论是为害还是为益。因此他极其恳切地命令我们砍断那些伤害我们的人，暗示那些引人绊倒的人。\par

你看到他如何消除了由绊脚石可能带来的祸害？他预言必定有绊脚石来到，这样就没有人处于懈怠状态，而是因期待而警醒。他指其祸害之大（因为他不会无端说\BibleQuote{因绊脚石的事，世界有祸了}，而要显示由此而来的祸害之大），又以更强烈的语气哀叹那引人绊脚石的人。因为\BibleQuote{但那人有祸了}这句话，表明惩罚极大；不仅如此，他还藉着添加的比较增加了恐惧。\par

然后他不仅满足于这些，还指出了人可避免绊脚石的道路。\par

但这道路是什么呢？他说：恶人，即使是你极其亲爱的朋友，也要从你的友谊中除掉。\par

祂又提出一個無可辯駁的理由。因為如果他們繼續作朋友，你並不能贏得他們，反而連自己也喪失；但若你與他們斷絕，至少你能獲得自己的救恩。所以，若任何人的友誼傷害你，就應與之斷絕。因為如果我們自己的肢體中，常有許多既無可救藥、又毀壞其餘部分時，我們尚且將其割除，那麼對待朋友更當如此。\par

但若邪惡出於本性，那麼這一切告誡與勸導便是多餘的，藉所提及的方法來預防也是多餘的。但若並非多餘——誠然不是多餘——則很明顯，邪惡是出於意志的。\par

\BibleQuote{你們小心，不要輕視這些小子中的一個；我告訴你們：他們的天使在天上，常見我在天之父的面。}\par

祂稱之為小子的，並非真正的小子，而是那些被眾人視為小子的，即窮人、受輕視者、無名之輩（因為那與全世界等價的人，怎能算為小子？那為天主所鍾愛的人，怎能算為小子？），而是那些在眾人心目中如此被看輕的人。\par

祂不僅談及許多人，甚至只談一個，藉此再次防範眾多絆腳石的傷害。因為正如逃避惡人有很大的益處，同樣尊敬善人也有極大的益處，且對留心者是一重雙保險：一是根除與犯罪者的友誼，二是因尊敬與榮耀這些聖者而得保護。\par

然後，祂又用另一種方式使他們成為可敬的，說：\BibleQuote{他們的天使在天上，常見我在天之父的面。}\par

由此可見，聖人們有天使，或者說所有人都有。因為宗徒論及婦女時也說：\BibleQuote{她為天使的緣故，應當在頭上有權威。} \BibleRef{格前 10:10} 梅瑟也說：\BibleQuote{祂按天主的眾天使的數目，劃定萬民的疆界。} \BibleRef{申 32:8}\par

但在這裡，祂不僅論及天使，更是指出那些比別的天使更大的。當祂說：\Quote{我在天之父的面}，祂不過是指他們更圓滿的信任與極大的尊榮。\par

\BibleQuote{因為人子來，是為拯救那迷失的。}\par

祂又提出了另一個比先前更強的理由，並以一個比喻與之相連，由此表明聖父也同樣渴望這些事。祂說：\BibleQuote{你們以為如何？若一個人有一百隻羊，其中一隻迷了路，他豈不撇下那九十九隻，往山裡去尋找那迷路的嗎？若找著了，他為這一隻歡喜，比為那沒有迷路的九十九隻歡喜更大。同樣，你們的天父也不願意這些小子中的一個喪亡。} \BibleRef{瑪 18:12}\par

你看，祂用了多少事情來敦促我們關心我們卑微的弟兄？那麼，不要說：\Quote{這人是鐵匠、鞋匠、農夫、愚笨人}，因而輕視他。為了使你不會如此感受，你看，祂用多少動機說服你操練謙卑，並督促你關心這些人。祂擺出一個小孩子，並說：\Quote{要像小孩子。} 又說：\Quote{凡接待這樣一個小孩子的，就是接待我；} 又說：\Quote{凡絆倒人的，} 必要受極重的刑罰。祂不僅以\Quote{磨石}的比喻為滿足，還加上\Quote{禍哉}，並命令我們與這樣的人斷絕，即使他們對我們如同手和眼一般。祂又藉著託管這些卑微弟兄的天使，使他們成為可敬的；並藉著自己的意願與苦難（因為當祂說：\Quote{人子來是為拯救那迷失的，} 祂暗示了十字架，正如保祿論及一位弟兄時說：\Quote{基督為他死了}）；又藉著聖父，因為連祂也不願一人喪亡；又藉著通常的習俗，因為牧人撇下安全的羊，去尋找那迷失的；當他找到那迷路的羊時，因尋獲與救回而極其喜悅。\par

5. 如果天主對被尋回的小子這樣喜悅，你怎能輕視那些天主所殷切眷顧的人呢？甚至為這些小子中的一個，人應當捨棄自己的性命。但他弱小卑微嗎？正因為此，我們更應當盡一切努力保全他。因為祂自己撇下九十九隻羊，去尋找這一隻，那麼多羊的安全並不能掩蓋這一隻的損失。但路加說，祂甚至把羊扛在肩上，並且\BibleQuote{對於一個罪人悔改，比對九十九個無須悔改的義人，歡喜更大。} \BibleRef{路 15:7} 祂為此撇下那些得救的羊，並對這一隻更感喜悅，顯出祂對這羊的殷切是多麼大。\par

因此，我們不要對這樣的靈魂漠不關心。因為這一切話都是為此目的。祂以警告說，若不成為小孩子，就絕不能進天國，並提到\Quote{磨石}，以此壓制驕傲者的傲慢；因為再沒有比驕傲更與愛德為敵的了。祂說\Quote{絆腳石是免不了的}，使我們警醒；又加上\Quote{但絆倒人的有禍了}，使每人努力不讓絆倒出於自己。祂命令割除那些絆倒的人，使救恩變得容易；祂又吩咐不要輕視他們——不僅是吩咐，而且帶著熱切（因為祂說：\Quote{你們小心，不要輕視這些小子中的一個}）；祂說：\Quote{他們的天使常見我在天之父的面}，以及\Quote{我為此而來}，和\Quote{我父願意這樣}，使那些應照顧他們的人更加殷勤。\par

你們看，祂在他們周圍設置了何等樣的圍牆，對那卑微而將滅亡的人又是何等切切的眷顧，既以不治之災威脅那使他們跌倒的人，又以極大的祝福應許那伺候他們、照管他們的人，並且再次從祂自己和聖父引出榜樣。\par

讓我們也效法祂，為了弟兄的緣故，不拒絕任何看似卑下艱難的任務；即使我們必須服務，即使被服務者渺小卑微，即使工作辛勞，即使我們必須翻山越嶺，也當為了弟兄的救恩把一切視為可忍受的。因為靈魂為天主所如此切切眷顧，以致\BibleQuote{祂沒有怜惜自己的兒子}。\BibleRef{羅 8:32}\par

因此我懇求，當清晨出現，我們一出家門，就要懷著這唯一目標，這超過一切的切切關懷：拯救那處於危險中的人；我指的不是僅憑感覺可知的危險，因為這根本不是危險，而是靈魂的危險，就是魔鬼加給人的危險。\par

因為商人為了增加財富，橫渡大海；工匠為了增添產業，做一切事。讓我們也不要以自己的得救為滿足，否則連這個也要失去。因為在戰爭中，在戰鬥中，只想靠逃跑自救的士兵也連同自己毀了其他人；而高尚的人，為保衛他人而持械站立的人，也與他人一起保全了自己。既然我們的狀態也是戰爭，而且是所有戰爭中最慘烈的，是戰鬥和搏鬥，就讓我們按照我們君王的命令在戰鬥中列陣，準備好面對屠殺、流血和死亡，為所有人的救恩著想，鼓勵站立的人，扶起跌倒的人。因為我們許多弟兄在這場衝突中倒下了，受了傷，在血泊中翻滾，卻沒有人醫治，沒有一個百姓，沒有一位司祭，沒有其他人，沒有人站在旁邊，沒有朋友，沒有兄弟，我們每個人只顧自己的事。\par

因此我們也損害了自己的利益。因為最大的信心和嘉許的途徑，就是不看自己的事。\par

所以我說，我們軟弱，容易被人和魔鬼制勝，是因為我們追求相反的事，沒有彼此鎖住盾牌，沒有以虔敬的愛堅固，反而為自己尋找其他友誼的動機：有的出於親屬關係，有的出於長期相識，有的出於共同利益，有的出於鄰里關係；我們與人為友，幾乎出於任何原因，卻不出於虔敬，而友誼本應唯獨建立於此。但現在做的正相反：我們有時與猶太人和希臘人為友，勝過與教會的兒女為友。\par

6. 是的，他說，因為一個是無用的，另一個卻是和善溫柔的。你說什麼？你竟稱你的兄弟為無用，而你是奉命連\OriginalTerm{拉加}{Raca}也不可稱他的？難道你不羞愧，不臉紅嗎？你竟暴露你的兄弟，你的同體，與你同受生育、共享同一餐桌的人？\par

但如果你有按肉體的兄弟，即使他做出萬般惡事，你仍竭力遮掩他，並在他受辱時認為自己也有份於羞恥；而對你的屬靈兄弟，當你本該為他解除誹謗時，你卻反倒羅列萬般指控？\par

\Quote{他為什麼是無用且難以忍受的？}你會說。不，正因如此，要成為他的朋友，好使你終止他這樣的人，使他悔改，領他回到德行。——\Quote{但他不服從，}你會說，\Quote{也不聽勸告。}——你從何知道？什麼，你曾勸戒他，試圖改正他嗎？——\Quote{我常常勸戒他，}你會說。多少次？——常常，既一次，又兩次。——哦！這算很多嗎？為什麼？如果你一直這樣做，豈不該厭倦而放棄？你難道沒看見天主如何時常藉著先知、宗徒、聖史勸戒我們嗎？那麼，我們都做到了嗎？我們在所有事上都服從了嗎？絕沒有。難道祂就停止勸戒了嗎？祂沉默了嗎？祂豈不是每天說：\BibleQuote{你們不能事奉天主，又事奉瑪門}\BibleRef{瑪 6:24}而對許多人來說，財富的過剩和暴政仍在增加？祂豈不是每天高呼：\BibleQuote{寬恕，你們也將獲得寬恕}\BibleRef{路 6:37}而我們卻更像野獸？祂豈不是不斷勸戒要克制慾望、控制邪惡情慾，而許多人在這罪中比豬更打滾？但祂仍然不停地說話。\par

那麼，我們為什麼不思考這些事，對自己說：天主也與我們講理，並不停止這樣做，儘管我們在許多事上不服從祂？\par

因此祂說：\Quote{得救的人少。}因為如果我們的德性本身不足以使我們得救，而且我們還必須帶著別人一起離世；當我們既沒有救自己，也沒有救別人時，我們將受什麼苦？我們還能從哪裡得到救恩的盼望？\par

但我為什麼為這些事責備人呢？因為我們連同住的人——妻子、兒女、僕人——也不關心，反而像醉酒的人一樣，偏重別的事：使僕人人數增多，殷勤服侍我們；使兒女從我們得到大量遺產；使妻子有金飾、貴重衣服和財富；我們不關心他們本身，只關心屬於他們的事物。因為我們既不關心自己的妻子，不為她著想，只關心屬於妻子的事物；也不關心兒女，只關心屬於兒女的事物。\par

我们也是这样，如同有人看见一所房子状况不佳、墙壁倒塌，却忽略去修葺它们，反而在外面筑起高大的篱笆；或者身体患病，却不照顾它，反而为它编织镀金的衣裳；或者女主人病了，却去注意女仆、织机和屋里的器皿，操心别的事，任凭她躺在那里呻吟。\par

时至今日，这也照样发生：当我们的灵魂处于邪恶与悲惨的境地——发怒、辱骂、不正当的贪欲、充满虚荣、争斗、被拖向尘世、被如此多的野兽撕裂时，我们却忽略从她身上驱除这些情欲，反而操心房屋和个人。如果一只熊偷偷逃出，我们会紧闭房屋，沿着狭窄的通道奔跑，以免碰上野兽；然而，此刻撕裂灵魂的不是一只野兽，而是许多这样的思想，我们却毫无感觉。在城市里，我们如此谨慎，将野兽关在偏僻之处和笼子里，既不在城市的元老院，也不在法庭，更不在王宫，而是在远处某个地方将它们锁起来；但在灵魂那里——那里是元老院、是王宫、是法庭——野兽却被释放，在思想本身和御座周围嚎叫喧嚣。因此，一切都颠倒了，一切都充满了混乱：内在的、外在的，我们与一座被蛮族蹂躏而陷入混乱的城市毫无区别；我们身上发生的事，就像一条蛇袭击一窝麻雀，麻雀发出微弱的叫声，惊慌失措地四处飞散，充满恐惧，却无处可去以结束它们的惊惶。\par

7. 因此，我恳求，让我们杀死那条蛇，让我们关住那些野兽，让我们扼杀它们，让我们消灭它们，让我们将这些邪恶的思想交于圣神的利剑，免得先知也以他向犹大所发的同样警告来威胁我们，说：\Quote{野驴要在那里跳舞，豪猪和蛇也在那里。}\par

因为确实有，甚至还有比野驴更坏的人，仿佛生活在旷野中，又踢又跳；的确，我们中间的年轻人多半如此。他们怀着狂野的欲望，就这样跳跃、踢踹，毫无约束地四处游荡，将精力耗费在毫无体统的事上。\par

而父亲们应受责备：他们强迫驯马师以极大的注意力调教马匹，不允许幼驹长期未经驯服，从一开始就给它们套上缰绳和一切；却忽视自己的年轻人，让他们长期不受约束、毫无节制地四处游荡，因淫乱、赌博和常留连邪恶剧场而玷污自己。他们本该在淫乱之前就给他娶一位妻子——一位贞洁且极富智慧的妻子；因为她既能将丈夫从最无序的生活中带出来，又能成为这匹幼驹的缰绳。\par

因为淫乱和通奸不是来自别的原因，正是来自年轻人不受约束。如果他有一位明智的妻子，他就会关心家庭、荣誉和名声。\Quote{但他还年轻，} 你说。我也知道。因为如果依撒格在四十岁娶妻，以童贞度过之前的所有岁月，那么承受恩宠的年轻人更应当实践这种克制。但唉，多么可悲！你们不肯关心他们的贞洁，反而忽略他们玷污、败坏自己，成为可咒骂的；好像你们不知道婚姻的益处是保持身体纯洁，如果做不到，婚姻就毫无益处。但你们却反其道而行：当他们沾满无数污点时，你们才毫无目的、毫无结果地将他们带入婚姻。\par

\Quote{为什么我必须等待，} 你会说，\Quote{使他得以认可，在国事中出人头地；} 但对灵魂，你们却毫不关心，视若弃儿。因此，一切都充满了混乱、无序和麻烦，因为这件事被放在次要地位，必要的事被忽略，不重要的事却得到许多关照。\par

你们不知道吗？你们能为年轻人所做的最大善事，莫过于使他保持纯洁、远离娼妓的污秽。因为没有什么比灵魂更重要。经上说：\BibleQuote{人纵使赚得全世界，} 祂说，\BibleQuote{却赔上自己的灵魂，有什么益处呢？} \BibleRef{玛 16:26} 但由于贪爱金钱倾覆并推倒了一切，推开了对天主的敬畏，如同叛军首领占据了堡垒一般占据了人的灵魂，因此我们既不顾儿女的救恩，也不顾自己的救恩，只着眼于一件事：变得富裕，好将财富留给别人，这些人再留给后人，如此传递下去，我们成了财产和金钱的传递者，而非主人。\par

因此，我们的愚昧深重；因此，自由人反不如奴隶被重视。因为奴隶，我们虽不为他们的缘故，也为我们自己的缘故而责备他们；但自由人连这种关心的益处也享受不到，在我们眼中比这些奴隶更卑贱。我为什么说比我们的奴隶更卑贱呢？因为我们的儿女还不如牲畜被重视；我们关心马和驴多于关心孩子。如果有人有一头骡子，他会非常焦急地寻找最好的饲养员，不苛刻、不欺诈、不酗酒、不熟悉自己的技艺；但若我们为孩子的心灵设立一位导师，我们却立刻随意拉来任何路人。然而，没有比这技艺更伟大的了。因为训练灵魂、塑造年轻人的心灵，有什么能与之相比呢？掌握这技艺的人，应当比任何画家和雕塑家更精确细致。\par

但我们毫不在意这一点，只注重一件事：那就是他自己的舌头得到训练。而为了金钱，我们又把努力用在这方面。因为他学习说话，不是为能言善辩，而是为赚取金钱；因为如果不用说话也能发财，我们甚至不会在乎这一点。\newline{}\par

你看见财富的暴政有多么巨大吗？它如何攫取一切，像捆绑奴隶或牲畜一样捆绑万物，随意拖曳？\newline{}\par

但这样指责财富对我们有什么益处呢？因为我们确实用言语向它射击，但它用行动征服我们。尽管如此，我们也绝不停止用舌头的言语向它射击。因为如果取得任何进展，我们和你们都受益；但如果你们停留在原状，至少我们的部分已经完成了。\newline{}\par

愿天主把你们从这种疾病中解救出来，并使我们因你们而夸耀，因为光荣与权能归于祂，永世无尽。阿们。\par

\SourceRecord{Translated by George Prevost and revised by M.B. Riddle.；Nicene and Post-Nicene Fathers, First Series；Vol. 10.；Edited by Philip Schaff.；Buffalo, NY: Christian Literature Publishing Co.,；1888.；<http://www.newadvent.org/fathers/200159.htm>.}

% structure: p0001 source=h1 target=section reason=page-title

\section{玛窦福音第六十篇讲道}

\Emph{\BibleRef{玛 18:15}}\par

\Quote{如果你的弟兄得罪了你，去，指出他的过失，\Emph{只在你和他之间。如果他听从了你，你就赢得了你的弟兄。}}\par

因为，既然祂严厉谴责那些引人跌倒的人，并从各方面使他们畏惧；为避免那被得罪的人反而因此松懈，既不把一切责任推到他人身上，而是被引向另一种恶，即自身软化，凡事要求被迁就，并陷入骄傲的暗礁，请看祂如何再次约束他们，命令只在两人之间指出过错，以免因多数人的见证而使控告更重，使对方被激怒而坚持不改。\par

因此祂说：\Quote{只在他和你之间，}又说：\Quote{如果他听从了你，你就赢得了你的弟兄。}什么是\Quote{如果他听从了你？}即如果他自责，如果他相信自己做错了。\par

\Quote{你就赢得了你的弟兄。}祂并没有说：你得到了足够的报复，而是说：\Quote{你就赢得了你的弟兄，}以表明仇恨对双方都是损失。因为祂不是说：\Quote{他仅仅赢得了自己，}而是说：\Quote{你也赢得了他的救恩。}由此祂表明，在此之前，两者都是失败者：一个失去了自己的弟兄，另一个失去了自己的救恩。\par

这一点，当祂坐在山上时也曾劝告过；一方面让那施加痛苦的人来到那受痛苦者面前，说：\Quote{与你的弟兄和好}，另一方面命令那受委屈的宽恕他的邻人。因为祂教导人们说：\Quote{求你宽免我们的罪债，如同我们也宽免得罪我们的人。} \BibleRef{玛 6:12}\par

但在此处，祂设计了另一种方式。因为\Emph{祂现在并没有召叫}那施加痛苦的人，而是将那受痛苦的人带到此人面前。因为那做错事的人出于羞愧和窘迫，不容易前来道歉，所以祂将另一人引到他那裏，并且不仅是这样，而且还是以一种纠正所发生之事的方式。祂并没有说：\Quote{控告}，或\Quote{责备他}，或\Quote{要求赔偿和交代}，而是说：\Quote{指出他的过失。}因为他因愤怒和羞愧而陷入一种麻木状态，如同被灌醉；而你，既然健康，就必须向那患病的人走去，使审判私下进行，并使补救措施易于接受。因为说\Quote{指出他的过失}，无非就是\Quote{提醒他的错误}，告诉他你在他手中所受的苦；若事情做得恰当，这本身就是为他辩护的部分，并热切地将他引向和好。\par

那么，如果他不听从，并执意固执下去呢？\Quote{你要另外带一两个人同去，好使一切事在两个人的口证下得以确立。}因为那人越是无耻和胆大，我们就越应当为治愈他而努力，而不是愤怒或愤慨。因为医生同样，当他看到疾病顽固时，并不放弃也不焦躁，反而做更充分的准备；这也是祂在此处命令我们做的。\par

因为既然你独自似乎过于软弱，就借此增加使自己更有力量。因为两个人足以定那犯罪者的罪。你看祂不单寻求那受苦者的益处，也寻求那施加痛苦者的益处。因为受伤害的正是这个被自己的情感俘虏的人，正是他患病、软弱和无力。因此祂常派另一人去，有时单独，有时带他人；但若他仍坚持，甚至带教会去。因为祂说：\Quote{告诉教会。} \BibleRef{玛 18:17}因为如果祂只寻求这一方的益处，祂就不会命令宽恕悔改的人七十个七次。祂也不会多次派这样多人去纠正他的情感；但如果他在第一次交谈后仍不改正，祂就会弃之不顾；而现在，祂命令尝试一次、两次、三次，有时单独，有时带两人，有时带更多人。\par

因此，对于外人，祂没有这样说，而是说：\Quote{如果有人掌击你的右颊，你把另一面也转给他。} \BibleRef{玛 5:39}但此处并非如此。因为保禄所说的\Quote{关於那些外人，那与我有关？} \BibleRef{格前 5:12}，但对弟兄，他命令指出他们的过失、躲避他们、并割断他们的联系，若不听从，好使他们羞愧；此处祂自己也这样做，为弟兄制定了这些法律，并设置两人为他作师傅和法官，在他醉酒时教导他所行的事。因为尽管是他说了、做了那些不合理的事，但他仍需要别人教导他这点，如同醉酒的人。因为愤怒和罪比任何醉酒更疯狂，更使灵魂烦乱。\par

例如，有谁比达味更智慧呢？然而，当他犯罪时，他却没有察觉，他的情欲压制了他的一切理智能力，如同烟雾充满他的灵魂。因此他需要从先知那里得到一盏灯，以及提醒他所作所为的话语。所以此处祂也把这些带到那犯罪者面前，为要与他论理他所做的事。\par

2. 但为什么祂命令此人指出他的过失，而不是另一人呢？因为此人他会更安静地忍受：这被冤屈、被痛苦、被虐待的人。因为被侮辱者本人指出过失，与其他人为侮辱者指出过失，受者的感受不同；尤其是当此人独自责备他时。因为当那理当要求严惩的人，甚至这人还关心他的救恩时，那世上再没有比这更能使他羞愧的了。\par

你看见这并非仅为惩罚，而是为改正吗？因此祂并不立刻命他带两个人同去，而是等他自己失败之后；即使那时，也不派遣大队去攻击他；只是加添到二人，或甚至一人；但当他连这些也轻视时，那时才把他带到教会。\par

祂如此恳切，以免我们邻人的罪被我们暴露。祂本来可以从一开始就命令这事，但为使这事不发生，祂没有命令，而是在第一次和第二次劝诫之后才指定这事。\par

但什么是\BibleQuote{凭两三个见证人的口，每句话都得以成立？}\BibleRef{玛 18:16}你已有足够的见证。祂的意思是，你已经尽了你的一切本分，没有忽略任何属于你应做的事。\par

\BibleQuote{若他拒绝听从他们，就告诉教会，}\BibleRef{玛 18:17}即告诉教会的首领；\BibleQuote{若他拒绝听从教会，就把他当作外邦人和税吏。}\BibleRef{玛 18:17}因为此后这样的人便无药可救。\par

但请你留意，祂如何处处把税吏当作最大邪恶的例证。因为上面祂也说：\BibleQuote{税吏不也是这样行吗？}\BibleRef{玛 5:46}后面又说：\BibleQuote{税吏和娼妓要在你们之先进入天国}\BibleRef{玛 21:31}即那些全然被弃绝和定罪的人。愿那些追求不义之财、计算利上加利的人听吧。\par

但为何祂把他与他们放在一起？为安慰受害的人，并惊吓他。这只是惩罚吗？不，且听下文：\BibleQuote{凡你们在地上捆绑的，在天上也要被捆绑。}\BibleRef{玛 18:18}祂没有对教会的首领说：\Quote{捆绑这人，}而是说\Quote{若你捆绑，}将整件事交给那受委屈的人自己，并且这捆绑是不可解除的。因此他将承受极大的祸患；但那把他带来受审的人没有责任，而是那不肯被说服的人。\par

你看见祂如何用双重约束束缚他吗？既有今世的报应，又有来世的惩罚？但祂说这些威胁，是为使这些情况不发生，而是因惧怕被逐出教会以及捆绑的危险和在天上被捆绑，他变得更温和。既知道这些，即使起初不，至少在许多法庭中他会放下他的愤怒。因此，我告诉你们，祂为他设了第一、第二和第三级法庭，使他即使不听从第一级，也会屈服于第二级；即使拒绝第二级，也会害怕第三级；即使不重视第三级，也会因将来的报应和从天主而来的判决而惊惶。\par

\BibleQuote{我又告诉你们：如果你们中有两人在地上同心合意地祈求任何事，我的天父必为他们成就。因为哪里有两三个人因我的名聚集，我就在他们中间。}\BibleRef{玛 18:19}\par

你看见祂如何以另一种动机消除我们的敌意，除去我们的小争吵，将我们彼此吸引，岂止是从所提到的惩罚，也是从由爱德产生的善果？因为祂在对抗争念中宣布了那些威胁之后，在这里赐下和谐的伟大赏报，至少同心合意的人能向父求得他们所求的事，并有基督在他们中间。\par

\Quote{那么，难道没有两人同心合意吗？}不，在许多地方，甚至可能到处都有。\Quote{那他们为何不获得一切呢？}因为有许多原因使他们失败。或因为他们常求无益的事。何必惊奇，这种事在别人身上如此，甚至在保禄身上也是如此，当他听到：\BibleQuote{我的恩宠为你够用，因为我的德能在软弱中才成全。}\BibleRef{格后 12:9}或因他们不配与听到这些话的人同列，也不贡献自己的部分；但祂寻找像他们一样的人；因此祂说\Quote{你们，}指有德之人，他们展现天使般的生活准则。或他们祈求反对那些得罪他们的人，寻求纠正和报复；这类事是被禁止的，因为祂说：\BibleQuote{为你们的仇人祈祷。}\BibleRef{玛 5:44}或因他们有不悔改的罪而祈求怜悯，这是不可能得到的，不仅他们自己求，即使别人对天主有极大的信心为他们祈求也一样，就像耶肋米亚为犹太人祈祷时所听到的：\BibleQuote{你不要为这百姓祈祷，因为我不听你。}\BibleRef{耶 11:14}\par

但如果一切条件都满足了，你祈求合宜的事，贡献你全部的本分，活出美德，并对你的邻人有和睦与爱，你将获得所求；因为主是爱人的。\par

3. 然后，因为祂说了\Quote{我的父，}为表明是祂自己赐予，而非仅生祂的那一位，祂加添说：\BibleQuote{因为无论何处有两三个人因我的名聚集，我就在他们中间。}\BibleRef{玛 18:20}\par

那么，难道没有两三个人因祂的名聚集吗？确实有，但稀少。因为祂不是仅仅谈论聚会，也不只要求这个；而是如我先前所说，还要求其他德行与此一起；此外，祂对此要求本身也是非常严格的。因为祂所说的是如此：\Quote{若有人把我作为他爱邻人的主要动机，我必与他同在，如果他在其他方面是有德的人。}\par

但现在我们看到大多数人具有其他的友谊动机。因为有人爱，是因为他被爱；有人爱，是因为他受到尊敬；第三人爱，是因为某人在其他世俗之事上对他有帮助；第四人爱，是因为其他类似的原因；但为了基督的缘故，要找到一个人真诚地爱他的邻人，如同他应当爱的那样，是一件困难的事。因为大多数人是被世俗事务彼此束缚的。但保禄并非这样爱，而是为了基督的缘故；因此，即便他没有得到像他所爱的那样被爱，他也没有停止他的爱，因为他为自己的情感种下了坚固的根；但我们的现状并非如此，然而如果我们调查，就会发现对大多数人来说，任何能产生友谊的东西都比这更常见。如果有人在这众多的人群中赐给我权力进行这项调查，我会展示出大多数人被世俗动机彼此束缚。\par

这一点从引发敌意的原因就可以明显看出。因为他们被这些暂时的动机彼此束缚，所以他们彼此既不热忱，也不持久，反而侮辱、金钱损失、嫉妒、虚荣之爱以及诸如此类的事临到他们，就会切断爱的纽带。因为爱找不到属灵的根。因为如果它确实是属灵的，就没有任何世俗之事能瓦解属灵之事。因为为了基督的缘故的爱是坚固的、不可破坏的、坚不可摧的，没有什么能把它撕裂——不是毁谤、不是危险、不是死亡，也不是任何这一类的事。因为那如此爱的人，即使遭受千般苦难，着眼于他爱的根基，他也不会停止。因为那因为被爱而爱的人，若遇到任何痛苦的事，就会结束他的爱；但那因这爱而被束缚的人，绝不会停止。\par

因此保禄也说：\Quote{爱德永不失败。} 你还有什么可说的呢？当他受尊敬时，他侮辱人？当他领受恩惠时，他想要杀害你？但即使这样，也促使你更去爱，如果你是为了基督的缘故而爱。因为在别人身上颠覆爱的事物，在这里却变得易于产生爱。如何？首先，因为这样的人对你是获得赏报的原因；其次，因为如此性情的人更需要救助和许多关注。因此我说，那如此爱的人不询问种族、国家、财富、他对自己是否爱，也不问任何其他这样的事，但尽管他被恨、被侮辱、被杀，他仍继续爱，以基督为他爱的充足理由；因此他也屹立不动，坚固，不可推翻，定睛于祂。\par

因为基督也是这样爱了他的仇敌，爱了那些固执的、有害的、亵渎的、恨祂的、连见都不愿见祂的人；那些宁可选树木石头而不选祂的人，并且是以最高之爱，再没有比这更大的了。祂说：\BibleQuote{再没有比这更大的爱情了，即人若为自己的朋友舍掉性命。}\BibleRef{若 15:13}\par

就连那些钉死祂、并在诸多事例中侮辱祂的人，请看祂是如何继续以仁慈对待。因为祂甚至为他们在圣父面前求情，说：\Quote{宽赦他们吧，因为他们不知道他们做的是什么。} 而且此后，祂还派遣了祂的门徒到他们那里去。\par

那么，让我们也效法这爱，让我们注目于此，好使我们作为基督的追随者，能藉着我们的主耶稣基督的恩宠和对人的爱，获得今世和来世的美善，愿光荣和德能归于祂，直到永远。阿们。\par

\SourceRecord{Translated by George Prevost and revised by M.B. Riddle.；Nicene and Post-Nicene Fathers, First Series；Vol. 10.；Edited by Philip Schaff.；Buffalo, NY: Christian Literature Publishing Co.,；1888.；<http://www.newadvent.org/fathers/200160.htm>.}

% structure: p0001 source=h1 target=section reason=page-title

\section{\WorkTitle{玛窦福音讲道六十一}}

\BibleRef{玛 18:21}\par

\BibleQuote{那时，伯多禄前来对祂说：\InnerQuote{主啊，我弟兄得罪我，我该宽恕他多少次？直到七次吗？}耶稣对他说：\InnerQuote{我不对你说直到七次，而是到七十个七次。}}\par

伯多禄自以为说了什么了不起的话，所以为了显大，他加上了\Quote{直到七次吗？}他说：\Quote{祢所命令要做的事，我该做多少次呢？因为如果他不断犯罪，但每当受到责备就悔改，祢命令我们要容忍这样的人到几时呢？}因为关于那个不悔改、也不承认自己过错的人，祢已经设定了界限，说：\Quote{让他对待你像外邦人和税吏一样}；但对他就不再是这样，祢却命令要接纳他。\par

那么，我该多少次容忍他，在他被告知过失并悔改之后？七次够吗？\par

那么，良善的、爱人如己的天主基督怎么说呢？\Quote{我不对你说直到七次，而是到七十个七次}，这不是设定一个数字，而是无限、永久和永远的。因为就如万次表示很多，这里也一样。因为经上说：\Quote{\InnerQuote{不生育的生了七子}} \BibleRef{撒上 2:5}，圣经是指多。所以祂并没有用数字限制宽恕，而是宣告它应当是永久的和永远的。\par

这一点至少祂通过后面的比喻指明了。因为为了不让任何人觉得所命令的事重大而难以承受，祂说了\Quote{七十个七次}之后，又加上了这个比喻，一方面引导他们遵行祂所说的话，另一方面压服那些以此自夸的人，表明这件事并不沉重，反而非常容易。因此，让我补充说，祂展示了自己对人的爱，以便通过比较，正如祂所说，你们可以学到：即使你宽恕了七十个七次，即使你不断地完全宽恕你的邻人所有的罪，那也只是一滴水与无边的海洋相比，或者说更多，你的爱人之心与天主那无尽的良善相比，仍然远远不及；而你正需要天主的良善，因为你要受审判，要交账。\par

因此祂接着说：\Quote{天国好比一个王，要同他的仆人算账。他开始算账的时候，有人带了一个欠他一万塔冷通的人来。因为他没有可还的，主人就下令把他和他妻子儿女，并一切所有的都卖掉偿还。}\par

后来这人得了恩惠后出去，\Quote{抓住一个欠他一百德纳的同伴，扼住他的喉咙} \BibleRef{玛 18:28}，因这些作为激怒了他的主人，使主人又把他投进监里，直到还清了全部债务。\par

你看见得罪人与得罪天主之间的差别有多大吗？就像一万塔冷通与一百德纳之间的差别，甚至更大。这既出于位格的差别，也出于我们不断犯罪的连续。因为当人看着我们时，我们就远离并躲避犯罪；但当天主每天看见我们时，我们并不停止，而是无所畏惧地做一切事、说一切话。\par

但不仅于此，而且从我们所分的恩惠和尊荣来看，我们的罪变得更严重了。\par

如果你愿意学习我们得罪祂的罪如何是一万塔冷通，甚至更多，我将试着简要说明。但我恐怕，对于那些倾向邪恶、喜欢不断犯罪的人，我会给他们提供更大的安全感，或者使温顺的人陷入绝望，他们会重复门徒们的那句话：\Quote{谁能得救？}\par

尽管如此，我还是要说，为使那些留心听的人更安全、更温顺。因为那些不治之症、麻木不仁的人，即使没有我这些话，也不会摆脱他们自己的疏忽和邪恶；即使他们因此得到更大的轻蔑的借口，过失不在于所说的话，而在于他们的麻木；因为所说的话确实足以约束那些留心的人，并刺痛他们的心；而温顺的人，一方面看到自己罪的重大，另一方面也学到悔改的力量，就会更加紧抓住它，因此有必要说。\par

那么我就要说，并陈明我们的罪，包括我们得罪天主和得罪人的地方，我不提出每个人自己的罪，而是共同的罪；但每个人自己的罪，之后由自己的良心加上。\par

我要这样做，首先陈明天主对我们的恩惠。祂的恩惠是什么呢？祂在虚无中创造了我们，并为我们造了所见的万物：天、海、空气，以及其中的一切，生物、植物、种子；因为我们必须简要地说一下祂作为的无边海洋。唯独我们这些地上的人，祂吹入了一口气，使我们拥有了现在这样的灵魂；祂栽种了乐园；祂赐予了配偶；祂使我们管理一切走兽；祂以光荣和尊贵为冠冕。\par

之后，当人对恩主忘恩负义时，祂又赐予了更大的恩赐。\par

2. 不要只看这一件，就是祂把人赶出了乐园，也要看见从中所得到的益处。因为把祂赶出乐园之后，又行了那无数善事，完成了种种不同的安排，甚至为了那些曾受恩惠却恨祂的人，派遣了祂自己的独生子，为我们打开了天堂，敞开了乐园本身，使我们这些敌人、忘恩负义的人成为儿子。\par

因此现在正适合说：\BibleQuote{啊，天主的智慧和知识的财富多么高深！} \BibleRef{罗 11:33}\par

祂也赐给了我们赦免罪的圣洗，免于惩罚，承受天国，并应许如果我们行正义，就有无数美善，祂伸出援手，将祂的圣神倾注在我们心中。\par

那么，在领受了如此众多且伟大的恩惠之后，我们应当怀有怎样的心境呢？难道即使我们每日为如此爱我们的主而死，就能做出相称的报偿，或者说，我们能否偿还债务的最小一部分？绝非如此，因为甚至这一点也转而成为我们的益处。\par

然而，我们这些本应怀有此种心境的人，如今又是怎样的呢？我们每日都侮辱祂的法律。但请你们不要生气，如果我对犯罪者放言指责，因为我不仅指控你们，也指控我自己。\par

那么，你们愿意我从何处开始呢？从奴隶，还是从自由人？从军中服役者，还是从平民？从统治者，还是从属民？从妇女，还是从男子？从老人，还是从青年？从哪个年龄？从哪个种族？从哪个阶层？从哪个行业？\par

那么，你们愿意我从服役的军人开始吗？这些人每日犯下何种罪过呢？侮辱、谩骂、狂乱、以他人之灾牟利、如同豺狼、从未远离过过犯——除非有人说大海也没有波浪。什么情欲不困扰他们？什么疾病不围攻他们的灵魂？\par

因为他们对同辈怀有嫉妒之心，他们嫉妒且追求虚荣；对属下，则贪得无厌；对那些前来诉讼、如投奔避风港般投奔他们的人，他们的行为却如同仇敌和作伪证者。他们中有多少掠夺！多少欺诈！多少诬告和卑劣行为！多少奴颜婢膝的谄媚！\par

那么，让我们在每种情形下应用基督的法律。\Quote{凡向弟兄说\InnerQuote{愚妄}的，应受地狱之罚。\BibleRef{玛 5:22} 凡注视妇女而有意贪恋她的，已在她心中犯了奸淫。\BibleRef{玛 5:28} 你们若不变成如同小孩一样，绝不能进天国。}\par

然而这些人甚至研习傲慢，对待那些臣服于他们、交付在他们手中、在他们面前战栗恐惧的人，比野兽还要凶猛；他们不为基督做任何事，却为肚腹、为金钱、为虚荣做一切。\par

人能用言语数算他们行为的过犯吗？该如何述说他们的判决、他们的欢笑、他们不合时宜的言谈、他们污秽的言语？至于贪财，简直无从谈起。因为正如山上的隐修士不知贪财为何物，这些亦然；但方式相反：前者因远离此病而不识此情欲，后者因深陷其中而甚至不觉此恶之巨大。因为此恶已如此排挤德行并施行暴政，以致在这些狂人眼中甚至不被视为重罪。\par

但你们愿意我们撇开这些人，转向其他更温和的人吗？那么，让我们来审查工匠和手艺人的群体。这些人似乎最以诚实的劳动和自身的汗水为生。然而，当他们不自警时，也从这里聚集了许多罪恶。因为他们将买卖中的不诚实带入诚实的劳动中，并常常在贪婪之上加上誓言、伪证和谎言，只沉溺于世俗之事，锚定于大地；他们为赚钱而做一切，却不太留意施舍给穷人，总想增加自己的财物。该如何述说涉及此类事务的谩骂、侮辱、借贷、高利贷、充满卑鄙交易的全盘交易、无耻的买卖呢？\par

3. 但你们愿意我们也撇开这些人，转向其他看似更公正的人吗？他们是谁？是拥有土地、收获大地所生财富的人。还有比他们更不义的吗？因为若有人查问他们如何对待那些可怜而劳累的工人，就会看到他们比野蛮人更残忍。对于因饥饿而憔悴、终年劳碌的人，他们既施加沉重且无法忍受的赋税，又加给劳苦的重担，如同对待驴子或骡子，甚或石头一样对待他们的身体，不允许他们稍得喘息；无论土地有收成还是没有收成，他们都同样折磨他们，毫不体恤。还有什么比这更可怜的呢？这些人经过整个冬天的劳作，被霜冻、雨水和彻夜不眠所消耗，却空手而归，甚至还负债，并且比饥荒和船难更畏惧和恐惧监工们的折磨、他们的拖曳、他们的要求、他们的监禁，以及任何哀求都无法解除的劳役！\par

又何必谈论他们从工人身上所做的买卖、所获的肮脏利益呢？他们用工人的劳动和汗水装满酒榨和酒缸，却不允许工人带回家哪怕一小杯，而是将所有果实榨入他们邪恶的酒桶中，并为此扔给他们一点钱。\par

他们还发明了新的高利贷方式，甚至按异教徒的法律也不合法，他们订立充满许多诅咒的借款契约。因为他们索取和逼取的不是款项的百分之一，而是款项的一半；而他们所逼迫的人有妻子、抚养子女、是一个人、并用自身的劳动充满他们的谷场和酒榨。\par

但这些人毫不思虑这些。因此，现在正适宜引出先知的话说：\Quote{惊惶吧，诸天！战栗恐怖吧，大地！}\BibleRef{耶 2:12} 人类的种族竟疯狂到何等巨大的残暴！\par

但我这样说，并非指责手艺、农耕或军役，而是指责我们自己。因为\Person{科尔乃略}也是百夫长，\Person{保禄}是制皮革的，传道后仍操此业，\Person{达味}是君王，\Person{约伯}拥有土地和大笔收入，但这些都未曾阻碍他们中任何人的德行。\par

牢记这一切，并思虑那一万塔冷通，至少让我们从现在起赶紧豁免我们近邻那些微末的债务。因为我们也有一笔账要交，就是我们所受托的诫命，而且我们无论做什么，都不能偿还一切。因此，天主赐给我们一条既现成又容易的偿还之路，能够勾销这一切，那就是不报复。\par

为使我们好好学会这一点，让我们按顺序听完整的比喻。经上说：\Quote{有人带了一个欠他一万塔冷通的来。因他没有可还的，主人就吩咐把他、他的妻子和儿女都卖掉。}请问，这是为什么？不是出于残忍，也不是出于不仁（因为损失还是落在自己身上，她也是奴隶），而是出于说不出的慈爱。\par

他的用意是用这种恐吓来惊吓他，好让他恳求，而不是真的要把他卖掉。因为如果祂真要这样做，就不会应允他的请求，也不会施恩给他。\par

那么，为什么祂不这样做，也不在算账之前就豁免债务呢？祂想教训他，使他从多少重担中被解救出来，这样至少能使他对自己的同伴仆人更温和。因为即使他晓得了自己债务的沉重和赦免的伟大之后，还继续扼住同伴的喉咙；如果祂没有先用这些药事先管教他，他还会残忍到什么地步呢？\par

另一个怎么说呢？\BibleQuote{求你容忍我，我必还清你。他的主人就动了心，释放了他，并免了他的债。}\BibleRef{玛 18:26}\par

你看到那超越的仁慈了吗？仆人只求延迟和宽限时日，而主人却赐下比他求的更多的：全部债务的豁免和赦免。因为主人从起初就愿意赐予，但他不愿意这恩赐只出于自己，也愿意出于这人的祈求，免得他空手而去。因为虽然这一切是出于主人，尽管这人俯伏在地祈求，赦免的动机已表明了，因为\Quote{动了心}就赦免了他。但即使如此，主人仍愿意那人也看似有所贡献，免得他极其羞愧，并使他因自己的灾祸受了教训，而对同伴仆人宽容。\par

4. 至此，这人还是善良可嘉的：因为他承认了，并答应还债，俯伏在地祈求，谴责自己的罪，并认识到债务的沉重。但后来的行为却与先前的行径不相称。因为他一出去，没有过多久，立刻就带着刚得到的恩惠，滥用主人赐予他的自由，把它用在了邪恶上。\par

因为他\BibleQuote{遇见一个同伴仆人，欠他一百德纳，就扼住他的喉咙，说：还你所欠的债。}\BibleRef{玛 18:28}\par

你看到主人的仁慈了吗？你看到仆人的残忍了吗？你们这些为钱财做这事的人，请听。因为如果为了罪我们也不该这样做，何况为了钱财呢？\par

另一个怎么说呢？\Quote{求你容忍我，我必还清你。}但他没有顾念那些曾使他得救的话语（因为他自己说了这话后，就从那一万塔冷通中被释放了），也没有认出那个使他免于沉船的海港；那恳求的姿态并没有唤起他对主人仁慈的记忆，反而因贪婪、残忍和报复心，他把这一切都抛在脑后，比任何野兽都更凶猛，扼住同伴的喉咙。\par

你做什么呢，人啊？你没有觉察到你在向自己索债，把刀刺向自己，废除了那判决和恩赐吗？但他毫不思考，既不记念自己的处境，也不让步，尽管那请求并非为同等事物。\par

因为一个恳求的是免去一万塔冷通，另一个却是一百德纳；一个是向同伴仆人恳求，另一个是向主人恳求；一个得到了完全的赦免，另一个只求延迟，甚至这一点他也没有给他，因为\Quote{他把他下在监里。}\par

\Quote{同伴们看见这事，就向主人报告了。}即是对人，这事也不讨喜悦，何况对天主呢？那些没有欠债的人，也分担了忧伤。\par

主人怎么说呢？\Quote{你这恶仆！因为你求了我，我就免了你所有的债；难道你不该也怜悯你的同伴仆人，如同我怜悯了你一样吗？}\par

再看主人的温柔。他与仆人理论，为自己辩护，即将撤回他的恩赐；或者更确切地说，不是主人撤回，而是那领受的人自己撤回了。因此祂说：\Quote{我免了你所有的债，因为你求了我；难道你不该也怜悯你的同伴仆人吗？}因为即使这事在你看来难以接受，但你应该看到那已经发生和将要发生的收益。即使这命令令人不快，你也该想到赏报；不是他令你伤心，而是你激怒了天主，而你仅凭祈祷就与祂和好了。但即使与伤你的人和好是件难事，落入地狱却更加痛苦；你若把这两者放在一起衡量，就会知道宽恕要轻松得多。\par

而当他欠那一万塔冷通时，主人没有称他为恶，也没有责备他，反而向他施恩；但当他变得对同伴严厉时，主人就说：\Quote{你这恶仆。}\par

让我们这些贪吝的人听吧，因为这话也是对我们说的。让我们这些无情残忍的人也听吧，因为我们残忍不是对别人，而是对自己。所以当你想要报复时，要想到你是在报复自己，而不是别人；你是在捆绑自己的罪，而非邻人的。因为无论你对这人做什么，你都是作为一个人、在此生做；但天主却不这样，祂会更有力地惩罚你，而且是在来世惩罚。\par

\BibleQuote{于是把他交出来，直到他还清所欠的债}——那就是，直到永远；因为他永远不会偿还。既然你因所受的恩待没有变得更好，那么唯有藉惩罚你才能被纠正。\par

然而，\BibleQuote{恩惠和恩赐是不会反悔的}；但邪恶却有如此力量，竟能废掉这律法。那么还有什么比报复更可悲的事呢？因为它竟能推翻天主如此伟大的恩赐。\par

他不仅\Quote{交出}他，而且\Quote{发怒}。因为当他下令将他卖掉时，那并不是愤怒的话语（因此他也没有这样做），而是一个极大的行善机会；但现在这判决却充满了义怒、报复和惩罚。\par

那么这比喻是什么意思？祂说：\BibleQuote{同样，我的圣父也要这样对待你们，如果你们不从心里饶恕你们的弟兄。}\par

祂不说\BibleQuote{你们的圣父}，而说\BibleQuote{我的圣父}。因为对这样邪恶恶毒的人，称天主为父是不合适的。\par

5. 因此祂在这里要求两件事：谴责自己的罪，以及宽恕他人；前者是为了后者，使后者更容易（因为想到自己罪过的人，对自己的同伴会更宽容）；而且不仅要口头上宽恕，还要从心里宽恕。\par

那么，我们不要因报复而把剑刺入自己。因为那个使你忧伤的人，他所加给你的痛苦，怎能比得上你因怀恨在心、招致天主的审判来定你的罪，而给自己带来的痛苦呢？如果你警醒自持，那恶就会落在他头上，受害的是他；但如果你继续愤怒、不满，那么你将受害，不是从他，而是从你自己。\par

所以不要说：他侮辱了我，诽谤了我，对我做了数不清的恶事。因为你越诉说，就越表明他是你的恩人。因为他给了你一个洗除罪恶的机会；所以他对你造成的伤害越大，就越成为你得更大赦罪的原因。\par

因为我们若愿意，没有人能伤害我们，连我们的仇敌也会大大地有益于我们。我说人做什么呢？有什么比魔鬼更邪恶呢？然而，即使从它那里，我们也有一个很好的自我证明的机会；约伯就是证明。但如果魔鬼能成为得冠冕的原因，你又何必害怕一个人作为仇敌呢？\par

那么请看，你若温和地承受敌人的恶行，会得到多少益处。第一，也是最大的，罪过的赦免；第二，坚毅和忍耐；第三，温柔和善心；因为那不懂得对得罪自己的人发怒的人，更会乐意服事爱自己的人。第四，持续地免于愤怒，这是无可比拟的。因为免于愤怒的人，显然也免于由此产生的沮丧，不会将生命耗费在徒劳的劳苦和悲伤中。因为不懂得恨的人，也不懂得忧伤，而会享受喜乐和千万祝福。所以我们因恨别人而惩罚自己，同样，我们因爱别人而造福自己。\par

除此之外，你甚至会成为你的仇敌——即使是魔鬼——所敬畏的对象；或者更确切地说，你若持有这样的心肠，你根本就不会有仇敌。\par

但比这一切更大、更首要的是，你获得天主的恩宠。你若犯了罪，必得赦免；若行了正义，必得更大的信赖。所以让我们做到不恨任何人，好使天主也爱我们，这样即使我们欠了万贯债务，祂也会怜悯我们。\par

你被他伤害了吗？那就可怜他，不要恨他；为他哭泣哀悼，不要避开他。因为不是你在冒犯天主，而是他；而你却因忍受而蒙悦纳。试想基督，当祂即将被钉十字架时，祂为自己而喜乐，却为那些钉死祂的人而哭泣。这也该是我们的心肠；我们受的伤害越大，就越该为那些伤害我们的人哀叹。因为对我们来说，这带来许多益处，对他们却相反。\par

他当着众人侮辱你、打你了吗？那么他就当着众人贬低羞辱了自己，开了无数控告者之口，为你编了更多的冠冕，并聚集了很多人传扬你的忍耐。\par

他曾在别人面前诽谤你吗？那又怎样？要算账的是天主，而不是那些听见的人。因为他是给自己增加了受罚的机会，不仅要为自己的罪交账，还要为他说你的话交账。他在人面前给你带来恶名，但他在天主面前却招致了恶名。\par

如果这些还不够，你要想到你的主也曾被魔鬼和人毁谤，甚至是被祂最亲爱的人毁谤；祂的独生子也同样。因此祂说：\BibleQuote{他们既然称家主为贝耳则步，何况他的家人呢？}\BibleRef{玛 10:25}\par

那恶鬼不仅毁谤祂，而且竟有人相信，并且不是在小事上毁谤祂，而是以最大的侮辱和指控。因为它声称祂被附魔，是个骗子，与天主为敌。\par

但你若行了善，却受了恶？不！为此更当哀恸怜悯那作恶的人，为你自己反倒应当喜乐，因为你已变得相似天主，\BibleQuote{祂使太阳上升，照耀恶人，也照耀善人。}\BibleRef{玛 5:45}\par

然而，若效法天主是力所不及——虽然对于警醒的人，连这也非难事——但尽管如此，如果你仍觉此事过高，那么让我们领你到你同仆那里：到那受尽无数苦难、却向弟兄们行善的若瑟那里；到那在众人无数次谋害他之后仍为他们祈祷的梅瑟那里；到那从他们所受之苦不可胜数、且甘愿为他们成为诅咒的有福的保禄那里；到那被人用石头砸死、却恳求赦免这罪的斯德望那里。你思考了这一切之后，要抛弃一切愤怒，好使天主也因我们的主耶稣基督的恩宠和爱人，赦免我们的一切过犯。愿光荣、权能、尊崇归于祂及圣父和圣神，自今至永远，世世无尽。阿们。\par

\SourceRecord{Translated by George Prevost and revised by M.B. Riddle.；Nicene and Post-Nicene Fathers, First Series；Vol. 10.；Edited by Philip Schaff.；Buffalo, NY: Christian Literature Publishing Co.,；1888.；<http://www.newadvent.org/fathers/200161.htm>.}

% structure: p0001 source=h1 target=section reason=page-title

\section{玛窦福音讲道集 第六十二讲}

\Emph{\BibleRef{玛 19:1}}\par

\Quote{耶稣讲完了这些话，就离开\Place{加里肋亚}，来到\Place{约但河}对岸的\Place{犹太}境内。}\par

祂由于那些人的嫉妒，经常离开犹太；从这时起，祂频繁地来到犹太，因为受难即将临近；然而，祂暂时不上耶路撒冷，而是来到\Quote{犹太境内}。\par

\Quote{并且，}祂来到后，\Quote{有许多群众跟随祂，祂便治好了他们。}\BibleRef{玛 19:2}\par

因为祂并非总是一味地以言语施教，也并非总是一味地行奇事，而是时而这样，时而那样，以不同的方式成就那些等待和追随祂的人的救恩，好使祂所说的话借着奇迹显得是一位值得信赖的师傅，又借着言语的教导增进奇迹的益处；这样，祂就亲手引导他们认识天主。\par

但是，请你留意这一点：宗徒们如何用一句话就略过了整群群众，没有逐一指名道姓地宣告被治愈的人。因为经上并没有说某某人和某某人，只说许多人，教导我们要不炫耀。但基督治愈他们，既造福了他们，也借他们造福了许多其他人。因为治愈这些人的疾病，为其他人认识天主奠定了基础。\par

但法利塞人却不如此，反倒因此更加凶悍，来到祂那里试探祂。因为他们无法阻止所行的奇事，便向祂提出问题。因为他们\Quote{来到祂跟前，试探祂说：\InnerQuote{人无论什么缘故都可以休妻吗？}}\BibleRef{玛 19:3}\par

真是愚蠢！他们以为能用问题使祂哑口无言，尽管他们已从祂身上获得了某种能力的证据。至少，当他们激烈争论关于安息日的事时，当他们说\Quote{祂亵渎了，}当他们说\Quote{祂附有魔鬼，}当他们责备祂的门徒在麦田里行走时，当他们争论关于不洗手的事时，每一次祂都封住了他们的口，堵住了他们无耻的舌头，这样打发走了他们。然而，即使这样，他们也不离开祂。因为邪恶和嫉妒就是如此无耻和胆大；即使被挫败一万次，一万次它也会再次进攻。\par

但是，也请你观察他们从问题形式中显露的诡诈。因为他们并没有对祂说：祢命令不可休妻（因为祂确实已经论述过这条法律），但他们却只字未提那些话；而是从这里找机会，以为能设下更大的圈套，并想把祂逼到必须与法律相抵触的境地，他们不说：祢为什么制定这个或那个？而是好像什么话也没有说过一样，问道：\Quote{可以休妻吗？}他们期望祂已经忘记自己说过的话；如果祂说\Quote{可以休妻，}他们就准备用祂自己说过的话来攻击祂，说：那么祢为什么以前断言相反的话？但若祂现在又说与以前相同的话，就用梅瑟的话来攻击祂。\par

那么祂说了什么？祂没有说：\Quote{你们试探我，伪君子们？}虽然祂后来这样说了，但此时祂没有这样说。这是为什么？为了连同祂的能力也展示祂的温和。因为祂并非总是沉默，免得他们以为自己是隐藏的；也并非总是斥责，为要教导我们以温和承受一切。\par

那么祂如何回答他们？\Quote{你们没有念过：那起初造人的，造男造女，并且说\InnerQuote{因此，人要离开父母，与妻子结合，二人成为一体}吗？这样，他们不再是两个，而是一体了。所以，天主所结合的，人不可拆散。}\BibleRef{玛 19:4}\par

请看师傅的智慧。就是说，当被问及\Quote{是否可以}时，祂没有立刻说\Quote{不可以}，免得他们骚动不安，而是在作出决定之前，先用论证证明了这一点，表明这本身就是祂父的诫命，祂吩咐这些事并非与梅瑟相反，而是完全与祂一致。\par

但是，请注意祂不仅从创造，而且从祂的命令出发，进行了强有力的论证。因为祂没有说只造了一男一女，而且还给了这条命令，要一男与一女结合。如果祂愿意这个人休掉这一个，再娶另一个，那么当祂造了一男之后，就会造许多女人了。\par

但实际上，祂既借着创造的方式，也借着颁布法律的方式，指明了一男必须与一女永久同居，永不分离。\par

请看祂如何说：\Quote{那起初造人的，造男造女，}就是说，他们出自同一根源，成为一体，\Quote{因为二人成为一体。}\par

此后，为使违犯这一法律成为可怕的事，并坚定法律，祂没有说：\Quote{因此不可分开，也不可离异，}而是说：\Quote{天主所结合的，人不可拆散。}\par

但是，如果你提出梅瑟，我就告诉你梅瑟的主；同时，我也依据时间。因为天主起初造男造女；这条法律更古老（尽管看似现在由我引入），并且极其郑重地确立。因为祂不仅将女人带给男人，还命令要离开父母。祂不仅制定法律要男人来到女人面前，而且\Quote{要与她结合，}通过语言的表达暗示他们不可分离。祂还不满足于此，还寻求另一种更大的结合：\Quote{因为二人，}祂说，\Quote{成为一体。}\par

然后，祂诵读了那古老的律法——这律法既由行为又由言语引进——并因那颁赐者的缘故而显明它是值得尊敬的；此后，祂自己以权威解释并颁布律法，说：\Quote{这样，他们不再是两个，而是一体了。}正如割裂肉身是可怕的事，同样，休妻也是不合法的。祂并没有停留于此，还引进了天主，说：\Quote{所以，天主所结合的，人不可拆散。}这表明此事既违背本性，又违背律法；违背本性，因为一体被割裂；违背律法，因为当天主已经结合并命令不可分开时，你却图谋此事。\par

2. 那么，他们此后应当做什么呢？难道不该保持沉默，称赞这话吗？难道不该惊叹祂的智慧吗？难道不该因祂与父的一致而惊异吗？但他们没有做这些事，反而好像在爲\Emph{律法}争辩，说：\Quote{那么，梅瑟爲什么吩咐给休书，便可以休她呢？}其实他们现在不该提出这事，反而祂该向他们提出；但祂并没有利用他们，也没有对他们说：\Quote{我现在不爲此所束缚。}而是也解决了这事。\par

的确，如果祂与旧盟约相悖，祂就不会爲梅瑟辩护，也不会从起初一次成就的事积极地论证；祂就不会努力表明自己的命令与旧命令一致。\par

事实上，梅瑟还颁布了许多其他诫命，包括关于食物和安息日的；那么，为什么他们不像在这里一样，在其他地方也提出他呢？因为他们想煽动众丈夫反对祂。因为这在犹太人中被视爲一件无所谓的事，而几乎所有人都这样做。因此，当山上说了这么多话之后，他们现在只记得这条诫命。\par

然而，那难以言喻的智慧甚至爲此辩护，说：\Quote{梅瑟因爲你们的心硬}才这样立法。祂甚至不让梅瑟受指责，因爲正是祂自己给了他律法；但祂使梅瑟免于责难，并把一切转回到他们头上，正如祂处处所做的那样。\par

因为，当祂的门徒掐麦穗而受责备时，祂就显明他们自己有罪；当他们因不洗手而指责门徒违反诫命时，祂就显明他们自己是违法者；关于安息日也是如此：处处如此，在这里也一样。\par

然后，因为这话难以承受，并给他们带来许多责难，祂迅速将话题转回到那古老的律法，如同祂以前说过的那样，说：\Quote{但在起初并不是这样。}这就是说，天主在起初的行为中规定了相反的事。为了使他们不能说：\Quote{从哪里显明梅瑟是\InnerQuote{因我们的心硬}才说这话？}祂再次封住他们的口。因为如果这是原始的律法，且对我们有益，那么起初就不会颁那另一个；天主在创造时就不会那样造，就不会说那样的话。\par

\Quote{我却对你们说：凡休妻的，除非爲淫乱的缘故，并另娶的，就是犯奸淫。}因为祂已封住他们的口，于是祂凭自己的权威颁布律法，如同关于食物，如同关于安息日一样。\par

因为关于食物同样如此：当祂胜过他们之后，那时——且仅在那时——祂才向群众宣布：\Quote{入口的不能污秽人；}\BibleRef{玛 15:11}关于安息日，当祂封住他们的口之后，祂说：\Quote{所以，在安息日行善是许可的；}\BibleRef{玛 12:12}在这里也正是同一回事。\par

但那里所发生的事，这里也发生了。因为当犹太人被堵住嘴后，门徒们感到困扰，与伯多禄一起来到祂面前说：\Quote{请给我们讲解这个比喻；}\BibleRef{玛 15:15}现在同样，他们也感到困扰，说：\Quote{如果男人这样与妻子相处，倒不如不结婚爲妙。}\par

因为他们现在比以前更明白这话。因此当时他们保持沉默，但现在，经过争辩、回答、提问和通过答复学习，而且律法显得更加清楚，他们便问祂。他们不敢公然反驳，而是提出一个似乎是痛苦且难堪的结果，说：\Quote{如果男人这样与妻子相处，倒不如不结婚爲妙。}因为拥有一位充满各种坏习气的妻子，常年像一头野兽一样被关在屋内，似乎是很难忍受的事。爲使你们知道这事大大困扰了他们，\Person{马尔谷}记载说：\BibleRef{谷 10:10}表明他们私下对祂说话。\par

3. 但\Quote{如果男人这样与妻子相处}是什么意思？也就是说，如果爲这目的他与她结合，即他们应成为一体，或者另一方面，如果人因这些事而招致责难，并因休妻而常常犯罪，那么与邪淫的妻子争战，还不如与自然欲望和自己争战更容易。\par

那么基督说了什么？祂没有说：\Quote{是的，这更容易，就这样做，}免得他们认为这是一条法律；而是补充说：\Quote{不是所有的人都领悟这话，只有那些得了恩赐的人，}\BibleRef{玛 19:11}提高了这事，并表明它是伟大的，从而吸引他们，敦促他们。\par

但请看其中的矛盾。主的确说这是大事，但他们却说这更容易。因为这两件事都应当成就，且应由主承认其为大事，以激励他们更踊跃；同时由他们自己的话显示其更容易，好使他们因此更选择童贞和节制。既然谈论童贞似乎令人痛苦，主便藉此法律的约束驱使他们产生这种渴望。然后，为显示其可能性，祂说：\BibleQuote{有些阉人，自从母胎就是这样生的；也有些阉人，是被人阉的；也有些阉人，是爲天國而自閹的。}\BibleRef{玛 19:12}这些话暗中引导他们选择此事，建立这种德行的可能性，并几乎在说：想想吧，如果你们生来如此，或从那些施加这种恶意伤害的人手中遭受了同样的事，你们会怎样做？被剥夺了享乐却没有赏报。因此现在感谢天主，因为你们带着赏报和冠冕忍受这些事，而那些人忍受时没有冠冕；或者甚至不是这样，而是轻省得多，因有希望和善工的觉悟支持，欲望不像波浪般在你们内翻腾。\par

因为切除肢体不能平息这样的波浪，不能像理性的缰绳那样带来平静；或者不如说，只有理性才能做到这一点。\par

为此祂引入那些其他人，甚至是为鼓励这些人，因为如果祂所建立的不是这个，那么祂论及其他阉人的话又意味着什么？但当祂说他们自阉时，祂指的并非切除肢体，绝非如此，而是摒弃邪恶的思想。因为事实上，自阉的人甚至遭受咒骂，如保禄所说：\Quote{我恨不得那些扰乱你们的人，甚至自阉了才好。}这很有道理。因为这样的人是冒险行凶，为毁谤天主造化的人提供口实，张开摩尼教徒的口，与希腊人中自阉的人犯了同样的不法行为。因为切除肢体从一开始就是魔鬼的作为和撒殚的诡计，为的是对天主的造化作出恶评，毁坏这个活物，使众人将一切归于肢体的本性而非选择，从而在安全中犯罪，不负责任；并且双倍地损害这活物，既通过切除肢体，又阻碍自由选择在善行上的踊跃。\par

这些是魔鬼的规条，除了我们已提及的事，还引入另一邪恶教义，并预先为有关命运和必然性的争论铺路，处处破坏天主赐予我们的自由，说服我们邪恶之事出于本性，从而暗中植入许多其他邪恶教义，尽管不公开。因为魔鬼的毒药就是这样。\par

因此，我恳求你们逃离这种不法之事。因为连同我所提及的，情欲的力量非但没有变得温和，反而更猛烈了。因为我们体内的种子有其来源来自另一个起源，其波浪的翻腾来自另一个原因。有人说这狂暴的冲动来自大脑，有人说来自腰部；但我认为，它无非来自无节制的意志和放任的心智：如果这心智节制，本性冲动就不会产生邪恶的结果。\par

在论及那些徒然无益地成为阉人的阉人（除非他们也在心智上实践节制），以及那些为天国而成为童贞的人之后，祂接着又说：\Quote{能领悟的，就领悟吧！}这同时使他们更加热切，显示这善工极其伟大，并因其无比的温良，不让这事受法律的强制。当祂显示这事极为可能时说了这话，好使自由选择的竞争更大。\par

如果这是自由选择的事，有人会说，祂又为何在开始时说：\Quote{这话不是人人所能领悟的，只有那些得了恩赐的人，才能领悟。}？好叫你们知道这争斗是大的，并非要你们怀疑任何强制的分配。因为恩赐是给那些愿意的人。\par

祂这样说，是为了显示进入这竞技场的人需要许多来自上天的助佑，而愿意的人必定会分得。因为祂惯于在善工伟大时使用这种说法，如祂说：\Quote{天国的奥秘，是给你们知道的。}\par

此事为真，从当前的例子也显而易见。因为如果只是来自上天的恩赐，而度童贞生活的人自己毫无贡献，那么祂许诺给他们天国，并将他们与其他阉人区别开来，就是徒然了。\par

但请看，我祈求，如何从一些人的恶行中，其他人获得益处。我的意思是，犹太人离去时一无所学，因为他们发问并非出于学习的意愿，但门徒们却甚至从这中间获益。\par

4. \BibleQuote{那时，有人带着一些小孩子来见耶稣，要祂给他们覆手祈祷，门徒却斥责他们。耶稣却说：你们让小孩子来吧！不要阻止他们到我跟前来，因为天国正是属于这样的人。耶稣给他们覆了手，就从那里走了。}\BibleRef{玛 19:13}\par

门徒们为什么阻止小孩子？是为尊严。那么祂做什么？祂教导他们谦卑，践踏世俗的骄傲，于是接了他们来，抱在怀里，并应许天国给这样的人；祂以前也说过类似的话。\BibleRef{玛 18:3}\par

那麼，我們若願意繼承天堂，也當以極大的熱忱掌握這德行。因為這是真正智慧的界限：明辨地保持純樸；這是天使的生命，是的，因為小孩子的靈魂純潔無染於一切情慾。他對惹惱他的人不懷怨恨，反而像對待朋友一樣走向他們，好像什麼事都沒有發生過；無論母親如何打他，他仍舊追隨她，將她置於一切之上。即使你給他看戴著王冠的王后，他也不會偏愛她勝過衣衫襤褸的母親，反而寧願看見母親穿著破衣，也不願看見王后華麗。因為他辨別屬於自己和屬於外人的東西，不是根據貧富，而是根據友誼。他所求的也不超過必需之物，只要吃飽母乳就停止吸吮。小孩子不為我們所憂慮的事憂慮，比如金錢的損失和類似之事；也不為我們所高興的事高興，即這些世俗之事；他不貪戀人的容貌。\par

因此祂說：\Quote{天國正是屬於這樣的人。}好叫我們憑著選擇去實踐小孩子們天生所做的事。因為法利塞人之所以行事，沒有比出於詭詐和驕傲更甚的，所以主處處囑咐門徒要心地純樸，既暗指那些人，也教導這些人。因為沒有什麼比權力和高位更能引人驕傲。既然門徒們要在普世享受極大的尊榮，主便預先佔據他們的心，不讓他們有屬人的感覺，也不讓他們向群眾要求榮譽，或讓人們為他們清道。\par

因為這些事看似微小，卻是許多大患的根源。法利塞人至少因這樣受訓，便陷入罪惡的頂峰，追求問候、首席、中位；從這些他們被拋到貪圖虛榮的狂瀾上，然後從那裡又陷入不敬。所以那些人因試探主而自招咒詛，而小孩子們卻得祝福，因他們脫離了這一切。\par

那麼，我們也當像小孩子一樣，\Quote{在惡事上作嬰孩。}\BibleRef{格前 14:20}因為一個人若不如此，便不可能看見天國；狡猾邪惡的人必定要投入地獄。\par

5. 而且在進入地獄之前，我們在此世也要遭受極大的苦難。經上說：\Quote{你若作惡，}經上說，\Quote{你獨自承受那惡；但若行善，則為你自己和你的鄰人。}請看，這事在從前的例子中也照樣發生了。因為沒有誰比撒烏耳更邪惡，也沒有人比達味更純樸和心地純潔。那麼，哪一個更強大呢？達味豈不是兩次將他拿在手裡，雖有能力殺他，卻仍忍耐？豈不是像網中囚犯一樣將他困住，卻饒恕了他？這是在別人催促他、他自己也能控訴他無數罪狀的時候；但他仍讓他平安離去。而另一人卻率領全軍追趕他，他卻帶著幾個絕望的流亡者四處漂泊遷徙；然而這流亡者勝過了國王，因為前者以純樸進入衝突，後者以邪惡進入。\par

因為有誰比那個人更邪惡呢？他率領軍隊，勝利地完成所有戰爭，承受勝利的勞苦和戰利品，卻將冠冕獻給那人，反倒圖謀殺害他？\par

6. 嫉妒的本性如此，它總是圖謀反對自己的榮譽，損耗擁有它的人，給他帶來無數災禍。例如，那可憐之人直到達味離開時，才發出那哀慟的哭聲，自怨自艾地說：\Quote{我很痛苦，培肋舍特人攻打我，上主已離棄了我。}\BibleRef{撒上 28:15}他與達味分開之前，沒有在戰爭中失敗，反而平安而榮耀；因為將軍的榮耀歸於國王。這人並無篡位之心，也沒有圖謀推翻對方，反而為他成就一切，並且忠誠於他；這從後來的事也明顯可知。當他確實服從他時，那些不仔細查考的人或許以為這是臣屬的常規；但當他離開撒烏耳的王國之後，還有什麼能阻止他，甚至殺死他呢？難道對方不是一次、兩次、多次惡待他嗎？豈不是在他受惠於對方之後？豈不是無可指責？撒烏耳的王國和安全對他自己豈不是危險和不確定嗎？他豈不是必須漂泊流浪，顫抖恐懼於極大的禍患，而對方卻活著並統治？然而，這一切沒有迫使他染指於血；當他看見對方睡著、被捆綁、孤身一人、在自己人中間，觸摸他的頭，並有許多人鼓動他，說這有利的時機是天主的審判時，他立刻斥責那些催促他的人，制止了殺戮，送他平安無恙地離去；他好像是他的保鏢和持盾者，而不是敵人，他指責軍隊對國王的背叛。\BibleRef{撒上 26:16}\par

有什麼能與這靈魂相比？有什麼能與那溫和相比？因為這甚至從已提及的事也可見，但更從現今所行的事可見。因為當我們思想自己的卑微時，我們會更完全地認識那些聖人的德行。所以我懇求你們努力效法他們。\par

因为如果你爱光荣，并为此图谋反对你的邻人，那么当你蔑视光荣并停止图谋时，你将更充分地享有它。正如致富与贪婪相反，同样，爱光荣也与获得光荣相反。如果你愿意，让我们探究每一件事。既然我们不惧怕地狱，也不十分关心天国，那么，就从眼前的事来引导你们吧。\par

请告诉我，谁是可笑的？岂不是那些为了从群众中得光荣而行事的人吗？谁又是受赞扬的？岂不是那些蔑视群众赞扬的人吗？因此，如果贪求虚光荣是应受责备的事，并且虚光荣的人爱它也无法隐藏，那么他必定成为受责备的对象，而爱光荣就成了他蒙受耻辱的原因。不仅在这方面他使自己蒙羞，而且他还被迫做出许多可耻之事，充满极大的耻辱。正如在利益方面，人往往因贪婪之病而遭受比任何事更大的损害（他们至少成为许多诡计的受害者，因小利而蒙受大损失，因此这句话甚至成为谚语）；同样，对于好色之人，其情欲也成为他享受快乐的障碍。那些极其沉溺于其中并成为女人奴隶的人，女人尤其会像对待仆人一样将他们牵来牵去，永远不屑于将他们当作男人对待，而是殴打、踢打、带领、到处拉着他们，摆架子，仅仅向他们发号施令。\par

同样，没有什么比那傲慢、疯狂追求光荣、自视甚高的人更卑鄙、更受羞辱了。因为人类是好争斗的种族，没有什么比反对一个自夸者、一个傲慢的人、一个光荣的奴隶更强烈的了。\par

而他自己，为了维持其骄傲的样式，却向普通人表现出奴隶的行为，谄媚、讨好他们，受着比用金钱买来的奴隶更痛苦的奴役。\par

既然知道这一切，让我们放下这些情欲，免得我们既在此世受罚，又在那里受无尽刑罚。让我们成为热爱美德的人。这样，即使在我们到达天国之前，我们也能在此世收获最大的益处，而当我们离开时，我们将分享永福；愿天主藉着我们的主耶稣基督的恩宠和爱人，恩赐我们众人获享这些，愿光荣和权能归于他，至于无穷之世。阿们。\par

\SourceRecord{Translated by George Prevost and revised by M.B. Riddle.；Nicene and Post-Nicene Fathers, First Series；Vol. 10.；Edited by Philip Schaff.；Buffalo, NY: Christian Literature Publishing Co.,；1888.；<http://www.newadvent.org/fathers/200162.htm>.}

% structure: p0001 source=h1 target=section reason=page-title

\section{《玛窦福音》讲道集第六十三篇}

\BibleRef{玛 19:16}\par

\BibleQuote{有一个人前来对祂说：善师，我该做什么，才能得永生？}\par

有些人确实指责这个青年，说他是虚伪和心怀恶意，带着试探来到耶稣面前；但我虽然不愿说他不贪爱钱财、不受财富束缚（因为基督确实指责他具有这样的品性），却绝不称他为虚伪之人，一则因为冒险论断不确定之事并不安全，尤其是在责备方面，二则因为马尔谷已经消除了这一嫌疑；因为他记述道：\BibleQuote{他跑来，跪在祂面前，求祂，}并说\BibleQuote{耶稣注视他，就喜爱他。}\BibleRef{谷 10:17}\par

然而财富的暴虐是何等巨大，由此可以显明；我的意思是，即使我们在其他方面品德高尚，这也会毁掉一切。保禄有理由断言它是一切邪恶的根源。他说：\BibleQuote{贪爱钱财是一切邪恶的根源。}\par

那么，基督为何如此回答他，说\BibleQuote{没有良善的？}因为他是作为一个普通人、一个凡夫、一个犹太教师来到祂面前；因此，基督也作为一个人与他交谈。事实上，在许多情况下，祂回应来到祂面前之人的隐秘思想；例如，当祂说\BibleQuote{我们朝拜我们所知道的，}\BibleRef{若 4:22}以及\BibleQuote{我若为自己作证，我的证验不真。}\BibleRef{若 5:31}因此，当祂说\BibleQuote{没有良善的}时，并非将自己排除在良善之外，绝非如此；因为祂并没有说\BibleQuote{你为什么称我为良善的？我不是良善的，}而是说\BibleQuote{没有良善的}，也就是说，在众人中没有。\par

祂说这话时，并非剥夺人的良善，而是相对于天主的良善而言。因此祂也补充道：\BibleQuote{惟独一位，就是天主；}祂并没有说\BibleQuote{惟独我的父，}好让你知道祂尚未向这个青年启示自己。同样，早先祂称人为恶，说：\BibleQuote{你们纵然不善，尚且知道把好东西给你们的儿女。}\BibleRef{玛 7:11}因为在那里，祂称他们为恶，并非谴责整个人类是恶的（因为祂所说的\Quote{你们}并非指\Quote{你们人类}），而是将人间的善与天主的善相比，如此称呼它；因此祂又补充说：\BibleQuote{何况你们的父，岂不更把好东西赐给求祂的人吗？}那么，有什么促使祂以这种方式回答，或有什么益处呢？祂逐步引导他，教导他远离一切奉承，将他从依赖万物中拉开，使他专注天主，劝勉他寻求永恒的事物，并认识那真正的善、万物的根源和泉源，将光荣归于天主。\par

同样，当祂说\BibleQuote{不要称呼地上的人为师傅}时，这是相对于祂自己而言的，好让他们知道万有的最高主权是谁。因为这个青年在此前已表现出不小的热忱，竟怀有如此渴望；当其他人有的试探，有的为医治自己或他人的疾病而来时，他却为永生而来见祂，并与祂交谈。因为土地肥沃而富饶，但荆棘丛生，窒息了种子。无论如何，请注意他是如何准备好遵守诫命的。因为他说：\BibleQuote{我该做什么，才能得永生？}他如此乐于遵行被告知的事。如果他来试探祂，福音作者也会向我们说明这一点，正如对其他人那样，例如对待法学士。即使他自己保持沉默，基督也不会容忍他隐藏，而是会明确地揭露他，或者至少暗示出来，以免他看似欺骗了祂并被隐藏，从而受到伤害。\par

如果他是来试探祂的，就不会因听到的话而忧愁地离开。法利塞人从未有过这种感受，反而在口被封住时变得恼怒。但这个人并非如此；他垂头丧气地离去，这并非微不足道的迹象，表明他来见祂并非心怀恶意，而是心志过于软弱，并且他确实渴望永生，却被另一种最强烈的感情所束缚。\par

因此，当基督说\BibleQuote{你若愿意进入生命，就该遵守诫命}时，他问道\BibleQuote{哪一条？}并非试探，绝无此意，而是以为除了律法之外还有其他诫命能为他赢得永生，这正是一个极其渴望之人的表现。然后，当耶稣提到律法中的诫命时，他说\BibleQuote{这一切我自幼就遵守了。}他并未止于此，而是又问\BibleQuote{我还缺少什么？}这本身再次表明他极其恳切。\par

那么基督说什么呢？由于祂将要吩咐一件大事，便先摆出赏报，说：\BibleQuote{你若愿意成为成全的，去变卖你所有的，施舍给穷人，你必有财宝在天上；然后来跟随我。}\par

2. 你看祂为这场赛跑预备了多少奖品、多少冠冕？如果他是来试探的，祂就不会告诉他这些事。但现在祂既说了，并且为了吸引他，还向他显示赏报何其大，且将一切完全交托给他自己的意愿，以此冲淡祂建议中看似难以接受的部分。因此，在提到争战和劳苦之前，祂先向他展示奖赏，说\BibleQuote{你若愿意成为成全的，}然后说\BibleQuote{变卖你所有的，施舍给穷人，}紧接着又是赏报：\BibleQuote{你必有财宝在天上；然后来跟随我。}因为跟随祂本身就是极大的赏报。\BibleQuote{你必有财宝在天上。}\par

因为祂的谈论涉及钱财，祂甚至劝他抛弃一切，表明他并非失去所有，而是增添财富，祂所赐的比他所舍弃的更多；不仅更多，而且大如天堂比地更大，甚至更大。\par

但祂称之为宝藏，表明报偿的丰盛、永久和稳妥，尽可能用人间的比喻向听者暗示。因此，不仅轻视财富是不够的，还必须周济穷人，并首先跟随基督，即遵行祂所命令的一切，随时准备殉道和每日的死亡。\BibleQuote{因为谁若愿意跟我来，就当否认自己，背起自己的十字架来跟随我。}\BibleRef{玛 16:24}因此，抛弃钱财比这最后一条诫命——甚至流血——要小得多；然而摆脱财富对此并非无益。\par

\BibleQuote{那少年人一听这话，就忧愁地走了。}\BibleRef{玛 19:22}之后，福音作者似乎要表明他不可能感受到的，说：\Quote{因为他拥有许多产业。}因为拥有少许的人不会同样被奴役，如同被巨大财富淹没的人那样，因为那时对财富的贪爱变得更加专制。我不断说的正是这一点：财富的增加更点燃火焰，使占有者更贫穷，因为它激起更大的欲望，并使他们更感到匮乏。\par

请看，就在这里，这种情欲显示了何等力量。那位曾带着喜乐和热忱来到祂面前的人，当基督命令他抛弃财富时，竟被如此压倒和拖累，甚至无法回答这些事，反而沉默、沮丧、生气地走开了。\par

那么基督说了什么？\Quote{富人进天国是多么难啊！}祂并非谴责财富，而是谴责那些被财富奴役的人。如果富人\Quote{难}，那么贪婪的人更难。因为若不捐出自己的财物成为进入天国的障碍，那么夺取别人的财物，想想会积聚多少火啊。\par

然而，为什么祂对门徒说\Quote{富人难以进入}，而他们是穷人，没有财产呢？祂教导他们不要因贫穷而羞愧，并好像为自己辩解，允许他们一无所有。\par

但祂说那是难的；接着祂表明甚至是不可能的，不仅是不可，而且是极其不可能；祂用骆驼和针眼的比喻表明这一点。\par

祂说：\BibleQuote{骆驼穿过针眼，比富人进天国还容易。}\BibleRef{玛 19:24}由此表明，那些富有并能自制的人，并非没有巨大的赏报。因此祂也断言这是天主的作为，为显示成就此事需要极大的恩宠。果然，当门徒们不安时，祂说：\BibleQuote{在人这是不能的，但在天主凡事都能。}\BibleRef{玛 19:26}\par

门徒们为何不安？他们本是贫穷，甚至极其贫穷。那么他们为何困惑？因为他们为其余人的救恩而担忧，对众人怀有极大的爱，并已具备教师般的慈悲心肠。他们因这宣告而对整个世界如此战栗恐惧，以至需要极大的安慰。\par

因此，祂首先\BibleQuote{注视他们，对他们说：在人不能的，在天主却能。}祂以温和仁慈的目光安抚他们颤抖的心灵，消除他们的痛苦（福音作者用\Quote{祂注视他们}表明这一点），然后祂也以言语安慰他们，向他们呈示天主的全能，从而使他们充满信心。\par

但若你也要明白其方式，以及不可能的事如何变为可能，请听。祂说\BibleQuote{在人不能的，在天主却能}，并非要你放弃并止步于不可能的事，而是要你思考这善工的伟大，从而热切地奔向它，并恳求天主在这些高贵竞赛中助佑你，从而获得永生。\par

3. 那么这如何成为可能？若你抛弃所有，若你腾空财富，若你戒除邪恶的欲望。为证明祂并非将此事完全归于天主，而是为让你知道这善工的伟大，请听下文。当伯多禄说：\BibleQuote{看，我们舍弃了一切，跟随了祢，}并问\BibleQuote{我们将得到什么呢？}时，祂为他们规定了赏报，并补充说：\BibleQuote{凡舍弃了房屋、田地、兄弟、姐妹、父亲或母亲的，必要领取百倍的赏报，并承受永生。}这样，不可能的变成可能。但有人会说：如何能做到舍弃这些呢？一个深陷财富欲望的人怎能自拔？若他开始腾空自己的财产，并割舍多余的。如此他便会更进一步，以后更容易前行。\par

因此，不要寻求一步登天，而要温和地、逐步地攀登这通向天堂的阶梯。因为如同患热病的人体内充满酸胆汁，若再进食饮，非但不能解渴，反而点燃火焰；同样，贪婪的人若将财富投入比胆汁更酸的邪恶欲望中，反而会加剧它。没有什么比暂时节制贪图之欲更能遏制它，如同酸胆汁通过禁食和排泄而得到抑制。\par

但这件事本身，要藉什么方法成就呢？有人会说。如果你深思，当你富有的时候，你将永不止息地渴求，并因贪求更多而憔悴；但若你摆脱了财产，你也能遏制这疾病。那么，不要用更多来包围自己，免得你追逐不可及的事物，变得不可救药，比众人更悲惨，如此疯狂。\par

因为请回答我，我们该说谁更受折磨和痛苦呢？是那贪求昂贵酒肉，却不能随心所欲享用的人，还是没有这种欲望的人？显然必须说，是那渴望却得不到所欲之物的人。因为渴望而不能享受，口渴而不能喝，是如此痛苦，以致基督要想形容地狱时，便如此描述，并引入那富人以这种方式受折磨。因为渴望一滴水而得不到，这乃是他的惩罚。所以，轻视财富的人平息了欲望，但渴望富有的人却把它煽得更旺，而且仍不停止；即使他得到一万塔冷通，他还渴望同样多；即使他得到了这些，他又瞄准双倍更多，继续下去，他甚至渴望山脉、大地、海洋，一切为他变成金子，被一种新的、可怕的疯狂所迷，而且这狂热永不能被这样扑灭。\par

为使你明白，这恶不是藉增添，而是藉除去才得止息；如果你曾经有过想飞行、在空中飘浮的荒唐欲望，你如何扑灭这无理欲望呢？是通过制造翅膀，准备其他器械，还是通过说服理性，明白它是在欲望不可能的事，且不应尝试其中任何一样？很明显，是通过说服理性。但你可能说，那是不可能做到的。然而，要为这欲望找到一个界限，更是难上加难。因为确实，人飞行比藉增添更多东西来使这贪欲止息更容易。因为当欲望的对象是可能的时候，人可以通过享受它们得到安抚；但当它们不可能时，人必须努力做一件事，就是使自己脱离欲望，否则至少不可能恢复灵魂。\par

因此，为使我们不致有冗余的忧愁，让我们抛弃那永远使人痛苦、永不安宁的贪爱钱财，让我们转向另一种爱，它既使我们幸福，又极为便利，让我们渴望天上的宝藏。因为这里的劳作并不那么大，收益亦不可言喻，凡是稍微警醒、清醒并轻视现世事物的人，都不可能得不到它们；正如另一方面，对于被现世事物奴役、完全沉迷其中的人，他必定不能获得那些更好的财富。\par

4. 因此，思量这一切，要除去对财富的邪恶意欲。因为你不能说这话：它虽剥夺我们未来的事物，却赐给我们现世的事物；即使如此，这也是极重的惩罚和报应。但现在连这也说不上。因为除了地狱，且在地狱之前，它在此世就已将你投入更惨烈的惩罚。因为许多家庭被这贪欲倾覆，激起了惨烈的战争，迫使人们以暴死结束生命；而在这些危险之前，它已毁坏了灵魂的高贵，常常使拥有它的人变得怯懦、无丈夫气、鲁莽、虚伪、诽谤、贪婪、欺压，以及所有最坏的事情。\par

但也许你看见银子的光芒、仆人的众多、房屋的美丽、市场中的奉承，就被这迷住了吗？那么，对这恶伤有何补救呢？如果你深思这些东西如何影响你的灵魂，它们使它变得多么黑暗、荒凉、污秽，多么丑陋；如果你计算这些东西是以多少邪恶得来的，以多少劳苦保存的，以多少危险——其实它们并不能保存到底，但当你逃过了所有人的攻击，死亡临到你时，常将这些东西移入你仇敌之手，然后它把你带走，使你赤贫，身后不留任何这些东西，唯独留下灵魂在离世前因这些东西所受的创伤和疮疤。当你看见任何人外表因衣服和众多随从而光彩夺目时，揭开他的良心，你会看见里面有许多蜘蛛网和许多灰尘。试想\Person{保禄}、\Person{伯多禄}。试想\Person{若望}、\Person{厄里亚}，或者更确切地说，\Person{天主子}自己，祂没有枕头的地方。要效法祂和祂的仆人，并为你自己想象他们不可言喻的财富。\par

但如果你藉这些获得了一点视力，却又变得昏暗，如同在船难中风暴来临之际，请听\Person{基督}的声明，祂断言：\BibleQuote{富人进入天国是不可能的。}你尽可以拿山岳、大地、海洋来对抗这声明；如果你愿意，就设想一切都是金子；但你将看到，没有什么能与你因此遭受的损失相比。你固然提及若干亩田地、十座二十座甚至更多的房屋、同样多的浴室、一千或两千个奴隶、以及镶银包金的战车；但我要说：如果你们中每个富人离开这贫穷（因为与我要说的相比，这些事就是贫穷），并拥有整个世界，每人都有现在遍布陆地和海洋的那么多人，每人都有一个海陆世界，到处是建筑、城市、民族，从四面八方，代替水、代替泉源，有金子为他涌流出来，我也不说这些如此富有的人，当他们被逐出天国时，价值三个铜板。\par

因为若现在他们瞄准那会朽坏的财富，而得不到时便痛苦不堪，若他们得以感知那说不尽的福乐，那么什么才能安慰他们呢？什么也不能。请不要对我讲他们财产的丰裕，而要想想贪爱这丰裕的人因此遭受多大的损失：因这些事物而失去天堂，且处于如同某人被逐出王宫的最尊荣之后，拥有一个粪堆，竟以此为荣。囤积金钱与此毫无区别，或可说那粪堆更好。因为粪堆可用于耕作、烧热浴池以及其他类似用途，而埋藏的金子却一无所用。但愿它仅仅是无用而已；但事实上，它还为拥有它的人点燃许多炉火，除非他正当地使用它；由此生出数不尽的祸患。\par

因此，那些外人常称贪财为万恶的堡垒；但蒙福的保禄说得更好、更生动，称它为\BibleQuote{万恶的根源}。\par

那么，既思量这一切，让我们效法那值得效法的事，不是华丽的建筑，不是昂贵的产业，而是对天主怀有极大信赖的人、那些拥有天上财富的人、那些宝藏的主人、真正富有的人、为基督的缘故而贫穷的人，好使我们能藉着我们的主耶稣基督对人的恩宠与慈爱，获致永恒的福乐；愿光荣、权能、尊崇归于祂，偕同圣父及圣神，自今而始，永远世世。阿们。\par

\SourceRecord{Translated by George Prevost and revised by M.B. Riddle.；Nicene and Post-Nicene Fathers, First Series；Vol. 10.；Edited by Philip Schaff.；Buffalo, NY: Christian Literature Publishing Co.,；1888.；<http://www.newadvent.org/fathers/200163.htm>.}

% structure: p0001 source=h1 target=section reason=page-title

\section{玛窦福音讲道集 第六十四篇}

\Emph{\BibleRef{玛 19:27}}\par

\BibleQuote{那时伯多禄回答祂说：看，我们舍弃了一切，跟随了祢；那么，我们将得到什么呢？}\par

这一切？哦，蒙福的伯多禄；那鱼竿？那鱼网？那船？那手艺？你就把这些告诉我，当作一切吗？是的，他说，但我这样说不是为了炫耀，而是为了通过这个问题把穷人的群众带进来。因为主曾说：\BibleQuote{你若愿意成为成全的，去变卖你所有的，施舍给穷人，你必在天上有财宝；} \BibleRef{玛 19:21} 免得穷人中有人说：那么，如果我没有财产，我就不能成为成全的吗？伯多禄问，为要叫你，穷人，学习到，你在这方面丝毫不被贬低：伯多禄问，为要叫你，不从伯多禄学习而怀疑（因为他那时还不成全，且没有圣神），而是从伯多禄的主那里领受声明，好叫你放心。\par

因为我们所做的一样（当谈论他人的事时，我们常把别人的事当作自己的事），宗徒也是这样，当他代表全世界向祂提出这个问题时。因为他至少肯定知道自己的一份，从先前所说的已显明出来；那位已经领受了天国钥匙的，对将来的事更可充满信心。\par

但也请注意，他的回答如何完全符合基督的要求。因为祂曾向那富贵人要求这两件事：把他所有的施舍给穷人，并跟随祂。因此，他也表达了这两件事：舍弃和跟随。他说：\Quote{看，我们舍弃了一切，}他说，\Quote{并跟随了祢。}因为舍弃是为了跟随，而跟随因舍弃变得更容易，使他们对于舍弃感到自信和喜乐。\par

那么祂说什么？祂说：\BibleQuote{我实在告诉你们：你们这些跟随我的人，在重生的时候，当人子坐在祂光荣的宝座上时，你们也要坐在十二个宝座上，审判以色列十二支派。} \BibleRef{玛 19:28} 那么，有人会说，犹大也坐在那里吗？绝不。那么，祂怎么说：\Quote{你们要坐在十二个宝座上？}那么这应许的条款如何实现？\par

请听这是如何以及基于什么原则。有一条天主颁布的法律，由先知耶肋米亚向犹太人诵读，话是这样说的：\BibleQuote{我何时论到一邦或一国，说：要拔除、摧毁、消灭；如果那一邦人从他们的邪恶中悔改，我也要后悔我本想给他们降的灾祸。我何时论到一邦或一国，说：要建立、要栽植；如果他们在我眼中行恶，不听从我的声音，我也要后悔我所说要赐给他们的福乐。} \BibleRef{耶 18:7}\par

对福乐之事我也遵守同样的习惯，祂说。因为虽然我说了要建立，如果他们显出自己配不上这应许，我就不再做了。就像在创造人时所做的那样，经上说：\BibleQuote{你们所惊恐和害怕的，要临到走兽，} \BibleRef{创 9:2} 但这事没有发生，因为人证明了自己不配统治，犹大也是如此。\par

因为为了避免人在惩罚的宣布中绝望而变得更刚硬，也为了避免因福乐的应许而毫无理由地变得更懈怠，祂用我先前提到的方法补救这两种恶，这样说：即使我威胁，也不要绝望；因为你能悔改，并翻转这宣布，像尼尼微人一样。即使我应许任何福乐，也不要因应许而懈怠。因为如果你显出自己配不上，我曾应许的事实不会使你受益，反而会带来惩罚。因为我应许时，你是配得上的。\par

因此，即使那时在祂与门徒的谈话中，祂不是简单地应许他们，因为祂并没有只说\Quote{你们，}而是加上了\Quote{那些跟随我的人，}为要既排除犹大，又吸引后来的人。因为这话不是只对他们说的，也不包括已经变为不配的犹大。\par

如今，祂向门徒应许将来的事，说：\BibleQuote{你们要坐在十二个宝座上，}因为他们现在已有了更高的品格，不再追求世上任何事物；但对其余的人，祂也应许今世的事。\par

因为祂说：\BibleQuote{凡为我的名舍弃了兄弟、姊妹、父亲、母亲、妻子、儿女、田地或房屋的，必要在今世得到百倍，并承受永生。}\par

因为免得有人听了\Quote{你们}这个词，就以为这是门徒独有的（我指的是在将来享受最大最先的尊荣），祂扩展了这话语，将应许散布到全地，并从今世之事建立将来之事。而且对门徒们，在起初他们还处于更不完全的状态时，祂也从今世之事推理。因为当祂从海边召叫他们，带他们离开职业，命令他们舍弃船只时，祂提及的不是天堂，不是宝座，而是此世的事，说：\BibleQuote{我要使你们成为渔人的渔夫；}但当祂已经陶成他们，使他们有更高的眼界时，此后祂就也讲论将来之事了。\par

2. 但什么是\Quote{审判以色列十二支派？}这就是\Quote{定他们的罪。}因为他们不是真的要坐审判席，而是如同祂说南方的女王要定那世代的罪，尼尼微人要定他们的罪一样，如今这些也要这样。因此祂没有说\Quote{外邦人}和\Quote{世界}，而是说\Quote{以色列的支派}。因为犹太人和宗徒们都在同样的法律、习俗和政制下长大；当犹太人说，他们因此不能信基督，因为法律禁止接受祂的诫命时，祂举出这些人，他们领受了同样的法律，却仍然相信了，祂就定所有这些人的罪；如同祂已经说过：\Quote{因此他们要作你们的审判官。}\BibleRef{玛 12:27}\par

可能有人说，如果尼尼微人和南方女王所得的，这些也要得，那祂给他们许诺了什么大事？首先，祂在此以前已经许诺了他们许多其他事，在此以后也许诺他们，并且这本身并非他们的赏报。\par

而且除此以外，祂在途中还暗示了比这些更多的事。因为对那些，祂只是说\Quote{尼尼微人要起来定这世代的罪。}\BibleRef{玛 12:41}，和\Quote{南方的女王要定它的罪。}；但关于这些，不只是这样，而是如何？\Quote{当人子坐在祂荣耀的宝座上时，你们也要坐在十二个宝座上，}祂说，表明他们也要与祂一同为王，分享那荣耀。因为经上说：\Quote{如果我们受苦，我们也要与祂一同为王。}\BibleRef{弟后 2:12} 因为宝座并非指审判的座位，因为唯独祂是那要坐并审判的，但宝座表明无法言说的尊荣和荣耀。\par

祂对这些说了这些事，但对其他所有人说了永生和今世的百倍赏报。但如果是对其他人，那么对这些更是如此，既得这些，也得今世的事物。\par

这确实成就了；因为他们一旦放下钓竿和渔网，就掌有所有人的财产、房屋和田地的价钱，以及信徒们的身体。因为他们常常甚至愿意为他们被杀，正如保禄也为许多人作证，他说：\Quote{如果可能，你们会挖出眼睛给我。}\BibleRef{迦 4:51} 但当祂说：\Quote{凡舍弃了妻子的，}祂不是说婚姻要无故破裂，而是如同祂论到人的生命时说：\Quote{凡为我丧失生命的，必得着生命。}\BibleRef{玛 10:39} 不是要我们毁灭自己，也不是要我们在此世使灵魂离开身体，而是要我们优先考虑虔诚胜过一切；这也是祂论到妻子和弟兄时所说的。\par

但我觉得祂在此也暗示了迫害。因为有很多例子：父亲催逼儿子行不敬虔之事，妻子催逼丈夫；当她们命令这些事时，祂说，既不要当作妻子，也不要是父母，正如保禄也说：\Quote{如果不信的人离开，就让他离开吧。}\BibleRef{格前 7:15}\par

当祂提升了众人的精神，使他们对自身和整个世界都有了信心后，祂加上说：\Quote{许多在先的将要在后，在后的将要在先。}但这虽然对许多其他人也同样没有区别地说，但也是对这些人和对不信的法利塞人说的，如同以前祂也曾说：\Quote{有许多人从东从西来，要与亚巴郎、依撒格和雅各伯一同坐席；但本国的子民将被驱逐出去。}\BibleRef{玛 8:11}\par

然后祂加上一个比喻，如同训练那些落后的人要有更大的积极性。\par

祂说：\Quote{天国好像一个家主，清早出去雇工人进他的葡萄园。他和工人讲定一天一个德纳，就派他们进葡萄园去。}\par

\Quote{约在第三时辰，他又看见另有些人闲站着，就对他们也说：\InnerQuote{你们也进葡萄园去，凡正义的，我必给你们。}约在第六和第九时辰，他也照样做了。约在第十一时辰，他看见还有人闲站着，就问他们说：\InnerQuote{你们为什么整天闲站着？}他们回答说：\InnerQuote{因为没有人雇我们。}他对他们说：\InnerQuote{你们也进我的葡萄园去，凡正义的，你们必得到。}}\par

\Quote{到了晚上，葡萄园的主人对他的管事的说：\InnerQuote{叫工人都来，给他们工钱，从最后的开始，到最先的。}那些约在第十一时辰雇来的人来了，每人领了一个德纳。那些最先来的人以为会多领，但他们也每人领了一个德纳。他们领了以后，就埋怨家主说：\InnerQuote{这些最后的人只工作了一个时辰，你却把他们与我们这整天承受劳苦和炎热的人同等看待。}他回答其中一人说：\InnerQuote{朋友，我没有亏待你；你我不是讲定了一个德纳吗？拿你的走吧，我愿意给这最后的人，也像给你一样。难道我不能拿我的财物做我所愿意的吗？还是因为我善良，你就眼红吗？}这样，在后的将要在先，在先的将要在后：因为被召的人多，被选的人少。}\par

3. 这个比喻对我们的用意是什么？因为开头与结尾说的并不协调，而是完全相反。因为在第一部分，祂显示所有人都享受同样的，没有一些人被驱逐，一些人被带入；但祂自己在比喻前后都说了相反的话。\Quote{在后的要在先，在先的要在后，}意思是在最前的人之前，那些不是始终保持在前而是变成在后的人。因为证明这就是祂的意思，祂加上了\Quote{被召的人多，被选的人少。}，这样既刺痛了前者，也安慰和激励了后者。\par

但比喻不是这样说，而是说他们将与那些被认可的、劳苦多时的人同等。他们说：\Quote{你使他们与我们同等，} 并说：\Quote{那些背负了全日重担和炎热的人。}\par

那么比喻的含义是什么？首先必须阐明这点，然后我们再澄清另一点。祂用葡萄园比喻天主的命令和诫命；用劳作的时间比喻现世生命；用工人比喻以不同方式被召来完成诫命的人；用清晨、第三、第九和第十一时辰比喻在不同年龄接近天主并证明自己的人。\par

但问题在于：那些光荣地证明了自己、讨天主喜悦、终日因劳作而闪耀的人，是否被最卑劣的恶习——嫉妒和猜忌——所占据？因为他们看见这些人享受同样的赏报，便说：\Quote{这些最后的人只做了一小时，你却使我们与他们这些背负了全日重担和炎热的人同等。} 说这些话时，他们虽未受伤害，自己的工钱也未减少，却因他人的好处而愤慨、极为不悦，这正是嫉妒和猜忌的证明。更甚的是，家主在为自己辩护、对说话的人解释时，揭露了他的邪恶和极端的嫉妒，说道：\Quote{你不是与我讲定了一个德纳吗？拿你的走吧！我愿意给这最后的人，如同给你一样。因为我为善，你就眼红吗？}\par

那么，通过这些事要确立什么呢？因为在其他比喻中也能看到同样的情况。那位被认可的儿子的表现，正与这里相似：当他看见那荡子兄弟享受甚至比自己更多的荣耀时，也感到同样的情绪。因为正如这些人因先获得而更快乐，那儿子也因所得之丰盛而更受尊荣；而被认可的儿子也见证了这些事。\par

那么我们可以说什么呢？在天国里，没有人会这样自我辩护或指责他人；想都不敢想！因为那地方是纯净的，没有嫉妒和猜忌。既然圣人们在世上都愿为罪人舍命，那么当他们看见罪人在那里享受这些恩惠时，更会喜乐，并视之为自己的福分。那么主为何如此讲述呢？这话是比喻，因此不必逐字细究比喻中所有细节，而应在了解其目的后收获益处，不必再费心其他。\par

那么，这个比喻为何如此构成？它的目的是什么？是要使那些在极晚年才悔改、变得更好的人更加努力，不让他们以为自己所得更少。为此，祂也引入他人对他们的福分不悦，并非表示那些人苦恼或烦恼——想都不敢想！而是教导我们：这些人所享受的尊荣之大，甚至足以引起他人的嫉妒。这也是我们常做的事，例如说：\Quote{某人责备我，因我认你配受大尊荣，} 实际上并未被责备，也不想诽谤那人，只是借此表明此人所得恩赐的伟大。\par

但为什么祂不一次召募所有人呢？就祂而言，祂召募了所有人；但若并非所有人都同时听从，那差异在于被召者的心态。因此，有些人清晨被召，有些在第三时辰，有些在第六，有些在第九，有些在第十一时辰，当他们愿意服从的时候。\par

保禄也宣告了这一点，他说：\BibleQuote{祂从母腹中就把我分别出来，并乐意} \BibleRef{迦 1:15} 祂何时乐意？当他准备服从的时候。因为祂从起初就愿意，但若保禄那时不屈服，祂便等到保禄也准备服从的时候。同样，祂召叫了强盗，虽然之前也能召叫他，但他不愿服从。因为若保禄起初不肯服从，那强盗更甚。\par

如果他们却说：\Quote{没有人雇我们，} 首先，如我所说，我们不必好奇比喻中所有细节；但在此处，家主并未说这话，而是他们说的；祂并未反驳他们，以免使他们困惑，而是为赢得他们。因为比喻表明，就祂而言，祂从起初就召叫了所有人，祂说：\Quote{他清晨出去雇人。}\par

从这一切我们清楚地看到，这个比喻是针对那些从幼年就抓住德行的人，以及那些在老年、较迟抓住德行的人说的；对于前者，使他们不致骄傲，也不可责备在第十一时辰被召的人；对于后者，使他们知道即使短时间内也能挽回一切。\par

因为主之前谈到热忱、舍弃财富、鄙视所有财物，但这需要极大的精神力量和青春热情；为了在他们心中点燃爱火，增强他们的意志，祂表明即使后来者也能得到全日的工钱。\par

但祂并不直接这样说，免得又使他们骄傲，而是表明这一切全凭祂对人的爱，正因如此，他们不会落空，反而会享受那不可言喻的福乐。\par

而祂主要就是要透过这个比喻来确立这一点。如果祂加上：\Quote{这样，最后的将成为最先的，最先的将成为最后的；因为被召的人多，被选的人少}，不要惊讶。因为祂说这话，并非从比喻推论出来，而是祂的意思是：就像这事发生一样，那事也要发生。因为在这里，最先的并没有成为最后的，而是所有人都出乎意料地得到了同样的报酬。但正如这个结果出乎意料地发生，后来者与先来看齐，同样也会发生比这更甚、更奇妙的事，我说的是：最后的甚至要排在最先的之前，而最先的要在他们之后。所以这是一回事，那是另一回事。\par

不过在我看来，祂说这些话，是在暗暗暗示犹太人，以及在信徒中那些起初发光、后来却忽略德行而退步的人；还有那些从恶行中兴起、超越许多人的其他人。因为我们在信仰和行为上都看到这样的变化。\par

因此我恳求你们，让我们用很大的勤勉，既坚守正确的信仰，又展现卓越的生活。因为如果我们不加上与信仰相称的生活，我们将遭受最严厉的惩罚。\par

蒙福的保禄从古时就指出了这一点，当他说：\Quote{他们都吃了同样的神粮，都饮了同样的神饮}，并补充说他们并未得救，\Quote{因为他们倒毙在旷野里}。基督也在福音中宣告了这一点，当祂带来一些曾经驱逐魔鬼、说过预言的人，却被引向惩罚。还有祂所有的比喻，如童女的比喻、网的比喻、荆棘的比喻、不结果子的树的比喻，都要求我们行为上的德行。因为关于教义祂很少谈论，因为那不需要费力；但关于生活则常常，或者可以说处处谈论，因为这场战争是持续的，所以劳作也是如此。\par

我何必说全部法典。因为即使忽略其中一部分，也会给人带来巨大的灾祸：例如，忽略施舍会使缺少施舍的人堕入地狱；然而这并非德行的全部，而只是其中一部分。但尽管如此，那些童女因没有施舍而受罚，那个富人因此受折磨，那些没有给饥饿者食物的人因此与魔鬼同被定罪。再者，不辱骂是其中很小的一部分，然而这也将没有做到的人驱逐出去。\Quote{因为凡向弟兄说\InnerQuote{愚妄}的，难免地狱的火。}\BibleRef{玛 5:22}同样，节制本身也是一部分，但尽管如此，没有它，没有人能见到主。因为经上说：\Quote{你们要追求与众人和睦，并要追求圣化；非圣化没有人能见主。}\BibleRef{希 12:14}同样，谦逊也是德行的一部分；但尽管如此，即使有人完成了其他善工，却没有达到这一点，他在天主面前也是不洁的。这从法利塞人身上很明显，他虽然有无数的善工，却因此失去了全部。\par

但是我还要说一些比这些更多的事。我是说，不仅忽略其中一项会关上天国之门，而且即使做了，若没有达到应有的完美和充足，也会产生同样的效果。\Quote{因为除非你们的义德超过经师和法利塞人的义德，你们决不能进入天国。}\BibleRef{玛 5:20}所以，即使你施舍，但若不超过他们，你也不能进去。\par

他们施舍了多少呢？有人可能会问。正是为这事，我现在想说，为了使不施舍的人被激励去施舍，施舍的人不自傲，反而增加他们的礼物。那么他们给了什么呢？他们所有财产的十分之一，再加上另一个十分之一，之后又第三个十分之一，这样他们几乎给出了三分之一，因为三个十分之一合起来就是这么多。此外还有初熟的果实、头胎的，以及其他东西，例如赎罪祭、洁净祭、节期的献仪、安息年的献仪、债务的豁免、仆人的释放、以及不取利息的借贷。但如果一个给出三分之一财产的人，或者一半（因为那些加起来就是一半），那么他给出了一半，还算不得什么大事，那么连十分之一都不肯给的人，又配得什么呢？祂有理由说：\Quote{得救的人很少。}\par

5. 那么，让我们不要轻视对自己生活的关心。因为如果其中一部分被轻视会带来如此大的毁灭，当我们四面都处于定罪判决之下时，我们如何逃避惩罚呢？我们将遭受怎样的刑罚呢？我们还有怎样的得救希望呢？有人会问，如果我们所数的每一件事都用地狱威胁我们？我也这么说；然而，如果我们留心，我们仍可得救，准备施舍的药膏，并照顾我们的伤口。\par

因为油不能如此强壮身体，正如仁爱立即强壮灵魂，使它对一切坚不可摧，对魔鬼固若金汤。因为无论他在哪里抓住我们，他的手都会滑脱，这油不允许他的抓握固定在我们的背上。\par

因此让我们不断地用这油膏抹自己。因为它是健康的原因、光的供应、喜乐的源泉。\Quote{但是这样的人，}你会说，\Quote{有那么多塔冷通的金子，却什么都不给。}那与你何干？因为这样，你在贫穷中比他更慷慨，就显得更值得赞叹。正是为此，保禄赞赏马其顿人，不是因为他们给予，而是因为尽管贫穷，他们仍然给予。\par

不要看这些，而要看众人共同的导师，祂\BibleQuote{没有枕头的地方}。\BibleRef{玛 8:20}你说，为什么某某人不这样做呢？不要判断别人，而要摆脱你自身的罪责。因为当你同时指责别人，自己却不做时，惩罚更大；当你判断别人时，你自己也再次落入同样的审判。因为连那些行正义的人，祂也不允许他们判断别人，更何况是犯罪的人呢？所以我们不要判断别人，也不要看那些安逸度日的他人，而要看耶稣，从祂那里汲取榜样。\par

怎么！我曾是你的恩人吗？怎么！我曾救赎了你，叫你仰望我吗？是另一位把这些赐给了你。你为什么放弃你的主，却看你的同仆？你没有听见祂说：\BibleQuote{你们跟我学，因为我是良善心谦的？}\BibleRef{玛 11:29}又说：\BibleQuote{你们中谁愿意为首，就该作众人的仆役}；又说：\BibleQuote{人子来不是受服事，而是服事人。}\BibleRef{玛 20:27}此外，为避免你因同仆中懈怠的人而跌倒，继续陷入轻慢，祂又说：\BibleQuote{我给你们立了榜样，叫你们也照样做。}\BibleRef{若 13:15}难道在你身边的人中没有一位德行的导师，能引导你做到这些吗？那么，即使没有导师供应，你成了可惊叹的人，你的赞美就会更大，称许也会更多。\par

这是可能的，不但可能，而且非常容易，只要我们愿意；那些首先妥善履行了这些事的人就证明了这一点，例如\Person{诺厄}、\Person{亚巴辣罕}、\Person{默基瑟德}、\Person{约伯}以及所有像他们的人。每天需要注目他们，而不是那些你不断效法、在集会中传讲他们名字的人。因为我到处只听到你们说这样的话：\Quote{某人买了多少亩地；某人富有，他在建造。}人啊，你为什么盯着外在的东西？你为什么看别人？如果你一心要看别人，就看看那些尽本分的人，那些自证可嘉的人，那些认真遵行法律的人，而不是那些犯罪、蒙羞的人。因为如果你看这些人，你会从中招致许多恶事，陷入懈怠、骄傲、论断他人；但如果你默想那些行义的人，你就会引导自己走向谦卑、勤勉、痛悔以及数不清的福分。\par

听一听法利塞人遭受了什么，因为他放过了行义的人，却注目那犯罪的人；听着，畏惧吧。\par

看达味如何成了可惊叹的人，因为他注目他那以德行著称的祖先。他说：\Quote{我原是旅客，是寄居的，如同我所有的祖先。}因为这个人以及所有像他的人，放过了犯罪的人，却想到了那些自证可嘉的人。\par

你也当这样做。因为你没有被设立去判断别人所犯的疏忽，也不去查问别人所犯的罪；你被要求审判自己，而不是别人。经上说：\BibleQuote{我们若先审判自己，就不致受审判；但我们受审判时，是为主所惩治。}\BibleRef{格前 11:31}但你却颠倒了次序，对自己不要求交代大罪小过，却对别人的过失严厉而好奇。\par

我们不要再这样做了，而要放弃这种混乱的方式，在我们自己内设立一个审判庭，审判我们自己犯的罪，使我们自己成为自己过错的控告者、审判者和执行者。\par

但如果你愿意忙碌于别人的事，那么你就忙于他们的善工，而不是他们的罪，这样，藉着追忆我们的疏忽，以及效法他们所行的善工，\Emph{并藉着将那个任何祈祷都不能解救的审判台摆在我们面前，每日被我们的良心像被刺棒扎伤一样}，我们就可以引导自己走向谦卑和更大的勤勉，并因我们的主耶稣基督的恩宠和祂对人的爱，获得将来的美好事物；愿光荣、权能、尊崇归于祂和父及圣神，从今直到永远，万世无疆。阿们。\par

\SourceRecord{Translated by George Prevost and revised by M.B. Riddle.；Nicene and Post-Nicene Fathers, First Series；Vol. 10.；Edited by Philip Schaff.；Buffalo, NY: Christian Literature Publishing Co.,；1888.；<http://www.newadvent.org/fathers/200164.htm>.}

% structure: p0001 source=h1 target=section reason=page-title

\section{\BibleBook{玛窦}{福音}讲道辞 第六十五篇}

\BibleRef{玛 20:17-19}\par

\BibleQuote{耶稣上耶路撒冷去的时候，在路上把十二门徒带到一边，对他们说：看，我们上耶路撒冷去，人子要被交给司祭长和经师，他们要定祂死罪，并把祂交给外邦人戏弄、鞭打、钉在十字架上；但第三天祂要复活。}\par

祂从\Place{加里肋亚}出来后，并没有立刻上耶路撒冷，而是先行了奇迹，堵住了法利塞人的口，并与门徒谈论了舍弃财产的事。因为祂说：\BibleQuote{你若愿意成为成全的，}祂说，\BibleQuote{去变卖你所有的：}\BibleRef{玛 19:21} 关于童贞，祂说：\BibleQuote{能领悟的，就领悟吧！}\BibleRef{玛 19:12} 关于谦卑：\BibleQuote{你们若不悔改，变成如同小孩一样，你们决不能进天国。}\BibleRef{玛 18:3} 关于今世的赏报：\BibleQuote{凡舍弃了房屋、或兄弟、或姊妹的，必要得到百倍的赏报。}\BibleRef{玛 19:29} 关于将来的赏报：\BibleQuote{他还要承受永生。} 然后祂就向京城进发，在即将上去的时候，再次谈论自己的苦难。因为门徒们不愿这事发生，很可能会忘记，祂便不断提醒他们，借着频繁的提醒锻炼他们的心志，减轻他们的痛苦。\par

但祂同他们\Quote{私下}说话，这是必要的；因为关于这些事的言论不宜让大众知道，也不宜说得太明白，因为这并无益处。如果门徒们听了这些事都感到困惑，那么百姓就更不用说了。\par

那么，难道没有告诉百姓吗？你可以说：确实也告诉了百姓，但没有这么明白。因为祂说：\BibleQuote{你们拆毁这座圣殿，三天之内我要把它重建起来。}\BibleRef{若 2:19} 又说：\BibleQuote{这一世代要求征兆，但除了约纳的征兆外，决不给它任何征兆。}\BibleRef{玛 12:39} 又说：\BibleQuote{我同你们在一起的时候不多了，你们要寻求我，却找不着我。}\BibleRef{若 7:33}\par

但对门徒却不是这样，正如祂把其他事更明白地告诉他们一样，这件事也是如此。那么，既然百姓不明白祂话中的力量，为什么还要说呢？为了让他们以后知道，祂是预见并甘愿受难的，并非出于无知或被迫。但对门徒们，祂预告这事不仅为此目的，而且，如我所说，是为了让他们借着预期的锻炼，能更轻松地承受苦难，免得因突然临到而慌乱。因此，开始时祂只提到自己的死亡，当门徒们练习并习惯于听到这事后，祂又加上了其他情节，例如：他们要把祂交给外邦人，要戏弄、鞭打祂。这一方面是为了让门徒们看到这些悲哀事件发生时，也能由此期待复活。因为那位没有向他们隐瞒令人痛苦和看似可耻的事，在好事上也理当被相信。\par

但请注意，祂在时间上也安排得很明智。起初并未告诉他们，免得使他们不安；也不是在事发当时才说，免得又使他们慌乱。而是在他们充分体验了祂的能力，并接受了关于永生的极大许诺之后，才引入关于这些事的话，并且一次又一次地将它们与祂的奇迹和教训交织在一起。\par

但另一位福音作者说，祂也引用了先知为证；\BibleRef{路 18:31} 另一位又说，连他们自己也不明白祂的话，这话对他们是隐藏的，他们跟着祂时都感到惊讶。\par

那么，有人会说，预言的益处就被取消了。因为如果他们不明白所听到的，就不能期待事件的发生；不期待，也就不能因预期而得到锻炼。\par

但我要说另一件更令人困惑的事：如果他们不知道，他们怎么会忧伤呢？因为另一位说，他们很忧伤。那么，如果他们不知道，怎么会忧伤呢？\Person{伯多禄}怎么又说：\BibleQuote{主，千万不可！这事绝不会临到祢身上。}\BibleRef{玛 16:22}\par

那么，我们可以说什么呢？他们确实知道祂要死，尽管并不清楚知道降生成人的奥秘。他们也不清楚知道复活的事，以及祂要成就什么；这些是向他们隐藏的。\par

为此，他们也感到痛苦。因为他们知道有些人曾被别人复活，但从未听说过有人能自己复活，并且这样复活后不再死亡。\par

这事他们虽然多次听说，却不明白；甚至对祂的这种死亡本身，他们也不清楚知道是什么，以及如何临到祂身上。因此，他们跟着祂时感到惊讶，但不仅为此；在我看来，祂谈论自己的苦难似乎也使他们惊讶。\par

2. 然而，这些事没有一样使他们鼓起勇气，尽管他们不断听到关于祂复活的事。因为连同祂的死亡，还有一件事特别困扰他们，即听到人们要\Quote{戏弄并鞭打祂}等。当他们想起祂的奇迹——祂所解救的附魔者、祂所复活的死者、祂所行的一切奇事——然后听到这些事时，他们感到惊奇：行这些事的那一位竟要如此受苦。因此他们甚至陷入困惑，时而信，时而疑，不能明白祂的话。至少他们是如此不明白祂所说的，以致载伯德的儿子们同时来到祂跟前，向祂谈论位次问题。据说：\Quote{我们愿意，一个坐在祢右边，一个坐在祢左边。}那么，这位福音作者怎么说他们的母亲来到祂跟前呢？很可能两件事都发生了。我的意思是，他们带着母亲一同前往，为要使他们的请求更有力，从而以此方式说服基督。\par

为证明我所说的这是实情，即请求实际上是他们自己的，并且他们因羞愧而推出他们的母亲，请看基督如何把祂的话直接指向他们。\par

但是，让我们先学习：他们求什么？怀着怎样的心意？是什么促使他们这样做？那么，是什么促使他们这样做呢？他们看到自己比其他门徒更受尊荣，便期望也能得到这个请求的应允。但他们所求的到底是什么呢？听另一位福音作者清楚说明这事。因为\Quote{祂临近耶路撒冷，而且因为他们以为天主的国即将显现}，\BibleRef{路 19:11}他们才求这些事。因为他们以为这国已在门口，是看得见的，而且他们以为一旦得到所求的，就不会遭受任何痛苦。因为他们寻求这个，不仅是为其本身，也好像是想要逃避艰难。\par

因此基督首先引导他们离开这些想法，命令他们等候杀戮、危险和极度的恐惧。因为祂说：\Quote{你们能喝我将要喝的杯吗？}\BibleRef{玛 20:22}\par

但任何人不要因宗徒们处于如此不完善的状态而困扰。因为那时十字架尚未完成，圣神的恩宠尚未赐下。但如果你愿意学习他们的美德，请注意这些事以后他们的情况，你将看到他们超乎一切情欲之上。因为祂揭示他们的不足，正是为此，使你日后知道他们因恩宠成了怎样的人。\par

由此可见，他们当时所求的其实没有什么是属灵的，也没有想到上头的国。但让我们也看看他们怎样来到祂跟前，以及他们说什么。据说：\Quote{我们愿意，我们无论向祢求什么，就给我们成就。}\BibleRef{谷 10:35}\par

基督对他们说：\Quote{你们要什么？}\BibleRef{谷 10:36}祂并非不知，而是为要迫使他们回答，使伤口暴露，然后敷上药物。但他们出于羞愧和窘迫，因为受人情欲驱使才来做这事，便私下将祂带离其他门徒，向祂请求。据说，他们走在前面，以免被别的门徒察觉，就这样说出了他们的愿望。因为他们曾听到\Quote{你们要坐在十二个宝座上}，于是想要在这些座位中占据首位。他们知道自己比其他门徒更有优势，但却惧怕伯多禄，于是说：\Quote{命令，一个坐在祢右边，一个坐在祢左边；}并且催促祂说：\Quote{命令。}\par

那么祂说什么呢？祂表明他们所求的毫无属灵之事，而且，如果他们知道所求的是什么，也不敢求这么大的事。祂说：\Quote{你们不知道你们求的是什么}——这所求是多么伟大、奇妙，甚至超越上界的能力。此后祂又加上：\Quote{你们能喝我将要喝的杯吗？并且受我即将受的洗吗？}你看见祂如何通过从相反的话题来构建祂的言论，立刻打消了他们的猜疑吗？祂说：你们跟我谈尊荣和冠冕，但我却对你们谈争战和劳苦。因为现在不是赏报的时候，我的荣耀也不会在此时显现，当下乃是杀戮、战争和危险的时候。\par

请看祂如何通过问话的形式，既催促又吸引他们。因为祂没有说：\Quote{你们能被杀吗？}\Quote{你们能倾流你们的血吗？}而是怎样说呢？祂说：\Quote{你们能喝这杯吗？}然后为吸引他们参与，祂说：\Quote{就是我将要喝的那杯。}这样，他们因与祂在此事上的共融而更加愿意。\par

祂又称之为一种洗礼，表明世界将从正在发生的事中得到极大的净化。\par

\Quote{他们对祂说：我们能。}\BibleRef{玛 20:22}由于他们的冒失，他们立即承担了这事，甚至不知道自己在说什么，只盼着听到他们所求的事。\par

那么祂说什么？\Quote{你们确实要喝我的杯，并且受我受的洗。}祂预言了极大的福分给他们。祂的意思是说：你们将被视为配得殉道，并要受我所受的苦；你们要以暴力的死亡结束生命，在这些事上你们将与我同受；\Quote{但坐在我的右边和左边，不是我可以赐予的，而是给那些为我父所预备的人。}\par

3. 祂先提升了他们的灵魂，使他们具有更高的品格，使他们成为哀伤不能征服的人，然后才责备他们的请求。\par

但这当前的话是什么意思呢？确实有两方面是许多人探究的问题：第一，是否任何人被预备坐在祂右边；第二，万有之主难道没有能力将之赐予祂所预备的人吗？\par

这句话是什么意思？如果我们解决了前一个问题，那么第二个问题对询问者也会变得清晰。那么这是什么？没有人会坐在祂右边或左边。因为那个宝座是所有人都不可及的，不仅是对人、圣人、宗徒，而且甚至是对天使、总领天使和所有天上的权势而言。\par

至少保禄将此作为独生子的独特特权，说：\Quote{祂曾对哪一个天使说过：\InnerQuote{你坐在我右边}？\BibleRef{希 1:13} 至于天使，祂说：\InnerQuote{祂使祂的天使为神风}，但论到子：\InnerQuote{天主啊，你的宝座。}} \BibleRef{希 1:7}\par

那么祂怎么说：\Quote{坐在我右边或左边，不是我可以给的，}好像有些人在那里坐？并非如此；祂是对那些求这恩宠之人的思想作答，迁就他们的理解。因为他们既不知道那崇高的宝座，也不知祂坐在圣父右边；他们怎么会知道呢？他们连比这些低得多的事，就是每日灌输给他们的，也不明白。但他们只寻求一件事：享受首位，站在其余人之前，且无人能胜过他们与祂在一起；正如我先前已经说过的，因为他们听到了十二个宝座，却不知道这句话是什么意思，便要求首位。\par

因此基督所说的就是这一点：\Quote{你们确实要为我而死，并为福音而被杀，且要与我一同有分，就受难而言；但这不足以使你们享有首位，并让你们占据第一个位置。因为如果有别人前来，虽同为殉道，却比你们更圆满地具备其他一切德行，那么不是因为我现在爱你们，且以你们为优于他人，我便会将那位因善工而卓越的人撇在一边，而把首位给你们。}\par

但祂没有这样说，以免使他们痛苦，而是含糊地暗示了同样的事，说：\Quote{你们确实要饮我的杯，并要受我所受的洗礼；但坐在我右边或左边，不是我可以给的，而是给那些预备了的人。}\par

但它是为谁预备的呢？为那些能因善工而卓越的人。因此祂没有说：不是我可以给的，而是我父给的，免得有人说祂太软弱或缺乏能力来回报他们；而是说：不是我可以给的，而是为那些预备了的人。\par

为了使我所说的更明白，让我们用一个比喻来说明。假设有一位竞赛的主持人，有许多出色的竞技者下场比赛，其中有两位与主持人最亲近的竞技者来对他说：\Quote{请使我们得冠冕并被宣布为优胜者，}他们信赖与他的友谊和善意；而他对他们说：\Quote{这不是我可以给的，而是给那些通过劳苦和辛劳为此预备了的人。}我们会因此断定他无能吗？绝不，我们会称赞他的公正和不徇情面。就像我们不会说他因缺乏能力而不给冠冕，而是因为他不想破坏竞赛的规则，也不想扰乱正义的秩序；同样，我现在要说，基督这样说，是为了从各个方面促使他们，在获得天主的恩宠之后，将救恩和赞许的希望寄托于自己善工的证明。\par

因此祂说：\Quote{给那些预备了的人。}祂说，如果别人显出比你们更好呢？如果他们做了更大的事呢？难道因为你们做了我的门徒，你们就会享受首位，如果你们自己显得不配被选呢？\par

因为祂对整个事物有权力，从祂拥有全部审判权就明显了。因为祂也对伯多禄这样说：\Quote{我要把天国的钥匙给你。}保禄也清楚地表明这一点，他说：\Quote{从今以后，正义的冠冕已为我预备好了，就是主，正义的审判者，在那一天要赏给我的；不但赏给我，也赏给所有爱慕祂显现的人。} \BibleRef{弟后 4:8}\par

如果祂把这些话说得有些隐晦，不必惊奇。因为祂想通过隐含的教训引导他们，使他们不再无目的、无理由地粗鲁要求首位（因为他们出于人的情欲而如此感觉），而又不想使他们痛苦，通过隐晦祂实现了这两个目的。\par

\Quote{于是那十个人就恼怒那两个人。}那时。什么时候？当祂责备了他们之后。只要审判是基督的，他们就不恼怒；但看到这两人被偏爱，他们就满意了，出于对师傅的尊敬和敬意而保持沉默。\par

即使他们心中烦恼，却不敢说出来。当他们因人的软弱而对伯多禄有些不满时——就是祂交双德拉克马的时候——他们并没有发怒，而只是问：\Quote{那么谁最大？}但既然此处是门徒们的请求，他们便恼怒起来。而且他们也不是在请求时就立刻恼怒，而是当基督责备了他们，并说除非他们显出自己配得这些，否则不会享有首位时。\par

4. 你们看到了吗？他们都处于不完美的状态，因为这两个人在那十人之上抬高自己，而那十人则嫉妒这两个人。但是，正如我所说，在这一切之后，你们看他们，就会发现他们摆脱了所有这些情欲。至少听听这个若望，就是现在为此事来见祂的那位，是如何在\WorkTitle{宗徒大事录}中处处把首位让给伯多禄，无论在向民众讲话，还是在行奇迹上。\par

他并不隐瞒\Person{伯多禄}的善行，反而叙述了当众人都沉默时他公开所作的宣认（\BibleRef{若 6:68}），以及他进入坟墓的事（\BibleRef{若 20:6}），并将这位宗徒放在自己之前。因为两人在祂被钉十字架时都陪伴着祂，为了除去自己受赞扬的根据，他说：\BibleQuote{那个门徒是大司祭所认识的。} \BibleRef{若 18:15}\par

但是\Person{雅各伯}没有活很长时间，而是从一开始就充满了极大的热忱，并如此抛弃了人类的一切事物，攀升到不可言喻的高处，以至立刻被杀害。因此，在各方面，他们以后都成了卓越的。\par

但是接着，\Quote{他们忿忿不平。} 基督于是说什么？\Quote{祂把他们叫到跟前，并说：外邦人的首领管辖他们。} 因为，他们既然感到不安和困扰，祂便透过在说话前先召唤他们、使他们亲近祂来安抚他们。原来那两人从十人的群体中分开，站得更靠近祂，为自己的利益恳求。因此，祂也藉着这个举动使他们亲近祂，并在其他人面前揭露和显明这事，安抚了这一个和那一个的情绪。\par

而且祂现在不像从前那样约束他们。因为从前祂把小孩子带到中间，命令人效法他们的单纯和谦卑；但在这里，祂以更严厉的方式从相反的一面责备他们，说：\Quote{外邦人的首领管辖他们，他们的有权势者对他们施行权威，但在你们中却不可这样；相反，你们中谁愿意成为大的，就让他作众人的仆人；谁愿意成为首位的，就让他作众人的末位。} 以此显示这样的情绪是外邦人的，我是指喜爱首位。因为这情欲是暴虐的，并不断阻碍即使是伟大的人；因此也需要更严厉的鞭策。因此，祂更是深深地刺入他们，藉着与外邦人的比较来羞辱他们燃烧的灵魂，并消除一方的嫉妒和另一方的傲慢，几乎是在说：\Quote{不要因为被侮辱而忿忿不平。因为谁以这种方式寻求首位，他们最伤害和羞辱自己，因为他们是在末位。因为我们的情况不像世间的情况。\InnerQuote{外邦人的首领管辖他们}，但在我这里，末位的人，甚至他才是首位。}\par

\Quote{为了证明我不是无缘无故说这些话，请从我所作和所受的事中接受我言语的证明。因为我亲自做了更甚的事。因为我是天上万军的君王，却甘愿成为人，并忍受被藐视和恶待。即使如此，我仍不满足，甚至一直来到死亡。因此，}祂说，\par

\BibleQuote{正如人子来不是为受服事，而是为服事人，并交出自己的生命，为大众作赎价。} \BibleRef{玛 20:28} \Quote{因为连这一步我也没有停住，}祂说，\Quote{反而献出了我的生命作为赎价；是为了谁？为了仇敌。但你若受辱，是为了你自己，而我却是为了你。}\par

那么，不要害怕，好像你的荣誉被剥夺了。因为无论你怎样谦卑自己，你都不能降得如同你的主一样低。然而祂的下降成了众人的上升，并使祂自己的荣耀彰显出来。因为祂在成为人之前，只在天使中被认识；但祂成为人并被钉十字架之后，那荣耀非但没有减少，反而获得了另外的，即来自世界认识的荣耀。\par

那么，不要害怕，好像你的荣誉被贬低，如果你谦卑自己，因为这样你的荣耀更被高举，这样它变得更大。这是王国的大门。那么我们不要走相反的路，也不要以自己为敌。因为如果我们想要显得伟大，我们就不会伟大，反而会成为众人中最受羞辱的。\par

你看见祂如何处处藉着相反的事物催促他们，赐给他们所渴望的吗？因为在前面的部分我们也在许多例子中表明了这一点，并在贪婪者和虚荣者的情形中，祂是这样做的。因为祂说：你为什么在人前施舍？是为了享受荣耀吗？那么你不应该这样做，你就必然享受它。你为什么积蓄财宝？是为了致富吗？那么你不应该积蓄财宝，你就会富有。同样在这里，你为什么一心追求首位？是为了超越别人吗？那么就选择末位，然后你就会得到首位。所以，如果你意愿成为伟大，不要寻求成为伟大，然后你就会伟大。因为相反的便是渺小。\par

5. 你看见祂如何使他们脱离这疾病，既向他们显示从那里会达不到目的，又显示从这里会有所得，好让他们逃避那一个，而追求这一个。\par

并且祂也为此提醒他们外邦人的事，好这样再次显示这事是可耻的并应被憎恶的。\par

因为骄傲的人必然是卑贱的，反之，谦卑的人是高尚的。因为这才是真实和真正的崇高，并不仅仅存在于名号上，也不存在于称呼的方式上。那来自外方的崇高出于必要和恐惧，但这崇高相似天主的崇高。这样的人，即使不被任何人钦佩，仍是高尚的；同样，那一个人，即使被所有人奉承，仍是众人中最卑贱的。一个是由必要而给予的荣誉，因此也容易消逝；但另一个是由原则而来的，因此也持久稳固。因为我们也因此钦佩圣人，他们比众人都伟大，却比众人都谦卑自己。所以直到今天，他们仍然是高尚的，甚至死亡也没有降低那崇高。\par

若你愿意，让我们也藉着推理来探究这事本身。任何人被称为高尚，或是因为他身材高大，或是因为他恰巧被安置在高处；卑贱也是如此，来自相反的情形。\par

那么让我们看看谁是这样的，是自夸者，还是那保持适度的人，好让你领悟，没有什么比谦卑的心更高，也没有什么比自夸更低。\par

因此，自夸者想要比所有人都大，并声称没有人配得上与他相等；无论他得到多少荣誉，他都渴望更多，要求更多，认为自己什么也没有得到，对他人极为藐视，却又追求从他们而来的荣誉——还有什么比这更不合理呢？这确实像一个谜。他轻视这些人，却渴望被他们荣耀。\par

你看见了吗？那想要被高举的人，反而跌倒在地上。因为他认为所有人相比他都算不得什么，他自己就表明这一点——这就是自夸。那么，你为什么要把自己投向那算不得什么的人呢？你为什么向他寻求尊荣呢？你为什么带着这么多群众跟随你呢？\par

你看见一个谦卑的人，处在卑微的位置。那么，让我们来查考那高尚的人。他知道人是什么，知道人是伟大的，而他自己是众人中最末的，因此无论他享受什么尊荣，他都视为重大，所以他是前后一致、高尚的，判断不动摇。因为他看重的人，从他们而来的尊荣他也看为重大，即使那尊荣可能微小，因为他认为赐予尊荣的人是伟大的。但自夸的人认为赐予尊荣的人算不得什么，却把所赐的尊荣看为重大。\par

再者，谦卑的人不被任何情欲控制，愤怒不能大大困扰他，虚荣、嫉妒、忌恨也不能。还有什么比摆脱了这些事物的灵魂更高尚呢？但自夸的人却被这一切所奴役，像爬行在泥淖中的虫子，因为忌妒、嫉妒和愤怒不断搅扰他的灵魂。\par

那么，哪一个是高尚的呢？是胜过情欲的人，还是情欲的奴隶？是战战兢兢害怕它们的人，还是不为所动、从未被它们俘获的人？我们该说哪一种鸟飞得更高？是飞得高于猎人的手和箭的，还是那因沿地飞行、从未能升高，而使猎人甚至不需要用箭的？傲慢的人不就是这样吗？因为每一种网罗都轻易地捉住他，因为他在地上爬行。\par

6. 但如果你愿意，甚至从那个邪恶的魔鬼身上也可以证明这一点。还有什么比魔鬼更卑下，因为他曾自高自大？还有什么比那自愿谦卑的人更高尚？前者在我们脚下爬行（因为主说：\Quote{你们践踏毒蛇和蝎子}），但后者与天使同在高处。\par

但如果你也愿意从骄傲之人的例子学习这一点，请想想那个率领如此大军的蛮族国王，他连众人皆知的事都不懂：例如，石头是石头，偶像是偶像；因此他甚至比这些东西还低下。然而虔敬而忠诚的人却被高举到太阳之上；还有什么比他们更高？他们上升到穹苍之上，超越天使，站在君王的宝座前。\par

为使你从另一角度明白他们的卑下：谁会被贬抑？是有天主作盟友的人，还是天主与他为敌的人？显然是与天主为敌的人。那么，请听圣经论及这两者所说的话。\BibleQuote{天主拒绝骄傲的人，却赐恩宠给谦卑的人。}\BibleRef{雅 4:6}\par

再者，我还要问你另一件事：谁更高尚？是向天主履行司祭职、献上祭品的人？还是远离祂的信任、不能坦然面对祂的人？谦卑的人献上什么祭品呢？可以这样说。请听达味的话：\BibleQuote{天主所喜悦的祭品是痛悔的精神；痛悔和谦卑的心，天主必不轻视。}\par

你看见这人的纯洁了吗？也请看另一人的不洁，因为\Quote{凡心里骄傲的，在天主面前都是不洁的。}此外，一人有天主住在内（因为祂说：\BibleQuote{我要垂顾的是谁呢？不是那谦卑、安静、敬畏我话语的人吗？}\BibleRef{依 66:2}），但另一人与魔鬼一同爬行，因为那心高气傲的人要受魔鬼的惩罚。因此保禄也说：\BibleQuote{免得他自高自大，陷于魔鬼所受的判决。}\BibleRef{弟前 3:6}\par

结果，事情与他所愿的正相反。他想要傲慢，以便被尊荣，但这种人却最受鄙视。因为他们最是笑柄，与众人为敌，最容易被敌人制服，轻易发怒，在天主面前不洁。\par

还有什么比这更糟呢？因为这是罪恶的极限。有什么比谦卑的人更甘甜、更有福呢？因为他们被天主渴慕和爱慕。而且，他们也最享受来自人的荣耀，所有人都尊敬他们如同父亲，拥抱他们如同兄弟，接纳他们如同自己的肢体。\par

那么，让我们变得谦卑，好使我们被高举。因为傲慢最彻底地贬抑人。傲慢贬抑了法郎。因为他说：\BibleQuote{我不认识上主，}\BibleRef{出 5:2}于是他变得比苍蝇、青蛙和蝗虫还不如，最后连同他的军队和战马都淹死在海中。与他对立的是亚巴郎，他说：\BibleQuote{我只是灰尘和泥土，}\BibleRef{创 18:27}却战胜了无数的蛮族，又陷入埃及人中间，却带着比先前更光荣的战利品回来，并且持守这美德，日益高超。因此，他处处受到歌颂和加冕；而法郎却成了尘土和灰烬，甚至比这些更卑下。因为天主最厌恶傲慢。为此，祂从起初就做了一切，为要根除这情欲。因为傲慢，我们变成了必死的，处于忧愁和哀哭之中；因为傲慢，我们劳苦、流汗、不断工作，并夹杂着痛苦。因为第一个人的确是出于傲慢而犯罪，想要与天主同等。因此，连他所拥有的也未能保持，甚至失去了这些。\par

傲慢就是这样，非但不能给我们的生活增添任何进步，反而连我们已有的也夺去；相反，谦卑非但不夺去我们已有的，反而把我们没有的也加给我们。\par

那么，让我们效法这美德，让我们追求它，好使我们既享有现世的荣誉，也获得将来的光荣，藉着我们的主\Person{耶稣基督}对人的恩宠和慈爱，愿光荣和德能归于\Person{圣父}，偕同\Person{圣神}，从今世直到永远。阿们。\par

\SourceRecord{Translated by George Prevost and revised by M.B. Riddle.；Nicene and Post-Nicene Fathers, First Series；Vol. 10.；Edited by Philip Schaff.；Buffalo, NY: Christian Literature Publishing Co.,；1888.；<http://www.newadvent.org/fathers/200165.htm>.}

% structure: p0001 source=h1 target=section reason=page-title

\section{玛窦福音讲道集 第66篇}

\Emph{玛 20:29-30.}\par

\Quote{他们从耶里哥出发，有许多群众跟随祂。看，有两个瞎子坐在路旁，听说耶稣经过，就喊叫说：\InnerQuote{主，达味之子，可怜我们吧！}}\par

请看祂从何处进入耶路撒冷，以及在此以前祂住在何处，由此我认为特别值得探究的是，为什么祂没有更早地从那里前往加里肋亚，而是经过撒玛黎雅。但这一点我们留给那些喜爱学习的人去研究。因为若有人愿意仔细考查这事，他会发现若望清楚地暗示了这一点，并说明了原因。\par

但我们还是专注于摆在我们面前的事情，让我们听听这些瞎子，他们比许多明眼人更好。因为他们既没有向导，也不能看见祂走近他们，却仍然努力来到祂面前，开始大声喊叫，当因说话受斥责时，他们喊得更厉害了。因为坚忍的灵魂的本性就是这样，正是那些阻碍它的事物，反而把它托起。\par

但基督容许他们受斥责，为的是使他们的热忱更能彰显出来，并让你知道，他们配得享受治愈的益处。所以祂没有像对许多人那样问他们说：\Quote{你们信吗？}因为他们的喊叫和来到祂面前，已足以显明他们的信德。\par

因此，亲爱的诸位，请学习这一点：即使我们极其卑贱和被遗弃，但若以热忱接近天主，我们自己也必能成就我们所求的一切。请看这些人，他们没有一位宗徒为他们求情，反而有许多人阻止他们开口，但他们却能越过障碍，来到耶稣自己面前。然而，圣史没有为他们作任何生活信德的见证，但热忱足以替代一切。\par

让我们也效法他们。即使天主延迟恩赐，即使有许多人劝阻我们，我们也不要停止祈求。因为这样我们最能把天主争取到我们这边。请看，至少在这里，贫穷、瞎眼、不被听见、受群众斥责，没有任何东西能阻止他们极度的热忱。这就是一个热切而劳苦的灵魂的本性。\par

那么，基督说什么呢？\Quote{祂叫了他们来，说：\InnerQuote{你们愿意我给你们做什么？}他们说：\InnerQuote{主，愿我们的眼睛被打开。}}（\BibleRef{玛 20:32}）祂为什么问他们？免得有人以为他们想接受一件东西，而祂却给了他们另一件。因为祂惯常每次首先将祂所医治之人的德行显明并展示给众人，然后才施行治愈；一方面是为了激励别人也效法，另一方面是为了表明他们配得享受这恩赐。例如，祂曾对客纳罕妇人这样做，也曾对百夫长这样做，又对血漏妇人这样做；更确切地说，那位了不起的妇人甚至预见了主的询问，但祂并没有忽略她，而是在治愈之后仍然使她显明出来。祂在每一个场合都如此热切地宣扬那些来到祂面前的人的好行为，并显明它们比实际更大，祂在这里也这样做。\par

然后，当他们说出自己的愿望时，祂动了怜悯之心，并触摸了他们。因为这就是他们得治愈的唯一原因，也是祂来到世界的原因。但尽管如此，虽然这是慈悲和恩宠，它仍寻找配得的人。\par

他们是配得的，这一点从他们所喊叫的，从他们在领受之后没有急忙离去（像许多人那样，受惠后忘恩负义）就显明了。不，他们并非如此，而是在领受之前坚持，在领受之后感恩，因为\Quote{他们跟随了祂}。\par

\Quote{当祂临近耶路撒冷，来到贝特法革，橄榄山那里，祂派遣了两个门徒，说：\InnerQuote{你们往对面的村庄去，会立刻看见一匹拴着的母驴，和同她在一起的小驴驹。解开它们，牵到我这里来。如果有人对你们说什么，你们就说：\InnerQuote{主需要它们。}那人就会立刻让它们来。}这事发生，是为应验先知匝加利亚的话，说：\InnerQuote{你们应向熙雍女子说：看，你的君王到你这里来，温和地骑着一匹驴，一匹小驴驹，即驴的崽子。}}\par

然而，祂以前常常进入耶路撒冷，但从未如此郑重其事。原因何在？那时是救恩工程的开始，祂不太为人所知，祂受难的时期也未临近；因此祂与人们相处时较为低调，更多隐藏自己。因为若祂这样显现，就不会受人钦佩，反而会激起他们更大的愤怒。但既已给了他们充足的能力证据，又因为十字架已近在眼前，祂便使自己更加引人注目，并更加郑重地做那些可能激怒他们的事。因为这在起初也是可能的，但那样做既无益处也不合时宜。\par

但是，我请求你注意，发生了多少奇迹，应验了多少预言。祂说：\Quote{你们会找到一匹驴；}祂预言没有人会阻止他们，反而所有人听见时都会保持沉默。\par

但这对犹太人来说并非轻微的谴责：那些素不相识、从未在祂面前出现的人，祂却说服他们放弃自己的财产，并不加反对，而且是通过祂的门徒做到的；而这些人（犹太人）在祂行奇迹时与祂同在，却不被说服。\par

2. 切莫以为所发生的事是小事。因为，当他们的财物被夺走，他们自己也许是穷苦的农夫，是谁说服了他们，使他们不加禁止呢？我岂是说禁止？甚至不发问，或是即使发问，也缄默不语，将之交出。因为这两件事同样令人惊奇：要么在牲口被拖走时一言不发，要么开口后听见\Quote{主需要它们}，便屈服、不抗拒，且此时他们看见的不是祂，而是祂的门徒。\par

祂借此教导他们：祂本有能力完全阻止犹太人，甚至在他们蓄意攻击祂时违背他们的意愿，使他们哑口无言，但祂不愿如此。\par

此外，祂同时还教导门徒另一件事：祂所求的任何事物，都要给予；即使祂要求他们舍弃自己的性命，也应交出，不得违抗。因为若连外人都向祂交出一切，他们岂不更该脱去所有？\par

除了我们所说过的，祂还成就了另一项预言，这预言是双重的：一部分在言语中，另一部分在行动中。在行动中的是骑驴；在言语中的是匝加利亚的预言，因为他曾说，君王要骑在驴上。祂骑上驴，成就了这预言，同时又给预言开了一个新的开端，借着祂所行的预先象征未来之事。\par

如何以及以何种方式？祂预先宣告了不洁外邦人的召叫，祂将安息于他们身上，他们将服从祂、跟随祂，预言接续预言。\par

但在我看来，祂骑驴似乎不仅为此目的，也为我们提供克己的准则。因为祂不仅成就预言，不仅栽种真理的教义，还借着这些事纠正我们的行为，处处为我们立下必要的规范，以各种方式改善我们的生活。\par

为此，我说，当祂将要诞生时，祂没有寻找华丽的房屋，也没有寻找富有而尊贵的母亲，而是一位贫穷的妇人，一位与木匠订婚的妇人；祂生于马厩，放在马槽里。拣选门徒时，祂没有挑选演说家、智者、富人、贵族，而是贫穷的人，出身贫寒、各方面卑微的人。预备饭食时，一次祂面前摆的是大麦饼，另一次又即时吩咐门徒去市场购买。铺床时，祂以草为席；穿衣时，祂穿着廉价、与常人无异的衣服；祂连一间房屋也没有。若须从一处往另一处，祂步行，且行至疲惫。坐时，祂不需宝座或枕头，而是坐在地上，有时在山中，有时在井旁，且不仅是在井旁，也是独自一人，与一位撒玛黎雅妇人谈话。\par

再者，为悲伤定下尺度，当祂需要哀恸时，祂适度哭泣，处处为我人立下准则，如我所说，定下当行多远，不可越过。因此现在也为这目的，既然有些人软弱，需要牲口驮负，祂也在此定下尺度，表明人不该套上马或骡来驮自己，而该用驴，不可逾矩，处处以需要为限度。\par

但我们也来看那预言——言语的预言与行动的预言。那么预言是什么？\Quote{看，你的君王到你这里来，温和地骑在驴上，和一头小驴驹上；}\BibleRef{匝 9:9}不像其他君王驾着战车，不索取贡物，不推挤人，不带卫兵，而是借此彰显祂极大的温和。\par

那么请问犹太人：有哪个君王骑驴进入耶路撒冷？他不能说出别的，唯独这一位。\par

但祂行这些事，如我所说，是预先象征未来之事。因为这里驴驹象征教会，象征新的人民——他们原是不洁的，但耶稣骑上他们后，就成了洁净的。请看这图像贯穿始终：我说，门徒解开驴驹。因为借着宗徒，他们和我们都被召叫；借着宗徒，我们得以亲近。但因着我们的接纳也激发了他们的竞争，所以驴跟随着驴驹。因为基督骑上外邦人之后，他们也将前来激发我们的竞争。保禄宣告这一点，说：\Quote{以色列有一部分硬心，直到外邦人的数目全满了，然后全以色列必要得救。}\BibleRef{罗 11:25}因为从所说的话明显知道，这是一项预言。否则先知何须如此精确地表达驴的年龄，除非真是如此。\par

但所说的事不仅象征这些，也象征宗徒将轻易带来他们。因为如同在这里，无人阻止他们牵驴，同样，关于外邦人，从前作他们主人的那些人，也不能阻止他们。\par

但祂不是坐在光秃的驴驹上，而是坐在宗徒的衣服上。因为牵了驴驹后，他们便交出一切，如保禄所说：\Quote{我极甘愿付出我自己，也为你们的灵魂而耗尽。}\BibleRef{格后 12:15}\par

但请看这驴驹多么驯服：它从未被驯服，从未认识缰绳，却不抗拒，而是有序前行。这件事本身便是未来的预言，象征外邦人的顺服，以及他们突然转向良好秩序。因为那句\Quote{解开它，带到我这里来}的话成就了一切，使不受管束的变为有序，不洁的从此成为洁净。\par

3. 但请看犹太人的卑劣。祂行了那么多奇迹，他们从未如此惊叹；但当他们见群众聚集，便惊叹了。\Quote{全城都骚动起来，说：这人是谁？群众却说：这是加里肋亚纳匝肋的先知耶稣。}当他们自以为说着了不起的话时，他们的思想仍是属地的、卑下的、拖在地上的。\par

但祂做这些事，并非为显耀排场，而是如我所说，既应验预言，又教导克己，同时也安慰那些为祂的死而忧伤的门徒，并向他们显示祂是甘心承受这一切。我请你留意先知如何准确预知了一切。有些事由\Person{达味}预告，有些则由\Person{匝加利亚}预告。我们也当照样行事，歌唱赞歌，将我们的外衣给予那背负祂的人。因为我们该当何罪？当有人给祂所骑的驴子披上衣服，甚至有人将衣服铺在驴蹄之下，而祂赤身露体时，我们见祂赤身，甚至未被命令脱衣，只需花费存积之物，却仍不慷慨？当他们在祂前后侍奉时，我们却当祂来到我们这里，将祂送走，推开祂，侮辱祂。\par

这些事该当何等严厉的惩罚、多大的报应！你的主来到你面前，身处困乏，你连倾听祂的恳求都不愿，反而指责、斥责祂，而且是在你听了这样的话之后。如果施舍一个饼、一点钱，你都这样吝啬、高傲、退缩；若是要你倾尽所有，你又会变成什么样子呢？\par

你难道没有看见那些在剧场中炫耀奢华的人，给妓女们多少馈赠？但你连一半也不给，甚至常常连最小部分也不给。魔鬼怂恿你任意施舍，为你招致地狱，你便施舍；而基督向穷人施舍，应许天国，你不但不给，反而侮辱他们。你宁可选服从魔鬼而受罚，也不愿服从基督而得救。\par

还有什么比这癫狂更糟的呢？一个招致地狱，另一个赐予天国，你却舍弃后者，奔向前者。当祂来到你面前时，你送走祂；当祂远离你时，你却呼求祂。你所做的，就好像一位国王身披王袍、献上冠冕，未能赢得你的选择，而一个强盗向你挥舞刀剑、威胁死亡，却赢得了它。\par

亲爱的，既想念这些事，让我们终久辨明真理，清醒过来。因为我现在羞于谈论施舍了，因为我多次谈论这个主题，却未产生与劝勉相称的效果。确实有些增长，但未如我所愿。因为我见你们播种，却非慷慨之手。因此我也害怕你们会\Quote{少收}。\BibleRef{格后 9:6}\par

为证明我们确实是少收，让我们查问，若以为合适，城中是穷人多还是富人多？那些既不穷也不富、处于中间地位的又是哪些？例如，十分之一是富人，十分之一是一无所有的穷人，其余的是中等阶层。\par

那么，让我们把全城众人分给穷人，你便会看到耻辱何其大。真正的富人很少，其次者很多；再者，穷人比这些少得多。尽管如此，虽有如此多能喂养饥饿者的人，仍有许多人饿着肚子睡觉，不是因为那些有财物者不能轻易救助他们，而是因为他们极其野蛮与不人道。因为若富人和他们附近的人都将需要衣食的穷人分给自己，几乎每个穷人都可摊到五十甚至一百人。然而，虽有如此众多的帮助者，他们每日仍在哀哭。为让你得知他人的不人道，教会虽拥有最穷富人之一而非大富者的进项，请思考她每日救助多少寡妇、多少贞女；她们的名单已达三千之数。此外，她还救助狱中之人、旅店病人、健康者、离乡者、身体残疾者、侍奉祭台者；并在衣食上救助每日偶然到来者；她的资财并未减少。所以，若只有十人愿意这样花费，便不会有穷人了。\par

4. 有人会说：\Quote{那我们的子女继承什么呢？}本金仍在，收入复又更为丰裕，财物在天上为他们积蓄了。\par

你不愿这样做吗？至少做一半，至少三分之一，至少四分之一，至少十分之一。因天主的恩宠，我们的城市本可养活十座城市的穷人。\par

你若愿意，让我们作些计算来证明这一点；或者毋需计算，事情本身就显出其容易。至少在公共场合，请看一个家庭往往如何不惜花费，甚至不曾感到开销之痛；若每个富人都愿意为穷人提供这样的服务，转瞬之间他便可夺取天堂。\par

那时有何托辞？有何辩护的阴影？当我们必定要分离的那些东西，在此处被取走后，我们连像别人向舞台人士那样慷慨地分施给穷人都不肯，尽管我们将从中获得如此多的益处？即使我们永远留在此处，也不该吝惜这种善用；但当我们不久之后要离开此处，赤身地被拖走，我们有什么辩护，连从收入中也不给饥饿困苦者？\par

我并非强迫你减少财产，并非因为我不愿意，而是因为我看你们十分退缩。我说的不是这个，而是：从你的果实中花费，不要从中积聚任何东西。收入之钱如泉水般涌流给你，这足够了；让穷人分享你的财物，成为天主赐予你之物的好管家。\par

但有人会说：我缴了税。为此你就藐视吗？因为在这种情况下无人向你索要？而那一位，无论大地出产与否，都强行勒索，你不敢反驳；而这一位如此温和，只在土地出产时才索要，你连一句话也不回答？谁将把你从那些不可忍受的刑罚中解救出来呢？没有人。因为如果在另一种情况下，由于不给予而会招致非常严厉的惩罚，因此你殷勤地支付；那么在这里也想想，有一种更严厉的：不是被捆绑，也不是被投进监牢，而是被投入永火之中。\par

因此，让我们出于一切理由首先缴纳这些税赋：因为方便甚大，赏报更大；收益更丰富，而我们若顽梗不化，惩罚也更重。因为降临到我们身上的惩罚，是没有终结的。\par

但若你对我说，士兵为你与蛮族作战，那么这里也有一个军营，即穷人的军营，以及一场战争，即穷人为你进行的战争。因为他们领受时，藉着祈祷使天主和蔼可亲；使祂和蔼可亲后，他们击退的，不是蛮族，而是魔鬼的袭击；他们不容恶者肆虐，也不容他持续攻击我们，反而削弱了他的力量。\par

5. 因此，既然每天都看见这些士兵藉着他们的恳求和祈祷为你与魔鬼作战，那么你当向自己索取这良好的捐献，即他们的供养。因为这位国王温和，没有指派任何人向你索要，却愿意你心甘情愿地给予；即使你一点一点地支付，祂也接受；即使你处境困难，很久以后才支付，祂也不催逼那没有的人。\par

那么，我们不要藐视祂的忍耐；让我们为自己积攒的，不是忿怒，而是救恩；不是死亡，而是生命；不是惩罚和报复，而是荣誉和冠冕。在这种情况下，无需支付运输捐献之物的运费；无需费力将其换成金钱。你若将它们交出，主亲自将它们运到天上；祂亲自使这交易为你更加有利。\par

此处无需找人来运送你所捐献的；只要你捐献，它立刻就升上去，不是为供养其他人作士兵，而是为你存留，获得大利。因为在这里，你所给出的一切，都不可能收回；但那里，你将带着极大的荣耀重新收到它们，并获致更大、更属神的收益。在这里，礼物是索债；在那里，是贷款，是有息的钱，是一笔债务。\par

不仅如此，天主还给了你债券。因为经上说：\BibleQuote{怜悯穷人的，就是借贷给上主。}\BibleRef{箴 19:17}祂还给了你定金和担保，而祂竟是天主！什么样的定金？就是今世的事物，可见的、属神的事物，以及未来事物的预尝。\par

那么，你为何迟延，为何退缩？你既已领受了这么多，又期待着这么多？\par

因为你所领受的是这些：祂亲自造了你的身体，祂亲自在你内放入了灵魂，祂唯独赐给你言语尊荣，超越地上万物，祂赐给你一切可见之物的使用，祂赏赐你认识祂的知识，祂为了你舍弃了祂的圣子，祂赐给你充满如此多恩宠的圣洗，祂赐给你圣桌，祂应许了天国，以及那不可言说的美善。\par

既然领受了如此多的美善，又将领受如此多，我再说一遍同样的话：你还在为那朽坏的财富斤斤计较，你还有什么借口呢？\par

但你全然顾念你的儿女？并为了他们而退缩？不，更当教导他们也赚取这样的收益。因为你若有一笔钱放贷取息，且有一个感恩的债户，你宁可一万次选择把债券留给你的孩子，而不是金子，好让他从中获得大笔收入，而不必四处奔走，寻找他人借贷。\par

如今，把这债券交给你的孩子，留下天主作他们的债户。你不卖掉土地留给儿女，而是留给他们，好让收入持续，使他们从那里获得更大的财富增长；但这债券比任何土地或收入都更丰产，结出如此多的果实，你却害怕留给他们？这是何等大的愚昧和疯狂！况且你明知，即使你留给他们，你自己也将再次带着它离开。\par

属神的事物就是如此；它们具有极大的慷慨。因此，让我们不要吝啬，不要对自己不仁慈和野蛮，而要从事那美好的交易；好让我们自己离去时带着它，也留给我们的儿女，并藉着我们的主耶稣基督的恩宠和爱人，达到未来的美善；愿光荣、威能、尊崇归于祂，与圣父及圣神，现在及永远，世世无尽。阿们。\par

\SourceRecord{Translated by George Prevost and revised by M.B. Riddle.；Nicene and Post-Nicene Fathers, First Series；Vol. 10.；Edited by Philip Schaff.；Buffalo, NY: Christian Literature Publishing Co.,；1888.；<http://www.newadvent.org/fathers/200166.htm>.}

% structure: p0001 source=h1 target=section reason=page-title

\section{玛窦福音讲道集 第67篇}

\BibleRef{玛 21:12-13}\par

\BibleQuote{耶稣进了圣殿，把一切在圣殿内作买卖的人赶出去，推倒了兑换银钱者的桌子和卖鸽者的凳子，对他们说：经上记载：\InnerQuote{我的殿宇应称为祈祷的殿}，你们竟使它成为贼窝了。}\par

\Person{若望}在他的福音开始也记载了这事，而这里（玛窦）是在末了。由此可知这事可能做了两次，且在不同的季节。\par

从时间和他们的回答都可以清楚看出。因为在那里（若望）他正好在逾越节来到，但这里（玛窦）则早得多。在那里犹太人说：\Quote{你给我们显什么神迹？}\BibleRef{若 2:18} 但这里他们保持沉默，虽然被责备，因为他在众人中备受赞叹。\par

这是对犹太人更重的指控，当他不仅一次，而是第二次这样做时，他们仍然做他们的买卖，并说他是天主的敌人，而他们本应从此事学到他对圣父的尊敬和他自己的大能。因为他确实也行了奇迹，他们看到他的话与他的行为一致。\par

但即使如此，他们也不被说服，反而\Quote{极其恼怒}，而他们听到先知大声呼喊，孩童们以超越年龄的方式宣告他。因此他也提出\Person{依撒意亚}作为控告者来反对他们，说：\BibleQuote{我的殿宇应称为祈祷的殿。}\BibleRef{依 56:7}\par

他不仅以这种方式显示他的权柄，还通过医治各种疾病。因为\BibleQuote{盲人和跛子来到他跟前，他治好了他们}，\BibleRef{玛 21:14} 这显示了他的能力和权柄。\par

但他们甚至这样也不被说服，反而在其余的神迹之外，听到孩童们宣告，就几乎气炸了，说：\Quote{你听见他们说什么吗？}而基督本应当对他们说：\Quote{你们没有听见这些说什么吗？}因为孩子们是向着他如同向天主一样歌唱。\par

那么他说了什么？既然他们反对明显的事，他以更责备的方式纠正他们，说：\Quote{你们没有念过：\BibleQuote{从婴孩和吃奶的口中，你完善了赞美}吗？}他又说：\Quote{从口中。} 因为所说的不是出于他们的理解，而是出于他的能力给他们尚未成熟的舌头以表达能力。\par

这也是外邦人呀呀学语，同时以理解和信心发出伟大事物的预表。\par

对宗徒来说，这也是一种不小的安慰。因为他们可能困惑，没有学问的他们如何能传扬福音，孩童们预先做到，并消除了他们所有的焦虑，教导他们，那使这些孩童歌唱的，也必赐给他们表达能力。\par

不仅如此，这奇迹还显示他是自然万物的创造者。那时孩童们，虽然年龄不成熟，却说出含义清晰且与天主一致的话，而成年人的话却充满疯狂和狂妄。因为邪恶的本性就是如此。\par

因为有很多事情激怒他们：从群众，从赶出买卖人，从奇迹，从孩童，他再次离开他们，给激情的爆发让出空间，不愿意开始教导，免得他们充满嫉妒，对他的话更加恼怒。\par

\Quote{早晨回城的时候，他饿了。}\BibleRef{玛 21:18} 早晨他怎么会饿？当他允许肉身时，肉身就显出其感觉。\Quote{他看见路旁一棵无花果树，就走到跟前，在树上什么也没找到，只有叶子。}\BibleRef{玛 21:19} 另一位福音作者说：\Quote{那时还不是\Emph{无花果}的季节。} 但如果不是季节，另一位福音作者为什么说：\Quote{他来到树前，或许能在上面找到果子。}由此可见，这属于他门徒的猜疑，他们那时还处于不太完备的状态。因为福音作者在许多地方记录了门徒的猜疑。\par

正如这是他们的猜疑，同样，认为它被诅咒是因为没有果子，也是猜疑。那么它为什么被诅咒？是为了门徒的缘故，好使他们能有信心。因为他在各处施恩，却不惩罚任何人；而他有必要给他们一个证明他也能惩罚的能力的示范，好使门徒和犹太人都学到：他能毁灭那些钉死他的人，却自愿顺从，不毁灭他们；他不愿意在人类身上显示这一点；却在植物上提供了他惩罚能力的证明。但当类似的事情发生在地方、植物或畜类身上时，不要好奇，也不要说：\Quote{如果还不是无花果的季节，无花果树怎么会被公正地枯萎呢？}因为说这话是最大的轻浮；而要注视奇迹，赞叹并荣耀行奇迹者。\par

因为关于被淹死的猪群，也有许多人说过这样的话，推敲公义的论证；但那里也不应留意，因为那些又是畜类，正如那是不结果实的植物一样。\par

那么为什么这个举动要具有这样的外表，以及这个诅咒的借口？正如我所说，这是门徒的猜疑。\par

但如果还不是时候，有些人徒然地说这里指的是法律。因为法律的果子是信德，那时正是这个果子的时节，而且它已经结了果子。因为他引用圣经说：\BibleQuote{因为田里的庄稼已经发白，可以收割了}\BibleRef{若 4:35}，又说：\BibleQuote{我派你们去收割你们没有劳苦的}\BibleRef{若 4:38}。\par

2. 因此，祂在此并非暗示这些事的任何一件，而是正如我所说的，祂显示了自己惩罚的权能，这藉著说\Quote{时候还没有到}而显明，表明祂是特意前往，并非为了饥饿，而是为了祂的门徒，他们确实极其惊奇，尽管已经行了许多更伟大的奇迹；但正如我所说，这事是奇特的，因为祂现在首次显出了祂报复的权能。因此，祂不在其他植物中，而在所有种植之物中最湿润的那一棵上行奇迹，为的是使奇迹由此显得更加伟大。\par

并且，为使你们知道，这是为他们而行的，为要训练他们满怀信心，请听祂随后说的话。祂说了什么？\Quote{你们若愿意相信，并因祈祷而满怀信心，也要行更大的事。}你们看到这一切都是为了他们的缘故吗？好使他们不因针对他们的阴谋而害怕和战栗。故此祂第二次又说了这话，为叫他们坚持祈祷和信德。\Quote{因为你们不仅能行这事，而且还能移山；并且，你们若在信德和祈祷上满怀信心，还能行更多的事。}\par

但骄傲自大的犹太人，想要打断祂的教导，来到祂跟前问说：\Quote{祢凭什么权柄做这些事？}\BibleRef{玛 21:23} 因为他们无法反对奇迹，便拿祂纠正殿中商贩的事来攻击祂。这事在若望福音中也出现他们前来询问，虽然措辞不同，但意图一样。因为在那里他们也说：\Quote{祢行这些事，给我们显什么征兆？}但祂在那里回答他们说：\Quote{你们拆毁这座殿，我三天之内要把它重建起来。}而在这里，祂却使他们陷入困境。由此显明，那时确实是奇迹的开端和序幕，而这里却是终结。\par

他们所说的是这样：你领受了教师的座席吗？你被祝圣为司祭了吗？以致你显示了这样的权威？这话是记载的。然而祂并未行任何傲慢之事，而是关切圣殿的秩序，但尽管如此，他们无话可说，却拿这事来反对祂。的确，当祂赶出他们时，他们因奇迹而不敢说什么；但当祂显明自己时，他们便指责祂。\par

那么，祂说了什么？祂没有直接回答他们，以表明他们若愿意看见祂的权柄，本是可以的；但祂反过来问他们说：\Quote{若翰的洗礼是从哪里来的？是从天上，还是从人来的？}\par

这推论是什么？确实是最重要的。因为若他们说是从天上来，祂会对他们说：\Quote{那么你们为什么不信他呢？}因为如果他们信了，就不会问这些事了。因为若翰曾论及祂说：\Quote{我连给祂解鞋带也不配}；又说：\Quote{看，天主的羔羊，除免世罪者}；并说：\Quote{这是天主子}；以及\Quote{那从天上来的，在万有之上}；\BibleRef{若 4:31} 还有\Quote{祂手里拿着簸箕，要彻底清理禾场。}\BibleRef{玛 3:12} 所以，如果他们信了他，就无任何阻碍使他们知道基督凭何权柄行这些事。\par

此后，因为他们狡猾地说：\Quote{我们不知道}，祂没有说\Quote{我也不知道}，而是说：\Quote{我也不告诉你们。}\BibleRef{玛 21:27} 因为若他们真是无知，理当受教导；但既然他们心存狡猾，祂有充分理由不回答他们。\par

他们为何不说那洗礼是从人来的？经上说：\Quote{他们害怕民众。}\BibleRef{玛 21:26} 你们看到悖逆的心吗？在每件事上他们都轻视天主，却为人的缘故而行一切。因为他们害怕这人是为民众的缘故，并非尊敬圣人，而是因着人；他们也不愿信从基督，是因着人；他们的一切灾祸皆由此而生。\par

此后，祂说：\Quote{你们以为怎样？一个人有两个儿子。他对第一个说：孩子，你今天到葡萄园去工作。他回答说：我不去；但后来后悔了，就去了。他又对第二个也说了同样的话。他回答说：主，我去；却没有去。这两个儿子中，哪一个遵行了父亲的意愿？}他们说：\Quote{第一个。}\par

祂又藉比喻定他们的罪，暗示他们不合理的顽梗，以及那些被他们完全定罪之人的顺服。因为这两个孩子说明了有关外邦人和犹太人的事。前者并未承诺服从，也未曾成为法律的听者，却在行为上显出了服从；后者说过：\Quote{凡上主所吩咐的，我们都要遵行，都要听从}\BibleRef{出 19:8}，却在行为上不服从。为此，我还要加上一点，好叫他们不以为法律会使他们受益，祂表明正是这件事定了他们的罪，如同保禄也说：\Quote{听法律的不是在天主面前为义人，行法律的才被成义。}为这目的，好叫他们甚至自我定罪，祂使审判由他们自己发出，正如祂在接下来的葡萄园比喻中所行的一样。\par

3. 为使这事成就，祂借着另一个人来考验控告。因为他们不愿直接承认，祂便藉比喻将他们引向祂所愿的。\par

但他们不明白祂的话，就作出了判断；之后祂便揭示隐藏的含义，说：\Quote{税吏和娼妓要在你们之前进入天国。因为若翰来到你们这里，显示正义之路，你们不信他；但税吏却信了他；你们见了以后，也没有后悔而相信他。}\BibleRef{玛 21:31}\par

因为若祂简单地说\Quote{娼妓要在你们之前进入}，这话在他们看来似乎冒犯；但现在，在他们的判断之后说出，就显得并不过于严厉。\par

因此祂也加上控诉。那么这是什么？\Quote{若翰来到}祂说，\Quote{你们这里}，不是到他们那里，不但如此，而且是\Quote{走在义的道路上}。\Quote{因为你们因此也不能指责他，说他是一个粗心大意、毫无益处的人；但他的生活无可指摘，他对你们的关心很大，而你们却不理会他。}\par

与此再加上另一项指控，就是税吏听从了；与此再加上又一项，就是\Quote{连在他们之后你们也没有做。因为你们本应在他们之前就做，但连在他们之后也不做，就丧失了一切借口}；于此，一方的赞美和另一方的指控都是无法言喻的。\Quote{他到你们这里来，你们不接受他；他没有到他们那里去，他们却接待他，而你们甚至连他们也不当作导师。}\par

看有多少事显明了那些人的称赞和这些人的指控。他到你们这里来，不是到他们那里。你们不相信，这没有触犯他们。他们相信了，这对你们无益。\par

但是\Quote{走在你们前面}这句话，并不是说这些人在跟随，而是说倘若他们愿意，就有希望。因为妒忌最能激发粗鄙之人。所以祂常说：\Quote{最先的将成为最后的，最后的将成为最先的。}因此祂引入了妓女和税吏，为要激起他们的妒忌。\par

因为这两个确实是主要的罪，由强烈的私欲所生，一个是性欲，另一个是贪财。祂指出，特别听从天主的法律就是相信若翰。因为妓女得以进入，不仅靠恩宠，也靠义德。因为她们并非继续作妓女而进入，而是服从了、相信了、被净化了、悔改了，如此才进入的。\par

你看见祂如何通过比喻，通过引入妓女，使祂的言论不那么刺耳而更具穿透力吗？因为祂没有立刻说：你们为什么不信若翰？而是更刺痛人的是，当祂提出了税吏和妓女之后，才加上这句话，按他们行为的次序，定了他们不可饶恕的罪，表明他们做事都是出于对人的恐惧和虚荣。因为他们怕被逐出会堂，就不承认基督；同样，他们不敢说若翰的坏话，这也不是出于敬畏，而是出于恐惧。这一切祂都用话语揭露，之后更严厉地打击，说道：\Quote{但你们，知道了以后，后来也不悔改，以致相信他。}\par

因为起初不选择善，固然是恶事，但甚至不能回转，则是更重的指控。因为这尤其使许多人变得邪恶，我见现在有些人因极端麻木不仁就是如此。\par

但不要让任何人如此；虽然一个人沉沦到邪恶的极致，也不要对向善的改变绝望。因为从邪恶的深渊中崛起是容易的事。\par

你们没有听说过那妓女吗？她在纵情声色上超过一切人，却在敬虔上闪耀超越众人。我不是指福音中的那个妓女，而是我们这一代的一个，来自腓尼基——那个最无法无天的城市。因为她曾是我们中的妓女，在舞台上享有最高的荣誉，她的名声四处传播，不仅在我们的城市，甚至远达基利家人和卡帕多细亚人。她毁了许多家产，败坏了许多孤儿；许多人还指控她行巫术，说她不仅以美貌，还以药物编织这样的罗网。这个妓女甚至一度赢得了皇后的兄弟，因为她的势力实在强大。\par

但忽然间，我不知道如何，或者说我确实知道，因为是由于这样心意转变、悔改，并将天主的恩宠引到自己身上，她鄙视了这一切，抛弃了魔鬼的伎俩，升到了天堂。\par

事实上，当她在舞台上时，没有什么比她更卑鄙的了；然而后来，她在极度节制方面超过了很多人，她穿上苦衣，所有时间都这样克苦自己。为了这个妇人，总督被煽动，士兵武装起来，但他们没有力量把她抢到舞台上去，也不能把她从接待她的贞女们手中带走。\par

这个女人被堪配那不可言传的奥迹，并表现出与（所赐的）恩宠相称的勤勉，如此结束了她的一生，通过恩宠洗尽了一切，并在领洗后显示了极大的节制。因为她甚至不允许那些曾经的爱慕者看一眼自己，当他们为此而来时，她把自己关闭起来，仿佛在监狱中度过了许多年。因此，\Quote{最后的将成为最先的，最先的将成为最后的；}因此我们在任何情况下都需要一个热切的灵魂，并且没有什么能阻止一个人成为伟大和可钦佩的：\par

所以，生活在恶习中的人，不要绝望；生活在美德中的人，不要昏睡。不要让这后者自信，因为妓女常常会超过他；也不要让前者绝望，因为他有可能超过那最先的。\par

你们听天主对耶路撒冷所说的话：\Quote{我说，在她行了这一切淫乱之后，你转向我，她却不回头。}\BibleRef{耶 3:7}当我们回归到天主切切的爱时，祂不记念先前的事。天主不像人，因为祂不拿过去责备我们，也不说：你为什么离开这么久？当我们悔改时；只让我们按应当的样式亲近祂。让我们切切地依附祂，将我们的心钉在敬畏祂上。\par

此类事不仅发生在新盟约之下，也发生在旧约之下。因为谁比默纳舍更恶？但他能平息天主的怒气。谁比撒罗满更有福？但当他昏睡时，他便跌倒了。倒不如说，我甚至能指明这两件事都发生在同一个人身上，就在这个人的父亲身上，因为同一个人在不同时期既好又坏。谁比犹达斯更有福？但他成了叛徒。谁比玛窦更不幸？但他成了福音作者。谁比保禄更坏？但他成了宗徒。谁比西满更令人羡慕？但他自己却成了所有人中最不幸的。\par

你会看到多少其他这样的变化，既在过去发生，也在每天发生？为此我说：让那靠杖之人不要绝望，让那在教会中的人不要自负。因为对后者说：\BibleQuote{那自以为站得稳的，要小心，免得跌倒。}\BibleRef{格前 10:12}；对前者说：\BibleQuote{跌倒的岂能不再起来吗？}\BibleRef{耶 8:4}，以及\Quote{抬起下垂的手，和无力的膝盖。}再者，对这些祂说：\Quote{警醒}；但对那些说：\BibleQuote{你这睡眠的，醒来吧，从死者中起来吧。}\BibleRef{弗 5:14}因为这些需要保存他们所拥有的，那些则需要成为他们所非是的；这些需要保持健康，那些则需要摆脱疾病，因为他们有病；但许多病人变得健康，许多健康人因懈怠而变得虚弱。\par

所以对这人祂说：\BibleQuote{看，你已经痊愈了，不要再犯罪，免得你遭遇更坏的事。}\BibleRef{若 5:14}；但对这些人说：\Quote{你愿意痊愈吗？起来，拿起你的床，回家去吧。}因为罪是一种可怕的、可怕的瘫痪，不，不仅是瘫痪，而且是更严重的事。因为这样的人不仅对善行无能，而且积极作恶。但尽管如此，即使你如此状态，并愿意稍加振作，一切恐怖就结束了。\par

即使你已如此\Quote{三十八年}，并且渴望痊愈，也没有人阻止你。基督现在也在，并说：\Quote{拿起你的床}，只愿你愿意振作，不要绝望。你没有一个人吗？但你有天主。没有人把你放进水池吗？但你有那使你不需要水池的。没有人把你放进去吗？但你有那命令你拿起床的。\par

你不能说：\BibleQuote{当我来到时，另一个人在我之前下去。}\BibleRef{若 5:7}因为如果你愿意下到泉源，没有人阻止你。恩宠不会被消耗，不会被耗尽；它是源源不断的泉源；因祂的丰满，我们灵魂和身体都得医治。那么让我们现在就来就它。因为辣哈布也曾是妓女，却得了救；盗贼是杀人犯，却成了天堂的公民；犹达斯与主同在却灭亡，盗贼在十字架上却成了门徒。这就是天主奇妙的作为。同样，贤士们证明了自己，税吏成了福音作者，亵渎者成了宗徒。\par

5. 看看这些事，永不绝望，但要永远自信，并振作自己。只要抓住通往那里的路，你就会快速前进。不要关闭门，不要封闭入口。今生的生命短暂，劳苦微小。即使它很大，也不该躲避。因为如果你不在悔改和德行这最光荣的劳苦中劳碌，你在世俗中必然以其他方式劳碌疲惫。但如果两者都有劳苦，为什么我们不选择那果实丰盛、赏报更大的呢？\par

然而这劳苦和那劳苦也不相同。因为在世俗追求中，有不断的危险、接连的损失、不确定的希望；有极大的奴役，以及财富、身体和灵魂的耗费；而果实的回报远低于我们的期望，如果它偶尔能生长的话。\par

因为世俗的劳苦并非处处都结出果实；不，即使它没有失败，甚至丰盛地产出果实，它持续的时间也很短暂。\par

因为当你年老，此后不再完美地感受到享受时，劳苦才带给你报酬。而劳苦是在身体强健时，果实和享受却是在年老衰弱时，那时时间已使感觉迟钝，即使没有迟钝，对结局的期待也不让我们感到快乐。\par

但在另一情形则不然：劳苦是在腐朽和必死的身体中，但冠冕却是不朽、不死、没有尽头的。劳苦既在先又短暂；但报酬既在后又无穷，此后你可以安心安息，不再期待任何不愉快的事。\par

因为你不再像这里一样担心变化或损失。那么，这些既不可靠、短暂、属世、未出现就消失、且以许多劳苦得来的福乐，是什么样的福乐呢？而哪些福乐能与那些不动摇、不衰老、无劳苦、甚至在争战之时就给你冠冕的福乐相比呢？\par

因为那轻视金钱的人，在这里就已经得了赏报，免于焦虑、竞争、诬告、阴谋、嫉妒。那节制、有序生活的人，甚至在离世之前，就得了冠冕，生活愉快，免于不雅、嘲笑、诬告的危险以及其他可怕的事。其余的一切德行也在这里就给我们回报。\par

因此，为了获得现在和将来的祝福，让我们逃避邪恶，选择德行。这样，我们既享受快乐，又获得将来的冠冕。愿天主恩赐我们众人达到这一切，因我们的主耶稣基督的恩宠和慈爱，愿光荣和权能归于祂，直到永远。阿们。\par

\SourceRecord{Translated by George Prevost and revised by M.B. Riddle.；Nicene and Post-Nicene Fathers, First Series；Vol. 10.；Edited by Philip Schaff.；Buffalo, NY: Christian Literature Publishing Co.,；1888.；<http://www.newadvent.org/fathers/200167.htm>.}

% structure: p0001 source=h1 target=section reason=page-title

\section{玛窦福音讲道集第六十八篇}

\Emph{\BibleRef{玛 21:33-44}}\par

\Emph{\BibleQuote{你们再听一个比喻。有一个家主，栽种了一个葡萄园，周围圈上篱笆，在里面挖了一个压酒池，盖了一座瞭望台，然后租给佃户，就往远方去了。到了收成的时节，他派仆人到佃户那里，去收取他的果子；佃户却拿住仆人，有的鞭打，有的杀死，有的用石头砸死。他又派其他的仆人，比先前的更多；他们照样对待了他们。最后，他派自己的儿子到他们那里去，说：\InnerQuote{他们必会尊敬我的儿子。}佃户看见儿子，就彼此说：\InnerQuote{这是继承人；来，我们杀他，得到他的产业。}于是他们捉住他，推出葡萄园外杀死了。那么，葡萄园的主人来到时，他要怎样处置那些佃户呢？他们对他说：\InnerQuote{他要凶恶地消灭那些恶人，把葡萄园租给别的佃户，就是那些在时节上把果子交给他的佃户。}耶稣对他们说：\InnerQuote{你们没有读过经上这句话吗？匠人弃而不用的石头，反而成了屋角的基石；} }}\par

祂借着这个比喻暗示了许多事：天主从起初就对他们施行的眷顾；他们从起初就有杀人的心性；凡是关乎谨慎照管他们的事，没有一样遗漏；即使先知被杀，祂也没有离弃他们，反而派遣了自己的圣子；新约和旧约的天主是同一的；祂的死将带来极大的祝福；他们因钉死祂的罪行要遭受极重的刑罚；外邦人的蒙召，犹太人的被弃。\par

因此，祂把这个比喻放在前一个比喻之后，为借此显示他们的罪责更大，且完全不可原谅。如何、以何种方式呢？就是虽然他们蒙受了如此多的眷顾，他们却比娼妓和税吏更坏，且坏得多。\par

也请注意祂极大的眷顾，和这些人极端的懈怠。因为凡属于佃户的事，祂自己都做了：周围圈上篱笆、栽种葡萄园、以及其他一切，祂只留下很少的事让他们去做：照管园中的一切，保存所交给他们的。没有一样被遗漏，一切都完成了；即便如此，他们也没有获益，尽管他们从祂那里享受了如此大的恩惠。因为当他们从埃及出来时，祂赐下了法律，设立了城邑，建造了圣殿，预备了祭台。\par

\Quote{就往远方去了；}这就是说，祂长久容忍他们，并非总是紧接着他们的罪过就施行惩罚；因为祂的往远方去，是指祂极大的忍耐。\par

而\Quote{祂派了祂的仆人，}即先知们，\Quote{去收取果子；}即他们的顺服，以及通过行为证明他们的顺服。但这些人甚至在此显出了他们的邪恶，不仅在接受如此大的眷顾后未能交出果子，这已经是懈怠的标志，而且还向那些来的人发怒。因为那些欠债不还的人，本不该愤慨，也不该发怒，而该恳求。但他们不仅愤慨，还双手沾满鲜血，自己本应受罚，却反而施加惩罚。\par

因此，祂又派了第二批和第三批人，既要显明这些人的邪恶，也要显明差遣者的慈爱。\par

为什么祂不立即派遣祂的儿子？为的是让他们因对待其他人的行为而自责，止息怒气，并在祂来临时尊敬祂。还有其他的原因，但眼下让我们继续看下文。可是，\Quote{他们必会尊敬，}是什么意思？这不是无知之语，千万不可这样想！而是要表明罪的严重性，且无可推诿。因为祂明知他们会杀祂，却还是派遣了。但祂说\Quote{他们必会尊敬，}是表明他们原本应当做的，即他们本该尊敬祂。因为别处祂也说：\Quote{也许他们会听从；}\BibleRef{则 2:5}这里祂也不是不知道，而是免得任何倔强的人会说，祂的预言迫使了他们的悖逆，因此祂用这样的措辞来说话，说：\Quote{他们是否会，}和\Quote{也许。}因为尽管他们曾对祂的仆人执拗，但他们对圣子的尊严仍当尊敬。\par

那么这些人做了什么呢？他们本该跑到祂面前，本该请求赦免他们的过犯，却反而更顽固地坚持先前的罪，继续增添他们的污秽，永远用后来的罪掩盖先前的过犯；正如祂自己所说的：\Quote{你们来充满你们祖宗的恶贯！}\BibleRef{玛 23:32}因为从一开始，先知就用这些话责备他们：\Quote{你们的手满是血；}\BibleRef{依 1:15}又：\Quote{他们以血接血；}\BibleRef{欧 4:2}又：\Quote{他们用血建造熙雍。}\BibleRef{米 3:10}\par

但他们没有学会节制，尽管他们首先领受了这条诫命：\Quote{不可杀人；}并因这条诫命被命令禁止无数其他事情，并通过多种方式被敦促遵守这条诫命。\par

尽管如此，他们却没有除去这一恶习；但当他们看见祂时，说了什么？\Quote{来，我们杀祂。}出于什么动机和原因？他们能提出什么指控，无论是大的还是小的？难道是因祂尊重你们，且身为天主子却为你们的缘故成为人，并行了无数奇迹？或是因祂赦免了你们的罪？或是因祂召叫你们进入天国？\par

但请注意，与他们的不虔敬相伴随的，是极大的愚昧，谋杀祂的理由充满了疯狂。因为经上说：\Quote{我们杀祂，产业就是我们的了。}\par

他们在哪里商议杀祂？\Quote{葡萄园外。}\par

2. 你看见祂甚至预言了自己将要被杀的地点吗？\Quote{他们把他扔出葡萄园，杀了他。}\par

路加确实说，祂宣告了这些人将要遭受什么；他们便说：\Quote{绝对不可}；祂又加上了圣经的证言。因为\BibleQuote{祂注视他们，说：\InnerQuote{那么经上写的这是什么？匠人所弃的石头，已成为屋角的基石；凡跌在这石头上的，必被摔碎。}} \BibleRef{路 20:17} 但玛窦说，是他们自己宣判了刑罚。但这并不矛盾。因为两件事都发生了：他们自己对自己宣判了刑罚；并且当他们意识到自己所说的话时，又加上了\Quote{绝对不可}；而祂又引用先知的话来反驳他们，说服他们这事一定会发生。\par

然而，即使如此，祂也没有明示外邦人，免得给他们留下把柄，只是隐晦地藉着说\Quote{祂要把葡萄园租给别的园户}来暗示。因此，祂用比喻说话，为的是让他们自己宣判刑罚，这在达味的事上也发生过，当他审判纳堂的比喻时。但请你留意，我恳求你，正是由此看出这判决是多么公正，因为将要受罚的人自己定了自己的罪。\par

然后，为使他们知道，不仅正义的本质要求这些事，而且圣神的恩宠从一开始就已预言了它们，天主也如此定断，祂便加上了一段预言，并以一种令他们羞愧的方式责备他们，说：\Quote{你们没有读过吗？匠人所弃的石头，已成为屋角的基石？这是上主所作的，在我们眼中神妙莫测。} 用这一切表明，他们因不信将被抛弃，而外邦人将被接进来。祂也藉着客纳罕妇人所暗示的正是这件事；藉着驴驹、百夫长和许多别的比喻也一样；现在也是如此。\par

为此，祂又加上了\Quote{这是上主所作的，在我们眼中神妙莫测}，预先宣告说，相信的外邦人和同样相信的犹太人将合而为一，尽管他们之间先前有如此大的区别。\par

然后，为使他们知道，所发生的一切没有一件与天主旨意相违，反而这事件极其可悦，超乎期望，令每个目睹者惊叹（因为奇迹实在无法言喻），祂又加上说：\Quote{这是上主所作的。} 祂以石头指自己，以工匠指犹太人的教师；正如厄则克耳也说：\Quote{他们修筑城墙，却用未泡透的灰抹墙。} \BibleRef{则 13:10} 但他们如何弃绝祂呢？他们说：\Quote{这个人不是从天主来的；} \BibleRef{若 9:16} \Quote{这个人迷惑百姓；} \BibleRef{若 7:12} 又说：\Quote{你是撒玛黎雅人，并且附有魔鬼。} \BibleRef{若 8:48}\par

然后，为使他们知道惩罚不仅限于被驱逐，祂又加上了刑罚，说：\Quote{凡跌在这石头上的，必被摔碎；但石头落在谁身上，就要把他压得粉碎。} 祂这里讲到两种毁灭方式：一种是由于绊倒和跌倒；这就是\Quote{凡跌在这石头上的}；另一种是由于被掳、灾祸和彻底毁灭，这也是祂清楚地预言的，说：\Quote{它要把他压得粉碎。} 藉着这些话，祂也隐晦地暗示了自己的复活。\par

依撒意亚先知说，祂责备葡萄园，但这里祂特别指控民众的首领。在那里祂说：\Quote{我应替我的葡萄园做什么，而没有做到呢？} \BibleRef{依 5:4} 还有其他地方说：\Quote{你们的祖先在我身上犯了什么不义？} \BibleRef{耶 2:5} 又说：\Quote{我的百姓，我对你做了什么？我在什么事上使你厌烦？} \BibleRef{米 6:3} 显示他们忘恩负义，当享受一切时，他们却以相反的行为回报；但这里祂表达得更有力。因为祂没有亲自辩护说：\Quote{我应做什么却没有做？} 而是让他们自己来判断，没有任何缺失，并审判自己。因为当他们说：\Quote{祂要凶恶地消灭那些恶人，并将葡萄园租给别的园户。} 他们所说的正是这个，以更大的力量宣告了自己的判决。\par

斯德望也正是以此责备他们，这最刺痛他们的是：他们一直享受了许多眷顾，却以相反的行为回报了恩主，这本身就是极大的记号：带来惩罚的原因不是惩罚者，而是受惩罚者自己。\par

这里也同样藉着比喻和预言表明了这一点。因为祂不仅满足于一个比喻，还加上了双重的预言：一个是达味的，另一个是祂自己的。\par

那么，他们听了这些事后应当做什么呢？他们难道不该朝拜、惊叹于那前后显示的慈爱眷顾吗？但如果这些事没有一样使他们变得更好，至少对惩罚的恐惧难道不该使他们更加节制吗？\par

但他们并没有如此；他们听了这些话以后做了什么？经上记载：\Quote{他们听了祂的比喻，明白是指着他们说的。他们就想要捉拿祂，但害怕群众，因为他们把祂当作先知。} 因为他们后来感觉到自己被暗示了。有时，当祂被捉拿时，祂从他们中间走过，不被看见；有时，当他们出现时，祂抑制了他们忙碌的急切；人们对此感到惊奇，说：\Quote{这不是耶稣吗？祂放胆讲话，他们竟一句话也不对祂说。} 但在这一事例中，因为他们被对群众的恐惧所约束，祂便满足了这一点，并没有像以前那样行奇迹从他们中间穿过并消失。因为祂并不希望以超人的方式做一切事，好使救恩计划被人相信。\par

但是他们，既不被群众所动，也不因所说的话而归正；他们不顾先知的见证，也不顾自己的判决，也不顾百姓的心意；权力之爱和虚荣之欲如此完全地蒙蔽了他们，连同对短暂之事的追求。\par

3. 因为没有什么比被这些腐朽之事牢牢钉住更能使人莽撞冲下悬崖，更能使人失去未来的事物。没有什么比将未来的事物置于一切之上更能确保人同时享受两者。因为基督说：\Quote{你们先寻求天主的国，这一切自会加给你们。} 确实，即使没有这个应许，我们也不应追求它们。但现在，在获得那些事物时，我们也可以得到这两样；即使如此，有些人仍不被说服，如同无感觉的石头，追逐快乐的影子。因为今世的事物中有什么是愉快的？有什么是令人愉悦的？因为今天我要更自由地向你们讲论；但请容忍，好使你们知道，这看似令人烦恼和厌倦的生活——我指的是隐修士们和那些被钉十字架者的生活——远比那看似安逸和精致的生活更甜美、更可羡慕。\par

对此，你们是见证人，你们在遭遇挫折和沮丧时，常常求死，并认为那些在山中、在洞穴中、没有结婚、过着超脱世俗生活的人是幸福的；你们这些从事技艺的、服兵役的、漫无目的生活的人，以及那些在剧院和乐池中度日的人。因为从这些事物中，尽管看似涌出无数快乐和欢笑的源泉，却也产生了无数更痛苦的利箭。\par

因为如果有人对那里跳舞的少女之一心生爱慕，他将承受比一万次行军、一万次离家旅行更痛苦的折磨，其悲惨境况超过任何被围困的城。\par

然而，目前不探究这些事，将它们留给那些被掳掠之人的良心，让我们来讲论普通人的生活，我们将发现这两种生活的差异如同港口与常被风浪击打的大海之间的差异一样巨大。\par

请从他们的退隐之处立即看出他们宁静的首要标志。因为他们逃离了市场、城市和人群的喧嚣，选择了山中的生活，那种与现世毫无共同之处、不受人类任何灾祸影响的生活——没有世俗的悲伤、没有忧愁、没有重大的挂虑、没有危险、没有阴谋、没有嫉妒、没有怨恨、没有无法无天的情欲，也没有任何这类事。\par

在这里，他们已经默想天国的事，与树林、山峦、泉水以及极大的宁静和孤独交谈，而首先是，与天主交谈。他们的斗室远离一切纷扰，他们的灵魂摆脱了每一情欲和疾病，纯净、轻盈，远胜于最纯净的空气。\par

他们的工作就是亚当起初、在他犯罪前的工作，那时他身披荣耀，自由地与天主交谈，住在充满大福乐的地方。因为他们在哪方面比他的境况更差呢？在他未悖逆之前，他被安置去耕种乐园。他没有世俗的挂虑吗？这些人也没有。他以纯洁的良心与天主交谈吗？这些人也如此；或者不如说，他们比他更有信心，因为他们借着圣神的倾注，享受了更大的恩宠。\par

你们本应通过目睹来领悟这些；但因为你们不愿意，反而在喧嚣和市场中度日，至少让我们用言语来教导你们，取他们生活方式的一部分（因为不可能穷尽他们的全部生活）。这些世界之光，在太阳升起时，或更确切地说，远在日出之前，就健康、清醒、警觉地从床上起身（因为没有忧愁、挂虑、头痛、劳苦、众多事务或任何其他这类事打扰他们，他们像天使一样生活在天上）；然后他们立刻从床上起来，愉快而欢喜，组成一个合唱队，以明亮的良心，同一声音，如同出于同一张口，向万有的天主歌唱赞美诗，尊崇他，感谢他的一切恩惠，无论是特殊的还是共同的。\par

所以，如果你们认为合适，让我们离开亚当，探究天使与这些在地上歌唱并说\Quote{光荣归于至高天主，在地上平安，良善归于人}的人群之间的差异。\par

他们的衣著与他们的刚毅相称。因为他们不像那些穿着拖地长袍、软弱无力、矫揉造作的人，而是像那些有福的天使、厄里亚、厄里叟、若翰、宗徒们一样穿着；他们的衣服是用山羊毛或骆驼毛制成的，有些人只用兽皮就够了，且这些衣服已穿了很久。\par

然后，在唱完那些诗歌后，他们屈膝，向那受他们赞美诗颂扬的天主祈求那些甚至连思想都不易达到的事。因为他们不祈求任何现世之物，他们对这些毫无挂虑，只求他们在独生子天主来审判生者死者时，能坦然面对那可怕的审判宝座，并且无人听见那可怕的声音说：\Quote{我不认识你们。} 以及他们能以纯洁的良心和许多善行度过这痛苦的生活，在顺风中驶过怒海。他们的父亲和领袖带领他们祈祷。\par

此后，当他们起身完成那些神圣而不断的祈祷后，太阳升起，他们各自去工作，从中为穷人收集大量的供给。\par

4. 如今那些投身于魔鬼的歌咏团、娼妓的歌曲、并坐在剧院里的人何在？我实在耻于提及他们，然而因你们的软弱，仍不得不如此。因为保禄也说：\Quote{你们从前怎样将肢体献给不洁作奴仆，现今也要照样将肢体献给正义作奴仆，归于圣洁。}\BibleRef{罗 6:19}\par

来吧，让我们也一同比较那由娼妓和舞台上卖淫的青年组成的团体，与这由蒙福者组成的同一团体，就快乐而言——大多数漫不经心的青年正是因此而被诱入陷阱。因为我们将发现这差别如同有人听见天使在天上唱出那全和谐之歌，而狗和猪在粪堆上嚎叫咕噜般巨大。因为藉着这些人的口，基督说话；藉着他们的舌头，魔鬼说话。\par

但是，当腮帮鼓起、琴弦拉到快要断裂时，笛声与无意义的噪声和令人不悦的表演结合在一起了吗？但在这里，圣神的恩宠发出声音，不是用笛子、竖琴或管子，而是用圣徒的嘴唇。\par

或者更确切地说，无论我们说什么，都不可能表达出其中的快乐，因为他们被钉在泥土和制砖上。因此我甚至希望抓住一个对这些事狂热的人，把他领到那里，向他展示那些圣徒的歌咏团，我就不再需要这些言语了。然而，尽管我们对泥泞的人说话，我们将试着，虽然通过言语，仍一点一点地，将他们从淤泥和沼泽中拉出来。因为在那里，听者立即接受了非法爱欲之火；因为就好像看到娼妓不足以点燃心灵，他们还用声音添加上祸害；但在这里，即使灵魂有此类事物，它也会立即放下。不仅是他们的声音，或是他们的面容，甚至他们的衣服也比这些更令观看者困惑。如果有一个粗俗而轻率的穷人，从这景象中他会千百次痛苦绝望地呼喊，并对自己说：\Quote{娼妓和卖淫的男童，厨师和鞋匠的子女，甚至常常是奴隶的子女，生活在如此精致中，而我一个自由人，由自由人所生，选择诚实的劳动，却连在梦中都想象不到这些事物；} 然后他离开时充满了不满之火。\par

但在隐修士的情况下，却没有这样的结果，反而完全相反。因为当他看到富人的子女和显贵祖先的后代，穿着连最穷的穷人都不如的衣服，并为此欢喜时，请想想他离开时会得到多大的对抗贫穷的安慰。如果他是个富人，他会清醒地回来，成为一个更好的人。同样在剧场里，当他们看到娼妓穿金戴银，而穷人会哀叹悲伤，看到自己的妻子毫无此类，富人则会因此景象轻视和蔑视他们的家庭伴侣。因为当娼妓向观看者展示服装、容貌、声音和步伐，一切都奢华时，他们离开时已被点燃，进入自己的家，从此成为俘虏。\par

因此产生侮辱和冒犯，因此产生敌意、战争和每日的死亡；因此对那些被俘虏的人，生活无法忍受，他们的家庭伴侣从此不再令人愉悦，孩子也不如以前那样被爱，家中一切变得颠倒，此后他们似乎被阳光本身搅乱。\par

但不会从这些歌咏团产生任何这样的不满，而是妻子会接纳她的丈夫安静而温顺，摆脱一切非法情欲，并会发现他比以前对她更温柔。那歌咏团产生这样的恶事，而这歌咏团产生善事，一个将羊变成狼，这一个将狼变成羔羊。但也许我们迄今尚未说到快乐之事。\par

还有什么比心灵不受困扰或悲伤，不沮丧不呻吟更令人愉快呢？然而，让我们进一步继续我们的论述，并考察这两种歌曲和表演的享受。我们会看到，一种确实持续到晚上，只要观众坐在剧场里，但之后却比任何刺更痛苦地折磨他；另一种则永远活跃在观看者的灵魂中。因为那些人的形象、地方的宜人、生活方式的甜美、规则的纯洁、以及那最美妙属灵歌曲的恩宠，他们已永远灌输给自己。至少那些持续享受那些港湾的人，从此就像逃避暴风雨一样，逃避人群的喧嚣。\par

但不仅在歌唱和祈祷时，而且在专注于书本时，他们也是观看者喜悦的景象。因为在他们结束歌咏之后，一位拿起\BibleBook{依撒意亚}{先知书}与他讨论，另一位与宗徒们交谈，另一位则查阅他人的劳苦，寻求关于天主、关于这个宇宙、关于可见之事、关于不可见之事、关于感官对象和理智对象、关于现世生命的卑贱和来世生命的伟大的智慧。\par

5. 他们以最上等的食物为养，不是摆上烹饪的兽肉，而是天主的圣言，超越蜂蜜和蜂巢，是一种奇妙的蜜，远胜过从前若翰在旷野所吃的。因为这种蜜不是野蜂采集、停在花朵上，也不是消化露水储存在蜂巢中，而是圣神的恩宠形成它，储存在圣徒的灵魂中，取代蜂巢、蜂房和管子，以便愿意的人可以安全地持续吃它。因此他们也在模仿这些蜜蜂，盘旋在那些圣书——圣经的蜂巢周围，从中收获极大的快乐。\par

如果你愿意了解他们的餐桌，就近前，你便会看到他们涌出如此温柔、甘饴、充满属神芬芳的事物。属神的口不能出污言秽语，没有愚妄的戏谑，没有粗暴之言，一切都是堪当天堂的。将那些在市场上爬行、痴迷世俗之人的口比作泥沟，不算为过；但这些人的嘴唇却如流淌着蜜、涌出清流的泉源。\par

但若有人因我将众人的口称为泥沟而不悦，让他知道，我这样说已是十分宽恕他们了。因为圣经并未用这样的尺度，而是用了更强烈的比喻。经上记载：\BibleQuote{蝮蛇的毒汁在他们的嘴唇下，他们的咽喉是敞开的坟墓。}但他们的却不是这样，而是充满芬芳。\par

他们在此世的境况如此，但将来如何，有什么言语能向我们陈明？有什么思想能领悟？是天使的份、不可言喻的福乐、难以述说的美善？\par

也许现在有些人心中发热，被激起对这等美好生活方式的渴望。但这有什么益处呢？当你还在此处时，你拥有这火；但当你离去时，你熄灭了火焰，这渴望便消失了。那么，如何才能避免呢？趁这渴望在你心中炽热时，你当奔赴那些天使那里，使它更燃旺。因为我们所讲的道理不能像亲眼目睹那样点起你的火。不要说：我要先与妻子说话，先料理我的事务。这迟延就是松懈的开端。请听，有人想向他家里的人告别，\BibleRef{列上 19:20}，先知不许他。我何必说告别呢？门徒想要埋葬他的父亲，\BibleRef{玛 8:21}，基督连这一点也不允许。试问，在你看来，有什么事情比父亲的葬礼更必需呢？但连这件事祂也不准许。\par

这是为什么呢？因为魔鬼正在一旁凶猛攻击，想要找些隐秘的缺口；即使只是小小的阻碍或迟延，它抓住了，就会造成巨大的松懈。因此有人劝诫说：\Quote{切勿一天又一天地拖延。} \BibleRef{德 5:7} 这样你便能在大多数事情上成功，你家里的事也会井井有条。因为经上说：\Quote{你们先该寻求天主的国，这一切自会加给你们。} \BibleRef{玛 6:33}\par

所以不要为你的利益思虑，把它们交托给天主。因为如果你为它们思虑，你只是以人的方式思虑；但如果天主供给，祂是以天主的方式供给。不要为它们思虑到忽略了更重要的事，因为那样祂就不会太多地供给它们了。因此，为让祂能完全供给它们，就把它们单单交给祂。因为如果你自己插手，放弃了属神的事，祂就不会为它们多多准备了。\par

为使这两件事都为你安排妥善，使你免于一切焦虑，你当紧守属神的事，轻视世俗的事；因为这样你便会拥有天地，获得将来的美善，借着我们的主耶稣基督的恩宠和慈爱，愿光荣与权能归于祂，至于无穷世。阿们。\par

\SourceRecord{Translated by George Prevost and revised by M.B. Riddle.；Nicene and Post-Nicene Fathers, First Series；Vol. 10.；Edited by Philip Schaff.；Buffalo, NY: Christian Literature Publishing Co.,；1888.；<http://www.newadvent.org/fathers/200168.htm>.}

% structure: p0001 source=h1 target=section reason=page-title

\section{玛窦福音释义第六十九篇}

\BibleRef{玛 22:1-14}\par

\BibleQuote{ \Emph{耶稣又开口用比喻对他们说：} \Emph{天国好比一位国王，为自己的儿子举办婚宴。} \Emph{他派仆人去叫那些被邀请的人来赴宴，但他们不愿意来。他又派别的仆人说：你们对被邀请的人说：我已经预备好了我的盛宴，我的公牛和肥畜已经宰了，一切都齐备了，请来赴婚宴吧！但他们却不理会，各自走了，一个到自己的田里，一个去做自己的生意；其余的竟拿住他的仆人，凌辱他们，把他们杀了。} }\par

你们在前一个比喻和这个比喻中，都看到了子与仆人之间的区别吗？你们同时看到了这两个比喻之间的极大相似之处，以及极大的不同吗？因为这也显示出天主的容忍，祂极大的眷顾，以及犹太人的忘恩负义。\par

但这个比喻还有比另一个比喻更多的东西。因为它预告了犹太人的被弃和外邦人的蒙召；并且它还同时指示了所要求的生活的严格，以及为疏忽者所设定的惩罚是多么巨大。\par

将这个比喻放在另一个比喻之后是恰当的。因为祂曾说过：\BibleQuote{它必赐给一个能结果实的外邦人。}，祂接着宣告那是怎样的外邦人；不仅如此，祂还再次揭示了祂对犹太人的眷顾，其程度难以言表。因为在那里，祂在受难之前召唤他们；但在这里，甚至在祂被杀之后，祂仍然敦促他们，努力争取他们。当他们理应受到最严厉的惩罚时，祂却既催促他们赴宴，又用最高的尊荣尊重他们。请注意，在那里，祂不是先召唤外邦人，而是犹太人，在这里也是如此。但是，正如在那里，当祂到来时，他们不愿意接待祂，反而杀了祂，祂就把葡萄园给了别人；同样，在这里，当他们不愿意来赴宴时，祂就召唤了别人。\par

那么，还有什么比他们更忘恩负义呢？被邀请参加婚宴，却匆忙离开？因为谁不会选择来参加婚宴呢，而且还是国王的婚宴，一位国王为他的儿子举办婚宴？\par

有人可能会问：为什么称它为婚宴？为使你们了解天主的温柔关怀，祂对我们的深情，事情的愉悦气氛，那里没有悲伤，也没有忧愁，一切都是属灵的喜乐。因此，若望也称祂为新郎，保禄也说：\BibleQuote{因为我曾把你们许配给一个丈夫。}\BibleRef{格后 11:2}，以及\BibleQuote{这是极大的奥迹，但我是指基督和教会说的。}\BibleRef{弗 5:32}\par

那么，为什么新娘不说许配给祂，而是许配给子呢？因为许配给子的，也就是许配给父。因为圣经中说这二者是等同的，并无区别，这是因为本质的同一性。\par

由此祂也宣告了复活。因为既然在前面祂提到了死亡，祂就表明，即使在死亡之后，仍有婚宴，仍有新郎。\par

但即便如此，这些人也没有变得更好或更温和，还有什么比这更坏的呢？因为这是第三项指控。第一项，他们杀害了先知；然后，是子；后来，即使在他们杀了祂之后，又被那被杀者本人邀请参加祂的婚宴，他们也不来，反而找借口，牛轭、田地、妻子。然而这些借口看似合理；但由此我们学到，即使妨碍我们的事情是必要的，也要将属灵的事置于一切之上。\par

并且祂不是突然召唤，而是提前很久。因为祂说：\BibleQuote{告诉那些被邀请的人。}，又说：\BibleQuote{叫那些被邀请的人来。}，这使对他们的指控加重。他们是在什么时候被邀请的呢？通过所有的先知；又通过若翰；因为他要把所有人都引向基督，说：\BibleQuote{祂必须兴盛，我必须衰微。}\BibleRef{若 3:30}；又通过子本人：\BibleQuote{凡劳苦和负重担的，你们都到我这里来，我要使你们安息。}，以及\BibleQuote{若是有人渴了，让他到我这里来喝。}\BibleRef{若 7:37}\par

但祂不仅用言语，也通过行动邀请他们：在祂升天后，通过伯多禄和与他一起的人。因为经上说：\BibleQuote{那在伯多禄身上运行，使他成为受割礼者的宗徒的，也在我身上运行，使我成为外邦人的宗徒。}\par

因为当他们看见子时，他们恼怒并杀死了祂，祂又通过祂的仆人邀请他们。他邀请他们做什么呢？是劳苦、辛劳和流汗吗？不，而是享乐。因为祂说：\BibleQuote{我的公牛和肥畜已经宰了。}请看祂的宴席多么丰盛，祂的慷慨多么巨大。\par

但就连这也没有使他们羞愧；祂越显示出容忍，他们就越是刚硬。因为他们不来，不是因为事务繁忙，而是由于\Quote{轻视。}\par

\Quote{那么，为什么一些人举出婚姻，另一些人举出牛轭呢？这些肯定是出于没有闲暇。}\par

绝非如此；因为当属灵的事召唤我们时，没有哪个事务繁忙具有必然性。\par

在我看来，他们似乎还利用这些借口，把这些作为他们疏忽的掩饰。不仅这一点是可悲的，即他们没有来，而且还有更暴力和疯狂的行为，即他们竟击打并凌辱那些来的人，还杀了他们；这比前者更严重。因为前者来到，是要收取产品和果实，却被杀害；而这些人，是要邀请他们参加那曾被他们杀害者的婚宴，却又被杀害。\par

有什么能与此疯狂相比呢？保禄也曾以此谴责他们，说：\BibleQuote{他们既杀了主，又杀了他们自己的先知，还迫害了我们。}\par

此外，免得他们说：\Quote{祂是反对天主的，所以我们不来。}，请听那些邀请他们的人所说的话：是父亲在举办婚宴，是祂在邀请他们。\par

那么在此之后祂又做了什么呢？既然他们不愿意来，甚至还杀了那些来到他们那里的人，祂便烧毁他们的城邑，派遣祂的军队并杀了他们。\par

祂说这些话，预先宣告了在\Person{韦斯帕芗}和\Person{提图斯}手下发生的事，以及他们因不信祂而激怒了圣父；无论如何，是圣父在报仇。\par

为此我补充说，并非在基督被杀害后立即被俘，而是在四十年之后，为要表明祂的忍耐，当时他们已杀害了\Person{斯德望}，处死了\Person{雅各伯}，并恶待了宗徒们。\par

你看到事件的真相及其迅速了吗？因为当\Person{若望}和许多其他与基督同在的人还在世时，这些事情就发生了，而听见这些话的人成了这些事件的见证人。\par

请看那难以言喻的眷顾。祂栽了一个葡萄园；祂做了一切事，并且完成了；当祂的仆人被处死后，祂又派了其他仆人；当那些人被杀后，祂派了儿子；而当祂被杀害时，祂邀请他们赴婚宴。他们不愿来。此后祂又派了其他仆人，他们也杀了这些人。\par

于是，祂便杀了他们，如同不治之症。因为他们患了不治之症，不仅由他们的行为得以证明，而且即使当娼妓和税吏信从时，他们也做这些事。因此，他们不仅因自己的罪行，也因别人能行善事而受谴责。\par

但若有人说，那时外邦人尚未被召，我是指在宗徒们挨打并遭受万般苦难之后，而是在复活之后（因为那时祂对他们说：\Quote{你们要去使万民成为门徒。}\BibleRef{玛 28:19}），我们会说，无论在受难之前还是之后，他们都首先向犹太人传道。因为在受难之前，祂对他们说：\Quote{你们往以色列家迷失的羊那里去。}\BibleRef{玛 10:6}；而在受难之后，祂不仅没有禁止，反而命令他们向犹太人传道。因为虽然祂说：\Quote{使万民成为门徒}，但当祂即将升天时，祂宣布他们应首先向这些人传道；因为祂说：\Quote{你们将领受圣神降临于你们身上的能力，并要在\Place{耶路撒冷}、全\Place{犹太}和直到地极为我作证人。}\BibleRef{宗 1:8}；\Person{保禄}又说：\Quote{那在\Person{伯多禄}身上为受割礼者之宗徒职分有效工作的，也在我身上为外邦人有效工作。}因此宗徒们也首先向犹太人传道，当他们长时间留在耶路撒冷，并被他们驱逐后，才分散到外邦人那里。\par

2. 你看，即使在此，祂也显出慷慨：\Quote{你们无论遇见谁，都请来赴婚宴。}因为在此之前，正如我所说，他们既向犹太人也向希腊人传道，大部分时间留在\Place{犹太}；但由于他们继续图谋害他们，请听\Person{保禄}解释这比喻并这样说：\Quote{天主的道本当先讲给你们听；但因你们弃绝自己，看为不配得永生，我们就转向外邦人。}\par

因此基督也说：\Quote{婚宴已经预备好了，但那些被请的人都不配。}\par

祂其实早已知道这点，但为不给他们留下无耻反驳的借口，虽然知道，祂仍先来到他们那里并派遣使者，既堵住他们的口，又教导我们要尽自己的本分，即使无人获益。\par

既然他们不配，祂说：\Quote{你们到大路上去，凡遇见的都请来。}无论平民还是被遗弃的人。因为祂曾说：\Quote{娼妓和税吏将继承天国}；以及\Quote{在先的将要在后，在后的将要在前。}祂表明这些事情是公义地发生的；这比任何事都更刺痛了犹太人，并比他们被推翻更严重地刺痛他们，因为他们看到外邦人被带入他们的特权之中，甚至得到比他们更大的特权。\par

然后，为了使甚至这些人也不仅依靠信德，祂也对他们论及要审判恶行；对尚未信从的人，论到因信德而来到祂面前；对已信从的人，论到对生活的关注。因为那衣服就是生命和实践。\par

然而，召叫是出于恩宠；那么为什么祂要严格追查呢？因为虽然被召和洗净是出于恩宠，但被召并穿上洁白的衣服后，继续保持洁净却是被召之人的勤奋。\par

被召并非出于功劳，而是出于恩宠。因此，理应回报恩宠，而不应在受此殊荣后显示如此大的邪恶。有人可能会说：\Quote{我没有享受像犹太人那么多的益处。}不，你享受了远大的益处。因为那些为他们预备了所有时间的东西，你立刻接受了，尽管你并不配。因此\Person{保禄}也说：\Quote{使外邦人因祂的怜悯而荣耀天主。}\BibleRef{罗 15:9}因为那些应归于他们的，你已领受。\par

因此，对于疏忽懈怠的人，也预备了重大的惩罚。因为正如他们因不来而侮辱，同样你也因以污秽的生活坐下而侮辱。因为穿着污秽的衣服进来，即是带着不纯洁的生命离开此世；因此他也哑口无言。\par

你是否看到，虽然事实如此明显，祂并没有立即惩罚，直到犯罪者自己宣判？因为无话可答，他定了自己的罪，然后被带去受那不可言喻的苦刑。\par

因为现在，当你听到\Quote{黑暗}时，不要以为他仅仅因被送入无光之处而受罚，而是那里\Quote{也有哀号和切齿}。\BibleRef{玛 22:13}祂这样说，表明那难以忍受的痛苦。\par

你们这些分享过奥秘、参赴了婚宴的人，要听：不要用污秽的行为装饰你们的灵魂。要听，你们是从哪里被召叫的。\newline{}\par

从大路上。你们是什么人？灵魂跛足和瘸腿的人——这比身体的残缺更为严重。要尊重那位召叫你们者的爱，不要让任何人继续穿着污秽的衣裳，但愿你们每个人都忙于你灵魂的衣裳。\newline{}\par

听吧，你们妇女；听吧，你们男人：我们不需要那些饰以黄金、\Emph{装饰我们外表的部分}的衣裳，而是需要那些装饰内在的衣裳。当我们拥有前者时，就很难穿上后者。不可能同时打扮灵魂和身体。不可能同时事奉\OriginalTerm{玛门}{\GreekText{μαμμωνάς}}又照我们应该的那样服从基督。\newline{}\par

因此，让我们摆脱这沉重的暴政。因为倘若有人用金帘子装饰你的房屋，却让你赤身裸体、衣衫褴褛地坐在那里，你也不会温顺地忍受。可是，看，你现在就是这样对待自己：用无数的帘子装饰你灵魂的房屋——我指的是身体——却让灵魂本身衣衫褴褛地坐着。难道你不知道君王应当比城池装饰得更华丽吗？因此，就如为城池准备了亚麻帷幔，为君王准备了紫袍和王冠，同样，你也应当用更简朴的衣服包裹身体，而用心智穿上紫袍、戴上冠冕，并把它放在高大显赫的战车上。因为你现在做的恰恰相反：用各种方式装饰城池，却让君王——心智——被粗暴的情欲捆绑拖曳。\newline{}\par

你难道不记得，你被邀请赴婚宴，就是天主的婚宴吗？你难道不想想，被邀请的灵魂应当如何穿着、佩戴金穗进入那些内室吗？\newline{}\par

3. 你们愿意我给你们看那些这样穿着、穿着婚宴礼服的人吗？\newline{}\par

请回想那些我最近向你们讲论的圣人：那些穿着粗毛衣服、住在旷野里的人。这些人首先是那婚宴礼服的穿着者；这从以下事实可以看出：即使你给他们多少件紫袍，他们也不愿接受；就如一个君王，若有人拿起乞丐的破衣劝他穿上，他会厌恶那衣服一样，那些圣人也厌恶他的紫袍。他们没有别的原因，只是因为他们知道自己衣服的美丽。因此，他们连那紫袍也像蜘蛛网一样弃绝。因为正是他们的苦衣教导了他们这些；事实上，他们远比那在位的君王更高贵、更荣耀。\newline{}\par

如果你能打开心智之门，看到他们的灵魂和里面的一切装饰，你一定会扑倒在地，因为你承受不住他们美的荣耀、那些衣服的光彩和良心的闪电般明亮。\newline{}\par

我们也可以讲述古代那些伟大而令人敬佩的人；但既然可见的榜样更能引导粗笨灵魂的人，所以我把你们甚至送到那些圣人的帐幕那里。因为他们没有任何忧伤，反而如同在天上扎营一样，远离今生的烦扰，安营对抗魔鬼；如同在唱诗班中一样，他们向他作战。因此我说，他们已搭好帐幕，逃离了城市、市场和房屋。因为作战的人不能坐在房屋里，而必须建造一种临时的住所，就像即将立刻转移一样，这样居住。所有那些圣人都如此，与我们相反。因为我们确实不是像在军营中生活，而是像在和平的城市中。\newline{}\par

因为在军营中，有谁曾立地基、为自己建造房屋，而他不久就要离开呢？没有这样的人；如果有人企图这样做，他会被作为叛徒处死。在军营中，有谁购买田地、为自己经营买卖呢？没有一个人，这是很合理的。\Quote{因为你们到这里来，}他们会说，\Quote{是来战斗，不是来做买卖；那么你为什么为这个你不久就要离开的地方烦恼呢？当我们离开这里回到我们的家乡时，再去做这些事情吧。}\newline{}\par

我现在也对你们说同样的话。当我们迁到上面的城市时，再做这些事；或者更确切地说，在那里你不再需要劳苦；之后君王会为你做一切。但在这里，只要挖一条壕沟、立一道栅栏就够了，没有必要建造房屋。\newline{}\par

听听斯基泰人的生活，他们住在四轮车上，正如人们所说，那是游牧部落的习惯。基督徒也应当这样生活：在世上漫游，与魔鬼作战，救出被他奴役的俘虏，并且不受一切世俗之物的束缚。\newline{}\par

人啊，你为什么预备房屋，好把你捆绑得更紧？你为什么埋藏财宝，邀请敌人来攻击你？你为什么用墙壁围困自己，为自己预备一座监狱？\newline{}\par

但如果这些事在你看来很难，就让我们去到那些人的帐幕那里，好让我们通过他们的行为学习其中的容易。因为他们搭起茅屋，如果必须离开这些，就像士兵一样离开，平静地离开他们的军营。因为他们也是这样安营的，或者甚至更美丽得多。\newline{}\par

因为看一个包含紧密相接的隐修士茅屋的旷野，比看士兵在军营里张起帆布、竖起长矛、在矛尖悬挂藏红花色的衣服、众多铜头的人、闪闪发光的盾牌中心凸起、全身披挂钢铁的人、匆忙建起的王庭、远远铲平的地面、以及人们吃喝吹笛，更令人愉快。因为连这个景象也不如我现在所讲的景象那样令人愉悦。\newline{}\par

因为如果我们进入旷野，看看基督士兵的帐幕，我们将看到的不是张起的帆布、矛尖，也不是构成王室营帐的金色衣服；而是仿佛有人在一片远比这大地更大、甚至无限的广阔土地上，张起了许多天空，那景象将是奇异而可畏的；在这里你也可以看到同样的情况。\newline{}\par

因为他们的居所在任何方面都不次于诸天；因为天使与他们同住，天使的主也与他们同住。如果他们来到\Person{亚伯拉罕}——一个有妻子、抚养儿女的人——面前，是因见他好客；那么，当他们找到更丰盛的德行，以及一个脱离了肉身、在肉身中无视肉身的人时，他们就更加逗留在他那里，举行相称的歌舞盛宴。因为在那里，桌子全然脱离了贪欲，充满了克己。\par

在他们中间没有血流成河，没有切割肉块，没有头脑沉重，没有精美烹调，也没有令人不快的肉味和刺鼻的烟雾，没有奔跑、喧闹、混乱和烦人的吵闹；只有面包和水——水来自纯洁的泉源，面包来自诚实的劳动。倘若他们有时想要更丰盛地宴饮，他们的丰盛在于果实，其中的愉悦远胜于王室的餐桌。那里没有恐惧或战栗；没有统治者控告，没有妻子激怒，没有孩子带来悲伤，没有无序的欢乐使人放荡，没有一群谄媚者使人自大；但餐桌是天使的餐桌，远离一切这样的纷扰。\par

他们的床榻只是身下的青草，如同基督在旷野设宴时一样。他们中的许多人甚至不在掩蔽之下，而以青天为屋顶，以月亮代替烛光，不需要油，也无需人照料；只有他们才配得上从上而来的光芒照耀。\par

4. 天使从天上看见这餐桌，就喜悦满足。因为如果为一个悔改的罪人而喜乐，那么为这么多效法他们的义人，他们岂不更加欢喜？那里没有主人和奴隶之分；全是仆人，全是自由人。不要以为这话是隐晦的箴言，因为他们彼此确实是仆人，又彼此是主人。\par

他们无需在夜幕降临之时忧伤，如同许多人那样，因回想白天的不幸而翻来覆去地焦虑不安。他们无需在晚餐后挂虑盗贼，关闭门户，加上门闩，也无需害怕其他灾祸，许多人因这些而恐惧，小心翼翼熄灭火烛，生怕一点火星引燃房屋。\par

他们的谈话同样充满了这种宁静。因为他们不谈我们所谈论的那些与我们无关的事：某人被任命为总督，某人已不再做总督；某人死了，另一个人继承了遗产，以及诸如此类。他们总是谈论将来的事，寻求智慧；如同住在另一个世界，如同已迁到天上，生活在那里，所以他们的全部谈话都是关于那边的事：关于\Person{亚伯拉罕}的怀抱，关于圣人们的冠冕，关于与基督的合唱；对眼前的事，他们既无记忆也无思想，正如我们不屑谈论蚂蚁在洞穴和裂缝中所做的事一样，他们也不屑谈论我们所做的事；而是谈论至高的君王，他们参与的战争，魔鬼的诡计，圣人们所成就的善行。\par

那么，与他们相比，我们与蚂蚁有何不同？因为我们关心身体的事，他们也关心；但但愿仅是这些而已：而现在甚至连更糟的事也关心。因为我们不仅像他们一样关心必需之物，还关心多余之物。因为那些昆虫从事的营生完全无可指责，而我们却追求一切贪欲，甚至不效法蚂蚁之道，反而效法豺狼之道，效法豹子之道，或者我们甚至比它们更坏。因为自然对它们的规定是那样觅食，但天主赐予我们言语和正义感，我们却变得比野兽更坏。\par

既然我们比野兽更坏，那些人却与天使同等，对世间之物如同客旅和寄居者；他们的一切都与我们不同：衣服、食物、房屋、鞋子、言语。如果有人听见他们谈话和我们谈话，他就会完全明白，他们确实是天上的公民，而我们连大地都不配。\par

因此，当任何身居官位的人来到他们中间，一切膨胀的骄傲就尽显虚妄。因为那里的劳动者和毫无世俗经验的人，坐在军队统帅的身旁——后者以自己的权威而自傲——他们坐在草地上，坐在低矮的垫子上。因为没有人颂扬他，没有人使他自大；但如同某人到金匠铺和玫瑰园，从金子和玫瑰中得到一些光彩；他们也是如此，从这些人的光辉中稍得照耀，就摆脱了昔日的傲慢。又如同某人登上高处，尽管极其矮小，也显得高大；同样，这些人来到他们崇高的心灵面前，仿佛和他们一样，只要留在那里；但当他们下去时，又从高处跌落，恢复卑下。\par

国王在他们中间算不得什么，总督算不得什么；如同我们看见孩童玩这些游戏时发笑，他们同样完全鄙弃那些在外招摇者的膨胀骄傲。这从以下事实可以证明：如果有人愿意给他们一个安稳的王国，他们绝不会接受；但他们却会接受，除非他们的思想专注于比王国更伟大的事，除非他们认为那事只是暂时的。\par

那么，我们岂不转投如此伟大的福乐？我们岂不来到这些天使面前，领取洁净的衣裳，参与这婚宴的礼仪；反而继续乞讨，与街上的穷人毫无差别，甚至更糟更苦？因为那些以恶道致富的人比这些穷人更糟；乞讨比掠夺好，因为前者有可原谅之处，后者则招致惩罚；乞丐丝毫不冒犯天主，而掠夺者却冒犯人和天主；他经历了掠夺的辛劳，而别人却常常收取一切好处。\par

那么，既然知道这些事，就让我们放下一切的贪婪，热切地渴望天上的事，\Quote{以猛力夺取天国。} \BibleRef{玛 11:12} 因为这是不可能的，绝不可能。凡懈怠的人，绝不能进入其中。\par

但愿天主使我们众人成为热忱警醒的人，得以达到那目标，藉着我们的主耶稣基督的恩宠和爱人之情，愿光荣与权能归于祂，于无穷世。阿们。\par

\SourceRecord{Translated by George Prevost and revised by M.B. Riddle.；Nicene and Post-Nicene Fathers, First Series；Vol. 10.；Edited by Philip Schaff.；Buffalo, NY: Christian Literature Publishing Co.,；1888.；<http://www.newadvent.org/fathers/200169.htm>.}

% structure: p0001 source=h1 target=section reason=page-title

\section{\WorkTitle{玛窦福音第七十篇讲道}}

\Emph{\BibleRef{玛 22:15}}\par

\Quote{然后法利塞人就去商议，怎样在谈话中圈套祂。}\par

然后。什么时候呢？当他们最应该感到痛悔，应当惊叹于祂对人的爱，应当害怕将来的事，并应当从过去的事相信将来的事时。因为所说的事确实在应验中大声宣告。我是指，税吏和娼妓相信了，先知和义人却被杀，从这些事他们不应反对自己的灭亡，反倒应当相信并清醒过来。\par

然而，即使如此，他们的恶行也没有停止，反而继续阵痛并推进。因为他们无法捉拿祂（因害怕群众），便采取另一种途径，意图使祂陷入危险，并使祂犯下危害国家的罪行。\par

因为\Quote{他们就派自己的门徒与黑落德党人一同到祂那里说：\InnerQuote{老师，我们知道你是诚实的，并且按真理教导天主的道，也不顾任何人，因为你不看人的情面。请告诉我们，你以为如何？纳税给凯撒，可以不可以？}}\BibleRef{玛 22:16}\par

因为他们那时已是纳贡者，国家已归于罗马人的统治。既然他们见\Person{特乌达}和\Person{犹达斯}\BibleRef{宗 5:36}及其同伙为此被处死，以为是准备叛乱，便想借这些话使祂也陷入这种嫌疑。因此他们派了自己的门徒和黑落德的兵士，在他们以为的两边都挖下悬崖，四面布下网罗，这样无论祂说什么，他们都能拿住话柄；如果祂的回答有利于黑落德党人，他们自己可以指责祂；如果有利于他们，别人就会控告祂。然而祂曾给过双德拉克马，但他们不知道。\par

他们确实指望在两方面都能拿住祂；但他们更希望祂说一些反对黑落德党人的话。所以他们派自己的门徒也去，用他们的在场催促祂，好将祂交给总督，当作篡权者。因为\BibleBook{路加}{福音}也暗示并表明这一点，说他们也当着群众的面问，为使见证更有效。\par

但结果完全相反；因为在更多的观众面前，他们反而证明了他们的愚蠢。\par

请看他们的谄媚和隐藏的诡计。\Quote{我们知道}，他们的话是，\Quote{您是诚实的。}但你们怎么曾说，\Quote{祂是骗子}，和\Quote{迷惑百姓}，和\Quote{有魔鬼}，和\Quote{不是从天主来的}？方才你们不是还图谋杀害祂吗？\par

但他们做一切事，无论诡计为他们暗示什么。因为既然不久之前他们刚愎自用地说过：\Quote{你凭什么权柄做这些事？}\BibleRef{玛 21:23}，而未能得到对问题的回答，他们便用谄媚来奉承祂，试图说服祂说一些反对既定法律、与现政权相对立的话。\par

因此他们也向祂证明真理，承认事实如此，然而并非出于正直的心，也不是自愿；并且加上说：\Quote{你不顾任何人。}请看他们多么明显地在敦促祂说那些话，使祂既得罪黑落德，又招致篡权者的嫌疑，仿佛对抗法律，以便他们能惩罚祂，当作煽动叛乱者和篡权者。因为说\Quote{你不顾任何人}和\Quote{你不看人的情面}，他们是在暗示黑落德和凯撒。\par

\Quote{请告诉我们，你以为如何？}现在你们尊敬祂，以祂为老师，而从前当祂讲论关于你们得救的事时，你们却多次藐视和侮辱祂。因此他们也成了同谋。\par

看他们的诡诈。他们不说：告诉我们，什么是好的，什么是合宜的，什么是合法的？而是说：\Quote{你以为如何？}他们如此只着眼于一个目的：背叛祂，使祂与统治者对立。而\BibleBook{马尔谷}{福音}表明这一点，更明白地揭露他们的刚愎和杀人的性情，肯定他们说过：\Quote{我们给凯撒纳税，还是不给呢？}\BibleRef{谷 12:15}因此他们充满怒气，孕育着针对祂的阴谋，却假装尊重。\par

那么祂说什么？\Quote{你们为什么试探我，假善人？}你们看祂如何以比平时更严厉的语气与他们说话？因为既然他们的邪恶现已完成并显露，祂便更深入，首先令他们难堪并哑口无言，揭露他们隐藏的念头，并向所有人显明他们怀着怎样的意图来见祂。\par

祂这样做，为要击退他们的邪恶，使他们不再试图同样的事而受害。然而他们的话充满许多尊重，因为他们既称祂为老师，又为祂的真理作证，说祂不徇情面；但祂既是天主，就不被这些事所欺骗。因此他们也应当猜到，这斥责不是出于猜测，而是祂知道他们隐秘思想的标记。\par

2. 祂却没有停留在斥责上，尽管足以揭露他们的意图并使他们的邪恶蒙羞；但祂不满足于此，而是以另一种方式封住他们的口；因为\Quote{拿一个税币给我看，}祂说。当他们拿给祂时，正如祂一贯所做的，祂通过他们的口引出决定，使他们判定这是合法的；这是一个清楚而明显的胜利。因此，当祂问时，不是由于无知而问，而是因为祂要使他们被自己的回答所束缚。因为当被问\Quote{这像是谁的？}他们回答\Quote{凯撒的}时，祂说：\Quote{那么，把凯撒的归给凯撒。}因为这不是给予而是归还，祂通过像和名号表明这一点。\par

为了不让他们说：\Quote{祢把我们从属于人。}祂接着说：\Quote{把属于天主的，归还于天主。}因为既能对人尽到他们的要求，也能向天主献上我们所亏欠于天主的。因此保禄也说：\BibleQuote{各人所当得的，就给他们；当纳粮的，就纳粮；当完税的，就完税；当敬畏的，就敬畏。}\BibleRef{罗 13:7}\par

但是你们，当你们听到：\Quote{把属于\Person{凯撒}的，归还给\Person{凯撒}。}要知道祂所说的，仅指那些无碍于虔敬之物；因为如果是这样的事，那就不再是凯撒的赋税，而是魔鬼的。\par

他们听了这些话，就哑口无言，并\Quote{惊奇}祂的智慧。难道他们不该相信，不该惊异吗？因为祂确实借着揭示他们心中的隐秘，给了他们祂天主性的证据，并以温和使他们无言。\par

那么怎样？他们信了吗？绝没有。反而他们\Quote{离开了祂，走了}；在他们之后，\Quote{撒杜塞人来到祂跟前。}\par

噢，愚昧！当别人都沉默时，这些人却发起攻击，而他们本应更退却才是。但鲁莽的本性就是这样：无耻、纠缠不休、企图做不可能的事。因此，圣史也惊异于他们的愚昧，指出了这一点，说：\Quote{那一天，他们来到祂那里。}哪一天？就是祂揭露了他们的诡诈并使她们蒙羞的那一天。这些人是谁？是犹太人中一个不同于法利塞人的派别，而且比他们更坏。这些人说：\BibleQuote{没有复活，也没有天使，也没有神灵。}\BibleRef{宗 23:8} 因为他们属于更粗俗的人，热衷于肉体的事。因为在犹太人中也有许多派别。因此保禄也说：\Quote{我是法利塞人，属于我们中最严格的派别。}\par

他们对复活并未直接说什么；却编造故事，捏造一个事例，依我看来，这根本从未发生过；他们想使祂为难，并渴望推翻两件事：即复活的存在，以及这种复活的存在。\par

这些人又装着有节制地攻击祂，说：\Quote{老师，梅瑟说：如果一个人死了，没有孩子，他的兄弟应娶他的妻子，为他兄弟立嗣。在我们这里有七个兄弟：第一个娶了妻，死了；没有子嗣，就把妻子留给了他的兄弟。同样，第二个、第三个，直到第七个。最后，那女人也死了。那么，在复活的时候，她应是七个中谁的妻子呢？}\par

请看祂像老师一样回答他们。他们虽是怀着狡诈到祂跟前来的，但他们的却问题出于无知。因此祂也没有对他们说：\Quote{你们这些假善人。}\par

此外，祂为了不责备他们，说：\Quote{为什么七个共有一个妻子？}他们却引用了梅瑟的权威；虽然，如我先前所说，至少在我看来，这是一个虚构。因为当第三个看到两个新郎都已死，他不会娶她；或即使第三个娶了，第四个或第五个也不会；如果连这些也娶了，那么第六个或第七个更不会接近那女人，反而会回避她。因为这就是犹太人的本性。如果现在许多人都有这种感觉，那么那时就更为如此；那时，即使没有这事，他们也常常避免这样结婚，即使法律强迫他们。无论如何，摩阿布女子卢德被推给比她亲属更远的人；塔玛尔也是如此被迫用计谋从她丈夫的亲属中得子。\par

他们为什么不编造两三个，而是七个？为了更充分地嘲讽他们所认为的复活。因此他们又说：\Quote{他们都曾拥有过她。}好像要把祂逼入困境。\par

那么基督怎么说？祂回答两者，不是针对言辞，而是针对意图，并每次揭示他们心中的隐秘；有时揭露他们，有时则把反驳他们的问题留到他们的良心。无论如何，请看祂在这里如何证明这两点：即会有复活，并且复活不是他们所怀疑的那种。\par

祂说什么？\BibleQuote{你们错了，既不明白圣经，也不明了天主的能力。}\BibleRef{玛 22:29} 因为既然他们好像知道圣经，就提出梅瑟和法律，祂指出这个问题是那些非常无知于圣经的人的问题。因为他们试探祂也是由此而来，即因他们无知于圣经，也没有按应有的方式认识天主的能力。\par

祂说：\Quote{那么这算什么奇迹呢？}祂说，\Quote{你们试探我，我对你们尚且陌生，而你们居然连天主的能力都不认识，你们对那能力曾有如此多经验，却既未从常识也未从圣经中认识它；}如果连常识也告诉我们，在天主一切都可能。祂首先回答所问的问题。因为这就是他们不信复活的原因：他们认为事物的秩序如此。祂医治原因，然后也是症状（因为疾病也从那里来），并显示了复活的方式。祂说：\Quote{因为在复活时，他们也不娶也不嫁，而是如同天主的天使在天堂中一样。}但路加说：\BibleQuote{作为天主的儿子。}\BibleRef{路 20:36}\par

如果那时他们不嫁娶，问题便是枉然。但他们之所以如同天使，不是因为他们不嫁娶，而是因为他们如同天使，因此他们不嫁娶。借此祂也消除了许多其他困难，这一切保禄用一句话暗示说：\BibleQuote{因为这世界的格局正在消逝。}\BibleRef{格前 7:31}\par

祂藉这些话表明了复活是何等伟大的事；不仅如此，祂还证明了复活的存在。其实在此之前所讲的也已同时表明了这一点，但祂仍在此之上，藉现在所说的话，将自己的话语又加以补充。因为祂并非只在他们发问时停下，而是在他们思想上止住。所以，当人们并非出于极大狡诈，而是出于无知询问时，祂甚至额外施教；但若只是出于恶意，祂连问题也不回答。\par

祂又藉梅瑟堵住他们的口，因为他们也引用了梅瑟。祂说：\Quote{论到死人复活，你们没有读过：\BibleQuote{我是亚巴郎的天主，以撒的天主，雅各伯的天主}吗？祂不是死人的天主，而是活人的天主。}祂的意思不是说他们不存在、已完全消灭、不再复活。因为祂没有说\Quote{我曾是}，而是说\Quote{我是}；那是指存在着、活着的人。就如亚当虽然在他吃树的那天还活着，却在判决中死去；同样，这些人虽然死了，却在复活的许诺中活着。\par

那么，祂为何在别处说：\BibleQuote{为叫祂做死人和活人的主？} \BibleRef{罗 14:9}但这与那并不矛盾。因为这里祂是指死人，而这些死人也要活着。此外，\Quote{我是亚巴郎的天主}与\BibleQuote{为叫祂做死人和活人的主}是不同的事。祂还知道另一种死亡，关于这死亡祂说：\BibleQuote{任凭死人去埋葬他们的死人。} \BibleRef{玛 8:22}\par

\BibleQuote{群众听了这话，就惊奇祂的教训。} \BibleRef{玛 22:33}然而连这时撒杜塞人也没有信服；他们败退而去，而无偏袒的群众则得了益处。\par

既然复活是这样的，来吧，让我们做一切事，好能在那里获得首位的荣誉。但若你们愿意，让我在复活之前就在这里向你们显示一些追求并收获这些福份的人，再次求助于旷野。因为我将再次进入同一篇讲道，因为我看见你们听得更愉快。\par

那么，让我们今天也观看属灵的营帐，观看他们那毫无恐惧的喜乐。因为他们不像士兵那样带着长矛扎营——这是我最近结束讲道时所谈到的——也没有盾牌和胸甲；但你会看见他们赤手空拳，却成就了甚至武装者也做不到的事。\par

若你能够观察，来，向我伸出手，让我们两人一同前去这场战争，看看他们的战阵。因为这些人每天也作战，击杀仇敌，征服一切阴谋陷害我们的情欲；你会看见这些情欲被摔在地上，甚至无力挣扎，而是以实际行动证实宗徒的这句话：\BibleQuote{属于基督耶稣的人，已把肉身同邪情和私欲钉在十字架上了。} \BibleRef{迦 5:24}\par

你看见那里有一大堆死人，被圣神的剑所杀吗？因此那地方没有醉酒和饕餮。他们的餐桌和其上树立的凯旋纪念碑就是证明。因为醉酒和饕餮被击倒，被饮水所击溃，尽管这种恶习是千变万化、多头怪物。好比传说中的斯库拉和海德拉，在醉酒中你可以看见许多头：一边是淫乱生长，另一边是忿怒；一边是懒惰，另一边是无法无天的情欲；但这一切都被除去了。然而，所有其他军队，即使赢得千万次战争，却被这些情欲所俘虏；而武器、长矛，无论还有什么，都不能抵挡这些方阵；那些巨人、英雄、无数勇敢行为的人，你会发现他们不被捆绑却被睡眠和醉酒所捆绑，没有杀戮或伤口却像受伤者一样躺卧，甚至更糟。因为至少那些人还会挣扎；但这些情欲根本不会，而是立刻投降。\par

你看见这军队更伟大、更令人赞叹吗？因为它仅凭意志就消灭了那些战胜他人的仇敌。他们如此削弱了所有邪恶之母，以致她不能再搅扰他们；首领被推翻，头部被除去，身体其余部分也就静止了。\par

我们可以看见每一个住在那里的隐修士都在赢得这场胜利。因为不像我们这些世俗战争，如果某个敌人被一个人打中，他就不再对另一个人构成威胁，因为他已经被打倒一次；但每个人都必须打击这头怪物；没有击中并打倒她的人，必定被她困扰。\par

你看见光荣的胜利了吗？因为全世界军队集合起来也不能建立的凯旋纪念碑，每个隐修士都能建立起来；所有从醉酒军队中混杂一起受伤的事物——疯狂的话语、妄念、不悦人的傲慢——都被击倒。他们效法他们自己的主，经上惊叹祂说：\BibleQuote{祂在路旁要喝溪水，因此祂必抬起头来。}\par

你们还想看到另一堆死人吗？让我们看看来自奢侈生活的欲望，那些被酱料制作者、厨师、筵席供应者、甜点师所珍爱的欲望。我实在羞于谈论所有一切；不过，我要讲讲来自法西斯的鸟，用各种东西混合的汤：湿的、干的菜肴，以及关于这些东西制定的法律。因为仿佛是规划一座城、调度军队一样，这些人也制定法律，规定某物第一，某物第二；有人先端上填满鱼的内馅的炭火烤鸟；另有人用其他材料开始这些违法筵席；在这些东西上——关于质量、顺序、数量——有很多竞争；他们以那些本应羞愧之事为荣：一些人说他们花了半天时间，一些人说花了整天，一些人说还加了黑夜。看哪，可怜的人啊，你肚腹的尺寸，并为你那无度的热忱而羞愧吧！\par

但在那些天使中并没有这样的事；所有这些欲望也都是死的。因为他们的饮食不是为了饱足和奢靡享乐，而是为了需要。那里没有捕鸟者，没有渔夫，只有饼和水。但这一切混乱、骚扰和动荡都从那里消除了，无论是从家中还是从身体上，那里是极大的避风港，而在这些人中却是极大的风暴。\par

现在，在思想上剖开那些以这些为食之人的肚子，你就会看到大量的残渣、污秽的通道和粉饰的坟墓。\BibleRef{玛 23:27}\par

但此后的事，我甚至羞于启齿：令人不快的嗳气、呕吐、从下往上的排泄物。\par

但你去看看，这些欲望在那里也是死的，还有那些由此而产生的更强烈的情欲；我指的是不洁的情欲。因为你会看到这些也都完全被推翻，连同它们的马匹和驮兽。因为不洁行为的驮兽、武器和马匹，就是不洁的言语。但你会看到这样的马匹和骑手一起倒毙，他们的武器也都被丢弃；而在这里却恰恰相反，灵魂被击倒而死。不仅在这些圣人的餐宴上，胜利是光荣的，在其他方面也是如此：在金钱、荣耀、嫉妒和灵魂的一切疾病上。\par

这岂不是比那一个更强盛，这宴席也更好吗？谁能否认呢？没有人，即使他们自己中的人，即使他非常疯狂。因为这一餐引导我们升天，那一餐将我们拖入地狱；这一餐是魔鬼摆设的，那一餐是基督摆设的；因为在这里，奢华和不节制制定法律，在那里，克己和清醒；在这里基督临在，在那里魔鬼。因为哪里有醉酒，魔鬼就在那里；哪里有污言秽语，哪里有暴饮暴食，那里就有魔鬼的合唱团。那个富人就有这样的宴席，因此他连一滴水也得不到。\BibleRef{路 16:24}\par

但这些圣人没有这样的宴席，他们已经在实践天使的生活方式。他们不娶不嫁，也不过度睡眠，也不奢侈享乐，除了少数事情外，他们甚至像是无身的。\BibleRef{玛 22:30}\par

那么有谁能像那个在用餐时立起凯旋纪念碑的人那样轻易地战胜敌人呢？因此先知也说：\Quote{在我敌人面前，你为我摆设了宴席。}\BibleRef{咏 23:5}关于这宴席，重复这句神谕是不错的。因为没有什么比混乱的情欲、奢华、醉酒以及由此产生的邪恶更能麻烦灵魂了；那些有经验的人深知这一点。\par

如果你也了解这宴席从何而来，那宴席从何而来，那么你就会清楚地看到两者的区别。这宴席来自何处？来自无数的眼泪，来自被欺骗的寡妇，来自被掠夺的孤儿；但那宴席来自诚实的劳动。这宴席好比一位美丽而端庄的妇人，不需要外在的装饰，其美来自天生；而那宴席好比一位丑陋而可憎的娼妓，涂了许多脂粉，却无法掩饰她的畸形，越是靠近，就越被揭穿。因为这宴席，当它接近赴宴者时，就更加显出其丑陋。因为，我告诉你，不要只看赴宴者入席之时，也要看他们离去之时，那时你就会看到它的丑陋。因为那宴席，如同自由人，不容许赴宴者说任何可耻的话；而这宴席，如同娼妓，受羞辱，不容许说出任何体面的话。这宴席寻求赴宴者的益处，那宴席寻求伤害。一个不允许得罪天主，另一个不允许我们不冒犯祂。\par

因此，让我们到那些人那里去。从那里我们将学到我们被多少绳索捆绑。从那里我们将学会为自己摆设一桌充满无数祝福、最甜蜜、无需花费、免除忧虑、免于嫉妒和一切疾病、充满美好希望并有许多凯旋的宴席。那里没有灵魂的骚动，没有悲伤，没有愤怒；一切都是平静，一切都是平安。\par

不要告诉我富人家中仆人的静默，而是赴宴者的喧哗；我不是指他们彼此之间的喧哗（因为这也值得嘲笑），而是指内心的喧哗，那灵魂中的喧哗，那给他们带来巨大俘虏的喧哗：思想的骚动、冰雹、黑暗、风暴，一切都混乱混杂，如同夜战。但在隐修士的帐篷中没有这样的事；那里是极大的平静，极大的安宁。那宴席之后是如同死亡的睡眠，而这宴席之后是清醒和警醒；那宴席导致惩罚，这宴席导致天国和不朽的赏报。\par

让我们效法这宴席，好使我们也能享受它的果实；愿天主赐予我们众人得以达到这境界，因我们的主耶稣基督的恩宠和慈爱，愿光荣与权能归于祂，直到永远。阿们。\par

\SourceRecord{Translated by George Prevost and revised by M.B. Riddle.；Nicene and Post-Nicene Fathers, First Series；Vol. 10.；Edited by Philip Schaff.；Buffalo, NY: Christian Literature Publishing Co.,；1888.；<http://www.newadvent.org/fathers/200170.htm>.}

% structure: p0001 source=h1 target=section reason=page-title

\section{\WorkTitle{玛窦福音讲道第七十一篇}}

\Emph{\BibleRef{玛 22:34-36}}\par

\Quote{但法利塞人听见祂使撒杜塞人闭口无言，就聚集在一起；其中有一个法学士，试探祂，问说：师傅，律法中哪一条诫命是最大的？}\par

圣史再次说明他们本应保持沉默的原因，也借此指出他们的放肆。如何？以何种方式？因为当那些人被驳得哑口无言时，这些人又来攻击祂。他们本应为此沉默，却竭力继续先前的努力，推出那位法学士，不是出于学习的意愿，而是试探祂，问说：\Quote{最大的诫命是什么？}\par

因为最大的诫命本是\BibleQuote{你当爱主、你的天主}，他们以为祂会给他们某种把柄，仿佛祂要修改这诫命，以显示祂自己也是天主，于是提出这问题。基督如何回答？祂指出他们之所以如此，是由于没有爱德，由于嫉妒焦心，由于被嫉妒所困，祂说：\BibleQuote{你当爱主、你的天主。这是第一条且最大的诫命。第二条也与此相似：你当爱邻人如你自己。} \BibleRef{玛 22:37}\par

但为何说\Quote{相似}？因为这条为那一条铺路，又由那一条得以坚固；\BibleQuote{凡作恶的，便恨光，也不来就光} \BibleRef{若 3:20} 再者，\Quote{愚妄人心里说：没有天主。} 由此导致什么？\Quote{他们败坏，行为可憎。} 再者，\BibleQuote{贪爱钱财是万恶的根源；有人贪求它，就离弃了信德} \BibleRef{弟前 6:10} 并且，\Quote{那爱我的，必遵守我的诫命。}\par

但祂的诫命以及它们的总纲是：\BibleQuote{你当爱主、你的天主，并爱邻人如你自己。} 因此，如果爱天主就是爱邻人——\BibleQuote{若你爱我}，祂说，\BibleQuote{伯多禄，喂养我的羊群} \BibleRef{若 21:16} 而爱邻人又能促成遵守诫命，所以祂有理地说：\BibleQuote{全律法和先知都系于这两条} \BibleRef{玛 22:40}\par

所以，祂先前所做的，此时也照样做。我的意思是，当被问及复活的方式时，祂教导了复活，指示他们超出所问的；同样，这里被问及第一条诫命时，祂也陈述了第二条，这一条并不比那一条逊色（因为虽然位居第二，却与它相似），藉此暗示他们，这问题出自仇恨。\BibleQuote{因为爱德不嫉妒} \BibleRef{格前 13:4} 借此，祂显示自己既顺服律法，也顺服先知。\par

但为何玛窦说他问话是试探祂，而马尔谷却相反？\BibleQuote{因为耶稣见他回答得明智，就对他说：你离天主的国不远了} \BibleRef{谷 12:34}\par

他们并非互相矛盾，而是完全一致。因为他起初确实出于试探而问，但后来因回答而得益处，受到称赞。因为祂并非一开始就称赞他，而是当他说了\Quote{爱邻人胜过全燔祭}之后，才说\Quote{你离天国不远了}；因为他看轻了低俗的事物，把握了德行的首要原则。的确，所有那些事物（如安息日等）都是为了这个目的。\par

即使这样，祂的称赞也不算完全，仍有所欠缺。因为祂说\Quote{你离得不远}，表明他仍有所欠缺，以便他追求那所缺的。\par

但如果当他说\Quote{天主是唯一的，除祂以外没有别的神}时，祂称赞了他，不必惊奇；由此也可见，祂如何按照来见祂的人的意见来回答。因为虽然人们对基督说了无数与祂的荣耀不相称的话，但他们至少不敢说祂根本不是天主。那么，为何祂赞扬那个说除了圣父以外没有别的神的人呢？\par

并非将祂自己排除于天主之外——千万不要这样想——而是因为披露祂天主性的时机尚未成熟，祂容他停留在先前的教理中，并称赞他深谙古训，使他适合于新约的教理，而这新约教理祂将在适当的时候带来。\par

此外，这句话\Quote{天主是唯一的，除祂以外没有别的神}，无论在旧约还是各处，都不是为了否定圣子，而是为了与偶像区分。所以，当祂称赞这样说话的人时，祂是出于这种心态称赞他。\par

然后，既然祂回答了，祂也反过来问：\Quote{你们对基督有何看法？祂是谁的儿子？他们对祂说：达味的儿子。}\par

请看，在这么多奇迹、这么多神迹、这么多问题、这么多言语和行为中与圣父一致的表现之后；在赞扬了那个说天主是唯一的人之后，祂提出这个问题，免得他们能说祂确实行了奇迹，却是律法的敌对者和天主的仇敌。\par

因此，在如此多事之后，祂问这些问题，暗中引导他们承认祂也是天主。祂先问门徒别人说什么，然后问他们自己；但对这些人却不然；因为他们肯定会说祂是骗子、恶人，因为祂毫不畏惧地说话。所以为此缘故，祂询问这些人自己的意见。\par

因为祂现在即将走向受难，祂提出那明确宣告祂是主的预言；并且不是要显得祂无缘无故这样做，也不是以此为宗旨，而是出于合理的缘由。\par

因为祂先问他们，既然他们没有回答关于祂的真理（因为他们说祂只是一个凡人），为了推翻他们的错误意见，祂就这样引入达味来宣告祂的天主性。因为他们确实以为祂只是一个凡人，所以他们也说：\Quote{达味之子}；但祂为纠正这一点，就带来先知证明祂是主，祂儿子身份的纯粹性，以及祂与圣父同等的尊荣。\par

并且祂并不止于此，而是为了促使他们畏惧，祂接着还说了以下的话：\Quote{直到我把祢的仇敌放在祢的脚下}；好让他们至少借此被争取过来。\par

为了不让他们说，祂是出于奉承才这样称呼祂，并且这是一个人的判断，请看祂说什么：\Quote{那么，达味怎么在圣神里称祂为主？}请看祂是多么谦卑地引出关于祂自己的言论和判断。首先，祂说过：\Quote{你们怎么想？祂是谁的儿子？}以问题方式引导他们回答。然后既然他们说：\Quote{达味之子}，祂没有说：\Quote{然而达味这样说这些话}，而是再次以问题的秩序说：\Quote{那么，达味怎么在圣神里称祂为主？}为的是不让这些话冒犯他们。因此祂也没有说：你们怎么看我，而是怎么看基督。为此，宗徒们也谦卑地推理说：\Quote{我们可坦然对圣祖达味说，他既死了，也埋葬了。}\BibleRef{宗 2:29}\par

祂自己也同样为此原因，以问题和推论的方式引入教义，说：\Quote{那么，达味怎么在圣神里称祂为主，说：主对我主说：祢坐在我右边，等我使祢的仇敌作祢的脚凳；}\BibleRef{玛 22:44} 又说：\Quote{达味既称祂为主，祂怎么又是他的儿子呢？}\BibleRef{玛 22:45} 这并不是否认祂是他的儿子，绝无此意；因为若如此，祂就不会为此责备伯多禄了，而是要纠正他们隐秘的想法。所以当祂说：\Quote{祂怎么是他的儿子？}时，祂的意思是：不像你们所说的那样。因为他们说，祂只是儿子，而不是主。这是在见证之后，然后谦卑地说：\Quote{达味既称祂为主，祂怎么又是他的儿子呢？}\par

然而，虽然他们听了这些，却什么也没有回答，因为他们不想学习任何必要的事。因此祂自己补充说：\Quote{祂是他们的主。}或者不如说，连这件事祂也不是没有支持而说的，而是带着先知，因为祂极其被他们不信任，且在他们中被恶名传扬。对此我们应当特别留意，如果祂说了任何卑微、顺从的话，不要被冒犯，因为原因就在这里，还有许多其他事情，祂以迁就的方式与他们交谈。\par

因此，现在祂也以问答的方式传授教义；但祂即使这样也隐约暗示了祂的尊威。因为被称为犹太人的主，不如被称为达味的主那么重要。\par

但是，我请你注意，这也是多么合时宜。因为当祂说：\Quote{只有一位主}时，然后祂就说自己是主，并借着预言证明这一点，不再只凭祂的作为。并且祂显示圣父亲自为祂向敌人复仇，因为祂说：\Quote{直到我使祢的仇敌作祢的脚凳}，并且借此显示生祂的那位对祂自己伟大的同心同德和尊荣。祂与他们的辩论以此崇高伟大的结局结束，足以封住他们的口。\par

因为从那时起他们就沉默了，不是自愿的，而是因为他们无话可说；他们受到如此致命的打击，以至于再也不敢尝试同样的事了。因为经上说：\Quote{从那天起，没有人敢再问祂什么。}\BibleRef{玛 22:46}\par

这对群众来说不是小的益处。因此，祂从此将祂的话转向他们，驱走了狼，击退了他们的阴谋。\par

因为那些人一无所获，被虚荣心所俘虏，落入了这可怕的情欲。因为这情欲是可怕的，而且多头，有些人为了虚荣而追求权力，有些人追求财富，有些人追求力量。但按顺序进行，它也延伸到施舍、禁食、祈祷和教导，这怪物的头很多。\par

然而，为那些其他事物而虚荣并不奇怪；但为禁食和祈祷而虚荣，这就奇怪而可悲了。\par

但为了让我们不只责备，来吧，让我们讲述避免它的方法。我们将先准备与谁对敌？是与那些为钱财虚荣的人，还是为衣着、为权力地位、为学问、为技艺、为容貌、为美丽、为装饰、为残忍、为人道和施舍、为邪恶、为死亡、为死后虚荣的人？因为正如我所说，这情欲有许多环节，且超越我们的生命。因为据说，有一个人死了，为了被钦佩，他命令做这样那样的事；因此某人贫穷，某人富有。\par

可悲的是，它甚至由相反的事物构成。\par

那么，我们将站立在谁面前，先列阵谁呢？因为同一个论述不足以对付所有。你们愿意先针对那些为施舍虚荣的人吗？\par

至少在我看来很好；因为我极爱这件事，且看到它被玷污，虚荣心像谄媚的乳母对待某个王室少女一样密谋反对它。因为她确实喂养她，却是为了耻辱和损害，出卖她并命令她藐视父亲；却打扮自己去取悦不洁和无耻的人；给她穿上一件陌生人所希望的衣服，可耻而不尊贵，而不是父亲所希望的。\par

來吧，現在讓我們瞄準這些；並且讓施捨大量地為向眾人展示而做。的確，首先虛榮領她離開了她的圣父的內室。然而，她的圣父連左手都不要知道（\BibleRef{瑪 6:3}），她卻把她顯示給奴僕和那些連認識她都未嘗的俗人。\par

你看見一個娼妓和拉皮條的人，把她投入愚昧男人的愛中，好讓他們如何要求，她就如何自處？你想看它如何使這樣的靈魂不僅成為娼妓，而且還變成瘋癲嗎？\par

那麼請留意她的心志。因為當她放開天堂，追逐逃奴和卑賤的奴僕，穿過街道和小巷追趕那些恨她的人、那些醜陋畸形的人、那些甚至不願看她的人、那些當她對他們燃燒愛火時卻恨她的人時，還有什麼比這更瘋癲的呢？因為眾人恨惡誰，也不如恨那些想要得到他們所能給予榮耀的人。至少他們對這些人提出無數指控，結果就像有人把國王的童貞女兒從御座上帶下來，強迫她賣身給那些厭惡她的角鬥士一樣。這些人，你越追逐他們，他們就越避開你；但天主，你若尋求從祂而來的榮耀，祂就越發吸引你到祂自己身邊，讚許你，賜給你巨大的賞報。\par

但若你樂意用另一種方式辨識其禍害，當你為了展示和炫耀而施捨時，想想那時臨於你的悲傷是多麼大，沮喪是多麼持久，而基督的聲音在你耳中響起：\Quote{你們已經失去了你們全部的賞報}（\BibleRef{瑪 6:1}）。因為在任何事上，虛榮確實是壞事，但在行善方面尤其如此，因為它是最殘酷的，拿他人的災難作秀，幾乎就是責備那些貧窮的人。因為若提到自己的善行就是責備，那麼你想，把它們甚至傳揚給許多其他人，又是什麼呢？\par

那麼我們將如何逃脫這危險？如果我們學會如何施捨，如果我們明白我們應當追求誰的好評。因為告訴我，誰擁有施捨的技巧？顯然，是天主，祂顯明了這事物，祂最知道它，並且無限量地實踐它。那麼怎樣？如果你在學習成為摔跤手，你仰望誰？或者你在摔跤學校裡向誰展示你的行為，是向賣菜賣魚的，還是向教練？然而他們人數眾多，而他只有一個。那麼，如果他讚賞你，而別人嘲笑你，你不會與他一同嘲笑他們嗎？\par

怎樣？如果你在學習拳擊，你不也會同樣地仰望那位知道如何教導此事的人嗎？如果你在練習演說，你不會接受修辭學教師的讚揚，而輕視其餘的人嗎？\par

那麼，在其他技藝上只仰望教師，而在此處卻反其道而行之，豈不是荒謬嗎？雖然損失並不相等。因為在那裡，你若按照群眾的意見摔跤，而不按照教師的意見，損失是在摔跤上；但在這裡卻是在永生上。你已變得相似天主，在施捨上；那麼你也要在不高調宣揚上相似祂。因為連祂醫治時也說，不要告訴任何人。\par

但你渴望在人們中間被稱為仁慈嗎？這有什麼益處？益處一無所有，損失卻無限。因為這些你喚來作見證的人，成了你天上寶藏的盜賊；或者更確切地說，不是他們，而是我們自己，毀壞了自己的產業，散盡了我們積攢在上面的東西。\par

哦，新的禍患！這奇怪的激情。哪裡蛀蟲不侵蝕，賊不挖洞，虛榮卻散盡。牠是那裡寶藏的蛀蟲；是我們天上財富的賊；牠偷走了那不能朽壞的財富；牠損壞並敗壞了一切。因為魔鬼看見那地方對賊和蟲子以及其它陰謀是無法攻破的，牠就藉著虛榮偷走了財富。\par

但你渴望榮耀嗎？那由領受者自己給予的，那來自我們恩慈的天主的，還不夠嗎？你卻還渴望從人而來的？當心，免得你遭受相反的後果，免得有人責備你不是在施恩，而是在作秀，尋求榮譽，拿他人的災難作秀。\par

因為施恩確實是一件奥迹。所以關上門，好叫沒有人看見那不虔敬展示的事。因為我們的奥迹，首要的就是顯示天主的仁慈和慈愛。祂按照祂的大仁慈，憐憫了我們這些悖逆的人。\par

第一個祈禱也充滿了仁慈，當我們為附魔者懇求時；第二個祈禱又是為其他在補贖中尋求大量仁慈的人；第三個祈禱是為我們自己，這祈禱提出百姓中無辜的兒童，為我們向天主祈求仁慈。因為既然我們為罪責備自己，我們就為那些犯了大罪、應受責備的人自己呼喊；但為我們自己，兒童呼喊；因為天國是為那些單純如兒童的人保留的。因為這圖像表明，那些像這些兒童一樣謙卑單純的人，正是最能以他們的祈禱拯救有罪之人的。\par

但那奥迹本身，充滿了多少仁慈、多少愛人之心，领洗者知道。\par

那麼，當你按你的能力向一個人施恩時，關上門，讓你施恩的對象只看見它；但如果可能，甚至連他也不要看見。但若你敞開門，你就是不敬地暴露了你的奥迹。\par

要想到，你所尋求其讚揚的那個人，連他自己也會責備你；如果他是朋友，他會在心裡責備你；如果是敵人，他甚至會向別人嘲笑你。你將遭受與你所願相反的事。因為你確實希望他稱你為仁慈的人；但他不會這樣稱你，而會稱你為虛榮的人、討好人的，以及其他比這些更嚴重的稱呼。\par

但若你将它隐藏起来，祂必以与此完全相反的名称来称呼你：仁慈的、良善的。因为天主不容许它被隐藏；但若你将它藏起，别人反而会将它显扬出来，而那钦佩会更大，收益也更丰富。所以，甚至为这受光荣的目的本身，炫耀对我们是有害的；因为我们最急于追求的那件事，正是在这一点上，这事对我们最为不利。因为我们非但得不到仁慈的名声，反而得到相反的名声，而且除此之外，我们所受的损失也很大。\par

因此，为了每一个动机，让我们戒绝此事，只将我们的爱寄托于天主的赞美。因为如此，我们既在此世获得尊荣，也享受永福，藉着我们的主耶稣基督的恩宠和爱人之心，愿光荣与权能归于祂，至于无穷之世。阿们。\par

\SourceRecord{Translated by George Prevost and revised by M.B. Riddle.；Nicene and Post-Nicene Fathers, First Series；Vol. 10.；Edited by Philip Schaff.；Buffalo, NY: Christian Literature Publishing Co.,；1888.；<http://www.newadvent.org/fathers/200171.htm>.}

% structure: p0001 source=h1 target=section reason=page-title

\section{玛窦福音讲道集 第七十二篇}

\Emph{\BibleRef{玛 23:1-3}}\par

\Quote{\Emph{那时，\Person{耶稣}对群众和祂的门徒说：经师和法利塞人坐在\Person{梅瑟}的座位上，因此，凡他们所吩咐你们的，你们都要遵守；}\Emph{但不要效法他们的行为。}}\par

那时。是什么时候？当祂说了这些话，当祂堵住了他们的口；当祂使他们不敢再试探祂；当祂表明他们的状况已不可救药。\par

既然祂提到了\Quote{主}和\Quote{我的主}，祂又回到法律。然而法律并未说这类事，而是说：\BibleQuote{上主你的天主是唯一的主}\BibleRef{申 6:4}。但圣经称整部\WorkTitle{旧约}为法律。\par

但祂说这些话，是为了藉一切显示祂与生祂的那位完全一致。因为如果祂是反对的，祂就会说关于法律相反的话；但现在祂命令对法律表示如此大的尊敬，以致即使教导它的人已败坏，祂仍命令他们持守它。\par

但此处祂正在论述他们的生活与道德，因为这是他们不信的主要原因——堕落的生命及爱慕虚荣。因此，为了改正祂的听众，祂首先强调最有助于救恩的事——不可轻视我们的教师，也不可反抗我们的司铎，祂以极其殷切的态度命令此事。但祂不仅命令，而且亲自实践。因为虽然他们败坏，祂并未废黜他们的尊位；这既加重了他们的定罪，也为祂的门徒不留下不顺服的借口。\par

我的意思是，免得有人会说：\Quote{因为我的老师不好，所以我变得更加松懈}，祂甚至消除了这个借口。尽管他们是坏人，祂仍如此牢固地确立他们的权威，以致在如此严厉的指控之后还说：\BibleQuote{凡他们所吩咐你们的，你们都要做。}因为他们所说的不是自己的话，而是天主的，是祂藉\Person{梅瑟}设立的法律。且看祂对\Person{梅瑟}表示了多么大的尊敬，再次显示祂与\WorkTitle{旧约}的一致；因为甚至藉此祂也使他们受人尊敬。因为祂说：\BibleQuote{他们坐在}，祂说，\BibleQuote{梅瑟的座位上。}因为祂不能凭着他们的生活使他们可信，祂便藉着向他们开放的理据——从他们的座位和他们的继承来自他——来做到这一点。但当你听到\Quote{所有}，不要理解为所有的法律，例如关于食物的规定、关于祭献的规定等；因为祂怎能对祂先前已废除的这些东西如此说呢？但祂指的是所有纠正道德原则、修正性情、与\WorkTitle{新约}法律一致，且不再使他们处于法律轭下的事物。\par

那么，为什么祂给予这些事神性权威，不是从恩宠的法律，而是从\Person{梅瑟}？因为还不是时候，在十字架之前，这些事不能清楚地被宣告。\par

但在我看来，除了已经说的，祂说这些话还为了另一个目的。因为祂即将要指控他们，免得在愚昧人看来祂似乎是为了这个权威而心存记恨，或出于仇恨而做这些事，祂首先消除了这个想法，并澄清自己不被怀疑，然后才开始指控。祂为什么定他们的罪，并长篇大论地攻击他们？为了警戒群众，使他们不陷入同样的罪。因为劝阻不像指出那些已犯过的人；正如推荐正确的事不像提出那些做得好的人。所以祂也预先说：\BibleQuote{不要效法他们的行为。}因为免得他们因听他们的话而认为也应该效法他们，祂就用这个纠正的方法，并将他们表面的尊严变成对他们的控诉。因为还有什么比一个老师更可怜，如果他的学生的保存就在于不注意他的生活？所以，他们表面的尊严是对他们的最沉重的指控，当他们被证明过着这样的生活，以至于效法者被毁灭。\par

因此，祂也抨击对他们的指控，但不仅为此，而且为表明他们先前的不信——他们不信，以及随后他们敢行的钉十字架——都不是对那位被钉和不信者的指控，而是对他们的邪僻的指控。\par

但看祂从哪里开始，以及祂从哪里加重对他们的指责。因为祂说：\BibleQuote{他们说却不行。}因为每个人在违反法律时都值得责备，但尤其担当教导权威的人，他应受双倍、三倍的定罪。原因之一，因为他违反法律；其二，因为他本应改正别人，自己却跛脚，因着他的尊位而应受双倍惩罚；第三，因为他甚至更加败坏别人，因为以教师的身份犯下这样的过犯。\par

和这些一起，祂还提到另一个对他们的指控：他们对自己管辖的人严厉。\par

\BibleQuote{他们把沉重难担的担子捆起来，放在人的肩上，自己却连一个指头也不肯动。}祂在这里提到双重邪恶：他们对属下要求极大且严格的生活，毫无包容，而对自己却允许极大自由；真正善牧应该相反：在自己身上，他要做不宽容和严厉的审判者，但在属下的问题上，要温和并乐于体谅；这些人的行为却相反。\par

2. 因为所有只在言语上实践克制的人都是这样的：不宽容，难以承受，如同没有体验过行动的困难。而这一点本身也不是小错，并且以不寻常的方式加重了先前的指控。\par

但是，我请求你注意，祂如何加重这项指控。因为祂没有说\Quote{他们不能}，而是说\Quote{他们不愿意}。祂没有说\Quote{承担}，而是说\Quote{用指头移动}，也就是说，连接近它们、触摸它们都不肯。\par

然而他们在什么事上热心而积极呢？在禁止的事上。因为祂说：\BibleQuote{他们所行的一切事，都是为叫人看见。}\BibleRef{玛 23:5} 祂说这些话，是指责他们贪图虚荣，而虚荣正是他们灭亡的原因。因为前面的事是冷酷和懈怠的标志，而这些事则是疯狂贪求荣耀的标志。虚荣使他们远离天主，使他们在别的观众面前争胜，并导致他们毁灭。无论一个人有哪种观众，既然他致力于取悦他们，他所表现出的竞争也是如此。在贵族中摔跤的人，他所进行的竞赛也是这样的；但在冷漠和懈怠的人中间，他自己也变得更加松懈。例如：有人有一个喜欢嘲笑的观众吗？他自己也变成引起嘲笑的人，以取悦观众；另一个人有一个严肃认真、操练自律的观众吗？他努力使自己成为那样的人，因为赞美他的人有那样的性情。\par

但请看，这里也是加重了指控。因为他们不是一些事这样做，另一些事那样做，而是全然如此。\par

然后，在责备他们贪图虚荣之后，祂指出他们甚至不是因重大而必要的事而自夸（因为他们也没有这些，反而缺乏善行），而是因无价值、无意义的事，以及那些明明证明他们卑贱的事：经匣和衣边。祂说：\BibleQuote{他们把自己的经匣放宽，把衣边加长。}\BibleRef{玛 23:5}\par

这些经匣和衣边是什么呢？因为犹太人不断忘记天主的恩惠，天主命令将祂奇妙的作为写在小小的牌子上，并挂在他们的手上（因此祂也说：\Quote{它们要在你们眼前固定不动}），他们称之为\OriginalTerm{经匣}{phylacteries}；就像我们现在的许多妇女把福音书挂在脖子上。为了通过另一件事提醒他们，正如许多人常做的，用一块亚麻布或线缠在手指上，因为他们容易忘记，天主命令他们如同孩子一般去做：\Quote{在衣服的边穗上缝一条蓝色的带子，垂在脚边，使他们看见并记起命令。}这些被称为\OriginalTerm{衣边}{borders}。\par

在这些事上，他们很殷勤，把牌子的条带加宽，把衣边加长；这是最极端虚浮的标志。因为你为什么自夸并把这些做得宽大呢？这难道是你们的善行吗？如果你们没有从中获得好处，这对你们有什么益处呢？因为天主所寻求的不是将这些加宽加大，而是我们记住祂的恩惠。但如果连施舍和祈祷，虽然劳苦，且是我们的善行，我们也不可追求虚荣，那么你，犹太人，怎能以这些最足以证明你懈怠的事自夸呢？\par

但不仅在这些事上，也在其他小事上，他们受到这种病的困扰。\par

因为祂说：\Quote{他们喜爱筵席上的首座，会堂里的高位，市场上的问候，以及被人称为\InnerQuote{辣彼}。}因为这些事，虽有人可能认为微小，却是大祸的根源。这些事曾倾覆了城市和教会。\par

现在甚至使我哭泣的是，当我听到首座、问候，并想到因此给天主的教会带来了多少祸害——这些我现在不必向你们公开；其实，凡是年长的人也无需从我们这里学习这些事。\par

但是，我请求你注意，虚荣心何等盛行：他们甚至在会堂里——他们进入那里本为教导别人——也被命令不可贪图虚荣。\par

因为在宴席上有这种感觉，即使程度很深，似乎也不是那么可怕的事；虽然即使在宴席上，教师也应受到尊敬，不仅是在教会里，而是一切的场所。如同一个人无论出现在哪里，都显然与野兽有所区别；同样，教师无论在说话还是沉默，在用餐还是做任何事时，都应当通过他的步态、容貌、服饰以及一切方面显得与众不同。但他们在各方面都成为笑柄，在一切事上羞辱自己，致力于追求他们本该逃避的事。经上说：他们喜爱这些；但如果喜爱它们本身已是可责备的，那么去做它们、追求它们又是多大的恶呢？\par

3. 其他事祂只是指控他们，认为这些事小而琐碎，似乎门徒们在这些事上完全不需要纠正；但是，那造成一切祸患的原因——野心和对教师席位的强夺，祂才提出来并殷勤地纠正，热切地、郑重地责备他们。\par

因为祂说什么？\BibleQuote{你们不要被称为\InnerQuote{辣彼}，因为你们的师傅只有一位，你们众人都是弟兄；}\BibleRef{玛 23:8} 并且一人不比另一人多什么，因为没有人靠自己知道什么。因此保禄也说：\BibleQuote{保禄算什么？阿颇罗算什么？不过是仆役。}\BibleRef{格前 3:5} 他没有说师傅。又说：\BibleQuote{不要称人为父，}\BibleRef{玛 23:9} 这并非不可称父，而是叫他们知道谁才是至高的父。因为正如师傅不主要是师傅，父也不主要是父。因为祂是万有的根源——包括师傅和父亲的根源。\par

祂又接着说：\Quote{也不要被称为导师，因为你们的导师只有一位，就是基督。}祂没有说：\Quote{我}。如先前祂说：\Quote{你们对基督有何看法？}而没有说：\Quote{对我}，此处亦然。\par

但我在此愿意问一问那些屡次将\OriginalTerm{独一}{One}一词单用于圣父、以排斥独生者的人：圣父是导师吗？人人都可宣告，无人能否认。然而祂说：\Quote{你们的导师只有一位，就是基督。}因为正如基督被称为唯一的导师，并不排除圣父为导师；同样，圣父被称为师父，也不排除圣子为师父。因为\Quote{唯一}\OriginalTerm{独一}{one}的说法是与世人及其余受造物相对而言的。\par

祂既警戒他们提防这严重的恶疾，并予以纠正，接着又教导他们如何藉谦卑逃避它。故祂又说：\Quote{你们中最大的，要做你们的仆人。因为凡自高的，必降为卑；凡自卑的，必被高举。}\par

因为没有什么能与谦逊的操练相比，所以祂屡次提醒他们这一德行，既在把孩童带到中间时，也在此处。在山上的真福中，祂也是从此开始。在此处，祂更是连根拔除骄傲，说：\Quote{自卑的，必被高举。}\par

你看见祂如何将听者完全转向相反的一面吗？祂不仅禁止他贪求首位，反而要求他追求末位。因为这样你才能得到所愿的，祂说。因此，那渴望首位的人，必须追求末位。\Quote{因为自卑的，必被高举。}\par

我们到哪里去寻找这种谦卑呢？你们愿意我们再去那座德行之城，圣人们的帐幕——我说的是山岭和树林吗？因为在那里，我们也将看到这种谦卑的高峰。\par

有些人因世间的地位而显赫，有些人因财富，但他们却在各方面贬抑自己：在衣着、住所、服务对象上，如同用文字一样，在所有事上刻写着谦卑。\par

那些激发傲慢的事物——如穿着华丽、建造辉煌的房屋、拥有众多仆役——这些常使人甚至违背自己的意愿而傲慢的事情，都被除去了。因为他们亲自点火，亲自劈柴，亲自烹饪，亲自服侍来客。\par

在那里，听不到有人侮辱人，也看不到有人受侮辱，没有人发号施令，也没有人接受命令；所有人都专心服侍那些被服侍的人，每个人都为陌生人洗脚，并且为此争相去做。他并不问那人是谁，是奴隶还是自由人，而是对每个人都履行这项服务。那里没有伟大与卑微之分。那么，会有混乱吗？绝无，而是最高的秩序。因为若有人卑微，那伟大的人并不看这一点，反而自认比那卑微者更低下，因此成为伟大的。\par

众人共用一张餐桌——服侍者与被服侍者相同；同样的食物，同样的衣服，同样的住所，同样的生活方式。在那里伟大的人，是积极抓住卑微工作的人。没有\Quote{我的}和\Quote{你的}，这个词已被消灭，它是无数战争的根源。\par

4. 如果众人的生活方式、餐桌和衣着都一样，又有什么可惊奇的呢？何况众人甚至只有一个灵魂，不仅在本性上（这也是众人共有的），而且在爱上。那么，它又如何能自相高抬呢？那里没有财富与贫穷，没有尊荣与羞辱；那么，骄傲和狂妄如何能进入呢？因为他们在德行上确有大小之分，但如我所说，无人注意到这一点。那小的不会因被轻视而痛苦，因为没有人轻视他；即使有人轻蔑他，他们所受的教导就是要被人轻视、被藐视、在言语和行为上被轻看。他们与贫穷和残疾的人交往，他们的餐桌充满了这样的客人；因此他们堪得天国。有人照料残废者的伤口，有人给瞎子引路，有人扶携瘸腿的人。\par

那里没有成群谄媚者或食客；或者说，他们甚至不知道什么是谄媚；那么，他们何时会自高呢？因为他们中间有很大的平等，因此也有很大的行德之便。\par

因为这样，那些次等的人比被迫让出首位时得到更好的教导。\par

正如那暴躁的人从被打的人那里得到教训并顺服；同样，那争名的人从不求荣耀、反鄙视荣耀的人那里得到教训。他们在此处丰盛地这样做，正如我们争抢首位那般激烈，他们争抢的却是不要首位，而是完全拒绝它；他们热切渴望的是谁能在尊敬别人、而非被尊敬上更胜一筹。\par

此外，连他们的工作也劝导他们操练节制，不要高傲。因为请问，谁在挖土、浇水、种植，或编篮子、织麻袋，或从事其他手工劳动时会骄傲呢？谁住在贫穷中，与饥饿搏斗，还会患上这种病呢？没有一人。因此谦卑对他们来说是容易的。而在这里，因为有许多人夸奖、钦佩我们，所以谦卑是艰难的；在那里却是极其容易的。\par

那人只关注旷野，看到鸟儿飞翔，树木摇曳，微风吹拂，溪流穿过幽谷。那么，住在如此孤寂中的人，如何会自高呢？\par

但这并不意味着我们可以因此有借口，因为我们身处人中间时便骄傲。因为的确，\Person{亚巴郎}身处客纳罕人中时说：\BibleQuote{我是尘埃和灰烬；}\BibleRef{创 18:27}，和\Person{达味}在军营中说：\Quote{我是虫，不是人；}宗徒在世界中说：\BibleQuote{我不配称为宗徒。}\BibleRef{格前 15:9}那时我们还有什么安慰？什么托词，当我们拥有如此伟大的榜样，却不实行节制？因为他们配得无数冠冕，首先走上了德行之路；同样，我们也该受无数惩罚，我们连那些已故的、在书中展示给我们的榜样，甚至那些活着、因他们的行为而受人钦佩的人，都不能被引领去效法。\par

因为你能说什么，为自己没有改正？你难道不识字，没有查阅圣经，好学习古人的美德？说实话，这本身就是可责备的，当教会常开不闭，你却不愿进入，分沾那神圣的活泉。\par

然而，即使你不通过圣经认识那些已故者，这些活着的人你总该看见。但难道没有人引导你吗？来找我，我指给你这些圣人的避难所；来向他们学习一些有益的东西。这些人在世界各地是明灯；如城墙般环绕着城市。为此他们占据了旷野，为要教导你轻视尘世间的喧嚣。\par

因为他们，如同强壮的人，即使在波涛汹涌中也能享受平静；但你，处处渗漏，需要安宁，并在接踵而来的波浪之后稍作喘息。所以，不断到那里去，好借着他们的祈祷和劝诫，涤除持久的污点，从而以最佳的方式度过今生，并到达将来的福乐，借着我们的主耶稣基督的恩宠和祂对人的爱，祂与圣父及圣神，同享光荣、权能、尊威，从今世直到永远，永世无尽。阿们。\par

\SourceRecord{Translated by George Prevost and revised by M.B. Riddle.；Nicene and Post-Nicene Fathers, First Series；Vol. 10.；Edited by Philip Schaff.；Buffalo, NY: Christian Literature Publishing Co.,；1888.；<http://www.newadvent.org/fathers/200172.htm>.}

% structure: p0001 source=h1 target=section reason=page-title

\section{玛窦福音论讲第七十三篇}

\Emph{\BibleRef{玛 23:14}}\par

\Emph{\Quote{祸哉，你们经师和法利塞假善人！因为你们吞没寡妇的家产，而以长久的祈祷作掩饰，为此，你们必受更重的惩罚。}}\par

此后，祂接着讥讽他们的贪食；而可悲的是，他们不是从富人的财物，而是从穷人那里满足自己的肚腹，加重了本应救济的穷人的贫困。因为他们不单是吃，而是吞没。\par

而且，他们欺诈的方式更为严重：\Quote{以长久的祈祷作掩饰。}\par

因为任何人作恶都该受罚；但若从虔诚中找寻作恶的理由，并用它作为恶行的掩饰，就更应受更严厉的惩罚。那么，祂为何没有废黜他们？因为时候尚未允许。因此，祂暂时容忍他们，但藉祂的言论确保百姓不受欺骗，免得因那些人的尊贵而被引向同样的效仿。\par

因为祂曾说过：\Quote{凡他们所吩咐你们的，你们要遵守}，祂便指出他们做了多少错事，免得愚昧人以为祂将一切托付给了他们。\par

\Quote{祸哉，你们！因为你们给人关闭了天国；你们自己不进去，也不让那想要进去的人进去。}但若对人无益已是罪过，更何况伤害和阻碍，那还有什么可推诿的呢？\Quote{想要进去的}是什么意思？是指那些适宜进去的人。因为他们吩咐别人时，使担子难以承受；但当他们自己要做所要求的事时，不但不去做，反而更加邪恶，甚至败坏别人。这些就是所谓的害群之马，以毁坏他人为业，与教师完全相反。因为教师的本分是拯救沉沦者，而毁灭将得救者则是毁灭者的作为。\par

之后，又是一项指控：\Quote{你们走遍海洋和陆地，为使一个人归依；及至他归依了，你们却使他成为比你们更坏的地狱之子。}也就是说，即使你们好不容易费尽千辛万苦才得到他，这也不使你们对他心存怜悯；通常对于辛苦获得的东西我们更为珍惜，但你们连这一点也不能使你们变得更温和。\par

这里祂指责他们两件事：一是他们无益于众人的救恩，且需很大努力才能赢得一人；二是他们疏于保守已赢得的人，甚至不仅疏忽，而且背叛，因他们生活的邪恶败坏了他，使他更糟。因为当门徒看到他的老师是这样的人，他会变得比老师更坏。他不停留在老师的邪恶上；正如当老师有德行时，门徒效仿他，同样，当老师邪恶时，门徒甚至超过他，因为我们容易倾向恶。\par

祂称他为\Quote{地狱之子}，即十分地狱。祂说\Quote{比你们还加倍}，为要惊吓那些人，也使这些人更感严重，因为他们是邪恶的教师。不仅如此，还因为他们致力向门徒灌输更大的邪恶，使他们变得比老师更腐化，这尤其显示出堕落的灵魂。\par

然后祂又因他们的愚昧讥讽他们，因为他们命人忽略更大的诫命。之前祂曾说过相反的话：\Quote{他们把沉重难负的担子捆起来，放在人的肩上。}但这些事他们再次做了，且为败坏属下之故，在小事上要求严格，而轻视大事。\par

\Quote{因为你们捐献十分之一的薄荷、茴香，却忽略了法律上更重要的：公义、仁爱与信义。这些固然该做，那些也不可忽略。}\par

所以祂自然在此处说这样的话，因为涉及十分之一和施舍；施舍何害之有？但不可忽略法律；因为经文并非如此说。因此，此处祂说：\Quote{这些固然该做}；但在谈论洁净与不洁净时，祂不再加上这话，而是作出区分，表明内心的纯洁必然带来外表的纯洁，但反之则不成立。\par

因为当涉及爱人的理由时，祂轻描淡写，正是为此，也因为尚未到明确废止法律之事的时机。但涉及身体净化的规矩时，祂更明确地推翻它们。\par

因此，关于施舍祂说：\Quote{这些固然该做，那些也不可忽略}，但关于洁净祂不说这话，而是说什么？\Quote{你们洗净杯盘的外面，里面却满是劫夺与不义。你们先洗净杯的里面，好使外面也干净。}祂从一件公认且明显的事，即杯盘说起。\par

2. 然后，为表明轻视身体洁净没有妨害，而不注重灵魂洁净（即德行）会有极大的惩罚，祂称这些为\Quote{蚊虫}，因为它们是微小而无足轻重的，但那些是\Quote{骆驼}，因为它们超乎人所能承担。因此祂说：\Quote{滤出蚊虫，却吞下骆驼。}因为前者乃是为后者（即仁爱与公义）而设立的；所以即使单独遵守前者也无益。因为既然小事是为大事而设，而大事被忽略，只在小事上劳苦，那么即使如此也毫无所得。因为大事不会随小事而来，但小事却必然随大事而来。\par

但祂说这些话是为表明，即使在恩宠来临之前，这些事也并非主要的事，也非人应劳苦追求的事，而是所需的事是另外的。但若在恩宠之前尚且如此，那么当崇高的诫命来临之时，这些事就更无益处，根本不宜去行了。\par

总之，无论在哪种情况下，恶习都是严重的事，尤其当它甚至不以为需要改正之时；而更严重的是，当它自以为足以改正他人时；基督为表明这一点，称他们为\Quote{盲目的向导}。因为盲人若不自知需要向导，已是极端的悲惨和不幸；而当他竟想引导他人时，请看那是何等深渊！\par

但祂说这些话，都是暗示他们对光荣的疯狂渴望，以及对此祸患的极端癫狂。因为这是他们一切邪恶的原因，即他们为炫耀而行一切事。这既使他们远离信德，又使他们忽视真正的德行，只专注于肉身的洁净，而忽略了灵魂的洁净。因此，为引导他们进入真正的德行和灵魂的洁净，祂提到了怜悯、公义和信德。因为正是这些事构成了我们的生命，这些事洁净灵魂：正义、爱人和真理。前者使我们倾向于宽恕，不让我们对犯罪者过分严厉和不宽容（因为这样我们便双重得益：既对人仁慈，也因此从万有的天主那里得到许多仁慈），并使我们同情受虐待者，帮助他们；后者则不容他们诡诈狡猾。\par

但祂说\Quote{这些是你们当行的，而那些也不可放下}，并非为引入遵守法律；绝无此意；祂说\Quote{先清洁杯盘里面，好叫外面也干净}时，也不是针对杯盘而论，而是完全相反，祂做这一切是为表明这些是多余的。因为祂不是说清洁外面，而是说里面，外面自然也就干净了。\par

此外，祂也不是在谈论杯盘，而是灵魂和身体，以外面指身体，以里面指灵魂。但若杯盘需要清洁里面，那么对于你们更是如此。\par

但你们做的正相反，祂说，你们注重琐细外表的事，却忽略了重大内在的事；由此产生极大的祸害，因为你们以为已尽全责，就轻视其他事；而轻视它们，你们便不努力、也不尝试去行它们。\par

此后，祂又因他们的虚荣而讥讽他们，称他们为\BibleQuote{粉饰的坟墓}\BibleRef{玛 23:27}，并对一切加上\Quote{假善人}；这是他们一切邪恶的原因和败坏的根源。祂不仅称他们为粉饰的坟墓，还说他们充满不洁和伪善。祂说这些话，是表明他们不信的原因，即因为他们充满伪善和不义。\par

然而不仅基督如此说，先知们也常指责他们，说他们掠夺，说他们的长官不按公义审判，处处可见祭献被拒绝，而这些都是所要求的。故此没有什么奇怪，没有什么新奇，无论在立法上，还是在指控上，甚至坟墓的比喻也是如此。因为先知也提到坟墓，且不单称他们为坟墓，而是\Quote{他们的咽喉是敞开的坟墓}。\par

如今许多人也是这样，外表装饰得漂亮，里面却充满不义。因为如今也有许多方式、许多对表面洁净的关心，但对灵魂的洁净却连一个也没有。若有人真的剖开每个人的良心，会发现许多虫子和腐烂，以及难以言喻的恶臭；我是指无理邪恶的情欲，它们比虫子更不洁。\par

3. 不过他们是这种人固然可怕（虽然确实可怕），但你们这些被尊为天主之殿的人，竟突然成为坟墓，带着同样恶臭，这才是极端的悲惨。基督住在其中、圣神运行、领受如此伟大奥迹的人，竟成为坟墓，这是何等悲惨！这需要何等的哀哭和悲叹！基督的肢体竟成为不洁的坟墓！想想你是如何出生的，你被赐予了什么，你领受了怎样的衣服，你是怎样被建造成无玷的圣殿！何等美丽！不是用金银装饰，而是用比这些更宝贵的圣神。\par

请想想，城中不会造坟墓，同样，你也不能出现在天上的城中。因为若在此地禁止，那里更不准。或者不如说，即使在此地，你也成为众人嘲笑的对象，带着死去的灵魂，不仅是嘲笑，也是躲避。请告诉我，若有人到处带着一具死尸，岂不人人都逃开？岂不都躲避？现在也请这样想。因为你到处带着比这更可怕的景象，一个被罪致死的灵魂，一个瘫痪的灵魂。\par

如今谁会怜悯这样的人呢？因为你不怜悯自己的灵魂，别人怎能怜悯那如此残忍、与自己为敌的人呢？若有人在你的卧室和饭厅埋了一具死尸，你会怎么办？但你却是埋葬一个死去的灵魂，不是在你的饭厅或卧室，而是在基督的肢体里；你岂不害怕千万雷电和霹雳从天上击打你的头吗？\par

你怎敢踏入天主的教会和神圣的殿宇，而身上却带着如此可憎的气味呢？因为如果有人带着死尸进入王宫并将其埋葬，必受极刑；而你竟敢踏入神圣的殿宇，使整个殿宇充满如此恶臭，试想你将受何等刑罚。\par

你当效法那位用香液傅抹基督双脚、使整个房屋充满香气的妓女；你对她主所行之事，却正与之相反！虽然你或许闻不到这恶臭，但正是这病最严重之处；因此你已病入膏肓，比那些身体残废、发臭的人更为严重。因为那种病既为患者所感知，且无过咎，反值得怜悯；而此病却招致憎恨和刑罚。\par

既然这一方面更为严重，且病人不按理感受它；那么，请你倾听我的话，使我能清楚地教导你这一祸害。\par

但先请听你在圣咏中所说的话：\BibleQuote{愿我的祷声如馨香上达于你面前。}既然如此，当你身上和你的行为中不是馨香，而是臭气升起时，你当受何等刑罚呢？\par

那么何为臭气呢？许多人进来，凝视妇女的美貌；另有人好奇地观看少年的青春。此后，你们岂不惊奇，为何没有闪电击下，一切事物不被连根拔起？因为所行之事既配受雷霆也配受地狱；但天主，祂是宽仁而大慈大悲的，暂时忍耐祂的义怒，召叫你们悔改和改善。\par

人哪，你在做什么？你好奇地观看妇女的美貌，难道不因这样侮辱天主的殿宇而战栗吗？难道教会对你而言就像一个妓院，比市场更不受尊敬？因为在市场上，你尚且害怕并羞于显得打量任何女人；但在天主的殿宇内，当天主亲自向你讲话，并警告你关于这些事时，你却在听闻不可行此事的同时，行奸淫和通奸。你难道不战栗，不惊愕吗？\par

这些事是纵情逸乐的景象所教导你的，是难以遏制的祸害，是恶毒的邪术，是愚昧者致命的陷阱，是放荡者乐在其中的毁灭。\par

因此，先知也曾责备你们说：\BibleQuote{你们的眼不好，你们的心也不正。}\par

这样的人倒不如瞎眼；倒不如患病，总比滥用双眼于这些事好。\par

诚然，你们本应在你们中间有一道隔墙将你们与妇女分开；但既然你们不如此想，我们的父辈认为有必要用这些木板将你们隔开；因为我从长辈听说，古时并没有这些分隔；\BibleQuote{因为在基督耶稣内，不再分男女性别。}在宗徒时代，男女也在一起。因为那时男人是男人，女人是女人；但现在完全相反：妇人已自甘于娼妓的作风，而男人也不比疯狂的牡马好多少。\par

你们难道没有听说，男人和女人曾在楼上聚集，那个聚会是配得上天堂的吗？这很合理。因为连妇人那时也多多克己，而男人则庄重而有节。例如，听那卖紫布的妇人说：\BibleQuote{如果你们认为我是忠于主的，就请进来，与我同住。}又听那些跟随宗徒的妇人，她们拥有男子的勇气，如普黎斯加、培尔息等；我们的妇人与她们相比，正如我们的男子与那些男子相比，相距甚远。\par

4. 因为那时，妇人们即使远行，也不致使自己蒙受恶名；但如今，即使被养在深闺，也难逃嫌疑。但这些事源于她们的装饰和奢侈。那时，那些妇人的事业是传播圣言；但现在，却是显得美丽、漂亮和容貌姣好。这就是她们的光荣，她们的救恩；而对于崇高伟大的事业，她们甚至连想都不想。\par

有哪个妇人努力使自己的丈夫更好？有哪个男子以此为己任来改造自己的妻子？一个也没有。妇人全副心思都在金饰、衣服和其他身体的装饰上，以及如何增加财产；而男子不仅关心这些，还有别的，但全都是世俗的。\par

有谁在将要结婚时打探那女子的品性和教养呢？无人；而是马上问及金钱、财产和各种各样的地产尺度，好像要买什么东西，或者要订立某个普通合同似的。\par

因此他们甚至这样称呼婚姻。因为我听许多人说，某人向某女子立约了，也就是说娶了她。他们侮辱天主的恩赐，像买卖一样娶嫁。\par

而且婚约所需的保障比买卖更大。你当学习古人如何婚配，并效仿他们。那么他们如何结婚呢？他们查考生活方式、品德和灵魂的美德。因此他们不需要婚约，也不需要羊皮纸和墨水的保证；因为新娘的品性对他们而言就足够了。\par

因此我也劝你们，不要寻求财富和富裕，而要寻求善良和温柔。寻求一位虔诚和克己的少女，这些将对你而言胜过无数宝藏。如果你寻求天主的事，这些其他的也会加给你；但如果你忽略那些，而急于这些，那么连这些也不会跟着来。\par

但有人会说：这样的人是靠妻子致富的！难道你不羞于举出这样的例子吗？我宁愿一万次成为穷人——正如我听许多人说过的那样——也不愿靠妻子获得财富。因为有什么比那财富更令人不快？有什么比富足更痛苦？有什么比因此出名更可耻，并且被所有人说\Quote{某人靠妻子致富}？至于家中因此必然产生的烦恼——妻子的骄傲、奴性、争吵、仆人的责骂——我都略过不提。\Quote{这个乞丐，}\Quote{这个衣衫褴褛的人，}\Quote{这个卑贱的、出身低微的人。}\Quote{他进来时有什么？}\Quote{难道一切不都属于我们的女主人吗？}但你根本不在意这些话，因为你也不是自由人。那些寄生虫同样听到比这更糟的话，却毫不痛苦，因此这些人也不痛苦，反而以自己的耻辱为荣。当我们告诉他们这些事时，他们中有人说：\Quote{让我有什么愉悦甜蜜的东西吧，就算噎死我。}唉！魔鬼把人世间什么样的谚语带来了，足以颠覆这些人的整个生活。至少看看这同一个魔鬼般有害的谚语，它充满了多大的毁灭。因为它的意思无非就是：不要顾及体面；不要顾及正义；抛弃一切，只寻求快乐。即使那东西使你窒息，也要选择它；即使所有遇见的人唾弃你，往你脸上抹泥，像赶狗一样赶你走，也要忍受一切。如果猪有声音，它们还会说什么？肮脏的狗还会说什么？但或许连它们也不会说这样的话，就像魔鬼说服人类发狂一样。\par

因此，我恳请你们意识到这类言语的愚昧，逃避这些谚语，并选择圣经中与它们相反的谚语。\par

这些是什么呢？经上说：\Quote{不要随从你的私欲，应克制你的贪欲。}\BibleRef{德 18:30}关于妓女，又有一段话与这个谚语相反：\Quote{不要留心坏女人，因为妓女的嘴唇滴下蜂蜜，在一段时间内你的喉咙感到甘甜；但此后你必发现它比苦胆更苦，比双刃剑更锋利。}因此，让我们听从后者而非前者。因为由此才生出我们的卑劣、我们的奴性思想，使人成为野兽，因为他们凡事都按这个谚语追求快乐——即使不用我们的论证，这谚语本身也是可笑的。因为一个人被噎死后，甘甜有何益处呢？\par

因此，请停止树立如此巨大的荒谬，停止点燃地狱和不灭的火；让我们（虽然为时已晚）按我们应当的那样坚定地展望未来的事，除去眼前的薄膜，以便我们既以极大的敬意和虔诚的敬畏度过今世，又借着我们的主耶稣基督的恩宠和对人的爱，获得未来的美善。愿荣耀归于祂，直到永远。阿们。\par

\SourceRecord{Translated by George Prevost and revised by M.B. Riddle.；Nicene and Post-Nicene Fathers, First Series；Vol. 10.；Edited by Philip Schaff.；Buffalo, NY: Christian Literature Publishing Co.,；1888.；<http://www.newadvent.org/fathers/200173.htm>.}

% structure: p0001 source=h1 target=section reason=page-title

\section{\WorkTitle{玛窦福音 第七十四篇讲道}}

\BibleQuote{\Emph{祸哉，你们！因为你们建造先知的墓，修饰他们的坟墓}，\Emph{并说：我们若在我们祖先的日子，必不与他们同流先知的血。}}\par

祂说\BibleQuote{祸哉}，并非因为他们建造，也并非因为他们责备别人，而是因为，他们既如此行事，又如此说话，假装谴责自己的祖先，却做出更坏的事。为证明这种谴责是假装，路加说：\BibleQuote{祸哉，你们！因为你们建造先知的墓，而你们的祖先杀了他们。你们确实作证，并赞同你们祖先的作为，因为他们杀了他们，你们却建造他们的墓。} \BibleRef{路 11:47} 因为此处祂责备他们的意图，他们建造的意图不是为尊崇被杀者，而是为了炫耀谋杀，并害怕坟墓因时间而湮灭后，这种胆大妄为的证明和记忆会消失，于是建立这些辉煌的建筑，如同某种凯旋纪念碑，并以此夸耀那些人的胆大行为，并展示它们。\par

因为你们现在所敢做的事，表明你们也是本着这种精神做这些事。因为，虽然你们说相反的话，祂说，如同谴责他们，例如：\BibleQuote{我们若在他们那时代，必不与他们同流；}然而你们说这些话的意向是明显的。因此，祂虽然隐约地，但仍将之表达出来而加以阐明。因为当祂说了：你们说\BibleQuote{我们若在我们祖先的日子，必不与他们同流先知的血}时，祂接着又说：\BibleQuote{所以你们自己证明自己是杀害先知者的子孙。}而如果一个人不分享父亲的心意，那么作为杀人犯的儿子又有什么罪责呢？没有。由此可见，祂提出这一点正是为此，暗示他们在邪恶上的相似。\par

这一点由后面的话也显然可见；祂至少接着说：\BibleQuote{你们这些蛇，毒蛇的种类。} \BibleRef{玛 23:33} 因为正如那些野兽在毒液的破坏性上与其父母相似，同样，你们在凶杀上也与你们的祖先相似。\par

然后，因为祂察验他们的心境，这对大多数人来说是隐晦的，祂便从他们即将犯下的、对所有人显然的事来确立祂的话。因为祂说了：\BibleQuote{所以你们自己证明自己是杀害先知者的子孙}，表明祂所说的是他们在邪恶上的相似，并且说\BibleQuote{我们必不与他们同流}是伪装，祂接着又说：\BibleQuote{你们就充满你们祖先的恶贯吧！} \BibleRef{玛 23:32} 这并不是命令，而是预言即将发生的事，即祂自己的被杀。\par

因此，祂引出了对他们的反驳，并表明他们为自己辩护的言语是伪装的，例如：\BibleQuote{我们必不与他们同流}（因为那些不放过主的人，怎能放过仆人呢？），在此之后，祂使祂的言语更具谴责性，称他们为\BibleQuote{蛇，毒蛇的种类}，并说：\BibleQuote{你们怎能逃脱地狱的刑罚？}你们既做出这样的事，又否认它们，并掩饰自己的意图。\par

然后，从另一个原因更严厉地责备他们，祂说：\BibleQuote{我派遣先知、贤士和经师到你们这里来，其中有的你们要杀死并钉死，有的你们要在会堂里鞭打。} 为了使他们不致说：\BibleQuote{我们虽然钉死了主，但如果那时我们在，我们必会放过仆人；}祂说：\BibleQuote{看，我也派遣仆人，即同样先知到你们这里来，你们也不会放过他们。} 但祂说这些话，是表明祂被这些儿子谋杀毫不奇怪，因为他们是凶杀和欺诈的，并且拥有许多诡诈，在暴行上超过他们的祖先。\par

此外，除了以上所说的，祂还表明他们极其虚荣。因为当他们说：\BibleQuote{我们若在我们祖先的日子，必不与他们同流}时，他们是出于虚荣而说话，只在言语上实践德行，但在行为上却相反。\par

你们这些蛇，毒蛇的种类，即恶人的恶子，比生养他们的更恶。因为祂表明他们犯下更大的罪行：既因为他们在这之后犯下这些罪行，又因为他们所做的比祖先更严重得多，而他们却声称自己绝不会陷入同样的事。因为他们添加了那既是终结又是他们恶行之冠冕的。因为祖先杀了那些来到葡萄园的人，但这些人杀了儿子以及召请他们赴婚宴的人。\par

但祂说这些话，是为了将他们与亚巴郎的亲属关系分开，并表明如果他们不追随他的行为，从那里毫无益处；因此祂接着说：\BibleQuote{你们怎能逃脱地狱的刑罚？}如果你们追随那些做出如此行为的人？\par

在这里祂使他们想起若翰的控告，因为他也曾以这名号称呼他们，并提醒他们将来的审判。然后，因为他们不被审判和地狱所惊吓，由于他们不相信这些，又因这事是未来的，祂以眼前的事恐吓他们，说：\BibleQuote{因此，看，我派遣先知和经师到你们这里来：其中有的你们要杀死并钉死，有的你们要在会堂里鞭打，为使流在地上的所有义人的血，从义人亚伯尔的血，直到你们在圣殿与祭台之间所杀的贝勒基雅的儿子匝加利亚的血，都归到你们身上。我实在告诉你们：这一切都要归到这一代人身上。} \BibleRef{玛 23:34}\par

2. 请看祂用多少事警告了他们。祂说：你们定你们的祖先有罪，因为你们说：\Quote{我们不会与他们同流合污}；这件事足以使他们羞愧。祂说：当你们定他们的罪时，你们自己却做了更坏的事；这足以使他们蒙羞。祂说：这些事不会没有惩罚；由此祂在他们心中植入了言语无法形容的恐惧。祂至少提醒了他们地狱。然后，因为那是将来的事，祂将恐惧带到他们面前，如同已经临到。祂说：\Quote{这一切都要临到这一代人。}\par

祂还在惩罚上加上了难以言表的严厉，说他们将遭受比所有人更痛苦的事；然而，这些事没有一件使他们变得更好。但若有人说：为什么他们比所有人受更重的苦？我们要说：因为他们首先犯下了比所有人更重的罪，并且因他们所行的事，没有一件使他们醒悟过来。\par

你没有听见\Person{拉默客}说：\Quote{从\Person{拉默客}要讨回七十七倍的罪}\BibleRef{创 4:24}吗？也就是说：\Quote{我比\Person{加音}更该受惩罚。}为什么会这样？他并没有杀他的兄弟；但是因为连他的榜样也没有使他变得更好。这正是天主在其他地方所说的：\Quote{对恨我的第三代和第四代子孙追讨祖先的罪债。}并不是说一个人要为别人所犯的罪受罚，而是说那些在许多人犯罪并受罚之后仍未变好、反而犯同样罪行的人，理当承受他们的惩罚。\par

但请看祂多么合时宜地提到了\Person{亚伯尔}，表明这谋杀同样出于嫉妒。那么你们有什么可说的呢？你们不知道\Person{加音}受了什么苦吗？天主对祂的作为保持沉默了吗？祂没有施行最严厉的惩罚吗？你们没有听见你们的祖先在杀害先知时受了什么苦吗？他们不是被交付给惩罚和无数次的报复吗？那么你们怎么没有变好呢？我为什么要提你们祖先的惩罚和所受的苦呢？你们自己定你们祖先的罪，怎么倒做了更坏的事？因为连你们自己也说过：\Quote{祂要凶恶地消灭那些恶人。}\BibleRef{玛 21:41}那么，你们在说了这样的话之后还做这样的事，还能得到什么恩惠呢？\par

这\Person{匝加利亚}是谁？有人说，是\Person{若翰}的父亲；有人说，是先知；有人说，是有两个不同名字的司祭，\WorkTitle{圣经}也称他为\Person{约雅达}的儿子。\par

但你要注意，这暴行是双重的。因为他们不仅杀害了圣洁的人，而且在圣洁的地方。祂说这些话时，不仅警告他们，也安慰祂的门徒，表明在他们之前的义人也遭受了这些事。但祂警告这些人，预言正如那些义人付出了代价，这些人也同样要受最极端的苦。因此祂称他们为\Quote{先知、明智者、经师}，这也剥夺了他们的一切借口。祂说：\Quote{你们不能说：\InnerQuote{你们是从外邦人中派来的，所以我们感到冒犯}；}但他们被引向谋杀的性情和嗜血。所以祂也事先说：\Quote{为此我派遣先知和经师。}先知们也以此责备他们，说：\Quote{他们把血与血混杂}，又说他们是血性之人。因此祂也命令将血献给自己，表明如果连牲畜的血都如此珍贵，何况是人呢？祂对\Person{诺厄}同样说：\Quote{凡流人血的，我也要追讨。}\BibleRef{创 9:5}人们还可以找到祂关于不杀人的成千上万其他命令；因此祂甚至命令他们不可吃勒死的牲畜。\par

哦，天主对人的爱！虽然祂预知他们不会得益，祂仍然尽祂的本分。祂说：因为我要派遣，并且知道他们会被杀害。所以连这一点也证明他们徒然说：\Quote{我们不会与我们的祖先同流合污。}因为这些人在他们的会堂里也杀害先知，既不尊重那地方，也不尊重那人的尊贵。因为他们杀害的不仅是普通人，而是先知和明智者，以致他们无可指责。祂指的是宗徒和之后的人，因为确实许多人曾说过预言。然后，为了加重他们的恐惧，祂说：\Quote{我实在告诉你们：这一切都要临到这一代人。}也就是说：我要把这一切都归在你们头上，并使惩罚更严厉。因为那知道许多人犯了罪，却没有醒悟，反而自己又犯同样的罪，而且不仅是同样的，甚至是更重的，他就理当受比他们更重的惩罚。因为正如他若愿意，本可以从他们的榜样中得到益处，同样，既然他仍不改正，他就该受更重的惩罚，因为他从先前犯罪受罚的人那里得到了更多的警告，却毫无所得。\par

3. 然后祂向\Place{耶路撒冷}城说话，以此意图纠正听众，并说：\Quote{\Place{耶路撒冷}，\Place{耶路撒冷}！}\BibleRef{玛 23:37}重复是什么意思？这是怜惜她、哀叹她、非常爱她的方式。因为就像对一位心爱的女子，她虽始终被爱，却轻视了爱她的人，因此即将受罚，祂在将要降罚时恳求她。先知们也这样说过：\Quote{我说：你归向我吧，她却没有回来。}\BibleRef{耶 3:7}\par

然后，祂召她前来，也指出了她血腥的行为，说：\BibleQuote{你杀害先知，用石头砸死那些被派到你这里来的人。我多少次愿意聚集你的子女，你们却不愿意。}祂藉此也为自己的作为辩护：即使有了这些事，你们也没有使我偏离，也没有把我对你们的深厚感情夺去；相反，我仍然渴望——不只是一次或两次，而是常常——把你们吸引到我身边来。因为祂说：\BibleQuote{我多少次愿意聚集你的子女，如同母鸡聚集小鸡在翅膀底下，你们却不愿意。}祂说这话，为表明他们常常因自己的罪恶而四散。祂用这个比喻来显示祂的情感；事实上，这种动物对幼雏怀着温暖的爱。在众先知中，也处处有这同一翅膀的意象，在梅瑟的歌和圣咏中也是如此，表明祂伟大的保护和关怀。\par

祂说：\BibleQuote{但你们不愿意。看，你们的房屋将留给你们，成为荒场。}\BibleRef{玛 23:38} 被剥夺了从我而来的援助。当然，正是那同一位的祂，从前保护他们、聚集他们、保存他们；正是祂常常责罚他们。祂制定了一种惩罚，就是他们一直极其惧怕的；因为这宣告了他们政体的彻底毁灭。\BibleQuote{因为我告诉你们：从今以后，你们再不会看见我，直到你们说：奉主名而来的，当受赞美！}\BibleRef{玛 23:39} 这是深爱着他们的人所说的话，热切地藉着将来的事吸引他们归向祂，而不仅仅是以过去的事来警告他们；因为祂在这里说的是祂第二次来临的未来日子。\par

那么，从那时起他们就没有再看见祂吗？但祂说\Quote{从今以后}，所指的不是那一时刻，而是直到祂被钉十字架的时候。\par

因为他们一直以此事控告祂，说祂是某种对立的神，与天主为敌；祂藉此来激发他们对祂的爱，即通过显示祂与父同心合意；祂表明自己就是那曾在先知们中间的同一位。因此祂也用了先知们所用的同样话语。\par

藉着这些话，祂预示了祂的复活和第二次来临，甚至对完全不信的人也明白地显示，那时他们必定会朝拜祂。祂如何使这事明白呢？祂说了许多首先将要发生的事：祂将派遣先知，他们将杀害他们；这事要在会堂里发生；他们将遭受极端的苦难；他们的房屋将被遗弃荒凉；他们将经历比任何人都更悲惨的事，是前所未有的。因为这一切足以向最愚钝、最好争辩的人提供明确的证据，证明祂来临时所将发生的事。\par

我要问他们：祂派遣了先知和智者吗？他们在会堂里杀害了他们吗？他们的房屋被遗弃荒凉了吗？所有的报应都临到了那一代人身上吗？显然是这样，没有人反驳。那么，正如这一切都发生了，将来那些事也必会发生，而且那时他们必定会降服。\par

但他们那时将不能从中获得任何辩护的益处，正如那些在人生道路上悔改的人也不能。\par

所以，让我们趁着还有时间，实践善行。因为正如他们那时从知识中不再获得益处，我们也不会从对罪恶的悔改中获益。因为对于舵手，当船因他的疏忽沉入大海时，他再也不能做什么；对于医生，当病人去世时，他也无能为力；每个人都必须在结束之前筹划并执行一切，以免陷入危险和羞耻；但此后，一切都是徒劳的。\par

因此，我们也要在生病时请来医生，花费金钱，坚持不懈地努力，以便从痛苦中起来，健康地离开这里。\par

我们为仆人身体生病时所操的心，也要同样为我们自己灵魂患病时表现出来。事实上，我们比仆人更亲近自己，我们的灵魂也比那些身体更为必要，但至少我们若能付出同等的努力就好。因为如果我们现在不这样做，一旦离世，此后我们再也得不到任何辩护的机会。\par

4. 有人会说：谁这样可怜，连这一点心思都不肯花呢？这正是令人惊奇的事：我们对自己如此轻视，以至于我们比对待仆人更加藐视自己。因为当我们的仆人生热病时，我们请来医生，在家中隔离，强迫他们遵守医术的规律；如果这些被忽视，我们就不高兴，并派人看守他们，即使他们想纵情，也不允许；如果照顾他们的人说必须花大价钱买药，我们也顺从；他们吩咐什么，我们都听从，并且为这些吩咐付给他们报酬。\par

\Quote{当我们患病时（更确切地说，没有什么时候我们不患病），我们甚至不请医生来，不动用钱财，反而好像涉及一个暴徒、仇敌和敌人那样，忽视我们的灵魂。我说这些话，不是指责我们对仆人的关心，而是认为至少应当同样关心我们的灵魂。那该怎么做呢？有人会说。将它呈给生病的\Person{保禄}；请来\Person{玛窦}；让\Person{若望}坐在它旁边。从他们那里听听，我们这些病人该做什么，他们一定会告诉我们，不会隐瞒。因为他们没有死，而是活着并说话。但灵魂因发烧而沉重，不听他们么？你要强迫它，唤醒它的理性。请来先知。无需付钱给这些医生，因为他们自己也不索要酬劳，他们配制的药物也不会迫使你承担花费，除非为了施舍；但在其他一切事上，他们甚至增加你的财产；例如，当他们要求你节制时，他们使你免于不合时宜且错误的花费；当他们告诉你要戒除醉酒时，他们使你更富裕。你看见这些医生的技艺了么？他们在健康之外，还供给你们财富。因此，坐在他们旁边，向他们学习你疾病的本质。例如，你爱财富和贪得无厌，如同发烧的人爱水？无论如何，听一听他们的劝诫。因为如同医生对你说：\Quote{如果你满足你的欲望，你将死亡，并遭受这个或那个；}同样，\Person{保禄}也说：\BibleQuote{那些想要致富的人，陷于诱惑，落入魔鬼的罗网，并且陷入许多愚蠢而有害的私欲，这私欲把人淹没在败坏和灭亡中。}\BibleRef{弟前 6:9}}\par

但你急躁吗？听他说：\BibleQuote{还有一点点时候，那要来的就来，必不迟延。\BibleRef{希 10:37} 主已经近了，什么都不要挂虑；}\BibleRef{斐 4:5} 又说：\BibleQuote{这世界的局面正在逝去。}\BibleRef{格前 7:31}\par

因为他不仅命令，而且安慰，如同医生应该做的。正如他们用其他东西代替冷饮，这个人也以另一种方式消除欲望。他说：你愿意致富吗？那就\Quote{在善工上}。你愿意积蓄财宝吗？我毫不禁止；只让它在天上。\par

正如医生所说，冷的东西对牙齿、神经、骨骼有害；同样，他——确实更简洁，而注意简洁，却更清晰有力地说：\BibleQuote{因为贪爱金钱是万恶的根源。}\BibleRef{弟前 6:10}\par

那么，应该使用什么呢？他也告诉了这个：用知足代替贪婪。他说：\BibleQuote{因为知足与虔诚结合，就是大利。}\BibleRef{弟前 6:6} 但如果你不满足，想要更多，还未能抛弃一切多余的东西，他也告诉那个这样患病的人，该怎样处理这些东西。\BibleQuote{那以财富为乐的，要像不乐的；那拥有的，要像没有的；那享用这世界的，要像不尽享用的。}\BibleRef{格前 7:30}\par

你看见他吩咐的是何等样的事吗？你还要请另一位医生来吗？至少在我看来很好。因为这些医生不像身体的医生，他们往往彼此竞争，把病人压垮。但这些医生并非如此，因为他们关心病人的健康，而不是自己的虚荣。那么，不要害怕他们的数量；一位导师在众人中说话，就是基督。\par

5. 看，例如，另一位又进来，说了关于这种疾病的严厉的话，更确切地说，是导师藉着他说话：\BibleQuote{因为你们不能服事天主而又服事\OriginalTerm{玛门}{Mammon}。}\BibleRef{玛 6:24} 他说：是的，但这些事如何成就？我们怎样停止欲望？由此我们也可以学到这个。我们怎样知道？听他说这个：\BibleQuote{不要在地上为自己积蓄财宝，那里有虫蛀和锈蚀，也有贼挖洞偷窃。}\BibleRef{玛 6:19}\par

你看见吗？藉着地方，藉着那里朽坏的东西，他吸引人离开此地的欲望，将他们锚定于天堂，因为那里万物坚固。因为如果你将你的财富迁移到那里，那里既无锈也不被虫蛀，也没有贼挖洞偷窃，你既能驱除这疾病，又能将你的灵魂安顿在最大的丰盈中。\par

此外，与我们所说的一起，他举了一个例子教你节制。正如医生为警告病人，说某某人因饮冷水而死；同样，他也引用那个富少年（\BibleRef{玛 19:16}），他确实劳苦，渴望生命和健康，却不能达到，因为他一心贪恋财产，反而空手离去。除了这个人，另一位福音作者又给你展示了另一个人，就是那个在痛苦中的人（\BibleRef{路 16:24}），他连一滴水也掌握不了。然后，为显示他的诫命是容易的，他说：\BibleQuote{你们看天空的飞鸟。}\BibleRef{玛 6:26} 但他出于怜悯，连富人也不让他绝望。他说：\BibleQuote{在人所不能的，在天主却是可能的。}\BibleRef{路 18:27} 因为即使你富有，医生也能治愈你。因为他所去掉的并不是财富，而是做财富的奴隶，做贪得无厌的人。\par

那么，富人怎能得救呢？与那些有需要的人共有他的财物，如同\Person{约伯}那样，并从他灵魂中根除更多的欲望，绝不超过真正的需要。\par

祂同时也向你展示了这位税吏，他曾深受贪婪热病的折磨，却迅速从中解脱。有什么比税吏更卑污的呢？然而，这人因遵守医师的法则，对财富变得漠不关心。因为祂所收为门徒的，正是这样的人，他们曾患有与我们相同的疾病，却迅速恢复了健康。祂向我们展示每一个，为使我们不致绝望。请看这位税吏。再看另一位，税吏长，他曾承诺凡勒索来的四倍偿还，并把他所有的一半，为了能接待\Person{耶稣}。\par

但你正燃烧着对财富的极度渴望吗？那就拥有全人类的财富，而不是你自己的。因为，祂说，我赐给你的远超过你所求的，为你打开了全世界富人的家门。\BibleQuote{凡舍弃父亲或母亲、土地或房屋的，必得到百倍赏报。}\BibleRef{玛 19:29} 这样，你不仅会享受更丰盛的财富，而且会彻底消除这痛苦的渴望，轻松承受一切，不但不再贪求更多，甚至常常不求必要之物。同样，\Person{保禄}忍饥挨饿，却比他在饮食时更受人尊敬。因为一个摔跤手在奋力拼搏、赢得桂冠时，也不愿放弃而休息；一个出发航海的商人也不愿后来无所事事。\par

因此，如果我们按所应得的品尝了属神的果实，从那时起，我们甚至不会把现世的事物当作一回事，而被对未来事物的渴望所攫取，如同被最高贵的陶醉所占据。\par

因此，让我们品尝这些，好使我们既从现世事物的纷扰中解脱，又能获得未来的美好事物，藉着\Person{我们的主耶稣基督}的恩宠与对人的慈爱，愿光荣与权能归于祂，从今世直到永远，世世无穷。阿们。\par

\SourceRecord{Translated by George Prevost and revised by M.B. Riddle.；Nicene and Post-Nicene Fathers, First Series；Vol. 10.；Edited by Philip Schaff.；Buffalo, NY: Christian Literature Publishing Co.,；1888.；<http://www.newadvent.org/fathers/200174.htm>.}

% structure: p0001 source=h1 target=section reason=page-title

\section{玛窦福音讲道第七十五篇}

\BibleRef{玛 24:1-2}\par

\Quote{\Emph{耶稣从圣殿里出来，就走了。} \Emph{门徒们前来，把圣殿的建筑指给他看。他回答他们说：\InnerQuote{你们不是看见这一切吗？我实在告诉你们：将来这里决没有一块石头留在另一块石头上，而不被拆毁的。}}}\par

因为祂说了：\Quote{你们的房屋将留给你们荒芜}，并曾预先警告他们许多痛苦的事；因此门徒们听了这些事，好像惊讶，来到祂跟前，把圣殿的华美指给祂看，心想，若如此华美、如此昂贵的材料、难以言喻的工艺将被毁灭；祂便不再仅仅谈论荒芜，而是预言彻底的毁灭。\Quote{你们不是看见这一切吗？}祂说，而你们惊讶、诧异吗？\Quote{将没有一块石头留在另一块石头上。}那么它怎么还留着呢？有人会说。但这算什么？因为预言并未落空。因为祂说这些话或是表示它的完全荒芜，或是就地而言。因为有一部分已被拆毁到地基。\par

此外，我们还愿附带说另一件事：根据已经发生的事，连最爱争辩的人也应当相信，关于遗留部分，它们终将被完全毁灭。\par

\Quote{当祂坐在橄榄山上时，门徒们私下前来对祂说：\InnerQuote{请告诉我们，什么时候要有这些事？你来临和世界终结的预兆是什么？}} \BibleRef{玛 24:3}\par

所以他们私下前来见祂，因为他们打算询问这样的事。他们切望知道祂来临的日子，因为他们急切渴望见到那荣耀，那荣耀是无数福佑的缘由。他们问祂两件事：这些事什么时候发生？即圣殿的倾覆；以及你来临的预兆是什么？但路加说，\BibleRef{路 21:6} 问的是关于耶路撒冷的事，好像他们以为那时就是祂的来临。而马尔谷说，并非所有门徒都问关于耶路撒冷的结局，只有伯多禄和若望，因为他们有更大的发言权。\par

那么祂说什么？\Quote{你们要谨慎，免得有人欺骗你们。因为将有许多人假冒我的名来说：我是基督，并且要欺骗许多人。你们要听见战争和战争的风声。你们要小心，不要惊慌；因为这一切都必须发生，但结局尚未来到。} \BibleRef{玛 24:4}\par

因为他们觉得，当听到将降临于耶路撒冷的事时，自己是在听别人受惩罚，似乎自己不受骚动，只梦想好事，并且期待这些好事立即降临；因此祂再次预言他们将遭遇痛苦的事，使他们警惕，并基于两个理由命令他们警醒，以免被欺骗者的诡计所迷惑，或被降临的祸患的暴力所压倒。\par

因为，祂说，战争将是双重的：欺骗者的战争与敌人的战争，但前者更为严重，因为它在混乱和骚动中临到他们，当人们惊恐不安时。因为那时风暴很大，罗马势力开始兴盛，城市被占领，营寨和武器被调动，许多人轻易相信。\par

但祂指的是耶路撒冷的战争；当然不是指外面和世界各地的战争；因为那些与他们何干？而且，如果祂谈论普世的灾难（这些灾难总是发生），祂就没有说出什么新事。因为在此之前就有战争、骚动和争斗；但祂指的是犹太人的战争，即将在不久后临到他们，因为罗马的武力从此令人忧虑。既然这些事也足以使他们困惑，祂便预言了一切。\par

然后，为了表明祂自己也与罗马人一同攻击犹太人，并与他们作战，祂不仅提到战争，还提到天主降下的灾祸：饥荒、瘟疫和地震，表明战争也是祂允许的，并且这些事的发生并非无缘无故，按照人类往常的常理，而是来自上天的愤怒。\par

因此祂说，它们不会自动或立刻来临，而是带着预兆。为免犹太人说，那些当时相信的人是这些灾祸的始作俑者，祂也告诉了他们灾祸临到他们身上的原因。因为祂先前曾说：\Quote{我实在告诉你们，}祂之前说过，\Quote{这一切都要归到这一代人身上}，提及他们手上的血债。\par

然后，为免他们听见祸患的倾盆大雨，就以为福音失败了，祂补充说：\Quote{你们要小心，不要惊慌，因为这一切都必须发生}，即我所预言的，而诱惑的临近不会使我的任何预言落空；但确实会有骚动和混乱，但没有什么能动摇我的预言。\par

然后，既然祂曾对犹太人说：\Quote{你们不得再见我，直到你们说：奉主名而来的，当受赞美}；门徒们以为毁灭也将是终结；为纠正他们的这种隐秘想法，祂说：\Quote{但结局尚未来到}。因为他们的确像我所说那样怀疑，你可以从他们的问题中得知。因为他们问了什么？这些事什么时候发生？即耶路撒冷何时被毁灭？以及你来临和世界终结的预兆是什么？\par

但祂没有直接回答这个问题，而是先说了其他紧急的事，那些是他们必须先学习的。因为祂既没有立刻谈到耶路撒冷，也没有谈到祂自己的第二次来临，而是谈到即将临到他们门口的祸患。因此祂也通过说\Quote{你们要谨慎，免得有人欺骗你们；因为将有许多人假冒我的名来说：我是基督}来鼓励他们努力。\par

此后，当祂激发他们留心这些事时（因为祂说：\Quote{你们要谨慎，免得有人迷惑你们}），并使他们充满活力，预备他们保持警醒，先讲了假基督的事，然后才讲耶路撒冷的灾难，藉着已经过去的事，尽管他们愚昧好争，仍使他们确信那些将要发生的事。\par

2. 但祂所说的\Quote{战争和战争的谣言}，是指我先前所说的，即将临到他们身上的苦难。此后，因为正如我已说过的，他们以为那场战争之后结局就会来到，请看祂如何警告他们，说：\Quote{但结局还未到。因为民族要起来反对民族，国家要起来反对国家。}\BibleRef{玛 24:7} 祂所讲的是犹太人灾难的前兆。\Quote{这一切是痛苦的开始}，即那些临到他们身上的痛苦。\Quote{那时，人要交出你们去受苦，并要杀害你们。}\par

祂适时地引入他们的灾难，使他们从共同的苦难中得到安慰；不仅这样，祂还加上说，这是\Quote{因我的名}。祂说：\Quote{你们要为我的名被众人恼恨。那时，许多人要跌倒，彼此出卖，彼此仇视；许多假基督和假先知要起来，迷惑许多人。由于罪恶增多，许多人的爱要冷淡；但是，那坚忍到底的，必然得救。}\par

更大的恶是当内部也有战争时，因为有许多假弟兄。你看见这战争是三重吗？来自欺骗者、来自仇敌、来自假弟兄。请看\Person{保禄}也为同样的事哀叹说：\Quote{外有争斗，内有恐惧}\BibleRef{格后 7:5}，以及\Quote{在假弟兄中的危险}\BibleRef{格后 11:26}，又说：\Quote{因为这样的人是假宗徒，是欺诈的工人，伪装成基督的宗徒。}\BibleRef{格后 11:13}\par

此后，比一切更严重的是，他们甚至得不到来自爱的安慰。然后，为表明这些事丝毫不能伤害那些高尚和坚定的人，祂说：不要害怕，也不要惊慌。因为如果你们显出应有的忍耐，危险就不会胜过你们。这有一个明显的证据：福音必定传遍天下各处，你们将超越所有令你们恐惧的事。为使他们不致说：那么我们怎么生活呢？祂更进一步说：你们不仅要生活，还要到处宣讲。因此祂又加上：\Quote{这天国的福音必先传遍天下，对万民作证，然后结局才来到}，即耶路撒冷倾覆的结局。\par

为证明祂所说的就是指这一点，并且在耶路撒冷被攻陷之前福音已经传开，请听\Person{保禄}所说：\Quote{他们的声音传遍了天下}\BibleRef{罗 10:18}；又说：\Quote{这福音是传与天下一切受造物的。}\BibleRef{哥 1:23} 你看见他从耶路撒冷一直跑到西班牙吗？如果一个人承担了这么大的部分，想想其余的宗徒也作了何等事业。\Person{保禄}在写给其他人的书信中又论及福音说：\Quote{这福音在天下一切受造物中结果实、增长。}\par

但是，\Quote{对万民作证}是什么意思呢？因为虽然福音传遍各地，却并未处处被人相信。祂说，这是为了给那些不信的人作证，即为了定罪、为了谴责、为了见证；因为那些相信的人将作证反对那些不信的人，并定他们的罪。为此缘故，福音传遍天下之后，耶路撒冷被毁灭，使他们连一丝固执己见的借口都没有。因为那些看见祂的能力在每个地方闪耀，并瞬间征服世界的人，还有什么借口继续固执己见呢？为证明当时福音已传遍各处，请听\Person{保禄}所说：\Quote{这福音是传与天下一切受造物的。}\BibleRef{哥 1:23}\par

这也是基督能力的一个极大的标记：在二十或至多三十年内，福音已达地极。\Quote{此后，}祂说，\Quote{耶路撒冷的结局就要来到。}因为祂暗示这一点，由下文显明了。\par

因为祂还引入了一个预言，以证实他们的荒凉，说：\Quote{当你们看见\Person{达尼尔}先知所说的\InnerQuote{招致荒凉的可憎之物}站在圣地时——读者应当明白。}祂引他们到达尼尔。所谓\Quote{可憎之物}是指那攻取城市的人所立的像，那使城市和圣殿荒凉的人将像立在圣殿内，因此基督称之为\Quote{招致荒凉的可憎之物}。并且，为使他们知道这些事将发生在他们之中有些人还活着的时候，祂因此说：\Quote{当你们看见招致荒凉的可憎之物时。}\par

3. 从这里，人最可以惊叹基督的能力和他们的勇气，因为他们在这样的时代宣讲——正是犹太国遭受战争的时候，正是人们视他们为叛乱煽动者的时候，当\Person{凯撒}命令将他们都驱逐出去。\BibleRef{宗 18:2} 结果就如同：当大海四面翻腾，黑暗充满天空，连续的船难发生，所有同船的水手在上方争斗，怪物从下方涌起，波涛吞噬水手，雷电坠落，有海盗，船上的人彼此暗算——这时，有人命令一群毫无航海经验、甚至没见过大海的人坐在舵位上，驾驶并作战，而一支庞大的舰队排列整齐地来攻击他们，他们却用一艘独木舟，在如此混乱的船员中，击沉并制服那舰队。因为在外邦人眼中，他们被恨恶如同犹太人；在犹太人眼中，他们被石头砸死，如同与律法作战；他们没有立足之地。\par

于是，一切都是悬崖、暗礁和岩石，城中的、田野的、屋内的，每一个人都与之争战；将军、统治者、平民，所有民族、所有人民，以及言语无法形容的混乱。因为犹太民族在罗马政府眼中极其可恶，给他们带来了无尽的麻烦；即便如此，圣言的宣讲并未受损；但城市被攻陷、被焚烧，使居民陷入无数灾难；然而从那里出来的宗徒，引入新的法律，甚至战胜了罗马人。\par

哦，奇异而奇妙的事实！罗马人当时征服了无数犹太民众，却没有战胜十二个赤手空拳与他们作战的人。什么语言能述说这个奇迹？因为教导者需要具备两样东西：可信赖，并被他们所教导的人所爱；此外，还要有容易接受的言论，以及没有麻烦和骚乱的时间。\par

但那时与这些完全相反。因为他们看似不可信赖，却使那些被他们欺骗的人离开了看似可信赖的人。他们非但不被爱，甚至被恨，并且将人从他们所爱的习惯、祖传风俗和法律中夺走。\par

再者，他们的诫命充满困难；而他们使人离开的事物却充满乐趣。而且他们和服从他们的人遭遇许多危险，许多死亡，此外，时间也给他们带来许多困难，充满了战争、骚乱、动荡，以至于即使没有我们提到的这些事，也足以使一切陷入混乱。\par

我们有充分理由说：\Quote{谁能述说主的大能，宣扬祂的一切赞美？}因为如果祂的百姓在神迹中，因着泥土和砖块就不听从\Person{梅瑟}，那么是谁说服了那些每日被殴打、被杀、遭受不治之症的人，离开平静的生活，而选择这充满血与死的生活，并且传道者对他们来说是陌生人，各方面都非常敌对？我不说对民族、城市和百姓，而是说，将家中所有人都憎恨的人带进一个小屋，通过他试图使他们离开他们所爱的人——父亲、妻子、孩子——他岂不是在开口之前就被撕碎了吗？如果再加上丈夫和妻子在屋内的争吵和冲突，他们岂不会在他踏上门槛之前就将他用石头砸死？如果他又是一个容易被人轻视的人，并且命令苛刻，命令那些生活在奢侈中的人实行克己，同时冲突的对象是人数远多于他且胜过他的人，那岂不是显而易见他将被彻底毁灭？然而，这在一家之中不可能做到的事，基督却在全世界成就了，尽管有悬崖、火炉、沟壑、岩石、陆地与海洋一起与祂争战，祂却将世界的医治者带了进来。\par

如果你愿意更清楚地了解这些事，我是说饥荒、瘟疫、地震和其他灾难，请阅读\Person{约瑟夫}所著的相关历史，你就会准确知道一切。因此祂自己也说：\BibleQuote{不要惊慌，因为这一切都必须发生；}以及\BibleQuote{那坚忍到底的，必然得救；}和\BibleQuote{福音必在全世界宣讲。}因为当他们因所说的恐惧而软弱无力时，祂鼓励他们说：即使发生一万件事，福音必须在世界各地宣讲，然后终结才会来临。\par

4. 你们看到当时的情形何等严峻，战争何等多样吗？这是开始，那时每件要成就的事最需要平静。那么当时的情形是怎样的？因为没有什么阻止我们再次重复同样的事。第一场战争是欺骗者的；\BibleQuote{因为将有，}祂说，\BibleQuote{假基督和假先知：}第二场是罗马人的，\BibleQuote{你们将听到，}祂说，\BibleQuote{战争：}第三场是带来\BibleQuote{饥荒：}第四场是\BibleQuote{瘟疫}和\BibleQuote{地震：}第五场是\BibleQuote{他们要把你们交给苦难：}第六场是\BibleQuote{你们要受到所有人的憎恨：}第七场是\BibleQuote{他们要彼此出卖，彼此憎恨}（这里祂揭示了内讧）；然后，\BibleQuote{假基督}和假弟兄；然后，\BibleQuote{大多数人的爱要冷淡，}这是一切恶事的原因。\par

你们看到无数种战争，新奇而奇怪吗？然而在这一切之中，甚至更多（因为内讧中还掺杂了亲属的战争），福音却传遍了全地。\BibleQuote{因为福音，}祂说，\BibleQuote{要在全世界宣讲。}\par

那么，那些认为占星术和时间的循环能对抗教会教义的人在哪里？因为谁曾记载另有一位基督出现？这样的事发生过吗？虽然他们虚假地肯定其他事，比如过去了十万年，但他们甚至连这一点也无法捏造。那么你们所说的循环是哪一种？因为再也没有另一个\Place{索多玛}，没有另一个\Place{哈摩辣}，也没有另一个洪水。你们还要戏弄到何时，谈论循环和诞生？\par

那么，怎么说他们说的许多事都应验了呢？因为你剥夺了天主赐予的帮助，出卖了自己，将自己置于祂的眷顾之外；因此恶灵随意转动和扭曲你的事务。\par

但在圣人中并非如此，或者不如说，在我们完全轻视它的罪人中也不是。因为虽然我们的行为难以忍受，但因天主的恩宠，我们严格持守真理的教义，就能超越恶灵的恶意。\par

总之，诞生是什么呢？无非是不义、混乱，以及一切事都在随意进行；或者更准确地说，不仅是随意，而且远远超过此，是愚昧。\par

\Quote{如果没有诞生，那么这般富有的人从何而来？这般贫穷的人又从何而来？}\par

我不知道。因为我要暂时这样同你辩论，教导你不要对一切事都好奇；也不要因此随意而轻率地行事。因为你即使不知道这些，也不应当虚构不存在的事。无知而正确，胜于学习错误。因为不知道原因的人，很快会找到正确的原因；但那个因为不知道真正的原因而虚构不真实原因的人，不容易接受真实的原因；他需要更多的辛劳和努力，才能除去前者。的确，一块平板，如果被擦拭光滑，任何人都可以随意在上面写字；但上面已写有字时，就不再一样了，我们必须先擦去写坏的字。同样，在医生中，什么都不用的人远比使用有害东西的人好；建筑不牢固的人，比根本不建筑的人更糟；正如不长任何东西的土地，远比长荆棘的土地好得多。\par

因此，我们不要急于了解一切，而要忍受对某些事无知，这样当我们找到老师时，不会给他增加双倍的辛劳。或者，许多人常常因为粗心大意地陷入错误意见而变得不可救药。因为先拔掉错误扎根的东西再播种，与在洁净的土地上种植，其辛劳是不同的。在前一种情况下，他必须先推翻，然后再放入其他东西；但在后一种情况下，听讲者已经准备好了。\par

那么，这般富有的人从何而来？我现在要说：许多人因天主的恩赐而致富，也因他的允许而致富。这就是简短而简单的解释。\par

那么，怎样呢？据说，难道他使嫖客、奸夫、与男子行淫者以及滥用自己财产的人致富吗？他不使他们致富，而是允许他们致富；这之间有极大的区别，且是无限的差别。但他为什么容忍呢？因为现在还不是审判的时候，好使每个人按自己的功过得到报应。\par

因为有什么比那个连一点碎屑也不肯给拉匝禄的富人更无价值的呢？然而他比所有人都更悲惨，因为他最终连一滴水也得不到；而这正是因为他虽然富有却残忍。因为如果有两个恶人，在此世没有相同的命运，一个富有，一个贫穷，他们在那里所受的惩罚不会相同，而是更富有的人受罚更重。\par

5. 你岂不见至少这个人，因为\Quote{已享用了他的福}而更可怕地受苦吗？因此，当你看到某个因不义而致富的人昌盛时，你要哀叹、哭泣；因为这份财富对他而言是惩罚的增加。正如那些犯罪多而又不愿悔改的人为自己积蓄忿怒的宝藏；同样，那些不仅不受惩罚反而享福的人，将要受更大的惩罚。\par

如果你愿意，我可以证明这一点，不仅从将来的事，也从今生的事。因为蒙福的达味，当他犯了巴特舍巴的罪而被先知谴责时，他之所以受更严厉的责备，正是因为即使他享有如此的安全，竟然如此。请听天主特别以此责备他：\Quote{我岂没有膏你为王，救你脱离撒乌耳的手，将你主人的一切、以及以色列和犹大的全家都给了你？如果这太少，我还会再加添；你为何行了邪恶在我眼中？} \BibleRef{撒下 12:7} 因为并不是所有的罪都有相同的惩罚，而是许多不同的惩罚，依照时间、依照人、依照他们的地位、依照他们的理解、及其他方面。为使我的话更清楚，我们以一个罪——奸淫为例，看我能找到多少不同的惩罚，不是出于我自己，而是出于圣经。有人在法律之前犯了奸淫，他受的惩罚不同；保禄指示了这一点，\Quote{凡没有法律而犯罪的，也必不按法律而丧亡。} \BibleRef{罗 2:12} 有人在法律之后犯了奸淫，他要受更严重的事：\Quote{凡在法律内犯罪的，也必按法律受审判。} 有人身为司铎而犯奸淫，他从自己的尊位获得惩罚的极大增加。因此，其他妇女因奸淫而被杀，但司铎的女儿们却被火焚；立法者更充分地表明，如果司铎犯此罪，等待他的是何等大的惩罚。因为如果他对女儿因她是司铎的女儿而施加更大的惩罚，更何况对那身负司铎职位的男人自己呢？若奸淫是强暴的，她甚至免于惩罚。若一人富有而犯奸淫，另一人贫穷而犯，这里也有区别。这从我们之前关于达味所说的可以显明。有人在基督来临后犯了奸淫，若他未受洗而死，他将受比所有那些人更重的惩罚。有人在洗礼后犯了奸淫，在这种情况下，这个罪再也得不到任何宽慰。保禄也亲自声明了这事，他说：\Quote{谁干犯梅瑟法律，凭两三个证人，尚且不得怜悯而死；那么，你们想，那践踏了天主子，将那使他成义的盟约之血视为俗物，并侮辱了恩宠之神的人，应受怎样更重的刑罚呢！} 如今有人身为司铎而犯奸淫？这更是罪恶的顶峰。\par

你们看到一条罪就有多少不同的形态吗？一种是在法律之前，另一种是在法律之后，另一种是身负司铎职务者；富妇的，贫妇的，\OriginalTerm{慕道者}{catechumen}的，信女的和司铎女儿的。\par

从知识上，也大有区别；\BibleQuote{因为知道主的旨意，而不遵行的，必受许多鞭打。}\BibleRef{路 12:47} 而在看到榜样后仍犯罪的，招致更大的惩罚。因此祂说：\BibleQuote{你们虽然看见了，后来却不悔改，}\BibleRef{玛 21:32} 虽然你们曾蒙受了极大的关怀。因此祂也同样责备\Place{耶路撒冷}说：\BibleQuote{我多少次愿意聚集你的子女，而你却不愿意！}\BibleRef{路 13:34}\par

至于在奢华中犯罪，这由\Person{拉匝禄}的事迹可见。而地点也使罪更为严重，祂亲自指明说：\BibleQuote{在圣殿和祭台之间。}\BibleRef{玛 23:35}\par

从罪行本身的相当性来看，\Quote{一个人因偷窃被捉，并不稀奇。}又说：\Quote{你杀了你的儿女，这超过了你的所有淫行和可憎之事。}\par

又从人的角度来看：\Quote{若一人得罪了另一人，他们可以为他祈祷；但若他得罪了天主，谁能为他求情呢？}\par

而当任何人比那些远逊于他的人更疏忽时，主在\BibleBook{厄则克耳}{先知书}中也以此谴责他们说：\BibleQuote{你们连列邦的法度也没有遵行。}\BibleRef{则 5:7}\par

而当一个人即使他人的榜样也不能使他醒悟时，经上说：\Quote{她看见她的姊妹，并称她为义。}\par

而当一个人获得了更丰厚的眷顾时，祂说：\BibleQuote{因为如果在\Place{提洛}和\Place{漆冬}所行了这些大能的事，他们早就悔改了；但是为\Place{提洛}和\Place{漆冬}，将比那座城更容易忍受。}\BibleRef{玛 11:21}\par

你们看到完全的精确性了吗？所有人因同样的罪所受的惩罚并不相同。因为当我们获得长期忍耐的恩惠却一无所用时，我们将遭受更糟的事。\Person{保禄}也指明这一点，他说：\BibleQuote{但因为你顽固不化和不肯悔改的心，你为自己积蓄忿怒。}\BibleRef{罗 2:5}\par

既然知道了这些事，让我们不要感到冒犯，也不要因所发生的事而慌乱，更不要让思想的风暴降临我们身上；而要顺从天主的眷顾，专注于德行，逃避恶行，好使我们也能获得将来的美物，借着我们的主耶稣基督的恩宠与对人的爱，愿光荣归于圣父，偕同圣神，从今直到永远，世世无尽。阿们。\par

\SourceRecord{Translated by George Prevost and revised by M.B. Riddle.；Nicene and Post-Nicene Fathers, First Series；Vol. 10.；Edited by Philip Schaff.；Buffalo, NY: Christian Literature Publishing Co.,；1888.；<http://www.newadvent.org/fathers/200175.htm>.}

% structure: p0001 source=h1 target=section reason=page-title

\section{玛窦福音讲道集 第七十六讲}

\Emph{\BibleRef{玛 24:16-18}}\par

\Emph{\Quote{那时，在犹太的，要逃到山上；在屋顶上的，不要下来取家里的东西；在田里的，也不要回去取衣裳。}}\par

论述了将要临到那城的灾祸、宗徒们的考验、以及他们必不被征服而要传遍全世界之后，他再次提及犹太人的灾难，表明当一方因教训了全世界而变得荣耀时，另一方却要处于灾难之中。\par

请看他是如何借看似微小的事描述这场战争，以显明其难以忍受。因为他说：\Quote{那时}，他说：\Quote{在犹太的，要逃到山上。}那时？何时？当这些事发生时，\Quote{当那造成荒凉的可憎之物站在圣地时。}由此，他似乎是在说军队。因此，他说：那时你们要逃跑，因为此后你们再没有得救的希望了。\par

因为过去在惨烈的战争中，他们常常得以恢复，如在山乃黑黎布手下，又如在安提约古手下（因为那时也有军队进入，圣殿已被占据，但玛加伯人集结起来扭转了局面）；为了避免他们现在也猜想会有这样的转变，他禁止一切此类想法。他说：今后能光着身子逃脱已是万幸。因此，他甚至不准在屋顶上的人进屋取衣服，以示灾祸不可避免，灾难无穷，卷入其中的人必死无疑。为此他又补充说：在田里的，也不要回去取衣裳。因为屋里的人尚且要逃，何况在屋外的人，更不该回屋避难。\par

\Quote{怀孕的和哺乳的有祸了}\BibleRef{玛 24:19}，对前者是因为她们更为迟钝，因身孕沉重而难以逃跑；对后者是因为她们被对孩子的感情所束缚，无法拯救乳儿。因为钱财容易轻视，也容易赚取和置办衣服，但天性的束缚谁能逃脱？孕妇怎能行动敏捷？哺乳的母亲怎能忽视自己所生的孩子？\par

接着，为再次显示灾难的深重，他说：\Quote{你们要祈祷，叫你们的逃亡不在冬天，也不在安息日。因为那时必有大灾难，从世界的起头直到如今没有过的，将来也必没有。}\BibleRef{玛 24:20}\par

你看见他的话是对犹太人说的，是在讲那将要临到他们的灾祸吗？因为宗徒们当然不必守安息日，也不必在韦斯巴芗做那些事的时候留在那里。事实上，他们中的大部分已经离世。即便有存留的，那时也住在世界其他地方。\par

但为什么\Quote{既不在冬天，也不在安息日？}不在冬天，是因季节的困难；不在安息日，是因律法的绝对权威。因为他们需要逃跑，而且是极速逃跑，但犹太人不敢在安息日逃跑（因律法禁止），而在冬天逃跑也不容易，因此他说：\Quote{你们要祈祷}；\Quote{因为那时必有大灾难，是从前没有的，将来也必没有。}\par

不要以为这话是夸张；要研读若瑟夫的著作，便知这话的真实。因为谁也不能说，他作为一个信徒，为了证实基督的话而夸大了悲惨的历史。他原是犹太人，而且是坚定的犹太人，非常热心，是生活在基督降临之后的人。\par

那么此人说了什么？他说那些恐怖超过了所有悲剧，从未有过的灾难临到那民族。饥荒之严重，以至于母亲们为了吃自己的孩子而争斗，甚至为此打仗；他还说许多人死后肚子被剖开。\par

因此，我很想问犹太人：为何从天主而来的、不可忍受的、比从前一切灾难更重的愤怒降在他们身上，不仅是在犹太，而且在全世界任何地方？难道不正是因十字架的行为和这拒绝？所有人都会这么说，而且事实本身在所有人和所有事之前已证明。\par

但请注意，我恳求你们，灾难的极大性：不仅与以往相比显得更重，而且与将来全部时间相比也是如此。因为在整个世界，无论过去还是将来，没有人能说有过这样的灾难。这很自然，因为无论过去还是将来，从未有人犯过如此邪恶可怕的罪行。因此他说：\Quote{那灾难是前所未有，以后也必没有的。}\par

\Quote{若不减少那些日子，凡有血气的总没有一个得救的；但为了被选的人，那些日子必会减少。}他以此表明他们应受比先前所说的更重的惩罚，现在他指的是战争和围困的日子。他所说的大意是：如果罗马人对那城的战争持续下去，所有的犹太人（这里\Quote{凡有血气的}指犹太血肉）都将灭亡，无论是在外的还是在家的。因为他们不仅与犹太地的人作战，而且因憎恨他们，还将各地散居的犹太人逐出并放逐。\par

2. 但祂在此所说的\Quote{被选者}是谁？是指他们中间被围困的信徒。为免犹太人会说，因着福音和敬拜基督，这些灾祸才发生，祂便表明，信徒绝非原因；若非他们，所有人都已完全灭亡了。因为，如果天主容许战争拖延，连犹太人的残存者也不会留下；但为使那些已成为信徒的犹太人不要与不信的犹太人一同灭亡，祂迅速平息了战事，结束了战争。因此祂说：\BibleQuote{但为了被选者的缘故，那些日子必会缩短。}祂说这些话，是为了鼓励那些被围困在他们中间的人，让他们得到喘息，免得他们恐惧，好像要与他们一同灭亡。如果在这里祂对他们如此关怀，以致为了他们的缘故，其他人也得救，并且为了基督徒的缘故，犹太人中留下了残存者，那么到了他们得冠冕的时候，他们的尊荣将会何等大呢？\par

借此祂也鼓励他们不要为自己的危险而苦恼，因为这些他人正遭受如此苦难，且毫无益处，只招致灾祸归于自己头上。\par

但祂不仅鼓励他们，还秘密且不引人注意地引导他们脱离犹太人的习俗。因为若此后没有改变，圣殿也不得存留，显然法律也必被废止。\par

然而，祂没有公开说这话，而是藉着他们的彻底毁灭暗中暗示。但祂没有公开说，免得在时机成熟前惊动他们。所以，祂起初也没有主动谈论这些事；而是先为城市哀哭，然后迫使他们向祂展示石头并询问祂，为的是好像在回答他们的问题，祂可能向他们预告一切将要发生的事。\par

但请你留意圣神的安排：\Person{若望}并未写下这些事中的任何一件，免得他似乎是依据所发生之事的历史而写（因为他确实在城陷后活了很久）；而他们，那些在城陷前去世、不曾看见这些事的人，却写下这些事，为使预言的能力从各方面清楚闪耀。\par

\BibleQuote{那时，若有人对你们说：看，基督在这里，或在那里，不要相信；因为将有假基督和假先知兴起，显大征兆和奇事，以致如果可能，连被选者也要被迷惑。看，我预先告诉了你们。所以，如果他们向你们说：看，祂在旷野里，你们不要出去；看，祂在内室里，你们不要相信。因为如同闪电从东方发出，直照到西方，人子的来临也要这样。因为无论那里有尸体，鹰也必聚集在那里。}\par

说完了关于\Place{耶路撒冷}的事，祂便转到自己的来临上，并述说其征兆，这不仅为他们的益处，也为我们以及所有以后来的人。\par

\Quote{那时。}何时？这里，正如我常说的，\Quote{那时}一词并不涉及时间上与前述之事的顺序关系。至少，当祂想要表达时间的联系时，祂加了\BibleQuote{那些日子的灾难一过，}\BibleRef{玛 24:29}但这里并非如此，\Quote{那时}的意思不是指这些事之后立即发生的事，而是指祂将要说的那些事要发生的时候。同样，当经上说\BibleQuote{那时，\Person{若翰洗者}出现，}他也不是指紧接着之后的时间，而是指许多年之后，以及祂将要说的那些事发生的时间。因为事实上，在讲述了\Person{耶稣}的诞生、贤士的来临和\Person{黑落德}的死亡之后，祂立刻说：\BibleQuote{那时，\Person{若翰洗者}出现，}尽管中间过了三十年。但这是圣经的习惯，我指的是使用这种叙述方式。所以这里也一样，在略过了从\Place{耶路撒冷}被毁直到末日预兆的所有中间时间之后，祂便谈到末日之前的时间。因此祂说：\BibleQuote{那时，若有人对你们说：看，基督在这里，或在那里，不要相信。}\par

祂一时以地点来保护他们，提及祂第二次来临的标记和迷惑者的迹象。因为不像祂第一次来临时出现在\Place{白冷}，在世界的一个小角落，起初无人认识祂那样，祂说那时也将如此；而是公开地、带着一切情况，无需任何人来告诉这些事。这是一个不小的迹象，表明祂不会秘密地来临。\par

但请注意，这里祂没有提到战争（因为祂正在解释关于祂来临的道理），而是提到那些试图欺骗的人。因为在宗徒时代，有人欺骗了群众，\BibleQuote{因为他们要来，}祂说，\BibleQuote{并要迷惑许多人；}\BibleRef{玛 24:11}在祂第二次来临之前，还有人会这样做，他们比前者更严厉。\BibleQuote{因为他们要显大征兆和奇事，}祂说，\BibleQuote{以致如果可能，连被选者也要被迷惑。}\BibleRef{玛 24:24}这里祂是在谈论假基督，并指出有些人也将为他服务。\Person{保禄}也这样谈到他。称他为\BibleQuote{罪恶之人}、\BibleQuote{丧亡之子}后，他又加上了\BibleQuote{他的来临，是靠撒殚的作为，具有各种德能、征兆和欺骗的奇事，以及各种不义的欺骗，在那些丧亡的人身上。}\BibleRef{得后 2:9}\par

请看祂如何保护他们；\BibleQuote{不要到旷野去，不要进入内室。}祂没有说：\Quote{去，不要相信；}而是说：\Quote{不要去，也不要往那里去。}因为那时迷惑很大，因为连欺骗性的奇迹也将被行出。\par

3. 祂在告诉了他们敌基督如何来临，例如，将在某个地方之后，又说祂自己如何来临。那么祂自己如何来临呢？\BibleQuote{闪电从东方发出，直照到西方，人子的来临也要这样。因为哪里有尸体，哪里就有鹰聚集。}\BibleRef{玛 24:27}\par

那么闪电怎样照耀呢？它不需要有人谈论，不需要传令官，而是在一瞬间，整个世界上，甚至对坐在屋里和在内室中的人都显现。那来临也将如此，因祂荣耀的光辉而同时处处显现。但祂还提到了另一个标记：\BibleQuote{哪里有尸体，哪里也有鹰聚集。}意指众天使、殉道者、诸圣人的群体。\par

然后祂讲述了可怕的异兆。这些异兆是什么？祂说：\BibleQuote{那些日子的灾难一过，太阳就要昏暗。}\BibleRef{玛 24:29}祂说的是哪些日子的灾难？是敌基督和假先知的吗？因为有许多欺骗者，将有大灾难。但这灾难不会延长时日。因为如果为了选民的缘故，犹太人的战争被缩短了，那么为了这些人的缘故，这诱惑更应该被限制。因此，祂没有说\Quote{灾难以后}，而是说\Emph{立刻}\BibleQuote{那些日子的灾难以后，太阳就要昏暗}，因为几乎所有的事都在同一时间发生。假先知和假基督会来造成混乱，然后祂自己立刻就到。因为那时世界上将有大动荡。\par

但祂怎样来临呢？受造界本身那时要改变，因为\BibleQuote{太阳要昏暗}，不是被消灭，而是被祂临在的光辉所胜过；星辰要坠落，因为从今以后，既然没有黑夜，它们还有什么用处呢？\BibleQuote{诸天的权势也要震动}，很可能，看到如此巨大的变化发生。因为如果当初星辰被造时，它们曾战栗惊奇（经上说：\Quote{当星辰被造时，所有天使都高声赞美我}）；那么，当看到万物都在改变，它们的同僚要交账，全世界要站在那可怕的审判宝座前，从亚当直至祂来临的人都必须为一切行为交账时，它们怎能不战栗震动呢？\par

\BibleQuote{那时，人子的记号要出现在天上；}\BibleRef{玛 24:30}也就是说，十字架比太阳更明亮，因为太阳将要昏暗并隐藏自己，而十字架将显现——如果它不比太阳的光芒更明亮得多，它就不会显现。但为什么这记号要出现呢？为了更充分地堵住犹太人的无耻。因为基督手持十字架作为最大的辩护，就这样来到审判宝座前，不仅显示祂的伤痕，也显示祂耻辱的死亡。\BibleQuote{那时，地上所有的支派都要哀号}，因为当他们看见十字架时，就不需要指控了；他们要哀号，因为他们的死亡对祂毫无益处；因为他们钉死了他们本该朝拜的那一位。\par

你看到祂多么可怕地描绘了祂的来临吗？祂怎样激起了门徒的心？为此，我还要说，祂先把悲伤的事放在前面，然后才说好事，这样祂也能安慰和鼓舞他们。祂提醒他们想起祂的苦难、复活，并以荣耀的光景提及祂的十字架，使他们不致羞愧或悲伤，因为祂那时要把它作为记号带出来。另一位经上说：\BibleQuote{他们要瞻望他们所刺透的那一位。}因此他们要哀号，当他们看见这就是祂时。\par

既然祂提到了十字架，祂又接着说：\BibleQuote{他们要看见人子来临，不再是在十字架上，而是乘着天上的云，带着威能和大光荣。}\BibleRef{玛 24:30}\par

因为，祂的意思是，不要因你听到十字架就以为又是可悲的事，因为祂要带着威能和大光荣来临。但祂带来十字架，是为了让他们的罪自我定罪，就像一个人被石头击中，他拿出石头或沾血的衣服一样。祂乘云而来，正如祂被接升天一样。众支派看见这些事就哀号。然而，恐惧在他们那里不会止于哀号；哀号是为了让他们从内心发出判决，定自己的罪。\par

然后又说：\BibleQuote{祂要派遣祂的天使，用大号角，从四方、从天这边到天那边，聚集祂的选民。}\BibleRef{玛 24:31}\par

当你听到这一点时，要想到那些留下来的人的惩罚。因为他们不仅承受先前的惩罚，还要承受这个。如同上面祂说，他们要说：\BibleQuote{奉主名而来的，当受赞颂；}\BibleRef{玛 22:39}这里却说他们要哀号。因为祂已向他们讲述了可怕的战争，好让他们知道，除了今世可怕的事，那里的折磨也在等待着他们；祂让他们哀号，与选民分离，并被交付给地狱；这又一次唤醒了门徒，并表明他们将从多少灾祸中被拯救，并享受多少福乐。\par

4. 为何现今祂藉天使召唤他们，如果祂如此公开地来临？这也是为了以此方式尊荣他们。但\Person{保禄}却说，他们\Quote{将被提到云中}。当祂论及复活时，也说了这话。因为经上说：\Quote{\BibleRef{得前 4:16}主亲自}，又经上说：\Quote{祂必带着号令，总领天使的声音，从天降下}。因此，当复活后，天使将聚集他们，聚集后，云彩将接他们上升；这一切都在一刹那，一瞬间完成。因为并非祂停留在高处呼召他们，而是祂亲自带着号声降临。号筒和声音有何含义？它们是为唤醒、为喜乐，以显明那时所发生之事的奇妙，并为留下的人带来悲伤。\par

我有祸了，为那可怕的日子！因为虽然我们听到这些事应当喜乐，我们却感到痛苦、沮丧，面容悲伤。或者只是我如此感觉，而你们听到这些事却喜乐？因为至少对我来说，当说到这些事时，有一种战栗临到我，我痛苦哀哭，从心底深深叹息。因为在这些事中我没有分，而是在随后的话中、对童女们、对埋藏所领受塔冷通的人、对恶仆所说的话中。因此我哭泣，想到我们将被从何等荣耀中逐出，从何等福禄的期盼中逐出，并且是永久地、永远地，只为节省一点小劳苦。因为如果这真是巨大的劳苦、严苛的律法，我们仍应做一切；尽管如此，许多松懈的人似乎至少有一些借口——虽是个糟糕的借口——但他们似乎有些理由，说劳苦巨大、时间无尽、负担难忍；但现在我们提不出这样的反对；那时，这情形将最甚地啃噬我们，不亚于地狱，当我们因缺乏一点努力、一点劳苦，而失去了天堂和说不尽的福乐。因为时间短暂，劳苦微小，但我们却软弱懈怠。你在世上奋斗，冠冕在天上；你受人惩罚，却受天主尊荣；赛程只有两天，赏报却永无穷尽；斗争是可朽坏的身体，赏报却是不朽坏的。\par

此外，我们也应思考另一件事：即使我们不愿为\Person{基督}忍受任何痛苦，我们在其他方面也必定会忍受。因为即使你不为\Person{基督}而死，你也不会不朽；即使你不为\Person{基督}抛弃财富，你也不会带着它们离开。祂向你要求的这些事，即使祂不要求，你也得放弃，因为你是必死的；祂愿你自愿去做那些你必然要做的事。祂只要求你加上这一点：为祂而做；因为这些事临到人并消逝，是出于自然的必然。你看，争战多么容易？你必然要忍受的，就为我而忍受吧；只加上这一点，我就满足。你打算借给别人的金子，请借给我，既得利更多，又更安全。你的身体，你将要为别人作战的身体，让它为我作战吧，因为我以极其丰厚的酬报超过你的劳苦。然而，在一切其他事情上，你宁愿选择那给你更多报酬的人，无论是在借贷、经商还是作战上；但唯独\Person{基督}，当祂赐予更多、远超一切时，你却不接纳。这如此深重的敌意是什么？这如此强烈的仇恨是什么？当你宁愿选择人而不选择人的理由，却不足以使你宁愿选择天主而不选择人时，你还有什么借口或辩护呢？\par

为何你将财宝托付给大地？祂说：\Quote{把它交在我手中。}大地的主岂不是比大地更值得信任吗？大地确实归还你所存放的，尽管常常连这也做不到，但祂为你保管还给你酬报？因为祂极其爱我们。因此，你若愿意借贷，祂已预备；若愿意播种，祂接受；若愿意建造，祂吸引你到祂那里，说：在我的地方建造。为何你跑向贫穷、乞讨的人，他们为微利给你带来大麻烦？然而，即使听到这些，我们也不下定决心，而是哪里有争斗、战争、激烈斗争、审判、诉讼、诬告，我们就奔向哪里。\par

5. 祂岂不应当公正地弃绝我们、惩罚我们吗？因为祂将自己一切都给了我们，我们却抗拒祂？这显然众人皆知。因为无论你愿装饰自己，祂说：\Quote{用我的装饰品}；或愿武装自己，\Quote{用我的武器}；或愿穿衣，\Quote{用我的衣服}；或愿进食，\Quote{在我的筵席上}；或愿旅行，\Quote{在我的路上}；或愿继承产业，\Quote{我的产业}；或愿进入一国，\Quote{我所建造营造的城}；或愿建造房屋，\Quote{在我的帐幕中}。\Quote{因为，我不仅不向你要求我所赐之物的报酬，甚至为使我自己欠你报酬，只要你愿意使用我所有的一切。}谁能比得上这样的慷慨？\Quote{我是父亲，我是兄弟，我是新郎，我是居所，我是食物，我是衣服，我是根，我是基础，凡你愿意的，我都是。}\Quote{你什么都不需要，我甚至要做仆人，因为我来了不是为受服事，而是为服事；我是朋友、肢体、头、兄弟、姊妹、母亲；我是一切，只要你紧紧依附我。我为你成了贫穷，为你成了流浪者，为你上了十字架，为你进入了坟墓；在上边我为你向圣父转求；在地上我为你而来，作为我圣父的使者。你对我是一切，兄弟、共同继承人、朋友、肢体。}你还要什么？为什么你要避开爱你的一位？为什么你要为世界劳苦？为什么你要把水倒在破裂的水池里？因为为现世生命劳苦正是如此。为什么你要把羊毛梳向火中？为什么你\Quote{打空气？}\BibleRef{格前 9:26}为什么你\Quote{空跑？}\BibleRef{迦 2:2}\par

难道每一种技艺都有其目的？这显然人人都知道。你也当显明你世俗热心的目的。但你做不到；因为\Quote{虚而又虚，万事皆虚。}\BibleRef{训 1:2}让我们到坟墓那里去；把你的父亲指给我看；把你的妻子指给我看。那个穿金戴银的人在哪儿？那个坐战车的人在哪儿？那个拥有军队、佩带腰带、有传令官的人在哪儿？那个杀戮这些、监禁那些的人？那个随意处死、随意释放的人？我只见骨头、虫子和蛛网；这一切都是尘土，一切都是虚构，一切都是梦影，都是空谈，都是图画，甚至还不如图画。因为图画我们至少看到相似，但这里连相似都没有。\par

但愿恶事就此止步。因为现世的荣誉、奢侈和显赫都以影子、言语告终；但它们的后果却不再限于影子和言语，而是持续，并随我们转到别处，且要显明给众人：抢劫、贪得、淫乱、奸淫、不可胜数的可怕之事；这些不是用比喻，也不在灰烬中，而是记载在天上，包括言语和行为。\par

那么我们将以怎样的眼睛瞻仰基督？因为如果有人因自觉得罪了父亲而不能忍受见父亲的面，那么对于那位在宽容上无限超越父亲的祂，我们将如何观看？我们将如何忍受？因为我们确实要站在基督的审判台前，那里将有一场对万事的严格审查。\par

但若有人不信将来的审判，让他看看这里的事：那些在监狱里的，在矿上的，在粪堆上的，附魔的，疯狂的人，与不治之症搏斗的，与持续贫穷抗争的，生活在饥荒中的，被无法补救的灾祸刺透的，被掳的。因为这些人不会在这里受这些苦，除非报应和惩罚也等待着所有其他犯有同样罪过的人。如果其余的人在这里未曾经历什么，你应当把这事实视为一个标记：在我们离开这里之后，必定还有事要发生。因为同一位万有的天主不会向一些人施行报复，而让其他犯有同样或更严重罪行的人不受惩罚，除非祂打算在那里给他们带来一些惩罚。\par

那么，借着这些论据和这些例证，让我们也谦卑自下；让那些顽固不信审判的人从此相信，并成为更好的人；使我们在此世活出相称于天国的生活，好能获得将来的美善，借着我们的主耶稣基督的恩宠和爱人，愿光荣归于祂，至于无穷之世。阿们。\par

\SourceRecord{Translated by George Prevost and revised by M.B. Riddle.；Nicene and Post-Nicene Fathers, First Series；Vol. 10.；Edited by Philip Schaff.；Buffalo, NY: Christian Literature Publishing Co.,；1888.；<http://www.newadvent.org/fathers/200176.htm>.}

% structure: p0001 source=h1 target=section reason=page-title

\section{玛窦福音讲道集 第七十七篇}

\BibleRef{玛 24:33-34}.\par

\BibleQuote{你们该由无花果树学个比喻：当它的枝条已经嫩了，长出叶子时，你们就知道夏天近了。同样，你们看见这一切事，也该知道它近了，就在门口。}\par

因为祂曾说过：\BibleQuote{那些日子的灾难一过，} 但门徒们寻求这事在多少时间以后发生，并特别想知道那日子，因此祂又加上无花果树的比喻，表明间隔不大了，祂的来临也会紧接而来。祂不仅用比喻，也用随后的话宣称：\BibleQuote{要知道它近了，就在门口。}\par

借此祂也预告另一件事，即一种属灵的夏天，以及那日（现世风暴之后）为义人预备的平静；但对罪人则相反，夏天之后是冬天，这在随后的话中表明，说那日要临到他们，当他们正生活奢侈时。\par

但祂提出这关于无花果树的比喻，不仅是为了表明间隔，因为也能用其他方式摆在他们面前；也是为了借此坚固祂的话，保证这事必定如此成就。因为正如这无花果树的事是必然的，那事也是必然的。因此，每当祂要谈论必定成就的事，祂都引用自然的必然律，不仅祂自己如此，蒙福的保禄也效法祂。为此，当祂论及自己的复活时，祂说：\BibleQuote{麦子若不落在地里死去，它仍然是一粒；但如果死去，就结出许多果实。} \BibleRef{若 12:24} 蒙福的保禄也受了教导，用同样的比喻说：\BibleQuote{愚昧的人哪！你所播的种子若不先死，就不会活过来。} \BibleRef{格前 15:36}\par

此后，为使门徒们不再立刻回到这个问题，问说：\Quote{何时？} 祂便提起已经说过的话，说：\BibleQuote{我实在告诉你们：这一世代决不会过去，直到这一切事都发生！} 这一切事？请问是什么？那些关于耶路撒冷的事，关于战争、饥荒、瘟疫、地震、假基督、假先知、福音传遍各处、叛乱、骚乱，以及我们说过直到祂来临前要发生的一切。那么，有人会问，祂为何说\Quote{这一世代？} 祂不是指当时的世代，而是指信徒的世代。因为祂惯常不仅以时间，也以宗教敬拜的方式和实践来区分一个世代，正如祂说：\BibleQuote{这是寻求主的世代。}\par

因为祂前面所说：\BibleQuote{这一切都必须发生，} \BibleRef{玛 24:6} 又说：\BibleQuote{福音必传遍，} \BibleRef{玛 24:14} 祂在此也宣告说：这一切事必将发生，而信者的世代将存留，不被上述任何事所剪除。因为耶路撒冷将毁灭，犹太人大部分将被消灭，但这一世代什么也不能克服它，无论是饥荒、瘟疫、地震、战争的骚乱、假基督、假先知、欺骗者、叛徒、引人跌倒者、假弟兄，或其他任何类似的诱惑。\par

为在信德上更进一步引导他们，祂说：\BibleQuote{天地要过去，但我的话决不会过去。} \BibleRef{玛 24:35} 意思是，这些稳固、固定、不动摇的形体被抹去，比我的话中任何一句落空更容易。谁若反驳这些事，让他来检验祂的话，当他发现这些话是真实的（他必定会发现），就让他从已发生的事相信那将要发生的事，并让他殷勤查考一切，他将看到实际事件为预言的真理作证。祂引用了元素（天地），一方面宣告教会比天地更尊贵，另一方面也借此表明祂自己是万有的创造者。因为祂在谈论终结，这是许多人所不信的，祂便引用天地，表明祂不可言喻的大能，并以极大的权威显示祂是万有的主，借此使祂的话即使对多疑的人也值得相信。\par

\BibleQuote{但那日子和时辰，没有人知道，连天上的天使也不知道，圣子也不知道，唯独圣父知道。}祂说连天使也不知道，以此堵住他们的口，使他们不再寻求学习天使所不知道的事；又说\BibleQuote{圣子也不知道}，不仅禁止他们学习，甚至连探问也不许。为证明祂这样说正是为此，请看祂复活后，见他们过于好奇时，是如何更坚决地堵住他们的口的。因为现在祂确实提到了无数确实的征兆；但那时祂仅仅说：\BibleQuote{不是你们可以知道时候或日期。}又为使他们不致说\Quote{我们被逼入困境，我们完全被鄙视，我们连这一点也不配知道}，祂说：\BibleQuote{那是圣父凭着自己的权柄所定的。}祂这样做，是因为祂极其关心尊荣他们，不愿向他们隐瞒任何事。因此祂将此归于圣父，既为使这事可畏，也为将祂所说的事排除在他们的探问之外。因为若不是这样，而是祂不知道，那么祂何时会知道呢？难道要和我们一起知道吗？但谁会这样说呢？祂清楚地认识圣父，如同清楚地认识圣子一样；祂岂会不知道日子？况且\BibleQuote{圣神洞察天主深奥的事}（\BibleRef{格前 2:10}），祂岂会连审判的时间也不知道？但祂知道该如何审判，对每个人的隐秘完全知晓；而比这平常得多的事，祂反而不知道？如果\BibleQuote{万物是借着祂造成的，凡受造的没有一样不是借着祂造的}，祂岂会不知道日子？因为祂创造了世界，显然也创造了时间；如果创造了时间，也就创造了那日子。那么祂怎么会不知道祂所创造的呢？\par

2. 你们确实说知道祂的本体，但圣子（祂常在圣父怀里）却连日子也不知道？然而祂的本体比日子大得多，甚至无限地大。那么，你们既然为自己选取了更大的事，为何连较小的也不允许圣子知道呢？\BibleQuote{在祂内蕴藏着智慧和知识的一切宝藏}（\BibleRef{哥 2:3}）。但你们既不知道天主的本体是什么（虽然你们千万次如此狂言），圣子也并非不知道那日子，而是完全确定地知道。\par

为此我说，当祂说尽了一切、包括时间和时期，并将之带到门前（\BibleQuote{近了，}祂说，\BibleQuote{就在门口}）时，祂就沉默不言那日子。因为若你们追问那日子和时辰，你们不会从我这里听见，祂说；但若问时期和预兆，我毫不隐瞒，必详细告诉你们一切。\par

因为事实上我并非不知道它，我已通过许多事表明：我提到了间隔、将要发生的一切、从现今到那日子本身是多么短暂（因为无花果树之喻正指明了这一点），并领你们直到门口；即使我不为你们开门，这也是为你们的好处。\par

为使你们也从另一件事知道，沉默并非出于祂的无知，请看，连同我们已提到的，祂还提出了另一个标记。\BibleQuote{就如诺厄的日子怎样，人们吃喝婚嫁，直到洪水来了，把他们全都冲去；人子的来临也要这样。}祂说这些话，是表明祂要突然、出乎意料地来临，而且多数人正过着奢侈的生活。因为保禄也这样说，他写道：\BibleQuote{当人们正说平安稳妥时，毁灭会突然临到他们；}为表明多么出乎意料，祂说，\BibleQuote{如同产痛临到孕妇}（\BibleRef{得前 5:3}）。那么祂怎么说\BibleQuote{那些日子一过，灾难就来了}呢？因为如果那时有奢侈、平安、稳妥，如保禄所说的，祂怎能说\BibleQuote{那些日子一过，灾难就来了}？如果有奢侈，怎能还有灾难？奢侈是对那些麻木不仁、自认为平安的人而言的。因此祂没有说\BibleQuote{当有平安时}，而是\BibleQuote{当人们正说平安稳妥时}，表明他们的麻木如同诺厄时代的人；因为在那般的灾祸中，他们仍过着奢侈的生活。\par

但义人却不如此，他们生活在苦难和沮丧中。由此祂表明，当假基督来时，罪人及那些对自身得救绝望的人，会更加热衷于不正当的享乐。那时必有贪食、宴乐和酗酒。因此祂特别举出一个与此相合的例证。因为正如当方舟建造时，他们不信，祂说；但当方舟立在他们中间，预先宣告将要来临的灾祸时，他们看见后仍纵情享乐，仿佛没有任何可怕的事要发生；同样，现在假基督将要出现，其后便是结局，以及结局中的刑罚和不可容忍的报复；但那些被邪恶之醉所控制的人，甚至不会察觉正要发生之事的可怕性质。因此保禄也说：\BibleQuote{如同产痛临到孕妇}，那些可怕而不可治愈的灾祸也要同样临到他们。\par

为什么祂不提索多玛的灾祸？祂愿意引入一个涵盖所有人的例子，一个在预言之后仍不被相信的例子。因此，正如大多数人对未来之事不信，祂用过去之事来证实它们，使他们内心恐惧。同时，祂也表明，过去之事也是祂所行的。接着，祂又设下另一个标记，通过这些事，祂表明祂并非不知道那日子。这标记是什么？\BibleQuote{那时，两个人在田里，一个被接去，一个被留下。两个女人在推磨，一个被接去，一个被留下。所以，你们要警醒，因为你们不知道你们的主什么时候来到。}这一切都证明祂知道，并旨在使他们不再追问。为此，祂也提到\Person{诺厄}的日子，为此祂也说：\BibleQuote{两个人躺在床上}，意指祂会如此出乎意料地降临，当人们毫无思想之时，以及\BibleQuote{两个女人在推磨}，这本身也不是深思熟虑之人的工作。\par

与此一起，祂宣告，仆人和主人都将被接去和被留下，无论是安闲的，还是劳碌的，从这一等级到另一等级；正如在旧约中祂说：\BibleQuote{从坐宝座的到在磨坊的俘虏女子。}因为祂曾说，富人很难得救，祂显示，连这些也并非完全丧亡，穷人也不是全都能得救，而是从这些人中，也从那些人中，有人得救，有人丧亡。\par

而在我看来，祂似乎宣告，降临将在夜间发生。因为\Person{路加}也这样说。\BibleRef{路 17:34}你们看，祂多么准确地知道一切！\par

此后，为使众人不再追问，祂又加上：\BibleQuote{所以，你们要警醒，因为你们不知道你们的主什么时候来到。}\BibleRef{玛 24:42}祂没有说：\Quote{我不知道}，而是说：\Quote{你们不知道。}因为当祂将他们带近那时刻，并安置在那里时，祂再次阻止他们追问，出于希望他们应当时常奋力。因此祂说：\Quote{警醒}，显示正是因为这一点，祂没有告诉。\par

\BibleQuote{但你们要知道，如果家主知道贼在什么时候来，他必会警醒，不容自己的房屋被挖开。所以你们也要准备好，因为在你们想不到的时候，人子就来了。}\par

为此祂不告诉他们，好使他们警醒，常常准备好；因此祂说：当你们不期待时，祂就来了，愿意他们迫切等候，并持续行善。\par

但祂的意思是这样：如果普通人知道自己何时会死，他们必定会在那一刻奋力拼搏。\par

3. 因此，为使他们在那一刻奋力，而不只是在那时，祂便不告诉他们共同的时间，也不告诉每个人各自的时间，愿意他们永远期盼这个，从而始终奋力。因此，祂使每个人的生命终结也变得不确定。\par

此后，祂公开称自己为主，此前从未如此明确说过。但在这里，在我看来，祂也使漫不经心者蒙羞，因为他们对灵魂的关注甚至不及期待盗贼的人对金钱的关注。的确，他们期待时便警醒，不容家中任何东西被夺走；但你们，虽然知道祂要来，而且必定要来，却仍不警醒，也不准备好，以免毫无准备地被带走。因此，那日子对沉睡者来说是毁灭。因为正如那人若知道就会逃脱，同样，你们若准备好，便可自由逃脱。\par

然后，当祂谈到审判时，便将话题转向教师，谈论惩罚与荣誉；先提到行善的人，最后以那些继续犯罪的人结束，使祂的讲道以警醒之言收尾。\par

因此祂首先说：\BibleQuote{那么，谁是那忠信而精明的仆人，主人将派他管理家仆，按时分粮给他们？当主人回来时，看见他如此行事，那仆人就有福了。我实在告诉你们，主人要派他管理一切财产。}\par

请告诉我，这也是无知者的话吗？因为如果因为祂说：\BibleQuote{圣子也不知道}，你们就说祂不知道；正如祂说：\BibleQuote{那么是谁呢？}你们会说什么？你们会说祂也不知道这个吗？绝对不可有此念！因为就连疯狂之人也不会这样说。然而在第一种情况中，或许可以指出一个原因；但在这里连这个也没有。而且当祂说：\BibleQuote{伯多禄，你爱我吗？}\BibleRef{若 21:16}问的时候，祂连这个也不知道吗？或者当祂说：\BibleQuote{你们把他放在哪里了？}\BibleRef{若 11:34}呢？\par

连圣父也要被发现有讲论类似的事。因为祂自己同样说：\BibleQuote{\Person{亚当}，你在哪里？}\BibleRef{创 3:9}以及\BibleQuote{\Place{索多玛}和\Place{哈摩辣}的喊声，在我面前越来越大。因此我要下去，看看他们的行为是否真的如同那达到我面前的喊声；如果不是，我也就知道了。}\BibleRef{创 18:20}另外，祂还说：\BibleQuote{他们或听，或不听，}\BibleRef{则 2:5}而在福音中也说：\Quote{他们或许会尊敬我的儿子。}——这都是无知的表达。但祂说这些话并非出于无知，而是为成就合乎祂身分的旨意：在\Person{亚当}的事上，为促使他替自己的罪过辩护；在\Place{索多玛}人的事上，为教导我们绝不可妄下断语，直到我们亲临现场；在先知的事上，为免得预言在愚人看来成为一种迫使人违命的束缚；在福音的比喻中，则为显示他们本应这样做，该当尊敬圣子。但在此处，既为使他们不再好奇或过度干预，也为表明这是珍稀宝贵的事。请看，这句话指示了多大的无知，如果祂连那被委派的人都不认识。祂的确祝福他，因为祂说：\Quote{那仆人有福了}；但祂没有说这是谁。因祂说：\Quote{那一个人是谁，他的主人要派他管理？}又说：\Quote{主人看见他如此做，那人就有福了。}\par

但这些话不仅指钱财，也包括言语、权力、神恩以及每个人受托的每一项管理职务。这个比喻也适用于国家中的统治者，因为人人都该善用自己所有以谋公共福利。若你有智慧、能力、财富或任何东西，不可用它来伤害同仆，也不可自取灭亡。因此，主要求他两样东西：智慧和忠信，因为罪也来自于愚昧。祂称他有智慧，因为他没有侵吞，也没有漫无目的或徒劳无益地浪费主人的财物；祂称他忠信，因为他知道如何按适当的方式分派所赐予他的。的确，两者我们都需具备：既不侵吞主人的财物，又能按适当方式分派。若缺其一，另一便跛足。若他忠信而不偷窃，却挥霍浪费在不相关的事上，那罪过就大了；若他善于分派，却侵吞，那对他的指控也不小。\par

我们这些有钱人也当听这些话。因为祂不仅对导师发言，也向富人发言。因为两者都受托了财富：那些教导的人受托了更必需的财富，你们受托了次等的财富。那么，当时导师们正在分散更大的财富，你们却连较小的也不愿慷慨施舍——或者不如说是诚实（因为你们给的是别人的东西）——你们还有什么借口呢？但现在，在我们看到行相反之事的人受罚之前，让我们先听一听那自证其仆人的荣耀。因为\Quote{我实在告诉你们，他要将全部财产交给他管理。}\par

有什么能与这荣耀相比？有什么言语能述说其尊严和福祉？当天的君王、拥有一切的那位，将要派一个人管理\Quote{他的全部财产}时。因此祂也称他为有智慧，因为他知道不为小事舍弃大事，在此世节制，从而到达了天堂。\par

4. 此后，祂一如既往，不仅以善者所存的荣耀，也以对恶者所威胁的惩罚来纠正听者。因此祂又加上了：\Quote{但是，如果那恶仆心里说：我的主人必会迟延；于是开始拷打同仆，并与醉汉一同吃喝饮食；那仆人的主人要在他想不到的日子，不知道的时候来到，把他斩成两半，并分派他与假善人一同：在那里要有哀号和切齿。}\par

但若有人说：你们看见因为日子不确定，他心里起了什么念头？他说：\Quote{我的主人必会迟延？}我们应肯定地说：不是因为日子不确定，而是因为仆人是恶的。否则，为何这念头没有进入那忠信而有智慧的仆人的心里？因为，即使主人迟延，你这可怜的人，你肯定指望他会来。那么你为何不谨慎呢？\par

因此我们学到了，祂并不真的迟延。因为这判断不是出于主，而是出于恶仆的心，所以他为此受责。为证明祂并不迟延，请听\Person{保禄}说：\BibleQuote{主已经近了，应当一无所挂；}\BibleRef{斐 4:5}以及\BibleQuote{那要来的一位必要来，决不迟延。}\BibleRef{希 10:37}。\par

但你也要听随后的话，学习祂如何不断提醒他们那日子是未知的，表明这对仆人们是有益的，适于唤醒并彻底激发他们。因为即使有人未曾由此获益，其他对他们有益的事也有些人未获益，然而祂仍不停做祂该做的部分。\par

那么，接下来的话是什么意思呢？\Quote{因为在他不期待的日子，在他不知道的时刻，他要来到}\BibleRef{玛 24:50}，并要对他施以极重的惩罚。你们看到，祂甚至各处都提出这一点，即他们的无知，表明这是有益的，并借此使他们始终勤奋用心吗？因为这是祂努力的重点，即我们应当时时警醒；既然我们总是在安逸中松懈，却在困苦中振奋，因此祂处处这样说：当有放松时，恐怖便来临。正如祂先前用诺厄的例子表明这一点，同样这里祂说，当那个仆人醉酒、殴打时，他的惩罚将是难以忍受的。\par

但我们不仅要注意为他所定的惩罚，也要留意这一点，免得我们自己也在不知不觉中做出同样的事。因为那些有钱却不周济穷人的，就像这个仆人。因为你们也是自己财产的管家，并不亚于那分发教会赈济的人。正如他无权随意挥霍你们为穷人而给的东西，因为那些东西是为维持穷人生计而给的；同样，你也不可挥霍你自己的财产。因为即使你从父亲那里继承了遗产，并以这种方式拥有了一切，这一切仍然是天主的。那么，你尚且希望你所给的东西被如此仔细地分发，难道你不认为天主会以更大的严格来向我们要求祂的财产，或者祂任凭它们被随意浪费吗？不是这样的，完全不是。因为祂把这些东西留在你手中，正是为了\Quote{按时给他们食粮}。但\Quote{按时}是什么意思？给有需要的人，给饥饿的人。因为正如你交给你的同仆去分发，同样主也愿意你把这些东西花费在必要的事上。因此，虽然祂能从你那里取走它们，祂却留下它们，使你有机会展示美德；通过使我们彼此需要，祂使我们对彼此的爱更加热切。\par

但你们，当你们接受了，非但不给，反而殴打。然而，如果不给已是过错，那么殴打还有什么借口呢？但在我看来，祂说这话是在暗示那些傲慢的人、贪婪的人，并指出他们的罪责沉重，当他们殴打那些他们奉命喂养的人时。\par

5. 祂似乎在这里也暗示那些生活奢侈的人，因为奢侈也存有大惩罚。经上说：\Quote{他同醉酒的人一起吃喝}，指向贪饕。因为你们接受财产不是为了这个目的，好叫你们把它用在奢侈上，而是为了把它用在施舍上。怎么！你们所拥有的难道是你们自己的吗？你们受托管理穷人的财物，即使你们是通过诚实劳动得到的，或者是从父亲那里继承的。难道天主不能从你们那里取走这些东西吗？但祂没有这样做，是为了给你们权力去慷慨地对待穷人。\par

但请你们注意，在所有比喻中，祂如何惩罚那些不把钱用在穷人身上的人。因为童女们没有抢夺别人的财物，而是没有拿出自己的；埋藏一个塔冷通的人没有盗用，而是没有倍增；那些忽视饥饿的人受惩罚，不是因为夺取别人的财产，而是因为不拿出自己的，像这个仆人一样。\par

让我们听从，我们这些取悦肚腹的人，我们这些把根本不属我们、而属于穷人的财富花费在昂贵宴会上的。因为不要因为出于极大的爱，你们被命令如同拿出自己的东西一样施舍，就以为这些东西真是你们自己的。祂把它们借给你们，使你们能证明自己。所以，不要把祂自己的东西归还给祂时，却以为它们是你们的。因为如果你们借给某人一些钱，让他去谋生，你们不会说那钱是他的。同样，天主也给了你们，使你们能为天国经商。因此，不要使祂爱的极大性成为忘恩负义的原因。\par

试想，在领洗后能找到一种方法来消除自己的罪过，这值得多么恳切的祈祷啊！如果祂没有说\Quote{施舍}，多少人会想说：\Quote{如果能用钱来摆脱未来的祸患就好了！}但既然这已成为可能，他们又变得懈怠了。\par

\Quote{但我施舍，}你说。这算什么呢？你还没有像那个投入两个小钱的妇女那样施舍；或者还不如她的一半，远不及她所给的极小部分。你把大部分花在无用的开支上，花在宴会、醉酒和极度的挥霍上；时而请客，时而被请；时而花钱，时而又怂恿别人花钱；以致惩罚对你甚至加倍，既来自你自己的行为，也来自你促使别人所行的。请看祂本人如何为此责备祂的仆人。因为祂说：\Quote{他同醉酒的人一起吃喝。}因为祂不仅惩罚醉酒的人，也惩罚与他们一起的人，非常恰当，因为他们（除了败坏自己）也轻视他人的得救。但没有什么比我们倾向于忽视邻人的事更惹天主发怒的了。因此，表示祂的愤怒，祂命令把他斩为两段。所以祂也肯定爱是祂门徒的标志，因为一个爱人的人必须关心他所爱之人的事。\par

那么，让我们持守这条道路，因为这是尤其通往天堂的道路，它使人成为基督的追随者，使人尽可能地相似天主。请看这些美德何等必要，它们沿着这条道路居住。如果你们愿意，让我们探究它们，并从天主的审判中引出论断。\par

那么，愿有两种最神圣的生活方式：一种确保实行者自身的善，另一种也确保邻人的善。让我们看看哪一种更蒙悦纳，并引领我们达到德性的顶峰。确实，那只寻求自己之事的人，甚至要从保禄那里受到无尽的责备——当我说从保禄，我意指从基督——但另一种则得到称赞和冠冕。这从何显明呢？请听祂对这一位说什么，对那一位说什么。\BibleQuote{谁也不要求自己的利益，但人人都要求别人的利益。}\BibleRef{格前 10:24} 你看见祂拒绝了前者，而引入了后者吗？又说：\BibleQuote{你们各人务要叫邻人喜悦，好使他得益处，得以建立。}随之而来的还有超越言词的赞美与劝诫：\BibleQuote{因为连基督也没有寻求自己的喜悦。}\BibleRef{罗 15:2}\par

这些判断本身足以显明胜利；但为使这事分外充足，让我们看看善工中哪些局限于我们自己，哪些也传递到他人。禁食、卧于硬地、守童贞、克己生活——这些事只给实行者本人带来益处；但从我们传递到邻人的则是施舍、教导、爱德。请听保禄在这事上也说：\BibleQuote{即使我把我的一切财物施舍给人，又舍身被人焚烧，若没有爱德，于我毫无益处。}\BibleRef{格前 13:3}\par

6. 你看见它本身被光荣地颂扬，并加冕吗？\par

但若你愿意，让我们也从第三个角度来比较：一人禁食、克己、殉道、被焚而死；另一人为邻人的益处而延迟殉道，不仅延迟，甚至没有殉道就离世；他们离世后谁更蒙悦纳？我们无需多言，也无需冗长的迂回。因为真福保禄就在眼前，作出判断说：\BibleQuote{我一心想离开世界，与基督同住，那是再好不过的；然而为了你们，我留在肉身内更为必要。}\BibleRef{斐 1:23} 他甚至将邻人的益处置于自己归向基督之上。因为与基督同住的最崇高意义，就是承行祂的旨意；而没有什么比促进邻人的益处更合乎祂的旨意了。\par

你愿意我也告诉你关于这些事的第四个证据吗？祂说：\BibleQuote{伯多禄，你爱我吗？你牧放我的羊。}\BibleRef{若 21:15} 三次询问后，祂宣示这是爱的确切证明。但这话不仅对司祭说，也对我们每个人说，因为我们也受托管理一小群羊。不要因它是小群就轻视，因为祂说：\BibleQuote{我父乐意将天国赐给你们。}\BibleRef{路 12:32} 我们每人有一只羊，要带领它到合适的牧场。人一起床就当寻求别的，只寻求如何做事或说话，使全家更虔诚。女人也当做好管家，但在此之前，她另有更急需的关怀：使全家人都行天国的事。因为在世务上，我们处理家务之前，先努力缴纳公赋，以免因不尽责而被拖到市场挨打，蒙受无数不体面的事；在属灵的事上更当如此，先向万有的君王天主尽本分，免得我们来到那\Quote{有咬牙切齿}的地方。\par

在这些德行之后，让我们寻求那些连同我们自己的救恩，能极大有益于邻人的善工。施舍即是如此，祈祷也是如此——或者更确切地说，后者因前者而有效并插上翅膀。经上说：\BibleQuote{你的祈祷和你的施舍已上达天主面前，获得纪念。}\BibleRef{宗 10:4} 不仅祈祷，禁食也由此得力。你若禁食而不施舍，这行为甚至算不得禁食；这样的人比贪食者和酒徒更糟，而且糟多少，就如残忍比奢侈更严重。我何必提禁食呢？即使你操练克己，即使你守童贞，若没有施舍，你也被摒于洞房之外。然而，有什么能与童贞相比呢？童贞在新约时代也未被法律强制，因其卓越。但尽管如此，缺少施舍时它也被抛弃。如果童女因没有充足施舍而被抛弃，那么没有施舍的人谁能获得宽恕呢？没有人，人若没有施舍，必定灭亡。\par

因为，在世俗事务上，没有人为自己活着：工匠、士兵、农夫、商人，都贡献于公共福祉和邻人的益处；在属灵的事上，我们更当如此。这才是真正的活着：因为那只为己活、忽视所有他人的人是无用的，甚至算不上人，也不是我们的同类。\par

你会说：那么，若我寻求他人的益处而忽略自己的利益，怎么办呢？寻求他人益处的人不可能忽略自己的利益；因为寻求他人益处的人不伤害任何人，却怜悯所有人，尽力帮助他们；他不抢劫，不贪恋他人财物，不偷盗，不作假见证；他远离一切恶行，致力于一切德行；他为仇敌祈祷，善待图谋反对他的人；他不毁谤任何人，不说任何人坏话，即使从他们那里听到无数恶言；他却说宗徒的话：\BibleQuote{有谁软弱，我不软弱呢？有谁跌倒，我不心焦呢？}\BibleRef{格后 11:29} 然而，当我们只着眼自己的益处时，他人的益处未必随之而来。\par

由这一切之事所劝服，既知若不关注公益，便无人能得救，又目睹那被腰斩之人与那埋藏塔冷通之人，让我们选择这条道路，以便我们也能获得永生。愿天主因我们的主\Person{耶稣基督}的恩宠与对人的慈爱，赏赐我们众人获得永生。愿光荣归于他，至于无穷之世。阿们。\par

\SourceRecord{Translated by George Prevost and revised by M.B. Riddle.；Nicene and Post-Nicene Fathers, First Series；Vol. 10.；Edited by Philip Schaff.；Buffalo, NY: Christian Literature Publishing Co.,；1888.；<http://www.newadvent.org/fathers/200177.htm>.}

% structure: p0001 source=h1 target=section reason=page-title

\section{玛窦福音释义第七十八篇}

\BibleRef{玛 25:1-30}\par

\Quote{那时，天国}祂说，\Quote{好比十个童女，拿着灯，出去迎接新郎。其中五个是明智的，五个是糊涂的，她们没有}祂说，\Quote{油。}\par

\Quote{新郎迟延的时候，她们都打盹睡着了。半夜有人喊说：看，新郎来了，你们出来迎接祂！五个童女起来，不知所措，对明智的说：请分点油给我们！明智的却不答应说：恐怕不够你我用的，不如你们到卖油的那里去买吧。}\par

\Quote{她们去买的时候，新郎到了；那预备好了的，就进去；这些后来也来了，说：主啊，主啊，给我们开门！他却回答说：我实在告诉你们，我不认识你们。所以你们要警醒，因为那日子，那时辰，你们不知道。}\par

\Quote{\Emph{祂又讲了另一个比喻：一个人要远行，叫了自己的仆人，把财产托付给他们；一个给了五个塔冷通，一个给了两个，一个给了一个，各按各的才能，然后走了。那两个领了塔冷通的，各自赚了一倍；那领了一个的，却原封不动地带来了，并受到责备，说：我知道你是严厉的人，在没有播种的地方也要收割，在没有散布的地方也要收集；而我} \Emph{我害怕，就去把你的塔冷通藏了起来；看，你的原物在这里。主人回答说：你这恶仆，你既知道我在没有播种的地方也要收割，在没有散布的地方也要收集，就该把我的钱交给钱庄，等我回来时，可以连本带利收回。所以把这个塔冷通夺过来，给那有十个塔冷通的。因为凡有的，还要给他，使他富足；凡没有的，连他所有的也要夺去。把这无用的仆人丢在外面的黑暗中，在那里必有哀号和切齿。}}\par

这些比喻与之前忠仆和浪费主人财物的恶仆的比喻相似。因为一共有四个比喻，从不同方面就同一主题警告我们，我指的是勤于施舍，以及用我们所能的一切方式帮助邻人，因为不如此就无法得救。但是，那里祂更一般地谈论应当给予邻人的一切帮助；而关于童女，祂特别谈论了施舍的仁慈，并且比前一个比喻更强烈。因为在那里，祂惩罚的是殴打、醉酒、浪费主人财物的人；而在这里，祂惩罚的甚至是不帮助、不慷慨地使用自己财物的人。因为她们确实有油，但不充足，因此她们也受了惩罚。\par

但是为什么祂以童女为比喻，而不是随便取一个人物？祂曾论及童贞说了伟大之言，说：\Quote{有些阉人，是为了天国而自阉的；}又说：\Quote{谁能接受，就让他接受。}\BibleRef{玛 19:12}祂也知道众人对童贞评价很高。因为这件事本性上确实伟大，并且从这样的事实可以看出来：在旧约时代，这些古圣并未实行它，在新约时代，它也不是律法的强制。因为祂没有命令这事，而是留给听众选择。因此保禄也说：\Quote{关于童贞，我没有主的命令。}\BibleRef{格前 7:25}\Quote{我虽然称赞达到这境界的人，却不勉强不愿的人，也不将此当作命令。}既然这件事本身是伟大的，并且在众人中享有极大的尊荣，以免有人达到此境界后，便以为已达到了全部，而对其他事漠不关心，祂便讲了这个比喻，足以说服他们：童贞即使拥有一切，但若缺乏由施舍而来的善行，就会被与娼妓一起驱逐出去，祂将冷酷无情的人与她们放在一起。这是非常合理的，因为一个是被肉欲的享乐所克服，而她们则是被金钱所克服。肉欲的享乐和金钱不等，肉欲享乐强烈得多，更专制。敌人越弱，被其克服的人越不可原谅。因此祂也称她们为糊涂的，因为她们承担了更大的辛劳，却因缺少较少的而辜负了一切。这里，祂用灯指着童贞本身，即圣洁的纯洁；用油指着\OriginalTerm{仁爱}{humanity}、施舍、帮助有需要的人。\par

\Quote{新郎迟延的时候，她们都打盹睡着了。}祂借此表明中间的时间不会短，使祂的门徒脱离祂的王国会立即显现的期望。因为门徒们确实希望这样，所以祂不断阻止他们抱这样的希望。同时祂也暗示死亡是一种睡眠。因为她们睡着了，祂说。\par

\Quote{约在半夜，有人喊说。} 祂或许是在继续比喻，或者又表明复活将在夜间。但保禄也指出那喊声，说：\Quote{随着一声呐喊，随着总领天神的声音，随着最后的号角，祂要从天上降下。} 号角是什么意思？那喊声说什么？\Quote{新郎来了。} 于是当她们整理好灯之后，愚拙的对明智的说：\Quote{请分些油给我们。} 祂再次称她们为愚拙的，表明没有什么比那些在此处富有、在那里赤身离去的人更愚拙了，在那里我们最需要仁慈，我们需要许多油。但不仅在这方面她们是愚拙的，还因为她们指望在那里得到它，并在不合时宜时寻求它；然而没有什么比那些童女更仁慈的了，她们正因此受到赞许。她们也没有寻求全部，因为她们说：\Quote{请分些油给我们；} 而且她们需要的紧迫性被表明：\Quote{因为我们的灯快要灭了。} 但即使如此，她们失败了，无论是他们所请求之人的仁慈，还是请求的容易，还是他们的急需和匮乏，都没有使她们得到。\par

但我们现在从这里学到什么？就是没有人在那里能保护我们，如果我们被自己的行为出卖了，不是因为他不愿意，而是因为他不能。因为这些人也以不可能为托辞。有福的亚巴郎也指出这一点，说：\Quote{在我们与你们之间，隔着一个巨大的深渊。} \BibleRef{路 16:26}，以致即使愿意，也不准他们越过它。\par

\Quote{但是去找那些贩卖的人，买油吧。} 那贩卖的是谁？是穷人。他们在哪里？在这里，那时她们本该寻找他们，而不是在那个时刻。\par

你看见从穷人那里给我们带来多大的益处吗？你若把他们除去，你就除去了我们救恩的大希望。因此我们必须在这里积攒油，以便在时间召唤我们时，它在那里对我们有用。因为那不是收集的时刻，而是现在。所以不要把你的财物浪费在奢侈和虚荣上。因为你在那里将需要许多油。\par

那些童女听见这些话就走了；但她们毫无所得。祂这样说，要么是继续并润色比喻，要么也是借这些事表明，即使我们在离世后变得仁慈，我们也无法从那里获得任何帮助以逃脱。因此，她们的积极对那童女毫无用处，因为她们去找那些贩卖的人，不是在这里，而是在那里；那富人也是如此，当他变得如此慈善，甚至为他的亲属担忧时。因为那个曾经过躺在门口的人，竟急于从危险和地狱中拯救那些他甚至没有见过的人，并恳求派人去告诉他们这些事。但尽管如此，他并未从中获益，正如那些童女一样。因为当她们听了这些话离去时，新郎来了，那些预备好了的同祂进去，其余的人却被关在门外。在她们许多劳苦之后，在她们无数辛劳之后，在那难以忍受的战斗之后，在她们制服自然欲望的疯狂所立下的那些战利品之后，她们蒙羞，灯也灭了，退下，脸伏于地。因为没有什么比没有怜悯的童贞更污秽的了；所以甚至民众也惯于称不仁慈的人为黑暗。那么童贞的益处在哪里呢，当她们见不到新郎时？甚至当她们敲门时，她们也没有得到，却听见那可怕的话：\Quote{离开我，我不认识你们。} \BibleRef{玛 25:12} 当祂说了这话，剩下的就只有地狱和那不能忍受的惩罚；或更好说，这话比地狱更痛苦。祂也对那些作恶的人说这话。\par

\Quote{所以你们要警醒，因为你们不知道那日子，也不知道那时辰。} \BibleRef{玛 25:13} 你看见祂如何不断地加上这句话，表明我们对于离世的时刻的无知是多么可怕吗？那些一生懈怠，但当我们责备他们时却说：\Quote{我临死的时候，会给穷人留下钱财。} 的人在哪里？让他们听这些话而改正吧。因为的确在那时，许多人未能做到，他们被突然夺去，甚至不被允许向亲人交代他们想做之事。\par

这个比喻是针对施舍方面的仁慈说的；但接下来的那个，是针对那些无论在钱财、言语、保护或任何其他事情上都不愿意帮助邻人，反而一概不给的人说的。\par

为什么这个比喻引进一位君王，而那个却引进一位新郎呢？为使你学习基督与那些舍弃自己财产的童女结合得多么密切；因为这正是童贞。所以保禄也以此作为该事的定义。他说：\Quote{未出嫁的妇女挂虑主的事；} \BibleRef{格前 7:34} 又说：\Quote{为那合宜的事，并使你们能无分心地事奉主。这些事我们劝告你们。}\par

如果在路加中关于塔冷通的比喻另有说法，那么应该说，这个比喻与那个确实不同。因为在那个比喻中，从同一本金产生了不同程度的增益，因为有人用一米纳赚了五米纳，另一个赚了十米纳；因此他们也没有得到同样的报酬；但这里相反，荣冠也因此相等。因为领了两个的给了两个，领了五个的也同样；但在那里，由于从相同的起点，有人增益较多，有人较少；正如所料，在赏报上，他们也不享受同样的。\par

但是，请看祂处处不立即要求归还。因为就葡萄园的比喻而言，祂将葡萄园租给佃户，就出国远行了；在这里，祂将塔冷通交付给他们，就起程了，为叫你学习祂的忍耐。在我看来，祂说这些话似乎是在暗示复活。但在这里，不再有葡萄园和佃户，而是所有的仆人。因为祂的话不是仅对统治者，也不是仅对犹太人，而是对所有人说的。那些向祂交账的人坦承，哪些是自己的，哪些是主人的。一人说：\Quote{主，祢给了我五个塔冷通。}另一人说：\Quote{两个，}指明他们从祂那里获得了得利的根源，而且非常感恩，将一切归于祂。\par

主人于是说什么？\BibleQuote{做得好，你这良善的}（因为这就是良善：顾及邻人）\BibleQuote{又忠信的仆人；你在小事上忠信，我要委派你管理大事：进入你主人的福乐吧！}\BibleRef{玛 25:23}祂用这句话意指一切福乐。\par

但另一个却不是如此，而是怎样？\BibleQuote{我知道你是个严厉的人；没下种的地方也收割，没有散播的地方也聚敛；我害怕了，就去把你的塔冷通藏起来；看，你原有的还你。}\BibleRef{玛 25:24}主人于是说什么？\BibleQuote{你应当把我的银子交给兑换银钱的人，}就是说，\Quote{应当说话，应当劝诫，应当忠告。}但若他们不肯听从呢？这与你不相干。\par

还有什么比这更温和的？因为人们确实不这样做，反而叫那放了高利贷的人负责去收回。但祂却不是这样；而是说：\Quote{你应当把它存到钱庄，把追讨的事交给我。我自会加上利息收回。}利息是指听道后的行为表现。你该做较容易的事，把更难的留给我。既然他没有这样做，祂就说：\BibleQuote{把塔冷通从他身上取去，给那有十个塔冷通的。因为凡有的，还要给他，使他富足；但凡没有的，连他所有的也要夺去。}\BibleRef{玛 25:28}那么这是什么意思？那拥有言语和教导的恩赐并以此获益却不去使用的人，也会失去这恩赐；但那勤勉的人，会为自己赢得更丰富的恩赐，正如另一人失去他所领受的。但懈怠者的惩罚不仅限于此，而且惩罚本身无法忍受，判决充满了严厉的控诉。因为祂说：\BibleQuote{把这无用的仆人丢到外面的黑暗里去；在那里要有哀号和切齿。}你看见没有？不仅掠夺者、贪得无厌者、作恶者，就连不行善的人，也要受极重的惩罚。\par

那么就让我们听从这些话吧。趁着有机会，我们当努力于自己的救恩，为自己的灯预备油，努力增加自己的塔冷通。因为若我们退缩，在此地虚度光阴，即使将来恸哭万遍，也无人再怜悯我们。身穿污秽衣服的人定自己的罪，毫无裨益。那领了一个塔冷通的，归还了所托付的，却仍被定罪。童女们一再恳求，到祂那里敲门，却都是徒然，毫无效果。\par

所以，知道这些事以后，让我们同样为邻人的益处贡献财富、勤勉、保护以及所有一切。因为这里的塔冷通是指各人的能力，不论是保护、金钱、教导，或任何此类的事。不要有人说：我只有一个塔冷通，什么都做不了。因为即使只有一个，你也能证明自己。因为你并不比那位寡妇贫穷；你并不比\Person{伯多禄}和\Person{若望}更无知，他们两人都是\BibleQuote{没有学识的平民}\BibleRef{宗 4:13}；但尽管如此，由于他们表现出了热忱，并为公众利益做了一切，他们获得了天国。因为再没有比为公共利益而活更中悦天主的了。\par

为此，天主赐给我们口舌、双手、双脚、体力、心智和领悟力，好让我们使用这一切，既为自身的救恩，也为邻人的益处。因为我们的言语不仅用于赞美和感恩，也用于教导和劝诫。如果我们为此目的使用它，就是在效法我们的主；但如果为相反的目的，就是效法魔鬼。因为\Person{伯多禄}在宣认基督时，是有福的，因为他说了圣父的话；但当他拒绝十字架并劝阻时，他受到了严厉的责备，因为他体会的是魔鬼的事。然而，如果在无知时说的话，责备尚且如此沉重，当我们故意犯许多罪时，还能得到什么恩惠呢？\par

那么，让我们说这样的话，使它们本身就是基督的话。因为不但如果我说\BibleQuote{起来，行走吧！}\BibleRef{玛 9:5}；也不但如果说\BibleQuote{塔彼达，起来！}\BibleRef{宗 9:40}那时我才说基督的话；而更当我在受辱骂时祝福，在受凌辱时为那凌辱我的人祈祷。近日我确实说过，我们的舌头是抓住天主脚的手；但现在我更要说，我们的舌头是模仿基督的舌头的舌头，如果它表现出合宜的严谨，如果我们说祂所愿意的话。但祂愿意我们说的是什么话？充满温和与良善的话，正如祂自己曾对侮辱祂的人说的：\BibleQuote{我没有附魔}；又说：\BibleQuote{我若说得不对，你指证那不对的地方。}\BibleRef{若 18:23}如果你也这样说；如果你为邻人的改进而说话，你将得到像那样的舌头。而天主亲自说：\BibleQuote{你若将珍宝从垃圾中取出，你就可做我的口舌。}\BibleRef{耶 15:19}这就是祂的话。\par

所以，当你的舌头成了基督的舌头，你的口成了圣父的口，你成了圣神的宫殿时，还有什么尊荣能与此相比？即使你的口是金子做的，甚至宝石做的，也不如现在这样，因着温良的装饰而发光。有什么比一个不知侮辱、惯于祝福\Emph{且出善言}的口更可爱呢？但若你不能忍受祝福那咒骂你的人，就暂且沉默，以此完成这一点；循序渐进，尽力奋斗，你也会达到那另一点，并拥有我们所说的那样的口。\par

4. 不要认为这话是冒失的。因为主是爱人者，恩赐来自祂的良善。拥有魔鬼的口，拥有恶鬼般的舌头，尤其是对于参与如此奥迹、领受主亲自的肉体的人而言，才是冒失的。故此，你当默思这些事，尽己所能变得像祂。那样，魔鬼再也不敢正视你，因为你已成了这样的人。他认出了君王的肖像，他知道基督的武器，就是基督借以击败他的武器。这些是什么？温和与温良。当基督在山上击溃并挫败那攻击祂的魔鬼时，祂不是通过显明自己是基督，而是用这些言语诱捕他，以温和擒获他，以温良使他逃遁。你也当如此行。若你看见一个人成了魔鬼，来攻击你，你也照样做。基督赐予你能力，使你尽己所能变得像祂。不要因听到这话而害怕。害怕才是不像祂。那么，按祂的方式说话，在这方面你就成了像祂那样的，以一个人所能达到的程度。\par

因此，如此说话的人比那说预言的人更大。因为说预言完全是恩赐，而在后者中也有你的劳苦和辛劳。教导你的灵魂为你塑造一个像基督的口。因为灵魂若愿意，就能创造这样的口；它知道这门技艺，只要不懈怠。这样的口如何造就？有人会问。通过哪种色彩？通过哪种材料？确实，不是通过色彩或材料，而是只通过德行、温良和谦卑。\par

我们也来看魔鬼的口是如何造就的，好让我们永不造就那样的口。那么它是如何造就的呢？通过咒骂、侮辱、嫉妒、伪誓。因为当一个人说出魔鬼的话时，他就取了魔鬼的舌头。那么我们将有什么借口？或者说，当我们忽略这舌头——我们被允许用它品尝主的肉体——却让它说出魔鬼的话时，我们将受怎样的惩罚？\par

我们不可忽略，而要尽一切努力，训练它效法它的主。因为若我们训练它如此，它就会使我们满怀信心地站在基督的审判台前。若有人不知如何这样说话，审判者甚至不会听他。就像法官若是罗马人，他不会听一个不知如何这样说话的人的辩护一样；同样，基督除非你按祂的方式说话，否则不会听你，也不会垂顾。\par

因此，让我们学会以我们的审判者惯于听闻的方式说话；让我们努力模仿那舌头。若你陷入悲伤，当心不要被沮丧的暴政扭曲你的舌头，而要像基督那样说话。因为祂也曾为拉匝禄和犹达斯悲伤。若你陷入恐惧，再寻求像祂那样说话。因为祂为你，就祂的人性而言，也曾陷入恐惧。另一群羊，为要指出一群的不结果实，因为山羊羔不会结果实；而另一群大有裨益，因为绵羊确实大有裨益，无论是奶、毛还是羔羊，而山羊羔在这些方面都缺乏。\par

\SourceRecord{Translated by George Prevost and revised by M.B. Riddle.；Nicene and Post-Nicene Fathers, First Series；Vol. 10.；Edited by Philip Schaff.；Buffalo, NY: Christian Literature Publishing Co.,；1888.；<http://www.newadvent.org/fathers/200178.htm>.}

% structure: p0001 source=h1 target=section reason=page-title

\section{\WorkTitle{玛窦福音释义第七十九篇}}

\BibleQuote{当人子在他的圣父的光荣中，与所有圣天使一同降来时，他要坐在他光荣的宝座上，把绵羊和山羊分开。}（\BibleRef{玛 25:31-46}）其中，前者因曾给饥饿的他食物、给口渴的他水喝、收留作客的他、给裸体的他衣穿、看望患病的他、到监狱中探望他而蒙接纳，他将赐予他们天国；后者则因相反的行为受责罚，被投入为魔鬼及其使者预备的永火中。\par

这最令人喜悦的圣经段落，我们不断回旋默想，现在当以全副热忱和痛悔之心来聆听；这恰是基督话语的终结，作为最后的话，合情合理，因为他深怀爱护人类与慈悲之心。因此，在前文中他已以不同方式论及此事；而此处则更为清晰、更为恳切地加以陈述，不是提出两个、三个或五个人，而是整个世界；尽管先前那些提到两个人的地方，确非指两个人，而是指人类的两部分：一部分是悖逆者，另一部分是服从者。但此处他更以令人敬畏的方式、更明了的光照来阐发这句话。因此他不再说\Quote{天国好比……}，而是公开显示自己，说\Quote{当人子在自己的光荣中降来时}。如今他是在屈辱中来临，在侮辱和责难中；但那时他必坐在他光荣的宝座上。\par

他不断提及光荣。因为十字架临近，看似是羞辱之事，为此他提升听众的心神，将审判座置于他们眼前，并召聚整个世界围绕四周。\par

不仅借此使他的话语可畏，更显示诸天敞开。他说，所有天使都与他同来，他们也要作证，见证他们曾多少次奉主差遣为人的救恩服务。\par

一切事物都使那日子变得可畏。他说，那时\BibleQuote{万民都要聚集}，即全体人类。\BibleQuote{他要将他们彼此分开，如同牧人分开绵羊和山羊一样。}如今他们并未分开，而是混杂在一起，但那时将精确地分开。暂时他按位置划分他们，使之显明；随后他藉名称指出各人的心性，称一些为山羊，另一些为绵羊，以表明山羊的不结果实，因为山羊不结果实；而绵羊则大有裨益，因为绵羊的奶、毛和羔羊都带来丰厚的收益，这些山羊一概没有。\par

牲畜的不结果实和结果实出于本性，而人则出于选择，因此一些人受罚，另一些人得冠冕。他并非不加申辩就惩罚他们；因此，当他将他们安置妥当后，便提出对他们的指控。他们以温顺的语气讲话，但此时已无济于事；这很合理，因为他们忽略了如此可嘉的工作。先知们随处都在说：\BibleQuote{我喜欢仁爱胜过祭献}（\BibleRef{欧 6:6}）；立法者以言以行始终敦促他们行此善工，自然本身也教导他们。\par

但请注意，他们并非缺少一两件事，而是缺少一切。因为他们不仅没有给饥饿者食物、给裸者衣穿，甚至连看望病人这样更容易的事也没有做。\par

请注意他的命令是多么轻松。他没有说：\BibleQuote{我在监里，你们释放了我；我患病，你们使我痊愈}，而是说\BibleQuote{你们看望了我}和\BibleQuote{你们到我这里来}。饥饿的要求也并不沉重。因为他并不要昂贵的筵席，只需必需品和必要的食物，而且他以乞求者的形象出现，所以一切都足以定他们的罪：要求的容易（不过是食物）、请求者的可怜（他是穷人）、本性的同情（他是人）、应许的吸引力（他应许天国）、惩罚的可怕（他威胁地狱）、受惠者的尊威（天主藉穷人接受）、他所屈尊给予的非凡荣誉、他所赋予的正当要求（他接受的是属于自己的）。但贪财一旦抓住人，便蒙蔽了他们，尽管如此大的威胁摆在面前。\par

再往前，他说，凡不接待这样一个小孩子的人，所受的比索多玛还重；此处他说：\BibleQuote{你们没有给这些最小兄弟中的一个做，就是没有给我做。}（\BibleRef{玛 25:45}）他是你的兄弟，为何称他们为最小的？正因如此，他们是兄弟：因为他们卑微、贫穷、遭遗弃。他正是邀请这样的人加入兄弟之谊：无名、可鄙的人；他不仅指隐修士和住在山上的人，而是指每一位信徒，即使他是个世俗人，只要他饥饿、赤贫、裸体、作客，他就应得到这一切照顾。因为圣洗使人成为兄弟，分享天主奥秘也使人成为兄弟。\par

2. 然后，为使你也从另一角度看到判决的公义，他先赞扬那些行善的人，说：\BibleQuote{来吧，我圣父所祝福的，承受自创世以来为你们预备的国度吧！因为我饿了，你们给了我吃的；我渴了，你们给了我喝的……}（\BibleRef{玛 25:34}）及以下的话。为使你们不能说\Quote{我们没有}，他藉同仆之口定他们的罪，如同童女被童女定罪，醉酒贪食的仆人被忠信仆人定罪，埋没塔冷通的人被带回两个塔冷通的人定罪，每个继续犯罪的人被行善的人定罪。\par

这番比较有时用于平等双方，正如这里以及关于童贞女的比喻中所见；有时用于一方占优势的情形，如祂说：\Quote{尼尼微人要起来，定这世代的罪，因为他们听了约纳的宣讲而悔改了；看，这里有一位大于约纳的！}以及\Quote{南方女王要定这世代的罪，因为她从地极而来，听撒罗满的智慧；} \BibleRef{玛 12:41} 又再次用于平等双方：\Quote{他们必要做你们的判官；} \BibleRef{玛 12:27} 以及用于一方占优势的情形：\Quote{你们不知道，我们要审判天使，更何况今生的事呢？} \BibleRef{格前 6:3}\par

但这里却是用于平等双方：祂将富人与富人相比，穷人与穷人相比。祂不仅藉此表明判决的公正，因他们的同仆在那相同处境下做了正直的事，而且更因他们在这些连贫穷也非障碍的事上（如给飢渴者饮水、看望被囚者、探视病人）并未顺服。当祂称赞那些做对了的人时，祂更显示祂对他们原本的爱有多么深重。因为祂说：\Quote{你们蒙我父降福的，来承受那从创世以来为你们预备的国度。} 这同一句话等于多少善事：蒙福，并蒙父降福？他们何以配得如此殊荣？原因何在？\Quote{我饿了，你们给了我吃的；我渴了，你们给了我喝的；}以及随后的话。\par

这是何等的尊荣，何等的福分！祂不是说\Quote{领受}，而是说\Quote{承受}，犹如自己的产业，犹如父的，犹如你们的，犹如从起初就应得的。因为祂说：在你们存在之前，这些就已为你们预备妥当，因为我预知你们会成为这样的。\par

因何回报他们得着这些？为一片屋顶的遮蔽，为一件衣服，为一块饼，为一口凉水，为一次探望，为一次探监。实际上，在每一种情形下都是为急需之物；有时甚至还不是。因为正如我曾说的，病人和被囚者所求的并不仅限于此：一个求释放，一个求痊愈。但祂仁慈，只要求我们能力所及之事，甚至比能力所及更少，留给我们余裕去行更多慷慨之事。\par

但祂对另一些人说：\Quote{可咒骂的，离开我，}（不再说是父的咒骂，因咒骂他们的不是祂，而是他们自己的行为）\Quote{到那为魔鬼和他的使者预备的永火里去。}关于天国，当祂说\Quote{来，承受国度}之后，加上了\Quote{从创世以来为你们预备的}；但关于火，则不再这样，而是\Quote{为魔鬼预备的}。祂说：我预备了天国给你们，但那火并不是为你们预备的，而是\Quote{为魔鬼和他的使者}；但既然你们自己投身其中，这要归咎于你们自己。不仅在此，祂还藉随后的话，如同为自己辩护一样，陈明了原因。\par

\Quote{因为我饿了，你们没有给我吃的。}即使来到你们面前的是你们的仇敌，难道祂所受的苦难还不足以感动和降服那毫无怜悯之心的人吗？飢渴、寒冷、锁链、赤身、疾病、到处漂泊无家？这些足以消除仇恨。但你们连对一个朋友——同时是朋友、恩主和主——也没有做这些。即使我们看到一条狗饿了，也常被触动；看到一只野兽，也心生怜悯；但看到主，你们却无动于衷？这些事有何可辩解的呢？\par

因为如果仅仅是这些（先不谈在全世界面前，从坐在圣父宝座上的祂口中听到这样的话并赢得国度），难道那行动本身还不够得赏报吗？但如今，即在全世界面前，在那不可言喻的荣耀显现时，祂宣扬你，为你加冕，承认你是祂的供养者和主人，不以说这些话为耻，为使你的冠冕更加光彩。\par

为此，一方面，那些人受罚是公义的；另一方面，这些人却蒙恩受冠。因为即使他们行了万千善事，这份慷慨仍是出于恩宠，为如此微小低廉的服务，竟回报以如此的天国、国度及如此大的荣耀。\par

\Quote{耶稣说完了这一切话，就对门徒说：你们知道，两天后就是逾越节，人子将被出卖，钉在十字架上。} \BibleRef{玛 26:1} 当祂提醒他们天国、那里的赏报以及不死的惩罚之后，祂适时地又谈到受难，仿佛在说：既然如此美好的事等待你们，你们为何害怕暂时的危险呢？\par

3. 但请注意，我恳求你，祂在最初的那些话中如何以新颖的方式处理并淡化了他们最痛苦的事。因为祂没有说\Quote{你们知道，两天后我要被出卖}，而是说\Quote{你们知道，两天后就是逾越节}，以表明所发生的是奥迹，是为世界的救恩而举行的节日和庆典，并且祂是预知一切而受苦的。因此，这足以安慰他们，祂甚至没有对他们再提复活的事；因为既然已经谈论了许多，再提便是多余。再者，正如我所说，祂藉逾越节提醒他们在埃及的古代恩惠，以此表明祂的受难本身正是无数邪恶的解脱。\par

\Quote{那时，司祭长、经师和民间长老聚集在大司祭该亚法的庭院里，商议如何用诡计捉拿耶稣，并杀害祂。但是他们说：不可在庆节期间，以免民间发生暴动。}\par

你们看见犹太国无可言喻的腐化了吗？他们企图行非法之事，来到大司祭那里，想要从本该找到阻碍的地方获得权柄。\par

那时有多少位大司祭？因为法律只许有一位，但那时却有许多位。由此可见，犹太人的体制已经开始瓦解。因为梅瑟，如我所说，命令只应有一位，当他死后另有一位，并且天主以这人的在世时期来量定无意杀人者的流放期。那么，那时怎么会有许多大司祭呢？后来他们是一年一任的。圣史在谈及匝加利亚时说明了这一点，说他属于阿彼雅班次。因此，他在这里称那些曾经作过大司祭的人为大司祭。\par

他们商议什么呢？是要秘密捉拿祂，还是置祂于死地？两者皆是；因为他们惧怕百姓。因此他们也等逾越节过去；因为\Quote{他们说：\InnerQuote{不可在庆节日子。}}魔鬼不愿这事发生在逾越节，免得受难成为显眼的事；但他们却恐怕发生骚乱。请看他们始终畏惧的，不是来自天主的灾祸，也不是怕因节期而招致更大的污秽，而总是来自人的灾祸。\par

尽管如此，他们仍怒火中烧，再次改变了主意。因为他们虽然说过\Quote{不可在庆节日子}，但一找到叛徒，便不等时机，就在庆节期间杀害了祂。但为什么那时捉拿祂呢？如我所说，他们怒气沸腾；他们期望那时找到祂，所做的一切都像瞎了眼。因为尽管祂自己最大程度地利用他们的恶行来完成祂的救恩计划，他们却并非因此无罪，反而因他们的心肠而该受无数惩罚。至少当众人，甚至罪人，都应得释放时，这些人却杀害了那无辜者，那曾赐给他们无数恩惠、并曾为他们的缘故暂缓眷顾外邦人的那一位。噢，仁爱啊！祂再次拯救这些如此败坏、\Emph{如此乖戾}、充满无数邪恶的人，并派遣宗徒为他们受死，且通过宗徒恳求他们。\Quote{因为我们是为基督作大使的。}\par

既然有这样的榜样，我不说我们应为敌人而死，因为连这事我们也应该做；但由于我们过于软弱，无法做到这点，我暂且说，至少不要以恶眼看我们的朋友，不要嫉妒我们的恩人。我不说暂且善待那恶待我们的人，因为我也渴望做到这点；但既然你们对此过于迟钝，至少不要报复自己。我们的处境像什么？像一场戏和表演？为什么你们直接对抗所吩咐的事呢？其他一切被记录下来，以及祂在十字架上所行足以召回他们的许多事，并非徒然；而是为使你们效法祂的良善，效法祂的仁爱。因为祂确实将他们击倒在地，又恢复了仆人的耳朵，且以忍耐讲论；当祂被举起时，行了伟大的奇迹：掩盖日光、使岩石崩裂、复活死者、以异梦惊吓审判祂的人的妻子；在审判时显示了完全的温柔（这比奇迹更有效力来赢得他们），在审判厅里预先警告了他们无数的事；在十字架上大声呼喊：\Quote{父啊，宽免他们的罪过。}当祂被埋葬后，为他们的得救显出了多少事？复活后，祂岂不立刻召叫了犹太人？岂不赐给他们罪的赦免？岂不将无数祝福摆在他们面前？还有什么比这更奇妙的呢？那些钉死祂、杀气腾腾的人，在钉死祂之后，竟成了天主的子女。\par

有什么能与这种温柔相比呢？听到这些，让我们掩面，因为我们离我们应该效法的那一位如此遥远。至少让我们看看这距离有多大，好使我们至少能自责，因为与那些基督为之舍命的人争战，且不愿与他们和好，而基督为和好甚至拒绝被杀；除非这也是一种花费和金钱支出，正是你们在施舍中所反对的。\par

4. 想一想你犯了多少过错；因此，你非但不迟迟不肯原谅那得罪你的人，反而会奔向你那叫你伤心的人，为使自己有得宽恕的理由，为能找到补救自己恶行的药方。\par

希腊人的子孙，他们并不追求什么伟大之事，却常在这方面表现出自制；而你，怀着如此希望即将离世，却畏缩不前，行动迟缓；时间所成就的事，你却不甘愿因天主的法律提前去做，反而情愿让这种情绪得不到酬报地熄灭，而不是为得酬报？因为即使这种情绪是因时间而起的，你也不会得到任何好处，反而要受重罚，因为时间所成就的事，天主的法律却没有说服你去做。\par

但若你说你因记得那侮辱而心如火焚，那么请回想那得罪你的人曾对你做过什么好事，以及你曾给他人造成多少祸害。\par

他曾恶言中伤你，使你蒙羞？也请想想你也曾这样论及别人。那么，你不宽恕别人，又如何能获得宽恕呢？但若你没有中伤过任何人？但你曾听人如此说而容许之。这也并非无罪。\par

你愿学习不记仇怨是何等美善，以及此事比一切更能悦乐天主吗？祂惩罚那些对天主亲自公正惩治的人幸灾乐祸的人。然而他们确是公正受罚，但你不可幸灾乐祸。因此先知提出许多控诉后，又加上这话说：\Quote{他们不为若瑟的苦难所动；}\BibleRef{亚 6:6} 又说：\Quote{住在厄南的，没有出来为她附近的场所哀哭。}然而，无论是若瑟（即从他的后裔所生的支派），还是这些人的邻舍，都是按照天主的旨意受罚；但祂仍愿意我们同情这些人。因为我们这些恶人，在惩罚仆人时，若见另一个同仆在笑，我们就更被激怒，把怒气转向他；何况天主更要惩罚那些对祂所惩治的人幸灾乐祸的人。如果对天主所惩治的人不可践踏，而应与之同悲，那么对得罪我们的人就更当如此。因为这是爱德的标记；天主喜爱爱德胜过一切。正如在朝廷的紫色袍中，那些构成这袍服的花朵和染料中，有些是珍贵的；同样，在这里，这些德行是珍贵的，它们保存爱德。但没有比不记念得罪我们的人更能维持爱德的了。\par

\Quote{难道天主没有保护另一方吗？难道祂没有把行不义的人赶向被亏负的人吗？难道祂没有打发他从祭台到那人那里，待和好之后再请他上桌吗？}但不要因此等候对方来，因为这样你就失去一切。因为祂特别为此指定给你不可言喻的赏报，使你抢先一步；因为如果你因他恳求而和好，这和睦就不再是出于天主的命令，而是对方的努力。因此你空手而去，他倒领受赏报。\par

你说什么？你有仇敌，还不羞愧吗？为什么魔鬼还不够我们对付，我们还要把同族的人也招惹来？但愿连他（魔鬼）也不曾有心攻打我们；但愿连他也不是魔鬼！\par

难道你不知道和好后的喜乐有多大吗？因为什么，尽管在仇视中它显得不大？因为爱那亏负我们的人比恨他更甘甜，在仇视消除后，你将能充分体会到。\par

5. 那么，为什么我们要效法疯人，互相吞噬，攻打自己的血肉呢？\par

你们即使在旧约下也听听，对这事有多重视：\Quote{报复的人的途径是死亡。一人对另一人怀怒，难道他还能寻求天主的治愈吗？}\BibleRef{德 18:3} \Quote{祂却允许了\InnerQuote{以眼还眼}和\InnerQuote{以牙还牙}，那么祂为何责备呢？}因为祂允许那些事，不是要我们彼此施行，而是藉着对受苦的恐惧，使我们避免犯罪。此外，那些行为是短暂愤怒的果实，而记仇是灵魂在恶中操练的表现。\par

但是你受了恶吗？然而没有比你自己因记仇而会加给自己的更大了。此外，一个善人几乎不可能遭受任何恶。因为假设有人，有子女有妻子，让他操练德行，又让他有许多被亏负的机会，以及大量财产、主权、许多朋友，又让他享受尊荣；只要让他操练德行——因为必须加上这一点——并且我们假设灾祸落在他身上。让一个恶人到他那里，使他受损失。那么，对那视金钱如无物的人，这算什么呢？让他杀他的孩子；对那在复活中得到智慧的人，这算什么呢？让他杀他的妻子；对那受教导不为长眠者悲伤的人，这算什么呢？让他使他受辱；对那视现世如草花的人，这算什么呢？如果你愿意，让他也折磨他的身体，把他下在监里；对那已经学会\Quote{我们的外在人虽然毁坏，内在人却日日更新}，并且\Quote{患难生称许？}的人，这算什么呢？\par

现在我曾说过他不会受伤害；但论述继续下去却显示他甚至得益，被更新，并成为可赞许的。\par

那么，我们不要因他人而烦恼，伤害自己，使我们的灵魂软弱。因为烦恼更多的是出于我们的软弱，而非邻人的邪恶。因此，若有人侮辱我们，我们就哭泣、皱眉；若有人抢夺我们，我们就像那些小孩子一样，他们同伴中较聪明的无缘无故招惹他们，为小事使他们忧愁；然而这些人，如果看见他们烦恼，就继续戏弄他们，但如果看见他们笑，反而停手。但我们比这些孩子更愚蠢：为那些我们本应一笑置之的事而哀哭。\par

因此我恳求，我们应放下这孩童心思，抓住天国。因为基督愿意我们成为成年人，完美的人。保禄也是这样命令的：\Quote{弟兄们，在见识上不要作孩子，}他说，\Quote{但在邪恶上却要作婴孩。}\par

因此，我们要在邪恶上作婴孩，逃避恶行，抓住德行，这样我们也能获得永生的美善，藉着我们的主耶稣基督的恩宠和爱人之心，愿光荣与权能归于祂，至于无穷之世。阿们。\par

\SourceRecord{Translated by George Prevost and revised by M.B. Riddle.；Nicene and Post-Nicene Fathers, First Series；Vol. 10.；Edited by Philip Schaff.；Buffalo, NY: Christian Literature Publishing Co.,；1888.；<http://www.newadvent.org/fathers/200179.htm>.}

% structure: p0001 source=h1 target=section reason=page-title

\section{玛窦福音释义 第八十讲}

\Emph{\BibleRef{玛 26:6-7}}。\par

\BibleQuote{耶稣正在伯达尼癞病人西满家里时，有一个女人拿着一玉瓶贵重的香液，来到祂面前，倒在祂的头上，当祂坐席的时候。}\par

这个女人似乎与所有圣史所记载的是同一人，但实际并非如此；因为尽管与前三圣史的记载在我看来似乎是同一人，但与若望的记载则不是，而是另一个人，一位非常令人钦佩的人，即拉匝禄的姐妹。\par

圣史提及西满的癞病并非没有目的，而是为显示那女人从何处获得信心而来到祂跟前。因为癞病似乎是一种最不洁净、最可憎恶的疾病，然而她看见耶稣不但治愈了那个人（否则祂不会选择与一个癞病人同住），而且还进了他的屋子；于是她信心倍增，相信祂也会轻易洗去她灵魂的不洁。祂也并非无故提及那城市伯达尼，而是为使你知道，祂是出于自己的意愿而走向受难。因为先前祂曾从他们中间逃开；那时，当他们的嫉妒最为炽烈时，祂却走近到大约十五弗隆以内；祂先前的退避，完全是安排的一环。\par

那女人看见祂，便从那里得到信心而来到祂跟前。因为那患血漏的妇人，虽然自己并不觉得有什么类似的罪，却因那种自然的不洁而战战兢兢地靠近祂；这个女人更可能因她的邪恶良心而迟疑退缩。因此，她是在许多妇人之后，即在撒玛黎雅妇人、客纳罕妇人、患血漏的妇人以及其他妇人之后，才来到祂跟前，因为她知道自己有许多不洁。而且，她不是在公开场合，而是在屋子里。而其他所有人都只是为了身体的治愈而来，她却是为了尊崇祂、为了灵魂的改善而来。因为她的身体根本没有受苦，所以这一点尤其令人惊叹。\par

她不是将祂当作一个普通人而来；因为那样她就不会用她的头发擦祂的脚，而是当作一位超乎人的存在。因此，她把全身最尊贵的部分，即她自己的头，放在了基督的脚下。\par

\Quote{当祂的门徒们看见时，就很不高兴——这是经上的话——说：\InnerQuote{为什么这样浪费？这香液原可卖得许多钱，施舍给穷人。}但耶稣知道了，便说：\InnerQuote{你们为什么叫这个女人难受？她在我身上做了一件善事。你们常有穷人同你们在一起，至于我，你们却不常有。她把这香液倒在我身上，是为安葬我而做的。我实在告诉你们：将来在普天之下，无论在什么地方传扬这福音，必要述说她所做的事，来纪念她。}}\par

那么他们的这种想法从何而来？他们常听到他们的主说：\BibleQuote{我喜欢仁爱，胜过祭献，}又责备犹太人，因为他们忽略了最重要的：公义、仁爱与信德，并在山上讲了有关施舍的许多话。从这些事，他们自己推断并认为，如果祂不接受全燔祭和古代敬礼，那么祂更不会接受香油的傅抹。\par

虽然他们这样想，但祂知道她的意图，就容许了她。因为她的虔诚确实伟大，她的热忱无以言表；因此，出于这种极大的屈尊，祂甚至容许香液倒在祂的头上。\par

因为如果祂不拒绝成为人，不拒绝在母胎中孕育，不拒绝吮吸母乳，你们又何必惊奇祂不彻底拒绝这事呢？因为正如圣父曾接受肉类的香气和烟雾，同样祂也接受这个罪妇，接纳（正如我已说过）她的心意。因为雅各伯也曾为天主傅油立石，祭献中也献过油，司祭们也都用香液受傅。\par

但是门徒们不知道她的目的，不合时宜地指责她，而他们对她提出的指控，反而显出这个女人的慷慨。因为他们说，这香液可以卖三百银币，这就表明这女人在香液上花了多少，显示了她多大的慷慨。因此祂也责备他们，说：\BibleQuote{你们为什么叫这个女人难受？}祂又加上一个理由，因为祂要再次使他们记起祂的受难：\BibleQuote{她做这事，是为安葬我。}还有一个理由：\BibleQuote{你们常有穷人同你们在一起，至于我，你们却不常有。}以及：\BibleQuote{将来在普天之下，无论在什么地方传扬这福音，必要述说她所做的事，来纪念她。}\par

你看见祂又是如何预先宣告福音传到外邦人吗？以此也安慰他们关于祂的死亡，因为十字架以后祂的能力如此彰显，以致福音要传到地极。\par

那么谁如此可怜，竟敢反对如此多的真理？因为看！祂所说的已经成就；无论你走到地球的哪个角落，你都会看到她被纪念。\par

但如此行事的这个人并非著名人物，所做之事也没有许多见证人，也不是在剧场中，而是在一所屋子里发生的，且这屋子是一个癞病人的家，只有门徒在场。\par

2. 那么，是谁宣告了它，并使它传扬开来？是正在说这些话的那一位的能力。尽管无数君王和将领的功业——甚至那些留有纪念的——都已消逝于沉默；他们攻陷城市，以城墙围困，树立战利品，奴役许多民族，却连传闻都不为人知，甚至名字也不被知晓，尽管他们立了雕像，制定了法律；然而，一个曾是妓女的女子，在某个癞病人的家里，在十个人面前倾倒了香液，这事全世界都在传颂；如此漫长的时间过去，那事的记忆仍未褪色，反而波斯人、印度人、西古提人、色雷斯人、萨尔马提人、摩尔人的种族，以及居住在不列颠群岛的人，都在传扬那由一个曾是妓女的女子在屋内秘密所作的事。\par

主的慈爱何其伟大。祂容忍一个妓女，一个妓女亲吻祂的脚，用香液抹湿，用她的头发擦拭，祂接纳了她，并斥责那些责怪她的人。因为如此大的热诚，那女子不应被驱入绝望。\par

但你也要注意这点：他们现在如何超脱于世，在施舍上如何前进。为什么祂不只是说\BibleQuote{她做了一件善事}，而是先说了\BibleQuote{你们为什么为难这女人？}？为使他们学到，起初不应向软弱者要求过高的原则。因此，祂也不仅仅孤立地审视这行为，而是考虑到这女子的身份。的确，如果祂是在立法，祂就不会提到这女子；但为使你们知道，是为了她的缘故才说这些话，以免他们破坏她初萌的信德，反倒要培育它，因此祂这样说，教导我们：任何人所作的善事，即使不完全，也要接纳它、鼓励它、促进它，而不是在起初就要求所有的完美。因为祂自己宁愿这样，这从祂没有枕头的地方却要求有一个钱囊的事实显明了。但那时的时间并不要求祂纠正这行为，而只是接纳它。因为正如如果有人在她未做这事时问祂，祂不会赞同；同样，在她做了之后，祂只关注一件事：她不要因门徒们的责备而陷入困惑，而是应从祂的关心中离开，变得更加愉快和更好。的确，在香液倾倒之后，他们的责备已不合时宜。\par

因此，你也要同样：如果你看见有人提供圣器并奉献，并乐意劳碌于教会其他装饰，如墙壁或地板；不要命令已造的变卖或推倒，以免挫伤他的热诚。但如果他在提供之前告诉你，命令他施舍给穷人；因为祂这样做也是为了不挫伤那女子的精神，祂所说的一切，都是为了安慰她。\par

因为祂曾说过\BibleQuote{她做这事是为我的安葬}；为使祂不因提到这样的事——即祂的埋葬和死亡——而使那女子困惑，请看祂如何借着下文安慰她，说：\BibleQuote{她所做的，将在全世界被传扬。}\par

这对门徒们既是安慰，对她也是安慰和赞扬。因为祂说，所有人将来都要赞扬她；现在她也预先宣告了我的受难，给我带来了葬礼所需之物，所以谁都不要责备她。因为我非但不谴责她做错了，不责备她没有做对，反而不会让所做之事隐藏，而是要让全世界知道在屋内秘密所作的事。因为这种行为确实来自一颗虔敬的心、炽热的信德和痛悔的灵魂。\par

为什么祂没有许诺那女子什么属灵的事物，而是永久的纪念？借此祂使她确信其他事情也必得报偿。因为如果她做了一件善事，显然她将得到相应的赏报。\par

\BibleQuote{那时，那十二人中的一个，名叫犹达斯依斯加略的，去见司祭长，对他们说：你们愿意给我什么，我就把祂交给你们？}那时。何时？当这些话被说出，当祂曾说过\BibleQuote{这是为我的安葬}，甚至这样也没有使他动心，当他听说福音将要被传扬到各处时，他也不害怕（尽管这是不可言说的大能之言），但当妇女们——且是曾为妓女的妇女——表现出如此大的尊敬时，他就行了魔鬼的作为。\par

但他们为什么提到他的绰号？因为还有另一个犹达斯。他们毫不避讳地说他是十二人之一；他们完全隐藏那些看似可耻的事。然而他们本可以只说他是门徒之一，因为还有别的门徒。但现在他们加上\Quote{十二人之一}，仿佛是说，从被选为最好的一群人中，从与伯多禄和若望在一起的那些人中。因为他们只关心一件事，那就是真理，而不是隐瞒所行的事。\par

因此，许多迹象他们略过不提，但那些看似可耻的事，他们毫不隐瞒；无论是话语、行为，还是任何这一类的，他们都满怀信心地宣扬。\par

3. 不仅这些，连那个宣讲更高深道理的人，若望本人，也是如此。因为他最多地告诉我们加于祂的侮辱和可耻之事。\par

请看犹达斯的恶意多么大：他主动去见他们，为钱财做这事，而且是那样一笔钱。\par

但路加说，他与首领们商议。因为在犹太人变得叛乱之后，罗马人委派了那些负责维护他们秩序的人。因为根据预言，他们的政府已经发生了变化。\par

於是他到他們那裏說：\Quote{你們願意給我什麼，我把他交給你們。}他們約定給他三十塊銀錢。從那時起，他便找機會出賣祂。原來他害怕群眾，想要單獨捉住祂。\par

哦，瘋狂！貪財如何完全蒙蔽了他！因為他屢次看見祂穿過人群而不被捉拿，並顯示了許多天主性與權能的證據，卻仍想捉拿祂；而祂還用許多可畏和溫柔的話如同咒語一般來對付他，以消除這邪惡的念頭。因為甚至在晚餐時，祂也沒有停止對他的關懷，直到最後一天還在談論這些事。但他毫無受益。然而主並沒有因此停止履行自己的職分。\par

知道這一點，讓我們也不要停止對那些犯罪和懈怠的人做一切可能的事：警告、教導、勸勉、告誡、建議，即使我們毫無益處。因為基督預知叛徒無可救藥，但祂仍不斷供應自己所能做的：既警告又威脅，還為他哀哭，且從不明說或公開，而是隱晦地。就在出賣的那一刻，祂甚至允許他親吻祂，但這對他毫無益處。貪財是何等大的惡事！它使他成了叛徒和褻聖的強盜。\par

請聽，你們這些貪財的人，你們這些患有猶達斯病的人；請聽，並警醒這災禍。因為那與基督同在、行奇蹟、受益於如此多教導的人，因未脫離這病而墮入如此深淵；何況你們，連聖經也不聽，終日緊貼於眼前事物，若無持續的警醒，豈不更容易成為這災禍的獵物？那個人天天與祂——連枕頭的地方都沒有的那一位——同在，天天受教於言行，不要有金子、銀子、兩件衣裳；但他仍學不會節制；你們若無認真關注和不懈努力，怎能指望逃脫這病？因為這怪物實在可怕，然而，只要你願意，你必能輕易勝過它。因為這慾望不是本性的；這從那些不受它困擾的人可以顯明。因為本性的事物是眾人所共有的；但這慾望僅源於懈怠；由此而生，由此而長，當它抓住那些貪婪注視它的人時，便使他們違背本性生活。因為當他們不顧同胞、朋友、兄弟，總之一切人，甚至不顧自己時，這便是違背本性生活。由此可見，貪財的惡習和疾病——猶達斯深陷其中而成了叛徒——是違背本性的。你或許會問：他既被基督召喚，如何成了這樣的人？因為天主的召喚不是強迫的，也不強迫那些不願選擇德行者的意志，而是告誡、勸導、做一切、安排一切，以說服人成為善；但如果有些人不忍受，祂並不強迫。但若你想知道他為何成了這樣的人，你會發現他是被貪財毀滅的。\par

有人會問：他如何被這災禍纏住？因為他懈怠了。因為這樣的改變源於懈怠，正如另一方面，向善的改變源於勤奮。有多少從前兇暴的人，如今比羔羊還溫順？有多少好色的人後來成了節制者？有多少從前貪財的人，如今甚至捨棄了自己的財產？相反的情形也因懈怠而發生。因為革哈齊也曾與聖人同住，他也因同樣的病而墮落。因為這災禍是最嚴重的一種。由此而產生盜墓者、殺人者、戰爭和鬥爭，以及你所能提及的一切邪惡。這樣的人在各個方面都無用，無論是帶領軍隊還是治理百姓；不僅在公共事務上，在私事上也是如此。若要娶妻，他不會娶賢德的女子，而是娶最卑賤的；若要買房，不會買適合自由人的，而是買能帶來高租金的；若要買奴隸或其他東西，他會挑最差的。\par

何必提及帶兵、治民、理家呢？因為即使他是一國之君，他也是所有人中最可憐的，是世間的禍害，而且是最窮困的。他會覺得自己與常人無異，不把眾人的財物視為己有，反而把自己當作眾人中的一員；當掠奪所有人的財物時，他認為自己擁有的比任何人都少。因為他以對未得之物的慾望去衡量眼前之物，便視前者為無物。因此有一位說：\Quote{沒有比貪財的人更邪惡的了。}（\BibleRef{德 10:9}）\par

4. 因为这样的人既出卖自己，又四处走动，成为世界的公敌；他因大地不产金子而产谷物、泉源不流金子而流清水、山岭不产金子而产石头而忧伤；他为季节的丰收而烦恼，为公共的恩惠而困扰；他回避一切不能获得钱财的途径，忍受一切能搜刮到哪怕两个小钱的事；他仇恨所有的人，穷人和富人：穷人，免得他们来向他乞讨；富人，因为他不拥有他们的财产。他认为所有人都拥有属于他的东西，仿佛被所有人伤害了，因而对所有人都感到不满。他不知满足，从未体验过饱足，他比任何人都更悲惨；另一方面，摆脱了这些事并实践节制的人，则是最令人羡慕的。因为有德之人，即使他是仆人、囚犯，也是最幸福的人。因为没有人能伤害他，即使全世界的所有人聚集起来，发动军队和营垒与他征战，也不能伤害他。但堕落卑劣、如我们所描述的那样的人，即使他是国王，戴着一千个王冠，也会遭受最极端的苦难，甚至出自普通人之手。恶习如此软弱，美德如此坚强。\par

那么，你为何因处于贫困而哀伤？你为何在庆祝节日时哭泣？因为这确实是节庆的场合。你为何哭泣，你这小孩子？因为我们应当称这样的人为小孩子。有人打了你？这算什么？他使你更能忍耐了？但他夺去了你的钱财？他移去了你大部分的重担。但他剥夺了你的荣誉？你又对我谈及另一种自由。连那些没有智慧教导的人也听见了这些事，并说：\Quote{你若不在意它，便未受损害。}但他夺去了你那有围墙环绕的大宅？看，整个大地都在你面前，公共建筑，无论你为享乐还是为实用。还有什么比穹苍更令人愉悦、更美丽呢？\par

你要贫穷困乏到几时呢？灵魂不富有的人，不可能富有；正如心中有贫困的人，不可能贫穷一样。因为灵魂比身体更高贵，较低贵的部分不能按照自身影响它；相反，高贵的部分将自己吸引过来，并改变那些不那么高贵的部分。因为心脏，当它受到任何伤害时，会相应地影响全身；若其状态失调，则会败坏一切；若其状态良好，则有益于一切。即使其余部分中有些变得败坏，只要心脏健全，它也能轻易摆脱它们中的邪恶。\par

为了使我所说的更加明白，我请问你，当根枯萎时，青翠的树枝有何用？而当根健全时，上面的叶子枯萎又有何害？同样，在这里，灵魂贫穷时，金钱毫无用处；灵魂富有时，贫穷也无害。或许有人说：灵魂既缺钱，怎能富有？此时最可能如此；因为那时它也往往是富有的。\par

因为，如我们常所示，轻看财富、一无所缺，乃是富足的确实证据；而贫穷则是有所欠缺；并且人在贫穷中比在财富中更容易轻看金钱。显然，处于贫困反倒使人富有。因为富人比穷人更看重金钱，这对每个人都是明摆着的；正如醉酒的人比喝够了的人更渴。因为他的欲望不会因过多而止息；反而，其本性是因此被点燃。火也一样，得到更多燃料时，便最为猛烈；财富的暴政，当你投入更多金子时，尤其加剧。\par

那么，若贪求更多是贫穷的标记，而拥有财富的人正是如此，他便尤其贫穷。你看见了吗？灵魂正是在富有时最贫穷，在贫穷时最富有？\par

若你愿意，让我们也在具体人物身上运用推理。设有两人，一人拥有一万塔冷通，另一人拥有十塔冷通，从两人都夺去这些。谁会更悲伤？失去一万的人。但他若非更爱它，便不会更悲伤；但若他爱得更深，便贪求更多；若贪求更多，便更贫穷。因为我们最渴望的，正是我们最欠缺的；因为渴望出于贫乏。因为满足之处，不可能有渴望。因为我们最渴时，正是最需要水时。\par

我说这一切，为表明若我们警醒，无人能伤害我们；伤害并非来自贫穷，而是来自我们自己。因此，我恳请你们尽心竭力，除去贪吝的祸害，好使我们在此既得富有，又享永生之美。愿天主恩赐我们众人达此境界，赖我们的主耶稣基督的恩宠与对人的慈爱，愿光荣归于他，永无穷尽。阿们。\par

\SourceRecord{Translated by George Prevost and revised by M.B. Riddle.；Nicene and Post-Nicene Fathers, First Series；Vol. 10.；Edited by Philip Schaff.；Buffalo, NY: Christian Literature Publishing Co.,；1888.；<http://www.newadvent.org/fathers/200180.htm>.}

% structure: p0001 source=h1 target=section reason=page-title

\section{第八十一讲：玛窦福音释义}

\Emph{\BibleRef{玛 26:17-18}}\par

\Quote{无酵节的第一天，门徒来问耶稣说：\InnerQuote{你愿我们往哪里去为你预备吃逾越节晚餐呢？}耶稣说：\InnerQuote{你们进城去见某人，对他说：\InnerQuote{老师说了，我的时候快到了，我要在你家里同我的门徒守逾越节。}}}\par

他所说的\Quote{无酵节的第一天}，是指该节日前的那一天；因为他们惯常从傍晚开始算日子，他所提及的正是逾越节羔羊必须在傍晚被宰杀的那一天（\BibleRef{若 13:1}）；因为在那一周的第五天，他们来到祂面前。这位（玛窦）称之为无酵节的前一日（\BibleRef{路 22:7}），指他们来到祂面前的时候；另一位（路加）则这样说：\Quote{那时无酵节的日子到了，逾越节的羔羊必须被宰杀；}他用\Quote{到了}这个词表示它临近了，就在门口，明显是指那个傍晚。因为他们从傍晚开始，所以每一位都补充说\Quote{当逾越节被宰杀的时候}。\par

他们问道：\Quote{你愿我们往哪里去为你预备吃逾越节晚餐呢？}由此也可明显看出，祂既无房子，也无寄居之处；我想他们也没有。因为他们若有，必定会请祂到那里去。但他们也没有，因为已经舍弃了一切。\par

但祂为什么守逾越节呢？为了一直到最后一天，透过一切事表明祂并不反对法律。\par

祂为何派遣他们到一个陌生人那里去呢？借此也要表明，祂本可避免受苦。因为祂竟然能使这人的心转变，以致接待他们，且是藉着言语；若是祂不愿意受苦，对那些钉死祂的人，祂还有什么不能作的呢？关于驴子的事，祂在这里也是如此行。因为在那里祂也说：\Quote{若有人对你们说什么，你们就说：主需要它们；}（\BibleRef{玛 21:3}）同样，在这里祂也说：\Quote{老师说了：我要在你家里守逾越节。}我不仅惊奇这人不认识祂却接待了祂，而且他料到自己会招致如此仇恨和无法和解的敌意，却藐视群众的敌意。\par

此后，因为他们不认识祂，祂便给他们一个记号，如同先知论到撒乌耳时所说：\Quote{你必遇见一个人上来，带着一瓶}（\BibleRef{撒上 10:3}）；在这里是\Quote{拿着一个水罐}。请看祂权能的再一次彰显。因为祂不仅说\Quote{我要守逾越节}，还加上了另一件事：\Quote{我的时候近了。}祂这样做，一方面不断提醒门徒受难的事，好让经常的预告使他们为将要发生的事做好准备；同时也要向他们本人、接待祂的人以及所有的犹太人（我常提及的）表明，祂并非不情愿地来受难。祂又加上了\Quote{同我的门徒}，以便预备充足，且使那人不会以为祂在隐藏自己。\par

\Quote{到了晚上，祂与十二门徒坐席。}哦，犹大之无耻！因为他也在那里，并前来领受奥迹和筵席，就在席上被定罪，纵使他是野兽，也会变得温顺。\par

为此，福音作者也表明，当他们正在吃的时候，基督讲论自己的背叛，好藉着时间和筵席，揭露叛徒的恶行。\par

当门徒照着耶稣所吩咐的作了之后，\Quote{到了晚上，祂与十二门徒坐席。正吃的时候，祂说：}我们读到，\Quote{我实在告诉你们，你们中间有一个人要出卖我。}在晚餐前，祂甚至洗了他们的脚。请看祂怎样宽恕叛徒。祂没有说\Quote{某人要出卖我}，而是说\Quote{你们中间有一个人}，这样藉着隐晦，再次给他悔改的机会。祂选择使所有人惊慌，为要拯救这个人。祂说：\Quote{你们十二人中，那常与我同在、我曾洗你们的脚、并应许你们许多事的。}\par

若望说他们\Quote{彼此相顾，猜疑}，各人虽然自己心里并无此事，却恐惧地问。但这位（玛窦）福音作者说：\Quote{他们就非常忧愁，一个一个地问祂说：\InnerQuote{主，是我吗？}祂回答说：\InnerQuote{同我蘸手在盘子里的，就是他要出卖我。}}（\BibleRef{玛 26:22}）\par

请注意祂在何时揭露了他。那是在祂愿意使其余的人脱离这困扰的时候，因为他们已吓得要死，所以急迫地追问。但这不仅是祂要使他们脱离痛苦，也是要改正叛徒。因为他在屡次泛泛地听到后仍执迷不悟，麻木不仁，祂便想使他更有所感觉，于是揭穿了他的面具。\par

当他们忧愁地开始问：\Quote{\InnerQuote{主，是我吗？}祂回答说：\InnerQuote{同我蘸手在盘子里的，就是他要出卖我。人子固然要按着指着祂所记载的而去，但出卖人子的那人却有祸了！那人若没有生下来，为他倒好。}}\par

现在有人说，他竟如此大胆，竟敢不尊敬他的老师，而与祂同蘸手；但在我看来，基督这样做也是为了更进一步羞辱他，使他归向更好的心态。因为这一举动还有更深的意义。\par

2. 但这些事我们不可随便放过，而应铭刻在心，愤怒就绝不会在任何时候有立足之地。\par

因为谁能记起那晚餐和坐在万有救主席上的叛徒，以及那将要被出卖的基督如此温良地论述，而不除去一切毒害的忿怒和怒气呢？请看祂对他是多么温良：\Quote{人子固然要按着指着祂所记载的而去。}\par

祂又说这些话，既是为安慰门徒，免得他们以为此事是软弱的迹象，也是为纠正那负卖者。\par

\Quote{但负卖人子的人有祸了！那人若没有生下来，为他倒好。}请看，在祂的责备中，又显出祂难以言说的温良。因为即使在此，祂也不是用谩骂，而是更以怜悯的方式说出祂的话，却又再次隐晦地表达；然而，不仅是他先前的麻木，还有他后来的无耻，都理应激起最大的愤慨。因为在他被如此定罪之后，他说：\Quote{主，是我吗？}哦，麻木！他明知自己做了这样的事，却还来询问。因为福音作者也惊叹于他的厚颜，才这样写。那么，最温和良善的耶稣说了什么？\Quote{你说的是。}其实祂本可以说：你这不洁的，全不洁的；该受诅咒的，亵渎的；你心怀恶念如此之久，走你的路，与撒殚订立契约，同意收受银钱，又被我揭穿，你还敢问吗？但祂一句这样的话也没有说，而是怎样？\Quote{你说的是？}为我们立定了忍耐的界限和规矩。\par

但有人会说：如果经上写着祂要受这些苦，那为何要责备犹达斯呢？因为他所做的事正是经上所写的。但并非出于此意，而是出于邪恶。因为如果你不去探究动机，你甚至会把魔鬼从对他的指控中释放出来。但这些事不是这样，绝不是。因为两者都该受无数惩罚，即使世界得到了救恩。因为拯救我们的，不是犹达斯的背叛，而是基督的智慧，以及祂巧妙利用他人的邪恶来为我们获益的巧妙安排。\par

\Quote{那么，}有人会说，\Quote{即使犹达斯没有出卖祂，难道别人不会出卖祂吗？}这与问题何干？\Quote{因为如果基督必须被钉十字架，那必然要借着某个人；如果借着某个人，那必定是像这样的人。但如果所有人都善良，为我们所做的救世工程就会受到阻碍。}并非如此。因为全知者知道如何成就我们的益处，即使发生这种事。因为祂的智慧富于筹划，难以理解。因此，为免有人以为犹达斯成了救世工程的工具，祂宣告了那人的不幸。但有人又会说：\Quote{如果那人没有生下来倒好，为何祂还允许这个人以及所有恶人来到世上？}当你本应责备恶人，因为他们本可以不成为那样，却成了恶人，你却放弃这一点，忙碌于探究天主的作为；虽然你知道没有人是因必然性而成为恶人的。\par

\Quote{但只应有善人生下来，}他会说，\Quote{就不需要地狱、惩罚、报应，也没有罪恶的痕迹；恶人要么根本不该出生，要么出生后就立刻死去。}\par

首先，最好是重复宗徒的话：\Quote{然而，人啊，你是谁，竟敢向天主抗辩？受造物岂能对造物主说：你为什么这样造了我？}\par

但如果你仍要求理由，我们可以说，善人在恶人中更受钦佩，因为那时他们的忍耐和伟大的自制最能显明。但你若说这些话，便剥夺了他们角力与争斗的机会。\Quote{那么，为了显明这些人是善良的，其他人就要受惩罚吗？}他说。绝不可这样！他们是因自己的邪恶而受罚。因为他们并非因来到世上而成为恶人，而是因自己的邪恶；因此他们受罚。因为怎么不该受罚呢？他们有这么多的德行导师，却一无所获。正如高尚善良的人应得双倍的尊荣，因为他们既成为善良的，又未从恶人那里受到伤害；同样，无价值的人也当受双倍的惩罚，因为他们本可成为善良的（那些成了善良的人证明了这一点），却成了恶人，并且从善人那里一无所获。\par

但让我们看这可怜的人，当他被师傅定罪时说了什么。那么他说了什么？\Quote{辣彼，是我吗？}\BibleRef{玛 26:25}为什么他不是一开始就问？他以为说\Quote{你们中间有一个人}就能逃避被认出；但当祂把他显明出来后，他又冒险再问，信赖师傅的仁慈，以为祂不会定他的罪。\par

3. 哦，盲目！它把他引到了哪里？贪婪就是这样，它使人愚蠢、麻木，甚至鲁莽，使人在为犬之后，变成魔鬼，甚至比狗更凶残。这个人至少在他图谋不轨时接纳了魔鬼，但耶稣即使对他行善，他也出卖了，他在意志上已经成了魔鬼。因为贪得无厌使人如此：失去理智，疯狂，完全沉溺于利益，就像犹达斯的情形一样。\par

但玛窦和其他福音作者怎么说？他们说，当他商议背叛之事时，魔鬼就抓住了他；但若望却说\Quote{那饼以后，撒殚进入了他。}\BibleRef{若 13:27}若望自己知道这一点，因为他在更早的地方说：\Quote{魔鬼已把出卖祂的念头放入犹达斯心中。}那么他怎么说\Quote{那饼以后，撒殚进入了他}？因为撒殚不是突然一下子进入，而是先多次试探；这里也是如此。因为他一开始试探他，悄悄攻击他，之后见他准备好接纳他，就完全将自己注入他里面，彻底控制了他。\par

但若他们在吃逾越节晚餐，他们是否违犯了法律？因为他们不应该坐着吃饭。\BibleRef{出 12:11}那么可以这样说：他们吃过逾越节晚餐后，这才坐下赴宴。\par

但另一位福音作者说，那日晚间祂不仅吃了逾越节晚餐，还说：\Quote{我渴望又渴望，同你们吃这一次逾越节晚餐。}\BibleRef{路 22:15} 即在那一年。为什麼呢？因为那时世界的救恩即将实现，奥迹要交付，悲伤的事要因祂的死亡而消除；十字架对祂而言如此可亲。但什麼也无法软化这野蛮的怪物，无法打动或羞辱他。主称他为有祸的，说：\Quote{那人是有祸的。}祂又警醒他，说：\Quote{那人若没有出生，为他更好。}祂羞辱他，说：\Quote{我蘸了饼递给谁，谁就是。}但没有一件能制止他；他被贪财所控制，如同被某种疯狂，或更确切地说，被一种更严重的疾病所控制。因为这种疯狂更为严重。\par

因为疯狂的人会做出这样的事吗？他不是从口中吐出泡沫，而是吐出了对其主的谋杀。他没有扭曲双手，却伸手索要宝血的代价。因此他的疯狂更大，因为他在健康时行了疯狂。\par

但他没有说出\Emph{你说}这样无意义的声音。而什麼比这话更无意义呢？\Quote{你们愿意给我什麼，我把他交给你们？}\BibleRef{玛 26:15} \Quote{我要交给你们，}魔鬼藉那口说话。但他没有用脚挣扎着踢地？不，这样挣扎比直立站著要好得多。但你会说，他没有用石头割自己？然而，做这样的事比起来又如何！\par

你们愿意我们把附魔者和贪财者带来，在二者之间作比较吗？但不要有人将所做的事视为对自己的羞辱。因为不是责难本性，而是哀叹行为。附魔者从不穿衣，用石头割自己，奔跑著，冲过崎岖的道路，被魔鬼驱赶。这些事岂不看来可怕吗？那麼，如果我证明贪财者对自己的灵魂做了更严重的事，而且严重得多，使这些附魔者的行为相比之下如同儿戏，又当如何？你们岂不真要避开这祸害吗？那么，让我们看看，他们在任何方面是否比附魔者处于更可容忍的状态。完全不，甚至处于更糟的状态；因为他们比一万个赤身者更为可耻。因为穿着贪财的果实四处走动，如同庆祝酒神狂欢节的人，远比赤身露体更为可耻。正如附魔者戴着疯人的面具和衣服，这些人也是如此。而且正如附魔者的赤身是由疯狂所致，这种衣服也由疯狂产生，并且这衣服比赤身更可怜。\par

我将借此尽力证明这一点。因为在疯人之中，我们应说谁更疯：是那割自己的，还是那既割自己又伤害遇见他的人？显然，是后者。那么，附魔者脱去自己的衣服，而这些家伙脱去所有遇见者的衣服。\Quote{但这些家伙撕裂别人的衣服。} 每一个受伤害的人何等愿意自己的衣服被撕裂，而不是被剥夺全部财产？\par

\Quote{但附魔者不打别人的脸。} 首先，贪财者确实这样做；即使不是全部，但所有贪财者都藉饥饿和贫困使肚腹遭受更剧烈的痛苦。\par

\Quote{但附魔者不用牙齿咬人。} 但愿只是用牙齿，而不是用比牙齿更锋利的贪财之箭。\Quote{因为他们的牙齿是武器和箭矢。} 谁会感到更痛苦：是被咬一次立刻痊愈的人，还是被贫困的牙齿永远啃噬的人？因为非自愿的贫困比火炉或野兽更残酷。\par

\Quote{但附魔者不像被魔鬼附身的人那样冲进荒野。} 但愿他们衝进的是荒野，而不是城市，那麼城中所有人都能享安全。因为现在，在这方面他们再次比所有疯人更不可容忍：他们在城市里做那些别人在荒野里做的事，使城市变成荒野，如同在无人阻拦的荒野中，掠夺所有人的财物。\par

\Quote{但附魔者不用石头扔遇见他们的人。} 这算什麼？石头容易提防；但藉由纸张和墨水对那些可怜穷人造成的创伤（写出满是无数打击的文件），遇见他们的人中，有谁能轻易提防呢？\par

4. 让我们也看看他们对自己做了什麼。他们赤身裸体地在城市里走来走去，因为他们没有德行这衣服。但如果这在他們看来不算耻辱，这又是他们特别疯狂的原因：他们不觉得不体面，反而为身体赤裸感到羞耻，却带着赤裸的灵魂而夸耀。如果你愿意，我也告诉你他们麻木的原因。是什麼原因呢？他们在一群同样赤裸的人中赤裸，因此不觉羞耻，正如我们在澡堂中也不觉羞耻。所以，如果确实有很多人穿着德行，那么他们的羞耻就会更明显。但现在，最值得流泪的是，因为坏人众多，坏事甚至不被视为耻辱。因为魔鬼除了其他事，还造成了这一点：不让他们获得对自己恶行的任何感觉，反而藉作恶者的众多，掩盖了他们的耻辱；因为如果这贪财者置身于一群克己者中间，他就会更清楚地看到自己的赤裸。\par

由这些事可知，他们比附魔者更为赤露；至于他们走入旷野，这一点也没有人能反驳。因为宽阔的大路比任何旷野都更为荒凉。虽然有许多人在其上行走，但那不是人，而是蛇、蝎子、狼、毒蛇和蝮蛇。行恶的人就是这样。这条路不仅荒凉，而且比\Emph{疯人的}路更为崎岖。这一点也很明显：石头、沟壑和悬崖不会像抢劫和贪婪那样伤害攀登者，它们伤害的是那行这些事的灵魂。\par

他们像附魔者那样住在坟墓旁，或者更确切地说，他们自己就是坟墓，这一点也很明显。坟墓是什么？是一块石头，里面躺着一具尸体。那么这些人的身体与那些石头有何不同？不如说，他们比石头更悲惨。因为不是石头装着一具尸体，而是比石头更麻木的身体带着一个死去的灵魂。因此，称他们为坟墓并不为过。我们的主也是这样称呼犹太人，尤其是因为这一点；祂接着说：\Quote{他们内里充满劫掠和贪婪。}\par

你们愿意我接下来指出他们如何用石头击打自己的头吗？那么，请问，你们想从哪里学到这一点？从今世的事，还是从将来的事？但他们不大关心将来的事；所以我们必须谈论今世的事。焦虑不是比许多石头更沉重吗？它们不击打头，却消耗灵魂。他们害怕那些不义得来的东西会公正地离开他们的家；他们为极端的祸患而战栗；他们发怒、激愤，攻击自己家里的人和外人；现在沮丧，现在恐惧，现在愤怒，接连不断；他们好像越过一个又一个悬崖，日日急切地追求尚未得到的东西。因此，他们不为已拥有的感到快乐，既因为对安全不放心，也因为全心全意地想着尚未攫取的东西。就像一个不断口渴的人，即使喝下一万个泉源，也不感到快乐，因为不满足；同样，这些人非但不感到快乐，反而越发痛苦，他们越聚敛越多，因为这种欲望永无止境。\par

今世的情况是这样；但让我们也谈谈将来的日子。虽然他们不留意，但我们还是必须说。在将来的日子，会看到这样的人到处受罚。因为当祂说：\Quote{我饿了，你们没有给我吃的；我渴了，你们没有给我喝的；}\BibleRef{玛 25:42}祂正是在惩罚这些人；当祂说：\Quote{离开这里，到那为魔鬼预备的永火里去吧，}祂正是把滥用财富的人送到那里。那个不把主人的财物分给同伴的恶仆，也属于这些人；还有埋了塔冷通的仆人，以及那五个童女。\par

无论你走到哪里，都会看到贪婪的人受罚。时而他们会听见：\Quote{在我们与你们之间有一个深渊；}时而：\Quote{离开我，到那预备的火里去。}\BibleRef{玛 25:41}他们被割裂后离去，到那有咬牙切齿的地方；从各处可见他们被驱赶，无处容身，唯独聚在地狱里。\par

5. 当我们听到这些事时，正确的信仰对我们的得救有何益处？在那里，有咬牙切齿、外边的黑暗、为魔鬼预备的火、被割裂、被驱逐；在这里，有仇恨、毁谤、诽谤、危险、忧虑、阴谋、被众人憎恨、被所有人厌弃，甚至包括那些似乎奉承我们的人。因为正如善人不仅被善人钦佩，甚至也被恶人钦佩；同样，恶人不仅被善人憎恨，也被卑鄙的人憎恨。为证明这是真的，我乐意问一问贪婪的人：他们彼此之间是否感到痛苦？是否把这些贪婪的人视为比那些曾大大伤害他们的人更大的敌人？他们是否也自责？如果有人以此责备他们，他们是否认为这是一种侮辱？因为确实这是一种极端的责备，是许多邪恶的确实证据；如果你不能忍受轻视财富，那么你能胜过什么？能胜过情欲、或对光荣的疯狂渴望、或愤怒、或暴怒吗？谁能被说服呢？因为至于情欲、愤怒和暴怒，许多人甚至归咎于肉体的气质，医学研究者也将它们的过度归因于此；他们断言具有更热、更怠惰气质的人更好色；但倾向于更干燥的恶劣气质的人则急切、易怒、暴躁。至于贪婪，从未听说他们说过这样的话。因此，这种灾祸完全是懈怠和麻木灵魂的结果。\par

因此，我恳求你们，让我们努力改正这一切，并在每一个年龄都给临到我们的情欲以相反的方向。但如果我们一生都在航行中错失德行的劳作，处处遭遇海难；那么当我们到达港口时，缺乏属灵的货载，我们将遭受极端的惩罚。因为今世的生活是一片广阔的海洋。如同在此海洋中，不同的海湾遭遇不同的风暴：\Place{爱琴海}因风而险恶，\Place{第勒尼安海峡}因狭窄而险恶，非洲附近的\Place{卡律布狄斯}因浅滩而险恶，\Place{普罗庞提斯}（黑海之外）因猛烈的海流而险恶，\Place{加的斯}以外的海域因荒凉、无路和未探明之地而险恶，其他部分则因其他原因而险恶。同样，在我们的生活中也是如此。\par

首先映入眼帘的是童年时期的海，它因愚昧、易变、不稳定而充满了风暴。因此，我们为它设置了向导和教师，以我们的勤勉补充自然的不足，就像那里用领航技术一样。\par

在此年龄之后，便是青年之海，其中的风暴如同在\Place{爱琴海}上一般猛烈，情欲在我们身上日益增长。这个年龄尤其缺乏修正；不仅因为他被更猛烈地围攻，而且他的过犯不受责备，因为教师和引导者在那之后都退去了。因此，当风暴更加猛烈地吹袭，舵手更加软弱，又没有帮手时，请思量这风暴的猛烈。\par

此后，又有另一个生命时期，即成年时期，那时家庭的忧虑压迫着我们，有妻子、婚姻、生养子女、管理家务，以及如雨般密集降下的忧虑。那时，贪婪和嫉妒尤其猖獗。\par

那么，当我们生命的每一部分都遭遇船难时，我们怎能应付现世的生活？我们怎能逃脱未来的惩罚？因为首先在最早的年龄，我们学不到任何有益的事物；然后在青年时期，我们不修习节制；待到成年，又不能战胜贪婪；到了老年，仿佛进入一个积满舱底污水的船舱，由于所有这些冲击，灵魂的小船已经朽坏，木板分离，我们将抵达那个港口，携带许多污秽而非属灵的货物，给魔鬼提供笑料，却给自己带来哀哭，并为自己招致不可忍受的惩罚。\par

为使这些事不致发生，让我们从各方面振作起来，抵挡我们的一切私欲偏情，摒弃对财富的贪恋，好能获得未来的美善，赖我们的主耶稣基督的恩宠和慈爱，愿光荣归于祂，至于无穷之世。阿们。\par

\SourceRecord{Translated by George Prevost and revised by M.B. Riddle.；Nicene and Post-Nicene Fathers, First Series；Vol. 10.；Edited by Philip Schaff.；Buffalo, NY: Christian Literature Publishing Co.,；1888.；<http://www.newadvent.org/fathers/200181.htm>.}

% structure: p0001 source=h1 target=section reason=page-title

\section{玛窦福音讲道 第八十二篇}

\BibleRef{玛 26:26-28}\par

\BibleQuote{\Emph{他们正吃的时候，耶稣拿起饼来，祝谢了，擘开，递给门徒说：\InnerQuote{你们拿去吃罢！这是我的身体。}}}\par

\BibleQuote{\Emph{他又拿起杯来，祝谢了，递给他们说：\InnerQuote{你们都由其中喝罢！因为这是我的血，新约的血，为大众倾流，以赦免罪过。}}}\par

唉！叛徒的盲目何其大！即使参与了奥迹，他仍然如故；被接纳到至圣的祭台前，他却没有改变。路加表明，此后撒殚进入了犹达斯，不是由于轻慢主的身体，而是从此嘲笑那叛徒的厚颜无耻。因为他的罪从两方面变得更大了：既以如此心态前来领受奥迹，又虽领受了却没有因恐惧、或恩惠、或尊荣而变得更好。但基督虽预知一切，却没有阻止他，好叫你学习祂不遗漏任何有助于改正的事。因此在这之前和之后，祂不断地劝诫他、制止他，借着行为与言语，借着恐惧与仁慈，借着威胁与尊荣。但这些事没有一样能使他脱离那严重的灾祸。\par

因此祂从此不再管他，而借着奥迹又提醒门徒祂将被杀，并在用餐中谈论十字架，通过不断重复预言，使祂的受难易于接受。因为如果已有如此多的事成就和预言，他们尚且困扰；若他们从未听到这些，他们会感到什么呢？\par

\Quote{\Emph{他们正吃的时候，祂拿起饼来，擘开了。}}为什么祂要在这个时候，即在逾越节时，订立这件圣事呢？好叫你从每件事上学习，祂既是旧约的立法者，而且旧约中的事是这些事的预象。因此，我说，哪里是预象，祂就把真理放在那里。\par

然而，傍晚是时期圆满的确切标记，表示诸事如今已到了尽头。\par

祂祝谢了，为教导我们应如何举行这件圣事，并表明祂并非不情愿地受难，也教导我们无论遭受什么，都要感恩地承受，由此也暗示了美好的希望。因为如果预象能带来那样的解脱，真理更将解放世界，而祂将为我们的种族而被交付。因此，我要补充说，祂之前并未设立这件圣事，而是在此后律法的礼仪将要停止的时候。这样，祂结束了节期的首脑，将它们转移到另一个最可畏的祭台，祂说：\Quote{你们拿去吃，这是我的身体，为许多人而擘开的。}\par

他们听到这个如何不感到困惑呢？因为祂之前已向他们讲了许多关乎此的大事情。因此祂不再建立什么，因为他们已充分听过了，但祂谈到祂受难的原因，即赦免罪恶。祂称它为\Quote{新约的血}，指的是保证、应许、新律法。因为祂从前也保证了这一点，而这就是包含在新律法中的约。如同旧约献祭羊羔和牛犊，这个有主的血。因此祂也表明祂快要死了，故也提到了约，并提醒他们前约，因为那也是用血奠定的。祂再次宣告祂死亡的原因：\Quote{为许多人倾流，以赦免罪过。}祂说：\Quote{你们要这样做，为纪念我。}你看到祂如何撤除并引导他们离开犹太习俗吗？就如你们在埃及奇迹的纪念中那样做，祂说，同样，你们也这样做以纪念我。那血是为保存长子而倾流，这血是为赦免全世界的罪。因为祂说：\Quote{这是我的血，为赦免罪过而倾流。}\par

祂这样说，以此指明祂的受难和十字架是一个奥迹，也借此再次安慰门徒。就如梅瑟说：\BibleQuote{这将是你们永远的纪念。}\BibleRef{出 12:14}同样，祂也说：\Quote{为纪念我}，直到我来。因此祂也说：\BibleQuote{我渴望又渴望，在这逾越节与你们同吃。}\BibleRef{路 22:15}即，要把新的礼仪交给你们，并给出一个逾越节，借此我要使你们成为属灵的。\par

祂自己也喝了。为了免得他们听到这个时说：\Quote{那么什么呢？我们喝血，吃肉？}然后感到困惑（因为当祂开始谈论这些时，甚至听了这话许多人就反感了），所以免得他们那时也困扰，祂先自己做了这事，引导他们平静地参与奥迹。因此祂喝了自己的血。那么，我们是否还必须遵守那古老的礼仪呢？有人会说。绝不。因为为此祂说：\Quote{你们这样做}，以撤除他们从另一个。因为如果这个带来罪赦，它确实带来，那另一个就是多余的了。\par

如同在犹太人中一样，祂在此也将恩惠的纪念与奥迹绑在一起，以此再次堵住异端者的口。因为当他们说：\Quote{从哪里看出基督被祭献了？}我们与其他论据一起，也从奥迹堵住他们的口。因为如果耶稣没有死，这些礼仪的象征是什么呢？\par

你看到用了多大努力，使祂为我们而死这件事常被记念吗？因为\OriginalTerm{马西昂派}{Marcionists}、\OriginalTerm{华伦提努派}{Valentinians}和\OriginalTerm{摩尼派}{Manichaeans}将要兴起，否认这一救恩计划，祂就借着奥迹不断提醒我们受难（使无人受欺骗），同时借着那神圣的祭台拯救和教导。因为这是恩惠的首端，因此保禄也处处强调这一点。\par

然後，當祂交付了之後，祂說：\Quote{從今以後，我不再喝這葡萄汁，直到我在我父的國裏同你們喝新的那日子。} 因為祂曾與他們談論受難和十字架，祂又引入祂要說的關於復活的事，在他們面前提起了國度、宴席，卻因醉酒而起立，其實理當感謝並以一首聖詠結束。\par

你們這些不等待奧跡最後祈禱的人，請聽這話，因為這是那事的象徵。祂在遞給門徒之前先感謝了，使我們也能感謝。祂在交付之後感謝並詠唱了聖詠，使我們也能做同樣的事。\par

但祂為何上山去？為了顯明自己，以便被捉拿，免得顯得是隱藏自己。因為祂急忙前往那地方，那地方猶達斯也知道。\par

然後\Quote{祂對他們說：你們都要為我的緣故而跌倒。}\BibleRef{瑪 26:31} 此後祂也提到一個預言：\Quote{因為經上記載：我要打擊牧人，羊群就四散。} 同時勸勉他們總要留意所記載的事，並顯明祂是按天主的旨意被釘死；藉這一切表明祂並不疏離舊約，也不疏離其中所傳揚的天主，而是所行的是天主的救恩計劃，並且先知們從起初就預先宣告了這一切涵蓋的事，好使他們也對更美好的事充滿信心。\par

祂教導我們認識門徒們在受難前是什麼樣子，受難後又是什麼樣子。因為確實，當祂被釘十字架時，他們連站都站不穩，這些人在祂死後卻變得剛強，比金剛石還堅硬。\par

這同一件事是祂死亡的明證——我指的是門徒們的驚慌和懦弱。因為如果當如此多的事情已經發生和說了之後，仍有一些無恥之徒說祂沒有被釘死；那麼如果這些事一件都沒有發生，他們會墮入何等程度的邪惡呢？所以，為此原因，祂不僅藉自己的苦難，也藉門徒們身上發生的事，證實了關於祂死亡的言語，並藉奧跡，從各方面使那些受馬爾西翁瘟疫感染的人困惑。為此，祂甚至讓首席宗徒否認祂。但如果祂沒有被捆綁也沒有被釘死，伯多祿和其他宗徒的恐懼從何而來呢？\par

然而，祂另一方面不讓他們等到憂傷的時候，而是說什麼？\Quote{但我復活以後，要在你們以先往加里肋亞去。} 因為祂不是立刻從天顯現，也不是要往任何遙遠的地方去，而是在同一民族中，就是在祂曾被釘死的同一地方附近，以便再次向他們保證被釘的就是復活的那一位，並這樣在憂傷時更豐盛地安慰他們。因此祂也說\Quote{在加里肋亞}，為使他們擺脫對猶太人的恐懼，從而相信祂的話。為此祂果然在那裡顯現。\par

\Quote{伯多祿回答說：\InnerQuote{即使眾人都為你的緣故而跌倒，我卻永不跌倒。}}\par

3. 你怎說，伯多祿？先知說：\Quote{羊群要四散}；基督已證實了這話，而你卻說\Quote{不}嗎？先前的事還不夠嗎？當你說：\Quote{這事絕不能臨到你}\BibleRef{瑪 16:22}，你的嘴就被堵住了嗎？因此祂讓他跌倒，藉此教導他凡事相信基督，並視祂的宣言比自己的良心更可信。而其餘的人也從他的否認中獲益匪淺，認識到人的軟弱和天主的真理。因為當祂預言任何事情時，我們不可再挑剔，也不可抬高自己超過常人。因為經上說：\Quote{你的誇耀要在你自身內，而不在別人身上。} 因為他本應祈禱說：幫助我們，使我們不被剪除，他卻自信地說：\Quote{即使眾人都為你的緣故而跌倒，我卻永不跌倒}；即使眾人都經歷這事，我也不會經歷，這使他逐漸變得自恃。基督於是為了壓制這一點，容許了他跌倒。因為他既不服從基督，也不服從先知（其實祂為此還引用了先知，使他們無法反駁），但既然他不服從祂的話，他就受事的教訓。\par

為證明祂容許這事是為要在他身上改正這一點，請聽祂說什麼：\Quote{我已為你祈求，使你的信德不致喪失。} 祂說這話是嚴厲責備他，表明他的跌倒比其他人更嚴重，需要更多幫助。因為可責備的有兩點：一是他反駁；二是他把自己置於他人之上；或者還有第三點，即他把一切都歸於自己。\par

為治癒這些事，祂便容許了跌倒，也為此緣故離開其他人，而懇切地向他說話。因為祂說：\Quote{西滿，西滿，看，撒旦已要求得著你們，好篩你們像篩麥子一樣；} 即是為了擾亂、困惑、試探你們；但 \Quote{我已為你祈求，使你的信德不致喪失。}\par

而如果撒旦想要所有的人，為何祂不說關於眾人的\Quote{我已為你祈求}呢？這豈不是很明顯，正如我先前提到的那樣，祂向他說話是為責備他，並表明他的跌倒比其他人更嚴重嗎？\par

祂为何不说：\InnerQuote{但我不容许}，反而说：\InnerQuote{我已祈求了？}此时祂正在通往受难的路上，说话开始谦卑，为要显明祂的人性。那位在伯多禄的宣认上建立了祂的教会，并如此坚固它，以致万般危险和死亡都不能胜过它的；那位赐给他天国的钥匙，并赋予他如此大的权柄，且完全不需要为此祈祷的（因为祂并没有说\InnerQuote{我已祈求}，而是凭自己的权柄说：\InnerQuote{我要建立我的教会，我要把天国的钥匙给你}），祂怎会需要祈祷来坚固一个单一灵魂的动摇呢？那么，祂为何如此说话？出于我所说的原因，也因他们的软弱，因为他们尚未拥有对祂合宜的看法。\par

那么，祂怎能否认呢？祂不是说：\InnerQuote{使你不否认}，而是说：\InnerQuote{使你的信德不丧失}，即你不会完全失落。因为这是来自祂的关怀。\par

因为恐惧确实驱走了一切，因为它超乎寻常，并且之所以变得超乎寻常，是由于天主极其剥夺了祂的帮助，而祂之所以极其剥夺，是因为在祂内过度存在着固执和反驳的激情。为要将其连根拔除，祂便容许恐怖临到祂身上。\par

为证明这种激情在他身上是痛苦的，他不满足于先前的话语，反驳先知和基督，而且在这些事情之后，当基督对他说：\InnerQuote{我实在告诉你，今夜鸡叫以前，你要三次否认我}，他回答说：\InnerQuote{即便我必须和祢同死，我也决不会否认祢。}路加还表示，基督越警告他，伯多禄就越极力反对。\par

这些事是什么意思，啊，伯多禄？当祂说：\InnerQuote{你们中有一人将要出卖我}时，你害怕自己可能是那个叛徒，便催促那位门徒去问，虽然你自己并不觉得有这种事；但现在，当祂明明呼喊说：\InnerQuote{众人都要跌倒}时，你却在反驳，而且不只一次，而是两次和多次？因为这就是路加所说的。\par

那么，这从何而来呢？来自极大的爱，来自极大的喜悦。我的意思是，自从他摆脱了关于出卖的那种令人痛苦的恐惧，知道了叛徒是谁，他便自信地说话，高抬自己超过其余的人，说：\InnerQuote{即使众人跌倒，我总不会跌倒。}在某种程度上，他的行为也源自嫉妒，因为晚餐时他们争论\InnerQuote{他们中哪一个最大}\BibleRef{路 22:24}，这激情竟如此困扰他们。因此，祂制止了他，并不是强迫他否认，万万不可！而是让他失去祂的帮助，揭露人性的软弱。\par

无论如何，请看此后他是如何被降伏的。复活后，当他说：\InnerQuote{主，这个人怎么样？}\BibleRef{若 21:21} 并受到斥责时，他不敢再像这里那样反驳，而是沉默不语。同样，在升天之时，当他听到：\InnerQuote{不是你们可以知道时间或时期}\BibleRef{宗 1:7}，他又沉默不语，没有反驳。此后，在房顶上，借着那块布，当他听到有声音对他说：\InnerQuote{天主所洁净的，你不可称为俗物}\BibleRef{宗 10:15}，即使当时他不知道这话是什么意思，他仍安静，没有争辩。\par

4. 这一切都是那次的跌倒造成的。在此之前，他把一切归功于自己，说：\InnerQuote{即使众人跌倒，我总不会跌倒}；以及\InnerQuote{即使我必须死，我也决不会否认祢}（他本应说：如果我得到祢的帮助），但在此之后，完全相反：\InnerQuote{你们为什么注视我们，好像我们凭自己的能力和虔诚使这人行走呢？}\par

我们由此学到一个伟大的教导：人的意愿是不够的，除非人得到从上而来的援助；同样，如果缺乏意愿，从上而来的援助也一无所得。这两点，犹达斯和伯多禄都表明了：前者虽得了许多帮助，却毫无益处，因为他没有意愿，也不贡献自己的部分；而后者，虽然心意已定，但因没有受到援助，便跌倒了。因为美德之网的确是由这两者编织而成的。\par

因此，我恳求你们：既不要将一切交托给天主而自己沉睡，也不要努力劳作，以为单凭自己的辛劳便可成就一切。因为天主不希望我们懈怠，所以并未独自完成一切；也不愿我们自夸，故未将一切交给我们；而是除掉了两者中有害的，留下对我们有益的。因此，祂甚至容许首席宗徒跌倒，既使他内心更加谦卑，也从此训练他拥有更大的爱。因为经上说：\InnerQuote{那得赦免多的，爱得也多。}\BibleRef{路 7:47}\par

那么，让我们在一切事上相信天主，不要在任何事上反驳祂，即使所说的似乎与我们的思想和感官相悖，但让祂的话语拥有高于推理和视力的权威。让我们在圣事中也这样做，不注视眼前的事物，而谨记祂的话语。\par

因为祂的话语不会欺骗，而我们的感官却容易受骗。那话语从未落空，而感官在多数事情上会出错。既然这话说：\InnerQuote{这是我的身体}，让我们既相信又信服，并以心灵的眼睛注视它。\par

因为基督没有赐下任何感官之物，虽然是在感官之物中，却全部要以心灵去领会。同样，在圣洗中，恩赐是通过感官之物（即水）授予的，但所成的事却由心灵领会，我指的是重生和更新。因为如果你是无形的，祂就会把无形的恩赐赤裸地交付给你；但由于灵魂被禁锢在身体内，祂便在感官之物中交付那些心灵所领会的事物。\par

如今有多少人说：我愿看见祂的形貌、印记、衣裳、鞋袜。看哪！你看见祂，你触摸祂，你吃祂。你固然渴望看见祂的衣裳，但祂却把自己赐给你，不仅为看见，也为触摸、吃下并接受到你内。\par

因此，不要有人冷漠地、不要有人怯懦地前来，却要怀着炽热的心、满腔热忱、全神贯注。因为犹太人站着，脚上穿着鞋，手中拿着杖，匆匆地吃，何况你更应该警醒。他们原是要前往巴勒斯坦，所以有行旅的装束，而你却是要升往天堂。\par

5. 因此，必须在各方面保持警醒，因为对于不相称地领受的人，确实有重大的惩罚。\par

试想你对那叛徒、对钉死祂的人是何等义愤。所以要当心，免得你自己也成了基督身体和血的有罪者。他们杀害了那至圣的身体，而你却在领受如此大的恩惠之后，以污秽的灵魂接受它。因为祂不仅满足于成为人、被击打、被杀害，而且还把自己与我们结合，不仅借着信德，更借着事实使我们成为祂的身体。那么，领受这祭献恩惠的人，岂不应该在纯洁上远超一切？那要切割这肉的手，那充满属灵之火的嘴，那被这至可怖的血染红的舌头，岂不应该比阳光更纯洁？请思量你被赋予何等尊荣，你参与的是何等筵席。天使看见它时都战栗，不敢因那从而来的光辉而无所敬畏地注视，我们却以此得喂养，以此与我们结合，与基督成为同一身体、同一血肉。\Quote{谁能述说上主的大能，宣扬祂的一切赞美？}有哪个牧人用自己的肢体喂养他的羊群？我说牧人干什么？常有母亲在生育之苦后，把孩子交给别的妇女哺乳；但祂却不这样做，而是亲自用自己的血喂养我们，并用各种方式把自己与我们交织在一起。\par

请注意，祂是由我们的本体所生。但你说，这与众人无关；虽然它确实涉及众人。因为如果祂来到我们的本性，显然是为众人；但若是为众人，便是为每一个人。那么你说，为什么不是所有人都从中获益？这并非由于祂的作为，祂的抉择是为众人做这事，而是那些不愿意之人的过犯。祂在奥秘中与每一位信徒结合，祂所生的人，祂亲自喂养，不托付给别人；这也再次说服你，祂已取了你的血肉。我们既蒙受了如此大的爱与尊荣，就不要懈怠。你们不见婴儿是如何渴望吸吮乳房？他们是如何迫切地把嘴唇贴在乳头上的？我们也当以同样的心情走近这筵席，走近这属灵之杯的乳头。或者更确切地说，让我们像吃奶的婴儿一样，更迫切地汲取圣神的恩宠，让我们唯一的忧愁就是不能领受这食物。摆在我们面前的作为并非出于人的能力。那在晚餐中行这些事的那一位，如今同样行这些事。我们占据仆人的位置。那祝圣并变化它们的，是同一位。所以不要让犹大在场，不要让贪财的人在场。如果谁不是门徒，就让他退去，这筵席不收留这样的人。因为祂说：\Quote{我同我的门徒守逾越节。}\BibleRef{玛 26:18}\par

这筵席与那晚餐是同一的，毫无逊色。因为并非基督行了那事，而人行了这事，而是祂也做这事。这就是那间楼房，他们当时所在之处；他们从这里出发往橄榄山去。\par

我们也当出去到穷人手中，因为这个地方就是橄榄山。因为穷人的众多就是栽在天主家里的橄榄树，滴下那对我们有益的油，就是五个贞女所有的，而其余的因没有油而灭亡。领受了这油，让我们进去，灯光明亮地迎接新郎；领受了这油，让我们从此出去。\par

不要让任何不人道的人在场，不要让任何残忍无情的人，不要让任何不洁的人在场。\par

6. 这些话我是对领受者说的，也是对服务者说的。因为我也必须对你们讲话，好使你们极其谨慎地分发这些恩赐。如果你们明知某人有任何邪恶，却仍容许他领受这筵席，你们就没有轻微的惩罚。\Quote{我必向你追讨他的血债。}\BibleRef{则 33:8}即使某人是将军、总督，甚至是戴冠冕者，若不相称地前来，也要阻止他；你们所拥有的权柄大于他的。如果你们受托为羊群看守一眼清澈的泉源，然后看见一只羊嘴上满是污泥，你们不会容许它低头弄浑水流；但现在你们受托看守的泉源不是水，而是血和神，如果你们看见任何人身上有罪——这比泥土和污泥更严重——走向它，难道你们不恼怒吗？不驱赶他们吗？你们还有什么可推诿的呢？\par

为此，天主给了你们这尊荣，叫你们辨别这些事。这就是你们的职责，你们的保障，你们全部的花冠，并非叫你们穿着白色闪亮的衣袍走来走去。\par

你也许会说：我怎能认识这个人或那个人呢？我不是说那些不认识的人，而是那些众所周知的。\par

我要说一件更可怕的事：附魔者在场，还不及像保禄所断言的那样践踏基督、\Quote{轻视那盟约的血，并侮慢圣神的恩宠}的人那么严重。因为陷入罪恶而走近的人，比附魔者更坏。附魔者因被附而不受惩罚，但那些人若不相称地走近，就被交付给不死的惩罚。所以我们不仅要把这些人赶走，而且要把所有我们看见不相称前来的人，无一例外地赶走。\par

\Quote{不属于门徒的人，不可领圣体。不要让\Person{犹达斯}领受，免得他遭受\Person{犹达斯}的命运。这群众也是基督的身体。因此，你主持圣事的人要小心，免得你惹动主，因为没有洁净这身体。别给人刀剑代替食物。}\par

\Quote{不，即便他是出于无知前来领受，也要禁止他，不要害怕。敬畏天主，不要怕人。你若怕人，连那人也要耻笑你；但你若敬畏天主，连人也会尊重你。}\par

\Quote{但若你不敢自己这样做，就把他带到我这里来；我不容许任何人胆敢做这些事。我宁愿舍弃性命，也不愿将主的血分给不配的人；我宁愿流自己的血，也不愿将如此可畏的血按不合宜的方式分给人。}\par

\Quote{但若经过多方查问，仍不认识那坏人，便无可责备。因为这些话是针对公开的罪人说的。因为如果我们改正了这些人，天主也会很快向我们揭示那些隐藏的；但若我们放任这些人不管，他为何要显明那些隐藏的呢？}\par

\Quote{但我说这些，不是要我们仅排斥他们，或与他们断绝，而是要我们改正他们，领他们回头，照顾他们。因为这样，我们既得天主垂顾，也会有许多人堪当领受；且因我们自己的勤勉和对他人的关怀，获得巨大赏报。愿天主因我主\Person{耶稣基督}的恩宠和爱人之心，使我们众人得以达到；愿光荣归于他，至于无穷世。阿们。}\par

\SourceRecord{Translated by George Prevost and revised by M.B. Riddle.；Nicene and Post-Nicene Fathers, First Series；Vol. 10.；Edited by Philip Schaff.；Buffalo, NY: Christian Literature Publishing Co.,；1888.；<http://www.newadvent.org/fathers/200182.htm>.}

% structure: p0001 source=h1 target=section reason=page-title

\section{玛窦福音讲道词第八十三篇}

\BibleRef{玛 26:36-38}\par

\BibleQuote{然后耶稣同他们来到一个名叫革责玛尼的地方，便对门徒说：\Quote{你们坐在这里，等我到那边去祈祷。}祂带了伯多禄和载伯德的两个儿子，开始忧愁痛苦，就对他们说：\Quote{我的心灵忧闷得要死，你们留在这里，同我一起醒寤。}}\par

因为他们紧紧跟随祂，所以祂说：\Quote{你们留在这里，等我离开去祈祷。}这是祂的习惯，分开祈祷。祂这样做，教导我们在祈祷时要为自己预备安静和极大的退隐。\par

祂带了那三个人，对他们说：\BibleQuote{我的心灵忧闷得要死。}为什么祂没有带所有人？为了不使他们沮丧；但祂带了那些曾见过祂荣耀的人。然而，连这些祂也遣散了：\BibleQuote{祂稍往前走，俯伏在地，祈祷说：\Quote{父啊！如果可能，就让这杯离开我！但不要照我，而要照你所愿意的。}祂来到门徒那里，见他们睡着了，便对伯多禄说：\Quote{你们竟不能同我醒寤一个时辰吗？醒寤祈祷吧！免陷于诱惑；心神固然切愿，但肉体却软弱。}}\par

祂并非无故最严厉地指责伯多禄，虽则其他人也睡着了；但也是为了使他通过这一点感受到我之前所说的原因。然后，因为其他人也说了同样的话（当伯多禄说（这是他的话）：\Quote{即便我必须同祢死，我决不会否认祢}，同时，\Quote{众门徒也都这样说了}），\BibleRef{玛 26:36}祂就对他们所有人说话，揭露他们的软弱。因为他们愿意同祂死，却连醒着同祂悲伤都不能，反而被睡意压倒。\par

祂恳切祈祷，为免此事看似表演。又有汗滴流下，同样是为了不让异端者说祂在演戏。因此有汗珠如血滴，并有天使显现加强祂，以及千百种恐惧的迹象，免得有人以为这些话是假装的。因此也有这个祈祷。当祂说：\Quote{如果可能，就让这杯离开我}，祂显示了祂的人性；但说：\Quote{但不要照我，而要照你所愿意的}，祂显示了祂的德行和自制，教导我们即使本性拉扯我们，也要跟随天主。因为为了愚昧的人，单单显示祂的面容是不够的，祂也用了言语。再者，言语本身也不够，还需要行动；祂也将这些与言语结合，即使那些极爱争论的人也可能相信，祂既成了人，也死了。因为即使这些事已然如此，有人仍然不信，那么若没有这些，更会不信。看祂通过多少事物来显示降生成人的真实性：通过祂所说的，通过祂所受的。然后祂来对伯多禄说，如经上所记：\Quote{怎么，你不能同我醒寤一个时辰吗？}\BibleRef{谷 14:37}所有人都睡着了，祂却责备伯多禄，在祂的话中暗示他。\Quote{同我}一词并非无故使用；似乎是说：你不能同我醒寤一个时辰，你还要为我舍命吗？随后的话也暗示了同一件事。因为祂说：\Quote{醒寤}，并\Quote{祈祷，免陷于诱惑}。看，祂又一次教导他们不要自负，而要有忏悔的心，谦卑，并将一切交托给天主。\par

祂有时对伯多禄说话，有时对众人说话。祂对伯多禄说：\Quote{西满，西满，撒殚想要得到你们，好筛你们像筛麦子一样，但我已为你祈求了}；祂对众人说：\Quote{祈祷，免陷于诱惑}；从各方面拔除他们的自以为是，使他们有认真之心。然后，为避免祂的话显得完全是指责，祂说：\Quote{心神固然切愿，但肉体却软弱}。因为即使你渴望轻看死亡，你也不能，除非天主伸出祂的手，因为肉性的心思把人往下拉。\par

祂又同样祈祷，说：\Quote{父啊！如果这杯不能离开我，非要我喝不可，就成就你的意愿吧！}，表明祂完全顺从天主的旨意，并且我们必须常常跟随这旨意，寻求它。\par

\Quote{祂来到，发现他们睡着了。}因为除了夜深，他们的眼睛也因沮丧而沉重。第三次祂来说了同样的话，证实了祂成为人的事实。因为在圣经中，第二第三次特别表示真理；如同若瑟也对法郎说：\Quote{天主使梦重复了两次，是因为这事已确实，且天主急促要实行。}\BibleRef{创 41:32}因此，祂也一次、两次、三次说了同样的事，为证实降生成人。\par

祂第二次来是为了什么？为了责备他们，因为他们如此沉浸在沮丧中，甚至对祂的临在毫无感觉。但祂并没有责备他们，而是离开他们稍远一点，显示他们不可言喻的软弱，即他们受到责备后还是无法忍受。但祂没有再次唤醒并责备他们，以免打击已经受打击的人，而是离开去祈祷，当祂又回来时，祂说：\Quote{现在你们睡吧，休息吧！}然而那时却需要醒寤，但为了显示他们连看见危险也承受不住，会因恐惧而逃跑并抛弃祂，并且祂不需要他们的援助，祂必须被交出，祂说：\Quote{现在你们睡吧，}祂说，\Quote{休息吧！看，时辰近了，人子被交在罪人手中。}\BibleRef{玛 26:45}\par

祂再次显示，所发生的事属于天主的安排。\par

2. 祂不仅仅这样做，而且还藉着说\Quote{交在罪人手中}，鼓舞他们的心神，表明这是由于他们的恶行，而非因祂有任何罪责。\par

\BibleQuote{起来，我们去吧；看，那出卖我的人近了。} \BibleRef{玛 26:46} 因为祂用各种方式教导他们，这事不是出于必然或软弱，而是出于某种奥秘的安排。因为，如我们所见的，祂预知犹达斯会来，非但没有逃走，反而去迎接他。无论如何，\BibleQuote{祂还说话的时候，看，犹达斯，十二宗徒之一，来了，并同他一起有许多带着刀剑棍棒的人，从司祭长和民间的长老那里来。} \BibleRef{玛 26:47} 司祭们的器具真是相称！他们带着\Quote{刀剑棍棒}来对付祂！并且犹达斯，据说，与他们一起，十二宗徒之一。他又称他为\Quote{十二宗徒之一}，并不羞愧。那出卖祂的人给了他们一个记号，说：\BibleQuote{我亲吻谁，谁就是祂，你们就抓住祂。} 哦！这叛徒的灵魂接受了何等堕落啊！因为他当时用什么样的眼睛看他的师父？用什么样的嘴亲吻祂？哦！可咒骂的企图；他想出了什么？他胆敢做什么？他给出什么样的背叛记号？他说：我亲吻谁。他被祂的温和所鼓动，这温和足以使他羞惭，并剥夺他所有的借口，因为他是在背叛如此温良的一位。\par

但祂为什么说这话？因为当祂被他们抓住时，常常从他们中间出去，他们竟不知道。然而，那时这事也会发生，若不是祂自己的旨意要祂被捉拿。至少为教导他们这一点，祂那时使他们眼瞎，并亲自问：\BibleQuote{你们找谁？} \BibleRef{若 18:4} 他们认不出祂，虽然带着灯笼火把，并有犹达斯同他们一起。后来，当他们说了\Quote{耶稣}时，祂说：\BibleQuote{我就是你们所找的}；而这里再次说：\Quote{朋友，你为何来？}\par

因为祂显示了自己的能力之后，便立刻把自己交出。但若望说，直到最后一刻祂仍继续责备他，说：\BibleQuote{犹达斯，你用亲吻来出卖人子吗？} \BibleRef{路 22:48} 祂说：你连背叛的形式都不感到羞耻吗？然而，因为连这也没能阻止他，祂任由他亲吻，并自愿被交出；他们就下手拿住祂，在吃过逾越节的那一夜捉住了祂，他们如此怒气沸腾，简直疯了。然而，他们不会有力量，除非祂自己允许。但这并没有使犹达斯逃脱不可忍受的惩罚，反而更严厉地定他的罪，因为虽然他接受了如此多的关于祂的能力、宽厚、温良和良善的证明，他却变得比任何野兽更凶猛。\par

既然如此，让我们逃避贪财。因为正是贪财那时使他疯狂；它使那些被它俘虏的人变得极端残忍和不人道。因为，当它使他们对自己的救恩绝望时，更使他们忽视其他人的救恩。这欲望是如此暴虐，有时竟能胜过最强烈的色欲。因此，我极其羞愧，因为为了节省钱财，他们可能确实克制了不贞，但为了对基督的敬畏，他们却不愿贞洁庄重地生活。\par

因此我说，让我们逃避它；因为我将永不停止地说这话。因为，人啊，你为什么积聚金子？为什么使你的束缚更苦？为什么使你的看守更痛苦？为什么使你的焦虑更痛苦？把那些埋藏在矿里的金属，那些在君王宫廷里的金属，都算作你自己的吧。因为如果你拥有那全部一堆，你只会保存它，而不会使用它。因为如果你现在没有使用你所拥有的东西，反而像它们属于别人一样禁戒不用，那么如果你拥有更多，情况更是如此。因为贪财者的做法是，他们越是在周围堆积，就越是吝惜。但你说：\Quote{但我知道这些是我的。} 那么占有只是假设，而非享用。但你说：\Quote{我应成为人所恐惧的对象。} 不，反而你会因此成为更容易的猎物，无论对富人或穷人，对强盗、诬告者、仆人，以及一般所有存心算计你的人。因为如果你想成为可惧的对象，就断除他们能抓住你并伤害你的机会，无论谁存心于此。难道你没有听到那个比喻吗？它说：\Quote{穷人和赤身露体的人，即使一百人聚集起来也从未能剥去他的衣服。} 因为他有他的贫穷作为最大的保护，这保护甚至连国王也永远不能征服和夺取。\par

3. 事实上，所有的人都来折磨贪财的人。我说人，何必呢？飞蛾和虫子也攻击这样的人。我为什么说飞蛾？光是时间本身，即使无人打扰他，也足以对这样的人造成最大的损害。\par

那么财富的快乐是什么？因为我看到它的不适，但请你告诉我它的快乐。它的不适是什么？你说：焦虑、阴谋、敌意、仇恨、恐惧；永远口渴和痛苦。\par

因为如果有人拥抱他所爱的少女，却不能满足他的欲望，他就遭受极大的折磨。富人也是如此。因为他拥有很多，并与她同在，却不能完全满足他的欲望；但同样的结果发生，正如某位智者所说：\Quote{阉人奸污童贞女的欲望}，以及\Quote{像阉人拥抱童贞女而叹息}；所有富人都是如此。\par

何必再提其他事？这样的人如何使所有人不悦：使他的仆人、工人、邻居，处理公务的人，受伤害的人，未受伤害的人，特别是他的妻子，以及超过一切他的儿女。因为他不是把人当作人来抚养，而是比仆人和买来的奴隶更悲惨。\par

他又为自己招致无数机会，惹来愤怒、烦恼、侮辱和嘲笑，成为众人共同的笑柄。这些苦恼便是如此，或许还不止于此；人无法用言辞一一列举，但经验能将它们呈现在我们面前。\par

但请告诉我由此而来的快乐。他说：\Quote{我显得富有，}且\Quote{被人视为富有。}被视为富有有何快乐？那是一个招人嫉妒的盛名。我说盛名，因为财富只是一个空洞的名号，并无实质。\par

他说：\Quote{然而那富有的人，}\Quote{沉溺于这种想法并以此自娱。}他因那些本应悲伤的事而自娱。\Quote{悲伤？为何？}他问。因为财富使他百无一用，面对流放和死亡都变得怯懦和缺乏男子气概，因为他加倍珍惜财富，贪爱钱财甚于光明。这样的人连天堂也无法使他喜悦，因为天堂不产黄金；太阳也不能使他喜悦，因为它不放射金色的光芒。\par

但有人说，有些人确实享受着他们所拥有的，过着奢侈、暴饮暴食、酗酒、挥霍无度的生活。你是在告诉我比前者更糟的人。因为后者尤其是那些毫无享受的人。至少前者因被一种爱束缚，而避免了其他邪恶；但后者比前者更坏，除了我们所说的，他们还为自己招来一群残酷的主人，每天都为肚腹、情欲、酗酒及其他各种放纵效劳，如同服侍许多残暴的暴君，豢养娼妓，摆设昂贵的宴席，收买食客和谄媚者，追求不自然的欲望，使他们的身体和灵魂因此陷入千百种疾病。\par

因为他们不是把财物花在必要之处，而是用来毁坏身体，同时也毁坏灵魂；他们所做的，就像一个装饰自身的人，以为是在为自己所需花钱一样。\par

因此，只有那为正当目的使用财富的人，才能真正享受快乐并主宰自己的财物；但这些人却是奴隶和俘虏，因为他们加剧了身体的激情和灵魂的疾病。这算什么样的享受？这里只有围困、战争和比海洋汹涌更可怕的风暴。因为如果财富遇到愚人，会使他们更愚昧；遇到放纵的人，会使他们更放纵。\par

你或许会说：智慧对穷人有什么用处？果然，你不知道；因为盲人也不懂得光明的好处。请听\Person{撒罗满}所说：\Quote{智慧胜过愚昧，如同光明胜过黑暗。}\BibleRef{训 2:13}\par

但我们如何教导那在黑暗中的人呢？因为贪财是黑暗，不容任何事物显示本相，而是呈现别样。譬如在黑暗中，即使看见金器、宝石、紫袍，也会以为它们算不得什么，因为他看不到它们的美；同样，贪财的人也不懂得那些值得关注之事的真正美善。所以，我恳请你驱散由这种情欲产生的迷雾，然后你就能看清事物的本性。\par

但这一切在贫穷中最为明显，那些看似是而非的事物，在克己中最能被揭穿。\par

4. 但愚昧的人啊！他们甚至咒骂穷人，说贫穷使房屋和生活都蒙羞，真是混淆一切。请问：什么是房屋的耻辱？既没有象牙床，也没有银器，只有陶器和木器。不，这恰恰是房屋最大的光荣和特点。因为对世俗事物淡然处之，常常使人将全部闲暇用于照顾自己的灵魂。\par

因此，当你看到对外在事物过分用心时，就应因这极大的不体面而感到羞耻。因为富人的房屋最缺乏体面。当你看到桌上铺着桌布，床上镶着银饰，如同剧场，如同舞台上的展示，还有什么比这更不体面呢？哪一种房屋最像舞台和舞台上的事物？富人的还是穷人的？显然不是富人的吗？因此，富人的房屋充满不体面。哪一种房屋最像\Person{保禄}或\Person{亚巴郎}的？显然是穷人的。因此，穷人的房屋最受装饰，最值得称赞。为了让你明白这尤其是房屋的装饰，请进入\Person{匝凯}的房屋，看看当\Person{基督}即将进去时，\Person{匝凯}是如何装饰它的。他没有跑到邻居那里去乞求帷幔、座椅和象牙椅子，也没有从壁橱里拿出拉科尼亚的挂毯；而是用一种适合\Person{基督}的装饰来装饰它。这是什么？他说：\Quote{我把一半的财物施舍给穷人；凡我欺骗过谁，就偿还四倍。}我们也应这样装饰我们的房屋，好让\Person{基督}也进入我们中间。这些才是美丽的帷幔，这些是在天上织成的，它们是在那里编织的。哪里有这些，哪里就有天堂之王。但如果你以另一种方式装饰，你就是邀请魔鬼和他的伙伴。\par

他也进入了税吏\Person{玛窦}的家。那么这个人又做了什么？他首先用自己的乐意、舍弃一切并跟随\Person{基督}来装饰自己。\par

同样，\Person{科尔乃略}用祈祷和施舍装饰了他的房屋；因此，直到今日，它比王宫更加光辉。因为房屋的低贱状态不在于器皿摆放杂乱，不在于床铺凌乱，也不在于墙壁被烟熏黑，而在于居住者的邪恶。\Person{基督}也显示了这一点：如果居住者德行高尚，祂并不羞于进入这样的房屋；但若进入那另一种房屋，即使它有金色的屋顶，祂也绝不进入。因此，这房屋比王宫更为辉煌，因为它接待了万有的主；而那有金色屋顶和柱子的房屋，则如同污浊的水沟和阴沟，因为它盛放着魔鬼的器皿。\par

但这些话我们所说的，不是针对那些为有用目的而富有的人，而是针对贪婪和贪得无厌的人。因为在他们中间，既没有对于必需之事的勤奋或关心，只有关于纵情于肚腹、醉酒和其他类似的不体面之事；但在另一些人那里，是关于自律。因此，基督从未进入华丽的房屋，而是进了税吏和首席税吏以及渔夫的家，离开了君王的宫殿和穿细软衣服的人。\par

因此，如果你也渴望邀请祂，要用施舍、祈祷、恳求、守夜来装饰你的房屋。这些是基督君王的装饰，而\OriginalTerm{玛门}{mammon}的装饰则是基督敌人的装饰。所以，没有人应为简陋的房屋感到羞耻，如果它有这些家具；让富人不以拥有昂贵的房屋而骄傲，反而应当掩面，寻求这另一种装饰，舍弃那一种，好使他在这里也能接待基督，并在那里享受永恒的居所，因我们的主耶稣基督的恩宠和对人的爱，愿光荣和权能归于祂，至于无穷之世。阿们。\par

\SourceRecord{Translated by George Prevost and revised by M.B. Riddle.；Nicene and Post-Nicene Fathers, First Series；Vol. 10.；Edited by Philip Schaff.；Buffalo, NY: Christian Literature Publishing Co.,；1888.；<http://www.newadvent.org/fathers/200183.htm>.}

% structure: p0001 source=h1 target=section reason=page-title

\section{玛窦福音讲道集第八十四篇}

\BibleRef{玛 26:51-54}\par

\Quote{看，同耶稣在一起的一个人，伸手拔出剑来，砍了大司祭的仆人一剑，削去了他的一个耳朵。}\par

\Quote{那时耶稣对他说：把你的剑放回原处，因为凡持剑的，必死于剑下。你想我不能祈求我父，他必立刻给我派遣十二个军团以上的天使吗？若这样，经上记载必须如此的事，怎能应验呢？}\par

这\Quote{一个人}是谁？若望说，他就是伯多禄。\BibleRef{若 18:10}因为这行为出于他的热心。\par

但另有一点值得探究：为什么他们佩带剑？因为他们佩带剑不仅从这里可以清楚看出，而且从他们被问时说\Quote{这里有两把}也可以得知。可是基督为什么容许他们持有剑呢？因为路加也证实了这事：他曾对他们说：\Quote{我以前派遣你们不带钱囊、不带口袋、不带鞋，你们缺少过什么吗？}他们回答说：\Quote{什么也没有。}他就说：\Quote{但现在，有钱囊的应当带上，有口袋的也一样；没有剑的，应当卖掉自己的外衣去买一把。}他们回答说：\Quote{这里有两把剑。}他就说：\Quote{够了。}\par

那么他为什么容许他们有剑呢？为使他们确知他要被出卖。因此他对他们说：\Quote{让他买一把剑。}这不是叫他们武装自己，绝不是；而是藉此表明他已被出卖。\par

为什么他也提到钱囊？他由此教训他们从此要清醒、警醒，并要自己多加勤勉。因为在起初，他因他们（没有经验）以诸多大能的表现来培育他们；但后来，如同把幼鸟带出巢穴，他命令他们运用自己的翅膀。然后，为了不让他们以为他让他们独自面对是因软弱，他在命令他们也尽本分时，提醒他们先前的事，说：\Quote{我以前派遣你们不带钱囊，你们缺少过什么吗？}好叫他们从这两件事上认识他的大能：既知道他曾保护他们，也知道他现在渐渐地让他们自谋其力。\par

但那里的剑是从哪里来的？他们是从晚餐和筵席上来的。因为有羔羊，很可能有剑；而且门徒们听说有人要来对付他，就取了剑作为防御，意思是为主人作战，但这只是他们的想法。因此伯多禄使用剑时也受到责备，且带着严厉的威胁。因为他正在抵抗来的仆人——确实激烈，但不是自卫，而是为主人做的。\par

然而基督没有容许任何伤害发生。因为他治好了那人，并显了一个大奇迹，足以同时表明他的容忍和能力，以及他门徒的情感和温良。那时他出于情感行事，现在则出于顺命。因为当他听见\Quote{把你的剑放回剑鞘}（\BibleRef{若 18:11}）后，他立即服从，此后也未再这样做。\par

但另一个人说，他们还问：\Quote{我们可以砍吗？}（\BibleRef{路 22:49}），但他禁止了，并治好了那人，又责备他的门徒，并加以威胁，为促使他服从。\Quote{因为凡持剑的，}他说，\Quote{必死于剑下。}\par

他又加上一个理由，说：\Quote{你想我不能祈求我父，他必立刻给我派遣十二个军团以上的天使吗？但为使经书得以应验。}（\BibleRef{玛 26:53}）他以此话平息了他们的愤怒，表明这事也蒙经书认可。因此他在那里祈祷，为使他们得知这事也是按天主的旨意成就时，能温顺地接受临到他身上的事。\par

他藉这两件事安慰了他们：一是那些图谋反对他的人所受的惩罚——\Quote{因为凡持剑的，}他说，\Quote{必死于剑下}；二是他并非不自愿地承受这些事——\Quote{因为我能祈求，}他说，\Quote{我的父。}\par

为什么他不说：\Quote{你想我不能消灭他们全部吗？}因为他所说的话更容易使人相信；那时他们对他的认识还不正确。而且不久之前，他曾说：\Quote{我的心灵忧闷得要死}，以及\Quote{父啊！愿这杯离开我吧}（\BibleRef{玛 26:38}）；他显出了痛苦和出汗，并由一位天使坚强了他。\par

因为他既然显示了许多人性的标记，若他说\Quote{你想我不能消灭他们}，似乎不太可能使人相信。因此他说：\Quote{你想我不能祈求我的父吗？}他又谦卑地说：\Quote{他必立刻给我派遣十二个军团的天使。}因为如果一个天使杀死了十八万五千名武装士兵（\BibleRef{列下 19:35}），那何需十二个军团来对付一千人呢？但他这样措辞是为了他们的恐惧和软弱，因为他们已吓得要死。因此他也引用经书反对他们，说：\Quote{那么经书怎能应验呢？}也以此使他们害怕。因为如果这事为经书所认可，你们岂能反对并与之争战吗？\par

2. 他对门徒说了这些话；但对其他人说：\Quote{你们带着刀剑棍棒出来拿我，如同对付强盗一样吗？我天天坐在圣殿里施教，你们并没有捉拿我。}\par

请看祂做了多少事来唤醒他们。祂将他们击倒在地，治好了仆人的耳朵，并警告他们将遭杀身之祸。\Quote{凡动刀的，}祂说，\Quote{必死于刀下。}藉着治耳，祂也证实了这些事；从各方面，既从当前的事，也从将来的事，彰显祂的大能，并表明这逮捕不是凭他们的力量。因此祂又补充说：\Quote{我天天与你们同在，坐着教导，你们没有拿住我。}这也表明这逮捕是出于祂的许可。祂略过奇迹，只提教导，以免显得自夸。\par

当我教导时，你们没有拿住我；当我沉默时，你们竟来攻击我？我在圣殿里，没有人拿住我，而现在你们在深夜带着刀棒来捉拿我？对于这位常与你们同在的，何需这些武器？藉着这些事教训他们，除非祂自愿投降，否则他们绝不可能成功。因为连那些曾将祂握在手中、却未能拿住祂的人，以及那些将祂围在中间却未能胜过祂的人，除非祂愿意，也绝不可能成功。\par

此后，祂也解决了为何祂那时愿意的难题。因为祂说：\Quote{这事成就，}祂说，\Quote{是为应验先知的话。}\BibleRef{玛 26:56}请看，直到最后一刻，甚至在背叛行为中，祂仍做一切事以使他们悔改，医治、预言、警告。祂说：\Quote{因为，}祂说，\Quote{他们必死于刀下。}为了表明祂是自愿受苦，祂说：\Quote{我天天与你们同在教导；}为了显明祂与父的合一，祂补充说：\Quote{为应验先知的话。}\par

但为何他们在圣殿里没有拿住祂？因为在圣殿里，他们因百姓的缘故不敢下手。因此祂也外出到外面，既藉着地点也藉着时间给予他们安全感，甚至到最后时刻还除去他们的借口。因为那为应验先知而连自己都舍弃的，怎么可能教导与他们相悖的事呢？\par

\Quote{那时，所有门徒，}据记载，\Quote{都离开祂逃走了。}因为当祂被捉拿时，他们还留下；但当祂对群众说了这些话后，他们就逃走了。因为从那时起，他们看出逃脱已不可能，当祂自愿将自己交给他们，并说这事是按圣经成就的。\par

当这些人逃走时，\Quote{他们把祂带到大司祭盖法那里；但伯多禄远远跟着，进入里面要看结局如何。}\par

这位门徒的热情多么大；他看见别人逃跑，自己却不逃跑，反而站定脚跟，与祂一同进去。就算若望也如此做了，但他却是\BibleQuote{为大司祭所认识的。}\BibleRef{若 18:15}\par

为何他们把祂带到他那里，那里众人聚集？为的是能经司祭长们同意做一切事。因为他当时是大司祭，众人都等候基督在那里，他们甚至彻夜不睡，为此放弃睡眠。因为他们当时并未吃逾越节晚餐，而是为这另一目的守望。因为若望说过\Quote{天还早}，又补充说\Quote{他们进了审判厅，免得被玷污，但为了吃逾越节晚餐。}\par

那么我们必须说什么呢？他们因急于行凶而在另一天吃逾越节晚餐，因此违犯了法律。因为基督不会逾越逾越节的时间，而他们却胆大包天，践踏千条法律。因为他们极其愤怒，多次试图捉拿祂却未能得逞；既然这次意外捉住了祂，他们甚至选择放弃逾越节，以满足他们的杀人欲望。\par

因此他们都聚集在一起，这是一群恶人的议会。他们问一些问题，想使这个阴谋具有法庭的外表。因为\Quote{他们的证词互不一致}，这个法庭是如此虚假，一切都充满混乱无序。\par

\Quote{但假见证人上前来说：这人说过：\InnerQuote{我能拆毁这座圣殿，并在三天内把它重建起来。}}\par

那么大司祭做什么？他想要逼祂申辩，以便借此捉拿祂，就说：\Quote{你没有听见这些人对你作证的事吗？但祂保持沉默。}\par

因为辩护的努力毫无益处，无人听他的。因为这只不过是法庭的虚饰，实际上是强盗的袭击，无缘无故地攻击祂，如同在洞穴或路上。\par

因此\Quote{祂保持沉默}，但那人继续说：\Quote{我指着永生的天主叫你起誓，告诉我们你是不是基督，永生天主之子。耶稣对他说：\InnerQuote{你说的是。但我告诉你们：从此你们要看见人子坐在大能者的右边，乘着云彩从天降来。}于是大司祭撕裂自己的衣服说：\InnerQuote{他说了亵渎的话。}}他做这事是为了加重指控，以行为夸大祂所说的话。因为那话使听者恐惧，他们曾对斯德望所做的，\BibleRef{宗 7:59}捂住耳朵，这位大司祭在这里也做同样的事。\par

3. 然而这究竟是怎样的亵渎呢？因为此前祂曾说过，当他们聚集时：\BibleQuote{主对我主说：你坐在我右边。}（\BibleRef{玛 22:43}）并解释了这话，他们却不敢说什么，而是沉默不语，从此不再质问祂。那么他们现在为什么称这话是亵渎呢？基督又为什么这样回答他们呢？是为了除去他们所有的借口，因为直到最后一天，祂教导说祂是基督，祂坐在圣父的右边，祂将再次降临审判世界，这是表明祂与圣父完全合一的话语。\par

因此，大司祭撕裂自己的衣服，说：\Quote{你们以为如何？}（\BibleRef{玛 26:66}）他不亲自宣判，而是邀请他们宣判，如同在承认的罪行和明显的亵渎案中一样。因为他们知道，如果这事被查究并仔细裁决，就会使祂免于一切责难，所以他们自己在他们中间定罪祂，抢先对听众说：\Quote{你们听见了这亵渎的话。}几乎是强迫并迫使人们宣判。那么他们说了什么？\Quote{祂是该死的。}这样，他们拿了祂当作已定罪的人，就能以此促使比拉多宣判。那些同谋者也说：\Quote{祂是该死的。}他们自己控告，自己审判，自己宣判，自己在那一刻充当一切。\par

但他们为什么不提出安息日的事呢？因为祂常使他们哑口无言；而且他们想要捉拿祂，并凭当时所说的话定祂的罪。\Emph{大司祭}抢先一步，以他们的名义宣判，并撕裂自己的衣服，把他们全都拉拢过来，然后把祂当作已被定罪的人带到比拉多面前，就这样做了。\par

至少在比拉多面前，他们没说这类话，而是说：\Quote{如果这人不是作恶的，我们便不会把祂交给你。}他们企图用政治控告置祂于死地。他们为什么不秘密杀害祂呢？他们也想要败坏祂的名声。因为现在有许多人听过祂，并赞赏祂、惊叹祂，所以他们竭力公开处死祂，在众人面前执行。\par

但基督没有阻止这事，反而善用他们的恶行来确立真理，使祂的死成为昭彰。结果却与他们的愿望相反。他们想要公开示众，以为这样能羞辱祂，但祂却因这些事而更加闪耀。正如他们说过：\Quote{让我们杀祂，免得罗马人来夺走我们的地土和民族。}结果在祂死后这事发生了；同样，在这里，他们的目的是公开钉死祂以损害祂的名声，结果却适得其反。\par

为了证明他们本来有能力自己在他们中间处死祂，请听比拉多说什么：\Quote{你们自己把祂带走，按你们的法律审判祂吧。}（\BibleRef{若 18:31}）但他们不愿如此，好让祂似乎是被当作犯法者、篡位者、煽动叛乱者处死的。因此他们还同祂一起钉了两个强盗；因此他们说：\Quote{不要写\InnerQuote{这是犹太人的王}，而要写\InnerQuote{祂说自己是犹太人的王}。}\par

但所有这些事都是为了真理而成就，使他们连任何无耻辩解的影子也没有。同样，在坟墓那里，封条和看守使真理更加显明；而嘲笑、讥讽、辱骂也产生了同样的效果。\par

因为错误的本质就是这样：它被那些它图谋毁灭的事物所毁灭；至少在这里也是这样，那些似乎得胜的人，他们最受羞辱、被击败、被毁灭；但那位似乎被击败的人，祂反而大放光芒，并大大得胜。\par

因此，我们不要处处寻求胜利，也不要处处回避失败。有时胜利带来伤害，失败却有益处。例如，在愤怒之人的情形中，那个极其暴怒的人似乎占了上风；但他恰恰是被最严重的激情所制服和伤害的人；而那勇敢忍受的人，反而占了上风和得胜。一个人没有能力胜过自己的疾病，另一个人却除去了别人的疾病；这个人被自己的激情压制，那个人却胜过了别人的激情；他不仅没有被烧毁，反而扑灭了别人高涨的火焰。如果他想要获得表面上的胜利，他自己也会被击败；并且激怒对方，使对方更严重地受苦；像女人一样，两者都会在愤怒中可耻而悲惨地倒下。但现在，那个自制的人既脱离了这个耻辱，又在自己和邻人身上为战胜愤怒树立了荣耀的纪念碑，借着他光荣的失败。\par

4. 因此，我们不要处处寻求胜利。因为欺骗者战胜了受屈者，但那是一种邪恶的胜利，给得胜者带来毁灭；而受屈者，如果他能克制自己忍受，反而正是那得到冠冕的人。因为常常被打败更好，这本身就是最佳的胜利方式。无论是欺骗、击打或嫉妒，那被打败而不进入冲突的人，正是得胜者。\par

为何我提及欺压与嫉妒？因为就连被拖去殉道的人，也这样藉着被捆绑、被鞭打、被残害、被杀害而获胜。在战争中是失败——即战士倒下——在我们这里却是胜利。因为我们绝非因行不义而获胜，而是处处因受不义而获胜。如此，胜利也变得更加光荣，当我们这些受苦者胜过那些行恶者时。由此表明，胜利是出于天主。事实上，它对外在的征服具有相反的性质，这又是一切之上绝无错谬的强健标记。同样，海中的岩石因受冲击而击碎波浪；同样，所有圣人都被传扬、被加冕、并竖立起他们荣耀的凯旋，赢得这平静的胜利。\Quote{你不可激动自己，}祂说，\Quote{也不可自寻烦扰。}天主赐给了你这力量，不是藉着冲突，而单单藉着忍耐来征服。不要像他那样也反对自己，你就得胜了；不要冲突，你就获得了冠冕。你为何羞辱自己？不可让他说，你是藉着冲突占了上风，而要让他惊讶并惊叹于你那不可战胜的力量，并向众人说，甚至未进入冲突你已得胜。\par

同样，蒙福的\Person{若瑟}获得了美名，处处因受不义而胜过那些行不义者。因为他的兄弟和那埃及女人都是谋害他的人，但这个人胜过了一切。因为不要告诉我这个人所住的监狱，或那个女人所居的王宫，而要向我指明谁被征服、谁失败、谁沮丧、谁欢喜。因为她不仅不能胜过正义者，甚至连自己的情欲也未能克制；而这个人既胜过了她，也胜过了那严重的病症。你若愿意，且听她亲口所说的话，你就必看见那凯旋。\Quote{你带到这里来的那个希伯来仆人，竟来戏弄我们。}\BibleRef{创 39:17}啊，可怜而不幸的女人，戏弄你的不是这个人，而是那告诉你你能击碎金刚石的魔鬼。带他到你这里的，不是\Emph{你的丈夫}，为谋害你，而是那邪恶的精灵带来了那不洁的色欲；他才是戏弄你的那一位。\par

那么\Person{若瑟}做了什么？他沉默不语，就这样被定罪，正如基督也被定罪一样。\par

因为所有那些事都是这些事的预像。他确实被捆绑，而她身处王宫。然而这算什么呢？因为他在捆绑中却比任何戴冠的胜利者更为荣耀，而她虽在王宫中却比任何囚犯更为悲惨。\par

但不只从此处可看出胜利与失败，也由其结局本身看出。哪一位成就了他的愿望？是囚犯，而非那出身高贵的女子？因为他努力保持贞洁，她却要破坏它。那么，谁成就了他所愿望的呢？是那受苦者，还是那行恶者？很明显，是那受苦者。确实，这就是那得胜者。\par

既然知道这些事，让我们追求这由受不义而得的胜利，逃避那由行不义而得的胜利。因为这样，我们既可今生全然安宁、极其平静地度日，又将获得未来的福乐，藉着我们的主耶稣基督的恩宠和慈爱，愿光荣与权能归于祂，世世无尽。阿们。\par

\SourceRecord{Translated by George Prevost and revised by M.B. Riddle.；Nicene and Post-Nicene Fathers, First Series；Vol. 10.；Edited by Philip Schaff.；Buffalo, NY: Christian Literature Publishing Co.,；1888.；<http://www.newadvent.org/fathers/200184.htm>.}

% structure: p0001 source=h1 target=section reason=page-title

\section{玛窦福音讲道辞 第八十五篇}

\BibleRef{玛 26:67-68}.\par

\BibleQuote{\Emph{那时，他们便向祂脸上吐唾沫，并用拳头打祂，另有人用手掌打祂}，\Emph{说：\InnerQuote{祢是基督，给我们说，是谁打了祢？}}}\par

他们既然要置祂于死地，为何还要做这些事？何需这样的嘲弄？为的是叫你们从一切事上看出他们放纵的性情；既把祂当作猎物捕获，就如此显出他们的醉态，任意发泄他们的疯狂；以此作为节日，快意地攻击祂，并展示他们凶杀的本性。\par

但我恳请你们赞美门徒们的自制，他们何等精确地记载了这些事。这清楚显明了他们热爱真理的心志，因为他们以全部真实记载了那些似乎可耻的事，毫不掩饰，也不以此为羞，反而视为极大的荣耀，事实也确是如此：宇宙的主宰竟为我们忍受了这样的事。这显示祂不可言喻的温柔，以及那些人的无可推诿的邪恶，他们竟忍心对如此温良、柔和、并用足以使狮子变为羔羊的话来感化他们的主，做出这样的事。因为祂在一切温柔上毫无欠缺，而他们在所言所行上，在无礼和残忍上也无欠缺。这一切，先知依撒意亚曾预先宣告，预先这样传报，用一句话就暗示了这一切无礼。因为他说：\Quote{许多人怎样因你而惊奇，你的容貌同样在人们面前被轻视，你的荣耀在人们之子面前也是如此。}\par

有什么能与这种无礼相比呢？面对那张海一见便敬畏、太阳在十字架上见到便收回光线的脸，他们竟吐唾沫，用手掌击打，有人还打祂的头，在各方面尽情放纵自己的疯狂。因为他们加给了祂最侮辱的打击：拳打、掌掴，又在这些打击上加上向祂吐唾沫的侮辱。又说出了充满讥讽的话：\Quote{祢是基督，给我们说，是谁打了祢？}因为群众称祂为先知。\par

但另一个人\BibleRef{路 22:64}说，他们用祂自己的衣服蒙住祂的脸，并做了这些事，好像他们中间抓了一个卑鄙无用的人。而且那时不仅自由人，连奴隶也向祂放纵了这种放纵。\par

我们要不断读这些事，要正确听这些事，要将这些事铭刻在我们心中，因为这就是我们的光荣。我引以为荣的，不仅是祂使数以千计的人复活，更是祂所忍受的苦难。保禄处处提出这些：十字架、死亡、苦难、辱骂、侮辱、讥讽。如今他说：\BibleQuote{我们到祂那里去，带着祂的凌辱。}\BibleRef{希 13:13}又说：\BibleQuote{祂因那摆在面前的喜乐，忍受了十字架，轻看了羞辱。}\BibleRef{希 12:2}\par

\BibleQuote{那时，伯多禄坐在外面的院子里；一个使女来到他跟前说：\InnerQuote{你也是和加里肋亚人耶稣一起的。}但他当着众人否认说：\InnerQuote{我不知道你说的什么。}他出去到了门廊，另一个使女看见他，就说：\InnerQuote{这个人也是和纳匝肋人耶稣一起的。}他又起誓否认。过了一会儿，旁边站着的人来对伯多禄说：\InnerQuote{你确是他们中的一个，因为你的口音把你暴露了。}于是他就开始咒骂起誓说：\InnerQuote{我不认识那个人。}立时鸡就叫了。伯多禄便想起耶稣所说的话：\InnerQuote{鸡叫以前，你要三次否认我。}他就出去痛哭。}\par

奇异而可叹的事！当他看见主被捕时，他是那样热切，以致拔出剑来，砍掉了那个人的耳朵；但当理当更加愤怒、激动、燃烧时，听到这般辱骂，他却成了否认者。因为那时所行的事，谁不会激怒到疯狂呢？然而门徒被恐惧胜过，不仅没有愤怒，反而否认，连一个可怜卑微的使女的威胁都承受不住，而且不是一次，而是第二次、第三次否认祂；时间短暂，且不是在审判官面前，因为是在外面：\BibleQuote{他出去到了门廊}，他们问他，他甚至没有立刻意识到自己的跌倒。路加记载说：\BibleRef{路 22:61}，即基督看了他一眼，表明他不仅否认了祂，而且内心也没有回忆起，尽管鸡已经叫了；但他还需要来自主的进一步提醒，主的眼光对他来说如同声音；他是那样充满了恐惧。\par

但马尔谷说，他第一次否认时，鸡就叫了；但第三次否认后，鸡叫了第二次；因为他更详细地陈述了门徒的软弱，以及他因恐惧而完全麻木；这些事他是从他的师傅本人\BibleRef{伯前 5:13}那里学来的，因为他是伯多禄的追随者。在此，最令人惊叹的是，他不但没有掩饰老师的过错，反而比其他人更清楚地陈述了这一点，正因他是他的门徒。\par

那么，玛窦所说的这话如何是真的？他说基督说：\BibleQuote{我实在告诉你们：鸡叫以前，你要三次否认我。}\BibleRef{玛 26:34}而马尔谷在第三次否认后说：\BibleQuote{鸡就叫了第二次。}\BibleRef{谷 14:72}不，二者都是真的，并且和谐一致。因为鸡每次啼叫，也常叫第三次和第四次；马尔谷为了表明连声音也没有阻止他、使他醒悟，就说了这话。所以两者都是真的。因为在鸡完成一次啼叫之前，他已经第三次否认了。并且即使被基督提醒了他的罪，他也不敢公开哭泣，免得眼泪暴露他，而是\BibleQuote{出去痛哭}。\par

\Quote{及至天亮，他们从盖法那里把耶稣解往比拉多那里。}因为他们想要处死祂，但因节期的缘故自己不能执行，便将祂解到总督那里。\par

但请你留意，这件事是如何被迫发生的，以至于发生在节期之中。因为从起初就是这样预表的。\par

\Quote{那时，出卖耶稣的犹达斯看见祂已被定罪，就后悔了，把那三十块银钱拿回来。}\par

这对犹达斯和那些人都是一个控告：对犹达斯，并非因为他后悔，而是因为他后悔得太迟、太缓慢，并且成了自我定罪（因为他自己承认了交出了祂）；对那些人，则因为他们本有能力挽回，却毫不后悔。\par

但请你留意，他是在什么时候感到懊悔的。是在他的罪完成并成就之后。因为魔鬼正是这样：它不让那些不警醒的人在罪成就之前看见罪恶，免得它所掳掠的人后悔。至少，当耶稣说了那么多话时，他并未受到影响；但当他的罪行完成后，懊悔才临到他；然而这懊悔并无益处。因为他谴责自己的行为、丢下银钱、不理会犹太民众，这些都是可接纳的；但上吊自尽，这又是不可饶恕的，是恶灵的工作。因为魔鬼过早地把他从懊悔中领出，使他不能从此获得任何果实；并借着一种最可耻的、众人皆知的死亡将他带走，诱使他自毁。\par

但请你留意，真理如何在各方面照耀出来，甚至借着仇敌的所作所为和所受的苦。因为事实上，连叛徒的结局本身也堵住了那些定祂罪之人的口，不容他们有任何丝毫的借口，那借口无疑是厚颜无耻的。因为他们还能说什么呢？当叛徒自己被显明对自己作出了这样的判决。\par

但让我们也看看所说的话：\Quote{他把那三十块银钱拿回来交给司祭长，说：\InnerQuote{我出卖了无辜的血，犯了罪了。}他们却说：\InnerQuote{这与我们何干？你自己看着办吧！}于是他把银钱扔在圣殿里，就离开出去，上吊自尽了。}\par

因为他再也无法忍受良心的鞭笞。但请你留意，犹太人也遭受了同样的事。这些人本应因所受的苦而改过，却不停手，直到他们完成了自己的罪。因为他的罪已经完成了，那是背叛；但他们的罪还没有。当他们完成了自己的罪，把祂钉在十字架上时，他们也乱了方寸：一会儿说：\Quote{不要写\InnerQuote{犹太人的君王}}（\BibleRef{若 19:21}）（然而你们为什么害怕？为什么被钉在十字架上的死尸所困扰？）；一会儿又看守祂，说：\Quote{免得祂的门徒来偷去，说祂复活了；那后来的迷惑就比先前的更坏了。}即使他们这样做了，这事如果为假，也会被驳斥。但他们怎敢这样说呢？他们连在祂被捉时都不敢站稳，他们的首领甚至三次否认祂，连一个使女的威吓都受不了。但正如我所说，司祭长现在乱了方寸；因为他们知道这行为是违法的，从他们的话\Quote{你自己看着办吧}就明显可见。\par

你们贪财的人要听，要想想那人的遭遇：他同时失去了钱财、犯了罪、也毁灭了自己的灵魂。\OriginalTerm{贪婪}{covetousness}的暴政就是如此。他既未享用钱财，也未享用今生或来世，而是瞬间失去了一切，甚至在那些人面前得了坏名声，于是上吊自尽。\par

但正如我所说，在行为之后，有些人就看得清楚了。至少请看这些人当时也不愿清楚地看清事实，却说：\Quote{你自己看着办吧}；这件事本身就是对他们最严厉的控告。因为这是他们为自身的胆大妄为和违法作证的话语，但他们被情欲冲昏了头，不愿停止撒旦般的企图，反而愚蠢地以假装无知的遮羞布包裹自己。\par

因为如果这些话是在钉十字架和祂被杀之后说的，那么即使那时这话也毫无合理之处，但它对他们的谴责不会那么严重；然而现在，你们既将祂握在自己手中，有权释放祂，你们怎能说出这些话呢？因为这种辩护将成为对你们最严厉的控告。如何？以何种方式？因为他们将全部责任推给叛徒（他们说\Quote{你自己看着办吧}），本可以使自己脱离这杀害基督的罪行，却撇下叛徒，甚至进一步推动罪行，将十字架加于背叛之上。因为当他们向他说\Quote{你自己看着办吧}时，什么阻止他们自己停止犯罪行为呢？但如今他们甚至做了相反的事，加上谋杀，并且在他们所做所说的一切中，将自己缠在不可避免的灾祸中。事实上，此后当比拉多将选择权交给他们时，他们选择释放强盗而不是耶稣；却杀害了那无罪且曾赐给他们如此多恩惠的祂。\par

3. 那么那人做了什么？当他看见自己徒劳无功，他们不肯收下银钱时，\Quote{他把银钱扔在圣殿里，就出去上吊自尽了。司祭长拾起那些银钱，说：\InnerQuote{不可放入献仪箱内，因为这是血价。}他们商议之后，就用那钱买了陶匠的田，作为埋葬外乡人之用。因此那块田直到今日还叫做\InnerQuote{血田}。这就应验了耶肋米亚先知所说的话：\InnerQuote{他们取了那三十块银钱，即被估价者之价，并用来买了陶匠的田，正如主所吩咐我的。}}\par

你岂不见他们又因自己的良心而自责吗？因为他们知道是自己买下了杀人的罪行，就不把银钱投入圣库，却买了田地埋葬外乡人。这也成了反对他们的见证，证明他们的叛逆。因为那地方的名字比号角更清晰地宣告了他们的血债。他们并非随意为之，而是商议之后，每一件事都同样如此，使无人能脱离干系，所有人都负罪。但这些事是预言自古以来就预告的。你岂不见不但宗徒们，连先知们也精确地宣布那些受责备的事，从各方面宣扬那苦难，且预先指明吗？\par

犹太人这样行，自己却不自觉。因为他们若将银钱投入圣库，事情就不会被如此清楚发现；但如今他们买了一块地，竟向后世之人显明了这一切。\par

你们这些想从凶杀中行善，并因人的生命收取报酬的人，要听！这些施舍是犹太式的，更确切说是撒殚式的。因为现在也有，正有许多人强取他人无数财物，却以为只要投下十枚或一百枚金币，就为一切作了辩解。\par

论到这些人，先知也说：\BibleQuote{你以泪水遮盖了我的祭台。} \BibleRef{拉 2:13}基督不愿被贪婪滋养，祂不接受这种食物。你为什么侮辱你的主，将不洁之物献给祂？宁可让人忍饥挨饿，也比从这些来源喂养他们更好。前者是残忍之人的行为，后者是既残忍又傲慢之人的行为。因为，告诉我，如果你看见两个人，一个赤身，一个有衣服，然后你剥了那有衣服之人的衣服，去遮盖那赤身的人，你岂不是行了不义？这显然是众所皆知的。但如果你将所夺的一切给了另一个人，你行了不义，而非显示仁慈；当你连所抢夺的极小部分也不给出，却称那行为为施舍，你将会受何等刑罚？因为人若献上瘸腿的牲畜尚且被责难，你若行更严重的事，又怎能得恩惠？因为若那首领先向物主本人归还，仍然是不义，且加四倍偿还不易涤除对自己的控告，这是在旧盟约之下；\BibleRef{出 22:1}那不偷窃，却强取，不但不还给被抢之人，反倒给了别人；且不还四倍，连一半也不给；何况不是活于旧规定之下，而是活于新规定之下；试想他是在自己头上堆积了何等大火！若他尚未受罚，我为此事更要哀哭他，因为除非悔改，他正为自己积蓄更重的义怒。因为祂说：\BibleQuote{你们以为，}祂说，\BibleQuote{那被塔倒压死的人就是罪更大的吗？我告诉你们：不是的，除非你们悔改，你们也要同样灭亡。}\par

那么，让我们悔改，并施与纯全不染贪婪的、极丰富的施舍。试想犹太人曾供养八千名肋未人，连同肋未人，还有寡妇和孤儿，他们还负担许多其他公共开支，同时还服兵役；但现在，在教会里，因着你们和你们的心硬，却有了田地、房屋、出租的寓所、车辆、骡夫、骡子，以及这一类的庞大排场。因为教会的这份储藏本应在你们之中，你们的慷慨本应成为她的收入；但现在发生了两件错事：你们继续不结果实，天主的司铎也不能履行他们应有的职责。\par

难道在宗徒时代，房屋和土地不能保留吗？那么他们为什么变卖并施散呢？因为这是更好的事。\par

4. 但现在，我们的先祖感到恐惧（是因为你们如此疯狂追求世俗之物，因为你们聚敛而不分散），恐怕那些寡妇、孤儿和童贞女会因饥饿而死亡；所以他们被迫提供了这些。因为他们本不愿投身于如此不体面的事；而是希望你们的善意成为他们的供给，他们从那里收取果实，而他们自己只应专务祈祷。\par

但如今你们迫使他们效仿管理公共事务者的家室；因而一切颠倒混乱。因为当你们和我们都被同样的事纠缠时，谁能为天主求情呢？因此，当教会的状态不比世俗之人更好时，我们无法开口。你们岂没听过，宗徒们连毫无困难地分派收集的款项也不肯同意吗？但如今我们的主教们在关心这些事上已经超过了代办、管家和小贩；当他们本应关心你们的灵魂时，却每天为这些事烦恼——那些客栈老板、税吏、会计和管家所关心的事。\par

我说这些事并非无故抱怨，而是为使有些悔改和转变，为使人们怜悯我们受着痛苦的奴役，为使你们成为教会的收入与储藏。\par

但若你们不愿意，看，穷人就在你们眼前；凡我们所能供养的，我们必不断喂养；但那些我们无法供养的，我们将留给你们，使你们在那可怕的日子听不到那对不仁慈和残忍之人所说的话：\BibleQuote{你看见我饿了，却没有给我吃。} \BibleRef{玛 25:42}\par

因为这种不人道，连同你们一起，使我们成了笑柄，因为我们放下祈祷、教导及其他圣洁的部分，却终日争斗，有的与酒商，有的与粮商，有的与零售其他商品的人。\par

因此产生了战斗、纷争、每日的辱骂、责备和嘲弄，并且每个司铎都被加上了更适用于世俗人家庭的称号；其实本应取别的名称来代替这些，应当以那些宗徒所任命的事情来命名：喂养饥饿者、保护受伤害者、照顾陌生人、救助受虐待者、供应孤儿、参与照顾寡妇、管理贞女；这些职务应当在我们中间分配，而不是管理土地和房屋。\par

这些才是教会的储备，这些才是相称于她的宝藏，它们在很大程度上既带给我们安逸，也带给你们利益；或者不如说带给你们安逸连同利益。因为我猜想，因天主的恩宠，在这里聚集的人数达到十万；如果每个人将一块饼施舍给某个穷人，所有人都将富足；但如果只给一个铜钱，就没有人会贫穷；我们也就不会因为对钱财的操心而遭受那么多的辱骂和嘲弄。的确，那句话：\BibleQuote{变卖你的财产，施舍给穷人，然后来跟随我}，若是针对教会财产向教会的主教们说，正是合时的。因为无论如何，若不能摆脱一切更粗俗和更世俗的操心，我们就不可能适当地跟随祂。\par

但现在，天主的司铎参与葡萄收成和收割，以及产品的买卖；而侍奉那影子的人虽然被托付了更属血肉的职务，却完全免于这些事务；我们这些被邀请到天上最内层的圣所，进入真正至圣所的人，却承担了商人和零售商的操心。\par

因此，大大忽略了圣经，祈祷懈怠，对其他一切职责漠不关心；因为人不可能以应有的热忱同时兼顾两者。因此我祈祷并恳求你们，愿许多泉源从各方涌向我们，愿你们的慷慨成为我们的打谷场和榨酒池。\par

因为这样，穷人将更容易得到供养，天主将受光荣，你们也将达致更伟大的爱人之德，并享受永恒的美善；愿天主恩赐我们都能获得这一切，藉着我们的主耶稣基督的恩宠和祂对人的爱，愿光荣归于祂，直到无穷之世。阿们。\par

\SourceRecord{Translated by George Prevost and revised by M.B. Riddle.；Nicene and Post-Nicene Fathers, First Series；Vol. 10.；Edited by Philip Schaff.；Buffalo, NY: Christian Literature Publishing Co.,；1888.；<http://www.newadvent.org/fathers/200185.htm>.}

% structure: p0001 source=h1 target=section reason=page-title

\section{玛窦福音讲道集 第八十六篇}

\BibleRef{玛 27:11-12}.\par

\Emph{\Quote{耶稣站在总督面前，总督问衪说：\InnerQuote{祢是犹太人的君王吗？}耶稣对他说：\InnerQuote{祢说了。}当衪被司祭长和长老控告时，衪什么也不回答。}}\par

你们看见祂首先被问的是什么？这正是他们不断以各种方式提出来的事吗？因为他们看见比拉多不理会法律方面的事，就把控告转向政治罪名。他们对宗徒们也是这样，常提出这些事，说他们到处宣扬一个名叫耶稣的君王，\BibleRef{宗 17:7}好像说的只是一个凡人，并使他们有篡位的嫌疑。\par

由此可见，撕裂衣服和惊愕都是假的。但他们制造一切，并竭力进行，以便把祂置于死地。\par

比拉多就是这样问的。那么基督说什么呢？\Quote{祢说了。}祂承认祂是一个君王，但属于天国的君王，这在别处祂也说得更清楚，回答比拉多说：\Quote{我的王国不属于这世界；}\BibleRef{若 18:36}免得他们和这人有机会控告祂这样的事。祂还给出了一个无可辩驳的理由，说：\Quote{如果我的王国属于这世界，我的臣仆必要战斗，使我不至于被交付。}为此，为了驳斥这种猜疑，祂既纳了税，\BibleRef{玛 22:17}又命别人纳税，而当人们要立祂为王时，祂就避开了。\BibleRef{若 6:15}\par

那么，为什么祂在彼时被控篡位时不提出这些事呢？因为从祂的行为中，祂的能力、温良、仁厚有无数的证据，但他们是故意眼瞎，行事不公，法庭也是腐败的。因此祂不再回答，而是保持沉默，但简短地回答（以免因持续沉默而获得傲慢的名声）当大司祭用圣经命令祂发誓时，当总督询问时，但对他们的控告祂不再说什么；因为祂现在不可能说服他们。正如先知古时预言此事所说：\Quote{在祂的谦抑中，祂的审判被夺去了。}\par

总督对这些事感到惊讶，的确值得赞叹地看到祂表现出如此大的忍耐，保持沉默，而祂本有无数的话可说。因为他们控告祂，并不是知道祂有什么恶事，而是出于嫉妒和嫉妒。至少当他们设了假见证后，为什么他们无话可说，仍然坚持己见呢？当他们看到犹达斯已死，比拉多洗手不干时，为什么他们不感到痛悔呢？因为在那一刻祂做了许多事，好使他们醒悟过来，但没有一个人改变。\par

那么比拉多说什么呢？\Quote{你没有听见他们作证控告你这么多的事吗？}他愿意祂为自己辩护并被释放，所以他也说了这些话；但由于祂什么也不回答，他又想出了另一个办法。\par

这是什么办法呢？当时有一个习惯，就是释放一个被定罪的人，他企图借此释放祂。因为如果你们不愿意释放无辜的祂，至少当作有罪的，因节日的缘故饶恕祂。\par

你们看到秩序颠倒了吗？因为为被定罪者求情本是民众的事，批准是官长的事；但现在相反，官长向民众求情；即使这样，他们也没有变得温和，反而变得更加凶残嗜血，被嫉妒的激情所驱动。因为他们没有理由控告祂，尽管祂沉默，但他们甚至在那时也被祂的众多义举所驳倒，而沉默的祂胜过了那说千万话并疯狂的人。\par

\Quote{当他坐在审判席上时，他的妻子派人到他那里说：\InnerQuote{你与这义人毫无关系，因为我今天在梦中因祂受了许多苦。}}你看又发生了一件足以使他们所有人悔悟的事。因为除了所发生的事的证据外，这梦也不是小事。为什么他自己没有看见呢？或者是因为她更配得，或者是因为他若看见了，就不会同样被相信；或者根本不会说出来。因此，命妻子看见，好使这事向所有人显明。而且她不仅看见了，还受了许多苦，好使这人因对妻子的感情而对这谋杀更加迟疑。时间也起了不小的作用，因为正是在那夜她看见了。\par

但有人说，释放祂是不安全的，因为他们说祂自称为王。那么他应该寻求证据和定罪，以及所有篡位的可靠标志，例如：祂是否招募军队，是否收集钱财，是否锻造武器，是否尝试任何其他类似的事。但他被随意引导，因此基督也没有为他开脱责任，说：\Quote{那把我交给你的，罪更大。}所以他是由于软弱才屈服，鞭打了祂，并交付了祂。\par

他于是无男子气概且软弱；但司祭长们邪恶且犯罪。因为他想出了一个办法，即节日的律法要求释放一个被定罪的人，他们用什么来对抗呢？据说，\Quote{他们说服了群众，要求巴拉巴。}\par

\newline{}2. 请看祂如何悉心为他们开脱罪责，他们又是何等费尽心机，以便不给自己留下一丝一毫的借口。究竟哪一个是合理的？是释放已定罪的罪犯，还是释放对其罪责尚有疑问的那一位？因为，如果对已认罪的犯人应当释放，那么对于有疑问的人更应如此。显然，这人在他们眼中并不比那些公认的杀人犯更坏。为此，经上不仅说他们有一个强盗，而且是著名的强盗，即臭名昭著的恶徒，犯下无数谋杀。然而即便如此，他们仍宁选此人而弃绝世界的救主，他们既不因时节神圣而敬畏，也不顾人道法则或任何此类事情，嫉妒已彻底蒙蔽了他们。除了他们自己的邪恶，他们还腐化民众，为的是连欺骗民众的罪也使他们在最严厉的惩罚中受苦。\par

因此，既然他们要求释放那一个，他说：\Quote{那么我该如何对待基督呢？}他这样说是想让他们羞愧，给他们选择的权力，好让他们至少因羞耻而要求释放祂，使整个事件都出于他们的慷慨。因为尽管说祂没有做错会让他们更加争辩，但出于人道要求祂得救，则带着无法反驳的说服和恳求。\par

但即便如此，他们仍说：\Quote{\InnerQuote{钉祂在十字架上。}但他说：\InnerQuote{为什么？祂做了什么恶事？}但他们越发喊叫：\InnerQuote{让祂被钉在十字架上。}他见毫无益处，就洗手说：\InnerQuote{我是无辜的。}}那么，你为什么把祂交出来？你为什么不像百夫长拯救保禄那样拯救祂？\EditorialNote{宗徒大事录第廿一章}因为那人也知道自己会取悦犹太人；而且因他发生了骚乱和喧哗，然而他仍然坚挺面对一切。但这个人并非如此，他极其懦弱无力，所有人都一起腐败了。因为这个人没有坚挺面对群众，群众也没有坚挺面对犹太人，在各方面他们的借口都被剥夺了。因为他们\Quote{越发喊叫}，即更大声地喊叫：\Quote{让祂被钉在十字架上。}因为他们不仅想要杀死祂，还想让祂被指控为恶人，尽管判官在反驳他们，他们仍继续喊叫同样的话。\par

你看见基督做了多少事好使他们回头吗？因为正如祂多次阻止犹达斯，同样祂也约束了这些人，无论是在祂的全部福音中，还是在祂受审的时刻。因为当他们看见总督和判官洗手说：\Quote{对这血我是无辜的。}他们本该因所言所行而受到痛悔，也当他们看见犹达斯上吊，以及看见比拉多亲自恳求他们另选一人代替祂。因为当控告者和叛徒定自己的罪，判官推卸罪责，而且当天夜里就出现了这样的异象，甚至祂如被定罪者恳求祂获释，他们还会有什么借口呢？因为他们既然不愿意祂是无辜的，至少也不该将祂置于一个强盗之下，一个公认的、臭名昭著的强盗。\par

那么他们做了什么？当他们看见判官洗手说：\Quote{我是无辜的。}他们就喊叫：\Quote{祂的血归在我们和我们的子孙身上。}\BibleRef{玛 27:25}最后，当他们自己定了自己的罪，他才让步让一切办成。\par

这里也看到他们极大的疯狂。因为情欲和邪恶的欲望正是这样。它们不让人看见任何正确的事。就算你们咒诅自己，为什么还要把咒诅引到你们的子孙身上呢？\par

然而，那爱人者虽然他们如此疯狂地对待自己和自己的子女，并没有确认对他们的子女的判决，甚至也没有确认对他们自己的判决，反而从两者中接纳了那些悔改的人，并视他们为配得无数美物。因为甚至连保禄也是他们中的一员，以及在耶路撒冷相信的成千上万人；因为\Quote{你看见经上说：\InnerQuote{兄弟，有多少万犹太人都信了。}}如果有些人继续\Emph{在他们的罪中}，就让他们把惩罚归咎于自己。\par

\Quote{于是他把巴拉巴释放给他们；至于耶稣，他鞭打了祂，就交给他们钉在十字架上。}\par

他为什么鞭打祂？要么是当作已定罪者，要么是出于审判的形式，要么是为了取悦他们。然而他本应抵抗他们。因为就连在此之前，他曾说：\Quote{你们把祂带走，按你们的法律审判祂。}\BibleRef{若 18:31}有许多事本可以阻止他和那些人，就是那些征兆和奇迹，以及那承受这一切者的极大忍耐，尤其是祂难以言喻的沉默。因为祂既藉着为自己辩护和祈祷显示了祂的人性，又藉着沉默和轻视所说的话显示了祂的高贵和性体的伟大；一切都在引导他们惊叹祂。但他们对这些事一概不让步。\par

\newline{}3. 因为一旦理性能力被醉酒或某种疯狂的癫狂所压倒，沉沦的灵魂就难以重新升起，除非它非常高贵。\par

因为让这些邪恶的情欲占据一席之地是可畏的，是可畏的，因此理当千方百计地防备和阻止它们进入。因为一旦它们抓住了灵魂并统治了它，就像火点燃了木材，它们就把火焰燃成烈火。\par

因此，我恳求你们做一切事来阻止它们的进入；不要用这种无情推理来安慰自己，从而把一切邪恶引到自己身上，说：这有什么？那有什么？因为无数祸患由此而生。因为魔鬼，既已堕落，就使用许多诡计、努力和自卑来毁灭人类，并开始以更微不足道之事攻击他们。\par

请注意，他想要把撒乌耳引入巫术的迷信中。但若他在一开始就建议这事，那人必不会听从；因为那甚至把巫术驱除的人，怎能听从呢？因此他温柔地、逐步地引导他走向此事。当撒乌耳违背了撒慕尔，并在撒慕尔不在场时献了全燔祭，因而受到责备时，他说：\Quote{敌人的逼迫太大}，而当他应当哀恸时，他却觉得自己没有做什么。\par

当天主又关于阿玛肋克人给他命令，他也违背了这些。从那里他继续走向关于达味的罪行，就这样轻易而逐步地滑落，不停歇，直到他来到毁灭的深坑，并将自己投了进去。同样，在加音的例子上，他没有立刻催促他杀害自己的兄弟，因为他不会说服他，而是首先使他奉献一些多少卑劣的祭品，说：\Quote{这不是罪。}其次，他燃起嫉妒和妒恨，说：这也没有什么。第三，他说服他杀人并否认他的谋杀；并且不离开他，直到他把恶的冠冕加在他身上。\par

因此，我们必须抵制开端。因为，即使初次的罪止于自身，也不应当轻视初次的罪；但如今，当心思疏忽时，它们还会继续发展成更大的罪。因此，我们应当尽一切努力消除它们的开端。\par

因为不要只看罪的本质，以为它小，而要看它被忽视时会成为大罪的根。因为若可以这样说一句奇妙的话：大罪并不需要像那些微小、无足轻重的罪那样多的热忱。因为前者，罪的本质使我们厌恶；但小罪却因这一点使我们陷入懈怠，不允许我们热心地振作起来除掉它们。因此，在我们沉睡时，它们迅速变得巨大。这一点在身体上也能看到。\par

同样，在犹达斯的例子中，那极大的邪恶就是由此而生。因为，若偷窃穷人的钱在他看来不是一件小事，他就不会走上这背叛之路。除非被虚荣所俘虏在犹太人看来不是小事，他们就不会撞上成为基督谋杀者的礁石。事实上，我们可以看到一切邪恶都由此而生。\par

因为没有人一下子突然冲向恶习。因为灵魂中确实植有羞耻之心和对正义之事的尊敬；它不会立刻变得如此无耻，以致在一次行为中抛弃一切，而是在疏忽时慢慢、逐步地消亡。同样，偶像崇拜也是这样进入的：人们过分尊崇生者和死者；同样，偶像也受敬拜；同样，淫乱和其他邪恶也盛行。\par

请看：一个人不合时宜地笑了；另一个人责备了他；第三个人消除了恐惧，说：这没有什么关系。\Quote{因为笑是什么？它能产生什么？}由此生出愚昧的戏谑；从那生出污秽的言谈；然后污秽的行为。\par

又如，另一个人因诽谤邻人、辱骂和造谣中伤而受责备，却轻视它，说：恶语算不了什么。由此他生出难以言喻的仇恨、无尽的辱骂；由辱骂生出殴打，由殴打常常生出谋杀。\par

4. 于是，从这些小事中，那邪恶的灵带来了大罪；又从大罪带来绝望；他发明了这另一个陷阱，其危害不亚于前者。因为犯罪不如绝望那样毁灭人。因为那犯了罪的人，若保持警惕，能通过悔改迅速弥补所犯的过犯；但那学会了沮丧而不悔改的人，因此未能通过悔改的补救措施来改正。\par

他还有第三个严重的陷阱；就是当他给罪披上虔诚的外表时。魔鬼曾在哪里得胜到如此地步，以致欺骗人到这种程度？请听，并提防他的诡计。基督藉着保禄命令：\Quote{妻子不可离开丈夫，\BibleRef{格前 7:10}，也不可彼此亏负，除非两相情愿；\BibleRef{格前 7:5}}但是有些人出于对节制的热爱，竟离开自己的丈夫，仿佛在做一件虔诚的事，却驱使他们犯奸淫。请想想，这是多大的邪恶：她们忍受如此多的劳苦，却应受谴责，犯下最大的不义，并遭受极重的惩罚，还把她们的丈夫推入毁灭的深坑。\par

另一些人，按照禁食的规定戒绝肉食，逐渐发展到厌恶它们；这本身就带来极大的惩罚。\par

但这事发生，当有人固执于自己的偏见，违背圣经所认可的事时。连格林多人中的一些人，认为不加区别地吃一切东西，甚至是禁忌之物，是成全的一部分；然而这不是成全，而是最大的不法。因此，保禄也严厉地责备他们，宣布他们应受极重的惩罚。另一些人又以为留长发是虔诚的标志。然而这是被禁止的事，且带有许多耻辱。\par

再者，另一些人追求为罪过度忧伤，以为这是有益的事；然而这也来自魔鬼的诡计，犹达斯就证明了这一点；至少因此他甚至上吊自缢。因此，保禄再次为那犯了奸淫的人担心，免得他遭遇这种事，就劝格林多人迅速释放他，\Quote{免得这样的人被过度的忧伤所吞没。}\BibleRef{格后 2:7}然后，表明这样的结果来自那恶者的陷阱，他说：\Quote{免得撒殚占我们的便宜，因为我们不是不知道他的诡计。}\BibleRef{格后 2:10}意思是他用许多狡猾攻击我们。因为若他清楚公开地与我们作战，胜利就会现成而容易；或者不如说，即使现在，只要我们保持警惕，胜利就会现成。因为事实上，天主已经用每一条这样的道路武装了我们。\par

為勸我們不可輕視這些小事，請聽祂怎樣警告我們說：\BibleQuote{凡向弟兄說\InnerQuote{愚妄的}，應受地獄的罰；} \BibleRef{瑪 5:22} 而且\BibleQuote{凡帶淫念注視的，已犯了姦淫。} \BibleRef{瑪 5:28} 祂也對歡笑的人宣佈禍災，處處除去邪惡的開端與種子，並說\BibleQuote{我們要為每一句空談交賬。} \BibleRef{瑪 12:36} 為此，約伯也為他子女的思念而施補救。\BibleRef{約 1:5}\par

關於不可絕望，經上說：\BibleQuote{人跌倒，豈不再起來？人轉離，豈不再回頭？} \BibleRef{耶 8:4} 又說：\BibleQuote{我不願惡人死亡，卻願他回頭而生活。} 又說：\BibleQuote{今天你們如果聽從祂的聲音。} 聖經中還有許多其他這類的話語與實例。為了不使我們在敬虔之懼的偽裝下滅亡，請聽保羅說：\BibleQuote{恐怕這樣的人被過度的憂愁所吞噬。}\par

故此，我們既知道這些事，就當在一切引人迷途的道路上，以聖經中汲取的智慧設起屏障。不要說：這有什麼關係，若我好奇地注視一個美貌女子？因為若你在心中犯了姦淫，不久你便會膽敢在肉體上犯之。不要說：這有什麼關係，若我從這窮人身邊走過？因為你若錯過這人，也會錯過下一個；若錯過他，便會錯過第三個。\par

再者，也不要說：這有什麼關係，若我貪戀鄰人的財物？因為這、這導致了阿哈布的滅亡；雖然他願付代價，卻是從不情願的人那裡強取。須知，人不可憑強力購買，而應以說服。但倘若他，願付公平價格者，因從不情願者強取而受罰，那麼，連這點也不做，而從不情願者暴力強奪，且是在恩寵之下生活的人，該當受怎樣的刑罰呢？\par

因此，為使我們不受罰，應使自己完全潔淨，遠離一切強暴與掠奪，並防範罪惡的根源及其本身，讓我們十分勤勉地專注於德行；這樣，藉著我們主耶穌基督的恩寵與對人的慈愛，我們也必得享永恆的福樂，願光榮歸於祂，直到永遠。阿們。\par

\SourceRecord{Translated by George Prevost and revised by M.B. Riddle.；Nicene and Post-Nicene Fathers, First Series；Vol. 10.；Edited by Philip Schaff.；Buffalo, NY: Christian Literature Publishing Co.,；1888.；<http://www.newadvent.org/fathers/200186.htm>.}

% structure: p0001 source=h1 target=section reason=page-title

\section{玛窦福音讲道 第八十七篇}

\Emph{\BibleRef{玛 27:27-29}}\par

\BibleQuote{\Emph{那时，总督的兵士把耶稣带进公堂，并召集了全队兵士围着祂；他们剥去祂的衣服，给祂披上一件紫红袍；又用荆棘编了一个冠冕，戴在祂头上，拿一根芦苇放在祂右手里；然后跪在祂面前，戏弄祂说：\InnerQuote{犹太人的君王，万岁！}}}\par

仿佛在某个信号下，魔鬼当时正凯旋进入一切。因为，且不说犹太人因嫉妒和猜忌而疯狂反对祂，至于那些兵士，他们又是从何而来，又是什么样的原因呢？岂不是显然魔鬼那时正疯狂地进入所有人的心吗？他们实在以侮辱祂为乐，是一群野蛮无情的人。我的意思是，他们本该感到敬畏，本该哭泣——就连百姓也这样做了——他们却没有，反而心狠手辣、傲慢无礼；或许他们自己也想讨好犹太人，又或者所做的一切都与他们邪恶的本性相符。\par

而且侮辱的方式各不相同，花样繁多。那神圣的头颅，一时被他们掌掴，一时被荆棘冠冕侮辱，一时又被芦苇击打——这些不洁而受诅咒的人！\par

在基督忍受了这些之后，我们因受伤而激动，还有什么可说的呢？因为所做的事是侮辱的极限。没有一个肢体，而是整个身体都成了被侮辱的对象：头部因冠冕、芦苇和掌掴；面部被吐唾沫；脸颊被手掌击打；全身因鞭打、被裹上袍子和虚假的敬拜；手因芦苇，他们给祂拿着代替权杖；口因被奉上酸醋。还有什么比这些更痛苦？还有什么比这更侮辱呢？\par

因为所做的事超出了言语所能表达。他们似乎害怕在罪行上有所欠缺——他们亲手杀死了先知，但这个人却由法官判决——所以他们每件事都亲力亲为；他们用自己的双手行这些事，并在他们当中和比拉多面前定罪、判刑，说：\BibleQuote{祂的血归于我们和我们的子孙}\BibleRef{玛 27:25}，然后侮辱祂、凌辱祂，亲自捆绑祂、带走祂，并成为兵士们所作恶事的发起者，将祂钉在十字架上，辱骂祂、向祂吐唾沫、嘲笑祂。因为比拉多在这件事上没有贡献什么，一切都是他们自己做的，他们成了控告者、法官、刽子手，身兼一切。\par

这些事情在我们当中宣读，当众人聚集时。为使外邦人不能说：你们向人民和民族展示的是光荣和显赫的事，比如征兆和奇迹，而隐藏这些被视为耻辱的事；圣神的恩宠促成了这事：在盛大的庆节，当众多男女和所有人，可以说，在逾越节前夕，所有这些事都被宣读；当全世界都在场时，所有这些行为都清晰响亮地宣告。这些事被宣读、被众人知晓，基督被相信是天主，并且，除了其他一切，祂还因这一点而被朝拜：祂竟屈尊为我们如此俯就，以至实际忍受了这些事，并教给我们一切美德。\par

所以让我们不断诵读这些事；因为实在有大收获，大益处。因为当你看见祂，无论通过姿态还是行为，被嘲笑、被如此戏弄地朝拜，被击打、忍受极端的侮辱，即使你是石头，你也会变得比蜡还软，并从你的灵魂中除去一切傲慢。\par

因此也听听随后的事。因为经上说：\BibleQuote{他们戏弄了祂之后，就带祂去钉十字架}，又，当他们剥去祂的衣服，就取了祂的衣裳，坐下看守祂，等祂死。他们分取祂的衣裳，这种事通常发生在最卑贱、最可鄙的罪犯身上，他们无人属己，完全孤独。\par

他们分了那衣裳，那曾行过如此伟大奇迹的衣裳。但那时奇迹一件也没有发生，基督抑制了祂不可言喻的能力。这又是侮辱的附加。因为对待一位卑贱可鄙的人，像我所说过的最卑贱的人，他们胆敢做一切事。至于那两盗贼，他们至少没有对他们做什么，但对基督，他们什么都敢做。他们把祂钉在他们中间，使祂分享他们的名声。\par

他们给祂苦胆喝，这也是为了侮辱祂，但祂不肯喝。但另一位说，祂尝了之后，说：\BibleQuote{完成了}\BibleRef{若 19:30}。\Quote{完成了}是什么意思？关于祂的预言应验了。经上说：\BibleQuote{他们拿苦胆给我作食物；我渴了，他们拿醋给我喝}。\par

然而，甚至在这里，他们的侮辱也没有停止，而是在剥去祂的衣服、钉祂十字架、献醋给祂之后，他们更进一步，看见祂被钉在十字架上，就辱骂祂，他们和路过的人都是；这比一切都更痛苦：因祂被指控为骗子、欺诈者而忍受这些，作为夸口者，徒然假装祂所说的。因此他们公开将祂钉死，为在众人眼前展示；也因此他们借兵士的手行此事，使这些事甚至由公共法庭施行，侮辱就更大了。\par

2. 然而，谁不会因跟随祂并哀悼祂的群众而感动呢？不，这些野兽却不会。因此，祂俯允回答群众，但对这些人却不然。因为在做了他们想做的事之后，他们还试图损害祂的尊荣，害怕祂的复活。于是他们公开说这些话，并与祂一同钉死强盗，希望证明祂是骗子，他们说：\Quote{你这拆毁圣殿，三日内建造起来的，从十字架上下来吧。} \BibleRef{玛 27:40} 因为当他们叫比拉多移除那控告（这写着\Quote{犹太人的君王}）时，未能得逞，但他坚持说：\Quote{我所写的，我已经写上了。} \BibleRef{若 19:22} 于是他们试图通过嘲笑祂来表明祂不是君王。\par

因此他们说了那些话，也说了这些。如果\Quote{祂是以色列的君王，现在就从十字架上下来吧。祂救了别人，却不能救自己。} \BibleRef{玛 27:42} 借此甚至要贬低祂以前的奇迹。又说：\Quote{如果祂是天主子，且祂愿意祂，就让祂救祂吧。}\par

啊，可憎！极为可憎！怎么，先知们不成为先知，义人们不成为义人，是因为天主没有将他们从危险中拯救出来吗？绝非如此，他们虽然遭受这些，仍然是先知和义人。那么，你们的愚昧还有什么能相比呢？因为如果危险临到他们身上，在你们眼中并未损害他们的尊荣，那么在这人身上，你们就更有理由不该被绊倒，因为祂所做的和所说的，总是在事先纠正你们的这种猜疑。\par

然而尽管如此，即使说了这些话、做了这些事，他们却没有成功，即使在当时也没有。无论如何，那个陷入如此邪恶、一生以杀人和偷窃为业的人，当这些话被说出时，却承认了祂，并提到一个王国，民众也为祂哀哭。然而，所发生的事在那些不认识神秘安排的人眼中，似乎证明了相反的情况，即祂软弱无力，但真理甚至借着相反的事而获胜。\par

听了这些，让我们武装自己，对抗一切狂怒，对抗一切愤怒。如果你发现自己的心在膨胀，就画十字印记你的胸膛。回想当时发生的一些事，你就能借回忆所发生的事，像尘土一样驱散一切愤怒。思考那些话语和行动；思考祂是主，你是仆人。祂为你受苦，你为你自己受苦；祂为那些曾受祂恩惠却将祂钉死的人受苦，你为你自己受苦；祂为那些曾虐待祂的人受苦，你却常常在那些受伤害的人手中受苦。祂是在全城眼前，或更好说在全体犹太人眼前（包括陌生人和本地人），在他们面前说了那些仁慈的话，而你却在少数人面前；对祂更侮辱的是，连祂的门徒也离弃了祂。因为那些先前关注祂的人已经抛弃了祂，但祂的敌人和仇敌，将祂置于十字架中间，从上下侮辱、谩骂、嘲笑、讥讽、嘲弄祂，犹太人和士兵从下面，盗贼从两边的上面：因为那两个盗贼都侮辱和谴责祂。那么路加怎么说其中一个\Quote{斥责}了呢？\BibleRef{路 23:40} 两件事都发生了，因为起初两人都谴责祂，但后来一人不再如此。为了让你不认为这事是出于某种约定，或者那盗贼不是盗贼，他借着辱骂向你显示，在十字架上他是一个盗贼和仇敌，并且立刻改变了。\par

那么，思考这一切之后，要控制自己。因为你所受的，有哪样与你的主所受的相似？你曾被公开侮辱吗？但不像这些。你曾被嘲笑吗？但你的全身体没有被这样鞭打和剥光。即使你被掌掴，也不像这样。\par

3. 此外，我请求你，思考这是由谁、为何、何时，以及祂是谁；而（最严重的是）这些事发生时，没有人责备，没有人指责所发生的事，反而所有人都赞同，并加入嘲笑和讥讽祂的行列；他们视祂为夸口者、骗子、欺诈者，不能在祂的作为中证实祂所说的话，就这样辱骂祂。但祂对一切都保持沉默，为我们准备了最有力的忍耐动机。\par

但我们，虽然每天听这些事，对仆人却没有那样忍耐，反而比野驴更冲动和踢踢，对于针对自己的伤害，变得凶残和无人性；但对针对天主的伤害却不怎么在意。对待朋友也是如此；如果有人惹我们烦恼，我们无法忍受；如果他侮辱我们，我们比野兽更凶残，我们这些每天读这些事的人。一个门徒出卖了祂，其余的人离弃祂逃跑了，那些曾受祂恩惠的人向祂吐唾沫，大司祭的仆人用手掌打祂，士兵们掌击祂；路过的人嘲笑和辱骂祂，盗贼控告祂；祂对任何人都没有说一句话，却以沉默胜过了一切；用祂的作为教导你，你越温和地忍受，就越能胜过那些恶待你的人，并在众人面前成为令人钦佩的对象。因为谁不会钦佩一个以容忍忍受那些恶待他之人的侮辱的人呢？正如虽有人公正地受苦，但若温和地忍受邪恶，多数人会认为他是不公正地受苦；同样，虽有人不公正地受苦，但若他暴力，就会引起公正受苦的嫌疑，并成为嘲笑的对象，因为被他的愤怒俘虏，失去了自己的高贵。对于这样的人，我们甚至不能称他为自由人，即使他是一万个仆人的主人。\par

但某人是否过分激怒了你？那又怎样？因为那时正应显出自制，既然无人刺激时，我们甚至看见野兽也是温顺的；因为它们并非总处于野蛮，而是当有人激怒它们时才如此。因此，如果我们只在无人触怒我们时才安静，我们比它们又有什么优点呢？因为它们也常常正当地忿怒，并且有许多借口：因受刺激和驱策而被激发，此外它们缺乏理性，且本性中有野性。\par

但是，请问，你能从哪里找到借口变得野蛮凶暴呢？你遭受了什么苦难？你被抢劫了吗？正是为了这个缘故，你应当忍受，以便获得更丰富的补偿。但你的名誉被剥夺了吗？这又算什么？如果你自制，你的状况绝不会因此变糟。但如果你没有受到任何委屈，你为什么向那既没有伤害你、反而对你有益的人生气呢？因为那些尊敬人的人，使不警惕的人更加虚荣；而那些侮辱和轻视人的人，使谨慎自守的人更加稳固。因为粗心大意的人受尊敬比受侮辱受害更大。如果我们清醒，前者（尊敬）成为我们自制的促进者，但后者（侮辱）激发我们的骄傲，使我们充满自夸、虚荣、愚妄，使我们的灵魂更软弱。\par

对此，父亲们可以作证：他们不常夸奖自己的子女，反而责备他们，恐怕他们因称赞而受到任何伤害；他们的教师也对他们采用同样的补救办法。所以，如果我们应当躲避什么人，那应当是那些奉承我们的人，而不是那些侮辱我们的人；因为对于那些不谨慎的人而言，这种诱饵比侮辱带来更大的伤害，而且控制这种情绪比控制侮辱更难。并且从那方面来的赏报也更丰厚，赞叹也更大。因为看见一个人受侮辱而不动，比看见一个人被打被击而不倒下更值得赞叹。\par

有人会说：\Quote{怎么可能不动心呢？}有人侮辱了你吗？在胸前划十字，默想那时所成就的一切，一切就都熄灭了。不要只考虑侮辱，也要考虑那侮辱你的人曾对你行过的任何善事，你立刻就会变得温和；或者最重要的是，先考虑对天主的敬畏，你很快就会变得柔和温良。\par

4. 此外，甚至从你自己的仆人身上，你可以学到关于这些事的功课；当你看见你自己在侮辱人，而你的仆人却保持沉默时，就要想到自制是可能的，并因你的暴戾而谴责自己；就在你正施行侮辱的时候，学会不侮辱；这样，即使你被侮辱，你也不会烦恼。要想到那傲慢的人已失去理智、癫狂了，你被侮辱时就不会忿怒，因为被附身的人击打我们，我们非但不生气，反而怜悯他们。你也要这样：怜悯那对你傲慢的人，因为他被一个可怕的怪物——暴怒——所控制，被一个凶恶的魔鬼——忿怒——所压制。你要解救他，因为他正被一个凶恶的魔鬼纠缠，正急速走向灭亡。因为这疾病如此严重，甚至不需要时间就能毁灭被它抓住的人。因此，有话说：\Quote{他的忿怒的权势是他的败落；}\BibleRef{德 1:22}这最清楚地显示它的暴虐：在短时间内它造成大恶，且无需与我们长久共存；所以，如果它除了力量之外还易于持久，那么与之抗争确实会是困难的。\par

我愿意展示侮辱人的人是什么样子，自制的人是什么样子，并将两者的灵魂赤裸地摆在你们面前。因为你们会看见一个像暴风雨中的大海，另一个像风平浪静的港湾。因为后者不被这些恶风所扰乱，反而轻易地使它们平息。因为那些侮辱人的人，所做的一切都是为了刺痛对方。当他们不能如愿时，他们自己从此也就平静了，并改过自新而离去。因为一个发怒的人不可能完全不谴责自己，正如另一方面，一个不发怒的人不可能自我谴责。因为即使必须报复，也可以不发怒地进行报复（这比带着怒气更容易、更明智），并且没有任何痛苦的感觉。因为如果我们愿意，美善之事将来自我们自己，并且在天主的恩宠下，我们将足以保全自己的平安和尊荣。\par

因为你为什么寻求从别人那里来的荣耀呢？你应当自己尊荣自己，就没有人能侮辱你；但如果你自己羞辱自己，即使所有人都尊荣你，你也不会受到尊荣。因为正如除非我们使自己处于邪恶状态，否则没有人能使我们如此；同样，除非我们侮辱自己，否则没有人能羞辱我们。\par

因为让一个人伟大而值得赞叹，而所有人却称他为奸夫、盗贼、掘墓者、杀人犯、强盗，但他既不恼怒也不忿恨，也不自觉有这些罪，那么他将因此蒙受什么耻辱呢？毫无。你会说：\Quote{那如果许多人对他有这种看法呢？}即便如此，他也没有受辱，而是他们使自己蒙羞，因为他们把并非如此的人当作如此。告诉我，如果有人以为太阳是黑暗的，他是给那太阳带来恶名，还是给自己？当然是给自己，使自己落得个瞎眼或癫狂的名声。同样，那些把恶人当作善人的人，以及那些犯相反错误的人，是羞辱自己。\par

因此，我们应当更加殷勤，保持良心清白，不给自己留下任何把柄，也不给他人恶意猜疑的借口；但如果别人在我们持这种态度时仍要发狂，我们就不必太在意，也不必忧伤。因为一个善人若被人看作恶人，并不会因此有损他实际为善的本性；但若有人毫无理由地妄加猜疑他人，自己反而受到最大的伤害。反之，一个恶人若被人看作相反的，也不会从中得到任何益处，反而要受更重的审判，并被引入更大的懈怠。因为一个恶人若被人察觉，或许会自卑而承认自己的罪；但若他逃脱了察觉，就会陷入麻木不仁的状态。如果罪人在众人指责时都难以被激发痛悔，那么当无人指责甚至有人称赞他们时，生活于恶习中的人又何时才能睁眼呢？你们听到\Person{保禄}也曾为此责备\Place{格林多}人，因为他们（非但不让那犯奸淫的人承认自己的罪）反而赞扬并尊敬他，从而促使他更加作恶。因此，我恳求我们不要在意众人的猜疑、侮辱或称赞，只应勤勉一事：就是自觉没有行任何恶事，也不侮辱自己。因为这样，在此世和来世，我们都要享受极大的光荣。愿我们的主\Person{耶稣基督}的恩宠和爱人之情，使我们众人得以获享这光荣。愿光荣归于祂，世世无穷。阿们。\par

\SourceRecord{Translated by George Prevost and revised by M.B. Riddle.；Nicene and Post-Nicene Fathers, First Series；Vol. 10.；Edited by Philip Schaff.；Buffalo, NY: Christian Literature Publishing Co.,；1888.；<http://www.newadvent.org/fathers/200187.htm>.}

% structure: p0001 source=h1 target=section reason=page-title

\section{论玛窦福音第八十八篇}

\BibleRef{玛 27:45-48}\par

\BibleQuote{从第六时辰起，黑暗笼罩了大地，直到第九时辰。大约第九时辰，耶稣大声呼喊说：\OriginalTerm{厄里，厄里，肋玛撒巴黑塔尼？}{Eli, Eli, lama sabachthani?}这就是说：我的天主，我的天主，你为什么弃舍了我？站在那里的人中，有几个听见了，就说：这人呼唤\Person{厄里亚}呢。其中有一个人立刻跑去，拿海绵蘸满了醋，绑在苇子上，送给他喝。}\par

这就是他先前答应赐给他们的征兆，当他们求征兆时，他说：\BibleQuote{一个邪恶淫乱的世代寻求征兆，除了先知约纳的征兆外，再没有征兆赐给它。}\BibleRef{玛 12:39}意指他的十字架、他的死亡、他的埋葬和他的复活。又用另一种方式宣告十字架的德能，他说：\BibleQuote{当你们高举了人子以后，你们便知道我是那一位。}\BibleRef{若 8:28}他所说的大意是：\Quote{当你们钉死了我，以为胜过了我时，那时你们尤其要认识我的大能。}\par

因为在十字架受刑之后，城市被毁灭，犹太国终结，他们失去了自己的政体和自由，福音兴盛，圣言传播到地极；海与陆地，居住之地和荒漠都不断宣扬其大能。他指的就是这些事，以及在十字架受刑时刻发生的那些事。因为当祂被钉在十字架上时，这些事的发生，确实比祂在地上行走时更为奇妙。奇妙不仅在此，还因为从天上也成就了他们所求的，并且这事遍及全世界，是前所未有的，只在埃及逾越节要完成时才发生过。因为那些事正是这些事的预象。\par

请注意这事发生的时间。在正午，使所有地上的居民都能知道，当时全世界都是白天；这足以使他们皈依，不仅因奇迹的伟大，也因其在适当的时候发生。因为在他们所有的侮辱和无法无天的嘲笑之后，在他们放下了愤怒，停止了嘲弄，对讥讽感到满足，说了他们想说的话之后，这才发生：他显示黑暗，为的是至少这样（他们发泄了怒气之后）能从这奇迹得益。因为他在十字架上能行这些事，比从十字架上下来更为奇妙。因为无论他们以为是他自己做的，他们本应相信并敬畏；或者不是他，而是圣父，也照样应该因此受到痛悔的感动，因为那黑暗是他对他们罪行的忿怒的标志。因为那既不是天蚀，而是忿怒和义愤，不仅由此显明，也由时间显明：它持续了三个时辰，而天蚀发生在瞬间，凡见过天蚀的都明白；实际上，在我们这一代也发生过天蚀。\par

你或许会说：为什么没有人全都惊奇而认他为天主呢？因为那时人类处于极大的漫不经心和恶习之中。这奇迹只是一个，发生后立即消失；没有人关心探究其原因，偏见和不虔敬的习惯很严重。他们不知道所发生之事的原因，他们可能认为这只是以天蚀或某种自然效应的方式发生。你对那些外面的人感到惊奇什么呢？他们一无所知，因极大的漠不关心而不加探究；而那些在\Place{犹太}境内的人，在这么多奇迹之后，仍然继续凌辱他，尽管他明明向他们显示是他自己做了这事。\par

为此，他在此后仍然说话，好让他们知道他还活着，并且是他自己行了这事，也使他们因此变得更为温和。他说：\BibleQuote{\Emph{\OriginalTerm{厄里，厄里，肋玛撒巴黑塔尼？}{Eli, Eli, lama sabachthani?}}}\BibleRef{玛 27:46}，以致直到他最后一口气，他们能看到他尊敬圣父，不是天主的仇敌。因此他也发出了先知的一声呼喊，甚至到最后一刻仍为旧约作证，而且不仅是先知的声音，还是用希伯来语，以便对他们清晰易懂，并且通过一切显示他与生他的父同心。\par

但是你们也要注意他们的放纵、无节制和愚妄。他们认为（经上记载）他所呼求的是\Person{厄里亚}，就立刻拿醋给他喝。\BibleRef{玛 27:48}另一个人来到他跟前，\BibleQuote{用枪刺透了他的肋旁}。还有什么比这些人更无法无天、更残忍呢？他们竟疯狂到如此地步，最后甚至对一具尸体加以侮辱。\par

但是，我请求你们注意，他如何利用他们的邪恶来成就我们的救恩。因为那一击之后，我们救恩的泉源就从那里涌出。\par

\BibleQuote{耶稣又大声呼喊，交付了灵魂。}这就是他所说的：\BibleQuote{我有权舍掉我的生命，也有权再取回来；}\BibleQuote{我是自己舍掉的。}\BibleRef{若 10:18}因此他大声呼喊，为要显示这行动是由权能完成的。马尔谷至少记载，\BibleQuote{比拉多惊奇他是否已经死了；}\BibleRef{谷 15:44}并且百夫长为此特别相信，因为他死得有权能。\BibleRef{谷 15:39}\par

这呼喊撕裂了帐幔，开启了坟墓，并使圣殿成为荒凉。他这样做，不是作为对圣殿的侮辱（因为他尚且说：\BibleQuote{不要使我父的殿成为商场；}\BibleRef{若 2:16}），而是声明他们甚至不配他住在那里；正如他把它交给巴比伦人时一样。但是这些事不仅为此而发生，而且所发生的是对即将来临的荒凉以及向更伟大更高境界转变的预言；也是他权能的\Emph{一个标志}。\par

连同这些事，祂也藉着随后发生的事——藉着复活死人——显示了自己。因为在\Person{厄里叟}的事例中，有人因触摸死尸而复活；但如今祂却藉着一句话使他们复活，祂的身体仍悬在十字架上。此外，那些事是这事的预象。为使这事得以相信，先前的一切才发生。他们不仅被复活，而且岩石崩裂，大地震动，为要叫他们知道，祂能使他们自己眼瞎，并将他们撕碎。因为那劈开岩石、使世界黑暗的那位，若祂愿意，更可对他们行这些事。但祂不愿意，却将祂的义愤倾泻在自然元素上，却愿以仁慈拯救他们。然而他们并未减少狂妄。嫉妒就是如此，猜忌就是如此，不易平息。那时，他们竟无耻地反对实际显现的事；后来，\Emph{甚至反对事物本身}——当祂身上加印封石，并有士兵看守而复活后，他们从守卫口中听到这些事，竟还花钱，为要败坏别人，并窃取复活的故事。\par

因此，即使此时他们仍然悖逆，不要惊奇，因为他们已完全准备好无耻地反对一切；但请留意另一点：祂行了何等伟大的征兆，有些来自天上，有些在地上，有些就在圣殿里，既表明祂的义愤，同时又显明那些不可接近的现在可以进入，天上将敞开，工作移至真正的至圣所。他们确实说：\BibleQuote{祂若是以色列的王，现在就从十字架上下来吧}\BibleRef{玛 27:42}，但祂显明祂是普世的王。而那些人说：\BibleQuote{你这拆毁圣殿，三日之内重建的}\BibleRef{玛 27:40}，祂显明这殿将永远荒凉。他们又说：\BibleQuote{祂救了别人，却不能救自己}\BibleRef{玛 27:42}，但祂留在十字架上，却在祂仆人的身体上最充分地证明了这一点。因为如果\Person{拉匝禄}在第四天复活是大事，那么所有早已长眠的人同时显现为活着的，更是何等大事，这是将来复活的征兆。因为经上说：\BibleQuote{许多长眠的圣人的身体起来了，进入圣城，显现给许多人。}为使所发生的不被当作幻像，他们甚至在城里显现给许多人。百夫长也光荣了天主，说：\BibleQuote{这人的确是义人。}\BibleRef{路 23:47}被钉者的力量如此伟大，以致在如此多的嘲弄、讥讽和嘲笑之后，百夫长和群众都痛悔了。有人说这位百夫长后来也殉道了，他在这些事后在信德中成长到成年。\par

\BibleQuote{有许多妇女在那里远远地观看，她们是从加利利跟随耶稣来服事祂的，其中有\Person{玛利亚玛达肋纳}，雅各伯和若瑟的母亲\Person{玛利亚}，并载伯德儿子的母亲。}\par

这些事妇女们看见了，她们是最倾向于同情祂、最哀悼祂的人。请注意她们的热心何等伟大。她们曾跟随祂服事祂，甚至临到危险时刻仍在那里。因此她们看到了一切：祂如何呼喊，如何交付出灵魂，岩石如何崩裂，以及其余一切。\par

这些妇女首先看见耶稣；那最受谴责的性别，却首先享受蒙福的看见，最显示其勇气。当门徒们逃散时，她们却在场。但她们是谁？祂的母亲（因她被称为雅各伯的母亲），以及其他。但另一位圣史\BibleRef{路 22:48}说，也有许多人哀悼所发生的事，并捶胸，这最显示犹太人的残酷，因为他们以别人哀悼的事为荣，既不感动于怜悯，也不畏惧害怕。因为所发生的事是极大的义愤，并且不仅仅是征兆，而全是愤怒的征兆：黑暗、裂开的岩石、中间裂开的帐幔、大地的震动，义愤的过甚极大。\par

\BibleQuote{但\Person{若瑟}去，求了身体。}这若瑟原是暗暗作门徒的，但在基督死后却变得非常勇敢。因为他既非常人，也非无名之辈，而是公会的成员，极有声望；尤其由此可看出他的勇气。因为他冒死，因对耶稣的爱而与众人为敌，既敢求取身体，在未得到之前决不罢休。他不仅取了身体，并以隆重的方式埋葬，还放在自己的新坟墓里，以此显示他的爱和勇气。这事并非无目的的安排，而是为免有丝毫猜疑，以为另一个人复活了。\par

\BibleQuote{有\Person{玛利亚玛达肋纳}和另一个\Person{玛利亚}，对着坟墓坐着。}她们为何在此等候？她们尚未知道任何关于祂的伟大、相称和高超的事，因此她们带了香料，在坟墓那里等候，以便若犹太人的狂怒稍缓，她们就去拥抱尸体。你们看见妇女的勇气吗？你们看见她们的情感吗？你们看见她们在钱财上的高尚精神吗？她们至死的高尚精神？\par

我们男人当效法那妇女们；在诱惑中不要离弃\Person{耶稣}。因为她们甚至为已死的祂花费那么多，并舍命；但我们（因为我又说同样的事）既不饿时喂养祂，也不赤身时给祂穿衣，反而看见祂乞讨，却走过。然而若你们看见祂本人，人人都会舍弃自己的一切财物。但现在也是一样。因为祂自己曾说过：\Quote{我就是。}那么你为什么不会舍弃一切呢？因为即使现在你也听见祂说：\Quote{你们为我做的。}并且你给这个人或给祂，没有区别；你并不比这些当时喂养祂的妇女有更少的，反而更多。但不要困惑！因为喂养祂显现的本体，虽足以打动铁石心肠，还不如因为祂的话而服事穷人、瘸子、弯腰的人。因为在前者，那显现者的面容和尊荣与你分功；但在此，赏报全归于你的慈爱；并且这是对祂更大敬畏的证明，当你因祂一句话而服事你的同仆，在一切事上使他得到安息。安慰他，并相信那接受的主说：\Quote{你们给了我。}因为除非你给了祂，祂不会算你配得天国。若你未曾离弃祂，祂就不会将你送入地狱，若你忽略了偶然遇见的人；但因为正是祂自己被轻视，因此责备是大的。\par

\Person{保禄}也是这样，在迫害属于祂的人时，迫害了祂；因此祂也说：\BibleQuote{你为什么迫害我？}\BibleRef{宗 9:4} 因此，当我们施与时，要觉得是施与基督自己。因为祂的话语确实比我们的眼见更确实。所以当你看见一个穷人时，要记得祂的话，祂藉此声明，是祂自己被喂养。因为那显现的虽然不是基督，但在这人的形像里，基督自己接受并乞讨。\par

但你听见基督乞讨，就羞愧吗？反而当你没有给那向你乞讨的祂时，你当羞愧。因为这是耻辱，这是刑罚和惩罚。因为祂乞讨是出于祂的良善，为此我们甚至应当夸耀；但你不给，是出于你的无人性。但若你现在不信，即因忽略一个信主的穷人而忽略了祂，到了那时，祂将把你领到中间并说：\BibleQuote{你们没有给这些最小中的一个做，就是没有给我做。}但愿天主不许我们这样学习，而宁可赐恩让我们现在相信，结果实，并听到那最蒙福的、领我们进入天国的声音。\par

但或许有人会说：\Quote{你每天都向我们谈论施舍和仁爱。}我也不会停止谈论这事。因为如果你们已经达到了，首先，即便如此我也不该停止，恐怕使你们更加懈怠；但如果你们达到了，我也许会稍加放松；但如果你们连一半也没达到，不要对我说这些话，而要对自己说。因为你们责备我，正好像一个小孩子，屡次听字母阿尔法，却学不会，反而责怪老师，因为他不断地、永远地提醒他。\par

这些讲论中有谁在施舍上变得更热忱了？谁放弃了金钱？谁献出了一半家产？谁献出了三分之一？没有人。那么，当你们没有学时，却要求我们停止教导，这不是荒谬吗？你们应当相反。即使我们有意停止，你们也应当阻止我们，并说：我们还没有学会这些事，你怎么就停止提醒我们了呢？假如有人眼睛患病，而我恰是医生，然后我给他敷药、涂油，并做了其他处理，但没有多大起色，于是就停止了；他难道不会到我诊所门口来，向我喊叫，指责我极大的疏忽，因为我自行撤走，而疾病仍在；而如果被责备时，我回答说，我敷了药，涂了油，他能忍受吗？绝不，反而立刻会说：\Quote{那有什么益处，如果我还疼痛呢？}对你们的灵魂也要这样推理。但假如我屡次热敷一只僵硬萎缩的手，却没有成功使它软化，难道我不会听到同样的话吗？即使现在，我们还在洗一只萎缩枯干的手，并且为此，直到我们完全把它伸展开，我们也不会停止。但愿你们也这样，在家里、在市场、在餐桌、在夜里，甚至在梦中，都不谈论别的事。因为如果我们白天常常留意这些事，就连在梦中也会从事它们。\par

你说什么？我永远在谈论施舍吗？我倒希望没有大需要向你们提出这个劝告，而是让我谈论与犹太人、外教人和异端者的争战；但当你们尚未痊愈时，谁能武装你们去作战呢？怎能带领你们上阵，而你们还有伤口和裂痕呢？因为如果我确实看见你们完全健康，我就会带领你们上那阵线，你们将藉着基督的恩宠看见成千上万的人倒毙，他们的头堆叠在一起。无论如何，在其他著作中，我们已谈论了许多关于这些事，但即便如此，我们也未能完全取得胜利，因为群众的懈怠。因为当我们在教义上打败他们一万次时，他们却用那些加入我们聚会的大众的生活、他们的伤口、他们灵魂的疾病来指责我们。\par

那时，我们若将你们置于战阵之中，又怎能满怀信心地向你们展示呢？你们反倒加害于我们，刚一被敌人击伤，就沦为笑柄。因为一个人的手患病萎缩，以致无法施舍；这样的人怎能持盾挡在身前，避免被残忍的讥讽所伤？另一些人脚跛，就是那些前往戏院和娼妓之所的人；他们又怎能立于战场，不被放荡的指控所伤？又有人患眼疾，目光不正，满心淫荡，攻击妇女的贞洁，颠覆婚姻；这样的人怎能直面敌人，挥舞长矛，投掷标枪，四面受讥讽的刺激？我们还可见许多人患腹疾，不亚于水肿病，被贪饕和酗酒所制伏；我又怎能带领这些醉汉去作战？又有人口舌腐烂，就是那些易怒、辱骂、亵渎的人；这样的人又怎能在战场上呐喊，成就任何伟大崇高之事？他同样沉醉于另一种醉态，给敌人带来许多笑柄。\par

因此，我每天在这营地里巡视，为你们包扎伤口，医治你们的疮痍。但若你们一旦振作起来，甚至能伤害他人，我必将这门战艺传授给你们，教导你们如何运用这些武器；或者更确切地说，你们自己的善工本身就是武器，众人会立刻顺服——如果你们变得仁慈、宽容、温和、忍耐，并彰显所有其他德行。若有人反驳，我们也将从我们这一方面提供明证，将你们引上前来；因为现在我们（至少在你们方面）倒是在这赛跑中受阻了。\par

请听。我们说基督成就了大事，使人变为天使；然后，当我们被要求交账，并需从这羊群中提供证据时，我们的口就被堵住了。因为我害怕，我带来的不是天使，而是从猪圈里牵出的猪，和因情欲而狂乱的马。\par

我知道你们感到痛苦，但这话并非针对你们所有人，而是针对有罪的人；或者更确切地说，甚至不是针对他们，若他们醒悟过来，而是为他们。因为如今一切都已败坏、毁灭，教会已变得不比牛棚、驴圈和骆驼圈好多少；我四处寻找一只羊，却找不到。众人都像马和野驴一样踢踹，用许多粪便充满这个地方——他们的言谈也是如此。若有人能看见每次集会中男女所说的话，就会看到他们的话比那些粪便更不洁。\par

因此，我恳求你们改掉这个恶习，好使教会散发馨香。但如今，我们在教会中存放感官的香料，却不用心清除并驱散头脑中的污秽。那有什么益处呢？因为我们把粪便带入教会，并不像彼此谈论这些事——谈论利益、商品、小买卖、与我们无关的事——那样使教会蒙羞；此处本应有天使的唱诗班，我们本应将教会变为天堂，除了恳切的祈祷和静默聆听之外，别无所知。\par

那么，让我们从此时起就这样做，好使我们净化自己的生活，并达到那应许的福分，藉着我们的主耶稣基督的恩宠和慈爱，愿光荣归于祂，世世无穷。阿们。\par

\SourceRecord{Translated by George Prevost and revised by M.B. Riddle.；Nicene and Post-Nicene Fathers, First Series；Vol. 10.；Edited by Philip Schaff.；Buffalo, NY: Christian Literature Publishing Co.,；1888.；<http://www.newadvent.org/fathers/200188.htm>.}

% structure: p0001 source=h1 target=section reason=page-title

\section{玛窦福音讲道集 第八十九篇}

\BibleRef{玛 27:62}\par

\Quote{\Emph{次日，就是预备日之后，司祭长和法利塞人一同来见比拉多，说：\Quote{大人，我们记得那骗子还活着的时候曾说：\InnerQuote{三天后我要复活。}所以请下令将坟墓把守妥当，直到第三天，免得祂的门徒来把祂偷去，然后对百姓说：\InnerQuote{祂从死人中复活了。}那最后的欺骗就比先前的更坏了。}}}\par

欺骗处处自食其果，并违背自己的意愿支持真理。请看。祂的死亡、复活和埋葬，这一切本应被相信，却由祂的仇敌促成了。无论如何，这些话正为这些事实作证。\Quote{我们记得}，正是这些话，\Quote{那骗子还活着的时候曾说}（因此祂现在已经死了），\Quote{三天后我要复活。} \Quote{所以请下令将坟墓封好}（因此祂已被埋葬），\Quote{免得祂的门徒来把祂偷去。} 因此，如果坟墓被封好，就不会有欺骗。因为不可能有。所以，通过你们所提出的，祂复活的证据已变得无可辩驳。因为坟墓被封了，就没有欺骗。但如果没有欺骗，而坟墓空了，显然祂复活了，明确而无可辩驳。你看，他们如何违背自己的意志为真理的证据争辩。\par

但请你留意门徒们对真理的爱，他们如何不向我们隐瞒仇敌所说的任何话，尽管那些话是辱骂的。无论如何，他们甚至称祂为骗子，而这些人对此并不沉默。\par

但这些事也显出他们的残忍（即甚至到祂死时，他们仍未放弃愤怒），以及这些人的纯朴和诚实的性情。\par

但值得探究那一点，即祂在哪儿说过\Quote{三天后我要复活}？因为人们不会发现这话如此明确地陈述，而是约纳的预像。因此，他们理解了祂的话，却故意行事不公。\par

那么比拉多说什么？\Quote{你们有看守的兵；尽你们所能去把守妥当。} \Quote{他们就去把守妥当，封了坟墓，并派兵看守。} 他不让士兵单独封墓，因为他已经得知有关基督的事，不再愿意与他们合作。但为了摆脱他们，他容忍了这事，并说：\Quote{你们随意去封吧，免得你们有机会责备他人。} 因为如果只有士兵封了墓，他们可能会说（尽管这种说法不可信且虚假，但正如在其他事上他们不顾羞耻，在这事上他们也可能会说）：士兵把尸体交出去让人偷走，给了祂的门徒机会编造关于祂复活的故事，但现在他们自己封了墓，就再也不能说这样的话了。\par

你看见他们如何违背自己的意志为真理效劳吗？因为他们自己来见比拉多，自己请求，自己封墓，并派兵看守，以便彼此指控和驳斥。而且，他们什么时候能偷走祂呢？在安息日？怎么可能？因为甚至出门也是不合法的。\BibleRef{出 16:29} 即使他们触犯法律，那些如此胆怯的人怎敢出来？他们又如何能说服群众呢？说什么？做什么？凭着什么样的热忱，他们能为死人辩护？期望什么样的报酬？什么样的回报？看见祂活着时仅仅被捕，他们就逃跑了；祂死后，除非祂已复活，否则他们怎敢为祂辩护？这些事合理吗？因为他们既不愿意也不能捏造一个没有发生的复活，从这一点可以清楚看出。祂曾多次向他们讲论复活，并不断地说，正如这些人所说的，\Quote{三天后我要复活。} 因此，如果祂没有复活，显然这些人（因祂的缘故受骗并成为整个民族的敌人，变得无家无城）会憎恨祂，不愿将这种荣耀归给祂，因为他们受骗，并因祂而陷入极大的危险。因为他们甚至不可能捏造复活，除非复活是真实的，这一点甚至无需推理。\par

因为他们凭什么自信？凭他们推理的机敏吗？不，他们是最无知的人。凭他们财富的丰富吗？不，他们既没有手杖也没有鞋。凭他们民族的杰出吗？不，他们是卑微的，出身卑微。凭他们国家的伟大吗？不，他们来自偏僻的地方。凭他们的人数众多吗？不，他们不超过十一人，而且分散各处。凭他们师傅的应许吗？什么样的应许？因为如果祂没有复活，那些应许也不会被他们相信。他们又如何能忍受一个疯狂的民族？因为如果他们的首领连一个看门的使女的话都受不了，如果其余的人看见祂被捆绑都逃散了，他们怎会想到跑到地极，去传播一个捏造的复活故事呢？因为如果那个人连一个女人的威吓都抵挡不住，他们连看见捆绑都受不了，他们又怎能抵挡君王、官长、和外邦人呢？那些地方有刀剑、烤架、火炉，以及成千上万的死亡，日复一日，除非他们得到复活之主的能力和恩宠？如此多的奇迹和如此大的事发生了，犹太人却一样也不在意，反而把行这些事的祂钉在十字架上，他们又怎会仅仅因为这些人关于复活的一番话就相信呢？这些事不可能，绝不可能，而是复活之主的权能成就了这一切。\par

2. 但是，请注意，我恳求你们，看看他们的诡计是多么荒谬。\Quote{我们记得，}这是他们的话，\Quote{那个骗子活着的时候曾说：三天后我要复活。}然而，如果他是骗子，狂妄地说谎，你们为什么害怕、跑来跑去，如此费心？我们害怕，有人回答说，恐怕门徒偷走他，欺骗群众。然而，这已经被证明毫无可能性。但恶意是一件好争辩且无耻的事，它试图做不合理的事。\par

他们吩咐严密看守三天，仿佛是为教义争辩，并有意证明在那之前他也是骗子，他们将恶意甚至延伸到他的坟墓。正因如此，他提前复活了，免得他们说他撒谎，并被偷走。他的提前复活是无懈可击的，而延迟则会充满猜疑。因为如果他当时没有在他们坐着看守时复活，而是在三天后他们撤走时才复活，他们就会有所说、有所反对，尽管是愚蠢的。因此，他提前了时间。因为复活应当在他们坐看守时发生。所以，它也应当在三天内发生，因为如果过了三天，那些人撤走了，这件事就会被怀疑。因此，他也允许他们随意封住坟墓，并派士兵围坐。\par

他们不在意做这些事并在安息日工作，只关注一个目的，即他们邪恶的意图，好像由此他们就能成功；这是极度愚蠢的标志，也是恐惧现在大大惊吓他们的标志。因为当他们活着时抓住他的人，现在却害怕死了的他。然而，如果他只是一个凡人，他们就有理由鼓起勇气。但为了让他们知道，当他活着时也是自愿承受他所承受的，看啊，既有封条、石头和看守，他们却无法留住他。但只有一个结果：埋葬被宣扬，复活也因此被证明。因为士兵确实坐在旁边，犹太人也在警戒。\par

\Quote{但在安息日将尽，向着一周的第一天破晓时，玛利亚玛达肋纳和另一个玛利亚来看坟墓。看，发生了大地震，因为主的一位天使从天上降下，前来把石头从墓门滚开，坐在上面。他的容貌好像闪电，衣服洁白如雪。}\par

复活后，天使来了。那么他为什么来并挪开石头呢？因为那些妇女，她们当时在坟墓里看见了他。所以，为了让他们相信他已复活，他们看见坟墓空了。为此他挪开石头，为此也发生了地震，好使他们彻底醒悟和惊醒。因为他们来是要给他抹油，这些事发生在夜里，可能有些人也昏昏欲睡。他为什么说\Quote{你们不要害怕？}首先他解除他们的恐惧，然后告诉他们复活的事。\Quote{你们}这个词显示对他们极大的尊荣，并表明那些胆敢做别人所做的事的人要面临极重的惩罚，除非他们悔改。因为害怕不是为你们，他意思是，而是为那些钉死他的人。\par

然后，他用言语和容貌解除了他们的恐惧（因为他显出明亮的样子，如同带来如此好消息），接着说：\Quote{我知道你们寻找被钉十字架的耶稣。}他不耻于称他为\Quote{被钉十字架的}，因为这是福泽的顶峰。\par

\Quote{他已经复活了。}\BibleRef{玛 28:6}从哪里是明显的呢？\Quote{照他所说的。}所以他似乎说：如果你们不肯相信我，请记住他的话，你们也就不会不信我。然后又一个证明：\Quote{来看他躺过的地方。}因为他已经挪开了石头，好叫他们从这点也能得到证明。\Quote{去告诉他的门徒，你们要在加里肋亚看见他。}他预备他们去传好消息给别人，这件事最使他们相信。他说\Quote{在加里肋亚}说得很好，使他们摆脱困扰和危险，免得恐惧阻碍他们的信德。\par

\Quote{她们带着恐惧和喜乐离开坟墓。}这是为什么呢？她们看见了一件令人惊奇、超出预料的事：坟墓空了，而之前她们看见他被放在那里。因此他也引导她们去看，使她们成为两件事的见证人：既是他的坟墓，也是他的复活。因为她们认为没有人能在他身边有这么多士兵看守时把他带走，除非他自己复活。因此她们也欢喜惊奇，并得到与他长久相处的回报：她们首先看见并喜乐地宣告，不仅是所说的话，还有她们所看见的。\par

3. 所以，当她们带着恐惧和喜乐离开后，\Quote{看，耶稣迎上她们，说：愿你们平安。}但\Quote{她们抱住他的脚}，以极大的喜乐和欢喜跑到他面前，通过触摸也得到复活的确凿证明和十足保证。\Quote{她们朝拜了他。}那么他说什么呢？\Quote{不要害怕。}他再次亲自驱除她们的恐惧，为信德开路：\Quote{去告诉我的弟兄们，叫他们往加里肋亚去，在那里他们要看见我。}请注意，他亲自通过这些妇女向他的门徒传报好消息，正如我常说的，将那最受轻视的性别提升到尊荣和善望，并医治那有病的。\par

或许你们有人愿意像她们那样，抱住耶稣的脚；你们现在就可以做到，不只抱住祂的脚和手，甚至抱住那神圣的头，以纯洁的良心领受可畏的奥迹。但是不仅在此，在那一天你们也要看见祂，带着那说不出的荣耀和众多天使来临，如果你们存心仁慈；你们不仅会听到\Quote{万福！}这些话，还会听到另一些话：\BibleQuote{我父所祝福的，你们来承受从创世以来为预备的国度吧！}\BibleRef{玛 25:34}。\par

所以，你们要仁慈，好能听到这些话。你们这些佩戴金饰的妇女，看到了这些妇女的奔跑，最后，虽然晚了，也要摒弃贪求金饰的疾病。如果你们效法这些妇女，就改变你们所佩戴的装饰，却以施舍作衣装。请问，这些宝石和金线绣花的衣裳有什么用处？你说：\Quote{我的灵魂，} 你说，\Quote{就喜欢这些，以此为乐。} 我问你益处，你却告诉我伤害。因为没有什么比沉迷于这些东西，以此为乐，被它们钉住更糟的了。因为这种痛苦的奴役更苦，当人甚至乐意作奴隶时。因为她何时会在属灵事务上按应有的样子勤勉？何时会按应有的样子嘲笑今世的事物？她竟把被金链锁住视为一件值得高兴的事！因为那留在监牢里并感到快乐的人，永远不愿意被释放；这妇女也一样；她如同成了这邪恶欲望的俘虏，再也不能以应有的渴望和勤奋去听属灵的话，更不用说从事这样的工作了。\par

那么这些装饰和奢华有何益处，我请问你们？你说：\Quote{我喜欢它们，} 你说。\Quote{但是，} 你说，\Quote{我也从观看者那里得到许多荣耀。} 这是什么？这是另一场毁灭的缘由，当你被高举到傲慢、自大时。来吧，既然你没有告诉我益处，请容忍我告诉你祸害。那么，由此产生的祸害是什么？焦虑，比快乐更大。因此，许多观看者，我是说那些较为粗俗的人，从这装饰中得到的快乐比佩戴金子的妇女更多。因为你确实带着焦虑打扮自己，而他们却不用焦虑，饱享眼福。\par

此外，还有其他事：灵魂的卑下，四面遭受嫉妒。因为邻近的妇女被刺痛，就武装起来对抗自己的丈夫，并激起严重的战争对付你。此外，所有闲暇和焦虑都花在这件事上，不认真从事属灵的成就；充满傲慢、自大和虚荣；被钉在地上，失去翅膀，不是鹰而是狗或猪。因为你放弃了仰望天堂和在那里飞翔，就像猪一样弯向大地，好奇于矿脉和洞穴，有一个不男子气和卑下的灵魂。但当你出现时，你难道吸引市场上的人的目光吗？好，正因如此，你不应佩戴金子，以免成为众矢之的，使许多指责者开口。因为那些看你的人中，没有人羡慕你，反而嘲笑你，说你爱打扮、自负、属肉体的女人。如果你进入教堂，你出去时一无所获，只有无数的斜眼、辱骂和咒诅，不仅来自观看者，也来自先知。因为那声音最洪亮的\Person{依撒意亚}，一看见你，就会喊道：\BibleQuote{主论到熙雍骄傲的女子这样说：因为他们昂首行走，媚眼传情，行走时拖曳衣裙，碎步而行，主必除掉她们的美饰，必有尘土代替馨香，必有绳索代替腰带。}\BibleRef{弟前 2:9}\par

这些话是对你的华丽服饰说的。因为这些话不仅对她们说，也对每个仿效她们的女人说。\Person{保禄}又同他一起作为控告者，吩咐\Person{弟茂德}命令妇女\BibleQuote{不要以编发、金子、珍珠或华服来装饰自己。} 所以，佩戴金子到处都有害，尤其当你进入圣堂、经过穷人时。因为如果你极其焦虑地要控告自己，就没有任何装饰比这副残忍和不人道的面孔更适合你了。\par

4. 试想，至少，你穿着这身打扮经过多少饥饿的肚腹，用这魔鬼般的展示经过多少赤身露体的身体。喂养饥饿的灵魂，比穿通你的耳垂，从上面悬挂无数穷人的食物，毫无目的或益处，要好得多。什么？富有是值得称赞的吗？什么？佩戴金子是荣耀吗？即使这些东西是从正当收入得来，这样你所做的已经是极重的罪责；但若更是从不当得利而来，试想其重大程度。\par

但是你爱慕赞美和荣誉吗？那么，脱去这可笑的衣服，这样人人都会钦佩你；那时你既能享有荣誉，又能享有纯粹的快乐。因为现在，你无论如何都被嘲弄淹没，因这些事给自己制造了许多烦恼的缘由。因为若其中任何一样缺失，想想由此生出的多少灾祸：多少婢女被打，多少人陷入困苦，多少人被处死，多少人被投入监牢。审判由此而生，诉讼亦然，无数的咒骂和控告：丈夫对妻子，妻子对朋友，灵魂对自己。\Quote{但是它们不会丢失。}首先，这不易确保，即使它们常常被保存，但保存本身也带来许多焦虑、挂虑和不适，并无益处。\par

因为这会给家庭带来什么益处？给穿戴它的女人自己带来什么好处？毫无益处，反而有许多不雅观，和来自各方的指责。你这样穿戴，怎能亲吻基督的脚，并依偎它们呢？这种装饰是他所厌恶的。为此，他屈尊降生于木匠之家，甚至不是在那家里，而是在一个棚屋和马槽中。那么，你怎能看见他呢？你在他眼中没有可羡慕的美，没有穿在他面前可爱的服饰，反而穿着可憎的。因为来到他跟前的人，不可用这样的衣服装饰自己，而要以德行作衣装。\par

请想想这些宝石终究是什么。不外乎尘土和灰烬。用水调和，它们就成了泥。想想吧，并羞愧于让泥土作你的主人，抛弃一切，依恋它，携带并背负它，甚至在你进入教堂时——那时你最应该远离它。因为教堂不是为此建造的，让你在里面展示这些财富，而是展示属灵的财富。但你，仿佛进入华丽的游行队伍，这样四面装饰自己，仿效舞台上的妇人，你也这样大量携带那可笑之物。\par

因此，我告诉你，你给许多人带来祸害，当会众散去后，在他们的家里，在餐桌上，你可以听到大多数人在描述这些事。因为他们已不说：先知如此说，宗徒如此说；他们却描述你衣服的昂贵，你宝石的大小，以及穿戴这些之人的所有其他不雅事。\par

这使你和你的丈夫在施舍上落后。因为你们中没有人会轻易同意打碎一件这样的饰物来供养穷人。既然你宁愿自己受苦也不愿看见这些东西被打碎，你又怎会以它们为代价喂养别人呢？\par

因为大多数妇女对这些东西的感受，如同对待某些活物，甚至不亚于对待她们的孩子。\Quote{天主不允许，}你说。那么，用你的行为向我证明吧，因为至少现在我看到了相反的事。那些完全被俘获的人中，有谁会熔掉这些东西来拯救一个孩子的灵魂免于死亡呢？我为什么说孩子呢？有谁借此赎回了自己的灵魂，当灵魂灭亡时？不，相反，大多数人甚至每天为这些东西出卖灵魂。若任何身体上的疾病发生，他们会做一切；但若他们看到自己的灵魂堕落，却不如此操心，反而对子女和自己的灵魂都漠不关心，为的是让这些东西能够留存到因时间生锈。\par

当你佩戴着价值万金的宝石时，基督的肢体却连必要的食物都未能享用。而万有的共同主将天堂和天上的事物以及属灵的筵席公平地分施给众人，你却连可朽坏的东西也不分施给他，为的是让你能永远被这些沉重的锁链束缚。\par

因此有无数的灾祸，因此男人们犯奸淫，因为你为他们准备了抛弃自制，你教导他们以那些娼妓用来装饰自己的东西为乐。为此，他们如此迅速地就被俘获。因为如果你曾教导他轻视这些东西，并喜爱贞洁、敬畏天主和谦卑，他就不会如此轻易地被奸淫之箭射中。因为娼妓能这样装饰自己，甚至更甚，但其他饰物则不能。那么，让他习惯于喜爱那些不能放在娼妓身上的饰物。你如何使他养成这个习惯呢？如果你脱下这些，穿上那些，那么你的丈夫就安全，你也有荣誉，天主会对你仁慈，所有人都会钦佩你，你将获得将来的美善，藉着我们的主耶稣基督的恩宠和慈爱，愿光荣和权能归于他，直到无穷之世。阿们。\par

\SourceRecord{Translated by George Prevost and revised by M.B. Riddle.；Nicene and Post-Nicene Fathers, First Series；Vol. 10.；Edited by Philip Schaff.；Buffalo, NY: Christian Literature Publishing Co.,；1888.；<http://www.newadvent.org/fathers/200189.htm>.}

% structure: p0001 source=h1 target=section reason=page-title

\section{论玛窦福音第九十篇讲道}

\BibleRef{玛 28:11-14}\par

\Quote{\Emph{他们去的时候，看，几个卫兵进了城，将所发生的事全都报告给司祭长。}\Emph{司祭长同长老聚会商议之后，就给了士兵许多钱，说：你们就说：他的门徒夜间来了，趁我们睡觉时把他偷去了。如果这事传到总督耳中，我们会说服他，保证你们无事。}}\par

为了这些士兵，那地震发生了，为要惊吓他们，并且使见证来自他们——结果确实如此。因为来自卫兵自己的报告，就没有嫌疑。因为那些征兆，有些是公开向世界显示的，有些是私下向在场的人显示的：公开的是黑暗，私下的是天使的出现和地震。当他们来报告时（因为真理照耀，藉着它的敌人而被宣扬），他们又给钱，要他们按所说的那样说：\Quote{他的门徒来把他偷去了。}\par

他们怎么偷的？你们这些最愚蠢的人啊！因为真理的清晰和显明，他们甚至无法编造谎言。因为他们所说的话极不可信，连表面合理性都没有。请问，门徒们——那些贫穷无学识、连露面都不敢的人——怎么偷他？什么？难道没有封条吗？什么？难道没有那么多卫兵、士兵和犹太人围在那里吗？什么？难道那些人没有预料到这一点，并思虑、不眠、为之焦虑吗？而且，他们偷他到底是为了什么？为了假装复活的道理？可是，怎么会有人想到假装这样的事——那些宁愿隐藏和活着的人？他们怎么能挪开那封牢的石头？他们怎么能躲过那么多人的眼目？即便他们不怕死，也不会无缘无故、徒劳地冒险去对抗那么多看守的人。而且，他们是胆怯的——他们先前所做的事清楚地表明了这一点：至少，当他们看见他被捕时，都从他身边逃跑了。如果那时他们看见他活着都不敢站住，那么当他死了以后，他们怎能不怕那么多士兵呢？什么？是要破门吗？是要不被发现吗？一块大石头放在上面，需要很多人才能挪动。\par

他们说：\Quote{最后的迷惑比先前的更坏}，\BibleRef{玛 27:64}——这话是针对他们自己说的，因为当他们经过那么多疯狂的行为后本应悔改，却反而企图超过先前的行为，编造荒谬的谎言。正如他活着时他们收买他的血，同样当他死了又复活了，他们又用金钱企图破坏他复活的证据。但请你注意，他们如何处处被自己的行为捉住。因为他们若不到比拉多那里去，也不要求卫兵，他们倒能更厚颜无耻地行事；但事实并非如此。事实上，他们好像竭力要堵住自己的嘴，所做的一切都是如此。因为如果门徒们没有力量与他一起警醒——而且是被他责备之后——他们怎么敢做这些事呢？而且，为什么他们不在这之前偷他，而是等到你们来了之后呢？如果他们有意这样做，他们会在第一夜——当坟墓还没有被看守时——就做，那时做起来没有危险且安全。因为他们在安息日来向比拉多求看守，并看守了坟墓；但第一夜，没有一个人在现场看守坟墓。\par

2. 而且，那些粘着没药的殓布又是什么意思？因为伯多禄看见这些布放在那里。如果他们有心偷盗，他们不会偷裸体，不仅因为不尊敬，而且为了不拖延时间、不耽误剥衣，不给那些可能醒来抓住他们的人机会。特别是没药是一种紧贴身体、粘着衣服的药，所以不易将衣服从身体上取下，做这事需要很多时间，因此从这一点来看，偷盗的说法也是不可信的。\par

什么？他们不知道犹太人的愤怒吗？不知道他们会把怒气发泄在他们身上吗？如果他没有复活，这对他们有什么益处呢？\par

这些人知道这一切都是他们编造的，就给了钱说：\Quote{你们就这样说，我们会说服总督。}因为他们希望这说法被传播，徒劳地与真理作战；他们越企图掩盖它，就越无意中使它显明。事实上，这甚至证实了复活——我指的是他们所说的门徒偷了他的话。因为这是那些承认身体不在那里的人的话。因此，当他们承认身体不在那里，而偷盗之说被证明是虚假和不可信的——有他们的看守、封条和门徒的胆怯为证——复活的证明即由此显明不可反驳。\par

然而，这些无耻狂妄的人，尽管有那么多事堵住他们的口，他们的话却是：\Quote{你们就说，我们会说服，并保证你们平安。}你看到一切都败坏了吗？比拉多（因为他被说服了）？士兵们？犹太民众？但不要惊讶，如果金钱胜过了士兵。因为如果它在门徒身上都显出如此大的力量，何况在这些士兵身上呢。\par

\Quote{这说法直到今天还在犹太人中间流传}，经上说。你再次看到门徒们热爱真理：他们并不羞于说出这样的话——这样的说法曾反对他们。\par

\Quote{随后，十一个门徒往加里肋亚去，有人见了就朝拜，有人却怀疑。}\par

在我看来，这似乎是祂在加里肋亚最后一次显现，那时祂派遣他们去施洗。若有人\Quote{怀疑}，在此又当钦佩他们的诚实，他们连自己的缺点也隐瞒到最后一刻。然而，就连这些人也因眼见而得到确信。\par

那么，当祂看见他们时，祂对他们说了什么？\Quote{天上地下的一切权柄都交给了我。}祂又用更合于人的方式对他们说话，因为他们尚未领受那能提升他们的圣神。\Quote{你们要去使万民成为门徒，因父、及子、及圣神之名给他们付洗，教训他们遵守我所吩咐你们的一切。}一方面以教训相托，另一方面以诫命相托。祂没有提及犹太人，也没有提出已经成就的事，更没有责备伯多禄的背主或其他人逃跑，而是将教义的概要——即洗礼经文所表达的——交在他们手中，命令他们向普世倾注。\par

此后，因为祂已吩咐他们大事，为鼓舞他们的勇气，祂说：\Quote{看！我同你们天天在一起，直到今世的终结。}（\BibleRef{玛 28:20}）你看到祂固有的权能了吗？你看到那些其他的话也是出于迁就了吗？祂不仅应许与那些人同在，也与所有后来信仰的人同在。因为显然宗徒们不会留到\Quote{今世的终结}，但祂是对信徒们如同对一个身体说话。因为，祂说，不要对我说事情的艰难，因为\Quote{我与你们同在}，我使一切变得容易。这在旧约中祂也常对先知们说，如对推辞自己年幼的耶肋米亚（\BibleRef{耶 1:6}），以及对畏缩使命的梅瑟和厄则克耳，祂说\Quote{我与你们同在}，在这里也对这些人说。请你留意他们的卓越之处，因为其他人被派到一个民族时，常常推辞，而这些人被派到整个世界，却毫无怨言。祂也提醒他们终结之事，为更多地吸引他们，使他们不仅看到眼前的危险，也看到那无尽的未来美物。\par

\Quote{因为，祂说，你们所忍受的厌烦之事都与现世生命一同结束，因为至少这世界本身也会终结，而你们将享受的美物却永存不朽，正如我以前常对你们说的。}如此，祂藉着那日的回忆，振奋并激发了他们的心灵，然后派遣了他们。因为那日对于活在善工中的人是可渴望的，反之，对于在罪中的人，则如同对受审判者一样可怕。\par

但我们不仅当恐惧战栗，也当趁有机会时悔改，从恶中奋起，因为只要我们愿意，我们就能。因为在恩宠之前许多人尚且如此，何况在恩宠之后呢？\par

3. 因为我们被吩咐了什么难事？要移开山岭？要飞上天空？还是要渡过托斯卡纳海？绝不是，而是一种如此容易的生活方式，几乎不需要任何工具，只需要灵魂和意志。因为这些宗徒成就了如此大事，用了什么工具呢？他们不是只穿一件外衣、赤着脚四处奔走吗？他们战胜了一切。\par

因为诫命中有什么难的呢？不要有仇敌。不要恨任何人。不要说任何人的坏话。不，相反的事才是更大的艰难。但你反驳说，祂说：舍弃你的钱财。这难道就是难事吗？首先，祂没有命令，而是劝告。但即使这是命令，抛弃负担和不及时的挂虑，又有什么难的呢？\par

哦，贪得无厌！一切都成了金钱；为此一切都被颠倒了。若有人称另一个人有福，便提到这一点；若说他可怜，也是由此得出可怜的描述。一切衡量都为此：这个人如何变富，那个人如何变穷。无论是从军、婚姻、贸易，还是任何人着手做的任何事，他并不专注于所提议之事，除非看到财富迅速涌入。此后，我们难道不该聚集商议如何驱除这场瘟疫吗？难道不该为我们祖先的善行感到羞愧吗？那三千、五千人，他们不是一切公用吗？\par

今生的利益有何用处，如果我们不将其用于将来的获益？你们要到何时才不奴役那奴役了你们的钱财？你们做钱财的奴隶要到何时？你们要到何时才不爱自由，不撕毁贪婪的契约？然而，如果你们成了人的奴隶，若有人承诺你们自由，你们便做一切事；但你们被贪婪俘虏，却甚至不考虑如何摆脱这痛苦的束缚。然而，前者并不可怕，后者却是最痛苦的暴政。\par

请想想基督为我们付出了多大的代价。祂倾流了自己的血；祂交付了自己。但你，即使在此之后，却变得懈怠；而最严重的是，你甚至以奴隶为乐，以耻辱为乐，那本应逃避的，竟成了你所渴望的。\par

但是，既然不仅要哀叹和责备，还要纠正，就让我们看看，这种激情和邪恶是从什么原因成为你们渴望的对象的。那么，它从何而来成为渴望的对象？因为，你说，它使我得到尊荣和安全。请问，是怎样的安全？是信心，相信不会遭受饥饿、寒冷、伤害和轻视。那么，如果我们向你许诺这种安全，你会放弃致富吗？如果财富正是为此而被渴望，而你可以不靠它们就获得安全，那么你还需要它们做什么？\Quote{这怎么可能呢？}你说，\Quote{一个不富有的人如何能获得这些？}不，怎么可能（我说相反的事），如果一个人富有？因为必须讨好许多人，包括统治者和臣民，恳求无数人，成为卑贱的奴隶，恐惧战栗，猜忌嫉妒者的眼神，害怕诬告者的舌头，以及别的贪婪之人的欲望。但贫穷却不是这样，而是完全相反。它是一个避难所和安全处，一个平静的港湾，一个角力场和练习自我克制的学校，是天使生活的模仿。\par

请听这些话，所有贫穷的人；或者更确切地说，所有渴望致富的人。可怕的不是贫穷，而是不愿意贫穷。把贫穷视作不可怕的，它对你就不再可怕。因为这种恐惧并不在事物的本性中，而是在懦弱之人的判断中。或者更确切地说，我甚至羞愧于有必要如此多地说论贫穷，以表明它并不可怕。因为如果你操练自我克制，它甚至是你无数祝福的泉源。如果有人向你提供主权、政治权力、财富和奢华，然后把贫穷放在它们对面，让你选择取哪个，你会立刻抓住贫穷，如果你确实知道它的美。\par

4. 我知道，当这些话被说出来时，许多人发笑；但我们并不困扰，我们只要求你们停留片刻，很快你们就会与我们一同判断。因为在我看来，贫穷好像一位美丽、匀称、可爱的少女，而贪婪则像一个怪物般的女人，一个\OriginalTerm{斯库拉}{Scylla}或\OriginalTerm{许德拉}{Hydra}，或其他一些寓言作家虚构的怪物。\par

因为，我请求你们，不要提出那些指责贫穷的人，而是提出那些因贫穷而闪耀的人。\Person{厄里亚}在贫穷中受养育，被接升天，获得那有福的升天。\Person{厄里叟}因贫穷而闪耀；\Person{若翰}也是；所有宗徒也是；而与之相反，\Person{阿哈布}、\Person{依则贝耳}、\Person{革哈齐}、\Person{犹达斯}、\Person{尼禄}、\Person{盖法}则被定罪。\par

但如果你们愿意，让我们不仅看那些因贫穷而光荣的人，也让我们观察这位少女本身的美。确实，她的眼睛清澈锐利，没有丝毫浑浊，不像贪婪的眼睛，时而充满愤怒，时而沉溺于享乐，时而被不节制所困扰。但贫穷的眼睛不是这样，而是温和、平静、慈祥地看着所有人，温柔、和蔼，不恨任何人，不躲避任何人。因为哪里有财富，哪里就有敌意和无尽的战争。贪婪的嘴则充满侮辱、傲慢、许多夸口、咒骂、欺骗；而贫穷的嘴和舌头是健康的，充满不断的感恩、祝福、温柔、慈爱、礼貌、赞美、称许的话语。如果你还想看她肢体的比例，她身材匀称高大，远高于财富。如果许多人逃避她，不必惊讶，因为愚人也逃避其他德行。\par

但你会说，穷人被富人侮辱。你又在向我宣告贫穷的赞美。因为，我请问，谁是有福的？侮辱者还是被侮辱者？显然是被侮辱者。但是，贪婪促使一个人侮辱他人；贫穷则劝人忍耐。\Quote{但穷人遭受饥饿，}你会说。\Person{保禄}也遭受过饥饿和饥荒。\Quote{他没有安息。}\BibleQuote{连人子都没有枕头的地方。}\BibleRef{玛 8:20}\par

你看到对贫穷的赞美到了何等程度，它把你放在何处，引你到什么人面前，使你成为主的追随者？如果拥有黄金是好的，基督——祂拥有不可言喻的祝福——就会把这赐给祂的门徒。但现在，祂不仅没有赐给，还禁止他们拥有。因此\Person{伯多禄}非但不以贫穷为耻，反而夸耀它，说：\BibleQuote{金银我都没有，但我把我所有的给你。}\BibleRef{宗 3:6}你们中谁不愿意说出这句话？我们所有人都极其愿意，也许有人会说。那么，扔掉你们的银子，扔掉你们的金子。\Quote{如果我扔掉，}你会说，\Quote{我会得到伯多禄的能力吗？}告诉我，是什么使伯多禄有福？难道是因为他使瘸子站起来吗？绝不是，而是没有这些财富，这为他赢得了天堂。因为行这些奇迹的人中，许多坠入了地狱，而做那些好事的人却获得了天国。这一点你甚至可以从伯多禄本人学到。他说了两件事：\BibleQuote{金银我都没有，}和\BibleQuote{因耶稣基督的名，起来行走！}\par

那么，是什么使得祂荣耀而蒙福呢？是使跛子起来，还是抛弃他的钱财？这事你们可从竞赛的主宰——祂自身——学到。那么，祂自己对寻求永生的富少年说了什么呢？祂并没有说：\Quote{使跛子起来}，而是说：\BibleQuote{变卖你的家产，分给穷人，然后来跟随我，你必有宝藏在天上。}\BibleRef{玛 19:21} \Person{伯多禄}也并没有说：\Quote{看，因你的名字我们驱逐魔鬼；}尽管他确实在驱逐魔鬼，而是说：\BibleQuote{看，我们舍弃了一切，跟随了你，那么我们将得到什么呢？}\BibleRef{玛 19:27} 基督在回答这位宗徒时，也并没有说：\Quote{若有人使跛子起来}，而是说：\BibleQuote{凡舍弃了房屋或田地的，必要在今生得到百倍，并承受永生。}\par

那么，让我们也效法这个人，使我们不致羞愧，而能满怀信心地站在基督的审判台前；使我们赢得祂与我们同在，如同祂曾与祂的门徒同在一样。因为祂将与我们同在，正如祂曾与他们同在，只要我们愿意跟随他们，并效法他们的生活和言行。因为天主由于这些事而加冕并称赞人，并不要求你复活死人，或医治跛子。因为使人像\Person{伯多禄}的，不是这些事，而是抛弃自己的财物，因为这是宗徒们的成就。\par

但是，你难道不觉得抛弃它们是可能的吗？首先，我说，这是可能的；但我并不强迫你，如果你不愿意，也不勉强你；但我恳求你，至少要拿出部分用于周济穷人，并且除了必需品之外，不要再为自己寻求什么。因为这样我们既能在此生无忧无虑、安稳度日，又能享受永生；愿天主恩赐我们众人得以获得永生，仰赖我主耶稣基督对人类恩宠与慈爱，愿光荣与权能归于祂，偕同圣父及圣神，从今世直到永远，世世无穷。阿们。\par

\SourceRecord{Translated by George Prevost and revised by M.B. Riddle.；Nicene and Post-Nicene Fathers, First Series；Vol. 10.；Edited by Philip Schaff.；Buffalo, NY: Christian Literature Publishing Co.,；1888.；<http://www.newadvent.org/fathers/200190.htm>.}
